\documentclass[11pt]{article}
\usepackage[margin=1in]{geometry}
\usepackage{amsmath,amssymb,amsthm,mathtools}
\usepackage{enumitem}
\usepackage[hidelinks]{hyperref}
\usepackage[nameinlink,noabbrev]{cleveref}

\newtheorem{theorem}{Theorem}[section]
\newtheorem{lemma}[theorem]{Lemma}
\newtheorem{proposition}[theorem]{Proposition}
\newtheorem{corollary}[theorem]{Corollary}
\newtheorem{definition}[theorem]{Definition}
\newtheorem{remark}[theorem]{Remark}
\newtheorem{assessment}[theorem]{Assessment}
\newtheorem{firewall}[theorem]{Firewall}
\newtheorem{programme}[theorem]{Programme}
\newtheorem{audit}[theorem]{Audit}
\newtheorem{decision}[theorem]{Decision}

\newcommand{\Tr}{\operatorname{Tr}}
\newcommand{\spec}{\operatorname{spec}}
\newcommand{\sgn}{\operatorname{sgn}}

\title{Renewal Standard Model--Einstein Continuum Research Notes\\
Ninth Series, V1}
\author{}
\date{\today}

\begin{document}
\maketitle
\tableofcontents

\section{Context, archive boundary, and purpose}
\label{sec:v9v1}
\label{sec:n9v1-context}

These notes continue a constructive programme whose goal is a
cutoff-independent continuum containing the Standard Model fields coupled
to Einstein gravity.  The target includes a coherent chiral fermion
measure, gauge and Higgs sectors, gravitational constraints and stress
response, regulator independence, a positive Euclidean state with the
required Ward and reflection properties, and eventual reconstruction or
Lorentzian interpretation.  The complete object has not been constructed.

The preceding active source reached its hundred-version boundary and has
been frozen byte for byte as
\begin{center}
\texttt{Renewal\_SM\_Einstein\_Continuum\_Research\_Notes\_V100\_f8.tex}.
\end{center}
It has \(40\,777\) logical source lines, \(1\,449\,798\) bytes, and
SHA--256
\begin{center}
\small\ttfamily
D8C6AE875DBD019AFB9CDA21F0B941F5F600CB8A20EC079C6FACFB67EED39502.
\end{center}
Earlier frozen SM files remain under their existing suffixes.  This compact
ninth-series V1 maps the proved interfaces and open gates; the full proofs
remain in the frozen sources.  Future versions append to this active file
and replace only the preceding active numeric version.  The newest
Yang--Mills and Navier--Stokes notes are audited every fifth version.  At
V100 this series will be frozen with the next unused suffix and the active
notes will again restart at V1.

A conditional theorem proves the implication from its displayed
hypotheses.  It does not assert that the physical regulator family supplies
those hypotheses.  Finite-regulator identities, analytic continuum cards,
and final existence statements remain separate.

\section{Audited summary of the major mathematical results}
\label{sec:n9v1-summary}

\subsection{Chiral geometry, Standard Model arithmetic, and Higgs blocks}

On an admissible finite lattice gauge sector, the archived analysis gives a
volume-uniform spectral window for the Wilson kernel under its stated
plaquette bound.  Functional calculus then gives exponential locality of
the overlap sign, Ginsparg--Wilson operator, and moving chiral projectors.
Riesz-projector and Schur--Berezin formulas retain zero and crossing modes
as finite Grassmann sectors before the complementary block is inverted.
At a simple transverse crossing, the least singular value is linear, the
determinant changes sign, and an unsaturated source limit depends on its
approach direction.

The finite determinant is a section of a determinant line.  The moving
projector gives its Berry connection and curvature; radial homotopy yields
local primitives, and a uniform finite-range gap gives exponentially local
curvature.  Local triviality does not prove global holonomy triviality.
For one Standard Model generation in the all-left-handed convention, the
finite representation arithmetic cancels the local
\(SU(3)^3\), \(SU(3)^2U(1)\), \(SU(2)^2U(1)\), \(U(1)^3\), and mixed
gravitational--\(U(1)\) anomalies, and the number of left-handed \(SU(2)\)
doublets is even.

On a Higgs chart separated from zero, the normalized doublet reduces the
electroweak group to \(U(1)_{\rm em}\), pairs the charged chiral
representations, and isolates the unmatched minimal-neutrino row.  The
improved Ginsparg--Wilson fields put kinetic and Yukawa corners into
massive-overlap blocks.  A positive mass gives a uniform inverse bound on
that chart.  The zero-Higgs region and global phase gluing remain open.

\subsection{Finite loaders, bosonic regulators, and cofinal response}

Exact relative determinant identities, finite-rank transfer, Schur
factorization, reduced pseudodeterminants, and trace-ideal factorizations
show that subtraction must be assembled before an absolute trace norm is
taken.  Positive-type compact-group plaquette weights supply finite
reflection-positive pure-gauge regulators.  Heat and symbol expansions
assign local divergences to declared counterterm owners.  Multipole
analysis of elliptic constraint substitution proves polynomial Coulomb
tails and identifies the monopole and dipole cancellations needed for
absolute summability.

Finite Fock and quasifree compressions, Nambu covariance identities,
Gibbs--Duhamel response, commutator expansions, replacement channels, and
Petz recovery maps have been derived at finite cutoff.  Conditional channel
theorems propagate boundary, recovery, metric, and refinement sources
through explicit inhomogeneous resolvents.  Summable adjacent observable
and tangent rows give a positive normalized state on a countable inductive
local algebra and a bounded derivative on the completed tangent core,
provided the stated carrier cards hold.

The weak Einstein residual is defined in a Banach quotient by allowed
finite local counterterm variations.  An absolutely summable lifted
cocycle gives a quotient residual limit.  A finite counterterm sets that
limit to zero exactly when its quotient class vanishes.  Regulator
independence still needs a separate common-refinement estimate.

\subsection{Coupled constraints, determinant, and coarea}

The latest frozen block gives a quantitative scalar--vector Einstein
constraint solve on the conformal-Killing quotient and an augmented
full-space system whose finite multipliers record compatibility defects.
Under uniform inverse and differentiability assumptions, a nonlinear
implicit-function theorem constructs a local coupled constraint chart and
computes the exact tangent multiplier ledger.

Two finite-dimensional coarea eliminations factor the augmented constraint
Jacobian and isolate the physical multiplier zero set.  Exact chart,
slow-basis, and adjacent-regulator cocycles prevent coordinates from being
mistaken for probability densities.  On a regular multiplier chart, a
uniform normal singular-value floor \(\beta>0\) gives a unique physical
graph, a coarea floor proportional to \(\beta^r\), controlled graph
transport, and a Lipschitz logarithmic density.  Summable adjacent
\(C^1\) changes give a summable coarea row.  Singular strata require a
quantitative exclusion or stratified analysis.

For the fourth-order normal operator, the finite-part heat construction
proves exact endpoint cocycles.  In dimension three its leading constraint
heat density is
\begin{equation}
 A_0(g)=C_{\rm con}\operatorname{Vol}_g(M),\qquad
 C_{\rm con}=
 {\Gamma(3/4)(17+3\sqrt3)\over64\pi^2}.               \label{eq:n9v1-Ccon}
\end{equation}
Fixed volume therefore cancels the leading relative divergence.  The odd
closed-manifold coefficients vanish by symbol parity.  The remaining
\(A_2\) density has an exact finite-jet symbol formula.  A conforming
Galerkin theorem reduces positive-time heat-trace convergence to explicit
uniform spectral and symbol hypotheses.

\subsection{The adjacent-density compiler}

After transport to a common local space, the adjacent logarithmic density
has the exact row structure
\begin{equation}
 \ell_{n,A}
 =
 \Delta S_{n,A}+\Delta D_{n,A}+\Delta C_{n,A}
 +\Delta T_{n,A}+\Delta K_{n,A},                      \label{eq:n9v1-ledger}
\end{equation}
where the rows are action, determinant, coarea, ambient transport, and
local counterterm contributions.  Field-independent constants disappear
under probability normalization.  The archived compiler proves adjacent
total-variation summability from either summable row oscillations or
summable centered \(L^2\) norms with a uniform combined exponential
moment.  One cofinal subsequence can serve countably many supports and rows.

\begin{theorem}[Audited dependency theorem]
\label{thm:n9v1-dependency}
Suppose one physical SM--Einstein regulator sequence satisfies, on one
cofinal subsequence:
\begin{enumerate}[label=\textup{(D\arabic*)},leftmargin=4.5em]
\item the coupled constraint chart and compatible chart cocycles;
\item the normal-heat convergence card, including its \(A_2\) rate and
matched slow-sector rule;
\item the regular multiplier-coarea card and a summable exceptional-set
debit;
\item summable action, ambient-transport, and counterterm rows in one of
the density compiler norms; and
\item projective consistency of all transported local marginals.
\end{enumerate}
Then the constrained marginals converge in total variation and define a
unique projective cylinder probability measure.  Closed physical support,
bounded cylinder Ward identities, and reflection-positive cylinder
inequalities valid at every regulator pass to the limit.
\end{theorem}

\begin{proof}
Assumptions (D1)--(D3) provide the chart, determinant, coarea, and
slow-sector terms in \eqref{eq:n9v1-ledger}.  Assumption (D4) places every
remaining centered row in a proved density compiler, hence
\(\sum_n\|\mu_{n+1,A}-\mu_{n,A}\|_{\rm TV}<\infty\)
for each local support \(A\).  Completeness in total variation gives the
local limits.  Assumption (D5) makes them projectively compatible.  Closed
support follows from the Portmanteau inequality; bounded Ward and
reflection expressions pass by total-variation continuity.
\end{proof}

\begin{firewall}[Scope of the dependency theorem]
\label{fire:n9v1-dependency}
The theorem does not establish (D1)--(D5) for the full physical regulator.
A cylinder measure is not a Radon measure on a chosen field topology, and
bounded Ward identities do not imply closability of unbounded generators.
\end{firewall}

\section{Open gates at the ninth-series restart}
\label{sec:n9v1-open}

\begin{theorem}[Current claim boundary]
\label{thm:n9v1-open}
The full SM--Einstein continuum remains unproved.  The main unpaid gates
are:
\begin{enumerate}[label=\textup{(G\arabic*)},leftmargin=4.5em]
\item a single physical regulator with its chiral determinant, zero-mode
patches, global determinant-line data, and reflection structure;
\item uniform verification of the coupled constraint and normal-heat cards;
\item uniform multiplier transversality or summable singular neighborhoods;
\item a summable centered adjacent action row for the interacting Einstein,
gauge, Higgs, and fermion action;
\item the remaining ambient field and base-measure transport Jacobians;
\item tightness in a declared distributional topology and regulator
independence under common refinement;
\item unbounded Ward generators, clustering, full reflection positivity,
and reconstruction or Lorentzian dynamics; and
\item positive nonlinear gravitational control of the Euclidean conformal
direction.
\end{enumerate}
\end{theorem}

\begin{proof}
Each item is an explicit hypothesis or firewall in the frozen sources.
Finite-cutoff identities and the conditional dependency theorem cannot
imply a missing item.  No full continuum-existence conclusion follows.
\end{proof}

\section{First new attack: the adjacent action row}
\label{sec:n9v1-action}

The V100 audit selected gate (G4).  The newest Yang--Mills note favors
bounded local insertions with an explicit truncation schedule, but its
operator tree theorem remains conditional and cannot be imported.  The
newest Navier--Stokes rank-one Oseen constants are kernel-specific and do
not change this acquisition order.

Let \(\nu_n\) be the transported reference probability for one support and
\begin{equation}
 u_n:=-\Delta S_{n,A}
 +\int\Delta S_{n,A}\,d\nu_n                         \label{eq:n9v1-action-center}
\end{equation}
be the centered action contribution.  We must control
\begin{equation}
 r_n={e^{u_n}\over\int e^{u_n}\,d\nu_n}.              \label{eq:n9v1-action-tilt}
\end{equation}

\begin{definition}[Martingale action card, MAC]
\label{def:n9v1-MAC}
MAC holds with budget \(V_n\) if
\[
 u_n=\sum_{j=1}^{m_n}X_{n,j},\qquad
 \mathbb E_{\nu_n}[X_{n,j}\mid\mathcal F_{n,j-1}]=0,
\]
and
\begin{equation}
 \mathbb E_{\nu_n}
 [e^{\lambda X_{n,j}}\mid\mathcal F_{n,j-1}]
 \leq e^{\lambda^2v_{n,j}/2}
 \quad (|\lambda|\leq2),\qquad
 V_n:=\sum_jv_{n,j}.                                  \label{eq:n9v1-MAC}
\end{equation}
\end{definition}

\begin{programme}[Ninth-series V2]
\label{prog:n9v1-V2}
V2 will prove from MAC that
\[
 \|r_n-1\|_{L^1(\nu_n)}
 \leq2\sqrt{2V_n}\,e^{V_n},
\]
so \(\sum_n\sqrt{V_n}<\infty\) closes the isolated action row.  It will
derive MAC from conditional bounded differences, record the exact
influence-square budget, and state the blockwise loader that later versions
must verify for the Einstein, gauge, Higgs, and fermion pieces.
\end{programme}

\begin{assessment}[Position after compact V1]
\label{ass:n9v1-status}
The large proof history is preserved in the \texttt{\_f8} archive.  The
active note now exposes the shortest rigorous route forward: determinant
and regular-coarea rows have conditional closure mechanisms, and V2 attacks
the first remaining density row without assuming its physical estimate.
\end{assessment}

\begin{firewall}[Claim boundary after ninth-series V1]
\label{fire:n9v1-boundary}
V1 is an audited handoff.  It does not assert MAC for the physical action,
close singular strata, prove tightness, or construct the full SM--Einstein
continuum.
\end{firewall}


% =============================================================================
% BEGIN APPENDED NINTH-SERIES V2 MATERIAL
% =============================================================================

\clearpage
\section{Ninth-series V2: a martingale compiler for the action row}
\label{sec:v9v2}

\subsection{Conditional moment propagation}
\label{sec:n9v2-mgf}

Fix one adjacent regulator step and suppress its index.  Let
\((\Omega,\mathcal F,\nu)\) be a probability space with filtration
\[
 \{\varnothing,\Omega\}=\mathcal F_0
 \subset\mathcal F_1\subset\cdots\subset\mathcal F_m=\mathcal F.
\]
Let \(X_j\) be \(\mathcal F_j\)-measurable and suppose
\begin{align}
 \mathbb E_\nu[X_j\mid\mathcal F_{j-1}]&=0,            \label{eq:n9v2-md}\\
 \mathbb E_\nu[e^{\lambda X_j}\mid\mathcal F_{j-1}]
 &\leq e^{\lambda^2v_j/2}
 \quad\text{a.s. for }|\lambda|\leq2.                 \label{eq:n9v2-cond-mgf}
\end{align}
Put
\begin{equation}
 u:=\sum_{j=1}^mX_j,\qquad V:=\sum_{j=1}^mv_j.         \label{eq:n9v2-uV}
\end{equation}

\begin{lemma}[Iteration of the conditional moment bound]
\label{lem:n9v2-mgf}
Under \eqref{eq:n9v2-md}--\eqref{eq:n9v2-uV},
\begin{equation}
 \mathbb E_\nu e^{\lambda u}
 \leq e^{\lambda^2V/2}\qquad(|\lambda|\leq2).          \label{eq:n9v2-mgf}
\end{equation}
In particular,
\begin{align}
 \mathbb E_\nu u&=0,                                  \label{eq:n9v2-mean}\\
 \mathbb E_\nu u^2&\leq V,                            \label{eq:n9v2-var}\\
 \mathbb E_\nu e^{2|u|}&\leq2e^{2V}.                  \label{eq:n9v2-abs-mgf}
\end{align}
\end{lemma}

\begin{proof}
For \(S_k=\sum_{j\leq k}X_j\), conditional expectation gives
\[
 \mathbb E e^{\lambda S_k}
 =
 \mathbb E\!\left[
 e^{\lambda S_{k-1}}
 \mathbb E(e^{\lambda X_k}\mid\mathcal F_{k-1})
 \right]
 \leq
 e^{\lambda^2v_k/2}\mathbb E e^{\lambda S_{k-1}}.
\]
Iteration proves \eqref{eq:n9v2-mgf}.  The martingale-difference property
gives \eqref{eq:n9v2-mean}.  Applying \eqref{eq:n9v2-mgf} at both signs,
\[
 \mathbb E\cosh(\lambda u)\leq e^{\lambda^2V/2}.
\]
Fatou's lemma applied to
\((\cosh(\lambda u)-1)/\lambda^2\) as \(\lambda\to0\) gives
\(\mathbb E u^2/2\leq V/2\), proving \eqref{eq:n9v2-var}.  Finally,
\[
 e^{2|u|}\leq e^{2u}+e^{-2u},
\]
and \eqref{eq:n9v2-mgf} at \(\lambda=\pm2\) proves
\eqref{eq:n9v2-abs-mgf}.
\end{proof}

\subsection{The normalized exponential tilt}
\label{sec:n9v2-tilt}

\begin{theorem}[MAC normalized-density bound]
\label{thm:n9v2-MAC}
Let \(u\) satisfy the martingale action card
\Cref{def:n9v1-MAC} with budget \(V\), and put
\[
 Z:=\mathbb E_\nu e^u,\qquad r:={e^u\over Z}.
\]
Then
\begin{equation}
 Z\geq1,\qquad
 \|r-1\|_{L^1(\nu)}
 \leq2\sqrt{2V}\,e^V.                                 \label{eq:n9v2-L1}
\end{equation}
Consequently, for a regulator sequence satisfying MAC,
\begin{equation}
 \sum_n\sqrt{V_n}<\infty
 \quad\Longrightarrow\quad
 \sum_n\|r_n-1\|_{L^1(\nu_n)}<\infty.                 \label{eq:n9v2-sum}
\end{equation}
\end{theorem}

\begin{proof}
Since \(\mathbb E_\nu u=0\), Jensen's inequality gives \(Z\geq1\).
Moreover,
\[
 \mathbb E|e^u-Z|
 \leq
 \mathbb E|e^u-1|+|Z-1|
 \leq2\mathbb E|e^u-1|.
\]
The elementary bound
\(|e^x-1|\leq |x|e^{|x|}\), Cauchy--Schwarz, and
\Cref{lem:n9v2-mgf} yield
\[
 \mathbb E|e^u-1|
 \leq
 (\mathbb E u^2)^{1/2}
 (\mathbb E e^{2|u|})^{1/2}
 \leq\sqrt{2V}\,e^V.
\]
Division by \(Z\geq1\) proves \eqref{eq:n9v2-L1}.  Under the hypothesis in
\eqref{eq:n9v2-sum}, \(V_n\to0\), so \(\sup_ne^{V_n}<\infty\) after
absorbing finitely many initial indices.  Summation proves the conclusion.
\end{proof}

\begin{corollary}[Total-variation debit of the isolated action row]
\label{cor:n9v2-TV}
If \(d\mu=r\,d\nu\), then
\begin{equation}
 \|\mu-\nu\|_{\rm TV}
 \leq\sqrt{2V}\,e^V.                                  \label{eq:n9v2-TV}
\end{equation}
\end{corollary}

\begin{proof}
For probabilities with a density,
\(\|\mu-\nu\|_{\rm TV}=\frac12\|r-1\|_{L^1(\nu)}\).
Apply \Cref{thm:n9v2-MAC}.
\end{proof}

\subsection{Conditional bounded differences}
\label{sec:n9v2-bd}

The preceding theorem needs a concrete loader.  The next lemma requires no
independence once a martingale difference has been constructed.

\begin{lemma}[Conditional Hoeffding loader]
\label{lem:n9v2-Hoeffding}
Suppose \(X_j\) satisfies \eqref{eq:n9v2-md} and, conditional on
\(\mathcal F_{j-1}\), lies almost surely in an interval of length \(c_j\),
where \(c_j\) is deterministic.  Then
\begin{equation}
 \mathbb E[e^{\lambda X_j}\mid\mathcal F_{j-1}]
 \leq\exp\!\left({\lambda^2c_j^2\over8}\right)
 \qquad(\lambda\in\mathbb R).                         \label{eq:n9v2-Hoeffding}
\end{equation}
Thus MAC holds with
\begin{equation}
 v_j={c_j^2\over4},\qquad
 V={1\over4}\sum_{j=1}^mc_j^2.                        \label{eq:n9v2-influence}
\end{equation}
\end{lemma}

\begin{proof}
Condition on \(\mathcal F_{j-1}\).  Hoeffding's convexity argument for a
mean-zero random variable in an interval of length \(c_j\) gives
\eqref{eq:n9v2-Hoeffding}; it is pointwise in the conditioning variable.
Comparing its exponent with \(\lambda^2v_j/2\) gives
\eqref{eq:n9v2-influence}.
\end{proof}

\begin{proposition}[Independent-innovation bounded differences]
\label{prop:n9v2-product}
Let \(Y_1,\ldots,Y_m\) be independent under \(\nu\), let
\(\mathcal F_j=\sigma(Y_1,\ldots,Y_j)\), and let
\(F=F(Y_1,\ldots,Y_m)\) be integrable.  Suppose replacement of coordinate
\(j\) changes \(F\) by at most \(c_j\):
\begin{equation}
 \sup_{y,y':\,y_i=y_i'\ (i\ne j)}
 |F(y)-F(y')|\leq c_j.                                \label{eq:n9v2-coordinate}
\end{equation}
Then \(u:=F-\mathbb E_\nu F\) satisfies MAC with the budget in
\eqref{eq:n9v2-influence}.
\end{proposition}

\begin{proof}
Let \(M_j=\mathbb E[F\mid\mathcal F_j]\) and
\(X_j=M_j-M_{j-1}\).  Then \(u=\sum_jX_j\) and the \(X_j\) are martingale
differences.  Conditional on the past, the function
\[
 y_j\longmapsto
 \mathbb E[F\mid Y_1,\ldots,Y_j=y_j]
\]
has oscillation at most \(c_j\): use the same law for the independent
future coordinates and integrate \eqref{eq:n9v2-coordinate}.  Subtracting
its conditional mean does not change the interval length.  Apply
\Cref{lem:n9v2-Hoeffding}.
\end{proof}

\begin{firewall}[Dependent interacting references]
\label{fire:n9v2-dependent}
The coordinate proof uses independent innovations, not independent terms
in the action.  For a dependent Gibbs or constrained reference, one must
construct an innovation representation or prove conditional ranges for its
Doob martingale directly.  Ordinary pointwise locality of the action does
not supply that property.
\end{firewall}

\subsection{Local incidence and sector allocation}
\label{sec:n9v2-incidence}

Let \(E\) be a finite family of local action terms.  Write
\begin{equation}
 F(y)=\sum_{e\in E}a_e\Psi_e(y_e),                     \label{eq:n9v2-local-action}
\end{equation}
where \(y_e\) is the set of coordinates used by the term.  Define the
single-coordinate oscillation
\[
 b_{e,j}:=
 \sup_{y,y':\,y_i=y_i'\ (i\ne j)}
 |\Psi_e(y_e)-\Psi_e(y'_e)|
\]
and set \(b_{e,j}=0\) when \(j\notin e\).

\begin{proposition}[Exact local influence ledger]
\label{prop:n9v2-incidence}
The coordinate influence of \eqref{eq:n9v2-local-action} obeys
\begin{equation}
 c_j\leq\sum_{e\ni j}|a_e|b_{e,j}.                    \label{eq:n9v2-cj}
\end{equation}
If every \(e\) contains at most \(q\) coordinates, every coordinate belongs
to at most \(\Delta\) terms, and \(b_{e,j}\leq b_e\), then
\begin{equation}
 V
 \leq{\Delta q\over4}\sum_{e\in E}|a_e|^2b_e^2.       \label{eq:n9v2-incidence-budget}
\end{equation}
\end{proposition}

\begin{proof}
Only terms incident to \(j\) change when coordinate \(j\) is replaced, so
the triangle inequality gives \eqref{eq:n9v2-cj}.  Cauchy--Schwarz at each
coordinate gives
\[
 c_j^2
 \leq
 \Delta\sum_{e\ni j}|a_e|^2b_e^2.
\]
Summing in \(j\), each \(e\) is counted at most \(q\) times.  Insert the
result into \eqref{eq:n9v2-influence}.
\end{proof}

\begin{corollary}[Four-sector influence allocation]
\label{cor:n9v2-sectors}
Suppose
\[
 F=F^{\rm Ein}+F^{\rm gauge}+F^{\rm Higgs}+F^{\rm ferm}
\]
and coordinate \(j\) has sector influences \(c_j^{(s)}\).  For arbitrary
weights \(\alpha_s>0\) with \(\sum_s\alpha_s=1\),
\begin{equation}
 V
 \leq{1\over4}
 \sum_s{1\over\alpha_s}\sum_j(c_j^{(s)})^2.           \label{eq:n9v2-sector-budget}
\end{equation}
\end{corollary}

\begin{proof}
The total influence is at most \(\sum_sc_j^{(s)}\).  Weighted
Cauchy--Schwarz gives
\[
 \left(\sum_sc_j^{(s)}\right)^2
 \leq\sum_s{(c_j^{(s)})^2\over\alpha_s}.
\]
Sum in \(j\) and use \eqref{eq:n9v2-influence}.
\end{proof}

\subsection{Truncation and the exceptional tail}
\label{sec:n9v2-truncation}

Bounded differences are inappropriate for an untruncated Higgs amplitude,
metric coefficient, or Gaussian auxiliary field.  The truncation error
must remain visible in the probability norm.

\begin{proposition}[Normalized truncation debit]
\label{prop:n9v2-truncation}
Let \(F,\widetilde F\) be measurable and let \(G\) be an event on which
\(F=\widetilde F\).  Put
\[
 Z=\int e^F\,d\nu,\qquad
 \widetilde Z=\int e^{\widetilde F}\,d\nu,
\]
and suppose \(Z,\widetilde Z>0\).  For the normalized tilted probabilities
\[
 d\mu={e^F\over Z}\,d\nu,\qquad
 d\widetilde\mu={e^{\widetilde F}\over\widetilde Z}\,d\nu,
\]
one has
\begin{equation}
 \|\mu-\widetilde\mu\|_{\rm TV}
 \leq
 {1\over\min(Z,\widetilde Z)}
 \int_{G^c}(e^F+e^{\widetilde F})\,d\nu.              \label{eq:n9v2-truncation}
\end{equation}
\end{proposition}

\begin{proof}
Let \(h=e^F\), \(\widetilde h=e^{\widetilde F}\).  Then
\[
 \left\|{h\over Z}-{\widetilde h\over\widetilde Z}\right\|_1
 \leq
 {\|h-\widetilde h\|_1\over Z}
 +{|Z-\widetilde Z|\over Z}
 \leq{2\|h-\widetilde h\|_1\over Z}.
\]
Interchanging the two tilts replaces \(Z\) by
\(\min(Z,\widetilde Z)\).  Since \(h=\widetilde h\) on \(G\),
\[
 \|h-\widetilde h\|_1
 \leq\int_{G^c}(h+\widetilde h)\,d\nu.
\]
Divide the \(L^1\) density bound by two to obtain total variation.
\end{proof}

\begin{corollary}[Summable truncated action route]
\label{cor:n9v2-truncated}
Suppose the truncated centered action at step \(n\) satisfies MAC with
budget \(V_n\), and the right side of
\eqref{eq:n9v2-truncation} is \(\delta_n\).  If
\begin{equation}
 \sum_n(\sqrt{V_n}+\delta_n)<\infty,                   \label{eq:n9v2-truncated-sum}
\end{equation}
then the full isolated action tilts have a summable total-variation debit
relative to the reference sequence.
\end{corollary}

\begin{proof}
Use the triangle inequality, \Cref{cor:n9v2-TV}, and
\Cref{prop:n9v2-truncation}, then sum.
\end{proof}

\subsection{The action card delivered to the density compiler}
\label{sec:n9v2-card}

\begin{definition}[Summable adjacent martingale action card, SAMAC]
\label{def:n9v2-SAMAC}
For every support in the countable local core, SAMAC records:
\begin{enumerate}[label=\textup{(A\arabic*)},leftmargin=4em]
\item the common transported reference probability and the centered
adjacent action increment;
\item an innovation filtration or a direct Doob martingale;
\item conditional range bounds \(c_{n,j}\), split by physical sector;
\item the budget \(V_n=\frac14\sum_jc_{n,j}^2\);
\item a truncation event and the normalized tail debit \(\delta_n\);
\item the summability condition
\(\sum_n(\sqrt{V_n}+\delta_n)<\infty\); and
\item compatibility of the action factorization with the determinant,
coarea, transport, and counterterm rows in the exact density ledger.
\end{enumerate}
\end{definition}

\begin{theorem}[Conditional closure of the isolated action row]
\label{thm:n9v2-closure}
SAMAC implies a summable adjacent total-variation debit for the isolated
action tilt.  If the other density rows are composed through one of the
proved joint compilers and item (A7) holds, the action contribution is no
longer an unpaid row of the audited dependency theorem
\Cref{thm:n9v1-dependency}.
\end{theorem}

\begin{proof}
The isolated statement is
\Cref{thm:n9v2-MAC,cor:n9v2-truncated}.  Item (A7) is needed because
separately normalized row estimates cannot be multiplied without tracking
the common reference and the normalization cocycle.  With that item, the
joint density compiler applies to the same exact logarithmic ledger
\eqref{eq:n9v1-ledger}.
\end{proof}

\begin{firewall}[The physical action estimate remains open]
\label{fire:n9v2-physical}
V2 proves the analytic compiler, not its physical loader.  In particular,
it does not yet bound the Einstein conformal mode, construct independent
innovations for a constrained interacting measure, control an unbounded
Higgs field without the tail debit, or show the required square-influence
summability under refinement.  Fermionic determinant terms already owned
by the determinant row must not be charged a second time as an action term.
\end{firewall}

\begin{assessment}[Progress at ninth-series V2]
\label{ass:n9v2-status}
The action row now has a quantitative target:
\[
 \sum_n
 \left[
 \left(\sum_jc_{n,j}^2\right)^{1/2}
 +\delta_n
 \right]<\infty.
\]
The compiler loses a square root of the quadratic influence budget, but it
does not lose the number of local terms linearly.  This is the useful
content imported from the audit direction, now proved directly in the
probability norm required by the constrained-density programme.
\end{assessment}

\begin{programme}[Ninth-series V3: gauge--Higgs influence test]
\label{prog:n9v2-V3}
V3 will apply the local incidence ledger to an explicit compact gauge
plaquette increment and a radially truncated Higgs hopping/potential
increment.  It will separate support growth from coefficient decay,
optimize the sector weights in \eqref{eq:n9v2-sector-budget}, and state the
tail estimate required to remove the Higgs truncation.  If the needed rate
is incompatible with the regulator schedule, the programme will move
upstream to a block or conditional-log-Sobolev loader.
\end{programme}

\begin{firewall}[Claim boundary after ninth-series V2]
\label{fire:n9v2-boundary}
V2 does not verify SAMAC for the physical SM--Einstein sequence, close the
full density, prove tightness, or construct the full continuum.
\end{firewall}

\begin{theorem}[Ninth-series V2 result]
\label{thm:n9v2-result}
Conditional sub-Gaussian martingale differences with total variance proxy
\(V_n\) give the exact normalized-density bound
\[
 \|r_n-1\|_1\leq2\sqrt{2V_n}\,e^{V_n}.
\]
Independent-coordinate bounded differences load this card with
\(V_n=\frac14\sum_jc_{n,j}^2\); local incidence, sector allocation, and
truncation have the explicit debits
\eqref{eq:n9v2-incidence-budget},
\eqref{eq:n9v2-sector-budget}, and
\eqref{eq:n9v2-truncation}.
\end{theorem}

\begin{proof}
Use \Cref{thm:n9v2-MAC,prop:n9v2-product,
prop:n9v2-incidence,cor:n9v2-sectors,
prop:n9v2-truncation}.
\end{proof}

\paragraph{Parent-version record.}
V2 was formed from
\texttt{Renewal\_SM\_Einstein\_Continuum\_Research\_Notes\_V1.tex}.
The V1 parent had \(310\) logical source lines, \(13\,908\) bytes, and
SHA--256
\[
 \texttt{14C8C4CE18BCB6D74D3077B69DDDF01CECFAA4C1A91896235B74218B10E6411B}.
\]
The next compulsory companion audit is V5.

% =============================================================================
% END APPENDED NINTH-SERIES V2 MATERIAL
% =============================================================================


% =============================================================================
% BEGIN APPENDED NINTH-SERIES V3 MATERIAL
% =============================================================================

\clearpage
\section{Ninth-series V3: the compact gauge--Higgs influence test}
\label{sec:v9v3}

\subsection{Finite local complex and adjacent action}
\label{sec:n9v3-setting}

Fix a transported local support \(A\).  At regulator \(n\), let
\(\mathcal V_n(A)\), \(\mathcal E_n(A)\), and \(\mathcal P_n(A)\) be its
vertices, oriented links, and plaquettes.  Let \(G\) be compact and let
\(\rho:G\to U(d_\rho)\) be unitary.  The variables are
\[
 U_e\in G,\qquad \phi_x\in\mathbb C^{d_H}.
\]
This version studies the adjacent gauge--Higgs logarithmic action on the
truncated event
\begin{equation}
 \mathcal G_n(R):=\{|\phi_x|\leq R\text{ for every }x\in\mathcal V_n(A)\}.
 \label{eq:n9v3-good}
\end{equation}
After all common-cell transports have been performed, write
\begin{align}
 F_{n,R}(U,\phi)
 &:=
 \sum_{p\in\mathcal P_n(A)}
 a_{n,p}\left[
 1-\frac1{d_\rho}\operatorname{Re}\Tr\rho(U_{\partial p})
 \right]                                             \label{eq:n9v3-Fg}\\
 &\quad+
 \sum_{e=(x,y)\in\mathcal E_n(A)}
 h_{n,e}\,|\phi_y-\rho(U_e)\phi_x|^2                 \label{eq:n9v3-Fhop}\\
 &\quad+
 \sum_{x\in\mathcal V_n(A)}
 \left[
 m_{n,x}|\phi_x|^2+\lambda_{n,x}|\phi_x|^4
 \right].                                            \label{eq:n9v3-Fpot}
\end{align}
The coefficients are adjacent differences, not the full bare couplings.
Their signs are unrestricted in the estimates below.

Let \(q_g\) bound the number of links in a plaquette boundary,
\(\Delta_g\) the number of plaquettes incident to one link, and
\(\Delta_H\) the number of hopping links incident to one vertex.

\subsection{Exact single-coordinate oscillations}
\label{sec:n9v3-osc}

\begin{lemma}[Compact plaquette oscillation]
\label{lem:n9v3-plaquette}
For
\[
 P_p(U):=
 1-\frac1{d_\rho}\operatorname{Re}\Tr\rho(U_{\partial p}),
\]
one has \(0\leq P_p\leq2\).  Consequently replacement of one link in
\(\partial p\) changes \(P_p\) by at most \(2\).
\end{lemma}

\begin{proof}
Every eigenvalue of the unitary matrix \(\rho(U_{\partial p})\) lies on the
unit circle, so
\(-d_\rho\leq\operatorname{Re}\Tr\rho(U_{\partial p})\leq d_\rho\).
The range of \(P_p\) is therefore contained in \([0,2]\), whose diameter is
\(2\).
\end{proof}

\begin{lemma}[Truncated hopping and potential oscillations]
\label{lem:n9v3-Higgs}
On \(\mathcal G_n(R)\), the hopping term
\[
 H_e(U,\phi):=|\phi_y-\rho(U_e)\phi_x|^2
\]
lies in \([0,4R^2]\).  Replacement of \(U_e\), \(\phi_x\), or \(\phi_y\)
therefore changes it by at most \(4R^2\).  The potential
\[
 W_x(\phi_x):=m_{n,x}|\phi_x|^2+\lambda_{n,x}|\phi_x|^4
\]
has oscillation at most
\begin{equation}
 2w_{n,x}(R),\qquad
 w_{n,x}(R):=|m_{n,x}|R^2+|\lambda_{n,x}|R^4.         \label{eq:n9v3-w}
\end{equation}
\end{lemma}

\begin{proof}
Unitarity and the triangle inequality give
\[
 |\phi_y-\rho(U_e)\phi_x|\leq|\phi_y|+|\phi_x|\leq2R.
\]
This proves the hopping range.  Also
\(|W_x(\phi_x)|\leq w_{n,x}(R)\); the diameter of a subset of
\([-w_{n,x}(R),w_{n,x}(R)]\) is at most \(2w_{n,x}(R)\).
\end{proof}

\begin{proposition}[Exact gauge--Higgs influence ledger]
\label{prop:n9v3-exact}
Treat every \(U_e\) and every \(\phi_x\) as one coordinate block.  On
\(\mathcal G_n(R)\), valid coordinate influences for \(F_{n,R}\) are
\begin{align}
 c^U_{n,e}(R)
 &:=
 2\sum_{p\ni e}|a_{n,p}|+4R^2|h_{n,e}|,              \label{eq:n9v3-cU}\\
 c^\phi_{n,x}(R)
 &:=
 4R^2\sum_{e\ni x}|h_{n,e}|+2w_{n,x}(R).             \label{eq:n9v3-cphi}
\end{align}
For a product reference on the coordinate blocks, the centered truncated
action satisfies MAC with
\begin{equation}
 V_{n,R}^{GH}
 :=
 {1\over4}\left[
 \sum_e(c^U_{n,e}(R))^2+
 \sum_x(c^\phi_{n,x}(R))^2
 \right].                                             \label{eq:n9v3-Vexact}
\end{equation}
\end{proposition}

\begin{proof}
Changing \(U_e\) affects exactly the plaquettes incident to \(e\) and the
hopping term on \(e\).  Apply
\Cref{lem:n9v3-plaquette,lem:n9v3-Higgs} and the triangle inequality to
obtain \eqref{eq:n9v3-cU}.  Changing \(\phi_x\) affects the hopping terms
incident to \(x\) and the potential at \(x\), giving
\eqref{eq:n9v3-cphi}.  The product bounded-difference loader
\Cref{prop:n9v2-product} gives \eqref{eq:n9v3-Vexact}.
\end{proof}

\subsection{A coarse incidence budget}
\label{sec:n9v3-budget}

\begin{theorem}[Gauge--Higgs quadratic budget]
\label{thm:n9v3-budget}
The exact budget satisfies
\begin{align}
 V_{n,R}^{GH}
 &\leq
 2\Delta_gq_g\sum_p|a_{n,p}|^2                       \label{eq:n9v3-budget}\\
 &\quad+
 (8+16\Delta_H)R^4\sum_e|h_{n,e}|^2
 +2\sum_xw_{n,x}(R)^2. \notag
\end{align}
Consequently the isolated truncated gauge--Higgs tilt obeys
\begin{equation}
 \|r_{n,R}^{GH}-1\|_1
 \leq2\sqrt{2V_{n,R}^{GH}}\,e^{V_{n,R}^{GH}}.         \label{eq:n9v3-L1}
\end{equation}
\end{theorem}

\begin{proof}
Use \((s+t)^2\leq2s^2+2t^2\) in \eqref{eq:n9v3-cU}:
\[
 \sum_e(c^U_{n,e})^2
 \leq
 8\sum_e\left(\sum_{p\ni e}|a_{n,p}|\right)^2
 +32R^4\sum_e|h_{n,e}|^2.
\]
At most \(\Delta_g\) plaquettes meet one link, so Cauchy--Schwarz and then
double counting give
\[
 \sum_e\left(\sum_{p\ni e}|a_{n,p}|\right)^2
 \leq\Delta_gq_g\sum_p|a_{n,p}|^2.
\]
Similarly,
\[
 \sum_x(c^\phi_{n,x})^2
 \leq
 32R^4\sum_x\left(\sum_{e\ni x}|h_{n,e}|\right)^2
 +8\sum_xw_{n,x}(R)^2.
\]
At most \(\Delta_H\) links meet \(x\), and each link has two endpoints, so
\[
 \sum_x\left(\sum_{e\ni x}|h_{n,e}|\right)^2
 \leq2\Delta_H\sum_e|h_{n,e}|^2.
\]
Add the estimates and divide by four as in
\eqref{eq:n9v3-Vexact}.  Equation \eqref{eq:n9v3-L1} is
\Cref{thm:n9v2-MAC}.
\end{proof}

\begin{corollary}[Uniform-coefficient form]
\label{cor:n9v3-uniform}
Let
\[
 N_{p,n}=|\mathcal P_n(A)|,\quad
 N_{e,n}=|\mathcal E_n(A)|,\quad
 N_{v,n}=|\mathcal V_n(A)|
\]
and suppose
\[
 |a_{n,p}|\leq A_n,\quad |h_{n,e}|\leq H_n,\quad
 |m_{n,x}|\leq M_n,\quad|\lambda_{n,x}|\leq L_n.
\]
Then
\begin{align}
 \sqrt{V_{n,R}^{GH}}
 &\leq
 C_g\sqrt{N_{p,n}}A_n
 +C_hR^2\sqrt{N_{e,n}}H_n                             \label{eq:n9v3-uniform}\\
 &\quad+
 \sqrt{2N_{v,n}}\,(M_nR^2+L_nR^4),\notag
\end{align}
where one may take
\[
 C_g=\sqrt{2\Delta_gq_g},\qquad
 C_h=\sqrt{8+16\Delta_H}.
\]
\end{corollary}

\begin{proof}
Apply \(\sqrt{x+y+z}\leq\sqrt x+\sqrt y+\sqrt z\) to
\eqref{eq:n9v3-budget} and insert the uniform coefficient bounds.
\end{proof}

\subsection{Mesh and truncation schedules}
\label{sec:n9v3-schedule}

Let \(\varepsilon_n\downarrow0\) be a mesh scale in a fixed physical
\(D\)-dimensional support.  Assume the bounded-incidence complex satisfies
\begin{equation}
 N_{p,n}+N_{e,n}+N_{v,n}\leq C_A\varepsilon_n^{-D}.
 \label{eq:n9v3-count}
\end{equation}

\begin{corollary}[Power-counted sufficient schedule]
\label{cor:n9v3-power}
Under \eqref{eq:n9v3-count}, a sufficient condition for summability of the
truncated action row is
\begin{equation}
 \sum_n\varepsilon_n^{-D/2}
 \left[
 A_n+R_n^2H_n+M_nR_n^2+L_nR_n^4
 \right]<\infty.                                     \label{eq:n9v3-power}
\end{equation}
\end{corollary}

\begin{proof}
Use \Cref{cor:n9v3-uniform} and absorb the fixed incidence and support
constants.  The displayed sum implies
\(\sum_n\sqrt{V_{n,R_n}^{GH}}<\infty\), so
\Cref{thm:n9v2-MAC} applies.
\end{proof}

\begin{remark}[Cell-volume interpretation]
\label{rem:n9v3-cell}
If an adjacent local coefficient has the form
\(A_n=\eta_n\varepsilon_n^D\), then its plaquette contribution to
\eqref{eq:n9v3-power} is
\(\eta_n\varepsilon_n^{D/2}\).  The quadratic compiler gains the square
root of the number of cells compared with a linear oscillation sum.  This
calculation is not a claim that the actual transported action has that
coefficient form.
\end{remark}

\begin{proposition}[A normalized weighted-tail schedule]
\label{prop:n9v3-tail}
Assume the normalized truncation debit of
\Cref{prop:n9v2-truncation} satisfies, for constants \(C,\kappa,q>0\),
\begin{equation}
 \delta_n(R)
 \leq C N_{v,n}e^{-\kappa R^q}.                       \label{eq:n9v3-tail}
\end{equation}
If
\begin{equation}
 R_n^q\geq {1\over\kappa}
 \left[
 \log(CN_{v,n})+2\log(n+1)
 \right],                                             \label{eq:n9v3-R}
\end{equation}
then \(\delta_n(R_n)\leq(n+1)^{-2}\), and hence
\(\sum_n\delta_n(R_n)<\infty\).
\end{proposition}

\begin{proof}
Substitute \eqref{eq:n9v3-R} into \eqref{eq:n9v3-tail}.
\end{proof}

\begin{corollary}[Geometric mesh costs only logarithmic truncation growth]
\label{cor:n9v3-geometric}
If \(\varepsilon_n=2^{-n}\) and
\(N_{v,n}\leq C_A2^{Dn}\), the choice in
\eqref{eq:n9v3-R} has \(R_n=O(n^{1/q})\).  Thus the Higgs terms in
\eqref{eq:n9v3-power} pay only the explicit polynomial factors
\(n^{2/q}\) and \(n^{4/q}\).
\end{corollary}

\begin{proof}
The logarithm of \(CN_{v,n}\) is \(O(n)\).  Take the \(q\)-th root in
\eqref{eq:n9v3-R}.
\end{proof}

\subsection{Where the product loader fails}
\label{sec:n9v3-dependence}

\begin{proposition}[Constraint conditioning destroys the product premise]
\label{prop:n9v3-conditioning}
Let \(\nu=\bigotimes_j\nu_j\) be a product reference and let a nonconstant
constraint map \(C(y_1,\ldots,y_m)\) depend on at least two coordinate
blocks.  In general, the conditional or coarea-reduced measure on
\(C^{-1}(0)\) is not a product measure in those blocks.  Therefore
\Cref{prop:n9v2-product} cannot be applied to the physical constrained
measure solely because the unreduced reference is a product.
\end{proposition}

\begin{proof}
Consider already two independent real coordinates with a positive product
density and the constraint \(C(y_1,y_2)=y_1+y_2\).  On \(C^{-1}(0)\) one
has \(y_2=-y_1\) almost surely.  If the two coordinates remained
independent and neither were deterministic, then the variance of
\(y_1+y_2\) would be the positive sum of their variances, a contradiction.
Thus a multi-coordinate constraint can create dependence.  Since the
Einstein and gauge constraints couple neighboring variables, product
structure after coarea reduction requires a separate theorem.
\end{proof}

\begin{firewall}[A product calculation is not the physical action card]
\label{fire:n9v3-product}
The bounds \eqref{eq:n9v3-Vexact}--\eqref{eq:n9v3-power} are exact for a
product innovation reference, and remain valid for any dependent reference
whose Doob ranges obey the same \(c_j\).  V3 has not proved those ranges for
the coarea-reduced SM--Einstein measure.  It also assumes the weighted tail
\eqref{eq:n9v3-tail}; an unweighted Higgs tail is insufficient.
\end{firewall}

\subsection{Gauge--Higgs action card}
\label{sec:n9v3-card}

\begin{theorem}[Conditional compact gauge--Higgs action closure]
\label{thm:n9v3-closure}
For every local support, suppose:
\begin{enumerate}[label=\textup{(GH\arabic*)},leftmargin=4.8em]
\item the transported adjacent gauge--Higgs action has the form
\eqref{eq:n9v3-Fg}--\eqref{eq:n9v3-Fpot};
\item its reference admits a Doob martingale with conditional ranges no
larger than \eqref{eq:n9v3-cU}--\eqref{eq:n9v3-cphi};
\item the coefficient and truncation schedules satisfy
\eqref{eq:n9v3-power} and \eqref{eq:n9v3-tail}--\eqref{eq:n9v3-R}; and
\item the exact density ledger assigns no determinant term twice.
\end{enumerate}
Then the gauge--Higgs part of SAMAC holds and its isolated adjacent
total-variation debits are summable.
\end{theorem}

\begin{proof}
Items (GH1)--(GH2) and \Cref{thm:n9v3-budget} give the MAC budget.
Item (GH3), \Cref{cor:n9v3-power,prop:n9v3-tail} give
\(\sum_n(\sqrt{V_{n,R_n}^{GH}}+\delta_n)<\infty\).
Apply \Cref{cor:n9v2-truncated}.  Item (GH4) is the ownership condition
needed for insertion into the full density ledger.
\end{proof}

\begin{assessment}[Progress at ninth-series V3]
\label{ass:n9v3-status}
The compact gauge and truncated Higgs action now have explicit local
influences, a quadratic incidence budget, and a mesh/tail schedule.  The
calculation shows that the number of cells enters through its square root.
The first new obstruction is also exact: the physical coarea reduction
creates dependence, so the independent-coordinate proof does not by itself
load SAMAC.
\end{assessment}

\begin{programme}[Ninth-series V4: dependent-reference loader]
\label{prog:n9v3-V4}
V4 will move upstream and derive a Doob-range bound from a quantitative
interdependence matrix.  It will prove the resolvent influence
\((I-C)^{-1}c\), insert it into the MAC budget, and identify the precise
loss as the Dobrushin margin tends to zero.  This is a loader theorem; the
physical constraint measure must still supply the matrix.
\end{programme}

\begin{firewall}[Claim boundary after ninth-series V3]
\label{fire:n9v3-boundary}
V3 does not prove the dependent-reference loader, the physical tail bound,
the Einstein action row, full density convergence, tightness, or the
SM--Einstein continuum.
\end{firewall}

\begin{theorem}[Ninth-series V3 result]
\label{thm:n9v3-result}
The compact gauge--Higgs adjacent action has the exact influence budget
\eqref{eq:n9v3-Vexact} and the incidence bound
\eqref{eq:n9v3-budget}.  The power-counted schedule
\eqref{eq:n9v3-power}, together with the normalized tail schedule
\eqref{eq:n9v3-R}, conditionally closes this sector.  Constraint coarea
reduction prevents an automatic product-measure application.
\end{theorem}

\begin{proof}
Use \Cref{prop:n9v3-exact,thm:n9v3-budget,
cor:n9v3-power,prop:n9v3-tail,prop:n9v3-conditioning,
thm:n9v3-closure}.
\end{proof}

\paragraph{Parent-version record.}
V3 was formed from
\texttt{Renewal\_SM\_Einstein\_Continuum\_Research\_Notes\_V2.tex}.
The V2 parent had \(761\) logical source lines, \(29\,203\) bytes, and
SHA--256
\[
 \texttt{1F009C8E05CB625A12018BA54BFA826F3D9A8E8D9952FC4AFFFE09D495F486E9}.
\]
The next compulsory companion audit is V5.

% =============================================================================
% END APPENDED NINTH-SERIES V3 MATERIAL
% =============================================================================


% =============================================================================
% BEGIN APPENDED NINTH-SERIES V4 MATERIAL
% =============================================================================

\clearpage
\section{Ninth-series V4: a dependent-reference Dobrushin loader}
\label{sec:v9v4}

\subsection{Coordinate oscillations and interdependence}
\label{sec:n9v4-setting}

Let
\[
 \Omega=\prod_{i=1}^m\Omega_i
\]
be a finite product of standard Borel spaces and let \(\nu\) be a
probability admitting regular one-site conditional kernels
\(\gamma_i(\,\cdot\,|x_{-i})\).  For a bounded measurable \(f\), define
\begin{equation}
 \delta_i(f):=
 \sup_{x,y:\,x_k=y_k\ (k\ne i)}|f(x)-f(y)|.           \label{eq:n9v4-delta}
\end{equation}

\begin{definition}[Dobrushin interdependence matrix]
\label{def:n9v4-C}
A nonnegative matrix \(C=(C_{ij})\), with \(C_{ii}=0\), is an
interdependence matrix for \(\nu\) if
\begin{equation}
 \left\|
 \gamma_i(\,\cdot\,|x_{-i})
 -\gamma_i(\,\cdot\,|y_{-i})
 \right\|_{\rm TV}
 \leq C_{ij}                                           \label{eq:n9v4-C}
\end{equation}
whenever \(x\) and \(y\) differ only at coordinate \(j\).
\end{definition}

The total-variation convention is
\(\|\mu-\eta\|_{\rm TV}=\sup_B|\mu(B)-\eta(B)|\).  Thus two laws admit a
maximal coupling whose disagreement probability is their total variation.

\subsection{Comparison under one boundary change}
\label{sec:n9v4-comparison}

\begin{lemma}[Finite Dobrushin comparison]
\label{lem:n9v4-comparison}
Assume the spectral radius of \(C\) is less than one and put
\begin{equation}
 D:=(I-C)^{-1}=\sum_{k=0}^{\infty}C^k.                \label{eq:n9v4-D}
\end{equation}
Fix a set of coordinates as boundary data, and consider two conditional
laws of the remaining coordinates whose boundary values differ only at
coordinate \(j\).  If \(f\) has finite coordinate oscillations
\(c_i=\delta_i(f)\), then
\begin{equation}
 |\mathbb E f-\mathbb E'f|
 \leq\sum_{i=1}^m c_iD_{ij}.                          \label{eq:n9v4-comparison}
\end{equation}
The same bound holds when further coordinates, equal under the two laws,
are fixed.
\end{lemma}

\begin{proof}
First suppose all spaces are finite.  Couple the two single-site heat-bath
chains by using maximal couplings at each update.  The differing boundary
coordinate \(j\) is held fixed.  If \(p_i\) is the stationary probability
that the two copies disagree at a free coordinate \(i\), the definition
\eqref{eq:n9v4-C} and successive replacement of the disagreeing boundary
coordinates give
\[
 p_i\leq C_{ij}+\sum_{k\ {\rm free}}C_{ik}p_k.
\]
Iteration from the zero vector, or multiplication by the nonnegative
resolvent of the free-coordinate submatrix, gives
\[
 p_i\leq\sum_{r\ {\rm free}}D_{ir}C_{rj}
 =D_{ij}-\mathbf1_{\{i=j\}}.
\]
For any coupling,
\[
 |f(X)-f(X')|
 \leq\sum_i c_i\mathbf1_{\{X_i\ne X_i'\}}.
\]
The boundary coordinate itself contributes \(c_j\); expectation and the
preceding disagreement bound give \eqref{eq:n9v4-comparison}.

For standard Borel spaces, replace each maximal coupling by a measurable
maximal coupling of the regular conditional kernels and apply the same
finite-coordinate heat-bath construction.  Cesaro averages of its coupled
laws have subsequential limits because only finitely many coordinates are
present; equivalently one may apply the standard finite-volume coupling
comparison argument to finite measurable partitions and then refine the
partitions.  The inequalities are preserved under the limit.  Fixing equal
additional boundary coordinates deletes rows and columns of \(C\); the
corresponding nonnegative resolvent is entrywise bounded by \(D\).
\end{proof}

\begin{remark}[The comparison premise is substantive]
\label{rem:n9v4-existence}
For noncompact coordinate spaces, the limiting-coupling step requires the
conditional laws to be tight.  This is automatic for a fixed finite
probability system, but uniformity in the regulator is a separate estimate.
One may instead include \eqref{eq:n9v4-comparison} itself as the loader
card; all later conclusions then use only that inequality.
\end{remark}

\subsection{Doob ranges from the resolvent influence}
\label{sec:n9v4-Doob}

Reveal the coordinate blocks in the order \(1,\ldots,m\), put
\(\mathcal F_j=\sigma(Y_1,\ldots,Y_j)\), and define
\[
 M_j:=\mathbb E_\nu[f\mid\mathcal F_j],\qquad
 X_j:=M_j-M_{j-1}.
\]

\begin{theorem}[Dobrushin--Doob range theorem]
\label{thm:n9v4-range}
Under the hypotheses of \Cref{lem:n9v4-comparison}, define
\begin{equation}
 c_i:=\delta_i(f),\qquad
 b_j:=\sum_{i=1}^mD_{ij}c_i=(D^Tc)_j.                \label{eq:n9v4-b}
\end{equation}
Conditional on \(\mathcal F_{j-1}\), the martingale difference \(X_j\)
lies in an interval of length at most \(b_j\).  Hence
\begin{equation}
 \mathbb E[e^{\lambda X_j}\mid\mathcal F_{j-1}]
 \leq e^{\lambda^2b_j^2/8},                           \label{eq:n9v4-Xmgf}
\end{equation}
and \(f-\mathbb E_\nu f\) satisfies MAC with
\begin{equation}
 V_{\rm Dob}(f)
 :={1\over4}\|D^Tc\|_2^2.                             \label{eq:n9v4-VD}
\end{equation}
\end{theorem}

\begin{proof}
Fix the revealed values through coordinate \(j-1\), and compare two
possible values \(y_j,y_j'\).  The remaining tail laws have identical
earlier boundary coordinates and differ only at \(j\).  Applying
\Cref{lem:n9v4-comparison} to \(f\) shows that the conditional expectation
\(M_j\), viewed as a function of \(y_j\), has oscillation at most \(b_j\).
Subtracting \(M_{j-1}\), which is constant after conditioning on
\(\mathcal F_{j-1}\), does not alter this interval length.  Conditional
Hoeffding, \Cref{lem:n9v2-Hoeffding}, gives
\eqref{eq:n9v4-Xmgf}.  Summing \(v_j=b_j^2/4\) gives
\eqref{eq:n9v4-VD}.
\end{proof}

\begin{corollary}[Spectral-margin form]
\label{cor:n9v4-margin}
If
\begin{equation}
 \|C\|_{2\to2}\leq\kappa<1,                            \label{eq:n9v4-kappa}
\end{equation}
then
\begin{equation}
 V_{\rm Dob}(f)
 \leq{\|c\|_2^2\over4(1-\kappa)^2}.                  \label{eq:n9v4-margin}
\end{equation}
The normalized tilt \(r=e^{f-\mathbb Ef}/\mathbb E e^{f-\mathbb Ef}\)
satisfies
\begin{equation}
 \|r-1\|_1
 \leq
 {\sqrt2\,\|c\|_2\over1-\kappa}
 \exp\!\left\{
 {\|c\|_2^2\over4(1-\kappa)^2}
 \right\}.                                            \label{eq:n9v4-density}
\end{equation}
\end{corollary}

\begin{proof}
The Neumann series gives
\(\|D\|_{2\to2}\leq(1-\kappa)^{-1}\).  Insert this into
\eqref{eq:n9v4-VD}, then use \Cref{thm:n9v2-MAC}.
\end{proof}

\begin{corollary}[Row--column criterion]
\label{cor:n9v4-row-column}
If
\[
 \alpha:=\max_i\sum_jC_{ij}<1,\qquad
 \beta:=\max_j\sum_iC_{ij}<1,
\]
then \(\|C\|_{2\to2}\leq\sqrt{\alpha\beta}<1\), so
\eqref{eq:n9v4-margin} holds with
\(\kappa=\sqrt{\alpha\beta}\).
\end{corollary}

\begin{proof}
The matrix norm inequality
\(\|C\|_2^2\leq\|C\|_1\|C\|_\infty\) gives the claim.
\end{proof}

\subsection{Block coordinates and the critical-margin loss}
\label{sec:n9v4-block}

\begin{proposition}[Block Dobrushin loader]
\label{prop:n9v4-block}
Partition the microscopic variables into blocks and define the block
conditional kernels, block oscillations, and block interdependence matrix
\(C^{\rm bl}\) by \eqref{eq:n9v4-delta}--\eqref{eq:n9v4-C}.  If
\(\|C^{\rm bl}\|_2\leq\kappa_{\rm bl}<1\), all conclusions of
\Cref{thm:n9v4-range,cor:n9v4-margin} hold with block influence vector
\(c^{\rm bl}\) and
\begin{equation}
 V_{\rm bl}
 \leq{\|c^{\rm bl}\|_2^2\over
 4(1-\kappa_{\rm bl})^2}.                             \label{eq:n9v4-block}
\end{equation}
\end{proposition}

\begin{proof}
The product decomposition by blocks is another finite product of standard
Borel spaces.  Apply the preceding results verbatim.
\end{proof}

\begin{theorem}[Quantitative collapse at critical dependence]
\label{thm:n9v4-collapse}
For a sequence with \(\kappa_n\uparrow1\), the Dobrushin loader requires
\begin{equation}
 \sum_n{\|c_n\|_2\over1-\kappa_n}<\infty              \label{eq:n9v4-collapse}
\end{equation}
as a sufficient small-budget schedule.  In particular,
\(\sum_n\|c_n\|_2<\infty\) alone does not give a uniform conclusion when
the interdependence margin collapses.
\end{theorem}

\begin{proof}
By \eqref{eq:n9v4-margin},
\(\sqrt{V_n}\leq\|c_n\|_2/[2(1-\kappa_n)]\).
The MAC summability theorem therefore gives
\eqref{eq:n9v4-collapse} as the sufficient schedule.  The resolvent norm
can attain order \((1-\kappa_n)^{-1}\), for example when
\(C_n=\kappa_nP\) and \(P\) has eigenvalue one.  Thus no estimate uniform in
\(\kappa_n\) follows from this route.
\end{proof}

\begin{firewall}[Critical Ward modes]
\label{fire:n9v4-Ward}
Earlier audited work shows that a Ward-critical Gaussian owner can force
an absolute one-site influence row to reach one.  In such a sector the
strict Dobrushin premise may fail for structural reasons.  Enlarging blocks
helps only if it produces a proved block margin; it cannot remove an exact
conserved mode by notation.  A quotient observable, finite-horizon, or
conditional covariance argument may be needed instead.
\end{firewall}

\subsection{Insertion of the gauge--Higgs influences}
\label{sec:n9v4-GH}

\begin{theorem}[Dependent gauge--Higgs action loader]
\label{thm:n9v4-GH}
Use the truncated gauge--Higgs influence vector
\[
 c_{n,R}^{GH}:=
 ((c^U_{n,e}(R))_e,(c^\phi_{n,x}(R))_x)
\]
from \eqref{eq:n9v3-cU}--\eqref{eq:n9v3-cphi}.  Suppose the transported,
coarea-reduced reference has an interdependence matrix \(C_{n,R}\) with
\[
 \|C_{n,R}\|_2\leq\kappa_{n,R}<1.
\]
Then the centered truncated action satisfies MAC with
\begin{equation}
 V_{n,R}^{GH,{\rm dep}}
 \leq
 {\|c_{n,R}^{GH}\|_2^2\over
 4(1-\kappa_{n,R})^2}.                                \label{eq:n9v4-GH}
\end{equation}
If, for some truncation schedule \(R_n\),
\begin{equation}
 \sum_n
 \left[
 {\|c_{n,R_n}^{GH}\|_2\over1-\kappa_{n,R_n}}
 +\delta_n(R_n)
 \right]<\infty,                                      \label{eq:n9v4-GH-sum}
\end{equation}
then the isolated dependent gauge--Higgs action row is summable in total
variation.
\end{theorem}

\begin{proof}
Apply \Cref{cor:n9v4-margin} to the V3 influence vector.  The first term in
\eqref{eq:n9v4-GH-sum} makes the square root of
\eqref{eq:n9v4-GH} summable.  The second is the normalized truncation
debit.  Use \Cref{cor:n9v2-truncated}.
\end{proof}

\begin{definition}[Dependent martingale action card, DMAC]
\label{def:n9v4-DMAC}
DMAC is SAMAC with the independent-innovation item replaced by:
\begin{enumerate}[label=\textup{(DM\arabic*)},leftmargin=4.8em]
\item a declared coordinate or block decomposition of the constrained
reference;
\item an interdependence matrix \(C_n\) proved from its conditional kernels;
\item the propagated influence \(b_n=(I-C_n)^{-T}c_n\); and
\item the summability
\(\sum_n(\|b_n\|_2+\delta_n)<\infty\).
\end{enumerate}
\end{definition}

\begin{corollary}[DMAC closes the isolated action row]
\label{cor:n9v4-DMAC}
DMAC implies the isolated action conclusion of
\Cref{thm:n9v2-closure}.
\end{corollary}

\begin{proof}
\Cref{thm:n9v4-range} supplies MAC with
\(\sqrt V=\frac12\|b\|_2\).  Apply the summability and truncation arguments
of V2.
\end{proof}

\begin{firewall}[The physical interdependence matrix is unpaid]
\label{fire:n9v4-physical}
V4 does not prove a Dobrushin margin for the constrained SM--Einstein
measure.  The coarea Jacobian may alter conditional kernels, elliptic
constraint substitution may generate long-range entries, and massless or
Ward modes may force the margin to zero.  The theorem records all such
losses in the resolvent \((I-C_n)^{-1}\).
\end{firewall}

\begin{assessment}[Progress at ninth-series V4]
\label{ass:n9v4-status}
The product-reference gap in V3 has been replaced by a precise dependent
loader.  A local influence vector \(c\) becomes
\[
 b=(I-C)^{-T}c,
\]
and the action variance proxy is exactly \(\|b\|_2^2/4\).  The price of a
spectral Dobrushin margin is at most \((1-\kappa)^{-2}\).  This advances the
action row while exposing the next physical question: whether the
coarea-reduced local marginals have a strict block margin at all.
\end{assessment}

\begin{programme}[Ninth-series V5: compulsory companion audit]
\label{prog:n9v4-V5}
V5 will inspect the newest Yang--Mills and Navier--Stokes files, compare
their treatment of slow or conserved modes with
\Cref{fire:n9v4-Ward}, and decide between a block Dobrushin acquisition and
a quotient/finite-horizon covariance loader.  The selected direction will
be continued in V6.
\end{programme}

\begin{firewall}[Claim boundary after ninth-series V4]
\label{fire:n9v4-boundary}
V4 proves a loader implication.  It does not verify DMAC for the physical
measure, close the Einstein action, prove density convergence or tightness,
or construct the SM--Einstein continuum.
\end{firewall}

\begin{theorem}[Ninth-series V4 result]
\label{thm:n9v4-result}
A dependent finite reference with Dobrushin matrix \(C\) loads the
martingale action card with
\[
 V={1\over4}\|(I-C)^{-T}c\|_2^2.
\]
Under a spectral margin \(\kappa<1\), this is bounded by
\(\|c\|_2^2/[4(1-\kappa)^2]\), and
\eqref{eq:n9v4-GH-sum} conditionally closes the dependent truncated
gauge--Higgs row.
\end{theorem}

\begin{proof}
Use \Cref{lem:n9v4-comparison,thm:n9v4-range,
cor:n9v4-margin,thm:n9v4-GH}.
\end{proof}

\paragraph{Parent-version record.}
V4 was formed from
\texttt{Renewal\_SM\_Einstein\_Continuum\_Research\_Notes\_V3.tex}.
The V3 parent had \(1\,173\) logical source lines, \(43\,269\) bytes, and
SHA--256
\[
 \texttt{B6752E1E806620DDA5528BEC887F29E4C2EC2FBAA9588AAE37AD4053D80A466B}.
\]
The next version is the compulsory V5 companion audit.

% =============================================================================
% END APPENDED NINTH-SERIES V4 MATERIAL
% =============================================================================


% =============================================================================
% BEGIN APPENDED NINTH-SERIES V5 MATERIAL
% =============================================================================

\clearpage
\section{Ninth-series V5: compulsory companion audit and slow-label decision}
\label{sec:v9v5}

\subsection{Files and internal version status}
\label{sec:n9v5-files}

The active companion files inspected at this checkpoint were
\begin{center}
\small\ttfamily
Renewal\_Yang\_Mills\_Mass\_Gap\_Research\_Notes\_V56.tex,\\
Renewal\_Yang\_Mills\_Mass\_Gap\_Research\_Notes\_V57.tex,\\
Renewal\_NS\_Stability\_Research\_Notes\_V26.tex.
\end{center}
The two Yang--Mills files are byte-identical.  Each has \(23\,843\) logical
lines, \(853\,375\) bytes, and SHA--256
\[
 \texttt{C88EAF4D16DBB3691769EAF2ED68C90C089E26FC05AF3876B62910A2689E2F8D}.
\]
Both end with the internal V56 record, so there is no distinct V57 theorem
to audit.  The NS file is unchanged: it has \(5\,319\) logical lines,
\(278\,283\) bytes, and SHA--256
\[
 \texttt{82C887F8C2C69E4F28FAD7C3DC8DE5B9099CB5A4BF431EBA1FFF3E91935978AD}.
\]

\subsection{New Yang--Mills content}
\label{sec:n9v5-YM}

The YM V55 schedule separates lattice size, occupation window, and chart
threshold.  It proves relative-form and Feshbach return rates for a
logarithmic occupation window, a fifth-order magnetic remainder on that
window, and an explicit local rate table.  The interacting exit row remains
open.

The genuinely new V56 result takes a different route.  It proves the exact
quasi-free Euclidean Weyl pair product, shows stability of the pair kernel,
and combines a Mayer identity with a stable tree-graph inequality to obtain
a dimension-free connected Weyl tree bound.  No compact occupation cutoff
is differentiated.  A finite-rank spatial correction appears as a toron
tail.  The remaining interface is the weighted Weyl-representation norm of
the compressed Wilson vertices.

\begin{proposition}[Type-correct V56 transfer]
\label{prop:n9v5-YM-transfer}
The YM V56 tree theorem does not imply DMAC or a constrained SM action
estimate.  It supports one type-correct change of architecture: retain a
finite-dimensional slow variable explicitly, condition on that variable,
and apply the fast-sector concentration or cumulant theorem uniformly in
the slow label.
\end{proposition}

\begin{proof}
The non-transfer follows because the YM theorem concerns quasi-free Weyl
moments and assumes a representation interface not present in the
coarea-reduced SM density.  The architectural conclusion follows from the
form of the proved YM estimate: the local tree bound is dimension-free in
the fast sector, whereas the toron contribution is retained as a separate
finite-rank term.  Conditioning is the probability-theoretic analogue of
that exact finite-rank separation.
\end{proof}

\begin{firewall}[The Weyl representation budget is still open]
\label{fire:n9v5-YM}
The YM note does not yet prove a polynomial weighted Weyl variation norm
for the required compressed vertices.  V5 imports neither that missing
budget nor the YM mass-gap conclusion.
\end{firewall}

\subsection{Navier--Stokes status and the neutral-mode analogy}
\label{sec:n9v5-NS}

NS V26 still gives pointwise rank-one norms for the first two Oseen-kernel
derivatives and plans certified integral constants.  These constants do
not enter the SM constraint measure.  Its earlier V23 result identifies the
time derivative of an ancient profile as a neutral mode of the linearized
equation.  That fact is not an SM estimate, but it reinforces the logical
rule that a genuine neutral direction must be retained or quotiented before
a strict contraction theorem is applied.

\begin{proposition}[Type-correct NS consequence]
\label{prop:n9v5-NS-transfer}
No NS theorem changes the quantitative gauge--Higgs influence bounds of
V3 or supplies the interdependence matrix of V4.  The only relevant
structural lesson is to declare neutral modes before testing coercivity or
contraction on their complement.
\end{proposition}

\begin{proof}
The Oseen kernel and ancient-profile generator differ from the elliptic
Einstein constraint system and its induced probability law.  Therefore no
constant transfers.  The statement about neutral modes is a logical
requirement common to any quotient estimate and introduces no numerical NS
claim into the SM programme.
\end{proof}

\subsection{Why a global Dobrushin test can give the wrong verdict}
\label{sec:n9v5-mixture}

\begin{proposition}[A slow mixture can be maximally dependent]
\label{prop:n9v5-example}
Let \(S\) be uniform on \(\{-1,1\}\), and conditional on \(S\) let
\(Y_1=\cdots=Y_m=S\).  Conditional on \(S\), the fast law is deterministic
and has zero internal interdependence.  After \(S\) is marginalized, for
every \(i\ne j\) the conditional law of \(Y_i\) changes by total variation
one when \(Y_j\) changes sign.  Thus the global Dobrushin matrix has
off-diagonal coefficients one.
\end{proposition}

\begin{proof}
Given \(S\), all \(Y_i\) are fixed, so there is no fast conditional
fluctuation.  Marginally, \(Y_i=Y_j\) almost surely.  Hence
\[
 \mathcal L(Y_i\mid Y_j=1)=\delta_1,\qquad
 \mathcal L(Y_i\mid Y_j=-1)=\delta_{-1},
\]
whose total variation is one.
\end{proof}

\begin{corollary}[Failure of a global margin does not rule out a conditional
fast margin]
\label{cor:n9v5-conditional}
A Dobrushin coefficient equal to one can be caused entirely by an explicit
finite slow label.  It is therefore invalid to infer failure of all
conditional concentration estimates from failure of the unconditioned
one-site margin.
\end{corollary}

\begin{proof}
\Cref{prop:n9v5-example} has global coefficient one and conditional fast
coefficient zero.
\end{proof}

\subsection{Audit decision}
\label{sec:n9v5-decision}

\begin{audit}[V5 companion conclusion]
\label{aud:n9v5}
The audit selects a slow-label conditional action compiler before any
attempt to prove a global block Dobrushin margin.  The reasons are:
\begin{enumerate}[label=\textup{(\roman*)},leftmargin=3.8em]
\item V4 records the exact critical loss of the global resolvent;
\item the archived SM work already has finite conformal-Killing and other
slow sectors that must be transported explicitly;
\item YM V56 succeeds by separating a finite toron correction from its
dimension-free fast tree estimate; and
\item the NS neutral-mode record supports the same quotient discipline.
\end{enumerate}
\end{audit}

\begin{decision}[V6 acquisition order]
\label{dec:n9v5-V6}
Disintegrate the transported reference as
\[
 \nu(ds\,dy)=\pi(ds)\nu_s(dy).
\]
V6 will decompose a centered action increment into a conditional-fast
fluctuation and a slow effective action, prove normalized exponential-tilt
bounds for both pieces, and show how their total-variation debits compose.
The fast budget may come from a uniform conditional Dobrushin matrix; the
slow budget will remain a finite-dimensional integral that must be
estimated directly.
\end{decision}

\begin{firewall}[Conditioning is not deletion]
\label{fire:n9v5-conditioning}
Separating the slow label does not discard its probability law or its
action dependence.  The induced slow partition function is part of the
effective action and must be controlled.  A conditional fast estimate
alone cannot close the full action row.
\end{firewall}

\begin{assessment}[Progress at ninth-series V5]
\label{ass:n9v5-status}
The audit prevents the programme from forcing a structurally impossible
global margin.  The next theorem will retain the exact slow distribution
and charge both its effective tilt and the uniformly conditioned fast
tilts.  This is compatible with the finite slow-sector transport already
used in the determinant and constraint charts.
\end{assessment}

\begin{programme}[Ninth-series V6: conditional-fast mixture compiler]
\label{prog:n9v5-V6}
V6 will prove the slow/fast factorization selected in
\Cref{dec:n9v5-V6}, including a direct total-variation disintegration
inequality and a uniform conditional MAC bound.  It will then state the
slow-labelled dependent action card for the physical loader.
\end{programme}

\begin{firewall}[Claim boundary after ninth-series V5]
\label{fire:n9v5-boundary}
V5 is an audit and architectural decision.  It does not prove the
slow-label action budgets, the full action row, density convergence,
tightness, or the SM--Einstein continuum.
\end{firewall}

\begin{theorem}[Ninth-series V5 result]
\label{thm:n9v5-result}
The current companion notes support a finite-slow-label conditional route,
not a direct theorem for the SM action.  A global Dobrushin failure can be
caused entirely by slow-label mixing, so V6 must control the fast
conditional law and the induced slow effective action separately.
\end{theorem}

\begin{proof}
Use \Cref{prop:n9v5-YM-transfer,prop:n9v5-NS-transfer,
prop:n9v5-example,cor:n9v5-conditional,aud:n9v5}.
\end{proof}

\paragraph{Parent-version record.}
V5 was formed from
\texttt{Renewal\_SM\_Einstein\_Continuum\_Research\_Notes\_V4.tex}.
The V4 parent had \(1\,559\) logical source lines, \(57\,078\) bytes, and
SHA--256
\[
 \texttt{926CE368843203AE61012E758B19EF98F2B1D24CAC897F9D146733310D475D29}.
\]
The next compulsory companion audit is V10.

% =============================================================================
% END APPENDED NINTH-SERIES V5 MATERIAL
% =============================================================================


% =============================================================================
% BEGIN APPENDED NINTH-SERIES V6 MATERIAL
% =============================================================================

\clearpage
\section{Ninth-series V6: a conditional-fast mixture compiler}
\label{sec:v9v6}

\subsection{Exact slow--fast factorization}
\label{sec:n9v6-factor}

Let \(S\) and \(Y\) be standard Borel spaces.  Disintegrate a transported
local reference probability as
\begin{equation}
 \nu(ds\,dy)=\pi(ds)\nu_s(dy).                         \label{eq:n9v6-disintegration}
\end{equation}
Let \(F:S\times Y\to\mathbb R\) be an adjacent action increment with
\(0<\int e^F\,d\nu<\infty\).  Define
\begin{align}
 Z_s&:=\int_Ye^{F(s,y)}\,\nu_s(dy),                   \label{eq:n9v6-Zs}\\
 Z&:=\int_SZ_s\,\pi(ds),                              \label{eq:n9v6-Z}\\
 \mu_s(dy)&:={e^{F(s,y)}\over Z_s}\nu_s(dy),          \label{eq:n9v6-mus}\\
 \widehat\pi(ds)&:={Z_s\over Z}\pi(ds).               \label{eq:n9v6-pihat}
\end{align}

\begin{theorem}[Exact tilted disintegration]
\label{thm:n9v6-factor}
The normalized action tilt
\[
 d\mu={e^F\over Z}\,d\nu
\]
has the exact disintegration
\begin{equation}
 \mu(ds\,dy)=\widehat\pi(ds)\mu_s(dy).                \label{eq:n9v6-factor}
\end{equation}
Moreover,
\begin{equation}
 \|\mu-\nu\|_{\rm TV}
 \leq
 \|\widehat\pi-\pi\|_{\rm TV}
 +\int_S\|\mu_s-\nu_s\|_{\rm TV}\,\widehat\pi(ds).
 \label{eq:n9v6-TV-disintegration}
\end{equation}
\end{theorem}

\begin{proof}
Multiplication of \eqref{eq:n9v6-mus} and \eqref{eq:n9v6-pihat} gives
\[
 \widehat\pi(ds)\mu_s(dy)
 ={e^{F(s,y)}\over Z}\pi(ds)\nu_s(dy),
\]
which proves \eqref{eq:n9v6-factor}.  Insert the intermediate probability
\(\widehat\nu(ds\,dy):=\widehat\pi(ds)\nu_s(dy)\).  The triangle inequality
gives
\[
 \|\mu-\nu\|_{\rm TV}
 \leq\|\mu-\widehat\nu\|_{\rm TV}
 +\|\widehat\nu-\nu\|_{\rm TV}.
\]
The second term is \(\|\widehat\pi-\pi\|_{\rm TV}\).  For the first, the
density relative to \(\widehat\nu\) is \(e^F/Z_s\); Tonelli's theorem gives
exactly the integral of the conditional total variations in
\eqref{eq:n9v6-TV-disintegration}.
\end{proof}

\subsection{Uniform conditional-fast control}
\label{sec:n9v6-fast}

Put
\begin{align}
 m(s)&:=\int_YF(s,y)\,\nu_s(dy),                      \label{eq:n9v6-m}\\
 u_s(y)&:=F(s,y)-m(s),                                \label{eq:n9v6-us}\\
 q(s)&:=\log\int_Ye^{u_s(y)}\,\nu_s(dy).              \label{eq:n9v6-q}
\end{align}
Then the slow effective action is
\begin{equation}
 A(s):=\log Z_s=m(s)+q(s).                            \label{eq:n9v6-A}
\end{equation}

\begin{theorem}[Conditional MAC decomposition]
\label{thm:n9v6-conditional-MAC}
Suppose that, for \(\pi\)-almost every \(s\), \(u_s\) satisfies MAC under
\(\nu_s\) with budget \(V(s)\), and
\begin{equation}
 \overline V:=\mathop{\rm ess\,sup}_{s\sim\pi}V(s)<\infty.
 \label{eq:n9v6-Vbar}
\end{equation}
Then
\begin{align}
 \mathop{\rm ess\,sup}_{s}
 \|\mu_s-\nu_s\|_{\rm TV}
 &\leq\sqrt{2\overline V}\,e^{\overline V},           \label{eq:n9v6-fast-TV}\\
 0\leq q(s)&\leq{V(s)\over2},                         \label{eq:n9v6-q-bound}\\
 \mathop{\rm ess\,osc}_{\pi}q
 &\leq{\overline V\over2}.                            \label{eq:n9v6-q-osc}
\end{align}
\end{theorem}

\begin{proof}
The first estimate is \Cref{cor:n9v2-TV} applied in each conditional
fiber.  Jensen's inequality and
\(\int u_s\,d\nu_s=0\) give \(q(s)\geq0\).  The conditional MAC moment bound
at \(\lambda=1\) gives
\[
 \int e^{u_s}\,d\nu_s\leq e^{V(s)/2},
\]
proving \eqref{eq:n9v6-q-bound}.  Since \(q\) takes values in
\([0,\overline V/2]\) almost surely, \eqref{eq:n9v6-q-osc} follows.
\end{proof}

\begin{firewall}[Uniformity is over the tilted slow support]
\label{fire:n9v6-uniform}
An exceptional set of slow labels cannot be ignored merely because it has
small \(\pi\)-measure: the factor \(Z_s\) may amplify it under
\(\widehat\pi\).  Either the conditional estimate must be uniform on the
relevant slow support or the weighted exceptional integral must be entered
as a separate normalized truncation debit.
\end{firewall}

\subsection{Two loaders for the slow effective action}
\label{sec:n9v6-slow}

\begin{lemma}[Bounded logarithmic correction]
\label{lem:n9v6-bounded-correction}
Let \(\rho_a\) and \(\rho_{a+q}\) be the normalized exponential tilts of a
probability \(\pi\) by \(a\) and \(a+q\).  If
\(\operatorname{osc}q\leq\eta\), then
\begin{equation}
 \|\rho_{a+q}-\rho_a\|_{\rm TV}
 \leq{e^\eta-1\over2}.                                \label{eq:n9v6-bounded-correction}
\end{equation}
\end{lemma}

\begin{proof}
Relative to \(\rho_a\), the density of \(\rho_{a+q}\) is
\[
 R={e^q\over\int e^q\,d\rho_a}.
\]
If \(q\in[\alpha,\alpha+\eta]\), then
\(e^{-\eta}\leq R\leq e^\eta\).  Hence
\(\int|R-1|\,d\rho_a\leq e^\eta-1\).  Divide by two.
\end{proof}

\begin{theorem}[Slow MAC loader]
\label{thm:n9v6-slow-MAC}
Assume the centered conditional mean
\[
 m-\int_Sm\,d\pi
\]
satisfies MAC under \(\pi\) with budget \(W\).  Under the conditional-fast
hypotheses of \Cref{thm:n9v6-conditional-MAC},
\begin{equation}
 \|\widehat\pi-\pi\|_{\rm TV}
 \leq
 \sqrt{2W}\,e^W
 {e^{\overline V/2}-1\over2}.                         \label{eq:n9v6-slow-MAC}
\end{equation}
\end{theorem}

\begin{proof}
Let \(\pi_m\) be the normalized tilt of \(\pi\) by \(m\).  The slow MAC
bound \Cref{cor:n9v2-TV} gives
\[
 \|\pi_m-\pi\|_{\rm TV}\leq\sqrt{2W}\,e^W.
\]
By \eqref{eq:n9v6-A}, \(\widehat\pi\) is the tilt by \(m+q\).  Use
\eqref{eq:n9v6-q-osc}, \Cref{lem:n9v6-bounded-correction}, and the triangle
inequality.
\end{proof}

\begin{theorem}[Compact Lipschitz slow loader]
\label{thm:n9v6-slow-Lip}
Suppose \(S=K\) is a compact metric space of diameter \(d_K\) and
\begin{equation}
 |m(s)-m(t)|\leq L\,d(s,t).                            \label{eq:n9v6-Lip}
\end{equation}
Then
\begin{equation}
 \|\widehat\pi-\pi\|_{\rm TV}
 \leq
 {1\over2}
 \left[
 \exp\!\left(d_KL+{\overline V\over2}\right)-1
 \right].                                             \label{eq:n9v6-slow-Lip}
\end{equation}
\end{theorem}

\begin{proof}
Equations \eqref{eq:n9v6-Lip} and \eqref{eq:n9v6-q-osc} give
\[
 \operatorname{osc}A
 \leq d_KL+\overline V/2.
\]
Apply \Cref{lem:n9v6-bounded-correction} with \(a=0\) and \(q=A\).
\end{proof}

\subsection{The complete mixture bounds}
\label{sec:n9v6-complete}

\begin{theorem}[Slow--fast action compiler]
\label{thm:n9v6-compiler}
Under the conditional-fast hypothesis with budget \(\overline V\), either
of the following bounds holds.
\begin{enumerate}[label=\textup{(\roman*)},leftmargin=3.8em]
\item If the conditional mean has slow MAC budget \(W\), then
\begin{equation}
 \|\mu-\nu\|_{\rm TV}
 \leq
 \sqrt{2\overline V}\,e^{\overline V}
 +\sqrt{2W}\,e^W
 +{e^{\overline V/2}-1\over2}.                        \label{eq:n9v6-total-MAC}
\end{equation}
\item If \(S=K\) and \(m\) has Lipschitz constant \(L\), then
\begin{equation}
 \|\mu-\nu\|_{\rm TV}
 \leq
 \sqrt{2\overline V}\,e^{\overline V}
 +{1\over2}
 \left[
 e^{d_KL+\overline V/2}-1
 \right].                                             \label{eq:n9v6-total-Lip}
\end{equation}
\end{enumerate}
\end{theorem}

\begin{proof}
Use the disintegration inequality
\eqref{eq:n9v6-TV-disintegration}.  Bound its conditional integral by
\eqref{eq:n9v6-fast-TV}.  Bound the slow term by
\Cref{thm:n9v6-slow-MAC} in case (i) and by
\Cref{thm:n9v6-slow-Lip} in case (ii).
\end{proof}

\begin{corollary}[Summable cofinal mixture schedules]
\label{cor:n9v6-sum}
For adjacent regulator steps \(n\), the action total-variation debits are
summable under either
\begin{align}
 \sum_n(\sqrt{\overline V_n}+\sqrt{W_n})&<\infty,
                                                              \label{eq:n9v6-sum-MAC}\\
 \sum_n(\sqrt{\overline V_n}+d_{K_n}L_n)&<\infty.       \label{eq:n9v6-sum-Lip}
\end{align}
\end{corollary}

\begin{proof}
Either assumption forces every term in its sum to zero.  For small
\(x\geq0\), \(e^x-1\leq2x\).  Apply this to
\eqref{eq:n9v6-total-MAC} or \eqref{eq:n9v6-total-Lip}; finitely many
initial indices contribute a finite amount.
\end{proof}

\subsection{Loading the conditional fast budget from V4}
\label{sec:n9v6-Dobrushin}

\begin{corollary}[Uniform conditional Dobrushin loader]
\label{cor:n9v6-Dobrushin}
Suppose that for \(\pi\)-almost every \(s\), the fast conditional law
\(\nu_s\) has Dobrushin matrix \(C_s\), and the conditional action
oscillation vector is \(c_s\).  If
\[
 \mathop{\rm ess\,sup}_s\|C_s\|_2\leq\kappa<1,\qquad
 \mathop{\rm ess\,sup}_s\|c_s\|_2\leq\chi,
\]
then
\begin{equation}
 \overline V\leq{\chi^2\over4(1-\kappa)^2}.           \label{eq:n9v6-Dobrushin}
\end{equation}
\end{corollary}

\begin{proof}
Apply \Cref{cor:n9v4-margin} in every fiber and take the essential
supremum.
\end{proof}

\begin{remark}[Finite rank does not mean bounded diameter]
\label{rem:n9v6-rank}
The slow space may have fixed dimension while its diameter grows with the
regulator.  The Lipschitz loader pays \(d_{K_n}L_n\), so compactness at each
cutoff is insufficient.  Alternatively, the slow MAC loader can use an
unbounded slow space if its moment card is uniform.
\end{remark}

\subsection{The slow-labelled action card}
\label{sec:n9v6-card}

\begin{definition}[Slow-labelled dependent action card, SLDAC]
\label{def:n9v6-SLDAC}
For every support in the countable local core, SLDAC records:
\begin{enumerate}[label=\textup{(S\arabic*)},leftmargin=4em]
\item a common measurable slow space and disintegration
\(\nu_n(ds\,dy)=\pi_n(ds)\nu_{n,s}(dy)\);
\item matched slow rank and an explicit transport of the slow coordinates;
\item a conditional fast action and a uniform conditional MAC budget
\(\overline V_n\);
\item the induced slow mean \(m_n(s)\) and either a slow MAC budget \(W_n\)
or a compact Lipschitz budget \(d_{K_n}L_n\);
\item every normalized truncation or exceptional-slow-label debit;
\item one of the summability schedules
\eqref{eq:n9v6-sum-MAC}--\eqref{eq:n9v6-sum-Lip}; and
\item compatibility with the exact determinant, coarea, transport, and
counterterm ownership ledger.
\end{enumerate}
\end{definition}

\begin{theorem}[Conditional SLDAC closure]
\label{thm:n9v6-closure}
SLDAC implies a summable total-variation debit for the isolated adjacent
action tilt.  Together with a proved joint density compiler for the other
rows, it discharges the action hypothesis in
\Cref{thm:n9v1-dependency}.
\end{theorem}

\begin{proof}
Apply \Cref{thm:n9v6-compiler,cor:n9v6-sum} at every adjacent step and add
the declared exceptional debits.  The ownership item permits insertion
into the exact joint density ledger without double counting.
\end{proof}

\begin{firewall}[Moving or branching slow sectors]
\label{fire:n9v6-moving}
SLDAC does not cover a changing slow dimension, a multiplier stratum that
branches, or a slow-coordinate map whose Jacobian lacks a summable
logarithmic bound.  Those cases return to the slow-sector transport and
singular-coarea problems; they cannot be hidden inside the conditional
kernel.
\end{firewall}

\begin{assessment}[Progress at ninth-series V6]
\label{ass:n9v6-status}
An exact conserved or toron-like label no longer has to satisfy a strict
global Dobrushin margin.  V6 conditions on it, proves the fast action tilt
uniformly, and charges the induced slow partition function through
\(A=m+q\).  The fast fluctuation contributes
\(O(\sqrt{\overline V})\); its feedback into the slow action is only
\(O(\overline V)\).  The remaining physical work is to identify the actual
slow coordinates and prove uniform conditional budgets.
\end{assessment}

\begin{programme}[Ninth-series V7: the Einstein slow-label interface]
\label{prog:n9v6-V7}
V7 will test SLDAC against the finite conformal-Killing multiplier sector
and matched slow eigenspaces of the constraint normal operator.  It will
derive the slow effective-action derivative by differentiating the
conditional partition function, state the covariance term exactly, and
identify which uniform moment bounds yield the Lipschitz or slow-MAC
budget.
\end{programme}

\begin{firewall}[Claim boundary after ninth-series V6]
\label{fire:n9v6-boundary}
V6 proves the mixture compiler, not SLDAC for the physical sequence.  It
does not control singular slow strata, the Euclidean conformal instability,
full density convergence, tightness, or the SM--Einstein continuum.
\end{firewall}

\begin{theorem}[Ninth-series V6 result]
\label{thm:n9v6-result}
For an exact slow--fast disintegration, uniform conditional MAC with budget
\(\overline V\) and either slow MAC budget \(W\) or compact Lipschitz budget
\(d_KL\) give the total-variation bounds
\eqref{eq:n9v6-total-MAC}--\eqref{eq:n9v6-total-Lip}.  Thus the corresponding
summable schedules close the isolated slow-labelled action row without a
global Dobrushin margin.
\end{theorem}

\begin{proof}
Use \Cref{thm:n9v6-factor,thm:n9v6-conditional-MAC,
thm:n9v6-slow-MAC,thm:n9v6-slow-Lip,
thm:n9v6-compiler,cor:n9v6-sum}.
\end{proof}

\paragraph{Parent-version record.}
V6 was formed from
\texttt{Renewal\_SM\_Einstein\_Continuum\_Research\_Notes\_V5.tex}.
The V5 parent had \(1\,787\) logical source lines, \(66\,510\) bytes, and
SHA--256
\[
 \texttt{05B1C08201B87AFB68012A0E649F63DC93EE2F376D28D03D4545F2D7016A5E29}.
\]
The next compulsory companion audit is V10.

% =============================================================================
% END APPENDED NINTH-SERIES V6 MATERIAL
% =============================================================================


% =============================================================================
% BEGIN APPENDED NINTH-SERIES V7 MATERIAL
% =============================================================================

\clearpage
\section{Ninth-series V7: the Einstein slow-label derivative interface}
\label{sec:v9v7}

\subsection{A dominated conditional Gibbs chart}
\label{sec:n9v7-setting}

Let \(K\subset\mathbb R^r\) be compact and convex.  Let \((Y,\lambda)\) be
a sigma-finite measured space, and suppose the fast conditional reference
has density
\begin{equation}
 \nu_s(dy)
 =
 {e^{-H(s,y)}\over N(s)}\,\lambda(dy),\qquad
 N(s):=\int_Ye^{-H(s,y)}\,\lambda(dy).                \label{eq:n9v7-nu}
\end{equation}
Let \(F(s,y)\) be the adjacent action increment, and define
\begin{align}
 Z_s&:=\int_Ye^{F(s,y)}\,\nu_s(dy),                   \label{eq:n9v7-Zs}\\
 A(s)&:=\log Z_s,                                     \label{eq:n9v7-A}\\
 \mu_s(dy)&:={e^{F(s,y)}\over Z_s}\,\nu_s(dy).        \label{eq:n9v7-mu}
\end{align}
Assume \(H\) and \(F\) are \(C^1\) in \(s\), and that on every compact
subchart the integrands obtained by differentiating
\eqref{eq:n9v7-nu}--\eqref{eq:n9v7-Zs} are dominated by one integrable
function.  This is the declared differentiation hypothesis; V7 does not
infer it from pointwise smoothness.

\subsection{The score and the exact derivative}
\label{sec:n9v7-score}

\begin{lemma}[Conditional Gibbs score]
\label{lem:n9v7-score}
For a coordinate direction \(a=1,\ldots,r\), the logarithmic score of
\(\nu_s\) is
\begin{equation}
 \sigma_a(s,y)
 :=
 \partial_a\log{d\nu_s\over d\lambda}(y)
 =
 -\partial_aH(s,y)
 +\mathbb E_{\nu_s}\partial_aH(s,\cdot),              \label{eq:n9v7-score}
\end{equation}
and
\begin{equation}
 \mathbb E_{\nu_s}\sigma_a=0.                         \label{eq:n9v7-score-zero}
\end{equation}
\end{lemma}

\begin{proof}
Differentiate
\(\log(d\nu_s/d\lambda)=-H(s,\cdot)-\log N(s)\).
Dominated differentiation gives
\[
 \partial_a\log N(s)
 =
 -{\int(\partial_aH)e^{-H}\,d\lambda\over N(s)}
 =-\mathbb E_{\nu_s}\partial_aH.
\]
This proves \eqref{eq:n9v7-score}; integration proves
\eqref{eq:n9v7-score-zero}.
\end{proof}

\begin{theorem}[Exact slow effective-action derivative]
\label{thm:n9v7-derivative}
Under the preceding assumptions,
\begin{align}
 \partial_aA(s)
 &=
 \mathbb E_{\mu_s}
 [\partial_aF(s,\cdot)+\sigma_a(s,\cdot)]             \label{eq:n9v7-dA-score}\\
 &=
 \mathbb E_{\mu_s}[\partial_aF-\partial_aH]
 +\mathbb E_{\nu_s}\partial_aH.                       \label{eq:n9v7-dA-Gibbs}
\end{align}
\end{theorem}

\begin{proof}
Differentiate \(Z_s=\int e^F\,d\nu_s\).  The derivative of the density is
\(\sigma_a\,d\nu_s\), so
\[
 \partial_aZ_s
 =
 \int e^F(\partial_aF+\sigma_a)\,d\nu_s.
\]
Division by \(Z_s\) gives \eqref{eq:n9v7-dA-score}.  Substitute
\eqref{eq:n9v7-score} to obtain \eqref{eq:n9v7-dA-Gibbs}.
\end{proof}

\begin{corollary}[Exact covariance derivative of the conditional mean]
\label{cor:n9v7-covariance}
For
\[
 m(s):=\mathbb E_{\nu_s}F(s,\cdot),
\]
one has
\begin{equation}
 \partial_am(s)
 =
 \mathbb E_{\nu_s}\partial_aF
 -\operatorname{Cov}_{\nu_s}(F,\partial_aH).          \label{eq:n9v7-dm}
\end{equation}
If \(q=A-m\), then
\begin{equation}
 \partial_aq
 =
 \mathbb E_{\mu_s}(\partial_aF-\partial_aH)
 -\mathbb E_{\nu_s}\partial_aF
 +\mathbb E_{\nu_s}\partial_aH
 +\operatorname{Cov}_{\nu_s}(F,\partial_aH).          \label{eq:n9v7-dq}
\end{equation}
\end{corollary}

\begin{proof}
Differentiate \(\int F\,d\nu_s\):
\[
 \partial_am
 =
 \mathbb E_{\nu_s}\partial_aF
 +\mathbb E_{\nu_s}(F\sigma_a).
\]
Use \eqref{eq:n9v7-score-zero} and \eqref{eq:n9v7-score} to identify the
second term as
\(-\operatorname{Cov}_{\nu_s}(F,\partial_aH)\).
Subtract \eqref{eq:n9v7-dm} from \eqref{eq:n9v7-dA-Gibbs}.
\end{proof}

\begin{firewall}[The score term cannot be dropped]
\label{fire:n9v7-score}
The formula
\(\partial_aA=\mathbb E_{\mu_s}\partial_aF\) is valid only when
\(\nu_s\) is independent of \(s\).  In a moving constraint, metric, or
coarea chart, the score term records the change of the conditional
reference.  Omitting it would assign a false slow Lipschitz budget.
\end{firewall}

\subsection{Two quantitative derivative bounds}
\label{sec:n9v7-bounds}

For a vector-valued derivative, use the Euclidean norm in the slow
coordinate.  Define
\begin{align}
 B_F&:=
 \sup_{s,y}\|\nabla_sF(s,y)\|,                        \label{eq:n9v7-BF}\\
 B_H&:=
 \sup_s\operatorname{osc}_{y}
 \nabla_sH(s,y),                                      \label{eq:n9v7-BH}
\end{align}
where the vector oscillation is
\[
 \operatorname{osc}_{y}G
 :=\sup_{y,y'}\|G(y)-G(y')\|.
\]

\begin{theorem}[Total-variation derivative bound]
\label{thm:n9v7-TV-bound}
Suppose \(B_F,B_H<\infty\).  If the conditional centered action satisfies
MAC with budget at most \(\overline V\), uniformly in \(s\), then
\begin{equation}
 \|\nabla_sA(s)\|
 \leq
 B_F+B_H\sqrt{2\overline V}\,e^{\overline V}
 \quad(s\in K).                                      \label{eq:n9v7-TV-bound}
\end{equation}
\end{theorem}

\begin{proof}
From \eqref{eq:n9v7-dA-Gibbs},
\[
 \nabla_sA
 =
 \mathbb E_{\mu_s}\nabla_sF
 -\left(
 \mathbb E_{\mu_s}\nabla_sH
 -\mathbb E_{\nu_s}\nabla_sH
 \right).
\]
The first term has norm at most \(B_F\).  For probabilities \(\alpha,\beta\)
and a bounded vector function \(G\),
\[
 \|\mathbb E_\alpha G-\mathbb E_\beta G\|
 \leq\operatorname{osc}(G)\|\alpha-\beta\|_{\rm TV};
\]
apply the scalar dual estimate after pairing with unit vectors.  The
conditional MAC bound \eqref{eq:n9v6-fast-TV} completes the proof.
\end{proof}

\begin{theorem}[Covariance derivative bound]
\label{thm:n9v7-cov-bound}
Suppose
\begin{align}
 \sup_s\mathbb E_{\nu_s}
 \|\nabla_sF(s,\cdot)\|&\leq B_1,                     \label{eq:n9v7-B1}\\
 \sup_s\mathbb E_{\nu_s}
 \|\nabla_sH-\mathbb E_{\nu_s}\nabla_sH\|^2
 &\leq Q_H^2.                                         \label{eq:n9v7-QH}
\end{align}
If \(F(s,\cdot)-m(s)\) satisfies conditional MAC with budget at most
\(\overline V\), then
\begin{equation}
 \|\nabla_sm(s)\|
 \leq B_1+Q_H\sqrt{\overline V}.                      \label{eq:n9v7-cov-bound}
\end{equation}
\end{theorem}

\begin{proof}
The vector form of \eqref{eq:n9v7-dm} is
\[
 \nabla_sm
 =
 \mathbb E_{\nu_s}\nabla_sF
 -\mathbb E_{\nu_s}
 [(F-m)(\nabla_sH-\mathbb E_{\nu_s}\nabla_sH)].
\]
Cauchy--Schwarz bounds the second term by
\[
 \sqrt{\operatorname{Var}_{\nu_s}F}\,
 \left(
 \mathbb E_{\nu_s}
 \|\nabla_sH-\mathbb E_{\nu_s}\nabla_sH\|^2
 \right)^{1/2}.
\]
The conditional MAC variance estimate
\eqref{eq:n9v2-var} gives
\(\operatorname{Var}_{\nu_s}F\leq\overline V\).
\end{proof}

\subsection{Loading the compact slow Lipschitz row}
\label{sec:n9v7-Lip}

\begin{corollary}[Effective slow Lipschitz budget]
\label{cor:n9v7-Lip}
Under \Cref{thm:n9v7-TV-bound}, \(A\) is Lipschitz on the convex chart
\(K\) with
\begin{equation}
 L_A
 \leq
 B_F+B_H\sqrt{2\overline V}\,e^{\overline V}.         \label{eq:n9v7-LA}
\end{equation}
Consequently,
\begin{equation}
 \|\widehat\pi-\pi\|_{\rm TV}
 \leq{1\over2}
 \left[
 \exp\!\left\{
 d_K\left(
 B_F+B_H\sqrt{2\overline V}\,e^{\overline V}
 \right)
 \right\}-1
 \right].                                             \label{eq:n9v7-slow-TV}
\end{equation}
\end{corollary}

\begin{proof}
Integrate the gradient bound along the segment joining two points of
\(K\).  Then apply the bounded-logarithmic-correction lemma
\Cref{lem:n9v6-bounded-correction} to the slow tilt by \(A\).
\end{proof}

\begin{theorem}[Summable derivative schedule]
\label{thm:n9v7-sum}
For adjacent steps \(n\), suppose
\(\overline V_n\to0\) and
\begin{equation}
 \sum_n\left[
 \sqrt{\overline V_n}
 +d_{K_n}B_{F,n}
 +d_{K_n}B_{H,n}\sqrt{\overline V_n}
 \right]<\infty.                                      \label{eq:n9v7-sum}
\end{equation}
Then the slow--fast action tilts have summable total-variation debits.
\end{theorem}

\begin{proof}
For large \(n\), \(e^{\overline V_n}\) is uniformly bounded and the
exponent in \eqref{eq:n9v7-slow-TV} tends to zero.  Use
\(e^x-1\leq2x\) for small \(x\), add the fast debit
\(\sqrt{2\overline V_n}e^{\overline V_n}\), and sum.  Finitely many initial
terms are harmless.
\end{proof}

\subsection{Relation to Einstein and gauge slow sectors}
\label{sec:n9v7-physical}

\begin{proposition}[Three different meanings of a slow sector]
\label{prop:n9v7-types}
The existing programme contains three mathematically distinct objects:
\begin{enumerate}[label=\textup{(\roman*)},leftmargin=3.8em]
\item a physical harmonic or toron coordinate that remains in the
probability space;
\item a finite kernel coordinate used to augment the Einstein constraint
solve; and
\item a matched finite spectral projector removed from a reduced
determinant.
\end{enumerate}
Only item (i) is automatically a random slow label of the disintegration
\eqref{eq:n9v6-disintegration}.  Item (ii) becomes such a label only if it
parametrizes surviving points of the physical multiplier locus.  Item
(iii) is determinant bookkeeping and cannot be integrated again as an
action variable without a separate identification.
\end{proposition}

\begin{proof}
The three objects live respectively in the field configuration, the
domain/codomain augmentation of the constraint map, and the spectral
factorization of the normal operator.  Their ambient spaces and measures
differ.  Equality of their finite ranks supplies no canonical
measure-preserving identification.
\end{proof}

\begin{definition}[Slow effective partition card, SEPC]
\label{def:n9v7-SEPC}
For a proposed Einstein or gauge slow chart, SEPC records:
\begin{enumerate}[label=\textup{(E\arabic*)},leftmargin=4em]
\item why the slow coordinate survives on the physical probability space;
\item a fixed dominating fast measure after all field transports;
\item the conditional reference energy \(H_n(s,y)\) and adjacent action
\(F_n(s,y)\);
\item dominated differentiability of both conditional partition functions;
\item the exact score \eqref{eq:n9v7-score};
\item uniform conditional MAC budget \(\overline V_n\);
\item the derivative budgets \(B_{F,n}\) and either \(B_{H,n}\) or
\(Q_{H,n}\); and
\item the summability schedule \eqref{eq:n9v7-sum}, plus every exceptional
chart debit.
\end{enumerate}
\end{definition}

\begin{theorem}[Conditional SEPC-to-SLDAC loader]
\label{thm:n9v7-loader}
SEPC with the \(B_H\) branch implies the compact-Lipschitz part of SLDAC
and therefore closes the isolated slow-labelled action row.
\end{theorem}

\begin{proof}
\Cref{cor:n9v7-Lip} supplies the slow effective-action Lipschitz constant.
The schedule \eqref{eq:n9v7-sum} and
\Cref{thm:n9v7-sum} give summable slow--fast debits.  The remaining SEPC
items are exactly the measure identification, differentiation, and
exceptional-set data required by SLDAC.
\end{proof}

\begin{firewall}[The conformal instability is not a score estimate]
\label{fire:n9v7-conformal}
A finite derivative bound on a compact slow chart does not control the
unbounded Euclidean conformal direction.  Removing the chart cutoff
requires a normalized weighted-tail estimate or a positive replacement
for the gravitational measure.  SEPC leaves that obstruction visible.
\end{firewall}

\begin{assessment}[Progress at ninth-series V7]
\label{ass:n9v7-status}
The induced slow action is no longer formal.  Its derivative contains the
conditional score, and the derivative of the conditional mean contains the
exact covariance
\(-\operatorname{Cov}(F,\nabla_sH)\).  Uniform conditional MAC turns this
into the explicit Lipschitz schedule \eqref{eq:n9v7-sum}.  The remaining
task is a type-correct identification and transport of the physical slow
coordinates.
\end{assessment}

\begin{programme}[Ninth-series V8: matched slow-chart transport]
\label{prog:n9v7-V8}
V8 will construct a common finite-dimensional slow coordinate from two
nearby spectral subspaces using the polar unitary, prove its Lipschitz
Jacobian and score transport, and state the gap/rank conditions needed to
preserve SEPC across adjacent regulators.  If the rank changes, it will
return to the singular-stratum debit rather than force an identification.
\end{programme}

\begin{firewall}[Claim boundary after ninth-series V7]
\label{fire:n9v7-boundary}
V7 does not verify SEPC for a physical slow chart, remove the conformal
cutoff, prove full density convergence or tightness, or construct the
SM--Einstein continuum.
\end{firewall}

\begin{theorem}[Ninth-series V7 result]
\label{thm:n9v7-result}
For a dominated conditional Gibbs chart, the effective slow action obeys
the exact score formula \eqref{eq:n9v7-dA-Gibbs}.  Conditional MAC gives
the quantitative derivative bound \eqref{eq:n9v7-TV-bound}, and the
summability schedule \eqref{eq:n9v7-sum} closes the isolated slow-labelled
action row under SEPC.
\end{theorem}

\begin{proof}
Use \Cref{lem:n9v7-score,thm:n9v7-derivative,
thm:n9v7-TV-bound,cor:n9v7-Lip,
thm:n9v7-sum,thm:n9v7-loader}.
\end{proof}

\paragraph{Parent-version record.}
V7 was formed from
\texttt{Renewal\_SM\_Einstein\_Continuum\_Research\_Notes\_V6.tex}.
The V6 parent had \(2\,173\) logical source lines, \(80\,011\) bytes, and
SHA--256
\[
 \texttt{83F6C34B95C2539A184A24CE0E47FC1C1AB4E7E6801A81B270E419808027DE1D}.
\]
The next compulsory companion audit is V10.

% =============================================================================
% END APPENDED NINTH-SERIES V7 MATERIAL
% =============================================================================


% =============================================================================
% BEGIN APPENDED NINTH-SERIES V8 MATERIAL
% =============================================================================

\clearpage
\section{Ninth-series V8: matched spectral slow-chart transport}
\label{sec:v9v8}

\subsection{Two nearby slow subspaces}
\label{sec:n9v8-projections}

Let \(\mathcal H\) be a real or complex Hilbert space and let \(P,Q\) be
orthogonal projections.  Put
\begin{equation}
 \delta:=\|P-Q\|.                                     \label{eq:n9v8-delta}
\end{equation}
The intended application takes \(P\) and \(Q\) to be adjacent slow spectral
projectors of a common-refinement realization of the constraint normal
operator.

\begin{lemma}[Compression floor]
\label{lem:n9v8-floor}
If \(\delta<1\), then on \(\operatorname{Ran}P\),
\begin{equation}
 PQP\geq(1-\delta^2)P.                                \label{eq:n9v8-floor}
\end{equation}
Consequently \(QP:\operatorname{Ran}P\to\operatorname{Ran}Q\) is injective
and has closed range.
\end{lemma}

\begin{proof}
For \(x=Px\),
\[
 \langle x,PQPx\rangle=\|Qx\|^2
 =\|x\|^2-\|(I-Q)x\|^2.
\]
Since \((I-Q)P=(P-Q)P\), its norm is at most \(\delta\).  This proves
\eqref{eq:n9v8-floor}.  The lower norm bound
\(\|Qx\|\geq\sqrt{1-\delta^2}\|x\|\) gives injectivity and closed range.
\end{proof}

\begin{theorem}[Canonical polar transport]
\label{thm:n9v8-polar}
Assume \(\delta<1\) and \(\operatorname{rank}P<\infty\).  Then
\(\operatorname{rank}Q=\operatorname{rank}P\), and
\begin{equation}
 W_{Q,P}
 :=
 QP(PQP)^{-1/2}:
 \operatorname{Ran}P\longrightarrow\operatorname{Ran}Q              \label{eq:n9v8-W}
\end{equation}
is unitary between the two ranges.  It is the polar factor of \(QP\) and
obeys
\begin{equation}
 \|W_{Q,P}-P\|_{\operatorname{Ran}P\to\mathcal H}
 \leq
 \delta+{\delta^2\over1-\delta^2}.                    \label{eq:n9v8-W-bound}
\end{equation}
In particular, if \(\delta\leq\frac12\), the right side is at most
\(2\delta\).
\end{theorem}

\begin{proof}
\Cref{lem:n9v8-floor} makes \(PQP\) positive and invertible on
\(\operatorname{Ran}P\).  Direct calculation gives
\[
 W_{Q,P}^*W_{Q,P}
 =(PQP)^{-1/2}PQ P(PQP)^{-1/2}=P.
\]
Thus \(W_{Q,P}\) is an isometry.  Interchanging \(P,Q\) shows that
\(PQ:\operatorname{Ran}Q\to\operatorname{Ran}P\) is also injective.
Finite dimensionality gives equality of the ranks, so the isometry is
onto \(\operatorname{Ran}Q\).

On \(\operatorname{Ran}P\),
\[
 W_{Q,P}-P
 =(Q-P)P(PQP)^{-1/2}
 +P\big[(PQP)^{-1/2}-P\big].
\]
The first term is bounded by
\(\delta(1-\delta^2)^{-1/2}\).  A slightly coarser decomposition
\[
 W_{Q,P}-P=(Q-P)P+QP[(PQP)^{-1/2}-P]
\]
gives
\[
 \|W_{Q,P}-P\|
 \leq\delta+
 \left[(1-\delta^2)^{-1/2}-1\right].
\]
For \(0\leq x<1\),
\((1-x)^{-1/2}-1\leq x/(1-x)\); use \(x=\delta^2\) to obtain
\eqref{eq:n9v8-W-bound}.  If \(\delta\leq1/2\), then
\(\delta^2/(1-\delta^2)\leq2\delta^2\leq\delta\).
\end{proof}

\subsection{Canonical frames and their Jacobian}
\label{sec:n9v8-frames}

\begin{corollary}[Frame transport has unit Jacobian]
\label{cor:n9v8-J}
Let \(E:\mathbb F^r\to\operatorname{Ran}P\) be an orthonormal frame, where
\(\mathbb F=\mathbb R\) or \(\mathbb C\), and define
\begin{equation}
 E':=W_{Q,P}E.                                        \label{eq:n9v8-frame}
\end{equation}
Then \(E'\) is an orthonormal frame of \(\operatorname{Ran}Q\).  The common
coordinate \(s\in\mathbb F^r\) represents adjacent slow vectors as
\[
 Es\longmapsto E's.
\]
This transport preserves Euclidean volume:
\begin{equation}
 |\det(E'^*W_{Q,P}E)|=1.                              \label{eq:n9v8-J}
\end{equation}
Hence it contributes zero to the logarithmic coordinate-Jacobian row.
\end{corollary}

\begin{proof}
Unitarity of \(W_{Q,P}\) on the ranges gives \(E'^*E'=I_r\).  Moreover,
\[
 E'^*W_{Q,P}E
 =E^*W_{Q,P}^*W_{Q,P}E=I_r,
\]
whose determinant is one.
\end{proof}

\begin{remark}[Orientation and complex phase]
\label{rem:n9v8-phase}
An independently chosen adjacent frame differs from \(E'\) by an
orthogonal or unitary \(r\times r\) matrix.  Its absolute real Jacobian is
one, although an orientation sign or complex determinant phase may remain.
The recursive polar choice fixes that ambiguity along a cofinal chain; it
does not prove trivial holonomy around a loop in parameter space.
\end{remark}

\subsection{Cofinal frame convergence}
\label{sec:n9v8-cofinal}

\begin{theorem}[Summable projector motion]
\label{thm:n9v8-sum}
Let \(P_n\) be finite-rank orthogonal projections of the same rank on a
common Hilbert space and suppose
\begin{equation}
 \delta_n:=\|P_{n+1}-P_n\|\leq{1\over2},\qquad
 \sum_n\delta_n<\infty.                               \label{eq:n9v8-sum}
\end{equation}
Starting from an orthonormal frame \(E_1\), define
\[
 E_{n+1}:=W_{P_{n+1},P_n}E_n.
\]
Then
\begin{equation}
 \|E_{n+1}-E_n\|\leq2\delta_n,                        \label{eq:n9v8-frame-step}
\end{equation}
so \(E_n\) converges in operator norm to an isometric frame \(E_\infty\).
\end{theorem}

\begin{proof}
Equation \eqref{eq:n9v8-frame-step} follows from
\eqref{eq:n9v8-W-bound}.  Thus \((E_n)\) is Cauchy in the Banach space
\(\mathcal B(\mathbb F^r,\mathcal H)\).  Its limit satisfies
\(E_\infty^*E_\infty=I_r\) by continuity.
\end{proof}

\subsection{A spectral-gap loader for projector motion}
\label{sec:n9v8-gap}

\begin{theorem}[Ordered-gap sine estimate]
\label{thm:n9v8-gap}
Let \(L,L'\) be bounded self-adjoint operators with spectral projections
\(P,Q\) onto finite low-energy clusters.  Suppose there are \(a\in\mathbb R\)
and \(\gamma>0\) such that
\begin{align}
 \sup\spec(L|_{\operatorname{Ran}P}),\ 
 \sup\spec(L'|_{\operatorname{Ran}Q})
 &\leq a,                                               \label{eq:n9v8-cross1}\\
 \inf\spec(L|_{\operatorname{Ran}(I-P)}),\ 
 \inf\spec(L'|_{\operatorname{Ran}(I-Q)})
 &\geq a+\gamma.                                       \label{eq:n9v8-cross2}
\end{align}
Then
\begin{equation}
 \|P-Q\|
 \leq{\|L'-L\|\over\gamma}.                           \label{eq:n9v8-gap}
\end{equation}
\end{theorem}

\begin{proof}
Set \(E=L'-L\).  The off-diagonal block
\(X=(I-Q)P\) satisfies
\[
 L'_{I-Q}X-XL_P=(I-Q)EP.
\]
For an equation \(BX-XA=Y\) with \(B\geq(a+\gamma)I\) and
\(A\leq aI\), the norm-convergent formula
\[
 X=\int_0^\infty e^{-tB}Ye^{tA}\,dt
\]
shows \(\|X\|\leq\|Y\|/\gamma\).  Hence
\(\|(I-Q)P\|\leq\|E\|/\gamma\).  Interchanging the two operators gives
\(\|(I-P)Q\|\leq\|E\|/\gamma\).  For orthogonal projections,
\[
 \|P-Q\|=\max\{\|(I-Q)P\|,\|(I-P)Q\|\}.
\]
This proves \eqref{eq:n9v8-gap}.
\end{proof}

\begin{corollary}[Summable operator-over-gap schedule]
\label{cor:n9v8-gap-sum}
If
\begin{equation}
 \sum_n{\|L_{n+1}-L_n\|\over\gamma_n}<\infty          \label{eq:n9v8-gap-sum}
\end{equation}
and the ratios are eventually at most \(1/2\), the polar frames satisfy
\eqref{eq:n9v8-sum}.
\end{corollary}

\begin{proof}
Apply \Cref{thm:n9v8-gap} at each adjacent step and then
\Cref{thm:n9v8-sum}.
\end{proof}

\subsection{Transport of the slow effective action and score}
\label{sec:n9v8-score}

Let \(K_R=\{s\in\mathbb F^r:|s|\leq R\}\).  In polar coordinates, define
the adjacent pulled-back functions
\[
 \widetilde F_{n+1}(s,y):=F_{n+1}(E_{n+1}s,y),\qquad
 \widetilde H_{n+1}(s,y):=H_{n+1}(E_{n+1}s,y).
\]

\begin{proposition}[Frame-motion debit]
\label{prop:n9v8-frame-debit}
Let \(G\) be \(C^1\) on the slow ball in \(\mathcal H\) and suppose
\(\|D_sG\|\leq L_G\).  If \(\delta_n\leq1/2\), then
\begin{equation}
 \sup_{|s|\leq R}
 |G(E_{n+1}s)-G(E_ns)|
 \leq2L_GR\delta_n.                                   \label{eq:n9v8-frame-debit}
\end{equation}
\end{proposition}

\begin{proof}
The mean-value theorem and \eqref{eq:n9v8-frame-step} give
\[
 |G(E_{n+1}s)-G(E_ns)|
 \leq L_G\|(E_{n+1}-E_n)s\|
 \leq2L_GR\delta_n.
\]
\end{proof}

\begin{proposition}[Score covariance under polar coordinates]
\label{prop:n9v8-score}
Assume the fast dominating measure does not depend on the choice of slow
frame.  In the common polar coordinate \(s\), the conditional Gibbs score
is
\begin{equation}
 \widetilde\sigma_{n,a}
 =
 -\partial_{s_a}\widetilde H_n
 +\mathbb E_{\widetilde\nu_{n,s}}
  \partial_{s_a}\widetilde H_n.                       \label{eq:n9v8-score}
\end{equation}
There is no additional basis Jacobian term.
\end{proposition}

\begin{proof}
\Cref{cor:n9v8-J} shows that the slow coordinate Jacobian is the constant
one.  Differentiate the pulled-back conditional density exactly as in
\Cref{lem:n9v7-score}.
\end{proof}

\begin{theorem}[Polar transport loader for SEPC]
\label{thm:n9v8-loader}
Suppose the slow spectral rank is fixed, the gap schedule
\eqref{eq:n9v8-gap-sum} holds, and the intrinsic adjacent changes of
\(F_n,H_n\) and their slow derivatives are summable after pullback to a
fixed frame.  If the frame derivative bounds obey
\begin{equation}
 \sum_nR_n(L_{F,n}+L_{H,n})\delta_n<\infty,            \label{eq:n9v8-loader}
\end{equation}
then the polar coordinate changes contribute a summable action/score
transport debit and zero logarithmic slow-Jacobian debit.  Thus slow-basis
motion does not obstruct SEPC.
\end{theorem}

\begin{proof}
The Jacobian statement is \Cref{cor:n9v8-J}.  Apply
\Cref{prop:n9v8-frame-debit} to \(F_n,H_n\) and to their declared
derivatives; add the intrinsic adjacent errors.  Their sums control the
functions entering the score formula \eqref{eq:n9v8-score} and hence the
SEPC derivative ledger.
\end{proof}

\begin{firewall}[Nonlocality and changing rank]
\label{fire:n9v8-rank}
Polar transport is finite-dimensional and need not be spatially local.
If \(\|P_{n+1}-P_n\|\geq1\), the ranks may differ or the subspaces may be
orthogonal; \(PQP\) then loses its floor and the construction fails.  Such
an index is a slow spectral crossing and must be assigned a finite-rank
crossing sector or a summable exceptional-set debit.
\end{firewall}

\begin{assessment}[Progress at ninth-series V8]
\label{ass:n9v8-status}
Matched slow spectral spaces now have a canonical adjacent identification.
Its coordinate Jacobian is exactly one, its frame motion is
\(O(\|P_{n+1}-P_n\|)\), and a cross-gap estimate turns operator convergence
into a summable frame schedule.  The result closes arbitrary slow-basis
motion inside SEPC while retaining spectral crossings as a separate
bottleneck.
\end{assessment}

\begin{programme}[Ninth-series V9: rank-change and crossing debit]
\label{prog:n9v8-V9}
V9 will analyze a simple slow eigenvalue crossing by a finite-dimensional
Lyapunov--Schmidt split.  It will prove the local determinant/coarea power,
derive the probability cost of a tubular crossing neighborhood, and state
the summability condition that allows the polar charts to cover its
complement.
\end{programme}

\begin{firewall}[Claim boundary after ninth-series V8]
\label{fire:n9v8-boundary}
V8 does not prove the physical spectral gap schedule, control rank-changing
crossings, remove the conformal cutoff, prove full density convergence or
tightness, or construct the SM--Einstein continuum.
\end{firewall}

\begin{theorem}[Ninth-series V8 result]
\label{thm:n9v8-result}
For matched finite-rank slow projectors with summable
\(\|P_{n+1}-P_n\|\), polar transport gives summable frame motion and unit
slow-coordinate Jacobian.  The cross-gap schedule
\eqref{eq:n9v8-gap-sum} is a sufficient operator-level loader, and
\eqref{eq:n9v8-loader} preserves the SEPC action/score budget.
\end{theorem}

\begin{proof}
Use \Cref{thm:n9v8-polar,cor:n9v8-J,
thm:n9v8-sum,thm:n9v8-gap,
cor:n9v8-gap-sum,thm:n9v8-loader}.
\end{proof}

\paragraph{Parent-version record.}
V8 was formed from
\texttt{Renewal\_SM\_Einstein\_Continuum\_Research\_Notes\_V7.tex}.
The V7 parent had \(2\,584\) logical source lines, \(93\,481\) bytes, and
SHA--256
\[
 \texttt{D00A8F1C3844C8AF27573B6A580833A213EF723DDC645E0F13334B88A3560570}.
\]
The next compulsory companion audit is V10.

% =============================================================================
% END APPENDED NINTH-SERIES V8 MATERIAL
% =============================================================================


% =============================================================================
% BEGIN APPENDED NINTH-SERIES V9 MATERIAL
% =============================================================================

\clearpage
\section{Ninth-series V9: a transverse slow-crossing debit}
\label{sec:v9v9}

\subsection{Uniform simple crossing chart}
\label{sec:n9v9-crossing}

Let \(K\subset\mathbb R^d\) be compact.  Let
\[
 A:K\times(-\rho,\rho)\longrightarrow
 \operatorname{Herm}(N)
\]
be \(C^2\).  Assume that for every \(x\in K\), \(A(x,0)\) has a
one-dimensional kernel with orthogonal projection \(P_0(x)\), and its
remaining spectrum satisfies
\begin{equation}
 \operatorname{dist}
 \bigl(0,\spec(A(x,0)|_{\operatorname{Ran}(I-P_0(x))})\bigr)
 \geq\gamma>0.                                        \label{eq:n9v9-gap}
\end{equation}
Assume \(P_0\) is continuous on \(K\).  For a unit vector
\(\psi_x\in\operatorname{Ran}P_0(x)\), suppose
\begin{equation}
 \left|
 \langle\psi_x,\partial_tA(x,0)\psi_x\rangle
 \right|\geq a_0>0.                                   \label{eq:n9v9-transverse}
\end{equation}
The expression is independent of the phase of \(\psi_x\).

\begin{theorem}[Uniform simple-eigenvalue factorization]
\label{thm:n9v9-factor}
After decreasing \(\rho\), there are a \(C^1\) eigenvalue
\(\lambda(x,t)\), a rank-one spectral projection \(P(x,t)\), and a
nonvanishing complementary determinant \(D_\perp(x,t)\) such that
\begin{align}
 A(x,t)P(x,t)&=\lambda(x,t)P(x,t),                    \label{eq:n9v9-eigen}\\
 \det A(x,t)&=\lambda(x,t)D_\perp(x,t),               \label{eq:n9v9-det}\\
 0<c_\perp\leq|D_\perp(x,t)|&\leq C_\perp,            \label{eq:n9v9-Dperp}\\
 {a_0\over2}|t|
 \leq|\lambda(x,t)|
 &\leq a_1|t|                                         \label{eq:n9v9-linear}
\end{align}
for uniform constants \(c_\perp,C_\perp,a_1>0\).
\end{theorem}

\begin{proof}
Choose a circle of radius less than \(\gamma/2\) about zero.  Uniform
continuity of \(A\) and compactness of \(K\) keep this circle in the
resolvent set for all sufficiently small \(|t|\).  The Riesz projection
\[
 P(x,t)={1\over2\pi i}\int_\Gamma
 (z-A(x,t))^{-1}\,dz
\]
is \(C^1\), has rank one, and selects a unique eigenvalue
\(\lambda(x,t)=\Tr(A(x,t)P(x,t))\).  The complementary determinant is the
determinant of the compression to \(\operatorname{Ran}(I-P(x,t))\), after
unitary transport to the \(t=0\) complement.  Its factors remain at least
\(\gamma/2\) from zero and are uniformly bounded above, proving
\eqref{eq:n9v9-Dperp} and the factorization
\eqref{eq:n9v9-det}.

The Feynman--Hellmann formula gives
\[
 \partial_t\lambda(x,0)
 =
 \langle\psi_x,\partial_tA(x,0)\psi_x\rangle.
\]
Its absolute value is at least \(a_0\).  By compactness and continuity,
\(|\partial_t\lambda(x,t)|\geq a_0/2\) with constant sign on each connected
crossing chart after decreasing \(\rho\).  The mean-value theorem gives the
lower bound in \eqref{eq:n9v9-linear}; a uniform upper bound for
\(|\partial_t\lambda|\) gives \(a_1\).
\end{proof}

\begin{corollary}[Normal-square crossing]
\label{cor:n9v9-square}
If \(B(x,t)\) has one singular value
\(\sigma(x,t)\asymp|t|\) and all complementary singular values have a
uniform positive floor, then
\[
 \det(B^*B)
 =\sigma(x,t)^2D_{\perp}^{\,2}(x,t)
 \asymp |t|^2.                                        \label{eq:n9v9-square}
\]
\end{corollary}

\begin{proof}
The eigenvalues of \(B^*B\) are the squared singular values of \(B\).
Separate the simple small singular value from the uniformly invertible
complement.
\end{proof}

\subsection{Determinant--coarea power balance}
\label{sec:n9v9-power}

Let a local unnormalized density on the crossing chart have the form
\begin{equation}
 w(x,t)
 =
 h(x,t)\,
 {|\lambda(x,t)|^p\over J(x,t)},                      \label{eq:n9v9-w}
\end{equation}
where \(p\geq0\) is the net determinant power retained in this row.  Assume
\begin{equation}
 0\leq h(x,t)\leq H,\qquad
 J(x,t)\geq j_0|t|^q                                 \label{eq:n9v9-J}
\end{equation}
for \(q\geq0\).  This covers a coarea loss of order \(q\); it does not
assert that the same physical eigenvalue controls both factors.

\begin{theorem}[Crossing-tube mass]
\label{thm:n9v9-tube}
Suppose \(p-q>-1\).  Let
\[
 \mathcal T_\varepsilon:=K\times(-\varepsilon,\varepsilon)
\]
with \(0<\varepsilon<\rho\).  If the full chart partition function obeys
\begin{equation}
 Z_{\rm chart}:=\int_{K\times(-\rho,\rho)}w(x,t)\,dx\,dt
 \geq z_0>0,                                          \label{eq:n9v9-Zfloor}
\end{equation}
then the normalized chart probability satisfies
\begin{equation}
 \mu(\mathcal T_\varepsilon)
 \leq
 {2H\,a_1^p\,\operatorname{Vol}(K)\over
 z_0j_0(p-q+1)}
 \varepsilon^{p-q+1}.                                 \label{eq:n9v9-tube}
\end{equation}
\end{theorem}

\begin{proof}
The upper bound in \eqref{eq:n9v9-linear} and
\eqref{eq:n9v9-J} give
\[
 w(x,t)\leq{Ha_1^p\over j_0}|t|^{p-q}.
\]
Because \(p-q>-1\), this power is integrable at zero.  Integrating over
\(K\times(-\varepsilon,\varepsilon)\) gives
\[
 \int_{\mathcal T_\varepsilon}w
 \leq
 {2Ha_1^p\operatorname{Vol}(K)\over j_0(p-q+1)}
 \varepsilon^{p-q+1}.
\]
Divide by \eqref{eq:n9v9-Zfloor}.
\end{proof}

\begin{firewall}[The partition floor is essential]
\label{fire:n9v9-Z}
An unnormalized tube integral does not bound probability unless the
normalizing partition function has a lower bound in the same chart or the
estimate is inserted into a globally normalized partition.  V9 records
\(z_0\) rather than silently treating it as one.
\end{firewall}

\begin{theorem}[Power-balance trichotomy]
\label{thm:n9v9-trichotomy}
For the local model \eqref{eq:n9v9-w}:
\begin{enumerate}[label=\textup{(\roman*)},leftmargin=3.8em]
\item if \(p-q>-1\), the crossing is locally integrable and its tube mass
has the power in \eqref{eq:n9v9-tube};
\item if \(p-q=-1\), the model upper envelope has a logarithmic
nonintegrability and no vanishing tube debit follows from these bounds; and
\item if \(p-q<-1\), the model power is nonintegrable at zero.
\end{enumerate}
\end{theorem}

\begin{proof}
This is the elementary integral criterion for
\(\int_0^\varepsilon t^{p-q}\,dt\).  In cases (ii)--(iii), an additional
cancellation or a sharper numerator estimate may still change the actual
density, but it is absent from the stated hypotheses.
\end{proof}

\begin{corollary}[Simple balances]
\label{cor:n9v9-balances}
A simple determinant zero with no coarea pole has \((p,q)=(1,0)\) and tube
mass \(O(\varepsilon^2)\).  A simple determinant zero cancelling a simple
coarea loss has \((p,q)=(1,1)\) and tube mass \(O(\varepsilon)\).  A normal
square zero against a simple coarea loss has \((p,q)=(2,1)\) and tube mass
\(O(\varepsilon^2)\).
\end{corollary}

\begin{proof}
Insert the three pairs into \(p-q+1\).
\end{proof}

\subsection{Summable exceptional neighborhoods}
\label{sec:n9v9-sum}

\begin{corollary}[Cofinal tube schedule]
\label{cor:n9v9-sum}
For adjacent measures \(\mu_n,\mu_{n+1}\) whose crossing-chart constants
are uniform and whose common exponent
\(\alpha:=p-q+1\) is positive, suppose
\begin{equation}
 \sum_n\varepsilon_n^\alpha<\infty.                   \label{eq:n9v9-eps-sum}
\end{equation}
Then
\begin{equation}
 \sum_n\left[
 \mu_n(\mathcal T_{\varepsilon_n})
 +\mu_{n+1}(\mathcal T_{\varepsilon_n})
 \right]<\infty.                                      \label{eq:n9v9-exceptional}
\end{equation}
\end{corollary}

\begin{proof}
Apply \Cref{thm:n9v9-tube} to both adjacent measures and sum.
\end{proof}

\begin{corollary}[Polynomial and geometric choices]
\label{cor:n9v9-choices}
The choice \(\varepsilon_n=(n+1)^{-r}\) satisfies
\eqref{eq:n9v9-eps-sum} when \(r\alpha>1\).  The choice
\(\varepsilon_n=2^{-n}\) satisfies it for every \(\alpha>0\).
\end{corollary}

\begin{proof}
Use the \(p\)-series and geometric-series criteria.
\end{proof}

\subsection{Regular polar charts outside the tube}
\label{sec:n9v9-polar}

\begin{proposition}[Crossing excision plus polar transport]
\label{prop:n9v9-polar}
Suppose that outside \(\mathcal T_{\varepsilon_n}\), the selected slow
projectors have fixed rank, adjacent distance at most \(1/2\), and satisfy
the summable polar-transport schedule of V8.  If
\eqref{eq:n9v9-exceptional} holds, then the regular complement contributes
the V8 summable transport/Jacobian row, while the crossing tubes contribute
a separate summable adjacent \(L^1\) debit.
\end{proposition}

\begin{proof}
Use \Cref{thm:n9v8-loader} on the regular complement.  Restricting the
difference of two adjacent probability measures to the exceptional tube
costs at most the sum of their probabilities there, which is summable by
\eqref{eq:n9v9-exceptional}.  Add the two debits.
\end{proof}

\begin{definition}[Transverse slow-crossing card, TSCC]
\label{def:n9v9-TSCC}
TSCC records:
\begin{enumerate}[label=\textup{(T\arabic*)},leftmargin=4em]
\item a finite-dimensional crossing chart and a complementary spectral
gap;
\item the nonzero Feynman--Hellmann derivative \(a_0\);
\item the determinant zero order \(p\) after ownership assignments;
\item the coarea loss order \(q\), proved for the same chart;
\item uniform bounds \(H,j_0,z_0\);
\item the integrability inequality \(p-q>-1\);
\item a tube schedule satisfying \eqref{eq:n9v9-eps-sum}; and
\item fixed-rank polar transport on the regular complement.
\end{enumerate}
\end{definition}

\begin{theorem}[Conditional crossing closure]
\label{thm:n9v9-closure}
TSCC implies a summable exceptional-set debit for the slow crossing and a
summable regular slow-chart transport row on its complement.
\end{theorem}

\begin{proof}
The tube conclusion is \Cref{cor:n9v9-sum}; the complement conclusion is
\Cref{prop:n9v9-polar}.
\end{proof}

\begin{firewall}[No automatic equality of determinant and coarea modes]
\label{fire:n9v9-mode}
The exponent balance is useful only after proving that the local factors
have the stated orders in one common transverse coordinate.  A zero mode of
the constraint normal operator, a zero of the chiral determinant, and a
loss of multiplier transversality are different objects in general.  Their
powers cannot be cancelled by matching names or dimensions.
\end{firewall}

\begin{assessment}[Progress at ninth-series V9]
\label{ass:n9v9-status}
A simple slow crossing now has a rigorous finite-dimensional loader.
Lyapunov--Schmidt/Riesz reduction isolates its linear eigenvalue, and the
net determinant--coarea exponent gives the exact tube probability power.
When that power is positive, the crossing can be charged as a summable
exceptional debit while polar transport handles the regular complement.
\end{assessment}

\begin{programme}[Ninth-series V10: compulsory companion audit]
\label{prog:n9v9-V10}
V10 will inspect the newest YM and NS notes, test whether they supply a
usable crossing, slow-representation, or tail estimate, and then choose
between verifying a physical TSCC instance and returning to the remaining
ambient-transport/tightness rows.  The selected branch will continue in
V11.
\end{programme}

\begin{firewall}[Claim boundary after ninth-series V9]
\label{fire:n9v9-boundary}
V9 does not prove TSCC for any physical SM--Einstein crossing, identify the
net physical exponent, remove the conformal cutoff, prove full density
convergence or tightness, or construct the continuum.
\end{firewall}

\begin{theorem}[Ninth-series V9 result]
\label{thm:n9v9-result}
A uniform simple crossing with determinant order \(p\), coarea loss order
\(q\), and partition floor has tube probability
\(O(\varepsilon^{p-q+1})\) when \(p-q>-1\).  A summable tube schedule and
V8 polar transport conditionally close the crossing and regular slow-chart
debits.
\end{theorem}

\begin{proof}
Use \Cref{thm:n9v9-factor,thm:n9v9-tube,
thm:n9v9-trichotomy,cor:n9v9-sum,
prop:n9v9-polar,thm:n9v9-closure}.
\end{proof}

\paragraph{Parent-version record.}
V9 was formed from
\texttt{Renewal\_SM\_Einstein\_Continuum\_Research\_Notes\_V8.tex}.
The V8 parent had \(2\,946\) logical source lines, \(105\,576\) bytes, and
SHA--256
\[
 \texttt{F3EA316EA4ADE97943C219A1B7ACC7FBBF16E3D54635631C31B4939484D5FFAB}.
\]
The next version is the compulsory V10 companion audit.

% =============================================================================
% END APPENDED NINTH-SERIES V9 MATERIAL
% =============================================================================


% =============================================================================
% BEGIN APPENDED NINTH-SERIES V10 MATERIAL
% =============================================================================

\clearpage
\section{Ninth-series V10: compulsory audit and the tightness branch}
\label{sec:v9v10}

\subsection{Files inspected}
\label{sec:n9v10-files}

The newest Yang--Mills files inspected were
\begin{center}
\small\ttfamily
Renewal\_Yang\_Mills\_Mass\_Gap\_Research\_Notes\_V57.tex,\\
Renewal\_Yang\_Mills\_Mass\_Gap\_Research\_Notes\_V58.tex.
\end{center}
They are distinct.  V57 has \(24\,303\) logical lines, \(868\,452\) bytes,
and SHA--256
\[
 \texttt{53E2B4F4EF02FC1A6ADCA56374B801B4B88D1C575905565B9482322575387CC9}.
\]
V58 has \(24\,741\) logical lines, \(883\,036\) bytes, and SHA--256
\[
 \texttt{A515FFEBF73F899269597201B576AF07D5920FBEEC45548AAD8D791407FDF4FC}.
\]
The Navier--Stokes file remains V26, with \(5\,319\) logical lines,
\(278\,283\) bytes, and SHA--256
\[
 \texttt{82C887F8C2C69E4F28FAD7C3DC8DE5B9099CB5A4BF431EBA1FFF3E91935978AD}.
\]

\subsection{Yang--Mills V57--V58}
\label{sec:n9v10-YM}

YM V57 proves Weyl inversion and Plancherel formulas, converts phase-space
powers to commutators, bounds weighted \(L^2\) characteristic moments, and
obtains absolute moment bounds for occupation-window matrix units.  A
degree-weighted labelled-tree count then closes the connected operator
series under the logarithmic occupation schedule.  It explicitly rejects
a stronger geometric-moment target that is false.

YM V58 constructs a dense finite-cylinder Weyl endpoint core in the
reference Fock space, proves that fixed Weyl endpoints pay no occupation
loss, and obtains uniform reference clustering on that core.  It also
proves stability of the quadratic reference edge under a weighted
self-energy perturbation and an integrated semigroup-response estimate.
Insertion of the actual occupation self-energy remains conditional on a
quadratic self-energy decomposition card; interacting OS density is not
proved.

\begin{proposition}[Type-correct consequence for the SM programme]
\label{prop:n9v10-YM}
The YM V57--V58 theorems do not verify TSCC, SLDAC, or the physical
SM--Einstein action row.  They support the following independent
acquisition: test the already constructed projective cylinder candidate on
a dense Weyl/cylinder core and prove tightness from regulator-uniform
field-space moments before attempting reconstruction.
\end{proposition}

\begin{proof}
The YM bounds use a quasi-free reference and an occupation representation
interface not supplied by the SM constrained measure, so no numerical
estimate transfers.  Their endpoint construction nevertheless shows the
correct logical order: a dense cylinder core is useful only together with
uniform control that survives completion.  In the SM programme, the
corresponding missing completion statement is tightness in a declared
distributional topology.
\end{proof}

\begin{firewall}[Reference clustering is not interacting tightness]
\label{fire:n9v10-YM}
Uniform clustering of the YM quasi-free reference on its Weyl core does not
imply tightness, clustering, or reflection positivity of the interacting
SM--Einstein measure.  The self-energy insertion needed even in the YM
programme remains conditional.
\end{firewall}

\subsection{Navier--Stokes status}
\label{sec:n9v10-NS}

The NS note has not changed since the V5 audit.  Its rank-one Oseen-kernel
constants and ancient-profile neutral mode neither supply a field-space
moment for the SM measure nor alter the crossing exponent of V9.

\begin{proposition}[NS non-transfer at V10]
\label{prop:n9v10-NS}
No current NS result loads the SM tightness row.
\end{proposition}

\begin{proof}
The NS estimates concern deterministic parabolic kernels and profile
constraints.  The SM tightness problem requires probability moments of
metric, connection, Higgs, and fermion-distribution coordinates under the
constrained Euclidean measures.  No identification between these objects
is present.
\end{proof}

\subsection{Why cylinder convergence is insufficient}
\label{sec:n9v10-example}

\begin{proposition}[Cylinder convergence without Hilbert-space tightness]
\label{prop:n9v10-example}
Let \((e_n)\) be the standard orthonormal basis of \(\ell^2\) and let
\(\mu_n=\delta_{e_n}\).  Every fixed finite-coordinate marginal converges
to the corresponding marginal of \(\delta_0\), but \((\mu_n)\) is not
tight on \(\ell^2\) with its norm topology.
\end{proposition}

\begin{proof}
For a projection onto the first \(k\) coordinates, the image of \(e_n\) is
zero when \(n>k\), proving finite-coordinate convergence.  If the measures
were tight, a compact set \(K\) would satisfy
\(\delta_{e_n}(K)\geq1/2\) for every \(n\); hence every \(e_n\) would lie in
\(K\).  This is impossible because
\(\|e_n-e_m\|=\sqrt2\) for \(n\ne m\), while every sequence in a compact
metric set has a convergent subsequence.
\end{proof}

\begin{corollary}[A new topology-specific estimate is necessary]
\label{cor:n9v10-necessary}
The total-variation convergence of every fixed local cylinder marginal in
\Cref{thm:n9v1-dependency} does not by itself imply tightness in a chosen
infinite-dimensional field topology.
\end{corollary}

\begin{proof}
\Cref{prop:n9v10-example} has eventual equality of every fixed coordinate
marginal and still lacks tightness.
\end{proof}

\subsection{Audit decision}
\label{sec:n9v10-decision}

\begin{audit}[V10 companion conclusion]
\label{aud:n9v10}
The action, slow-label, and crossing cards V2--V9 remain the correct local
density architecture, but their physical loaders are not supplied by the
companions.  The next noncircular acquisition is the topology upgrade:
derive a compact-containment criterion from stronger negative-Sobolev
moments, first for one bosonic field and then for the product of metric,
gauge, and Higgs fields.
\end{audit}

\begin{decision}[V11 acquisition order]
\label{dec:n9v10-V11}
V11 will prove that a uniform \(H^{-s_0}\) moment gives tightness in
\(H^{-s}\) for \(s>s_0\) on a compact base, using compact Sobolev
embedding and Markov's inequality.  It will give a mode-by-mode loader,
finite-product assembly, and a compatibility theorem upgrading the
projective cylinder candidate to a Radon probability on the declared
negative-Sobolev product space.
\end{decision}

\begin{assessment}[Progress at ninth-series V10]
\label{ass:n9v10-status}
The audit prevents a false completion from local total-variation limits.
The next theorem will identify the exact extra derivative of compactness
needed for a global field-space measure.  The action and crossing cards
remain active hypotheses and will later have to be combined with this
tightness row.
\end{assessment}

\begin{programme}[Ninth-series V11: negative-Sobolev tightness]
\label{prog:n9v10-V11}
Carry out \Cref{dec:n9v10-V11}, retaining separate regularity indices and
moment budgets for every bosonic field species.  Fermionic Grassmann data
will remain in the observable algebra rather than being treated as an
ordinary positive random distribution.
\end{programme}

\begin{firewall}[Claim boundary after ninth-series V10]
\label{fire:n9v10-boundary}
V10 is an audit.  It does not prove the required moments, tightness, the
physical density cards, reconstruction, or the SM--Einstein continuum.
\end{firewall}

\begin{theorem}[Ninth-series V10 result]
\label{thm:n9v10-result}
The current companion work motivates a dense cylinder endpoint core but
does not load the SM measure.  Cylinder convergence alone can fail
Hilbert-space tightness, so the audited next gate is a uniform
negative-Sobolev moment and compact-embedding theorem.
\end{theorem}

\begin{proof}
Use \Cref{prop:n9v10-YM,prop:n9v10-NS,
prop:n9v10-example,cor:n9v10-necessary,aud:n9v10}.
\end{proof}

\paragraph{Parent-version record.}
V10 was formed from
\texttt{Renewal\_SM\_Einstein\_Continuum\_Research\_Notes\_V9.tex}.
The V9 parent had \(3\,284\) logical source lines, \(117\,948\) bytes, and
SHA--256
\[
 \texttt{482DE9B8EE5C4B3C1CA3A7250D7CFE8885B8D0D3D04F06C61F5110DE0D2A2722}.
\]
The next compulsory companion audit is V15.

% =============================================================================
% END APPENDED NINTH-SERIES V10 MATERIAL
% =============================================================================


% =============================================================================
% BEGIN APPENDED NINTH-SERIES V11 MATERIAL
% =============================================================================

\clearpage
\section{Ninth-series V11: negative-Sobolev tightness and Radon landing}
\label{sec:v9v11}

\subsection{Sobolev scale on a compact base}
\label{sec:n9v11-scale}

Let \(M\) be a compact \(d\)-dimensional Riemannian manifold and \(E\to M\)
a finite-rank Euclidean or Hermitian bundle.  Choose a positive
self-adjoint Laplace-type operator \(\Delta_E\) with eigenpairs
\((\lambda_k,e_k)\).  For \(t\in\mathbb R\), use the Hilbert norm
\begin{equation}
 \|u\|_{H^t(E)}^2
 :=
 \sum_{k=1}^{\infty}(1+\lambda_k)^t
 |\langle u,e_k\rangle|^2.                            \label{eq:n9v11-Hs}
\end{equation}
Every finite-regulator field is assumed to have been embedded measurably
into this common distribution space before a tightness claim is made.

\begin{lemma}[Compact loss of Sobolev order]
\label{lem:n9v11-compact}
If \(s>s_0\), the inclusion
\begin{equation}
 H^{-s_0}(E)\hookrightarrow H^{-s}(E)                 \label{eq:n9v11-embedding}
\end{equation}
is compact.
\end{lemma}

\begin{proof}
In the eigenbasis, conjugate the inclusion by the isometries from the two
Sobolev spaces to \(\ell^2\).  It becomes the diagonal operator with
singular values
\[
 (1+\lambda_k)^{-(s-s_0)/2}.
\]
They tend to zero because \(\lambda_k\to\infty\).  A diagonal operator on
\(\ell^2\) whose diagonal tends to zero is the operator-norm limit of its
finite-rank truncations, hence compact.
\end{proof}

\begin{firewall}[The regularity loss is necessary for this argument]
\label{fire:n9v11-loss}
The closed unit ball of an infinite-dimensional Hilbert space is not
compact in its own norm topology.  A uniform \(H^{-s_0}\) moment therefore
does not prove tightness in \(H^{-s_0}\) by the ball argument.  V11 uses the
strict loss \(s>s_0\).
\end{firewall}

\subsection{Moment-to-tightness theorem}
\label{sec:n9v11-moment}

\begin{theorem}[Negative-Sobolev moment criterion]
\label{thm:n9v11-tight}
Let \(\mu_n\) be Borel probabilities on \(H^{-s}(E)\).  Suppose, for some
\(p>0\), \(s>s_0\), and \(C<\infty\),
\begin{equation}
 \sup_n\int\|u\|_{H^{-s_0}(E)}^p\,\mu_n(du)\leq C.     \label{eq:n9v11-moment}
\end{equation}
Then \((\mu_n)\) is tight on \(H^{-s}(E)\).  More precisely, if
\(K_R\) is the closure in \(H^{-s}\) of the \(H^{-s_0}\) ball of radius
\(R\), then \(K_R\) is compact and
\begin{equation}
 \sup_n\mu_n(K_R^c)\leq {C\over R^p}.                 \label{eq:n9v11-Markov}
\end{equation}
\end{theorem}

\begin{proof}
\Cref{lem:n9v11-compact} shows that \(K_R\) is compact in \(H^{-s}\).
Markov's inequality gives
\[
 \mu_n(\|u\|_{H^{-s_0}}>R)
 \leq R^{-p}\int\|u\|_{H^{-s_0}}^p\,d\mu_n.
\]
The \(H^{-s_0}\) ball is contained in \(K_R\), so its complement contains
\(K_R^c\).  This proves \eqref{eq:n9v11-Markov}.  Given
\(\varepsilon>0\), choose \(R>(C/\varepsilon)^{1/p}\).
\end{proof}

\subsection{Mode-by-mode loaders}
\label{sec:n9v11-modes}

\begin{proposition}[Exact second-moment mode loader]
\label{prop:n9v11-modes}
Suppose \(u\) is a random distribution whose spectral coordinates are
square integrable and
\begin{equation}
 \sum_{k=1}^{\infty}
 (1+\lambda_k)^{-s_0}
 \mathbb E|\langle u,e_k\rangle|^2
 \leq C.                                               \label{eq:n9v11-mode-sum}
\end{equation}
Then
\begin{equation}
 \mathbb E\|u\|_{H^{-s_0}}^2\leq C.                  \label{eq:n9v11-mode-moment}
\end{equation}
\end{proposition}

\begin{proof}
Apply monotone convergence to the nonnegative partial sums in
\eqref{eq:n9v11-Hs}.
\end{proof}

\begin{corollary}[Weyl-counted polynomial covariance]
\label{cor:n9v11-Weyl}
Assume
\begin{equation}
 \mathbb E|\langle u_n,e_k\rangle|^2
 \leq C_0(1+\lambda_k)^\alpha                         \label{eq:n9v11-cov}
\end{equation}
uniformly in \(n,k\).  If
\begin{equation}
 s_0>\alpha+{d\over2},                                \label{eq:n9v11-index}
\end{equation}
then \eqref{eq:n9v11-moment} holds with \(p=2\).
\end{corollary}

\begin{proof}
Weyl's eigenvalue bound on a compact \(d\)-manifold gives
\(\lambda_k\geq c k^{2/d}-C_1\) for large \(k\).  Hence
\[
 \sum_k(1+\lambda_k)^{\alpha-s_0}
\]
converges precisely under the sufficient inequality
\((2/d)(s_0-\alpha)>1\).  Use
\Cref{prop:n9v11-modes}.
\end{proof}

\begin{remark}[Means and covariances]
\label{rem:n9v11-mean}
The loader uses the full second moment, not only the centered covariance.
A diverging deterministic mean can destroy tightness even when every
centered covariance is zero.  One may split
\(\mathbb E|u_k|^2=|\mathbb Eu_k|^2+\operatorname{Var}(u_k)\), but both
sums must be controlled.
\end{remark}

\subsection{Several bosonic species}
\label{sec:n9v11-product}

\begin{theorem}[Finite-product tightness]
\label{thm:n9v11-product}
Let
\[
 \mathcal X:=\prod_{a=1}^JH^{-s_a}(E_a),
\]
with \(J<\infty\).  Suppose that for each species \(a\), some
\(s_a>s_{0,a}\), \(p_a>0\), and \(C_a<\infty\) satisfy
\begin{equation}
 \sup_n\int
 \|u_a\|_{H^{-s_{0,a}}(E_a)}^{p_a}\,d\mu_n
 \leq C_a.                                             \label{eq:n9v11-product-moment}
\end{equation}
Then \((\mu_n)\) is tight on \(\mathcal X\), without any independence
assumption between the species.
\end{theorem}

\begin{proof}
Given \(\varepsilon>0\), use \Cref{thm:n9v11-tight} to choose compact
\(K_a\subset H^{-s_a}(E_a)\) such that
\[
 \sup_n\mu_n(u_a\notin K_a)<{\varepsilon\over J}.
\]
The finite product \(K=\prod_aK_a\) is compact.  The union bound gives
\[
 \sup_n\mu_n(K^c)
 \leq\sum_a\sup_n\mu_n(u_a\notin K_a)<\varepsilon.
\]
\end{proof}

\begin{definition}[Bosonic negative-Sobolev moment card, BNSMC]
\label{def:n9v11-BNSMC}
For a compact physical base, BNSMC records:
\begin{enumerate}[label=\textup{(B\arabic*)},leftmargin=4em]
\item common measurable embeddings of the metric-coordinate, connection,
and Higgs regulators into fixed tensor-bundle distribution spaces;
\item for each species, indices \(s_a>s_{0,a}\) and a positive moment order
\(p_a\);
\item a uniform bound \eqref{eq:n9v11-product-moment}, proved under the full
normalized constrained measure;
\item compatibility of local trivializations and gauge/coordinate chart
transports; and
\item a separate support statement describing which generalized metrics
and connections the limiting topology permits.
\end{enumerate}
\end{definition}

\begin{corollary}[BNSMC gives bosonic tightness]
\label{cor:n9v11-BNSMC}
BNSMC implies tightness of the bosonic constrained measures on the declared
finite product \(\mathcal X\).
\end{corollary}

\begin{proof}
Apply \Cref{thm:n9v11-product}.
\end{proof}

\begin{firewall}[Fermions are not positive random coordinates]
\label{fire:n9v11-fermion}
Grassmann fields are not ordinary random variables in a positive Borel
probability space.  V11 places the bosonic variables in \(\mathcal X\);
fermionic information remains in determinant sections and the graded
observable algebra.  Tightness of bosonic measures does not construct the
fermionic continuum field.
\end{firewall}

\subsection{From the cylinder candidate to a Radon measure}
\label{sec:n9v11-Radon}

\begin{theorem}[Tight cylinder-to-Radon landing]
\label{thm:n9v11-Radon}
Let \(\mathcal X\) be the separable Hilbert product in
\Cref{thm:n9v11-product}.  Suppose:
\begin{enumerate}[label=\textup{(\roman*)},leftmargin=3.8em]
\item \(\mu_n\) are Borel probabilities on \(\mathcal X\) and are tight;
\item for a countable total family of smooth test sections, every finite
collection of pairings has a limiting joint law; and
\item those finite-dimensional laws are projectively consistent.
\end{enumerate}
Then there is a unique Borel probability \(\mu\) on \(\mathcal X\) with
those cylinder laws, and
\begin{equation}
 \mu_n\Longrightarrow\mu\quad\text{weakly on }\mathcal X.             \label{eq:n9v11-weak}
\end{equation}
The measure \(\mu\) is Radon.
\end{theorem}

\begin{proof}
By Prokhorov's theorem, every subsequence has a weakly convergent
subsequence, say to \(\widetilde\mu\).  Pairing with a smooth test section
is a continuous linear functional on each negative Sobolev factor.
Therefore the continuous mapping theorem identifies every prescribed
finite-dimensional marginal of \(\widetilde\mu\) with the given cylinder
limit.

Choose the total test family so that its linear span is dense in the dual
coordinate family.  In a separable Hilbert space, the Borel sigma-algebra
is generated by the coordinates of one countable orthonormal basis.
Hence two Borel probabilities with all these finite-dimensional marginals
are equal.  Every subsequential weak limit is therefore the same measure
\(\mu\).  Tightness then implies convergence of the whole sequence:
otherwise a subsequence staying outside a weak neighborhood of \(\mu\)
would have a further subsequence converging to \(\mu\), a contradiction.
Every Borel probability on the Polish space \(\mathcal X\) is Radon.
\end{proof}

\begin{corollary}[Upgrade of the audited dependency theorem]
\label{cor:n9v11-upgrade}
If the hypotheses of \Cref{thm:n9v1-dependency} hold, its cylinder limits
are realized by measures \(\mu_n\) on \(\mathcal X\), and BNSMC holds,
then the bosonic cylinder candidate extends uniquely to the Radon
probability of \Cref{thm:n9v11-Radon}.
\end{corollary}

\begin{proof}
The audited dependency theorem supplies the consistent finite-dimensional
limits.  BNSMC supplies tightness.  Apply
\Cref{thm:n9v11-Radon}.
\end{proof}

\begin{firewall}[Support and nonlinear geometry remain separate]
\label{fire:n9v11-support}
A Radon measure on a negative-Sobolev space may be supported on
distributions that are not classical metrics or connections.  Positivity
of a metric, invertibility, bundle topology, and the nonlinear Einstein
equations require closed-support or capacity estimates compatible with the
chosen topology.  Tightness alone proves none of them.
\end{firewall}

\begin{assessment}[Progress at ninth-series V11]
\label{ass:n9v11-status}
The cylinder-to-field-space gap now has a precise loader.  One stronger
negative-Sobolev moment per bosonic species supplies compact containment
after an arbitrarily small loss of regularity.  Together with the existing
projective cylinder limits, this gives a unique Radon probability on the
declared product distribution space.  The moment bounds must still be
proved under the full constrained density.
\end{assessment}

\begin{programme}[Ninth-series V12: transport moments through the density]
\label{prog:n9v11-V12}
V12 will prove a Hölder/entropy loader that transfers reference
negative-Sobolev moments through the normalized action, determinant,
coarea, and transport density.  It will identify the required \(L^q\)
bound for the full Radon--Nikodym derivative and show why local
total-variation convergence alone does not preserve an unbounded Sobolev
moment.
\end{programme}

\begin{firewall}[Claim boundary after ninth-series V11]
\label{fire:n9v11-boundary}
V11 does not prove BNSMC for the physical measures, construct the
fermionic field, establish nonlinear geometric support, reconstruction, or
the full SM--Einstein continuum.
\end{firewall}

\begin{theorem}[Ninth-series V11 result]
\label{thm:n9v11-result}
Uniform \(H^{-s_0}\) moments imply tightness in every \(H^{-s}\) with
\(s>s_0\).  The mode loader \eqref{eq:n9v11-mode-sum} and finite-product
assembly give BNSMC, and consistent cylinder convergence plus BNSMC yields
a unique Radon bosonic probability on the declared negative-Sobolev product
space.
\end{theorem}

\begin{proof}
Use \Cref{lem:n9v11-compact,thm:n9v11-tight,
prop:n9v11-modes,thm:n9v11-product,
thm:n9v11-Radon,cor:n9v11-upgrade}.
\end{proof}

\paragraph{Parent-version record.}
V11 was formed from
\texttt{Renewal\_SM\_Einstein\_Continuum\_Research\_Notes\_V10.tex}.
The V10 parent had \(3\,489\) logical source lines, \(126\,297\) bytes, and
SHA--256
\[
 \texttt{A0614F08CF59A87ECD06197DE9EB336044964E33B4E4973C70867676B1631FF3}.
\]
The next compulsory companion audit is V15.

% =============================================================================
% END APPENDED NINTH-SERIES V11 MATERIAL
% =============================================================================


% =============================================================================
% BEGIN APPENDED NINTH-SERIES V12 MATERIAL
% =============================================================================

\clearpage
\section{Ninth-series V12: transporting moments through the full density}
\label{sec:v9v12}

\subsection{The Hölder loader}
\label{sec:n9v12-Holder}

Let \(\nu\) be a reference probability, let \(\mu\ll\nu\), and write
\[
 r:={d\mu\over d\nu}.
\]

\begin{theorem}[Density-to-moment transfer]
\label{thm:n9v12-Holder}
Let \(q>1\), \(q'=q/(q-1)\), \(p>0\), and let \(X\geq0\) be measurable.
Then
\begin{equation}
 \int X^p\,d\mu
 \leq
 \|r\|_{L^q(\nu)}
 \left(\int X^{pq'}\,d\nu\right)^{1/q'}.              \label{eq:n9v12-Holder}
\end{equation}
Consequently, uniform \(L^q\) control of the full density and a uniform
reference \(pq'\)-moment imply a uniform target \(p\)-moment.
\end{theorem}

\begin{proof}
Since \(d\mu=r\,d\nu\), Holder's inequality gives
\[
 \int X^pr\,d\nu
 \leq
 \left(\int r^q\,d\nu\right)^{1/q}
 \left(\int X^{pq'}\,d\nu\right)^{1/q'}.
\]
\end{proof}

\begin{corollary}[The \(L^2\)--fourth-moment route]
\label{cor:n9v12-L2}
If
\[
 \sup_n\|r_n\|_{L^2(\nu_n)}\leq R,\qquad
 \sup_n\int\|u\|_{H^{-s_0}}^{2p}\,d\nu_n\leq C,
\]
then
\begin{equation}
 \sup_n\int\|u\|_{H^{-s_0}}^p\,d\mu_n
 \leq R\sqrt C.                                       \label{eq:n9v12-L2}
\end{equation}
\end{corollary}

\begin{proof}
Take \(q=q'=2\) and
\(X=\|u\|_{H^{-s_0}}\) in \Cref{thm:n9v12-Holder}.
\end{proof}

\subsection{Loading the density norm from a logarithmic moment}
\label{sec:n9v12-log}

Let
\[
 r={e^\ell\over\int e^\ell\,d\nu},\qquad
 u:=\ell-\int\ell\,d\nu.
\]
Centering does not change \(r\).

\begin{theorem}[Centered log-density loader]
\label{thm:n9v12-log}
Suppose, for some \(q>1\) and \(V\geq0\),
\begin{equation}
 \int e^{q u}\,d\nu\leq e^{q^2V/2}.                  \label{eq:n9v12-log-mgf}
\end{equation}
Then
\begin{equation}
 \|r\|_{L^q(\nu)}\leq e^{qV/2}.                       \label{eq:n9v12-rq}
\end{equation}
In particular, MAC with moment range containing \(\lambda=q\) implies
\eqref{eq:n9v12-rq}.
\end{theorem}

\begin{proof}
Jensen's inequality gives
\(\int e^u\,d\nu\geq e^{\int u\,d\nu}=1\).  Hence
\[
 \|r\|_q
 =
 {\left(\int e^{qu}\,d\nu\right)^{1/q}
  \over\int e^u\,d\nu}
 \leq e^{qV/2}.
\]
\end{proof}

\begin{corollary}[MAC \(L^2\) density]
\label{cor:n9v12-MAC-L2}
Under the V2 martingale action card with budget \(V\),
\begin{equation}
 \int r^2\,d\nu\leq e^{2V},\qquad
 \|r\|_2\leq e^V.                                     \label{eq:n9v12-MAC-L2}
\end{equation}
\end{corollary}

\begin{proof}
Use \Cref{lem:n9v2-mgf} at \(\lambda=2\) and
\Cref{thm:n9v12-log}.
\end{proof}

\begin{firewall}[The full logarithmic density is required]
\label{fire:n9v12-full}
An \(L^2\) estimate for the isolated action factor does not bound the
density after determinant, coarea, transport, and counterterm factors are
multiplied and renormalized.  Equation \eqref{eq:n9v12-log-mgf} must hold
for the full owned logarithmic density or be obtained from a proved
sequential factorization whose reference measure is tracked at every step.
\end{firewall}

\subsection{An entropy alternative}
\label{sec:n9v12-entropy}

\begin{lemma}[Entropy variational bound]
\label{lem:n9v12-entropy}
Let
\[
 \operatorname{Ent}(\mu\mid\nu):=\int r\log r\,d\nu.
\]
For every measurable \(Y\) and every \(\theta>0\) with
\(\int e^{\theta Y}\,d\nu<\infty\),
\begin{equation}
 \int Y\,d\mu
 \leq
 {1\over\theta}
 \left[
 \operatorname{Ent}(\mu\mid\nu)
 +\log\int e^{\theta Y}\,d\nu
 \right].                                             \label{eq:n9v12-entropy}
\end{equation}
\end{lemma}

\begin{proof}
Define the exponential tilt
\[
 d\nu_{\theta,Y}
 ={e^{\theta Y}\over\int e^{\theta Y}\,d\nu}\,d\nu.
\]
Nonnegativity of
\(\operatorname{Ent}(\mu\mid\nu_{\theta,Y})\) expands to
\[
 0\leq
 \operatorname{Ent}(\mu\mid\nu)
 -\theta\int Y\,d\mu
 +\log\int e^{\theta Y}\,d\nu.
\]
Rearrange.
\end{proof}

\begin{proposition}[Renyi control of entropy]
\label{prop:n9v12-Renyi}
For \(q>1\),
\begin{equation}
 \operatorname{Ent}(\mu\mid\nu)
 \leq{1\over q-1}\log\int r^q\,d\nu.                 \label{eq:n9v12-Renyi}
\end{equation}
Under MAC, this gives
\begin{equation}
 \operatorname{Ent}(\mu\mid\nu)\leq2V                \label{eq:n9v12-MAC-entropy}
\end{equation}
by taking \(q=2\).
\end{proposition}

\begin{proof}
Convexity of
\[
 \varphi(q):=\log\int r^q\,d\nu
\]
and \(\varphi(1)=0\) imply
\(\varphi'(1)\leq[\varphi(q)-\varphi(1)]/(q-1)\).
But \(\varphi'(1)=\int r\log r\,d\nu\), by truncation if necessary.
The MAC conclusion follows from \eqref{eq:n9v12-MAC-L2}.
\end{proof}

\begin{corollary}[Exponential-reference moment transfer]
\label{cor:n9v12-exp}
If
\[
 \sup_n\operatorname{Ent}(\mu_n\mid\nu_n)\leq E,\qquad
 \sup_n\int
 e^{\theta\|u\|_{H^{-s_0}}^p}\,d\nu_n\leq M,
\]
then
\begin{equation}
 \sup_n\int\|u\|_{H^{-s_0}}^p\,d\mu_n
 \leq {E+\log M\over\theta}.                          \label{eq:n9v12-exp}
\end{equation}
\end{corollary}

\begin{proof}
Apply \Cref{lem:n9v12-entropy} with
\(Y=\|u\|_{H^{-s_0}}^p\).
\end{proof}

\subsection{Why total variation is insufficient for moments}
\label{sec:n9v12-counterexample}

\begin{proposition}[Vanishing total variation with exploding moments]
\label{prop:n9v12-counterexample}
On \(\mathbb R\), let
\[
 \nu=\delta_0,\qquad
 \mu_n=\left(1-{1\over n}\right)\delta_0
       +{1\over n}\delta_{n^2}.
\]
Then
\begin{equation}
 \|\mu_n-\nu\|_{\rm TV}={1\over n}\longrightarrow0,   \label{eq:n9v12-TV-zero}
\end{equation}
while
\begin{equation}
 \int x^2\,d\mu_n(x)=n^3\longrightarrow\infty.        \label{eq:n9v12-moment-explode}
\end{equation}
\end{proposition}

\begin{proof}
The measures differ by mass \(1/n\) moved from zero to \(n^2\), proving
\eqref{eq:n9v12-TV-zero}.  Direct calculation gives
\((1/n)(n^2)^2=n^3\).
\end{proof}

\begin{corollary}[The V98 compiler does not imply BNSMC]
\label{cor:n9v12-TV}
Summable adjacent total-variation changes of bounded cylinder marginals do
not by themselves supply the unbounded Sobolev moment required by BNSMC.
\end{corollary}

\begin{proof}
The example can be placed in one coordinate of a Hilbert space and made
stationary in all other coordinates.
\end{proof}

\subsection{Modewise fourth-moment route}
\label{sec:n9v12-modes}

\begin{theorem}[Density-loaded spectral moment]
\label{thm:n9v12-modes}
Let \(u_n\) be random distributions under reference probabilities
\(\nu_n\), and let \(\mu_n=r_n\nu_n\).  Suppose
\begin{align}
 \sup_n\|r_n\|_{L^2(\nu_n)}&\leq R,                   \label{eq:n9v12-mode-r}\\
 \int|\langle u_n,e_k\rangle|^4\,d\nu_n
 &\leq C_4(1+\lambda_k)^{2\alpha}                    \label{eq:n9v12-mode-four}
\end{align}
for all \(n,k\).  If \(s_0>\alpha+d/2\), then
\begin{equation}
 \sup_n\int\|u_n\|_{H^{-s_0}}^2\,d\mu_n<\infty.       \label{eq:n9v12-mode-target}
\end{equation}
\end{theorem}

\begin{proof}
Holder's inequality gives
\[
 \int|\langle u_n,e_k\rangle|^2\,d\mu_n
 \leq
 \|r_n\|_2
 \left(
 \int|\langle u_n,e_k\rangle|^4\,d\nu_n
 \right)^{1/2}
 \leq R\sqrt{C_4}(1+\lambda_k)^\alpha.
\]
The Weyl-counted loader \Cref{cor:n9v11-Weyl} then proves
\eqref{eq:n9v12-mode-target}.
\end{proof}

\begin{corollary}[Gaussian reference form]
\label{cor:n9v12-Gaussian}
If each real reference mode is centered Gaussian with variance at most
\(C_2(1+\lambda_k)^\alpha\), then
\eqref{eq:n9v12-mode-four} holds with \(C_4=3C_2^2\).
\end{corollary}

\begin{proof}
A centered real Gaussian variable of variance \(\sigma^2\) has fourth
moment \(3\sigma^4\).
\end{proof}

\subsection{Cumulative bounded rows}
\label{sec:n9v12-cumulative}

\begin{theorem}[Summable oscillation gives a uniform cumulative density]
\label{thm:n9v12-cumulative}
Suppose all measures are transported to one measurable space and
\[
 {d\mu_{n+1}\over d\mu_n}
 ={e^{\ell_n}\over\int e^{\ell_n}\,d\mu_n},
 \qquad
 \operatorname{osc}\ell_n\leq\varepsilon_n.
\]
If \(\sum_n\varepsilon_n<\infty\), then for every \(N\),
\begin{equation}
 \left\|{d\mu_N\over d\mu_1}\right\|_{L^\infty(\mu_1)}
 \leq
 \exp\!\left(\sum_{n=1}^{N-1}\varepsilon_n\right)
 \leq e^{\sum_n\varepsilon_n}.                        \label{eq:n9v12-cumulative}
\end{equation}
\end{theorem}

\begin{proof}
For a function of oscillation at most \(\varepsilon\), its normalized
exponential lies pointwise between \(e^{-\varepsilon}\) and
\(e^\varepsilon\).  The chain rule for Radon--Nikodym derivatives gives
\[
 {d\mu_N\over d\mu_1}
 =\prod_{n=1}^{N-1}{d\mu_{n+1}\over d\mu_n}.
\]
Multiply the upper bounds.
\end{proof}

\begin{corollary}[Moment propagation along the oscillation branch]
\label{cor:n9v12-cumulative}
Under \Cref{thm:n9v12-cumulative}, every nonnegative \(X\) satisfies
\begin{equation}
 \sup_N\int X\,d\mu_N
 \leq
 e^{\sum_n\varepsilon_n}\int X\,d\mu_1.              \label{eq:n9v12-cumulative-moment}
\end{equation}
\end{corollary}

\begin{proof}
Integrate \(X\) against the density in
\eqref{eq:n9v12-cumulative}.
\end{proof}

\begin{firewall}[Changing spaces require exact transports]
\label{fire:n9v12-transport}
The cumulative product is meaningful only after all adjacent measures are
represented on one common space and the Radon--Nikodym chain rule is exact.
For changing lattices, this requires the ambient transport and
common-refinement row that remains open in the main programme.
\end{firewall}

\subsection{The density-loaded moment card}
\label{sec:n9v12-card}

\begin{definition}[Density-loaded bosonic moment card, DLBMC]
\label{def:n9v12-DLBMC}
DLBMC consists of BNSMC together with, for each bosonic species, either:
\begin{enumerate}[label=\textup{(M\arabic*)},leftmargin=4em]
\item a full-density \(L^q\) bound and the corresponding higher reference
Sobolev moment in \eqref{eq:n9v12-Holder};
\item a full relative-entropy bound and the reference exponential moment in
\eqref{eq:n9v12-exp}; or
\item an exact common-space cumulative oscillation bound as in
\eqref{eq:n9v12-cumulative}.
\end{enumerate}
The card records the complete owned density; no factor may be omitted or
charged twice.
\end{definition}

\begin{theorem}[DLBMC loads Radon tightness]
\label{thm:n9v12-closure}
DLBMC supplies the full-measure moments required by BNSMC.  Hence, under
the cylinder hypotheses of \Cref{thm:n9v11-Radon}, it gives a unique Radon
bosonic limit on the declared negative-Sobolev product space.
\end{theorem}

\begin{proof}
Use \Cref{thm:n9v12-Holder}, \Cref{cor:n9v12-exp}, or
\Cref{cor:n9v12-cumulative} according to the selected branch.  Then apply
\Cref{cor:n9v11-BNSMC,thm:n9v11-Radon}.
\end{proof}

\begin{assessment}[Progress at ninth-series V12]
\label{ass:n9v12-status}
The tightness row is now attached to the density ledger.  A full-density
\(L^q\) norm transfers higher reference moments by Holder; relative entropy
transfers reference exponential moments; and a summable oscillation branch
gives a uniform cumulative density.  The remaining physical work is to
obtain one of these bounds for the complete constrained density, not only
for its action factor.
\end{assessment}

\begin{programme}[Ninth-series V13: reference Gaussian field moments]
\label{prog:n9v12-V13}
V13 will compute the spectral fourth-moment loader for massive scalar and
gauge-fixed one-form Gaussian references, state the zero-mode subtraction
exactly, and determine the negative-Sobolev indices delivered to DLBMC.
The Euclidean gravitational conformal mode will remain outside this
positive Gaussian calculation.
\end{programme}

\begin{firewall}[Claim boundary after ninth-series V12]
\label{fire:n9v12-boundary}
V12 does not prove a full-density \(L^q\) or entropy bound, BNSMC for the
physical measures, nonlinear geometric support, fermionic reconstruction,
or the SM--Einstein continuum.
\end{firewall}

\begin{theorem}[Ninth-series V12 result]
\label{thm:n9v12-result}
Uniform \(L^q\) control of the full density transports higher reference
Sobolev moments to the constrained measure by
\eqref{eq:n9v12-Holder}.  MAC supplies such an \(L^2\) bound for one
complete logarithmic factor, entropy gives an exponential-moment
alternative, and DLBMC combines either route with V11 to yield conditional
Radon tightness.
\end{theorem}

\begin{proof}
Use \Cref{thm:n9v12-Holder,thm:n9v12-log,
lem:n9v12-entropy,thm:n9v12-modes,
thm:n9v12-cumulative,thm:n9v12-closure}.
\end{proof}

\paragraph{Parent-version record.}
V12 was formed from
\texttt{Renewal\_SM\_Einstein\_Continuum\_Research\_Notes\_V11.tex}.
The V11 parent had \(3\,819\) logical source lines, \(138\,567\) bytes, and
SHA--256
\[
 \texttt{1DEF47D780EF93AFEE8A91864185CC341C2329BD1745A3BCBF2383B982D4FAB9}.
\]
The next compulsory companion audit is V15.

% =============================================================================
% END APPENDED NINTH-SERIES V12 MATERIAL
% =============================================================================


% =============================================================================
% BEGIN APPENDED NINTH-SERIES V13 MATERIAL
% =============================================================================

\clearpage
\section{Ninth-series V13: Gaussian reference Sobolev moments}
\label{sec:v9v13}

\subsection{A uniformly dominated spectral Gaussian}
\label{sec:n9v13-setting}

Let \(\Delta_Ee_k=\lambda_ke_k\) be the Laplace-type spectral resolution
from V11.  For each regulator \(n\), let
\begin{equation}
 X_n=\sum_{k=1}^{K_n}\sqrt{c_{n,k}}\,\xi_{n,k}e_k,     \label{eq:n9v13-X}
\end{equation}
where the \(\xi_{n,k}\) are independent standard real Gaussians.  The
cutoff \(K_n\) may tend to infinity.  Assume, for constants \(C>0\) and
\(\beta\geq0\),
\begin{equation}
 0\leq c_{n,k}\leq C(1+\lambda_k)^{-\beta}            \label{eq:n9v13-cov}
\end{equation}
uniformly in \(n,k\).

\begin{theorem}[Uniform Gaussian Sobolev moments]
\label{thm:n9v13-moments}
If
\begin{equation}
 s_0>{d\over2}-\beta,                                 \label{eq:n9v13-index}
\end{equation}
then
\begin{align}
 T_{s_0}
 &:=
 C\sum_{k=1}^{\infty}
 (1+\lambda_k)^{-s_0-\beta}<\infty,                  \label{eq:n9v13-T}\\
 \sup_n\mathbb E\|X_n\|_{H^{-s_0}}^2
 &\leq T_{s_0},                                       \label{eq:n9v13-second}\\
 \sup_n\mathbb E\|X_n\|_{H^{-s_0}}^4
 &\leq3T_{s_0}^2.                                     \label{eq:n9v13-fourth}
\end{align}
\end{theorem}

\begin{proof}
Put
\[
 a_{n,k}:=(1+\lambda_k)^{-s_0}c_{n,k}.
\]
Then
\[
 \|X_n\|_{H^{-s_0}}^2
 =\sum_{k\leq K_n}a_{n,k}\xi_{n,k}^2.
\]
Weyl's eigenvalue bound shows that the series in \eqref{eq:n9v13-T}
converges exactly under the sufficient inequality
\(s_0+\beta>d/2\).  Taking expectation gives
\eqref{eq:n9v13-second}.  Independence and
\(\mathbb E\xi_k^4=3\) yield
\begin{align*}
 \mathbb E\|X_n\|_{H^{-s_0}}^4
 &=
 \sum_k3a_{n,k}^2
 +2\sum_{k<j}a_{n,k}a_{n,j}\\
 &=
 \left(\sum_ka_{n,k}\right)^2
 +2\sum_ka_{n,k}^2\\
 &\leq3\left(\sum_ka_{n,k}\right)^2
 \leq3T_{s_0}^2.
\end{align*}
\end{proof}

\begin{corollary}[Tightness under a full \(L^2\) density]
\label{cor:n9v13-density}
Let \(\nu_n\) be the laws of \(X_n\), and let
\(\mu_n=r_n\nu_n\) with
\begin{equation}
 \sup_n\|r_n\|_{L^2(\nu_n)}\leq R.                    \label{eq:n9v13-r}
\end{equation}
Then
\begin{equation}
 \sup_n\int\|u\|_{H^{-s_0}}^2\,d\mu_n(u)
 \leq\sqrt3\,R\,T_{s_0}.                              \label{eq:n9v13-target}
\end{equation}
Consequently \((\mu_n)\) is tight in \(H^{-s}(E)\) for every
\(s>s_0\).
\end{corollary}

\begin{proof}
Apply \Cref{cor:n9v12-L2} with \(p=2\) and
\eqref{eq:n9v13-fourth}, then apply
\Cref{thm:n9v11-tight}.
\end{proof}

\begin{corollary}[Sharp range supplied by this loader]
\label{cor:n9v13-range}
For every
\begin{equation}
 s>{d\over2}-\beta,                                   \label{eq:n9v13-range}
\end{equation}
one may choose
\[
 {d\over2}-\beta<s_0<s
\]
and obtain tightness in \(H^{-s}\).
\end{corollary}

\begin{proof}
Choose the intermediate index and use
\Cref{cor:n9v13-density}.
\end{proof}

\subsection{Massive and massless scalar fields}
\label{sec:n9v13-scalar}

\begin{proposition}[Massive scalar covariance]
\label{prop:n9v13-massive}
For \(m>0\), the covariance
\[
 C_m=(\Delta+m^2)^{-1}
\]
has spectral variances
\[
 c_k=(\lambda_k+m^2)^{-1}
 \leq C_m'(1+\lambda_k)^{-1}.
\]
Thus \(\beta=1\), and the Gaussian scalar cutoffs satisfy the V13 moment
card for every
\begin{equation}
 s_0>{d\over2}-1.                                     \label{eq:n9v13-scalar-index}
\end{equation}
\end{proposition}

\begin{proof}
The ratio
\((1+\lambda)/(\lambda+m^2)\) is bounded on \([0,\infty)\).  Apply
\Cref{thm:n9v13-moments} with \(\beta=1\).
\end{proof}

\begin{proposition}[Massless scalar with the constant mode removed]
\label{prop:n9v13-massless}
Assume \(M\) is connected.  On the mean-zero subspace
\[
 L_0^2(M):=\left\{f:\int_Mf=0\right\},
\]
the covariance \(\Delta^{-1}\) satisfies
\[
 c_k=\lambda_k^{-1}\leq C(1+\lambda_k)^{-1}
\]
because the first nonzero eigenvalue has a positive floor.  Hence the same
index \eqref{eq:n9v13-scalar-index} holds.
\end{proposition}

\begin{proof}
The zero eigenspace consists of constants.  On its orthogonal complement,
\(\lambda_k\geq\lambda_1^+>0\), so
\((1+\lambda_k)/\lambda_k\leq1+1/\lambda_1^+\).
Apply \Cref{thm:n9v13-moments}.
\end{proof}

\begin{firewall}[Removing a zero mode is a measure decision]
\label{fire:n9v13-zero}
The massless covariance is not defined on constants.  One must fix the
constant mode, quotient it, or supply a separate probability law for it.
Deleting it from a formula without declaring which of these operations was
performed would change the theory.
\end{firewall}

\subsection{Gauge-fixed one-forms}
\label{sec:n9v13-oneform}

Let
\[
 L^2\Omega^1
 =
 \overline{d\Omega^0}
 \oplus\overline{d^*\Omega^2}
 \oplus\mathcal H^1
\]
be the Hodge decomposition.  The coexact sector
\(\overline{d^*\Omega^2}\) is a linear gauge-fixed transverse space.

\begin{theorem}[Transverse one-form Gaussian]
\label{thm:n9v13-oneform}
On the coexact one-form sector, the covariance
\[
 C_{1,m}=(\Delta_1+m^2)^{-1}
\]
for \(m\geq0\) has \(\beta=1\); when \(m=0\), harmonic one-forms are
excluded.  Uniform spectral cutoffs therefore satisfy
\eqref{eq:n9v13-second}--\eqref{eq:n9v13-fourth} whenever
\[
 s_0>{d\over2}-1.
\]
\end{theorem}

\begin{proof}
The Hodge Laplacian is Laplace type and has the same Weyl exponent \(d/2\).
For \(m>0\), use the massive argument.  For \(m=0\), the coexact sector is
orthogonal to the finite-dimensional harmonic kernel and its first
positive eigenvalue has a floor.  Apply
\Cref{thm:n9v13-moments}.
\end{proof}

\begin{corollary}[Finite harmonic slow law]
\label{cor:n9v13-harmonic}
Let \(h_1,\ldots,h_b\) be an orthonormal basis of \(\mathcal H^1\), and
suppose a slow vector \(Z_n\in\mathbb R^b\) satisfies
\[
 \sup_n\mathbb E|Z_n|^4<\infty.
\]
Then adding
\[
 \sum_{a=1}^bZ_{n,a}h_a
\]
to the transverse Gaussian preserves the uniform fourth
\(H^{-s_0}\) moment.
\end{corollary}

\begin{proof}
All Sobolev norms are equivalent on the fixed finite-dimensional harmonic
space.  Use
\(\|x+y\|^4\leq8(\|x\|^4+\|y\|^4)\) and
\Cref{thm:n9v13-oneform}.
\end{proof}

\begin{firewall}[Linear gauge fixing is not the nonabelian quotient]
\label{fire:n9v13-gauge}
The coexact Gaussian is a linear reference calculation.  A nonabelian
connection space has Gribov and bundle-chart issues, and its interacting
coarea density is not the Gaussian measure above.  V13 supplies a
reference moment only after a physical regulator proves the required chart
and full-density \(L^2\) loader.
\end{firewall}

\subsection{General smoothing order and field species}
\label{sec:n9v13-species}

\begin{proposition}[Covariance order table]
\label{prop:n9v13-order}
If a positive Gaussian reference covariance on a finite-rank bundle obeys
\[
 C\preceq C_0(1+\Delta_E)^{-\beta},
\]
then its spectral cutoffs satisfy the V13 fourth-moment card for
\[
 s_0>{d\over2}-\beta.
\]
In particular:
\begin{center}
\begin{tabular}{c|c|c}
reference covariance & \(\beta\) & sufficient \(s_0\)\\ \hline
white noise & \(0\) & \(s_0>d/2\)\\
second-order massive/free field & \(1\) & \(s_0>d/2-1\)\\
fourth-order positive reference & \(2\) & \(s_0>d/2-2\)
\end{tabular}
\end{center}
\end{proposition}

\begin{proof}
The quadratic-form domination bounds the variance of every covariance
eigenmode after simultaneous diagonalization when the covariance is a
function of \(\Delta_E\); in the general noncommuting case, apply the trace
bound
\[
 \Tr[(1+\Delta_E)^{-s_0}C]
 \leq C_0\Tr(1+\Delta_E)^{-s_0-\beta}.
\]
The Gaussian fourth Hilbert-norm identity is bounded by three times the
square of this trace.  Weyl summability gives the stated threshold.
\end{proof}

\begin{firewall}[Positive fourth-order reference versus gravity]
\label{fire:n9v13-gravity}
A positive fourth-order Gaussian reference may be useful for domination,
but it is not the Euclidean Einstein measure.  The real conformal mode of
the Einstein action is unbounded in the wrong direction.  V13 does not
repair that sign or justify an analytic continuation.
\end{firewall}

\subsection{Gaussian reference moment card}
\label{sec:n9v13-card}

\begin{definition}[Gaussian reference Sobolev card, GRSC]
\label{def:n9v13-GRSC}
For each bosonic species, GRSC records:
\begin{enumerate}[label=\textup{(R\arabic*)},leftmargin=4em]
\item a positive Gaussian reference on a fixed bundle distribution space;
\item covariance domination exponent \(\beta\) and constant \(C\);
\item every removed kernel and the separate law of each retained harmonic
or slow coordinate;
\item a common spectral embedding of all regulator cutoffs;
\item an index \(s_0>d/2-\beta\); and
\item a full constrained-density \(L^2\) bound relative to that reference.
\end{enumerate}
\end{definition}

\begin{theorem}[GRSC-to-DLBMC loader]
\label{thm:n9v13-loader}
GRSC gives the uniform full-measure \(H^{-s_0}\) second moment required by
DLBMC.  For any \(s>s_0\), the corresponding species is tight in
\(H^{-s}\).
\end{theorem}

\begin{proof}
\Cref{thm:n9v13-moments} supplies the reference fourth moment.  The density
item and \Cref{cor:n9v13-density} supply the target second moment.
\Cref{thm:n9v11-tight} gives tightness.
\end{proof}

\begin{assessment}[Progress at ninth-series V13]
\label{ass:n9v13-status}
The reference side of the bosonic tightness card is explicit.  A covariance
of smoothing order \(\beta\) gives uniform fourth moments above the
threshold \(s_0>d/2-\beta\); a full-density \(L^2\) bound then transfers a
second moment and yields tightness after an arbitrarily small additional
loss.  Scalar and transverse one-form references have \(\beta=1\), with
zero and harmonic modes kept as declared slow sectors.
\end{assessment}

\begin{programme}[Ninth-series V14: local-to-global distribution spaces]
\label{prog:n9v13-V14}
V14 will extend the compact-base theorem to an exhaustion and prove
tightness in \(H^{-s}_{\rm loc}\) from weighted local moment budgets.  It
will also state the common-trivialization cocycle needed for connection and
metric tensor charts.  The result will prepare the compulsory V15 audit.
\end{programme}

\begin{firewall}[Claim boundary after ninth-series V13]
\label{fire:n9v13-boundary}
V13 does not prove the full-density \(L^2\) bound, the interacting
nonabelian or gravitational measure, nonlinear support, reconstruction, or
the SM--Einstein continuum.
\end{firewall}

\begin{theorem}[Ninth-series V13 result]
\label{thm:n9v13-result}
Uniform Gaussian covariance domination
\(C_n\preceq C(1+\Delta)^{-\beta}\) gives uniform second and fourth
\(H^{-s_0}\) moments for \(s_0>d/2-\beta\).  With a full-density \(L^2\)
bound, GRSC loads DLBMC and gives tightness in every
\(H^{-s}\) with \(s>s_0\).
\end{theorem}

\begin{proof}
Use \Cref{thm:n9v13-moments,cor:n9v13-density,
prop:n9v13-massive,thm:n9v13-oneform,
prop:n9v13-order,thm:n9v13-loader}.
\end{proof}

\paragraph{Parent-version record.}
V13 was formed from
\texttt{Renewal\_SM\_Einstein\_Continuum\_Research\_Notes\_V12.tex}.
The V12 parent had \(4\,247\) logical source lines, \(151\,482\) bytes, and
SHA--256
\[
 \texttt{90EA3DA54A89CCA0FB37F0C8A44566B14BD753CE856208E52536DD7874A2AE95}.
\]
The next compulsory companion audit is V15.

% =============================================================================
% END APPENDED NINTH-SERIES V13 MATERIAL
% =============================================================================


% =============================================================================
% BEGIN APPENDED NINTH-SERIES V14 MATERIAL
% =============================================================================

\clearpage
\section{Ninth-series V14: local-to-global distribution-space tightness}
\label{sec:v9v14}

\subsection{An abstract projective Polish space}
\label{sec:n9v14-projective}

Let \(X_j\) be Polish spaces with continuous restriction maps
\[
 r_{j+1,j}:X_{j+1}\longrightarrow X_j,
\qquad
 r_{k,j}=r_{j+1,j}\circ\cdots\circ r_{k,k-1}.
\]
Define the inverse-limit space
\begin{equation}
 X:=
 \left\{
 (x_j)\in\prod_{j=1}^{\infty}X_j:
 r_{j+1,j}x_{j+1}=x_j\ \text{for every }j
 \right\}.                                             \label{eq:n9v14-X}
\end{equation}
Give \(X\) the subspace topology from the countable product.

\begin{lemma}[The inverse limit is Polish]
\label{lem:n9v14-Polish}
The space \(X\) in \eqref{eq:n9v14-X} is a closed subspace of the Polish
product \(\prod_jX_j\), hence is Polish.
\end{lemma}

\begin{proof}
For every \(j\), the map
\[
 (x_k)\longmapsto(r_{j+1,j}x_{j+1},x_j)
\]
is continuous, and the diagonal of \(X_j\times X_j\) is closed.  The
compatibility locus is the countable intersection of the inverse images of
these diagonals, so it is closed.  A countable product of Polish spaces is
Polish, as is a closed subspace.
\end{proof}

\begin{theorem}[Projective local tightness]
\label{thm:n9v14-projective}
Let \(\mu_n\) be Borel probabilities on \(X\), and let
\(\mu_{n,j}\) be their \(X_j\)-marginals.  If
\begin{equation}
 \{\mu_{n,j}:n\geq1\}\ \text{is tight on }X_j
 \quad\text{for every }j,                             \label{eq:n9v14-local-tight}
\end{equation}
then \((\mu_n)\) is tight on \(X\).
\end{theorem}

\begin{proof}
Fix \(\varepsilon>0\).  By \eqref{eq:n9v14-local-tight}, choose compact
\(K_j\subset X_j\) such that
\[
 \sup_n\mu_{n,j}(K_j^c)<\varepsilon2^{-j}.
\]
The product \(\prod_jK_j\) is compact by Tychonoff's theorem; because the
product is countable and metrizable, this is ordinary sequential
compactness as well.  The set
\[
 K:=X\cap\prod_jK_j
\]
is closed in that product and hence compact in \(X\).  The union bound gives
\[
 \sup_n\mu_n(K^c)
 \leq\sum_j\sup_n\mu_{n,j}(K_j^c)
 <\varepsilon.
\]
\end{proof}

\subsection{Local negative-Sobolev realization}
\label{sec:n9v14-local-Sobolev}

Let \(M\) be a sigma-compact manifold with a smooth exhaustion
\[
 K_1\Subset K_2\Subset\cdots,\qquad \bigcup_jK_j=M.
\]
Choose cutoffs \(\chi_j\in C_c^\infty(M)\) with
\(\chi_j=1\) on \(K_j\).  For a bundle \(E\to M\), define the local
distribution topology by the seminorms
\begin{equation}
 p_j(u):=\|\chi_ju\|_{H^{-s}(E)}.                     \label{eq:n9v14-seminorm}
\end{equation}
Equivalent choices of exhaustion and cutoffs give the usual
\(H^{-s}_{\rm loc}(E)\) topology.

\begin{theorem}[Local Sobolev moment criterion]
\label{thm:n9v14-local}
Suppose \(s>s_0\), and for every \(j\) there are \(p_j>0\) and
\(C_j<\infty\) such that
\begin{equation}
 \sup_n\int
 \|\chi_ju\|_{H^{-s_0}}^{p_j}\,d\mu_n(u)
 \leq C_j.                                             \label{eq:n9v14-local-moment}
\end{equation}
Then \((\mu_n)\) is tight on \(H^{-s}_{\rm loc}(E)\).
\end{theorem}

\begin{proof}
For fixed \(j\), \Cref{thm:n9v11-tight} gives tightness of the laws of
\(\chi_ju\) in \(H^{-s}\).  The local space is the inverse limit of the
compatible restrictions to the exhaustion, with topology generated by
\eqref{eq:n9v14-seminorm}.  Apply
\Cref{thm:n9v14-projective}.
\end{proof}

\begin{corollary}[Several local bosonic species]
\label{cor:n9v14-species}
If \eqref{eq:n9v14-local-moment} holds for every exhaustion index and every
member of a finite family of metric-coordinate, connection, and Higgs
bundle fields, then their joint laws are tight on the product of the
corresponding \(H^{-s_a}_{\rm loc}\) spaces.
\end{corollary}

\begin{proof}
Use finite-product tightness at each exhaustion index, followed by
\Cref{thm:n9v14-projective}.  Independence is not required.
\end{proof}

\subsection{Smooth trivialization cocycles}
\label{sec:n9v14-cocycle}

The local Sobolev norm is computed in a bundle trivialization.  The
transition functions must act continuously and with controlled constants.
For a nonnegative integer \(m\), write
\[
 \|G\|_{C^m(K)}
 :=\max_{|\alpha|\leq m}\sup_K|\partial^\alpha G|.
\]

\begin{lemma}[Multiplication on negative integer Sobolev space]
\label{lem:n9v14-multiplier}
Let \(K\) lie in one coordinate chart and let \(G\) be a smooth
matrix-valued function.  Multiplication \(u\mapsto Gu\) satisfies
\begin{equation}
 \|Gu\|_{H^{-m}(K)}
 \leq C_{K,m}\|G\|_{C^m(K)}\|u\|_{H^{-m}(K)}.         \label{eq:n9v14-multiplier}
\end{equation}
\end{lemma}

\begin{proof}
By duality,
\[
 \|Gu\|_{H^{-m}}
 =
 \sup_{\|\varphi\|_{H^m}\leq1}
 |\langle u,G^*\varphi\rangle|.
\]
Leibniz's rule gives
\[
 \|G^*\varphi\|_{H^m}
 \leq C_{K,m}\|G\|_{C^m}\|\varphi\|_{H^m}.
\]
Insert this estimate in the dual norm.
\end{proof}

\begin{proposition}[Tensor-chart moment transport]
\label{prop:n9v14-tensor}
Suppose two local tensor frames are related on \(K\) by \(u'=G u\), with
\[
 \|G\|_{C^m(K)}\leq L.
\]
Then every \(p>0\) obeys
\begin{equation}
 \int\|u'\|_{H^{-m}}^p\,d\mu
 \leq(C_{K,m}L)^p
 \int\|u\|_{H^{-m}}^p\,d\mu.                         \label{eq:n9v14-tensor}
\end{equation}
\end{proposition}

\begin{proof}
Raise \eqref{eq:n9v14-multiplier} to the power \(p\) and integrate.
\end{proof}

\begin{proposition}[Affine connection-chart transport]
\label{prop:n9v14-connection}
In two gauge trivializations related by a fixed smooth
\(g:K\to G\), a connection one-form transforms as
\[
 A'=g^{-1}Ag+g^{-1}dg.
\]
If \(g,g^{-1}\) have bounded \(C^{m+1}(K)\) norms, then
\begin{equation}
 \|A'\|_{H^{-m}}
 \leq C_g\bigl(1+\|A\|_{H^{-m}}\bigr).                \label{eq:n9v14-connection}
\end{equation}
Consequently a uniform \(p\)-moment for \(A\) gives one for \(A'\).
\end{proposition}

\begin{proof}
Apply \Cref{lem:n9v14-multiplier} twice to \(g^{-1}Ag\).  The smooth
inhomogeneous term \(g^{-1}dg\) has finite \(H^{-m}\) norm controlled by
the declared \(C^{m+1}\) bounds.  Use
\((1+x)^p\leq C_p(1+x^p)\).
\end{proof}

\begin{firewall}[Field-dependent gauge changes]
\label{fire:n9v14-field-gauge}
The preceding proposition treats a fixed smooth transition function.
A gauge-fixing map \(g(A)\) depending on the random field can have singular
derivatives and Gribov discontinuities.  Its moment transport requires a
separate measurable chart theorem and cannot be inferred from
\eqref{eq:n9v14-connection}.
\end{firewall}

\subsection{Local Radon landing}
\label{sec:n9v14-Radon}

\begin{theorem}[Local cylinder-to-Radon theorem]
\label{thm:n9v14-Radon}
Let \(\mathcal X_{\rm loc}\) be a finite product of the local Sobolev spaces
above.  Suppose:
\begin{enumerate}[label=\textup{(\roman*)},leftmargin=3.8em]
\item the regulated bosonic measures are Borel probabilities on
\(\mathcal X_{\rm loc}\);
\item the local moment card \eqref{eq:n9v14-local-moment} holds in compatible
trivializations; and
\item all finite smooth-test cylinder laws converge to one projectively
consistent family.
\end{enumerate}
Then the regulated measures converge weakly to a unique Radon probability
on \(\mathcal X_{\rm loc}\) with the prescribed cylinder laws.
\end{theorem}

\begin{proof}
\Cref{thm:n9v14-local,cor:n9v14-species} give tightness.
\Cref{lem:n9v14-Polish} shows the target is Polish.  Repeat the Prokhorov
and countable-coordinate uniqueness proof of
\Cref{thm:n9v11-Radon}.
\end{proof}

\begin{definition}[Local distribution landing card, LDLC]
\label{def:n9v14-LDLC}
LDLC consists of:
\begin{enumerate}[label=\textup{(L\arabic*)},leftmargin=4em]
\item a sigma-compact base, exhaustion, and finite product of bundle
\(H^{-s_a}_{\rm loc}\) spaces;
\item common measurable embeddings of all regulators;
\item density-loaded local moments \eqref{eq:n9v14-local-moment};
\item uniformly controlled fixed tensor and gauge transition cocycles;
\item projectively consistent smooth-test cylinder limits; and
\item separate nonlinear support and field-dependent gauge-chart cards.
\end{enumerate}
\end{definition}

\begin{corollary}[LDLC closes the bosonic topology row]
\label{cor:n9v14-LDLC}
LDLC yields the Radon bosonic probability of
\Cref{thm:n9v14-Radon}.
\end{corollary}

\begin{proof}
Items (L1)--(L5) are the hypotheses of
\Cref{thm:n9v14-Radon}; item (L6) records conclusions not supplied by
tightness.
\end{proof}

\begin{firewall}[Local tightness does not give asymptotic geometry]
\label{fire:n9v14-infinity}
On a noncompact base, \(H^{-s}_{\rm loc}\) tightness controls every compact
region but imposes no decay, finite energy, or boundary condition at
infinity.  Those require weighted global spaces or a separate asymptotic
card.
\end{firewall}

\begin{assessment}[Progress at ninth-series V14]
\label{ass:n9v14-status}
The bosonic topology loader now works on an exhaustion.  Local moments give
projective compact containment, smooth tensor and fixed gauge cocycles
preserve those moments, and cylinder convergence then yields a local Radon
measure.  Field-dependent gauge charts, nonlinear metric support, and
asymptotic control remain separate.
\end{assessment}

\begin{programme}[Ninth-series V15: compulsory companion audit]
\label{prog:n9v14-V15}
V15 will inspect the newest YM and NS files, compare their endpoint and
locality developments with LDLC, and choose the next acquisition among
full-density moment control, nonlinear support, and common-refinement
ambient transport.  The selected branch will continue in V16.
\end{programme}

\begin{firewall}[Claim boundary after ninth-series V14]
\label{fire:n9v14-boundary}
V14 does not verify LDLC for the physical measures, control
field-dependent gauge charts, prove nonlinear Einstein support,
reconstruction, or the SM--Einstein continuum.
\end{firewall}

\begin{theorem}[Ninth-series V14 result]
\label{thm:n9v14-result}
Tightness of every compatible local restriction implies tightness on the
inverse-limit local distribution space.  Density-loaded negative-Sobolev
moments provide those local restrictions, and controlled smooth
trivialization cocycles preserve them.  LDLC therefore conditionally lands
the cylinder candidate as a local Radon bosonic measure.
\end{theorem}

\begin{proof}
Use \Cref{thm:n9v14-projective,thm:n9v14-local,
lem:n9v14-multiplier,prop:n9v14-connection,
thm:n9v14-Radon,cor:n9v14-LDLC}.
\end{proof}

\paragraph{Parent-version record.}
V14 was formed from
\texttt{Renewal\_SM\_Einstein\_Continuum\_Research\_Notes\_V13.tex}.
The V13 parent had \(4\,611\) logical source lines, \(163\,016\) bytes, and
SHA--256
\[
 \texttt{C33B5654C9D7B50A0083F1F34ABD00B43F7AAC2F917A2CD15D13099C40860EF1}.
\]
The next version is the compulsory V15 companion audit.

% =============================================================================
% END APPENDED NINTH-SERIES V14 MATERIAL
% =============================================================================


% =============================================================================
% BEGIN APPENDED NINTH-SERIES V15 MATERIAL
% =============================================================================

\clearpage
\section{Ninth-series V15: compulsory audit and form-stable moments}
\label{sec:v9v15}

\subsection{Files and version integrity}
\label{sec:n9v15-files}

The newest Yang--Mills files inspected were
\begin{center}
\small\ttfamily
Renewal\_Yang\_Mills\_Mass\_Gap\_Research\_Notes\_V59.tex,\\
Renewal\_Yang\_Mills\_Mass\_Gap\_Research\_Notes\_V60.tex.
\end{center}
They are byte-identical.  Each has \(25\,326\) logical lines, \(903\,648\)
bytes, and SHA--256
\[
 \texttt{7845079A5A5D987B4777EDEE7CEC541B066DF861C4A2E29E1FB03CD006117B96}.
\]
Both end internally at V59, so no distinct V60 theorem is counted.  The
Navier--Stokes file remains V26, with \(5\,319\) logical lines,
\(278\,283\) bytes, and SHA--256
\[
 \texttt{82C887F8C2C69E4F28FAD7C3DC8DE5B9099CB5A4BF431EBA1FFF3E91935978AD}.
\]

\subsection{Yang--Mills V59}
\label{sec:n9v15-YM}

YM V59 corrects the preceding self-energy plan by proving a number-band
lemma and a one-crossing boundary identity.  For a polynomial interaction,
the Feshbach self-energy is supported at the artificial occupation
boundary; it gives no polynomial occupation correction to sufficiently low
vacuum sectors.  The boundary projection is controlled by number energy,
which yields relative-form smallness and an explicit logarithmic-window
rate.  Mixed and analytic-tail returns remain separate.  A
support-resolved high-occupation resolvent locality card is still open.

\begin{proposition}[Type-correct V59 transfer]
\label{prop:n9v15-YM}
The YM boundary theorem does not provide the full SM density \(L^2\) bound
or any negative-Sobolev moment.  It does identify a reusable analytic
interface: a positive Gaussian reference whose inverse covariance receives
a relatively form-small quadratic correction retains the same covariance
smoothing order, up to an explicit margin loss.
\end{proposition}

\begin{proof}
The non-transfer is forced by the different Hamiltonian, occupation
projection, and interacting measure.  The stated interface is purely
operator-theoretic: relative quadratic-form comparison implies inverse
operator comparison.  It can therefore be proved independently in V16 and
then loaded with whichever SM quadratic correction actually satisfies its
hypothesis.
\end{proof}

\begin{firewall}[Relative form does not give spatial locality]
\label{fire:n9v15-locality}
YM V59 explicitly leaves support-resolved resolvent locality open.
Likewise, an operator inequality for a covariance can preserve Sobolev
moments without proving exponential kernel decay, clustering, or an
ambient-transport estimate.
\end{firewall}

\subsection{Navier--Stokes status}
\label{sec:n9v15-NS}

NS V26 remains unchanged.  Its pointwise Oseen derivative norms neither
give a quadratic-form comparison for the SM Gaussian reference nor supply
the full-density moment loader.

\begin{proposition}[NS non-transfer at V15]
\label{prop:n9v15-NS}
No current NS result changes GRSC, DLBMC, or LDLC.
\end{proposition}

\begin{proof}
Those cards concern covariance domination and probability moments on
bundle distribution spaces.  The NS constants control a deterministic
parabolic kernel in a different problem.
\end{proof}

\subsection{Audit decision}
\label{sec:n9v15-decision}

\begin{audit}[V15 companion conclusion]
\label{aud:n9v15}
The next acquisition will remain on the reference-moment side of the
tightness programme.  It will prove covariance domination under a
relatively form-small perturbation, include finite-rank slow corrections,
and expose the exact loss as the form margin approaches one.  Analytic
tails and spatial locality will remain separate cards.
\end{audit}

\begin{decision}[V16 acquisition order]
\label{dec:n9v15-V16}
Let \(L>0\) be a reference inverse covariance and let \(B\) be a symmetric
form perturbation satisfying
\[
 |\langle u,Bu\rangle|
 \leq\eta\langle u,Lu\rangle,\qquad \eta<1.
\]
V16 will prove
\[
 (L+B)^{-1}\preceq(1-\eta)^{-1}L^{-1},
\]
propagate the V13 second and fourth Sobolev moments, add a finite-rank slow
covariance with a trace budget, and state the corresponding form-stable
Gaussian reference card.
\end{decision}

\begin{firewall}[This route does not solve the conformal sign]
\label{fire:n9v15-conformal}
Relative smallness with \(\eta<1\) presupposes a positive reference form.
An unbounded negative Euclidean conformal direction is not a small
perturbation of a positive form on the full field space.
\end{firewall}

\begin{assessment}[Progress at ninth-series V15]
\label{ass:n9v15-status}
The audit keeps the new tightness branch type-correct.  V16 will strengthen
the Gaussian reference card without claiming an interacting density or
spatial locality theorem.  This helps accommodate gauge/Higgs
self-energies and other positive quadratic corrections while leaving the
gravitational sign obstruction explicit.
\end{assessment}

\begin{programme}[Ninth-series V16: form-stable covariance moments]
\label{prog:n9v15-V16}
Carry out \Cref{dec:n9v15-V16} and insert the result into GRSC and DLBMC.
\end{programme}

\begin{firewall}[Claim boundary after ninth-series V15]
\label{fire:n9v15-boundary}
V15 is an audit.  It does not prove the form hypothesis for a physical
sector, the full-density \(L^2\) bound, nonlinear support, reconstruction,
or the SM--Einstein continuum.
\end{firewall}

\begin{theorem}[Ninth-series V15 result]
\label{thm:n9v15-result}
The new YM result motivates, but does not supply, a form-stability loader
for Gaussian reference moments.  The selected V16 theorem will preserve
covariance smoothing under a strict positive form margin while keeping
finite-rank slow and nonlocal analytic tails separate.
\end{theorem}

\begin{proof}
Use \Cref{prop:n9v15-YM,prop:n9v15-NS,aud:n9v15}.
\end{proof}

\paragraph{Parent-version record.}
V15 was formed from
\texttt{Renewal\_SM\_Einstein\_Continuum\_Research\_Notes\_V14.tex}.
The V14 parent had \(4\,935\) logical source lines, \(174\,062\) bytes, and
SHA--256
\[
 \texttt{0EF95D84793C63165A1C330026E91F0103D7C175A6A04C3D441AD2219F5B997D}.
\]
The next compulsory companion audit is V20.

% =============================================================================
% END APPENDED NINTH-SERIES V15 MATERIAL
% =============================================================================


% =============================================================================
% BEGIN APPENDED NINTH-SERIES V16 MATERIAL
% =============================================================================

\clearpage
\section{Ninth-series V16: form-stable Gaussian covariance moments}
\label{sec:v9v16}

\subsection{Closed positive form perturbation}
\label{sec:n9v16-form}

Let \(\mathcal H\) be a Hilbert space and let \(L\) be a strictly positive
self-adjoint operator with closed quadratic form
\[
 \mathfrak l[u]:=\|L^{1/2}u\|^2.
\]
Let \(\mathfrak b\) be a symmetric form on
\(\operatorname{Dom}\mathfrak l\) satisfying
\begin{equation}
 |\mathfrak b[u]|
 \leq\eta\,\mathfrak l[u],
 \qquad 0\leq\eta<1.                                  \label{eq:n9v16-relative}
\end{equation}
Define
\[
 \mathfrak a[u]:=\mathfrak l[u]+\mathfrak b[u].
\]

\begin{theorem}[Positive form and inverse comparison]
\label{thm:n9v16-inverse}
The form \(\mathfrak a\) is closed and strictly positive, with associated
self-adjoint operator \(A\), and
\begin{equation}
 (1-\eta)L\preceq A\preceq(1+\eta)L                  \label{eq:n9v16-form-order}
\end{equation}
in quadratic-form order.  Its inverse satisfies
\begin{equation}
 0\preceq A^{-1}
 \preceq{1\over1-\eta}L^{-1}.                         \label{eq:n9v16-inverse}
\end{equation}
\end{theorem}

\begin{proof}
Equation \eqref{eq:n9v16-relative} gives
\[
 (1-\eta)\mathfrak l[u]
 \leq\mathfrak a[u]
 \leq(1+\eta)\mathfrak l[u].
\]
Thus the form norms of \(\mathfrak a\) and \(\mathfrak l\) are equivalent;
closedness of \(\mathfrak l\) implies closedness of \(\mathfrak a\), and
the representation theorem gives \(A\).  This proves
\eqref{eq:n9v16-form-order}.

For a strictly positive closed form \(\mathfrak q\) with operator \(Q\),
\[
 \langle f,Q^{-1}f\rangle
 =
 \sup_{u\in\operatorname{Dom}\mathfrak q}
 \left[2\operatorname{Re}\langle f,u\rangle-\mathfrak q[u]\right].
\]
Since \(\mathfrak a\geq(1-\eta)\mathfrak l\), the variational formula gives
\[
 \langle f,A^{-1}f\rangle
 \leq
 \sup_u\left[
 2\operatorname{Re}\langle f,u\rangle
 -(1-\eta)\mathfrak l[u]
 \right]
 ={1\over1-\eta}\langle f,L^{-1}f\rangle.
\]
\end{proof}

\begin{corollary}[Additive lower-order debit]
\label{cor:n9v16-additive}
Suppose \(L\geq m_0I\) and
\begin{equation}
 |\mathfrak b[u]|
 \leq\eta\,\mathfrak l[u]+a\|u\|^2.                  \label{eq:n9v16-additive}
\end{equation}
If
\begin{equation}
 \eta_{\rm eff}:=\eta+{a\over m_0}<1,                \label{eq:n9v16-eta-eff}
\end{equation}
then \eqref{eq:n9v16-inverse} holds with
\(\eta_{\rm eff}\) in place of \(\eta\).
\end{corollary}

\begin{proof}
The spectral floor gives
\(\|u\|^2\leq m_0^{-1}\mathfrak l[u]\).  Insert this into
\eqref{eq:n9v16-additive} and apply
\Cref{thm:n9v16-inverse}.
\end{proof}

\begin{firewall}[A gap or quotient is needed]
\label{fire:n9v16-gap}
If \(L\) has a kernel, the inverse comparison is made only after removing,
fixing, or separately integrating that kernel.  The additive debit in
\Cref{cor:n9v16-additive} cannot be absorbed without a positive floor.
This is the same slow-sector ownership requirement used in V6--V9.
\end{firewall}

\subsection{Preservation of covariance smoothing}
\label{sec:n9v16-smoothing}

Let \(\Delta_E\) be a fixed Laplace-type operator on a compact
\(d\)-manifold.  Suppose the reference covariance satisfies
\begin{equation}
 L^{-1}\preceq C_0(1+\Delta_E)^{-\beta}.              \label{eq:n9v16-reference}
\end{equation}

\begin{theorem}[Form-stable smoothing order]
\label{thm:n9v16-smoothing}
Under \eqref{eq:n9v16-relative} and
\eqref{eq:n9v16-reference},
\begin{equation}
 A^{-1}
 \preceq {C_0\over1-\eta}(1+\Delta_E)^{-\beta}.       \label{eq:n9v16-smoothing}
\end{equation}
Thus the smoothing exponent \(\beta\) is unchanged; only its constant is
multiplied by \((1-\eta)^{-1}\).
\end{theorem}

\begin{proof}
Combine \eqref{eq:n9v16-inverse} and
\eqref{eq:n9v16-reference} using transitivity of quadratic-form order.
\end{proof}

\begin{corollary}[Critical form-margin loss]
\label{cor:n9v16-collapse}
As \(\eta\uparrow1\), the covariance domination constant supplied by this
route can diverge like \((1-\eta)^{-1}\), and the Gaussian second-moment
trace has the same possible loss.  No uniform tightness estimate follows
from form smallness if the margin is allowed to reach one.
\end{corollary}

\begin{proof}
This is the constant in \eqref{eq:n9v16-smoothing}.  It is attained in the
scalar example \(A=(1-\eta)L\).
\end{proof}

\subsection{Second and fourth Sobolev moments}
\label{sec:n9v16-moments}

Let \(X_A\) be a centered Gaussian with covariance \(A^{-1}\), or any
finite spectral compression of that covariance.  Define
\begin{equation}
 T_A(s_0):=
 \Tr\!\left[
 (1+\Delta_E)^{-s_0/2}
 A^{-1}
 (1+\Delta_E)^{-s_0/2}
 \right].                                             \label{eq:n9v16-TA}
\end{equation}

\begin{theorem}[Form-stable Sobolev moment]
\label{thm:n9v16-moments}
If
\begin{equation}
 s_0>{d\over2}-\beta,                                 \label{eq:n9v16-index}
\end{equation}
then
\begin{align}
 T_A(s_0)
 &\leq
 {C_0\over1-\eta}
 \Tr(1+\Delta_E)^{-s_0-\beta}
 =:\mathcal T_A<\infty,                               \label{eq:n9v16-trace}\\
 \mathbb E\|X_A\|_{H^{-s_0}}^2
 &\leq\mathcal T_A,                                   \label{eq:n9v16-second}\\
 \mathbb E\|X_A\|_{H^{-s_0}}^4
 &\leq3\mathcal T_A^2.                                \label{eq:n9v16-fourth}
\end{align}
\end{theorem}

\begin{proof}
Sandwich \eqref{eq:n9v16-smoothing} by
\((1+\Delta_E)^{-s_0/2}\) and take the trace.  Weyl summability and
\eqref{eq:n9v16-index} make the right side finite.  In the
\(H^{-s_0}\) Hilbert space, the Gaussian covariance has trace at most
\(\mathcal T_A\).  If its eigenvalues are \(\tau_k\), then
\[
 \mathbb E\|X_A\|^2=\sum_k\tau_k,\qquad
 \mathbb E\|X_A\|^4
 =\left(\sum_k\tau_k\right)^2+2\sum_k\tau_k^2
 \leq3\left(\sum_k\tau_k\right)^2.
\]
\end{proof}

\subsection{Finite-rank slow covariance}
\label{sec:n9v16-slow}

\begin{theorem}[Finite-rank Gaussian addition]
\label{thm:n9v16-slow}
Let \(R\geq0\) be finite rank and define
\begin{equation}
 \tau_R(s_0):=
 \Tr\!\left[
 (1+\Delta_E)^{-s_0/2}
 R
 (1+\Delta_E)^{-s_0/2}
 \right].                                             \label{eq:n9v16-tauR}
\end{equation}
If \(X_R\) is an independent centered Gaussian of covariance \(R\), then
\(X_A+X_R\) is centered Gaussian of covariance \(A^{-1}+R\) and
\begin{align}
 \mathbb E\|X_A+X_R\|_{H^{-s_0}}^2
 &\leq\mathcal T_A+\tau_R(s_0),                       \label{eq:n9v16-slow-second}\\
 \mathbb E\|X_A+X_R\|_{H^{-s_0}}^4
 &\leq3[\mathcal T_A+\tau_R(s_0)]^2.                 \label{eq:n9v16-slow-fourth}
\end{align}
\end{theorem}

\begin{proof}
Independence adds the covariances.  The transformed covariance trace is
\(T_A(s_0)+\tau_R(s_0)\).  Repeat the Gaussian Hilbert-norm calculation in
\Cref{thm:n9v16-moments}.
\end{proof}

\begin{corollary}[Uniform moving slow sectors]
\label{cor:n9v16-slow}
For a regulator sequence, a uniform rank bound alone is insufficient.
A sufficient slow condition is
\begin{equation}
 \sup_n\tau_{R_n}(s_0)<\infty.                        \label{eq:n9v16-slow-uniform}
\end{equation}
Together with \(\sup_n\eta_n<1\), it gives uniform second and fourth
Sobolev moments.
\end{corollary}

\begin{proof}
Use \eqref{eq:n9v16-slow-second}--\eqref{eq:n9v16-slow-fourth}.  A
fixed-rank covariance can still have eigenvalues diverging to infinity, so
the trace condition cannot be replaced by rank alone.
\end{proof}

\subsection{Density loading and tightness}
\label{sec:n9v16-density}

\begin{theorem}[Form-stable GRSC loader]
\label{thm:n9v16-loader}
Let \(\nu_n\) be Gaussian references with inverse covariances
\(A_n=L_n+B_n\).  Suppose uniformly in \(n\):
\begin{enumerate}[label=\textup{(\roman*)},leftmargin=3.8em]
\item \(L_n^{-1}\preceq C_0(1+\Delta_E)^{-\beta}\);
\item \(|\mathfrak b_n[u]|\leq\eta_*\mathfrak l_n[u]\) with
\(\eta_*<1\);
\item every retained finite-rank slow covariance satisfies
\eqref{eq:n9v16-slow-uniform}; and
\item the full constrained density satisfies
\(\|d\mu_n/d\nu_n\|_{L^2(\nu_n)}\leq R\).
\end{enumerate}
Then, for every \(s_0>d/2-\beta\),
\begin{equation}
 \sup_n\int\|u\|_{H^{-s_0}}^2\,d\mu_n(u)<\infty,      \label{eq:n9v16-target}
\end{equation}
and \((\mu_n)\) is tight in every \(H^{-s}\) with \(s>s_0\).
\end{theorem}

\begin{proof}
\Cref{thm:n9v16-moments,thm:n9v16-slow} give a uniform reference fourth
moment.  The \(L^2\) density loader \Cref{cor:n9v12-L2} gives
\eqref{eq:n9v16-target}.  Apply \Cref{thm:n9v11-tight}.
\end{proof}

\begin{definition}[Form-stable Gaussian reference card, FSGRC]
\label{def:n9v16-FSGRC}
FSGRC augments GRSC by recording:
\begin{enumerate}[label=\textup{(F\arabic*)},leftmargin=4em]
\item the closed positive reference form and its removed kernel;
\item every quadratic correction and its relative bound \(\eta_*<1\);
\item every additive lower-order debit and spectral floor used to absorb it;
\item every finite-rank slow covariance and its transformed trace;
\item the unchanged smoothing exponent and the factor
\((1-\eta_*)^{-1}\); and
\item analytic or nonlocal tails not included in the quadratic form.
\end{enumerate}
\end{definition}

\begin{corollary}[FSGRC loads DLBMC conditionally]
\label{cor:n9v16-FSGRC}
FSGRC plus a uniform full-density \(L^2\) bound supplies the moment part of
DLBMC and hence the corresponding bosonic tightness row.
\end{corollary}

\begin{proof}
This is \Cref{thm:n9v16-loader}, followed by the V11/V14 local or global
landing theorem appropriate to the base.
\end{proof}

\begin{firewall}[What relative form control does not provide]
\label{fire:n9v16-locality}
FSGRC does not imply a weighted resolvent kernel, exponential clustering,
a determinant trace ideal, or support-resolved locality.  It also does not
control a nonquadratic analytic tail unless that tail is included in the
full-density \(L^2\) estimate or assigned a separate normalized debit.
\end{firewall}

\begin{firewall}[Euclidean conformal direction]
\label{fire:n9v16-conformal}
The theorem requires \(A_n>0\).  The negative Euclidean conformal mode
cannot satisfy \eqref{eq:n9v16-relative} with \(\eta<1\) on the full real
linear space.  A contour prescription, nonlinear positive measure, or
different gravitational construction remains necessary.
\end{firewall}

\begin{assessment}[Progress at ninth-series V16]
\label{ass:n9v16-status}
The Gaussian tightness loader is stable under every quadratic correction
whose relative form margin stays below one.  The covariance smoothing
order is unchanged, the exact loss is \((1-\eta)^{-1}\), and finite slow
covariances are controlled by a transformed trace rather than rank alone.
This is the type-correct consequence of the V15 audit.
\end{assessment}

\begin{programme}[Ninth-series V17: full-density \(L^2\) composition]
\label{prog:n9v16-V17}
V17 will return from the reference to the full density.  It will derive a
sequential \(L^2\) composition rule for action, determinant, coarea,
transport, and counterterm tilts under changing intermediate measures, and
state the exact higher-exponent loss needed to avoid an invalid product of
separate \(L^2\) bounds.
\end{programme}

\begin{firewall}[Claim boundary after ninth-series V16]
\label{fire:n9v16-boundary}
V16 does not prove the physical relative-form margin or full-density
\(L^2\) bound, repair the gravitational conformal sign, establish
nonlinear support, reconstruction, or the SM--Einstein continuum.
\end{firewall}

\begin{theorem}[Ninth-series V16 result]
\label{thm:n9v16-result}
A symmetric perturbation of relative form size \(\eta<1\) preserves
Gaussian covariance smoothing with the explicit factor
\((1-\eta)^{-1}\).  Uniform finite-rank trace and full-density \(L^2\)
cards then give the Sobolev moments and tightness required by DLBMC.
\end{theorem}

\begin{proof}
Use \Cref{thm:n9v16-inverse,thm:n9v16-smoothing,
thm:n9v16-moments,thm:n9v16-slow,
thm:n9v16-loader,cor:n9v16-FSGRC}.
\end{proof}

\paragraph{Parent-version record.}
V16 was formed from
\texttt{Renewal\_SM\_Einstein\_Continuum\_Research\_Notes\_V15.tex}.
The V15 parent had \(5\,104\) logical source lines, \(180\,528\) bytes, and
SHA--256
\[
 \texttt{3EC76B9C6151E763F736F8AC1A88A6C02805DC0490EAA7734C0ADA94194166A8}.
\]
The next compulsory companion audit is V20.

% =============================================================================
% END APPENDED NINTH-SERIES V16 MATERIAL
% =============================================================================


% =============================================================================
% BEGIN APPENDED NINTH-SERIES V17 MATERIAL
% =============================================================================

\clearpage
\section{Ninth-series V17: sequential composition of the full density}
\label{sec:v9v17}

\subsection{Sequential normalized tilts}
\label{sec:n9v17-setting}

Let \(\nu_0\) be a probability on a common transported measurable space.
For \(j=1,\ldots,m\), define
\begin{equation}
 d\nu_j=r_j\,d\nu_{j-1},\qquad
 r_j:={e^{\ell_j}\over\int e^{\ell_j}\,d\nu_{j-1}}.   \label{eq:n9v17-step}
\end{equation}
Put
\begin{equation}
 R_j:={d\nu_j\over d\nu_0}=\prod_{i=1}^jr_i.          \label{eq:n9v17-R}
\end{equation}
The five physical rows may be inserted in any declared order, provided
every \(\ell_j\) is evaluated after the preceding transports and ownership
assignments.

\subsection{Why separate \(L^2\) rows are insufficient}
\label{sec:n9v17-counterexample}

\begin{proposition}[Two sequential \(L^2\) factors can have a non-\(L^2\)
product]
\label{prop:n9v17-counterexample}
There are probabilities \(\nu_0,\nu_1,\nu_2\) such that
\[
 {d\nu_1\over d\nu_0}\in L^2(\nu_0),\qquad
 {d\nu_2\over d\nu_1}\in L^2(\nu_1),
\]
but \(d\nu_2/d\nu_0\notin L^2(\nu_0)\).
\end{proposition}

\begin{proof}
Take \(\nu_0\) to be Lebesgue probability on \((0,1)\).  Let
\[
 r_1(x)=c_1x^{-3/10},\qquad
 r_2(x)=c_2x^{-1/5},
\]
where \(c_1\) normalizes \(r_1\) under \(\nu_0\) and \(c_2\) normalizes
\(r_2\) under \(\nu_1=r_1\nu_0\).  Then
\[
 \int r_1^2\,d\nu_0
 =c_1^2\int_0^1x^{-3/5}\,dx<\infty,
\]
and
\[
 \int r_2^2\,d\nu_1
 =c_1c_2^2\int_0^1x^{-7/10}\,dx<\infty.
\]
But
\[
 \int(r_1r_2)^2\,d\nu_0
 =c_1^2c_2^2\int_0^1x^{-1}\,dx=\infty.
\]
\end{proof}

\begin{firewall}[No multiplication of rowwise \(L^2\) constants]
\label{fire:n9v17-row-L2}
The reference changes after every normalized tilt.  Thus the product of
the rowwise \(L^2(\nu_{j-1})\) norms is not a bound for
\(\|R_m\|_{L^2(\nu_0)}\).  A valid composition must pay higher exponents or
use pointwise bounds for all but a controlled integrable row.
\end{firewall}

\subsection{General escalating-exponent composition}
\label{sec:n9v17-general}

For \(p>1\), define
\[
 M_j(p):=\int R_j^p\,d\nu_0,\qquad
 L_j(p):=\int r_j^p\,d\nu_{j-1}.
\]

\begin{lemma}[One-step exponent transfer]
\label{lem:n9v17-one-step}
Let \(p_j>1\), \(A_j>1\), and
\[
 B_j={A_j\over A_j-1},\qquad
 p_{j-1}:=1+(p_j-1)A_j,\qquad
 s_j:=p_jB_j.
\]
Then
\begin{equation}
 M_j(p_j)
 \leq
 M_{j-1}(p_{j-1})^{1/A_j}
 L_j(s_j)^{1/B_j}.                                   \label{eq:n9v17-one-step}
\end{equation}
\end{lemma}

\begin{proof}
Using \(d\nu_{j-1}=R_{j-1}\,d\nu_0\),
\begin{align*}
 M_j(p_j)
 &=\int R_{j-1}^{p_j}r_j^{p_j}\,d\nu_0\\
 &=\int R_{j-1}^{p_j-1}r_j^{p_j}\,d\nu_{j-1}.
\end{align*}
Holder's inequality with conjugate exponents \(A_j,B_j\) gives
\[
 M_j(p_j)
 \leq
 \left(
 \int R_{j-1}^{(p_j-1)A_j}\,d\nu_{j-1}
 \right)^{1/A_j}
 L_j(p_jB_j)^{1/B_j}.
\]
The first integral equals
\[
 \int R_{j-1}^{(p_j-1)A_j+1}\,d\nu_0
 =M_{j-1}(p_{j-1}).
\]
\end{proof}

\begin{theorem}[Dyadic sequential \(L^2\) compiler]
\label{thm:n9v17-dyadic}
Set \(p_m=2\) and recursively
\begin{equation}
 p_{j-1}:=2p_j-1
 \quad(2\leq j\leq m).                                \label{eq:n9v17-p}
\end{equation}
Thus
\[
 p_j=2^{m-j}+1.
\]
If
\begin{align}
 L_1(p_1)&<\infty,                                    \label{eq:n9v17-row1}\\
 L_j(2p_j)&<\infty\quad(2\leq j\leq m),               \label{eq:n9v17-rowj}
\end{align}
then \(R_m\in L^2(\nu_0)\), with
\begin{equation}
 M_m(2)
 \leq
 L_1(p_1)^{\,2^{-(m-1)}}
 \prod_{j=2}^m
 L_j(2p_j)^{\,2^{-(m-j+1)}}.                          \label{eq:n9v17-dyadic}
\end{equation}
\end{theorem}

\begin{proof}
Apply \Cref{lem:n9v17-one-step} at every step with
\(A_j=B_j=2\).  This gives
\[
 M_j(p_j)
 \leq M_{j-1}(p_{j-1})^{1/2}L_j(2p_j)^{1/2}.
\]
Iterating to \(j=1\), where \(M_1(p_1)=L_1(p_1)\), gives
\eqref{eq:n9v17-dyadic}.
\end{proof}

\begin{corollary}[The five-row exponent ledger]
\label{cor:n9v17-five}
For \(m=5\), the dyadic compiler requires row moments, in insertion order,
\begin{equation}
 17,\quad18,\quad10,\quad6,\quad4.                    \label{eq:n9v17-five}
\end{equation}
Rows with bounded logarithmic oscillation should therefore be assigned the
largest required exponents, while a row with limited exponential moments
should be placed as late as the exact ownership ledger permits.
\end{corollary}

\begin{proof}
The sequence \(p_j=2^{5-j}+1\) is
\((17,9,5,3,2)\).  Row one uses \(p_1=17\), and subsequent rows use
\(2p_j=(18,10,6,4)\).
\end{proof}

\subsection{Logarithmic moment loader at higher exponent}
\label{sec:n9v17-log}

\begin{proposition}[Row moment from a centered logarithmic MGF]
\label{prop:n9v17-log}
Let
\[
 r={e^\ell\over\int e^\ell\,d\nu},\qquad
 u:=\ell-\int\ell\,d\nu.
\]
If, for \(s>1\),
\begin{equation}
 \int e^{su}\,d\nu\leq e^{s^2V/2},                   \label{eq:n9v17-log}
\end{equation}
then
\begin{equation}
 \int r^s\,d\nu\leq e^{s^2V/2}.                      \label{eq:n9v17-rs}
\end{equation}
\end{proposition}

\begin{proof}
Centering leaves \(r\) unchanged, and Jensen gives
\(\int e^u\,d\nu\geq1\).  Hence
\[
 \int r^s\,d\nu
 =
 {\int e^{su}\,d\nu\over(\int e^u\,d\nu)^s}
 \leq e^{s^2V/2}.
\]
\end{proof}

\begin{corollary}[Higher-moment martingale row]
\label{cor:n9v17-martingale}
If the conditional martingale MGF bound of V2 holds for
\(|\lambda|\leq s\), then \eqref{eq:n9v17-rs} holds.  Conditional bounded
differences supply this for every finite \(s\).
\end{corollary}

\begin{proof}
The tower iteration in \Cref{lem:n9v2-mgf} is unchanged on the larger
moment interval.  The conditional Hoeffding lemma
\Cref{lem:n9v2-Hoeffding} holds for every real \(\lambda\).
\end{proof}

\begin{firewall}[The original MAC range stops at two]
\label{fire:n9v17-MAC-range}
The basic MAC definition in V1 assumes the conditional MGF only for
\(|\lambda|\leq2\).  It gives the isolated \(L^2\) density but does not
automatically give the exponents in \eqref{eq:n9v17-five}.  A dyadic
five-row application must use bounded differences, bounded oscillation, or
an explicitly strengthened MGF interval.
\end{firewall}

\subsection{One integrable row and bounded remaining rows}
\label{sec:n9v17-mixed}

\begin{lemma}[Oscillation gives a pointwise normalized bound]
\label{lem:n9v17-osc}
If
\[
 r={e^\ell\over\int e^\ell\,d\nu},
 \qquad \operatorname{osc}\ell\leq\varepsilon,
\]
then
\begin{equation}
 e^{-\varepsilon}\leq r\leq e^\varepsilon
 \quad \nu\text{-almost surely}.                      \label{eq:n9v17-osc}
\end{equation}
\end{lemma}

\begin{proof}
If \(a=\mathop{\rm ess\,inf}\ell\), then
\(e^a\leq\int e^\ell\,d\nu\leq e^{a+\varepsilon}\).
Divide \(e^\ell\) by these two bounds.
\end{proof}

\begin{theorem}[Mixed sequential \(L^2\) compiler]
\label{thm:n9v17-mixed}
Fix one row \(k\).  Suppose
\begin{equation}
 \|r_k\|_{L^2(\nu_{k-1})}\leq L_k,                    \label{eq:n9v17-one-L2}
\end{equation}
and for every \(j\ne k\),
\begin{equation}
 \|r_j\|_{L^\infty(\nu_{j-1})}\leq B_j.              \label{eq:n9v17-bounded}
\end{equation}
Then
\begin{equation}
 \|R_m\|_{L^2(\nu_0)}
 \leq
 L_k
 \left(\prod_{j<k}B_j\right)^{1/2}
 \prod_{j>k}B_j.                                      \label{eq:n9v17-mixed}
\end{equation}
\end{theorem}

\begin{proof}
Let \(R_{k-1}=\prod_{j<k}r_j\).  Pointwise,
\[
 R_{k-1}\leq\prod_{j<k}B_j.
\]
Therefore
\begin{align*}
 \|R_{k-1}r_k\|_{L^2(\nu_0)}^2
 &=
 \int R_{k-1}r_k^2\,d\nu_{k-1}\\
 &\leq
 \left(\prod_{j<k}B_j\right)
 \int r_k^2\,d\nu_{k-1}\\
 &\leq
 \left(\prod_{j<k}B_j\right)L_k^2.
\end{align*}
Multiplication by each later \(r_j\) costs at most \(B_j\) pointwise.
\end{proof}

\begin{corollary}[One MAC row plus four oscillation rows]
\label{cor:n9v17-physical}
For the five-row ledger, suppose the centered action row satisfies MAC with
budget \(V_A\) under the intermediate measure at which it is inserted, and
the other four rows have oscillations \(\varepsilon_j\).  Then the full
density satisfies
\begin{equation}
 \|R_5\|_2
 \leq
 e^{V_A}
 \exp\!\left\{
 {1\over2}\sum_{j<k}\varepsilon_j
 +\sum_{j>k}\varepsilon_j
 \right\},                                             \label{eq:n9v17-physical}
\end{equation}
where \(k\) is the action-row position.
\end{corollary}

\begin{proof}
\Cref{cor:n9v12-MAC-L2} gives \(L_k\leq e^{V_A}\).
\Cref{lem:n9v17-osc} gives \(B_j\leq e^{\varepsilon_j}\).
Insert both in \eqref{eq:n9v17-mixed}.
\end{proof}

\begin{remark}[Order optimization]
\label{rem:n9v17-order}
In \eqref{eq:n9v17-physical}, a bounded row inserted before the \(L^2\) row
costs only half its logarithmic oscillation, while one inserted after costs
the full oscillation.  This favors putting the integrable row last, but the
order may be constrained by which reference measure makes its MAC estimate
valid.
\end{remark}

\subsection{Uniform cofinal density card}
\label{sec:n9v17-card}

\begin{definition}[Sequential full-density \(L^2\) card, SFD2]
\label{def:n9v17-SFD2}
For every local or global support used in the tightness theorem, SFD2
records:
\begin{enumerate}[label=\textup{(Q\arabic*)},leftmargin=4em]
\item one common transported measurable space and the exact five-row
factorization \eqref{eq:n9v17-step};
\item the insertion order and the intermediate probability for each row;
\item either all higher moments required by
\eqref{eq:n9v17-row1}--\eqref{eq:n9v17-rowj}, or one \(L^2\) row and
pointwise bounds \eqref{eq:n9v17-bounded} for all others;
\item uniformity of the resulting full-density \(L^2\) bound in the
regulator;
\item every truncation and exceptional-chart debit; and
\item exact action, determinant, coarea, transport, and counterterm
ownership.
\end{enumerate}
\end{definition}

\begin{theorem}[SFD2 loads the tightness programme]
\label{thm:n9v17-closure}
SFD2 plus FSGRC gives the full-measure Sobolev moments of
\Cref{thm:n9v16-loader}.  With LDLC, the bosonic constrained measures are
tight and their consistent cylinder limit lands as a Radon probability on
the declared local distribution space.
\end{theorem}

\begin{proof}
\Cref{thm:n9v17-dyadic} or \Cref{thm:n9v17-mixed} gives the uniform full
density \(L^2\) bound.  FSGRC and
\Cref{thm:n9v16-loader} give the density-loaded moments.  Apply
\Cref{thm:n9v14-Radon}.
\end{proof}

\begin{firewall}[Physical rows are still conditional]
\label{fire:n9v17-physical}
V17 proves composition, not the row estimates.  The determinant and
regular coarea rows have conditional oscillation cards, but the ambient
transport, counterterm, singular-chart, and untruncated action bounds must
all be proved on the same intermediate measures before SFD2 is loaded.
\end{firewall}

\begin{assessment}[Progress at ninth-series V17]
\label{ass:n9v17-status}
The full-density \(L^2\) requirement in V12--V16 now has two rigorous
compilers.  Arbitrary sequential rows pay the explicit escalating exponents
\((17,18,10,6,4)\) in the five-row case.  If four rows are bounded in
logarithmic oscillation, a single MAC action row suffices with the sharper
bound \eqref{eq:n9v17-physical}.  The invalid product of separate \(L^2\)
norms has been ruled out by an explicit counterexample.
\end{assessment}

\begin{programme}[Ninth-series V18: ambient transport row]
\label{prog:n9v17-V18}
V18 will attack the bounded-row premise that is least developed: ambient
transport.  It will derive the finite-dimensional Jacobian for a near-
identity common-refinement map, separate volume growth from local support,
and prove a trace/log-determinant bound whose oscillation is summable under
an explicit nuclear or finite-rank schedule.
\end{programme}

\begin{firewall}[Claim boundary after ninth-series V17]
\label{fire:n9v17-boundary}
V17 does not prove SFD2 for the physical regulator, the ambient transport
bound, full density convergence, nonlinear support, reconstruction, or the
SM--Einstein continuum.
\end{firewall}

\begin{theorem}[Ninth-series V17 result]
\label{thm:n9v17-result}
Sequential rowwise \(L^2\) bounds alone do not control the full density.
The dyadic higher-exponent compiler \eqref{eq:n9v17-dyadic} and the mixed
one-\(L^2\)-plus-bounded compiler \eqref{eq:n9v17-mixed} are valid
replacements.  Either one loads SFD2 and hence the density side of the
bosonic tightness theorem.
\end{theorem}

\begin{proof}
Use \Cref{prop:n9v17-counterexample,lem:n9v17-one-step,
thm:n9v17-dyadic,thm:n9v17-mixed,
thm:n9v17-closure}.
\end{proof}

\paragraph{Parent-version record.}
V17 was formed from
\texttt{Renewal\_SM\_Einstein\_Continuum\_Research\_Notes\_V16.tex}.
The V16 parent had \(5\,465\) logical source lines, \(193\,080\) bytes, and
SHA--256
\[
 \texttt{A5D2C22D82B287D94C6E9290EE76EA89B64249CAD975894958D0FD40E321FF2E}.
\]
The next compulsory companion audit is V20.

% =============================================================================
% END APPENDED NINTH-SERIES V17 MATERIAL
% =============================================================================


% =============================================================================
% BEGIN APPENDED NINTH-SERIES V18 MATERIAL
% =============================================================================

\clearpage
\section{Ninth-series V18: the near-identity ambient transport row}
\label{sec:v9v18}

\subsection{A finite-regulator common chart}
\label{sec:n9v18-setting}

Let the two adjacent finite regulators be represented, after spectator
completion, on \(\mathbb R^N\).  Let
\begin{equation}
 T(x)=x+v(x),\qquad A(x):=Dv(x),                      \label{eq:n9v18-T}
\end{equation}
where \(v\in C^1(\mathbb R^N;\mathbb R^N)\) and
\begin{equation}
 \sup_x\|A(x)\|_{\rm op}\leq\kappa<1.                 \label{eq:n9v18-kappa}
\end{equation}
The dimension \(N\) may depend on the regulator.  All estimates below are
stated in norms whose dependence on \(N\) remains visible.

\begin{lemma}[Global near-identity diffeomorphism]
\label{lem:n9v18-diffeomorphism}
Under \eqref{eq:n9v18-kappa}, \(T\) is a global \(C^1\) diffeomorphism of
\(\mathbb R^N\), and
\begin{align}
 (1-\kappa)|x-y|
 &\leq|T(x)-T(y)|
 \leq(1+\kappa)|x-y|,                                \label{eq:n9v18-biLip}\\
 \|DT^{-1}\|_{\rm op}
 &\leq(1-\kappa)^{-1}.                                \label{eq:n9v18-inverse}
\end{align}
\end{lemma}

\begin{proof}
The mean-value formula gives
\(|v(x)-v(y)|\leq\kappa|x-y|\), hence
\eqref{eq:n9v18-biLip} and injectivity.  Given \(z\), the map
\[
 \Phi_z(x):=z-v(x)
\]
is a contraction of the complete space \(\mathbb R^N\), so it has a unique
fixed point \(x\), which satisfies \(T(x)=z\).  Thus \(T\) is onto.  Since
\(DT=I+A\) and \(\|A\|<1\), the inverse function theorem and the Neumann
series give \eqref{eq:n9v18-inverse}.
\end{proof}

\subsection{The finite-dimensional logarithmic Jacobian}
\label{sec:n9v18-J}

\begin{lemma}[Nuclear logarithmic determinant bound]
\label{lem:n9v18-logdet}
For a real \(N\times N\) matrix \(A\) with
\(\|A\|_{\rm op}\leq\kappa<1\),
\begin{equation}
 \log\det(I+A)
 =
 \sum_{m=1}^{\infty}{(-1)^{m+1}\over m}\Tr(A^m),      \label{eq:n9v18-series}
\end{equation}
and
\begin{equation}
 |\log\det(I+A)|
 \leq{\|A\|_1\over1-\kappa},                          \label{eq:n9v18-logdet}
\end{equation}
where \(\|\cdot\|_1\) is the trace norm.
\end{lemma}

\begin{proof}
The spectrum of \(I+A\) lies in the disk centered at one of radius
\(\kappa\), so its real determinant is positive and the principal
logarithm is defined.  The matrix logarithm has the absolutely convergent
operator series in \eqref{eq:n9v18-series}.  Moreover,
\[
 |\Tr A^m|
 \leq\|A^m\|_1
 \leq\|A\|_1\|A\|_{\rm op}^{m-1}.
\]
Sum the coarser geometric series
\(\sum_{m\geq1}\kappa^{m-1}\).
\end{proof}

\begin{lemma}[Lipschitz logarithmic determinant]
\label{lem:n9v18-logdet-difference}
If \(\|A\|_{\rm op},\|B\|_{\rm op}\leq\kappa<1\), then
\begin{equation}
 |\log\det(I+A)-\log\det(I+B)|
 \leq{\|A-B\|_1\over1-\kappa}.                        \label{eq:n9v18-logdet-difference}
\end{equation}
\end{lemma}

\begin{proof}
Set \(A_t=B+t(A-B)\).  Then \(\|A_t\|_{\rm op}\leq\kappa\), and
\[
 {d\over dt}\log\det(I+A_t)
 =\Tr[(I+A_t)^{-1}(A-B)].
\]
The inverse norm is at most \((1-\kappa)^{-1}\).  Trace duality gives
\[
 |\Tr[(I+A_t)^{-1}(A-B)]|
 \leq{1\over1-\kappa}\|A-B\|_1.
\]
Integrate in \(t\).
\end{proof}

\begin{corollary}[Jacobian oscillation]
\label{cor:n9v18-J-osc}
If
\begin{equation}
 \omega_A:=
 \sup_{x,y}\|A(x)-A(y)\|_1<\infty,                   \label{eq:n9v18-omegaA}
\end{equation}
then
\begin{equation}
 \operatorname{osc}_x\log\det DT(x)
 \leq{\omega_A\over1-\kappa}.                         \label{eq:n9v18-J-osc}
\end{equation}
\end{corollary}

\begin{proof}
Apply \Cref{lem:n9v18-logdet-difference} to \(A(x),A(y)\) and take the
supremum.
\end{proof}

\begin{firewall}[Operator norm does not control volume growth]
\label{fire:n9v18-dimension}
The bound \(\|A\|_{\rm op}\leq\kappa\) alone permits
\(\log\det(I+A)=N\log(1+\kappa)\).  A cofinal Jacobian estimate therefore
needs trace-norm, finite-rank, or centered-oscillation control.  It is
invalid to suppress the regulator dimension in an operator-norm constant.
\end{firewall}

\begin{corollary}[Finite-rank variation]
\label{cor:n9v18-rank}
If every difference \(A(x)-A(y)\) has rank at most \(r\) and operator norm
at most \(\delta\), then
\begin{equation}
 \operatorname{osc}\log\det DT
 \leq{r\delta\over1-\kappa}.                           \label{eq:n9v18-rank}
\end{equation}
\end{corollary}

\begin{proof}
A rank-\(r\) matrix obeys \(\|B\|_1\leq r\|B\|_{\rm op}\).  Use
\Cref{cor:n9v18-J-osc}.
\end{proof}

\subsection{Base-density transport}
\label{sec:n9v18-base}

Let adjacent reference measures have positive \(C^1\) densities
\[
 d\lambda_i(x)=q_i(x)\,dx,\qquad i=0,1.
\]
Pulling \(\lambda_1\) back by \(T\) gives the logarithmic transport row
\begin{equation}
 \tau(x)
 :=
 \log q_1(Tx)-\log q_0(x)+\log\det DT(x).             \label{eq:n9v18-tau}
\end{equation}
Put
\[
 g(x):=\log q_1(x)-\log q_0(x).
\]

\begin{theorem}[Ambient transport oscillation]
\label{thm:n9v18-transport}
Suppose
\begin{align}
 \operatorname{osc}g&\leq\varepsilon_{\rm base},      \label{eq:n9v18-base-osc}\\
 \sup_x\|\nabla\log q_1(x)\|&\leq L_q,                \label{eq:n9v18-Lq}\\
 \sup_x|v(x)|&\leq\rho.                               \label{eq:n9v18-rho}
\end{align}
Then
\begin{equation}
 \operatorname{osc}\tau
 \leq
 \varepsilon_{\rm base}
 +2L_q\rho
 +{\omega_A\over1-\kappa}.                            \label{eq:n9v18-transport}
\end{equation}
\end{theorem}

\begin{proof}
Decompose
\[
 \log q_1(Tx)-\log q_0(x)
 =
 g(x)+[\log q_1(Tx)-\log q_1(x)].
\]
The bracket has absolute value at most \(L_q\rho\), hence oscillation at
most \(2L_q\rho\).  Add \eqref{eq:n9v18-base-osc} and the Jacobian
oscillation \eqref{eq:n9v18-J-osc}.
\end{proof}

\begin{corollary}[Linear transports have no centered Jacobian row]
\label{cor:n9v18-linear}
If \(T(x)=Mx+b\) with \(M\) constant, then \(\log|\det M|\) is
field-independent and disappears under probability normalization.
Only the transported base-density term in \eqref{eq:n9v18-tau} can
contribute to its oscillation.
\end{corollary}

\begin{proof}
The derivative \(DT=M\) is constant, so its log determinant has zero
oscillation.  Adding a constant to a logarithmic density does not change
its normalized exponential.
\end{proof}

\subsection{Spectator completion}
\label{sec:n9v18-spectator}

\begin{proposition}[Normalized spectators do not alter a marginal]
\label{prop:n9v18-spectator}
Let \(X,Z\) be measurable spaces and let
\[
 d\Lambda(x,z)=d\lambda(x)\,k_x(dz),
\]
where \(k_x\) is a probability kernel.  Then the \(X\)-marginal of
\(\Lambda\) is exactly \(\lambda\).  If an unnormalized kernel
\(\widetilde k_x\) is used instead, its mass
\[
 Z_{\rm sp}(x):=\widetilde k_x(Z)
\]
contributes the effective logarithmic row \(\log Z_{\rm sp}(x)\).
\end{proposition}

\begin{proof}
For measurable \(B\subset X\),
\[
 \Lambda(B\times Z)
 =\int_Bk_x(Z)\,\lambda(dx)=\lambda(B).
\]
For an unnormalized kernel, disintegrate
\(\widetilde k_x=Z_{\rm sp}(x)k_x\); the factor remains in the marginal
density.
\end{proof}

\begin{firewall}[Spectator normalization may be field-dependent]
\label{fire:n9v18-spectator}
Adding refinement variables is neutral only when their conditional law is
normalized for every old field configuration.  A field-dependent
partition function is an effective action and must be assigned to the
action or transport ledger.
\end{firewall}

\subsection{Cofinal transport card}
\label{sec:n9v18-card}

\begin{theorem}[Summable near-identity transport row]
\label{thm:n9v18-sum}
For adjacent transports \(T_n=I+v_n\), suppose one common
\(\kappa<1\) satisfies \eqref{eq:n9v18-kappa}, and
\begin{equation}
 \sum_n\left[
 \varepsilon_{{\rm base},n}
 +L_{q,n}\rho_n
 +{\omega_{A,n}\over1-\kappa}
 \right]<\infty.                                      \label{eq:n9v18-sum}
\end{equation}
Then the ambient transport rows have summable logarithmic oscillation and
may serve as bounded rows in the mixed full-density compiler
\Cref{thm:n9v17-mixed}.
\end{theorem}

\begin{proof}
Sum \eqref{eq:n9v18-transport}.  The normalized transport factor is
pointwise bounded by the exponential of its oscillation by
\Cref{lem:n9v17-osc}.
\end{proof}

\begin{definition}[Near-identity ambient transport card, NIATC]
\label{def:n9v18-NIATC}
NIATC records:
\begin{enumerate}[label=\textup{(N\arabic*)},leftmargin=4em]
\item the common-refinement space and every normalized spectator kernel;
\item the map \(T_n=I+v_n\) and a uniform derivative margin
\(\|Dv_n\|_{\rm op}\leq\kappa<1\);
\item the trace-norm derivative oscillation \(\omega_{A,n}\);
\item the base log-density oscillation and displacement-gradient budget;
\item the summability condition \eqref{eq:n9v18-sum};
\item chart-exit and exceptional-set probabilities; and
\item compatibility with the physical-graph transport already owned by
the coarea row.
\end{enumerate}
\end{definition}

\begin{theorem}[Conditional NIATC closure]
\label{thm:n9v18-closure}
NIATC supplies a bounded summable ambient-transport row for SFD2.  If the
action is the single \(L^2\) row and the determinant, regular coarea, and
counterterm rows also have summable oscillations, the full-density bound
\eqref{eq:n9v17-physical} is uniform.
\end{theorem}

\begin{proof}
\Cref{thm:n9v18-sum} gives the transport oscillation.  Add the declared
chart-exit debit to the V98 exceptional-set compiler.  Apply
\Cref{cor:n9v17-physical} to the five owned rows.
\end{proof}

\begin{firewall}[No infinite-dimensional Lebesgue Jacobian]
\label{fire:n9v18-infinite}
V18 is a finite-regulator theorem.  There is no translation-invariant
Lebesgue measure in infinite dimension.  A continuum nonlinear Gaussian
change of variables would require a Cameron--Martin/Ramer-type
quasi-invariance theorem and a renormalized determinant.  NIATC instead
requires the finite Jacobian oscillations to be summable before the limit.
\end{firewall}

\begin{assessment}[Progress at ninth-series V18]
\label{ass:n9v18-status}
The ambient transport row now has an explicit bounded compiler.  The
nonlinear Jacobian is controlled by trace-norm variation, linear
field-independent determinants disappear after centering, and normalized
spectators are exactly neutral.  The remaining physical work is the
construction of common-refinement maps satisfying the nuclear schedule
\eqref{eq:n9v18-sum}.
\end{assessment}

\begin{programme}[Ninth-series V19: common-refinement composition]
\label{prog:n9v18-V19}
V19 will compose two common-refinement transports, prove the exact
Jacobian cocycle, and derive a path-independence criterion from summable
two-route defects.  This will connect NIATC to the regulator-independence
gate before the compulsory V20 audit.
\end{programme}

\begin{firewall}[Claim boundary after ninth-series V18]
\label{fire:n9v18-boundary}
V18 does not verify NIATC for the physical discretizations, prove
regulator independence, nonlinear support, reconstruction, or the full
SM--Einstein continuum.
\end{firewall}

\begin{theorem}[Ninth-series V18 result]
\label{thm:n9v18-result}
A finite near-identity common-refinement map has logarithmic Jacobian
oscillation bounded by
\(\omega_A/(1-\kappa)\).  Including base-density and displacement terms
gives \eqref{eq:n9v18-transport}; its summability loads the bounded ambient
transport row in SFD2.
\end{theorem}

\begin{proof}
Use \Cref{lem:n9v18-diffeomorphism,
lem:n9v18-logdet-difference,cor:n9v18-J-osc,
thm:n9v18-transport,thm:n9v18-sum,
thm:n9v18-closure}.
\end{proof}

\paragraph{Parent-version record.}
V18 was formed from
\texttt{Renewal\_SM\_Einstein\_Continuum\_Research\_Notes\_V17.tex}.
The V17 parent had \(5\,888\) logical source lines, \(206\,106\) bytes, and
SHA--256
\[
 \texttt{C47D321E156AD6C7783A471033FF0A3BA15179048945C8F4C5148CCEBA98D649}.
\]
The next compulsory companion audit is V20.

% =============================================================================
% END APPENDED NINTH-SERIES V18 MATERIAL
% =============================================================================


% =============================================================================
% BEGIN APPENDED NINTH-SERIES V19 MATERIAL
% =============================================================================

\clearpage
\section{Ninth-series V19: common-refinement cocycles and path independence}
\label{sec:v9v19}

\subsection{Exact composition of transport densities}
\label{sec:n9v19-cocycle}

Let \(X_a,X_b,X_c\) be finite-dimensional regulator charts with positive
base densities
\[
 d\lambda_i=q_i(x_i)\,dx_i.
\]
Let \(T_{b,a}:X_a\to X_b\) and \(T_{c,b}:X_b\to X_c\) be
orientation-preserving \(C^1\) diffeomorphisms, and define
\[
 T_{c,a}:=T_{c,b}\circ T_{b,a}.
\]
The logarithmic pullback row is
\begin{equation}
 \tau_{b,a}(x)
 :=
 \log q_b(T_{b,a}x)-\log q_a(x)
 +\log\det DT_{b,a}(x).                               \label{eq:n9v19-tau}
\end{equation}

\begin{theorem}[Exact ambient transport cocycle]
\label{thm:n9v19-cocycle}
The rows obey
\begin{equation}
 \tau_{c,a}
 =
 \tau_{b,a}
 +\tau_{c,b}\circ T_{b,a}.                            \label{eq:n9v19-cocycle}
\end{equation}
\end{theorem}

\begin{proof}
The chain rule gives
\[
 DT_{c,a}(x)
 =
 DT_{c,b}(T_{b,a}x)\,DT_{b,a}(x),
\]
hence
\[
 \log\det DT_{c,a}
 =
 \log\det DT_{c,b}\circ T_{b,a}
 +\log\det DT_{b,a}.
\]
Also,
\[
 \log q_c(T_{c,a}x)-\log q_a(x)
 =
 [\log q_c(T_{c,b}T_{b,a}x)-\log q_b(T_{b,a}x)]
 +[\log q_b(T_{b,a}x)-\log q_a(x)].
\]
Add the two identities.
\end{proof}

\begin{corollary}[Normalized density cocycle]
\label{cor:n9v19-density-cocycle}
If every action, determinant, coarea, transport, and counterterm row has an
exact endpoint cocycle under the same maps, their sum has an exact cocycle.
Its normalized exponential gives the same transported probability whether
the step \(a\to c\) is taken directly or through \(b\).
\end{corollary}

\begin{proof}
Add the row cocycles.  Equality of the unnormalized pullback densities is
preserved by division by their common total mass.
\end{proof}

\begin{firewall}[Row cocycles must use the same transport]
\label{fire:n9v19-same-map}
An endpoint counterterm cocycle expressed in one coordinate transport
cannot be combined with a determinant or coarea cocycle expressed in
another without inserting the coordinate-change row.  Equality of endpoint
labels is not equality of pulled-back densities.
\end{firewall}

\subsection{A quantitative refinement diamond}
\label{sec:n9v19-diamond}

Consider a diamond
\[
 a\longrightarrow b\longrightarrow d,
 \qquad
 a\longrightarrow c\longrightarrow d.
\]
Let
\begin{equation}
 U:=T_{d,b}\circ T_{b,a},\qquad
 V:=T_{d,c}\circ T_{c,a}.                             \label{eq:n9v19-UV}
\end{equation}
Both maps run from \(X_a\) to \(X_d\).  Define the terminal route defect
\begin{equation}
 R:=U\circ V^{-1}:X_d\longrightarrow X_d.             \label{eq:n9v19-R}
\end{equation}

\begin{theorem}[Near-identity diamond defect]
\label{thm:n9v19-diamond}
Suppose \(R=I+w\) satisfies the V18 near-identity hypotheses on \(X_d\):
\begin{align}
 \|Dw\|_{\rm op}&\leq\kappa<1,                        \label{eq:n9v19-kappa}\\
 \|w\|_\infty&\leq\rho,                               \label{eq:n9v19-rho}\\
 \sup_{x,y}\|Dw(x)-Dw(y)\|_1&\leq\omega,              \label{eq:n9v19-omega}\\
 \|\nabla\log q_d\|_\infty&\leq L_d.                 \label{eq:n9v19-L}
\end{align}
Then the two pullbacks \(U^*\lambda_d\) and \(V^*\lambda_d\), represented
on \(X_a\), have logarithmic density ratio of oscillation at most
\begin{equation}
 \eta_\diamond
 :=
 2L_d\rho+{\omega\over1-\kappa}.                      \label{eq:n9v19-eta}
\end{equation}
After normalization to probabilities \(\widehat\lambda_U,\widehat\lambda_V\),
\begin{equation}
 \|\widehat\lambda_U-\widehat\lambda_V\|_{\rm TV}
 \leq{e^{\eta_\diamond}-1\over2}.                     \label{eq:n9v19-diamond-TV}
\end{equation}
\end{theorem}

\begin{proof}
Since \(U=R\circ V\),
\[
 U^*\lambda_d=V^*(R^*\lambda_d).
\]
The logarithmic ratio between \(R^*\lambda_d\) and \(\lambda_d\) is the V18
transport row with identical initial and terminal base densities.  Thus
the intrinsic base difference vanishes, and
\Cref{thm:n9v18-transport} gives the bound
\eqref{eq:n9v19-eta}.  Composition with the bijection \(V\) does not change
oscillation.  The normalized bounded-correction estimate
\Cref{lem:n9v6-bounded-correction} gives
\eqref{eq:n9v19-diamond-TV}.
\end{proof}

\begin{remark}[Exact commuting diamond]
\label{rem:n9v19-exact}
If \(U=V\), then \(R=I\), all three defect parameters vanish, and the two
routes agree exactly.
\end{remark}

\subsection{Full-density diamond defects}
\label{sec:n9v19-full}

Let \(\Lambda_d=e^{\Phi_d}\lambda_d\) be the full owned terminal measure,
where \(\Phi_d\) includes the action, determinant, coarea, and counterterm
rows not already included in \(q_d\).

\begin{proposition}[Full endpoint route defect]
\label{prop:n9v19-full}
Under the hypotheses of \Cref{thm:n9v19-diamond}, suppose additionally
\begin{equation}
 \operatorname{osc}_x[\Phi_d(Rx)-\Phi_d(x)]
 \leq\eta_\Phi.                                       \label{eq:n9v19-Phi}
\end{equation}
Then the normalized full pullbacks along \(U\) and \(V\) differ in total
variation by at most
\begin{equation}
 {e^{\eta_\diamond+\eta_\Phi}-1\over2}.               \label{eq:n9v19-full-TV}
\end{equation}
\end{proposition}

\begin{proof}
Relative to \(\Lambda_d\), the pullback by \(R\) adds the base/Jacobian
row of \Cref{thm:n9v19-diamond} and the endpoint action difference
\(\Phi_d\circ R-\Phi_d\).  Oscillations add.  Apply
\Cref{lem:n9v6-bounded-correction}.
\end{proof}

\begin{firewall}[Endpoint regularity is a separate debit]
\label{fire:n9v19-endpoint}
Small coordinate displacement does not imply
\eqref{eq:n9v19-Phi} for an unbounded or highly oscillatory full action.
The endpoint density must have a proved response bound, or the route defect
must be estimated directly in an \(L^2\) or entropy norm.
\end{firewall}

\subsection{Summable homotopy of refinement paths}
\label{sec:n9v19-path}

\begin{theorem}[Summable diamond path independence]
\label{thm:n9v19-path}
Suppose two cofinal refinement paths are related beyond level \(n\) by a
sequence of elementary diamonds whose normalized full-density debits are
\(d_k\), and
\begin{equation}
 \sum_{k=1}^{\infty}d_k<\infty.                       \label{eq:n9v19-diamond-sum}
\end{equation}
Then the total-variation distance between the two transported endpoint
measures after level \(n\) is at most
\begin{equation}
 \sum_{k\geq n}d_k,                                   \label{eq:n9v19-tail}
\end{equation}
and hence tends to zero.
\end{theorem}

\begin{proof}
Replace one path by the other one diamond at a time.  Total variation is
invariant under a common measurable bijective transport and satisfies the
triangle inequality.  Summing the elementary diamond bounds gives
\eqref{eq:n9v19-tail}.  The tail of a convergent nonnegative series tends
to zero.
\end{proof}

\begin{corollary}[Bound from near-identity diamond parameters]
\label{cor:n9v19-path}
If the \(k\)-th diamond satisfies
\[
 d_k\leq{1\over2}
 \left[
 \exp(\eta_{\diamond,k}+\eta_{\Phi,k})-1
 \right]
\]
and
\begin{equation}
 \sum_k(\eta_{\diamond,k}+\eta_{\Phi,k})<\infty,      \label{eq:n9v19-eta-sum}
\end{equation}
then \eqref{eq:n9v19-diamond-sum} holds.
\end{corollary}

\begin{proof}
The exponents tend to zero.  For all sufficiently large \(k\),
\(e^x-1\leq2x\); finitely many initial terms are harmless.
\end{proof}

\subsection{Common-refinement equality of two regulator limits}
\label{sec:n9v19-regulator}

\begin{theorem}[Common-refinement bridge]
\label{thm:n9v19-bridge}
Let \(\mu_n\to\mu\) and \(\nu_n\to\nu\) in total variation on a common
local observable space.  Suppose there are cofinal subsequences
\(n_k,m_k\) and common-refinement probabilities \(\rho_k\) such that
\begin{equation}
 \|\mu_{n_k}-\rho_k\|_{\rm TV}
 +\|\nu_{m_k}-\rho_k\|_{\rm TV}
 \longrightarrow0.                                   \label{eq:n9v19-bridge}
\end{equation}
Then
\begin{equation}
 \mu=\nu.                                              \label{eq:n9v19-equality}
\end{equation}
\end{theorem}

\begin{proof}
The triangle inequality gives
\begin{align*}
 \|\mu-\nu\|_{\rm TV}
 &\leq
 \|\mu-\mu_{n_k}\|_{\rm TV}
 +\|\mu_{n_k}-\rho_k\|_{\rm TV}\\
 &\quad+
 \|\rho_k-\nu_{m_k}\|_{\rm TV}
 +\|\nu_{m_k}-\nu\|_{\rm TV}.
\end{align*}
Every term tends to zero.
\end{proof}

\begin{corollary}[Equality of Radon landings]
\label{cor:n9v19-Radon}
If two tight regulator schemes satisfy LDLC and their local cylinder
limits are joined by \Cref{thm:n9v19-bridge} on a countable total cylinder
family, their Radon limits on the declared separable local distribution
space are equal.
\end{corollary}

\begin{proof}
The bridge identifies every finite-dimensional cylinder law.  The
countable-coordinate uniqueness argument in
\Cref{thm:n9v11-Radon,thm:n9v14-Radon} identifies the Borel probabilities.
\end{proof}

\subsection{Common-refinement path card}
\label{sec:n9v19-card}

\begin{definition}[Common-refinement path card, CRPC]
\label{def:n9v19-CRPC}
CRPC records:
\begin{enumerate}[label=\textup{(C\arabic*)},leftmargin=4em]
\item a directed common-refinement graph and common transported local
spaces;
\item exact cocycles for every owned density row along each edge;
\item for noncommuting diamonds, the terminal defect map
\eqref{eq:n9v19-R} and its NIATC parameters;
\item the full endpoint response debit \eqref{eq:n9v19-Phi};
\item the summable diamond condition \eqref{eq:n9v19-eta-sum};
\item summable chart-exit and singular-stratum probabilities; and
\item cofinal common refinements joining every pair of regulator schemes
whose limits are to be identified.
\end{enumerate}
\end{definition}

\begin{theorem}[Conditional regulator independence]
\label{thm:n9v19-closure}
Assume ACDC/SFD2 gives local total-variation limits, LDLC gives their Radon
landings, and CRPC holds.  Then the limit is independent of refinement
path and of every regulator scheme joined by the declared common
refinements.
\end{theorem}

\begin{proof}
\Cref{thm:n9v19-path} gives path independence.
\Cref{thm:n9v19-bridge} identifies local limits from different schemes.
\Cref{cor:n9v19-Radon} identifies the Radon measures.
\end{proof}

\begin{firewall}[Existence of common refinements is not proved]
\label{fire:n9v19-physical}
V19 proves the comparison theorem once compatible common refinements and
their quantitative diamond bounds exist.  It does not construct those
refinements for arbitrary lattice, triangulation, gauge bundle, spin
structure, or metric discretizations.
\end{firewall}

\begin{assessment}[Progress at ninth-series V19]
\label{ass:n9v19-status}
The ambient transport row now has an exact cocycle and a quantitative
path-independence theorem.  A noncommuting refinement diamond is reduced to
a terminal near-identity map plus a full endpoint response.  Summable
diamond debits make refinement paths asymptotically identical, and a
common-refinement bridge identifies regulator limits.
\end{assessment}

\begin{programme}[Ninth-series V20: compulsory companion audit]
\label{prog:n9v19-V20}
V20 will inspect the newest YM and NS notes, compare their new locality or
common-refinement results with NIATC and CRPC, and choose the next
acquisition among physical common-refinement construction, nonlinear
metric support, and the full-density moment row.  The chosen branch will
continue in V21.
\end{programme}

\begin{firewall}[Claim boundary after ninth-series V19]
\label{fire:n9v19-boundary}
V19 does not verify CRPC for physical regulator schemes, prove the full
density cards, nonlinear support, reconstruction, or the SM--Einstein
continuum.
\end{firewall}

\begin{theorem}[Ninth-series V19 result]
\label{thm:n9v19-result}
Exact transports obey the cocycle \eqref{eq:n9v19-cocycle}.  A
near-identity diamond has the total-variation debit
\eqref{eq:n9v19-full-TV}; summability gives path independence, and the
common-refinement bridge \eqref{eq:n9v19-bridge} identifies the resulting
local and Radon limits.
\end{theorem}

\begin{proof}
Use \Cref{thm:n9v19-cocycle,thm:n9v19-diamond,
prop:n9v19-full,thm:n9v19-path,
thm:n9v19-bridge,thm:n9v19-closure}.
\end{proof}

\paragraph{Parent-version record.}
V19 was formed from
\texttt{Renewal\_SM\_Einstein\_Continuum\_Research\_Notes\_V18.tex}.
The V18 parent had \(6\,255\) logical source lines, \(218\,230\) bytes, and
SHA--256
\[
 \texttt{0C12B64A014F4EBC43B6DA7A5F7A696E90DDF0937289997900B326FA7BE6E56A}.
\]
The next version is the compulsory V20 companion audit.

% =============================================================================
% END APPENDED NINTH-SERIES V19 MATERIAL
% =============================================================================


% =============================================================================
% BEGIN APPENDED NINTH-SERIES V20 MATERIAL
% =============================================================================

\clearpage
\section{Ninth-series V20: compulsory audit and normalized chart exit}
\label{sec:v9v20}

\subsection{Files and internal version status}
\label{sec:n9v20-files}

The newest Yang--Mills files inspected were
\begin{center}
\small\ttfamily
Renewal\_Yang\_Mills\_Mass\_Gap\_Research\_Notes\_V62.tex,\\
Renewal\_Yang\_Mills\_Mass\_Gap\_Research\_Notes\_V63.tex.
\end{center}
They are byte-identical.  Each has \(26\,916\) logical lines, \(961\,508\)
bytes, and SHA--256
\[
 \texttt{D559406640DA883813024610CA08DC1AE7328059280E12CE2444C13B44013261}.
\]
Both end internally at V62; no distinct V63 theorem is counted.  The
Navier--Stokes file remains V26, with \(5\,319\) lines, \(278\,283\) bytes,
and SHA--256
\[
 \texttt{82C887F8C2C69E4F28FAD7C3DC8DE5B9099CB5A4BF431EBA1FFF3E91935978AD}.
\]

\subsection{Yang--Mills V60--V62}
\label{sec:n9v20-YM}

YM V60 proves stability of a positive spectral gap and a form-resolvent
comparison under a small positive relative-form subtraction.  It also
proves exact number-distance shielding and superexponentially small
occupation-shell matrix elements between fixed Weyl vectors.  Temporal
gap persistence is kept separate from spatial clustering.

V61 proves cumulant shielding and an anchored distinguished-insertion tree
bound, which loads its local tail-disappearance card.  It also proves a
no-go theorem: positivity and top-shell support alone do not imply spatial
locality.  A separate anchored boundary-decay card remains necessary.

V62 proves stabilization of connected coefficients below the reachable
occupation number, a two-cutoff comparison, diagonal occupation-regulator
independence, and local restoration of the analytic Wilson tail.  Its dense
finite-particle core still needs interacting endpoint continuity.  The next
YM acquisition is precisely a normalized interacting chart-exit bound
from a full-density \(L^q\) estimate.

\begin{proposition}[Type-correct V20 transfer]
\label{prop:n9v20-YM}
The YM V60--V62 results do not prove CRPC, SFD2, or chart exit for the
SM--Einstein measure.  They support two already compatible conclusions:
\begin{enumerate}[label=\textup{(\roman*)},leftmargin=3.8em]
\item relative form and cutoff-support statements must remain separate
from spatial locality; and
\item an exit probability under the interacting measure must be bounded
using the normalized full density, not the Gaussian reference tail alone.
\end{enumerate}
\end{proposition}

\begin{proof}
The YM theorems concern an occupation Feshbach problem and a quasi-free
Weyl core, so their constants do not identify SM fields or constraint
charts.  Their logical separations are measure-theoretic and operator-
theoretic, however.  Item (i) agrees with the distinct FSGRC and NIATC
cards.  Item (ii) is exactly the missing exceptional-set input in UMCC,
TSCC, NIATC, and LDLC.
\end{proof}

\begin{firewall}[Diagonal cutoff independence is programme-specific]
\label{fire:n9v20-YM}
The YM two-cutoff comparison proves independence of its occupation schedule
inside a fixed programme.  It does not construct common refinements between
SM lattice, metric, spin, or bundle regulators and cannot be substituted
for CRPC.
\end{firewall}

\subsection{Navier--Stokes status}
\label{sec:n9v20-NS}

NS V26 has not changed.  Its deterministic Oseen-kernel derivative norms
do not supply a normalized density or a reference tail for any SM chart.

\begin{proposition}[NS non-transfer at V20]
\label{prop:n9v20-NS}
No current NS result loads a chart-exit row of the SM programme.
\end{proposition}

\begin{proof}
An interacting chart-exit estimate requires two probabilities and their
Radon--Nikodym derivative.  These objects are absent from the NS result.
\end{proof}

\subsection{Audit decision}
\label{sec:n9v20-decision}

\begin{audit}[V20 companion conclusion]
\label{aud:n9v20}
The next acquisition is a normalized full-density exit compiler.  It will
turn a reference tail and an \(L^q\), entropy, or oscillation bound for the
complete density into the summable adjacent exceptional-set debit required
by the coarea, crossing, action-truncation, and ambient-transport cards.
\end{audit}

\begin{decision}[V21 acquisition order]
\label{dec:n9v20-V21}
V21 will prove:
\[
 \mu(E)\leq
 \left\|{d\mu\over d\nu}\right\|_{L^q(\nu)}
 \nu(E)^{1-1/q},
\]
derive the \(q=2\) consequence of SFD2, give entropy and bounded-oscillation
alternatives, and state a diagonal cutoff-selection theorem with no hidden
normalization assumption.
\end{decision}

\begin{assessment}[Progress at ninth-series V20]
\label{ass:n9v20-status}
The audit selects a theorem needed simultaneously by four existing cards.
It also prevents the reference Gaussian tail from being mistaken for the
physical constrained tail.  V21 will make the density-norm loss and cutoff
schedule explicit.
\end{assessment}

\begin{programme}[Ninth-series V21: full-density chart exit]
\label{prog:n9v20-V21}
Carry out \Cref{dec:n9v20-V21} and insert the resulting exit card into the
conditional continuum ledger.
\end{programme}

\begin{firewall}[Claim boundary after ninth-series V20]
\label{fire:n9v20-boundary}
V20 is an audit.  It does not prove the full-density norm, any physical
exit tail, CRPC, nonlinear support, reconstruction, or the SM--Einstein
continuum.
\end{firewall}

\begin{theorem}[Ninth-series V20 result]
\label{thm:n9v20-result}
The companion audit confirms that normalized chart exit, not a reference
tail alone, is the next common bottleneck.  V21 will load it from the full
density and keep spatial locality and common-refinement independence as
separate requirements.
\end{theorem}

\begin{proof}
Use \Cref{prop:n9v20-YM,prop:n9v20-NS,aud:n9v20}.
\end{proof}

\paragraph{Parent-version record.}
V20 was formed from
\texttt{Renewal\_SM\_Einstein\_Continuum\_Research\_Notes\_V19.tex}.
The V19 parent had \(6\,631\) logical source lines, \(230\,978\) bytes, and
SHA--256
\[
 \texttt{4FC7DE00186317D150ECFB6EDF6AD3540B2FB8251DFBDDB6D8686036D7F92F40}.
\]
The next compulsory companion audit is V25.

% =============================================================================
% END APPENDED NINTH-SERIES V20 MATERIAL
% =============================================================================


% =============================================================================
% BEGIN APPENDED NINTH-SERIES V21 MATERIAL
% =============================================================================

\clearpage
\section{Ninth-series V21: normalized full-density chart exit}
\label{sec:v9v21}

\subsection{The \(L^q\) exit inequality}
\label{sec:n9v21-Lq}

Let \(\nu\) be a reference probability, let \(\mu\ll\nu\), and put
\[
 r={d\mu\over d\nu}.
\]

\begin{theorem}[Full-density exit bound]
\label{thm:n9v21-Lq}
For \(q>1\), \(q'=q/(q-1)\), and every measurable event \(E\),
\begin{equation}
 \mu(E)
 \leq
 \|r\|_{L^q(\nu)}\,\nu(E)^{1/q'}.                    \label{eq:n9v21-Lq}
\end{equation}
\end{theorem}

\begin{proof}
Apply Holder's inequality to
\[
 \mu(E)=\int r\,\mathbf1_E\,d\nu.
\]
The \(L^{q'}(\nu)\) norm of \(\mathbf1_E\) is
\(\nu(E)^{1/q'}\).
\end{proof}

\begin{corollary}[SFD2 square-root tail]
\label{cor:n9v21-L2}
If SFD2 gives
\(\|r\|_{L^2(\nu)}\leq R\), then
\begin{equation}
 \mu(E)\leq R\sqrt{\nu(E)}.                            \label{eq:n9v21-L2}
\end{equation}
\end{corollary}

\begin{proof}
Take \(q=q'=2\) in \Cref{thm:n9v21-Lq}.
\end{proof}

\begin{corollary}[Centered logarithmic MGF exit]
\label{cor:n9v21-MGF}
Let
\[
 r={e^\ell\over\int e^\ell\,d\nu},\qquad
 u=\ell-\int\ell\,d\nu.
\]
If
\[
 \int e^{qu}\,d\nu\leq e^{q^2V/2},
\]
then
\begin{equation}
 \mu(E)
 \leq
 e^{qV/2}\nu(E)^{1/q'}.                               \label{eq:n9v21-MGF}
\end{equation}
\end{corollary}

\begin{proof}
\Cref{thm:n9v12-log} gives \(\|r\|_q\leq e^{qV/2}\).
Apply \Cref{thm:n9v21-Lq}.
\end{proof}

\begin{firewall}[The density is the complete normalized density]
\label{fire:n9v21-full}
The norm in \eqref{eq:n9v21-Lq} must belong to the full constrained
probability relative to the reference used for the tail.  An action-only or
determinant-only density norm does not control the exit probability after
the other rows are multiplied.
\end{firewall}

\subsection{Oscillation and entropy alternatives}
\label{sec:n9v21-alternatives}

\begin{proposition}[Bounded logarithmic density exit]
\label{prop:n9v21-osc}
If
\[
 r={e^\ell\over\int e^\ell\,d\nu},
 \qquad\operatorname{osc}\ell\leq\varepsilon,
\]
then
\begin{equation}
 \mu(E)\leq e^\varepsilon\nu(E).                      \label{eq:n9v21-osc}
\end{equation}
\end{proposition}

\begin{proof}
\Cref{lem:n9v17-osc} gives \(r\leq e^\varepsilon\) almost surely.
Integrate over \(E\).
\end{proof}

\begin{theorem}[Relative-entropy exit bound]
\label{thm:n9v21-entropy}
Let
\[
 H:=\operatorname{Ent}(\mu\mid\nu)<\infty.
\]
If \(0<\delta:=\nu(E)<1\), then
\begin{equation}
 \mu(E)
 \leq {H+\log2\over\log(1/\delta)}.                   \label{eq:n9v21-entropy}
\end{equation}
If \(\delta=0\), then \(\mu(E)=0\).
\end{theorem}

\begin{proof}
Apply the entropy variational bound \Cref{lem:n9v12-entropy} to
\(Y=\mathbf1_E\):
\[
 \mu(E)
 \leq{1\over t}
 \left[
 H+\log(1-\delta+\delta e^t)
 \right].
\]
Choose \(t=\log(1/\delta)\).  Then
\[
 1-\delta+\delta e^t=2-\delta\leq2,
\]
which proves \eqref{eq:n9v21-entropy}.  If \(\nu(E)=0\), absolute
continuity gives \(\mu(E)=0\).
\end{proof}

\begin{remark}[Entropy is weaker for rare events]
\label{rem:n9v21-entropy}
A uniform positive entropy bound decays only like
\(1/\log(1/\nu(E))\).  The \(L^q\) route gives a power of the reference
tail and is usually better for summable chart exits.  Entropy becomes
competitive when \(H_n\to0\) with the regulator.
\end{remark}

\subsection{Adjacent exceptional-set debit}
\label{sec:n9v21-adjacent}

\begin{theorem}[Summable adjacent exit criterion]
\label{thm:n9v21-adjacent}
For every adjacent step \(n\), let \(E_n\) be a bad set on the common
transported space.  Suppose
\[
 d\mu_{i}=r_i\,d\nu_i\qquad(i=n,n+1)
\]
and, for one \(q>1\),
\begin{equation}
 \sum_n\left[
 \|r_n\|_{L^q(\nu_n)}
 \nu_n(E_n)^{1/q'}
 +
 \|r_{n+1}\|_{L^q(\nu_{n+1})}
 \nu_{n+1}(E_n)^{1/q'}
 \right]<\infty.                                      \label{eq:n9v21-adjacent}
\end{equation}
Then
\begin{equation}
 \sum_n[\mu_n(E_n)+\mu_{n+1}(E_n)]<\infty.            \label{eq:n9v21-adjacent-target}
\end{equation}
\end{theorem}

\begin{proof}
Apply \Cref{thm:n9v21-Lq} to the two probabilities at every step and sum.
\end{proof}

\begin{corollary}[Uniform density norm]
\label{cor:n9v21-uniform}
If both adjacent full densities have \(L^q\) norm at most \(R\), it is
sufficient that
\begin{equation}
 \sum_n[
 \nu_n(E_n)^{1/q'}+
 \nu_{n+1}(E_n)^{1/q'}]<\infty.                       \label{eq:n9v21-uniform}
\end{equation}
\end{corollary}

\begin{proof}
Factor \(R\) out of \eqref{eq:n9v21-adjacent}.
\end{proof}

\subsection{Diagonal cutoff selection}
\label{sec:n9v21-diagonal}

\begin{theorem}[Noncircular exit cutoff]
\label{thm:n9v21-diagonal}
For each \(n\), let \(E_{n}(R)\) decrease as \(R\to\infty\).  Suppose the
full, untruncated densities obey
\begin{equation}
 \|r_n\|_{L^q(\nu_n)}\leq K_n<\infty                 \label{eq:n9v21-K}
\end{equation}
independently of \(R\), and
\begin{equation}
 \nu_n(E_n(R))\longrightarrow0.                       \label{eq:n9v21-tail-zero}
\end{equation}
Given any positive summable sequence \(a_n\), one can choose \(R_n\) so
that
\begin{equation}
 \mu_n(E_n(R_n))\leq a_n.                             \label{eq:n9v21-diagonal}
\end{equation}
\end{theorem}

\begin{proof}
By \eqref{eq:n9v21-tail-zero}, choose \(R_n\) with
\[
 \nu_n(E_n(R_n))
 \leq\left({a_n\over K_n}\right)^{q'}.
\]
Then \Cref{thm:n9v21-Lq} gives \eqref{eq:n9v21-diagonal}.
\end{proof}

\begin{corollary}[Two-sided adjacent cutoff]
\label{cor:n9v21-two-sided}
If a common bad set \(E_n(R)\) has vanishing reference probability under
both \(\nu_n\) and \(\nu_{n+1}\), and their full density norms
\(K_n,K_{n+1}\) are finite independently of \(R\), then \(R_n\) can be
chosen so that
\begin{equation}
 \mu_n(E_n(R_n))+\mu_{n+1}(E_n(R_n))\leq a_n.         \label{eq:n9v21-two-sided}
\end{equation}
\end{corollary}

\begin{proof}
Choose \(R_n\) large enough that the two reference probabilities are at
most
\[
 (a_n/(2K_n))^{q'},\qquad
 (a_n/(2K_{n+1}))^{q'},
\]
respectively, and apply \Cref{thm:n9v21-Lq}.
\end{proof}

\begin{firewall}[Truncation-dependent density norms can be circular]
\label{fire:n9v21-circular}
If only a truncated density norm \(K_n(R)\) is known, the diagonal theorem
requires the actual product
\[
 K_n(R)\nu_n(E_n(R))^{1/q'}
\]
to tend to zero.  Reference-tail decay alone cannot be used while ignoring
growth of the normalized interacting density.
\end{firewall}

\subsection{Explicit tail schedules}
\label{sec:n9v21-schedules}

\begin{proposition}[Exponential reference tail]
\label{prop:n9v21-exponential}
Suppose
\[
 \nu_n(E_n(R))\leq C_ne^{-\kappa R^\alpha},
 \qquad
 \|r_n\|_q\leq e^{b_n}.
\]
Then \(\mu_n(E_n(R_n))\leq a_n\) whenever
\begin{equation}
 R_n^\alpha
 \geq
 {q'\over\kappa}
 \left[b_n+\log(1/a_n)\right]
 +{1\over\kappa}\log C_n.                             \label{eq:n9v21-exponential}
\end{equation}
\end{proposition}

\begin{proof}
\Cref{thm:n9v21-Lq} gives
\[
 \mu_n(E_n(R))
 \leq
 e^{b_n}C_n^{1/q'}
 e^{-\kappa R^\alpha/q'}.
\]
The displayed lower bound on \(R_n^\alpha\) makes the right side at most
\(a_n\).
\end{proof}

\begin{proposition}[Polynomial reference tail]
\label{prop:n9v21-polynomial}
If
\[
 \nu_n(E_n(R))\leq C_nR^{-\gamma},
 \qquad \|r_n\|_q\leq K_n,
\]
then \(\mu_n(E_n(R_n))\leq a_n\) whenever
\begin{equation}
 R_n\geq
 \left({K_nC_n^{1/q'}\over a_n}\right)^{q'/\gamma}.
 \label{eq:n9v21-polynomial}
\end{equation}
\end{proposition}

\begin{proof}
Insert the polynomial tail into \eqref{eq:n9v21-Lq} and solve for \(R_n\).
\end{proof}

\subsection{Union of physical bad sets}
\label{sec:n9v21-union}

At an adjacent step, let
\begin{equation}
 E_n
 =
 E_n^{\rm Higgs}
 \cup E_n^{\rm coa}
 \cup E_n^{\rm cross}
 \cup E_n^{\rm tr}
 \cup E_n^{\rm slow},                                 \label{eq:n9v21-union}
\end{equation}
represent respectively Higgs-amplitude truncation, coarea
nontransversality, slow spectral crossing, ambient-transport chart exit,
and exceptional slow labels.

\begin{proposition}[One full-density exit ledger]
\label{prop:n9v21-union}
If
\[
 \nu_n(E_n^{a})\leq\delta_n^a
\]
for the five bad-set types, then
\begin{equation}
 \mu_n(E_n)
 \leq
 \|r_n\|_q
 \left(\sum_a\delta_n^a\right)^{1/q'}.                \label{eq:n9v21-union-bound}
\end{equation}
\end{proposition}

\begin{proof}
The union bound gives
\(\nu_n(E_n)\leq\sum_a\delta_n^a\).  Apply
\Cref{thm:n9v21-Lq}.
\end{proof}

\begin{definition}[Normalized chart-exit card, NCXC]
\label{def:n9v21-NCXC}
NCXC records:
\begin{enumerate}[label=\textup{(X\arabic*)},leftmargin=4em]
\item every bad set on the common transported adjacent space;
\item its reference tail under both adjacent reference probabilities;
\item an \(L^q\), entropy, or oscillation bound for the complete normalized
density relative to those references;
\item the density-weighted summability condition
\eqref{eq:n9v21-adjacent};
\item any cutoff dependence of the density norm; and
\item ownership of the resulting exceptional debit among UMCC, TSCC,
SAMAC/SLDAC, NIATC, and CRPC.
\end{enumerate}
\end{definition}

\begin{theorem}[Conditional closure of all declared chart exits]
\label{thm:n9v21-closure}
NCXC implies the summable adjacent exceptional-set condition required by
the regular coarea, slow crossing, action truncation, ambient transport,
and common-refinement theorems.  If their regular-complement estimates and
SFD2 also hold, no declared chart exit remains unpaid in the local
total-variation compiler.
\end{theorem}

\begin{proof}
\Cref{thm:n9v21-adjacent} supplies the common exceptional debit.
\Cref{prop:n9v21-union} permits the physical bad sets to be combined before
applying the density norm.  The cited regular-complement cards then cover
the complement, and the V98/V17 compiler adds the two parts.
\end{proof}

\begin{firewall}[Positivity is required]
\label{fire:n9v21-positive}
The event inequalities use positive probabilities.  A complex chiral
weight or signed gauge-fixed density must first be represented by a
positive measure with the phase retained in observables, or be controlled
by a total-variation measure whose normalization and moments are proved.
\end{firewall}

\begin{assessment}[Progress at ninth-series V21]
\label{ass:n9v21-status}
Every exceptional chart in the current programme now has one normalized
loader.  A full \(L^q\) density turns a reference tail into the physical
tail with exponent \(1-1/q\); SFD2 gives the square-root case.  Entropy and
bounded-oscillation alternatives are explicit, and the diagonal theorem
shows exactly when cutoff selection is noncircular.
\end{assessment}

\begin{programme}[Ninth-series V22: nonlinear metric support]
\label{prog:n9v21-V22}
V22 will attack the next topology gap: preservation of positive-definite
metric support under negative-Sobolev landing.  It will prove what can be
obtained from a uniform ellipticity event in a stronger continuous
topology, give the interpolation/moment condition needed for that event,
and state why distributional tightness alone cannot preserve pointwise
positivity.
\end{programme}

\begin{firewall}[Claim boundary after ninth-series V21]
\label{fire:n9v21-boundary}
V21 does not prove NCXC or SFD2 for the physical density, positivity of the
chiral measure, nonlinear metric support, reconstruction, or the
SM--Einstein continuum.
\end{firewall}

\begin{theorem}[Ninth-series V21 result]
\label{thm:n9v21-result}
For a complete density \(r\in L^q(\nu)\), every chart exit satisfies
\(\mu(E)\leq\|r\|_q\nu(E)^{1-1/q}\).  The adjacent summability theorem,
noncircular diagonal cutoff, and unified bad-set ledger conditionally close
the exceptional parts of UMCC, TSCC, the action cards, NIATC, and CRPC.
\end{theorem}

\begin{proof}
Use \Cref{thm:n9v21-Lq,thm:n9v21-adjacent,
thm:n9v21-diagonal,prop:n9v21-union,
thm:n9v21-closure}.
\end{proof}

\paragraph{Parent-version record.}
V21 was formed from
\texttt{Renewal\_SM\_Einstein\_Continuum\_Research\_Notes\_V20.tex}.
The V20 parent had \(6\,798\) logical source lines, \(237\,412\) bytes, and
SHA--256
\[
 \texttt{C0BDE6158CF43428BFD923E7A8EC356C5904C32B8CD2CEB197D687F591D5A3AB}.
\]
The next compulsory companion audit is V25.

% =============================================================================
% END APPENDED NINTH-SERIES V21 MATERIAL
% =============================================================================


% =============================================================================
% BEGIN APPENDED NINTH-SERIES V22 MATERIAL
% =============================================================================

\clearpage
\section{Ninth-series V22: distributional and classical metric support}
\label{sec:v9v22}

\subsection{The closed distributional positive cone}
\label{sec:n9v22-cone}

Let \(M\) be a smooth manifold with a fixed smooth background metric
\(g_*\).  Pair symmetric covariant two-tensor distributions with smooth
compactly supported contravariant test tensors using the fixed volume form
of \(g_*\).  For \(\lambda\geq0\), define
\begin{equation}
 \mathcal P_\lambda
 :=
 \left\{
 g\in\mathcal D'(S^2T^*M):
 \langle g-\lambda g_*,\varphi X\otimes X\rangle\geq0
 \text{ for all }\varphi\geq0,\ X
 \right\},                                             \label{eq:n9v22-cone}
\end{equation}
where \(\varphi\in C_c^\infty(M)\) and
\(X\in C_c^\infty(TM)\).

\begin{theorem}[Closed distributional ellipticity cone]
\label{thm:n9v22-cone}
The set \(\mathcal P_\lambda\) is convex and closed in the distribution
topology.  If probabilities \(\mu_n\) converge weakly on a Polish
distribution space continuously embedded in
\(\mathcal D'(S^2T^*M)\), and
\begin{equation}
 \mu_n(\mathcal P_\lambda)\longrightarrow1,           \label{eq:n9v22-cone-prob}
\end{equation}
then the weak limit \(\mu\) satisfies
\begin{equation}
 \mu(\mathcal P_\lambda)=1.                           \label{eq:n9v22-cone-limit}
\end{equation}
\end{theorem}

\begin{proof}
Every inequality in \eqref{eq:n9v22-cone} defines a closed affine
half-space because evaluation on a fixed test tensor is continuous in
\(\mathcal D'\).  Their intersection is closed and convex.  Its inverse
image in the chosen Polish distribution space is closed.  The closed-set
Portmanteau inequality gives
\[
 1=\limsup_n\mu_n(\mathcal P_\lambda)
 \leq\mu(\mathcal P_\lambda),
\]
which proves \eqref{eq:n9v22-cone-limit}.
\end{proof}

\begin{corollary}[A countable test core suffices]
\label{cor:n9v22-countable}
On a second-countable manifold, one may choose countable dense families of
nonnegative test functions and compactly supported vector fields whose
inequalities determine \(\mathcal P_\lambda\).
\end{corollary}

\begin{proof}
Choose countable dense subsets in the usual test-function Frechet spaces
on a countable compact exhaustion.  Distributional evaluation is
continuous, and every nonnegative test function can be approximated by
nonnegative members after smoothing and a vanishing positive adjustment.
Polarization and density of vector fields then extend the inequalities to
the full test family.
\end{proof}

\begin{firewall}[A positive tensor distribution need not be a metric]
\label{fire:n9v22-distribution}
Membership in \(\mathcal P_\lambda\) is meaningful distributional
positivity.  It does not imply that \(g\) is a continuous tensor, that
\(\det g\) is a function, or that \(g^{-1}\), the Levi--Civita connection,
and curvature exist.
\end{firewall}

\subsection{Two sharp counterexamples}
\label{sec:n9v22-counterexamples}

\begin{proposition}[Strict metrics can degenerate without a uniform floor]
\label{prop:n9v22-degenerate}
The smooth metrics
\[
 g_n={1\over n}g_*
\]
are positive definite for every \(n\) and converge in every distribution
topology to the zero tensor, which is not a metric.
\end{proposition}

\begin{proof}
Pointwise positive definiteness is immediate.  Pairing with any compactly
supported smooth test tensor gives a factor \(1/n\), hence convergence to
zero.
\end{proof}

\begin{proposition}[A uniform floor can still have a singular limit]
\label{prop:n9v22-singular}
Let \(M\) be compact, let \(x_0\in M\), and let
\(\rho_n\geq0\) be smooth approximate identities converging to
\(\delta_{x_0}\).  Then
\[
 g_n=(1+\rho_n)g_*\geq g_*
\]
are smooth uniformly elliptic metrics, but
\[
 g_n\longrightarrow g_*+\delta_{x_0}g_*
\]
as tensor distributions.  The limit is not a continuous metric.
\end{proposition}

\begin{proof}
The lower bound is pointwise.  For every smooth compactly supported test
tensor \(T\),
\[
 \langle\rho_ng_*,T\rangle
 \longrightarrow\langle\delta_{x_0}g_*,T\rangle
\]
by the defining property of an approximate identity.  The limit contains
a point mass and therefore is not represented by a continuous tensor.
\end{proof}

\begin{corollary}[Negative-Sobolev tightness does not give classical
geometry]
\label{cor:n9v22-negative}
The LDLC landing in a negative-Sobolev space, even together with a uniform
distributional ellipticity floor, does not by itself construct a
classical metric.
\end{corollary}

\begin{proof}
For \(s>d/2\), a point mass lies in \(H^{-s}\).  Thus
\Cref{prop:n9v22-singular} can converge in a negative-Sobolev topology
while its limit remains nonclassical.
\end{proof}

\subsection{Positive-Sobolev compactness}
\label{sec:n9v22-positive-Sobolev}

\begin{theorem}[Moment criterion for \(C^k\) tightness]
\label{thm:n9v22-Ck-tight}
Let \(M\) be compact, \(k\) a nonnegative integer, and
\begin{equation}
 r>k+{d\over2}.                                       \label{eq:n9v22-index}
\end{equation}
If Borel probabilities \(\mu_n\) on \(C^k(S^2T^*M)\) satisfy, for some
\(p>0\),
\begin{equation}
 \sup_n\int\|g\|_{H^r}^p\,d\mu_n(g)\leq C,            \label{eq:n9v22-Hr-moment}
\end{equation}
then \((\mu_n)\) is tight in \(C^k(S^2T^*M)\).
\end{theorem}

\begin{proof}
Choose \(r'\) with
\[
 k+d/2<r'<r.
\]
Rellich compactness gives a compact inclusion \(H^r\hookrightarrow H^{r'}\),
and Sobolev embedding gives a continuous inclusion
\(H^{r'}\hookrightarrow C^k\).  Therefore the closure in \(C^k\) of every
bounded \(H^r\) ball is compact.  Markov's inequality applied to
\eqref{eq:n9v22-Hr-moment} gives compact containment exactly as in
\Cref{thm:n9v11-tight}.
\end{proof}

\begin{corollary}[Density-loaded \(C^k\) tightness]
\label{cor:n9v22-density}
Let \(\mu_n=r_n\nu_n\).  If
\[
 \sup_n\|r_n\|_{L^q(\nu_n)}\leq R
\]
and the reference laws have a uniform \(H^r\) moment of order \(pq'\),
then \eqref{eq:n9v22-Hr-moment} holds and the target laws are tight in
\(C^k\).
\end{corollary}

\begin{proof}
Use the density-to-moment theorem \Cref{thm:n9v12-Holder} with
\(X=\|g\|_{H^r}\), then apply
\Cref{thm:n9v22-Ck-tight}.
\end{proof}

\begin{firewall}[Ultraviolet field measures rarely have positive
regularity]
\label{fire:n9v22-UV}
The positive-Sobolev hypothesis is much stronger than the negative
regularity delivered by a free quantum field.  It is a valid classical
metric route, but V22 does not assume that an ultraviolet gravitational
measure has such moments.  If it does not, generalized geometry or a
different basic variable is required.
\end{firewall}

\subsection{Closed support on uniformly elliptic continuous metrics}
\label{sec:n9v22-classical}

Define
\begin{equation}
 \mathcal M_\lambda^0
 :=
 \{g\in C^0(S^2T^*M):g\geq\lambda g_*\},
 \qquad\lambda>0.                                     \label{eq:n9v22-Mlambda}
\end{equation}

\begin{theorem}[Classical metric support]
\label{thm:n9v22-classical}
The set \(\mathcal M_\lambda^0\) is closed in \(C^0\).  If
\(\mu_n\Rightarrow\mu\) weakly in \(C^0\) and
\begin{equation}
 \mu_n(\mathcal M_\lambda^0)\longrightarrow1,         \label{eq:n9v22-ell-prob}
\end{equation}
then \(\mu(\mathcal M_\lambda^0)=1\).  On this set the inverse map is
continuous and satisfies
\begin{equation}
 \|g^{-1}-h^{-1}\|_{C^0}
 \leq\lambda^{-2}\|g-h\|_{C^0}                       \label{eq:n9v22-inverse}
\end{equation}
in background-orthonormal operator norms.
\end{theorem}

\begin{proof}
If \(g_j\to g\) uniformly and \(g_j\geq\lambda g_*\), passage to the limit
in \(g_j(x)(v,v)\geq\lambda g_*(x)(v,v)\) proves that
\(g\in\mathcal M_\lambda^0\).  Portmanteau and
\eqref{eq:n9v22-ell-prob} give full limiting support.  Finally,
\[
 g^{-1}-h^{-1}=g^{-1}(h-g)h^{-1}.
\]
Both inverse norms are at most \(\lambda^{-1}\), giving
\eqref{eq:n9v22-inverse}.
\end{proof}

\begin{corollary}[Volume density continuity]
\label{cor:n9v22-volume}
On \(\mathcal M_\lambda^0\), the maps
\[
 g\longmapsto g^{-1},\qquad
 g\longmapsto d{\rm vol}_g
\]
are continuous in \(C^0\).  Thus bounded continuous cylinder functionals
of the inverse metric and volume density pass under weak \(C^0\)
convergence.
\end{corollary}

\begin{proof}
The inverse statement is \eqref{eq:n9v22-inverse}.  On each bounded
matrix set with eigenvalues at least \(\lambda\), the square root of the
determinant is continuous uniformly; a uniformly convergent sequence is
bounded, so the volume densities converge uniformly.
\end{proof}

\subsection{A reference-tail loader for ellipticity}
\label{sec:n9v22-tail}

Let \(\bar g\geq\lambda_0g_*\) be a smooth background and write
\[
 g=\bar g+h.
\]
Fix \(0<\lambda<\lambda_0\), and let \(C_{\rm emb}\) be the norm of the
embedding \(H^r\hookrightarrow C^0\), with \(r>d/2\).

\begin{proposition}[Reference ellipticity exit]
\label{prop:n9v22-exit}
For every reference probability \(\nu\),
\begin{equation}
 \nu(g\notin\mathcal M_\lambda^0)
 \leq
 \left({C_{\rm emb}\over\lambda_0-\lambda}\right)^p
 \int\|h\|_{H^r}^p\,d\nu(h).                         \label{eq:n9v22-exit}
\end{equation}
\end{proposition}

\begin{proof}
If
\[
 \|h\|_{H^r}<{\lambda_0-\lambda\over C_{\rm emb}},
\]
then \(\|h\|_{C^0}<\lambda_0-\lambda\), so
\[
 g\geq[\lambda_0-\|h\|_{C^0}]g_*\geq\lambda g_*,
\]
after using background-orthonormal tensor norms.  The bad event is
therefore contained in the complementary Sobolev-norm event.  Apply
Markov's inequality.
\end{proof}

\begin{corollary}[Full-density ellipticity exit]
\label{cor:n9v22-exit}
If \(\mu=r\nu\) and \(r\in L^q(\nu)\), then
\begin{equation}
 \mu(g\notin\mathcal M_\lambda^0)
 \leq
 \|r\|_q
 \left[
 \left({C_{\rm emb}\over\lambda_0-\lambda}\right)^p
 \int\|h\|_{H^r}^p\,d\nu
 \right]^{1/q'}.                                      \label{eq:n9v22-full-exit}
\end{equation}
\end{corollary}

\begin{proof}
Combine \Cref{prop:n9v22-exit} with the normalized exit theorem
\Cref{thm:n9v21-Lq}.
\end{proof}

\begin{firewall}[A bounded moment need not make exit vanish]
\label{fire:n9v22-tail}
A uniform positive-Sobolev moment makes the laws tight but does not force
the right side of \eqref{eq:n9v22-full-exit} to tend to zero.  Classical
metric support needs exact regulated support or a vanishing/summable
ellipticity-exit budget.
\end{firewall}

\subsection{Regularity ladder for Einstein geometry}
\label{sec:n9v22-ladder}

\begin{proposition}[Classical regularity ladder]
\label{prop:n9v22-ladder}
On uniformly elliptic metrics:
\begin{enumerate}[label=\textup{(\roman*)},leftmargin=3.8em]
\item \(C^0\) convergence gives continuous inverse metrics and volume
densities;
\item \(C^1\) convergence gives convergence of Levi--Civita connection
coefficients in \(C^0\); and
\item \(C^2\) convergence gives convergence of Riemann, Ricci, and scalar
curvature in \(C^0\).
\end{enumerate}
\end{proposition}

\begin{proof}
Item (i) is \Cref{thm:n9v22-classical,cor:n9v22-volume}.  In local
coordinates,
\[
 \Gamma^k_{ij}
 ={1\over2}g^{k\ell}
 (\partial_ig_{j\ell}+\partial_jg_{i\ell}-\partial_\ell g_{ij}),
\]
so item (ii) follows from uniform inverse control and \(C^1\) convergence.
Curvature is a sum of first derivatives of \(\Gamma\) and quadratic
\(\Gamma\)-terms, proving item (iii).
\end{proof}

\begin{definition}[Classical metric support card, CMSC]
\label{def:n9v22-CMSC}
CMSC records:
\begin{enumerate}[label=\textup{(M\arabic*)},leftmargin=4em]
\item a common \(C^k\) tensor realization of the metric regulators;
\item a density-loaded \(H^r\) moment with \(r>k+d/2\);
\item a fixed ellipticity floor \(\lambda>0\);
\item exact support on \(\mathcal M_\lambda^0\) or a vanishing normalized
exit bound \eqref{eq:n9v22-full-exit};
\item compatibility of metric chart transitions; and
\item the regularity level \(k=0,1,2\) required by the claimed geometric
observable.
\end{enumerate}
\end{definition}

\begin{theorem}[Conditional classical metric landing]
\label{thm:n9v22-closure}
CMSC yields a Radon limiting probability supported on uniformly elliptic
\(C^k\) metrics.  It supports inverse metric and volume observables for
\(k=0\), Levi--Civita connection observables for \(k\geq1\), and classical
curvature observables for \(k\geq2\).
\end{theorem}

\begin{proof}
\Cref{thm:n9v22-Ck-tight} gives tightness.  Cylinder uniqueness gives the
Radon landing as in V11.  The closed support theorem
\Cref{thm:n9v22-classical} gives the ellipticity floor.  Apply
\Cref{prop:n9v22-ladder}.
\end{proof}

\begin{firewall}[Quantum gravity may require generalized geometry]
\label{fire:n9v22-generalized}
CMSC is a sufficient classical route, not a necessary description of a
quantum gravitational continuum.  If only negative-Sobolev moments are
available, the limit may live in \(\mathcal P_\lambda\) without admitting
classical curvature.  A generalized curvature construction would then be
a separate programme.
\end{firewall}

\begin{assessment}[Progress at ninth-series V22]
\label{ass:n9v22-status}
Metric support is now split correctly.  Negative-Sobolev landing preserves
a closed distributional positivity cone but not a classical inverse.
Classical \(C^k\) geometry follows from density-loaded positive-Sobolev
moments, a uniform ellipticity floor, and vanishing normalized exit.
The regularity needed for inverse, connection, and curvature observables is
explicit.
\end{assessment}

\begin{programme}[Ninth-series V23: weak Einstein residual on the Radon
landing]
\label{prog:n9v22-V23}
V23 will avoid imposing the excessive \(C^2\) route when possible.  It will
formulate the Einstein tensor in weak divergence form for uniformly
elliptic \(H^1\) metrics, identify the quadratic connection term and its
integrability, and prove a stability theorem under strong \(H^1\) plus
uniform \(L^\infty\) ellipticity control.
\end{programme}

\begin{firewall}[Claim boundary after ninth-series V22]
\label{fire:n9v22-boundary}
V22 does not prove CMSC for the physical measures, define curvature for the
negative-Sobolev limit, repair the conformal sign, establish reconstruction,
or construct the SM--Einstein continuum.
\end{firewall}

\begin{theorem}[Ninth-series V22 result]
\label{thm:n9v22-result}
Uniform distributional ellipticity is closed under negative-Sobolev
landing, but it does not produce a classical metric.  Density-loaded
\(H^r\) moments with \(r>k+d/2\), together with a uniform ellipticity
floor and vanishing normalized exit, give tightness and support on
classical \(C^k\) metrics as recorded by CMSC.
\end{theorem}

\begin{proof}
Use \Cref{thm:n9v22-cone,prop:n9v22-singular,
thm:n9v22-Ck-tight,thm:n9v22-classical,
cor:n9v22-exit,thm:n9v22-closure}.
\end{proof}

\paragraph{Parent-version record.}
V22 was formed from
\texttt{Renewal\_SM\_Einstein\_Continuum\_Research\_Notes\_V21.tex}.
The V21 parent had \(7\,231\) logical source lines, \(249\,888\) bytes, and
SHA--256
\[
 \texttt{850568FA26C59AD3973D32FBEAEAB1700E0434EF864A7B6A960E47B584320236}.
\]
The next compulsory companion audit is V25.

% =============================================================================
% END APPENDED NINTH-SERIES V22 MATERIAL
% =============================================================================


% =============================================================================
% BEGIN APPENDED NINTH-SERIES V23 MATERIAL
% =============================================================================

\clearpage
\section{Ninth-series V23: a weak Einstein--Hilbert density at \(H^1\)}
\label{sec:v9v23}

\subsection{Uniformly elliptic \(H^1\) metrics}
\label{sec:n9v23-setting}

Work first in a relatively compact coordinate chart \(U\).  Let \(g\) be a
symmetric matrix field satisfying
\begin{equation}
 g\in H^1(U)\cap L^\infty(U),\qquad
 \lambda I\leq g(x)\leq\Lambda I\quad\text{a.e.}      \label{eq:n9v23-elliptic}
\end{equation}
for fixed \(0<\lambda\leq\Lambda<\infty\).  Then
\(g^{-1},\sqrt{\det g}\in L^\infty\), and the weak Christoffel symbols
\begin{equation}
 \Gamma^k_{ij}(g)
 :={1\over2}g^{k\ell}
 (\partial_ig_{j\ell}+\partial_jg_{i\ell}
 -\partial_\ell g_{ij})                              \label{eq:n9v23-Gamma}
\end{equation}
belong to \(L^2(U)\).

\begin{lemma}[Strong stability under bounded elliptic approximation]
\label{lem:n9v23-Gamma}
Let \(g_n,g\) satisfy the same bounds \eqref{eq:n9v23-elliptic}.  If
\begin{equation}
 g_n\longrightarrow g\quad\text{strongly in }H^1(U),
 \qquad g_n(x)\longrightarrow g(x)\quad\text{a.e.}, \label{eq:n9v23-convergence}
\end{equation}
then
\begin{align}
 g_n^{-1}&\to g^{-1}\quad\text{a.e. and in every finite }L^p,
                                                               \label{eq:n9v23-inverse}\\
 \sqrt{\det g_n}&\to\sqrt{\det g}\quad\text{a.e. and in every finite }L^p,
                                                               \label{eq:n9v23-volume}\\
 \Gamma(g_n)&\to\Gamma(g)\quad\text{in }L^2.          \label{eq:n9v23-Gamma-conv}
\end{align}
\end{lemma}

\begin{proof}
Continuity of matrix inversion and determinant on the compact ellipticity
set gives pointwise convergence, while the common bounds give convergence
in every finite \(L^p\) by dominated convergence.  Write schematically
\[
 \Gamma(g_n)-\Gamma(g)
 =g_n^{-1}\partial(g_n-g)+(g_n^{-1}-g^{-1})\partial g.
\]
The first term tends to zero in \(L^2\).  For the second, the coefficient
tends to zero a.e., is uniformly bounded, and multiplies the fixed
\(L^2\) function \(\partial g\).  Dominated convergence proves
\eqref{eq:n9v23-Gamma-conv}.
\end{proof}

\subsection{The \(\Gamma\Gamma+\partial B\) formula}
\label{sec:n9v23-density}

Define
\begin{align}
 Q(g)
 &:=
 \sqrt{\det g}\,g^{ij}
 \left(
 \Gamma^k_{i\ell}\Gamma^\ell_{jk}
 -\Gamma^k_{ij}\Gamma^\ell_{k\ell}
 \right),                                             \label{eq:n9v23-Q}\\
 B^k(g)
 &:=
 \sqrt{\det g}
 \left(
 g^{ij}\Gamma^k_{ij}
 -g^{ik}\Gamma^j_{ij}
 \right).                                             \label{eq:n9v23-B}
\end{align}
For \eqref{eq:n9v23-elliptic}, one has
\[
 Q(g)\in L^1(U),\qquad B(g)\in L^2(U).
\]

\begin{definition}[Weak scalar-curvature density]
\label{def:n9v23-R}
For \(\varphi\in C_c^\infty(U)\), define
\begin{equation}
 \langle\mathcal R(g),\varphi\rangle
 :=
 \int_UQ(g)\varphi\,dx
 -\int_UB^k(g)\partial_k\varphi\,dx.                  \label{eq:n9v23-R}
\end{equation}
\end{definition}

\begin{proposition}[Agreement with smooth scalar curvature]
\label{prop:n9v23-smooth}
If \(g\) is \(C^2\), then
\begin{equation}
 \mathcal R(g)=\sqrt{\det g}\,R(g)                    \label{eq:n9v23-smooth}
\end{equation}
as distributions.
\end{proposition}

\begin{proof}
The coordinate identity for the Einstein--Hilbert density is
\[
 \sqrt{\det g}\,R(g)=Q(g)+\partial_kB^k(g).
\]
Pairing with \(\varphi\) and integrating the divergence by parts gives
\eqref{eq:n9v23-R}.
\end{proof}

\begin{theorem}[Strong \(H^1\) stability of the weak curvature density]
\label{thm:n9v23-stability}
Under the hypotheses of \Cref{lem:n9v23-Gamma},
\begin{equation}
 Q(g_n)\to Q(g)\text{ in }L^1,\qquad
 B(g_n)\to B(g)\text{ in }L^2,                        \label{eq:n9v23-QB}
\end{equation}
and consequently
\begin{equation}
 \mathcal R(g_n)\longrightarrow\mathcal R(g)
 \quad\text{in }\mathcal D'(U).                       \label{eq:n9v23-R-conv}
\end{equation}
\end{theorem}

\begin{proof}
All coefficient factors are uniformly bounded and converge a.e., while
\(\Gamma(g_n)\to\Gamma(g)\) in \(L^2\).  The product estimate
\[
 \|ab-a'b'\|_1
 \leq\|a-a'\|_2\|b\|_2+\|a'\|_2\|b-b'\|_2
\]
first handles the change of Christoffel products.  The remaining fixed
quadratic product is multiplied by a bounded coefficient converging a.e.;
dominated convergence handles it.  This proves the first part of
\eqref{eq:n9v23-QB}.  The same decomposition, with a fixed \(L^2\)
Christoffel factor and dominated convergence for its squared norm, gives
the \(L^2\) convergence of \(B\).  Insert both limits in
\eqref{eq:n9v23-R}.
\end{proof}

\begin{corollary}[Coordinate-independent weak extension]
\label{cor:n9v23-coordinate}
On a smooth manifold, the local definitions \eqref{eq:n9v23-R} agree on
overlaps and define a global scalar-density distribution for uniformly
elliptic \(H^1\cap L^\infty\) metrics.
\end{corollary}

\begin{proof}
After shrinking a compact subchart, extend and mollify the coefficient
matrix.  Convexity of
\(\{\lambda I\leq A\leq\Lambda I\}\) preserves the same bounds away from
the artificial boundary, and standard mollification gives strong
\(H^1\) and a.e. convergence along a subsequence.  For smooth metrics the
expression is the invariant density \(\sqrt{\det g}R(g)\).
\Cref{thm:n9v23-stability} gives the same distributional limit in every
smooth coordinate chart.  Exhausting the chart by compact subcharts shows
that the local limits agree on overlaps.
\end{proof}

\begin{firewall}[This is the scalar density, not the full Einstein tensor]
\label{fire:n9v23-tensor}
V23 defines the Einstein--Hilbert scalar density at \(H^1\) regularity.
The tensor-valued first variation, Bianchi identity, boundary terms, and
coupling to a rough stress tensor require additional differentiability and
integrability.  They are not consequences of
\eqref{eq:n9v23-R-conv}.
\end{firewall}

\subsection{Quantitative bounds and expectations}
\label{sec:n9v23-probability}

\begin{lemma}[Tested curvature growth]
\label{lem:n9v23-growth}
For every \(\varphi\in C_c^\infty(U)\), metrics satisfying
\eqref{eq:n9v23-elliptic} obey
\begin{equation}
 |\langle\mathcal R(g),\varphi\rangle|
 \leq
 C_{\lambda,\Lambda,\varphi}
 \left(
 1+\|\nabla g\|_{L^2(U)}^2
 \right).                                             \label{eq:n9v23-growth}
\end{equation}
\end{lemma}

\begin{proof}
The ellipticity bounds control every undifferentiated coefficient in
\eqref{eq:n9v23-Q}--\eqref{eq:n9v23-B}.  Equation
\eqref{eq:n9v23-Gamma} gives
\(\|\Gamma\|_2\leq C_\lambda\|\nabla g\|_2\).  Therefore
\[
 \|Q\|_1\leq C\|\nabla g\|_2^2,\qquad
 \|B\|_2\leq C\|\nabla g\|_2.
\]
Use Cauchy--Schwarz in the \(B\)-term and absorb the linear norm by
\(x\leq1+x^2\).
\end{proof}

\begin{theorem}[Passage of expected weak curvature]
\label{thm:n9v23-expectation}
Let \(\mu_n\Rightarrow\mu\) on a Polish space of uniformly bounded,
uniformly elliptic metrics whose topology implies strong local \(H^1\)
convergence, and suppose all measures are supported on the common
ellipticity set \eqref{eq:n9v23-elliptic}.  If, for
some \(\varepsilon>0\),
\begin{equation}
 \sup_n\int
 \|\nabla g\|_2^{\,2+\varepsilon}\,d\mu_n(g)<\infty,  \label{eq:n9v23-UI}
\end{equation}
then for every test function \(\varphi\),
\begin{equation}
 \int\langle\mathcal R(g),\varphi\rangle\,d\mu_n(g)
 \longrightarrow
 \int\langle\mathcal R(g),\varphi\rangle\,d\mu(g).    \label{eq:n9v23-expectation}
\end{equation}
\end{theorem}

\begin{proof}
Strong \(H^1\) convergence has an a.e.-convergent subsequence.  Applying
\Cref{thm:n9v23-stability} to every such subsequence, and using the
subsequence principle, makes the tested curvature functional continuous
in the declared topology.  The growth bound
\eqref{eq:n9v23-growth} and \eqref{eq:n9v23-UI} make these random variables
uniformly integrable.  Truncate the continuous functional at height \(R\),
pass to the weak limit for the bounded truncation, and then let \(R\to
\infty\) uniformly using uniform integrability.
\end{proof}

\begin{firewall}[Weak convergence alone does not pass curvature means]
\label{fire:n9v23-UI}
The tested curvature is quadratic in first derivatives and unbounded.
Weak convergence of metric laws, even in strong \(H^1\) topology, does not
pass its expectation without uniform integrability such as
\eqref{eq:n9v23-UI}.
\end{firewall}

\subsection{Weak scalar Einstein--Hilbert card}
\label{sec:n9v23-card}

\begin{definition}[Weak scalar-curvature landing card, WSCLC]
\label{def:n9v23-WSCLC}
WSCLC records:
\begin{enumerate}[label=\textup{(W\arabic*)},leftmargin=4em]
\item tightness in a topology that implies strong local \(H^1\)
convergence on the bounded ellipticity sets;
\item common lower and upper ellipticity bounds, or summable exits from
them;
\item strong \(H^1\) chart-transition compatibility with common
coefficient bounds;
\item the moment \eqref{eq:n9v23-UI} under the full constrained density;
\item boundary conditions needed to assemble the local divergence terms;
and
\item separate matter-stress and counterterm residual cards.
\end{enumerate}
\end{definition}

\begin{theorem}[Conditional weak scalar-curvature landing]
\label{thm:n9v23-closure}
WSCLC gives a well-defined limiting scalar-curvature density and convergence
of all its tested expectations.  It supplies the geometric scalar part of
the weak Einstein residual, but not the tensor equation.
\end{theorem}

\begin{proof}
Use \Cref{cor:n9v23-coordinate,thm:n9v23-expectation}; the final limitation
is \Cref{fire:n9v23-tensor}.
\end{proof}

\begin{assessment}[Progress at ninth-series V23]
\label{ass:n9v23-status}
The classical \(C^2\) requirement has been lowered for one important
observable.  Uniformly elliptic \(H^1\cap L^\infty\) metrics carry a stable
Einstein--Hilbert scalar density in divergence form, with stability under
strong \(H^1\) approximation inside a common ellipticity box.  Its expected weak
limit requires exactly a \(2+\varepsilon\) first-derivative moment.  The
tensor Einstein equation and rough matter coupling remain open.
\end{assessment}

\begin{programme}[Ninth-series V24: tensor first variation]
\label{prog:n9v23-V24}
V24 will differentiate the weak Einstein--Hilbert functional along smooth
compactly supported metric variations, isolate the bulk tensor and
boundary terms, and determine the additional regularity required for a
closable weak Einstein residual.
\end{programme}

\begin{firewall}[Claim boundary after ninth-series V23]
\label{fire:n9v23-boundary}
V23 does not prove WSCLC for the physical measure, construct the full
Einstein tensor or stress coupling, repair the conformal sign,
reconstruction, or the SM--Einstein continuum.
\end{firewall}

\begin{theorem}[Ninth-series V23 result]
\label{thm:n9v23-result}
Uniformly elliptic \(H^1\cap L^\infty\) metrics possess the stable weak
scalar-curvature density \eqref{eq:n9v23-R}.  Strong convergence in that
topology gives distributional curvature convergence, and the full-density
moment \eqref{eq:n9v23-UI} passes tested expectations to the Radon limit.
\end{theorem}

\begin{proof}
Use \Cref{lem:n9v23-Gamma,prop:n9v23-smooth,
thm:n9v23-stability,lem:n9v23-growth,
thm:n9v23-expectation,thm:n9v23-closure}.
\end{proof}

\paragraph{Parent-version record.}
V23 was formed from
\texttt{Renewal\_SM\_Einstein\_Continuum\_Research\_Notes\_V22.tex}.
The V22 parent had \(7\,671\) logical source lines, \(265\,160\) bytes, and
SHA--256
\[
 \texttt{4097241C1D1F354B043E57CC7C6A653AA3448274A0F6CEC75CA3570444BAE948}.
\]
The next compulsory companion audit is V25.

% =============================================================================
% END APPENDED NINTH-SERIES V23 MATERIAL
% =============================================================================


% =============================================================================
% BEGIN APPENDED NINTH-SERIES V24 MATERIAL
% =============================================================================

\clearpage
\section{Ninth-series V24: the tensor Einstein variation at \(H^2\)}
\label{sec:v9v24}

\subsection{Why one more derivative is imposed}
\label{sec:n9v24-setting}

V23 constructed the scalar-curvature density at uniformly elliptic
\(H^1\) regularity.  A tensor equation requires the Ricci components
themselves and their contraction against arbitrary symmetric variations.
For this purpose fix a relatively compact chart \(U\) and assume
\begin{equation}
 g\in H^2(U)\cap W^{1,\infty}(U),\qquad
 \lambda I\leq g\leq\Lambda I,qquad
 \|\nabla g\|_\infty\leq K.                         \label{eq:n9v24-class}
\end{equation}
The constants \(\lambda,\Lambda,K\) are common throughout each limiting
family considered below.

\begin{lemma}[\(H^1\) connection and \(L^2\) curvature]
\label{lem:n9v24-regularity}
For \eqref{eq:n9v24-class},
\begin{equation}
 \Gamma(g)\in H^1(U)\cap L^\infty(U),\qquad
 \Ric(g),R(g),G(g)\in L^2(U),                        \label{eq:n9v24-reg}
\end{equation}
where \(G_{ij}:=\Ric_{ij}-\frac12Rg_{ij}\).  Moreover
\begin{equation}
 \|\Gamma\|_\infty+\|\nabla\Gamma\|_2
 +\|\Ric\|_2+\|R\|_2+\|G\|_2
 \leq C_{\lambda,\Lambda,K,U}
 \bigl(1+\|g\|_{H^2(U)}\bigr).                      \label{eq:n9v24-growth}
\end{equation}
\end{lemma}

\begin{proof}
The formula \(\Gamma=g^{-1}*\partial g\) gives the \(L^\infty\) bound.
Differentiation gives
\[
 \partial\Gamma
 =g^{-1}*\partial^2g+g^{-1}*g^{-1}*\partial g*\partial g.
\]
The first term is in \(L^2\); the second is in \(L^2\) because one first
derivative is bounded and the other is square integrable.  In coordinates,
\begin{equation}
 \Ric_{ij}
 =\partial_k\Gamma^k_{ij}-\partial_j\Gamma^k_{ik}
 +\Gamma^k_{k\ell}\Gamma^\ell_{ij}
 -\Gamma^k_{j\ell}\Gamma^\ell_{ik}.                \label{eq:n9v24-Ricci}
\end{equation}
Thus \(\Ric\in L^2\), and bounded contraction coefficients give the claims
for \(R\) and \(G\).  The displayed estimates give
\eqref{eq:n9v24-growth}.
\end{proof}

\subsection{The covariant first variation}
\label{sec:n9v24-variation}

For a smooth compactly supported symmetric covariant tensor \(h\), put
\(g_t=g+th\).  For small \(|t|\), this remains uniformly elliptic.  If
\(\chi\in C_c^\infty(U)\) equals one on a neighbourhood of
\(\supp h\), define the localized action difference
\begin{equation}
 \Delta_h\mathcal S_g(t)
 :=\langle\mathcal R(g_t)-\mathcal R(g),\chi\rangle. \label{eq:n9v24-action}
\end{equation}
Because \(g_t=g\) off \(\supp h\), the right side is independent of the
choice of such \(\chi\).

\begin{theorem}[First variation of the weak Einstein--Hilbert density]
\label{thm:n9v24-firstvariation}
For \(g\) satisfying \eqref{eq:n9v24-class} and
\(h\in C_c^\infty(U;\Sym^2T^*U)\), the map
\(t\mapsto\Delta_h\mathcal S_g(t)\) is differentiable at zero and
\begin{equation}
 {d\over dt}\bigg|_{t=0}\Delta_h\mathcal S_g(t)
 =-\int_U\sqrt{\det g}\,G^{ij}(g)h_{ij}\,dx.        \label{eq:n9v24-firstvariation}
\end{equation}
Here indices on \(G\) are raised with \(g^{-1}\).  Thus the sign in
\eqref{eq:n9v24-firstvariation} corresponds to varying the covariant
metric \(g_{ij}\).
\end{theorem}

\begin{proof}
For a smooth metric the pointwise variation formula is
\begin{align}
 {d\over dt}\bigg|_{0}\bigl(\sqrt{\det g_t}R(g_t)\bigr)
 &=-\sqrt{\det g}\,G^{ij}h_{ij}+\partial_k\Theta^k(g,h),
                                                               \label{eq:n9v24-pointwise}\\
 \Theta^k(g,h)
 &:=\sqrt{\det g}\left(\nabla_i h^{ik}
              -\nabla^k\tr_g h\right).              \label{eq:n9v24-Theta}
\end{align}
The divergence integrates to zero because \(h\) is compactly supported.

For the stated rough metric, extend it slightly past each compact subset
and mollify.  Convexity of the ellipticity box preserves the common
matrix bounds; mollification preserves the gradient bound and converges
strongly in \(H^2\) after restriction.  The calculation in
\Cref{lem:n9v24-regularity} and dominated convergence show that the
Christoffels converge in \(H^1\), while
\(\Ric,R,G\) converge in \(L^2\).  Apply the smooth formula to the
approximants.  The bulk term converges by \(L^2\)-convergence.  In the
localized formula the variation of \(\partial_k\Theta^k\) pairs with
\(\chi\), whose derivative vanishes on the support of \(\Theta(g,h)\).
The V23 stability theorem passes the left side.  Equivalently, direct
difference quotients in the polynomial expressions
\eqref{eq:n9v23-Q}--\eqref{eq:n9v23-B} are dominated in \(L^1+H^{-1}\)
by \eqref{eq:n9v24-class}.  This proves differentiability and
\eqref{eq:n9v24-firstvariation}.
\end{proof}

\begin{corollary}[Cosmological first variation]
\label{cor:n9v24-Lambda}
For
\(
 \mathcal S_\Lambda(g)=\int(R(g)-2\Lambda_c)\,d\vol_g
\)
in the same localized sense,
\begin{equation}
 D\mathcal S_\Lambda(g)[h]
 =-\int_U\sqrt{\det g}\,
 \bigl(G^{ij}+\Lambda_c g^{ij}\bigr)h_{ij}\,dx.     \label{eq:n9v24-Lambda}
\end{equation}
\end{corollary}

\begin{proof}
Use \eqref{eq:n9v24-firstvariation} and
\(
 D\sqrt{\det g}[h]=\frac12\sqrt{\det g}\,g^{ij}h_{ij}
\).
\end{proof}

\subsection{Strong closure and weak Bianchi identity}
\label{sec:n9v24-closure}

Define the contravariant Einstein density
\begin{equation}
 \mathfrak G^{ij}(g):=\sqrt{\det g}\,G^{ij}(g).      \label{eq:n9v24-density}
\end{equation}

\begin{theorem}[Strong \(H^2\) closure inside the bounded class]
\label{thm:n9v24-stability}
Let \(g_n,g\) satisfy the same bounds \eqref{eq:n9v24-class}, assume
\(g_n\to g\) strongly in \(H^2(U)\), and, after passage to a harmless
subsequence, assume pointwise convergence of \(g_n\) and \(\nabla g_n\).
Then
\begin{equation}
 \Gamma(g_n)\to\Gamma(g)\text{ in }H^1,qquad
 \mathfrak G(g_n)\to\mathfrak G(g)\text{ in }L^2.  \label{eq:n9v24-stability}
\end{equation}
Consequently the same conclusion holds for the full sequence without
writing the pointwise subsequence assumption.
\end{theorem}

\begin{proof}
The zeroth-order connection difference is handled as in V23.  On
differentiating \(\Gamma=g^{-1}*\partial g\), the second-derivative
terms converge in \(L^2\).  A typical quadratic difference satisfies
\[
 \|\partial g_n*\partial g_n-\partial g*\partial g\|_2
 \leq 2K\|\partial g_n-\partial g\|_2.
\]
Coefficient differences multiplying the fixed \(L^2\) function
\(\partial^2g\) converge by dominated convergence.  Hence
\(\Gamma_n\to\Gamma\) in \(H^1\).  Formula
\eqref{eq:n9v24-Ricci} gives \(L^2\)-convergence of Ricci; bounded
pointwise-convergent inverse and volume coefficients give the assertion
for \(\mathfrak G\).  Every subsequence of a strongly \(H^2\)-convergent
sequence has a further pointwise-convergent subsequence, so the
subsequence principle removes the auxiliary assumption.
\end{proof}

\begin{theorem}[Contracted Bianchi identity at the closure regularity]
\label{thm:n9v24-Bianchi}
Every metric satisfying \eqref{eq:n9v24-class} obeys
\begin{equation}
 \partial_i\mathfrak G^{ij}
 +\Gamma^j_{ik}\mathfrak G^{ik}=0
 \quad\text{in }\mathcal D'(U).                     \label{eq:n9v24-Bianchi}
\end{equation}
\end{theorem}

\begin{proof}
For smooth metrics this is
\(\nabla_iG^{ij}=0\) written in density form.  Use the mollified metrics
from the proof of \Cref{thm:n9v24-firstvariation}.  The first term passes
because \(\mathfrak G_n\to\mathfrak G\) in \(L^2\).  The product in the
second term passes in \(L^1\), since both \(\Gamma_n\) and
\(\mathfrak G_n\) converge strongly in \(L^2\).
\end{proof}

\begin{firewall}[Bianchi is an identity, not a solved Einstein equation]
\label{fire:n9v24-Bianchi}
Equation \eqref{eq:n9v24-Bianchi} is geometric compatibility of the
limiting tensor.  It neither constructs a stress tensor nor proves its
covariant conservation.  A proposed matter limit must satisfy the same
compatibility after counterterms and anomalies are included.
\end{firewall}

\subsection{Expected tensor residuals}
\label{sec:n9v24-expectation}

Let \(\mathfrak T^{ij}\) be a contravariant matter stress density, already
including the volume factor.  For a test tensor \(h\), define
\begin{equation}
 \mathscr E_h(g,\mathfrak T)
 :=\int_U
 \left[
 \mathfrak G^{ij}(g)+\Lambda_c\sqrt{\det g}\,g^{ij}
 -\kappa\mathfrak T^{ij}
 \right]h_{ij}\,dx.                                 \label{eq:n9v24-residual}
\end{equation}

\begin{theorem}[Passage of expected Einstein residuals]
\label{thm:n9v24-expectation}
Suppose laws \(\mu_n\Rightarrow\mu\) are supported on a common class
\eqref{eq:n9v24-class}, their topology implies strong local \(H^2\)
convergence of \(g\) and strong local \(L^2\) convergence of
\(\mathfrak T\), and for some \(\varepsilon>0\),
\begin{equation}
 \sup_n\int
 \left(\|g\|_{H^2(U)}+\|\mathfrak T\|_{L^2(U)}\right)^{1+\varepsilon}
 d\mu_n<\infty.                                     \label{eq:n9v24-moment}
\end{equation}
Then
\begin{equation}
 \int\mathscr E_h\,d\mu_n\longrightarrow
 \int\mathscr E_h\,d\mu                            \label{eq:n9v24-expected}
\end{equation}
for every compactly supported smooth symmetric \(h\).
\end{theorem}

\begin{proof}
\Cref{thm:n9v24-stability} gives continuity of the geometric term; the
assumed topology gives continuity of the matter term.  By
\eqref{eq:n9v24-growth} and Cauchy--Schwarz,
\[
 |\mathscr E_h|
 \leq C_h\left(1+\|g\|_{H^2}+\|\mathfrak T\|_2\right).
\]
The moment \eqref{eq:n9v24-moment} makes these variables uniformly
integrable.  Bounded truncation followed by removal of the truncation
proves \eqref{eq:n9v24-expected}.
\end{proof}

\begin{definition}[Tensor Einstein landing card, TELC]
\label{def:n9v24-TELC}
TELC records:
\begin{enumerate}[label=\textup{(T\arabic*)},leftmargin=4em]
\item a Radon limit in a topology implying strong local \(H^2\) metric
and strong local \(L^2\) stress-density convergence;
\item common local ellipticity and \(W^{1,\infty}\) bounds, or summable
exit probabilities from a diagonal sequence of such bounds;
\item the uniform-integrability moment \eqref{eq:n9v24-moment};
\item a constructed renormalized stress density \(\mathfrak T\);
\item convergence of cosmological and gravitational counterterms; and
\item a vanishing tested residual \eqref{eq:n9v24-residual} on a dense
countable family of test tensors.
\end{enumerate}
\end{definition}

\begin{theorem}[Conditional tensor Einstein landing]
\label{thm:n9v24-landing}
If TELC holds and the finite-cutoff expected residuals vanish on its dense
test family, then the limiting law satisfies
\begin{equation}
 \int\mathscr E_h\,d\mu=0                            \label{eq:n9v24-limit-equation}
\end{equation}
for every compactly supported smooth symmetric test tensor \(h\).
\end{theorem}

\begin{proof}
Pass each member of the countable family by
\Cref{thm:n9v24-expectation}.  The growth estimate makes
\(h\mapsto\int\mathscr E_h,d\mu\) continuous in the test-function
topology.  Density extends the equality to all \(h\).
\end{proof}

\begin{firewall}[The new hypotheses are obligations]
\label{fire:n9v24-obligations}
V24 identifies a sufficient tensor closure class; it does not prove that
the physical SM--Einstein measures have uniform \(W^{1,\infty}\) bounds,
strong \(H^2\) tightness, or a strongly convergent renormalized stress
density.  The scalar V23 construction survives below this tensor class,
but the present proof does not.
\end{firewall}

\begin{assessment}[Progress at ninth-series V24]
\label{ass:n9v24-status}
The geometric side of a tested Einstein equation is now closable in the
explicit class \eqref{eq:n9v24-class}.  The first variation, tensor
stability, and weak contracted Bianchi identity are proved there.  The
main unresolved analytic load has moved to physical-measure estimates and
construction of the renormalized matter stress density.
\end{assessment}

\begin{programme}[Ninth-series V25: compulsory companion audit]
\label{prog:n9v24-V25}
V25 will inspect the newest Yang--Mills and Navier--Stokes notes before
choosing between a stress-density construction, an exit-probability
loader for TELC, and an upstream repair of the regularity demand.
\end{programme}

\begin{firewall}[Claim boundary after ninth-series V24]
\label{fire:n9v24-boundary}
V24 does not establish TELC for the physical measure, construct the
renormalized Standard Model stress tensor, prove anomaly cancellation or
reflection positivity, repair the conformal sign, perform reconstruction,
or complete the SM--Einstein continuum.
\end{firewall}

\begin{theorem}[Ninth-series V24 result]
\label{thm:n9v24-result}
Within common uniformly elliptic \(H^2\cap W^{1,\infty}\) bounds, the weak
Einstein--Hilbert action has the covariant first variation
\eqref{eq:n9v24-firstvariation}; its densitized Einstein tensor is strongly
\(L^2\)-stable and satisfies the distributional contracted Bianchi
identity.  The moment \eqref{eq:n9v24-moment} passes expected tested
Einstein residuals to a Radon limit.
\end{theorem}

\begin{proof}
Use \Cref{lem:n9v24-regularity,thm:n9v24-firstvariation,
thm:n9v24-stability,thm:n9v24-Bianchi,
thm:n9v24-expectation,thm:n9v24-landing}.
\end{proof}

\paragraph{Parent-version record.}
V24 was formed from
\texttt{Renewal\_SM\_Einstein\_Continuum\_Research\_Notes\_V23.tex}.
The V23 parent had \(7\,995\) logical source lines, \(277\,384\) bytes, and
SHA--256
\[
 \texttt{78BEC9A76122752870F98BFCB0D1CC0098BA9C0E1CB990E14706049D9840892B}.
\]
The next version is the compulsory V25 companion audit.

% =============================================================================
% END APPENDED NINTH-SERIES V24 MATERIAL
% =============================================================================


% =============================================================================
% BEGIN APPENDED NINTH-SERIES V25 MATERIAL
% =============================================================================

\clearpage
\section{Ninth-series V25: companion audit and local Einstein exits}
\label{sec:v9v25}

\subsection{Files inspected and audit boundary}
\label{sec:n9v25-files}

The compulsory five-version audit inspected the newest active companion
notes read-only:
\begin{center}
\small
\begin{tabular}{@{}p{0.46\textwidth}r r@{}}
\toprule
file & source lines & bytes\\
\midrule
\texttt{Renewal\_Yang\_Mills\_Mass\_Gap\_Research\_Notes\_V66.tex}
 & \(28\,877\) & \(1\,031\,339\)\\
\texttt{Renewal\_NS\_Stability\_Research\_Notes\_V26.tex}
 & \(5\,319\) & \(278\,283\)\\
\bottomrule
\end{tabular}
\end{center}
Their SHA--256 fingerprints are respectively
\begin{align}
 &\texttt{8DE9AF7186FFE4B293BD024C491D1EF72C5344445A90E644C03E215AA9C94DE9},
                                                               \label{eq:n9v25-YMhash}\\
 &\texttt{82C887F8C2C69E4F28FAD7C3DC8DE5B9099CB5A4BF431EBA1FFF3E91935978AD}.
                                                               \label{eq:n9v25-NShash}
\end{align}
Only proved statements and explicit firewalls are compared below; their
internal programme instructions are not instructions for this manuscript.

\subsection{Yang--Mills V63--V66: the transferable mechanism}
\label{sec:n9v25-YM}

The YM notes first prove that an all-block single-chart event has
probability tending to zero at positive block-exit density.  They replace
that invalid global target by local marginal density transfer, a replica
pressure defect, compact Morse--Bott one-block tails, uniform conditional
block barriers, finite-range coloring, and a mixed good/bad polymer
expansion.  At V66 the last physical loaders remain the uniform
conditional transverse barrier and a mixed boundary product majorant.

The exact transfer to the Einstein problem is structural: the global
event that every metric block obeys the V24 bounds should not be demanded.
One needs uniform conditional exit estimates and local bad-polymer
summability.  Wilson compactness, gauge slices, oscillator activities,
and the YM replica covariance do not transfer to a noncompact metric
field.

\subsection{Navier--Stokes V26: a precise non-transfer}
\label{sec:n9v25-NS}

NS V26 computes closed-form pointwise rank-one norms for first and second
derivatives of the Oseen kernel.  This sharpens a deterministic singular
integral constant in that programme.  The note contains no compact
configuration probability, conditional Gibbs barrier, metric
ellipticity estimate, or stress-density construction.

\begin{proposition}[Companion-audit decision]
\label{prop:n9v25-decision}
The YM finite-color conditional-exit mechanism transfers to the local
regularity events required by TELC.  No theorem specific to the Wilson
action or the NS Oseen kernel is imported.  The next SM task is therefore
a local conditional Einstein-regularity loader, followed by a construction
of the renormalized stress density.
\end{proposition}

\begin{proof}
The V24 tensor and residual are local distributions tested against compactly
supported tensors.  Hence their closure only uses finitely many metric and
matter blocks at a time.  Conditional exits and coloring control precisely
such finite families without demanding an infinite-volume all-good event.
The rejected YM and NS ingredients depend on field-specific algebra absent
from the hypotheses of V24.
\end{proof}

\subsection{Finite-color domination for local regularity exits}
\label{sec:n9v25-color}

Let \({\cal B}\) be a block graph.  Associate to each block \(b\) a finite
buffer \(b^+\).  The conflict graph joins \(b,c\) whenever their interior
variables cannot be conditionally separated after the exterior is fixed;
assume its maximum degree is \(\Delta_+\).  It has a proper coloring
\begin{equation}
 {cal B}={\cal B}_1\sqcup\cdots\sqcup{\cal B}_{K_c},
 \qquad K_c\leq\Delta_++1.                           \label{eq:n9v25-color}
\end{equation}

For thresholds \(A=(\lambda_A,\Lambda_A,K_A,L_A)\), let \(E_b(A)\)
be the event that on \(b^+\) at least one of
\begin{equation}
 \lambda_A I\leq g\leq\Lambda_A I,qquad
 \|\nabla g\|_\infty\leq K_A,qquad
 \|g\|_{H^2}+\|\mathfrak T\|_2\leq L_A             \label{eq:n9v25-good}
\end{equation}
fails.

\begin{definition}[Local conditional Einstein regularity card, LCERC]
\label{def:n9v25-LCERC}
LCERC at threshold \(A\) and probability \(p_A\in(0,1)\) consists of:
\begin{enumerate}[label=\textup{(L\arabic*)},leftmargin=4em]
\item a finite-range block specification with coloring
\eqref{eq:n9v25-color};
\item for every color, conditional factorization of disjoint block
interiors after all complementary variables are fixed; and
\item the uniform one-block estimate
\begin{equation}
 \mu\bigl(E_b(A)\mid{\cal F}_{b^c}\bigr)\leq p_A
 \quad\text{a.s. for every }b.                       \label{eq:n9v25-oneblock}
\end{equation}
\end{enumerate}
\end{definition}

\begin{lemma}[Same-color multiplication]
\label{lem:n9v25-samecolor}
Under LCERC, for finite \(S\subset{cal B}_j\),
\begin{equation}
 \mu\left(\bigcap_{b\in S}E_b(A)\right)
 \leq p_A^{|S|}.                                     \label{eq:n9v25-samecolor}
\end{equation}
\end{lemma}

\begin{proof}
Condition on the variables outside the mutually separated interiors.
Conditional factorization turns the conditional probability of the
intersection into the product of its one-block conditional probabilities.
Each factor is at most \(p_A\) by \eqref{eq:n9v25-oneblock}; integrate the
result over the exterior variables.
\end{proof}

\begin{theorem}[Arbitrary-set Einstein exit domination]
\label{thm:n9v25-domination}
For every finite \(S\subset{cal B}\), LCERC implies
\begin{equation}
 \mu\left(\bigcap_{b\in S}E_b(A)\right)
 \leq p_A^{|S|/K_c}.                                 \label{eq:n9v25-domination}
\end{equation}
\end{theorem}

\begin{proof}
Partition \(S=\bigsqcup_jS_j\).  Some \(S_{j_*}\) has cardinality at
least \(|S|/K_c\).  The full intersection is contained in the
\(S_{j_*}\) intersection, so \Cref{lem:n9v25-samecolor} gives
\[
 \mu(\cap_{b\in S}E_b)
 \leq p_A^{|S_{j_*}|}
 \leq p_A^{|S|/K_c}.
\]
\end{proof}

\begin{corollary}[Bad-component summability]
\label{cor:n9v25-polymer}
Suppose the block graph has maximum degree \(\Delta\), and let
\(a_m(b)\) be the number of connected \(m\)-block sets containing \(b\).
Using the coarse bound \(a_m(b)\leq\Delta^{2(m-1)}\), for every
\(\alpha\geq0\) such that
\begin{equation}
 \rho_A:=\Delta^2e^\alpha p_A^{1/K_c}<1,             \label{eq:n9v25-rho}
\end{equation}
one has
\begin{equation}
 \sum_{\substack{\Gamma\ni b\\\Gamma\ {
m connected}}}
 e^{\alpha|\Gamma|}
 \mu\left(\bigcap_{c\in\Gamma}E_c(A)\right)
 \leq {e^\alpha p_A^{1/K_c}\over1-\rho_A}.          \label{eq:n9v25-polymer}
\end{equation}
\end{corollary}

\begin{proof}
Apply \Cref{thm:n9v25-domination} to each \(\Gamma\), group by
\(m=|\Gamma|\), and sum the resulting geometric series.
\end{proof}

\begin{firewall}[Local domination is conditional on a physical loader]
\label{fire:n9v25-loader}
LCERC is not proved by coloring.  Its hard input is the uniform
conditional estimate \eqref{eq:n9v25-oneblock}, including ellipticity,
\(W^{1,\infty}\), \(H^2\), and stress-density exits under the complete
normalized SM--Einstein law.  The raw Euclidean Einstein--Hilbert action
does not supply that estimate merely from the formal first variation.
\end{firewall}

\subsection{Local closure without an all-good universe}
\label{sec:n9v25-localclosure}

Let \(h\) be supported in a fixed finite union of blocks \(S_h\), enlarged
once to contain every buffer entering the residual
\eqref{eq:n9v24-residual}.  Put
\begin{equation}
 F_A(h):=\bigcap_{b\in S_h}E_b(A)^c.                 \label{eq:n9v25-localgood}
\end{equation}

\begin{lemma}[Local bad-event probability]
\label{lem:n9v25-union}
Under LCERC,
\begin{equation}
 \mu(F_A(h)^c)\leq |S_h|p_A.                         \label{eq:n9v25-union}
\end{equation}
\end{lemma}

\begin{proof}
The one-block conditional estimate implies
\(\mu(E_b(A))\leq p_A\).  Apply the union bound over \(S_h\).
\end{proof}

\begin{theorem}[Local-exit completion of the V24 passage]
\label{thm:n9v25-localclosure}
Let \((g_n,\mathfrak T_n)\) converge in law locally in the topology of
\Cref{thm:n9v24-expectation}.  Suppose thresholds \(A_n\) satisfy LCERC
with \(p_{A_n}\to0\), and for some \(\varepsilon>0\),
\begin{equation}
 \sup_n\int|\mathscr E_h(g_n,\mathfrak T_n)|^{1+\varepsilon}
 \,d\mu_n<\infty.                                   \label{eq:n9v25-residual-UI}
\end{equation}
If the V24 closure and finite-cutoff residual identity hold on
\(F_{A_n}(h)\), then the bad sectors contribute zero and
\begin{equation}
 \int\mathscr E_h\,d\mu_n\longrightarrow
 \int\mathscr E_h\,d\mu.                            \label{eq:n9v25-localpass}
\end{equation}
In particular, no event requiring every block in the full volume to be
good is needed.
\end{theorem}

\begin{proof}
Write the expectation as its parts on \(F_{A_n}(h)\) and its complement.
The good part passes by the V24 strong closure hypothesis.  Holder's
inequality and \Cref{lem:n9v25-union} give
\begin{align*}
 \int_{F_{A_n}(h)^c}|\mathscr E_h|\,d\mu_n
 &\leq
 \left(\int|\mathscr E_h|^{1+\varepsilon}d\mu_n\right)^{1/(1+\varepsilon)}
 \mu_n(F_{A_n}(h)^c)^{\varepsilon/(1+\varepsilon)}\\
 &\leq C_h\bigl(|S_h|p_{A_n}\bigr)^{\varepsilon/(1+\varepsilon)}
 \longrightarrow0.
\end{align*}
The same truncation argument identifies the limit expectation.  If the
finite-cutoff good residual vanishes, only the displayed vanishing error
remains.
\end{proof}

\begin{firewall}[A growing support needs the polymer estimate]
\label{fire:n9v25-growing}
The union bound in \Cref{lem:n9v25-union} is sufficient for each fixed
compactly supported test tensor.  Uniform constructions with growing
supports, connected expansions, or thermodynamic pressures require the
stronger summability \eqref{eq:n9v25-polymer} and compatible bounds on
mixed good/bad activities.
\end{firewall}

\begin{assessment}[Progress at ninth-series V25]
\label{ass:n9v25-status}
The audit replaces an impossible global regularity target by a local one.
Finite coloring turns a uniform conditional metric-and-stress exit bound
into arbitrary-set domination and summable bad components.  For every
fixed Einstein test tensor, a vanishing local exit probability plus a
\(1+\varepsilon\) residual moment is enough to remove the bad sector.
The next bottleneck is the conditional one-block estimate itself.
\end{assessment}

\begin{programme}[Ninth-series V26: a coercive conditional loader]
\label{prog:n9v25-V26}
V26 will derive an explicit conditional Sobolev-exit bound from a local
coercive regulator, including the normalization denominator and boundary
oscillation.  It will then state exactly what uniformity is required when
that regulator is removed.
\end{programme}

\begin{firewall}[Claim boundary after ninth-series V25]
\label{fire:n9v25-boundary}
V25 does not prove LCERC for the physical law, UCTBB for Yang--Mills,
Navier--Stokes regularity, the renormalized Standard Model stress tensor,
conformal stability, reconstruction, or the full SM--Einstein continuum.
\end{firewall}

\begin{theorem}[Ninth-series V25 result]
\label{thm:n9v25-result}
The compulsory audit is complete.  Under LCERC, metric-and-stress exits
obey \eqref{eq:n9v25-domination}, their connected components satisfy
\eqref{eq:n9v25-polymer}, and the local bad contribution to each tested
Einstein residual vanishes by \eqref{eq:n9v25-residual-UI}.  Thus the V24
tensor closure needs local conditional control rather than a global
all-good event.
\end{theorem}

\begin{proof}
Use \Cref{prop:n9v25-decision,thm:n9v25-domination,
cor:n9v25-polymer,thm:n9v25-localclosure}.
\end{proof}

\paragraph{Parent-version record.}
V25 was formed from
\texttt{Renewal\_SM\_Einstein\_Continuum\_Research\_Notes\_V24.tex}.
The V24 parent had \(8\,342\) logical source lines, \(291\,123\) bytes, and
SHA--256
\[
 \texttt{36DE2D142D9195AC11412189C860CED5C9B5BCC191A237BD73C08FE2DCEB8255}.
\]
The compulsory audit was completed before selecting V26, and V26 follows
it as required by the five-version protocol.

% =============================================================================
% END APPENDED NINTH-SERIES V25 MATERIAL
% =============================================================================


% =============================================================================
% BEGIN APPENDED NINTH-SERIES V26 MATERIAL
% =============================================================================

\clearpage
\section{Ninth-series V26: a normalized coercive loader for local exits}
\label{sec:v9v26}

\subsection{Conditional block law and its normalization}
\label{sec:n9v26-law}

Fix one finite-cutoff block with interior configuration space \(X_a\),
reference probability \(m_a\), and boundary label \(\eta\).  The subscript
\(a\) may include lattice spacing, volume, gauge-fixing labels, and finite
mode cutoffs.  Let \({\cal N}_a:X_a\to[0,\infty)\) be a measurable
high-regularity size.  Consider the conditional law
\begin{equation}
 d\mu_{a,\eta}(x)
 ={e^{-H_a(x)+u_{a,\eta}(x)}\over Z_{a,\eta}},dm_a(x),
 qquad
 Z_{a,\eta}:=\int e^{-H_a+u_{a,\eta}},dm_a.          \label{eq:n9v26-law}
\end{equation}

\begin{definition}[Normalized coercive block data, NCBD]
\label{def:n9v26-NCBD}
NCBD with exponent \(p>0\) consists of numbers
\(c_a>0,b_a,B_a,D_a\geq0\), a measurable core \(C_a\subset X_a\), and
\(m_{0,a}>0\), such that uniformly in the admissible boundary label
\(\eta\),
\begin{align}
 H_a(x)&\geq c_a{\cal N}_a(x)^p-b_a
       &&\text{for }m_a\text{-a.e. }x,               \label{eq:n9v26-coercive}\\
 H_a(x)&\leq B_a &&\text{for }m_a\text{-a.e. }x\in C_a,              \label{eq:n9v26-core-energy}\\
 m_a(C_a)&\geq m_{0,a},                              \label{eq:n9v26-core-mass}\\
 \operatorname*{ess\,sup}_{x,y\in X_a}
 |u_{a,\eta}(x)-u_{a,\eta}(y)|&\leq D_a.             \label{eq:n9v26-osc}
\end{align}
No absolute bound on \(u_{a,\eta}\) is required; its additive constant
cancels from \eqref{eq:n9v26-law}.
\end{definition}

\begin{lemma}[Oscillation comparison]
\label{lem:n9v26-osc}
Let
\(
 d\nu_a=Z_a^{-1}e^{-H_a}dm_a
\).
Under \eqref{eq:n9v26-osc}, every measurable \(E\subset X_a\) satisfies
\begin{equation}
 \mu_{a,\eta}(E)\leq e^{D_a}\nu_a(E).                \label{eq:n9v26-comparison}
\end{equation}
\end{lemma}

\begin{proof}
Let \(u^-\) and \(u^+\) be the essential infimum and supremum of
\(u_{a,\eta}\).  Bound the numerator of \(\mu(E)\) by
\(e^{u^+}\int_Ee^{-H_a}dm_a\), and its denominator from below by
\(e^{u^-}\int e^{-H_a}dm_a\).  Their ratio costs
\(e^{u^+-u^-}\leq e^{D_a}\).
\end{proof}

\begin{theorem}[Normalized conditional Sobolev tail]
\label{thm:n9v26-tail}
NCBD implies, for every \(R\geq0\),
\begin{equation}
 \mu_{a,\eta}({\cal N}_a>R)
 \leq
 \exp\!\left[
 \Omega_a-c_aR^p
 \right],
 \qquad
 \Omega_a:=D_a+B_a+b_a+\log m_{0,a}^{-1}.            \label{eq:n9v26-tail}
\end{equation}
The estimate is uniform in \(\eta\) and contains the complete conditional
normalization.
\end{theorem}

\begin{proof}
The unperturbed partition function has the lower bound
\begin{equation}
 Z_a=\int e^{-H_a}dm_a
 \geq\int_{C_a}e^{-H_a}dm_a
 \geq m_{0,a}e^{-B_a}.                               \label{eq:n9v26-Zlower}
\end{equation}
Because \(m_a\) is a probability and \eqref{eq:n9v26-coercive} holds,
\[
 \int_{\{{\cal N}_a>R\}}e^{-H_a}dm_a
 \leq e^{b_a-c_aR^p}.
\]
Divide by \eqref{eq:n9v26-Zlower}, then apply
\Cref{lem:n9v26-osc}.
\end{proof}

\begin{firewall}[A numerator estimate is insufficient]
\label{fire:n9v26-normalization}
The coercive numerator \(e^{b_a-c_aR^p}\) is not a probability bound
until the denominator has been controlled.  The terms
\(B_a+\log m_{0,a}^{-1}\) in \(\Omega_a\) record the exact normalization
cost.  They may diverge with a local mode cutoff and cannot be silently
absorbed into a uniform constant.
\end{firewall}

\subsection{Loading the geometric good event}
\label{sec:n9v26-geometry}

Let \(s>\max\{2,d/2+1\}\) on a \(d\)-dimensional block.  Suppose the fundamental
block variables \(z\) determine a metric fluctuation \(q(z)\), a positive
metric
\begin{equation}
 g(z)=\bar g^{1/2}\exp(2q(z))\bar g^{1/2},            \label{eq:n9v26-expmetric}
\end{equation}
and a renormalized stress candidate \(\mathfrak T(z)\).  Assume
\begin{equation}
 \|q(z)\|_{H^s}+\|\mathfrak T(z)\|_{L^2}
 \leq F_a({\cal N}_a(z))                             \label{eq:n9v26-map}
\end{equation}
for an increasing deterministic function \(F_a\).

\begin{lemma}[Deterministic threshold map]
\label{lem:n9v26-threshold}
For every \(R<\infty\), the event \({\cal N}_a\leq R\) implies V24
bounds of the form
\begin{equation}
 \lambda_a(R)I\leq g\leq\Lambda_a(R)I,qquad
 \|\nabla g\|_\infty\leq K_a(R),qquad
 \|g\|_{H^2}+\|\mathfrak T\|_2\leq L_a(R),          \label{eq:n9v26-threshold}
\end{equation}
where all four threshold functions are finite, with
\(\lambda_a(R)>0\).
\end{lemma}

\begin{proof}
Sobolev embedding gives
\(\|q\|_{W^{1,\infty}}\leq C_sF_a(R)\).  The matrix exponential in
\eqref{eq:n9v26-expmetric} has eigenvalues between
\(e^{-2C_sF_a(R)}\) and \(e^{2C_sF_a(R)}\).  Combining these with the
fixed background ellipticity gives positive finite choices of
\(\lambda_a(R)\) and \(\Lambda_a(R)\).  The chain rule, the product rule,
and the smooth functional calculus for the matrix exponential give finite
\(W^{1,\infty}\) and \(H^2\) bounds depending only on \(F_a(R)\), the
background metric, and the block.  The stress bound is already contained
in \eqref{eq:n9v26-map}.
\end{proof}

\begin{corollary}[NCBD loads the LCERC one-block row]
\label{cor:n9v26-LCERC}
Let \(A_a(R)\) denote the thresholds in
\eqref{eq:n9v26-threshold}.  Under NCBD and
\eqref{eq:n9v26-map},
\begin{equation}
 \mu_{a,\eta}\bigl(E_b(A_a(R))\bigr)
 \leq p_a(R):=\min\{1,e^{\Omega_a-c_aR^p}\}.         \label{eq:n9v26-LCERC}
\end{equation}
If the finite-range conditional factorization row of LCERC also holds,
then all conclusions of V25 apply with \(p_A=p_a(R)\).
\end{corollary}

\begin{proof}
By \Cref{lem:n9v26-threshold}, failure of the geometric good event implies
\({\cal N}_a>R\).  Apply \Cref{thm:n9v26-tail}; the remaining statement is
\Cref{thm:n9v25-domination,cor:n9v25-polymer}.
\end{proof}

\begin{firewall}[The exponential metric is a positivity device]
\label{fire:n9v26-chart}
Formula \eqref{eq:n9v26-expmetric} is a sufficient global parametrization
of positive metrics, not a derived gauge choice for the physical path
integral.  Any other parametrization may replace it if it supplies the
deterministic implication \eqref{eq:n9v26-threshold} with controlled
threshold growth.
\end{firewall}

\subsection{Diagonal removal of a coercive regulator}
\label{sec:n9v26-diagonal}

The useful case has \(c_a\downarrow0\) as an auxiliary high-derivative
regulator is removed.  The radius must then grow.  Write the downstream
bad-sector loss as \({\cal L}_a(R)\); it includes ellipticity constants,
test norms, finite-color factors, and any mixed good/bad activity estimate.

\begin{theorem}[Explicit regulator-removal balance]
\label{thm:n9v26-diagonal}
Assume \(p>q\geq0\), \(\theta>0\), and
\begin{equation}
 p_a(R)\leq e^{\Omega_a-c_aR^p},qquad
 {cal L}_a(R)\leq e^{d_a(1+R^q)},                   \label{eq:n9v26-loss}
\end{equation}
where \(c_a>0\) and \(d_a,\Omega_a<\infty\) at every finite cutoff.
There exists a diagonal choice \(R_a\to\infty\) such that
\begin{equation}
 p_a(R_a)\to0,qquad
 {cal L}_a(R_a)p_a(R_a)^\theta\to0.                 \label{eq:n9v26-balance}
\end{equation}
The conclusion allows \(c_a\downarrow0\), provided the finite-cutoff
coercivity is strictly positive.
\end{theorem}

\begin{proof}
For each cutoff choose \(R_a\geq a^{-1}\) so large that
\begin{equation}
 \theta c_aR_a^p
 -\theta\Omega_a-d_a(1+R_a^q)\geq a^{-1}.           \label{eq:n9v26-choice}
\end{equation}
Such a choice exists because \(p>q\) and \(c_a>0\).  Then
\[
 \log\bigl({\cal L}_a(R_a)p_a(R_a)^\theta\bigr)
 \leq-a^{-1}\longrightarrow-\infty.
\]
Dropping the nonnegative loss term in \eqref{eq:n9v26-choice} also gives
\(\theta(c_aR_a^p-\Omega_a)\to\infty\), hence
\(p_a(R_a)\to0\).
\end{proof}

The notation \(a^{-1}\to\infty\) in the preceding theorem means that the
cutoff parameter is indexed along a sequence tending to zero.  For a
general sequence index \(n\), replace every occurrence of \(a^{-1}\) in
the choice by \(n\).

\begin{corollary}[Polynomial and exponential threshold losses]
\label{cor:n9v26-exp-loss}
If the deterministic metric map has at most exponential loss
\({\cal L}_a(R)\leq e^{d_a(1+R)}\) and the regulator is quadratic
\((p=2)\), then \Cref{thm:n9v26-diagonal} applies with \(q=1\).
In particular, the exponential ellipticity constants generated by
\eqref{eq:n9v26-expmetric} can be dominated by a sufficiently large
quadratic-regulator radius at each finite cutoff.
\end{corollary}

\begin{firewall}[Exit balance does not create continuum compactness]
\label{fire:n9v26-compactness}
The diagonal theorem makes bad-sector errors vanish.  If
\(c_a\downarrow0\), it does not give a cutoff-uniform moment of
\({\cal N}_a\), strong \(H^2\) tightness, or convergence of
\(\mathfrak T_a\).  Those are separate good-sector obligations in TELC.
Allowing \(R_a\) to grow merely moves the boundary between the good and
bad sectors.
\end{firewall}

\subsection{A moment consequence when coercivity is uniform}
\label{sec:n9v26-moment}

\begin{lemma}[Uniform stretched-exponential moment]
\label{lem:n9v26-moment}
Suppose \(c_a\geq c_*>0\) and \(\Omega_a\leq\Omega_*<\infty\).  Then for
every \(0<\tau<c_*\),
\begin{equation}
 \sup_{a,\eta}
 \int e^{\tau{cal N}_a^p}\,d\mu_{a,\eta}
 \leq
 1+{\tau e^{\Omega_*}\over c_*-\tau}.               \label{eq:n9v26-expmoment}
\end{equation}
Consequently every polynomial moment of \({\cal N}_a\) is uniform.
\end{lemma}

\begin{proof}
For \(Y={\cal N}_a^p\), the layer-cake identity gives
\[
 \mathbb E e^{\tau Y}
 =1+\tau\int_0^\infty e^{\tau t}\mathbb P(Y>t)\,dt.
\]
Use \eqref{eq:n9v26-tail} with the bound
\(\mathbb P(Y>t)\leq e^{\Omega_*-c_*t}\) and integrate.  Polynomial
moments follow from domination by a stretched exponential.
\end{proof}

\begin{assessment}[Progress at ninth-series V26]
\label{ass:n9v26-status}
The V25 conditional-exit card now has a precise analytic loader.  Local
coercivity, an explicit core lower bound for the partition function, and
uniform boundary oscillation give the normalized tail
\eqref{eq:n9v26-tail}.  A high Sobolev control map converts it into the
metric, ellipticity, and stress exits required by V24.  Quadratic
coercivity can beat exponential threshold losses along a diagonal, but
good-sector compactness and the physical stress construction remain open.
\end{assessment}

\begin{programme}[Ninth-series V27: renormalized matter stress]
\label{prog:n9v26-V27}
V27 will define the finite-cutoff Standard Model stress density by metric
variation, separate improvement and counterterm contributions, prove its
diffeomorphism Ward identity under anomaly-free regularization, and state
the exact topology and moments needed for the V24--V26 closure chain.
\end{programme}

\begin{firewall}[Claim boundary after ninth-series V26]
\label{fire:n9v26-boundary}
V26 does not prove NCBD for the unregulated Einstein action, uniform
coercivity after regulator removal, the stress control map
\eqref{eq:n9v26-map}, TELC for the physical measure, reflection positivity,
reconstruction, or the full SM--Einstein continuum.
\end{firewall}

\begin{theorem}[Ninth-series V26 result]
\label{thm:n9v26-result}
NCBD gives the fully normalized conditional tail
\eqref{eq:n9v26-tail}; a deterministic Sobolev-to-geometry map loads the
LCERC one-block probability \eqref{eq:n9v26-LCERC}; and whenever the
coercive exponent dominates the logarithmic downstream loss,
\Cref{thm:n9v26-diagonal} removes the bad-sector error along a diagonal.
Uniform positive coercivity additionally supplies all polynomial moments.
\end{theorem}

\begin{proof}
Use \Cref{thm:n9v26-tail,lem:n9v26-threshold,
cor:n9v26-LCERC,thm:n9v26-diagonal,lem:n9v26-moment}.
\end{proof}

\paragraph{Parent-version record.}
V26 was formed from
\texttt{Renewal\_SM\_Einstein\_Continuum\_Research\_Notes\_V25.tex}.
The V25 parent had \(8\,654\) logical source lines, \(303\,632\) bytes, and
SHA--256
\[
 \texttt{493C54481AED0DB8BBFE1BB00E96FAC09C3DEACD47E12AB4C9247B970726CD61}.
\]
V26 is the required post-audit continuation.

% =============================================================================
% END APPENDED NINTH-SERIES V26 MATERIAL
% =============================================================================


% =============================================================================
% BEGIN APPENDED NINTH-SERIES V27 MATERIAL
% =============================================================================

\clearpage
\section{Ninth-series V27: measure stress and insertion-complete Einstein landing}
\label{sec:v9v27}

\subsection{Nonduplication audit of the stress programme}
\label{sec:n9v27-audit}

The frozen first-series archive already contains the classical minimal and
Higgs-improved Hilbert stress, the variation of
\(\xi_HR H^\dagger H\), the assembled diffeomorphism Ward identity, the
critical Higgs concentration stress, the Einstein-frame diagonalization,
the state-evaluation bridge, the one-loop proper-time stress telescope, and
the invariant local RG projection of the first interacting shell
derivative.  In particular, the exact first-cumulant row
\[
 \delta_g\langle I\rangle_0
 =\langle\delta_gI\rangle_0
 -\hbar^{-1}\operatorname{Cov}_0(I,\delta_gS_0)
\]
and the critical Higgs defect decomposition were proved there.  Those
calculations are not repeated.

The new issue created by V24 is a type mismatch.  Its strong tensor closure
uses an \(L^2\) stress density, whereas critical weak compactness in the
archive can produce a tensor-valued Radon measure.  V27 determines what the
Einstein equation itself says about that measure and distinguishes a mean
residual identity from an almost-sure distributional equation.

\subsection{Deterministic landing with a measure-valued source}
\label{sec:n9v27-measure}

Let \(U\Subset\mathbb R^d\), and let
\(\mathcal M(U;\Sym^2)\) denote finite symmetric contravariant
tensor-valued Radon measures, with componentwise weak-star convergence.
The measures below are densitized, so no additional volume factor is
attached to them.  Put
\begin{equation}
 \mathfrak E^{ij}_{\Lambda_c}(g)
 :=\sqrt{\det g}\bigl(G^{ij}(g)+\Lambda_cg^{ij}\bigr).              \label{eq:n9v27-geometric}
\end{equation}

\begin{theorem}[Measure-source closure]
\label{thm:n9v27-measure-closure}
Let \(g_n,g\) obey one common V24 class
\eqref{eq:n9v24-class} and let \(g_n\to g\) strongly in \(H^2(U)\).
Suppose
\begin{equation}
 \tau_n\overset{*}{\rightharpoonup}\tau
 \quad\text{in }\mathcal M(U;\Sym^2),               \label{eq:n9v27-measure-conv}
\end{equation}
and, for every \(h\in C_c^\infty(U;\Sym^2T^*U)\),
\begin{equation}
 2\chi_R\int_U\mathfrak E^{ij}_{\Lambda_c}(g_n)h_{ij}\,dx
 -\int_Uh_{ij}\,d\tau_n^{ij}\longrightarrow0.       \label{eq:n9v27-residual}
\end{equation}
Then
\begin{equation}
 \tau^{ij}
 =2\chi_R\mathfrak E^{ij}_{\Lambda_c}(g)\,dx         \label{eq:n9v27-limit-equation}
\end{equation}
as tensor distributions.  In particular, if \(\chi_R\) and \(\Lambda_c\)
are finite, \(\tau\) is absolutely continuous and has an \(L^2\) density.
\end{theorem}

\begin{proof}
By \Cref{thm:n9v24-stability},
\(\mathfrak G(g_n)\to\mathfrak G(g)\) in \(L^2\).  The determinant and
inverse coefficients converge boundedly and pointwise along every relevant
subsequence, so the cosmological density converges in every finite
\(L^p\).  Thus the geometric pairing in
\eqref{eq:n9v27-residual} converges.  The source pairing converges by
\eqref{eq:n9v27-measure-conv}.  Passing to the limit gives
\eqref{eq:n9v27-limit-equation}.  Its right side belongs to \(L^2dx\), so
uniqueness of the Lebesgue decomposition of each component of \(\tau\)
removes its singular part.
\end{proof}

\begin{corollary}[Only the total singular stress is excluded]
\label{cor:n9v27-total-defect}
Suppose the archived stress decomposition has a regular part and defect
measures
\begin{equation}
 \tau=\tau_{\rm reg},dx
 +\mathfrak K-\lambda_Hg\,\nu_H-2\xi_H\mathfrak I+\mathfrak S_{\rm other},
                                                               \label{eq:n9v27-defects}
\end{equation}
where \(\nu_H\geq0\) is the quartic Higgs concentration measure.  Under
\Cref{thm:n9v27-measure-closure},
\begin{equation}
 \bigl(mathfrak K-\lambda_Hg\,\nu_H
 -2\xi_H\mathfrak I+\mathfrak S_{\rm other}\bigr)^{\rm sing}=0.       \label{eq:n9v27-total-singular}
\end{equation}
This does not imply that the individual measures in
\eqref{eq:n9v27-defects} vanish.
\end{corollary}

\begin{proof}
The total source has no singular part by
\Cref{thm:n9v27-measure-closure}; subtract the absolutely continuous
regular part.  Individual vanishing does not follow because two nonzero
tensor measures can cancel: for any nonzero tensor measure \(\sigma\), the
pair \((\sigma,-\sigma)\) has zero sum.
\end{proof}

\begin{firewall}[No sectorwise positivity is available]
\label{fire:n9v27-cancellation}
Although \(\nu_H\) is nonnegative, it is multiplied by the tensor
\(-\lambda_Hg\) and is combined with kinetic and improvement defects.
Stress has no general positive ordering that forbids cancellation in
\eqref{eq:n9v27-total-singular}.  Removing \(\nu_H\) still requires a
compactness input such as the archived local Orlicz or monotonicity gate,
or a new sign-separating theorem.
\end{firewall}

\subsection{Mean equations do not give samplewise equations}
\label{sec:n9v27-mean}

Let \((\Omega,\mathcal F,\mathbb P)\) carry random \(g\) and \(\tau\).
For a test tensor \(h\), define the scalar residual
\begin{equation}
 X_h
 :=2\chi_R\int_U\mathfrak E^{ij}_{\Lambda_c}(g)h_{ij}\,dx
 -\int_Uh_{ij}\,d\tau^{ij}.                          \label{eq:n9v27-Xh}
\end{equation}

\begin{proposition}[Failure of the mean-to-pathwise inference]
\label{prop:n9v27-mean-failure}
The identities
\begin{equation}
 \mathbb E X_h=0\qquad\text{for every }h              \label{eq:n9v27-mean}
\end{equation}
do not imply \(X_h=0\) almost surely.
\end{proposition}

\begin{proof}
Fix a nonzero deterministic tensor measure \(\sigma\) and a Rademacher
random variable \(\epsilon\) with equal probabilities at \(\pm1\).  Keep
the metric deterministic and replace a solution source by
\(\tau+\epsilon\sigma\).  The extra residual has zero expectation for
every \(h\), but is nonzero with probability one for some \(h\).
\end{proof}

The correct probabilistic strengthening uses insertions.  Let
\(\mathcal A\) be a unital algebra of bounded cylinder observables which
generates \(\mathcal F\).

\begin{definition}[Insertion-complete Einstein residual, ICER]
\label{def:n9v27-ICER}
ICER holds on a countable distribution-determining family
\(\{h_m\}_{m\geq1}\subset C_c^\infty(U;\Sym^2T^*U)\) if
\begin{equation}
 X_{h_m}\in L^1(\mathbb P),
 qquad
 \mathbb E[F X_{h_m}]=0
 \quad\text{for every }F\in\mathcal A, m\geq1.      \label{eq:n9v27-ICER}
\end{equation}
The family is distribution determining when a distribution vanishing on
all \(h_m\) vanishes on every test tensor.
\end{definition}

\begin{theorem}[Insertion completeness gives the pathwise equation]
\label{thm:n9v27-ICER}
If ICER holds, then with probability one
\begin{equation}
 X_h=0
 \quad\text{for every }h\in C_c^\infty(U;\Sym^2T^*U).               \label{eq:n9v27-pathwise}
\end{equation}
Consequently, on the event that \(g\) belongs to the V24 class and
\(\tau\) is locally finite,
\begin{equation}
 \tau=2\chi_R\mathfrak E_{\Lambda_c}(g)\,dx
 \quad\text{and}\quad \tau^{\rm sing}=0              \label{eq:n9v27-pathwise-ac}
\end{equation}
almost surely.
\end{theorem}

\begin{proof}
Fix \(m\).  The class of bounded random variables \(F\) satisfying
\(\mathbb E[FX_{h_m}]=0\) is a monotone class containing the generating
algebra \(\mathcal A\).  The functional monotone-class theorem therefore
extends the identity to every bounded \(\mathcal F\)-measurable \(F\).
Taking bounded approximations to
\(\operatorname{sgn}(X_{h_m})\) gives
\(\mathbb E|X_{h_m}|=0\).  Intersect the resulting probability-one events
over the countable set of \(m\).  On this intersection the random residual
distribution vanishes on the determining family and hence on every test
tensor.  This proves \eqref{eq:n9v27-pathwise}; uniqueness of the Lebesgue
decomposition then proves \eqref{eq:n9v27-pathwise-ac} sample by sample.
\end{proof}

\begin{remark}[Why the countable family is enough]
\label{rem:n9v27-countable}
The test-function space is separable.  On each compact exhaustion choose a
countable dense set in every seminorm and take their union.  A distribution
is continuous in the test-function topology, so vanishing on this family
determines it.
\end{remark}

\subsection{Passing the inserted identities from finite cutoff}
\label{sec:n9v27-passage}

\begin{theorem}[Uniform-integrability passage for ICER]
\label{thm:n9v27-passage}
Let \((g_n,\tau_n)\Rightarrow(g,\tau)\) in a topology for which the metric
convergence is the V24 strong local \(H^2\) convergence on local good sets
and the source convergence is local weak-star convergence of measures.
Let \(F\) be a bounded continuous cylinder observable and \(h=h_m\).  If
\begin{equation}
 \sup_n\mathbb E|X_{n,h}|^{1+\varepsilon}<\infty       \label{eq:n9v27-UI}
\end{equation}
for some \(\varepsilon>0\), and local exits satisfy the V25--V26 vanishing
debit, then
\begin{equation}
 \mathbb E[F X_{n,h}]\longrightarrow\mathbb E[F X_h].               \label{eq:n9v27-inserted-pass}
\end{equation}
Hence finite-cutoff inserted residual identities pass to every observable
in a countable convergence-determining cylinder algebra.
\end{theorem}

\begin{proof}
On the local good event, V24 gives continuity of the geometric pairing and
weak-star convergence gives continuity of the source pairing.  Multiplying
by bounded continuous \(F\) preserves convergence.  Condition
\eqref{eq:n9v27-UI} gives uniform integrability.  The bad-event term is
bounded by Holder exactly as in \Cref{thm:n9v25-localclosure}; its
probability tends to zero by V26.  Truncate, pass the bounded continuous
part under weak convergence, and remove the truncation uniformly.
\end{proof}

\begin{corollary}[Conditional almost-sure Einstein landing]
\label{cor:n9v27-landing}
Assume the hypotheses of \Cref{thm:n9v27-passage} simultaneously for the
countable test family and a countable generating cylinder algebra.  If the
finite-cutoff inserted residuals vanish, then the limiting law satisfies
the Einstein equation \eqref{eq:n9v27-pathwise-ac} almost surely.
\end{corollary}

\begin{proof}
Pass every pair \((F,h_m)\) by
\Cref{thm:n9v27-passage}; countability places all identities on one event.
Apply \Cref{thm:n9v27-ICER}.
\end{proof}

\begin{firewall}[Ward identities and Einstein insertions are distinct]
\label{fire:n9v27-Ward}
The archived diffeomorphism Ward identity constrains the covariant
divergence of the complete assembled source.  ICER asserts the full tested
Einstein residual with arbitrary bounded cylinder insertions.  Neither the
uninserted mean equation nor the Ward identity implies ICER.  Moreover,
before absolute continuity is known, multiplying a singular stress measure
by an almost-everywhere-defined \(L^\infty\) Christoffel symbol is not
canonical; a measure-level Ward identity needs a \(C^1\) metric,
absolute continuity, or a separately defined product.
\end{firewall}

\subsection{Updated stress landing card}
\label{sec:n9v27-card}

\begin{definition}[Measure and insertion Einstein landing card, MIELC]
\label{def:n9v27-MIELC}
MIELC records:
\begin{enumerate}[label=\textup{(M\arabic*)},leftmargin=4em]
\item V24 strong local metric closure on the V25 good blocks;
\item V26 normalized conditional exits with a vanishing diagonal debit;
\item local weak-star tightness of the complete assembled densitized
stress as one tensor-valued measure;
\item the moment \eqref{eq:n9v27-UI} for a countable test family;
\item finite-cutoff Einstein residual identities with every insertion in a
countable generating cylinder algebra; and
\item convergence of the complete stress after improvement, BV source,
state evaluation, and invariant counterterm projection have been assembled.
\end{enumerate}
\end{definition}

\begin{theorem}[MIELC closes the singular-source interface]
\label{thm:n9v27-MIELC}
MIELC implies an almost-sure distributional Einstein equation whose total
stress is an \(L^2\) density and has no singular Radon part.  It does not by
itself identify or eliminate cancellations among the sectorwise defect
measures.
\end{theorem}

\begin{proof}
Items (M1)--(M5) give
\Cref{cor:n9v27-landing}; item (M6) identifies its source as the complete
renormalized SM stress.  Apply \Cref{cor:n9v27-total-defect} for the final
limitation.
\end{proof}

\begin{assessment}[Progress at ninth-series V27]
\label{ass:n9v27-status}
The archived stress calculations have been reused at their actual endpoint
rather than repeated.  V27 proves that V24 regularity excludes a singular
total stress once the Einstein equation holds pathwise.  It also isolates
the missing logical strength: an uninserted expectation equation controls
only a barycenter, while a countable insertion-complete hierarchy forces
the residual distribution to vanish almost surely.
\end{assessment}

\begin{programme}[Ninth-series V28: acquiring inserted residuals]
\label{prog:n9v27-V28}
V28 will derive a finite-dimensional Schwinger--Dyson integration-by-parts
identity for the metric variables with bounded cylinder insertions,
including the reference-measure score, constraint/coarea score, and
boundary terms.  It will determine whether that identity is the Einstein
residual required by ICER or contains an additional regulator score that
must be shown to vanish.
\end{programme}

\begin{firewall}[Claim boundary after ninth-series V27]
\label{fire:n9v27-boundary}
V27 does not prove MIELC, finite-cutoff inserted Einstein identities,
stress-measure tightness, absence of sectorwise concentration, interacting
shell summability, conformal stability, reconstruction, or the full
SM--Einstein continuum.
\end{firewall}

\begin{theorem}[Ninth-series V27 result]
\label{thm:n9v27-result}
Strong V24 metric convergence and weak-star convergence of the complete
stress measure pass the tested Einstein equation.  Any pathwise limit
source is then an \(L^2\) density.  ICER is sufficient to upgrade the full
inserted expectation hierarchy to that pathwise equation, and the
V25--V26 exit mechanism passes ICER from finite cutoff under
\eqref{eq:n9v27-UI}.
\end{theorem}

\begin{proof}
Use \Cref{thm:n9v27-measure-closure,thm:n9v27-ICER,
thm:n9v27-passage,thm:n9v27-MIELC}.
\end{proof}

\paragraph{Parent-version record.}
V27 was formed from
\texttt{Renewal\_SM\_Einstein\_Continuum\_Research\_Notes\_V26.tex}.
The V26 parent had \(8\,976\) logical source lines, \(316\,302\) bytes, and
SHA--256
\[
 \texttt{8AE5F3ECCC74B4CEDBBEE7019E0C8141589CAAA3414BF332E884B77D57F23B1F}.
\]

% =============================================================================
% END APPENDED NINTH-SERIES V27 MATERIAL
% =============================================================================


% =============================================================================
% BEGIN APPENDED NINTH-SERIES V28 MATERIAL
% =============================================================================

\clearpage
\section{Ninth-series V28: the Einstein Schwinger--Dyson score}
\label{sec:v9v28}

\subsection{Finite-dimensional integration by parts}
\label{sec:n9v28-IBP}

Let \(X\) be an oriented finite-dimensional manifold, possibly with
piecewise smooth boundary, let \(m\) be a positive smooth volume density,
and let
\begin{equation}
 d\mu=Z^{-1}e^{-S/\hbar},dm,
 qquad 0<Z<\infty,quad\hbar>0.                     \label{eq:n9v28-law}
\end{equation}
Let \(V\) be a smooth vector field for which the following integrals are
finite.  Define its divergence relative to \(m\) by
\begin{equation}
 {cal L}_Vm=(\div_mV)m                              \label{eq:n9v28-div}
\end{equation}
and its assembled Schwinger--Dyson score by
\begin{equation}
 Q_V:=VS-\hbar\div_mV.                               \label{eq:n9v28-score}
\end{equation}

\begin{theorem}[Exact inserted score identity]
\label{thm:n9v28-SD}
For every smooth bounded \(F\) with the required integrability,
\begin{equation}
 \mathbb E_\mu[FQ_V]
 =\hbar\mathbb E_\mu[VF]-\hbar\mathfrak b_V(F),      \label{eq:n9v28-SD}
\end{equation}
where the normalized outward flux is
\begin{equation}
 \mathfrak b_V(F)
 :={1\over Z}\int_{\partial X}F e^{-S/\hbar}\iota_Vm.               \label{eq:n9v28-flux}
\end{equation}
If \(X\) has no boundary, or \(V\) is tangent to it, the flux vanishes.
\end{theorem}

\begin{proof}
Cartan's formula and Stokes' theorem give
\begin{align*}
 \int_{\partial X}Fe^{-S/\hbar}\iota_Vm
 &=\int_X{cal L}_V(Fe^{-S/\hbar}m)\\
 &=\int_Xe^{-S/\hbar}
 \left(VF+F\div_mV-\hbar^{-1}FVS\right)dm.
\end{align*}
Divide by \(Z\), multiply by \(\hbar\), and rearrange.
\end{proof}

In coordinates with \(dm=\rho(x)dx\),
\begin{equation}
 \div_mV=\partial_AV^A+V\log\rho.                   \label{eq:n9v28-coordinate-div}
\end{equation}
Thus the reference-measure score is part of \(Q_V\); it is not an optional
normalization constant.

\subsection{Independence from the action--measure split}
\label{sec:n9v28-split}

The finite regulator can place determinants, coarea factors, gauge-fixing
Jacobians, or local counterterms either in an effective action or in the
reference density.  The assembled score is invariant under this bookkeeping.

\begin{proposition}[Score partition invariance]
\label{prop:n9v28-partition}
Let \(A\) be smooth and set
\begin{equation}
 S':=S+A,qquad m':=e^{A/\hbar}m.                    \label{eq:n9v28-resplit}
\end{equation}
Then \(e^{-S'/\hbar}m'=e^{-S/\hbar}m\) and
\begin{equation}
 VS'-\hbar\div_{m'}V
 =VS-\hbar\div_mV.                                  \label{eq:n9v28-invariant}
\end{equation}
\end{proposition}

\begin{proof}
The equality of unnormalized measures is immediate.  Since
\(m'=e^{A/\hbar}m\),
\[
 \div_{m'}V=\div_mV+\hbar^{-1}VA.
\]
Substitution cancels \(VA\) exactly.
\end{proof}

\begin{corollary}[Coarea and determinant rows must be assembled]
\label{cor:n9v28-assembled}
An Einstein score computed after constraint elimination is unchanged if a
positive determinant or coarea factor is moved between the effective action
and reference density.  Estimating its two appearances separately can lose
their exact cancellation and does not estimate the invariant score
\eqref{eq:n9v28-score}.
\end{corollary}

\begin{proof}
Apply \Cref{prop:n9v28-partition} successively to the logarithm of every
positive factor.  Phase and signed factors require their already-declared
line-bundle or finite-sector treatment before this real logarithmic argument
applies.
\end{proof}

\subsection{Constraint surfaces and tangent metric variations}
\label{sec:n9v28-constraint}

Let \(C=\{\Phi=0\}\subset X\) be a regular constraint surface, and let its
reduced reference density be
\begin{equation}
 dm_C=\rho_CJ_\Phi^{-1},d\vol_C,                    \label{eq:n9v28-coarea}
\end{equation}
where \(J_\Phi>0\) is the normal coarea Jacobian.  If \(V\) is tangent to
\(C\), then in intrinsic coordinates
\begin{equation}
 \div_{m_C}V
 =\div_CV+V\log\rho_C-V\log J_\Phi.                 \label{eq:n9v28-coarea-score}
\end{equation}

\begin{lemma}[Reduced constrained score]
\label{lem:n9v28-constrained}
For a tangent vector field \(V\) on \(C\),
\Cref{thm:n9v28-SD} holds on \(C\) with
\begin{equation}
 Q_V^C
 =VS-\hbar\div_CV-\hbar V\log\rho_C
 +\hbar V\log J_\Phi.                               \label{eq:n9v28-constrained-score}
\end{equation}
\end{lemma}

\begin{proof}
Apply \Cref{thm:n9v28-SD} to the manifold \(C\) with volume density
\eqref{eq:n9v28-coarea}, and insert
\eqref{eq:n9v28-coarea-score}.
\end{proof}

\begin{firewall}[A free metric variation need not be tangent]
\label{fire:n9v28-tangent}
After Hamiltonian and momentum constraints are eliminated, an arbitrary
ambient metric variation generally moves \(C\).  The constrained identity
requires the tangent lift supplied by the coupled implicit-function chart,
including its lapse--shift multiplier response.  Applying
\eqref{eq:n9v28-constrained-score} to an ambient constant vector without
that lift omits a shape derivative and is invalid.
\end{firewall}

\subsection{Relation to the classical Einstein residual}
\label{sec:n9v28-Einstein}

Use the archived covariant-metric convention
\begin{equation}
 \delta S_{\rm matter}[h]
 ={1\over2}\int\mathfrak T^{ij}h_{ij},qquad
 \delta S_{\rm grav}[h]
 =-\chi_R\int\mathfrak E_{\Lambda_c}^{ij}(g)h_{ij}.  \label{eq:n9v28-convention}
\end{equation}
For the complete finite-cutoff action, after every improvement and
counterterm has been included, define
\begin{equation}
 \mathcal E_h
 :=\int
 \left(2\chi_R\mathfrak E_{\Lambda_c}^{ij}(g),dx
       -d\tau^{ij}\right)h_{ij}.                    \label{eq:n9v28-classical-residual}
\end{equation}
Then \(V_hS=-\frac12\mathcal E_h\), provided \(V_h\) is the complete
constraint-compatible metric variation and no row has been omitted.

\begin{theorem}[Quantum Einstein identity with contact term]
\label{thm:n9v28-quantum-Einstein}
Under the preceding convention and with vanishing boundary flux,
\begin{equation}
 \boxed{
 \mathbb E\left[
 F\left(\mathcal E_h+2\hbar\div_mV_h\right)
 \right]
 =-2\hbar\mathbb E[V_hF].}                           \label{eq:n9v28-quantum-Einstein}
\end{equation}
With a boundary, the right side gains
\(2\hbar\mathfrak b_{V_h}(F)\).
\end{theorem}

\begin{proof}
Insert \(V_hS=-\frac12\mathcal E_h\) into
\eqref{eq:n9v28-score} and then into \eqref{eq:n9v28-SD}.  Multiply the
result by \(-2\).
\end{proof}

\begin{assessment}[The divergence row is a quantum source]
\label{ass:n9v28-divergence}
The term \(2\hbar\div_mV_h\) contains the finite-dimensional base-measure,
gauge-fixing, determinant, coarea, chart, and constraint-lift responses not
already placed in the convention-fixed action.  By
\Cref{prop:n9v28-partition}, its sum with \(\mathcal E_h\) is independent
of the bookkeeping split.  A continuum stress definition must therefore
be attached to the assembled score, not to one selected density factor.
\end{assessment}

\subsection{Why insertion-complete classical residuals fail at fixed
\texorpdfstring{\(\hbar\)}{hbar}}
\label{sec:n9v28-no-go}

\begin{proposition}[ICER is not the fixed-\(\hbar\) path-integral identity]
\label{prop:n9v28-ICER-no-go}
Suppose \(X=\mathbb R^N\), \(m=dx\), \(V\) is a nonzero constant vector,
the boundary flux vanishes at infinity, and \(\mu\) is a probability with
a positive smooth density.  The score identities and
\begin{equation}
 \mathbb E[FQ_V]=0\quad\text{for every }F\in C_b^1(\mathbb R^N)       \label{eq:n9v28-false-ICER}
\end{equation}
cannot both hold.
\end{proposition}

\begin{proof}
Together with \Cref{thm:n9v28-SD},
\eqref{eq:n9v28-false-ICER} would give
\(\int VF,d\mu=0\) for every bounded \(C^1\) test function.  Hence
\(\div(V\mu)=0\) distributionally, so \(\mu\) is invariant under
translations along \(V\).  Choose a bounded slab of positive \(\mu\)-mass
whose translates by sufficiently separated multiples of \(V\) are
disjoint.  Translation invariance assigns them all the same positive mass,
contradicting finiteness of \(\mu\).
\end{proof}

\begin{corollary}[Correction to the V27 acquisition route]
\label{cor:n9v28-correction}
At fixed nonzero \(\hbar\), ordinary Schwinger--Dyson integration by parts
does not load the V27 insertion-complete classical residual.  It loads
\eqref{eq:n9v28-quantum-Einstein}, whose contact term is intrinsic.
MIELC remains a valid conditional theorem for a measure supported on the
classical Einstein locus, but it is not the generic quantum path-integral
landing card.
\end{corollary}

\begin{proof}
The contact term in \eqref{eq:n9v28-quantum-Einstein} is nonzero for general
metric-dependent insertions.  The no-go proposition shows that deleting it
for a generating family would contradict a nondegenerate finite probability
in a flat metric direction.
\end{proof}

\subsection{Two legitimate continuum targets}
\label{sec:n9v28-targets}

\begin{theorem}[Zero-noise route to the V27 pathwise equation]
\label{thm:n9v28-semiclassical}
Let \(\hbar_n\downarrow0\).  Suppose the finite-cutoff identity
\eqref{eq:n9v28-quantum-Einstein} holds and, for every pair in the V27
countable insertion/test family,
\begin{align}
 \sup_n\mathbb E|V_{n,h}F|&<\infty,                 \label{eq:n9v28-contact-bound}\\
 \hbar_n\mathbb E|F\div_{m_n}V_{n,h}|&\longrightarrow0,             \label{eq:n9v28-div-vanish}\\
 \hbar_n|\mathfrak b_{V_{n,h}}(F)|&\longrightarrow0.                \label{eq:n9v28-flux-vanish}
\end{align}
If the residuals also satisfy the V27 passage hypotheses, then ICER holds
in the limit and the V27 almost-sure classical Einstein equation follows.
\end{theorem}

\begin{proof}
The right side of \eqref{eq:n9v28-quantum-Einstein} tends to zero by
\eqref{eq:n9v28-contact-bound}; the divergence and flux rows tend to zero
by \eqref{eq:n9v28-div-vanish}--\eqref{eq:n9v28-flux-vanish}.  Pass the
remaining classical residual by \Cref{thm:n9v27-passage}, then apply
\Cref{thm:n9v27-ICER}.
\end{proof}

\begin{definition}[Quantum Einstein Schwinger--Dyson card, QESDC]
\label{def:n9v28-QESDC}
At fixed \(\hbar>0\), QESDC records:
\begin{enumerate}[label=\textup{(Q\arabic*)},leftmargin=4em]
\item a common countable cylinder core for observables and metric
variations;
\item the complete constraint-compatible vector fields \(V_h\);
\item the split-invariant scores \(Q_{V_h}\), with all determinant,
coarea, chart, BV, and counterterm rows assembled;
\item convergence and uniform integrability of \(FQ_{V_h}\) and
\(V_hF\);
\item convergence or vanishing of the boundary flux; and
\item closability of the limiting metric-derivative operator on the
cylinder core.
\end{enumerate}
\end{definition}

\begin{theorem}[Conditional fixed-\(\hbar\) landing]
\label{thm:n9v28-QESDC}
If QESDC holds and every finite regulator satisfies
\eqref{eq:n9v28-SD}, then the limiting measure satisfies
\begin{equation}
 \mathbb E_\mu[FQ_h]
 =\hbar\mathbb E_\mu[V_hF]-\hbar\mathfrak b_h(F)     \label{eq:n9v28-limit-SD}
\end{equation}
on the declared cylinder core.  This is the fixed-\(\hbar\) quantum
Einstein equation supplied by finite-dimensional path-integral integration
by parts.
\end{theorem}

\begin{proof}
Pass the three terms of the finite identity using items (Q3)--(Q5).
Item (Q6) identifies the limiting derivative as an operator rather than
only a family of unrelated cylinder formulas.
\end{proof}

\begin{firewall}[A Schwinger--Dyson hierarchy is not reconstruction]
\label{fire:n9v28-reconstruction}
QESDC supplies integration-by-parts identities for a Euclidean measure.
It does not prove reflection positivity, clustering, essential
self-adjointness of reconstructed generators, a Lorentzian Einstein
equation, or existence of the measure itself.
\end{firewall}

\begin{assessment}[Progress at ninth-series V28]
\label{ass:n9v28-status}
The correct finite-cutoff quantum metric identity is now fixed.  Its score
is invariant under transfers between action and measure, includes the
coarea and base-density responses, and carries the unavoidable insertion
contact term.  This corrects the proposed V27 acquisition route: pathwise
Einstein landing is appropriate for a zero-noise or stationary-support
limit, while the fixed quantum continuum requires QESDC.
\end{assessment}

\begin{programme}[Ninth-series V29: closability of the metric derivative]
\label{prog:n9v28-V29}
V29 will construct the cylinder gradient and its logarithmic derivative
from the score identity, prove an \(L^2\) closability criterion, and derive
the associated symmetric Dirichlet form.  This attacks item (Q6) without
assuming the continuum measure already has a differentiable density.
\end{programme}

\begin{firewall}[Claim boundary after ninth-series V28]
\label{fire:n9v28-boundary}
V28 does not prove QESDC for the physical regulator, convergence of the
assembled metric score, vanishing boundary flux, closability, conformal
stability, reflection positivity, reconstruction, or the full
SM--Einstein continuum.
\end{firewall}

\begin{theorem}[Ninth-series V28 result]
\label{thm:n9v28-result}
Finite-dimensional metric integration by parts gives the split-invariant
score identity \eqref{eq:n9v28-SD} and its Einstein form
\eqref{eq:n9v28-quantum-Einstein}.  The contact term prevents generic
fixed-\(\hbar\) ICER, while the controlled zero-noise limit
\eqref{eq:n9v28-contact-bound}--\eqref{eq:n9v28-flux-vanish} recovers it.
At fixed \(\hbar\), QESDC is the type-correct continuum target.
\end{theorem}

\begin{proof}
Use \Cref{thm:n9v28-SD,prop:n9v28-partition,
lem:n9v28-constrained,thm:n9v28-quantum-Einstein,
prop:n9v28-ICER-no-go,thm:n9v28-semiclassical,
thm:n9v28-QESDC}.
\end{proof}

\paragraph{Parent-version record.}
V28 was formed from
\texttt{Renewal\_SM\_Einstein\_Continuum\_Research\_Notes\_V27.tex}.
The V27 parent had \(9\,326\) logical source lines, \(331\,064\) bytes, and
SHA--256
\[
 \texttt{00362F0E48C608897BEE4B37F06105D56A35E066AEBED9E1A88A4D4DF5BC1D07}.
\]

% =============================================================================
% END APPENDED NINTH-SERIES V28 MATERIAL
% =============================================================================


% =============================================================================
% BEGIN APPENDED NINTH-SERIES V29 MATERIAL
% =============================================================================

\clearpage
\section{Ninth-series V29: closability of the metric cylinder gradient}
\label{sec:v9v29}

\subsection{Cylinder directions and logarithmic derivatives}
\label{sec:n9v29-gradient}

Let \(E\) be a separable field space with probability \(\mu\), and let
\({\cal C}\subset L^2(\mu)\) be a unital algebra of bounded smooth cylinder
functions, dense in \(L^2(\mu)\).  Let \(H\) be a separable Hilbert space
of admissible constraint-compatible metric variations, with orthonormal
basis \((e_k)_{k\geq1}\).  For \(F\in{\cal C}\), let \(D_kF\) be its
directional derivative along \(e_k\); only finitely many coordinates are
used by each cylinder function.  Assume the product and chain rules.

\begin{definition}[\(L^2\) Einstein logarithmic derivative card, ELDC]
\label{def:n9v29-ELDC}
ELDC consists of functions \(\beta_k\in L^2(\mu)\) such that
\begin{equation}
 \int_E D_kF\,d\mu=\int_EF\beta_k\,d\mu
 \quad(F\in{\cal C}, k\geq1).                      \label{eq:n9v29-IBP}
\end{equation}
For the flux-free QESDC convention of V28,
\begin{equation}
 \beta_k=\hbar^{-1}Q_{e_k}.                          \label{eq:n9v29-beta-score}
\end{equation}
\end{definition}

\begin{remark}[The score is already assembled]
\label{rem:n9v29-assembled}
By V28 score partition invariance, \(\beta_k\) contains the action,
reference-volume, determinant, coarea, chart, and constraint-lift
variations in their invariant sum.  ELDC is not a separate square-integrable
bound on each of those summands.
\end{remark}

\subsection{One-direction closability}
\label{sec:n9v29-one}

\begin{lemma}[Formal adjoint on the cylinder core]
\label{lem:n9v29-adjoint}
Under ELDC, for bounded \(F,G\in{\cal C}\),
\begin{equation}
 \int(D_kF)G\,d\mu
 =\int F\bigl(-D_kG+\beta_kG\bigr)\,d\mu.            \label{eq:n9v29-adjoint}
\end{equation}
Thus on its cylinder domain
\begin{equation}
 D_k^*G=-D_kG+\beta_kG.                              \label{eq:n9v29-adjoint-formula}
\end{equation}
\end{lemma}

\begin{proof}
Apply \eqref{eq:n9v29-IBP} to \(FG\) and use
\(D_k(FG)=(D_kF)G+F(D_kG)\).  Rearrangement gives
\eqref{eq:n9v29-adjoint}.  Since \(G\) and \(D_kG\) are bounded and
\(\beta_k\in L^2\), the displayed adjoint lies in \(L^2(\mu)\).
\end{proof}

\begin{theorem}[Closability of every cylinder direction]
\label{thm:n9v29-one-closable}
Under ELDC, the operator
\begin{equation}
 D_k:{\cal C}\subset L^2(\mu)\longrightarrow L^2(\mu)             \label{eq:n9v29-Dk}
\end{equation}
is closable for every \(k\).
\end{theorem}

\begin{proof}
Let \(F_n\to0\) in \(L^2(\mu)\) and
\(D_kF_n\to Y\) in \(L^2(\mu)\).  For every bounded cylinder \(G\),
\Cref{lem:n9v29-adjoint} gives
\[
 \int YG\,d\mu
 =\lim_n\int F_n(-D_kG+\beta_kG)\,d\mu=0.
\]
The cylinder core is dense, hence \(Y=0\).  This is the graph criterion for
closability.
\end{proof}

\subsection{The full gradient and its Dirichlet form}
\label{sec:n9v29-form}

Define on \({\cal C}\)
\begin{equation}
 DF:=\sum_{k\geq1}(D_kF)e_k,                          \label{eq:n9v29-D}
\end{equation}
where the sum is finite for each cylinder function.

\begin{theorem}[Closability of the Hilbert-valued gradient]
\label{thm:n9v29-gradient-closable}
Under ELDC,
\begin{equation}
 D:{\cal C}\subset L^2(\mu)\longrightarrow L^2(\mu;H)             \label{eq:n9v29-gradient-operator}
\end{equation}
is closable.
\end{theorem}

\begin{proof}
Suppose \(F_n\to0\) in \(L^2(\mu)\) and \(DF_n\to Y\) in
\(L^2(\mu;H)\).  Taking the \(e_k\)-component gives
\(D_kF_n\to\langle Y,e_k\rangle_H\) in \(L^2(\mu)\).  By
\Cref{thm:n9v29-one-closable}, every component vanishes.  Completeness of
the orthonormal basis gives \(Y=0\).
\end{proof}

Let \(\overline D\) be the closure and define
\begin{equation}
 \mathcal E(F,G)
 :=\int_E\langle\overline DF,\overline DG\rangle_H,d\mu,
 qquad F,G\in\operatorname{Dom}(\overline D).        \label{eq:n9v29-form}
\end{equation}

\begin{theorem}[Closed symmetric Markov form]
\label{thm:n9v29-Dirichlet}
The form \((\mathcal E,\operatorname{Dom}(\overline D))\) is densely
defined, closed, symmetric, nonnegative, and Markovian.  Therefore it has a
unique nonnegative self-adjoint operator \(L\) such that
\begin{equation}
 \mathcal E(F,G)=\langle L^{1/2}F,L^{1/2}G\rangle_{L^2(\mu)}.        \label{eq:n9v29-generator}
\end{equation}
For a cylinder function involving only the first \(m\) directions, its
formal generator is
\begin{equation}
 LF=\sum_{k=1}^mD_k^*D_kF
 =\sum_{k=1}^m\left(-D_k^2F+\beta_kD_kF\right),       \label{eq:n9v29-cylinder-generator}
\end{equation}
whenever the right side belongs to \(L^2(\mu)\).
\end{theorem}

\begin{proof}
Density follows from density of \({\cal C}\); closedness is the graph
closedness of \(\overline D\).  Symmetry and nonnegativity are immediate.
For a smooth normal contraction \(\eta:\mathbb R\to\mathbb R\) with
\(|\eta'|\leq1\), the cylinder chain rule gives
\[
 \mathcal E(\eta(F),\eta(F))
 =\int|\eta'(F)|^2|DF|_H^2d\mu
 \leq\mathcal E(F,F).
\]
Approximation by smooth normal contractions and closedness extend this to
the standard Markov contraction; hence the form is Dirichlet.  The closed
form representation theorem gives \(L\).  Finally use
\Cref{lem:n9v29-adjoint} componentwise to obtain
\eqref{eq:n9v29-cylinder-generator}.
\end{proof}

\begin{firewall}[No diffusion has yet been constructed]
\label{fire:n9v29-process}
A closed Dirichlet form has an \(L^2\) semigroup.  A field-valued Hunt
process requires quasi-regularity or an equivalent topological theorem,
which ELDC does not provide.  The form also acts in configuration space;
its positivity is not Osterwalder--Schrader reflection positivity in
spacetime.
\end{firewall}

\subsection{Weighted score vectors}
\label{sec:n9v29-weighted}

Individual \(L^2\) logarithmic derivatives suffice for closability.  A
genuine Hilbert-valued score needs a summability row.  For weights
\(w_k>0\), let \(H_w\) be the completion of the algebraic span with
\(\langle e_j,e_k\rangle_{H_w}=w_k\delta_{jk}\), so that
\(w_k^{-1/2}e_k\) is orthonormal and the weighted gradient is
\begin{equation}
 D^{(w)}F:=\sum_k w_k^{-1}(D_kF)e_k,qquad
 \|D^{(w)}F\|_{H_w}^2=\sum_kw_k^{-1}|D_kF|^2.         \label{eq:n9v29-weighted-gradient}
\end{equation}

\begin{proposition}[Hilbert-valued logarithmic derivative]
\label{prop:n9v29-score-vector}
If
\begin{equation}
 \sum_{k=1}^\infty w_k^{-1}\|\beta_k\|_{L^2(\mu)}^2<\infty,        \label{eq:n9v29-score-sum}
\end{equation}
then
\begin{equation}
 B:=\sum_{k=1}^\infty w_k^{-1}\beta_ke_k
 \quad\text{belongs to }L^2(\mu;H_w),               \label{eq:n9v29-score-vector}
\end{equation}
and \(\langle B,e_k\rangle_{H_w}=\beta_k\).  The weighted gradient is
closable and has the corresponding weighted Dirichlet form.
\end{proposition}

\begin{proof}
The squared \(L^2(\mu;H_w)\) norm of the partial sums is exactly the partial
sum in \eqref{eq:n9v29-score-sum}; hence they are Cauchy.  The component
identity follows from the definition of the \(H_w\) inner product.  The
closability proof of \Cref{thm:n9v29-gradient-closable} applies after the
same coordinate rescaling.
\end{proof}

\begin{firewall}[A strong tangent norm can make the score divergent]
\label{fire:n9v29-white}
In a continuum field theory the logarithmic derivative often has negative
Sobolev regularity.  Then \eqref{eq:n9v29-score-sum} fails for an unweighted
\(L^2\) tangent basis but can hold after choosing a sufficiently weak
weighted tangent space.  This choice changes the configuration-space
Dirichlet form and must be compatible with the geometric metric variations
used in QESDC.
\end{firewall}

\subsection{Extraction of a limiting score from finite cutoffs}
\label{sec:n9v29-limit}

Let \(E\) now be Polish.  Suppose \(\mu_n\Rightarrow\mu\), the common
cylinder algebra consists of bounded continuous functions, and
\(D_kF\) is bounded continuous whenever \(F\in{\cal C}\).

\begin{theorem}[Uniform \(L^2\) score compactness]
\label{thm:n9v29-score-compactness}
Fix \(k\).  Suppose finite-cutoff logarithmic derivatives satisfy
\begin{equation}
 \int D_kF\,d\mu_n=\int F\beta_{n,k}\,d\mu_n,qquad
 \sup_n\int|\beta_{n,k}|^2d\mu_n\leq C_k<\infty.     \label{eq:n9v29-finite-IBP}
\end{equation}
Then there is a unique \(\beta_k\in L^2(\mu)\), with
\(\|\beta_k\|_2^2\leq C_k\), such that
\begin{equation}
 \int D_kF\,d\mu=\int F\beta_k\,d\mu               \label{eq:n9v29-limit-IBP}
\end{equation}
for every \(F\in{\cal C}\).  After passage to a subsequence, the signed
score measures satisfy
\begin{equation}
 \beta_{n,k}\mu_n\Rightarrow\beta_k\mu.              \label{eq:n9v29-score-measures}
\end{equation}
If every subsequential score limit acts on a dense cylinder core, uniqueness
makes \eqref{eq:n9v29-score-measures} valid for the full sequence.
\end{theorem}

\begin{proof}
Put \(d\sigma_n=\beta_{n,k}d\mu_n\).  Cauchy--Schwarz gives, for every
Borel set \(A\),
\begin{equation}
 |\sigma_n|(A)\leq C_k^{1/2}\mu_n(A)^{1/2}.           \label{eq:n9v29-tight}
\end{equation}
Weak convergence of probabilities gives uniform tightness of \(\mu_n\),
so \eqref{eq:n9v29-tight} gives uniform tightness and bounded total
variation of \(\sigma_n\).  Extract a weakly convergent signed-measure
subsequence \(\sigma_n\Rightarrow\sigma\).  For bounded continuous \(F\),
\[
 \left|\int Fd\sigma_n\right|
 \leq C_k^{1/2}\left(\int F^2d\mu_n\right)^{1/2}.
\]
Passing to the limit shows that \(F\mapsto\int Fd\sigma\) extends to a
bounded functional on \(L^2(\mu)\).  Riesz representation gives
\(d\sigma=\beta_kd\mu\) with the stated norm bound.  Pass the left side of
\eqref{eq:n9v29-finite-IBP} by bounded continuity of \(D_kF\) to obtain
\eqref{eq:n9v29-limit-IBP}.  If two \(L^2\) functions satisfy that identity,
their difference pairs to zero with the dense cylinder core and is zero;
this proves uniqueness and the final subsequence assertion.
\end{proof}

\begin{corollary}[Countable QESDC gradient]
\label{cor:n9v29-countable}
If \eqref{eq:n9v29-finite-IBP} holds for every \(k\) with constants
\(C_k\), one diagonal subsequence supplies ELDC for all directions.  If in
addition
\begin{equation}
 \sum_kw_k^{-1}C_k<\infty,                            \label{eq:n9v29-uniform-weighted}
\end{equation}
then the limiting logarithmic derivative is \(H_w\)-valued and obeys
\eqref{eq:n9v29-score-sum}.
\end{corollary}

\begin{proof}
Use \Cref{thm:n9v29-score-compactness} successively and diagonalize over the
countable set of \(k\).  Lower semicontinuity of every finite partial sum,
followed by monotone convergence and
\eqref{eq:n9v29-uniform-weighted}, proves the weighted estimate.
\end{proof}

\begin{firewall}[Boundary flux must be represented or removed]
\label{fire:n9v29-boundary}
If the V28 boundary functional \(\mathfrak b_{e_k}(F)\) is nonzero, it need
not equal integration against an \(L^2(\mu)\) bulk function.  Then ELDC in
the present form does not follow.  One must impose a tangent/no-flux domain,
adjoin a boundary state space, or prove an \(L^2\) representation of the
flux.
\end{firewall}

\subsection{Metric-gradient closure card}
\label{sec:n9v29-card}

\begin{definition}[Metric score closability card, MSCC]
\label{def:n9v29-MSCC}
MSCC records:
\begin{enumerate}[label=\textup{(C\arabic*)},leftmargin=4em]
\item a common dense smooth cylinder algebra and countable Hilbert tangent
core;
\item constraint-compatible finite-cutoff derivatives satisfying the exact
V28 score identity;
\item vanishing flux or an explicitly represented boundary row;
\item the coordinate score bounds in \eqref{eq:n9v29-finite-IBP};
\item weak convergence of the transported probabilities and continuity of
cylinder derivatives; and
\item, when a score vector is required, the weighted sum
\eqref{eq:n9v29-uniform-weighted}.
\end{enumerate}
\end{definition}

\begin{theorem}[MSCC supplies QESDC closability]
\label{thm:n9v29-MSCC}
MSCC gives a limiting \(L^2\) logarithmic derivative in every declared
metric direction, a closable cylinder gradient, and the closed symmetric
Dirichlet form \eqref{eq:n9v29-form}.  Thus it discharges QESDC item (Q6)
and passes its integration-by-parts identity on the common cylinder core.
\end{theorem}

\begin{proof}
Use \Cref{thm:n9v29-score-compactness,cor:n9v29-countable} to obtain ELDC,
then \Cref{thm:n9v29-gradient-closable,thm:n9v29-Dirichlet}.
\end{proof}

\begin{assessment}[Progress at ninth-series V29]
\label{ass:n9v29-status}
The unbounded metric derivative now has a precise closure mechanism.  A
uniform \(L^2\) bound on each complete finite-cutoff Einstein score yields
the limiting logarithmic derivative and closable cylinder gradient; a
weighted square sum yields a Hilbert-valued score.  The physical work has
therefore been reduced to complete-score estimates and boundary-flux
control, rather than formal differentiation of individual density rows.
\end{assessment}

\begin{programme}[Ninth-series V30: compulsory companion audit]
\label{prog:n9v29-V30}
V30 will inspect the newest YM and NS notes before selecting the first
physical loader for the MSCC score bounds.  It will compare their current
conditional/polymer or kernel estimates only at type-compatible interfaces,
then continue to V31 in the selected direction.
\end{programme}

\begin{firewall}[Claim boundary after ninth-series V29]
\label{fire:n9v29-claim}
V29 does not prove the uniform score bounds for the physical SM--Einstein
measure, eliminate boundary flux, establish quasi-regularity, reflection
positivity, reconstruction, conformal stability, or the full continuum.
\end{firewall}

\begin{theorem}[Ninth-series V29 result]
\label{thm:n9v29-result}
ELDC makes every metric cylinder derivative and the full Hilbert-valued
gradient closable and produces a closed symmetric Dirichlet form.  Uniform
finite-cutoff \(L^2\) score bounds yield ELDC after weak convergence, and
the weighted estimate \eqref{eq:n9v29-uniform-weighted} yields an
\(H_w\)-valued limiting Einstein score.
\end{theorem}

\begin{proof}
Use \Cref{thm:n9v29-one-closable,thm:n9v29-gradient-closable,
thm:n9v29-Dirichlet,prop:n9v29-score-vector,
thm:n9v29-score-compactness,thm:n9v29-MSCC}.
\end{proof}

\paragraph{Parent-version record.}
V29 was formed from
\texttt{Renewal\_SM\_Einstein\_Continuum\_Research\_Notes\_V28.tex}.
The V28 parent had \(9\,694\) logical source lines, \(345\,361\) bytes, and
SHA--256
\[
 \texttt{55D81CEA593A216C5A34F835C58655853275C038D6CD1E4E604532CB9A122173}.
\]

% =============================================================================
% END APPENDED NINTH-SERIES V29 MATERIAL
% =============================================================================


% =============================================================================
% BEGIN APPENDED NINTH-SERIES V30 MATERIAL
% =============================================================================

\clearpage
\section{Ninth-series V30: companion audit and regular/singular score gluing}
\label{sec:v9v30}

\subsection{Compulsory file audit}
\label{sec:n9v30-files}

The newest companion files were inspected read-only.  Their exact records
are
\begin{center}
\small
\begin{tabular}{@{}p{0.47\textwidth}r r@{}}
\toprule
file & source lines & bytes\\
\midrule
\texttt{Renewal\_Yang\_Mills\_Mass\_Gap\_Research\_Notes\_V68.tex}
 & \(29\,958\) & \(1\,072\,290\)\\
\texttt{Renewal\_Yang\_Mills\_Mass\_Gap\_Research\_Notes\_V69.tex}
 & \(29\,958\) & \(1\,072\,290\)\\
\texttt{Renewal\_NS\_Stability\_Research\_Notes\_V26.tex}
 & \(5\,319\) & \(278\,283\)\\
\bottomrule
\end{tabular}
\end{center}
The two YM files are byte-identical, with SHA--256
\begin{equation}
 \texttt{C05180F2C249C152BC49137AEA4A06E7FE18D59D5C4D68C880EB8C88948BE8FD}.
                                                               \label{eq:n9v30-YMhash}
\end{equation}
The V69 file ends with V68 material and contains no substantive V69
appendix, so V68 is the newest completed YM result.  The NS hash remains
\begin{equation}
 \texttt{82C887F8C2C69E4F28FAD7C3DC8DE5B9099CB5A4BF431EBA1FFF3E91935978AD}.
                                                               \label{eq:n9v30-NShash}
\end{equation}
No companion file is modified or removed.

\subsection{Type-compatible companion conclusions}
\label{sec:n9v30-digest}

YM V67 proves the flat Wilson Hessian is a covariant cochain square,
identifies its physical kernel with twisted first cohomology, and replaces
an all-label uniform barrier by regular/singular conditional domination.
YM V68 proves a bounded-incidence owner compiler, exhibits an absolute
energy-repair route to product domination, and gives a decisive type
counterexample: a uniform Dirichlet-form gap does not control the
probability of a specified event.  It leaves the physical Wilson repair or
reflection-positive product estimate open.

NS V26 is unchanged.  Its pointwise rank-one Oseen derivative norms sharpen
a deterministic kernel constant but do not provide a probability event,
Einstein logarithmic derivative, Gibbs integration-by-parts score, or
metric tangent form.  No NS estimate transfers to MSCC at this audit.

\begin{proposition}[V30 audit decision]
\label{prop:n9v30-decision}
The SM programme retains V26 absolute-energy tails for local metric exits
and V29 \(L^2\) score bounds as distinct inputs.  The former cannot be
deduced from a Dirichlet-form gap, and the latter cannot be deduced from
event rarity alone.  They can be joined by a higher-moment estimate on the
singular sector.  V31 will seek the regular score estimate through an
integration-by-parts Hessian identity, not through a spectral-gap-to-event
inference.
\end{proposition}

\begin{proof}
The first non-implication is the Bernoulli projection counterexample proved
in YM V68.  The converse and the valid gluing estimate are proved below.
The NS objects have different domains and codomains from the logarithmic
metric score in V29.
\end{proof}

\subsection{Rare events do not bound a score}
\label{sec:n9v30-no-go}

\begin{proposition}[Exit-rarity no-go]
\label{prop:n9v30-rarity-no-score}
There are probabilities \(\mu_n\), events \(E_n\), and centered random
variables \(\beta_n\) such that
\begin{equation}
 \mu_n(E_n)=p_n\longrightarrow0,qquad
 \int\beta_n,d\mu_n=0,qquad
 \int|\beta_n|^2d\mu_n\longrightarrow\infty.         \label{eq:n9v30-no-go}
\end{equation}
\end{proposition}

\begin{proof}
On the two-point space with \(\mu_n(E_n)=p_n\), put
\begin{equation}
 \beta_n=p_n^{-1}{\bf1}_{E_n}
 -(1-p_n)^{-1}{\bf1}_{E_n^c}.                        \label{eq:n9v30-example}
\end{equation}
Its mean is zero, while
\[
 \int\beta_n^2d\mu_n=p_n^{-1}+(1-p_n)^{-1}	o\infty.
\]
\end{proof}

\begin{firewall}[Product domination and Fisher information are different]
\label{fire:n9v30-type}
Even product domination of all intersections of bad blocks controls only
their indicator algebra.  An Einstein score can grow arbitrarily rapidly
on those events.  Conversely, an \(L^2\) score or a Dirichlet gap does not
identify a chosen event as energetically costly.  Neither row may be used
as a substitute for the other.
\end{firewall}

\subsection{The valid two-norm gluing estimate}
\label{sec:n9v30-gluing}

\begin{lemma}[Regular/singular \(L^2\) gluing]
\label{lem:n9v30-gluing}
Let \(\beta\) be measurable, let \(G\) be an event, and let
\(\delta>0\).  If
\begin{align}
 \int_G|\beta|^2d\mu&\leq A,                         \label{eq:n9v30-regular}\\
 \int|\beta|^{2+\delta}d\mu&\leq M,                 \label{eq:n9v30-high}\\
 \mu(G^c)&\leq p,                                   \label{eq:n9v30-rare}
\end{align}
then
\begin{equation}
 \int|\beta|^2d\mu
 \leq A+M^{2/(2+\delta)}p^{\delta/(2+\delta)}.        \label{eq:n9v30-gluing}
\end{equation}
\end{lemma}

\begin{proof}
Split the integral over \(G\) and \(G^c\).  Holder's inequality with
exponents \((2+\delta)/2\) and \((2+\delta)/\delta\) gives
\[
 \int_{G^c}|\beta|^2d\mu
 \leq
 \left(\int|\beta|^{2+\delta}d\mu\right)^{2/(2+\delta)}
 \mu(G^c)^{\delta/(2+\delta)}.
\]
Insert \eqref{eq:n9v30-high}--\eqref{eq:n9v30-rare}.
\end{proof}

\begin{corollary}[Vanishing singular score]
\label{cor:n9v30-vanishing}
If \(M\) is uniform and \(p_n\to0\), then
\begin{equation}
 \|\beta_n{\bf1}_{G_n^c}\|_{L^2(\mu_n)}\longrightarrow0.           \label{eq:n9v30-vanishing}
\end{equation}
Thus a uniform regular-sector \(L^2\) bound supplies the full MSCC
coordinate bound.
\end{corollary}

\begin{proof}
Use the singular-sector term in the proof of
\Cref{lem:n9v30-gluing} and let \(p_n\to0\).
\end{proof}

\begin{theorem}[Weighted regular/singular MSCC loader]
\label{thm:n9v30-weighted}
For each metric direction \(e_k\), let \(G_{n,k}\) be a local regular
event and \(\beta_{n,k}\) the complete Einstein logarithmic derivative.
Suppose for a common \(\delta>0\),
\begin{align}
 \int_{G_{n,k}}|\beta_{n,k}|^2d\mu_n&\leq A_k,        \label{eq:n9v30-Ak}\\
 \int|\beta_{n,k}|^{2+\delta}d\mu_n&\leq M_k,        \label{eq:n9v30-Mk}\\
 \mu_n(G_{n,k}^c)&\leq p_{n,k}.                      \label{eq:n9v30-pnk}
\end{align}
If
\begin{equation}
 \sum_kw_k^{-1}left[
 A_k+M_k^{2/(2+\delta)}sup_np_{n,k}^{\delta/(2+\delta)}
 \right]<\infty,                                    \label{eq:n9v30-weighted}
\end{equation}
then the V29 weighted finite-cutoff score bound
\eqref{eq:n9v29-uniform-weighted} holds.  If moreover
\(p_{n,k}\to0\) for every \(k\), every limiting coordinate score is the
\(L^2\) limit of its regular-sector truncations along any coupling where
the regular parts converge.
\end{theorem}

\begin{proof}
Apply \Cref{lem:n9v30-gluing} for every \(n,k\), multiply by
\(w_k^{-1}\), and sum.  The asserted full bound is
\eqref{eq:n9v30-weighted}.  Coordinatewise vanishing of the singular part
is \Cref{cor:n9v30-vanishing}; adding convergence of the regular part gives
the last statement by the triangle inequality.
\end{proof}

\begin{firewall}[The higher moment is an analytic obligation]
\label{fire:n9v30-high}
The \((2+\delta)\)-moment in \eqref{eq:n9v30-Mk} cannot be replaced by the
same \(L^2\) norm one is trying to prove.  It must come from coercivity,
an exponential score moment, a Hessian identity with a higher-order bound,
or a separate truncation estimate.  Otherwise the gluing argument is
circular.
\end{firewall}

\subsection{Conditional regular and singular labels}
\label{sec:n9v30-labels}

The YM regular/singular mechanism has the following type-correct use here.
For a metric direction supported in block \(b\), let \(\Sigma_b\) be the
event that the boundary label fails the hypotheses of a regular score
calculation.  On \(\Sigma_b^c\), prove
\eqref{eq:n9v30-Ak}.  Use V26 energy repair or another normalized density
argument to prove product domination for \(\Sigma_b\); this supplies
\eqref{eq:n9v30-pnk}.  Finally prove the high score moment
\eqref{eq:n9v30-Mk}.  Finite coloring or a bounded-incidence owner map can
then compile multi-block singular labels without changing the one-score
Holder estimate.

\begin{proposition}[No virtual-minimum substitution]
\label{prop:n9v30-virtual}
A score bound evaluated at a boundary-dependent conditional minimizer does
not imply \eqref{eq:n9v30-Ak} for the sampled field unless a deterministic
or probabilistic comparison between the minimizer and the sampled interior
is proved.
\end{proposition}

\begin{proof}
The conditional minimizer and sampled interior are different points of the
block configuration space.  A nonconstant score may lie below a threshold
at one and above it at the other.  No event inclusion follows without a
comparison map.  This is the same logical distinction identified for the
Wilson certificate in YM V68, now applied to the Einstein score.
\end{proof}

\subsection{Regular/singular metric score card}
\label{sec:n9v30-card}

\begin{definition}[Regular/singular Einstein score card, RSESC]
\label{def:n9v30-RSESC}
RSESC records:
\begin{enumerate}[label=\textup{(R\arabic*)},leftmargin=4em]
\item the complete split-invariant score \(\beta_{n,k}\) of V28;
\item a sampled-field regular event \(G_{n,k}\), not merely a virtual
conditional-minimum certificate;
\item the regular \(L^2\) estimate \eqref{eq:n9v30-Ak};
\item a normalized probability or repair proof of
\eqref{eq:n9v30-pnk};
\item the noncircular higher moment \eqref{eq:n9v30-Mk}; and
\item the weighted summability \eqref{eq:n9v30-weighted} on the countable
metric tangent core.
\end{enumerate}
\end{definition}

\begin{theorem}[RSESC implies MSCC's analytic score row]
\label{thm:n9v30-RSESC}
RSESC gives the uniform weighted \(L^2\) score estimate required by V29.
Together with the remaining algebra, flux, and convergence rows of MSCC,
it yields a closable limiting metric gradient and the fixed-\(\hbar\)
Einstein Schwinger--Dyson identity.
\end{theorem}

\begin{proof}
The analytic implication is \Cref{thm:n9v30-weighted}.  Then apply
\Cref{thm:n9v29-MSCC,thm:n9v28-QESDC}.
\end{proof}

\begin{assessment}[Progress at ninth-series V30]
\label{ass:n9v30-status}
The compulsory audit has kept event probability, Dirichlet form, and
Einstein-score integrability in their correct spaces.  V30 proves the
exact bridge: rare singular labels become harmless for an \(L^2\) score
only with a uniform higher score moment.  The next acquisition is the
regular-sector score estimate and, if possible, its higher-moment upgrade.
\end{assessment}

\begin{programme}[Ninth-series V31: Fisher--Hessian score identity]
\label{prog:n9v30-V31}
V31 will apply the exact V28 integration-by-parts formula to the score
itself.  Under a justified core approximation it will derive
\(\mathbb E\beta_h^2=\mathbb E D_h\beta_h\), identify the complete second
metric variation and divergence correction, and state the weighted trace
condition that loads the V29--V30 score bounds.
\end{programme}

\begin{firewall}[Claim boundary after ninth-series V30]
\label{fire:n9v30-boundary}
V30 does not prove RSESC for the physical measure, a physical Einstein
score moment, a Wilson repair, Navier--Stokes regularity, reflection
positivity, conformal stability, reconstruction, or the full continuum.
\end{firewall}

\begin{theorem}[Ninth-series V30 result]
\label{thm:n9v30-result}
The companion audit is complete.  Event rarity and score integrability are
logically independent, while
\eqref{eq:n9v30-gluing} is a valid noncircular bridge under a
\((2+\delta)\)-moment.  The weighted version
\eqref{eq:n9v30-weighted} supplies the analytic MSCC score row and fixes the
V31 acquisition target.
\end{theorem}

\begin{proof}
Use \Cref{prop:n9v30-decision,prop:n9v30-rarity-no-score,
lem:n9v30-gluing,thm:n9v30-weighted,thm:n9v30-RSESC}.
\end{proof}

\paragraph{Parent-version record.}
V30 was formed from
\texttt{Renewal\_SM\_Einstein\_Continuum\_Research\_Notes\_V29.tex}.
The V29 parent had \(10\,066\) logical source lines, \(360\,500\) bytes,
and SHA--256
\[
 \texttt{CAB3DB39F6498AE274619391763FC5F64D47C7A8E3C67F5BDF5BED9FD9D87786}.
\]
V31 follows this audit before the next five-version checkpoint.

% =============================================================================
% END APPENDED NINTH-SERIES V30 MATERIAL
% =============================================================================


% =============================================================================
% BEGIN APPENDED NINTH-SERIES V31 MATERIAL
% =============================================================================

\clearpage
\section{Ninth-series V31: Fisher--Hessian identities for the Einstein score}
\label{sec:v9v31}

\subsection{Using the logarithmic derivative as a test}
\label{sec:n9v31-Fisher}

Fix one finite regulator, one constraint-compatible metric vector field
\(V\), and write
\begin{equation}
 D_VF:=VF,qquad
 \beta_V:=\hbar^{-1}Q_V
 =\hbar^{-1}VS-\div_mV.                              \label{eq:n9v31-beta}
\end{equation}
Assume the V28 boundary flux vanishes on the functions used below, so
\begin{equation}
 \int D_VF\,d\mu=\int F\beta_V\,d\mu.                \label{eq:n9v31-IBP}
\end{equation}

\begin{lemma}[Graph-core extension]
\label{lem:n9v31-core}
Suppose \(\beta_V\in L^2(\mu)\) and there are cylinder functions \(F_j\)
such that
\begin{equation}
 F_j\to\beta_V,qquad D_VF_j\to D_V\beta_V
 \quad\text{in }L^2(\mu).                           \label{eq:n9v31-graph}
\end{equation}
Then \eqref{eq:n9v31-IBP} extends to \(F=\beta_V\).
\end{lemma}

\begin{proof}
The right sides converge because
\[
 \left|\int(F_j-\beta_V)\beta_Vd\mu\right|
 \leq\|F_j-\beta_V\|_2\|\beta_V\|_2.
\]
The left sides converge because the constant function belongs to
\(L^2(\mu)\) and \(D_VF_j\to D_V\beta_V\) in \(L^2\).
\end{proof}

\begin{theorem}[Noncircular Fisher--Hessian acquisition]
\label{thm:n9v31-Fisher}
Assume \(D_V\beta_V\in L^1(\mu)\).  Suppose the bounded Lipschitz
truncations
\[
 T_R(\beta_V):=\max\{-R,\min\{\beta_V,R\}\}
\]
are admissible in \eqref{eq:n9v31-IBP}, by smooth approximation if
necessary, and obey the Sobolev chain rule.  Then
\(\beta_V\in L^2(\mu)\) and
\begin{equation}
 \boxed{
 \int\beta_V^2d\mu
 =\int D_V\beta_Vd\mu
 ={1\over\hbar}\int V(VS)d\mu
  -\int V(\div_mV)d\mu.}                            \label{eq:n9v31-Fisher}
\end{equation}
The expression \(V(VS)\) is the second derivative along the flow of
\(V\); it includes the derivative of \(V\) when the constraint lift is
field dependent.
\end{theorem}

\begin{proof}
Put \(X_R:=\int\beta_VT_R(\beta_V)d\mu\).  The integrand is nonnegative
and increases pointwise to \(\beta_V^2\).  Integration by parts and the
chain rule give
\[
 X_R=\int{\bf1}_{\{|\beta_V|<R\}}D_V\beta_V\,d\mu,
\]
with the level set handled by smooth truncation.  Hence
\(0\leq X_R\leq\int|D_V\beta_V|d\mu\), uniformly in \(R\).
Monotone convergence proves \(\beta_V\in L^2\).  Dominated convergence on
the right then proves
\(\int\beta_V^2d\mu=\int D_V\beta_Vd\mu\).
Finally differentiate \eqref{eq:n9v31-beta} along \(V\).
The graph-core argument of \Cref{lem:n9v31-core} gives the same identity
when that stronger approximation is already available.
\end{proof}

\begin{corollary}[Flat constant-direction form]
\label{cor:n9v31-flat}
If \(X=\mathbb R^N\), \(m=dx\), and \(V\) is constant, then
\begin{equation}
 {1\over\hbar^2}\int(VS)^2d\mu
 ={1\over\hbar}\int V^2S\,d\mu.                    \label{eq:n9v31-flat}
\end{equation}
\end{corollary}

\begin{proof}
Here \(\div_mV=0\) and \(\beta_V=\hbar^{-1}VS\).  Insert these facts in
\eqref{eq:n9v31-Fisher}.
\end{proof}

\begin{firewall}[A lower Hessian bound is not the required upper estimate]
\label{fire:n9v31-lower}
Coercivity or a spectral gap gives a lower bound on a quadratic form.  The
right side of \eqref{eq:n9v31-Fisher} requires an integrable upper bound on
the complete second metric variation and divergence correction.  A lower
bound alone cannot supply the MSCC score norm.
\end{firewall}

\subsection{Cross-direction Fisher matrix}
\label{sec:n9v31-matrix}

Let \(V_1,\ldots,V_m\) be declared metric directions with scores
\(\beta_1,\ldots,\beta_m\).  Assume the graph-core hypothesis for their
linear combinations.

\begin{proposition}[Fisher matrix identity]
\label{prop:n9v31-matrix}
For all \(j,k\),
\begin{equation}
 \int\beta_j\beta_kd\mu
 =\int D_{V_j}\beta_kd\mu.                           \label{eq:n9v31-matrix}
\end{equation}
Consequently the matrix on the right is symmetric and positive
semidefinite after integration, even when the pointwise derivative matrix
is neither.
\end{proposition}

\begin{proof}
Apply the integration-by-parts identity for direction \(V_j\) with
\(F=\beta_k\).  Symmetry and positivity follow because the left side is
the Gram matrix of the score variables in \(L^2(\mu)\).
\end{proof}

\begin{corollary}[Finite weighted trace identity]
\label{cor:n9v31-trace}
For positive weights \(w_k\) and every \(m<\infty\),
\begin{equation}
 \sum_{k=1}^mw_k^{-1}\|\beta_k\|_2^2
 =\sum_{k=1}^mw_k^{-1}
 \int\left[
 \hbar^{-1}V_k(V_kS)-V_k(\div_mV_k)
 \right]d\mu.                                       \label{eq:n9v31-trace}
\end{equation}
\end{corollary}

\begin{proof}
Sum \eqref{eq:n9v31-Fisher} over the finite family.
\end{proof}

\subsection{Higher score moments from differentiated scores}
\label{sec:n9v31-higher}

The V30 regular/singular bridge needs an exponent above two.  That exponent
can also be obtained by integration by parts, provided the differentiated
score has the corresponding moment.

\begin{theorem}[Noncircular higher Fisher identity and bound]
\label{thm:n9v31-higher}
Let \(r>2\), assume
\begin{equation}
 D_V\beta_V\in L^{r/2}(\mu),                         \label{eq:n9v31-Dbeta}
\end{equation}
and suppose bounded Lipschitz functions of \(\beta_V\) are admissible in
\eqref{eq:n9v31-IBP} with the Sobolev chain rule.  Then
\(\beta_V\in L^r(\mu)\) and
\begin{align}
 \int|\beta_V|^r d\mu
 &=(r-1)\int|\beta_V|^{r-2}D_V\beta_V\,d\mu,          \label{eq:n9v31-higher-identity}\\
 \int|\beta_V|^r d\mu
 &\leq(r-1)^{r/2}
       \int|D_V\beta_V|^{r/2}d\mu.                  \label{eq:n9v31-higher-bound}
\end{align}
\end{theorem}

\begin{proof}
Let
\[
 \theta_R(t):=\sgn(t)\min\{|t|,R\}^{r-1},
 \qquad
 X_R:=\int\beta_V\theta_R(\beta_V)d\mu.
\]
Use smooth approximations at the two corners.  The integrand defining
\(X_R\) is nonnegative and increases to \(|\beta_V|^r\), while
\[
 |\theta_R'(t)|
 \leq(r-1){\bf1}_{\{|t|<R\}}|t|^{r-2}.
\]
The integration-by-parts and chain-rule identities, followed by
Holder's inequality, give
\begin{align*}
 X_R
 &\leq(r-1)\int_{\{|\beta_V|<R\}}
       |\beta_V|^{r-2}|D_V\beta_V|d\mu\\
 &\leq(r-1)X_R^{(r-2)/r}
       \left(\int|D_V\beta_V|^{r/2}d\mu\right)^{2/r}.
\end{align*}
Thus \(X_R\leq(r-1)^{r/2}
\int|D_V\beta_V|^{r/2}d\mu\), uniformly in \(R\).
Monotone convergence proves \(\beta_V\in L^r\) and
\eqref{eq:n9v31-higher-bound}.  Now
\(|\beta_V|^{r-2}|D_V\beta_V|\in L^1\) by Holder.
Dominated convergence in the differentiated truncation identity gives
\eqref{eq:n9v31-higher-identity}.
\end{proof}

\begin{corollary}[The V30 moment with \(r=2+\delta\)]
\label{cor:n9v31-V30}
For \(r=2+\delta\), a uniform bound
\begin{equation}
 \int|D_{V_k}\beta_k|^{1+\delta/2}d\mu_n\leq Y_k     \label{eq:n9v31-Yk}
\end{equation}
implies the V30 score moment with
\begin{equation}
 M_k=(1+\delta)^{1+\delta/2}Y_k.                     \label{eq:n9v31-Mk}
\end{equation}
\end{corollary}

\begin{proof}
Insert \(r=2+\delta\) in \eqref{eq:n9v31-higher-bound}.
\end{proof}

\begin{firewall}[Differentiating an incomplete score is invalid]
\label{fire:n9v31-incomplete}
The derivative \(D_V\beta_V\) must be formed after the action,
reference-volume, determinant, coarea, constraint-lift, BV, and counterterm
rows have been assembled into the split-invariant V28 score.  Bounds on
the absolute values of the separately differentiated rows can destroy
cancellations and do not prove \eqref{eq:n9v31-Yk}.
\end{firewall}

\subsection{Invariance under action--measure repartition}
\label{sec:n9v31-invariance}

\begin{proposition}[The complete Hessian score is split invariant]
\label{prop:n9v31-invariance}
Under the V28 repartition
\(S'=S+A\), \(m'=e^{A/\hbar}m\),
\begin{equation}
 \beta'_V=\beta_V,qquad
 D_V\beta'_V=D_V\beta_V.                            \label{eq:n9v31-invariance}
\end{equation}
In expanded form, the apparent term \(\hbar^{-1}V(VA)\) in the action
Hessian is cancelled exactly by the same term in
\(V(\div_{m'}V)\).
\end{proposition}

\begin{proof}
The first equality is V28 score partition invariance divided by
\(\hbar\).  Apply \(D_V\) to it for the second.  Alternatively, use
\(\div_{m'}V=\div_mV+\hbar^{-1}VA\) and differentiate once more.
\end{proof}

\subsection{Boundary correction}
\label{sec:n9v31-boundary}

If the flux does not vanish, V28 gives
\begin{equation}
 \int D_VF\,d\mu
 =\int F\beta_Vd\mu+\mathfrak b_V(F).                \label{eq:n9v31-boundary-IBP}
\end{equation}
Therefore the Fisher identity becomes
\begin{equation}
 \int\beta_V^2d\mu
 =\int D_V\beta_Vd\mu-\mathfrak b_V(\beta_V),        \label{eq:n9v31-boundary-Fisher}
\end{equation}
provided the boundary evaluation is defined.

\begin{firewall}[The boundary score cannot be dropped by sign]
\label{fire:n9v31-boundary}
The flux \(\mathfrak b_V(\beta_V)\) has no automatic favorable sign and
may not be represented by a bulk \(L^2\) function.  A physical Fisher
estimate must use a tangent/no-flux metric domain, prove a quantitative
trace bound, or enlarge the state space to include the boundary port.
\end{firewall}

\subsection{A finite-cutoff Hessian loader}
\label{sec:n9v31-loader}

For direction \(V_{n,k}\), define the complete differentiated score
\begin{equation}
 {cal H}_{n,k}
 :=D_{V_{n,k}}\beta_{n,k}
 ={1\over\hbar}V_{n,k}(V_{n,k}S_n)
  -V_{n,k}(\div_{m_n}V_{n,k}).                       \label{eq:n9v31-Hnk}
\end{equation}

\begin{definition}[Complete Hessian moment card, CHMC]
\label{def:n9v31-CHMC}
CHMC with exponent \(r>2\) records:
\begin{enumerate}[label=\textup{(H\arabic*)},leftmargin=4em]
\item the exact split-invariant score and its complete derivative
\eqref{eq:n9v31-Hnk};
\item a cylinder/Sobolev core justifying bounded score truncations and
their chain rules, without assuming the target score moment;
\item vanishing flux or the explicit correction
\eqref{eq:n9v31-boundary-Fisher};
\item bounds
\begin{equation}
 \int|{\cal H}_{n,k}|^{r/2}d\mu_n\leq Y_k;           \label{eq:n9v31-CHMC-bound}
\end{equation}
\item the V30 regular/singular event and probability rows; and
\item weighted summability after inserting
\(M_k=(r-1)^{r/2}Y_k\) in
\eqref{eq:n9v30-weighted}.
\end{enumerate}
\end{definition}

\begin{theorem}[CHMC loads the V30 and V29 analytic cards]
\label{thm:n9v31-CHMC}
CHMC implies the \(L^r\) score moment required by RSESC and the weighted
\(L^2\) score estimate required by MSCC.  With the remaining convergence
and flux rows, the limiting metric gradient is closable and satisfies the
fixed-\(\hbar\) Einstein Schwinger--Dyson identity.
\end{theorem}

\begin{proof}
Use \Cref{thm:n9v31-higher} to obtain
\(M_k=(r-1)^{r/2}Y_k\), then
\Cref{thm:n9v30-weighted,thm:n9v30-RSESC,
thm:n9v29-MSCC,thm:n9v28-QESDC}.
\end{proof}

\begin{firewall}[CHMC is an estimate, not an identity alone]
\label{fire:n9v31-physical}
The Fisher formula is exact once integration by parts is justified, but it
does not bound \({\cal H}_{n,k}\).  The physical programme still needs a
uniform moment estimate for the complete second metric variation.  In
particular, the Euclidean conformal direction, principal-symbol metric
variations, and interacting BV covariance can make that estimate fail.
\end{firewall}

\begin{assessment}[Progress at ninth-series V31]
\label{ass:n9v31-status}
The score norm and the V30 higher moment have both been reduced to a
type-correct object: the complete differentiated Einstein score.  The
identities are invariant under action--measure repartition and expose the
boundary correction.  The remaining analytic gate is now the weighted
\(L^{1+\delta/2}\) moment of \({\cal H}_{n,k}\), not an unjustified
spectral-gap or rare-event inference.
\end{assessment}

\begin{programme}[Ninth-series V32: Gaussian principal-score model]
\label{prog:n9v31-V32}
V32 will compute CHMC for a finite Gaussian metric block, diagonalize the
weighted trace, and determine the exact ultraviolet summability condition
on the tangent weights.  This solvable model will show which part of the
physical score estimate is mere dimension counting and which part belongs
to interactions and the conformal sign.
\end{programme}

\begin{firewall}[Claim boundary after ninth-series V31]
\label{fire:n9v31-claim}
V31 does not prove CHMC for the physical SM--Einstein regulator, control
the conformal score, close boundary flux, construct the continuum measure,
prove reflection positivity or reconstruction, or complete the continuum.
\end{firewall}

\begin{theorem}[Ninth-series V31 result]
\label{thm:n9v31-result}
For an admissible complete Einstein logarithmic derivative, the Fisher
identity \eqref{eq:n9v31-Fisher} gives its \(L^2\) norm from the complete
second metric variation.  The higher identity
\eqref{eq:n9v31-higher-identity} gives the V30 \(2+\delta\) moment from
\eqref{eq:n9v31-Yk}.  CHMC therefore loads the analytic score rows of
RSESC and MSCC, subject to its explicit core and flux conditions.
\end{theorem}

\begin{proof}
Use \Cref{thm:n9v31-Fisher,prop:n9v31-matrix,
thm:n9v31-higher,prop:n9v31-invariance,thm:n9v31-CHMC}.
\end{proof}

\paragraph{Parent-version record.}
V31 was formed from
\texttt{Renewal\_SM\_Einstein\_Continuum\_Research\_Notes\_V30.tex}.
The V30 parent had \(10\,380\) logical source lines, \(372\,979\) bytes,
and SHA--256
\[
 \texttt{03CF2AD04BC386984C5C28175832FEC0651082EE2328ADDA78646BF460CFB910}.
\]

% =============================================================================
% END APPENDED NINTH-SERIES V31 MATERIAL
% =============================================================================


% =============================================================================
% BEGIN APPENDED NINTH-SERIES V32 MATERIAL
% =============================================================================

\clearpage
\section{Ninth-series V32: the Gaussian Einstein-score trace}
\label{sec:v9v32}

\subsection{Finite positive Gaussian block}
\label{sec:n9v32-finite}

Let \(A\) be a positive definite symmetric operator on \(\mathbb R^N\),
with orthonormal eigenbasis \(e_k\) and eigenvalues \(a_k>0\).  Put
\begin{equation}
 S_0(x)={1\over2}\langle x,Ax\rangle,qquad
 d\mu_A(x)=Z_A^{-1}e^{-S_0(x)/\hbar}dx.              \label{eq:n9v32-Gaussian}
\end{equation}
For the constant direction \(e_k\), the logarithmic derivative is
\begin{equation}
 \beta_k(x)={a_kx_k\over\hbar},qquad
 D_k\beta_k={a_k\over\hbar}.                        \label{eq:n9v32-score}
\end{equation}

\begin{theorem}[Exact Gaussian score moments]
\label{thm:n9v32-moments}
For every \(r>0\),
\begin{equation}
 \int|\beta_k|^r d\mu_A
 =m_r\left({a_k\over\hbar}\right)^{r/2},qquad
 m_r:={2^{r/2}\Gamma((r+1)/2)\over\sqrt\pi}.         \label{eq:n9v32-moments}
\end{equation}
In particular,
\begin{equation}
 \|\beta_k\|_2^2={a_k\over\hbar}
 =\int D_k\beta_kd\mu_A,                             \label{eq:n9v32-Fisher}
\end{equation}
and the V31 higher Fisher identity is exact because
\begin{equation}
 m_r=(r-1)m_{r-2}.                                   \label{eq:n9v32-recurrence}
\end{equation}
\end{theorem}

\begin{proof}
The coordinates are independent centered Gaussians with
\(x_k\sim N(0,\hbar/a_k)\).  Hence
\(\beta_k\sim N(0,a_k/\hbar)\).  The absolute moment of a standard normal
is \(m_r\), which proves \eqref{eq:n9v32-moments}.  The case \(r=2\)
gives \eqref{eq:n9v32-Fisher}.  The gamma recurrence
\(\Gamma((r+1)/2)=((r-1)/2)\Gamma((r-1)/2)\) gives
\eqref{eq:n9v32-recurrence}; substitution in
\eqref{eq:n9v31-higher-identity} verifies the higher identity.
\end{proof}

\begin{corollary}[The CHMC bound is dimension free by direction]
\label{cor:n9v32-CHMC}
For each fixed direction, CHMC holds for every finite \(r>2\), with
\begin{equation}
 Y_k=left({a_k\over\hbar}\right)^{r/2},qquad
 M_k=m_r\left({a_k\over\hbar}\right)^{r/2}.          \label{eq:n9v32-YM}
\end{equation}
The dependence on the number of modes enters only when the direction
moments are summed.
\end{corollary}

\begin{proof}
The differentiated score is constant by \eqref{eq:n9v32-score}, and the
score moment is \eqref{eq:n9v32-moments}.
\end{proof}

\subsection{Weighted continuum trace}
\label{sec:n9v32-trace}

Let \(P\) be a positive elliptic operator of order \(2\sigma>0\) on a
closed \(d\)-dimensional manifold, after its finite-dimensional kernel has
been quotiented or retained separately.  Let
\(\lambda_k\) be eigenvalues of a fixed positive Laplace-type operator and
assume the two-sided principal comparison
\begin{equation}
 c(1+\lambda_k)^\sigma
 \leq a_k\leq C(1+\lambda_k)^\sigma.                 \label{eq:n9v32-elliptic}
\end{equation}
For tangent weight \(t\in\mathbb R\), put
\begin{equation}
 w_k=(1+\lambda_k)^t.                                \label{eq:n9v32-weight}
\end{equation}

\begin{theorem}[Sharp Gaussian weighted-score threshold]
\label{thm:n9v32-threshold}
The Gaussian logarithmic derivative is an \(H_t\)-valued score in the
sense of V29 if and only if
\begin{equation}
 \sum_{k=1}^\infty {a_k\over\hbar w_k}<\infty.        \label{eq:n9v32-series}
\end{equation}
Under \eqref{eq:n9v32-elliptic}, Weyl asymptotics give the equivalent
condition
\begin{equation}
 \boxed{t>\sigma+{d\over2}.}                         \label{eq:n9v32-threshold}
\end{equation}
At equality the series diverges logarithmically.
\end{theorem}

\begin{proof}
V29 gives
\[
 \mathbb E\|B\|_{H_t}^2
 =\sum_kw_k^{-1}\|\beta_k\|_2^2
 ={1\over\hbar}\sum_k{a_k\over w_k},
\]
so \eqref{eq:n9v32-series} is necessary and sufficient.  Weyl's law gives
\(\lambda_k\asymp k^{2/d}\).  The summand is therefore comparable to
\(k^{2(\sigma-t)/d}\), whose series converges exactly when
\(2(\sigma-t)/d<-1\).  This is
\eqref{eq:n9v32-threshold}; exponent \(-1\) gives logarithmic divergence.
\end{proof}

\begin{corollary}[All finite score moments have the same trace threshold]
\label{cor:n9v32-all-r}
For any fixed \(r>2\), the weighted V30 moment contribution is
\begin{equation}
 \sum_kw_k^{-1}M_k^{2/r}
 =m_r^{2/r}{1\over\hbar}\sum_k{a_k\over w_k}.        \label{eq:n9v32-rtrace}
\end{equation}
Thus the threshold remains \(t>\sigma+d/2\); asking for the
\((2+\delta)\)-moment does not impose an additional ultraviolet derivative
in the Gaussian model.
\end{corollary}

\begin{proof}
Raise the exact moment in \eqref{eq:n9v32-YM} to the power \(2/r\), then
apply \Cref{thm:n9v32-threshold}.
\end{proof}

\begin{proposition}[Unweighted score-trace divergence]
\label{prop:n9v32-unweighted}
For \(t=0\), the score-vector norm diverges for every
positive-order elliptic \(P\).  With a spectral cutoff
\(\lambda_k\leq\Lambda^2\), its leading size is
\begin{equation}
 \sum_{\lambda_k\leq\Lambda^2}\|\beta_k\|_2^2
 \asymp\hbar^{-1}\Lambda^{d+2\sigma}.               \label{eq:n9v32-divergence}
\end{equation}
\end{proposition}

\begin{proof}
The eigenvalue counting density has order
\(\lambda^{d/2-1}d\lambda\).  Integrating
\(\lambda^\sigma\) up to \(\Lambda^2\) gives order
\(\Lambda^{d+2\sigma}\).  The lower comparison in
\eqref{eq:n9v32-elliptic} proves divergence.
\end{proof}

\begin{firewall}[The tangent space must be declared]
\label{fire:n9v32-tangent}
The condition \(t>\sigma+d/2\) means the admissible metric variations are
sufficiently smooth and the score is correspondingly distributional in
the dual tangent scale.  It does not prove that arbitrary physical metric
variations lie in that Hilbert space.  Changing \(t\) to force convergence
changes the closed Dirichlet form and the QESDC domain.
\end{firewall}

\subsection{A controlled interaction perturbation}
\label{sec:n9v32-interaction}

At finite cutoff let
\begin{equation}
 S=S_0+I,qquad
 d\mu=Z^{-1}e^{-S/\hbar}dx,                          \label{eq:n9v32-interacting}
\end{equation}
and assume this is a probability for which the V31 truncation argument is
valid.  Then
\begin{equation}
 \beta_k={1\over\hbar}(a_kx_k+\partial_kI),qquad
 D_k\beta_k={1\over\hbar}(a_k+\partial_k^2I).         \label{eq:n9v32-interacting-score}
\end{equation}

\begin{theorem}[Interacting diagonal-Hessian loader]
\label{thm:n9v32-interacting}
Let \(r>2\) and suppose
\begin{equation}
 Z_{k,r}:=int|\partial_k^2I|^{r/2}d\mu<\infty.       \label{eq:n9v32-Zkr}
\end{equation}
Then
\begin{align}
 \|\beta_k\|_2^2
 &={1\over\hbar}
 \left(a_k+\int\partial_k^2I\,d\mu\right),          \label{eq:n9v32-interacting-Fisher}\\
 \int|\beta_k|^rd\mu
 &\leq{(r-1)^{r/2}2^{r/2-1}\over\hbar^{r/2}}
 \left(a_k^{r/2}+Z_{k,r}\right).                    \label{eq:n9v32-interacting-r}
\end{align}
Consequently the weighted CHMC and RSESC moment rows hold if
\begin{equation}
 \sum_kw_k^{-1}a_k<\infty,qquad
 \sum_kw_k^{-1}Z_{k,r}^{2/r}<\infty.                 \label{eq:n9v32-interacting-sum}
\end{equation}
\end{theorem}

\begin{proof}
The V31 Fisher identity and
\eqref{eq:n9v32-interacting-score} give
\eqref{eq:n9v32-interacting-Fisher}.  The higher Fisher bound gives
\[
 \int|\beta_k|^rd\mu
 \leq(r-1)^{r/2}\hbar^{-r/2}
 \int|a_k+\partial_k^2I|^{r/2}d\mu.
\]
Use \(|u+v|^q\leq2^{q-1}(|u|^q+|v|^q)\) for
\(q=r/2\), proving \eqref{eq:n9v32-interacting-r}.  Raising that bound to
\(2/r\), using subadditivity up to a fixed \(r\)-constant, and summing
with \(w_k^{-1}\) proves the final assertion.
\end{proof}

\begin{firewall}[Diagonal Hessians must come from the complete score]
\label{fire:n9v32-complete}
In the physical theory, \(I\) in
\eqref{eq:n9v32-interacting-score} must denote the complete assembled
interaction after constraint, determinant, BV, counterterm, and
reference-density repartition.  Sectorwise Hessian moments do not prove
\eqref{eq:n9v32-Zkr} if their signed sum contains essential cancellations.
\end{firewall}

\subsection{The conformal sign obstruction in the solvable model}
\label{sec:n9v32-conformal}

\begin{proposition}[Negative Gaussian direction is not normalizable]
\label{prop:n9v32-negative}
If \(A\) has an eigenvector \(e\) with eigenvalue \(a<0\), then
\begin{equation}
 \int_{\mathbb R^N}e^{-\langle x,Ax\rangle/(2\hbar)}dx=\infty.       \label{eq:n9v32-negative}
\end{equation}
Hence no positive Gaussian probability, ELDC identity, or Fisher trace is
defined on that real contour.
\end{proposition}

\begin{proof}
Restrict the integral to a cylinder with bounded coordinates orthogonal to
\(e\).  Along \(x=se+y\), the exponential contains
\(e^{|a|s^2/(2\hbar)}\) times a positive bounded-
\(y\) factor.  Its integral over \(s\in\mathbb R\) diverges.
\end{proof}

\begin{corollary}[What a positive regulator does and does not supply]
\label{cor:n9v32-regulator}
If an auxiliary positive operator \(R_a\) makes
\(A_a=A+R_a>0\) at every finite cutoff, then V32 applies with eigenvalues
of \(A_a\).  Removing \(R_a\) while a negative conformal eigenvalue
reappears cannot yield a positive Gaussian limit by this argument.
\end{corollary}

\begin{proof}
Positivity gives the finite Gaussian construction.  If the limiting real
quadratic form is negative in one direction, the preceding proposition
reproduces the normalization obstruction; the finite score estimates alone
cannot change that sign.
\end{proof}

\begin{firewall}[Contour rotation changes the problem]
\label{fire:n9v32-contour}
Rotating a negative conformal coordinate can make an oscillatory or complex
Gaussian integral meaningful, but it no longer defines the positive
probability used by V28--V31.  Such a contour requires a separate reality,
Ward, and reflection analysis and cannot be inserted into QESDC by notation.
\end{firewall}

\subsection{Gaussian principal-score card}
\label{sec:n9v32-card}

\begin{definition}[Gaussian principal-score card, GPSC]
\label{def:n9v32-GPSC}
GPSC records:
\begin{enumerate}[label=\textup{(G\arabic*)},leftmargin=4em]
\item positivity of the quotient principal operator at every finite cutoff;
\item elliptic comparison \eqref{eq:n9v32-elliptic};
\item tangent weight \(t>\sigma+d/2\);
\item the complete interaction Hessian moments
\eqref{eq:n9v32-interacting-sum};
\item treatment of zero modes as gauge or retained slow variables; and
\item a separate conformal-sign mechanism stable under regulator removal.
\end{enumerate}
\end{definition}

\begin{theorem}[GPSC loads the principal CHMC trace]
\label{thm:n9v32-GPSC}
Items (G1)--(G5) give the weighted Gaussian and interacting score moments
needed by CHMC on the positive quotient sector.  Item (G6) is logically
independent and is necessary before those finite-cutoff estimates can
belong to a positive unregulated Einstein measure.
\end{theorem}

\begin{proof}
Use \Cref{thm:n9v32-threshold,cor:n9v32-all-r,
thm:n9v32-interacting}.  Independence of (G6) is
\Cref{prop:n9v32-negative,cor:n9v32-regulator}.
\end{proof}

\begin{assessment}[Progress at ninth-series V32]
\label{ass:n9v32-status}
The dimension-counting part of MSCC is now exact.  For an order
\(2\sigma\) positive Gaussian principal operator in dimension \(d\), the
Einstein score is Hilbert valued precisely on tangent weights
\(t>\sigma+d/2\), and all finite higher score moments use the same
threshold.  Interactions reduce to weighted moments of the complete
diagonal Hessian.  The real conformal sign remains an obstruction before
normalization, not a missing trace estimate.
\end{assessment}

\begin{programme}[Ninth-series V33: relative Hessian perturbations]
\label{prog:n9v32-V33}
V33 will replace the coordinatewise interaction moment
\eqref{eq:n9v32-interacting-sum} by basis-independent Schatten and relative
form conditions.  It will determine which trace ideals preserve the
weighted Fisher bound under interacting Hessian perturbations and isolate
the part still requiring the archived derivative-faithful BV shell estimate.
\end{programme}

\begin{firewall}[Claim boundary after ninth-series V32]
\label{fire:n9v32-boundary}
V32 does not prove GPSC for the physical regulator, positivity of the
conformal sector, the complete interaction Hessian moment, boundary-flux
closure, reflection positivity, reconstruction, or the full continuum.
\end{firewall}

\begin{theorem}[Ninth-series V32 result]
\label{thm:n9v32-result}
The finite positive Gaussian Einstein score has exact moments
\eqref{eq:n9v32-moments}.  Its continuum weighted trace is finite exactly
when \(t>\sigma+d/2\); the same condition controls every V30 higher-moment
row.  An interacting perturbation is controlled by
\eqref{eq:n9v32-interacting-sum}, while any negative real conformal
eigenvalue destroys the underlying Gaussian probability.
\end{theorem}

\begin{proof}
Use \Cref{thm:n9v32-moments,thm:n9v32-threshold,
cor:n9v32-all-r,thm:n9v32-interacting,
prop:n9v32-negative,thm:n9v32-GPSC}.
\end{proof}

\paragraph{Parent-version record.}
V32 was formed from
\texttt{Renewal\_SM\_Einstein\_Continuum\_Research\_Notes\_V31.tex}.
The V31 parent had \(10\,763\) logical source lines, \(387\,066\) bytes,
and SHA--256
\[
 \texttt{1071AF6A03DC5734AC4BA708D35C26BF162914A297BF50ACF05B3DFF05C94DFE}.
\]

% =============================================================================
% END APPENDED NINTH-SERIES V32 MATERIAL
% =============================================================================


% =============================================================================
% BEGIN APPENDED NINTH-SERIES V33 MATERIAL
% =============================================================================

\clearpage
\section{Ninth-series V33: relative Hessian ideals and the mixed-moment obstruction}
\label{sec:v9v33}

\subsection{Operator form of the weighted Fisher trace}
\label{sec:n9v33-trace}

Work first at finite cutoff.  Let \(A>0\) be the Gaussian principal
Hessian, let \(W>0\) be the tangent-weight operator, and let
\(K(x)=D^2I(x)\) be the self-adjoint Hessian of the complete assembled
interaction.  In a common orthonormal basis \((e_k)\) diagonalizing
\(A\) and \(W\), write
\begin{equation}
 Ae_k=a_ke_k,qquad We_k=w_ke_k.                     \label{eq:n9v33-AW}
\end{equation}
Assume the complete law
\(
d\mu=Z^{-1}e^{-(S_0+I)/\hbar}dx
\)
is normalizable and satisfies the V31 truncation and no-flux hypotheses.
The complete differentiated score is
\begin{equation}
 D_k\beta_k={1\over\hbar}
 \langle e_k,(A+K)e_k\rangle.                       \label{eq:n9v33-Hdiag}
\end{equation}

\begin{proposition}[Trace-ideal loader for the \(L^2\) score]
\label{prop:n9v33-S1}
Suppose
\begin{equation}
 W^{-1/2}AW^{-1/2}\in\mathfrak S_1,qquad
 W^{-1/2}K(x)W^{-1/2}\in\mathfrak S_1               \label{eq:n9v33-S1}
\end{equation}
almost surely, and
\begin{equation}
 \mathbb E\|W^{-1/2}KW^{-1/2}\|_{\mathfrak S_1}<\infty.            \label{eq:n9v33-S1moment}
\end{equation}
Then
\begin{align}
 \sum_kw_k^{-1}\|\beta_k\|_2^2
 &={1\over\hbar}\Tr(W^{-1/2}AW^{-1/2})
 \notag\\
 &\quad+{1\over\hbar}\mathbb E
       \Tr(W^{-1/2}KW^{-1/2}),                       \label{eq:n9v33-trace}\
 \sum_kw_k^{-1}\|\beta_k\|_2^2
 &\leq{1\over\hbar}\left[
 \|W^{-1/2}AW^{-1/2}\|_{\mathfrak S_1}
 +\mathbb E\|W^{-1/2}KW^{-1/2}\|_{\mathfrak S_1}
 \right].                                            \label{eq:n9v33-trace-bound}
\end{align}
\end{proposition}

\begin{proof}
At finite rank, sum the V31 Fisher identity with weights \(w_k^{-1}\).
The principal sum is the first trace.  The interaction sum is
\[
 \sum_kw_k^{-1}\langle e_k,Ke_k\rangle
 =\Tr(W^{-1/2}KW^{-1/2}).
\]
Absolute integrability follows from \eqref{eq:n9v33-S1moment}, so expectation
and trace may be interchanged.  The trace is bounded by the trace norm,
giving \eqref{eq:n9v33-trace-bound}.  Uniform finite-rank approximation
then gives the stated trace-class version.
\end{proof}

\begin{remark}[Cancellations remain inside the complete trace]
\label{rem:n9v33-cancellation}
Equation \eqref{eq:n9v33-trace} retains signed interaction cancellations.
The sufficient bound \eqref{eq:n9v33-trace-bound} spends them by taking one
trace norm, but only after all physical sectors and counterterm owners have
been assembled into \(K\).
\end{remark}

\subsection{Why a random trace norm does not load the higher moments}
\label{sec:n9v33-no-go}

Put \(q=r/2>1\).  The V30--V31 higher-moment trace contains
\begin{equation}
 \sum_kw_k^{-1}
 \left(\mathbb E|langle e_k,Ke_k\rangle|^q\right)^{1/q}.           \label{eq:n9v33-mixed}
\end{equation}
This is an \(\ell^1(L^q)\) mixed norm.  It is not controlled by an
\(L^q(\mathfrak S_1)\) norm.

\begin{theorem}[Random-mode Schatten no-go]
\label{thm:n9v33-no-go}
For every \(N\), there is a random positive rank-one operator \(K_N\) on
\(\mathbb R^N\) such that
\begin{equation}
 \|K_N\|_{\mathfrak S_1}=1\quad\text{almost surely},                \label{eq:n9v33-fixed-S1}
\end{equation}
but, for the standard basis and every \(q>1\),
\begin{equation}
 \sum_{k=1}^N
 \left(\mathbb E|\langle e_k,K_Ne_k\rangle|^q\right)^{1/q}
 =N^{1-1/q}.                                         \label{eq:n9v33-growing-mixed}
\end{equation}
Hence no dimension-free bound of \eqref{eq:n9v33-mixed} follows from
\(\mathbb E\|K\|_{\mathfrak S_1}^q\).
\end{theorem}

\begin{proof}
Let \(J\) be uniform on \(\{1,\ldots,N\}\) and put
\(K_N=|e_J\rangle\langle e_J|\).  It is a positive rank-one projection,
so \eqref{eq:n9v33-fixed-S1} holds.  Its \(k\)-th diagonal entry is
\({\bf1}_{\{J=k\}}\), whose \(L^q\) norm is \(N^{-1/q}\).  Summing over
\(k\) gives \eqref{eq:n9v33-growing-mixed}.
\end{proof}

\begin{firewall}[The order of the two norms matters]
\label{fire:n9v33-mixed}
Minkowski gives
\(
 \|\sum_k|K_{kk}|\|_{L^q}
 \leq\sum_k\|K_{kk}\|_{L^q}
\), in the wrong direction for controlling
\eqref{eq:n9v33-mixed} by a random trace norm.  Interchanging the
\(\ell^1\) and \(L^q\) norms without an additional domination hypothesis
would erase the random-mode counterexample.
\end{firewall}

\subsection{A basis-independent sufficient condition}
\label{sec:n9v33-relative}

\begin{definition}[Random relative Hessian bound]
\label{def:n9v33-relative}
The interaction Hessian is \(A\)-form dominated by a random variable
\(\kappa\geq0\) if
\begin{equation}
 |\langle h,K(x)h\rangle|
 \leq\kappa(x)\langle h,Ah\rangle
 \quad\text{for every }h                            \label{eq:n9v33-relative}
\end{equation}
outside one common null set.
\end{definition}

\begin{theorem}[Relative-form higher-moment loader]
\label{thm:n9v33-relative}
Let \(q=r/2>1\), assume \eqref{eq:n9v33-relative}, and suppose
\(\kappa\in L^q(\mu)\).  Then
\begin{align}
 \left(
 \mathbb E|\langle e_k,(A+K)e_k\rangle|^q
 \right)^{1/q}
 &\leq a_k\bigl(1+\|\kappa\|_{L^q}\bigr),           \label{eq:n9v33-diag-bound}\\
 \sum_kw_k^{-1}
 \left(
 \mathbb E|D_k\beta_k|^q
 \right)^{1/q}
 &\leq{1+\|\kappa\|_{L^q}\over\hbar}
       \sum_k{a_k\over w_k}.                         \label{eq:n9v33-relative-trace}
\end{align}
Consequently the V32 threshold \(t>\sigma+d/2\) loads CHMC and the V30
higher score moment with no additional ultraviolet derivative loss.
\end{theorem}

\begin{proof}
Set \(h=e_k\) in \eqref{eq:n9v33-relative}.  Then
\[
 |\langle e_k,(A+K)e_k\rangle|
 \leq(1+\kappa)a_k.
\]
Take the \(L^q\) norm to obtain \eqref{eq:n9v33-diag-bound}.  Divide by
\(\hbar\), multiply by \(w_k^{-1}\), and sum.  The final statement follows
from \Cref{thm:n9v32-threshold,cor:n9v31-V30}.
\end{proof}

\begin{corollary}[Deterministic relative bound]
\label{cor:n9v33-deterministic}
If \(-\varepsilon A\leq K(x)\leq\varepsilon A\) as quadratic forms with a
deterministic \(\varepsilon<\infty\), then
\eqref{eq:n9v33-relative-trace} holds with
\(\|\kappa\|_{L^q}=\varepsilon\) for every \(q\).  If
\(\varepsilon<1\), the complete Hessian also has the pointwise lower bound
\begin{equation}
 A+K\geq(1-\varepsilon)A.                            \label{eq:n9v33-positive}
\end{equation}
\end{corollary}

\begin{proof}
The two-sided form inequality is exactly
\eqref{eq:n9v33-relative} with constant \(\kappa=\varepsilon\).  Adding
the lower inequality \(K\geq-\varepsilon A\) to \(A\) gives
\eqref{eq:n9v33-positive}.
\end{proof}

\begin{firewall}[Relative Hessian positivity is not action coercivity]
\label{fire:n9v33-coercivity}
Even \eqref{eq:n9v33-positive} is a local or pointwise Hessian statement.
Global normalizability also depends on the topology of field space,
zero/slow modes, boundary behavior, and growth of the action along long
curves.  Conversely, a nonconvex but coercive action can define a
probability without satisfying a uniform \(\varepsilon<1\) form bound.
\end{firewall}

\subsection{The exact mixed diagonal ideal}
\label{sec:n9v33-ideal}

For a fixed weighted tangent structure and probability, define
\begin{equation}
 \|K\|_{\mathfrak D_q(W;\mu)}
 :=\sum_kw_k^{-1}
 \|\langle e_k,Ke_k\rangle\|_{L^q(\mu)}.             \label{eq:n9v33-Dq}
\end{equation}

\begin{proposition}[Minimal diagonal loader]
\label{prop:n9v33-Dq}
If
\begin{equation}
 \Tr(W^{-1/2}AW^{-1/2})<\infty,qquad
 \|K\|_{\mathfrak D_q(W;\mu)}<\infty,               \label{eq:n9v33-Dq-condition}
\end{equation}
then the weighted CHMC moment trace is finite.  The relative-form condition
\eqref{eq:n9v33-relative} implies
\begin{equation}
 \|K\|_{\mathfrak D_q(W;\mu)}
 \leq\|\kappa\|_{L^q}
 \Tr(W^{-1/2}AW^{-1/2}).                             \label{eq:n9v33-Dq-relative}
\end{equation}
\end{proposition}

\begin{proof}
Use the triangle inequality in \(L^q\) on each diagonal entry of
\(A+K\), then sum with weights.  This is exactly the trace required in
\eqref{eq:n9v33-relative-trace}.  The second assertion follows by applying
\eqref{eq:n9v33-relative} to every \(e_k\).
\end{proof}

\begin{remark}[Basis dependence is meaningful here]
\label{rem:n9v33-basis}
The ideal \(\mathfrak D_q\) records moments of the declared cylinder
directions, so it depends on that tangent frame.  Relative-form domination
is the invariant sufficient replacement.  Without such a replacement,
the frame dependence cannot be removed by a Schatten norm because of
\Cref{thm:n9v33-no-go}.
\end{remark}

\subsection{Adjacent interacting shell Hessians}
\label{sec:n9v33-shells}

Let the invariant local RG projection of the frozen archive be performed
before the following decomposition.  On one transported probability
space, write the irrelevant complete Hessian increment of shell \(n\) as
\(K_n\).

\begin{theorem}[Summable relative-Hessian shell compiler]
\label{thm:n9v33-shell}
Suppose for some \(q>1\) there are nonnegative random variables
\(\kappa_n\) such that
\begin{equation}
 |\langle h,K_nh\rangle|
 \leq\kappa_n\langle h,Ah\rangle,qquad
 \sum_n\|\kappa_n\|_{L^q}<\infty.                  \label{eq:n9v33-shell-bound}
\end{equation}
Then \(\kappa:=\sum_n\kappa_n\) belongs to \(L^q\), the quadratic-form
series \(K:=\sum_nK_n\) converges absolutely relative to \(A\), and
\begin{equation}
 |\langle h,Kh\rangle|
 \leq\kappa\langle h,Ah\rangle.                    \label{eq:n9v33-shell-limit}
\end{equation}
Consequently the complete irrelevant interaction satisfies the V33
relative-form higher-moment loader.
\end{theorem}

\begin{proof}
Minkowski's inequality gives
\(\|\sum_n\kappa_n\|_{L^q}\leq\sum_n\|\kappa_n\|_{L^q}\), so the sum is
finite almost surely and in \(L^q\).  For every \(h\),
\[
 \sum_n|\langle h,K_nh\rangle|
 \leq\left(\sum_n\kappa_n\right)\langle h,Ah\rangle.
\]
This proves absolute form convergence and \eqref{eq:n9v33-shell-limit}.
Apply \Cref{thm:n9v33-relative}.
\end{proof}

\begin{firewall}[One more metric derivative is required]
\label{fire:n9v33-archive}
The frozen interacting-shell programme controls a first metric derivative
only after the connected covariance and invariant local RG projection are
assembled.  Equation \eqref{eq:n9v33-shell-bound} is a bound on the next
metric derivative in relative form and does not follow formally from that
first-derivative estimate.  Coincident insertions and contact counterterms
must be included before its norm is taken.
\end{firewall}

\subsection{Relative Hessian score card}
\label{sec:n9v33-card}

\begin{definition}[Relative Hessian score card, RHSC]
\label{def:n9v33-RHSC}
RHSC records:
\begin{enumerate}[label=\textup{(R\arabic*)},leftmargin=4em]
\item the positive quotient principal Hessian \(A\) and tangent weight
\(W\) satisfying \(W^{-1/2}AW^{-1/2}\in\mathfrak S_1\);
\item complete interaction Hessians after invariant local RG subtraction;
\item either the relative bound \eqref{eq:n9v33-relative} in \(L^q\), or
the exact mixed ideal \eqref{eq:n9v33-Dq-condition};
\item for a shell construction, the summability
\eqref{eq:n9v33-shell-bound};
\item the V31 core and boundary-flux rows; and
\item a separate positive treatment of conformal and retained slow modes.
\end{enumerate}
\end{definition}

\begin{theorem}[RHSC loads CHMC without a basiswise hypothesis]
\label{thm:n9v33-RHSC}
RHSC gives the weighted \(L^2\) and higher score traces required by CHMC,
RSESC, and MSCC on the positive quotient sector.  It does not obtain the
conformal row or prove the relative interacting-shell estimate.
\end{theorem}

\begin{proof}
The \(L^2\) assertion is \Cref{prop:n9v33-S1}; the higher assertion is
\Cref{thm:n9v33-relative,prop:n9v33-Dq}.  The shell alternative is
\Cref{thm:n9v33-shell}.  Apply \Cref{thm:n9v31-CHMC} and retain item (R6)
as a separate hypothesis by \Cref{prop:n9v32-negative}.
\end{proof}

\begin{assessment}[Progress at ninth-series V33]
\label{ass:n9v33-status}
The interaction Hessian gate is now basis independent under a random
relative-form bound.  Trace class control alone closes the Fisher
\(L^2\) trace but provably fails for the mixed higher moments; the
random-mode counterexample quantifies the failure.  A summable relative
bound on complete RG-projected shell Hessians is the exact new analytic
target, and it requires one more derivative than the archived first-shell
stress estimate.
\end{assessment}

\begin{programme}[Ninth-series V34: the conformal direction upstream]
\label{prog:n9v33-V34}
V34 will return to the conformal mode at the action level.  It will compute
the Euclidean Einstein--Hilbert conformal transformation, prove the
negative-gradient direction, and test which local positive additions can
make a finite-cutoff real measure coercive without claiming that their
removal preserves positivity.
\end{programme}

\begin{firewall}[Claim boundary after ninth-series V33]
\label{fire:n9v33-boundary}
V33 does not prove RHSC for the interacting BV shells, control coincident
second metric insertions, stabilize the conformal direction, close boundary
flux, reconstruct a Lorentzian theory, or complete the continuum.
\end{firewall}

\begin{theorem}[Ninth-series V33 result]
\label{thm:n9v33-result}
An \(L^1\) trace-ideal moment of the complete relative Hessian loads the
weighted Fisher \(L^2\) trace.  For the higher score moments, no random
Schatten bound suffices in general; the exact mixed ideal
\eqref{eq:n9v33-Dq} or the invariant relative-form domination
\eqref{eq:n9v33-relative} does.  Summability of the shell bounds
\eqref{eq:n9v33-shell-bound} passes that domination to the interacting
irrelevant limit.
\end{theorem}

\begin{proof}
Use \Cref{prop:n9v33-S1,thm:n9v33-no-go,
thm:n9v33-relative,prop:n9v33-Dq,
thm:n9v33-shell,thm:n9v33-RHSC}.
\end{proof}

\paragraph{Parent-version record.}
V33 was formed from
\texttt{Renewal\_SM\_Einstein\_Continuum\_Research\_Notes\_V32.tex}.
The V32 parent had \(11\,115\) logical source lines, \(400\,513\) bytes,
and SHA--256
\[
 \texttt{818AF24A8375D6E271ACC2CBA534B65FCF35C790C0CB688354956EEC693ABD58}.
\]

% =============================================================================
% END APPENDED NINTH-SERIES V33 MATERIAL
% =============================================================================


% =============================================================================
% BEGIN APPENDED NINTH-SERIES V34 MATERIAL
% =============================================================================

\clearpage
\section{Ninth-series V34: constraint pullback of the Einstein Fisher Hessian}
\label{sec:v9v34}

\subsection{Why the conformal calculation is not repeated}
\label{sec:n9v34-audit}

The frozen SM archives already prove the Euclidean conformal-gradient sign,
the failure of a naive real Gaussian, the exact fourth-order/low-mode
positivity threshold, the cosmological-Hessian conflict, and the
linearized ADM statement that nonzero-momentum scalar configuration modes
are constraint or gauge variables.  They also prove that a sourced scalar
returns as an inverse-Laplacian response and identify the norm-one
divergence-source estimate in the gradient-energy topology.

The missing interface is different.  V31--V33 estimate the differentiated
Einstein score on the actual probability carrier.  If the conformal
coordinate is removed by a constraint chart, the relevant Hessian is the
pullback Hessian, the tangent vector contains the constraint response, and
the reduced coarea density contributes its own logarithmic second
derivative.  V34 derives that complete object.

\subsection{A nonlinear reduced chart}
\label{sec:n9v34-chart}

Let \(P\subset\mathbb R^m\) and \(Y\subset\mathbb R^r\) be open, and let
\begin{equation}
 \iota(p):=(p,\varphi(p)),qquad
 C:=\iota(P)\subset P\times Y                       \label{eq:n9v34-iota}
\end{equation}
be a \(C^2\) constraint chart.  Let the reduced positive density be
\begin{equation}
 dm_C=\rho_C(p)J_\Phi(p)^{-1}sqrt{\det G(p)},dp
 =:r_C(p)\,dp,                                      \label{eq:n9v34-density}
\end{equation}
where \(G=(D\iota)^*D\iota\) is the induced Gram matrix and
\(J_\Phi\) is the normal coarea Jacobian.  Put
\begin{equation}
 \bar S(p):=S(\iota(p)),qquad
 d\bar\mu=\bar Z^{-1}e^{-\bar S/\hbar}r_C(p)dp.       \label{eq:n9v34-law}
\end{equation}

For a constant reduced direction \(h\in\mathbb R^m\), define
\begin{equation}
 \bar V_h:=h\cdot\nabla_p,qquad
 v_h(p):=D\iota(p)h=(h,D\varphi(p)h).                \label{eq:n9v34-lift}
\end{equation}

\begin{lemma}[Reduced logarithmic score]
\label{lem:n9v34-score}
Assuming zero boundary flux in direction \(h\), the reduced logarithmic
derivative is
\begin{equation}
 \bar\beta_h
 ={1\over\hbar}D\bar S[h]-D\log r_C[h],              \label{eq:n9v34-score}
\end{equation}
and
\begin{equation}
 D_h\bar\beta_h
 ={1\over\hbar}D^2\bar S[h,h]-D^2\log r_C[h,h].      \label{eq:n9v34-Dscore}
\end{equation}
\end{lemma}

\begin{proof}
The divergence of the constant coordinate field relative to
\(r_Cdp\) is \(D\log r_C[h]\).  Insert this in the V28 score formula and
differentiate once more along the constant direction.
\end{proof}

\begin{theorem}[Exact pullback Hessian]
\label{thm:n9v34-pullback}
For every \(p\in P\) and \(h\in\mathbb R^m\),
\begin{equation}
 \boxed{
 D^2\bar S(p)[h,h]
 =D^2S(\iota(p))[v_h,v_h]
 +DS(\iota(p))[D^2\iota(p)[h,h]].}                  \label{eq:n9v34-pullback}
\end{equation}
The second term cannot in general be discarded merely because \(p\) is a
stationary point of the restricted action.
\end{theorem}

\begin{proof}
Differentiate \(\bar S=S\circ\iota\) once:
\(D\bar S[h]=DS[D\iota h]\).  Differentiating the product of the ambient
differential with the tangent lift gives the two terms in
\eqref{eq:n9v34-pullback}.  Restricted stationarity says only that
\(DS[D\iota h]=0\); it does not imply \(DS=0\) on ambient normal vectors,
including \(D^2\iota[h,h]\).
\end{proof}

\begin{corollary}[Complete differentiated reduced score]
\label{cor:n9v34-complete}
Combining \Cref{lem:n9v34-score,thm:n9v34-pullback},
\begin{equation}
 \boxed{
 \hbar D_h\bar\beta_h
 =D^2S[v_h,v_h]
 +DS[D^2\iota[h,h]]
 -\hbar D^2\log r_C[h,h].}                          \label{eq:n9v34-complete}
\end{equation}
Thus the CHMC Hessian has ambient, constraint-curvature, and coarea-density
rows.
\end{corollary}

\begin{proof}
Substitute \eqref{eq:n9v34-pullback} into
\eqref{eq:n9v34-Dscore} and multiply by \(\hbar\).
\end{proof}

\subsection{The constrained stationary Hessian is a Lagrangian Hessian}
\label{sec:n9v34-Lagrange}

Suppose \(C=\{\Phi=0\}\) for a submersion
\(\Phi:P\times Y\to\mathbb R^r\).  At a constrained stationary point
\(x=\iota(p)\), use the convention
\begin{equation}
 DS(x)=D\Phi(x)^*\lambda                             \label{eq:n9v34-multiplier}
\end{equation}
for a multiplier \(\lambda\).

\begin{theorem}[Bordered Hessian identity]
\label{thm:n9v34-bordered}
At such a point,
\begin{equation}
 D^2\bar S[h,h]
 =D^2\bigl(S-\langle\lambda,\Phi\rangle\bigr)
   [v_h,v_h].                                       \label{eq:n9v34-bordered}
\end{equation}
\end{theorem}

\begin{proof}
Differentiate \(\Phi\circ\iota=0\) twice:
\begin{equation}
 D\Phi[D^2\iota[h,h]]+D^2\Phi[v_h,v_h]=0.           \label{eq:n9v34-Phi2}
\end{equation}
By \eqref{eq:n9v34-multiplier},
\[
 DS[D^2\iota[h,h]]
 =\langle\lambda,D\Phi[D^2\iota[h,h]]\rangle
 =-\langle\lambda,D^2\Phi[v_h,v_h]\rangle.
\]
Insert this in \eqref{eq:n9v34-pullback}.
\end{proof}

\begin{firewall}[The ambient Hessian is not the physical Hessian]
\label{fire:n9v34-ambient}
At a constrained stationary point, the physical quadratic form is the
Hessian of the Lagrangian \(S-\langle\lambda,\Phi\rangle\) restricted to
the tangent space.  Using \(D^2S\) alone loses the multiplier curvature
term.  This term is also present in the differentiated Einstein score and
must be assigned to the constraint/coarea ledger.
\end{firewall}

\subsection{Linear constraints and the fate of a negative normal mode}
\label{sec:n9v34-linear}

Let \(C=\{(p,y):y=Bp\}\) for a linear map \(B\), so
\(\iota(p)=(p,Bp)\) and \(J=(I,B)\).  For a quadratic ambient action
\(S(x)=\frac12\langle x,Hx\rangle\),
\begin{equation}
 \bar S(p)={1\over2}\langle p,H_{
m red}p\rangle,qquad
 H_{\rm red}:=J^*HJ.                                \label{eq:n9v34-Hred}
\end{equation}

\begin{proposition}[Restriction can remove an ambient negative direction]
\label{prop:n9v34-negative-normal}
The reduced real Gaussian is normalizable exactly when
\(H_{\rm red}>0\), regardless of the sign of \(H\) on vectors transverse
to \(\operatorname{Ran}J\).  In particular, an ambient negative conformal
direction is harmless for this Gaussian only if the physical measure is
actually supported on a constraint surface whose tangent excludes it.
\end{proposition}

\begin{proof}
The reduced density is proportional to
\(e^{-\langle p,H_{\rm red}p\rangle/(2\hbar)}dp\) times a constant linear
coarea factor.  It is integrable exactly when every eigenvalue of
\(H_{\rm red}\) is positive.  Values of \(H\) off
\(\operatorname{Ran}J\) never occur in the restricted integral.  If the
ambient coordinate is integrated rather than constrained, V32's negative
Gaussian no-go applies instead.
\end{proof}

\begin{corollary}[Constraint imposition and contour repair are distinct]
\label{cor:n9v34-distinct}
Replacing an ambient integral by a coarea-supported constraint measure can
remove a negative normal mode while retaining positivity.  Rotating that
mode on a complex contour and integrating it is a different construction
with different reality and reflection requirements.
\end{corollary}

\begin{proof}
The first statement is \Cref{prop:n9v34-negative-normal}; the second follows
because a complex contour is not the positive measure
\eqref{eq:n9v34-law} on the real constraint manifold.
\end{proof}

\subsection{Relative-form control after nonlinear pullback}
\label{sec:n9v34-relative}

Let \(A_P>0\) be a declared principal quadratic form on reduced directions.
Suppose there are nonnegative random variables
\(\kappa_1,\kappa_2,\kappa_3\) such that for every reduced \(h\),
\begin{align}
 |D^2S[v_h,v_h]|&\leq\kappa_1\langle h,A_Ph\rangle,                 \label{eq:n9v34-k1}\\
 |DS[D^2\iota[h,h]]|&\leq\kappa_2\langle h,A_Ph\rangle,           \label{eq:n9v34-k2}\\
 \hbar|D^2\log r_C[h,h]|&\leq\kappa_3\langle h,A_Ph\rangle.      \label{eq:n9v34-k3}
\end{align}

\begin{theorem}[Constraint-pullback RHSC loader]
\label{thm:n9v34-relative}
If \(q>1\) and \(\kappa_j\in L^q(\bar\mu)\), then the complete reduced
differentiated score satisfies
\begin{equation}
 \hbar|D_h\bar\beta_h|
 \leq\kappa\langle h,A_Ph\rangle,qquad
 \kappa:=\kappa_1+\kappa_2+\kappa_3\in L^q.          \label{eq:n9v34-relative}
\end{equation}
Consequently V33 relative-form control and the V31 higher Fisher moment
hold on the reduced carrier whenever
\begin{equation}
 \Tr(W^{-1/2}A_PW^{-1/2})<\infty.                   \label{eq:n9v34-trace}
\end{equation}
\end{theorem}

\begin{proof}
Take absolute values in \eqref{eq:n9v34-complete} and apply
\eqref{eq:n9v34-k1}--\eqref{eq:n9v34-k3}.  Minkowski gives
\(\|\kappa\|_q\leq\sum_j\|\kappa_j\|_q\).  The remaining assertions are
\Cref{thm:n9v33-relative,thm:n9v31-higher} applied to the reduced
directions.
\end{proof}

\begin{firewall}[A smooth implicit chart does not give uniform constants]
\label{fire:n9v34-uniform}
The implicit-function theorem gives \(D\varphi\) and \(D^2\varphi\) at one
finite regulator, but the moments \(\kappa_2,\kappa_3\) may diverge when the
constraint inverse loses its normal singular-value floor or the coarea
Jacobian approaches zero.  Uniform RHSC requires the quantitative
constraint and singular-neighborhood cards from the current programme.
\end{firewall}

\subsection{Source response and the tangent-lift norm}
\label{sec:n9v34-source}

The archived linearized ADM calculation has the schematic constraint
response
\begin{equation}
 y=(-\Delta)^{-1}\rho.                               \label{eq:n9v34-response}
\end{equation}
The tangent lift \(v_h=(h,D\varphi h)\) is useful in an energy topology
only if this response is uniformly bounded there.

\begin{proposition}[Imported divergence-source row loads the lift]
\label{prop:n9v34-divergence}
Assume the nonzero-mode source factors as
\(\rho=\operatorname{div}J\).  In the gradient-energy norm,
\begin{equation}
 \|\nabla(-\Delta)^{-1}\operatorname{div}J\|_2
 \leq\|J\|_2.                                       \label{eq:n9v34-Riesz}
\end{equation}
Therefore a cofinally bounded source-to-\(J\) map gives a cofinally bounded
linear tangent lift.  Without the divergence factorization, the operator
from \(\rho\in L^2\) to the gradient response has norm
\(|k_{\min}|^{-1}\) on a periodic box and is not volume uniform.
\end{proposition}

\begin{proof}[Archived proof and present use]
The multiplier estimate and the least-mode counterexample were proved in
the frozen eighth-series V49 section entitled
\emph{Linearized ADM reduction and the transferred scalar source}.  V34
does not repeat that calculation.  Applying its norm-one source-to-gradient
operator to the second component of \(v_h=(h,D\varphi h)\) gives the stated
uniform tangent-lift bound; the archived least-mode example gives the
failure without factorization.
\end{proof}

\begin{firewall}[Stress conservation is not the factorization]
\label{fire:n9v34-conservation}
Covariant conservation couples spatial and temporal stress components and
zero-mode charges.  It does not automatically express the energy source
on one slice as the spatial divergence of a uniformly bounded \(L^2\)
vector.  The source factorization in
\Cref{prop:n9v34-divergence} remains a separate constraint-response row.
\end{firewall}

\subsection{Constraint Fisher score card}
\label{sec:n9v34-card}

\begin{definition}[Constraint Fisher pullback card, CFPC]
\label{def:n9v34-CFPC}
CFPC records:
\begin{enumerate}[label=\textup{(C\arabic*)},leftmargin=4em]
\item a positive reduced measure of the form \eqref{eq:n9v34-law};
\item the complete nonlinear tangent lift and multiplier Hessian
\eqref{eq:n9v34-bordered};
\item the coarea-density derivative in
\eqref{eq:n9v34-complete};
\item the relative moment bounds
\eqref{eq:n9v34-k1}--\eqref{eq:n9v34-k3};
\item a trace-class reduced tangent weight
\eqref{eq:n9v34-trace};
\item divergence/charge-sector control of long-range sourced scalar
responses; and
\item separate treatment of zero modes, boundary flux, and singular
constraint charts.
\end{enumerate}
\end{definition}

\begin{theorem}[CFPC loads the constraint-reduced CHMC]
\label{thm:n9v34-CFPC}
CFPC gives the relative-Hessian and higher Fisher-score bounds required by
RHSC and CHMC on the physical reduced carrier.  An ambient negative
conformal direction is thereby removed only when the support and tangent
statements in items (C1)--(C2) are actually proved.
\end{theorem}

\begin{proof}
Use \Cref{thm:n9v34-relative,prop:n9v34-divergence}, followed by
\Cref{thm:n9v33-RHSC,thm:n9v31-CHMC}.  The final qualification is
\Cref{prop:n9v34-negative-normal}.
\end{proof}

\begin{assessment}[Progress at ninth-series V34]
\label{ass:n9v34-status}
The archived conformal and ADM results are now connected to the current
Einstein-score programme.  Constraint reduction replaces the ambient
Hessian by the exact Lagrangian pullback Hessian and adds a coarea-density
second derivative.  These rows admit a single relative-form moment bound,
while the sourced tangent lift is volume uniform only in the declared
divergence/charge topology.
\end{assessment}

\begin{programme}[Ninth-series V35: compulsory companion audit]
\label{prog:n9v34-V35}
V35 will inspect the newest YM and NS notes before selecting between the
constraint-coarea Hessian moment, the boundary-flux row, and the
derivative-faithful interacting shell Hessian.  It will continue to V36
after the audit.
\end{programme}

\begin{firewall}[Claim boundary after ninth-series V34]
\label{fire:n9v34-boundary}
V34 does not prove CFPC for the physical nonlinear constraint measure,
the uniform multiplier/coarea Hessian moments, the matter divergence row,
zero-mode control, reflection positivity, reconstruction, or the full
continuum.
\end{firewall}

\begin{theorem}[Ninth-series V34 result]
\label{thm:n9v34-result}
On a nonlinear constraint chart the complete differentiated metric score
is \eqref{eq:n9v34-complete}; at a constrained stationary point its action
row is the restricted Lagrangian Hessian
\eqref{eq:n9v34-bordered}.  The relative moment conditions
\eqref{eq:n9v34-k1}--\eqref{eq:n9v34-k3} load RHSC and CHMC on the reduced
carrier, while a negative ambient conformal mode is excluded only by an
actual support restriction.
\end{theorem}

\begin{proof}
Use \Cref{thm:n9v34-pullback,thm:n9v34-bordered,
prop:n9v34-negative-normal,thm:n9v34-relative,
prop:n9v34-divergence,thm:n9v34-CFPC}.
\end{proof}

\paragraph{Parent-version record.}
V34 was formed from
\texttt{Renewal\_SM\_Einstein\_Continuum\_Research\_Notes\_V33.tex}.
The V33 parent had \(11\,496\) logical source lines, \(415\,332\) bytes,
and SHA--256
\[
 \texttt{AE44D72E783B4E59F2997A0DB0E40BB8732CAC1920DD5DF0FD08313DE0DE7328}.
\]

% =============================================================================
% END APPENDED NINTH-SERIES V34 MATERIAL
% =============================================================================

% =============================================================================
% BEGIN APPENDED NINTH-SERIES V35 MATERIAL
% =============================================================================

\clearpage
\section{Ninth-series V35: normalized chart labels, their Fisher debit,
and the local rare-sector interface}
\label{sec:v9v35}

\subsection{Compulsory companion audit and direction}

This is the compulsory five-version audit announced in
\Cref{prog:n9v34-V35}.  The newest substantive companion files inspected
were
\begin{align*}
 &\texttt{Renewal\_Yang\_Mills\_Mass\_Gap\_Research\_Notes\_V73.tex},
 &&32\,200\text{ lines},\quad 1\,150\,509\text{ bytes},\\
 &\texttt{Renewal\_NS\_Stability\_Research\_Notes\_V26.tex},
 &&5\,319\text{ lines},\quad 278\,283\text{ bytes}.
\end{align*}
Their respective SHA--256 digests were
\begin{align*}
 &\texttt{E5ACD30866839956E854C7A10D4E68A05258329975B8B68DF95B668BCABF82A7},\\
 &\texttt{82C887F8C2C69E4F28FAD7C3DC8DE5B9099CB5A4BF431EBA1FFF3E91935978AD}.
\end{align*}
The Yang--Mills V69--V73 increment supplies a normalized-label
contraction, a local $L^2$--probability compiler, and an atlas IMS debit.
It also warns that a full-volume $L^2$ density norm is generally the
wrong object.  The Navier--Stokes file supplies no new type-compatible
input for the constraint-coarea Fisher row.  The useful transfer is
therefore structural: introduce normalized auxiliary chart labels,
derive their exact score, and retain the derivative cost of the labels.
The latter is the score counterpart of the atlas IMS debit.

The derivation below is self-contained because a careless use of
conditional Jensen would lose an essential term.  Normalization removes
overlap multiplicity from conditional averaging, but a moving label
kernel has its own Fisher information.

\subsection{The exact score of a normalized label lift}

Let $(X,\mathcal F,\mu)$ be a probability space equipped with a
derivation $D$ on an algebra $\mathcal A$.  Assume that $\mu$ has score
$\beta$ in the convention
\begin{equation}
 \int_X Df\,d\mu=\int_X f\beta\,d\mu,
 \qquad f\in\mathcal A.
 \label{eq:n9v35-base-ibp}
\end{equation}
Let $I$ be a finite label set and let
$\eta_\alpha\in\mathcal A$, $\alpha\in I$, be real functions satisfying
\begin{equation}
 \kappa_\alpha:=\eta_\alpha^2\ge0,
 \qquad
 \sum_{\alpha\in I}\kappa_\alpha=1.
 \label{eq:n9v35-pou}
\end{equation}
The lifted probability measure on $X\times I$ is
\begin{equation}
 \widehat\mu(dx,d\alpha)=\mu(dx)\kappa_\alpha(x).
 \label{eq:n9v35-lift}
\end{equation}
Its $X$-marginal is exactly $\mu$.

\begin{theorem}[Normalized-label score and Fisher identity]
\label{thm:n9v35-label-score}
Suppose \eqref{eq:n9v35-base-ibp} holds and the products used below
belong to the integration-by-parts domain.  On
$\{\kappa_\alpha>0\}$ define
\begin{equation}
 \widehat\beta_\alpha
 :=\beta-D\log\kappa_\alpha
 =\beta-2\frac{D\eta_\alpha}{\eta_\alpha}.
 \label{eq:n9v35-joint-score}
\end{equation}
Then $\widehat\beta$ is the score of the lifted measure:
\begin{equation}
 \sum_\alpha\int_X \kappa_\alpha Df_\alpha\,d\mu
 =\sum_\alpha\int_X
      \kappa_\alpha f_\alpha\widehat\beta_\alpha\,d\mu.
 \label{eq:n9v35-joint-ibp}
\end{equation}
Moreover, the conditional mean and conditional second moment are
\begin{align}
 \sum_\alpha\kappa_\alpha\widehat\beta_\alpha
 &=\beta,
 \label{eq:n9v35-cond-mean}\\
 \sum_\alpha\kappa_\alpha|\widehat\beta_\alpha|^2
 &=|\beta|^2+\mathcal I_\kappa[D],
 \label{eq:n9v35-fisher-pythagoras}\\
 \mathcal I_\kappa[D]
 &:=\sum_{\alpha:\kappa_\alpha>0}
       \frac{|D\kappa_\alpha|^2}{\kappa_\alpha}
   =4\sum_\alpha|D\eta_\alpha|^2.
 \label{eq:n9v35-label-fisher}
\end{align}
The last expression defines the zero-set extension.  Thus the label
lift never lowers Fisher information and its exact debit is
$\mathcal I_\kappa[D]$.
\end{theorem}

\begin{proof}
Apply \eqref{eq:n9v35-base-ibp} to $f_\alpha\kappa_\alpha$:
\[
 \int\kappa_\alpha Df_\alpha\,d\mu
 =\int f_\alpha(\kappa_\alpha\beta-D\kappa_\alpha)\,d\mu.
\]
This is \eqref{eq:n9v35-joint-ibp}.  Differentiating
$\sum_\alpha\kappa_\alpha=1$ gives
$\sum_\alpha D\kappa_\alpha=0$, which proves
\eqref{eq:n9v35-cond-mean}.  Expanding the square in
\eqref{eq:n9v35-joint-score}, the cross term is
\[
 -2\beta\sum_\alpha D\kappa_\alpha=0.
\]
The remaining positive term is
$\sum_\alpha|D\kappa_\alpha|^2/\kappa_\alpha$.
Since $D\kappa_\alpha=2\eta_\alpha D\eta_\alpha$, it equals the
right-hand side of \eqref{eq:n9v35-label-fisher} wherever
$\eta_\alpha\ne0$.  The square-root expression is finite and supplies
the canonical extension across zeros.
\end{proof}

\begin{corollary}[Fixed labels are free; moving labels are not]
\label{cor:n9v35-fixed-label}
If $D\kappa_\alpha=0$ for every $\alpha$, the lift preserves the Fisher
second moment.  Conversely, equality in
\eqref{eq:n9v35-fisher-pythagoras} holds if and only if
$D\eta_\alpha=0$ for every $\alpha$, modulo $\mu$-null sets.
\end{corollary}

\begin{proof}
This follows from the nonnegative sum on the right-hand side of
\eqref{eq:n9v35-label-fisher}.
\end{proof}

\begin{remark}[Score-splitting invariance and label dependence]
\label{rem:n9v35-split}
The sign in \eqref{eq:n9v35-joint-score} is forced by
\eqref{eq:n9v35-base-ibp}.  It is the same bookkeeping principle as the
score-splitting invariance of V28: moving $\kappa_\alpha$ into the
reference density subtracts $D\log\kappa_\alpha$ from the score.  Dropping
this term would change the measure and invalidate the inserted
Schwinger--Dyson identity.
\end{remark}

\subsection{Weighted tangent trace and the atlas IMS debit}

Let $(e_k)_{k\ge1}$ be tangent directions and let $w_k>0$ be the Hilbert
weights of V29--V33.  Write $D_k:=D_{e_k}$ and assume the series below
are measurable.  Define
\begin{align}
 \|\beta\|_{w,*}^2
 &:=\sum_k w_k^{-1}|\beta_k|^2,
 \label{eq:n9v35-base-trace}\\
 \mathfrak I_w(\kappa)
 &:=4\sum_k w_k^{-1}\sum_\alpha|D_k\eta_\alpha|^2.
 \label{eq:n9v35-atlas-debit}
\end{align}

\begin{theorem}[Weighted label-Fisher trace identity]
\label{thm:n9v35-weighted-trace}
For every finite mode cutoff $K$,
\begin{equation}
 \sum_{k\le K}w_k^{-1}
 \sum_\alpha\kappa_\alpha|\widehat\beta_{k,\alpha}|^2
 =\sum_{k\le K}w_k^{-1}|\beta_k|^2
  +4\sum_{k\le K}w_k^{-1}
       \sum_\alpha|D_k\eta_\alpha|^2.
 \label{eq:n9v35-finite-trace}
\end{equation}
Consequently, monotone convergence gives
\begin{equation}
 \mathbb E_{\widehat\mu}
       \|\widehat\beta\|_{w,*}^2
 =\mathbb E_\mu\|\beta\|_{w,*}^2
  +\mathbb E_\mu\mathfrak I_w(\kappa),
 \label{eq:n9v35-full-trace}
\end{equation}
with equality in $[0,\infty]$.  In particular, the lifted score belongs
to $L^2(\widehat\mu;H_w)$ exactly when both terms on the right are
finite.
\end{theorem}

\begin{proof}
Apply \Cref{thm:n9v35-label-score} direction by direction, multiply by
$w_k^{-1}$, and sum to $K$.  All summands are nonnegative, so monotone
convergence permits $K\to\infty$ without an unproved interchange of
conditionally convergent series.
\end{proof}

\begin{proposition}[Coordinate-free finite-rank form]
\label{prop:n9v35-coordinate-free}
Let $P$ be a finite-rank orthogonal projection in the tangent Hilbert
space and let $\nabla_P$ denote the corresponding gradient.  Then
\begin{equation}
 \mathbb E_{\widehat\mu}\|P\widehat\beta\|^2
 =\mathbb E_\mu\|P\beta\|^2
  +4\mathbb E_\mu\sum_\alpha\|\nabla_P\eta_\alpha\|^2.
 \label{eq:n9v35-projector-identity}
\end{equation}
The added term is independent of the orthonormal basis of
$\operatorname{Ran}P$.
\end{proposition}

\begin{proof}
Use \eqref{eq:n9v35-finite-trace} with unit weights in an orthonormal
basis of $\operatorname{Ran}P$.  Both squared norms are invariant under
orthogonal changes of that basis.
\end{proof}

\begin{remark}[Exact relation to an IMS cost]
\label{rem:n9v35-ims}
For a square partition of unity $\sum_\alpha\eta_\alpha^2=1$, the
standard IMS localization error is a positive sum of
$|\nabla\eta_\alpha|^2$.  Equation
\eqref{eq:n9v35-projector-identity} shows that the same quantity, with
the exact factor four fixed by logarithmic differentiation, is the
Fisher debit of adjoining the chart label.  Conditional normalization
therefore removes a crude number-of-charts factor but does not erase the
transition-gradient cost.
\end{remark}

\subsection{Local $L^2$--probability gluing}

The next lemma records precisely the part of the Yang--Mills V73
mechanism that is type-compatible with the metric-score programme.  It
is stated for a local insertion rather than a full-volume density.

\begin{lemma}[Normalized-label Cauchy--Schwarz compiler]
\label{lem:n9v35-local-cs}
Let $E\in\mathcal F$ and let $G_\alpha$ be square-integrable under the
lifted law.  Then
\begin{equation}
 \left|
  \mathbb E_\mu\!left[
   \mathbf1_E\sum_\alpha\kappa_\alpha G_\alpha
  \right]
 \right|
 \le
 \left(
  \mathbb E_\mu\sum_\alpha\kappa_\alpha|G_\alpha|^2
 \right)^{1/2}
 \mu(E)^{1/2}.
 \label{eq:n9v35-local-cs}
\end{equation}
No factor $|I|$ occurs.
\end{lemma}

\begin{proof}
Conditional Jensen gives
$|\sum_\alpha\kappa_\alpha G_\alpha|^2
\le\sum_\alpha\kappa_\alpha|G_\alpha|^2$.
Cauchy--Schwarz between this conditional mean and $\mathbf1_E$ proves
\eqref{eq:n9v35-local-cs}.
\end{proof}

\begin{theorem}[Good-chart/bad-chart score insertion bound]
\label{thm:n9v35-good-bad}
Let $B$ be a bad-chart event with $\mu(B)\le\pi$.  Let
$H_\alpha$ be a local differentiated-score or coarea-Hessian insertion
such that
\begin{equation}
 \mathbb E_\mu\sum_\alpha\kappa_\alpha|H_\alpha|^2\le A^2.
 \label{eq:n9v35-H-L2}
\end{equation}
Then
\begin{equation}
 \left|
  \mathbb E_\mu\!left[
   \mathbf1_B\sum_\alpha\kappa_\alpha H_\alpha
  \right]
 \right|
 \le A\pi^{1/2}.
 \label{eq:n9v35-bad-H}
\end{equation}
If, on $B^c$,
\begin{equation}
 \left|\sum_\alpha\kappa_\alpha H_\alpha\right|\le C
 \quad\text{almost surely},
 \label{eq:n9v35-good-H}
\end{equation}
then
\begin{equation}
 \left|
  \mathbb E_\mu\sum_\alpha\kappa_\alpha H_\alpha
 \right|
 \le C+A\pi^{1/2}.
 \label{eq:n9v35-total-H}
\end{equation}
The estimate remains local: $A$ is the norm of the insertion on its
declared support, not the $L^2$ norm of a full-volume Radon--Nikodym
density.
\end{theorem}

\begin{proof}
Equation \eqref{eq:n9v35-bad-H} is
\Cref{lem:n9v35-local-cs} with $E=B$.  On $B^c$, use
\eqref{eq:n9v35-good-H} and $\mu(B^c)\le1$, then add the two pieces.
\end{proof}

\begin{corollary}[Vanishing bad sector]
\label{cor:n9v35-vanishing-bad}
For a regulator sequence $n$, suppose
$A_n\pi_n^{1/2}\to0$ and the good-sector expectations converge.  Then
the full local insertion expectations have the same limit as their
good-sector parts.
\end{corollary}

\begin{proof}
Subtract the good-sector expectation and use
\eqref{eq:n9v35-bad-H}.
\end{proof}

\begin{remark}[Why rarity alone still does not load Fisher control]
\label{rem:n9v35-rarity}
The factor $A$ in \eqref{eq:n9v35-bad-H} is indispensable.  A random
variable may be arbitrarily large on an event of arbitrarily small
probability.  Thus this theorem respects the V30 no-go result: rare exits
do not by themselves imply an $L^2$ score bound.  Here rarity is paired
with a local $L^2$ bound for the differentiated insertion.
\end{remark}

\subsection{Insertion into the nonlinear constraint programme}

On a regular constraint chart of V34, let $\bar\beta_{k,\alpha}$ be the
complete reduced score, including the pullback action, Lagrangian
constraint-curvature term, and coarea-density term.  Introduce smooth
square labels $\eta_\alpha$ subordinate to the regular charts.  The
joint chart-labelled score is then
\begin{equation}
 \widehat{\bar\beta}_{k,\alpha}
 =\bar\beta_k-2\frac{D_k\eta_\alpha}{\eta_\alpha},
 \label{eq:n9v35-reduced-joint-score}
\end{equation}
where $\bar\beta_k$ denotes the chart-independent score after the
cocycle compatibility of the local descriptions has been proved.

\begin{definition}[Normalized labelled Fisher chart condition (NLFCC)]
\label{def:n9v35-NLFCC}
A cofinal family of reduced constraint measures satisfies NLFCC on a
local support $Q$ when:
\begin{enumerate}
\item the square labels obey \eqref{eq:n9v35-pou} and their chart
cocycles identify the same reduced score on overlaps;
\item the base reduced weighted score trace is uniformly integrable;
\item the label debit
$\mathbb E\mathfrak I_w(\kappa)$ is uniformly finite and its unwanted
high-shell or transition part tends to zero in the declared limit;
\item the bad-chart probability $\pi_n(Q)$ and local insertion norm
$A_n(Q)$ satisfy $A_n(Q)\pi_n(Q)^{1/2}\to0$; and
\item all constants are uniform in the exterior condition required by
the local-to-global cofinal construction.
\end{enumerate}
\end{definition}

\begin{theorem}[NLFCC closes the labelled chart interface]
\label{thm:n9v35-NLFCC}
Assume NLFCC.  Then adjoining the chart label preserves the reduced
configuration marginal, produces a uniformly $L^2$ joint score, and
does not change the limiting local differentiated-score identity.  The
only new Fisher term is the explicit positive debit
\eqref{eq:n9v35-atlas-debit}; the discarded bad-chart contribution
vanishes at rate $A_n(Q)\pi_n(Q)^{1/2}$.
\end{theorem}

\begin{proof}
Marginal preservation is \eqref{eq:n9v35-lift}.  The exact score and its
$L^2$ norm are given by \Cref{thm:n9v35-label-score,
thm:n9v35-weighted-trace}.  NLFCC items 2--3 give uniform joint Fisher
control.  Item 4 and \Cref{cor:n9v35-vanishing-bad} remove the local bad
sector.  The overlap average has no chart-count loss by
\Cref{lem:n9v35-local-cs}.  Uniformity in the exterior condition is
exactly what permits this local conclusion to enter a cofinal
construction.
\end{proof}

\begin{firewall}[Cocycle compatibility is prior to averaging]
\label{fire:n9v35-cocycle}
Equations \eqref{eq:n9v35-cond-mean} and
\eqref{eq:n9v35-reduced-joint-score} do not repair incompatible local
scores.  Before normalized averaging, the chart scores must transform
as descriptions of one logarithmic derivative of one reduced measure.
Otherwise the average is merely a new chosen field, not the score of the
physical marginal.
\end{firewall}

\begin{assessment}[Progress at ninth-series V35]
\label{ass:n9v35-status}
The compulsory companion audit has yielded an exact interface theorem
for constraint-chart localization.  The normalized label kernel removes
overlap multiplicity, while its derivative contributes the positive
Fisher/IMS debit $4\sum|D\eta|^2$.  A local $L^2$ differentiated-score
bound combined with bad-chart probability suppresses the singular
sector with the sharp Cauchy--Schwarz exponent $1/2$.
\end{assessment}

\begin{programme}[Ninth-series V36: quantitative label schedule]
\label{prog:n9v35-V36}
V36 will convert NLFCC item 3 into a quantitative cofinal atlas
condition.  It will determine how transition width, active tangent-mode
trace, chart multiplicity, and bad-chart probability must scale so that
the label Fisher debit and singular-chart insertion both vanish without
using a global full-volume density norm.
\end{programme}

\begin{firewall}[Claim boundary after ninth-series V35]
\label{fire:n9v35-boundary}
V35 does not prove the required physical constraint atlas, cocycle
compatibility, uniform label-gradient bounds, bad-chart rarity, local
coarea-Hessian moment, OS positivity, reconstruction, or the full
SM--Einstein continuum.  It proves the exact accounting identity and the
local interface estimate that those inputs would load.
\end{firewall}

\begin{theorem}[Ninth-series V35 result]
\label{thm:n9v35-result}
For a normalized square chart kernel, the exact labelled score is
\eqref{eq:n9v35-joint-score} and its weighted Fisher trace equals the
unlabelled trace plus the positive atlas debit
\eqref{eq:n9v35-atlas-debit}.  Under NLFCC, normalized conditional
averaging has no chart-count loss and the bad-chart differentiated-score
insertion is at most $A\pi^{1/2}$.
\end{theorem}

\begin{proof}
Use \Cref{thm:n9v35-label-score,thm:n9v35-weighted-trace,
thm:n9v35-good-bad,thm:n9v35-NLFCC}.
\end{proof}

\paragraph{Parent-version record.}
V35 was formed from
\texttt{Renewal\_SM\_Einstein\_Continuum\_Research\_Notes\_V34.tex}.
The V34 parent had $11\,882$ logical source lines, $430\,833$ bytes, and
SHA--256
\[
 \texttt{D69E88B6D183562F04B47A3CC3A265FD8250EB8B3FD5EA74A9FA45A1EE6C3689}.
\]

% =============================================================================
% END APPENDED NINTH-SERIES V35 MATERIAL
% =============================================================================

% =============================================================================
% BEGIN APPENDED NINTH-SERIES V36 MATERIAL
% =============================================================================

\clearpage
\section{Ninth-series V36: quantitative cofinal atlas schedules and
transition-sector alternatives}
\label{sec:v9v36}

\subsection{Purpose and scale variables}

V35 identified the exact label Fisher debit but left its cofinal control
as NLFCC item 3.  This note supplies a quantitative sufficient condition
and an obstruction.  There are two legitimate ways to make a transition
debit small: flatten the partition of unity over a widening transition,
or confine its gradients to a transition sector whose probability tends
to zero.  The two effects may also be combined.  A growing number of
available charts is harmless only when the number simultaneously active
at one configuration remains controlled.

For regulator $n$, let $(e_{n,k})_{1\le k\le K_n}$ be the retained
tangent directions with weights $w_{n,k}>0$.  Let
$\eta_{n,\alpha}$ be a finite square partition of unity,
\begin{equation}
 \sum_{\alpha\in I_n}\eta_{n,\alpha}^2=1,
 \qquad
 \kappa_{n,\alpha}=\eta_{n,\alpha}^2,
 \label{eq:n9v36-square-pou}
\end{equation}
and write $D_{n,k}=D_{e_{n,k}}$.  Its cutoff label debit is
\begin{equation}
 \mathfrak I_{w,n}
 :=4\sum_{k=1}^{K_n}w_{n,k}^{-1}
       \sum_{\alpha\in I_n}|D_{n,k}\eta_{n,\alpha}|^2.
 \label{eq:n9v36-debit}
\end{equation}

\subsection{Active multiplicity rather than atlas cardinality}

\begin{lemma}[Pointwise active-chart bound]
\label{lem:n9v36-active}
Suppose that at each $x$ at most $m_n$ of the functions
$D_{n,k}\eta_{n,\alpha}(x)$ are nonzero.  If, on a measurable transition
set $T_n$,
\begin{equation}
 |D_{n,k}\eta_{n,\alpha}(x)|
 \le \frac{C_n c_{n,k}}{\omega_n}
 \quad (x\in T_n)
 \label{eq:n9v36-gradient-bound}
\end{equation}
and all these derivatives vanish outside $T_n$, then
\begin{equation}
 \mathfrak I_{w,n}(x)
 \le
 \frac{4m_nC_n^2}{\omega_n^2}\,
 \mathcal T_n\mathbf1_{T_n}(x),
 \qquad
 \mathcal T_n:=\sum_{k=1}^{K_n}w_{n,k}^{-1}c_{n,k}^2.
 \label{eq:n9v36-pointwise-debit}
\end{equation}
The bound is independent of $|I_n|$.
\end{lemma}

\begin{proof}
For fixed $k$ and $x$, at most $m_n$ nonzero squares occur in the
inner sum of \eqref{eq:n9v36-debit}.  Bound each by
$C_n^2c_{n,k}^2\omega_n^{-2}$, multiply by $4w_{n,k}^{-1}$, and sum in
$k$.
\end{proof}

\begin{theorem}[Widening/rarity atlas compiler]
\label{thm:n9v36-atlas-compiler}
Assume the hypotheses of \Cref{lem:n9v36-active} and
\begin{equation}
 \mu_n(T_n)\le\tau_n.
 \label{eq:n9v36-transition-prob}
\end{equation}
Then
\begin{equation}
 \mathbb E_{\mu_n}\mathfrak I_{w,n}
 \le
 \Delta_n^{\rm lab}:=
 \frac{4m_nC_n^2\mathcal T_n}{\omega_n^2}\tau_n.
 \label{eq:n9v36-expected-debit}
\end{equation}
Consequently:
\begin{enumerate}
\item uniform boundedness of $\Delta_n^{\rm lab}$ gives uniform control
of the additional labelled Fisher trace;
\item $\Delta_n^{\rm lab}\to0$ makes the label lift asymptotically
Fisher-neutral; and
\item neither $|I_n|\to\infty$ nor $K_n\to\infty$ is by itself an
obstruction, provided the displayed product has the required bound.
\end{enumerate}
\end{theorem}

\begin{proof}
Integrate \eqref{eq:n9v36-pointwise-debit} and use
\eqref{eq:n9v36-transition-prob}.  The conclusions follow from the
weighted trace identity \eqref{eq:n9v35-full-trace}.
\end{proof}

\begin{corollary}[Three sufficient scaling regimes]
\label{cor:n9v36-three-regimes}
Each of the following implies asymptotic Fisher neutrality:
\begin{align}
 &\frac{m_nC_n^2\mathcal T_n}{\omega_n^2}\longrightarrow0
 &&\text{with }\sup_n\tau_n\le1,
 \label{eq:n9v36-wide}\\
 &m_nC_n^2\mathcal T_n=O(1),\quad \tau_n\longrightarrow0
 &&\text{with }\inf_n\omega_n>0,
 \label{eq:n9v36-rare}\\
 &\log(m_nC_n^2\mathcal T_n)-2\log\omega_n+log\tau_n
 \longrightarrow-\infty.
 \label{eq:n9v36-mixed}
\end{align}
The third condition is the general mixed regime.
\end{corollary}

\begin{proof}
Each condition makes the right-hand side of
\eqref{eq:n9v36-expected-debit} tend to zero.
\end{proof}

\begin{remark}[High-shell remainder]
\label{rem:n9v36-high-shell}
If the tangent family is infinite, apply
\Cref{thm:n9v36-atlas-compiler} through $K_n$ and retain
\[
 r_n:=4\mathbb E\sum_{k>K_n}w_k^{-1}
       \sum_\alpha|D_k\eta_{n,\alpha}|^2.
\]
The correct conclusion is then
$\mathbb E\mathfrak I_{w,n}\le\Delta_n^{\rm lab}+r_n$.
Thus $r_n\to0$ is a separate ultraviolet requirement; it cannot be
deduced from a finite-mode transition bound.
\end{remark}

\subsection{A transition-width obstruction}

The widening alternative has a genuine lower bound.  It is not enough
to rename a sharp transition as a smooth chart overlap.

\begin{lemma}[One-dimensional transition lower bound]
\label{lem:n9v36-lower}
Let $\eta\in H^1(0,\omega)$ have traces
$\eta(0)=1$ and $\eta(\omega)=0$.  Then
\begin{equation}
 \int_0^\omega|\eta'(t)|^2\,dt\ge\frac1\omega.
 \label{eq:n9v36-lower}
\end{equation}
Equality holds for the affine transition $\eta(t)=1-t/\omega$.
\end{lemma}

\begin{proof}
The fundamental theorem for $H^1$ functions and Cauchy--Schwarz give
\[
 1=\left|\int_0^\omega\eta'(t)\,dt\right|
 \le\omega^{1/2}
      \left(\int_0^\omega|\eta'(t)|^2\,dt\right)^{1/2}.
\]
Squaring proves the bound.  The affine function saturates it.
\end{proof}

\begin{corollary}[Fixed-width nonvanishing criterion]
\label{cor:n9v36-fixed-width}
Suppose a family of transition fibres carries conditional measures whose
density with respect to $dt$ is bounded below by $\rho>0$, and on each
fibre a label changes from one to zero over length at most
$\omega_*<\infty$.  Then that label contributes at least
$4\rho/\omega_*$ to the integrated one-direction Fisher debit per unit
mass of the family.  In particular, a positive-mass family of such
fibres prevents asymptotic Fisher neutrality.
\end{corollary}

\begin{proof}
Multiply \eqref{eq:n9v36-lower} by the density lower bound and by the
factor four in \eqref{eq:n9v35-label-fisher}, then integrate over the
fibre family.
\end{proof}

\begin{firewall}[What the lower bound does and does not say]
\label{fire:n9v36-lower}
The obstruction concerns a positive-mass family with a uniform density
lower bound.  It does not exclude fixed-width transitions whose
probability vanishes, nor transitions placed in directions assigned
vanishing dual weight.  Those are precisely the rarity and weighted
routes in \eqref{eq:n9v36-expected-debit}.
\end{firewall}

\subsection{Coupling to the bad-chart tail}

Let $B_n(R_n)$ be the singular or bad-chart event at threshold $R_n$.
V26 supplied the coercive conditional tail form
\begin{equation}
 \mu_n(B_n(R_n))
 \le \exp\{\Omega_n-c_nR_n^p\},
 \qquad p>0,quad c_n>0.
 \label{eq:n9v36-coercive-tail}
\end{equation}
Let $A_n(R_n)$ be the local $L^2$ norm of the differentiated-score or
coarea-Hessian insertion from \eqref{eq:n9v35-H-L2}.

\begin{theorem}[Quantitative singular-chart compiler]
\label{thm:n9v36-singular-compiler}
Under \eqref{eq:n9v36-coercive-tail}, the bad-chart insertion is bounded
by
\begin{equation}
 \Delta_n^{\rm bad}
 :=A_n(R_n)
   \exp\!\left\{\frac{\Omega_n-c_nR_n^p}{2}\right\}.
 \label{eq:n9v36-bad-debit}
\end{equation}
It vanishes whenever
\begin{equation}
 \log A_n(R_n)+\frac{\Omega_n}{2}
 -\frac{c_nR_n^p}{2}\longrightarrow-\infty.
 \label{eq:n9v36-bad-condition}
\end{equation}
In particular, if
\begin{equation}
 A_n(R)\le C_nR^s,
 \label{eq:n9v36-polynomial-A}
\end{equation}
it suffices that
\begin{equation}
 \log C_n+s\log R_n+\frac{\Omega_n}{2}
 -\frac{c_nR_n^p}{2}\longrightarrow-\infty.
 \label{eq:n9v36-polynomial-condition}
\end{equation}
More generally, if
$A_n(R)\le C_n\exp\{b_nR^q\}$ with $0\le q<p$, it suffices that
\begin{equation}
 \log C_n+b_nR_n^q+\frac{\Omega_n}{2}
 -\frac{c_nR_n^p}{2}\longrightarrow-\infty.
 \label{eq:n9v36-subcoercive-condition}
\end{equation}
\end{theorem}

\begin{proof}
Insert \eqref{eq:n9v36-coercive-tail} into
\eqref{eq:n9v35-bad-H}.  This gives
\eqref{eq:n9v36-bad-debit}.  Taking logarithms gives
\eqref{eq:n9v36-bad-condition}; the last two statements follow by their
respective upper bounds on $A_n$.
\end{proof}

\begin{remark}[Uniformity cannot be suppressed]
\label{rem:n9v36-uniformity}
The condition $q<p$ alone is insufficient when $b_n/c_n$ or
$\Omega_n/c_n$ diverges too quickly.  The actual diagonal condition is
\eqref{eq:n9v36-subcoercive-condition}.  This is the same uniformity
issue already isolated by the normalized conditional block loader of
V26.
\end{remark}

\subsection{The combined cofinal criterion}

\begin{definition}[Cofinal atlas scale condition (CASC)]
\label{def:n9v36-CASC}
A cofinal constraint atlas satisfies CASC on a declared local support
$Q$ if:
\begin{enumerate}
\item its chart scores satisfy the cocycle requirement of
\Cref{fire:n9v35-cocycle};
\item the unlabelled reduced score has the uniform weighted $L^2$ bound
required by V29--V34;
\item the finite-mode atlas variables satisfy
$\Delta_n^{\rm lab}(Q)+r_n(Q)\to0$;
\item the singular-chart variables satisfy
$\Delta_n^{\rm bad}(Q)\to0$; and
\item the bounds are uniform in admissible exterior data and in the
finite set of cylinder insertions used at that stage of the diagonal.
\end{enumerate}
\end{definition}

\begin{theorem}[CASC loads NLFCC and preserves the local score limit]
\label{thm:n9v36-CASC}
Assume CASC.  Then NLFCC holds on $Q$.  The chart-labelled and
unlabelled local integration-by-parts identities have the same cofinal
limit, and their weighted Fisher traces differ by a quantity tending to
zero.  More precisely, the two additional error scales are bounded by
\begin{equation}
 \mathbb E\mathfrak I_{w,n}
 \le\Delta_n^{\rm lab}+r_n,
 \qquad
 |\mathcal E_n^{\rm bad}|
 \le\Delta_n^{\rm bad}.
 \label{eq:n9v36-two-errors}
\end{equation}
\end{theorem}

\begin{proof}
The first inequality is \Cref{thm:n9v36-atlas-compiler} together with
\Cref{rem:n9v36-high-shell}.  The second is
\Cref{thm:n9v36-singular-compiler}.  CASC items 3--4 make both errors
vanish.  Items 1--2 provide NLFCC items 1--2, and item 5 provides its
required conditional uniformity.  Apply \Cref{thm:n9v35-NLFCC}.
\end{proof}

\begin{corollary}[Diagonal selection]
\label{cor:n9v36-diagonal}
Suppose that for each fixed local support and finite cylinder family one
can choose parameters $(K,\omega,R)$ so that the right-hand sides of
\eqref{eq:n9v36-two-errors} are arbitrarily small, with the estimates
uniform beyond a regulator index depending on those choices.  Then a
diagonal sequence exists on which CASC items 3--5 hold for a countable
exhaustion of local supports and cylinder insertions.
\end{corollary}

\begin{proof}
Enumerate the countable requirements.  At stage $j$, choose the auxiliary
parameters so the first $j$ error pairs are at most $2^{-j}$, and then
choose the regulator index beyond all corresponding uniformity
thresholds.  Increase these indices strictly.  Every fixed requirement
then has errors tending to zero along the selected subsequence.
\end{proof}

\begin{assessment}[Progress at ninth-series V36]
\label{ass:n9v36-status}
The atlas part of the constraint-reduced Fisher programme now has an
explicit scale law.  The controlling label quantity is
$m_nC_n^2\mathcal T_n\tau_n/\omega_n^2$, plus a separately declared
high-shell remainder.  The singular chart is controlled by the actual
competition between its local insertion norm and half of the coercive
tail exponent.  A fixed-width positive-mass transition has a rigorous
lower bound and therefore cannot be declared negligible.
\end{assessment}

\begin{programme}[Ninth-series V37: chart-cocycle construction]
\label{prog:n9v36-V37}
The next bottleneck is upstream of the numerical scale balance: CASC
assumes that local constraint scores are cocycle compatible.  V37 will
derive the score transformation law under a change of reduced chart,
including the coordinate Jacobian, coarea density, and square-label
term, and will state the exact obstruction at singular chart changes.
\end{programme}

\begin{firewall}[Claim boundary after ninth-series V36]
\label{fire:n9v36-boundary}
V36 proves sufficient scale conditions and a transition lower bound.  It
does not prove that the physical SM--Einstein constraint atlas has those
scales, that its chart cocycles close, or that the coercive tail and
coarea-Hessian norm are uniform.  It also does not prove OS positivity,
reconstruction, or the full continuum.
\end{firewall}

\begin{theorem}[Ninth-series V36 result]
\label{thm:n9v36-result}
The cofinal chart-label Fisher debit is bounded by
\eqref{eq:n9v36-expected-debit}, while the bad-chart insertion is bounded
by \eqref{eq:n9v36-bad-debit}.  CASC makes both vanish and thereby loads
NLFCC.  Conversely, a positive-mass family of fixed-width transitions
with a uniform density lower bound has a nonzero label Fisher debit.
\end{theorem}

\begin{proof}
Use \Cref{thm:n9v36-atlas-compiler,lem:n9v36-lower,
thm:n9v36-singular-compiler,thm:n9v36-CASC}.
\end{proof}

\paragraph{Parent-version record.}
V36 was formed from
\texttt{Renewal\_SM\_Einstein\_Continuum\_Research\_Notes\_V35.tex}.
The V35 parent had $12\,325$ logical source lines, $447\,771$ bytes, and
SHA--256
\[
 \texttt{F524BEC7CF475CE54994A83FE7FBFA6B8C254672B82CF40D2CF9FB7043D0600E}.
\]

% =============================================================================
% END APPENDED NINTH-SERIES V36 MATERIAL
% =============================================================================
% =============================================================================
% BEGIN APPENDED NINTH-SERIES V37 MATERIAL
% =============================================================================

\clearpage
\section{Ninth-series V37: exact reduced-chart score cocycles and the
singular-change obstruction}
\label{sec:v9v37}

\subsection{Two descriptions of one constraint carrier}

V36 made cocycle compatibility the first CASC hypothesis.  We now derive
that compatibility rather than postulate it.  Let $C$ be an
$m$-dimensional regular constraint carrier with two $C^2$ embeddings
\begin{equation}
 \iota_\alpha:P_\alpha\longrightarrow C,
 \qquad
 \iota_\gamma:P_\gamma\longrightarrow C.
 \label{eq:n9v37-two-charts}
\end{equation}
On an overlap suppose
\begin{equation}
 q=F(p),\qquad
 \iota_\alpha(p)=\iota_\gamma(F(p)),
 \label{eq:n9v37-transition}
\end{equation}
where $F$ is a $C^2$ diffeomorphism between the coordinate overlap
domains.  Let the geometric constrained reference measure be
\begin{equation}
 dm_C=\rho_CJ_\Phi^{-1}\,d\operatorname{vol}_C.
 \label{eq:n9v37-geometric-measure}
\end{equation}
In the two charts write
\begin{align}
 dm_C&=r_\alpha(p)\,dp=r_\gamma(q)\,dq,
 \label{eq:n9v37-chart-densities}\\
 r_\alpha&=(\rho_CJ_\Phi^{-1})\circ\iota_\alpha
             \sqrt{\det G_\alpha},
 \qquad
 G_\alpha=(D\iota_\alpha)^*D\iota_\alpha,
 \label{eq:n9v37-ra}
\end{align}
and analogously for $\gamma$.  Let
$S_\alpha=S\circ\iota_\alpha$ and
$S_\gamma=S\circ\iota_\gamma$.

\begin{lemma}[Coarea-coordinate density cocycle]
\label{lem:n9v37-density-cocycle}
On the overlap,
\begin{align}
 G_\alpha(p)
 &=DF(p)^T G_\gamma(F(p))DF(p),
 \label{eq:n9v37-Gram-cocycle}\\
 r_\alpha(p)
 &=r_\gamma(F(p))\,|\det DF(p)|,
 \label{eq:n9v37-r-cocycle}\\
 S_\alpha(p)&=S_\gamma(F(p)).
 \label{eq:n9v37-S-cocycle}
\end{align}
Thus the coordinate Jacobian is already part of the reduced density
cocycle; it must not be inserted a second time into the physical coarea
factor.
\end{lemma}

\begin{proof}
Differentiate \eqref{eq:n9v37-transition} to obtain
$D\iota_\alpha=(D\iota_\gamma\circ F)DF$.  Taking Gram matrices gives
\eqref{eq:n9v37-Gram-cocycle}, hence
\[
 \sqrt{\det G_\alpha}
 =\sqrt{\det G_\gamma\circ F}\,|\det DF|.
\]
The functions $\rho_C$ and $J_\Phi$ are geometric scalars on the same
point of $C$, which proves \eqref{eq:n9v37-r-cocycle}.  The action is also
a scalar on $C$, giving \eqref{eq:n9v37-S-cocycle}.
\end{proof}

\subsection{The invariant score is indexed by a vector field}

For a $C^1$ reduced vector field $V_\alpha=v_\alpha^i\partial_{p^i}$
with compact support in the overlap, set
\begin{equation}
 \operatorname{div}_{r_\alpha}V_\alpha
 :=r_\alpha^{-1}\partial_{p^i}(r_\alpha v_\alpha^i),
 \qquad
 \beta_{\alpha,V}
 :=\hbar^{-1}V_\alpha S_\alpha
   -\operatorname{div}_{r_\alpha}V_\alpha.
 \label{eq:n9v37-general-score}
\end{equation}
Its pushforward is
\begin{equation}
 V_\gamma=F_*V_\alpha,
 \qquad
 v_\gamma(F(p))=DF(p)v_\alpha(p).
 \label{eq:n9v37-pushforward}
\end{equation}

\begin{theorem}[Reduced-score cocycle]
\label{thm:n9v37-score-cocycle}
Under \eqref{eq:n9v37-r-cocycle}--\eqref{eq:n9v37-pushforward},
\begin{align}
 (\operatorname{div}_{r_\gamma}V_\gamma)\circ F
 &=\operatorname{div}_{r_\alpha}V_\alpha,
 \label{eq:n9v37-div-cocycle}\\
 \beta_{\gamma,F_*V}\circ F&=\beta_{\alpha,V}.
 \label{eq:n9v37-score-cocycle}
\end{align}
Consequently the integration-by-parts identity
\begin{equation}
 \int V f\,d\mu=\int f\beta_V\,d\mu,
 \qquad
 d\mu=Z^{-1}e^{-S/\hbar}dm_C,
 \label{eq:n9v37-invariant-ibp}
\end{equation}
is independent of the regular reduced chart.
\end{theorem}

\begin{proof}
For a compactly supported test function $g$ on the overlap,
change variables using \eqref{eq:n9v37-r-cocycle}:
\[
 \int g\circ F\,
      \operatorname{div}_{r_\alpha}V_\alpha\,r_\alpha dp
 =-\int V_\alpha(g\circ F)r_\alpha dp
 =-\int V_\gamma g\,r_\gamma dq.
\]
The last expression equals
$\int g\operatorname{div}_{r_\gamma}V_\gamma r_\gamma dq$.
Since both divergences are continuous, their pullbacks agree pointwise,
proving \eqref{eq:n9v37-div-cocycle}.  The chain rule and
\eqref{eq:n9v37-S-cocycle} give
$(V_\gamma S_\gamma)\circ F=V_\alpha S_\alpha$; insert both identities
in \eqref{eq:n9v37-general-score}.  Finally, ordinary integration by
parts against $e^{-S/\hbar}r\,dp$ proves
\eqref{eq:n9v37-invariant-ibp} in either chart.
\end{proof}

\begin{corollary}[Why constant coordinate scores do not transform as
constant coordinate scores]
\label{cor:n9v37-constant}
Let $V_\alpha=h^i\partial_{p^i}$ with constant $h$.  In the $q$ chart its
components are
\begin{equation}
 u(q)=DF(F^{-1}(q))h,
 \label{eq:n9v37-variable-u}
\end{equation}
which are generally nonconstant.  The transformed score is
\begin{equation}
 \beta_{\gamma,u}
 =\hbar^{-1}u^a\partial_{q^a}S_\gamma
  -u^a\partial_{q^a}\log r_\gamma
  -\partial_{q^a}u^a.
 \label{eq:n9v37-u-score}
\end{equation}
The last term is required for the cocycle.  It vanishes for every
constant $h$ only when the relevant pushforward fields have zero
coordinate divergence, as for an affine volume-preserving transition.
\end{corollary}

\begin{proof}
Equation \eqref{eq:n9v37-variable-u} is
\eqref{eq:n9v37-pushforward}.  Expand the weighted divergence in
\eqref{eq:n9v37-general-score}.
\end{proof}

\begin{remark}[The V34 constant-direction formula]
\label{rem:n9v37-v34}
The score \eqref{eq:n9v34-score} is correct in its declared $p$ chart
because the coordinate vector $h$ is constant there.  Under a nonlinear
chart change it must be compared with \eqref{eq:n9v37-u-score}, not with
the constant-vector formula obtained by freezing $u$ at one point.
Freezing would drop $\partial_a u^a$ and manufacture a false cocycle
defect.
\end{remark}

\subsection{The differentiated score and the geometric acceleration row}

Define the complete same-field differentiated score by
\begin{equation}
 \mathcal H_{\alpha,V}:=V_\alpha\beta_{\alpha,V}.
 \label{eq:n9v37-complete-H}
\end{equation}

\begin{theorem}[Differentiated-score cocycle]
\label{thm:n9v37-H-cocycle}
For corresponding vector fields,
\begin{equation}
 \mathcal H_{\gamma,F_*V}\circ F
 =\mathcal H_{\alpha,V}.
 \label{eq:n9v37-H-cocycle}
\end{equation}
In the $q$ chart, its action part is
\begin{equation}
 \hbar^{-1}u(uS_\gamma)
 =\hbar^{-1}\left[
  D^2S_\gamma[u,u]+DS_\gamma[D_u u]
 \right].
 \label{eq:n9v37-acceleration}
\end{equation}
Thus the acceleration term $DS_\gamma[D_u u]$ is the
change-of-coordinate counterpart of the constraint-curvature term in
the pullback Hessian \eqref{eq:n9v34-pullback}.
\end{theorem}

\begin{proof}
Equation \eqref{eq:n9v37-score-cocycle} is an equality of scalar
functions under pullback.  Apply $V_\alpha$ to it and use
$V_\alpha(f\circ F)=(F_*V_\alpha f)\circ F$ to obtain
\eqref{eq:n9v37-H-cocycle}.  Expanding $u(u^a\partial_aS)$ gives
\eqref{eq:n9v37-acceleration}.
\end{proof}

\begin{firewall}[Coordinate Hessian is not the complete differentiated
score]
\label{fire:n9v37-coordinate-Hessian}
The matrix of second partial derivatives of $S_\gamma$ alone is not a
chart tensor.  Equation \eqref{eq:n9v37-acceleration}, together with the
derivative of $\operatorname{div}_{r_\gamma}u$, supplies the terms needed
for the invariant object \eqref{eq:n9v37-complete-H}.  Any relative
Hessian estimate used in CHMC must either estimate this complete object
or be expressed through a declared connection for which the omitted
terms are separately controlled.
\end{firewall}

\subsection{Square labels and coordinate changes commute}

Let a physical label kernel be a family of scalar functions on $C$.
Thus on the overlap
\begin{equation}
 \kappa_{\alpha,\ell}(p)
 =\kappa_{\gamma,\ell}(F(p)),
 \qquad
 \sum_\ell\kappa_{\alpha,\ell}=1.
 \label{eq:n9v37-label-cocycle}
\end{equation}
Define the labelled scores as in V35,
\begin{equation}
 \widehat\beta_{\alpha,V,\ell}
 =\beta_{\alpha,V}-V_\alpha\log\kappa_{\alpha,\ell}.
 \label{eq:n9v37-labelled-score}
\end{equation}

\begin{proposition}[Labelled-score cocycle]
\label{prop:n9v37-labelled-cocycle}
On every positive-label overlap,
\begin{equation}
 \widehat\beta_{\gamma,F_*V,\ell}\circ F
 =\widehat\beta_{\alpha,V,\ell}.
 \label{eq:n9v37-labelled-cocycle}
\end{equation}
The conditional mean over $\ell$ is the unlabelled invariant score, and
the Fisher debit is the scalar quadratic form
\begin{equation}
 \mathcal I_\kappa(V)
 =4\sum_\ell|V\eta_\ell|^2,
 \qquad \kappa_\ell=\eta_\ell^2.
 \label{eq:n9v37-invariant-label-debit}
\end{equation}
\end{proposition}

\begin{proof}
The chain rule applied to \eqref{eq:n9v37-label-cocycle} gives
$(F_*V)\log\kappa_{\gamma,\ell}\circ F
=V\log\kappa_{\alpha,\ell}$.  Combine this with
\eqref{eq:n9v37-score-cocycle}.  The last two statements are
\eqref{eq:n9v35-cond-mean} and \eqref{eq:n9v35-label-fisher}, now written
for the geometric vector field $V$.
\end{proof}

\subsection{Tangent norms and boundary flux}

The weighted trace used in V29--V36 is geometric only if its tangent
quadratic form is transported with the chart.  Let $a_\alpha$ and
$a_\gamma$ be positive cotangent quadratic forms satisfying
\begin{equation}
 a_\alpha(p)(\xi,\xi)
 =a_\gamma(F(p))\bigl(DF(p)^{-T}\xi,DF(p)^{-T}\xi\bigr).
 \label{eq:n9v37-cotangent-metric}
\end{equation}

\begin{proposition}[Fisher norm cocycle]
\label{prop:n9v37-norm-cocycle}
If the score one-form and the cotangent metric are transformed by their
natural pullback laws, then the pointwise squared Fisher norm and its
integral are chart independent.  Assigning unrelated diagonal weights
in the two coordinate systems need not preserve either quantity.
\end{proposition}

\begin{proof}
The score is a linear functional of $V$, hence a cotangent vector.
Equation \eqref{eq:n9v37-cotangent-metric} is precisely the congruence
law that preserves its quadratic norm.  The measures agree by
\eqref{eq:n9v37-r-cocycle}, so the integrals agree.  Unrelated diagonal
matrices do not in general satisfy the congruence law.
\end{proof}

\begin{proposition}[Artificial overlap boundaries]
\label{prop:n9v37-boundary}
If the vector field or the labelled test function is not compactly
supported in a chart overlap, the local integration-by-parts formula has
the boundary contribution
\begin{equation}
 \int_{\partial P_\alpha}
 f\,e^{-S_\alpha/\hbar}r_\alpha
 (v_\alpha\cdot n_\alpha)\,d\sigma_\alpha.
 \label{eq:n9v37-boundary-flux}
\end{equation}
For a subordinate label with $\kappa_{\alpha,\ell}$ vanishing at that
artificial boundary, its labelled flux vanishes.  Otherwise the flux
must be retained and shown to cancel with the adjacent chart under the
orientation and density cocycles.
\end{proposition}

\begin{proof}
This is the divergence theorem applied to
$f e^{-S_\alpha/\hbar}r_\alpha v_\alpha$.  Multiplication by a label that
vanishes at the boundary kills the trace.  Without that vanishing, a
chartwise proof has a nonzero boundary term; cancellation is the only
other way to recover the boundary-free geometric identity.
\end{proof}

\subsection{The singular transition obstruction}

\begin{theorem}[Regular cocycle criterion and singular obstruction]
\label{thm:n9v37-regular-criterion}
The chartwise reduced score, complete differentiated score, square-label
debit, and Fisher norm define common geometric objects on an overlap if:
\begin{enumerate}
\item $F$ is a $C^2$ diffeomorphism with nonzero Jacobian;
\item the reduced density obeys \eqref{eq:n9v37-r-cocycle};
\item vector fields, labels, and tangent quadratic forms obey
\eqref{eq:n9v37-pushforward}, \eqref{eq:n9v37-label-cocycle}, and
\eqref{eq:n9v37-cotangent-metric}; and
\item artificial boundary fluxes vanish or cancel.
\end{enumerate}
If $\det DF=0$, the pointwise logarithmic transformation contains
$D\log|\det DF|$ and the inverse transformation of tangent/cotangent
objects fails.  The regular cocycle proof therefore gives no score
extension through that set.  Such a set must instead be removed by a
tube estimate, resolved by a different regular atlas, or handled by a
separately proved weak logarithmic-derivative theorem.
\end{theorem}

\begin{proof}
Under items 1--4, apply
\Cref{lem:n9v37-density-cocycle,thm:n9v37-score-cocycle,
thm:n9v37-H-cocycle,prop:n9v37-labelled-cocycle,
prop:n9v37-norm-cocycle,prop:n9v37-boundary}.  When $\det DF=0$, the
density cocycle is not invertible and neither the logarithm of its
Jacobian nor $DF^{-T}$ is defined there.  Hence the displayed regular
formulas cannot supply an extension.  The listed alternatives are
additional hypotheses, not consequences of the regular calculation.
\end{proof}

\begin{definition}[Reduced score cocycle card (RSCC)]
\label{def:n9v37-RSCC}
RSCC consists of the four hypotheses of
\Cref{thm:n9v37-regular-criterion}, uniform on the cofinal regular atlas,
together with a declared treatment of the singular transition set.
\end{definition}

\begin{corollary}[RSCC supplies CASC item 1]
\label{cor:n9v37-RSCC-CASC}
RSCC implies the chart-cocycle hypothesis in CASC and NLFCC.  It also
identifies the correct complete differentiated score to which the V33
relative-form estimate must be applied.
\end{corollary}

\begin{proof}
Use \Cref{thm:n9v37-regular-criterion} and the definition of CASC in
\Cref{def:n9v36-CASC}.
\end{proof}

\begin{assessment}[Progress at ninth-series V37]
\label{ass:n9v37-status}
The regular chart-cocycle premise has now been reduced to geometric
identities.  The coordinate Jacobian is exactly the Gram-density
cocycle, a constant direction becomes a variable vector field, and its
coordinate divergence and acceleration rows are mandatory.  Square
labels commute with chart changes when they are scalar kernels on the
constraint carrier.  Singular transitions remain outside this regular
argument.
\end{assessment}

\begin{programme}[Ninth-series V38: singular transition tubes]
\label{prog:n9v37-V38}
V38 will go upstream to the singular-change set.  It will combine the
determinant--coarea tube powers already proved earlier in this ninth
series with the derivative order of
$D\log|\det DF|$ and the local $L^2$--probability compiler, deriving the
precise integrability threshold rather than assuming that small tube
mass controls a singular score.
\end{programme}

\begin{firewall}[Claim boundary after ninth-series V37]
\label{fire:n9v37-boundary}
V37 proves regular chart covariance.  It does not construct the physical
constraint atlas, extend a score through a rank-changing chart, bound
the singular Jacobian derivative, establish the cofinal scale conditions
of V36, or prove OS positivity, reconstruction, or the full continuum.
\end{firewall}

\begin{theorem}[Ninth-series V37 result]
\label{thm:n9v37-result}
On every regular overlap, the reduced score and its complete same-field
derivative obey the exact cocycles
\eqref{eq:n9v37-score-cocycle} and \eqref{eq:n9v37-H-cocycle}.  The
coordinate Jacobian, vector-field divergence, acceleration, square-label
derivative, tangent norm, and boundary flux are the complete ledger.
RSCC loads the cocycle premise of CASC; a rank-changing transition
requires a separate singular theorem.
\end{theorem}

\begin{proof}
Use \Cref{lem:n9v37-density-cocycle,thm:n9v37-score-cocycle,
thm:n9v37-H-cocycle,prop:n9v37-labelled-cocycle,
prop:n9v37-norm-cocycle,thm:n9v37-regular-criterion,
cor:n9v37-RSCC-CASC}.
\end{proof}

\paragraph{Parent-version record.}
V37 was formed from
\texttt{Renewal\_SM\_Einstein\_Continuum\_Research\_Notes\_V36.tex}.
The V36 parent had $12\,698$ logical source lines, $461\,280$ bytes, and
SHA--256
\[
 \texttt{7A45D614DF7B19E14D89B6560E2FEE165CF70BF485A5184EB820F2B4C49DDA75}.
\]

% =============================================================================
% END APPENDED NINTH-SERIES V37 MATERIAL
% =============================================================================
% =============================================================================
% BEGIN APPENDED NINTH-SERIES V38 MATERIAL
% =============================================================================

\clearpage
\section{Ninth-series V38: singular chart-score powers and the sharp
tube threshold}
\label{sec:v9v38}

\subsection{The punctured transverse model}

V37 proved the score cocycle only where the chart transition is a
diffeomorphism.  Near a rank-changing set, the logarithmic Jacobian can
have a pole.  The earlier V9 tube theorem controls the mass of a
transverse tube, but mass alone does not control that pole.  We now join
the two powers.

Let $K$ be compact and let $0<|t|<\rho$.  Assume that the normalized
local law has density $w(x,t)$ satisfying the two-sided power bound
\begin{equation}
 c_0|t|^a\le w(x,t)\le C_0|t|^a,
 \qquad a>-1,
 \label{eq:n9v38-density-power}
\end{equation}
for almost every $(x,t)\in K\times(-\rho,\rho)$, with
$0<c_0\le C_0<\infty$.  In the determinant--coarea notation of V9,
$a=p-q$ and the tube-mass exponent is
\begin{equation}
 \lambda:=a+1=p-q+1.
 \label{eq:n9v38-lambda}
\end{equation}
The lower bound in \eqref{eq:n9v38-density-power} is needed only for the
necessity statements; the upper bound suffices for all loaders.

Let $J_F$ denote the coordinate-transition Jacobian, distinct from the
normal coarea Jacobian $J_\Phi$.  Suppose on one punctured side that
\begin{equation}
 J_F(x,t)=|t|^d b(x,t),
 \qquad d>0,qquad 0<b_-\le b\le b_+,
 \label{eq:n9v38-J-factor}
\end{equation}
where $b$ has bounded first and second logarithmic derivatives in $t$.

\begin{lemma}[Logarithmic Jacobian orders]
\label{lem:n9v38-log-orders}
Under \eqref{eq:n9v38-J-factor},
\begin{align}
 \partial_t\log J_F
 &=\frac d t+g_0(x,t),
 \label{eq:n9v38-first-pole}\\
 \partial_t^2\log J_F
 &=-\frac d{t^2}+g_1(x,t),
 \label{eq:n9v38-second-pole}
\end{align}
where $g_0,g_1$ are bounded.  The order $d$ changes the coefficient, not
the pole order.
\end{lemma}

\begin{proof}
On each punctured side,
$\log J_F=d\log|t|+\log b$.  Differentiate once and twice and use the
assumed logarithmic derivative bounds for $b$.
\end{proof}

The complete chart score may contain cancellations between the
coordinate Jacobian, vector-field divergence, coarea density, action,
and label rows.  We therefore state the result for the net leading
coefficient.  Let
\begin{align}
 \beta(x,t)&=\frac c t+B(x,t),
 \label{eq:n9v38-net-score}\\
 H(x,t)&=-\frac c{t^2}+R(x,t),
 \label{eq:n9v38-net-H}
\end{align}
where $B,R$ are bounded and $c$ is the uncancelled coefficient.  For an
exact differentiated score, $H=\partial_t\beta$ and bounded
$\partial_tB$ gives this form.

\subsection{Sharp moment matching}

\begin{theorem}[Singular score moment threshold]
\label{thm:n9v38-score-threshold}
Let $r>0$.  Under the upper bound in
\eqref{eq:n9v38-density-power},
\begin{equation}
 \beta\in L^r_{\rm loc}(\mu)quad\text{whenever}\quad
 \lambda>r.
 \label{eq:n9v38-score-sufficient}
\end{equation}
If $c\ne0$ and the lower density bound also holds, this condition is
sharp:
\begin{equation}
 \beta\in L^r_{\rm loc}(\mu)quad\Longleftrightarrow\quad
 \lambda>r.
 \label{eq:n9v38-score-iff}
\end{equation}
At $\lambda=r$ the singular moment diverges logarithmically.
\end{theorem}

\begin{proof}
Boundedness of $B$ gives
$|\beta|^r\le C_r(|c|^r|t|^{-r}+1)$ for a constant $C_r$ valid also
when $0<r<1$.  Multiplication by the upper density bound reduces local
integrability to
\[
 \int_0^\rho t^{a-r}\,dt<\infty,
\]
which is equivalent to $a-r>-1$, or $\lambda>r$.  If $c\ne0$, choose
$\rho_0$ so small that $|B|\le |c|/(2|t|)$ for
$0<|t|<\rho_0$.  Then $|\beta|\ge |c|/(2|t|)$.  The lower density bound
reduces necessity to the same power integral.  Equality of the exponent
to $-1$ gives logarithmic divergence.
\end{proof}

\begin{theorem}[Differentiated-score moment threshold]
\label{thm:n9v38-H-threshold}
Let $r>0$.  Under the upper density bound,
\begin{equation}
 H\in L^{r/2}_{\rm loc}(\mu)quad\text{whenever}\quad
 \lambda>r.
 \label{eq:n9v38-H-sufficient}
\end{equation}
If $c\ne0$ and the lower density bound holds, the condition is sharp.
In particular, the requirements
\begin{equation}
 \beta\in L^r,qquad H\in L^{r/2}
 \label{eq:n9v38-matched-requirements}
\end{equation}
have the same transverse threshold $\lambda>r$.
\end{theorem}

\begin{proof}
The leading power of $|H|^{r/2}$ is $|t|^{-r}$.  Repeat the proof of
\Cref{thm:n9v38-score-threshold} with $H$ and exponent $r/2$.
\end{proof}

\begin{corollary}[Fisher and CHMC thresholds]
\label{cor:n9v38-Fisher-CHMC}
An uncancelled simple transverse score pole has:
\begin{enumerate}
\item finite Fisher moment $\beta\in L^2$ and integrable differentiated
score $H\in L^1$ exactly when $\lambda>2$;
\item the V31 higher score/Hessian pair
$\beta\in L^r$, $H\in L^{r/2}$, $r>2$, exactly when $\lambda>r$; and
\item logarithmic failure at the respective equality threshold.
\end{enumerate}
For the V9 power balance this means $p-q+1>r$.
\end{corollary}

\begin{proof}
Set $r=2$ or the declared $r>2$ in
\Cref{thm:n9v38-score-threshold,thm:n9v38-H-threshold} and use
\eqref{eq:n9v38-lambda}.
\end{proof}

\begin{remark}[Mass integrability is much weaker]
\label{rem:n9v38-mass-weaker}
V9 requires only $\lambda>0$ for finite tube mass.  An $L^2$ score with
an uncancelled simple pole requires $\lambda>2$.  Therefore a crossing
may be harmless for normalization and still be fatal for Fisher
closability.
\end{remark}

\subsection{Quantitative excision}

Let $\mathcal T_\varepsilon=K\times\{|t|<\varepsilon\}$.

\begin{proposition}[Sharp tube-moment debit]
\label{prop:n9v38-tube-moment}
If $\lambda>r$ and only the upper hypotheses above are assumed, then for
$0<\varepsilon<\rho$,
\begin{align}
 \int_{\mathcal T_\varepsilon}|\beta|^r\,d\mu
 &\le C\bigl(\varepsilon^{\lambda-r}
              +\varepsilon^\lambda\bigr),
 \label{eq:n9v38-beta-tube}\\
 \int_{\mathcal T_\varepsilon}|H|^{r/2}\,d\mu
 &\le C\bigl(\varepsilon^{\lambda-r}
              +\varepsilon^\lambda\bigr),
 \label{eq:n9v38-H-tube}
\end{align}
where $C$ is independent of $\varepsilon$.  Hence both singular moment
debits vanish at leading rate $O(\varepsilon^{\lambda-r})$.
\end{proposition}

\begin{proof}
The proof of the two threshold theorems gives integrands bounded by a
constant times $|t|^{a-r}+|t|^a$.  Integrating from zero to
$\varepsilon$ yields the two powers because
$a-r+1=\lambda-r$ and $a+1=\lambda$.
\end{proof}

\begin{corollary}[Summable singular-score tube schedule]
\label{cor:n9v38-summable}
If $\lambda>r$ and
\begin{equation}
 \sum_n\varepsilon_n^{\lambda-r}<\infty,
 \label{eq:n9v38-summable-condition}
\end{equation}
with uniform local constants, then the adjacent singular score and
differentiated-score moment debits are summable.  For
$\varepsilon_n=(n+1)^{-s}$ it suffices that
$s(\lambda-r)>1$; a geometric schedule works for every
$\lambda>r$.
\end{corollary}

\begin{proof}
Apply \Cref{prop:n9v38-tube-moment} to the two adjacent measures and
sum.  The elementary series criteria give the final statement.
\end{proof}

\begin{proposition}[Direct power integration versus the local
$L^2$--probability compiler]
\label{prop:n9v38-direct-vs-cs}
For the Fisher identity one needs the $L^1$ tube debit of $H$.  Direct
power integration requires only $\lambda>2$.  Using the V35
Cauchy--Schwarz compiler with $G=H$ would require $H\in L^2$, hence
$\lambda>4$, unless an independent bounded local $L^2$ insertion norm is
available.  Thus the V35 compiler is robust under abstract information,
whereas the present direct calculation is sharper when the transverse
power law is proved.
\end{proposition}

\begin{proof}
The direct claim is \Cref{thm:n9v38-H-threshold} with $r=2$.
The $L^2$ condition for $H\sim |t|^{-2}$ is
$\int_0^\rho t^{a-4}dt<\infty$, equivalent to
$a+1=\lambda>4$.
\end{proof}

\subsection{Higher codimension}

The same calculation exposes the role of codimension.  Let
$z\in\mathbb R^c$, $r_z=|z|$, and assume a local density comparable to
$r_z^a$ with $a+c>0$.

\begin{proposition}[Codimension-$c$ threshold]
\label{prop:n9v38-codimension}
Suppose the net score has leading order $r_z^{-\ell}$ and consider its
$s$th moment.  Then the local moment is finite exactly when
\begin{equation}
 a+c>\ell s,
 \label{eq:n9v38-codim-threshold}
\end{equation}
under corresponding two-sided bounds and a nonzero leading coefficient.
The tube-mass exponent is $\lambda_c=a+c$, and the tube-moment debit is
$O(\varepsilon^{\lambda_c-\ell s})$ above threshold.
\end{proposition}

\begin{proof}
Polar coordinates supply volume element
$r_z^{c-1}dr_zd\theta$.  The radial power is
$r_z^{a+c-1-\ell s}$, which is integrable at zero exactly when its
exponent exceeds $-1$.  Integration to $\varepsilon$ gives the stated
debit.
\end{proof}

\subsection{Cancellation and normalized labels}

\begin{proposition}[Net-coefficient rule]
\label{prop:n9v38-net-coefficient}
Suppose finitely many rows of the complete score have expansions
\begin{equation}
 \beta_j(x,t)=\frac{c_j}{t}+B_j(x,t),
 \qquad B_j\in L^\infty.
 \label{eq:n9v38-row-expansion}
\end{equation}
Then the threshold theorems apply with
$c=\sum_jc_j$.  If $\sum_jc_j=0$, they yield no obstruction; the next
nonzero asymptotic order must be computed.  Verifying only one row, such
as $D\log J_F$, cannot establish the net pole.
\end{proposition}

\begin{proof}
Sum \eqref{eq:n9v38-row-expansion}.  The bounded remainders remain
bounded and the leading coefficients add.
\end{proof}

\begin{proposition}[Labels cannot cure a non-$L^2$ marginal score]
\label{prop:n9v38-label-no-cure}
For every normalized square-label lift,
\begin{equation}
 \mathbb E_\mu|\beta|^2
 \le\mathbb E_{\widehat\mu}|\widehat\beta|^2.
 \label{eq:n9v38-label-monotone}
\end{equation}
Hence if an uncancelled score pole has $\lambda\le2$, no choice of
normalized labels can make the joint score square-integrable while
preserving the same marginal.
\end{proposition}

\begin{proof}
This is the nonnegativity of the label debit in
\eqref{eq:n9v35-fisher-pythagoras}.  Combine it with
\Cref{thm:n9v38-score-threshold} at $r=2$.
\end{proof}

\begin{definition}[Singular chart threshold card (SCTC)]
\label{def:n9v38-SCTC}
For target score moment $r\ge2$, SCTC records:
\begin{enumerate}
\item a common transverse chart and a uniform density upper power
$a=p-q$;
\item the complete score ledger and its net first-pole coefficient;
\item either cancellation of that coefficient with a proved next-order
bound, or the strict threshold $\lambda=p-q+1>r$;
\item the matching differentiated-score bound of order $|t|^{-2}$;
\item a tube schedule satisfying
\eqref{eq:n9v38-summable-condition}; and
\item regular-complement RSCC and CASC estimates uniform in the tube
radius.
\end{enumerate}
\end{definition}

\begin{theorem}[SCTC closes the singular chart-score interface]
\label{thm:n9v38-SCTC}
SCTC supplies summable $L^r$ score and $L^{r/2}$ differentiated-score
debits on the excised tubes and the regular score cocycle on their
complements.  It therefore closes the singular-chart row of RSCC for the
declared local insertions and target moment.
\end{theorem}

\begin{proof}
Use \Cref{prop:n9v38-net-coefficient,prop:n9v38-tube-moment,
cor:n9v38-summable} on the tube.  Use V37 RSCC on the regular complement.
Uniformity in SCTC item 6 permits the two regions to be combined.
\end{proof}

\begin{assessment}[Progress at ninth-series V38]
\label{ass:n9v38-status}
The singular chart problem now has a sharp exponent test.  A tube-mass
power $\lambda$ controls an uncancelled simple score pole in $L^r$ and
its differentiated pole in $L^{r/2}$ exactly when $\lambda>r$.
Normalization needs only $\lambda>0$, so it is not a substitute for
Fisher control.  Cancellations must be checked in the complete score
ledger before this obstruction is applied.
\end{assessment}

\begin{programme}[Ninth-series V39: physical net-pole ledger]
\label{prog:n9v38-V39}
V39 will compute the net transverse coefficient for a constrained
finite-dimensional model with action, Gram determinant, normal coarea
Jacobian, coordinate-transition Jacobian, and square label all present.
The aim is to identify which cancellations are geometric identities and
which require a physical determinant or counterterm input.
\end{programme}

\begin{firewall}[Claim boundary after ninth-series V38]
\label{fire:n9v38-boundary}
V38 proves the abstract sharp threshold.  It does not identify the
physical SM--Einstein transverse exponent, prove a cancellation of its
net pole, establish uniform two-sided chart density bounds, or complete
RSCC/CASC.  OS positivity, reconstruction, and the full continuum remain
open.
\end{firewall}

\begin{theorem}[Ninth-series V38 result]
\label{thm:n9v38-result}
If a rank-changing tube has mass exponent $\lambda=p-q+1$ and the
complete reduced score has a nonzero simple transverse pole, then
$\beta\in L^r$ and $D\beta\in L^{r/2}$ hold precisely above the strict
threshold $\lambda>r$.  Above threshold, the excised moment debit is
$O(\varepsilon^{\lambda-r})$ and admits a summable cofinal schedule.
\end{theorem}

\begin{proof}
Use \Cref{thm:n9v38-score-threshold,thm:n9v38-H-threshold,
prop:n9v38-tube-moment,cor:n9v38-summable,
prop:n9v38-net-coefficient}.
\end{proof}

\paragraph{Parent-version record.}
V38 was formed from
\texttt{Renewal\_SM\_Einstein\_Continuum\_Research\_Notes\_V37.tex}.
The V37 parent had $13\,131$ logical source lines, $477\,083$ bytes, and
SHA--256
\[
 \texttt{2825A500121B467F7608214148ADBDBE3FF1D7440276DF37D1ACD826AF22F975}.
\]

% =============================================================================
% END APPENDED NINTH-SERIES V38 MATERIAL
% =============================================================================
% =============================================================================
% BEGIN APPENDED NINTH-SERIES V39 MATERIAL
% =============================================================================

\clearpage
\section{Ninth-series V39: the complete transverse pole ledger and
vector-field order}
\label{sec:v9v39}

\subsection{A factorized constrained density}

V38 treated a net simple score pole.  To use that result one must first
compute the net coefficient without double-counting the coordinate
Jacobian.  Consider a punctured transverse chart
$K\times\{0<|t|<\rho\}$ and a reduced law
\begin{equation}
 d\mu=Z^{-1}e^{-S(t,x)/\hbar}r(t,x)\,dt\,dx.
 \label{eq:n9v39-law}
\end{equation}
Suppose the positive reference density has the factorization
\begin{equation}
 r=\rho_{\rm amb}\,
    \mathcal D\,
    \sqrt{\det G}\,
    J_\Phi^{-1}\,
    \mathcal C,
 \label{eq:n9v39-factorization}
\end{equation}
where $\mathcal D$ denotes the owned positive determinant modulus,
$G$ the reduced Gram matrix, $J_\Phi$ the normal coarea Jacobian, and
$\mathcal C$ any separately owned positive counterterm/reference factor.
Assume uniform punctured asymptotics
\begin{align}
 \rho_{\rm amb}&=|t|^u\rho_0,
 &\mathcal D&=|t|^pD_0,
 &\sqrt{\det G}&=|t|^gG_0,
 \label{eq:n9v39-numerator-orders}\\
 J_\Phi&=|t|^qJ_0,
 &\mathcal C&=|t|^kC_0,
 &S&=s\log|t|+S_0,
 \label{eq:n9v39-denominator-action}
\end{align}
where all displayed zero-order factors are bounded above and below by
positive constants and have bounded first two logarithmic transverse
derivatives.  Set
\begin{equation}
 a:=u+p+g-q+k,
 \qquad
 \lambda:=a+1.
 \label{eq:n9v39-a-lambda}
\end{equation}
Local normalization requires $\lambda>0$ when the action logarithm is
kept in $S$.  Equivalently, the actual total density has exponent
$a-s/\hbar$; below we retain the action row separately because that is
the score convention used throughout V28--V38.

\subsection{General transverse vector fields}

Let
\begin{equation}
 V=v(t,x)\partial_t,
 \qquad
 v(t,x)=v_0(x)|t|^\nu(1+O(|t|)),
 \label{eq:n9v39-vector-order}
\end{equation}
on one punctured side, with $v_0$ bounded away from zero and with the
corresponding differentiated asymptotics.  The constant-coordinate case
of V34 and V38 is $\nu=0$.

\begin{theorem}[Complete first-pole coefficient]
\label{thm:n9v39-first-coefficient}
For the score convention
\begin{equation}
 \beta_V=\hbar^{-1}VS-\operatorname{div}_rV,
 \qquad
 \operatorname{div}_rV=v\partial_t\log r+\partial_tv,
 \label{eq:n9v39-score}
\end{equation}
one has
\begin{equation}
 \beta_V
 =v_0(x)c_V|t|^{\nu-1}+O(|t|^\nu)+O(|t|^{\nu-1}|t|),
 \label{eq:n9v39-score-asymptotic}
\end{equation}
on either fixed punctured side, where
\begin{equation}
 \boxed{c_V=\frac{s}{\hbar}-(u+p+g-q+k)-\nu
       =\frac{s}{\hbar}-a-\nu.}
 \label{eq:n9v39-net-coefficient}
\end{equation}
Thus the vector-field divergence contributes the indispensable row
$-\nu$.
\end{theorem}

\begin{proof}
The factorization and logarithmic derivative hypotheses give
\[
 \partial_t\log r=\frac a t+O(1),
 \qquad
 \partial_tS=\frac s t+O(1),
 \qquad
 \partial_tv=\frac\nu t v+O(|t|^\nu).
\]
Insert these three expansions into \eqref{eq:n9v39-score}.  The
coefficient multiplying $v/t$ is $s/\hbar-a-\nu$, which proves the
claim.
\end{proof}

\begin{corollary}[Row-by-row ownership table]
\label{cor:n9v39-ownership-table}
The leading coefficient in \eqref{eq:n9v39-net-coefficient} has the
following additive rows:
\begin{center}
\begin{tabular}{c|c}
row & contribution to $c_V$\\ \hline
action logarithm $s\log|t|$ & $s/\hbar$\\
ambient/reference density $|t|^u$ & $-u$\\
owned positive determinant $|t|^p$ & $-p$\\
Gram density $|t|^g$ & $-g$\\
normal coarea denominator $|t|^{-q}$ & $+q$\\
positive counterterm factor $|t|^k$ & $-k$\\
vector-field order $|t|^\nu$ & $-\nu$
\end{tabular}
\end{center}
Every physical factor must appear exactly once, either in the action or
in the reference density.
\end{corollary}

\begin{proof}
Expand \eqref{eq:n9v39-net-coefficient}.
\end{proof}

\begin{proposition}[Score-splitting invariance at a crossing]
\label{prop:n9v39-split-invariance}
Move a factor $|t|^z f_0(t,x)$ from the reference density into the
effective action by
\begin{equation}
 r_{\rm new}=r(|t|^zf_0)^{-1},
 \qquad
 S_{\rm new}=S-\hbar\log(|t|^zf_0).
 \label{eq:n9v39-split-move}
\end{equation}
Then the probability law, the complete score, and the coefficient
$c_V$ are unchanged.
\end{proposition}

\begin{proof}
The density identity
$e^{-S_{\rm new}/\hbar}r_{\rm new}=e^{-S/\hbar}r$ is immediate.
The action coefficient changes from $s/\hbar$ to
$s/\hbar-z$, while the reference exponent changes from $a$ to $a-z$.
Their difference in \eqref{eq:n9v39-net-coefficient} is unchanged.  The
bounded factor $f_0$ cancels by the same logarithmic differentiation.
\end{proof}

\subsection{The vector-order moment theorem}

For moment statements it is cleaner to absorb the action logarithm into
the actual probability exponent.  Put
\begin{equation}
 a_{\rm prob}:=a-\frac{s}{\hbar},
 \qquad
 \lambda_{\rm prob}:=a_{\rm prob}+1,
 \label{eq:n9v39-probability-exponent}
\end{equation}
and assume $\lambda_{\rm prob}>0$.  Notice that
$c_V=1-\lambda_{\rm prob}-\nu$.

\begin{theorem}[Vector-order score threshold]
\label{thm:n9v39-vector-threshold}
Assume $c_V\ne0$ and two-sided power bounds for the actual probability
density $|t|^{a_{\rm prob}}$.  For every $r>0$,
\begin{equation}
 \beta_V\in L^r_{\rm loc}(\mu)
 \quad\Longleftrightarrow\quad
 \boxed{\lambda_{\rm prob}>r(1-\nu).}
 \label{eq:n9v39-vector-threshold}
\end{equation}
At equality the moment diverges logarithmically.  If $\nu\ge1$, the
right-hand side is nonpositive and local normalization already implies
every finite score moment allowed by the bounded remainder hypotheses.
\end{theorem}

\begin{proof}
The leading score magnitude is comparable to $|t|^{\nu-1}$.  Its
$r$th moment has transverse power
$a_{\rm prob}+r(\nu-1)$.  Integrability is equivalent to this exponent
being greater than $-1$, which rearranges to
\eqref{eq:n9v39-vector-threshold}.  Equality gives the power $-1$.
\end{proof}

\begin{theorem}[Matched differentiated-score threshold]
\label{thm:n9v39-vector-H}
Suppose in addition that the leading derivative does not vanish, namely
$c_V(\nu-1)\ne0$.  Then
\begin{equation}
 V\beta_V
 =v_0(x)^2c_V(\nu-1)|t|^{2\nu-2}\bigl(1+o(1)\bigr),
 \label{eq:n9v39-H-asymptotic}
\end{equation}
and
\begin{equation}
 V\beta_V\in L^{r/2}_{\rm loc}(\mu)
 \quad\Longleftrightarrow\quad
 \lambda_{\rm prob}>r(1-\nu).
 \label{eq:n9v39-H-threshold}
\end{equation}
Thus the $L^r$ score and $L^{r/2}$ differentiated-score thresholds still
match for a general vector-field order.  When $\nu=1$ or $c_V=0$, the
displayed leading differentiated row cancels and the next order must be
computed.
\end{theorem}

\begin{proof}
Apply $V=v\partial_t$ to the leading term
$v_0c_V|t|^{\nu-1}$.  This gives
\eqref{eq:n9v39-H-asymptotic}.  Raising it to power $r/2$ again produces
$|t|^{r(\nu-1)}$, so the same one-dimensional power test applies.
\end{proof}

\begin{corollary}[Constant and linearly vanishing fields]
\label{cor:n9v39-special-fields}
For $\nu=0$, the V38 condition
$\lambda_{\rm prob}>r$ is recovered.  For $\nu=1$, the score is bounded
to leading order and its leading same-field derivative vanishes; local
normalization $\lambda_{\rm prob}>0$ is sufficient at this order.
\end{corollary}

\begin{proof}
Insert $\nu=0$ and $\nu=1$ in
\eqref{eq:n9v39-vector-threshold} and
\eqref{eq:n9v39-H-asymptotic}.
\end{proof}

\subsection{Where the coordinate Jacobian belongs}

Suppose $q=F(t)$ on the punctured overlap and
\begin{equation}
 F'(t)=f_0|t|^d(1+O(|t|)),
 \qquad f_0\ne0.
 \label{eq:n9v39-coordinate-order}
\end{equation}
If the same constraint embedding is written as
$\iota_t(t)=\iota_q(F(t))$, then
\begin{equation}
 \sqrt{\det G_t(t)}
 =\sqrt{\det G_q(F(t))}\,|F'(t)|
 \label{eq:n9v39-Gram-coordinate}
\end{equation}
in the one-dimensional transverse block.

\begin{proposition}[No Gram/Jacobian double count]
\label{prop:n9v39-no-double-count}
The coordinate-transition order $d$ in
\eqref{eq:n9v39-coordinate-order} is already included in the Gram order
$g_t$ through \eqref{eq:n9v39-Gram-coordinate}.  Adding a further
$|F'|$ factor to \eqref{eq:n9v39-factorization} would shift the score
coefficient by a spurious $-d$.
\end{proposition}

\begin{proof}
Equation \eqref{eq:n9v39-Gram-coordinate} is the one-dimensional
determinant form of \eqref{eq:n9v37-Gram-cocycle}.  Hence the coordinate
Lebesgue Jacobian is precisely the factor by which the Gram volume
density changes.  Counting it again adds $d$ to $a$ and therefore
subtracts $d$ a second time in \eqref{eq:n9v39-net-coefficient}.
\end{proof}

\begin{proposition}[Punctured cocycle check]
\label{prop:n9v39-punctured-check}
Assume the $q$-coordinate density and action are regular and nonzero at
$q=0$, while \eqref{eq:n9v39-coordinate-order} holds.  For the constant
$t$-field $\partial_t$, the pushed field is
$F'(t)\partial_q$.  Its $q$-coordinate divergence has leading term
$d/t$, and therefore
\begin{equation}
 \beta_{q,F'\partial_q}\circ F
 =-\frac d t+O(1)
 =\beta_{t,\partial_t}.
 \label{eq:n9v39-punctured-cocycle}
\end{equation}
The pole is the same geometric score evaluated on the same vector field;
it is not a contradiction between charts.
\end{proposition}

\begin{proof}
In the $t$ chart the density contains $|F'|$, so its logarithmic
derivative is $d/t+O(1)$.  In the $q$ chart,
$\partial_qF'(t(q))=d/t+O(1)$ because division by $F'$ cancels the
$t^d$ factor.  Apply \eqref{eq:n9v37-u-score}.
\end{proof}

\subsection{Square labels at a singular point}

Suppose a labelled kernel has on one side
\begin{equation}
 \kappa_\ell(t,x)=|t|^{2b_\ell}k_{0,\ell}(t,x),
 \qquad b_\ell\ge0,
 \label{eq:n9v39-label-order}
\end{equation}
with a positive zero-order factor whenever the label is present.

\begin{proposition}[Labelled exponent and the nonvanishing-label row]
\label{prop:n9v39-label-orders}
For label $\ell$, the joint probability exponent is
$a_{\rm prob}+2b_\ell$ and the labelled score coefficient is
\begin{equation}
 c_{V,\ell}=c_V-2b_\ell.
 \label{eq:n9v39-labelled-coefficient}
\end{equation}
If $\sum_\ell\kappa_\ell=1$ up to $t=0$, at least one active label has
$b_\ell=0$.  Thus vanishing labels can improve their own joint
integrability, but normalized labelling cannot remove the marginal
threshold of an uncancelled base score.
\end{proposition}

\begin{proof}
Multiplication by $\kappa_\ell$ adds $2b_\ell$ to the joint density
power.  The labelled score subtracts
$V\log\kappa_\ell=2b_\ell v/t+O(v)$, proving
\eqref{eq:n9v39-labelled-coefficient}.  If every $b_\ell>0$, every label
tends to zero at $t=0$, contradicting their normalized sum.  The final
claim also follows from \eqref{eq:n9v38-label-monotone}.
\end{proof}

\subsection{The complete pole ledger card}

\begin{definition}[Transverse pole ownership card (TPOC)]
\label{def:n9v39-TPOC}
For each physical rank-changing sector, TPOC records:
\begin{enumerate}
\item one common transverse variable and the orders
$(u,p,g,q,k,s)$ in \eqref{eq:n9v39-numerator-orders}--
\eqref{eq:n9v39-denominator-action};
\item a proof that every positive factor is owned exactly once, with the
Gram/coordinate rule of \Cref{prop:n9v39-no-double-count};
\item the order $\nu$ of every tested geometric tangent field;
\item the complete coefficient $c_V$ in
\eqref{eq:n9v39-net-coefficient} and every claimed cancellation;
\item the actual probability exponent $\lambda_{\rm prob}$ and the
strict moment condition \eqref{eq:n9v39-vector-threshold}; and
\item uniform zero-order bounds and a summable tube schedule for the
first nonzero asymptotic row.
\end{enumerate}
\end{definition}

\begin{theorem}[TPOC loads SCTC]
\label{thm:n9v39-TPOC}
If TPOC holds for target moment $r\ge2$, then the singular score and
differentiated-score rows satisfy SCTC for the declared vector fields.
If $c_V\ne0$, the required strict condition is
\begin{equation}
 \lambda_{\rm prob}>r(1-\nu).
 \label{eq:n9v39-TPOC-threshold}
\end{equation}
If $c_V=0$, TPOC must provide the next nonzero order before SCTC can be
claimed.
\end{theorem}

\begin{proof}
Use \Cref{thm:n9v39-first-coefficient,
thm:n9v39-vector-threshold,thm:n9v39-vector-H} to identify the power and
moment threshold, followed by the V38 tube integration and summability
argument.  The no-double-count and uniformity clauses provide the
ownership and constant hypotheses of SCTC.
\end{proof}

\begin{assessment}[Progress at ninth-series V39]
\label{ass:n9v39-status}
The singular threshold can now be computed from a complete owned ledger.
The action, determinant, Gram volume, coarea denominator, counterterm,
and vector-field divergence contribute additively to the net pole.
The vector-field order changes the moment threshold from
$\lambda>r$ to $\lambda>r(1-\nu)$, and score splitting leaves the result
unchanged.  Coordinate Jacobians are already present in the Gram density.
\end{assessment}

\begin{programme}[Ninth-series V40: compulsory companion audit]
\label{prog:n9v39-V40}
V40 will inspect the newest Yang--Mills and Navier--Stokes notes before
choosing the next physical TPOC row.  It will continue to V41 after the
audit, as required by the five-version protocol.
\end{programme}

\begin{firewall}[Claim boundary after ninth-series V39]
\label{fire:n9v39-boundary}
V39 proves the finite-dimensional pole ledger and its vector-order
threshold.  It does not identify the physical orders for the
SM--Einstein constraint map, prove a net cancellation, establish the
uniform tube hypotheses, or complete RSCC/CASC.  OS positivity,
reconstruction, and the full continuum remain open.
\end{firewall}

\begin{theorem}[Ninth-series V39 result]
\label{thm:n9v39-result}
For the factorized constrained density
\eqref{eq:n9v39-factorization} and a transverse field of order $\nu$,
the complete pole coefficient is \eqref{eq:n9v39-net-coefficient} and
the sharp score/Hessian moment threshold is
\eqref{eq:n9v39-vector-threshold}.  This ledger is invariant under
action/reference score splitting and counts the coordinate Jacobian
exactly once through the Gram density.
\end{theorem}

\begin{proof}
Use \Cref{thm:n9v39-first-coefficient,prop:n9v39-split-invariance,
thm:n9v39-vector-threshold,thm:n9v39-vector-H,
prop:n9v39-no-double-count,thm:n9v39-TPOC}.
\end{proof}

\paragraph{Parent-version record.}
V39 was formed from
\texttt{Renewal\_SM\_Einstein\_Continuum\_Research\_Notes\_V38.tex}.
The V38 parent had $13\,514$ logical source lines, $490\,855$ bytes, and
SHA--256
\[
 \texttt{D3DCCE5BCBF2EB81C3EC62020FD28A36064C9289FA1C7FBC087B6A4410A48243}.
\]

% =============================================================================
% END APPENDED NINTH-SERIES V39 MATERIAL
% =============================================================================
% =============================================================================
% BEGIN APPENDED NINTH-SERIES V40 MATERIAL
% =============================================================================

\clearpage
\section{Ninth-series V40: compulsory companion audit and the
repeated-insertion obstruction}
\label{sec:v9v40}

\subsection{Audit record}

This is the compulsory audit announced in
\Cref{prog:n9v39-V40}.  The newest Yang--Mills files were
\begin{align*}
 &\texttt{Renewal\_Yang\_Mills\_Mass\_Gap\_Research\_Notes\_V77.tex},\\
 &\texttt{Renewal\_Yang\_Mills\_Mass\_Gap\_Research\_Notes\_V78.tex}.
\end{align*}
They were byte-identical and contained substantive material only through
V77.  Each had $34\,181$ logical source lines, $1\,220\,556$ bytes, and
SHA--256
\[
 \texttt{18E18847220951E2AC2FEE7F27823402D8D0E0AA8E6D35BB3A21A2B2DC77BB30}.
\]
Thus V78 was not counted as independent evidence.  The newest
Navier--Stokes file remained
\texttt{Renewal\_NS\_Stability\_Research\_Notes\_V26.tex}, with
$5\,319$ lines, $278\,283$ bytes, and SHA--256
\[
 \texttt{82C887F8C2C69E4F28FAD7C3DC8DE5B9099CB5A4BF431EBA1FFF3E91935978AD}.
\]

The substantive Yang--Mills V74--V77 increment establishes the following
structural points.
\begin{enumerate}
\item V74 transfers a local reference fourth moment to a physical mixed
$L^2$ activity and keeps the estimate support-local.
\item V75 proves compact Haar Fisher--Hessian identities and polynomial
cutoff bounds for every fixed finite derivative word.
\item V76 enumerates the local vertex types and isolates a defect: the
second-order nature of one generator does not bound the number of its
occurrences in a Duhamel expansion.
\item V77 repairs that multiplicity at the resolvent level by
energy-normalizing a relatively bounded symmetric form and resumming it
through an inverse.  Weighted locality still requires a separate
energy-factor estimate.
\end{enumerate}
The Navier--Stokes increment provides no new type-compatible input for
the present constraint-score row.

The transferable conclusion is not a Yang--Mills estimate.  It is the
logical distinction between a fixed score moment and arbitrary repeated
insertions.  V31--V39 provide finite Fisher/Hessian moments under explicit
conditions; those moments alone do not justify an infinite Duhamel or
resolvent series.

\subsection{Finite moments do not control multiplicity}

\begin{proposition}[Any fixed moment leaves a higher insertion
uncontrolled]
\label{prop:n9v40-fixed-moment}
For every $r>0$ and every integer $N>r$, there is a nonnegative random
variable $X$ such that
\begin{equation}
 \mathbb E X^r<\infty,
 \qquad
 \mathbb E X^N=\infty.
 \label{eq:n9v40-fixed-counterexample}
\end{equation}
Consequently an $L^r$ score estimate cannot by itself bound a Duhamel
word containing $N$ copies of that score.
\end{proposition}

\begin{proof}
Choose $\alpha$ with $r<\alpha<N$ and let $X$ have Pareto density
$\alpha x^{-\alpha-1}\mathbf1_{[1,\infty)}(x)$.  Direct integration gives
$\mathbb E X^s=\alpha/(\alpha-s)$ for $s<\alpha$ and infinity for
$s\ge\alpha$.
\end{proof}

\begin{theorem}[All polynomial moments still need not sum]
\label{thm:n9v40-lognormal}
There exists a positive random variable with every finite polynomial
moment for which
\begin{equation}
 \sum_{n=0}^\infty\frac{t^n}{n!}\mathbb E X^n=\infty
 \qquad\text{for every }t>0.
 \label{eq:n9v40-divergent-series}
\end{equation}
Hence even a hierarchy of finite moments, without quantitative growth
control, does not justify exchanging expectation with an exponential or
arbitrarily long insertion series.
\end{theorem}

\begin{proof}
Let $X=e^Y$ with $Y\sim N(0,\sigma^2)$ and $\sigma>0$.  Then
\begin{equation}
 \mathbb E X^n=e^{\sigma^2n^2/2}.
 \label{eq:n9v40-lognormal-moments}
\end{equation}
For the $n$th summand $a_n=t^ne^{\sigma^2n^2/2}/n!$,
\[
 \frac{a_{n+1}}{a_n}
 =\frac{t}{n+1}e^{\sigma^2(2n+1)/2}\longrightarrow\infty.
\]
The summands therefore fail even to tend to zero.
\end{proof}

\begin{corollary}[CHMC is a fixed-order loader]
\label{cor:n9v40-CHMC-fixed}
The CHMC conclusion of V31 and its singular-chart refinements in
V38--V39 control the declared finite score/Hessian moments.  They do not,
without an additional growth or resummation theorem, control arbitrary
same-block multiplicity.
\end{corollary}

\begin{proof}
The statements of CHMC fix a moment exponent $r$.  Apply
\Cref{prop:n9v40-fixed-moment}; even estimates for all fixed exponents do
not supply uniform growth by \Cref{thm:n9v40-lognormal}.
\end{proof}

\subsection{Two valid scalar repairs}

The scalar analogues clarify the required input before the operator
resummation is developed in V41.

\begin{proposition}[Factorial moment criterion]
\label{prop:n9v40-factorial}
Suppose constants $C,R>0$ satisfy
\begin{equation}
 \mathbb E|X|^n\le C n!R^n
 \qquad(n\ge0).
 \label{eq:n9v40-factorial-bound}
\end{equation}
Then for $|t|<R^{-1}$,
\begin{equation}
 \sum_{n=0}^\infty\frac{|t|^n}{n!}\mathbb E|X|^n
 \le\frac{C}{1-|t|R}.
 \label{eq:n9v40-factorial-sum}
\end{equation}
In particular Tonelli justifies termwise expectation of the absolute
exponential series.
\end{proposition}

\begin{proof}
Insert \eqref{eq:n9v40-factorial-bound}; the resulting majorant is the
geometric series $C\sum_n(|t|R)^n$.
\end{proof}

\begin{proposition}[Exponential-integrability criterion]
\label{prop:n9v40-exponential}
If $\mathbb E e^{t_0|X|}<\infty$ for some $t_0>0$, then for
$0\le t<t_0$,
\begin{equation}
 \sum_{n=0}^\infty\frac{t^n}{n!}\mathbb E|X|^n
 =\mathbb E e^{t|X|}<\infty.
 \label{eq:n9v40-exponential-sum}
\end{equation}
Moreover
$\mathbb E|X|^n\le n!t_0^{-n}\mathbb E e^{t_0|X|}$.
\end{proposition}

\begin{proof}
Tonelli applies to the nonnegative exponential series.  The moment bound
follows from $(t_0|X|)^n/n!\le e^{t_0|X|}$.
\end{proof}

\begin{remark}[Why the operator repair is preferable here]
\label{rem:n9v40-operator-repair}
For an unbounded Einstein/constraint Hessian, exponential integrability
of its pointwise norm is substantially stronger than the relative-form
control already isolated in V33--V34.  Energy-normalized form
resummation uses that existing type of information directly and avoids
assigning one independent probability moment to every occurrence.
\end{remark}

\subsection{Audit decision}

\begin{definition}[Repeated Einstein-insertion interface (REII)]
\label{def:n9v40-REII}
REII asks for one of the following, on the actual reduced carrier:
\begin{enumerate}
\item a quantitative moment-growth estimate strong enough for the
intended insertion series; or
\item a symmetric relative-form representation whose energy-normalized
operator has norm strictly below one, together with the locality weight
needed by the cofinal construction.
\end{enumerate}
Fixed CHMC moments alone do not constitute REII.
\end{definition}

\begin{theorem}[V40 audit verdict]
\label{thm:n9v40-audit}
The V35--V39 constraint-score programme should next pursue REII item 2.
The V33 relative-Hessian card is type-compatible with an
energy-normalized form inverse, whereas the companion notes provide no
physical SM--Einstein exponential-moment input for REII item 1.  The
transfer is conditional on proving the relative form on the
constraint-reduced tangent carrier and does not import the Yang--Mills
weighted locality estimate.
\end{theorem}

\begin{proof}
The insufficiency of fixed moments is
\Cref{prop:n9v40-fixed-moment,thm:n9v40-lognormal}.  The available SM
input is the relative-form architecture of V33--V34.  The audited V77
argument identifies the compatible functional-analytic repair, while
its own firewall leaves its weighted energy-factor estimate unproved for
the full Yang--Mills form; it therefore cannot be copied as a physical
SM estimate.
\end{proof}

\begin{assessment}[Progress at ninth-series V40]
\label{ass:n9v40-status}
The compulsory audit is complete.  It has exposed a new scope boundary:
finite Fisher/Hessian moments load fixed differentiated identities but
not arbitrary repeated insertions.  A Pareto example defeats any one
fixed moment, and a lognormal example defeats unquantified all-order
moments.  The selected repair is energy-normalized relative-form
resummation on the reduced carrier.
\end{assessment}

\begin{programme}[Ninth-series V41: reduced form-resolvent normal form]
\label{prog:n9v40-V41}
V41 will construct the self-adjoint reduced Einstein/constraint form
sum from a positive reference form and a strict relative perturbation,
derive its exact energy-normalized resolvent, and show how arbitrary
insertion multiplicity is replaced by one inverse factor.  It will keep
the weighted locality estimate as a separate loader.
\end{programme}

\begin{firewall}[Claim boundary after ninth-series V40]
\label{fire:n9v40-boundary}
V40 proves a multiplicity no-go and selects a repair.  It does not prove
the physical reduced relative-form bound, a weighted locality estimate,
semigroup or OS reconstruction, or the full SM--Einstein continuum.
\end{firewall}

\begin{theorem}[Ninth-series V40 result]
\label{thm:n9v40-result}
Fixed $L^r$ score/Hessian control cannot justify arbitrary same-block
insertion multiplicity, and even finiteness of every polynomial moment
does not suffice without growth control.  REII identifies factorial or
exponential moment growth and strict energy-normalized form resummation
as valid repairs; the V40 audit selects the latter for the next
SM--Einstein step.
\end{theorem}

\begin{proof}
Use \Cref{prop:n9v40-fixed-moment,thm:n9v40-lognormal,
prop:n9v40-factorial,prop:n9v40-exponential,thm:n9v40-audit}.
\end{proof}

\paragraph{Parent-version record.}
V40 was formed from
\texttt{Renewal\_SM\_Einstein\_Continuum\_Research\_Notes\_V39.tex}.
The V39 parent had $13\,936$ logical source lines, $505\,736$ bytes, and
SHA--256
\[
 \texttt{43F7A0ADC32CF5A6695F1E35E2180DA3E7631B60358A4A0527690EF2ECFB9CCC}.
\]

% =============================================================================
% END APPENDED NINTH-SERIES V40 MATERIAL
% =============================================================================
% =============================================================================
% BEGIN APPENDED NINTH-SERIES V41 MATERIAL
% =============================================================================

\clearpage
\section{Ninth-series V41: constraint-reduced relative-form resummation
and the essential-supremum barrier}
\label{sec:v9v41}

\subsection{Energy normalization of a symmetric form}

Let $\mathcal K$ be a complex Hilbert space.  Let $A_0$ be a nonnegative
self-adjoint operator with closed quadratic form
\begin{equation}
 a_0[u,v]=\langle A_0^{1/2}u,A_0^{1/2}v\rangle,
 \qquad
 \mathcal Q:=D(A_0^{1/2}).
 \label{eq:n9v41-reference-form}
\end{equation}
Let $k$ be a symmetric sesquilinear form on $\mathcal Q$ satisfying
\begin{equation}
 |k[u,u]|\le\eta a_0[u,u]+c\|u\|^2,
 \qquad 0\le\eta<1,quad c\ge0.
 \label{eq:n9v41-relative-bound}
\end{equation}
For $\sigma>0$, put
\begin{equation}
 A_\sigma:=A_0+\sigma I,
 \qquad
 a_\sigma[u,v]:=a_0[u,v]+\sigma\langle u,v\rangle.
 \label{eq:n9v41-shift}
\end{equation}

\begin{lemma}[Strict shifted relative coefficient]
\label{lem:n9v41-shifted-bound}
For every $u\in\mathcal Q$,
\begin{equation}
 |k[u,u]|\le\theta_\sigma a_\sigma[u,u],
 \qquad
 \theta_\sigma:=\max\{\eta,c/\sigma\}.
 \label{eq:n9v41-theta}
\end{equation}
In particular $\theta_\sigma<1$ whenever $\sigma>c$.
\end{lemma}

\begin{proof}
Use
$\eta a_0[u,u]+c\|u\|^2
\le\max\{\eta,c/\sigma\}
(a_0[u,u]+\sigma\|u\|^2)$.
\end{proof}

\begin{theorem}[Energy-normalized form representation]
\label{thm:n9v41-energy-normalization}
Assume $\sigma>c$.  There is a unique bounded self-adjoint operator
$T_\sigma$ on $\mathcal K$ such that
\begin{equation}
 k[u,v]
 =\left\langle A_\sigma^{1/2}u,
 T_\sigma A_\sigma^{1/2}v\right\rangle,
 \qquad u,v\in\mathcal Q,
 \label{eq:n9v41-T-representation}
\end{equation}
and
\begin{equation}
 \|T_\sigma\|\le\theta_\sigma<1.
 \label{eq:n9v41-T-bound}
\end{equation}
The form $h=a_0+k$ is closed and lower semibounded.  If $H$ is its
associated self-adjoint operator, then in form sense
\begin{equation}
 H+\sigma I
 =A_\sigma^{1/2}(I+T_\sigma)A_\sigma^{1/2}.
 \label{eq:n9v41-form-factorization}
\end{equation}
\end{theorem}

\begin{proof}
Equip $\mathcal Q$ with the energy inner product $a_\sigma$.  The
diagonal estimate \eqref{eq:n9v41-theta}, polarization, and symmetry show
that $k$ is a bounded symmetric form on this Hilbert space with norm at
most $\theta_\sigma$.  The map
$U_\sigma:\mathcal Q\to\mathcal K$,
$U_\sigma u=A_\sigma^{1/2}u$, is unitary because
$A_\sigma\ge\sigma I$ makes $A_\sigma^{-1/2}$ bounded and onto
$\mathcal Q$.  Riesz representation after transport by $U_\sigma$
gives the unique self-adjoint $T_\sigma$ and
\eqref{eq:n9v41-T-bound}.  Since
$I+T_\sigma\ge(1-\theta_\sigma)I$, the sum form is closed and lower
semibounded.  The first representation theorem gives $H$ and
\eqref{eq:n9v41-form-factorization}.
\end{proof}

\begin{corollary}[Exact resolvent and spectral comparison]
\label{cor:n9v41-resolvent}
Under the hypotheses of \Cref{thm:n9v41-energy-normalization},
\begin{align}
 (H+\sigma I)^{-1}
 &=A_\sigma^{-1/2}(I+T_\sigma)^{-1}A_\sigma^{-1/2},
 \label{eq:n9v41-resolvent}\\
 \|(I+T_\sigma)^{-1}\|
 &\le(1-\theta_\sigma)^{-1},
 \label{eq:n9v41-inverse-bound}\\
 (1-\theta_\sigma)a_\sigma[u,u]
 &\le h[u,u]+\sigma\|u\|^2
 \le(1+\theta_\sigma)a_\sigma[u,u].
 \label{eq:n9v41-form-comparison}
\end{align}
\end{corollary}

\begin{proof}
The strict norm bound makes $I+T_\sigma$ positive and boundedly
invertible.  Invert \eqref{eq:n9v41-form-factorization} through the
bounded maps $A_\sigma^{-1/2}$.  The inverse norm and form comparison
follow from the spectral interval
$\sigma(T_\sigma)\subset[-\theta_\sigma,\theta_\sigma]$.
\end{proof}

\subsection{Arbitrary multiplicity is one inverse}

\begin{theorem}[Norm-convergent insertion resummation]
\label{thm:n9v41-Neumann}
The inverse in \eqref{eq:n9v41-resolvent} has the norm-convergent series
\begin{equation}
 (I+T_\sigma)^{-1}
 =\sum_{n=0}^\infty(-T_\sigma)^n,
 \label{eq:n9v41-Neumann}
\end{equation}
and the remainder after order $N$ satisfies
\begin{equation}
 \left\|(I+T_\sigma)^{-1}
 -\sum_{n=0}^N(-T_\sigma)^n\right\|
 \le\frac{\theta_\sigma^{N+1}}{1-\theta_\sigma}.
 \label{eq:n9v41-remainder}
\end{equation}
Thus arbitrarily many occurrences of the symmetric perturbation cost
one factor $(1-\theta_\sigma)^{-1}$ rather than one new probability
moment per occurrence.
\end{theorem}

\begin{proof}
This is the Banach-space geometric series because
$\|T_\sigma\|<1$.  Summing the tail of the norm majorant gives
\eqref{eq:n9v41-remainder}.
\end{proof}

\begin{proposition}[Sandwiched observable bound]
\label{prop:n9v41-sandwich}
Let the middle product below denote the bounded extension supplied by
\eqref{eq:n9v41-resolvent}.  For bounded operators $B,C$ on $\mathcal K$,
\begin{equation}
 \|BA_\sigma^{1/2}(H+\sigma)^{-1}A_\sigma^{1/2}C\|
 \le\frac{\|B\|\|C\|}{1-\theta_\sigma}.
 \label{eq:n9v41-sandwich}
\end{equation}
More generally, if $BA_\sigma^{-1/2}$ and
$A_\sigma^{-1/2}C$ are bounded, the unsandwiched resolvent obeys the
corresponding product bound.
\end{proposition}

\begin{proof}
The middle energy-normalized resolvent is
$(I+T_\sigma)^{-1}$ by \eqref{eq:n9v41-resolvent}.  Apply
\eqref{eq:n9v41-inverse-bound} and submultiplicativity.
\end{proof}

\subsection{Application to the constraint-reduced Einstein Hessian}

Let $\mathcal K=L^2(\bar\mu;\mathcal H_P)$ be the direct-integral
reduced tangent space from V29--V39.  Let $A_P(x)$ be the positive
reference tangent operator and let $K_C(x)$ denote the complete
constraint-reduced Hessian, including the ambient, constraint-curvature,
coarea-density, chart-divergence, and label rows identified in
V34--V39.  On a common form core define
\begin{align}
 a_0[U]&:=\int\langle U(x),A_P(x)U(x)\rangle,d\bar\mu(x),
 \label{eq:n9v41-direct-a}\\
 k_C[U]&:=\int\langle U(x),K_C(x)U(x)\rangle,d\bar\mu(x).
 \label{eq:n9v41-direct-k}
\end{align}

\begin{definition}[Constraint-reduced form-resummation card (CRFRC)]
\label{def:n9v41-CRFRC}
CRFRC consists of:
\begin{enumerate}
\item the physical support and reduced tangent space of CFPC;
\item a positive closed reference form $a_0$;
\item a complete symmetric form $k_C$ containing every V34--V39 row;
\item constants $\eta<1$ and $c<\infty$, uniform in the regulator and
admissible exterior data, for which
\eqref{eq:n9v41-relative-bound} holds;
\item a shift $\sigma>c$ uniform in the same parameters; and
\item a separate weighted locality estimate for the energy-normalized
operator $T_\sigma$ whenever spatial decay is required.
\end{enumerate}
\end{definition}

\begin{theorem}[CRFRC closes the repeated symmetric insertion row]
\label{thm:n9v41-CRFRC}
Under CRFRC, the constraint-reduced form sum defines a self-adjoint
operator and its shifted resolvent is
\eqref{eq:n9v41-resolvent}.  All repeated occurrences of $k_C$ in the
energy-normalized resolvent are resummed with uniform cost at most
$(1-\theta_\sigma)^{-1}$.  This proves REII item 2 at the unweighted
resolvent level.
\end{theorem}

\begin{proof}
Apply \Cref{thm:n9v41-energy-normalization,
cor:n9v41-resolvent,thm:n9v41-Neumann} to
\eqref{eq:n9v41-direct-a}--\eqref{eq:n9v41-direct-k}.  Uniformity of the
CRFRC constants gives the regulator-uniform inverse bound.  The statement
is unweighted because CRFRC item 6 has not yet been loaded.
\end{proof}

\begin{firewall}[Resolvent resummation is not a semigroup or OS theorem]
\label{fire:n9v41-resolvent-only}
The identity \eqref{eq:n9v41-resolvent} controls a shifted resolvent.
Passing it to time-local semigroups, normalized Euclidean correlations,
or Osterwalder--Schrader reconstruction requires additional functional
calculus, locality, normalization, and reflection-positivity inputs.
None follows from the Neumann series alone.
\end{firewall}

\subsection{Why the V33 moment card is not yet CRFRC}

V33 used the pointwise relative inequality
\begin{equation}
 |\langle h,K_C(x)h\rangle|
 \le\kappa(x)\langle h,A_P(x)h\rangle
 \label{eq:n9v41-random-relative}
\end{equation}
with finite moments of the random coefficient $\kappa$.  This is enough
for the fixed-moment loaders proved there, but strict form resummation
asks a different question.

\begin{theorem}[Direct-integral essential-supremum criterion]
\label{thm:n9v41-essential-sup}
Suppose \eqref{eq:n9v41-random-relative} holds.  Then
\begin{equation}
 |k_C[U]|
 \le\|\kappa\|_{L^\infty(\bar\mu)}a_0[U]
 \label{eq:n9v41-Linfty-sufficient}
\end{equation}
whenever $\kappa\in L^\infty$.  Conversely, in the scalar model
$A_P=I$ and $K_C=\kappa I$ with $\kappa\ge0$, the optimal relative-form
constant is exactly
\begin{equation}
 \sup_{U\ne0}\frac{k_C[U]}{a_0[U]}
 =\operatorname*{ess\,sup}\kappa.
 \label{eq:n9v41-optimal-essential}
\end{equation}
Thus no finite $L^q$ moment, by itself, implies a strict relative-form
coefficient.
\end{theorem}

\begin{proof}
Integrate \eqref{eq:n9v41-random-relative} and take the essential
supremum to obtain \eqref{eq:n9v41-Linfty-sufficient}.  In the scalar
model the upper bound is immediate.  For every
$M<\operatorname*{ess\,sup}\kappa$, the set
$E_M=\{\kappa>M\}$ has positive measure.  Take
$U=\mathbf1_{E_M}/\bar\mu(E_M)^{1/2}$.  Then $a_0[U]=1$ and
$k_C[U]>M$.  Let $M$ increase to the essential supremum.
\end{proof}

\begin{corollary}[All finite moments do not give a form bound]
\label{cor:n9v41-all-moments}
There exists $\kappa$ with
$\kappa\in\bigcap_{q<\infty}L^q(\bar\mu)$ for which the scalar form above
has no finite relative bound.  One example is a lognormal random
variable.
\end{corollary}

\begin{proof}
A lognormal variable has all finite moments by
\eqref{eq:n9v40-lognormal-moments} but has infinite essential supremum.
Apply \eqref{eq:n9v41-optimal-essential}.
\end{proof}

\begin{proposition}[Rare bad sets do not improve operator norm]
\label{prop:n9v41-rarity-no-operator}
Let $E$ have arbitrarily small positive probability and let
$\kappa\ge M$ on $E$.  In the scalar direct-integral model, the
restriction of the multiplication form to functions supported on $E$
has relative constant at least $M$, independently of $\bar\mu(E)$.
Therefore the V35 $L^2$--probability compiler cannot replace a strict
relative-form bound.
\end{proposition}

\begin{proof}
Use the normalized test function
$U=\mathbf1_E/\bar\mu(E)^{1/2}$.  The probability cancels from the form
ratio, leaving a value at least $M$.
\end{proof}

\begin{remark}[Available repair choices]
\label{rem:n9v41-repairs}
The essential-supremum barrier can be addressed only by additional
structure: prove a uniform strict bound for the complete physical form;
absorb the large positive or indefinite row into a new reference
operator; prove a different operator-ideal resummation; or restrict the
claimed result to a genuinely invariant subspace.  Merely excising a
rare configuration set inside an $L^2$ operator norm is insufficient.
\end{remark}

\subsection{The next loader}

\begin{definition}[Uniform reduced relative-form card (URRFC)]
\label{def:n9v41-URRFC}
URRFC is CRFRC items 1--5 with a proved strict coefficient
$\theta_\sigma<1$ for the complete reduced form on the actual
direct-integral carrier.  Random $L^q$ control of a pointwise coefficient
does not count as URRFC unless an additional theorem converts it to that
form bound.
\end{definition}

\begin{corollary}[URRFC loads unweighted REII]
\label{cor:n9v41-URRFC}
URRFC proves the unweighted repeated-insertion interface, with inverse
cost $(1-\theta_\sigma)^{-1}$ and remainder
\eqref{eq:n9v41-remainder}.
\end{corollary}

\begin{proof}
This is \Cref{thm:n9v41-CRFRC} without its still-separate locality row.
\end{proof}

\begin{assessment}[Progress at ninth-series V41]
\label{ass:n9v41-status}
The repeated symmetric Hessian row now has an exact functional-analytic
loader.  A strict shifted relative form produces a self-adjoint sum and
resums all multiplicities into one energy-normalized inverse.  The
analysis also exposes a new upstream gap: the V33 random relative-form
moments do not imply the required direct-integral operator bound, even
when all finite moments exist.
\end{assessment}

\begin{programme}[Ninth-series V42: reference absorption and sign split]
\label{prog:n9v41-V42}
V42 will attack URRFC upstream.  It will split the complete reduced
Hessian into a positive absorbable part and a residual signed part,
prove the exact new-reference coefficient, and identify when a large
positive constraint/coarea row improves the reference rather than being
treated as a perturbation.  It will retain the conformal and singular
support firewalls.
\end{programme}

\begin{firewall}[Claim boundary after ninth-series V41]
\label{fire:n9v41-boundary}
V41 proves the abstract form-resolvent theorem and the
essential-supremum obstruction.  It does not prove URRFC for the physical
SM--Einstein Hessian, weighted locality, semigroup transfer, reflection
positivity, reconstruction, or the full continuum.
\end{firewall}

\begin{theorem}[Ninth-series V41 result]
\label{thm:n9v41-result}
A symmetric perturbation with shifted relative coefficient
$\theta_\sigma<1$ admits the exact energy-normalized resolvent
\eqref{eq:n9v41-resolvent} and resums arbitrary multiplicity with cost
$(1-\theta_\sigma)^{-1}$.  On the reduced direct-integral carrier,
finite moments of a random pointwise relative coefficient do not imply
this strict bound; the scalar optimal constant is its essential
supremum.
\end{theorem}

\begin{proof}
Use \Cref{thm:n9v41-energy-normalization,cor:n9v41-resolvent,
thm:n9v41-Neumann,thm:n9v41-essential-sup,
prop:n9v41-rarity-no-operator}.
\end{proof}

\paragraph{Parent-version record.}
V41 was formed from
\texttt{Renewal\_SM\_Einstein\_Continuum\_Research\_Notes\_V40.tex}.
The V40 parent had $14\,201$ logical source lines, $515\,816$ bytes, and
SHA--256
\[
 \texttt{B5DBEE59196C4E11BB531FC5BE467AF0D14B89EC8FBD269F665BC3D2557AD87E}.
\]

% =============================================================================
% END APPENDED NINTH-SERIES V41 MATERIAL
% =============================================================================
% =============================================================================
% BEGIN APPENDED NINTH-SERIES V42 MATERIAL
% =============================================================================

\clearpage
\section{Ninth-series V42: positive-row absorption and the residual
relative-form criterion}
\label{sec:v9v42}

\subsection{Abstract absorption theorem}

V41 showed that probability moments do not provide the strict form bound
needed for resolvent resummation.  A large Hessian row is nevertheless
harmless when it is nonnegative and is owned by the reference energy.
The correct coefficient then measures the residual relative to the
enlarged reference, not the complete Hessian relative to the old one.

Let $a$ and $p$ be densely defined closed nonnegative forms on a Hilbert
space $\mathcal K$.  Assume their common domain
\begin{equation}
 \mathcal Q_b:=D(a)\cap D(p)
 \label{eq:n9v42-common-domain}
\end{equation}
is dense, and put
\begin{equation}
 b[u,v]:=a[u,v]+p[u,v].
 \label{eq:n9v42-absorbed-reference}
\end{equation}
The form $b$ is closed on $\mathcal Q_b$.  Let $r$ be a symmetric form on
$\mathcal Q_b$ satisfying
\begin{equation}
 |r[u,u]|
 \le\alpha a[u,u]+\beta p[u,u]+c\|u\|^2,
 \qquad 0\le\alpha,\beta<1,quad c\ge0.
 \label{eq:n9v42-two-row-bound}
\end{equation}

\begin{theorem}[Positive-row absorption]
\label{thm:n9v42-absorption}
For $\sigma>0$, define
\begin{equation}
 b_\sigma[u,u]=b[u,u]+\sigma\|u\|^2,
 \qquad
 \theta_\sigma:=\max\{\alpha,\beta,c/\sigma\}.
 \label{eq:n9v42-absorbed-theta}
\end{equation}
Then
\begin{equation}
 |r[u,u]|\le\theta_\sigma b_\sigma[u,u].
 \label{eq:n9v42-residual-relative}
\end{equation}
If $\sigma>c$, then $\theta_\sigma<1$, the form
$h=b+r$ is closed and lower semibounded, and its shifted resolvent admits
the V41 energy-normalized representation with $b_\sigma$ as reference.
All residual insertion multiplicities are resummed with cost at most
$(1-\theta_\sigma)^{-1}$.
\end{theorem}

\begin{proof}
The intersection domain is complete in the norm
$b[u,u]+\|u\|^2$: a sequence Cauchy in this norm is Cauchy in both closed
form norms, and their limits agree in $\mathcal K$.  Thus $b$ is closed.
Furthermore,
\[
 \alpha a[u,u]+\beta p[u,u]+c\|u\|^2
 \le\theta_\sigma
      \bigl(a[u,u]+p[u,u]+\sigma\|u\|^2\bigr),
\]
which proves \eqref{eq:n9v42-residual-relative}.  Apply
\Cref{thm:n9v41-energy-normalization,thm:n9v41-Neumann} with $b$ in
place of $a_0$.
\end{proof}

\begin{corollary}[Positive-minus-negative split]
\label{cor:n9v42-positive-negative}
Let the perturbation have the form
\begin{equation}
 k=p-n,
 \label{eq:n9v42-sign-split}
\end{equation}
where $p,n$ are nonnegative forms.  If
\begin{equation}
 n[u,u]\le\theta\bigl(a[u,u]+p[u,u]\bigr)+c\|u\|^2,
 \qquad 0\le\theta<1,
 \label{eq:n9v42-negative-relative}
\end{equation}
then absorbing $p$ into the reference yields a strict residual form after
any shift $\sigma>c/(1-\theta)$; more precisely the shifted coefficient
is at most $\theta+c/\sigma<1$.
\end{corollary}

\begin{proof}
Use $r=-n$.  From \eqref{eq:n9v42-negative-relative},
\[
 |r[u,u]|\le(\theta+c/\sigma)b_\sigma[u,u].
\]
The declared inequality on $\sigma$ makes the coefficient less than one.
Apply \Cref{thm:n9v42-absorption}.
\end{proof}

\begin{remark}[Two legitimate shifted coefficients]
\label{rem:n9v42-shift-coefficients}
The estimate in \Cref{cor:n9v42-positive-negative} uses the additive
coefficient $\theta+c/\sigma$.  When a rowwise bound of the stronger form
$n\le\alpha a+\beta p+c\|\cdot\|^2$ is available,
\Cref{thm:n9v42-absorption} gives the sharper maximum coefficient
$\max\{\alpha,\beta,c/\sigma\}$.  These bounds should not be conflated.
\end{remark}

\subsection{Exact scalar and fibrewise ratios}

The absorption mechanism is already visible in a scalar fibre.  Let
$A(x)>0$, $P(x)\ge0$, and let $R(x)$ be real.

\begin{proposition}[Scalar optimal absorbed coefficient]
\label{prop:n9v42-scalar-optimal}
On $L^2(\mu)$ define
\begin{align}
 b[U]&=\int(A+P)|U|^2,d\mu,
 \label{eq:n9v42-scalar-b}\\
 r[U]&=\int R|U|^2,d\mu.
 \label{eq:n9v42-scalar-r}
\end{align}
Then the optimal relative-form coefficient is
\begin{equation}
 \sup_{U\ne0}\frac{|r[U]|}{b[U]}
 =\operatorname*{ess\,sup}_x
   \frac{|R(x)|}{A(x)+P(x)}.
 \label{eq:n9v42-scalar-ratio}
\end{equation}
In particular $P/A$ may be unbounded while the residual coefficient is
strictly less than one.
\end{proposition}

\begin{proof}
The essential supremum gives the upper bound after integration.  For the
reverse inequality, choose a positive-measure set on which the ratio is
within $\varepsilon$ of its essential supremum and choose $U$ supported
there, with its phase fixed so $R$ has a constant sign after splitting
the set into $\{R\ge0\}$ and $\{R<0\}$.  Normalize in the $b$ norm and
let $\varepsilon\downarrow0$.
\end{proof}

\begin{corollary}[Large positive rows can improve the margin]
\label{cor:n9v42-positive-improves}
If $R=-N$ with $N\ge0$, the absorbed coefficient is
$\operatorname*{ess\,sup}N/(A+P)$.  Increasing $P$ pointwise can only
decrease this coefficient, even if $P$ has no finite essential supremum
relative to $A$.
\end{corollary}

\begin{proof}
This is immediate from \eqref{eq:n9v42-scalar-ratio}.
\end{proof}

For finite-dimensional tangent fibres, let $A(x)>0$ and $P(x)\ge0$ be
self-adjoint matrices and put $B(x)=A(x)+P(x)>0$.

\begin{proposition}[Fibrewise normalized residual]
\label{prop:n9v42-fibrewise}
If $R(x)=R(x)^*$ is measurable and
\begin{equation}
 \operatorname*{ess\,sup}_x
 \|B(x)^{-1/2}R(x)B(x)^{-1/2}\|\le\theta,
 \label{eq:n9v42-fibre-bound}
\end{equation}
then the direct-integral forms satisfy
\begin{equation}
 |r[U,U]|\le\theta b[U,U].
 \label{eq:n9v42-fibre-form}
\end{equation}
For scalar fibres, \eqref{eq:n9v42-fibre-bound} is also necessary and
optimal by \Cref{prop:n9v42-scalar-optimal}.
\end{proposition}

\begin{proof}
Pointwise,
\[
 |\langle U,RU\rangle|
 \le\theta\langle U,BU\rangle.
\]
Integrate.  The scalar necessity is
\eqref{eq:n9v42-scalar-ratio}.
\end{proof}

\begin{firewall}[Moment bounds still do not replace the normalized
essential supremum]
\label{fire:n9v42-moment}
Finite moments of
$\|B^{-1/2}RB^{-1/2}\|$ do not imply
\eqref{eq:n9v42-fibre-bound}; the V41 localization counterexample applies
with the enlarged reference $B$.  Absorption changes the ratio that must
be bounded, but it does not turn a probability estimate into an operator
norm estimate.
\end{firewall}

\subsection{Ownership and maximal positive absorption}

Suppose a complete symmetric reduced Hessian form has a declared
decomposition
\begin{equation}
 k_C=p_C+r_C,
 \qquad p_C\ge0.
 \label{eq:n9v42-complete-split}
\end{equation}
The decomposition is useful only when it preserves the actual form
domain and the locality needed later.

\begin{lemma}[Monotonicity under additional positive absorption]
\label{lem:n9v42-monotonicity}
Let $0\le p_1\le p_2$ as forms on a common domain, and let $n\ge0$.
If
\begin{equation}
 n\le\theta(a+p_1)
 \label{eq:n9v42-first-margin}
\end{equation}
for some $\theta$, then
\begin{equation}
 n\le\theta(a+p_2).
 \label{eq:n9v42-second-margin}
\end{equation}
Thus absorbing an additional nonnegative row cannot worsen a one-sided
negative-part relative bound.
\end{lemma}

\begin{proof}
The form inequality $a+p_1\le a+p_2$ gives the result.
\end{proof}

\begin{proposition}[Partial absorption ledger]
\label{prop:n9v42-partial}
Let $0\le\zeta\le1$ and rewrite
\begin{equation}
 a+k_C=(a+\zeta p_C)+\bigl((1-\zeta)p_C+r_C\bigr).
 \label{eq:n9v42-partial-split}
\end{equation}
This identity is exact for every $\zeta$.  If $r_C=-n_C$ with
$n_C\ge0$, full absorption $\zeta=1$ minimizes the negative-part ratio
$n_C/(a+\zeta p_C)$ in the scalar fibre model.  For a general signed
$r_C$, no monotonic conclusion for the absolute residual norm follows
without additional commutation or form inequalities.
\end{proposition}

\begin{proof}
The algebraic identity is immediate.  In the scalar negative case the
denominator increases with $\zeta$.  For a signed residual, the numerator
in \eqref{eq:n9v42-partial-split} also changes, so the scalar monotonicity
argument no longer applies.
\end{proof}

\begin{firewall}[Spectral sign splitting may destroy locality]
\label{fire:n9v42-sign-locality}
Every finite-dimensional self-adjoint fibre has positive and negative
spectral parts.  This does not automatically provide a usable physical
decomposition: spectral projectors can lose differentiability at
crossings, can enlarge support for nonlocal operators, and can fail the
weighted locality card.  A positive row used in the reference must be
proved closed, measurable, correctly owned, and compatible with the
cofinal local structure.
\end{firewall}

\subsection{Constraint-reduced absorption card}

For the complete V34--V39 Hessian ledger, possible rows include the
ambient action Hessian, the Lagrange-multiplier constraint curvature,
the second logarithmic derivative of the coarea density, the chart
acceleration/divergence row, and the label Fisher/IMS debit.  Only the
last of these is manifestly nonnegative in its Fisher-trace role; the
others require a separate sign theorem before absorption.

\begin{definition}[Constraint positive-reference absorption card (CPRAC)]
\label{def:n9v42-CPRAC}
CPRAC consists of:
\begin{enumerate}
\item the actual constraint-reduced carrier and complete symmetric
Hessian form $k_C$;
\item an owned decomposition $k_C=p_C+r_C$ with $p_C$ a densely defined
closed nonnegative form;
\item a closed enlarged reference $b=a+p_C$ on a dense common domain;
\item a uniform residual estimate of the form
\eqref{eq:n9v42-two-row-bound}, or equivalently a uniform strict
energy-normalized residual bound after a declared shift;
\item preservation of RSCC/CASC chart covariance and support under the
decomposition; and
\item the separate weighted locality estimate for the enlarged
reference resolvent.
\end{enumerate}
\end{definition}

\begin{theorem}[CPRAC loads URRFC]
\label{thm:n9v42-CPRAC}
CPRAC items 1--4 load the unweighted URRFC of V41 with reference
$b=a+p_C$.  The inverse cost is
$(1-\theta_\sigma)^{-1}$ with $\theta_\sigma$ from
\eqref{eq:n9v42-absorbed-theta}.  Items 5--6 retain the chart and
locality interfaces required for cofinal use.
\end{theorem}

\begin{proof}
Apply \Cref{thm:n9v42-absorption} and then
\Cref{cor:n9v41-URRFC}.  The last statement is exactly the content of
CPRAC items 5--6; no locality estimate is inferred from form positivity.
\end{proof}

\begin{firewall}[Conformal and coarea sign rows]
\label{fire:n9v42-conformal}
An ambient negative conformal mode cannot be placed in $p_C$.  It must be
removed by an actual constraint support theorem as in V34 or controlled
inside the residual negative part.  Likewise
$-\hbar D^2\log r_C$ has no automatic sign.  Calling either row
``stabilizing'' without a proved form inequality would invalidate CPRAC.
\end{firewall}

\begin{assessment}[Progress at ninth-series V42]
\label{ass:n9v42-status}
The strict-form bottleneck has been sharpened.  Large nonnegative rows may
be absorbed into the reference even when they are unbounded relative to
the old energy; only the residual normalized by the enlarged reference
must have essential norm below one.  This can turn an unusable full
Hessian ratio into a small negative-part ratio, but it requires a genuine
sign, domain, ownership, and locality theorem.
\end{assessment}

\begin{programme}[Ninth-series V43: logarithmic-density sign calculus]
\label{prog:n9v42-V43}
The next upstream question is which constraint/coarea rows are
nonnegative.  V43 will derive an exact covariance formula for the
expected Hessian of a logarithmic density, distinguish pointwise sign
from integrated Fisher positivity, and test whether the label Fisher
debit can be absorbed without assuming a false pointwise sign for
$-D^2\log r_C$.
\end{programme}

\begin{firewall}[Claim boundary after ninth-series V42]
\label{fire:n9v42-boundary}
V42 proves the abstract absorption theorem and its direct-integral
criterion.  It does not prove CPRAC for the physical SM--Einstein
Hessian, a sign for the coarea row, weighted locality, semigroup
transfer, reflection positivity, reconstruction, or the full continuum.
\end{firewall}

\begin{theorem}[Ninth-series V42 result]
\label{thm:n9v42-result}
If a complete reduced Hessian decomposes as a closed nonnegative row plus
a residual whose energy-normalized essential norm is strictly below one,
then the positive row may be absorbed and V41 resums every residual
insertion.  The optimal scalar coefficient is
$\operatorname*{ess\,sup}|R|/(A+P)$; probability moments of that ratio do
not replace the essential bound.
\end{theorem}

\begin{proof}
Use \Cref{thm:n9v42-absorption,prop:n9v42-scalar-optimal,
prop:n9v42-fibrewise,thm:n9v42-CPRAC}.
\end{proof}

\paragraph{Parent-version record.}
V42 was formed from
\texttt{Renewal\_SM\_Einstein\_Continuum\_Research\_Notes\_V41.tex}.
The V41 parent had $14\,586$ logical source lines, $529\,943$ bytes, and
SHA--256
\[
 \texttt{B0915676F931373B8A9F27B577F1DFD849A96F2ADE282384351C98961EEAA887}.
\]

% =============================================================================
% END APPENDED NINTH-SERIES V42 MATERIAL
% =============================================================================
% =============================================================================
% BEGIN APPENDED NINTH-SERIES V43 MATERIAL
% =============================================================================

\clearpage
\section{Ninth-series V43: integrated Fisher positivity, logarithmic
density signs, and label compression}
\label{sec:v9v43}

\subsection{Coordinate-free expected Hessian identity}

V42 requires an actual nonnegative form before a row can be absorbed.
The differentiated score has a positive expectation, but this does not
make each of its action, coarea, chart, or label pieces pointwise
nonnegative.  We first state the invariant identity that separates these
notions.

Let $\mu$ be a probability measure with an integration-by-parts algebra
$\mathcal A$ and vector fields $V_1,\dots,V_m$ satisfying
\begin{equation}
 \int V_i f\,d\mu=\int f\beta_i\,d\mu,
 \qquad f\in\mathcal A.
 \label{eq:n9v43-ibp}
\end{equation}
Assume all products and derivatives below are integrable and that no
boundary flux remains.

\begin{theorem}[Expected differentiated-score matrix]
\label{thm:n9v43-expected-H}
The matrix
\begin{equation}
 M_{ij}:=\mathbb E_\mu[V_i\beta_j]
 \label{eq:n9v43-M}
\end{equation}
satisfies
\begin{equation}
 \boxed{M_{ij}=\mathbb E_\mu[\beta_i\beta_j].}
 \label{eq:n9v43-Fisher-matrix}
\end{equation}
Consequently $M$ is symmetric positive semidefinite and, since
$\mathbb E\beta_i=0$, it is the covariance matrix of the score vector.
This is an integrated statement; $V_i\beta_j(x)$ need not be symmetric
or positive at a fixed configuration.
\end{theorem}

\begin{proof}
Apply \eqref{eq:n9v43-ibp} with $f=\beta_j$ to obtain
\eqref{eq:n9v43-Fisher-matrix}.  Applying it with $f=1$ gives
$\mathbb E\beta_i=0$.  For $z\in\mathbb R^m$,
\[
 z^TMz=\mathbb E\left|\sum_i z_i\beta_i\right|^2\ge0.
\]
Symmetry follows from the same Gram representation.
\end{proof}

\begin{corollary}[Constant reduced directions]
\label{cor:n9v43-constant}
For the V34 reduced density
$d\bar\mu\propto e^{-\bar S/\hbar}r_Cdp$ and constant coordinate
directions, write
\begin{equation}
 a_i:=\hbar^{-1}\partial_i\bar S,
 \qquad
 \ell_i:=\partial_i\log r_C,
 \qquad
 \bar\beta_i=a_i-\ell_i.
 \label{eq:n9v43-a-l-beta}
\end{equation}
Then
\begin{equation}
 \mathbb E\left[
  \hbar^{-1}\partial_i\partial_j\bar S
  -\partial_i\partial_j\log r_C
 \right]
 =\mathbb E[\bar\beta_i\bar\beta_j]\succeq0.
 \label{eq:n9v43-complete-positive}
\end{equation}
\end{corollary}

\begin{proof}
Use \eqref{eq:n9v34-Dscore} and
\Cref{thm:n9v43-expected-H}.
\end{proof}

\begin{remark}[Relation to V31]
\label{rem:n9v43-v31}
For one direction, \eqref{eq:n9v43-Fisher-matrix} is the Fisher--Hessian
identity proved by bounded truncation in V31.  V43 uses it as an input
and adds the matrix, label-compression, and row-sign consequences needed
by the V42 absorption problem.
\end{remark}

\subsection{The coarea/log-density row is not positive by itself}

\begin{proposition}[Exact expected row decomposition]
\label{prop:n9v43-row-decomposition}
With \eqref{eq:n9v43-a-l-beta},
\begin{align}
 \mathbb E[\partial_i a_j]
 &=\mathbb E[\bar\beta_i a_j],
 \label{eq:n9v43-action-row}\\
 \mathbb E[-\partial_i\ell_j]
 &=\mathbb E[\ell_i\ell_j]
   -\mathbb E[a_i\ell_j],
 \label{eq:n9v43-density-row}\\
 \mathbb E[\partial_i a_j-\partial_i\ell_j]
 &=\mathbb E[\bar\beta_i\bar\beta_j].
 \label{eq:n9v43-row-sum}
\end{align}
Neither \eqref{eq:n9v43-action-row} nor
\eqref{eq:n9v43-density-row} has a fixed sign without extra hypotheses;
their sum is positive semidefinite.
\end{proposition}

\begin{proof}
Apply integration by parts with $f=a_j$ and $f=\ell_j$:
\[
 \mathbb E[\partial_i a_j]=\mathbb E[\bar\beta_i a_j],
 \qquad
 \mathbb E[\partial_i\ell_j]=\mathbb E[\bar\beta_i\ell_j].
\]
Insert $\bar\beta_i=a_i-\ell_i$ into the second identity.  Adding the
rows gives \eqref{eq:n9v43-row-sum}, also directly from
\Cref{thm:n9v43-expected-H}.
\end{proof}

\begin{example}[A changing-sign logarithmic Hessian]
\label{ex:n9v43-circle}
On the circle let $r(x)=Z^{-1}e^{b\cos x}$ with $b\ne0$ and take no
action.  Then
\begin{equation}
 -\partial_x^2\log r=b\cos x,
 \label{eq:n9v43-circle-pointwise}
\end{equation}
which changes sign.  Nevertheless periodic integration by parts gives
\begin{equation}
 \mathbb E_r[-\partial_x^2\log r]
 =\mathbb E_r[(\partial_x\log r)^2]ge0.
 \label{eq:n9v43-circle-integrated}
\end{equation}
Thus integrated Fisher positivity does not justify pointwise absorption
of the logarithmic-density Hessian.
\end{example}

\begin{proof}
Differentiate $b\cos x$ twice for
\eqref{eq:n9v43-circle-pointwise}.  The periodic boundary term vanishes,
and integration of
$\partial_x(r\partial_x\log r)$ gives
\eqref{eq:n9v43-circle-integrated}.
\end{proof}

\begin{firewall}[Expectation positivity is not form positivity]
\label{fire:n9v43-expectation}
The finite matrix
$\mathbb E[V_i\beta_j]$ being positive semidefinite does not imply that
the multiplication form
$\int\langle U(x),(V\beta)(x)U(x)\rangle d\mu$ is nonnegative for
configuration-dependent $U$.  Such a field can concentrate where the
pointwise matrix is negative, exactly as in the V41
essential-supremum argument.
\end{firewall}

\subsection{Normalized labels: a pointwise positive compressed row}

Let $\kappa_\alpha=\eta_\alpha^2$ be the normalized square labels of V35,
and define on their positive sets
\begin{equation}
 L_{i,\alpha}:=V_i\log\kappa_\alpha.
 \label{eq:n9v43-label-log-gradient}
\end{equation}
The labelled scores are
$\widehat\beta_{i,\alpha}=\beta_i-L_{i,\alpha}$.

\begin{lemma}[Pointwise label Hessian identity]
\label{lem:n9v43-label-Hessian}
For every $i,j$,
\begin{align}
 \sum_\alpha\kappa_\alpha L_{i,\alpha}&=0,
 \label{eq:n9v43-label-mean-zero}\\
 -\sum_\alpha\kappa_\alpha V_iL_{j,\alpha}
 &=\sum_\alpha\kappa_\alpha
       L_{i,\alpha}L_{j,\alpha}
 =4\sum_\alpha(V_i\eta_\alpha)(V_j\eta_\alpha).
 \label{eq:n9v43-label-Hessian-identity}
\end{align}
The matrix on the right is pointwise positive semidefinite.
\end{lemma}

\begin{proof}
The first identity is
$\sum_\alpha V_i\kappa_\alpha=V_i1=0$.  Next,
\[
 \kappa_\alpha V_iV_j\log\kappa_\alpha
 =V_iV_j\kappa_\alpha
  -\frac{(V_i\kappa_\alpha)(V_j\kappa_\alpha)}
        {\kappa_\alpha}.
\]
Sum in $\alpha$.  The first terms vanish because
$\sum_\alpha\kappa_\alpha=1$; move the remaining term to the other side.
Finally use $V_i\kappa_\alpha=2\eta_\alpha V_i\eta_\alpha$ and extend
through zeros by the square-root expression.  Positivity follows because
for $z\in\mathbb R^m$ the quadratic form is
$4\sum_\alpha|\sum_i z_iV_i\eta_\alpha|^2$.
\end{proof}

\begin{theorem}[Labelled Fisher/Hessian compression]
\label{thm:n9v43-label-compression}
At each configuration,
\begin{align}
 \sum_\alpha\kappa_\alpha
  \widehat\beta_{i,\alpha}widehat\beta_{j,\alpha}
 &=\beta_i\beta_j+\mathcal I^\kappa_{ij},
 \label{eq:n9v43-labelled-Fisher}\\
 \sum_\alpha\kappa_\alpha
  V_i\widehat\beta_{j,\alpha}
 &=V_i\beta_j+\mathcal I^\kappa_{ij},
 \label{eq:n9v43-labelled-Hessian}\\
 \mathcal I^\kappa_{ij}
 &:=\sum_\alpha\kappa_\alpha
       L_{i,\alpha}L_{j,\alpha}succeq0.
 \label{eq:n9v43-label-matrix}
\end{align}
After expectation, both sides of
\eqref{eq:n9v43-labelled-Hessian} equal the labelled Fisher matrix in
\eqref{eq:n9v43-labelled-Fisher}.  Hence the exact additional row in the
label-independent compression is the pointwise positive matrix
$\mathcal I^\kappa$.
\end{theorem}

\begin{proof}
Expand the product of labelled scores.  The cross terms vanish by
\eqref{eq:n9v43-label-mean-zero}, proving
\eqref{eq:n9v43-labelled-Fisher}.  Next,
$V_i\widehat\beta_{j,\alpha}=V_i\beta_j-V_iL_{j,\alpha}$; average and use
\eqref{eq:n9v43-label-Hessian-identity}.  Apply
\Cref{thm:n9v43-expected-H} to the joint labelled integration-by-parts
law to identify the expectations.
\end{proof}

\begin{corollary}[Recovery of the V35 trace debit]
\label{cor:n9v43-v35-trace}
Taking the weighted trace of \eqref{eq:n9v43-label-matrix} gives
\begin{equation}
 \sum_iw_i^{-1}\mathcal I^\kappa_{ii}
 =4\sum_iw_i^{-1}\sum_\alpha|V_i\eta_\alpha|^2,
 \label{eq:n9v43-trace-debit}
\end{equation}
which is exactly the label Fisher/IMS debit of V35.
\end{corollary}

\begin{proof}
Use the last equality in
\eqref{eq:n9v43-label-Hessian-identity} and sum the diagonal.
\end{proof}

\subsection{Compression is weaker than a full labelled operator identity}

Define the isometry
\begin{equation}
 J:L^2(\mu;\mathbb C^m)longrightarrow
 L^2(\widehat\mu;\mathbb C^m),
 \qquad (JU)(x,\alpha)=U(x),
 \label{eq:n9v43-J}
\end{equation}
and its adjoint
\begin{equation}
 (J^*F)(x)=\sum_\alpha\kappa_\alpha(x)F(x,\alpha).
 \label{eq:n9v43-Jstar}
\end{equation}

\begin{proposition}[Exact physical-subspace compression]
\label{prop:n9v43-physical-compression}
Let $K_\alpha(x)$ be the labelled differentiated-score matrix with
entries $V_i\widehat\beta_{j,\alpha}$.  As multiplication forms on their
common core,
\begin{equation}
 J^*KJ=K_0+\mathcal I^\kappa,
 \label{eq:n9v43-compressed-operator}
\end{equation}
where $(K_0)_{ij}=V_i\beta_j$.  The added form
\begin{equation}
 p_\kappa[U]
 :=\int\langle U,\mathcal I^\kappa U\rangle,d\mu
 \label{eq:n9v43-positive-label-form}
\end{equation}
is nonnegative.
\end{proposition}

\begin{proof}
Equations \eqref{eq:n9v43-J}--\eqref{eq:n9v43-Jstar} show that
$J^*KJ$ is the conditional label average of $K_\alpha$.  Apply
\eqref{eq:n9v43-labelled-Hessian}.  Pointwise positivity of
$\mathcal I^\kappa$ proves the last statement.
\end{proof}

\begin{firewall}[Compression does not commute with inversion
automatically]
\label{fire:n9v43-compression-inverse}
In general
\[
 J^*(H_{\rm lab}+\sigma)^{-1}J
 \ne\bigl(J^*(H_{\rm lab}+\sigma)J\bigr)^{-1}.
\]
Equality requires the label-independent range of $J$ to reduce the full
labelled operator, or a separate Schur-complement theorem controlling
its coupling to label-dependent modes.  Therefore
\eqref{eq:n9v43-positive-label-form} is immediately absorbable only in a
proof formulated on the compressed physical carrier.
\end{firewall}

\begin{proposition}[Reducing-subspace criterion]
\label{prop:n9v43-reducing}
Let $P=JJ^*$ be the orthogonal projection onto label-independent fields.
If a self-adjoint labelled operator $H_{\rm lab}$ satisfies
$PH_{\rm lab}\subset H_{\rm lab}P$, then
\begin{equation}
 J^*(H_{\rm lab}+\sigma)^{-1}J
 =\bigl(J^*H_{\rm lab}J+\sigma\bigr)^{-1}
 \label{eq:n9v43-reducing-resolvent}
\end{equation}
for every common resolvent point $-\sigma$.
\end{proposition}

\begin{proof}
The commutation makes $\operatorname{Ran}J$ reducing, so the resolvent is
block diagonal with respect to
$\operatorname{Ran}P\oplus\ker P$.  The restriction to
$\operatorname{Ran}J$ is unitarily equivalent through $J$ to
$J^*H_{\rm lab}J$.
\end{proof}

\subsection{A valid label-absorption interface}

\begin{definition}[Label Fisher absorption card (LFAC)]
\label{def:n9v43-LFAC}
LFAC requires:
\begin{enumerate}
\item a normalized square-label kernel satisfying the regularity and
CASC estimates of V35--V37;
\item the compressed physical carrier of
\eqref{eq:n9v43-J}, or a proved reducing/Schur-complement replacement;
\item a densely defined closed positive form $p_\kappa$ from
\eqref{eq:n9v43-positive-label-form};
\item a CPRAC residual bound relative to $a+p_\kappa$; and
\item retention of the high-shell and singular-chart debits from
V36--V39.
\end{enumerate}
\end{definition}

\begin{theorem}[LFAC loads the label row of CPRAC]
\label{thm:n9v43-LFAC}
Under LFAC, the label Fisher matrix is a legitimately nonnegative owned
row and may be absorbed into the V42 enlarged reference.  This statement
does not assign a sign to the unaveraged chart Hessians
$-V_iV_j\log\kappa_\alpha$ or to the physical coarea Hessian
$-D^2\log r_C$.
\end{theorem}

\begin{proof}
The pointwise positive compression is
\Cref{thm:n9v43-label-compression,
prop:n9v43-physical-compression}.  Closedness and the residual estimate
are LFAC items 3--4, so \Cref{thm:n9v42-CPRAC} applies.  The final
qualification follows from \Cref{ex:n9v43-circle} and from the fact that
\eqref{eq:n9v43-label-Hessian-identity} holds only after conditional
label averaging.
\end{proof}

\begin{assessment}[Progress at ninth-series V43]
\label{ass:n9v43-status}
The sign issue is now exact.  The complete expected differentiated score
is the Fisher covariance, while the action and coarea/log-density rows
need not be positive separately.  Normalized label averaging is special:
it produces the pointwise positive matrix
$4\sum_\alpha d\eta_\alpha\otimes d\eta_\alpha$ on the compressed
physical carrier.  Full labelled resolvents still require reduction or a
Schur complement.
\end{assessment}

\begin{programme}[Ninth-series V44: label-sector Schur complement]
\label{prog:n9v43-V44}
V44 will remove the compression caveat by decomposing the labelled
Hilbert space into physical conditional means and mean-zero label modes.
It will derive the exact block resolvent and a sufficient Schur margin
for integrating out label-dependent modes without asserting that the
physical subspace is automatically reducing.
\end{programme}

\begin{firewall}[Claim boundary after ninth-series V43]
\label{fire:n9v43-boundary}
V43 proves integrated Fisher positivity and pointwise positivity of the
compressed label row.  It does not prove a pointwise sign for the coarea
Hessian, reduction of the full labelled operator, LFAC/CPRAC for the
physical theory, weighted locality, OS reconstruction, or the full
continuum.
\end{firewall}

\begin{theorem}[Ninth-series V43 result]
\label{thm:n9v43-result}
The expected complete differentiated-score matrix is the positive Fisher
covariance \eqref{eq:n9v43-Fisher-matrix}, but its action and logarithmic
density rows need not have separate signs.  Conditional square-label
compression adds the pointwise positive matrix
\eqref{eq:n9v43-label-matrix}; under LFAC this row may be absorbed into
the V42 reference form.
\end{theorem}

\begin{proof}
Use \Cref{thm:n9v43-expected-H,prop:n9v43-row-decomposition,
ex:n9v43-circle,thm:n9v43-label-compression,
prop:n9v43-physical-compression,thm:n9v43-LFAC}.
\end{proof}

\paragraph{Parent-version record.}
V43 was formed from
\texttt{Renewal\_SM\_Einstein\_Continuum\_Research\_Notes\_V42.tex}.
The V42 parent had $14\,954$ logical source lines, $543\,255$ bytes, and
SHA--256
\[
 \texttt{7C664E186C9F0C38B1AD76DE578BFADF5AF3C368017205E4ADAB9F60A90D8069}.
\]

% =============================================================================
% END APPENDED NINTH-SERIES V43 MATERIAL
% =============================================================================
% =============================================================================
% BEGIN APPENDED NINTH-SERIES V44 MATERIAL
% =============================================================================

\clearpage
\section{Ninth-series V44: exact elimination of mean-zero label modes}
\label{sec:v9v44}

\subsection{Physical and auxiliary label sectors}

Let $\widehat\mu(dx,d\alpha)=\mu(dx)\kappa_\alpha(x)$ be the normalized
label lift of V35 and let
\begin{equation}
 \widehat{\mathcal K}=L^2(\widehat\mu;\mathcal H),
 \qquad
 \mathcal K=L^2(\mu;\mathcal H).
 \label{eq:n9v44-Hilbert-spaces}
\end{equation}
The isometry $J:\mathcal K\to\widehat{\mathcal K}$ and conditional mean
$J^*$ are given by \eqref{eq:n9v43-J}--\eqref{eq:n9v43-Jstar}.  Put
\begin{equation}
 P:=JJ^*,
 \qquad Q:=I-P.
 \label{eq:n9v44-PQ}
\end{equation}
Then $P$ projects onto label-independent fields and
\begin{equation}
 Q\widehat{\mathcal K}
 =\left\{F:\sum_\alpha\kappa_\alpha(x)F(x,\alpha)=0
           \text{ for almost every }x\right\}.
 \label{eq:n9v44-meanzero}
\end{equation}

The aim is to compute the physical resolvent even when $P$ is not a
reducing projection.

\subsection{A form-level two-sector theorem}

Let $\mathcal K_P$ and $\mathcal K_Q$ be Hilbert spaces.  Let $A>0$ and
$D>0$ be positive self-adjoint operators on them, with bounded inverses.
Suppose a sesquilinear cross form $b$ on
$D(A^{1/2})\times D(D^{1/2})$ satisfies
\begin{equation}
 |b[p,q]|
 \le\gamma\|A^{1/2}p\|\,\|D^{1/2}q\|,
 \qquad 0\le\gamma<1.
 \label{eq:n9v44-cross-bound}
\end{equation}

\begin{lemma}[Energy-normalized cross operator]
\label{lem:n9v44-cross-operator}
There is a unique bounded operator
$C:\mathcal K_Q\to\mathcal K_P$ with
\begin{equation}
 b[p,q]=\langle A^{1/2}p,CD^{1/2}q\rangle,
 \qquad \|C\|\le\gamma.
 \label{eq:n9v44-C}
\end{equation}
\end{lemma}

\begin{proof}
Transport $b$ by the unitaries from the form domains equipped with their
energy norms to $\mathcal K_P$ and $\mathcal K_Q$.  The transported
bounded sesquilinear form has norm at most $\gamma$, so Riesz
representation gives $C$.
\end{proof}

On $D(A^{1/2})\oplus D(D^{1/2})$, define
\begin{equation}
 \mathfrak l[p\oplus q]
 :=\|A^{1/2}p\|^2+\|D^{1/2}q\|^2
   +2\operatorname{Re}
      \langle A^{1/2}p,CD^{1/2}q\rangle.
 \label{eq:n9v44-block-form}
\end{equation}

\begin{theorem}[Exact label-sector Schur complement]
\label{thm:n9v44-Schur}
The form \eqref{eq:n9v44-block-form} is closed and strictly positive and
defines a self-adjoint operator $L$.  Its physical Schur form is
\begin{equation}
 \mathfrak s[p]
 =\left\langle A^{1/2}p,
       (I-CC^*)A^{1/2}p\right\rangle,
 \label{eq:n9v44-Schur-form}
\end{equation}
with associated operator denoted $S$.  One has
\begin{align}
 P L^{-1}P\big|_{\mathcal K_P}&=S^{-1},
 \label{eq:n9v44-compressed-resolvent}\\
 \mathfrak s[p]&\ge(1-\gamma^2)\|A^{1/2}p\|^2,
 \label{eq:n9v44-Schur-margin}\\
 \|A^{1/2}S^{-1}A^{1/2}\|&\le(1-\gamma^2)^{-1},
 \label{eq:n9v44-Schur-inverse-bound}
\end{align}
where the last product denotes its bounded extension.  Formally, and
exactly whenever the operator blocks have a common core,
\begin{equation}
 S=A-BD^{-1}B^*,
 \qquad B=A^{1/2}CD^{1/2}.
 \label{eq:n9v44-formal-Schur}
\end{equation}
\end{theorem}

\begin{proof}
Complete the square:
\begin{align*}
 \mathfrak l[p\oplus q]
 &=\|D^{1/2}q+C^*A^{1/2}p\|^2\\
 &\quad+
 \langle A^{1/2}p,(I-CC^*)A^{1/2}p\rangle.
\end{align*}
Since $I-CC^*\ge(1-\gamma^2)I$, this form is equivalent to the direct-sum
energy norm and is therefore closed and strictly positive.  For fixed
$p$, its minimizing $Q$-energy coordinate is
$D^{1/2}q=-C^*A^{1/2}p$, and the minimum is
\eqref{eq:n9v44-Schur-form}.  Solving the weak equation
$L(p\oplus q)=f\oplus0$ by this elimination shows that the $P$ component
is $S^{-1}f$, proving \eqref{eq:n9v44-compressed-resolvent}.  The margin
and inverse bound follow from the spectral bound for $CC^*$.  Expanding
$C=A^{-1/2}BD^{-1/2}$ gives \eqref{eq:n9v44-formal-Schur} on a common
operator core; the form formula remains the primary statement.
\end{proof}

\begin{corollary}[Error of naive physical compression]
\label{cor:n9v44-naive-error}
The exact energy-normalized physical resolvent is
\begin{equation}
 A^{1/2}S^{-1}A^{1/2}=(I-CC^*)^{-1}.
 \label{eq:n9v44-exact-physical}
\end{equation}
Consequently
\begin{equation}
 \left\|A^{1/2}(S^{-1}-A^{-1})A^{1/2}\right\|
 \le\frac{\gamma^2}{1-\gamma^2}.
 \label{eq:n9v44-naive-error}
\end{equation}
If $\gamma_n\to0$, inversion after conditional compression is
asymptotically valid in the energy-normalized norm; a merely uniform
$\gamma<1$ gives boundedness but not disappearance of the correction.
\end{corollary}

\begin{proof}
Equation \eqref{eq:n9v44-Schur-form} is the form factorization
$S=A^{1/2}(I-CC^*)A^{1/2}$.  Invert it as in V41.  Then
\[
 (I-CC^*)^{-1}-I=(I-CC^*)^{-1}CC^*,
\]
whose norm is at most $\gamma^2/(1-\gamma^2)$.
\end{proof}

\begin{corollary}[Reducing case]
\label{cor:n9v44-reducing}
If the physical and mean-zero sectors decouple, then $C=0$ and
$S=A$.  Equation \eqref{eq:n9v44-compressed-resolvent} reduces to the
V43 reducing-subspace identity \eqref{eq:n9v43-reducing-resolvent}.
\end{corollary}

\begin{proof}
Decoupling means $b=0$, hence $C=0$ by
\Cref{lem:n9v44-cross-operator}.
\end{proof}

\subsection{The induced negative row}

The Schur correction has a definite sign and therefore fits the V42
positive-minus-negative architecture.

\begin{proposition}[Mean-zero elimination produces a controlled negative
form]
\label{prop:n9v44-negative-row}
Relative to the compressed positive block $A$, eliminating the mean-zero
label modes produces the nonnegative subtraction
\begin{equation}
 n_{\rm lab}[p]
 :=\|C^*A^{1/2}p\|^2,
 \qquad
 0\le n_{\rm lab}[p]
 \le\gamma^2\|A^{1/2}p\|^2.
 \label{eq:n9v44-negative-row}
\end{equation}
Thus the effective physical form is
\begin{equation}
 \mathfrak s[p]=a_P[p]-n_{\rm lab}[p]
 \label{eq:n9v44-positive-minus-negative}
\end{equation}
with strict negative-part relative coefficient at most $\gamma^2$.
\end{proposition}

\begin{proof}
Subtract \eqref{eq:n9v44-Schur-form} from
$\|A^{1/2}p\|^2$ and use $\|C^*\|\le\gamma$.
\end{proof}

\begin{corollary}[Loading the V42 absorption theorem]
\label{cor:n9v44-load-V42}
If $A$ already contains the closed positive reference and the compressed
label Fisher form $p_\kappa$ of V43, then
\eqref{eq:n9v44-negative-row} is a valid residual negative row for
\Cref{cor:n9v42-positive-negative}.  No pointwise sign is required for
the off-diagonal block itself.
\end{corollary}

\begin{proof}
The Schur correction is nonnegative by construction and is relatively
bounded by $\gamma^2<1$ against the complete $A$ block.  Apply
\Cref{cor:n9v42-positive-negative} with $c=0$.
\end{proof}

\subsection{A checkable gap/coupling estimate}

The abstract coefficient $\gamma$ separates two tasks: a lower bound in
the mean-zero sector and an off-diagonal coupling bound.

\begin{proposition}[Absolute cross-bound loader]
\label{prop:n9v44-gap-loader}
Suppose
\begin{equation}
 A\ge a_*I,
 \qquad D\ge d_*I,
 \qquad \|B\|\le b_*,
 \label{eq:n9v44-absolute-data}
\end{equation}
for positive $a_*,d_*$ and a bounded operator block $B$.  Then
\begin{equation}
 \|A^{-1/2}BD^{-1/2}\|
 \le\frac{b_*}{\sqrt{a_*d_*}}.
 \label{eq:n9v44-gamma-absolute}
\end{equation}
Hence $b_*^2<a_*d_*$ supplies a strict Schur margin.
\end{proposition}

\begin{proof}
The spectral lower bounds give
$\|A^{-1/2}\|\le a_*^{-1/2}$ and
$\|D^{-1/2}\|\le d_*^{-1/2}$.  Use submultiplicativity.
\end{proof}

\begin{proposition}[Form cross-bound loader]
\label{prop:n9v44-form-loader}
Suppose on the form domains
\begin{equation}
 |b[p,q]|^2\le\delta^2 a_P[p]d_Q[q]
 \label{eq:n9v44-form-cross}
\end{equation}
with $0\le\delta<1$.  Then the energy-normalized cross operator exists
with $\gamma\le\delta$, and the physical Schur margin is at least
$1-\delta^2$.
\end{proposition}

\begin{proof}
This is \Cref{lem:n9v44-cross-operator,thm:n9v44-Schur} with
$\gamma=\delta$.
\end{proof}

\begin{firewall}[A label count is not a mean-zero gap]
\label{fire:n9v44-label-gap}
The finiteness of the label set and the normalization
$\sum_\alpha\kappa_\alpha=1$ do not imply $D\ge d_*I$ on the mean-zero
sector.  A gap must come from the actual labelled operator or from an
explicit auxiliary label penalty.  Likewise the V35 conditional Jensen
inequality does not bound the off-diagonal block $B$.
\end{firewall}

\subsection{Weighted locality and cofinal errors}

For spatially indexed blocks, introduce a weighted operator norm
$\|\cdot\|_\mu$ that is submultiplicative and stable under adjoints.  No
particular kernel realization is assumed here.

\begin{proposition}[Weighted Schur series]
\label{prop:n9v44-weighted}
Suppose
\begin{equation}
 \|C\|_\mu\le\gamma_\mu<1
 \label{eq:n9v44-weighted-C}
\end{equation}
in a Banach $*$-algebra norm.  Then
\begin{align}
 \|(I-CC^*)^{-1}\|_\mu
 &\le\frac1{1-\gamma_\mu^2},
 \label{eq:n9v44-weighted-inverse}\\
 \|(I-CC^*)^{-1}-I\|_\mu
 &\le\frac{\gamma_\mu^2}{1-\gamma_\mu^2}.
 \label{eq:n9v44-weighted-error}
\end{align}
\end{proposition}

\begin{proof}
Submultiplicativity and adjoint stability give
$\|CC^*\|_\mu\le\gamma_\mu^2<1$.  Sum the geometric series in the Banach
algebra and its tail after the identity term.
\end{proof}

\begin{firewall}[Unweighted margin does not imply weighted locality]
\label{fire:n9v44-weighted}
The Hilbert operator bound $\|C\|<1$ does not imply
\eqref{eq:n9v44-weighted-C}.  Spatial decay of both energy normalizers
and the physical/label coupling must be proved.  This is the same type
boundary retained in the V40 audit of the Yang--Mills WEFC result.
\end{firewall}

\subsection{Label-sector Schur card}

\begin{definition}[Label-sector Schur card (LSSC)]
\label{def:n9v44-LSSC}
LSSC consists of:
\begin{enumerate}
\item the exact conditional-mean decomposition
$\widehat{\mathcal K}=P\widehat{\mathcal K}\oplus Q\widehat{\mathcal K}$;
\item a strictly positive physical block $A$ containing every owned
positive reference row, including $p_\kappa$ when LFAC applies;
\item a strictly positive mean-zero label block $D$;
\item a form cross-bound \eqref{eq:n9v44-cross-bound} with a uniform
$\gamma<1$;
\item $\gamma_n\to0$ if the naive compressed resolvent, rather than the
exact Schur resolvent, is used in the cofinal limit; and
\item the weighted version \eqref{eq:n9v44-weighted-C} whenever spatial
locality is claimed.
\end{enumerate}
\end{definition}

\begin{theorem}[LSSC closes the auxiliary-label inversion interface]
\label{thm:n9v44-LSSC}
Under LSSC, the physical labelled resolvent is exactly the inverse Schur
form \eqref{eq:n9v44-compressed-resolvent}.  The mean-zero label modes
produce a negative residual of relative size at most $\gamma^2$, so
V42--V43 positive absorption and V41 resummation apply.  If the weighted
row is present, the same conclusion holds in the declared locality norm.
\end{theorem}

\begin{proof}
Use \Cref{thm:n9v44-Schur,prop:n9v44-negative-row,
cor:n9v44-load-V42}.  The locality conclusion is
\Cref{prop:n9v44-weighted}.
\end{proof}

\begin{assessment}[Progress at ninth-series V44]
\label{ass:n9v44-status}
The label-compression caveat is now an exact block problem.  Integrating
out conditional-mean-zero labels subtracts the positive Schur row
$BD^{-1}B^*$ from the physical block.  Its relative size is the square
of the energy-normalized cross coupling.  A strict coupling margin gives
a valid absorbed resolvent; asymptotic equivalence to naive compression
requires the stronger condition $\gamma_n\to0$.
\end{assessment}

\begin{programme}[Ninth-series V45: compulsory companion audit]
\label{prog:n9v44-V45}
V45 will perform the next compulsory Yang--Mills/Navier--Stokes audit,
test for a type-compatible energy-normalizer locality or label-gap
estimate, and continue to V46 after the audit.
\end{programme}

\begin{firewall}[Claim boundary after ninth-series V44]
\label{fire:n9v44-boundary}
V44 proves the exact label-sector Schur theorem.  It does not prove the
physical mean-zero label gap, cross-coupling margin, weighted locality,
CPRAC for the SM--Einstein Hessian, semigroup transfer, OS
reconstruction, or the full continuum.
\end{firewall}

\begin{theorem}[Ninth-series V44 result]
\label{thm:n9v44-result}
Without a reducing label projection, the physical resolvent is the
inverse Schur form \eqref{eq:n9v44-Schur-form}.  A cross coefficient
$\gamma<1$ gives physical margin $1-\gamma^2$ and a negative residual
row of relative size at most $\gamma^2$; the naive compressed inverse is
accurate only when $\gamma\to0$.
\end{theorem}

\begin{proof}
Use \Cref{thm:n9v44-Schur,cor:n9v44-naive-error,
prop:n9v44-negative-row,thm:n9v44-LSSC}.
\end{proof}

\paragraph{Parent-version record.}
V44 was formed from
\texttt{Renewal\_SM\_Einstein\_Continuum\_Research\_Notes\_V43.tex}.
The V43 parent had $15\,380$ logical source lines, $557\,864$ bytes, and
SHA--256
\[
 \texttt{B08DDDC593152FAF54974A22595022A3874B90D7779403ED905B90F5098E03B0}.
\]

% =============================================================================
% END APPENDED NINTH-SERIES V44 MATERIAL
% =============================================================================
% =============================================================================
% BEGIN APPENDED NINTH-SERIES V45 MATERIAL
% =============================================================================

\clearpage
\section{Ninth-series V45: compulsory companion audit, sector margins,
and Fock amplification}
\label{sec:v9v45}

\subsection{Audit record and transfer decision}

The newest Yang--Mills files at this audit were
\begin{align*}
 &\texttt{Renewal\_Yang\_Mills\_Mass\_Gap\_Research\_Notes\_V80.tex},\\
 &\texttt{Renewal\_Yang\_Mills\_Mass\_Gap\_Research\_Notes\_V81.tex}.
\end{align*}
They were byte-identical, so V81 was not counted as independent
mathematical evidence.  Each had $35\,570$ logical source lines,
$1\,270\,763$ bytes, and SHA--256
\[
 \texttt{C14E143490EEC3EE796B5D4F76B194D9547DA1B4FB7E456C8B67686056BAFAA8}.
\]
The newest Navier--Stokes file remained
\texttt{Renewal\_NS\_Stability\_Research\_Notes\_V26.tex}, with
$5\,319$ lines, $278\,283$ bytes, and SHA--256
\[
 \texttt{82C887F8C2C69E4F28FAD7C3DC8DE5B9099CB5A4BF431EBA1FFF3E91935978AD}.
\]

The substantive Yang--Mills V78--V80 increment changes the interpretation
of the V44 Schur programme in two ways.
\begin{enumerate}
\item V78 proves a weighted estimate for a quadratic one-particle
operator but gives a spectator counterexample to an overstrong
full-Fock operator transfer.  It replaces one universal treatment by
separate quadratic, good nonlinear, and bad-cluster resummations.
\item V79 saturates bad clusters and turns boundary interactions into a
distinct-edge forest, eliminating arbitrary boundary multiplicity.
\item V80 proves that a compact Gauss projection yields normalized
positive boundary-sector weights.  Uniform sector bounds therefore
have no representation-cardinality loss, but label-dependent constants
still require a supremum or a state-weighted summability theorem.
\end{enumerate}
The Navier--Stokes note supplies no new type-compatible input for the
label-sector operator margin.

The transferable lesson is a scope distinction.  The V44 Schur theorem
is stable on its declared Hilbert space, but a one-particle or one-sector
margin does not automatically survive an unrestricted Fock lift or a
direct sum of sectors.  We derive both obstructions and their valid
repairs below rather than importing a Yang--Mills estimate.

\subsection{Direct sums of label sectors}

Let
\begin{equation}
 \mathcal K=\bigoplus_{\lambda\in\Lambda}\mathcal K_\lambda
 \label{eq:n9v45-direct-sum}
\end{equation}
and suppose sector $\lambda$ has a V44 energy-normalized Schur coupling
$C_\lambda$ with
\begin{equation}
 \gamma_\lambda:=\|C_\lambda\|<1.
 \label{eq:n9v45-sector-gamma}
\end{equation}

\begin{theorem}[Uniform direct-sum Schur criterion]
\label{thm:n9v45-direct-sum}
On the orthogonal direct sum,
\begin{align}
 \left\|\bigoplus_\lambda
  (I-C_\lambda C_\lambda^*)^{-1}\right\|
 &=\sup_\lambda
   \|(I-C_\lambda C_\lambda^*)^{-1}\|,
 \label{eq:n9v45-direct-inverse}\\
 &\le\sup_\lambda(1-\gamma_\lambda^2)^{-1}.
 \label{eq:n9v45-direct-bound}
\end{align}
Hence the sectorwise Schur inverses are uniformly bounded whenever
\begin{equation}
 \sup_\lambda\gamma_\lambda<1.
 \label{eq:n9v45-uniform-margin}
\end{equation}
Conversely, the facts $\gamma_\lambda<1$ for every fixed $\lambda$ and
$|\Lambda|<\infty$ at every regulator do not yield a regulator-uniform
bound if the retained sector set grows and its largest $\gamma_\lambda$
tends to one.
\end{theorem}

\begin{proof}
The operator norm of an orthogonal direct sum is the supremum of the
component norms.  Apply \eqref{eq:n9v44-Schur-inverse-bound} in each
sector.  For the last statement take a growing finite set containing
sectors with $\gamma_\lambda\uparrow1$.
\end{proof}

\begin{proposition}[Normalized state mixture]
\label{prop:n9v45-state-mixture}
Let $w_\lambda\ge0$, $\sum_\lambda w_\lambda=1$, and let
$\omega_\lambda$ be normalized positive sector states.  For any sector
family $F_\lambda$,
\begin{align}
 \left|\sum_\lambda w_\lambda
   \omega_\lambda(F_\lambda)\right|
 &\le\sum_\lambda w_\lambda
   |\omega_\lambda(F_\lambda)|,
 \label{eq:n9v45-mixture-bound}\\
 &\le\sup_\lambda\|F_\lambda\|.
 \label{eq:n9v45-mixture-sup}
\end{align}
More generally, a label-dependent bound
$|\omega_\lambda(F_\lambda)|\le M_\lambda$ is sufficient when
\begin{equation}
 \sum_\lambda w_\lambda M_\lambda<\infty.
 \label{eq:n9v45-weighted-sector-sum}
\end{equation}
No factor $|\Lambda|$ appears.
\end{proposition}

\begin{proof}
Use the triangle inequality, positivity and normalization of the
weights, and $|\omega_\lambda(F)|\le\|F\|$.
\end{proof}

\begin{corollary}[State-level Schur criterion]
\label{cor:n9v45-state-Schur}
For a normalized sector mixture of energy-normalized physical
resolvents, it is sufficient that
\begin{equation}
 \sum_\lambda
 \frac{w_\lambda}{1-\gamma_\lambda^2}<\infty.
 \label{eq:n9v45-weighted-Schur}
\end{equation}
This can hold even when the direct-sum operator norm is infinite.
\end{corollary}

\begin{proof}
Use \Cref{prop:n9v45-state-mixture} with
$M_\lambda=(1-\gamma_\lambda^2)^{-1}$.  For strict separation of the two
notions, take
$1-\gamma_n^2=2^{-n}$ and normalized weights proportional to $2^{-2n}$.
Then the weighted series converges, whereas
$\sup_n(1-\gamma_n^2)^{-1}=\infty$.
\end{proof}

\begin{firewall}[Normalized labels remove cardinality, not bad margins]
\label{fire:n9v45-cardinality}
A convex sector mixture removes an unweighted label-count factor.  It
does not make $(1-\gamma_\lambda^2)^{-1}$ uniform.  An operator claim
needs \eqref{eq:n9v45-uniform-margin}; a specified state expectation may
instead use \eqref{eq:n9v45-weighted-Schur}.  These are different
conclusions and cannot be interchanged.
\end{firewall}

\subsection{Fock amplification of a one-particle bound}

Let $\mathfrak h$ be a one-particle Hilbert space and let
$c=c^*$ be bounded.  On the bosonic or fermionic $N$-particle sector,
define
\begin{equation}
 d\Gamma_N(c)
 :=\sum_{j=1}^N
 I\otimes\cdots\otimes c\otimes\cdots\otimes I
 \label{eq:n9v45-second-quantization}
\end{equation}
restricted to the symmetric or antisymmetric tensor product.

\begin{theorem}[Spectator amplification]
\label{thm:n9v45-amplification}
For every $N$,
\begin{equation}
 \|d\Gamma_N(c)\|\le N\|c\|.
 \label{eq:n9v45-N-bound}
\end{equation}
The factor $N$ is sharp: if $c=\gamma I$ on a nonzero one-particle
space, then
\begin{equation}
 \|d\Gamma_N(c)\|=N|\gamma|.
 \label{eq:n9v45-N-sharp}
\end{equation}
Consequently a nonzero one-particle operator gives an unbounded
$d\Gamma(c)$ on unrestricted Fock space unless it is normalized by an
energy that grows with particle number.
\end{theorem}

\begin{proof}
The triangle inequality gives \eqref{eq:n9v45-N-bound}.  When
$c=\gamma I$, every summand is $\gamma I$ and their sum is
$N\gamma I$ on the $N$-particle sector.  Taking the orthogonal sum over
all $N$ makes these norms unbounded.
\end{proof}

\begin{corollary}[Occupation-cutoff condition]
\label{cor:n9v45-cutoff}
On the truncated Fock space $\mathcal F_{\le M}$,
\begin{equation}
 \|d\Gamma(c)\|\le M\|c\|.
 \label{eq:n9v45-cutoff-bound}
\end{equation}
Thus a V44 Schur argument applied to such an absolute cross operator
requires $M\|c\|<1$; cutoff removal needs at least
$M_n\|c_n\|$ bounded strictly below one.
\end{corollary}

\begin{proof}
Take the maximum of \eqref{eq:n9v45-N-bound} over $0\le N\le M$.
\end{proof}

\subsection{Energy-relative second quantization is stable}

The correct repair compares the one-particle form with an extensive
one-particle energy before lifting it.

\begin{theorem}[Fock-stable relative form]
\label{thm:n9v45-Fock-relative}
Let $a\ge0$ be self-adjoint on $\mathfrak h$, and let $c$ be a symmetric
form on $D(a^{1/2})$ satisfying
\begin{equation}
 |c[f,f]|\le\theta\langle f,af\rangle,
 \qquad 0\le\theta<1.
 \label{eq:n9v45-one-relative}
\end{equation}
Then on every bosonic or fermionic finite-particle sector,
\begin{equation}
 |d\Gamma_N(c)[\Psi,\Psi]|
 \le\theta\,
 d\Gamma_N(a)[\Psi,\Psi].
 \label{eq:n9v45-Fock-relative}
\end{equation}
The coefficient is independent of $N$.  The direct-sum inequality holds
on the natural finite-energy form domain of Fock space.
\end{theorem}

\begin{proof}
First take finite sums of simple tensors in the algebraic form core.
Apply \eqref{eq:n9v45-one-relative} to the $j$th variable while holding
all spectator variables fixed, and integrate over the spectators.  Sum
the resulting inequalities over $j=1,\dots,N$.  Symmetric and
antisymmetric restrictions preserve the inequality.  Close the forms on
each sector, then sum the nonnegative sector energies on the Fock form
domain.
\end{proof}

\begin{corollary}[Energy normalization removes spectator growth]
\label{cor:n9v45-energy-normalization}
Under \eqref{eq:n9v45-one-relative}, the Fock energy-normalized form
operator has norm at most $\theta$ on every nonvacuum sector and on their
orthogonal sum.  Therefore V41 resummation has cost
$(1-\theta)^{-1}$ independent of occupation number.
\end{corollary}

\begin{proof}
Equation \eqref{eq:n9v45-Fock-relative} is precisely the uniform form
bound required by \Cref{thm:n9v41-energy-normalization}; apply that
theorem sectorwise and take the direct sum.  The vacuum sector has zero
second-quantized form and is handled separately or after a positive
shift.
\end{proof}

\begin{firewall}[Nonlinear interactions are not second quantizations]
\label{fire:n9v45-nonlinear}
The theorem applies to $d\Gamma(c)$, a sum of one-particle forms.
Pairing terms, quartic interactions, gauge constraints coupling several
particles, and field-dependent chart maps require their own many-body
relative bounds or a different resummation.  One-particle locality does
not prove those estimates.
\end{firewall}

\subsection{Audit-adjusted SM interface}

\begin{definition}[Sector/Fock stability card (SFSC)]
\label{def:n9v45-SFSC}
SFSC requires, for every gauge/label/Fock component used in the reduced
SM--Einstein construction:
\begin{enumerate}
\item either a uniform sector Schur margin
\eqref{eq:n9v45-uniform-margin} for operator claims, or the explicit
state-weighted sum \eqref{eq:n9v45-weighted-Schur} for the specified
state observable;
\item an occupation-uniform many-body relative-form estimate, obtained
from \eqref{eq:n9v45-Fock-relative} only for genuine second
quantizations;
\item separate ownership of nonlinear, pairing, bad-cluster, and
boundary rows; and
\item compatibility with the V35--V44 chart, singular-tube, absorption,
and Schur cards.
\end{enumerate}
\end{definition}

\begin{theorem}[SFSC protects LSSC from sector amplification]
\label{thm:n9v45-SFSC}
If SFSC holds, the LSSC Schur margin and the V41 inverse bound persist
under the declared sector sum and Fock lift.  Normalized sector weights
introduce no cardinality factor, and genuine one-particle relative forms
introduce no occupation factor.  All non-second-quantized rows remain
separate loaders.
\end{theorem}

\begin{proof}
Use \Cref{thm:n9v45-direct-sum,prop:n9v45-state-mixture} for the sector
decomposition and \Cref{thm:n9v45-Fock-relative,
cor:n9v45-energy-normalization} for the Fock lift.  The final statement
is \Cref{fire:n9v45-nonlinear}.
\end{proof}

\begin{assessment}[Progress at ninth-series V45]
\label{ass:n9v45-status}
The compulsory audit has corrected the scope of the label Schur
programme.  Normalized sectors remove cardinality but not a closing
Schur margin.  Absolute one-particle bounds amplify linearly with
occupation, whereas energy-relative one-particle forms lift with the
same coefficient to all Fock sectors.  This identifies the exact form
of the next Standard Model matter/gauge loader.
\end{assessment}

\begin{programme}[Ninth-series V46: Fock-stable SM kinetic reference]
\label{prog:n9v45-V46}
V46 will build the free Standard Model matter/gauge Fock reference as a
direct sum of second-quantized positive one-particle energies, prove the
multi-species relative-form compiler, and isolate Yukawa, pairing,
quartic Higgs, gauge, and Einstein-source rows that are not covered by
second quantization.
\end{programme}

\begin{firewall}[Claim boundary after ninth-series V45]
\label{fire:n9v45-boundary}
V45 proves sector and Fock stability criteria.  It does not verify SFSC
for the physical Standard Model, control nonlinear interactions or
cutoff removal, prove weighted locality or OS reconstruction, or
construct the full SM--Einstein continuum.
\end{firewall}

\begin{theorem}[Ninth-series V45 result]
\label{thm:n9v45-result}
Sectorwise Schur bounds yield a uniform direct-sum resolvent only under a
uniform margin, while normalized state mixtures may instead use a
weighted sector sum.  A one-particle absolute norm amplifies by $N$ on
the $N$-particle sector, but a one-particle relative-form coefficient
$\theta$ lifts unchanged to the complete Fock energy.
\end{theorem}

\begin{proof}
Use \Cref{thm:n9v45-direct-sum,cor:n9v45-state-Schur,
thm:n9v45-amplification,thm:n9v45-Fock-relative,
thm:n9v45-SFSC}.
\end{proof}

\paragraph{Parent-version record.}
V45 was formed from
\texttt{Renewal\_SM\_Einstein\_Continuum\_Research\_Notes\_V44.tex}.
The V44 parent had $15\,768$ logical source lines, $571\,086$ bytes, and
SHA--256
\[
 \texttt{D5FCE806FABC59D8903631640CD6781055A863B97818B29564433F052BB9FC63}.
\]

% =============================================================================
% END APPENDED NINTH-SERIES V45 MATERIAL
% =============================================================================
% =============================================================================
% BEGIN APPENDED NINTH-SERIES V46 MATERIAL
% =============================================================================

\clearpage
\section{Ninth-series V46: a multi-species Fock-stable kinetic
reference and the nonlinear interaction boundary}
\label{sec:v9v46}

\subsection{The reduced finite-regulator species space}

At a finite regulator, let $\mathfrak S_B$ and $\mathfrak S_F$ be finite
sets of bosonic and fermionic species.  They may include physical
transverse gauge polarizations, Higgs components, quark and lepton
flavours, and any finite auxiliary sector that has not yet been removed.
For species $s$, let $\mathfrak h_s$ be its one-particle Hilbert space and
let $a_s\ge0$ be a positive self-adjoint one-particle energy.  Define
\begin{equation}
 \mathcal F_s=
 \begin{cases}
  \Gamma_{\rm sym}(\mathfrak h_s),&s\in\mathfrak S_B,\\
  \Gamma_{\rm asym}(\mathfrak h_s),&s\in\mathfrak S_F,
 \end{cases}
 \qquad
 \mathcal F=\widehat\bigotimes_{s\in\mathfrak S_B\cup\mathfrak S_F}
             \mathcal F_s.
 \label{eq:n9v46-Fock-space}
\end{equation}
Let $N_s=d\Gamma_s(I)$ and introduce nonnegative one-particle reference
shifts $\tau_s$:
\begin{equation}
 A_{s,\tau}:=d\Gamma_s(a_s)+\tau_sN_s,
 \qquad
 A_\tau:=\sum_s A_{s,\tau}.
 \label{eq:n9v46-reference}
\end{equation}
The sums are form sums on the natural finite-energy domain.  A global
resolvent shift $\sigma I$ may later be added to control the common
vacuum.

\subsection{One-particle lower terms become number terms}

Let $c_s$ be a symmetric one-particle form on $D(a_s^{1/2})$ satisfying
\begin{equation}
 |c_s[f,f]|
 \le\eta_s\langle f,a_sf\rangle+b_s\|f\|^2,
 \qquad 0\le\eta_s<1,quad b_s\ge0.
 \label{eq:n9v46-one-particle-bound}
\end{equation}
Write $C_s=d\Gamma_s(c_s)$ for its form second quantization and
$C=\sum_s C_s$.

\begin{lemma}[Shifted one-particle coefficient]
\label{lem:n9v46-one-shift}
If $\tau_s>0$, then
\begin{equation}
 |c_s[f,f]|
 \le\theta_s\langle f,(a_s+\tau_sI)f\rangle,
 \qquad
 \theta_s:=\max\{\eta_s,b_s/\tau_s\}.
 \label{eq:n9v46-theta-s}
\end{equation}
Thus $\theta_s<1$ whenever $\tau_s>b_s$.
\end{lemma}

\begin{proof}
This is the one-particle version of
\Cref{lem:n9v41-shifted-bound}.
\end{proof}

\begin{theorem}[Multi-species Fock relative-form compiler]
\label{thm:n9v46-multispecies}
Assume $\tau_s>b_s$ for every species and put
\begin{equation}
 \theta_*:=\max_s\theta_s<1.
 \label{eq:n9v46-theta-star}
\end{equation}
Then on the finite-particle form core, and hence on the closed
$A_\tau$ form domain,
\begin{equation}
 |C[\Psi,\Psi]|
 \le\theta_* A_\tau[\Psi,\Psi].
 \label{eq:n9v46-multispecies-bound}
\end{equation}
The coefficient is independent of every bosonic occupation number,
fermionic particle number, number of retained one-particle modes, and
finite species-sector decomposition.
\end{theorem}

\begin{proof}
Apply the Fock lift \eqref{eq:n9v45-Fock-relative} to
\eqref{eq:n9v46-theta-s} in each species:
\[
 |C_s[\Psi,\Psi]|
 \le\theta_s A_{s,\tau}[\Psi,\Psi].
\]
Sum absolute values and use
$\theta_s\le\theta_*$.  Closure gives the result on the full form
domain.  No dimension or occupation factor appears in the form
inequality.
\end{proof}

\begin{corollary}[Common vacuum and shifted resolvent]
\label{cor:n9v46-vacuum}
Both $A_\tau$ and $C$ vanish on the common Fock vacuum.  For
$\sigma>0$, the shifted reference
\begin{equation}
 A_{\tau,\sigma}:=A_\tau+\sigma I
 \label{eq:n9v46-global-shift}
\end{equation}
is strictly positive and $C$ remains relatively bounded with coefficient
at most $\theta_*$.  V41 therefore gives a resolvent inverse with cost
$(1-\theta_*)^{-1}$.
\end{corollary}

\begin{proof}
Second quantizations vanish in the zero-particle sector.  Adding a
positive scalar form only enlarges the reference.  Apply
\Cref{thm:n9v41-energy-normalization}.
\end{proof}

\begin{firewall}[A global shift does not replace a number shift]
\label{fire:n9v46-global-vs-number}
The one-particle lower term $b_s\|f\|^2$ lifts to $b_sN_s$, not to a
constant multiple of the identity.  On unrestricted bosonic Fock space,
$N_s$ is unbounded and cannot be controlled by a fixed global
$\sigma I$.  The per-particle reference term $\tau_sN_s$, a physical
one-particle gap, or another extensive positive energy is needed before
\eqref{eq:n9v46-multispecies-bound} follows.
\end{firewall}

\subsection{Gauge and constraint reduction}

Let $P_{\rm phys}$ be an orthogonal finite-regulator Gauss/constraint
projection on $\mathcal F$.

\begin{proposition}[Restriction to a reducing physical sector]
\label{prop:n9v46-physical-restriction}
Suppose $P_{\rm phys}$ reduces both $A_\tau$ and the form $C$.  Then on
$P_{\rm phys}\mathcal F$,
\begin{equation}
 |C[\Psi,\Psi]|
 \le\theta_*A_\tau[\Psi,\Psi]
 \label{eq:n9v46-physical-bound}
\end{equation}
with the same coefficient.  If the projection does not reduce the
operator, the physical resolvent is instead a Schur complement and must
be treated by V44.
\end{proposition}

\begin{proof}
The first statement is simply the restriction of
\eqref{eq:n9v46-multispecies-bound} to the reducing subspace.  The second
is \Cref{thm:n9v44-Schur} applied to the physical and orthogonal sectors.
\end{proof}

\begin{proposition}[Compact product-group sector labels]
\label{prop:n9v46-product-group}
At finite representation cutoff, a compact product gauge group
\begin{equation}
 K=K_1\times\cdots\times K_J
 \label{eq:n9v46-product-K}
\end{equation}
has irreducible labels
$\lambda=(\lambda_1,\dots,\lambda_J)$.  A diagonal cut projection pairs
a left label only with the componentwise dual
$\lambda^*=(\lambda_1^*,\dots,\lambda_J^*)$.  For a $U(1)$ factor this
means opposite integer charge.  Normalized physical traces are therefore
convex sums over the product labels, and uniform component estimates
carry no factor equal to the number or dimensions of retained sectors.
\end{proposition}

\begin{proof}
Irreducible unitary representations of a compact direct product are
tensor products of irreducibles of the factors.  The invariant subspace
of $V_\lambda\otimes V_\mu$ is nonzero exactly when
$\mu=\lambda^*$ componentwise, by Schur orthogonality.  Haar averaging
is the orthogonal projection onto these invariant lines.  Normalizing
the trace divides the nonnegative sector contributions by their sum, so
the weights form a probability vector.  Apply
\Cref{prop:n9v45-state-mixture}.
\end{proof}

\begin{remark}[Standard Model gauge labels]
\label{rem:n9v46-SM-labels}
For the compact finite-regulator gauge group
$SU(3)\times SU(2)\times U(1)$, the paired boundary label has conjugate
$SU(3)$ representation, dual $SU(2)$ representation, and opposite
hypercharge.  This finite-cut statement controls label cardinality; it
does not prove anomaly cancellation, the continuum Gauss law, or a
uniform sector gap.
\end{remark}

\subsection{Metric/source variations of quadratic rows}

Let $g$ denote a finite-regulator metric variable and let
$a_s(g)$ be the one-particle kinetic energy.  For a metric direction $h$,
suppose its first variation is a symmetric form
\begin{equation}
 c_{s,h}:=D_ga_s[h]
 \label{eq:n9v46-metric-variation}
\end{equation}
satisfying \eqref{eq:n9v46-one-particle-bound} with constants uniform in
the declared metric chart and regulator.

\begin{theorem}[Quadratic metric-source lift]
\label{thm:n9v46-metric-lift}
Under the uniform one-particle bounds, the complete quadratic
matter/gauge metric variation
\begin{equation}
 C_h=\sum_s d\Gamma_s(c_{s,h})
 \label{eq:n9v46-source-lift}
\end{equation}
obeys \eqref{eq:n9v46-multispecies-bound}.  Its repeated symmetric form
insertions are therefore resummed by V41 with no occupation or
species-cardinality loss.
\end{theorem}

\begin{proof}
Apply \Cref{thm:n9v46-multispecies} to the family $c_{s,h}$ and then
\Cref{thm:n9v41-Neumann}.
\end{proof}

\begin{firewall}[A metric derivative must include the changing physical
subspace]
\label{fire:n9v46-moving-projection}
If $P_{\rm phys}(g)$ depends on the metric or chart, differentiating the
compressed kinetic energy also produces
$D_gP_{\rm phys}[h]$ rows.  These are not contained in
\eqref{eq:n9v46-source-lift}; they require the chart/projector response
and Schur analysis of V35--V44.
\end{firewall}

\subsection{Why the nonlinear Standard Model rows are separate}

The finite-mode occupation basis gives a sharp obstruction.  In one
bosonic mode let $a,a^*$ be the annihilation and creation operators and
$N=a^*a$.

\begin{proposition}[Quartic occupation obstruction]
\label{prop:n9v46-quartic}
On the number vector $|n\rangle$,
\begin{equation}
 (a^*)^2a^2|n\rangle=n(n-1)|n\rangle.
 \label{eq:n9v46-quartic-number}
\end{equation}
Hence the positive quartic form $(a^*)^2a^2=N(N-1)$ is not relatively
bounded by $N$ on full bosonic Fock space.  On the cutoff sector
$N\le M$ it satisfies
\begin{equation}
 0\le N(N-1)\le(M-1)N.
 \label{eq:n9v46-quartic-cutoff}
\end{equation}
\end{proposition}

\begin{proof}
Use $a|n\rangle=\sqrt n|n-1\rangle$ twice and then apply $a^*$ twice.
The ratio $n(n-1)/n=n-1$ is unbounded.  The cutoff inequality holds on
each number eigenvalue $0\le n\le M$.
\end{proof}

\begin{corollary}[Positive quartic absorption]
\label{cor:n9v46-quartic-absorption}
A positive Higgs quartic row may be included in the V42 reference form,
provided its closedness, ownership, gauge compatibility, and metric
dependence are proved.  Treating it as a small perturbation of a linear
number energy is not occupation-uniform.
\end{corollary}

\begin{proof}
The first statement is the positive-row absorption theorem
\Cref{thm:n9v42-absorption}; the second is
\Cref{prop:n9v46-quartic}.
\end{proof}

\begin{lemma}[Creation and annihilation growth]
\label{lem:n9v46-creation}
For $f\in\mathfrak h$ and vectors in the natural domains,
\begin{align}
 \|a(f)\Psi\|&\le\|f\|\,\|N^{1/2}\Psi\|,
 \label{eq:n9v46-annihilation}\\
 \|a^*(f)\Psi\|&\le\|f\|\,\|(N+I)^{1/2}\Psi\|.
 \label{eq:n9v46-creation}
\end{align}
\end{lemma}

\begin{proof}
Decompose $\Psi$ into particle-number sectors.  On the $n$-particle
sector the annihilation norm is at most $\sqrt n\|f\|$, and the creation
norm at most $\sqrt{n+1}\|f\|$.  Square and sum the orthogonal sector
contributions.
\end{proof}

\begin{proposition}[Fixed-cutoff polynomial degree]
\label{prop:n9v46-degree}
At a finite mode regulator, a Wick monomial containing $d$ bosonic
creation/annihilation factors with a bounded coefficient tensor has norm
on $N\le M$ bounded by
\begin{equation}
 C_d(M+d)^{d/2}.
 \label{eq:n9v46-degree-bound}
\end{equation}
Thus, relative to a number energy of order $M$, a degree-$d$ interaction
can carry a worst-case cutoff ratio of order
$(M+d)^{d/2}/M$.  A coupling or stability estimate must compensate this
growth when $d>2$.
\end{proposition}

\begin{proof}
Apply \Cref{lem:n9v46-creation} successively.  At the $j$th step the
intermediate occupation is at most $M+j$, giving a product of $d$ square
roots bounded by $(M+d)^{d/2}$.  Contracting with the finite-mode
coefficient tensor contributes the constant $C_d$.  Divide by the
linear number scale for the last statement.
\end{proof}

\begin{firewall}[Interaction ledger]
\label{fire:n9v46-interactions}
Yukawa vertices, non-Abelian cubic and quartic gauge vertices, the Higgs
quartic, pairing terms, and metric derivatives of these interactions are
not sums of one-particle forms.  They are therefore not controlled by
\Cref{thm:n9v46-multispecies}.  Positive stable rows may be absorbed;
the remaining rows require cutoff-explicit relative, cluster, or
state-moment estimates.
\end{firewall}

\subsection{Fock-stable kinetic reference card}

\begin{definition}[Fock-stable kinetic reference card (FSKRC)]
\label{def:n9v46-FSKRC}
FSKRC requires:
\begin{enumerate}
\item the reduced one-particle species spaces and positive energies
$a_s$ on the actual finite-regulator physical carrier;
\item uniform bounds \eqref{eq:n9v46-one-particle-bound} for every
quadratic metric/source row;
\item an owned extensive positive term controlling every $b_sN_s$ and a
uniform $\theta_*<1$;
\item a reducing physical projection or an LSSC Schur margin;
\item a separate nonlinear interaction ledger satisfying the firewall
above; and
\item locality, chart-cocycle, singular-tube, and cutoff estimates
uniform in the same parameters.
\end{enumerate}
\end{definition}

\begin{theorem}[FSKRC loads the quadratic part of SFSC]
\label{thm:n9v46-FSKRC}
FSKRC gives an occupation- and species-uniform relative-form coefficient
for all quadratic Standard Model kinetic and metric-source rows.  These
rows load the V41 resolvent resummation and the quadratic part of SFSC.
The nonlinear ledger remains an independent interface.
\end{theorem}

\begin{proof}
Use \Cref{thm:n9v46-multispecies,cor:n9v46-vacuum,
prop:n9v46-physical-restriction,thm:n9v46-metric-lift}.  The scope
restriction is \Cref{fire:n9v46-interactions}.
\end{proof}

\begin{assessment}[Progress at ninth-series V46]
\label{ass:n9v46-status}
The free multi-species Standard Model sector now has a precise
occupation-stable form compiler.  Lower one-particle constants become
number operators and require an extensive reference term; a global
resolvent shift cannot replace it.  Quadratic metric-source variations
lift uniformly.  Quartic and other nonlinear interactions have been
separated with their explicit occupation-growth obstruction.
\end{assessment}

\begin{programme}[Ninth-series V47: positive Higgs quartic reference]
\label{prog:n9v46-V47}
V47 will analyze the stable finite-regulator Higgs potential as an
absorbable positive reference row after vacuum-energy normalization.  It
will derive a coercive number/field-amplitude bound and isolate the
Yukawa and gauge couplings relative to that enlarged reference without
assuming that a quartic field term is a one-particle operator.
\end{programme}

\begin{firewall}[Claim boundary after ninth-series V46]
\label{fire:n9v46-boundary}
V46 proves the abstract multi-species kinetic compiler.  It does not
construct the physical reduced one-particle spaces, prove FSKRC, control
the nonlinear Standard Model interactions, remove the auxiliary number
shifts, establish OS reconstruction, or construct the full continuum.
\end{firewall}

\begin{theorem}[Ninth-series V46 result]
\label{thm:n9v46-result}
Uniform one-particle bounds
$|c_s|\le\eta_sa_s+b_sI$ lift to the full multi-species Fock space with
coefficient
$\max_s\max\{\eta_s,b_s/\tau_s\}<1$ relative to
$\sum_s[d\Gamma(a_s)+\tau_sN_s]$.  Quadratic metric-source rows are
therefore occupation stable, while the nonlinear Standard Model
vertices require a separate ledger.
\end{theorem}

\begin{proof}
Use \Cref{lem:n9v46-one-shift,thm:n9v46-multispecies,
thm:n9v46-metric-lift,prop:n9v46-quartic,
fire:n9v46-interactions}.
\end{proof}

\paragraph{Parent-version record.}
V46 was formed from
\texttt{Renewal\_SM\_Einstein\_Continuum\_Research\_Notes\_V45.tex}.
The V45 parent had $16\,127$ logical source lines, $584\,595$ bytes, and
SHA--256
\[
 \texttt{77B38DA57C32BDB8F28F3F3AE242CD092E04CFBE7320B4EF960076BD0DC2C9C0}.
\]

% =============================================================================
% END APPENDED NINTH-SERIES V46 MATERIAL
% =============================================================================
% =============================================================================
% BEGIN APPENDED NINTH-SERIES V47 MATERIAL
% =============================================================================

\clearpage
\section{Ninth-series V47: positive Higgs reference, Yukawa form bounds,
and vacuum-energy ownership}
\label{sec:v9v47}

\subsection{Finite-regulator Higgs potential}

Let $\Lambda$ be a finite set of regulator cells with positive weights
$w_x$.  At each $x$, let $\Phi_x\in\mathbb R^d$ denote the real
components of the regulated Higgs multiplet.  Consider
\begin{equation}
 V_{H,x}(\Phi_x)
 =\lambda_x|\Phi_x|^4-m_x^2|\Phi_x|^2,
 \qquad \lambda_x>0,quad m_x^2\ge0.
 \label{eq:n9v47-original-potential}
\end{equation}
Put
\begin{equation}
 v_x^2:=\frac{m_x^2}{2\lambda_x},
 \qquad
 e_x:=\frac{m_x^4}{4\lambda_x}.
 \label{eq:n9v47-v-e}
\end{equation}

\begin{lemma}[Exact square completion]
\label{lem:n9v47-square}
For every $x$,
\begin{equation}
 V_{H,x}(\Phi_x)
 =\lambda_x\bigl(|\Phi_x|^2-v_x^2\bigr)^2-e_x.
 \label{eq:n9v47-square}
\end{equation}
Thus the finite-volume potential decomposes exactly as
\begin{equation}
 \sum_xw_xV_{H,x}=P_H-E_{\rm vac},
 \quad
 P_H:=\sum_xw_x\lambda_x(|\Phi_x|^2-v_x^2)^2\ge0,
 \quad
 E_{\rm vac}:=\sum_xw_xe_x.
 \label{eq:n9v47-PH-E}
\end{equation}
\end{lemma}

\begin{proof}
Expand the square and use
$2\lambda_xv_x^2=m_x^2$ and
$\lambda_xv_x^4=e_x$.
\end{proof}

\begin{proposition}[Quartic amplitude coercivity]
\label{prop:n9v47-coercivity}
For every $r,v\ge0$,
\begin{equation}
 (r^2-v^2)^2\ge\frac12r^4-v^4.
 \label{eq:n9v47-scalar-coercivity}
\end{equation}
Consequently
\begin{equation}
 P_H\ge
 \frac12\sum_xw_x\lambda_x|\Phi_x|^4
 -\sum_xw_x\lambda_xv_x^4.
 \label{eq:n9v47-PH-coercivity}
\end{equation}
If $\lambda_x\ge\lambda_*>0$, this controls the weighted fourth-power
amplitude up to the displayed finite volume constant.
\end{proposition}

\begin{proof}
Young's inequality
$2v^2r^2\le\frac12r^4+2v^4$ gives
$r^4-2v^2r^2+v^4\ge\frac12r^4-v^4$.
Multiply by $w_x\lambda_x$ and sum.
\end{proof}

\subsection{Closed positive absorption}

Let $T_H$ be a nonnegative closed Higgs kinetic form on
$L^2(\mathbb R^{d|\Lambda|})$, possibly tensored with the other finite
regulator sectors.  Regard $P_H$ as a nonnegative multiplication form.

\begin{theorem}[Closed Higgs quartic reference]
\label{thm:n9v47-closed-reference}
If the common form domain
\begin{equation}
 \mathcal Q_H:=D(T_H)\cap L^2(P_H,d\Phi)
 \label{eq:n9v47-domain}
\end{equation}
is dense, then
\begin{equation}
 B_H:=T_H+P_H
 \label{eq:n9v47-BH}
\end{equation}
is a densely defined closed nonnegative form.  It is therefore a valid
positive row for the V42 absorption theorem.
\end{theorem}

\begin{proof}
Both $T_H$ and the nonnegative multiplication form $P_H$ are closed.
A sequence Cauchy in the sum form norm is Cauchy in each form norm; the
two limits agree in the ambient $L^2$ space.  Hence the intersection is
complete in the sum norm.  Density is assumed.  Apply
\Cref{thm:n9v42-absorption}.
\end{proof}

\begin{remark}[Gauge-covariant kinetic energy should be kept whole]
\label{rem:n9v47-covariant-whole}
At finite lattice regulator the gauge--Higgs kinetic row has the form
\begin{equation}
 P_{gH}=\sum_{(x,y)}w_{xy}
 \left|U_{xy}\Phi_y-\Phi_x\right|^2\ge0.
 \label{eq:n9v47-gauge-Higgs}
\end{equation}
Its expansion contains signed cross terms, but the complete covariant
square is nonnegative.  Absorbing the complete owned row preserves this
positivity; assigning its expanded pieces to unrelated ledgers can lose
the sign and double-count the link interaction.
\end{remark}

\begin{corollary}[Positive Higgs/gauge reference]
\label{cor:n9v47-positive-reference}
If $P_{gH}$ is closed on a dense common domain with $B_H$, then
\begin{equation}
 B_{Hg}:=T_H+P_H+P_{gH}
 \label{eq:n9v47-BHg}
\end{equation}
is a closed nonnegative reference form and may absorb the stable Higgs
potential and gauge-covariant Higgs gradient simultaneously.
\end{corollary}

\begin{proof}
Repeat the nonnegative form-sum argument in
\Cref{thm:n9v47-closed-reference}.
\end{proof}

\subsection{Lower-degree fields are infinitesimal relative to the
quartic}

\begin{lemma}[Sharp scalar quartic Young bounds]
\label{lem:n9v47-Young}
For $r\ge0$ and $\delta>0$,
\begin{align}
 r^2&\le\delta r^4+\frac1{4\delta},
 \label{eq:n9v47-r2}\\
 r^3&\le\delta r^4+\frac{27}{256\delta^3},
 \label{eq:n9v47-r3}\\
 br&\le\delta r^4+
 \frac{3}{4^{4/3}}b^{4/3}\delta^{-1/3}
 \qquad(b\ge0).
 \label{eq:n9v47-br}
\end{align}
\end{lemma}

\begin{proof}
Maximize respectively
$r^2-\delta r^4$, $r^3-\delta r^4$, and
$br-\delta r^4$ over $r\ge0$.  The critical points are
$(2\delta)^{-1/2}$, $3/(4\delta)$, and
$(b/(4\delta))^{1/3}$, giving the displayed maxima.
\end{proof}

\begin{proposition}[Shifted-well lower-degree bound]
\label{prop:n9v47-shifted-Young}
Fix $v\ge0$, $\lambda>0$, and $0<\varepsilon<1$.  For every $b\ge0$
there is an explicit finite constant $C(\varepsilon,b,\lambda,v)$ such
that
\begin{equation}
 br\le\varepsilon\lambda(r^2-v^2)^2
       +C(\varepsilon,b,\lambda,v)
 \qquad(r\ge0).
 \label{eq:n9v47-shifted-linear}
\end{equation}
One valid choice is
\begin{equation}
 C=\varepsilon\lambda v^4
 +\frac{3}{4^{4/3}}b^{4/3}
   (\varepsilon\lambda/2)^{-1/3}.
 \label{eq:n9v47-shifted-constant}
\end{equation}
\end{proposition}

\begin{proof}
By \eqref{eq:n9v47-scalar-coercivity},
\[
 \varepsilon\lambda(r^2-v^2)^2
 \ge(\varepsilon\lambda/2)r^4-\varepsilon\lambda v^4.
\]
Apply \eqref{eq:n9v47-br} with
$\delta=\varepsilon\lambda/2$ and rearrange.
\end{proof}

\subsection{Finite-regulator Yukawa row}

Let $\mathcal F_F$ be the finite-regulator fermion Fock space.  Suppose
$Y_x=Y_x^*$ is the regulated fermion bilinear appearing in the local
Yukawa interaction and
\begin{equation}
 \|Y_x\|\le y_x<\infty.
 \label{eq:n9v47-Y-bound}
\end{equation}
Let $M_x(\Phi_x)$ be a real linear component of the Higgs multiplet, so
$|M_x(\Phi_x)|\le|\Phi_x|$.  Define the symmetric multiplication/operator
form
\begin{equation}
 V_Y=\sum_xw_xM_x(\Phi_x)Y_x.
 \label{eq:n9v47-Yukawa}
\end{equation}

\begin{theorem}[Local Yukawa--quartic form bound]
\label{thm:n9v47-Yukawa}
For every $0<\varepsilon<1$,
\begin{equation}
 |V_Y[\Psi,\Psi]|
 \le\varepsilon P_H[\Psi,\Psi]
 +C_{Y,\varepsilon}\|\Psi\|^2,
 \label{eq:n9v47-Yukawa-bound}
\end{equation}
where one may take
\begin{equation}
 C_{Y,\varepsilon}
 :=\sum_xw_x
 C(\varepsilon,y_x,\lambda_x,v_x)
 \label{eq:n9v47-Yukawa-constant}
\end{equation}
with $C$ from \eqref{eq:n9v47-shifted-constant}.  Thus the regulated
Yukawa row is infinitesimally form bounded relative to the positive Higgs
quartic.  The lower constant retains its volume and fermion-cutoff
dependence explicitly.
\end{theorem}

\begin{proof}
Pointwise in $\Phi$ and the other coordinates,
\[
 |\langle\Psi,M_x(\Phi_x)Y_x\Psi\rangle_{\mathcal F_F}|
 \le y_x|\Phi_x|\,\|\Psi\|_{\mathcal F_F}^2.
\]
Apply \eqref{eq:n9v47-shifted-linear} at each cell, multiply by $w_x$,
sum, and integrate over the remaining variables.
\end{proof}

\begin{corollary}[Shifted Yukawa resolvent]
\label{cor:n9v47-Yukawa-resolvent}
At fixed regulator, $B_H+V_Y$ is closed and lower semibounded.  After a
shift
\begin{equation}
 \sigma>\frac{C_{Y,\varepsilon}}{1-\varepsilon},
 \label{eq:n9v47-Y-shift}
\end{equation}
the Yukawa residual has strict coefficient
$\varepsilon+C_{Y,\varepsilon}/\sigma<1$ relative to
$B_H+\sigma I$, and V41 resums its symmetric form insertions.
\end{corollary}

\begin{proof}
Apply \Cref{cor:n9v42-positive-negative} to the absolute form bound
\eqref{eq:n9v47-Yukawa-bound}, or directly use the additive shifted
coefficient in \Cref{rem:n9v42-shift-coefficients}.
\end{proof}

\begin{firewall}[The fixed-regulator Yukawa constant is not uniform]
\label{fire:n9v47-Y-uniform}
The norm $y_x$ can grow with the fermion mode or representation cutoff,
and $C_{Y,\varepsilon}$ is extensive in the number of cells.  A global
resolvent shift growing with volume is not a cofinal uniform estimate.
The physical programme must obtain a local cluster/form bound, absorb an
extensive positive reference, or prove a uniform density estimate before
using \Cref{cor:n9v47-Yukawa-resolvent} in the continuum limit.
\end{firewall}

\subsection{Vacuum energy is an Einstein source}

For a fixed background metric, subtracting $E_{\rm vac}$ in
\eqref{eq:n9v47-PH-E} merely multiplies an unnormalized partition
function by a scalar.  When the metric is dynamical, the cell weights
approximate the volume density and the same row varies with the metric.

\begin{proposition}[Metric variation of the completed-square constant]
\label{prop:n9v47-vacuum-variation}
In continuum notation let
\begin{equation}
 S_{\rm vac}[g]=-e\int_M d\operatorname{vol}_g
 \label{eq:n9v47-vacuum-action}
\end{equation}
with $e$ independent of $g$.  For a covariant metric variation $h$, using
$\delta d\operatorname{vol}_g=\frac12\operatorname{tr}_g(h)
d\operatorname{vol}_g$,
\begin{equation}
 DS_{\rm vac}[h]
 =-\frac e2\int_M\operatorname{tr}_g(h)
       \,d\operatorname{vol}_g.
 \label{eq:n9v47-vacuum-first}
\end{equation}
Thus the square-completion constant contributes a stress tensor
proportional to $g^{-1}$ and is algebraically a cosmological-constant
row.  It cannot be discarded from the Einstein variation without an
explicit cosmological renormalization condition.
\end{proposition}

\begin{proof}
Differentiate \eqref{eq:n9v47-vacuum-action} using the standard
determinant identity
$D\sqrt{\det g}[h]=\frac12\sqrt{\det g}\,
\operatorname{tr}(g^{-1}h)$.
\end{proof}

\begin{remark}[Sign convention]
\label{rem:n9v47-vacuum-sign}
The tensor sign inferred from \eqref{eq:n9v47-vacuum-first} depends on
whether the Einstein variation is taken with respect to $g_{\mu\nu}$ or
$g^{\mu\nu}$ and on the stress-tensor convention.  The invariant content
is that the variation is proportional to the metric and shifts the
cosmological row.  V24's covariant-metric convention fixes the final sign
when this term is inserted there.
\end{remark}

\begin{firewall}[Ground-energy normalization versus gravity]
\label{fire:n9v47-ground-energy}
Ground-energy subtraction can make a matter Hamiltonian nonnegative and
can cancel inside a normalized fixed-background state.  In a coupled
metric path integral, a regulator- and metric-dependent ground energy is
part of the gravitational effective action.  Its volume, boundary, and
metric-derivative terms must be owned by counterterms and cannot be
silently canceled between numerator and denominator before the Einstein
source is derived.
\end{firewall}

\subsection{Higgs positive-reference card}

\begin{definition}[Higgs positive-reference card (HPRC)]
\label{def:n9v47-HPRC}
HPRC requires:
\begin{enumerate}
\item a regulator-uniform positive quartic floor and the exact square
completion \eqref{eq:n9v47-PH-E};
\item closedness and gauge compatibility of the complete positive form
$B_{Hg}$;
\item explicit local Yukawa bilinear bounds $y_x$ and a cofinal treatment
of the extensive constant in \eqref{eq:n9v47-Yukawa-bound};
\item ownership of $E_{\rm vac}$ and its metric variations in the
cosmological/counterterm ledger;
\item separate bounds for gauge self-interactions and other nonlinear
rows not contained in $P_{gH}$; and
\item compatibility with FSKRC, CPRAC, RSCC, and the singular-tube cards.
\end{enumerate}
\end{definition}

\begin{theorem}[HPRC loads the stable Higgs/Yukawa form interface]
\label{thm:n9v47-HPRC}
Under HPRC, the positive Higgs potential and gauge-covariant gradient are
valid absorbed reference rows.  The regulated Yukawa form is a strict
residual after the declared local/extensive treatment, and its repeated
symmetric insertions are controlled by V41.  The induced vacuum energy
is retained as an Einstein cosmological source.
\end{theorem}

\begin{proof}
Use \Cref{thm:n9v47-closed-reference,
cor:n9v47-positive-reference,thm:n9v47-Yukawa,
cor:n9v47-Yukawa-resolvent,prop:n9v47-vacuum-variation}.  HPRC item 3 is
needed because \Cref{fire:n9v47-Y-uniform} prevents a volume-uniform
claim from the finite-regulator estimate alone.
\end{proof}

\begin{assessment}[Progress at ninth-series V47]
\label{ass:n9v47-status}
The stable Higgs sector now has a correct positive-reference
architecture.  Completing the square yields a closed quartic form that
controls finite-regulator Yukawa multiplication infinitesimally.  The
result exposes two cofinal debts rather than hiding them: cutoff/volume
growth of the Yukawa lower constant and the metric-dependent vacuum
energy that shifts the Einstein cosmological source.
\end{assessment}

\begin{programme}[Ninth-series V48: local Yukawa block compiler]
\label{prog:n9v47-V48}
V48 will replace the extensive Yukawa lower constant by a local block
estimate.  It will derive a conditional tensorization/ownership rule,
show how blockwise quartic coercivity pays the Higgs amplitude, and keep
fermion bilinear cutoff growth as an explicit activity tested against a
cofinal local rarity or coupling schedule.
\end{programme}

\begin{firewall}[Claim boundary after ninth-series V47]
\label{fire:n9v47-boundary}
V47 proves finite-regulator positive absorption and a local algebraic
Yukawa bound.  It does not prove HPRC uniformly in volume and cutoff,
renormalize the cosmological term, control the remaining gauge
interactions, establish OS reconstruction, or construct the full
SM--Einstein continuum.
\end{firewall}

\begin{theorem}[Ninth-series V47 result]
\label{thm:n9v47-result}
The finite-regulator Higgs potential splits exactly into a closed
nonnegative quartic well and a vacuum-energy constant.  The quartic well
infinitesimally controls every bounded regulated Yukawa bilinear through
\eqref{eq:n9v47-Yukawa-bound}; the constant is a metric-dependent
cosmological source and must remain in the Einstein ledger.
\end{theorem}

\begin{proof}
Use \Cref{lem:n9v47-square,thm:n9v47-closed-reference,
thm:n9v47-Yukawa,prop:n9v47-vacuum-variation,
thm:n9v47-HPRC}.
\end{proof}

\paragraph{Parent-version record.}
V47 was formed from
\texttt{Renewal\_SM\_Einstein\_Continuum\_Research\_Notes\_V46.tex}.
The V46 parent had $16\,543$ logical source lines, $599\,926$ bytes, and
SHA--256
\[
 \texttt{728ACC1496E6CD49D7B507F3B517F7A3309E8FB0118AACE4C6A2B285012D14F9}.
\]

% =============================================================================
% END APPENDED NINTH-SERIES V47 MATERIAL
% =============================================================================
% =============================================================================
% BEGIN APPENDED NINTH-SERIES V48 MATERIAL
% =============================================================================

\clearpage
\section{Ninth-series V48: local Yukawa ownership and ground-normalized
block resummation}
\label{sec:v9v48}

\subsection{Owned finite-range blocks}

Let $\mathcal B_n$ be a finite family of blocks covering the regulator
lattice.  To each block $B$ associate a fixed-radius ownership collar
$B^+$.  Every local Higgs, Yukawa, and gauge--Higgs term is assigned to
exactly one block whose collar contains its full support.  Write
\begin{equation}
 P_H=\sum_{B\in\mathcal B_n}P_B,
 \qquad
 V_Y=\sum_{B\in\mathcal B_n}Y_B,
 \label{eq:n9v48-owned-sum}
\end{equation}
with $P_B\ge0$ and no term repeated between blocks.

For fixed local block size, the V47 estimate gives, for every
$0<\varepsilon<1$,
\begin{equation}
 |Y_B[\Psi,\Psi]|
 \le\varepsilon P_B[\Psi,\Psi]
 +C_{B,n}(\varepsilon)\|\Psi\|^2,
 \label{eq:n9v48-local-Yukawa}
\end{equation}
where $C_{B,n}$ retains the local fermion-cutoff and metric dependence.
The crucial change from \eqref{eq:n9v47-Yukawa-bound} is that
$C_{B,n}$ is a local constant, not a sum over the full volume.

\begin{definition}[Exact ownership]
\label{def:n9v48-ownership}
The decomposition \eqref{eq:n9v48-owned-sum} is exactly owned when every
elementary interaction support occurs in one and only one pair
$(B,B^+)$.  Terms meeting two collars are assigned by a fixed
regulator-independent ordering rule and remain supported in the chosen
collar.
\end{definition}

\begin{lemma}[Ownership prevents multiplicity inflation]
\label{lem:n9v48-ownership}
Under exact ownership,
\begin{equation}
 \sum_B P_B=P_H,
 \qquad
 \sum_B Y_B=V_Y
 \label{eq:n9v48-exact-sums}
\end{equation}
as forms on the common finite-regulator core, and the number of block
labels never multiplies the coefficient of one elementary term.
\end{lemma}

\begin{proof}
Partition the finite elementary interaction index set according to its
unique owner.  Reordering finite form sums proves
\eqref{eq:n9v48-exact-sums}; disjointness of the index classes proves the
last statement.
\end{proof}

\subsection{Local lower bounds and ground normalization}

Let $T_B\ge0$ be the owned local kinetic/covariant-gradient form and set
\begin{equation}
 H_B:=T_B+P_B+Y_B.
 \label{eq:n9v48-block-H}
\end{equation}

\begin{theorem}[Local Yukawa block lower bound]
\label{thm:n9v48-local-lower}
Assume \eqref{eq:n9v48-local-Yukawa}.  Then $H_B$ is closed and lower
semibounded on $D(T_B+P_B)$, and
\begin{equation}
 H_B\ge T_B+(1-\varepsilon)P_B-C_{B,n}(\varepsilon)I
 \ge-C_{B,n}(\varepsilon)I
 \label{eq:n9v48-block-lower}
\end{equation}
in form sense.  In particular
\begin{equation}
 E_B:=\inf\operatorname{spec}H_B
 \ge-C_{B,n}(\varepsilon).
 \label{eq:n9v48-ground-lower}
\end{equation}
\end{theorem}

\begin{proof}
The absolute form bound in \eqref{eq:n9v48-local-Yukawa} has relative
coefficient $\varepsilon<1$ against $T_B+P_B$ and a finite lower term.
The KLMN conclusion follows from \Cref{thm:n9v41-energy-normalization},
and subtracting the absolute bound gives
\eqref{eq:n9v48-block-lower}.  The spectral lower bound is the
variational characterization of the infimum of the spectrum.
\end{proof}

\begin{corollary}[Exact internal contraction]
\label{cor:n9v48-contraction}
Define
\begin{equation}
 \widehat H_B:=H_B-E_BI\ge0.
 \label{eq:n9v48-ground-normalized}
\end{equation}
Then for every $t\ge0$,
\begin{equation}
 \|e^{-t\widehat H_B}\|\le1.
 \label{eq:n9v48-contraction}
\end{equation}
All internal Yukawa and positive Higgs occurrences in that block are
already resummed inside this contraction; no separate factor
$C_{B,n}^N$ is charged for their Duhamel multiplicity.
\end{corollary}

\begin{proof}
The spectral theorem gives
$\|e^{-t\widehat H_B}\|
=\sup_{\lambda\in\operatorname{spec}\widehat H_B}e^{-t\lambda}\le1$.
The second statement follows because $\widehat H_B$ is the self-adjoint
operator associated with the complete block form, rather than a
truncated expansion of $Y_B$.
\end{proof}

\begin{firewall}[Contraction is not small activity]
\label{fire:n9v48-contraction-small}
Equation \eqref{eq:n9v48-contraction} bounds arbitrary internal
multiplicity, but it does not make the block contribution close to the
free block, nor does it make a cross-boundary interaction small.  A
cluster expansion still needs a local activity or boundary estimate and
the ground-energy normalization must cancel or be owned correctly.
\end{firewall}

\subsection{Exact tensorization for disjoint collars}

Let $\mathcal C\subset\mathcal B_n$ be a family whose collars $B^+$ are
pairwise disjoint.  Assume the extended local Hilbert space factorizes as
\begin{equation}
 \mathcal H_{\mathcal C}
 =\widehat\bigotimes_{B\in\mathcal C}\mathcal H_B
 \label{eq:n9v48-tensor-space}
\end{equation}
and the owned block operators act on their corresponding factors.

\begin{theorem}[Same-family tensorized contraction]
\label{thm:n9v48-tensor-contraction}
For a finite disjoint family,
\begin{align}
 H_{\mathcal C}&:=\sum_{B\in\mathcal C}H_B,
 &E_{\mathcal C}&:=\sum_{B\in\mathcal C}E_B,
 \label{eq:n9v48-sum-H-E}\\
 e^{-t(H_{\mathcal C}-E_{\mathcal C})}
 &=\widehat\bigotimes_{B\in\mathcal C}
    e^{-t(H_B-E_B)},
 \label{eq:n9v48-tensor-semigroup}\\
 \|e^{-t(H_{\mathcal C}-E_{\mathcal C})}\|&\le1.
 \label{eq:n9v48-tensor-norm}
\end{align}
The contraction has no factor exponential in $|\mathcal C|$.
\end{theorem}

\begin{proof}
Operators on distinct tensor factors strongly commute.  The functional
calculus therefore factorizes their semigroup, giving
\eqref{eq:n9v48-tensor-semigroup}.  Tensor-product operator norms
multiply, and every factor has norm at most one by
\eqref{eq:n9v48-contraction}.
\end{proof}

\begin{firewall}[Gauss constraints can spoil literal tensorization]
\label{fire:n9v48-Gauss}
On the physical gauge-invariant Hilbert space, boundary Gauss constraints
can couple otherwise disjoint collars.  Literal tensorization is then
replaced by a normalized sector disintegration as in
\Cref{prop:n9v46-product-group}, or by a physical/auxiliary Schur
complement as in V44.  Extended-space factorization alone is not a proof
on the physical carrier.
\end{firewall}

\subsection{A finite coloring of ownership collars}

Form the collar-intersection graph with vertex set $\mathcal B_n$ and an
edge between $B$ and $B'$ when $B^+\cap(B')^+\ne\varnothing$.

\begin{lemma}[Finite collar coloring]
\label{lem:n9v48-coloring}
If the collar-intersection graph has maximum degree at most $\Delta_*$,
then it admits a coloring with at most $\Delta_*+1$ colors.  Collars of
blocks having the same color are pairwise disjoint.
\end{lemma}

\begin{proof}
Order the finite vertices arbitrarily and color them greedily.  When a
vertex is reached, at most $\Delta_*$ colors are forbidden by previously
colored neighbors, so one of $\Delta_*+1$ colors remains.  Equal colors
cannot be adjacent, which is exactly collar disjointness.
\end{proof}

\begin{corollary}[Colorwise internal resummation]
\label{cor:n9v48-color-resummation}
If the regulator geometry gives a volume-independent $\Delta_*$, then
the complete block family splits into a fixed number of disjoint-color
families.  Under the tensorization or physical-sector hypotheses of
\Cref{thm:n9v48-tensor-contraction,fire:n9v48-Gauss}, every color has an
internal contraction bound one.
\end{corollary}

\begin{proof}
Apply \Cref{lem:n9v48-coloring} and then
\Cref{thm:n9v48-tensor-contraction} to each color.
\end{proof}

\subsection{Boundary terms and distinct ownership}

Let $W_{BB'}$ be an interaction whose support meets two differently
colored collars after all internal terms have been assigned.  Let
$\mathcal E_n$ be the finite set of such owned boundary interactions.

\begin{definition}[Distinct boundary word]
\label{def:n9v48-distinct-boundary}
A boundary word is distinct when each $e\in\mathcal E_n$ occurs at most
once.  Repeated internal dynamics between its insertions is retained in
the exact ground-normalized block semigroups.
\end{definition}

\begin{proposition}[Bounded distinct-boundary estimate]
\label{prop:n9v48-boundary}
Let $S_j$ be contractions made from ground-normalized internal block
semigroups and let $W_{e_1},\dots,W_{e_k}$ be bounded, distinctly owned
boundary operators.  Then
\begin{equation}
 \|S_0W_{e_1}S_1\cdots W_{e_k}S_k\|
 \le\prod_{j=1}^k\|W_{e_j}\|.
 \label{eq:n9v48-boundary-word}
\end{equation}
No internal Yukawa multiplicity appears on the right.
\end{proposition}

\begin{proof}
Apply submultiplicativity and $\|S_j\|\le1$.
\end{proof}

\begin{firewall}[Distinctness needs an expansion theorem]
\label{fire:n9v48-distinctness}
Exact ownership does not itself ensure that a time-ordered expansion
uses each boundary interaction only once.  A forest interpolation,
polymer derivative, or another exact combinatorial representation must
prove distinctness.  Without it, repeated unbounded boundary operators
return the V40 multiplicity problem.
\end{firewall}

\subsection{Local ground-energy source ledger}

The block normalization introduces the scalar
$E_{\mathcal C}=\sum_BE_B$.  It cancels from a normalized fixed-metric
expectation only when numerator and denominator contain the same owned
blocks.  Its metric dependence remains relevant to Einstein variation.

\begin{proposition}[Variational Lipschitz control of a block ground
energy]
\label{prop:n9v48-ground-Lipschitz}
Let $H_B(g)$ and $H_B(g')$ be self-adjoint lower-semibounded block
operators on a common form domain.  Suppose
\begin{equation}
 |(H_B(g)-H_B(g'))[\Psi,\Psi]|
 \le L_Bd(g,g')\|\Psi\|^2
 \label{eq:n9v48-operator-Lipschitz}
\end{equation}
for every normalized form-domain vector.  Then
\begin{equation}
 |E_B(g)-E_B(g')|\le L_Bd(g,g').
 \label{eq:n9v48-ground-Lipschitz}
\end{equation}
\end{proposition}

\begin{proof}
The variational principle and
\eqref{eq:n9v48-operator-Lipschitz} give
$E_B(g)\le E_B(g')+L_Bd(g,g')$ by taking an approximating ground-state
sequence for $g'$.  Interchange $g$ and $g'$ for the reverse inequality.
\end{proof}

\begin{corollary}[Ground-source ownership]
\label{cor:n9v48-ground-source}
If the constants $L_B$ are local and uniformly bounded and the block
ownership has bounded incidence, then the ground-energy functional
$\sum_BE_B(g)$ has an extensive local Lipschitz bound.  It may be assigned
to a local cosmological/effective-action counterterm ledger, but it is
not a volume-uniform scalar resolvent shift.
\end{corollary}

\begin{proof}
Sum \eqref{eq:n9v48-ground-Lipschitz}; bounded incidence controls how many
block metrics depend on one local metric variation.  The sum still
scales with the number of varied blocks.
\end{proof}

\begin{firewall}[Lipschitz is not differentiability]
\label{fire:n9v48-ground-derivative}
Ground-state crossings can make $E_B(g)$ nondifferentiable.  The
Lipschitz estimate gives an almost-everywhere or subdifferential control,
not a classical Einstein stress tensor.  A differentiable source needs
an isolated spectral sector, a trace-state replacement, or a proved
weak variation theorem.
\end{firewall}

\subsection{Local Yukawa compiler}

\begin{definition}[Local Yukawa ownership card (LYOC)]
\label{def:n9v48-LYOC}
LYOC requires:
\begin{enumerate}
\item exact finite-range ownership with a volume-independent collar
degree;
\item a uniform local bound \eqref{eq:n9v48-local-Yukawa}, keeping
$C_{B,n}$ explicit;
\item closed local block Hamiltonians and exact ground-normalized
contractions;
\item physical tensorization by normalized gauge-sector disintegration
or a proved Schur margin;
\item an exact distinct-boundary expansion and summable boundary
activities; and
\item ownership and weak/differentiable control of the local ground
energy in the Einstein source ledger.
\end{enumerate}
\end{definition}

\begin{theorem}[LYOC closes internal Yukawa multiplicity]
\label{thm:n9v48-LYOC}
Under LYOC, arbitrarily many Higgs/Yukawa occurrences internal to an
owned block or disjoint color family are resummed into contractions of
norm at most one.  The remaining costs are exactly the distinct boundary
activities and the owned ground-energy source; no extensive sum of local
lower constants is used as one global resolvent shift.
\end{theorem}

\begin{proof}
Use \Cref{thm:n9v48-local-lower,cor:n9v48-contraction,
thm:n9v48-tensor-contraction,cor:n9v48-color-resummation} for internal
dynamics.  Use \Cref{prop:n9v48-boundary} and LYOC item 5 for boundary
terms.  Ground-energy ownership is item 6 and
\Cref{cor:n9v48-ground-source}.
\end{proof}

\begin{assessment}[Progress at ninth-series V48]
\label{ass:n9v48-status}
The extensive Yukawa lower shift has been replaced by a local
architecture.  Each owned block is closed and ground-normalized, so all
internal multiplicity costs one contraction.  Bounded collar coloring
keeps the number of tensorized families fixed.  The remaining physical
inputs are a distinct-boundary expansion, cutoff-uniform boundary
activity, gauge-sector factorization, and an Einstein-compatible
ground-energy source.
\end{assessment}

\begin{programme}[Ninth-series V49: local Higgs-amplitude tails]
\label{prog:n9v48-V49}
V49 will derive a uniform conditional tail from the block quartic under
an explicitly bounded exterior linear source.  It will test whether the
fermion-cutoff growth in $C_{B,n}$ can be paid by a large-amplitude bad
event without using rarity as an operator-norm substitute.
\end{programme}

\begin{firewall}[Claim boundary after ninth-series V48]
\label{fire:n9v48-boundary}
V48 proves the local ownership and contraction mechanism.  It does not
prove LYOC for the physical regulator, a distinct-boundary expansion,
uniform boundary activities, differentiability of the ground energy,
OS reconstruction, or the full SM--Einstein continuum.
\end{firewall}

\begin{theorem}[Ninth-series V48 result]
\label{thm:n9v48-result}
Uniform local quartic/Yukawa bounds make every owned block
lower-semibounded; subtracting its true ground energy resums arbitrary
internal multiplicity into a contraction.  A bounded-degree collar
coloring tensorizes disjoint blocks, while all cross-boundary and
ground-energy source terms remain explicitly owned interfaces.
\end{theorem}

\begin{proof}
Use \Cref{thm:n9v48-local-lower,cor:n9v48-contraction,
thm:n9v48-tensor-contraction,lem:n9v48-coloring,
prop:n9v48-boundary,thm:n9v48-LYOC}.
\end{proof}

\paragraph{Parent-version record.}
V48 was formed from
\texttt{Renewal\_SM\_Einstein\_Continuum\_Research\_Notes\_V47.tex}.
The V47 parent had $16\,964$ logical source lines, $614\,424$ bytes, and
SHA--256
\[
 \texttt{452B8C7D57CAB7E94EEAAFD46C9AE8F307937A64C892A1B3866840C831155691}.
\]

% =============================================================================
% END APPENDED NINTH-SERIES V48 MATERIAL
% =============================================================================
% =============================================================================
% BEGIN APPENDED NINTH-SERIES V49 MATERIAL
% =============================================================================

\clearpage
\section{Ninth-series V49: conditional Higgs quartic tails and the
good/bad boundary schedule}
\label{sec:v9v49}

\subsection{A normalized conditional block law}

Let $B$ be one fixed-size regulator block with Higgs coordinate
$\phi\in\mathbb R^D$ and exterior data $\xi$.  Assume a positive
conditional law
\begin{equation}
 d\mu_B(\phi\mid\xi)
 =Z_B(\xi)^{-1}e^{-U_B(\phi\mid\xi)/\hbar}\,d\phi,
 \qquad \hbar>0.
 \label{eq:n9v49-conditional-law}
\end{equation}
Suppose constants $c_4>0$, $J_*\ge0$, and $C_0<\infty$, independent of
the regulator and admissible $\xi$, satisfy
\begin{equation}
 U_B(\phi\mid\xi)
 \ge c_4|\phi|^4-J_*|\phi|-C_0.
 \label{eq:n9v49-conditional-coercivity}
\end{equation}
This is the conditional form of the V47 quartic bound after all owned
cross-boundary linear sources have been bounded by $J_*$.  Also assume a
local normalization floor: there exist $r_0>0$ and $M_0<\infty$ such
that
\begin{equation}
 U_B(\phi\mid\xi)\le M_0
 \quad\text{for }|\phi|\le r_0.
 \label{eq:n9v49-local-upper}
\end{equation}

\begin{lemma}[Uniform partition floor]
\label{lem:n9v49-Z-floor}
Under \eqref{eq:n9v49-local-upper},
\begin{equation}
 Z_B(\xi)
 \ge z_0:=|B_D(0,r_0)|e^{-M_0/\hbar}>0
 \label{eq:n9v49-Z-floor}
\end{equation}
uniformly in admissible exterior data.
\end{lemma}

\begin{proof}
Restrict the defining integral for $Z_B$ to the ball
$B_D(0,r_0)$ and use \eqref{eq:n9v49-local-upper}.
\end{proof}

\begin{lemma}[Absorption of the exterior linear source]
\label{lem:n9v49-source-absorption}
For every $r\ge0$,
\begin{equation}
 J_*r\le\frac{c_4}{2}r^4+C_J,
 \qquad
 C_J:=\frac{3}{4^{4/3}}J_*^{4/3}(c_4/2)^{-1/3}.
 \label{eq:n9v49-source-absorption}
\end{equation}
Consequently
\begin{equation}
 U_B(\phi\mid\xi)
 \ge a_4|\phi|^4-C_*,
 \qquad
 a_4:=c_4/2,\quad C_*:=C_0+C_J.
 \label{eq:n9v49-pure-quartic}
\end{equation}
\end{lemma}

\begin{proof}
Use \eqref{eq:n9v47-br} with $b=J_*$ and
$\delta=c_4/2$, then insert the result in
\eqref{eq:n9v49-conditional-coercivity}.
\end{proof}

\subsection{Uniform quartic tail}

Define
\begin{equation}
 I_{D,a}:=\int_{\mathbb R^D}e^{-a|\phi|^4}\,d\phi<\infty.
 \label{eq:n9v49-quartic-integral}
\end{equation}

\begin{theorem}[Conditional quartic tail]
\label{thm:n9v49-tail}
Under \eqref{eq:n9v49-conditional-law}--
\eqref{eq:n9v49-local-upper}, for every $R\ge0$,
\begin{equation}
 \mu_B\{\,|\phi|>R\mid\xi\}
 \le A_*e^{-c_*R^4},
 \label{eq:n9v49-tail}
\end{equation}
where
\begin{equation}
 c_*:=\frac{a_4}{2\hbar},
 \qquad
 A_*:=z_0^{-1}e^{C_*/\hbar}I_{D,a_4/(2\hbar)}.
 \label{eq:n9v49-tail-constants}
\end{equation}
All constants are uniform in the admissible exterior condition.
\end{theorem}

\begin{proof}
By \eqref{eq:n9v49-pure-quartic}, on $|\phi|>R$,
\begin{align*}
 e^{-U_B/\hbar}
 &\le e^{C_*/\hbar}e^{-a_4|\phi|^4/\hbar}\\
 &\le e^{C_*/\hbar}e^{-a_4R^4/(2\hbar)}
       e^{-a_4|\phi|^4/(2\hbar)}.
\end{align*}
Integrate and divide by the floor \eqref{eq:n9v49-Z-floor}.
Finiteness of \eqref{eq:n9v49-quartic-integral} follows from polar
coordinates.
\end{proof}

\begin{corollary}[Uniform polynomial moments]
\label{cor:n9v49-moments}
For every $m>0$ there is a finite constant $K_m$, depending only on the
data in \eqref{eq:n9v49-tail-constants}, such that
\begin{equation}
 \sup_\xi\mathbb E_{\mu_B(\cdot\mid\xi)}|\phi|^m\le K_m.
 \label{eq:n9v49-moments}
\end{equation}
One may take
\begin{equation}
 K_m=A_*c_*^{-m/4}\Gamma(1+m/4)
 \label{eq:n9v49-moment-constant}
\end{equation}
after replacing $A_*$ by $\max\{1,A_*\}$.
\end{corollary}

\begin{proof}
Use the layer-cake identity
$\mathbb E X^m=m\int_0^\infty R^{m-1}\mathbb P(X>R)dR$ and
\eqref{eq:n9v49-tail}.  The change of variables $u=c_*R^4$ gives
$m\int_0^\infty R^{m-1}e^{-c_*R^4}dR
=c_*^{-m/4}\Gamma(1+m/4)$.
\end{proof}

\begin{corollary}[Uniform subquartic exponential moments]
\label{cor:n9v49-exponential}
For every $0\le\theta<a_4/\hbar$,
\begin{equation}
 \sup_\xi\mathbb E_{\mu_B(\cdot\mid\xi)}
 e^{\theta|\phi|^4}<\infty.
 \label{eq:n9v49-exponential}
\end{equation}
\end{corollary}

\begin{proof}
Insert $e^{\theta|\phi|^4}$ in the conditional integral and use
\eqref{eq:n9v49-pure-quartic}.  The remaining integral has decay
$e^{-(a_4/\hbar-\theta)|\phi|^4}$ and the same partition floor.
\end{proof}

\begin{firewall}[The normalization floor is not optional]
\label{fire:n9v49-normalization}
The coercive lower bound controls the unnormalized tail integral.  A
conditional probability bound also needs the uniform floor
\eqref{eq:n9v49-Z-floor}.  If the exterior source drives the action high
on every fixed ball, $Z_B(\xi)$ can be arbitrarily small and the stated
uniform conditional tail does not follow.
\end{firewall}

\subsection{Good/bad boundary splitting}

Let
\begin{equation}
 G_R:=\{|\phi|\le R\},
 \qquad B_R:=\{|\phi|>R\}.
 \label{eq:n9v49-good-bad}
\end{equation}
Suppose a local boundary Yukawa insertion $W_n$ has on $G_R$ the
deterministic bound
\begin{equation}
 \|W_n\mathbf1_{G_R}\|\le q_n(R):=g_ny_nR,
 \label{eq:n9v49-good-activity}
\end{equation}
where $g_n$ is its owned coupling and $y_n$ is the regulated local
fermion-bilinear norm.  Suppose its bad-sector expectation is tested
through a local insertion $G_n$ with
\begin{equation}
 \left(
  \mathbb E\sum_\alpha\kappa_\alpha|G_{n,\alpha}|^2
 \right)^{1/2}\le B_n.
 \label{eq:n9v49-bad-L2}
\end{equation}

\begin{theorem}[Quartic-tail boundary compiler]
\label{thm:n9v49-boundary-compiler}
Under the uniform conditional tail,
\begin{equation}
 \left|
  \mathbb E\!\left[
   \mathbf1_{B_R}\sum_\alpha\kappa_\alpha G_{n,\alpha}
  \right]
 \right|
 \le B_nA_*^{1/2}e^{-c_*R^4/2}.
 \label{eq:n9v49-bad-error}
\end{equation}
Thus a threshold sequence $R_n$ closes both the good and bad boundary
rows whenever
\begin{align}
 g_ny_nR_n&\longrightarrow0,
 \label{eq:n9v49-good-condition}\\
 \log B_n+\tfrac12\log A_*-\tfrac12c_*R_n^4
 &\longrightarrow-\infty.
 \label{eq:n9v49-bad-condition}
\end{align}
\end{theorem}

\begin{proof}
Apply the normalized-label Cauchy--Schwarz compiler
\eqref{eq:n9v35-local-cs} with $E=B_R$, followed by
\eqref{eq:n9v49-tail}.  The first condition makes the good activity
\eqref{eq:n9v49-good-activity} vanish and the second makes
\eqref{eq:n9v49-bad-error} vanish.
\end{proof}

\begin{corollary}[Explicit logarithmic threshold]
\label{cor:n9v49-log-threshold}
Let $s_n\to\infty$ and choose
\begin{equation}
 R_n^4=\frac{2}{c_*}\bigl(\log_+B_n+s_n\bigr),
 \qquad \log_+x:=\max\{0,\log x\}.
 \label{eq:n9v49-R-choice}
\end{equation}
Then the bad error is at most a constant times $e^{-s_n}$.  The good
condition becomes
\begin{equation}
 g_ny_n\bigl(\log_+B_n+s_n\bigr)^{1/4}\longrightarrow0.
 \label{eq:n9v49-good-log}
\end{equation}
If $B_n\le C M_n^b$, one may choose $s_n=o(\log M_n)$ with
$s_n\to\infty$; it then suffices that
\begin{equation}
 g_ny_n(\log M_n)^{1/4}\longrightarrow0.
 \label{eq:n9v49-polynomial-schedule}
\end{equation}
\end{corollary}

\begin{proof}
Substitution of \eqref{eq:n9v49-R-choice} into
\eqref{eq:n9v49-bad-error} gives the exponential factor $e^{-s_n}$,
up to $A_*^{1/2}$ and the harmless case $B_n<1$.  Taking fourth roots
gives \eqref{eq:n9v49-good-log}.  Polynomial growth makes
$\log_+B_n=O(\log M_n)$.
\end{proof}

\begin{firewall}[Rarity pays only the bad expectation]
\label{fire:n9v49-rarity}
The tail factor in \eqref{eq:n9v49-bad-error} does not improve the
operator norm on $G_R$.  If $g_ny_n$ fails to decay fast enough for
\eqref{eq:n9v49-good-condition}, increasing $R_n$ makes the good bound
worse.  Quartic rarity therefore cannot hide an uncontrolled regulated
fermion bilinear; it works only when the two displayed scale conditions
have a common threshold.
\end{firewall}

\begin{remark}[Internal and boundary Yukawa rows differ]
\label{rem:n9v49-internal-boundary}
An internal owned Yukawa row is already resummed in the V48
ground-normalized block contraction and need not satisfy
\eqref{eq:n9v49-good-condition}.  The good/bad small-activity schedule is
required for boundary or polymer rows that remain outside that exact
block Hamiltonian.
\end{remark}

\subsection{Exterior-field and positivity obstructions}

For a gauge-covariant nearest-neighbor square, the conditional block
action contains boundary terms of the schematic form
\begin{equation}
 -2\operatorname{Re}\langle\phi_x,U_{xy}\xi_y\rangle.
 \label{eq:n9v49-boundary-source}
\end{equation}
Compactness of $U_{xy}$ preserves $|\xi_y|$ but does not bound it.

\begin{proposition}[Nested-source alternative]
\label{prop:n9v49-nested-source}
If $|\xi_y|\le R_{\rm ext}$ on an exterior-good event and each block has
a fixed number of boundary edges, then
\eqref{eq:n9v49-conditional-coercivity} holds there with
$J_*=C_BR_{\rm ext}$ for a geometric constant $C_B$.  On the exterior-bad
event, a second quartic tail or a joint enlarged-block estimate is
required.  No exterior-uniform $J_*$ follows from compact gauge links
alone.
\end{proposition}

\begin{proof}
The Cauchy--Schwarz inequality and unitarity of $U_{xy}$ bound
\eqref{eq:n9v49-boundary-source} by
$2|\phi_x||\xi_y|$.  Sum over the fixed boundary incidence.  The last
statement follows because $|U_{xy}\xi_y|=|\xi_y|$ can be arbitrarily
large.
\end{proof}

\begin{firewall}[Fermion positivity]
\label{fire:n9v49-fermion-positivity}
The probability law \eqref{eq:n9v49-conditional-law} presupposes a
positive conditional density.  After integrating chiral fermions, the
determinant or Pfaffian can have a phase.  A modulus tail bound does not
control phase cancellations or define a positive conditional
probability.  The physical Standard Model programme needs a separate
determinant-phase, doubled-representation, or constructive positivity
input before applying this conditional law to the full fermionic
measure.
\end{firewall}

\subsection{Local amplitude-tail card}

\begin{definition}[Local Higgs amplitude-tail card (LHATC)]
\label{def:n9v49-LHATC}
LHATC requires:
\begin{enumerate}
\item the positive normalized conditional block law and uniform
partition floor;
\item the quartic coercive bound with all exterior sources controlled
directly or by a nested/enlarged-block tail;
\item uniform constants $A_*,c_*$ in \eqref{eq:n9v49-tail};
\item a boundary good/bad decomposition with
\eqref{eq:n9v49-good-condition}--\eqref{eq:n9v49-bad-condition};
\item exact separation from the internally resummed V48 rows; and
\item a separate treatment of any fermion phase and of metric derivatives
of the conditional normalization.
\end{enumerate}
\end{definition}

\begin{theorem}[LHATC closes the Higgs-amplitude boundary row]
\label{thm:n9v49-LHATC}
Under LHATC, large Higgs amplitudes contribute a vanishing local
expectation error and the regular-amplitude boundary activity also tends
to zero.  The conclusion is support-local and compatible with the V48
block contraction; it makes no operator-norm claim on the rare sector.
\end{theorem}

\begin{proof}
Use \Cref{thm:n9v49-tail,thm:n9v49-boundary-compiler}.  LYOC/LHATC
ownership distinguishes internal from boundary rows, and
\Cref{fire:n9v49-rarity} fixes the scope of the result.
\end{proof}

\begin{assessment}[Progress at ninth-series V49]
\label{ass:n9v49-status}
The local Higgs quartic now yields a normalized conditional
$e^{-cR^4}$ tail with an explicit partition floor and exterior-source
hypothesis.  A logarithmic threshold balances polynomial bad-insertion
growth against a good boundary activity.  The analysis also exposes two
physical obstructions: unbounded exterior Higgs data require a nested
tail, and a chiral fermion phase prevents direct use of a positive
conditional probability.
\end{assessment}

\begin{programme}[Ninth-series V50: compulsory companion audit]
\label{prog:n9v49-V50}
V50 will inspect the newest Yang--Mills and Navier--Stokes notes, test for
a type-compatible nested boundary tail, phase, or state-level
factorization input, and then continue to V51.
\end{programme}

\begin{firewall}[Claim boundary after ninth-series V49]
\label{fire:n9v49-boundary}
V49 proves the positive conditional quartic-tail compiler.  It does not
prove the exterior-uniform hypotheses or fermion positivity for the full
Standard Model, close the boundary activity schedule, renormalize the
metric-dependent normalization, establish OS reconstruction, or
construct the full continuum.
\end{firewall}

\begin{theorem}[Ninth-series V49 result]
\label{thm:n9v49-result}
Uniform conditional quartic coercivity plus a local partition floor gives
the tail \eqref{eq:n9v49-tail}.  A boundary row closes when one threshold
$R_n$ satisfies the simultaneous good and bad conditions
\eqref{eq:n9v49-good-condition}--\eqref{eq:n9v49-bad-condition}; for
polynomial bad-insertion growth this reduces to the logarithmic schedule
\eqref{eq:n9v49-polynomial-schedule}.
\end{theorem}

\begin{proof}
Use \Cref{lem:n9v49-Z-floor,lem:n9v49-source-absorption,
thm:n9v49-tail,thm:n9v49-boundary-compiler,
cor:n9v49-log-threshold,thm:n9v49-LHATC}.
\end{proof}

\paragraph{Parent-version record.}
V49 was formed from
\texttt{Renewal\_SM\_Einstein\_Continuum\_Research\_Notes\_V48.tex}.
The V48 parent had $17\,362$ logical source lines, $629\,393$ bytes, and
SHA--256
\[
 \texttt{7422214EFEABEA2623C1600040F6177303542F430219CFB1B0D32C50E097CB40}.
\]

% =============================================================================
% END APPENDED NINTH-SERIES V49 MATERIAL
% =============================================================================
% =============================================================================
% BEGIN APPENDED NINTH-SERIES V50 MATERIAL
% =============================================================================

\clearpage
\section{Ninth-series V50: compulsory audit and the worst-exterior
Higgs influence obstruction}
\label{sec:v9v50}

\subsection{Compulsory companion audit}

The newest Yang--Mills note inspected was
\[
 \texttt{Renewal\_Yang\_Mills\_Mass\_Gap\_Research\_Notes\_V84.tex},
\]
with $37\,464$ logical source lines, $1\,337\,700$ bytes, and SHA--256
\[
 \texttt{6FB9582D4AAB62A939596684A9680A9C541A441E2CC2DD6F8BB5C0B6C8990C1F}.
\]
The newest Navier--Stokes note remained
\[
 \texttt{Renewal\_NS\_Stability\_Research\_Notes\_V26.tex},
\]
with $5\,319$ lines, $278\,283$ bytes, and SHA--256
\[
 \texttt{82C887F8C2C69E4F28FAD7C3DC8DE5B9099CB5A4BF431EBA1FFF3E91935978AD}.
\]

The substantive Yang--Mills V81--V84 increment contributes the following
structural information.
\begin{enumerate}
\item V81 proves that an infrared-relative perturbation of a massless
quadratic mode does not create a gap; a genuinely nonperturbative
terminal input is required.
\item V82 proves a terminal strong-coupling Dobrushin/OS gap for a compact
Wilson measure, while retaining the ultraviolet-to-terminal RG matching
as a separate problem.
\item V83 replaces a bare-action target by an open weighted Dobrushin ball
for complete quasi-local effective interactions.
\item V84 derives the corrected marginal influence recursion, including
the eliminated-sector feedback denominator, and proves that rare bad
blocks do not control worst-exterior total-variation influence.
\end{enumerate}
The Navier--Stokes file provides no new conditional-measure input for the
Higgs boundary problem.

The type-compatible conclusion is decisive for V49.  Its quartic tail
may participate in an averaged disagreement theorem, but it cannot be
used to assert a worst-exterior Dobrushin coefficient.  The obstruction
is even stronger for an unbounded Higgs field: a nonzero linear boundary
coupling has total-variation influence one in the worst exterior limit.

\subsection{An oscillation bound for normalized conditionals}

Let $P$ and $Q$ be probability measures with positive densities on a
common space.  Put $L=dP/dQ$.

\begin{lemma}[Likelihood-ratio oscillation bound]
\label{lem:n9v50-likelihood}
If
\begin{equation}
 0<m\le L\le M<\infty,
 \qquad \frac{M}{m}\le R,
 \label{eq:n9v50-likelihood-range}
\end{equation}
then
\begin{equation}
 \|P-Q\|_{\rm TV}
 \le\frac{\sqrt R-1}{\sqrt R+1}
 =\tanh\!\left(\frac14\log R\right).
 \label{eq:n9v50-TV-ratio}
\end{equation}
\end{lemma}

\begin{proof}
Because $\mathbb E_Q L=1$, necessarily $m\le1\le M$.  Among random
variables in $[m,M]$ with mean one, the convex functional
$\mathbb E|L-1|$ is maximized by a variable supported at the endpoints.
Its total variation is then
\[
 \frac12\mathbb E|L-1|
 =\frac{(M-1)(1-m)}{M-m}.
\]
For fixed ratio $M/m\le R$, this expression is maximized at
$m=R^{-1/2}$, $M=R^{1/2}$ and equals the right side of
\eqref{eq:n9v50-TV-ratio}.
\end{proof}

\begin{corollary}[Potential-oscillation bound]
\label{cor:n9v50-potential}
Let
\begin{equation}
 dP=Z_U^{-1}e^{-U/\hbar}d\nu,
 \qquad
 dQ=Z_V^{-1}e^{-V/\hbar}d\nu.
 \label{eq:n9v50-two-Gibbs}
\end{equation}
If
\begin{equation}
 \operatorname{osc}(U-V)
 :=\operatorname*{ess\,sup}(U-V)
  -\operatorname*{ess\,inf}(U-V)
 \le\Delta,
 \label{eq:n9v50-osc}
\end{equation}
then
\begin{equation}
 \|P-Q\|_{\rm TV}\le\tanh\!\left(\frac{\Delta}{4\hbar}\right).
 \label{eq:n9v50-Gibbs-TV}
\end{equation}
\end{corollary}

\begin{proof}
The logarithm of $dP/dQ$ is
$-(U-V)/\hbar+\log Z_V-\log Z_U$.  Its oscillation is at most
$\Delta/\hbar$, so the likelihood-ratio range is at most
$e^{\Delta/\hbar}$.  Apply \Cref{lem:n9v50-likelihood}.
\end{proof}

\subsection{Truncation and its exact probability debit}

\begin{lemma}[Good-set comparison]
\label{lem:n9v50-good-set}
Let $G$ satisfy $P(G),Q(G)>0$, and let $P_G,Q_G$ be the corresponding
conditional laws.  Then
\begin{equation}
 \|P-Q\|_{\rm TV}
 \le P(G^c)+Q(G^c)+\|P_G-Q_G\|_{\rm TV}.
 \label{eq:n9v50-good-set}
\end{equation}
\end{lemma}

\begin{proof}
The triangle inequality gives
\[
 \|P-Q\|_{\rm TV}
 \le\|P-P_G\|_{\rm TV}
   +\|P_G-Q_G\|_{\rm TV}
   +\|Q_G-Q\|_{\rm TV}.
\]
For conditioning on an event,
$\|P-P_G\|_{\rm TV}=P(G^c)$, and similarly for $Q$.
\end{proof}

\begin{theorem}[Truncated conditional influence]
\label{thm:n9v50-truncated}
Consider two block conditionals of the form
\eqref{eq:n9v50-two-Gibbs} and a common good set
$G_R=\{|\phi|\le R\}$.  Suppose
\begin{align}
 P(G_R^c)+Q(G_R^c)&\le2\pi_R,
 \label{eq:n9v50-two-tails}\\
 \operatorname{osc}_{G_R}(U-V)&\le\Delta_R.
 \label{eq:n9v50-good-osc}
\end{align}
Then
\begin{equation}
 \|P-Q\|_{\rm TV}
 \le2\pi_R+\tanh\!\left(\frac{\Delta_R}{4\hbar}\right).
 \label{eq:n9v50-truncated-TV}
\end{equation}
\end{theorem}

\begin{proof}
Apply \Cref{lem:n9v50-good-set} and then
\Cref{cor:n9v50-potential} to the normalized restrictions on $G_R$.
\end{proof}

\begin{corollary}[Bounded Higgs boundary source]
\label{cor:n9v50-bounded-source}
For an owned interaction
\begin{equation}
 W(\phi,\xi)=-2\kappa\operatorname{Re}
   \langle\phi,U\xi\rangle
 \label{eq:n9v50-linear-boundary}
\end{equation}
with $U$ unitary, changing $\xi$ to $\xi'$ gives on $|\phi|\le R$
\begin{equation}
 \operatorname{osc}_\phi
 [W(\phi,\xi)-W(\phi,\xi')]
 \le4|\kappa|R|\xi-\xi'|.
 \label{eq:n9v50-linear-osc}
\end{equation}
If $|\xi|,|\xi'|\le R_{\rm ext}$ and both conditional tails obey V49,
then
\begin{equation}
 c_{B\leftarrow\partial B}
 \le2A_*e^{-c_*R^4}
 +\tanh\!\left(
   \frac{2|\kappa|RR_{\rm ext}}{\hbar}
 \right).
 \label{eq:n9v50-bounded-influence}
\end{equation}
\end{corollary}

\begin{proof}
The difference in \eqref{eq:n9v50-linear-boundary} is a real linear
functional of $\phi$ of norm at most
$2|\kappa||\xi-\xi'|$.  Its range over the radius-$R$ ball has length at
most twice the radius times that norm, proving
\eqref{eq:n9v50-linear-osc}.  Use
$|\xi-\xi'|\le2R_{\rm ext}$ in
\Cref{thm:n9v50-truncated} and insert \eqref{eq:n9v49-tail}.
\end{proof}

\begin{remark}[The two cutoffs compete]
\label{rem:n9v50-compete}
Increasing $R$ suppresses the tail term but increases the oscillation
term.  Increasing $R_{\rm ext}$ worsens the worst-good-exterior bound.
A useful contraction therefore needs a coupling/scale schedule and a
separate probability debit for exterior-bad configurations.
\end{remark}

\subsection{Worst-exterior influence is one}

The following elementary model rules out a uniform total-variation
Dobrushin coefficient for an unbounded linearly coupled Higgs spin.

\begin{theorem}[Quartic linear-source no-go]
\label{thm:n9v50-quartic-no-go}
For $\lambda,\hbar>0$, let $P_J$ be the probability law on $\mathbb R$
with density
\begin{equation}
 p_J(x)=Z_J^{-1}
 \exp\!\left\{-\frac{\lambda x^4-Jx}{\hbar}\right\}.
 \label{eq:n9v50-pJ}
\end{equation}
Then
\begin{equation}
 \lim_{J\to+\infty}P_J\{x>0\}=1,
 \qquad
 \lim_{J\to+\infty}P_{-J}\{x>0\}=0,
 \label{eq:n9v50-opposite-limits}
\end{equation}
and consequently
\begin{equation}
 \sup_{J,J'\in\mathbb R}\|P_J-P_{J'}\|_{\rm TV}=1.
 \label{eq:n9v50-TV-one}
\end{equation}
\end{theorem}

\begin{proof}
For $J>0$, the numerator of $P_J\{x\le0\}$ is bounded by
$\int_{-\infty}^0e^{-\lambda x^4/\hbar}dx=:C_0$.  Put
$x_J=(J/(4\lambda))^{1/3}$.  The exponent
$F_J(x)=(Jx-\lambda x^4)/\hbar$ is maximized at $x_J$, with
\[
 F_J(x_J)=\frac{3Jx_J}{4\hbar}.
\]
On the interval
$[x_J,x_J+x_J^{-1}]$, Taylor's theorem and
$|F_J''(x)|=12\lambda x^2/\hbar$ show
$F_J(x)\ge F_J(x_J)-C$ for all sufficiently large $J$, with $C$
independent of $J$.  Hence
\[
 Z_J\ge x_J^{-1}e^{F_J(x_J)-C}\longrightarrow\infty
\]
faster than any power, and $P_J\{x\le0\}\le C_0/Z_J\to0$.
The law $P_{-J}$ is the reflection of $P_J$, proving the second limit.
Total variation dominates the difference of the probabilities of
$\{x>0\}$, so its supremum is one.
\end{proof}

\begin{corollary}[Unbounded exterior Higgs Dobrushin coefficient]
\label{cor:n9v50-Higgs-no-go}
If a scalar component of a block Higgs field has a conditional quartic
law whose exterior neighbor enters through a nonzero linear source
$J=\kappa\xi$, and $\xi$ ranges over all of $\mathbb R$, then the
worst-exterior total-variation influence coefficient equals one for
every $\kappa\ne0$.
\end{corollary}

\begin{proof}
Choose exterior values $\xi=J/\kappa$ and $\xi'=-J/\kappa$ and let
$J\to\infty$.  Apply \Cref{thm:n9v50-quartic-no-go}.
\end{proof}

\begin{firewall}[Quartic tails do not contradict the no-go]
\label{fire:n9v50-tail-no-go}
For each fixed bounded exterior source, V49 gives a strong quartic tail.
The Dobrushin coefficient takes a supremum over all exterior values
before averaging their probability.  The sources used in
\Cref{thm:n9v50-quartic-no-go} move the conditional mass to arbitrarily
large amplitudes, so the fixed-source tail constants are not uniform in
that supremum.
\end{firewall}

\subsection{The feedback denominator}

For a retained/exterior decomposition, let nonnegative influence
majorants be arranged as
\begin{equation}
 C=
 \begin{pmatrix}
 C_{RR}&C_{RE}\\
 C_{ER}&C_{EE}
 \end{pmatrix}.
 \label{eq:n9v50-influence-blocks}
\end{equation}

\begin{proposition}[Conditional-influence return paths]
\label{prop:n9v50-return-path}
If a chosen influence norm is submultiplicative and
$\|C_{EE}\|<1$, the eliminated-sector return majorant is
\begin{equation}
 A=C_{RE}(I-C_{EE})^{-1}C_{ER}.
 \label{eq:n9v50-return}
\end{equation}
A retained target with self-return $A_{ii}$ has the safe row bound
\begin{equation}
 c^{\rm eff}_{ij}
 \le\frac{c_{ij}+A_{ij}}{1-A_{ii}}
 \label{eq:n9v50-feedback}
\end{equation}
whenever $A_{ii}<1$.  Omitting the denominator assumes away the
retained--eliminated--retained feedback loop.
\end{proposition}

\begin{proof}
The Neumann series for $(I-C_{EE})^{-1}$ sums all paths confined to the
eliminated sector.  A disagreement at retained target $i$ can return
through $A_{ii}$ repeatedly, giving
\[
 p_i\le c_{ij}+A_{ij}+A_{ii}p_i.
\]
Solving this scalar inequality gives \eqref{eq:n9v50-feedback}.
\end{proof}

\begin{remark}[Audit provenance]
\label{rem:n9v50-provenance}
The form of \eqref{eq:n9v50-feedback} is the type-compatible conclusion
of the audited Yang--Mills V84 marginal-coupling theorem.  The present
note does not import its physical Yang--Mills constants; it records the
abstract feedback algebra needed by any future Higgs averaged-influence
recursion.
\end{remark}

\subsection{Audit-adjusted terminal interface}

\begin{definition}[Higgs influence alternatives card (HIAC)]
\label{def:n9v50-HIAC}
HIAC requires one of:
\begin{enumerate}
\item a hard or dynamically invariant bounded-amplitude carrier on which
the uniform truncated coefficient \eqref{eq:n9v50-bounded-influence}
has weighted row sum below one; or
\item an averaged disagreement-percolation theorem that charges
exterior-bad amplitudes by their probability and uses the good-set
coefficient without taking a supremum over those bad values.
\end{enumerate}
In either case, marginalization must retain the return paths and
feedback denominator of \eqref{eq:n9v50-feedback}.
\end{definition}

\begin{theorem}[V50 audit verdict]
\label{thm:n9v50-audit}
For the unbounded physical Higgs carrier with a nonzero linear
nearest-neighbor coupling, HIAC item 1 is unavailable without an
additional invariant hard cutoff.  The type-compatible continuation is
item 2: an averaged disagreement process combining the V49 quartic tail
with the good-set oscillation bound and the corrected feedback recursion.
\end{theorem}

\begin{proof}
\Cref{cor:n9v50-Higgs-no-go} rules out a strict worst-exterior
total-variation coefficient on the unbounded carrier.  V49 supplies the
probability tail and \Cref{thm:n9v50-truncated} supplies the good-set
coefficient.  \Cref{prop:n9v50-return-path} gives the required
marginalization algebra.
\end{proof}

\begin{assessment}[Progress at ninth-series V50]
\label{ass:n9v50-status}
The compulsory audit has prevented a false terminal claim.  A quartic
tail does not load a worst-exterior Dobrushin norm, and the scalar
quartic conditional proves that this norm is exactly one under an
unbounded linear source.  The valid data now available are a finite
good-set oscillation coefficient, an explicit bad probability, and a
feedback-corrected path recursion.
\end{assessment}

\begin{programme}[Ninth-series V51: averaged disagreement percolation]
\label{prog:n9v50-V51}
V51 will construct a two-type disagreement process with regular
transmission and rare amplitude-defect sites.  It will prove exponential
decay when the sum of the good transmission kernel and the defect
branching kernel has weighted norm below one, retaining the defect
probability rather than replacing it by a worst-case conditional
coefficient.
\end{programme}

\begin{firewall}[Claim boundary after ninth-series V50]
\label{fire:n9v50-boundary}
V50 proves the worst-exterior no-go and selects an averaged route.  It
does not prove an averaged disagreement coupling for the physical
SM--Einstein measure, a terminal clustering theorem, reflection
positivity, OS reconstruction, or the full continuum.
\end{firewall}

\begin{theorem}[Ninth-series V50 result]
\label{thm:n9v50-result}
Potential oscillation gives the truncated influence
\eqref{eq:n9v50-truncated-TV}, but an unbounded quartic spin with any
nonzero linear exterior coupling has worst-exterior total-variation
coefficient one.  Hence quartic rarity can enter only an averaged
disagreement criterion, whose marginal recursion must include the
feedback denominator \eqref{eq:n9v50-feedback}.
\end{theorem}

\begin{proof}
Use \Cref{lem:n9v50-likelihood,thm:n9v50-truncated,
thm:n9v50-quartic-no-go,cor:n9v50-Higgs-no-go,
prop:n9v50-return-path,thm:n9v50-audit}.
\end{proof}

\paragraph{Parent-version record.}
V50 was formed from
\texttt{Renewal\_SM\_Einstein\_Continuum\_Research\_Notes\_V49.tex}.
The V49 parent had $17\,762$ logical source lines, $643\,054$ bytes, and
SHA--256
\[
 \texttt{13353EDC7A50BA52C1A26E0013354522E31F4D43E022825E1D6A3260719AEC87}.
\]

% =============================================================================
% END APPENDED NINTH-SERIES V50 MATERIAL
% =============================================================================
% =============================================================================
% BEGIN APPENDED NINTH-SERIES V51 MATERIAL
% =============================================================================

\clearpage
\section{Ninth-series V51: averaged regular/defect disagreement
percolation}
\label{sec:v9v51}

\subsection{Why connected propagation is the correct object}

V50 rules out a worst-exterior total-variation contraction for an
unbounded Higgs field.  Replacing that coefficient by an additive defect
probability would also be insufficient: an error floor at every site
does not decay with distance.  A rare defect can contribute to distant
boundary sensitivity only when it belongs to a connected propagation
from the source of the boundary disagreement.  We therefore formulate a
support-resolved path theorem.

Let $\mathcal G=(V,E)$ be a locally finite graph with graph distance
$d$ and maximum degree at most $\Delta\ge2$.  Let $A\subset V$ be the
finite source set where two boundary conditions differ.  A disagreement
coupling supplies events $D_i$ that site $i$ disagrees and a connected
path representation:
\begin{equation}
 D_i\subset
 \bigcup_{\gamma:i\leadsto A}\ \bigcup_{F\subset\gamma^\circ}
 \mathcal P(\gamma,F).
 \label{eq:n9v51-path-representation}
\end{equation}
Here $\gamma$ is a self-avoiding path from $i$ to $A$,
$\gamma^\circ$ is the set of propagation vertices after the source, and
$F$ marks the vertices at which propagation uses an amplitude-defect
mechanism; the other vertices use regular good-field transmission.

\begin{definition}[Mixed averaged path bound]
\label{def:n9v51-mixed-path}
The coupling has parameters $q,b\ge0$ when every length-$n$ path and
marking $F\subset\gamma^\circ$ satisfies
\begin{equation}
 \mathbb P\bigl(\mathcal P(\gamma,F)\bigr)
 \le q^{\,n-|F|}b^{\,|F|}.
 \label{eq:n9v51-mixed-path}
\end{equation}
This is an annealed support estimate.  It is not a supremum of
conditionals on defect configurations.
\end{definition}

\subsection{Exponential path summation}

\begin{theorem}[Two-type disagreement-path compiler]
\label{thm:n9v51-path-compiler}
Assume \eqref{eq:n9v51-path-representation} and
\eqref{eq:n9v51-mixed-path}.  If
\begin{equation}
 \zeta:=\Delta e^\mu(q+b)<1
 \label{eq:n9v51-zeta}
\end{equation}
for some $\mu>0$, then
\begin{equation}
 \mathbb P(D_i)
 \le
 \frac{e^{-\mu d(i,A)}}{1-\zeta}.
 \label{eq:n9v51-disagreement-decay}
\end{equation}
More precisely, the numerator may be multiplied by
$[\Delta(q+b)e^\mu]^{d(i,A)}$, but the simpler displayed bound is
uniform.
\end{theorem}

\begin{proof}
For a fixed length-$n$ path, summing
\eqref{eq:n9v51-mixed-path} over all markings gives
\[
 \sum_{F\subset\gamma^\circ}
 q^{n-|F|}b^{|F|}=(q+b)^n.
\]
The number of self-avoiding paths of length $n$ beginning at $i$ is at
most $\Delta^n$.  Every path reaching $A$ has
$n\ge d(i,A)$.  The union bound therefore gives
\[
 \mathbb P(D_i)
 \le\sum_{n\ge d(i,A)}[\Delta(q+b)]^n
 =\sum_{n\ge d(i,A)}e^{-\mu n}\zeta^n
 \le\frac{e^{-\mu d(i,A)}}{1-\zeta}.
\]
\end{proof}

\begin{corollary}[Anisotropic kernel form]
\label{cor:n9v51-kernel}
Let $K=(K_{ij})$ be a nonnegative propagation kernel and assume the
coupling yields
\begin{equation}
 p_i\le s_i+\sum_jK_{ij}p_j,
 \qquad
 \|K\|_\mu:=
 \sup_i\sum_je^{\mu d(i,j)}K_{ij}<1,
 \label{eq:n9v51-kernel-recursion}
\end{equation}
where $s$ is supported on $A$.  Then
\begin{equation}
 \sup_i e^{\mu d(i,A)}p_i
 \le
 \frac{\sup_i e^{\mu d(i,A)}s_i}{1-\|K\|_\mu}.
 \label{eq:n9v51-kernel-decay}
\end{equation}
The scalar path theorem is the special case in which the sum of regular
and defect step kernels has weighted row norm at most $\zeta$.
\end{corollary}

\begin{proof}
For the weighted supremum norm
$\|p\|_{\infty,\mu,A}=\sup_ie^{\mu d(i,A)}p_i$, the triangle inequality
for $d$ gives
\[
 e^{\mu d(i,A)}(Kp)_i
 \le\sum_je^{\mu d(i,j)}K_{ij}
       e^{\mu d(j,A)}p_j.
\]
Thus $\|Kp\|_{\infty,\mu,A}\le\|K\|_\mu
\|p\|_{\infty,\mu,A}$.  Iterate the nonnegative recursion and sum the
geometric series.
\end{proof}

\subsection{Loading the defect weight from a quartic tail}

For a threshold $R$, let the block defect be
\begin{equation}
 \mathcal D_B(R):=\{|\phi_B|>R\}.
 \label{eq:n9v51-defect}
\end{equation}

\begin{lemma}[Uniform conditional tails imply product domination]
\label{lem:n9v51-product}
Suppose that for every block $B$, every admissible value of all other
blocks, and one common $\pi_R$,
\begin{equation}
 \mathbb P(\mathcal D_B(R)\mid\text{all other blocks})\le\pi_R.
 \label{eq:n9v51-single-conditional}
\end{equation}
Then for every finite set $\mathcal F$ of blocks,
\begin{equation}
 \mathbb P\!\left(\bigcap_{B\in\mathcal F}\mathcal D_B(R)\right)
 \le\pi_R^{|\mathcal F|}.
 \label{eq:n9v51-product-domination}
\end{equation}
\end{lemma}

\begin{proof}
Enumerate $\mathcal F$.  Integrate the last block using its conditional
law given every other block and apply
\eqref{eq:n9v51-single-conditional}.  Repeat for the remaining blocks.
Each step contributes one factor $\pi_R$.
\end{proof}

\begin{firewall}[Good-exterior tails do not imply product domination]
\label{fire:n9v51-product}
If \eqref{eq:n9v51-single-conditional} holds only when the exterior is
itself good, the induction can enter an exterior-bad configuration and
fails.  The nested-source issue of V49 must be resolved by enlarged
blocks, a multilevel marking, or a direct joint estimate before
\eqref{eq:n9v51-product-domination} is used.
\end{firewall}

\begin{lemma}[Local $L^2$ conversion of a defect marking]
\label{lem:n9v51-L2-defect}
Let $\mathcal P(\gamma,F)$ carry a regular local amplitude
$G_{\gamma,F}$ supported on the path, and suppose
\begin{align}
 \mathbb E|G_{\gamma,F}|^2
 &\le q^{\,2(n-|F|)},
 \label{eq:n9v51-path-L2}\\
 \mathbb P\!\left(
  \bigcap_{v\in F}\mathcal D_v(R)
 \right)&\le\pi_R^{|F|}.
 \label{eq:n9v51-path-defects}
\end{align}
Then
\begin{equation}
 \mathbb E\!\left[
  |G_{\gamma,F}|
  \mathbf1_{\cap_{v\in F}\mathcal D_v(R)}
 \right]
 \le q^{\,n-|F|}b_R^{\,|F|},
 \qquad b_R:=\pi_R^{1/2}.
 \label{eq:n9v51-b-sqrt}
\end{equation}
\end{lemma}

\begin{proof}
Cauchy--Schwarz gives the product of the square roots of
\eqref{eq:n9v51-path-L2} and
\eqref{eq:n9v51-path-defects}.
\end{proof}

\begin{corollary}[Quartic defect activity]
\label{cor:n9v51-quartic-b}
If the V49 tail is uniform in the strong conditional sense
\eqref{eq:n9v51-single-conditional}, then
\begin{equation}
 \pi_R=A_*e^{-c_*R^4},
 \qquad
 b_R=A_*^{1/2}e^{-c_*R^4/2}
 \label{eq:n9v51-quartic-b}
\end{equation}
are valid product and local-$L^2$ defect weights.
\end{corollary}

\begin{proof}
Use \Cref{thm:n9v49-tail,lem:n9v51-product,
lem:n9v51-L2-defect}.
\end{proof}

\subsection{Regular transmission from the good-set oscillation}

For adjacent blocks with both Higgs amplitudes bounded by $R$, V50 gives
a total-variation transmission coefficient.  With interaction strength
$\kappa_n$ and fixed boundary incidence constant $c_B$, set
\begin{equation}
 q_n(R):=
 \tanh\!\left(\frac{c_B|\kappa_n|R^2}{\hbar}\right).
 \label{eq:n9v51-qR}
\end{equation}
This is a worst-case bound only within the double-good sector.

\begin{theorem}[Explicit quartic averaged criterion]
\label{thm:n9v51-explicit}
Assume:
\begin{enumerate}
\item an exact disagreement-path representation
\eqref{eq:n9v51-path-representation};
\item the regular good-field step is bounded by $q_n(R)$ in
\eqref{eq:n9v51-qR};
\item defect markings satisfy the mixed bound with
$b_R$ in \eqref{eq:n9v51-quartic-b}; and
\item the block graph has degree at most $\Delta$.
\end{enumerate}
Then boundary sensitivity decays as in
\eqref{eq:n9v51-disagreement-decay} whenever
\begin{equation}
 \Delta e^\mu\left[
 \tanh\!\left(\frac{c_B|\kappa_n|R^2}{\hbar}\right)
 +A_*^{1/2}e^{-c_*R^4/2}
 \right]<1.
 \label{eq:n9v51-explicit-criterion}
\end{equation}
\end{theorem}

\begin{proof}
The two step types give \eqref{eq:n9v51-mixed-path} with
$q=q_n(R)$ and $b=b_R$.  Apply
\Cref{thm:n9v51-path-compiler}.
\end{proof}

\begin{corollary}[A small-coupling threshold choice]
\label{cor:n9v51-small-coupling}
Suppose $A_*,c_*,c_B,\hbar,\Delta$ are uniform and
$|\kappa_n|\to0$.  Choose
\begin{equation}
 R_n^4=\frac{2}{c_*}\log|\kappa_n|^{-1}
 \label{eq:n9v51-Rkappa}
\end{equation}
for sufficiently large $n$.  Then
\begin{align}
 b_{R_n}&=A_*^{1/2}|\kappa_n|,
 \label{eq:n9v51-b-rate}\\
 q_n(R_n)&=
 O\!\left(
 |\kappa_n|\sqrt{\log|\kappa_n|^{-1}}
 \right),
 \label{eq:n9v51-q-rate}
\end{align}
so \eqref{eq:n9v51-explicit-criterion} holds for any fixed sufficiently
small $\mu>0$ at all large $n$.
\end{corollary}

\begin{proof}
The first identity follows by substitution.  Since
$R_n^2=[2c_*^{-1}\log|\kappa_n|^{-1}]^{1/2}$ and
$\tanh x\le x$ for $x\ge0$, the second bound follows.  Both terms tend to
zero, so their product with fixed $\Delta e^\mu$ is eventually below
one.
\end{proof}

\begin{firewall}[A fixed physical coupling is not covered by the
small-coupling corollary]
\label{fire:n9v51-fixed-coupling}
The Higgs nearest-neighbor coupling need not tend to zero in a physical
continuum scaling.  For fixed $\kappa$, the criterion
\eqref{eq:n9v51-explicit-criterion} may or may not have a usable
threshold $R$; this requires the actual dimensionless block coupling and
quartic constants.  The corollary is a conditional scaling result, not a
claim about Standard Model parameters.
\end{firewall}

\subsection{From disagreement decay to local covariance decay}

For a bounded local observable $F$, define its site oscillation
\begin{equation}
 \delta_i(F):=
 \sup\{|F(x)-F(x')|:x_j=x'_j\text{ for }j\ne i\}.
 \label{eq:n9v51-observable-oscillation}
\end{equation}

\begin{lemma}[Coupling comparison]
\label{lem:n9v51-coupling-comparison}
If $X,X'$ are coupled configurations, then
\begin{equation}
 |\mathbb EF(X)-\mathbb EF(X')|
 \le\sum_{i\in\operatorname{supp}F}
 \delta_i(F)\,\mathbb P(X_i\ne X'_i).
 \label{eq:n9v51-coupling-comparison}
\end{equation}
\end{lemma}

\begin{proof}
Change the coordinates of $X$ to those of $X'$ one at a time and
telescope.  Each changed coordinate contributes at most
$\delta_i(F)$ and only on its disagreement event.  Take expectations.
\end{proof}

\begin{theorem}[Averaged boundary mixing]
\label{thm:n9v51-boundary-mixing}
Assume the path compiler uniformly for any two admissible boundary
conditions differing on $A$.  Then every bounded local $F$ satisfies
\begin{equation}
 |\mathbb E^\eta F-\mathbb E^{\eta'}F|
 \le
 \frac{e^{-\mu d(\operatorname{supp}F,A)}}{1-\zeta}
 \sum_{i\in\operatorname{supp}F}\delta_i(F).
 \label{eq:n9v51-boundary-mixing}
\end{equation}
If the infinite-volume specification exists and this comparison applies
after conditioning on the support of a bounded local observable $G$,
then
\begin{equation}
 |\operatorname{Cov}_\mu(F,G)|
 \le
 \frac{\|G\|_\infty}{1-\zeta}
 e^{-\mu d(\operatorname{supp}F,\operatorname{supp}G)}
 \sum_{i\in\operatorname{supp}F}\delta_i(F).
 \label{eq:n9v51-covariance}
\end{equation}
\end{theorem}

\begin{proof}
Combine \Cref{lem:n9v51-coupling-comparison} with
\eqref{eq:n9v51-disagreement-decay}.  For covariance, condition on the
configuration in the support of $G$.  The conditional expectation of
$F$ has oscillation bounded by the first estimate with
$A=\operatorname{supp}G$.  Subtract its mean and multiply by $G$;
the expectation is bounded by $\|G\|_\infty$ times that oscillation.
\end{proof}

\begin{firewall}[Covariance decay is not yet an OS mass theorem]
\label{fire:n9v51-OS}
Turning \eqref{eq:n9v51-covariance} into a Hamiltonian mass gap requires
a reflection-positive Euclidean measure, a time-reflection-compatible
cofinal limit, density of the physical observable algebra, and the
correct physical conversion of graph distance.  The averaged
disagreement theorem supplies clustering only.
\end{firewall}

\subsection{Marginalization and feedback}

Let $K=G+B$ be the nonnegative averaged propagation kernel.  Under a
retained/eliminated split, define
\begin{equation}
 A_K:=K_{RE}(I-K_{EE})^{-1}K_{ER}.
 \label{eq:n9v51-K-return}
\end{equation}

\begin{proposition}[Averaged Schur-path recursion]
\label{prop:n9v51-averaged-Schur}
If $\|K_{EE}\|_\mu<1$ and
$\tau:=\|K_{RE}\|_\mu\|K_{ER}\|_\mu/
(1-\|K_{EE}\|_\mu)<1$, then the retained averaged propagation has the
safe weighted bound
\begin{equation}
 \|K^{\rm eff}\|_\mu
 \le
 \frac{\|K_{RR}\|_\mu+\tau}{1-\tau}.
 \label{eq:n9v51-averaged-recursion}
\end{equation}
The denominator retains repeated target--eliminated--target feedback.
\end{proposition}

\begin{proof}
The weighted Neumann series bounds
$\|(I-K_{EE})^{-1}\|_\mu\le
(1-\|K_{EE}\|_\mu)^{-1}$ and hence
$\|A_K\|_\mu\le\tau$.  Apply the scalar feedback solution
\eqref{eq:n9v50-feedback} rowwise and sum in the weighted norm.
\end{proof}

\subsection{Averaged amplitude-disagreement card}

\begin{definition}[Averaged amplitude-disagreement card (AADC)]
\label{def:n9v51-AADC}
AADC requires:
\begin{enumerate}
\item an exact coupling/path representation
\eqref{eq:n9v51-path-representation} for the normalized physical block
measure;
\item uniform conditional defect product domination or a direct mixed
path estimate;
\item support-local $L^2$ regular amplitudes loading
\eqref{eq:n9v51-mixed-path};
\item the strict weighted branching condition \eqref{eq:n9v51-zeta};
\item the feedback-corrected marginal recursion
\eqref{eq:n9v51-averaged-recursion} at every elimination step; and
\item reflection, gauge/constraint, metric-uniformity, and cofinal
compatibility needed by the intended physical conclusion.
\end{enumerate}
\end{definition}

\begin{theorem}[AADC gives averaged exponential clustering]
\label{thm:n9v51-AADC}
Under AADC, boundary sensitivity and bounded local covariances decay
exponentially with rate at least the declared $\mu$, uniformly in the
finite regulator.  The conclusion survives each controlled
marginalization step whose effective weighted norm remains below one.
\end{theorem}

\begin{proof}
Use \Cref{thm:n9v51-path-compiler,
thm:n9v51-boundary-mixing} and iterate
\Cref{prop:n9v51-averaged-Schur}.  AADC items 1--3 provide the physical
coupling rather than merely the algebraic path sum.
\end{proof}

\begin{assessment}[Progress at ninth-series V51]
\label{ass:n9v51-status}
The valid replacement for worst-exterior Higgs Dobrushin contraction is
now explicit.  Regular good-field propagation and rare amplitude
defects enter one connected path weight $q+b$.  Uniform conditional
quartic tails give $b=\sqrt\pi$ when paired with a local $L^2$ path
amplitude.  A strict weighted branching inequality yields exponential
boundary and covariance decay while retaining the marginal feedback
denominator.
\end{assessment}

\begin{programme}[Ninth-series V52: nested exterior defects]
\label{prog:n9v51-V52}
V52 will attack the missing product-domination hypothesis when a block's
quartic tail is uniform only under good exterior amplitudes.  It will
construct an enlarged connected defect closure and derive the resulting
effective branching weight, or record the exact failure if nested
sources proliferate without a subcritical bound.
\end{programme}

\begin{firewall}[Claim boundary after ninth-series V51]
\label{fire:n9v51-boundary}
V51 proves an abstract averaged disagreement/path theorem.  It does not
construct the required physical coupling, prove uniform defect product
domination for the unbounded Higgs exterior, verify the Standard Model
branching constant, establish reflection positivity or OS
reconstruction, or construct the full continuum.
\end{firewall}

\begin{theorem}[Ninth-series V51 result]
\label{thm:n9v51-result}
If a physical disagreement admits connected mixed paths with regular
weight $q$ and defect weight $b$, then
$\Delta e^\mu(q+b)<1$ yields exponential boundary sensitivity and local
covariance decay.  A uniform conditional quartic tail and support-local
$L^2$ amplitude load $b=\pi^{1/2}$; controlled marginalization retains
the feedback recursion \eqref{eq:n9v51-averaged-recursion}.
\end{theorem}

\begin{proof}
Use \Cref{thm:n9v51-path-compiler,lem:n9v51-product,
lem:n9v51-L2-defect,thm:n9v51-boundary-mixing,
prop:n9v51-averaged-Schur,thm:n9v51-AADC}.
\end{proof}

\paragraph{Parent-version record.}
V51 was formed from
\texttt{Renewal\_SM\_Einstein\_Continuum\_Research\_Notes\_V50.tex}.
The V50 parent had $18\,189$ logical source lines, $657\,459$ bytes, and
SHA--256
\[
 \texttt{8A10C3F2ADAAF8C6C560EDBB23CC5D303CDAE2AE82AE20BB1208C4F8559D4469}.
\]

% =============================================================================
% END APPENDED NINTH-SERIES V51 MATERIAL
% =============================================================================
% =============================================================================
% BEGIN APPENDED NINTH-SERIES V52 MATERIAL
% =============================================================================

\clearpage
\section{Ninth-series V52: maximal Higgs-defect closure and joint
quartic component weights}
\label{sec:v9v52}

\subsection{Maximal components remove the unbounded exterior source}

Let $\mathcal G=(V,E)$ be the finite regulator block graph with degree at
most $\Delta$.  Fix $R>0$ and call a block bad when
\begin{equation}
 |\phi_i|>R.
 \label{eq:n9v52-bad}
\end{equation}
The bad blocks split uniquely into maximal connected components.  If
$C$ is such a component, every block in its exterior vertex boundary
\begin{equation}
 \partial C:=\{j\notin C:\exists i\in C,\ (i,j)\in E\}
 \label{eq:n9v52-boundary}
\end{equation}
satisfies $|\phi_j|\le R$.  Thus the unbounded-source problem of V49 is
converted into a joint conditional estimate on $C$ with a bounded
boundary source.  The cost is that the whole component, rather than one
site, must be integrated at once.

\subsection{A joint component coercivity hypothesis}

Let $\phi_C=(\phi_i)_{i\in C}\in(\mathbb R^D)^C$.  Assume that, conditional
on exterior data $\xi$ with $|\xi_j|\le R$ for $j\in\partial C$, the
positive component law has density
\begin{equation}
 d\mu_C(\phi_C\mid\xi)
 =Z_C(\xi)^{-1}e^{-U_C(\phi_C\mid\xi)/\hbar}\,d\phi_C.
 \label{eq:n9v52-component-law}
\end{equation}
Suppose constants $a>0$ and $A,J\ge0$, independent of $C$, $R$, the
regulator, and admissible $\xi$, give
\begin{equation}
 U_C(\phi_C\mid\xi)
 \ge
 a\sum_{i\in C}|\phi_i|^4
 -A|C|-JR\sum_{i\in C}|\phi_i|.
 \label{eq:n9v52-joint-coercivity}
\end{equation}
The last sum safely overcounts the boundary incidence; a fixed degree is
absorbed in $J$.  Also suppose that on the unit product ball
\begin{equation}
 Q_C:=\{|\phi_i|\le r_0\text{ for every }i\in C\}
 \label{eq:n9v52-product-ball}
\end{equation}
one has
\begin{equation}
 U_C(\phi_C\mid\xi)
 \le M(1+R)|C|
 \label{eq:n9v52-joint-upper}
\end{equation}
for fixed $r_0,M>0$.

\begin{lemma}[Joint source absorption]
\label{lem:n9v52-source}
There is a constant
\begin{equation}
 C_J:=\frac{3}{4^{4/3}}J^{4/3}(a/2)^{-1/3}
 \label{eq:n9v52-CJ}
\end{equation}
such that
\begin{equation}
 U_C(\phi_C\mid\xi)
 \ge
 \frac a2\sum_{i\in C}|\phi_i|^4
 -\bigl[A+C_JR^{4/3}\bigr]|C|.
 \label{eq:n9v52-absorbed-coercivity}
\end{equation}
\end{lemma}

\begin{proof}
Apply \eqref{eq:n9v47-br} at each site with
$b=JR$ and $\delta=a/2$, then sum.
\end{proof}

\begin{lemma}[Component partition floor]
\label{lem:n9v52-component-floor}
Let $v_0=|B_D(0,r_0)|$.  Then
\begin{equation}
 Z_C(\xi)\ge
 \left[v_0e^{-M(1+R)/\hbar}\right]^{|C|}.
 \label{eq:n9v52-component-floor}
\end{equation}
\end{lemma}

\begin{proof}
Restrict the integral to $Q_C$, whose volume is $v_0^{|C|}$, and use
\eqref{eq:n9v52-joint-upper}.
\end{proof}

\subsection{Exponential component activity}

Let
\begin{equation}
 I_*:=\int_{\mathbb R^D}
 e^{-a|\phi|^4/(4\hbar)}\,d\phi<\infty.
 \label{eq:n9v52-Istar}
\end{equation}

\begin{theorem}[Maximal-component quartic bound]
\label{thm:n9v52-component}
For every finite connected $C$ and every exterior configuration with
$|\xi_j|\le R$ on $\partial C$,
\begin{equation}
 \mu_C\{|\phi_i|>R\text{ for all }i\in C\mid\xi\}
 \le \pi_R^{|C|},
 \label{eq:n9v52-component-probability}
\end{equation}
where
\begin{equation}
 \pi_R:=
 K_0\exp\!\left[
 -\frac{aR^4}{4\hbar}
 +\frac{A+C_JR^{4/3}+M(1+R)}{\hbar}
 \right],
 \qquad K_0:=I_*/v_0.
 \label{eq:n9v52-piR}
\end{equation}
In particular $\pi_R\to0$ with leading exponent $-aR^4/(4\hbar)$.
\end{theorem}

\begin{proof}
On the event in question, use
\eqref{eq:n9v52-absorbed-coercivity} and split
\[
 e^{-a|\phi_i|^4/(2\hbar)}
 \le
 e^{-aR^4/(4\hbar)}
 e^{-a|\phi_i|^4/(4\hbar)}
\]
for every $i\in C$.  The unnormalized numerator is therefore at most
\[
 \left[
 e^{[A+C_JR^{4/3}]/\hbar}
 e^{-aR^4/(4\hbar)}I_*
 \right]^{|C|}.
\]
Divide by \eqref{eq:n9v52-component-floor}.
\end{proof}

\begin{corollary}[Probability of a prescribed maximal component]
\label{cor:n9v52-maximal}
Let $\mathcal M_C(R)$ be the event that $C$ is a maximal connected
component of the bad set.  Then
\begin{equation}
 \mathbb P(\mathcal M_C(R))\le\pi_R^{|C|}.
 \label{eq:n9v52-maximal}
\end{equation}
\end{corollary}

\begin{proof}
On $\mathcal M_C(R)$, every site of $C$ is bad and every site of
$\partial C$ is good.  Condition on all variables outside $C$.  The
boundary values obey the hypotheses of
\Cref{thm:n9v52-component}; integrate its uniform bound.
\end{proof}

\begin{firewall}[Positive component law]
\label{fire:n9v52-positive}
The component estimate uses a positive normalized conditional law.
The V49 fermion-phase firewall remains: a complex chiral determinant
cannot be inserted into \eqref{eq:n9v52-component-law} as a probability
density without a separate phase theorem.
\end{firewall}

\subsection{Why the joint hypotheses are natural but not automatic}

\begin{lemma}[Absorption of finite-range bilinear Higgs couplings]
\label{lem:n9v52-bilinear}
For $u,v\ge0$, $\kappa\ge0$, and $\delta>0$,
\begin{equation}
 \kappa uv
 \le\frac{\delta}{2}(u^4+v^4)
 +\frac{\kappa^2}{4\delta}.
 \label{eq:n9v52-bilinear}
\end{equation}
Consequently, on a degree-$\Delta$ graph, all internal bilinear
nearest-neighbor magnitudes can be absorbed into a quartic floor by
choosing $\delta$ below the one-site quartic coefficient divided by
$\Delta$, at the cost of a constant proportional to $|C|$.
\end{lemma}

\begin{proof}
First $\kappa uv\le\frac\kappa2(u^2+v^2)$.  Apply
$\frac\kappa2r^2\le\frac\delta2r^4+\kappa^2/(8\delta)$ to $u$ and $v$.
Each site occurs in at most $\Delta$ internal edges, giving the last
statement after summation.
\end{proof}

\begin{proposition}[Loading the joint coercive form]
\label{prop:n9v52-load-coercivity}
Suppose each site has a uniform positive quartic coefficient, all
negative quadratic terms are uniformly bounded, the internal
Higgs/gauge link interaction is finite range with bounded coefficients,
and every boundary neighbor of a maximal component has amplitude at most
$R$.  Then repeated use of
\Cref{lem:n9v52-bilinear,lem:n9v49-source-absorption} yields
\eqref{eq:n9v52-joint-coercivity} with constants independent of the
component size.  A component-size-independent upper bound
\eqref{eq:n9v52-joint-upper} additionally requires uniform bounds on all
determinant, coarea, metric, and reference-density rows on $Q_C$.
\end{proposition}

\begin{proof}
Absorb negative quadratic and internal bilinear rows into a fixed
fraction of the quartic site sum.  Each site is charged at most a bounded
number of times because the degree is bounded.  Boundary bilinear terms
have one factor at most $R$ and give the linear source in
\eqref{eq:n9v52-joint-coercivity}.  On $Q_C$, the elementary bosonic
terms are bounded per site and edge; the final qualification lists the
additional rows needed for the full density upper bound.
\end{proof}

\begin{firewall}[Component partition floors include every density row]
\label{fire:n9v52-floor}
Bosonic quartic coercivity alone does not prove
\eqref{eq:n9v52-joint-upper}.  A determinant zero, coarea pole, metric
degeneration, or singular chart density can destroy the product-ball
floor.  The V9 and V34--V39 cards must be imposed on the same component
chart.
\end{firewall}

\subsection{Rooted component summation}

\begin{lemma}[Connected-set count]
\label{lem:n9v52-connected-count}
The number of connected vertex sets of size $m\ge1$ containing a fixed
root is at most
\begin{equation}
 \Delta^{2(m-1)}.
 \label{eq:n9v52-connected-count}
\end{equation}
\end{lemma}

\begin{proof}
Every finite connected set has a spanning tree rooted at the chosen
vertex.  A depth-first traversal of that tree has length
$2(m-1)$.  Recording the successive neighbor choices gives at most
$\Delta^{2(m-1)}$ words.  This encoding may count one set many times,
which is harmless for an upper bound.
\end{proof}

\begin{corollary}[Rooted defect-polymer activity]
\label{cor:n9v52-rooted}
Let a component of size $m$ carry an $L^2$-converted activity at most
$\pi_R^{m/2}$.  If
\begin{equation}
 \Delta^2\sqrt{\pi_R}<1,
 \label{eq:n9v52-rooted-small}
\end{equation}
then the total activity of components containing a fixed root is at most
\begin{equation}
 b_R^{\rm comp}
 :=
 \frac{\sqrt{\pi_R}}
 {1-\Delta^2\sqrt{\pi_R}}.
 \label{eq:n9v52-rooted-activity}
\end{equation}
\end{corollary}

\begin{proof}
Use \Cref{lem:n9v52-connected-count} and sum
\[
 \sum_{m\ge1}\Delta^{2(m-1)}\pi_R^{m/2}
 =\frac{\sqrt{\pi_R}}
 {1-\Delta^2\sqrt{\pi_R}}.
\]
\end{proof}

\begin{remark}[Component activity replaces single-site product
domination]
\label{rem:n9v52-replacement}
V51 used $b_R=\sqrt{\pi_R}$ under uniform single-site conditional product
domination.  When that hypothesis fails because exterior Higgs
amplitudes are unbounded, the maximal-component construction supplies
the larger but still vanishing rooted weight
$b_R^{\rm comp}$.  The connected-set denominator is the price of nested
defect closure.
\end{remark}

\subsection{Insertion into the disagreement compiler}

Assume an exact coupling expansion in which each maximal bad component
encountered by a disagreement path is collapsed to a connected defect
polymer, and the remaining steps use the regular good-field coefficient
$q_R$.

\begin{theorem}[Nested-defect disagreement criterion]
\label{thm:n9v52-nested-criterion}
Suppose:
\begin{enumerate}
\item maximal components obey \eqref{eq:n9v52-maximal};
\item their support-local insertion amplitudes admit the square-root
conversion used in \Cref{cor:n9v52-rooted};
\item the coupling/path expansion counts every encountered component
once and regular steps have weight at most $q_R$; and
\item the effective component-incidence degree is at most
$\Delta_{\rm eff}$.
\end{enumerate}
If
\begin{equation}
 \Delta_{\rm eff}e^\mu
 \left(q_R+b_R^{\rm comp}\right)<1,
 \label{eq:n9v52-nested-criterion}
\end{equation}
then the V51 boundary-sensitivity and covariance conclusions hold with
$b=b_R^{\rm comp}$.
\end{theorem}

\begin{proof}
Sum all connected component shapes rooted at a path encounter using
\Cref{cor:n9v52-rooted}; this replaces the component by one effective
defect step of weight $b_R^{\rm comp}$.  Exact component ownership avoids
recounting.  The remaining expansion satisfies the V51 mixed path bound
with step weights $q_R$ and $b_R^{\rm comp}$.  Apply
\Cref{thm:n9v51-path-compiler,thm:n9v51-boundary-mixing}.
\end{proof}

\begin{corollary}[Large-threshold behavior]
\label{cor:n9v52-large-R}
Under uniform component constants,
\begin{equation}
 b_R^{\rm comp}
 =
 O\!\left(
 \exp\{-aR^4/(8\hbar)+O(R^{4/3})\}
 \right)
 \label{eq:n9v52-large-R}
\end{equation}
as $R\to\infty$.  Hence the nested component correction retains quartic
decay, with half the probability exponent because of the local
$L^2$ conversion.
\end{corollary}

\begin{proof}
Take the square root of \eqref{eq:n9v52-piR}.  Its denominator in
\eqref{eq:n9v52-rooted-activity} tends to one.
\end{proof}

\begin{firewall}[Exact component encounter is a required coupling row]
\label{fire:n9v52-encounter}
The component probability and rooted sum do not by themselves construct
the disagreement coupling.  The physical expansion must show that a
path entering one maximal component can be charged to that component
once, with bounded exit incidence.  Otherwise repeated visits can
reintroduce uncontrolled multiplicity.
\end{firewall}

\subsection{Maximal Higgs-defect closure card}

\begin{definition}[Maximal Higgs-defect closure card (MHDCC)]
\label{def:n9v52-MHDCC}
MHDCC requires:
\begin{enumerate}
\item a positive normalized joint component law with
\eqref{eq:n9v52-joint-coercivity} and
\eqref{eq:n9v52-joint-upper};
\item exact maximal-component ownership and bounded graph degree;
\item the uniform component probability \eqref{eq:n9v52-maximal};
\item a support-local $L^2$ insertion bound and rooted smallness
\eqref{eq:n9v52-rooted-small};
\item an exact once-per-component disagreement expansion with bounded
exit incidence; and
\item chart, coarea, determinant, metric, gauge, and fermion-phase
compatibility on the same enlarged component.
\end{enumerate}
\end{definition}

\begin{theorem}[MHDCC loads the nested row of AADC]
\label{thm:n9v52-MHDCC}
MHDCC replaces the unavailable uniform single-site exterior tail by the
rooted component weight \eqref{eq:n9v52-rooted-activity}.  If
\eqref{eq:n9v52-nested-criterion} holds, it loads the defect row of AADC
and yields uniform averaged exponential clustering for the declared
bounded local observables.
\end{theorem}

\begin{proof}
Use \Cref{thm:n9v52-component,cor:n9v52-maximal,
cor:n9v52-rooted,thm:n9v52-nested-criterion,
thm:n9v51-AADC}.
\end{proof}

\begin{assessment}[Progress at ninth-series V52]
\label{ass:n9v52-status}
The nested exterior-amplitude problem now has a concrete repair.
Maximal bad components have automatically good boundaries; joint
quartic coercivity gives an activity exponential in component size and
$R^4$.  Connected-set summation turns that activity into a rooted defect
weight.  The remaining physical rows are the complete component
partition floor, fermion positivity, and an exact once-per-component
disagreement expansion.
\end{assessment}

\begin{programme}[Ninth-series V53: fermion-phase transfer]
\label{prog:n9v52-V53}
V53 will attack the positivity firewall.  It will derive the exact
reweighting denominator for a complex fermion determinant, prove that
modulus-measure bounds transfer only with an inverse average-phase cost,
and identify sufficient nonvanishing or paired-representation
conditions without assuming that anomaly cancellation implies
positivity.
\end{programme}

\begin{firewall}[Claim boundary after ninth-series V52]
\label{fire:n9v52-boundary}
V52 proves the positive-measure maximal-component theorem.  It does not
prove MHDCC for the full Standard Model, control the chiral determinant
phase, construct the disagreement coupling, establish reflection
positivity or OS reconstruction, or construct the full continuum.
\end{firewall}

\begin{theorem}[Ninth-series V52 result]
\label{thm:n9v52-result}
A maximal Higgs-amplitude component has a good exterior boundary.
Under uniform joint quartic coercivity and a component partition floor,
its probability is at most $\pi_R^{|C|}$ with
$\log\pi_R=-aR^4/(4\hbar)+O(R^{4/3})$.  Rooted connected summation yields
the effective defect weight \eqref{eq:n9v52-rooted-activity}, which loads
the V51 averaged path theorem under
\eqref{eq:n9v52-nested-criterion}.
\end{theorem}

\begin{proof}
Use \Cref{lem:n9v52-source,lem:n9v52-component-floor,
thm:n9v52-component,cor:n9v52-maximal,
lem:n9v52-connected-count,cor:n9v52-rooted,
thm:n9v52-nested-criterion}.
\end{proof}

\paragraph{Parent-version record.}
V52 was formed from
\texttt{Renewal\_SM\_Einstein\_Continuum\_Research\_Notes\_V51.tex}.
The V51 parent had $18\,677$ logical source lines, $674\,103$ bytes, and
SHA--256
\[
 \texttt{5E308CBD574675A1748CE246F2437988CD80884AA15D442AD5D9659ECBC26381}.
\]

% =============================================================================
% END APPENDED NINTH-SERIES V52 MATERIAL
% =============================================================================
% =============================================================================
% BEGIN APPENDED NINTH-SERIES V53 MATERIAL
% =============================================================================

\clearpage
\section{Ninth-series V53: fermion-phase reweighting, positivity
alternatives, and the Einstein phase source}
\label{sec:v9v53}

\subsection{Exact phase-quenched representation}

At a finite regulator, let $\lambda$ be a positive bosonic reference
measure and let the integrated fermion factor be
\begin{equation}
 \mathcal D(\Phi)=|\mathcal D(\Phi)|e^{i\Theta(\Phi)}.
 \label{eq:n9v53-polar}
\end{equation}
Assume
\begin{equation}
 0<Z_{\rm abs}:=\int|\mathcal D|\,d\lambda<\infty
 \label{eq:n9v53-Zabs}
\end{equation}
and define the positive phase-quenched probability
\begin{equation}
 d\nu:=Z_{\rm abs}^{-1}|\mathcal D|\,d\lambda.
 \label{eq:n9v53-nu}
\end{equation}
The average phase and physical partition function are
\begin{equation}
 \zeta:=\mathbb E_\nu e^{i\Theta},
 \qquad
 Z=Z_{\rm abs}\zeta.
 \label{eq:n9v53-zeta}
\end{equation}
When $\zeta\ne0$, the normalized physical functional is
\begin{equation}
 \langle F\rangle_{\rm phys}
 =\frac{\mathbb E_\nu[Fe^{i\Theta}]}{\zeta}.
 \label{eq:n9v53-reweighting}
\end{equation}

\begin{theorem}[Phase reweighting norm]
\label{thm:n9v53-reweighting}
For $1\le p\le\infty$ and conjugate exponent $p'$, every
$F\in L^p(\nu)$ satisfies
\begin{equation}
 |\langle F\rangle_{\rm phys}|
 \le\frac{\|F\|_{L^p(\nu)}
              \|e^{i\Theta}\|_{L^{p'}(\nu)}}{|\zeta|}
 =\frac{\|F\|_{L^p(\nu)}}{|\zeta|}.
 \label{eq:n9v53-reweighting-bound}
\end{equation}
The factor $|\zeta|^{-1}$ is unavoidable for this general transfer.
\end{theorem}

\begin{proof}
Apply Hölder to the numerator of \eqref{eq:n9v53-reweighting} and use
$|e^{i\Theta}|=1$.  To see that no smaller universal denominator factor
is available, take $F=e^{-i\Theta}$; its physical numerator equals one
and the normalized value is $\zeta^{-1}$.
\end{proof}

\begin{corollary}[Transfer of a rare-set insertion]
\label{cor:n9v53-rare}
For an event $E$ and $F\in L^2(\nu)$,
\begin{equation}
 |\langle F\mathbf1_E\rangle_{\rm phys}|
 \le\frac{\|F\|_{L^2(\nu)}\,\nu(E)^{1/2}}{|\zeta|}.
 \label{eq:n9v53-rare-transfer}
\end{equation}
Thus a V49--V52 defect weight transfers algebraically only if its decay
beats the inverse average-phase cost.
\end{corollary}

\begin{proof}
Apply \Cref{thm:n9v53-reweighting} to $F\mathbf1_E$ and then
Cauchy--Schwarz.
\end{proof}

\begin{corollary}[Volume-dependent phase cost]
\label{cor:n9v53-volume-cost}
If along a cofinal sequence
\begin{equation}
 |\zeta_n|\ge e^{-\chi_n}
 \label{eq:n9v53-zeta-lower}
\end{equation}
and a local defect numerator is bounded by $\delta_n$, then its physical
bound is at most $e^{\chi_n}\delta_n$.  It vanishes only under the actual
balance
\begin{equation}
 \log\delta_n+\chi_n\longrightarrow-\infty.
 \label{eq:n9v53-phase-balance}
\end{equation}
In particular an average phase exponentially small in total volume can
destroy a volume-uniform local modulus-measure estimate.
\end{corollary}

\begin{proof}
Insert \eqref{eq:n9v53-zeta-lower} in
\eqref{eq:n9v53-reweighting-bound} and take logarithms.
\end{proof}

\subsection{Elementary sufficient nonvanishing conditions}

\begin{proposition}[Small total phase]
\label{prop:n9v53-small-phase}
If a real representative of $\Theta$ satisfies
\begin{equation}
 |\Theta|\le\delta<\frac{\pi}{2}
 \quad \nu\text{-almost surely},
 \label{eq:n9v53-uniform-phase}
\end{equation}
then
\begin{equation}
 |\zeta|\ge\operatorname{Re}\zeta\ge\cos\delta>0.
 \label{eq:n9v53-cos-bound}
\end{equation}
More generally, if
\begin{equation}
 \mathbb E_\nu|1-e^{i\Theta}|\le\eta<1,
 \label{eq:n9v53-L1-phase}
\end{equation}
then
\begin{equation}
 |\zeta|\ge1-\eta.
 \label{eq:n9v53-L1-lower}
\end{equation}
\end{proposition}

\begin{proof}
Under \eqref{eq:n9v53-uniform-phase},
$\cos\Theta\ge\cos\delta$.  For the second assertion,
\[
 |\zeta|=|1-\mathbb E_\nu(1-e^{i\Theta})|
 \ge1-\mathbb E_\nu|1-e^{i\Theta}|.
\]
\end{proof}

\begin{corollary}[An $L^1$ phase representative]
\label{cor:n9v53-L1-theta}
If $\mathbb E_\nu|\Theta|\le\eta<1$, then
$|\zeta|\ge1-\eta$.
\end{corollary}

\begin{proof}
Use $|1-e^{it}|\le|t|$ in
\eqref{eq:n9v53-L1-phase}.
\end{proof}

\begin{firewall}[Extensive small local phases need not give a small total
phase]
\label{fire:n9v53-extensive}
If $\Theta=\sum_X\Theta_X$, a small bound on each $\Theta_X$ can still
produce an extensive total phase.  Neither
\eqref{eq:n9v53-uniform-phase} nor
\eqref{eq:n9v53-L1-phase} follows from per-block smallness without a
summable connected cancellation or cluster theorem.
\end{firewall}

\subsection{Exact conjugate pairing}

\begin{theorem}[Conjugate-determinant positivity]
\label{thm:n9v53-pairing}
Suppose the regulated fermion species can be paired so that for every
bosonic configuration,
\begin{equation}
 \det M_{\bar s}(\Phi)=\overline{\det M_s(\Phi)}.
 \label{eq:n9v53-conjugate-pair}
\end{equation}
Then each pair contributes
\begin{equation}
 \det M_s\,\det M_{\bar s}
 =|\det M_s|^2\ge0.
 \label{eq:n9v53-positive-pair}
\end{equation}
If all remaining fermion factors are nonnegative, the complete fermion
weight is positive, $\Theta=0$ modulo $2\pi$ off its zero set, and
$\zeta=1$.  The positive-measure tail and coupling theorems of
V49--V52 then apply without a phase denominator.
\end{theorem}

\begin{proof}
Multiply \eqref{eq:n9v53-conjugate-pair} by $\det M_s$.  Products of
nonnegative factors remain nonnegative.  The phase-quenched and physical
weights then coincide after normalization.
\end{proof}

\begin{remark}[Zero modes]
\label{rem:n9v53-zeros}
At $\det M_s=0$ the phase is undefined, but the paired product
\eqref{eq:n9v53-positive-pair} is well defined and zero.  Integrability
near that set remains the determinant--coarea and score-threshold problem
of V9 and V38--V39.
\end{remark}

\begin{firewall}[The chiral Standard Model is not automatically
conjugate-paired]
\label{fire:n9v53-SM-pairing}
The physical Standard Model assigns different left- and right-handed
gauge representations and includes chiral weak couplings.  Cancellation
of perturbative gauge anomalies is a representation-theoretic identity;
it does not exhibit a configurationwise partner satisfying
\eqref{eq:n9v53-conjugate-pair}.  Positivity must be proved for the chosen
regulator rather than inferred from anomaly cancellation.
\end{firewall}

\subsection{Anomaly cancellation does not imply a nonzero average phase}

\begin{proposition}[Gauge-invariant phase with zero average]
\label{prop:n9v53-anomaly-no}
There is a positive gauge-invariant probability space and a
gauge-invariant phase whose infinitesimal gauge variation vanishes
identically but whose average phase is zero.
\end{proposition}

\begin{proof}
Take a two-point probability space with equal weights and trivial gauge
action.  Set $e^{i\Theta}=i$ at one point and
$e^{i\Theta}=-i$ at the other.  The phase is gauge invariant, so every
gauge variation is zero, while
$\mathbb E e^{i\Theta}=0$.
\end{proof}

\begin{corollary}[Separate anomaly and sign cards]
\label{cor:n9v53-separate-cards}
A regulator proof of anomaly cancellation or gauge invariance does not
load any positive lower bound for $|\zeta|$.  Gauge consistency and
phase noncancellation are logically separate interfaces.
\end{corollary}

\begin{proof}
Apply \Cref{prop:n9v53-anomaly-no}.
\end{proof}

\subsection{Metric differentiation of the average phase}

Let $g$ be a finite-dimensional regulator metric coordinate.  Write the
unnormalized phase-quenched weight as $w_{\rm abs}(g,\Phi)$ and define
\begin{equation}
 A_h:=D_g\log w_{\rm abs}[h],
 \qquad
 \dot\Theta_h:=D_g\Theta[h].
 \label{eq:n9v53-A-theta-dot}
\end{equation}
Assume differentiation under the integral is justified.

\begin{theorem}[Einstein average-phase response]
\label{thm:n9v53-phase-response}
The logarithmic derivative of the physical partition function splits as
\begin{equation}
 D_g\log Z[h]
 =D_g\log Z_{\rm abs}[h]+D_g\log\zeta[h],
 \label{eq:n9v53-logZ-split}
\end{equation}
where
\begin{equation}
 D_g\log\zeta[h]
 =
 \frac{
  \mathbb E_\nu[e^{i\Theta}i\dot\Theta_h]
  +\operatorname{Cov}_\nu(e^{i\Theta},A_h)
 }{\zeta}.
 \label{eq:n9v53-phase-response}
\end{equation}
Thus the fermion phase contributes an explicit Einstein source row and
inherits the inverse average-phase denominator.
\end{theorem}

\begin{proof}
Differentiate $Z=Z_{\rm abs}\zeta$.  For a normalized parameter-dependent
law,
\[
 D_g\mathbb E_\nu F[h]
 =\mathbb E_\nu[D_gF[h]]
  +\operatorname{Cov}_\nu(F,A_h).
\]
Set $F=e^{i\Theta}$, for which
$D_gF[h]=ie^{i\Theta}\dot\Theta_h$, and divide by $\zeta$.
\end{proof}

\begin{corollary}[A sufficient phase-source bound]
\label{cor:n9v53-phase-source}
If $|\zeta|\ge z_*>0$ and
$\dot\Theta_h,A_h\in L^2(\nu)$, then
\begin{equation}
 |D_g\log\zeta[h]|
 \le z_*^{-1}
 \left(
  \|\dot\Theta_h\|_{L^1(\nu)}
  +\|A_h-\mathbb EA_h\|_{L^2(\nu)}
 \right).
 \label{eq:n9v53-phase-source-bound}
\end{equation}
\end{corollary}

\begin{proof}
Use $|e^{i\Theta}|=1$, the triangle inequality, and Cauchy--Schwarz in
\eqref{eq:n9v53-phase-response}.
\end{proof}

\begin{firewall}[Phase derivatives at determinant zeros]
\label{fire:n9v53-phase-zeros}
The polar derivative $\dot\Theta_h$ can be singular where
$\mathcal D=0$.  The $L^2$ premise in
\Cref{cor:n9v53-phase-source} must be proved using the transverse
determinant powers and the V38--V39 threshold; it is not implied by
$|\mathcal D|$ being integrable.
\end{firewall}

\subsection{Why reweighting is not a probability coupling}

\begin{proposition}[Complex physical functional]
\label{prop:n9v53-not-probability}
Unless $e^{i\Theta}/\zeta$ is nonnegative real almost surely, the
functional \eqref{eq:n9v53-reweighting} is not integration against a
positive probability measure.  Consequently total-variation couplings,
conditional probabilities, and disagreement events used in AADC cannot
be constructed directly from it.
\end{proposition}

\begin{proof}
If a normalized density with respect to $\nu$ defines a positive
probability, that density must be nonnegative real almost surely.  Here
the density is $e^{i\Theta}/\zeta$.  A nonconstant phase violates this
condition.
\end{proof}

\begin{firewall}[A modulus coupling plus a denominator is not yet
physical clustering]
\label{fire:n9v53-clustering}
One may couple the positive modulus law $\nu$ and use
\eqref{eq:n9v53-rare-transfer} for individual numerators.  Physical
connected correlations also contain the global phase in both numerator
and normalization.  A nonlocal phase can correlate distant observables
even when $\nu$ clusters.  Transferring AADC clustering therefore needs
conjugate positivity or a separate quasi-local phase/cumulant theorem,
not merely $|\zeta|>0$.
\end{firewall}

\subsection{Fermion-phase transfer card}

\begin{definition}[Fermion-phase transfer card (FPTC)]
\label{def:n9v53-FPTC}
FPTC requires one of:
\begin{enumerate}
\item configurationwise conjugate pairing and positivity as in
\Cref{thm:n9v53-pairing}; or
\item a phase-quenched representation with a quantitative average-phase
lower bound, quasi-local phase control sufficient for the intended
correlations, and the balance \eqref{eq:n9v53-phase-balance}.
\end{enumerate}
It also requires integrable metric phase derivatives, determinant-zero
thresholds, gauge consistency, and a separate reflection-positivity
proof.
\end{definition}

\begin{theorem}[FPTC transfer]
\label{thm:n9v53-FPTC}
Under FPTC item 1, the V49--V52 positive-measure tail, component, and
averaged-disagreement machinery applies to the full paired fermion
weight.  Under item 2, modulus-measure local insertion estimates transfer
with the explicit factor $|\zeta|^{-1}$, while clustering transfers only
through the additional quasi-local phase theorem.  In both cases the
Einstein source includes \eqref{eq:n9v53-phase-response}.
\end{theorem}

\begin{proof}
The positive case is \Cref{thm:n9v53-pairing}.  The reweighted insertion
case is \Cref{thm:n9v53-reweighting,cor:n9v53-rare,
cor:n9v53-volume-cost}.  The source identity is
\Cref{thm:n9v53-phase-response}; the clustering qualification is
\Cref{fire:n9v53-clustering}.
\end{proof}

\begin{assessment}[Progress at ninth-series V53]
\label{ass:n9v53-status}
The fermion positivity bottleneck now has an exact ledger.  Every
phase-quenched estimate costs the inverse average phase; a volume-small
average can defeat local quartic rarity.  Configurationwise conjugate
pairing removes the cost, whereas anomaly cancellation alone does not.
Metric differentiation introduces a separate phase source with the same
denominator and possible determinant-zero singularities.
\end{assessment}

\begin{programme}[Ninth-series V54: quasi-local phase clusters]
\label{prog:n9v53-V54}
V54 will derive a finite-volume nonvanishing and local-correlation
transfer theorem when the phase factor has an absolutely summable
connected polymer expansion.  It will distinguish a local phase pressure,
which may be extensive, from the connected phase insertions needed for
local clustering.
\end{programme}

\begin{firewall}[Claim boundary after ninth-series V53]
\label{fire:n9v53-boundary}
V53 proves exact reweighting and sufficient positivity alternatives.  It
does not prove FPTC for the chiral Standard Model regulator, a uniform
average-phase lower bound, quasi-local phase control, reflection
positivity, OS reconstruction, or the full continuum.
\end{firewall}

\begin{theorem}[Ninth-series V53 result]
\label{thm:n9v53-result}
A complex fermion determinant transfers a phase-quenched $L^p$ estimate
with the unavoidable cost $|\zeta|^{-1}$.  Exact conjugate determinant
pairs give a positive weight and $\zeta=1$; anomaly cancellation does
not.  The Einstein partition derivative contains the phase response
\eqref{eq:n9v53-phase-response}, including the same noncancellation
denominator.
\end{theorem}

\begin{proof}
Use \Cref{thm:n9v53-reweighting,thm:n9v53-pairing,
prop:n9v53-anomaly-no,thm:n9v53-phase-response,
thm:n9v53-FPTC}.
\end{proof}

\paragraph{Parent-version record.}
V53 was formed from
\texttt{Renewal\_SM\_Einstein\_Continuum\_Research\_Notes\_V52.tex}.
The V52 parent had $19\,127$ logical source lines, $689\,443$ bytes, and
SHA--256
\[
 \texttt{DE8119EDD035EB70C78C74167C00FFE19D6136031CE922771FC867D78FD90190}.
\]

% =============================================================================
% END APPENDED NINTH-SERIES V53 MATERIAL
% =============================================================================
% =============================================================================
% BEGIN APPENDED NINTH-SERIES V54 MATERIAL
% =============================================================================

\clearpage
\section{Ninth-series V54: local phase cancellation and a connected
complex-cluster criterion}
\label{sec:v9v54}

\subsection{Global average phase can be the wrong local norm}

V53 gives the general reweighting bound $|\zeta|^{-1}$.  That bound is
sharp for arbitrary observables, but it can be grossly pessimistic for a
fixed local observable when phase factors outside its support cancel
between numerator and denominator.

Let
\begin{equation}
 (X,\nu)=\prod_{x\in\Lambda}(X_x,\nu_x),
 \qquad
 e^{i\Theta(X)}=\prod_{x\in\Lambda}u_x(X_x),
 \qquad |u_x|=1.
 \label{eq:n9v54-product-phase}
\end{equation}
Put $\zeta_x=\int u_x\,d\nu_x$.

\begin{theorem}[Exact exterior phase cancellation]
\label{thm:n9v54-product-cancellation}
If every $\zeta_x\ne0$, then
\begin{equation}
 \zeta_\Lambda=\prod_{x\in\Lambda}\zeta_x.
 \label{eq:n9v54-product-zeta}
\end{equation}
For every observable $F$ supported in $S\subset\Lambda$,
\begin{equation}
 \langle F\rangle_{\rm phys}
 =
 \frac{
  \displaystyle\int F(X_S)\prod_{x\in S}u_x(X_x)
      \prod_{x\in S}d\nu_x
 }{
  \displaystyle\prod_{x\in S}\zeta_x
 }.
 \label{eq:n9v54-local-cancellation}
\end{equation}
Thus if $|\zeta_x|\ge z_*>0$, the local bound costs at most
$z_*^{-|S|}$, even though $|\zeta_\Lambda|$ may decay exponentially in
$|\Lambda|$.  Physical connected correlations between disjoint site
sets vanish exactly in this product model.
\end{theorem}

\begin{proof}
Fubini gives \eqref{eq:n9v54-product-zeta}.  In the numerator containing
$F$, integration over every $x\notin S$ produces $\zeta_x$; these factors
cancel the identical factors in the denominator, proving
\eqref{eq:n9v54-local-cancellation}.  The remaining normalized complex
functional is a product over sites, so observables on disjoint sets
factor.
\end{proof}

\begin{corollary}[Global phase lower bounds are sufficient, not
necessary, for local control]
\label{cor:n9v54-global-not-necessary}
A volume-uniform lower bound on $|\zeta_\Lambda|$ is not necessary for
uniform fixed-support expectations or clustering.  What is necessary is
control of the phase factors connected to the observable support.
\end{corollary}

\begin{proof}
Take a common $0<|\zeta_x|<1$.  Then
$|\zeta_\Lambda|=|\zeta_x|^{|\Lambda|}\to0$, while
\eqref{eq:n9v54-local-cancellation} has a fixed cost for fixed $S$ and
all disjoint connected correlations vanish.
\end{proof}

\subsection{Connected generating functional}

Let $F_1,\dots,F_m$ be bounded local insertions with supports
$A_1,\dots,A_m$.  For a finite regulator define
\begin{equation}
 Z_\Lambda(J)
 :=
 \mathbb E_{\nu_\Lambda}
 \exp\!\left\{
  i\Theta_\Lambda+\sum_{a=1}^mJ_aF_a
 \right\}.
 \label{eq:n9v54-source-partition}
\end{equation}
Assume $Z_\Lambda(J)\ne0$ in a polydisc around $J=0$ and choose the
analytic branch
\begin{equation}
 W_\Lambda(J):=\log Z_\Lambda(J).
 \label{eq:n9v54-W}
\end{equation}

\begin{lemma}[Connected derivatives]
\label{lem:n9v54-connected-derivatives}
At $J=0$,
\begin{align}
 \partial_{J_a}W_\Lambda
 &=\langle F_a\rangle_{\rm phys},
 \label{eq:n9v54-first-derivative}\\
 \partial_{J_a}\partial_{J_b}W_\Lambda
 &=\langle F_aF_b\rangle_{\rm phys}
   -\langle F_a\rangle_{\rm phys}
    \langle F_b\rangle_{\rm phys}.
 \label{eq:n9v54-second-derivative}
\end{align}
\end{lemma}

\begin{proof}
Differentiate $\log Z=Z'/Z$ once and twice.  The resulting ratios are
exactly the reweighted physical functionals.
\end{proof}

\subsection{An abstract absolutely convergent phase-cluster expansion}

Let $\mathcal X_\Lambda$ be the finite connected subsets of the regulator
block graph.  Suppose
\begin{equation}
 W_\Lambda(J)=\sum_{X\in\mathcal X_\Lambda}w_X(J)
 \label{eq:n9v54-cluster-expansion}
\end{equation}
absolutely and locally uniformly in the source polydisc, where $w_X$
depends on $J_a$ only if $X\cap A_a\ne\varnothing$.

\begin{definition}[Connected phase norm]
\label{def:n9v54-phase-norm}
For two source insertions, let $c_X(F,G)$ satisfy
\begin{equation}
 |\partial_F\partial_Gw_X(0)|
 \le c_X\|F\|_\infty\|G\|_\infty.
 \label{eq:n9v54-cX}
\end{equation}
The expansion has connected norm $(C,\mu)$ when
\begin{equation}
 \sup_{i\in\Lambda}
 \sum_{X\ni i}e^{\mu\operatorname{diam}X}c_X
 \le C
 \label{eq:n9v54-connected-norm}
\end{equation}
uniformly in the regulator.
\end{definition}

\begin{theorem}[Complex phase-cluster covariance decay]
\label{thm:n9v54-cluster-decay}
Assume \eqref{eq:n9v54-cluster-expansion} is locally uniformly absolutely
convergent and has connected norm $(C,\mu)$.  Then bounded local
observables $F,G$ with supports $A,B$ satisfy
\begin{equation}
 |\langle FG\rangle_{\rm phys}
 -\langle F\rangle_{\rm phys}\langle G\rangle_{\rm phys}|
 \le
 C|A|\|F\|_\infty\|G\|_\infty
 e^{-\mu d(A,B)}.
 \label{eq:n9v54-cluster-decay}
\end{equation}
The bound contains no inverse of the full-volume average phase.
\end{theorem}

\begin{proof}
Local uniform absolute convergence permits two source derivatives to be
taken termwise.  A cluster derivative
$\partial_F\partial_Gw_X$ vanishes unless $X$ meets both $A$ and $B$.
Every such connected $X$ has
$\operatorname{diam}X\ge d(A,B)$.  Therefore, using
\Cref{lem:n9v54-connected-derivatives},
\begin{align*}
 |\operatorname{Cov}_{\rm phys}(F,G)|
 &\le\|F\|_\infty\|G\|_\infty
 \sum_{\substack{X:X\cap A\ne\varnothing\\X\cap B\ne\varnothing}}c_X\\
 &\le e^{-\mu d(A,B)}
 \|F\|_\infty\|G\|_\infty
 \sum_{i\in A}\sum_{X\ni i}e^{\mu\operatorname{diam}X}c_X,
\end{align*}
which is \eqref{eq:n9v54-cluster-decay}.
\end{proof}

\begin{corollary}[Extensive phase pressure cancels from connected
correlators]
\label{cor:n9v54-pressure}
The source-free terms $w_X(0)$ may sum to an extensive complex pressure
and make $|Z_\Lambda(0)|$ exponentially small.  They do not enter
\eqref{eq:n9v54-cluster-decay}; only clusters meeting both source
supports survive the connected derivatives.
\end{corollary}

\begin{proof}
Source-free constants vanish after either source derivative.  The
remaining statement is the support argument in
\Cref{thm:n9v54-cluster-decay}.
\end{proof}

\begin{firewall}[Writing a formal logarithm is not a nonvanishing proof]
\label{fire:n9v54-log}
The expansion \eqref{eq:n9v54-cluster-expansion} must be constructed in a
domain where $Z_\Lambda(J)\ne0$ and the logarithm is analytic.  A formal
cumulant series or termwise diagram list does not prove nonvanishing,
absolute convergence, or the uniform connected norm.
\end{firewall}

\subsection{Local observables and one-source clusters}

Assume additionally that
\begin{equation}
 |\partial_Fw_X(0)|
 \le a_X\|F\|_\infty,
 \qquad
 \sup_i\sum_{X\ni i}a_X\le A_1.
 \label{eq:n9v54-one-source-norm}
\end{equation}

\begin{proposition}[Volume-uniform local expectation]
\label{prop:n9v54-local-expectation}
Under \eqref{eq:n9v54-one-source-norm},
\begin{equation}
 |\langle F\rangle_{\rm phys}|
 \le A_1|\operatorname{supp}F|\|F\|_\infty.
 \label{eq:n9v54-local-expectation}
\end{equation}
\end{proposition}

\begin{proof}
Differentiate \eqref{eq:n9v54-cluster-expansion} once.  Only clusters
meeting the support of $F$ remain.  Sum them by first choosing one
support site.
\end{proof}

\begin{remark}[Comparison with the V53 denominator bound]
\label{rem:n9v54-denominator}
The V53 estimate is the correct black-box bound when no phase locality
is known.  The connected expansion is stronger information: it
organizes the cancellation of all phase clusters disjoint from the
insertions before taking absolute values.
\end{remark}

\subsection{Einstein source density from phase clusters}

Let $h$ be a compactly supported metric variation with support $A_h$.
Suppose each cluster coefficient is differentiable in the metric and
\begin{equation}
 |D_gw_X(0)[h]|
 \le d_X\|h\|_{\mathcal H},
 \qquad
 \sup_i\sum_{X\ni i}e^{\mu\operatorname{diam}X}d_X
 \le C_g.
 \label{eq:n9v54-metric-cluster}
\end{equation}

\begin{theorem}[Local phase-source bound]
\label{thm:n9v54-phase-source}
Under the convergent cluster expansion and
\eqref{eq:n9v54-metric-cluster},
\begin{equation}
 |D_gW_\Lambda(0)[h]|
 \le C_g|A_h|\|h\|_{\mathcal H}
 \label{eq:n9v54-phase-source}
\end{equation}
uniformly in total volume.  If the local cluster derivatives converge
cofinally, they define a local weak Einstein phase-source functional.
\end{theorem}

\begin{proof}
Differentiate \eqref{eq:n9v54-cluster-expansion} in the metric.  Locality
of the variation makes the derivative vanish for clusters disjoint from
$A_h$.  Sum the remaining bounds by a site of $A_h$ as in
\Cref{prop:n9v54-local-expectation}.  Cofinal convergence on a dense
local test class defines the asserted weak functional.
\end{proof}

\begin{firewall}[A weak phase source is not an $L^2$ stress tensor]
\label{fire:n9v54-source-regularity}
The bound \eqref{eq:n9v54-phase-source} gives a local continuous
functional in the declared test norm.  Identifying an $L^2$ tensor
density requires the corresponding Riesz/duality hypothesis and the
V23--V27 Einstein landing argument.  Cluster locality alone does not
raise regularity.
\end{firewall}

\subsection{Reflection positivity remains separate}

\begin{proposition}[Nonvanishing and clustering do not imply positivity]
\label{prop:n9v54-no-positivity}
There are normalized complex product functionals satisfying exact
factorization, nonzero local phase averages, and zero connected
correlations at nonzero separation, yet failing positivity on
$F^*F$.
\end{proposition}

\begin{proof}
On a one-site two-point space with equal positive reference weights, take
phase values $1$ and $i$.  Its normalized complex functional assigns
nonreal values to some nonnegative indicator functions, so it is not
positive.  Tensor products have exact factorization and hence zero
separated connected correlations, while retaining the one-site failure.
\end{proof}

\begin{corollary}[OS row cannot be inferred from the phase cluster]
\label{cor:n9v54-OS-separate}
Absolute complex cluster convergence and exponential connected decay do
not imply Osterwalder--Schrader reflection positivity.  A paired positive
representation or a separate reflection-positive fermion construction
is still required.
\end{corollary}

\begin{proof}
Apply \Cref{prop:n9v54-no-positivity}; ordinary positivity is already a
necessary part of the OS inner-product construction.
\end{proof}

\subsection{Quasi-local phase cluster card}

\begin{definition}[Quasi-local phase cluster card (QLPCC)]
\label{def:n9v54-QLPCC}
QLPCC requires:
\begin{enumerate}
\item finite-volume nonvanishing in a uniform source polydisc;
\item a locally uniformly absolutely convergent connected expansion
\eqref{eq:n9v54-cluster-expansion};
\item the one- and two-source norms
\eqref{eq:n9v54-connected-norm} and
\eqref{eq:n9v54-one-source-norm};
\item the metric derivative norm
\eqref{eq:n9v54-metric-cluster} and determinant-zero control;
\item compatibility with the modulus-measure Higgs/Yukawa component
expansion; and
\item a separate positive/reflection-positive physical reconstruction.
\end{enumerate}
\end{definition}

\begin{theorem}[QLPCC transfers local phase observables]
\label{thm:n9v54-QLPCC}
QLPCC gives volume-uniform local physical expectations, exponential
connected correlation decay, and a local weak phase contribution to the
Einstein source, even when the full-volume average phase is exponentially
small.  It does not by itself give a positive Hilbert-space
reconstruction.
\end{theorem}

\begin{proof}
Use \Cref{prop:n9v54-local-expectation,
thm:n9v54-cluster-decay,thm:n9v54-phase-source,
cor:n9v54-OS-separate}.
\end{proof}

\begin{assessment}[Progress at ninth-series V54]
\label{ass:n9v54-status}
The phase problem has been localized.  The inverse global average phase
is necessary for black-box reweighting but can cancel from fixed-support
observables.  An absolutely convergent connected expansion controls only
clusters meeting the insertions, yielding volume-uniform local
expectations, exponential complex clustering, and a local Einstein phase
source.  Reflection positivity remains a distinct unsolved row.
\end{assessment}

\begin{programme}[Ninth-series V55: compulsory companion audit]
\label{prog:n9v54-V55}
V55 will inspect the newest Yang--Mills and Navier--Stokes notes, test for
a type-compatible connected phase/nonvanishing or averaged-influence
input, and continue to V56 after the audit.
\end{programme}

\begin{firewall}[Claim boundary after ninth-series V54]
\label{fire:n9v54-boundary}
V54 proves consequences of QLPCC but does not construct its polymer
expansion for a chiral Standard Model regulator, prove nonvanishing or
reflection positivity, establish tensor stress regularity, reconstruct
the OS theory, or construct the full continuum.
\end{firewall}

\begin{theorem}[Ninth-series V54 result]
\label{thm:n9v54-result}
Exterior phase factors cancel exactly for product-local observables.
More generally, a nonvanishing absolutely convergent connected phase
expansion with norm \eqref{eq:n9v54-connected-norm} yields the exponential
physical covariance bound \eqref{eq:n9v54-cluster-decay} and the local
Einstein phase-source bound \eqref{eq:n9v54-phase-source}, independently
of the full-volume average-phase magnitude.
\end{theorem}

\begin{proof}
Use \Cref{thm:n9v54-product-cancellation,
thm:n9v54-cluster-decay,prop:n9v54-local-expectation,
thm:n9v54-phase-source,thm:n9v54-QLPCC}.
\end{proof}

\paragraph{Parent-version record.}
V54 was formed from
\texttt{Renewal\_SM\_Einstein\_Continuum\_Research\_Notes\_V53.tex}.
The V53 parent had $19\,556$ logical source lines, $704\,029$ bytes, and
SHA--256
\[
 \texttt{AF5CA23656C9B51107A2C201E88F29C8A17551ABD488A7E22AA1D902CEEA8967}.
\]

% =============================================================================
% END APPENDED NINTH-SERIES V54 MATERIAL
% =============================================================================
% =============================================================================
% BEGIN APPENDED NINTH-SERIES V55 MATERIAL
% =============================================================================

\clearpage
\section{Ninth-series V55: compulsory audit and the slow-carrier
covariance ledger}
\label{sec:v9v55}

\subsection{Compulsory companion audit}

The newest Yang--Mills files inspected were
\begin{align*}
 &\texttt{Renewal\_Yang\_Mills\_Mass\_Gap\_Research\_Notes\_V86.tex},\\
 &\texttt{Renewal\_Yang\_Mills\_Mass\_Gap\_Research\_Notes\_V87.tex}.
\end{align*}
They were byte-identical, so V87 was not counted as independent
evidence.  Each had $38\,688$ logical source lines, $1\,382\,380$ bytes,
and SHA--256
\[
 \texttt{B5D4170727DBCB0AF75A047303373EC0A19D1F7F8A0CC2CC39C884AC1B8D3A70}.
\]
The newest Navier--Stokes file remained
\texttt{Renewal\_NS\_Stability\_Research\_Notes\_V26.tex}, with
$5\,319$ lines, $278\,283$ bytes, and SHA--256
\[
 \texttt{82C887F8C2C69E4F28FAD7C3DC8DE5B9099CB5A4BF431EBA1FFF3E91935978AD}.
\]

Yang--Mills V85 proves a selected-path rare-defect criterion with a
regular-transmission entropy threshold.  This is structurally consistent
with V51--V52 and adds no new Higgs constant.  V86 proves a bridge
resolvent for an already clustering regular kernel, then isolates the
slow conditional-mean covariance.  It also gives a counterexample showing
that small normalized local $L^2$ density defects can coexist with a
nondecaying global carrier.  The Navier--Stokes note again supplies no
new type-compatible input.

The audit changes the SM--Einstein target.  Exponential decay may be
appropriate for a massive conditional matter sector, but the complete
theory contains slow metric and possibly massless gauge carriers.  The
correct theorem must state which sector is expected to be gapped and
must retain the slow covariance instead of calling it an error.

\subsection{Exact slow--fast covariance split}

Let $(\Omega,\mu)$ be a positive probability space and let
$\mathcal S$ be the sigma algebra generated by the declared slow fields:
the metric, retained gauge/constraint variables, zero modes, and any
coarse collective coordinates.

\begin{theorem}[Law of total covariance]
\label{thm:n9v55-total-covariance}
For $F,G\in L^2(\mu)$,
\begin{align}
 \operatorname{Cov}_\mu(F,G)
 &=
 \mathbb E_\mu[
  \operatorname{Cov}_\mu(F,G\mid\mathcal S)]
 \notag\\
 &\quad+
 \operatorname{Cov}_\mu(
  \mathbb E_\mu[F\mid\mathcal S],
  \mathbb E_\mu[G\mid\mathcal S]).
 \label{eq:n9v55-total-covariance}
\end{align}
\end{theorem}

\begin{proof}
Write $F=F_0+\mathbb E[F\mid\mathcal S]$ and similarly for $G$.
The conditional means of $F_0,G_0$ vanish, so the cross terms vanish by
the tower property.  The two remaining terms are exactly those in
\eqref{eq:n9v55-total-covariance}.
\end{proof}

\begin{corollary}[Conditional decay does not imply full decay]
\label{cor:n9v55-fast-only}
If $F,G$ are $\mathcal S$-measurable, the first term in
\eqref{eq:n9v55-total-covariance} is zero and the full covariance is the
slow covariance.  Hence no estimate of conditional fast fluctuations
alone can control such observables.
\end{corollary}

\begin{proof}
Conditioning fixes $F$ and $G$, so their conditional covariance
vanishes.
\end{proof}

\subsection{A sector-aware decay theorem}

Let $r=d(A,B)$ be the separation of the supports of bounded local
observables $F,G$.

\begin{theorem}[Decay inheritance]
\label{thm:n9v55-decay-inheritance}
Suppose
\begin{align}
 \left|
 \mathbb E[
  \operatorname{Cov}(F,G\mid\mathcal S)]
 \right|
 &\le C_{\rm fast}(F,G)e^{-m_{\rm fast}r},
 \label{eq:n9v55-fast-decay}\\
 \left|
 \operatorname{Cov}(
  \mathbb E[F\mid\mathcal S],
  \mathbb E[G\mid\mathcal S])
 \right|
 &\le C_{\rm slow}(F,G)D_{\rm slow}(r).
 \label{eq:n9v55-slow-decay}
\end{align}
Then
\begin{equation}
 |\operatorname{Cov}(F,G)|
 \le C_{\rm fast}e^{-m_{\rm fast}r}
    +C_{\rm slow}D_{\rm slow}(r).
 \label{eq:n9v55-combined-decay}
\end{equation}
If $D_{\rm slow}(r)=(1+r)^{-\alpha}$, the full estimate is polynomial at
large distance unless the slow conditional-mean covariance vanishes or
has faster decay.
\end{theorem}

\begin{proof}
Take absolute values in \eqref{eq:n9v55-total-covariance} and insert the
two hypotheses.  A nonzero polynomial term eventually dominates an
exponential term.
\end{proof}

\begin{corollary}[Centered fast observables]
\label{cor:n9v55-centered}
If
\begin{equation}
 \mathbb E[F\mid\mathcal S]=0
 \quad\text{or}\quad
 \mathbb E[G\mid\mathcal S]=0,
 \label{eq:n9v55-centered}
\end{equation}
then the slow term vanishes and the conditional fast decay is the full
decay estimate.
\end{corollary}

\begin{proof}
One argument of the second covariance in
\eqref{eq:n9v55-total-covariance} is zero.
\end{proof}

\begin{remark}[Conditional centering changes the observable]
\label{rem:n9v55-centering}
Replacing $F$ by $F-\mathbb E[F\mid\mathcal S]$ isolates its fast
fluctuation, but it is generally a nonlocal functional of the slow
field.  Locality and reflection support of the centered observable must
be proved before using it in an OS argument.
\end{remark}

\subsection{A massless slow carrier forbids a global exponential target}

\begin{proposition}[Explicit massless Gaussian diagnostic]
\label{prop:n9v55-massless}
In Euclidean dimension $d>2$, the translation-invariant massless scalar
Gaussian covariance is
\begin{equation}
 C_0(x)=\frac{\Gamma(d/2-1)}{4\pi^{d/2}}|x|^{2-d}
 \qquad(x\ne0).
 \label{eq:n9v55-massless-covariance}
\end{equation}
It obeys no bound
$C_0(x)\le Ce^{-m|x|}$ with fixed $C<\infty$ and $m>0$ for all large
$|x|$.
\end{proposition}

\begin{proof}
The displayed kernel is the fundamental solution of $-\Delta$ in
$d>2$.  If the exponential bound held, multiplying by
$e^{m|x|}|x|^{2-d}$ would remain bounded; instead it tends to infinity.
\end{proof}

\begin{corollary}[No universal mass-gap requirement for the full
SM--Einstein field algebra]
\label{cor:n9v55-no-global-gap}
A continuum target retaining a physical massless carrier cannot demand
uniform exponential decay for every observable having a nonzero
projection onto that carrier.  Massive conditional sectors and massless
slow sectors require separate spectral/decay statements.
\end{corollary}

\begin{proof}
Any observable whose slow conditional mean contains the linear Gaussian
field of \Cref{prop:n9v55-massless} inherits its polynomial covariance
through \eqref{eq:n9v55-total-covariance}.
\end{proof}

\begin{firewall}[Gauge-fixed components versus physical observables]
\label{fire:n9v55-gauge}
The scalar diagnostic illustrates the infrared type; it does not prove
that a particular gauge-invariant graviton or photon observable has
exactly \eqref{eq:n9v55-massless-covariance}.  The physical programme
must identify its gauge/constraint-reduced observables and their actual
infrared powers.
\end{firewall}

\subsection{Local-limit closure}

\begin{lemma}[Uniform decay passes to a local limit]
\label{lem:n9v55-limit}
Let $\mu_n$ converge locally to $\mu$.  If for fixed bounded local
$F,G$ and every translate $\tau_xG$,
\begin{equation}
 |\operatorname{Cov}_{\mu_n}(F,\tau_xG)|
 \le C_{F,G}D(x)
 \label{eq:n9v55-uniform-limit-bound}
\end{equation}
with a regulator-independent right side, then the same bound holds for
$\mu$.
\end{lemma}

\begin{proof}
For fixed $x$, local convergence passes the three bounded local
expectations defining covariance to the limit.  Take $n\to\infty$ in
\eqref{eq:n9v55-uniform-limit-bound}.
\end{proof}

\begin{corollary}[Massless-limit obstruction]
\label{cor:n9v55-massless-limit}
If a locally convergent subsequence has a limiting observable covariance
with no positive exponential rate, then the approximating covariances
cannot have a common positive exponential rate and common local
prefactor.
\end{corollary}

\begin{proof}
Take the contrapositive of \Cref{lem:n9v55-limit}.
\end{proof}

\subsection{Positivity scope}

\begin{firewall}[Complex phase functionals]
\label{fire:n9v55-complex}
Conditional expectation and the Hilbert-space proof of
\eqref{eq:n9v55-total-covariance} require a positive probability law.
They apply directly to a conjugate-paired fermion regulator or the
positive phase-quenched law.  For a genuinely complex chiral functional,
V54's connected generating-functional expansion is the appropriate
replacement; one cannot condition a complex weight as a probability.
\end{firewall}

\subsection{Sector-aware continuum correlation card}

\begin{definition}[Sector-aware continuum correlation card (SACCC)]
\label{def:n9v55-SACCC}
SACCC requires:
\begin{enumerate}
\item a positive physical measure or a QLPCC complex connected
functional with the corresponding decomposition theorem;
\item a declared slow sigma algebra and exact conditional marginal;
\item uniform decay for conditional fast fluctuations;
\item a separate infrared bound for the covariance of slow conditional
means, with exponential decay required only for sectors claimed massive;
\item cofinal local convergence preserving both bounds; and
\item reflection positivity and spectral reconstruction stated
sectorwise, including any massless subspace.
\end{enumerate}
\end{definition}

\begin{theorem}[SACCC gives the complete covariance ledger]
\label{thm:n9v55-SACCC}
Under SACCC, every declared local covariance is bounded by the sum of its
fast and slow estimates as in \eqref{eq:n9v55-combined-decay}, and the
same sectorwise bounds pass to the local continuum limit.  No artificial
global mass gap is inferred across a retained massless carrier.
\end{theorem}

\begin{proof}
Use \Cref{thm:n9v55-decay-inheritance,lem:n9v55-limit}; item 1 selects
the positive conditional proof or the V54 connected-functional
replacement.
\end{proof}

\begin{assessment}[Progress at ninth-series V55]
\label{ass:n9v55-status}
The compulsory audit has corrected the terminal correlation target.
Rare-defect and phase-cluster estimates may control fast conditional
matter fluctuations, but the metric/gauge slow conditional means remain
as an exact covariance term.  A coupled SM--Einstein continuum should
state massive and massless sectors separately rather than force all
observables into an exponential clustering ball.
\end{assessment}

\begin{programme}[Ninth-series V56: effective slow Einstein action]
\label{prog:n9v55-V56}
V56 will integrate the fast normalized sector at finite regulator,
derive the first and second metric variations of the exact slow
effective action, and identify the induced stress expectation and stress
covariance rows.  It will keep the fermion phase and reference-measure
responses explicit.
\end{programme}

\begin{firewall}[Claim boundary after ninth-series V55]
\label{fire:n9v55-boundary}
V55 proves the sector-aware covariance architecture.  It does not
construct the physical slow marginal, establish its infrared decay,
prove SACCC, reflection positivity, OS reconstruction, or the full
SM--Einstein continuum.
\end{firewall}

\begin{theorem}[Ninth-series V55 result]
\label{thm:n9v55-result}
The full covariance is exactly the sum of conditional fast covariance
and slow conditional-mean covariance.  Exponential fast decay therefore
combines with, rather than replaces, the slow infrared law.  SACCC
preserves this sectorwise ledger under local continuum limits and avoids
a false global mass-gap target.
\end{theorem}

\begin{proof}
Use \Cref{thm:n9v55-total-covariance,
thm:n9v55-decay-inheritance,prop:n9v55-massless,
lem:n9v55-limit,thm:n9v55-SACCC}.
\end{proof}

\paragraph{Parent-version record.}
V55 was formed from
\texttt{Renewal\_SM\_Einstein\_Continuum\_Research\_Notes\_V54.tex}.
The V54 parent had $19\,954$ logical source lines, $718\,067$ bytes, and
SHA--256
\[
 \texttt{C4632DA3BC28D343357814E74CB96F1D9C84BD6244ECDB24A8049B84C96521D2}.
\]

% =============================================================================
% END APPENDED NINTH-SERIES V55 MATERIAL
% =============================================================================
% =============================================================================
% BEGIN APPENDED NINTH-SERIES V56 MATERIAL
% =============================================================================

\clearpage
\section{Ninth-series V56: the exact fast-integrated Einstein action
and its stress covariance}
\label{sec:v9v56}

\subsection{Finite-regulator slow/fast disintegration}

Let $s$ denote the finite-dimensional slow variables, including the
regulated metric and retained gauge/constraint coordinates, and let $y$
denote the fast Standard Model variables.  Write the positive fast
weight as
\begin{equation}
 w_f(s,y)
 :=e^{-S_f(s,y)/\hbar}r_f(s,y),
 \qquad
 Z_f(s):=\int w_f(s,y)\,m(dy).
 \label{eq:n9v56-fast-weight}
\end{equation}
Assume $0<Z_f(s)<\infty$ and define
\begin{equation}
 d\nu_s(y):=Z_f(s)^{-1}w_f(s,y)m(dy).
 \label{eq:n9v56-fast-conditional}
\end{equation}
The exact slow effective action is
\begin{equation}
 S_{\rm eff}(s)
 :=S_{\rm slow}(s)-\hbar\log Z_f(s).
 \label{eq:n9v56-Seff}
\end{equation}
This definition includes all normalized determinant, coarea,
gauge-fixing, chart, counterterm, and vacuum-energy factors owned by the
fast integral.

For a constant slow direction $h$, define the fast logarithmic response
\begin{equation}
 A_h(s,y):=D_s\log w_f(s,y)[h]
 =-\hbar^{-1}D_sS_f[h]+D_s\log r_f[h].
 \label{eq:n9v56-Ah}
\end{equation}

\subsection{Differentiation theorem}

\begin{definition}[Fast differentiation domain]
\label{def:n9v56-differentiation}
On a slow chart $K$, the fast weight has a second-order differentiation
domain when $w_f$ is twice directionally differentiable and, for every
pair of directions in the declared finite test family, there are
$m$-integrable functions $G_0,G_1,G_2$ such that uniformly for
$s\in K$,
\begin{equation}
 |w_f|\le G_0,\qquad
 |D_sw_f[h]|\le G_1,\qquad
 |D_s^2w_f[h,k]|\le G_2,
 \label{eq:n9v56-dominators}
\end{equation}
and $\inf_{s\in K}Z_f(s)>0$.
\end{definition}

\begin{lemma}[Derivative of a normalized fast expectation]
\label{lem:n9v56-expectation-derivative}
Under the differentiation domain, for a differentiable insertion $F$,
\begin{equation}
 D_s\mathbb E_{\nu_s}F[k]
 =
 \mathbb E_{\nu_s}[D_sF[k]]
 +\operatorname{Cov}_{\nu_s}(F,A_k).
 \label{eq:n9v56-expectation-derivative}
\end{equation}
\end{lemma}

\begin{proof}
Differentiate the quotient
$Z_f^{-1}\int Fw_f\,dm$ under the integral.  Since
$D_kw_f=w_fA_k$ and
$D_kZ_f=Z_f\mathbb EA_k$, the quotient rule gives the covariance term.
\end{proof}

\begin{theorem}[Exact first and second slow variations]
\label{thm:n9v56-variations}
Under the second-order differentiation domain,
\begin{align}
 D_s\log Z_f[h]
 &=\mathbb E_{\nu_s}A_h,
 \label{eq:n9v56-DlogZ}\\
 D_sS_{\rm eff}[h]
 &=D_sS_{\rm slow}[h]
   -\hbar\mathbb E_{\nu_s}A_h,
 \label{eq:n9v56-first}\\
 D_s^2S_{\rm eff}[h,k]
 &=D_s^2S_{\rm slow}[h,k]
   -\hbar\mathbb E_{\nu_s}[D_kA_h]
   -\hbar\operatorname{Cov}_{\nu_s}(A_h,A_k).
 \label{eq:n9v56-second}
\end{align}
Equivalently, using \eqref{eq:n9v56-Ah},
\begin{align}
 D_sS_{\rm eff}[h]
 &=D_sS_{\rm slow}[h]
   +\mathbb E D_sS_f[h]
   -\hbar\mathbb E D_s\log r_f[h],
 \label{eq:n9v56-first-expanded}\\
 D_s^2S_{\rm eff}[h,k]
 &=D_s^2S_{\rm slow}[h,k]
   +\mathbb E D_s^2S_f[h,k]
   -\hbar\mathbb E D_s^2\log r_f[h,k]
 \notag\\
 &\quad-\hbar\operatorname{Cov}(A_h,A_k).
 \label{eq:n9v56-second-expanded}
\end{align}
\end{theorem}

\begin{proof}
The first identity follows by differentiating
$Z_f=\int w_fdm$.  Insert it into \eqref{eq:n9v56-Seff} for the first
variation.  Differentiate once more and apply
\Cref{lem:n9v56-expectation-derivative} with $F=A_h$.  Expanding $A_h$
and $D_kA_h$ gives the final two formulas.
\end{proof}

\begin{corollary}[Induced covariance is a negative Hessian row]
\label{cor:n9v56-negative-covariance}
For real directions and a positive fast law,
\begin{equation}
 \mathcal C_s(h,k):=
 \operatorname{Cov}_{\nu_s}(A_h,A_k)
 \label{eq:n9v56-C}
\end{equation}
is positive semidefinite.  Therefore the induced row
$-\hbar\mathcal C_s$ in \eqref{eq:n9v56-second} is negative semidefinite.
It must be included in the residual negative part of the V42 absorption
ledger unless it cancels another explicitly computed row.
\end{corollary}

\begin{proof}
Linearity of $A_h$ in $h$ gives
$\mathcal C_s(h,h)=
\mathbb E|A_h-\mathbb EA_h|^2\ge0$.  Multiply by $-\hbar$.
\end{proof}

\begin{firewall}[Mean Hessian is not the Hessian of the effective action]
\label{fire:n9v56-mean-Hessian}
The first three terms in \eqref{eq:n9v56-second-expanded} are the bare
slow Hessian plus mean fast Hessian rows.  Omitting the covariance term
does not differentiate the normalized fast integral correctly.  The
omission is precisely the same error as replacing a Schur complement by
its upper-left block.
\end{firewall}

\subsection{Gaussian Schur check}

\begin{proposition}[Exact quadratic model]
\label{prop:n9v56-Gaussian}
Let
\begin{equation}
 S(s,y)=\frac12s^TAs+s^TBy+\frac12y^TCy,
 \qquad C>0,
 \label{eq:n9v56-quadratic}
\end{equation}
with Lebesgue reference density.  Integrating $y$ gives, up to an
$s$-independent constant,
\begin{equation}
 S_{\rm eff}(s)
 =\frac12s^T(A-BC^{-1}B^T)s.
 \label{eq:n9v56-Gaussian-Seff}
\end{equation}
The covariance row in \eqref{eq:n9v56-second} is exactly
$-BC^{-1}B^T$.
\end{proposition}

\begin{proof}
Complete the square:
\[
 \frac12(y+C^{-1}B^Ts)^TC(y+C^{-1}B^Ts)
 +\frac12s^T(A-BC^{-1}B^T)s.
\]
The conditional covariance of $y$ is $\hbar C^{-1}$.  Its logarithmic
response is
$A_h=-\hbar^{-1}h^TBy$ after assigning the pure slow term to
$S_{\rm slow}$.  Hence
\[
 -\hbar\operatorname{Cov}(A_h,A_k)
 =-h^TBC^{-1}B^Tk,
\]
which is the Schur row.
\end{proof}

\begin{corollary}[Quantum covariance and operator Schur complements have
the same sign]
\label{cor:n9v56-Schur-sign}
The V44 label elimination and the exact fast Gaussian integration both
subtract a positive response covariance from the retained block.  A
large fast gap can make this row small, but positivity of the fast block
alone does not make it vanish.
\end{corollary}

\begin{proof}
Compare \eqref{eq:n9v44-formal-Schur} with
\eqref{eq:n9v56-Gaussian-Seff}.
\end{proof}

\subsection{Locality of the induced Einstein Hessian}

Suppose slow variations decompose locally,
\begin{equation}
 A_h=\sum_x h_x\mathcal O_x,
 \label{eq:n9v56-local-response}
\end{equation}
where $h_x$ is supported at block $x$.  Assume the conditional fast
covariance satisfies
\begin{equation}
 |\operatorname{Cov}_{\nu_s}(\mathcal O_x,\mathcal O_y)|
 \le C_{\mathcal O}e^{-m_fd(x,y)}
 \label{eq:n9v56-fast-covariance}
\end{equation}
uniformly in $s$ and the regulator.

\begin{theorem}[Quasi-local induced Hessian]
\label{thm:n9v56-quasilocal}
The covariance part of the effective Hessian has kernel
\begin{equation}
 K^{\rm cov}_{xy}
 =-\hbar\operatorname{Cov}_{\nu_s}(\mathcal O_x,\mathcal O_y)
 \label{eq:n9v56-cov-kernel}
\end{equation}
and obeys
\begin{equation}
 |K^{\rm cov}_{xy}|
 \le\hbar C_{\mathcal O}e^{-m_fd(x,y)}.
 \label{eq:n9v56-cov-decay}
\end{equation}
If distance spheres grow at most like $C_de^{\eta r}$ with
$m_f>\eta+\mu$, then its weighted row norm is finite:
\begin{equation}
 \sup_x\sum_y e^{\mu d(x,y)}|K^{\rm cov}_{xy}|
 \le
 \frac{\hbar C_{\mathcal O}C_d}
 {1-e^{-(m_f-\eta-\mu)}}.
 \label{eq:n9v56-cov-row}
\end{equation}
\end{theorem}

\begin{proof}
Insert \eqref{eq:n9v56-local-response} in
\eqref{eq:n9v56-C}; the bilinear coefficient is
\eqref{eq:n9v56-cov-kernel}.  Use
\eqref{eq:n9v56-fast-covariance} and sum over distance spheres for the
weighted bound.
\end{proof}

\begin{firewall}[Fast locality does not give a strict relative margin]
\label{fire:n9v56-locality-margin}
The finite weighted row norm in \eqref{eq:n9v56-cov-row} is a locality
statement.  CPRAC/URRFC additionally needs the negative covariance form
to be strictly smaller than the enlarged positive slow reference.  That
comparison requires the actual slow kinetic/constraint energy and is
not implied by exponential off-diagonal decay.
\end{firewall}

\subsection{Gauge and diffeomorphism covariance}

Let a finite-dimensional group element $\gamma$ act on slow and fast
variables by $(s,y)\mapsto(\gamma s,\gamma y)$.

\begin{proposition}[Effective-action invariance]
\label{prop:n9v56-invariance}
Suppose $S_{\rm slow}(\gamma s)=S_{\rm slow}(s)$ and the fast weighted
measure is equivariant:
\begin{equation}
 w_f(\gamma s,\gamma y)m(d(\gamma y))
 =w_f(s,y)m(dy).
 \label{eq:n9v56-equivariance}
\end{equation}
Then
\begin{equation}
 Z_f(\gamma s)=Z_f(s),
 \qquad
 S_{\rm eff}(\gamma s)=S_{\rm eff}(s).
 \label{eq:n9v56-Seff-invariance}
\end{equation}
Every infinitesimal gauge/diffeomorphism orbit direction therefore
annihilates $S_{\rm eff}$, subject to the finite-regulator action being
defined.
\end{proposition}

\begin{proof}
Change variables $y\mapsto\gamma y$ in the fast integral and use
\eqref{eq:n9v56-equivariance}.  Add slow invariance and differentiate the
resulting identity along a one-parameter group.
\end{proof}

\begin{firewall}[Anomalous Jacobian]
\label{fire:n9v56-anomaly}
If the fast change of variables has a determinant phase or Jacobian,
\eqref{eq:n9v56-equivariance} fails by that cocycle.  Perturbative anomaly
cancellation must be realized by the regulator measure itself before
\Cref{prop:n9v56-invariance} applies.  Configurationwise positivity and
reflection positivity remain separate, as in V53--V54.
\end{firewall}

\subsection{Complex fermion phase}

When the fast fermion factor is complex, write
\begin{equation}
 Z_f=Z_{\rm abs}\zeta,
 \qquad
 S_{\rm eff}
 =S_{\rm slow}-\hbar\log Z_{\rm abs}
              -\hbar\log\zeta
 \label{eq:n9v56-complex-Seff}
\end{equation}
on a nonvanishing branch.

\begin{proposition}[Phase row in the effective Einstein action]
\label{prop:n9v56-phase-row}
The first phase contribution is
\begin{equation}
 D_sS_{\rm eff}^{\rm phase}[h]
 =-\hbar D_s\log\zeta[h],
 \label{eq:n9v56-phase-first}
\end{equation}
with $D_s\log\zeta$ given by
\eqref{eq:n9v53-phase-response}.  Its second variation is
$-\hbar D_s^2\log\zeta$ whenever the QLPCC metric cluster can be
differentiated twice.  Neither row is a covariance of a positive
physical probability unless conjugate pairing makes the weight
positive.
\end{proposition}

\begin{proof}
Differentiate \eqref{eq:n9v56-complex-Seff}.  The explicit first
derivative is V53.  A second differentiation is valid under the stated
analytic QLPCC hypothesis.  The final statement is
\Cref{prop:n9v53-not-probability}.
\end{proof}

\subsection{Vacuum and ground-energy rows}

\begin{proposition}[Fast normalization owns local vacuum energies]
\label{prop:n9v56-vacuum}
If a local fast block Hamiltonian is shifted by a metric-dependent scalar
$E_B(s)$, then
\begin{equation}
 Z_f(s)\mapsto e^{E_B(s)/\hbar}Z_f(s)
 \quad\Longrightarrow\quad
 S_{\rm eff}(s)\mapsto S_{\rm eff}(s)-E_B(s).
 \label{eq:n9v56-ground-shift}
\end{equation}
Thus the V47--V48 Higgs vacuum and ground-energy normalizations reappear
exactly in the slow Einstein action and its metric derivative.
\end{proposition}

\begin{proof}
Insert the scalar factor in \eqref{eq:n9v56-Seff} and take its logarithm.
\end{proof}

\subsection{Effective slow Einstein action card}

\begin{definition}[Effective slow Einstein action card (ESEAC)]
\label{def:n9v56-ESEAC}
ESEAC requires:
\begin{enumerate}
\item a common slow chart and exact normalized fast disintegration;
\item a uniform positive partition floor and the second-order
differentiation domain;
\item the complete logarithmic response $A_h$, including reference,
coarea, determinant, chart, vacuum, and counterterm rows;
\item a paired positive fermion law or the nonvanishing analytic QLPCC
phase branch;
\item quasi-local first and second response kernels, including the stress
covariance \eqref{eq:n9v56-cov-kernel};
\item a strict CPRAC relative bound for the induced negative rows;
\item regulator-level gauge/diffeomorphism covariance without an
unowned anomaly cocycle; and
\item convergence of the first variation in the V23--V27 weak tensor
source topology.
\end{enumerate}
\end{definition}

\begin{theorem}[ESEAC loads the fast-to-Einstein handoff]
\label{thm:n9v56-ESEAC}
Under ESEAC, integrating out the fast Standard Model sector produces a
well-defined twice differentiable finite-regulator slow action with
exact source and Hessian \eqref{eq:n9v56-first}--
\eqref{eq:n9v56-second}.  Its induced covariance is an owned quasi-local
negative form controlled by CPRAC, its phase and vacuum rows remain
explicit, and its first variation lands in the weak Einstein source
topology.
\end{theorem}

\begin{proof}
Use \Cref{thm:n9v56-variations,cor:n9v56-negative-covariance,
thm:n9v56-quasilocal,prop:n9v56-invariance,
prop:n9v56-phase-row,prop:n9v56-vacuum}.  ESEAC items 6 and 8 load the
V42 form resummation and the V23--V27 tensor landing respectively.
\end{proof}

\begin{assessment}[Progress at ninth-series V56]
\label{ass:n9v56-status}
The fast-to-slow Einstein handoff now has an exact finite-regulator
formula.  The mean fast stress appears in the first variation, while the
second variation contains a negative stress-response covariance, the
quantum analogue of a Schur complement.  Conditional fast clustering
makes that row quasi-local but does not by itself make it small.  Fermion
phase, vacuum energy, anomaly, and weak tensor landing remain explicitly
owned.
\end{assessment}

\begin{programme}[Ninth-series V57: induced covariance relative margin]
\label{prog:n9v56-V57}
V57 will compare the quasi-local negative covariance kernel with the
positive constraint-reduced metric reference.  It will derive a
weighted Schur/diagonal-dominance criterion giving the strict CPRAC
margin and isolate the massless slow modes on which no such gap-based
comparison is available.
\end{programme}

\begin{firewall}[Claim boundary after ninth-series V56]
\label{fire:n9v56-boundary}
V56 proves the exact effective-action calculus.  It does not prove
ESEAC for the physical chiral regulator, the strict induced-covariance
margin, slow-sector infrared control, reflection positivity, OS
reconstruction, or the full SM--Einstein continuum.
\end{firewall}

\begin{theorem}[Ninth-series V56 result]
\label{thm:n9v56-result}
The exact fast-integrated action is
$S_{\rm slow}-\hbar\log Z_f$.  Its first variation is the complete mean
fast source \eqref{eq:n9v56-first-expanded}; its Hessian is the mean
complete fast Hessian minus the positive response covariance in
\eqref{eq:n9v56-second-expanded}.  Under ESEAC these rows are quasi-local,
symmetry compatible, and land in the declared weak Einstein source
topology.
\end{theorem}

\begin{proof}
Use \Cref{thm:n9v56-variations,prop:n9v56-Gaussian,
thm:n9v56-quasilocal,thm:n9v56-ESEAC}.
\end{proof}

\paragraph{Parent-version record.}
V56 was formed from
\texttt{Renewal\_SM\_Einstein\_Continuum\_Research\_Notes\_V55.tex}.
The V55 parent had $20\,286$ logical source lines, $730\,064$ bytes, and
SHA--256
\[
 \texttt{166B3F8E29F08DF5A5D8F5F54138CB221E1ED0C29932D35F52D1294056609A2C}.
\]

% =============================================================================
% END APPENDED NINTH-SERIES V56 MATERIAL
% =============================================================================
% =============================================================================
% BEGIN APPENDED NINTH-SERIES V57 MATERIAL
% =============================================================================

\section{Ninth-series V57: induced covariance margin and the massless-mode test}
\label{sec:v9v57}

\subsection{Purpose and finite-regulator setting}

V56 identified the negative covariance row in the Hessian of the
fast-integrated Einstein action.  Quasi-locality controls its spatial
tail, but it does not compare its size with the positive
constraint-reduced metric reference.  This note supplies that comparison
at finite regulator.  It also separates two logically different routes:
a massive reference admits a weighted block Schur test, whereas a
massless reference requires a zero-momentum cancellation and a moment
bound.  The latter cancellation is an additional Ward/counterterm datum;
it is not a consequence of exponential clustering alone.

Let $X$ be the finite set of slow metric blocks, let ${\cal H}_x$ be the
real tangent fibre over $x$, and put
\begin{equation}
 {cal H}=\bigoplus_{x\in X}{\cal H}_x,
 \qquad
 {mathfrak a}(u)=\sum_{x\in X}
       \langle u_x,A_xu_x\rangle,
 \label{eq:n9v57-reference}
\end{equation}
where every $A_x$ is positive definite.  The induced covariance form is
written
\begin{equation}
 {mathfrak c}(u,v)=\sum_{x,y\in X}
       \langle u_x,C_{xy}v_y\rangle,
 \qquad C_{yx}=C_{xy}^{*}.
 \label{eq:n9v57-covariance-blocks}
\end{equation}
For the V56 logarithmic responses this is the positive form
$\operatorname{Cov}_{\nu_s}(A_u,A_v)$.  Thus the effective Hessian
contains $-\hbar {\mathfrak c}$.

\subsection{The normalized block Schur test}

Define
\begin{equation}
 K_{xy}=A_x^{-1/2}C_{xy}A_y^{-1/2},
 \qquad k_{xy}=\lVert K_{xy}\rVert.
 \label{eq:n9v57-normalized-block}
\end{equation}

\begin{lemma}[Weighted block Schur estimate]
\label{lem:n9v57-Schur}
Suppose that positive numbers $w_x$ and constants $R_+,R_-<\infty$
satisfy
\begin{equation}
 \sup_x\frac1{w_x}\sum_y k_{xy}w_y\le R_+,
 \qquad
 \sup_y\frac1{w_y}\sum_x k_{xy}w_x\le R_-.
 \label{eq:n9v57-Schur-hyp}
\end{equation}
Then
\begin{equation}
 |{\mathfrak c}(u,v)|
 \le \sqrt{R_+R_-}\,
      {\mathfrak a}(u)^{1/2}{\mathfrak a}(v)^{1/2}.
 \label{eq:n9v57-Schur-bound}
\end{equation}
If $k_{xy}=k_{yx}$, one may take $R_+=R_-=R$ from either weighted row
bound.
\end{lemma}

\begin{proof}
Put $f_x=\lVert A_x^{1/2}u_x\rVert$ and
$g_y=\lVert A_y^{1/2}v_y\rVert$.  Then
\begin{equation}
 |{\mathfrak c}(u,v)|\le\sum_{x,y}k_{xy}f_xg_y.
 \label{eq:n9v57-majorization}
\end{equation}
Insert $1=(w_y/w_x)^{1/2}(w_x/w_y)^{1/2}$ into each summand and apply
Cauchy--Schwarz over $(x,y)$.  The first squared factor is at most
$R_+\sum_x f_x^2$ and the second at most $R_-\sum_y g_y^2$ by
\eqref{eq:n9v57-Schur-hyp}.  Equation
\eqref{eq:n9v57-reference} gives the result.  Symmetry identifies the two
row constants.
\end{proof}

\begin{theorem}[Strict induced-covariance margin]
\label{thm:n9v57-margin}
Let ${\mathfrak r}$ denote all other negative Hessian rows and suppose
\begin{equation}
 |{\mathfrak r}(u,u)|\le\delta {\mathfrak a}(u),
 \qquad 0\le\delta<1.
 \label{eq:n9v57-other-row}
\end{equation}
Under \Cref{lem:n9v57-Schur}, the full quadratic form
\begin{equation}
 {\mathfrak h}(u)
 ={\mathfrak a}(u)+{\mathfrak r}(u,u)-\hbar {\mathfrak c}(u,u)
 \label{eq:n9v57-full-Hessian}
\end{equation}
satisfies
\begin{equation}
 {\mathfrak h}(u)\ge
 \bigl(1-\delta-\hbar\sqrt{R_+R_-}\bigr){\mathfrak a}(u).
 \label{eq:n9v57-strict-margin}
\end{equation}
Consequently the covariance row has a strict CPRAC relative margin if
\begin{equation}
 \delta+\hbar\sqrt{R_+R_-}<1.
 \label{eq:n9v57-margin-condition}
\end{equation}
\end{theorem}

\begin{proof}
Apply \eqref{eq:n9v57-Schur-bound} with $u=v$, use positivity of
${\mathfrak c}$, and combine the result with
\eqref{eq:n9v57-other-row}.
\end{proof}

\begin{remark}[The normalization carries the comparison]
\label{rem:n9v57-normalization}
An unnormalized decay estimate for $C_{xy}$ is insufficient.  The
quantities entering \eqref{eq:n9v57-margin-condition} are the blocks
$A_x^{-1/2}C_{xy}A_y^{-1/2}$.  Small eigenvalues of the metric reference
therefore amplify the same covariance kernel.
\end{remark}

\subsection{From fast clustering to a computable row constant}

Assume that the local response decomposes as
$A_u=\sum_x L_xu_x$, with $L_xu_x$ centred in
$L^2(\nu_s)$.  Suppose the conditional fast law obeys
\begin{equation}
 |\operatorname{Cov}_{\nu_s}(L_xu_x,L_yv_y)|
 \le b_xb_y e^{-m d(x,y)}
 \lVert A_x^{1/2}u_x\rVert
 \lVert A_y^{1/2}v_y\rVert.
 \label{eq:n9v57-cluster-response}
\end{equation}

\begin{proposition}[Clustering-to-margin compiler]
\label{prop:n9v57-cluster-compiler}
Under \eqref{eq:n9v57-cluster-response},
\begin{equation}
 k_{xy}\le b_xb_y e^{-m d(x,y)}.
 \label{eq:n9v57-k-decay}
\end{equation}
If $b_x\le b$ and the interaction graph has maximum degree $\Delta$,
then, when $(\Delta-1)e^{-m}<1$,
\begin{equation}
 R\le b^2\left(
 1+\frac{\Delta e^{-m}}{1-(\Delta-1)e^{-m}}
 \right).
 \label{eq:n9v57-geometric-row}
\end{equation}
Thus \eqref{eq:n9v57-margin-condition} follows from the same expression
with its right-hand side substituted for $R$.
\end{proposition}

\begin{proof}
Taking suprema over unit $A_x^{1/2}u_x$ and
$A_y^{1/2}v_y$ gives \eqref{eq:n9v57-k-decay}.  At distance $r\ge1$
there are at most $\Delta(\Delta-1)^{r-1}$ vertices.  Summing the
resulting geometric series, including the vertex at $r=0$, gives
\eqref{eq:n9v57-geometric-row}.  Apply
\Cref{lem:n9v57-Schur,thm:n9v57-margin}.
\end{proof}

\begin{firewall}[Decay is not smallness]
\label{fire:n9v57-decay-smallness}
Increasing $m$ suppresses off-diagonal blocks but leaves the on-site
term $b^2$.  Hence arbitrarily fast clustering does not establish the
strict margin when the local normalized response is too large.  The
missing input must be a larger owned positive row, a smaller response,
or a cancellation that changes the infrared order.
\end{firewall}

\subsection{Constraint pullback without a locality shortcut}

Let ${\cal G}$ be the reduced tangent Hilbert space and let
$T:{\cal G}\to{\cal H}$ be the exact linearized constraint/chart map.
The pulled-back forms are
\begin{equation}
 {\mathfrak a}_{\rm red}(q)={\mathfrak a}(Tq),
 \qquad
 {\mathfrak c}_{\rm red}(q,r)={\mathfrak c}(Tq,Tr).
 \label{eq:n9v57-pullback}
\end{equation}

\begin{proposition}[Relative bounds survive exact pullback]
\label{prop:n9v57-pullback}
If $|{\mathfrak c}(u,v)|\le R
{\mathfrak a}(u)^{1/2}{\mathfrak a}(v)^{1/2}$ on
$\operatorname{Ran}T$, then
\begin{equation}
 |{\mathfrak c}_{\rm red}(q,r)|
 \le R {\mathfrak a}_{\rm red}(q)^{1/2}
          {\mathfrak a}_{\rm red}(r)^{1/2}.
 \label{eq:n9v57-pullback-bound}
\end{equation}
No locality property of $T$ is needed for this conclusion.
\end{proposition}

\begin{proof}
Substitute $u=Tq$ and $v=Tr$ in the assumed bound.
\end{proof}

\begin{firewall}[Projection may be nonlocal]
\label{fire:n9v57-nonlocal-projection}
The proposition transports an already proved relative bound.  It does
not transport the block-decay hypotheses used to prove that bound.
Elliptic gauge projections may contain an inverse Laplacian, so applying
the Schur test only after a global projection can silently replace an
exponential kernel by a long-range one.  One must prove the comparison
on $\operatorname{Ran}T$ before projection, or separately establish a
weighted kernel bound for $T$.
\end{firewall}

\subsection{The massless obstruction}

The preceding local normalization used invertible $A_x$.  The continuum
metric reference instead has massless directions.  The following
elementary test is decisive.

\begin{lemma}[Kernel compatibility is necessary]
\label{lem:n9v57-kernel}
Let $A\ge0$ and $C\ge0$ be quadratic forms on a finite-dimensional real
space.  If
\begin{equation}
 0\le C(u,u)\le R\langle u,Au\rangle
 \label{eq:n9v57-semidefinite-bound}
\end{equation}
for every $u$, then $C(z,v)=0$ for every $z\in\ker A$ and every $v$.
\end{lemma}

\begin{proof}
For $z\in\ker A$, \eqref{eq:n9v57-semidefinite-bound} gives
$C(z,z)=0$.  Cauchy--Schwarz for the positive semidefinite form $C$
then yields $|C(z,v)|^2\le C(z,z)C(v,v)=0$.
\end{proof}

\begin{proposition}[Fourier ratio criterion and infrared no-go]
\label{prop:n9v57-Fourier-no-go}
On a translation-invariant finite torus, let $\widehat A(p)$ and
$\widehat C(p)$ be the matrix symbols on the physical fibre.  The least
relative coefficient is
\begin{equation}
 R_*=sup_{p}\ sup_{v:\,\langle v,\widehat A(p)v\rangle>0}
 \frac{\langle v,\widehat C(p)v\rangle}
      {\langle v,\widehat A(p)v\rangle},
 \label{eq:n9v57-Fourier-ratio}
\end{equation}
provided $\widehat C(p)$ annihilates $\ker\widehat A(p)$.  If along
momenta $p_j\to0$ and unit vectors $v_j$ one has
\begin{equation}
 \langle v_j,\widehat A(p_j)v_j\rangle\le a|p_j|^2,
 \qquad
 \langle v_j,\widehat C(p_j)v_j\rangle\ge c>0,
 \label{eq:n9v57-massless-no-go-hyp}
\end{equation}
then $R_*=\infty$ in the continuum/volume limit.
\end{proposition}

\begin{proof}
Discrete Fourier transformation diagonalizes both forms, so the
Rayleigh quotient gives \eqref{eq:n9v57-Fourier-ratio}.  Under
\eqref{eq:n9v57-massless-no-go-hyp} the quotient is at least
$c/(a|p_j|^2)$, which diverges.
\end{proof}

\begin{corollary}[A correlation length does not create a graviton mass]
\label{cor:n9v57-no-fake-gap}
Exponential decay of the fast covariance makes $\widehat C(p)$ analytic
near $p=0$, but it does not force $\widehat C(0)=0$.  Consequently a
finite fast correlation length cannot by itself yield a relative bound
against a massless $p^2$ metric reference.
\end{corollary}

\begin{proof}
Exponential decay gives absolute summability and analyticity of the
Fourier series.  Its value at zero is the generally nonzero sum of the
position-space kernel.  Apply \Cref{prop:n9v57-Fourier-no-go}.
\end{proof}

\subsection{Zero-mode cancellation repairs the comparison}

We now state a sufficient condition that has the correct derivative
order.  Let $C(z)=C(-z)^*$ be a translation-invariant matrix kernel on
$\mathbb Z^d$ whose convolution form is positive semidefinite, with
\begin{equation}
 \sum_z C(z)=0,
 \qquad
 M_2:=\sum_z |z|^2\lVert C(z)\rVert<\infty.
 \label{eq:n9v57-zero-second-moment}
\end{equation}

\begin{theorem}[Cancellation-to-gradient relative bound]
\label{thm:n9v57-cancellation}
Under \eqref{eq:n9v57-zero-second-moment}, the symbol satisfies
\begin{equation}
 \lVert\widehat C(p)\rVert
 \le \frac12 M_2 |p|^2.
 \label{eq:n9v57-symbol-p2}
\end{equation}
If the physical reference obeys
\begin{equation}
 \widehat A(p)\ge a_0 |p|^2 I
 \label{eq:n9v57-A-p2}
\end{equation}
on the reduced fibre, then
\begin{equation}
 |{\mathfrak c}(u,u)|
 \le \frac{M_2}{2a_0}{\mathfrak a}(u).
 \label{eq:n9v57-gradient-relative}
\end{equation}
The induced covariance has a strict margin, including the massless
sector, whenever
\begin{equation}
 \delta+\frac{\hbar M_2}{2a_0}<1.
 \label{eq:n9v57-gradient-margin}
\end{equation}
\end{theorem}

\begin{proof}
Positivity gives $\widehat C(p)\ge0$.  For every fibre vector $v$, the
scalar function
$f_v(p)=\langle v,\widehat C(p)v\rangle$ is nonnegative and, by the
zero-sum identity, satisfies $f_v(0)=0$.  Its first derivative at the
minimum $p=0$ therefore vanishes.  Polarization gives
$D\widehat C(0)=0$.  Consequently Taylor's formula gives
\begin{equation}
 \widehat C(p)=
 -\int_0^1(1-t)\sum_z C(z)(p\cdot z)^2e^{-itp\cdot z}\,dt.
 \label{eq:n9v57-symbol-subtract}
\end{equation}
Taking the operator norm and using
$|p\cdot z|^2\le |p|^2|z|^2$ proves
\eqref{eq:n9v57-symbol-p2}.  Compare it with
\eqref{eq:n9v57-A-p2} mode by mode and sum by Parseval.  The last claim
follows as in \Cref{thm:n9v57-margin}.
\end{proof}

\begin{remark}[Ownership of the cancellation]
\label{rem:n9v57-cancellation-ownership}
The identity $\sum_zC(z)=0$ says that a constant physical metric
variation has zero fast response variance after all vacuum,
cosmological, reference-density, coarea, and counterterm rows are placed
in their declared owners.  A diffeomorphism Ward identity controls gauge
orbit directions, but it need not impose this identity on transverse
physical directions.  V57 therefore records
\eqref{eq:n9v57-zero-second-moment} as a condition to be derived from the
renormalized regulator, not as an automatic symmetry consequence.
\end{remark}

\subsection{An honest infrared split}

Let $P_{<\mu}$ project onto physical momenta $|p|<\mu$ and
$P_{\ge\mu}=I-P_{<\mu}$.

\begin{proposition}[High-mode comparison]
\label{prop:n9v57-high-mode}
Suppose $\widehat A(p)\ge a_0|p|^2I$ and
$\sup_p\lVert\widehat C(p)\rVert\le C_0$.  Then
\begin{equation}
 {\mathfrak c}(P_{\ge\mu}u,P_{\ge\mu}u)
 \le \frac{C_0}{a_0\mu^2}
       {\mathfrak a}(P_{\ge\mu}u).
 \label{eq:n9v57-high-mode-bound}
\end{equation}
No corresponding uniform estimate follows on $P_{<\mu}$ as
$\mu\downarrow0$ without the cancellation of
\Cref{thm:n9v57-cancellation} or a separate infrared construction.
\end{proposition}

\begin{proof}
For $|p|\ge\mu$, one has
$C_0\le C_0|p|^2/\mu^2$.  Compare Fourier symbols and apply Parseval.
The coefficient diverges as $\mu\downarrow0$, proving the final
statement at the level of this estimate; the no-go example in
\Cref{prop:n9v57-Fourier-no-go} shows the divergence can be genuine.
\end{proof}

\begin{definition}[Induced covariance relative-margin card (ICRMC)]
\label{def:n9v57-ICRMC}
ICRMC requires:
\begin{enumerate}
\item the exact complete response covariance from V56;
\item an owned positive constraint-reduced metric reference;
\item a normalized weighted Schur bound on every uniformly coercive
sector;
\item the aggregate strict inequality
$\delta+\hbar R<1$ after all negative rows are counted once;
\item exact pullback to the physical tangent space without assuming a
nonlocal projection preserves kernel decay;
\item on every massless physical sector, either the zero-mode and
second-moment conditions \eqref{eq:n9v57-zero-second-moment} or an
explicit infrared block retained outside the relative-form inversion;
and
\item regulator-uniform constants in the chosen continuum scaling.
\end{enumerate}
\end{definition}

\begin{theorem}[ICRMC loads the covariance row into CPRAC]
\label{thm:n9v57-ICRMC}
Under ICRMC, the negative fast stress covariance is relatively form
bounded by the positive physical metric reference with aggregate
coefficient strictly below one.  Massive/coercive blocks are controlled
by \Cref{thm:n9v57-margin}; massless blocks are controlled by
\Cref{thm:n9v57-cancellation} or are left in the declared infrared
effective theory.  Hence the V56 Hessian inversion does not use a false
mass gap.
\end{theorem}

\begin{proof}
Apply the weighted Schur estimate on the coercive direct summands and the
gradient comparison on the cancelled massless summands.  Orthogonal
summation preserves the maximum relative coefficient.  ICRMC item 4
places the combined coefficient below one.  If an infrared block is
retained, no inversion is asserted there.
\end{proof}

\begin{assessment}[Progress at ninth-series V57]
\label{ass:n9v57-status}
The induced covariance bottleneck has been reduced to two quantitative
tests.  On coercive blocks the test is the weighted row sum of the
reference-normalized covariance kernel.  On massless blocks the test is
the vanishing zero mode plus a uniform second moment, or else an explicit
infrared theory.  This exposes the exact datum that clustering cannot
supply.
\end{assessment}

\begin{programme}[Ninth-series V58: renormalized zero-mode ownership]
\label{prog:n9v57-V58}
V58 will decompose the zero-momentum metric response into vacuum energy,
cosmological, local counterterm, anomaly, and connected stress rows.  It
will determine which subtraction is fixed by the background Einstein
equation, which is fixed by Ward identities, and which remains a genuine
renormalization condition needed for ICRMC.
\end{programme}

\begin{firewall}[Claim boundary after ninth-series V57]
\label{fire:n9v57-boundary}
V57 proves the relative-margin criteria.  It does not prove their
regulator-uniform constants for the physical chiral Standard Model, the
required transverse zero-mode cancellation, the slow infrared theory,
reflection positivity, OS reconstruction, or the full SM--Einstein
continuum.
\end{firewall}

\begin{theorem}[Ninth-series V57 result]
\label{thm:n9v57-result}
The fast-induced negative Hessian row has a strict physical relative
margin if the reference-normalized covariance passes the weighted Schur
test on coercive sectors and the aggregate coefficient satisfies
\eqref{eq:n9v57-margin-condition}.  A massless $p^2$ sector additionally
requires covariance cancellation at zero momentum with a finite second
moment, yielding \eqref{eq:n9v57-gradient-relative}, or it must remain
outside the gap-based inversion as an explicit infrared block.
\end{theorem}

\begin{proof}
Combine \Cref{thm:n9v57-margin,prop:n9v57-pullback,
prop:n9v57-Fourier-no-go,thm:n9v57-cancellation,
thm:n9v57-ICRMC}.
\end{proof}

\paragraph{Parent-version record.}
V57 was formed from
\texttt{Renewal\_SM\_Einstein\_Continuum\_Research\_Notes\_V56.tex}.
The V56 parent had $20\,742$ logical source lines, $745\,315$ bytes, and
SHA--256
\[
 \texttt{82123934C5E9DF581B5C6F6629249673E28852825A69DFA8AB640BEAE7627B6B}.
\]

% =============================================================================
% END APPENDED NINTH-SERIES V57 MATERIAL
% =============================================================================
% =============================================================================
% BEGIN APPENDED NINTH-SERIES V58 MATERIAL
% =============================================================================

\section{Ninth-series V58: renormalized zero-mode ownership}
\label{sec:v9v58}

\subsection{Why the combined row is the physical object}

V57 gave a sufficient massless comparison when the covariance kernel
itself has vanishing zero mode.  That hypothesis is stronger than the
effective Einstein equation requires.  The Hessian seen by the slow
metric is the sum of the mean fast Hessian, the negative response
covariance, vacuum and ground-energy rows, and local counterterms.  Only
this combined row is invariant under a redistribution of local terms
between the bare action and the fast normalization.  V58 makes that
ownership explicit.

At finite regulator write
\begin{equation}
 \Gamma_n(g)=S_{{\rm E},n}(g)+S_{{\rm ct},n}(g)
              -\hbar\log Z_{f,n}(g).
 \label{eq:n9v58-effective-action}
\end{equation}
At a background $\bar g$, V56 gives
\begin{align}
 D^2\Gamma_n(\bar g)[h,k]
 &=D^2S_{{\rm E},n}[h,k]+{\mathfrak j}_n(h,k),
 \label{eq:n9v58-Hessian-split}\\
 {\mathfrak j}_n(h,k)
 &=D^2S_{{\rm ct},n}[h,k]
   -\hbar\,\mathbb E_{\nu_{\bar g}}D A_h[k]
   -\hbar\,\operatorname{Cov}_{\nu_{\bar g}}(A_h,A_k).
 \label{eq:n9v58-combined-row}
\end{align}
Equivalently, after expanding $A_h$ as in V56, the middle term contains
the mean fast action Hessian and every reference-density Hessian.  Metric
dependent block ground energies from V47--V48 are included either in
$S_{{\rm ct},n}$ or in $Z_{f,n}$, but never in both.

\begin{proposition}[Local redistribution invariance]
\label{prop:n9v58-redistribution}
Let $L_n(g)$ be any twice differentiable local functional and perform
the exact redistribution
\begin{equation}
 S_{{\rm ct},n}\mapsto S_{{\rm ct},n}+L_n,
 \qquad
 Z_{f,n}\mapsto e^{L_n/\hbar}Z_{f,n}.
 \label{eq:n9v58-redistribution-map}
\end{equation}
Then $\Gamma_n$, $D\Gamma_n$, $D^2\Gamma_n$, and the combined row
${\mathfrak j}_n$ are unchanged.  The separate mean-Hessian and
counterterm rows are not invariant.
\end{proposition}

\begin{proof}
The two changes in \eqref{eq:n9v58-effective-action} cancel exactly:
$L_n-\hbar\log(e^{L_n/\hbar})=0$.  Differentiate this identity twice.
\end{proof}

\begin{corollary}[Corrected massless ownership]
\label{cor:n9v58-corrected-ownership}
The V57 cancellation-to-gradient theorem may be applied to the total
renormalized induced kernel representing ${\mathfrak j}_n$, rather than
to the positive covariance kernel alone.  A nonzero covariance zero
mode is admissible if it is cancelled by an owned mean-Hessian,
vacuum-energy, or counterterm zero mode in
\eqref{eq:n9v58-combined-row}.
\end{corollary}

\begin{proof}
\Cref{thm:n9v57-cancellation} uses self-adjointness, the zero moment,
and a second-moment bound.  Positivity was used there only to derive the
vanishing first moment from the covariance zero mode.  For a general
combined self-adjoint kernel, impose evenness or the vanishing first
moment explicitly and repeat the Taylor estimate.  Redistribution
invariance identifies ${\mathfrak j}_n$ as the physical kernel.
\end{proof}

\subsection{What the Ward identity actually fixes}

Let $V_X(g)={\cal L}_Xg$ be the infinitesimal diffeomorphism vector field
on metric space.

\begin{proposition}[Off-shell and on-shell Hessian Ward identities]
\label{prop:n9v58-Ward-Hessian}
If $\Gamma_n(\phi^*g)=\Gamma_n(g)$ for regulator diffeomorphisms $\phi$,
then
\begin{align}
 D\Gamma_n(g)[V_X(g)]&=0,
 \label{eq:n9v58-first-Ward}\\
 D^2\Gamma_n(g)[k,V_X(g)]
 +D\Gamma_n(g)[{\cal L}_Xk]&=0.
 \label{eq:n9v58-second-Ward}
\end{align}
If $\bar g$ is stationary, the Hessian annihilates every gauge orbit
direction:
\begin{equation}
 D^2\Gamma_n(\bar g)[k,{\cal L}_X\bar g]=0.
 \label{eq:n9v58-onshell-gauge-null}
\end{equation}
\end{proposition}

\begin{proof}
Differentiate $\Gamma_n(\phi_t^*g)=\Gamma_n(g)$ at $t=0$ to obtain
\eqref{eq:n9v58-first-Ward}.  Differentiate that identity with respect
to $g$ in direction $k$.  Since $D_gV_X[k]={\cal L}_Xk$, this gives
\eqref{eq:n9v58-second-Ward}.  At a stationary background the second
term vanishes.
\end{proof}

\begin{firewall}[Ward identity versus transverse zero mode]
\label{fire:n9v58-Ward-transverse}
Equation \eqref{eq:n9v58-onshell-gauge-null} fixes gauge orbit
directions.  It does not set the zero-momentum Hessian to zero on every
transverse physical metric direction.  That stronger conclusion needs
the classification of local zero-derivative invariants, the background
equation, and control of nonlocal remainders.
\end{firewall}

\begin{firewall}[Anomaly row]
\label{fire:n9v58-anomaly}
If the regulated fermion measure transforms by a cocycle
$\Gamma_n(\phi^*g)-\Gamma_n(g)={\cal A}_n(\phi,g)$, then the right-hand
sides of \eqref{eq:n9v58-first-Ward}--
\eqref{eq:n9v58-second-Ward} acquire derivatives of ${\cal A}_n$.
A cohomologically nontrivial anomaly cannot be repaired by choosing a
cosmological or curvature counterterm.  Regulator-level anomaly
cancellation is therefore prior to the zero-mode subtraction below.
\end{firewall}

\subsection{The unique algebraic metric density}

\begin{lemma}[Zero-derivative invariant density]
\label{lem:n9v58-algebraic-density}
Let $F(g)$ be a continuous local density depending on a positive
definite metric matrix $g$ but on no derivatives, background tensors, or
matter fields.  If
\begin{equation}
 F(M^TgM)=|\det M|F(g)
 \label{eq:n9v58-density-covariance}
\end{equation}
for every invertible linear coordinate change $M$, then
\begin{equation}
 F(g)=\lambda\sqrt{\det g}
 \label{eq:n9v58-only-volume}
\end{equation}
for a constant $\lambda$.
\end{lemma}

\begin{proof}
Choose $M=g^{-1/2}$ in \eqref{eq:n9v58-density-covariance}.  Then
$F(I)=(\det g)^{-1/2}F(g)$, which is
\eqref{eq:n9v58-only-volume} with $\lambda=F(I)$.
\end{proof}

\begin{proposition}[First and second cosmological rows]
\label{prop:n9v58-volume-variations}
For $V(g)=\int\sqrt{\det g}\,dx$,
\begin{align}
 DV_g[h]
 &=\frac12\int \operatorname{tr}_g(h)\,dV_g,
 \label{eq:n9v58-volume-first}\\
 D^2V_g[h,k]
 &=\int\left(
   \frac14\operatorname{tr}_g(h)\operatorname{tr}_g(k)
  -\frac12\operatorname{tr}_g(hg^{-1}k)
 \right)dV_g.
 \label{eq:n9v58-volume-second}
\end{align}
Thus a cosmological coefficient contributes an algebraic, momentum-zero
Hessian row even on trace-free variations.
\end{proposition}

\begin{proof}
Use $D\log\det g[h]=\operatorname{tr}(g^{-1}h)$ and
$Dg^{-1}[k]=-g^{-1}kg^{-1}$, then differentiate once more.
\end{proof}

\begin{theorem}[Flat stationary background removes the total algebraic row]
\label{thm:n9v58-flat-algebraic}
Assume:
\begin{enumerate}
\item $\bar g$ is flat and translation invariant;
\item the regulator effective action is diffeomorphism invariant and
its local derivative expansion contains no background tensor;
\item all nonlocal quadratic remainders are continuous at $p=0$ and
have been separated from the local algebraic coefficient; and
\item $D\Gamma_n(\bar g)=0$ for arbitrary constant conformal
variations.
\end{enumerate}
Then the coefficient $\lambda_{\rm eff,n}$ of the total
zero-derivative metric density vanishes.  Consequently its contribution
to the full physical Hessian at $p=0$ vanishes.
\end{theorem}

\begin{proof}
By \Cref{lem:n9v58-algebraic-density}, the complete local
zero-derivative row is $\lambda_{\rm eff,n}V(g)$.  Curvature terms vanish
to first order on constant conformal variations of a flat metric.
Using $h=\bar g$ in \eqref{eq:n9v58-volume-first}, stationarity gives
$0=(d/2)\lambda_{\rm eff,n}V(\bar g)$.  Hence
$\lambda_{\rm eff,n}=0$, and \eqref{eq:n9v58-volume-second} removes its
Hessian row.
\end{proof}

\begin{remark}[Curved backgrounds]
\label{rem:n9v58-curved}
On an Einstein background, curvature and cosmological terms mix in the
background equation and the Lichnerowicz Hessian contains curvature
potentials.  The correct comparison is then with that covariant
reference operator.  Setting a coordinate Fourier zero mode to zero is
neither invariant nor required.
\end{remark}

\subsection{Taylor ownership of local and nonlocal rows}

Let $J_n(z)=J_n(-z)=J_n(z)^*$ be the translation-invariant physical
kernel of the combined row ${\mathfrak j}_n$ after gauge reduction.  Put
\begin{align}
 J_{0,n}&=\sum_zJ_n(z),
 \label{eq:n9v58-J0}\\
 Q_{2,n}(p)&=-\frac12\sum_zJ_n(z)(p\cdot z)^2.
 \label{eq:n9v58-Q2}
\end{align}

\begin{lemma}[Second-order local/nonlocal decomposition]
\label{lem:n9v58-Taylor-decomposition}
If
\begin{equation}
 M_{4,n}:=\sum_z|z|^4\lVert J_n(z)\rVert<\infty,
 \label{eq:n9v58-fourth-moment}
\end{equation}
then
\begin{equation}
 \widehat J_n(p)=J_{0,n}+Q_{2,n}(p)+R_{4,n}(p),
 \qquad
 \lVert R_{4,n}(p)\rVert\le\frac{M_{4,n}}{24}|p|^4.
 \label{eq:n9v58-Taylor-bound}
\end{equation}
\end{lemma}

\begin{proof}
Evenness gives
$\widehat J_n(p)=\sum_zJ_n(z)\cos(p\cdot z)$.  Apply
$|\cos t-1+t^2/2|\le |t|^4/24$, sum the operator norms, and use
$|p\cdot z|\le|p||z|$.
\end{proof}

\begin{definition}[Renormalized derivative ownership]
\label{def:n9v58-derivative-ownership}
The following allocation counts each Taylor row once:
\begin{enumerate}
\item $J_{0,n}$ joins the cosmological/vacuum coefficient and is fixed
by the chosen background Einstein equation;
\item the diffeomorphism-invariant part of $Q_{2,n}$ joins the
Einstein--Hilbert row and fixes the renormalized Newton coefficient;
\item allowed higher local Taylor polynomials join their declared
curvature counterterms; and
\item $R_{4,n}$ is the connected nonlocal remainder.
\end{enumerate}
\end{definition}

\begin{proposition}[Infrared smallness after two Taylor owners]
\label{prop:n9v58-IR-smallness}
Suppose the renormalized positive physical reference satisfies
$\widehat A_n(p)\ge a_{0,n}|p|^2I$, the $J_{0,n}$ and $Q_{2,n}$ rows have
been assigned as in \Cref{def:n9v58-derivative-ownership}, and
\eqref{eq:n9v58-fourth-moment} holds.  On $0<|p|\le\mu$,
\begin{equation}
 \lVert R_{4,n}(p)\rVert
 \le \frac{M_{4,n}\mu^2}{24a_{0,n}}
       \lambda_{\min}(\widehat A_n(p)).
 \label{eq:n9v58-R4-relative}
\end{equation}
Thus the infrared relative coefficient tends to zero with $\mu$ if
$M_{4,n}/a_{0,n}$ is regulator-uniformly bounded.
\end{proposition}

\begin{proof}
By \eqref{eq:n9v58-Taylor-bound},
$\lVert R_{4,n}(p)\rVert\le M_{4,n}|p|^4/24$.  Since
$|p|^4\le\mu^2|p|^2$ and
$\widehat A_n(p)\ge a_{0,n}|p|^2I$, the stated comparison follows.
\end{proof}

\begin{firewall}[A subtraction is not a prediction]
\label{fire:n9v58-subtraction}
Locality identifies the form of the $J_0$ and $Q_2$ owners, but it does
not determine their finite coefficients.  The cosmological coefficient
is fixed by a background/measurement condition and the Newton
coefficient by a gravitational normalization condition.  Calling these
rows counterterms does not prove their required continuum tuning.
\end{firewall}

\subsection{Renormalized zero-mode ownership card}

\begin{definition}[Renormalized zero-mode ownership card (RZMOC)]
\label{def:n9v58-RZMOC}
RZMOC requires:
\begin{enumerate}
\item the invariant combined Hessian row
\eqref{eq:n9v58-combined-row}, with every block energy counted once;
\item regulator-level anomaly cancellation and the on-shell Ward
identity \eqref{eq:n9v58-onshell-gauge-null};
\item a declared background Einstein condition fixing the total
cosmological/vacuum row;
\item a declared Newton normalization fixing the physical $p^2$ row;
\item evenness or explicit first-moment cancellation for the connected
remainder;
\item regulator-uniform second or fourth spatial moments, according to
whether the $p^2$ row is retained or subtracted; and
\item the V57 strict aggregate relative margin after the owned local
rows have been moved into the reference.
\end{enumerate}
\end{definition}

\begin{theorem}[RZMOC supplies the massless input to ICRMC]
\label{thm:n9v58-RZMOC}
Under RZMOC, the zero-momentum obstruction in V57 is tested on the total
renormalized effective Hessian rather than on an arbitrarily separated
covariance row.  The algebraic term is fixed by the background Einstein
condition, the two-derivative term renormalizes the positive metric
reference, and the connected remainder is relatively small in the
infrared under the uniform moment hypothesis.
\end{theorem}

\begin{proof}
Use \Cref{prop:n9v58-redistribution,prop:n9v58-Ward-Hessian,
thm:n9v58-flat-algebraic,lem:n9v58-Taylor-decomposition,
prop:n9v58-IR-smallness}.  Apply the V57 Schur test outside the infrared
window and combine the two bounds through the declared momentum split.
\end{proof}

\begin{assessment}[Progress at ninth-series V58]
\label{ass:n9v58-status}
The apparent demand that the stress covariance vanish at zero momentum
has been replaced by the invariant condition on the complete effective
Hessian.  Ward identities remove gauge directions on shell; the
background equation owns the algebraic gravitational row; Newton
normalization owns the quadratic row; and a moment estimate controls the
remaining nonlocal kernel.  These are distinct obligations.
\end{assessment}

\begin{programme}[Ninth-series V59: regulator-uniform moment transport]
\label{prog:n9v58-V59}
V59 will derive second- and fourth-moment bounds from the sector-aware
fast clustering obtained in V51--V55.  It will track lattice spacing,
block diameter, correlation length, and response normalization so that
RZMOC does not conceal a divergent continuum constant.
\end{programme}

\begin{firewall}[Claim boundary after ninth-series V58]
\label{fire:n9v58-boundary}
V58 proves the ownership and Taylor estimates.  It does not establish
the physical chiral regulator's anomaly cancellation, the needed
uniform moments, the numerical strict margin, the slow infrared theory,
reflection positivity, OS reconstruction, or the full SM--Einstein
continuum.
\end{firewall}

\begin{theorem}[Ninth-series V58 result]
\label{thm:n9v58-result}
The massless comparison must be imposed on the redistribution-invariant
combined Hessian row \eqref{eq:n9v58-combined-row}.  On a flat stationary
background its local algebraic owner vanishes; its quadratic Taylor row
renormalizes the Einstein--Hilbert reference; and, after these owners are
removed, a finite fourth moment gives the infrared coefficient
$M_{4,n}\mu^2/(24a_{0,n})$.
\end{theorem}

\begin{proof}
Combine \Cref{prop:n9v58-redistribution,thm:n9v58-flat-algebraic,
lem:n9v58-Taylor-decomposition,prop:n9v58-IR-smallness,
thm:n9v58-RZMOC}.
\end{proof}

\paragraph{Parent-version record.}
V58 was formed from
\texttt{Renewal\_SM\_Einstein\_Continuum\_Research\_Notes\_V57.tex}.
The V57 parent had $21\,219$ logical source lines, $762\,876$ bytes, and
SHA--256
\[
 \texttt{03E453E827D3E05DFE3CE36BDDBB70F3FAAD55D1D977ED8843C5FE444C334695}.
\]

% =============================================================================
% END APPENDED NINTH-SERIES V58 MATERIAL
% =============================================================================
% =============================================================================
% BEGIN APPENDED NINTH-SERIES V59 MATERIAL
% =============================================================================

\section{Ninth-series V59: regulator-uniform moment transport}
\label{sec:v9v59}

\subsection{Physical scaling behind a lattice decay estimate}

V58 reduced the infrared comparison to spatial moments of the complete
renormalized Hessian kernel.  A decay rate in graph distance is not yet
a continuum estimate: the number of lattice blocks per physical volume
diverges as the mesh is removed.  This note records the amplitude and
length scaling that makes the second and fourth moments uniform.

Let $X_n$ be a quasi-uniform $d$-dimensional block lattice of mesh
$a_n$.  Write $d_n(x,y)$ for graph distance and assume the shell bound
\begin{equation}
 \#\{y:r\le d_n(x,y)<r+1\}
 \le S_d(1+r)^{d-1}
 \label{eq:n9v59-shell-count}
\end{equation}
uniformly in $n$ and $x$.  If $k_n(x,y)$ is the operator norm of a
reference-normalized kernel block, define its physical row moment
\begin{equation}
 {cal M}_{q,n}(k)=
 \sup_x\sum_y(a_nd_n(x,y))^q k_n(x,y).
 \label{eq:n9v59-row-moment}
\end{equation}

\begin{lemma}[Exponential lattice sum]
\label{lem:n9v59-exp-sum}
For $q\ge0$, $c>0$, and $0<\epsilon\le1$, there is a constant
$C_{d,q,c}$ such that
\begin{equation}
 \epsilon^{d+q}
 \sum_{r=0}^{\infty}(1+r)^{d-1}r^qe^{-c\epsilon r}
 \le C_{d,q,c}.
 \label{eq:n9v59-exp-sum}
\end{equation}
\end{lemma}

\begin{proof}
For $r\ge1$, compare each summand with a constant multiple of the
integral of $(1+t)^{d-1+q}e^{-c\epsilon t/2}$ over
$[r-1,r]$.  The change of variables $s=\epsilon t$ bounds the resulting
integral by a constant times $\epsilon^{-d-q}$.  The $r=0$ term is
bounded separately.
\end{proof}

Let $\xi_n$ be a physical correlation length and set
\begin{equation}
 \ell_n=\max\{a_n,\xi_n\},
 \qquad \epsilon_n=a_n/\ell_n.
 \label{eq:n9v59-saturated-length}
\end{equation}

\begin{theorem}[Uniform exponential moment compiler]
\label{thm:n9v59-exp-compiler}
Suppose
\begin{equation}
 k_n(x,y)\le B_n\left(\frac{a_n}{\ell_n}\right)^d
 \exp\left[-c\frac{a_nd_n(x,y)}{\ell_n}\right].
 \label{eq:n9v59-scaled-envelope}
\end{equation}
Then for every $q\ge0$,
\begin{equation}
 {cal M}_{q,n}(k)
 \le S_dC_{d,q,c}B_n\ell_n^q.
 \label{eq:n9v59-uniform-moment}
\end{equation}
In particular, the row sum and the second and fourth moments are
regulator-uniform when $B_n$ and $\ell_n$ are uniformly bounded.
\end{theorem}

\begin{proof}
Use \eqref{eq:n9v59-shell-count} and
\eqref{eq:n9v59-scaled-envelope} in
\eqref{eq:n9v59-row-moment}.  Since
$a_nr=\ell_n\epsilon_nr$, the remaining series is precisely
\eqref{eq:n9v59-exp-sum} times $B_n\ell_n^q$.
\end{proof}

\begin{firewall}[The cell-volume factor is indispensable]
\label{fire:n9v59-cell-factor}
The factor $(a_n/\ell_n)^d$ in
\eqref{eq:n9v59-scaled-envelope} is the discrete cell-volume
normalization of an integral kernel at range $\ell_n$.  If the block
amplitude stays of order one while $\ell_n/a_n\to\infty$, the row sum
grows like $(\ell_n/a_n)^d$.  Exponential shape alone therefore does not
give a continuum Schur constant.
\end{firewall}

\subsection{Response normalization}

Let $A_{x,n}$ be the local positive metric reference and let
$L_{x,n}u_x$ be the centred fast logarithmic response.  Define the
normalized covariance blocks by
\begin{equation}
 k_n(x,y)=
 \sup_{u_x,v_y\ne0}
 \frac{|\operatorname{Cov}(L_{x,n}u_x,L_{y,n}v_y)|}
 {\lVert A_{x,n}^{1/2}u_x\rVert
  \lVert A_{y,n}^{1/2}v_y\rVert}.
 \label{eq:n9v59-normalized-response}
\end{equation}

\begin{definition}[Continuum response envelope (CRE)]
\label{def:n9v59-CRE}
CRE at range $\ell_n$ is the estimate
\eqref{eq:n9v59-scaled-envelope} for the complete normalized response
kernel, with $B_n$ independent of the volume.  Uniform CRE means
\begin{equation}
 \sup_n B_n<\infty,
 \qquad
 \sup_n\ell_n<\infty.
 \label{eq:n9v59-uniform-CRE}
\end{equation}
\end{definition}

\begin{proposition}[AADC/MHDCC rate converted to physical length]
\label{prop:n9v59-path-length}
Suppose the connected-path argument of V51--V52 yields
\begin{equation}
 k_n(x,y)\le\beta_n e^{-\mu_nd_n(x,y)},
 \qquad
 \mu_n=-\log\bigl(\Delta_n(q_n+b_n)\bigr)>0.
 \label{eq:n9v59-path-decay}
\end{equation}
Its physical correlation length is
\begin{equation}
 \xi_n^{\rm path}=a_n/\mu_n.
 \label{eq:n9v59-path-xi}
\end{equation}
The path estimate implies CRE only if its amplitude also satisfies
\begin{equation}
 \beta_n\le B_n
 \left(\frac{a_n}{\ell_n}\right)^d,
 \qquad
 \ell_n\ge\max\{a_n,\xi_n^{\rm path}\},
 \label{eq:n9v59-path-amplitude}
\end{equation}
up to a fixed change of the exponential constant.
\end{proposition}

\begin{proof}
Equation \eqref{eq:n9v59-path-xi} rewrites
$e^{-\mu_nd_n}$ as
$e^{-a_nd_n/\xi_n^{\rm path}}$.  Increasing the range from
$\xi_n^{\rm path}$ to $\ell_n$ only weakens the exponential.  The
remaining comparison with \eqref{eq:n9v59-scaled-envelope} is exactly
\eqref{eq:n9v59-path-amplitude}.
\end{proof}

\begin{assessment}[What V51--V52 do and do not provide]
\label{ass:n9v59-path-assessment}
The disagreement expansion supplies a decay exponent and an explicit
defect dependence.  A continuum moment theorem additionally needs the
reference-normalized response amplitude in
\eqref{eq:n9v59-path-amplitude}.  This amplitude is a local
renormalization estimate, not a combinatorial consequence of path
counting.
\end{assessment}

\subsection{Polynomial tails and sharp moment thresholds}

\begin{lemma}[Polynomial moment threshold]
\label{lem:n9v59-poly-moment}
Suppose instead that
\begin{equation}
 k_n(x,y)\le B_n\left(\frac{a_n}{\ell_n}\right)^d
 \left(1+\frac{a_nd_n(x,y)}{\ell_n}\right)^{-\sigma}.
 \label{eq:n9v59-poly-envelope}
\end{equation}
Then ${\cal M}_{q,n}(k)\le C B_n\ell_n^q$ when
\begin{equation}
 \sigma>d+q.
 \label{eq:n9v59-poly-threshold}
\end{equation}
For a system of physical diameter $L_n$, the same bound acquires a
factor of order
\begin{equation}
 \begin{cases}
  1,&\sigma>d+q,\\
  1+\log(1+L_n/\ell_n),&\sigma=d+q,\\
  (1+L_n/\ell_n)^{d+q-\sigma},&\sigma<d+q.
 \end{cases}
 \label{eq:n9v59-poly-three-regimes}
\end{equation}
\end{lemma}

\begin{proof}
After the shell bound and the substitution
$t=a_nr/\ell_n$, the relevant integral is a constant multiple of
\begin{equation}
 B_n\ell_n^q\int_0^{L_n/\ell_n}
 t^{d-1+q}(1+t)^{-\sigma}\,dt.
 \label{eq:n9v59-poly-integral}
\end{equation}
Its large-$t$ exponent is $d-1+q-\sigma$, giving the three cases.
\end{proof}

\begin{corollary}[Four-dimensional stress diagnostic]
\label{cor:n9v59-4d-stress}
In $d=4$, a scale-invariant stress kernel with tail $|x|^{-8}$ has a
finite second moment but a logarithmically divergent fourth moment.  It
can satisfy the V57 $O(p^2)$ comparison after zero-mode ownership, but it
does not satisfy the finite-fourth-moment hypothesis of V58.
\end{corollary}

\begin{proof}
Set $d=4$, $\sigma=8$.  For $q=2$ one has $8>6$; for $q=4$ one has the
borderline equality $8=8$ in
\eqref{eq:n9v59-poly-three-regimes}.
\end{proof}

\subsection{The borderline logarithm is still infrared subordinate}

The fourth-moment condition in V58 was sufficient, not necessary.  The
critical tail admits a logarithmic replacement.

\begin{theorem}[Borderline fourth-order Taylor bound]
\label{thm:n9v59-borderline}
Let $J(z)=J(-z)=J(z)^*$ on a $d$-dimensional lattice, suppose its zeroth
and quadratic Taylor owners have been subtracted, and assume for
$|z|\ge\ell$
\begin{equation}
 \lVert J(z)\rVert\le B|z|^{-d-4},
 \label{eq:n9v59-critical-tail}
\end{equation}
with a finite corresponding short-distance fourth moment.  Then for
$0<|p|\ell\le1$,
\begin{equation}
 \lVert\widehat J(p)\rVert
 \le C B |p|^4
 \bigl(1+|\log(|p|\ell)|\bigr),
 \label{eq:n9v59-p4log}
\end{equation}
where $C$ also absorbs the short-distance moment.  Relative to
$\widehat A(p)\ge a_0|p|^2I$, the coefficient is bounded by
\begin{equation}
 \frac{CB}{a_0}|p|^2
 \bigl(1+|\log(|p|\ell)|\bigr)\longrightarrow0.
 \label{eq:n9v59-p4log-relative}
\end{equation}
\end{theorem}

\begin{proof}
By evenness and the subtraction of the zeroth and quadratic owners,
\begin{equation}
 \widehat J(p)=\sum_zJ(z)
 \left(\cos(p\cdot z)-1+\frac12(p\cdot z)^2\right).
 \label{eq:n9v59-subtracted-symbol}
\end{equation}
Split the sum at $|z|=|p|^{-1}$.  On the inner part use
$|\cos t-1+t^2/2|\le |t|^4/24$.  The shell sum of
$|z|^4\lVert J(z)\rVert$ grows logarithmically, giving the first factor
in \eqref{eq:n9v59-p4log}.  On the outer part use
$|\cos t-1+t^2/2|\le 2+t^2/2$.  The two tail integrals are of orders
$|p|^4$ and $|p|^4$, respectively.  This proves
\eqref{eq:n9v59-p4log}; division by $a_0|p|^2$ proves
\eqref{eq:n9v59-p4log-relative}.
\end{proof}

\begin{remark}[Running rather than a missing cancellation]
\label{rem:n9v59-running}
The logarithm in \eqref{eq:n9v59-p4log} is the momentum-space signature
of a borderline curvature-squared row.  It is compatible with
infrared subordination to the Einstein $p^2$ reference, but it prevents
one from replacing the kernel by a regulator-uniform analytic fourth
order Taylor coefficient.
\end{remark}

\subsection{Sector-aware integration}

Let the fast response be conditioned on retained slow variables, as in
V55--V56.  The conditional covariance entering the fast-integrated
Hessian is
\begin{equation}
 {\mathfrak c}_{f|s}(u,v)
 =\operatorname{Cov}_{\nu_s}(A_u^f,A_v^f).
 \label{eq:n9v59-conditional-fast-cov}
\end{equation}

\begin{proposition}[Conditional integration prevents a false moment]
\label{prop:n9v59-sector-split}
Assume the fast conditional kernel obeys uniform CRE while a retained
slow sector has only the polynomial envelope
\eqref{eq:n9v59-poly-envelope}.  Integrating only the fast variables
produces the uniform moments \eqref{eq:n9v59-uniform-moment}; the slow
tail does not enter \eqref{eq:n9v59-conditional-fast-cov}.  After the
slow variables are integrated, their contribution reappears through the
law of total covariance and must be estimated with
\Cref{lem:n9v59-poly-moment}.
\end{proposition}

\begin{proof}
The V56 differentiation formula is conditional on the fixed slow
configuration, so its covariance is precisely
\eqref{eq:n9v59-conditional-fast-cov}.  The law of total covariance from
V55 adds the covariance of conditional means only when the outer slow
law is taken.  Apply the exponential and polynomial moment lemmas to
the two terms separately.
\end{proof}

\begin{firewall}[Cross-sector labels do not erase the slow tail]
\label{fire:n9v59-cross-sector}
A conditional label or finite-range block decomposition can localize the
fast covariance, but the covariance of the conditional mean retains the
slow sector's infrared law.  Assigning that outer covariance to the fast
row would produce a spurious uniform correlation length.
\end{firewall}

\subsection{Uniform moment transport card}

\begin{definition}[Uniform moment transport card (UMTC)]
\label{def:n9v59-UMTC}
UMTC requires:
\begin{enumerate}
\item a quasi-uniform physical embedding and shell count
\eqref{eq:n9v59-shell-count};
\item the complete reference-normalized response
\eqref{eq:n9v59-normalized-response};
\item the cell-volume amplitude and physical range in uniform CRE;
\item a conditional fast/retained slow split before invoking an
exponential moment;
\item uniform second moments for the zero-mode comparison;
\item either uniform fourth moments or an explicit borderline
$p^4\log p$ estimate after the quadratic owner is removed;
\item separate tracking of defect-path decay and local response
amplitude; and
\item constants uniform in volume, mesh, chart, and accepted boundary
sector.
\end{enumerate}
\end{definition}

\begin{theorem}[UMTC supplies the quantitative RZMOC moments]
\label{thm:n9v59-UMTC}
Under UMTC, the fast integrated kernel has regulator-uniform Schur,
second-moment, and, away from the declared borderline, fourth-moment
bounds.  A $|x|^{-d-4}$ sector instead has the explicit
$p^4\log p$ remainder \eqref{eq:n9v59-p4log}, which remains infrared
subordinate to a positive Einstein $p^2$ reference.
\end{theorem}

\begin{proof}
Use \Cref{thm:n9v59-exp-compiler,lem:n9v59-poly-moment,
thm:n9v59-borderline,prop:n9v59-sector-split}.  The $q=0$ estimate is the
Schur row sum, and $q=2,4$ give the RZMOC moments.
\end{proof}

\begin{assessment}[Progress at ninth-series V59]
\label{ass:n9v59-status}
The moment hypothesis is now expressed in physical units.  Uniformity
requires both a bounded physical correlation length and the correct
cell-volume response amplitude.  Massless retained sectors obey their
own power-law thresholds; the four-dimensional stress tail is exactly
fourth-moment borderline but still gives an infrared-small
$p^4\log p$ remainder after local ownership.
\end{assessment}

\begin{programme}[Ninth-series V60: cross-program audit and scaling ledger]
\label{prog:n9v59-V60}
V60 is the compulsory five-version audit.  It will inspect the newest
Yang--Mills and Navier--Stokes notes, compare their contraction and
terminal-source mechanisms with UMTC, and produce a single scaling
ledger for the next SM--Einstein step.
\end{programme}

\begin{firewall}[Claim boundary after ninth-series V59]
\label{fire:n9v59-boundary}
V59 proves the scaling compilers.  It does not establish uniform CRE for
the physical chiral determinant, the required local amplitude, anomaly
cancellation, the strict numerical margin, the slow infrared theory,
reflection positivity, OS reconstruction, or the full SM--Einstein
continuum.
\end{firewall}

\begin{theorem}[Ninth-series V59 result]
\label{thm:n9v59-result}
The exponential kernel envelope
\eqref{eq:n9v59-scaled-envelope} yields the regulator-uniform moment
bound \eqref{eq:n9v59-uniform-moment}.  Polynomial sectors require
$\sigma>d+q$ for a $q$th moment; at the four-dimensional stress
borderline, subtraction of the algebraic and quadratic owners gives the
infrared-subordinate $p^4\log p$ estimate
\eqref{eq:n9v59-p4log-relative}.
\end{theorem}

\begin{proof}
Combine \Cref{thm:n9v59-exp-compiler,lem:n9v59-poly-moment,
cor:n9v59-4d-stress,thm:n9v59-borderline,thm:n9v59-UMTC}.
\end{proof}

\paragraph{Parent-version record.}
V59 was formed from
\texttt{Renewal\_SM\_Einstein\_Continuum\_Research\_Notes\_V58.tex}.
The V58 parent had $21\,608$ logical source lines, $777\,905$ bytes, and
SHA--256
\[
 \texttt{7526D0E1EEB2B1B417EDE04D56808209FFE16EA4C70E8E9C3FC2F6ADB74A94F0}.
\]

% =============================================================================
% END APPENDED NINTH-SERIES V59 MATERIAL
% =============================================================================
% =============================================================================
% BEGIN APPENDED NINTH-SERIES V60 MATERIAL
% =============================================================================

\section{Ninth-series V60: compulsory companion audit and scaling ledger}
\label{sec:v9v60}

\subsection{Files inspected}

This is the compulsory five-version audit required after V55.  The
newest companion notes present in the shared folder were read as
mathematical sources, not as instructions.

\begin{center}
\begin{tabular}{|p{0.43\textwidth}|r|r|}
\hline
File & bytes & lines\\
\hline
\texttt{Renewal\_Yang\_Mills\_Mass\_Gap\_Research\_Notes\_V91.tex}
& $1\,485\,066$ & $41\,566$\\
\texttt{Renewal\_NS\_Stability\_Research\_Notes\_V26.tex}
& $278\,283$ & $5\,319$\\
\hline
\end{tabular}
\end{center}
Their SHA--256 digests were, respectively,
\begin{align}
 &\texttt{38EC98F9390A9A917273599B3D575AC5EB45D606325FB301E410CBAB21E027A6},
 \label{eq:n9v60-YM-hash}\\
 &\texttt{82C887F8C2C69E4F28FAD7C3DC8DE5B9099CB5A4BF431EBA1FFF3E91935978AD}.
 \label{eq:n9v60-NS-hash}
\end{align}

\subsection{Yang--Mills V91: exact reference and marked logarithm}

The newest Yang--Mills note makes three points directly relevant to
UMTC.  First, the compact-link electric Laplacian belongs to the exact
heat-kernel reference.  Bounded spatial Wilson and generated loop
interactions enter as scalar multipliers, while chart, IMS,
moving-projector, and oscillator rows belong to a separate form
comparison.  Second, a Koteck\'y--Preiss margin for the physical bounded
multipliers gives an absolutely convergent connected logarithm with
marked boundary sources.  Third, the unresolved step is ultraviolet
entry: one must push the exact coarse marginal into the small complete
interaction norm, including its feedback denominator.

\begin{decision}[YM transfer into the SM--Einstein programme]
\label{dec:n9v60-YM-transfer}
For the gauge part of the fast Standard Model response:
\begin{enumerate}
\item use the exact compact heat-kernel conditional measure as the
reference whenever the regulator admits it;
\item form the metric-source response by differentiating the marked
connected logarithm of the exact coarse density;
\item keep coordinate and electric-generator comparison rows out of the
bounded physical polymer activity; and
\item treat entry of the exact coarse marginal into a uniform rooted
norm as an open quantitative step.
\end{enumerate}
\end{decision}

The connection with V59 can be stated without importing any unproved
YM conclusion.

\begin{proposition}[Rooted marked clusters imply a response envelope]
\label{prop:n9v60-rooted-to-envelope}
Let $c_{n,X}^{(2)}:{\cal H}_x\times{\cal H}_y\to\mathbb R$ be connected
two-source coefficients supported on finite block sets $X$.  Suppose
that every contribution to the normalized kernel $k_n(x,y)$ contains
$x,y\in X$ and that
\begin{equation}
 \sup_x\sum_{X\ni x}
 e^{\mu_n\operatorname{diam}_nX}\lVert c_{n,X}^{(2)}\rVert
 \le R_n.
 \label{eq:n9v60-rooted-norm}
\end{equation}
Then
\begin{equation}
 k_n(x,y)\le R_ne^{-\mu_nd_n(x,y)}.
 \label{eq:n9v60-cluster-decay}
\end{equation}
Consequently it satisfies the V59 continuum response envelope at
$\ell_n=\max\{a_n,a_n/\mu_n\}$ if
\begin{equation}
 R_n\le B_n(a_n/\ell_n)^d.
 \label{eq:n9v60-rooted-amplitude}
\end{equation}
\end{proposition}

\begin{proof}
Any $X$ containing $x$ and $y$ has
$\operatorname{diam}_nX\ge d_n(x,y)$.  Hence
\begin{align}
 k_n(x,y)
 &\le\sum_{X\ni x,y}\lVert c_{n,X}^{(2)}\rVert\\
 &\le e^{-\mu_nd_n(x,y)}
       \sum_{X\ni x}e^{\mu_n\operatorname{diam}_nX}
       \lVert c_{n,X}^{(2)}\rVert,
 \end{align}
which proves \eqref{eq:n9v60-cluster-decay}.  The final statement is
\Cref{prop:n9v59-path-length}.
\end{proof}

\begin{firewall}[KP convergence is not yet continuum normalization]
\label{fire:n9v60-KP-normalization}
A volume-uniform bound $R_n=O(1)$ proves clustering in lattice units.
When $\mu_n\to0$, CRE requires the sharper density scaling
$R_n=O(\mu_n^d)$ in \eqref{eq:n9v60-rooted-amplitude}.  The marked
cluster construction and the continuum normalization are therefore
separate estimates.
\end{firewall}

\subsection{Navier--Stokes V26: norms on the actual input class}

The newest Navier--Stokes note computes exact pointwise norms of first
and second Oseen-kernel derivatives on rank-one inputs.  Its first
derivative gains nothing over the full symmetric trace-free norm; its
second derivative is strictly smaller on an intermediate radial range,
with a pointwise gain up to about eleven percent and a smaller integrated
gain.  The transferable principle is to normalize on the actual input
class and to certify whether a gain exists, rather than assuming that a
restriction is helpful.

Let $T_x:{\cal G}_x\to{\cal H}_x$ describe an exact block-local class of
admissible physical variations.  Set
\begin{equation}
 A_x^{\rm red}=T_x^*A_xT_x,
 \qquad
 C_{xy}^{\rm red}=T_x^*C_{xy}T_y.
 \label{eq:n9v60-local-restriction}
\end{equation}

\begin{proposition}[Restricted normalized block norm]
\label{prop:n9v60-restricted-norm}
Assume $A_x^{\rm red}>0$ and define
\begin{equation}
 U_x=A_x^{1/2}T_x(A_x^{\rm red})^{-1/2}.
 \label{eq:n9v60-isometry}
\end{equation}
Then $U_x^*U_x=I$ and
\begin{equation}
 (A_x^{\rm red})^{-1/2}C_{xy}^{\rm red}
 (A_y^{\rm red})^{-1/2}
 =U_x^*K_{xy}U_y.
 \label{eq:n9v60-compressed-block}
\end{equation}
In particular,
\begin{equation}
 k_{xy}^{\rm red}\le k_{xy}.
 \label{eq:n9v60-restricted-less}
\end{equation}
The inequality can be an equality.
\end{proposition}

\begin{proof}
The definition of $A_x^{\rm red}$ gives
$U_x^*U_x=(A_x^{\rm red})^{-1/2}T_x^*A_xT_x
(A_x^{\rm red})^{-1/2}=I$.  Insert
$C_{xy}=A_x^{1/2}K_{xy}A_y^{1/2}$ into
\eqref{eq:n9v60-local-restriction} to obtain
\eqref{eq:n9v60-compressed-block}.  Isometries have norm one, which
proves \eqref{eq:n9v60-restricted-less}.  Equality occurs when a top
singular pair of $K_{xy}$ lies in the two isometry ranges.
\end{proof}

\begin{decision}[NS transfer into the SM--Einstein programme]
\label{dec:n9v60-NS-transfer}
The local response amplitude $B_n$ in UMTC should be computed on the
exact constraint-reduced metric/source tangent class when that class has
a block-local parametrization.  The ambient norm remains the valid
fallback.  No gain is booked until an identity such as
\eqref{eq:n9v60-compressed-block} and a certified restricted norm are
available.
\end{decision}

\begin{firewall}[Local compression versus elliptic projection]
\label{fire:n9v60-local-compression}
\Cref{prop:n9v60-restricted-norm} is block-local.  It does not override
the V57 warning for a global elliptic constraint projector, whose
inverse-Laplacian tail may alter the kernel decay.  In that case the
relative form should be restricted before projecting, or the projector
kernel must be estimated separately.
\end{firewall}

\subsection{The unified scaling ledger}

The audit yields the following regulator ledger.  A row marked
``proved conditionally'' means that the implication has been proved in
these SM notes, while its physical hypotheses remain acquisition tasks.

\begin{center}
\small
\begin{tabular}{|p{0.18\textwidth}|p{0.24\textwidth}|p{0.25\textwidth}|p{0.20\textwidth}|}
\hline
Quantity & Exact owner & Required scaling & Present status\\
\hline
Fast gauge reference & compact-link heat kernel/Gauss sector
& normalized conditional probability
& exact finite-cutoff route identified\\
\hline
Physical activities & bounded Wilson, Higgs, and generated local
multipliers
& uniform complete KP norm below its feedback threshold
& coarse-entry estimate open\\
\hline
Metric response & marked derivatives of the connected coarse logarithm
& rooted norm \eqref{eq:n9v60-rooted-norm}
& cluster implication proved conditionally\\
\hline
Continuum amplitude & reference-normalized two-source coefficient
& $R_n\le B_n(a_n/\ell_n)^d$
& local amplitude estimate open\\
\hline
Constraint reduction & exact physical tangent map
& local compression or controlled projector tail
& compression theorem proved conditionally\\
\hline
Zero and quadratic rows & cosmological/vacuum and Newton owners
& RZMOC background and normalization conditions
& ownership proved; tuning open\\
\hline
Nonlocal remainder & connected marked clusters
& uniform $M_2$, $M_4$, or $p^4\log p$
& UMTC compiler proved conditionally\\
\hline
Strict inversion & complete positive metric reference
& aggregate CPRAC coefficient below one
& numerical uniform margin open\\
\hline
Slow infrared sector & retained massless fields
& sector-aware power-law effective theory
& not yet constructed\\
\hline
\end{tabular}
\end{center}

\begin{theorem}[Audit integration theorem]
\label{thm:n9v60-audit-integration}
Suppose an exact coarse SM conditional density admits a marked connected
expansion satisfying \eqref{eq:n9v60-rooted-norm} and
\eqref{eq:n9v60-rooted-amplitude}, with physical activities separated
from heat-kernel and coordinate-form owners.  Suppose also that the
exact local physical tangent compression satisfies
\eqref{eq:n9v60-local-restriction}.  Then its fast metric-response
kernel satisfies CRE with an amplitude no larger than the ambient
amplitude, and hence loads the V59 moment compiler.
\end{theorem}

\begin{proof}
Use \Cref{prop:n9v60-rooted-to-envelope} to obtain CRE.  Apply
\Cref{prop:n9v60-restricted-norm} to every marked two-source block; the
rooted sum cannot increase.  Then invoke
\Cref{thm:n9v59-exp-compiler}.
\end{proof}

\begin{assessment}[Direction after the V60 audit]
\label{ass:n9v60-direction}
The common bottleneck is no longer generic clustering.  It is the exact
coarse pushforward with its complete interaction norm and correctly
scaled marked response amplitude.  YM V91 supplies the appropriate
compact reference and connected-logarithm architecture.  NS V26
justifies an exact restricted-tangent norm calculation only after that
kernel has been constructed.
\end{assessment}

\begin{programme}[Ninth-series V61: one exact marked coarse pushforward]
\label{prog:n9v60-V61}
V61 will disintegrate one finite compact gauge--Higgs block map, express
the coarse density as a conditional partition function, and derive exact
first and second metric-source derivatives of its connected logarithm.
It will display the feedback denominator and identify the additional
estimate needed to make the rooted coefficient obey
\eqref{eq:n9v60-rooted-amplitude}.
\end{programme}

\begin{firewall}[Claim boundary after ninth-series V60]
\label{fire:n9v60-boundary}
The audit imports no companion programme conclusion as a proved SM
hypothesis.  V60 does not prove coarse KP entry, the continuum response
amplitude, the physical chiral phase control, anomaly cancellation, the
strict margin, the slow infrared theory, reflection positivity, OS
reconstruction, or the full SM--Einstein continuum.
\end{firewall}

\begin{theorem}[Ninth-series V60 result]
\label{thm:n9v60-result}
The companion audit selects the exact compact heat-kernel coarse
marginal and its marked connected logarithm as the source of UMTC.  A
rooted marked coefficient bound gives the decay
\eqref{eq:n9v60-cluster-decay}, but continuum moments additionally
require \eqref{eq:n9v60-rooted-amplitude}.  Exact block-local physical
tangent compression cannot increase the normalized kernel norm, though
it need not improve it strictly.
\end{theorem}

\begin{proof}
Combine \Cref{dec:n9v60-YM-transfer,
prop:n9v60-rooted-to-envelope,prop:n9v60-restricted-norm,
dec:n9v60-NS-transfer,thm:n9v60-audit-integration}.
\end{proof}

\paragraph{Parent-version record.}
V60 was formed from
\texttt{Renewal\_SM\_Einstein\_Continuum\_Research\_Notes\_V59.tex}.
The V59 parent had $22\,014$ logical source lines, $792\,553$ bytes, and
SHA--256
\[
 \texttt{956227CDF8F4A95E9AC56601C4325A0F864D9858A5C9D10EF5D9AF5A52A1AC9B}.
\]

% =============================================================================
% END APPENDED NINTH-SERIES V60 MATERIAL
% =============================================================================
% =============================================================================
% BEGIN APPENDED NINTH-SERIES V61 MATERIAL
% =============================================================================

\section{Ninth-series V61: one exact marked coarse pushforward}
\label{sec:v9v61}

\subsection{Exact disintegration of a finite gauge--Higgs block}

Fix one finite regulator and a finite family of compact gauge links and
Higgs variables.  Let $(\Omega,{\cal F})$ be the fine configuration
space, let $(\Sigma,{\cal S})$ be the coarse block space, and let
\begin{equation}
 B:\Omega\longrightarrow\Sigma
 \label{eq:n9v61-block-map}
\end{equation}
be a measurable, gauge-equivariant block map.  At slow metric $g$, let
$\mu_{0,g}$ be the exact positive reference probability.  It contains
the compact heat-kernel/electric reference, the positive Higgs quartic
reference, and the exact Gauss-sector normalization assigned in V47,
V56, and V60.  Standard Borel disintegration gives
\begin{equation}
 \mu_{0,g}(d\omega)=\kappa_g(d\omega\mid s)q_{0,g}(ds),
 \qquad s=B\omega,
 \label{eq:n9v61-disintegration}
\end{equation}
where $q_{0,g}=B_\#\mu_{0,g}$.

Let $V_g(\omega)$ be the remaining real interaction and assume
$e^{-V_g/\hbar}$ is integrable.  Define
\begin{align}
 R_g(s)&=\int e^{-V_g(\omega)/\hbar}
                 \kappa_g(d\omega\mid s),
 \label{eq:n9v61-conditional-partition}\\
 U_g(s)&=-\hbar\log R_g(s).
 \label{eq:n9v61-coarse-potential}
\end{align}

\begin{theorem}[Exact one-step coarse density]
\label{thm:n9v61-exact-pushforward}
Let
\begin{equation}
 Z_g=\int e^{-V_g/\hbar},d\mu_{0,g}.
 \label{eq:n9v61-full-Z}
\end{equation}
Then the coarse pushforward of the interacting probability is
\begin{equation}
 q_g(ds)=Z_g^{-1}R_g(s)q_{0,g}(ds)
        =Z_g^{-1}e^{-U_g(s)/\hbar}q_{0,g}(ds).
 \label{eq:n9v61-exact-coarse-law}
\end{equation}
No truncation or locality assertion is used in this identity.
\end{theorem}

\begin{proof}
For bounded measurable $F$ on $\Sigma$, disintegration gives
\begin{align}
 \int F(B\omega)e^{-V_g(\omega)/\hbar}\mu_{0,g}(d\omega)
 &=\int F(s)R_g(s)q_{0,g}(ds).
 \end{align}
Divide by \eqref{eq:n9v61-full-Z}.  Taking $F=1$ also proves the
normalization.
\end{proof}

\begin{remark}[Additive coarse-action gauge]
\label{rem:n9v61-additive-gauge}
Multiplying $R_g$ by a metric-dependent constant shifts $U_g$ by a
constant and is compensated by $Z_g$.  Vacuum and ground-energy rows
must follow the V58 redistribution rule so that this harmless action
gauge is not counted as a second physical source.
\end{remark}

\subsection{Exact marked metric derivatives}

Fix a background $\bar g$.  Assume the conditional reference has a
twice differentiable Radon--Nikodym density in the metric directions
under consideration.  Its conditional score is
\begin{equation}
 b_h(\omega\mid s)
 =D_g\log\kappa_g(d\omega\mid s)\big|_{\bar g}[h],
 \qquad
 \mathbb E_{\kappa_{\bar g}(\cdot\mid s)}b_h=0.
 \label{eq:n9v61-conditional-score}
\end{equation}
Define the complete marked insertion
\begin{equation}
 X_h(\omega\mid s)
 =-\hbar^{-1}D_gV_g(\omega)\big|_{\bar g}[h]
   +b_h(\omega\mid s).
 \label{eq:n9v61-complete-mark}
\end{equation}
Let $\nu_s$ be the tilted conditional law
\begin{equation}
 \nu_s(d\omega)=R_{\bar g}(s)^{-1}
 e^{-V_{\bar g}(\omega)/\hbar}\kappa_{\bar g}(d\omega\mid s).
 \label{eq:n9v61-tilted-conditional}
\end{equation}

\begin{theorem}[First and second marked coarse responses]
\label{thm:n9v61-marked-derivatives}
Under an integrable second-order differentiation domain,
\begin{align}
 D_g\log R_g(s)\big|_{\bar g}[h]
 &=\mathbb E_{\nu_s}X_h,
 \label{eq:n9v61-logR-first}\\
 D_g^2\log R_g(s)\big|_{\bar g}[h,k]
 &=\mathbb E_{\nu_s}D_gX_h[k]
   +\operatorname{Cov}_{\nu_s}(X_h,X_k),
 \label{eq:n9v61-logR-second}\\
 D_gU_g(s)[h]
 &=-\hbar\mathbb E_{\nu_s}X_h,
 \label{eq:n9v61-U-first}\\
 D_g^2U_g(s)[h,k]
 &=-\hbar\mathbb E_{\nu_s}D_gX_h[k]
   -\hbar\operatorname{Cov}_{\nu_s}(X_h,X_k).
 \label{eq:n9v61-U-second}
\end{align}
\end{theorem}

\begin{proof}
Differentiate \eqref{eq:n9v61-conditional-partition}.  The logarithmic
derivative of its complete integrand is
\eqref{eq:n9v61-complete-mark}, proving
\eqref{eq:n9v61-logR-first}.  Differentiating a normalized expectation
gives the mean derivative plus covariance, exactly as in V56, proving
\eqref{eq:n9v61-logR-second}.  Multiply by $-\hbar$.
\end{proof}

\begin{corollary}[Coarse and reference sources are separate owners]
\label{cor:n9v61-reference-owner}
The logarithmic metric derivative of the unnormalized coarse density
$R_g(s)q_{0,g}(ds)$ is
\begin{equation}
 \mathbb E_{\nu_s}X_h
 +D_g\log q_{0,g}(ds)[h].
 \label{eq:n9v61-two-source-owners}
\end{equation}
The conditional reference score in $X_h$ and the coarse reference score
in the second term arise from the exact disintegration and must not be
reinserted as physical Wilson/Higgs activities.
\end{corollary}

\begin{proof}
Apply the product rule to $R_gq_{0,g}$ and use
\eqref{eq:n9v61-logR-first}.
\end{proof}

\begin{firewall}[Reference form marks]
\label{fire:n9v61-reference-form}
For a heat-kernel reference, $b_h$ may be a Duhamel or quadratic-form
insertion rather than a bounded scalar multiplier.  It can be used as a
marked observable under an integrable form bound, but it cannot be
placed among the unmarked bounded polymer activities merely to invoke a
scalar KP theorem.  This is the type separation imported at V60.
\end{firewall}

\subsection{When the conditional partition is a polymer gas}

Assume the coarse variable contains a skeleton sufficient to make the
fine reference interiors conditionally independent over elementary
cells.  Let the remaining bounded interaction factors be
\begin{equation}
 e^{-V_{\bar g}/\hbar}=\prod_{b\in{\cal B}}(1+f_b),
 \qquad
 f_b=e^{-\Phi_b/\hbar}-1.
 \label{eq:n9v61-factor-expansion}
\end{equation}
The Higgs quartic and electric generator remain in the reference so
that the $f_b$ are bounded in the positive paired sector.

\begin{proposition}[Exact conditional polymer representation]
\label{prop:n9v61-polymer-representation}
Suppose conditionally independent reference cells factor whenever their
dependency supports are disjoint.  Then
\begin{equation}
 R_{\bar g}(s)=
 \sum_{\substack{\Gamma\ {
m finite}\\
                  \gamma\sim\gamma'\ {
m for}\ \gamma\ne\gamma'}}
 \prod_{\gamma\in\Gamma}z_s(\gamma),
 \label{eq:n9v61-polymer-gas}
\end{equation}
where $\gamma$ ranges over connected sets of interaction labels,
$\sim$ means disjoint dependency supports, and
\begin{equation}
 z_s(\gamma)=
 \mathbb E_{\kappa_{\bar g}(\cdot\mid s)}
 \prod_{b\in\gamma}f_b
 \label{eq:n9v61-polymer-activity}
\end{equation}
with all conditionally independent components removed in the usual
connected grouping.
\end{proposition}

\begin{proof}
Expand the finite product in \eqref{eq:n9v61-factor-expansion} over
subsets of ${\cal B}$.  Decompose each subset into connected dependency
components.  Conditional factorization turns the expectation into the
product of its component expectations.  Summing compatible components
gives \eqref{eq:n9v61-polymer-gas}.
\end{proof}

\begin{firewall}[Insufficient coarse skeleton]
\label{fire:n9v61-skeleton}
If conditioning on $s$ couples disjoint fine interiors through an
unresolved Gauss constraint, determinant, or block-map fibre, the
factorization used in \Cref{prop:n9v61-polymer-representation} fails.
One must enlarge the coarse skeleton, use an exact normalized sector
disintegration, or place the residual coupling in the activity.  It is
not legitimate to delete it.
\end{firewall}

\subsection{Marked cluster scaling}

Let $t_x$ be a local metric-source parameter and let
$z_s(\gamma;t)$ denote the corresponding activity.  Write
$\partial_x$ for a unit derivative normalized by the local positive
metric reference.  Assume a KP margin with contraction parameter
$0<\theta<1$ and exponential size/diameter weight.

\begin{definition}[Half-density mark bound]
\label{def:n9v61-half-density}
At physical range $\ell_n$, the one- and two-mark activity bounds have
half-density size if, uniformly in the accepted coarse boundary $s$,
\begin{align}
 |\partial_x z_s(\gamma;0)|
 &\le m_n\mathbf 1_{\{x\in\operatorname{supp}\gamma\}}
       \zeta_s(\gamma),
 \label{eq:n9v61-one-mark}\\
 |\partial_x\partial_y z_s(\gamma;0)|
 &\le m_n^2\mathbf 1_{\{x,y\in\operatorname{supp}\gamma\}}
       \zeta_s(\gamma),
 \label{eq:n9v61-two-mark}\\
 m_n&\le M(a_n/\ell_n)^{d/2},
 \label{eq:n9v61-half-density-scale}
\end{align}
where $\zeta_s$ satisfies the same KP majorant as $|z_s|$.
\end{definition}

\begin{theorem}[Two marked sources acquire cell-volume size]
\label{thm:n9v61-marked-KP}
Under a uniform KP margin and
\Cref{def:n9v61-half-density}, the connected logarithm
$\log R_{\bar g}(s;t)$ is twice differentiable at $t=0$.  Its normalized
two-source kernel has a rooted bound
\begin{equation}
 k_n(x,y)\le C_{\rm KP}(\theta) m_n^2
 e^{-\mu_nd_n(x,y)}.
 \label{eq:n9v61-marked-decay}
\end{equation}
Consequently
\begin{equation}
 R_n\le C_{\rm KP}(\theta)M^2
       (a_n/\ell_n)^d,
 \label{eq:n9v61-cell-volume-output}
\end{equation}
which is the V60 rooted-amplitude condition.
\end{theorem}

\begin{proof}
The KP tree majorant makes the Ursell series for $\log R$ locally
uniformly absolutely convergent.  Differentiate it twice.  A term in
the mixed derivative has either two singly marked activities or one
doubly marked activity.  In both cases
\eqref{eq:n9v61-one-mark}--\eqref{eq:n9v61-two-mark} contribute at most
$m_n^2$ times the same positive tree majorant.  A connected cluster
carrying marks at $x$ and $y$ has diameter at least $d_n(x,y)$, so the
diameter weight supplies $e^{-\mu_nd_n(x,y)}$.  Rooting the remaining
tree sum at the first mark gives a finite constant
$C_{\rm KP}(\theta)$, uniform when the KP margin is uniform.
Finally use \eqref{eq:n9v61-half-density-scale}.
\end{proof}

\begin{remark}[Why two marks matter]
\label{rem:n9v61-two-marks}
The continuum response kernel needs a cell-volume factor
$(a_n/\ell_n)^d$.  It is produced as the product of two normalized
one-source factors of size $(a_n/\ell_n)^{d/2}$.  A volume-uniform KP
constant alone cannot manufacture this scaling.
\end{remark}

\subsection{The exact feedback denominator}

Coarse re-exponentiation and normalization often make the complete
coarse interaction $u_g$ satisfy an implicit equation in a Banach space
${\mathbb B}$,
\begin{equation}
 u_g=u_g^{(0)}+{\cal P}_g(u_g).
 \label{eq:n9v61-feedback-equation}
\end{equation}
Here ${\cal P}_g$ is the connected pushforward, including the accepted
boundary-sector normalization.

\begin{proposition}[Feedback resolvent]
\label{prop:n9v61-feedback}
Suppose ${\cal P}_g$ is differentiable and
\begin{equation}
 \lVert D_u{\cal P}_g(u_g)\rVert\le\theta<1.
 \label{eq:n9v61-feedback-margin}
\end{equation}
Then
\begin{equation}
 D_gu_g[h]
 =(I-D_u{\cal P}_g(u_g))^{-1}
 \left(D_gu_g^{(0)}[h]+D_g{\cal P}_g(u_g)[h]\right),
 \label{eq:n9v61-feedback-derivative}
\end{equation}
and
\begin{equation}
 \lVert D_gu_g[h]\rVert
 \le\frac{
 \lVert D_gu_g^{(0)}[h]\rVert+
 \lVert D_g{\cal P}_g(u_g)[h]\rVert}{1-\theta}.
 \label{eq:n9v61-feedback-denominator}
\end{equation}
The same inverse appears in the second derivative, whose right-hand side
also contains $D_u^2{\cal P}[Du,Du]$ and the two mixed derivatives.
\end{proposition}

\begin{proof}
Differentiate \eqref{eq:n9v61-feedback-equation} and move
$D_u{\cal P}\,Du$ to the left.  The Neumann series gives
$\lVert(I-D_u{\cal P})^{-1}\rVert\le(1-\theta)^{-1}$.  A second
differentiation gives the stated terms and the same left-hand operator.
\end{proof}

\begin{corollary}[Uniform feedback preserves the density exponent]
\label{cor:n9v61-feedback-scaling}
If the unresummed one-mark source has the half-density scale
\eqref{eq:n9v61-half-density-scale}, all second derivatives of
${\cal P}$ are uniformly bounded in the corresponding marked norms, and
$\sup_n\theta_n<1$, then feedback changes only the constant in
\eqref{eq:n9v61-cell-volume-output}.  If $\theta_n\uparrow1$, the factor
$(1-\theta_n)^{-1}$ can destroy the continuum amplitude bound.
\end{corollary}

\begin{proof}
Apply \eqref{eq:n9v61-feedback-denominator} to each mark.  In the second
derivative every term contains either one two-mark input or a product of
two one-mark inputs, hence carries $(a_n/\ell_n)^d$.  Uniform resolvent
and derivative constants preserve that exponent.
\end{proof}

\subsection{Scope: positivity and the chiral phase}

\begin{firewall}[Positive paired sector only]
\label{fire:n9v61-chiral}
The probability disintegration and absolute KP majorant above apply to
a positive gauge--Higgs or conjugate-paired fermion sector.  For the
physical chiral determinant, the local phase must satisfy the V53--V56
nonvanishing analytic branch and marked connected-norm conditions.
Replacing the complex conditional weight by its modulus changes
$U_g$ and does not prove the physical coarse response.
\end{firewall}

\begin{definition}[Exact marked coarse pushforward card (EMCPC)]
\label{def:n9v61-EMCPC}
EMCPC requires:
\begin{enumerate}
\item the exact disintegration \eqref{eq:n9v61-disintegration} and
coarse density \eqref{eq:n9v61-exact-coarse-law};
\item complete conditional and coarse reference scores with the type
separation of \Cref{cor:n9v61-reference-owner};
\item an exact conditional factorization or an enlarged sector skeleton;
\item a uniform KP margin for all bounded physical activities;
\item the half-density one- and two-mark estimates;
\item a feedback derivative norm uniformly below one;
\item uniformity over volume, mesh, chart, and accepted coarse boundary;
and
\item a positive paired law or a separately controlled analytic chiral
phase.
\end{enumerate}
\end{definition}

\begin{theorem}[EMCPC loads the marked coarse kernel into UMTC]
\label{thm:n9v61-EMCPC}
Under EMCPC, one exact coarse step produces the conditional effective
potential \eqref{eq:n9v61-coarse-potential} with the exact response
\eqref{eq:n9v61-U-first}--\eqref{eq:n9v61-U-second}.  Its connected
two-source kernel obeys the cell-volume rooted bound
\eqref{eq:n9v61-cell-volume-output}; uniform feedback changes only its
constant.  Hence the step supplies the continuum response envelope
needed by UMTC.
\end{theorem}

\begin{proof}
Use \Cref{thm:n9v61-exact-pushforward,
thm:n9v61-marked-derivatives,prop:n9v61-polymer-representation,
thm:n9v61-marked-KP,cor:n9v61-feedback-scaling}.  Apply
\Cref{prop:n9v60-rooted-to-envelope}.
\end{proof}

\begin{assessment}[Progress at ninth-series V61]
\label{ass:n9v61-status}
The one-step coarse law and its metric-source derivatives are now exact.
The missing continuum estimate has been isolated to a local
half-density mark bound plus a uniform feedback margin.  This is sharper
than asking only for KP convergence: the two marked sources must carry
the cell-volume exponent required by V59.
\end{assessment}

\begin{programme}[Ninth-series V62: heat-kernel metric mark]
\label{prog:n9v61-V62}
V62 will compute the conditional metric derivative of a compact
heat-kernel reference by Duhamel differentiation, separate its centred
score from the coarse pushforward score, and test the half-density bound
in the reference energy norm.  This goes upstream to the local estimate
that V61 leaves open.
\end{programme}

\begin{firewall}[Claim boundary after ninth-series V61]
\label{fire:n9v61-boundary}
V61 proves exact identities and conditional compilers.  It does not
prove the physical block map has the required factorization, the KP or
feedback margins, the half-density mark bound, the chiral phase control,
anomaly cancellation, reflection positivity through the coarse step,
OS reconstruction, or the full SM--Einstein continuum.
\end{firewall}

\begin{theorem}[Ninth-series V61 result]
\label{thm:n9v61-result}
For one exact finite gauge--Higgs coarse step, the coarse interaction is
the conditional logarithm \eqref{eq:n9v61-coarse-potential}, with exact
metric mean and covariance responses
\eqref{eq:n9v61-U-first}--\eqref{eq:n9v61-U-second}.  Under a uniform
marked KP margin, two half-density metric marks yield the required
cell-volume coefficient \eqref{eq:n9v61-cell-volume-output}; a uniform
feedback resolvent preserves that exponent.
\end{theorem}

\begin{proof}
Combine \Cref{thm:n9v61-exact-pushforward,
thm:n9v61-marked-derivatives,thm:n9v61-marked-KP,
prop:n9v61-feedback,thm:n9v61-EMCPC}.
\end{proof}

\paragraph{Parent-version record.}
V61 was formed from
\texttt{Renewal\_SM\_Einstein\_Continuum\_Research\_Notes\_V60.tex}.
The V60 parent had $22\,327$ logical source lines, $804\,577$ bytes, and
SHA--256
\[
 \texttt{9B0349BA0B4AD870CAF87BDF67068F7B2557C232CD5A3AEF257EEBDE12B05E50}.
\]

% =============================================================================
% END APPENDED NINTH-SERIES V61 MATERIAL
% =============================================================================
% =============================================================================
% BEGIN APPENDED NINTH-SERIES V62 MATERIAL
% =============================================================================

\section{Ninth-series V62: heat-kernel metric marks and the ultraviolet Fisher row}
\label{sec:v9v62}

\subsection{Exact Duhamel differentiation of the reference}

V61 left the half-density estimate for the conditional reference score
open.  A compact heat-kernel reference depends on the metric through its
principal electric generator, so its derivative is a quadratic-form
insertion rather than a bounded scalar activity.  This note computes the
exact derivative and shows why its raw microscopic Fisher row is not
small under time subdivision.

Let ${\cal K}$ be a Hilbert space and let $H_g\ge0$ be the self-adjoint
electric/heat-kernel generator at metric $g$.  At $\bar g$ put
\begin{equation}
 H=H_{\bar g},
 \qquad B_h=D_gH_g\big|_{\bar g}[h].
 \label{eq:n9v62-generator-mark}
\end{equation}

\begin{theorem}[Semigroup Duhamel derivative]
\label{thm:n9v62-Duhamel}
If $B_h$ is a symmetric form on ${\cal D}((H+\sigma)^{1/2})$ with
\begin{equation}
 |B_h(u,v)|\le b_h
 \lVert(H+\sigma)^{1/2}u\rVert
 \lVert(H+\sigma)^{1/2}v\rVert,
 \label{eq:n9v62-form-bound}
\end{equation}
then, in the weak form sense,
\begin{equation}
 D_ge^{-tH_g}\big|_{\bar g}[h]
 =-\int_0^t e^{-(t-r)H}B_he^{-rH}\,dr.
 \label{eq:n9v62-Duhamel}
\end{equation}
Moreover,
\begin{equation}
 \left\lVert D_ge^{-tH_g}[h]\right\rVert
 \le b_h\left{
 \frac{\pi}{2e}+4\sqrt{\frac{\sigma t}{2e}}+\sigma t
 \right}.
 \label{eq:n9v62-Duhamel-bound}
\end{equation}
\end{theorem}

\begin{proof}
Write $A=H+\sigma$ and represent the form as
$B_h=A^{1/2}T_hA^{1/2}$ with $\lVert T_h\rVert\le b_h$.  The standard
form perturbation identity gives \eqref{eq:n9v62-Duhamel}.  Spectral
calculus gives, for $r>0$,
\begin{equation}
 \lVert A^{1/2}e^{-rH}\rVert
 \le (2er)^{-1/2}+\sqrt\sigma.
 \label{eq:n9v62-smoothing}
\end{equation}
Insert the form factorization in the integrand.  The product of the two
right-hand sides of \eqref{eq:n9v62-smoothing} integrates to
$\pi/(2e)+4\sqrt{\sigma t/(2e)}+\sigma t$.
\end{proof}

\begin{corollary}[Principal and bounded marks scale differently]
\label{cor:n9v62-principal-bounded}
A bounded scalar perturbation $B_h$ obeys
\begin{equation}
 \lVert D e^{-tH}[h]\rVert\le t\lVert B_h\rVert,
 \label{eq:n9v62-bounded-t}
\end{equation}
whereas the general principal-form estimate
\eqref{eq:n9v62-Duhamel-bound} has a nonzero $t\downarrow0$ bound.
Thus the electric metric mark cannot be assigned the small-time factor
of a bounded Wilson multiplier.
\end{corollary}

\begin{proof}
For bounded $B_h$, take norms directly in
\eqref{eq:n9v62-Duhamel}.  The comparison with
\eqref{eq:n9v62-Duhamel-bound} is immediate.
\end{proof}

\subsection{Kernel and bridge scores}

Let $G$ be a compact connected Lie group and let $K_t^g(x,y)>0$ be the
heat kernel of $H_g$.  Kernel evaluation of
\eqref{eq:n9v62-Duhamel} gives
\begin{equation}
 D_gK_t^g(x,y)\big|_{\bar g}[h]
 =-\int_0^t
 \bigl(e^{-(t-r)H}B_he^{-rH}\bigr)(x,y),dr.
 \label{eq:n9v62-kernel-Duhamel}
\end{equation}
Hence the endpoint mark is
\begin{equation}
 q_{t,h}(x,y)=D_g\log K_t^g(x,y)\big|_{\bar g}[h]
 =-\frac1{K_t(x,y)}\int_0^t
 \bigl(e^{-(t-r)H}B_he^{-rH}\bigr)(x,y),dr.
 \label{eq:n9v62-endpoint-mark}
\end{equation}

Divide $t$ into $m$ equal pieces, $\delta=t/m$, and condition the
heat-kernel path on endpoints $x_0,x_m$.  Its intermediate-point density
is
\begin{equation}
 P_g(dx_1\cdots dx_{m-1}\mid x_0,x_m)
 =\frac{\prod_{j=1}^mK_\delta^g(x_{j-1},x_j)}
 {K_t^g(x_0,x_m)}\prod_{j=1}^{m-1}dx_j.
 \label{eq:n9v62-discrete-bridge}
\end{equation}

\begin{proposition}[Exact centred bridge score]
\label{prop:n9v62-bridge-score}
The conditional reference score in V61 is
\begin{equation}
 b_h^{(m)}=
 \sum_{j=1}^mq_{\delta,h}(x_{j-1},x_j)
 -q_{t,h}(x_0,x_m),
 \label{eq:n9v62-bridge-score}
\end{equation}
and
\begin{equation}
 \mathbb E_P b_h^{(m)}=0.
 \label{eq:n9v62-bridge-centred}
\end{equation}
The subtracted endpoint score is precisely the coarse-reference owner
in the disintegration.
\end{proposition}

\begin{proof}
Differentiate the logarithm of
\eqref{eq:n9v62-discrete-bridge}.  This gives
\eqref{eq:n9v62-bridge-score}.  Differentiating the identity
$\int dP_g=1$ gives \eqref{eq:n9v62-bridge-centred}.
\end{proof}

\begin{firewall}[Continuous-time Radon--Nikodym shortcut]
\label{fire:n9v62-singular-paths}
Changing the principal diffusion metric can make the corresponding
continuous-time path measures mutually singular.  The finite-$m$ score
\eqref{eq:n9v62-bridge-score} is exact, but one may not assume it has an
$L^2$ limit as $m\to\infty$.  The form-Duhamel identity remains the safe
continuum object.
\end{firewall}

\subsection{Exact Gaussian bridge calculation}

The tangent Gaussian model already decides the ultraviolet scaling.  In
$\mathbb R^D$, take independent increments
\begin{equation}
 \Delta_j\sim N(0,2\kappa\delta I_D),
 \qquad j=1,\ldots,m,
 \label{eq:n9v62-Gaussian-increments}
\end{equation}
and condition on $\sum_j\Delta_j=Y$.  Use the logarithmic diffusion
parameter $\theta=\log\kappa$.

\begin{theorem}[Gaussian bridge Fisher identity]
\label{thm:n9v62-Gaussian-Fisher}
The conditional score is
\begin{equation}
 S_m=-\frac{D(m-1)}2
 +\frac1{4\kappa\delta}
  \sum_{j=1}^m\left|\Delta_j-\frac Ym\right|^2,
 \label{eq:n9v62-Gaussian-score}
\end{equation}
and satisfies
\begin{equation}
 \mathbb E(S_m\mid Y)=0,
 \qquad
 \operatorname{Var}(S_m\mid Y)=\frac{D(m-1)}2.
 \label{eq:n9v62-Gaussian-Fisher}
\end{equation}
In particular, the conditional Fisher information grows linearly with
the number of time subdivisions.
\end{theorem}

\begin{proof}
On the affine constraint $\sum_j\Delta_j=Y$, the residual vector
$(\Delta_j-Y/m)_{j=1}^m$ is Gaussian on a subspace of dimension
$\nu=D(m-1)$ with covariance $2\kappa\delta$ times the identity on that
subspace.  Differentiating its density with respect to
$\theta=\log\kappa$ gives \eqref{eq:n9v62-Gaussian-score}.  Moreover,
\begin{equation}
 \frac1{2\kappa\delta}
 \sum_j\left|\Delta_j-\frac Ym\right|^2\sim\chi^2_\nu.
 \end{equation}
The mean and variance of $\chi^2_\nu$ are $\nu$ and $2\nu$, proving
\eqref{eq:n9v62-Gaussian-Fisher}.
\end{proof}

\begin{corollary}[Endpoint owner plus conditional owner]
\label{cor:n9v62-Gaussian-split}
If the endpoint $Y$ is also sampled from
$N(0,2\kappa tI_D)$, its score has variance $D/2$ and is independent of
the residual bridge score.  Hence the full $m$-increment score has
variance $Dm/2$, split exactly into the coarse endpoint owner $D/2$ and
the conditional owner $D(m-1)/2$.
\end{corollary}

\begin{proof}
The Gaussian sample mean and residual subspace are independent.  The
endpoint score is $-D/2+|Y|^2/(4\kappa t)$ and has variance $D/2$.
Add the independent variances.
\end{proof}

\begin{proposition}[Compact-group small-time diagnostic]
\label{prop:n9v62-compact-diagnostic}
For a smooth uniformly elliptic heat kernel on a compact
$D$-dimensional group, the one-step logarithmic diffusion score has
Fisher information tending to $D/2$ as $t\downarrow0$, locally away
from the cut locus.  A bridge with $m$ microscopic substeps therefore
has the Gaussian leading term $D(m-1)/2$.
\end{proposition}

\begin{proof}
In exponential coordinates the heat-kernel parametrix is
\begin{equation}
 K_t(e,\exp X)
 =(4\pi\kappa t)^{-D/2}
 e^{-|X|^2/(4\kappa t)}
 \bigl(j(X)^{-1/2}+O(t)\bigr).
 \label{eq:n9v62-parametrix}
\end{equation}
Under the heat-kernel law, $X=O_{L^q}(\sqrt t)$ for every fixed $q$;
the complement of a normal neighbourhood is exponentially small.
Differentiating the parametrix in $\log\kappa$ gives the Gaussian score
plus an $L^2$ error tending to zero.  Apply the same argument on the
$D(m-1)$-dimensional local bridge subspace for fixed $m$.
\end{proof}

\begin{firewall}[Order of limits]
\label{fire:n9v62-order-limits}
The last proposition is uniform for fixed $m$ as $t\downarrow0$; it is
not a uniform $m\to\infty$ theorem.  The exact Gaussian identity shows
that no such uniform Fisher bound should be expected for a principal
diffusion variation.
\end{firewall}

\subsection{Failure of the raw half-density estimate}

Let a physical correlation block contain
\begin{equation}
 N_n\asymp(\ell_n/a_n)^d
 \label{eq:n9v62-block-count}
\end{equation}
microscopic reference cells.

\begin{proposition}[Microscopic Fisher no-go]
\label{prop:n9v62-half-density-no-go}
Suppose a unit normalized microscopic metric mark has Fisher information
at least $c_0>0$, uniformly in $n$.  If $\ell_n/a_n\to\infty$, that raw
mark cannot satisfy the V61 half-density estimate
\begin{equation}
 \operatorname{Var}(b_h)le C(a_n/\ell_n)^d.
 \label{eq:n9v62-half-density-target}
\end{equation}
The Gaussian heat-kernel scale mark satisfies the hypothesis with
$c_0=D/2$ for one increment.
\end{proposition}

\begin{proof}
The right-hand side of \eqref{eq:n9v62-half-density-target} tends to
zero, contradicting the uniform lower bound $c_0$.  The last statement
is \Cref{thm:n9v62-Gaussian-Fisher} with one unconditioned increment.
\end{proof}

\begin{assessment}[Meaning of the no-go]
\label{ass:n9v62-no-go-meaning}
The raw principal heat-kernel score contains an ultraviolet local Fisher
row.  It must be assigned to local gravitational owners and
renormalized, or the coarse mark must be normalized as a block average.
KP decay of physical multipliers cannot remove this reference row.
\end{assessment}

\subsection{Block averaging recovers the exponent}

The required exponent does arise from a normalized average.  Let
$S_i$, $i=1,\ldots,N$, be centred cell scores with covariance row sum
\begin{equation}
 \sup_i\sum_j|\operatorname{Cov}(S_i,S_j)|\le C_0.
 \label{eq:n9v62-score-row}
\end{equation}

\begin{lemma}[Averaged score half-density]
\label{lem:n9v62-average}
For $\overline S_N=N^{-1}\sum_{i=1}^NS_i$,
\begin{equation}
 \operatorname{Var}(\overline S_N)\le\frac{C_0}{N}.
 \label{eq:n9v62-average-variance}
\end{equation}
With \eqref{eq:n9v62-block-count}, this is precisely
$O((a_n/\ell_n)^d)$.
\end{lemma}

\begin{proof}
\begin{align}
 \operatorname{Var}(\overline S_N)
 &\le N^{-2}\sum_{i,j}|\operatorname{Cov}(S_i,S_j)|
 \le N^{-2}\sum_iC_0=C_0/N.
\end{align}
\end{proof}

\begin{firewall}[Average versus total source]
\label{fire:n9v62-average-total}
The derivative of an action with respect to a constant block metric is
normally a sum of cell stresses, not their arithmetic average.  The
factor $N^{-1}$ in \Cref{lem:n9v62-average} is legitimate only if it is
produced by the Riesz normalization against the coarse positive metric
reference or by the declared block variable.  Inserting it by hand
would change the Einstein source.
\end{firewall}

\subsection{Renormalized reference-mark card}

\begin{definition}[Heat-kernel metric-mark card (HKMMC)]
\label{def:n9v62-HKMMC}
HKMMC requires:
\begin{enumerate}
\item the exact Duhamel form derivative
\eqref{eq:n9v62-Duhamel};
\item the exact conditional/coarse score split
\eqref{eq:n9v62-bridge-score};
\item a local heat-kernel expansion assigning ultraviolet Fisher rows to
the cosmological, Newton, and allowed curvature owners of RZMOC;
\item an $L^2$ or form estimate for the renormalized nonlocal marked
remainder;
\item an exact Riesz normalization or coarse block variable producing
the half-density scale, if block averaging is used;
\item a uniform covariance row sum for those normalized marks; and
\item no continuous-path Radon--Nikodym assumption for a varying
principal diffusion metric.
\end{enumerate}
\end{definition}

\begin{theorem}[HKMMC supplies the reference part of EMCPC]
\label{thm:n9v62-HKMMC}
Under HKMMC, the heat-kernel metric derivative remains in its exact form
owner, its local ultraviolet Fisher contribution is counted once in the
RZMOC counterterm ledger, and its normalized connected remainder obeys
the V61 half-density mark estimate.  Hence it can be combined with the
bounded physical activity marks without treating a principal generator
as a scalar polymer multiplier.
\end{theorem}

\begin{proof}
The first two claims are \Cref{thm:n9v62-Duhamel,
prop:n9v62-bridge-score}.  HKMMC item 3 implements the V58 ownership.
Items 4--6 and \Cref{lem:n9v62-average} give the half-density estimate.
The V61 marked KP theorem then combines this marked observable with the
bounded activity cluster expansion.
\end{proof}

\begin{assessment}[Progress at ninth-series V62]
\label{ass:n9v62-status}
The reference-mark bottleneck is no longer hidden inside a KP constant.
A principal heat-kernel metric variation has order-one Fisher
information per microscopic increment, and a conditional bridge score
grows under time subdivision.  The viable route is Duhamel/form
differentiation, local gravitational ownership of ultraviolet rows, and
an exact coarse Riesz normalization for the remaining averaged mark.
\end{assessment}

\begin{programme}[Ninth-series V63: coarse Riesz normalization]
\label{prog:n9v62-V63}
V63 will compute the block scaling of the positive Einstein metric
reference and its dual stress-source norm.  It will decide when the
Riesz representative of a block-averaged stress really supplies the
$N^{-1/2}$ one-mark factor, and will expose any dimensional obstruction
for four-dimensional metric fluctuations.
\end{programme}

\begin{firewall}[Claim boundary after ninth-series V62]
\label{fire:n9v62-boundary}
V62 proves the Duhamel and Gaussian Fisher statements and identifies the
needed renormalized averaging.  It does not construct the compact
gauge--Higgs local counterterms, prove the coarse Riesz scaling, control
the chiral phase or anomaly, establish the strict margin, construct the
slow infrared theory, prove reflection positivity, OS reconstruction,
or the full SM--Einstein continuum.
\end{firewall}

\begin{theorem}[Ninth-series V62 result]
\label{thm:n9v62-result}
The exact heat-kernel reference mark is the Duhamel form insertion
\eqref{eq:n9v62-Duhamel}, with conditional bridge score
\eqref{eq:n9v62-bridge-score}.  In the Gaussian tangent model its
conditional Fisher information is exactly $D(m-1)/2$, so raw
microscopic marks cannot satisfy the V61 half-density estimate at a
macroscopic correlation range.  A normalized block average recovers the
required factor under the covariance row bound
\eqref{eq:n9v62-score-row}.
\end{theorem}

\begin{proof}
Combine \Cref{thm:n9v62-Duhamel,prop:n9v62-bridge-score,
thm:n9v62-Gaussian-Fisher,prop:n9v62-half-density-no-go,
lem:n9v62-average,thm:n9v62-HKMMC}.
\end{proof}

\paragraph{Parent-version record.}
V62 was formed from
\texttt{Renewal\_SM\_Einstein\_Continuum\_Research\_Notes\_V61.tex}.
The V61 parent had $22\,776$ logical source lines, $821\,499$ bytes, and
SHA--256
\[
 \texttt{1F450A0E0CF60E71E1F7D7192B0D6DC6C22330C31181E9BB65BCE6FB0EC59D62}.
\]

% =============================================================================
% END APPENDED NINTH-SERIES V62 MATERIAL
% =============================================================================
% =============================================================================
% BEGIN APPENDED NINTH-SERIES V63 MATERIAL
% =============================================================================

\section{Ninth-series V63: coarse Riesz normalization for a massless Einstein reference}
\label{sec:v9v63}

\subsection{The reference is a gradient seminorm}

V62 showed that an averaged microscopic score can have
$N^{-1/2}$ fluctuation size.  That fact does not by itself provide the
V61 normalized metric mark: the physical Einstein reference is massless
and controls derivatives of the metric perturbation, not its constant
block value.  The correct source norm is therefore an $H^{-1}$ norm
after gauge and zero-mode ownership.

On a flat $d$-torus with mesh $a$, let ${\cal H}_a$ be a physical
gauge-reduced tensor fibre and define
\begin{equation}
 {\mathfrak a}_a(h)
 =Z_a a^d\sum_{e}|\nabla_ah(e)|^2,
 \qquad
 \nabla_ah(e)=\frac{h(e_+)-h(e_-)}a,
 \label{eq:n9v63-Einstein-reference}
\end{equation}
where $Z_a>0$ is the renormalized Einstein coefficient.  Constants lie
in the kernel and are treated by the RZMOC background condition.

\begin{definition}[Massless Riesz dual norm]
\label{def:n9v63-dual}
For a linear source $F$ annihilating $\ker{\mathfrak a}_a$, set
\begin{equation}
 \lVert F\rVert_{-1,a}^2
 =\sup_{h:{\mathfrak a}_a(h)>0}
 \frac{|F(h)|^2}{{\mathfrak a}_a(h)}.
 \label{eq:n9v63-dual-norm}
\end{equation}
\end{definition}

\begin{lemma}[Constant-mode compatibility]
\label{lem:n9v63-constant-mode}
If a linear source $F$ does not annihilate a constant physical tensor
$h_0$, then it has no continuous extension to the Hilbert completion of
the quotient defined by \eqref{eq:n9v63-Einstein-reference}.  In
particular, its massless Riesz norm is infinite or undefined.
\end{lemma}

\begin{proof}
The constant $h_0$ represents the zero vector in the reference
seminorm, while $F(h_0)\ne0$.  Thus $F$ is not well-defined on the
quotient.  Equivalently, adding arbitrary multiples of $h_0$ changes
$F$ without changing ${\mathfrak a}_a$.
\end{proof}

\begin{firewall}[No local inverse for the constant metric]
\label{fire:n9v63-no-local-inverse}
The block matrices $A_x^{-1/2}$ used in the coercive V57 Schur test do
not exist for the pure massless constant mode.  A collar Poincar\'e
condition, global mean-zero condition, curvature potential, or explicit
infrared block must be declared before a block Riesz representative is
formed.
\end{firewall}

\subsection{Mean-zero sources and the Poincar\'e scale}

Let $Q_\ell$ be a cube of side $\ell$ and let $h$ have zero mean on
$Q_\ell$ or vanish on its boundary collar.  The discrete Poincar\'e
inequality reads
\begin{equation}
 a^d\sum_{x\in Q_\ell}|h(x)|^2
 \le C_P\ell^2 a^d\sum_{e\subset Q_\ell}|\nabla_ah(e)|^2.
 \label{eq:n9v63-Poincare}
\end{equation}

\begin{proposition}[Local $L^2$ stress to massless dual norm]
\label{prop:n9v63-L2-dual}
Suppose a centred random stress $T_a(x)$ satisfies
\begin{equation}
 \operatorname{Var}\left(a^d\sum_{x\in Q_\ell}
 T_a(x):h(x)\right)
 \le \rho_a a^d\sum_{x\in Q_\ell}|h(x)|^2.
 \label{eq:n9v63-white-row}
\end{equation}
Then on the mean-zero or collar subspace,
\begin{equation}
 \operatorname{Var}\left(a^d\sum_xT_a(x):h(x)\right)
 \le\frac{C_P\rho_a\ell^2}{Z_a}{\mathfrak a}_a(h).
 \label{eq:n9v63-Hminus1-bound}
\end{equation}
\end{proposition}

\begin{proof}
Apply \eqref{eq:n9v63-white-row}, then
\eqref{eq:n9v63-Poincare}, and identify
\eqref{eq:n9v63-Einstein-reference}.
\end{proof}

\begin{remark}[The physical scaling coefficient]
\label{rem:n9v63-physical-scaling}
The massless relative coefficient is
$C_P\rho_a\ell^2/Z_a$.  The factor $(a/\ell)^d$ from averaging is only
one possible input into $\rho_a$; it is not the complete dimensionless
Einstein comparison.  The Newton coefficient and the Poincar\'e length
must also be tracked.
\end{remark}

\subsection{Averaging does not remove momentum zero}

Let $w_\ell$ be a normalized block-average kernel,
$\sum_xa^dw_\ell(x)=1$, and let
\begin{equation}
 \widehat w_\ell(p)=a^d\sum_xw_\ell(x)e^{-ip\cdot x}.
 \label{eq:n9v63-window-symbol}
\end{equation}

\begin{proposition}[Block averaging preserves the zero mode]
\label{prop:n9v63-average-zero}
For every normalized block average,
\begin{equation}
 \widehat w_\ell(0)=1.
 \label{eq:n9v63-window-zero}
\end{equation}
If a translation-invariant source covariance has symbol
$\widehat C_a(p)$, the averaged covariance has symbol
\begin{equation}
 \widehat C_a^{\rm av}(p)
 =|\widehat w_\ell(p)|^2\widehat C_a(p),
 \qquad
 \widehat C_a^{\rm av}(0)=\widehat C_a(0).
 \label{eq:n9v63-averaged-symbol}
\end{equation}
Thus averaging cannot repair a nonzero covariance mass term relative to
the Einstein $p^2$ reference.
\end{proposition}

\begin{proof}
The normalization gives \eqref{eq:n9v63-window-zero}.  Convolution by
$w_\ell$ is multiplication by $\widehat w_\ell$ in Fourier space, so a
covariance acquires its modulus squared.  Evaluate at $p=0$ and apply
the V57 Fourier no-go when the result is nonzero.
\end{proof}

\begin{corollary}[Correct order of operations]
\label{cor:n9v63-order}
Zero-mode ownership or cancellation must precede the use of the
$N^{-1/2}$ block-average gain.  The averaging estimate in V62 controls
fluctuation density after that compatibility condition; it does not
establish the condition.
\end{corollary}

\subsection{Gradient-factorized sources}

A source with an exact discrete divergence automatically annihilates
constants.  Let $J_a(e)$ be a centred random edge tensor and define
\begin{equation}
 F_{J,a}(h)=a^d\sum_eJ_a(e):\nabla_ah(e).
 \label{eq:n9v63-gradient-source}
\end{equation}

\begin{theorem}[Gradient-source Riesz bound]
\label{thm:n9v63-gradient-source}
Suppose the edge covariance obeys
\begin{equation}
 \operatorname{Var}\left(a^d\sum_eJ_a(e):v(e)\right)
 \le\chi_a a^d\sum_e|v(e)|^2
 \label{eq:n9v63-edge-row}
\end{equation}
for every edge field $v$.  Then $F_{J,a}$ annihilates constants and
\begin{equation}
 \mathbb E\lVert F_{J,a}\rVert_{-1,a}^2
 \le\frac{\chi_a}{Z_a}.
 \label{eq:n9v63-gradient-dual}
\end{equation}
Equivalently,
\begin{equation}
 \operatorname{Var}(F_{J,a}(h))
 \le\frac{\chi_a}{Z_a}{\mathfrak a}_a(h).
 \label{eq:n9v63-gradient-relative}
\end{equation}
\end{theorem}

\begin{proof}
For constant $h$, $\nabla_ah=0$.  Use
\eqref{eq:n9v63-edge-row} with $v=\nabla_ah$ and compare with
\eqref{eq:n9v63-Einstein-reference}.  Taking the supremum in
\eqref{eq:n9v63-dual-norm} gives the result.
\end{proof}

\begin{proposition}[Fourier factorization criterion]
\label{prop:n9v63-Fourier-factor}
Let $C_a\ge0$ be a translation-invariant covariance on the physical
fibre.  If
\begin{equation}
 0\le\widehat C_a(p)\le\chi_a|\widehat\nabla_a(p)|^2I,
 \label{eq:n9v63-symbol-gradient}
\end{equation}
then there is a bounded Fourier multiplier $M_a(p)$ with
\begin{equation}
 \widehat C_a(p)^{1/2}=M_a(p)\widehat\nabla_a(p),
 \qquad
 \lVert M_a(p)\rVert\le\sqrt{\chi_a},
 \label{eq:n9v63-square-root-factor}
\end{equation}
on the nonzero momentum fibres, extended by zero at $p=0$.  Hence the
covariance is representable as a gradient-factorized Gaussian source
for purposes of the quadratic-form bound.
\end{proposition}

\begin{proof}
For $p\ne0$, define
$M_a(p)=\widehat C_a(p)^{1/2}
\widehat\nabla_a(p)^*/|\widehat\nabla_a(p)|^2$ on the scalar gradient
direction, with the corresponding tensor contraction.  Inequality
\eqref{eq:n9v63-symbol-gradient} bounds its norm by
$\sqrt{\chi_a}$.  At $p=0$ the inequality gives
$\widehat C_a(0)=0$, so the zero extension is consistent.  Multiplying
by the gradient symbol proves \eqref{eq:n9v63-square-root-factor}.
\end{proof}

\begin{firewall}[Conservation is not the factorization]
\label{fire:n9v63-conservation}
A stress Ward identity such as $p^\mu\widehat C_{\mu\nu,\rho\sigma}=0$
annihilates longitudinal gauge directions.  It does not by itself imply
the full physical estimate \eqref{eq:n9v63-symbol-gradient}.  The latter
also requires the transverse zero mode and its $p^2$ coefficient to be
controlled after RZMOC ownership.
\end{firewall}

\subsection{Block fluctuation scaling with the Einstein norm}

Let $Q_\ell$ contain $N\asymp(\ell/a)^d$ microscopic cells.  Suppose
the renormalized stress density obeys \eqref{eq:n9v63-white-row} with
\begin{equation}
 \rho_a\le\rho_0(a/\ell)^d.
 \label{eq:n9v63-density-gain}
\end{equation}

\begin{corollary}[Averaged massless coefficient]
\label{cor:n9v63-averaged-massless}
On a mean-zero/collar block,
\begin{equation}
 \operatorname{Var}(F(h))
 \le
 \frac{C_P\rho_0}{Z_a}
 a^d\ell^{2-d}{\mathfrak a}_a(h).
 \label{eq:n9v63-dimensional-coefficient}
\end{equation}
In $d=4$ the coefficient is
$C_P\rho_0a^4/(Z_a\ell^2)$.  Uniform smallness therefore depends on the
renormalized scaling of $Z_a$ and on the relation between $a$ and
$\ell$; the central-limit factor alone is not the final criterion.
\end{corollary}

\begin{proof}
Insert \eqref{eq:n9v63-density-gain} into
\eqref{eq:n9v63-Hminus1-bound} and simplify.
\end{proof}

\begin{remark}[Dimensionless rescaling]
\label{rem:n9v63-dimensionless}
If coarse coordinates and $Z_a$ are rescaled so that the block Einstein
form is dimensionless and of order one, the expression in
\eqref{eq:n9v63-dimensional-coefficient} may reduce to the V61
cell-volume factor.  Such a rescaling must be written explicitly; it is
not contained in the phrase ``unit metric mark.''
\end{remark}

\subsection{Coarse Einstein Riesz card}

\begin{definition}[Massless Einstein Riesz scaling card (MERSC)]
\label{def:n9v63-MERSC}
MERSC requires:
\begin{enumerate}
\item the exact physical gauge-reduced Einstein form and its coefficient
$Z_a$;
\item removal or explicit retention of every constant/infrared kernel
mode;
\item the source covariance in the corresponding $H^{-1}$ dual norm;
\item either a Poincar\'e/collar estimate with the physical length
$\ell$, or the gradient-symbol estimate
\eqref{eq:n9v63-symbol-gradient};
\item separate accounting of block averaging, Newton scaling, and the
Poincar\'e factor;
\item a controlled nonlocal constraint projector as required by V57;
and
\item a regulator-uniform aggregate coefficient below the CPRAC margin.
\end{enumerate}
\end{definition}

\begin{theorem}[MERSC replaces naive half-density normalization on the massless sector]
\label{thm:n9v63-MERSC}
Under MERSC, the marked fast source is a continuous functional of the
massless physical metric energy.  Its relative coefficient is given by
\eqref{eq:n9v63-Hminus1-bound} on collar blocks or
\eqref{eq:n9v63-gradient-relative} under exact gradient factorization.
Block averaging may improve the stress-density constant, but it is used
only after constant-mode compatibility.
\end{theorem}

\begin{proof}
Use \Cref{lem:n9v63-constant-mode,prop:n9v63-L2-dual,
prop:n9v63-average-zero,thm:n9v63-gradient-source,
cor:n9v63-averaged-massless}.  MERSC item 7 inserts the resulting bound
into the V57 aggregate margin.
\end{proof}

\begin{assessment}[Progress at ninth-series V63]
\label{ass:n9v63-status}
The half-density target has been split correctly.  It remains a useful
sufficient condition on coercive local blocks, but the massless Einstein
sector requires an $H^{-1}$ comparison.  Averaging reduces fluctuation
density while preserving momentum zero; cancellation or explicit
infrared ownership must come first.
\end{assessment}

\begin{programme}[Ninth-series V64: multiscale massless covariance]
\label{prog:n9v63-V64}
V64 will construct a finite-range or annular scale decomposition of the
massless physical metric covariance.  It will keep the true infrared
block explicit and apply the coercive Schur/marked-cluster machinery only
on individual momentum shells, where the reference has a scale-dependent
gap.
\end{programme}

\begin{firewall}[Claim boundary after ninth-series V63]
\label{fire:n9v63-boundary}
V63 proves the Riesz and averaging criteria.  It does not prove the
physical source satisfies the required $H^{-1}$ estimate, construct the
constraint-reduced covariance, tune $Z_a$, control the chiral phase or
anomaly, prove the strict aggregate margin, establish reflection
positivity, OS reconstruction, or the full SM--Einstein continuum.
\end{firewall}

\begin{theorem}[Ninth-series V63 result]
\label{thm:n9v63-result}
A massless Einstein reference normalizes sources in $H^{-1}$, not by a
blockwise inverse.  A nonzero constant source is incompatible with that
norm, and normalized averaging does not alter its zero mode.  After
RZMOC cancellation, the Poincar\'e estimate
\eqref{eq:n9v63-Hminus1-bound} or the gradient estimate
\eqref{eq:n9v63-gradient-relative} gives the correct relative
coefficient, with block averaging, $Z_a$, and the physical length tracked
separately.
\end{theorem}

\begin{proof}
Combine \Cref{lem:n9v63-constant-mode,prop:n9v63-L2-dual,
prop:n9v63-average-zero,thm:n9v63-gradient-source,
prop:n9v63-Fourier-factor,thm:n9v63-MERSC}.
\end{proof}

\paragraph{Parent-version record.}
V63 was formed from
\texttt{Renewal\_SM\_Einstein\_Continuum\_Research\_Notes\_V62.tex}.
The V62 parent had $23\,189$ logical source lines, $836\,565$ bytes, and
SHA--256
\[
 \texttt{D313C9CB0BE50DF3382B161D67F7BE824DCEB557CB3932350C9F0BE64797B8BB}.
\]

% =============================================================================
% END APPENDED NINTH-SERIES V63 MATERIAL
% =============================================================================
% =============================================================================
% BEGIN APPENDED NINTH-SERIES V64 MATERIAL
% =============================================================================

\section{Ninth-series V64: annular decomposition of the massless Einstein covariance}
\label{sec:v9v64}

\subsection{Exact spectral separation of the infrared block}

V63 replaced a nonexistent blockwise inverse by the physical
$H^{-1}$ source norm.  The remaining task is to use the coercive
machinery without inserting a false global mass.  At finite regulator,
let ${\cal H}_{\rm phys}$ be the constraint-reduced metric tangent
space and let $A\ge0$ be the self-adjoint positive Einstein reference on
that space.  Its kernel has already been assigned to gauge, background,
or global-modulus owners by RZMOC and MERSC.

Fix $0<\lambda_{
m IR}<\lambda_{
m UV}$ and define
\begin{equation}
 P_{<}=\mathbf 1_{(0,\lambda_{
m IR})}(A),
 \qquad
 P_{\ge}=\mathbf 1_{[\lambda_{
m IR},\lambda_{
m UV}]}(A).
 \label{eq:n9v64-IR-UV-projections}
\end{equation}
The zero eigenspace is omitted from the pseudoinverse and retained in
its declared owner.

\begin{proposition}[Exact covariance split]
\label{prop:n9v64-covariance-split}
Let $\chi_j:[0,\infty)\to[0,1]$, $j=0,\ldots,J$, be a finite smooth
partition satisfying
\begin{equation}
 \sum_{j=0}^J\chi_j(\lambda)=1
 \quad\hbox{on }[\lambda_{
m IR},\lambda_{
m UV}],
 \label{eq:n9v64-partition-unity}
\end{equation}
with $\chi_j$ supported in an annulus
$[\lambda_j/2,2\lambda_j]$.  Then
\begin{equation}
 A^{-1}P_{\ge}=\sum_{j=0}^JC_j,
 \qquad
 C_j=A^{-1}\chi_j(A)\ge0,
 \label{eq:n9v64-covariance-shells}
\end{equation}
and
\begin{equation}
 \lVert C_j\rVert\le 2/\lambda_j.
 \label{eq:n9v64-shell-covariance-norm}
\end{equation}
Independent centred Gaussian fields with covariances $C_j$ sum to the
cutoff physical Gaussian field with covariance $A^{-1}P_{\ge}$.
\end{proposition}

\begin{proof}
Functional calculus and \eqref{eq:n9v64-partition-unity} give the sum.
Every multiplier $\lambda^{-1}\chi_j(\lambda)$ is nonnegative and is at
most $2/\lambda_j$ on its support.  Covariances add under sums of
independent Gaussian variables.
\end{proof}

\begin{firewall}[The true infrared block remains dynamical]
\label{fire:n9v64-IR-retained}
The covariance $A^{-1}P_{<}$ has norm at least of order
$\lambda_{
m IR}^{-1}$, and the lowest nonzero eigenvalue tends to zero
with volume.  V64 does not integrate this block with a uniform gap.  It
remains in the slow effective theory while the ultraviolet and
intermediate annuli are controlled.
\end{firewall}

\subsection{Coercivity on sharp annuli}

For the relative-form comparison use disjoint spectral annuli.  Put
\begin{equation}
 \Lambda_j=4^j\lambda_{
m IR},
 \qquad
 P_j=\mathbf 1_{[\Lambda_j,4\Lambda_j)}(A),
 \qquad j=0,\ldots,J-1,
 \label{eq:n9v64-sharp-annuli}
\end{equation}
with the last interval truncated at $\lambda_{
m UV}$.  Write
$u_j=P_ju$ and
\begin{equation}
 {\mathfrak a}_j(u_j)=\langle u_j,Au_j\rangle.
 \label{eq:n9v64-shell-energy}
\end{equation}
Then
\begin{equation}
 \Lambda_j\lVert u_j\rVert^2
 \le {\mathfrak a}_j(u_j)
 <4\Lambda_j\lVert u_j\rVert^2.
 \label{eq:n9v64-shell-gap}
\end{equation}

Let $R=R^*$ be the complete renormalized fast-induced Hessian row after
the V58 local owners are assigned.  Define
\begin{equation}
 K_{ij}=(A|_{P_i})^{-1/2}P_iRP_j(A|_{P_j})^{-1/2},
 \qquad r_{ij}=\lVert K_{ij}\rVert.
 \label{eq:n9v64-scale-blocks}
\end{equation}

\begin{theorem}[Scale-index Schur margin]
\label{thm:n9v64-scale-Schur}
Suppose positive weights $w_j$ satisfy
\begin{equation}
 \sup_i\frac1{w_i}\sum_jr_{ij}w_j\le\rho<1.
 \label{eq:n9v64-scale-row}
\end{equation}
Then on $P_{\ge}{\cal H}_{\rm phys}$,
\begin{equation}
 |\langle u,Ru\rangle|
 \le\rho\langle u,Au\rangle,
 \qquad
 \langle u,(A+R)u\rangle
 \ge(1-\rho)\langle u,Au\rangle.
 \label{eq:n9v64-scale-margin}
\end{equation}
The estimate is uniform in the number of annuli whenever $\rho$ is.
\end{theorem}

\begin{proof}
The annuli are orthogonal and $A$ commutes with every $P_j$.  Apply the
weighted block Schur argument of V57 to the Hilbert fibres
$P_j{\cal H}_{\rm phys}$ with reference forms
\eqref{eq:n9v64-shell-energy}.  Self-adjointness makes the row and column
conditions identical.
\end{proof}

\begin{proposition}[Exponentially decoupled scale rows]
\label{prop:n9v64-scale-decay}
If
\begin{equation}
 r_{ij}\le c_0q^{|i-j|},
 \qquad 0<q<1,
 \label{eq:n9v64-scale-decay}
\end{equation}
then the unweighted scale Schur constant is
\begin{equation}
 \rho\le c_0\frac{1+q}{1-q}.
 \label{eq:n9v64-scale-geometric}
\end{equation}
Thus $c_0(1+q)<1-q$ is a sufficient annulus-uniform margin.
\end{proposition}

\begin{proof}
For every $i$,
$\sum_jq^{|i-j|}\le1+2\sum_{m\ge1}q^m=(1+q)/(1-q)$.
Apply \Cref{thm:n9v64-scale-Schur} with $w_j=1$.
\end{proof}

\subsection{Translation-invariant diagnostic}

On a flat torus, suppose $A$ and $R$ are simultaneous Fourier
multipliers on the physical fibre.  Then the sharp-annulus off-diagonal
blocks vanish.

\begin{proposition}[Massless and borderline shell coefficients]
\label{prop:n9v64-symbol-shell}
Assume
\begin{equation}
 \widehat A(p)\ge Z|p|^2I.
 \label{eq:n9v64-A-symbol}
\end{equation}
If
\begin{equation}
 \lVert\widehat R(p)\rVert\le c_2|p|^2,
 \label{eq:n9v64-R-p2}
\end{equation}
then every nonzero annulus has relative coefficient at most $c_2/Z$.
If, after the local quadratic owner is removed,
\begin{equation}
 \lVert\widehat R_{
m rem}(p)\rVert
 \le c_4|p|^4(1+|\log(|p|\ell)|),
 \label{eq:n9v64-R-p4log}
\end{equation}
then on an annulus $\kappa/2\le|p|\le2\kappa$ with
$2\kappa\ell\le1$ its relative coefficient is at most
\begin{equation}
 \rho_{\rm rem}(\kappa)
 \le\frac{4c_4}{Z}\kappa^2
 \bigl(1+|\log(2\kappa\ell)|+\log4\bigr),
 \label{eq:n9v64-shell-p4log}
\end{equation}
which tends to zero as $\kappa\downarrow0$.
\end{proposition}

\begin{proof}
Divide \eqref{eq:n9v64-R-p2} by
\eqref{eq:n9v64-A-symbol}.  For the second statement, divide
\eqref{eq:n9v64-R-p4log} by $Z|p|^2$, use $|p|^2\le4\kappa^2$, and bound
the logarithm throughout the annulus.  Simultaneous Fourier
diagonalization removes all $i\ne j$ blocks.
\end{proof}

\begin{assessment}[What the annuli accomplish]
\label{ass:n9v64-annuli}
The massless $p^2$ reference has no global spectral floor, but each
annulus has the exact floor \eqref{eq:n9v64-shell-gap}.  A correctly
owned $O(p^2)$ response has the same relative coefficient on every
annulus, while the $p^4\log p$ remainder becomes increasingly harmless
in the infrared.  The zero mode is still a separate compatibility
condition.
\end{assessment}

\subsection{Integrating high shells and the infrared Schur complement}

Let $P_L=P_{<}$ and $P_H=P_{\ge}$.  Write
\begin{equation}
 H_{\rm tot}=A+R
 =\begin{pmatrix}
 H_{LL}&H_{LH}\\
 H_{HL}&H_{HH}
 \end{pmatrix}
 \label{eq:n9v64-two-scale-Hessian}
\end{equation}
on $P_L{\cal H}\oplus P_H{\cal H}$.

\begin{theorem}[High-shell Feshbach debit]
\label{thm:n9v64-Feshbach}
Assume
\begin{align}
 \langle v_H,H_{HH}v_H\rangle
 &\ge(1-\rho_H){\mathfrak a}_H(v_H),
 \qquad \rho_H<1,
 \label{eq:n9v64-HH-margin}\\
 |\langle u_L,H_{LH}v_H\rangle|
 &\le\gamma
 {\mathfrak a}_L(u_L)^{1/2}{\mathfrak a}_H(v_H)^{1/2}.
 \label{eq:n9v64-cross-bound}
\end{align}
Then $H_{HH}$ is invertible on the high space and its Feshbach
correction satisfies
\begin{equation}
 0\le
 \langle u_L,H_{LH}H_{HH}^{-1}H_{HL}u_L\rangle
 \le\frac{\gamma^2}{1-\rho_H}{\mathfrak a}_L(u_L).
 \label{eq:n9v64-Feshbach-debit}
\end{equation}
Consequently the exact high-shell effective Hessian is
\begin{equation}
 H_{\rm eff,L}=H_{LL}-H_{LH}H_{HH}^{-1}H_{HL},
 \label{eq:n9v64-effective-low-Hessian}
\end{equation}
with an additional negative relative debit at most
$\gamma^2/(1-\rho_H)$.
\end{theorem}

\begin{proof}
The first inequality makes $H_{HH}$ positive and invertible at finite
regulator.  Let $z=H_{HH}^{-1}H_{HL}u_L$.  Then
\begin{align}
 \langle u_L,H_{LH}H_{HH}^{-1}H_{HL}u_L\rangle
 &=\langle z,H_{HH}z\rangle\\
 &\le \gamma {\mathfrak a}_L(u_L)^{1/2}
             {\mathfrak a}_H(z)^{1/2}
 \end{align}
by \eqref{eq:n9v64-cross-bound}, while
\eqref{eq:n9v64-HH-margin} bounds the left side below by
$(1-\rho_H){\mathfrak a}_H(z)$.  Divide and square.  Completing the
square in the high variable gives \eqref{eq:n9v64-effective-low-Hessian}.
\end{proof}

\begin{corollary}[Aggregate low-sector margin]
\label{cor:n9v64-low-margin}
If
\begin{equation}
 \langle u_L,H_{LL}u_L\rangle
 \ge(1-\rho_L){\mathfrak a}_L(u_L),
 \label{eq:n9v64-LL-margin}
\end{equation}
then
\begin{equation}
 \langle u_L,H_{\rm eff,L}u_L\rangle
 \ge\left(1-\rho_L-\frac{\gamma^2}{1-\rho_H}\right)
 {\mathfrak a}_L(u_L).
 \label{eq:n9v64-low-aggregate}
\end{equation}
The low block remains coercive only if the coefficient in parentheses
is positive.
\end{corollary}

\begin{proof}
Subtract \eqref{eq:n9v64-Feshbach-debit} from
\eqref{eq:n9v64-LL-margin}.
\end{proof}

\begin{firewall}[A retained infrared block still receives a debit]
\label{fire:n9v64-IR-debit}
Leaving $P_L$ unintegrated avoids inverting its small eigenvalues, but
integrating high shells changes its Hessian through
\eqref{eq:n9v64-effective-low-Hessian}.  That negative Schur row must be
kept in the slow action and cannot be discarded as an ultraviolet
normalization constant.
\end{firewall}

\subsection{Locality and reconstruction cautions}

\begin{proposition}[Projection after the local estimate]
\label{prop:n9v64-project-after}
Suppose a physical-space form $R$ already satisfies
\begin{equation}
 |\langle u,Rv\rangle|
 \le\rho\langle u,Au\rangle^{1/2}
             \langle v,Av\rangle^{1/2}.
 \label{eq:n9v64-global-relative}
\end{equation}
Then every spectral compression $P_iRP_j$ satisfies the corresponding
energy-normalized bound with coefficient at most $\rho$.  No spatial
locality of $P_i$ is required.
\end{proposition}

\begin{proof}
Apply \eqref{eq:n9v64-global-relative} to $P_iu$ and $P_jv$ and use
$AP_j=P_jA$.
\end{proof}

\begin{firewall}[Sharp annuli are an analysis device]
\label{fire:n9v64-sharp-locality}
Sharp spectral projectors generally have long position-space kernels
and need not preserve reflection positivity.  V64 uses them to prove
relative-form estimates after the local marked-cluster construction.
A local RG transformation requires smooth quasi-local multipliers or a
separately proved finite-range covariance decomposition.
\end{firewall}

\begin{firewall}[Gauge reduction and spectral calculus]
\label{fire:n9v64-gauge-reduction}
The operator $A$ must already act on the exact physical tangent space or
on a gauge-fixed space with controlled ghosts and Jacobians.  Applying
spectral calculus to an unconstrained Hessian and deleting its gauge
kernel afterwards does not prove the same covariance or determinant.
\end{firewall}

\subsection{Massless annular covariance card}

\begin{definition}[Massless annular covariance card (MACC)]
\label{def:n9v64-MACC}
MACC requires:
\begin{enumerate}
\item an exact positive constraint-reduced Einstein operator $A$ and
declared zero-mode owner;
\item the positive covariance partition
\eqref{eq:n9v64-covariance-shells};
\item a true infrared block retained without a uniform gap claim;
\item a scale-index Schur bound for the complete renormalized Hessian
row on all integrated annuli;
\item a high/low cross bound and the Feshbach debit
\eqref{eq:n9v64-Feshbach-debit};
\item local marked-cluster estimates proved before nonlocal spectral
compression;
\item regulator-uniform constants as the number of annuli grows; and
\item a separate locality and reflection-positivity argument for any
actual multiscale integration map.
\end{enumerate}
\end{definition}

\begin{theorem}[MACC integrates all gapped annuli without a false mass]
\label{thm:n9v64-MACC}
Under MACC, every annulus above $\lambda_{
m IR}$ has a uniform
scale-relative inverse controlled by
\Cref{thm:n9v64-scale-Schur}.  Integrating those modes produces the
explicit low-sector Schur correction
\eqref{eq:n9v64-effective-low-Hessian}, while modes below
$\lambda_{
m IR}$ remain in the slow theory.  No estimate uses a
volume-uniform gap for the full Einstein operator.
\end{theorem}

\begin{proof}
Use \Cref{prop:n9v64-covariance-split,
thm:n9v64-scale-Schur,thm:n9v64-Feshbach,
cor:n9v64-low-margin}.  The retained block is never inverted.
\end{proof}

\begin{assessment}[Progress at ninth-series V64]
\label{ass:n9v64-status}
The massless covariance is now separated into honest gapped annuli and a
true infrared block.  The price of high-shell integration is explicit:
each annulus must pass a scale Schur test, and high/low mixing produces a
negative Feshbach debit.  The remaining construction problem is a local,
symmetry-compatible scale map with the same bounds.
\end{assessment}

\begin{programme}[Ninth-series V65: compulsory companion audit]
\label{prog:n9v64-V65}
V65 will inspect the newest Yang--Mills and Navier--Stokes notes, compare
their current coarse-transfer and scale-localization results with MACC,
and choose the next local multiscale acquisition.  V66 will continue
after that audit.
\end{programme}

\begin{firewall}[Claim boundary after ninth-series V64]
\label{fire:n9v64-boundary}
V64 proves the annular spectral and Feshbach estimates.  It does not
construct a local finite-range physical metric covariance, establish
the needed uniform scale rows for the chiral Standard Model, control the
anomaly or phase, prove reflection positivity of the scale map, close
the true infrared sector, reconstruct the OS theory, or prove the full
SM--Einstein continuum.
\end{firewall}

\begin{theorem}[Ninth-series V64 result]
\label{thm:n9v64-result}
The constraint-reduced massless Einstein covariance admits the exact
positive annular decomposition \eqref{eq:n9v64-covariance-shells} above
an explicit infrared threshold.  The integrated annuli are controlled
by the scale Schur condition \eqref{eq:n9v64-scale-row}; their exact
elimination adds the low-sector debit
$\gamma^2/(1-\rho_H)$.  The infrared block remains dynamical, so no
volume-uniform graviton mass is assumed.
\end{theorem}

\begin{proof}
Combine \Cref{prop:n9v64-covariance-split,
thm:n9v64-scale-Schur,prop:n9v64-symbol-shell,
thm:n9v64-Feshbach,thm:n9v64-MACC}.
\end{proof}

\paragraph{Parent-version record.}
V64 was formed from
\texttt{Renewal\_SM\_Einstein\_Continuum\_Research\_Notes\_V63.tex}.
The V63 parent had $23\,552$ logical source lines, $850\,000$ bytes, and
SHA--256
\[
 \texttt{4D5AB0CEE301A02A6D972E6FB6021CA3D2C94F5EAFC2BCC422F76118AB9350CA}.
\]

% =============================================================================
% END APPENDED NINTH-SERIES V64 MATERIAL
% =============================================================================
% =============================================================================
% BEGIN APPENDED NINTH-SERIES V65 MATERIAL
% =============================================================================

\section{Ninth-series V65: compulsory audit and the eliminated determinant row}
\label{sec:v9v65}

\subsection{Companion files inspected}

This is the compulsory audit five versions after V60.  The newest
Yang--Mills files were V96 and V97; they were byte-identical, so V96 is
the substantive endpoint.  The Navier--Stokes file remained V26.

\begin{center}
\begin{tabular}{|p{0.43\textwidth}|r|r|}
\hline
File & bytes & lines\\
\hline
\texttt{Renewal\_Yang\_Mills\_Mass\_Gap\_Research\_Notes\_V97.tex}
& $1\,560\,740$ & $43\,695$\\
\texttt{Renewal\_NS\_Stability\_Research\_Notes\_V26.tex}
& $278\,283$ & $5\,319$\\
\hline
\end{tabular}
\end{center}
Their SHA--256 digests were
\begin{align}
 &\texttt{C14B8DBD6CF49735BC2B62D12994C3673FEBBE76D47B8F82D81CBF61D4EEB8DA},
 \label{eq:n9v65-YM-hash}\\
 &\texttt{82C887F8C2C69E4F28FAD7C3DC8DE5B9099CB5A4BF431EBA1FFF3E91935978AD}.
 \label{eq:n9v65-NS-hash}
\end{align}

\subsection{Audit findings}

The substantive Yang--Mills V96 result concerns a Gaussian determinant
formed only on a gapped eliminated physical subspace.  It differentiates
$\frac12\log\det L$, derives exponential locality from a finite-range
gap, and obtains a response row of order
$g^2(P_2(M)+P_1(M)^2)$ under an explicit weak-field derivative scaling.
It also shows that a massless four-dimensional inverse gives a
logarithmically divergent absolute response row.  These conclusions
match the division in V64 between gapped integrated modes and a retained
infrared block.

The Navier--Stokes V26 content and hash have not changed.  Its applicable
lesson remains V60's restricted-input norm: use an exact physical
tangent compression when available, but do not assume a strict gain.

\begin{decision}[Transfer from the V65 audit]
\label{dec:n9v65-transfer}
The SM--Einstein programme will:
\begin{enumerate}
\item form Gaussian determinants only after all gauge, toron, and
retained infrared modes have been removed from the eliminated Hessian;
\item use exact trace derivatives and finite-range inverse locality to
bound the determinant metric-response row;
\item keep weak-field derivative scaling as a physical hypothesis to be
proved for each gauge--Higgs--fermion block; and
\item replace sharp spectral annuli by positive resolvent differences
for the next candidate local scale covariance.
\end{enumerate}
\end{decision}

\subsection{Self-contained determinant derivative calculus}

Let $L(s)>0$ be a finite-dimensional eliminated Hessian depending
twice differentiably on retained block variables $s_x$.  Put
\begin{equation}
 G=L^{-1},
 \qquad
 B_x=D_xL,
 \qquad
 C_{xy}=D_yD_xL.
 \label{eq:n9v65-L-derivatives}
\end{equation}

\begin{theorem}[Exact eliminated determinant responses]
\label{thm:n9v65-logdet-derivatives}
For
\begin{equation}
 \Gamma_{\rm det}(s)=\frac\hbar2\log\det L(s),
 \label{eq:n9v65-logdet}
\end{equation}
one has
\begin{align}
 D_x\Gamma_{\rm det}
 &=\frac\hbar2\operatorname{Tr}(GB_x),
 \label{eq:n9v65-logdet-first}\\
 D_yD_x\Gamma_{\rm det}
 &=\frac\hbar2\operatorname{Tr}
 \bigl(GC_{xy}-GB_yGB_x\bigr).
 \label{eq:n9v65-logdet-second}
\end{align}
The second term is the negative Gaussian response covariance/Feshbach
row; the direct $C_{xy}$ row is a separate local owner.
\end{theorem}

\begin{proof}
Jacobi's formula gives
$D\log\det L=\operatorname{Tr}(L^{-1}DL)$.  Differentiate once more and
use $D_yG=-GB_yG$.
\end{proof}

\begin{firewall}[Pseudo-determinant with a closing floor]
\label{fire:n9v65-pseudodet}
Replacing $\det L$ by a pseudo-determinant does not make
\eqref{eq:n9v65-logdet-second} uniform if its smallest retained
eigenvalue tends to zero.  Both copies of $G$ in the covariance row
detect the closing floor.  Such modes belong in the V64 infrared block.
\end{firewall}

\subsection{Finite-range inverse locality}

Let the eliminated variables be decomposed into finite-rank block
fibres with projections $P_x$.  Suppose $L\ge\lambda I$, its off-diagonal
interaction range is $R_0$, and
\begin{equation}
 \sup_x\sum_{y\ne x}\lVert P_xLP_y\rVert\le J.
 \label{eq:n9v65-offdiag-row}
\end{equation}

\begin{lemma}[Weighted inverse estimate]
\label{lem:n9v65-weighted-inverse}
For
\begin{equation}
 0<\nu\le R_0^{-1}\log\left(1+\frac\lambda{2J}\right),
 \label{eq:n9v65-nu}
\end{equation}
one has
\begin{equation}
 \lVert P_xGP_y\rVert
 \le\frac2\lambda e^{-\nu d(x,y)}.
 \label{eq:n9v65-inverse-decay}
\end{equation}
\end{lemma}

\begin{proof}
Fix $y$ and conjugate by the diagonal weight
$W_\nu(x)=e^{\nu d(x,y)}$.  Finite range and
\eqref{eq:n9v65-offdiag-row} give
\begin{equation}
 \lVert W_\nu LW_\nu^{-1}-L\rVert
 \le J(e^{\nu R_0}-1)\le\lambda/2.
 \end{equation}
Hence $W_\nu LW_\nu^{-1}$ is invertible with inverse norm at most
$2/\lambda$.  Taking its $(x,y)$ block yields
\eqref{eq:n9v65-inverse-decay}.
\end{proof}

Assume $B_x$ is supported on a collar of radius $R_1$ around $x$, with
rank at most $r_*$ and norm at most $b_1$.  Assume $C_{xy}=0$ for
$d(x,y)>R_2$ and $\lVert C_{xy}\rVert\le b_2$ on its support.

\begin{theorem}[Quasi-local determinant Hessian]
\label{thm:n9v65-det-locality}
Under these hypotheses,
\begin{equation}
 |D_yD_x\Gamma_{\rm det}|
 \le\frac{\hbar r_*}{2}left[
 \frac{b_2}{\lambda}\mathbf1_{\{d(x,y)\le R_2\}}
 +\frac{4b_1^2}{\lambda^2}
 e^{-2\nu(d(x,y)-2R_1)_+}
 \right].
 \label{eq:n9v65-det-kernel-bound}
\end{equation}
On a polynomial-growth block lattice, every weighted row with exponent
$0<\mu<2\nu$ is finite uniformly in total volume.
\end{theorem}

\begin{proof}
The direct trace is supported where $C_{xy}\ne0$ and is bounded by rank
times $\lVert G\rVert\lVert C_{xy}\rVert\le r_*b_2/\lambda$.
For the response trace, insert the two derivative collars:
\begin{equation}
 \operatorname{Tr}(GB_yGB_x)
 =\operatorname{Tr}
 (P_xGP_yB_yP_yGP_xB_x),
 \end{equation}
with the projections enlarged by $R_1$.  Two applications of
\eqref{eq:n9v65-inverse-decay}, the rank bound, and the two $b_1$ bounds
give the second term.  Polynomial shell growth is dominated by
$e^{-(2\nu-\mu)d}$.
\end{proof}

\begin{corollary}[Weak-field determinant debit]
\label{cor:n9v65-weak-field}
If
\begin{equation}
 b_1\le C_1gP_1(M),
 \qquad
 b_2\le C_2g^2P_2(M),
 \label{eq:n9v65-weak-derivatives}
\end{equation}
while $\lambda,R_0,J,R_1,R_2,r_*$ are uniform, then
\begin{equation}
 \sup_x\sum_y e^{\mu d(x,y)}
 |D_yD_x\Gamma_{\rm det}|
 \le C\hbar g^2\bigl(P_2(M)+P_1(M)^2\bigr).
 \label{eq:n9v65-weak-row}
\end{equation}
For $g=\beta^{-1/2}$ and fixed logarithmic-polynomial $P_i(M)$ with
$M=(\log\beta)^\chi$, the row tends to zero as $\beta\to\infty$.
\end{corollary}

\begin{proof}
Insert \eqref{eq:n9v65-weak-derivatives} in
\eqref{eq:n9v65-det-kernel-bound}, sum the weighted row, and use that
every fixed power of $\log\beta$ is $o(\beta)$.
\end{proof}

\begin{firewall}[Weak field does not follow from proximity]
\label{fire:n9v65-weak-field}
A norm estimate $\lVert L(s)-L(0)\rVert=o(1)$ does not imply the first
and second derivative estimates
\eqref{eq:n9v65-weak-derivatives}.  Moving gauge projectors, chart
Jacobians, Higgs backgrounds, and chiral determinant rows must be
differentiated explicitly and counted once.
\end{firewall}

\subsection{Why the audit changes the local scale map}

The sharp annuli in V64 are exact for spectral bookkeeping, but their
projectors are nonlocal.  The inverse estimate above suggests a rational
scale partition built from gapped resolvents.  If
$m_0>m_1>\cdots>m_J>0$, then
\begin{equation}
 (A+m_j^2)^{-1}-(A+m_{j-1}^2)^{-1}
 =(m_{j-1}^2-m_j^2)
 (A+m_j^2)^{-1}(A+m_{j-1}^2)^{-1}ge0.
 \label{eq:n9v65-resolvent-difference}
\end{equation}
Each factor has a genuine mass floor and therefore admits a weighted
inverse estimate when $A$ is local.

\begin{proposition}[Positive telescoping preview]
\label{prop:n9v65-telescoping-preview}
The positive covariances
\begin{equation}
 C_j=(A+m_j^2)^{-1}-(A+m_{j-1}^2)^{-1}
 \label{eq:n9v65-Cj-preview}
\end{equation}
satisfy
\begin{equation}
 \sum_{j=1}^JC_j
 =(A+m_J^2)^{-1}-(A+m_0^2)^{-1}.
 \label{eq:n9v65-telescope-preview}
\end{equation}
Thus they provide positive quasi-local covariance slices between the UV
and infrared mass scales without a sharp projector.
\end{proposition}

\begin{proof}
The resolvent identity gives positivity in
\eqref{eq:n9v65-resolvent-difference}; the sum telescopes.
\end{proof}

\begin{firewall}[A resolvent slice is not yet a complete RG step]
\label{fire:n9v65-resolvent-preview}
Positivity and exponential kernel decay do not by themselves prove
finite range, reflection positivity, gauge covariance of a nonlinear
block map, or locality of the induced interaction.  These are separate
V66 obligations.
\end{firewall}

\subsection{Eliminated determinant locality card}

\begin{definition}[SM eliminated determinant card (SMEDC)]
\label{def:n9v65-SMEDC}
SMEDC requires:
\begin{enumerate}
\item an exact eliminated physical Hessian with all gauge and retained
slow modes removed;
\item a uniform positive floor, finite interaction range, and
off-diagonal row bound;
\item local first and second retained-variable derivatives;
\item explicit moving-projector, chart, measure, and phase owners;
\item the trace identities \eqref{eq:n9v65-logdet-first}--
\eqref{eq:n9v65-logdet-second};
\item a weighted determinant response row strictly inside the aggregate
CPRAC budget; and
\item compatibility with the exact coarse logarithm and cutoff limits.
\end{enumerate}
\end{definition}

\begin{theorem}[SMEDC loads the Gaussian determinant into MACC]
\label{thm:n9v65-SMEDC}
Under SMEDC, the eliminated Gaussian determinant contributes a
quasi-local metric Hessian row given exactly by
\eqref{eq:n9v65-logdet-second}.  If the weak-field scaling
\eqref{eq:n9v65-weak-derivatives} holds, its weighted row tends to zero
along the stated schedule and can be included in the V64 scale margin.
\end{theorem}

\begin{proof}
Use \Cref{thm:n9v65-logdet-derivatives,
lem:n9v65-weighted-inverse,thm:n9v65-det-locality,
cor:n9v65-weak-field}.  SMEDC item 6 allocates the resulting constant in
the MACC Schur ledger.
\end{proof}

\begin{assessment}[Direction after the V65 audit]
\label{ass:n9v65-direction}
The audit reinforces the retained-infrared decision and supplies an
exact determinant row for every genuinely gapped eliminated block.  It
also selects positive resolvent differences as the next local
mass-scale covariance candidate.  The unchanged NS note supplies no new
direction beyond exact restriction to the physical input class.
\end{assessment}

\begin{programme}[Ninth-series V66: positive resolvent scale slices]
\label{prog:n9v65-V66}
V66 will construct the full telescoping resolvent covariance, bound each
slice in weighted position-space norms, identify the ultraviolet and
infrared remainders, and test scale summability.  It will also state
precisely which locality and reflection properties do not follow from
the resolvent algebra.
\end{programme}

\begin{firewall}[Claim boundary after ninth-series V65]
\label{fire:n9v65-boundary}
V65 proves the determinant derivative and locality compilers under an
eliminated gap.  It does not prove the physical weak-field derivative
scaling, a uniform gap for the actual SM block, a local reflection-
positive scale map, chiral phase or anomaly control, closure of the
infrared theory, OS reconstruction, or the full SM--Einstein continuum.
\end{firewall}

\begin{theorem}[Ninth-series V65 result]
\label{thm:n9v65-result}
The V65 audit transfers the exact eliminated determinant Hessian
\eqref{eq:n9v65-logdet-second} and the gapped inverse estimate
\eqref{eq:n9v65-inverse-decay} into the SM scaling ledger.  Under local
weak-field derivatives its weighted response is $O(\hbar g^2)$ up to
the declared logarithmic polynomials.  Positive resolvent differences
\eqref{eq:n9v65-Cj-preview} are selected for the next quasi-local scale
construction.
\end{theorem}

\begin{proof}
Combine \Cref{dec:n9v65-transfer,
thm:n9v65-logdet-derivatives,thm:n9v65-det-locality,
cor:n9v65-weak-field,prop:n9v65-telescoping-preview,
thm:n9v65-SMEDC}.
\end{proof}

\paragraph{Parent-version record.}
V65 was formed from
\texttt{Renewal\_SM\_Einstein\_Continuum\_Research\_Notes\_V64.tex}.
The V64 parent had $23\,994$ logical source lines, $864\,780$ bytes, and
SHA--256
\[
 \texttt{1364C83F729A103AA1085DB5B7A9E9066C03D5B29683A84F7968E3180D49207A}.
\]

% =============================================================================
% END APPENDED NINTH-SERIES V65 MATERIAL
% =============================================================================
% =============================================================================
% BEGIN APPENDED NINTH-SERIES V66 MATERIAL
% =============================================================================

\section{Ninth-series V66: positive resolvent covariance slices}
\label{sec:v9v66}

\subsection{Exact telescoping construction}

V64 used spectral annuli as an exact analysis device, while the V65
audit selected massive resolvent differences as a quasi-local scale
candidate.  Let $A\ge0$ be the constraint-reduced Einstein reference on
the orthogonal complement $Q{\cal H}$ of its declared kernel.  At finite
volume $A_Q=QAQ$ is positive.  Choose
\begin{equation}
 m_0>m_1>\cdots>m_J>0
 \label{eq:n9v66-mass-sequence}
\end{equation}
and set
\begin{align}
 R_j&=(A_Q+m_j^2)^{-1},
 \label{eq:n9v66-resolvent}\\
 C_j&=R_j-R_{j-1},\qquad j=1,\ldots,J.
 \label{eq:n9v66-slice}
\end{align}

\begin{theorem}[Positive resolvent telescope]
\label{thm:n9v66-telescope}
Every $C_j$ is positive and
\begin{align}
 C_j&=(m_{j-1}^2-m_j^2)R_jR_{j-1},
 \label{eq:n9v66-resolvent-identity}\\
 A_Q^{-1}
 &=C_{\rm UV}+\sum_{j=1}^JC_j+C_{{\rm IR},J},
 \label{eq:n9v66-full-telescope}\\
 C_{\rm UV}&=R_0,
 \qquad
 C_{{\rm IR},J}=A_Q^{-1}-R_J
 =m_J^2A_Q^{-1}R_J\ge0.
 \label{eq:n9v66-remainders}
\end{align}
Thus independent Gaussian fields with these covariances sum exactly to
the finite-volume massless physical Gaussian field.
\end{theorem}

\begin{proof}
All factors are functions of $A_Q$ and commute.  The scalar identity
\begin{equation}
 \frac1{\lambda+m_j^2}-\frac1{\lambda+m_{j-1}^2}
 =\frac{m_{j-1}^2-m_j^2}
 { (\lambda+m_j^2)(\lambda+m_{j-1}^2)}
 \end{equation}
is nonnegative for $\lambda>0$, proving
\eqref{eq:n9v66-resolvent-identity}.  Sum the differences and insert
$A_Q^{-1}-R_J$; the second formula in
\eqref{eq:n9v66-remainders} is the same resolvent identity.  Independent
Gaussian covariances add.
\end{proof}

\begin{corollary}[No kernel regularization is smuggled in]
\label{cor:n9v66-kernel-owner}
The construction takes place on $Q{\cal H}$.  If it were applied to a
zero eigenvector, the partial sum would contain
$m_J^{-2}-m_0^{-2}$ and diverge as $m_J\downarrow0$.  Kernel removal or
explicit integration is therefore prior to the telescope.
\end{corollary}

\begin{proof}
Evaluate \eqref{eq:n9v66-slice} at $\lambda=0$.
\end{proof}

\subsection{Energy size of one scale}

For a geometric schedule
\begin{equation}
 m_j=q,m_{j-1},
 \qquad 0<q<1,
 \label{eq:n9v66-geometric-mass}
\end{equation}
the energy-normalized slice has a sharp universal bound.

\begin{proposition}[Exact energy-normalized slice norm]
\label{prop:n9v66-slice-norm}
For every $j$,
\begin{equation}
 \lVert A_Q^{1/2}C_jA_Q^{1/2}\rVert
 \le\frac{m_{j-1}-m_j}{m_{j-1}+m_j}
 =\frac{1-q}{1+q}.
 \label{eq:n9v66-exact-slice-norm}
\end{equation}
The supremum is attained spectrally at
$\lambda=m_jm_{j-1}$ when that value belongs to the spectrum, and is
otherwise the corresponding upper bound.
\end{proposition}

\begin{proof}
Functional calculus reduces the norm to the supremum of
\begin{equation}
 f(\lambda)=
 \frac{\lambda(m_{j-1}^2-m_j^2)}
 { (\lambda+m_j^2)(\lambda+m_{j-1}^2)}.
 \end{equation}
Differentiation gives the unique maximum at
$\lambda=m_jm_{j-1}$.  Substitution yields
$(m_{j-1}-m_j)/(m_{j-1}+m_j)$.  On a discrete spectrum this is an upper
bound, with equality under the stated spectral condition.
\end{proof}

\begin{firewall}[Small slices can have a large total]
\label{fire:n9v66-small-slices}
Taking $q\uparrow1$ makes each norm in
\eqref{eq:n9v66-exact-slice-norm} small, but increases the number of
slices.  Triangle summation is not the covariance identity and may
produce a logarithmic scale-volume loss.
\end{firewall}

\begin{proposition}[Logarithmic scale count]
\label{prop:n9v66-log-count}
If $m_J=q^Jm_0$, then
\begin{equation}
 J\frac{1-q}{1+q}
 =\frac{1-q}{(1+q)|\log q|}
   \log\frac{m_0}{m_J}.
 \label{eq:n9v66-log-count}
\end{equation}
As $q\uparrow1$, the prefactor tends to $1/2$.  Therefore a nonlinear
debit proportional only to the per-slice norm need not be summable as
$m_J\downarrow0$.
\end{proposition}

\begin{proof}
The identity $J=\log(m_0/m_J)/|\log q|$ gives the formula.  Apply
l'H\^opital's rule to $(1-q)/|\log q|$.
\end{proof}

\subsection{Continuous mass representation}

\begin{proposition}[Positive continuous resolvent decomposition]
\label{prop:n9v66-continuous}
For every $M>0$,
\begin{equation}
 A_Q^{-1}=(A_Q+M^2)^{-1}
 +\int_0^{M^2}(A_Q+s)^{-2}\,ds.
 \label{eq:n9v66-continuous-decomposition}
\end{equation}
In logarithmic mass $s=m^2$, the energy-normalized density is
\begin{equation}
 2m^2A_Q(A_Q+m^2)^{-2},d\log m,
 \label{eq:n9v66-log-density}
\end{equation}
whose operator norm is at most $1/2$ per unit $d\log m$.
\end{proposition}

\begin{proof}
Integrate
$-\partial_s(A_Q+s)^{-1}=(A_Q+s)^{-2}$ from $0$ to $M^2$.
For the second claim maximize
$2m^2\lambda/(\lambda+m^2)^2$, whose maximum $1/2$ occurs at
$\lambda=m^2$.
\end{proof}

\begin{remark}[Running estimates are required]
\label{rem:n9v66-running}
Equation \eqref{eq:n9v66-log-density} makes the scale measure explicit.
A per-log-scale interaction row must decay, cancel, or be absorbed into
running local owners to remain integrable over an unbounded infrared
scale interval.
\end{remark}

\subsection{Quasi-locality of massive resolvents}

The general V65 weighted-conjugation estimate gives a decay rate of
order $m^2$ from only finite range and an off-diagonal row.  A second-
order Einstein operator has a stronger energy commutator estimate.  Let
$d_a$ denote physical lattice distance and let
$W_\nu=e^{\nu d_a(\cdot,y)}$.

\begin{definition}[Quadratic weighted coercivity]
\label{def:n9v66-QWC}
The local operator $A_Q$ has QWC constant $c_A$ if
\begin{equation}
 \operatorname{Re}\langle u,W_\nu A_QW_\nu^{-1}u\rangle
 \ge\langle u,A_Qu\rangle-c_A\nu^2\lVert u\rVert^2
 \label{eq:n9v66-QWC}
\end{equation}
for all sufficiently small $\nu$ and all $u$ in the form domain.
\end{definition}

\begin{lemma}[Mass-length inverse locality]
\label{lem:n9v66-mass-locality}
Under QWC, if
\begin{equation}
 0<\nu\le m/\sqrt{2c_A},
 \label{eq:n9v66-mass-rate}
\end{equation}
then
\begin{equation}
 \lVert P_x(A_Q+m^2)^{-1}P_y\rVert
 \le\frac2{m^2}e^{-\nu d_a(x,y)}.
 \label{eq:n9v66-mass-resolvent-decay}
\end{equation}
\end{lemma}

\begin{proof}
By \eqref{eq:n9v66-QWC}, the real part of
$W_\nu(A_Q+m^2)W_\nu^{-1}$ is at least
$m^2-c_A\nu^2\ge m^2/2$.  Its inverse norm is at most $2/m^2$.
Taking the $(x,y)$ block after undoing the weight proves the estimate.
\end{proof}

\begin{corollary}[Slice quasi-locality]
\label{cor:n9v66-slice-locality}
Each $C_j$ obeys
\begin{equation}
 \lVert P_xC_jP_y\rVert
 \le 2(m_j^{-2}+m_{j-1}^{-2})
 e^{-c m_jd_a(x,y)},
 \label{eq:n9v66-slice-decay}
\end{equation}
with $c=(2c_A)^{-1/2}$, within the QWC range.  Thus the physical
localization length is $O(m_j^{-1})$.
\end{corollary}

\begin{proof}
Use $C_j=R_j-R_{j-1}$, the triangle inequality, and
\Cref{lem:n9v66-mass-locality}; the smaller mass gives the slower decay.
\end{proof}

\begin{firewall}[Operator locality versus absolute kernel rows]
\label{fire:n9v66-operator-row}
The weighted $\ell^2$ operator estimate
\eqref{eq:n9v66-mass-resolvent-decay} is sufficient for relative forms
and Feshbach inversion.  It does not automatically give the optimally
scaled absolute block-row norm required by a polymer expansion.  A
heat-kernel representation, positivity-preserving kernel, or sharper
off-diagonal summation estimate is needed for that conversion.
\end{firewall}

\subsection{Metric derivatives without moving spectral projectors}

Let $A=A(g)$ on a locally trivialized physical tangent bundle and put
\begin{equation}
 B_h=D_gA[h],
 \qquad
 D_{h,k}=D_g^2A[h,k].
 \label{eq:n9v66-A-derivatives}
\end{equation}

\begin{proposition}[Exact massive resolvent derivatives]
\label{prop:n9v66-resolvent-derivatives}
For $R_m=(A+m^2)^{-1}$,
\begin{align}
 D_gR_m[h]
 &=-R_mB_hR_m,
 \label{eq:n9v66-R-first}\\
 D_g^2R_m[h,k]
 &=R_mB_kR_mB_hR_m+R_mB_hR_mB_kR_m
   -R_mD_{h,k}R_m.
 \label{eq:n9v66-R-second}
\end{align}
The derivatives of $C_j$ are the differences of these formulas at
$m_j$ and $m_{j-1}$.
\end{proposition}

\begin{proof}
Differentiate $(A+m^2)R_m=I$ once and twice.  The mass is fixed and
there is no derivative of a spectral window.
\end{proof}

\begin{corollary}[Marked quasi-locality under local form derivatives]
\label{cor:n9v66-marked-locality}
If $B_h$ and $D_{h,k}$ have fixed local supports and uniform form bounds,
then every term in \eqref{eq:n9v66-R-first}--
\eqref{eq:n9v66-R-second} is a finite product of massive quasi-local
resolvents and local insertions.  Its kernel therefore decays
exponentially at a rate controlled by the smallest participating mass,
with explicit powers of the corresponding inverse floor.
\end{corollary}

\begin{proof}
Insert local support projections around each $B$ or $D$ and apply
\eqref{eq:n9v66-mass-resolvent-decay} to the intervening resolvents.
Sum over the finitely many insertion collars.
\end{proof}

\begin{firewall}[Moving physical subspace]
\label{fire:n9v66-moving-physical}
The formulas assume a fixed local trivialization of the physical tangent
bundle.  If the gauge/constraint projector depends on $g$, its first and
second derivatives enter $B_h$ and $D_{h,k}$.  Resolvent calculus avoids
moving annular projectors; it does not erase moving physical-projector
rows.
\end{firewall}

\subsection{Ultraviolet and infrared remainders}

\begin{proposition}[Remainder bounds]
\label{prop:n9v66-remainder-bounds}
The ultraviolet covariance satisfies
\begin{equation}
 0\le C_{\rm UV}\le m_0^{-2}I.
 \label{eq:n9v66-UV-bound}
\end{equation}
For every spectral threshold $\Lambda>0$, the infrared remainder obeys
\begin{equation}
 \left\lVert
 A_Q^{1/2}C_{{\rm IR},J}A_Q^{1/2}
 \mathbf1_{[\Lambda,\infty)}(A_Q)
 \right\rVert
 \le\frac{m_J^2}{\Lambda+m_J^2}.
 \label{eq:n9v66-IR-high-bound}
\end{equation}
It is concentrated spectrally on $\lambda\lesssim m_J^2$ and must be
retained as $m_J\downarrow0$ unless a separate infrared construction is
available.
\end{proposition}

\begin{proof}
The first claim is spectral monotonicity.  The energy-normalized infrared
symbol is
\begin{equation}
 \lambda\left(\lambda^{-1}-(\lambda+m_J^2)^{-1}\right)
 =\frac{m_J^2}{\lambda+m_J^2},
 \end{equation}
whose supremum on $\lambda\ge\Lambda$ is the right side of
\eqref{eq:n9v66-IR-high-bound}.
\end{proof}

\subsection{Symmetry and reflection boundaries}

\begin{proposition}[Equivariance of resolvent slices]
\label{prop:n9v66-equivariance}
If a unitary regulator symmetry $U$ commutes with $A_Q$ and preserves
$Q{\cal H}$, then it commutes with every $R_j$, $C_j$, and both
remainders in \eqref{eq:n9v66-remainders}.
\end{proposition}

\begin{proof}
Resolvents and their differences are functions of $A_Q$ by functional
calculus.
\end{proof}

\begin{firewall}[Positive covariance is not reflection positivity]
\label{fire:n9v66-reflection}
Operator positivity of $C_j$ constructs an ordinary Gaussian measure.
Reflection positivity requires positivity of the reflected cross-form,
which is not preserved by subtracting two reflection-positive
covariances in general.  Each slice or the integrated scale map needs a
separate reflection argument before OS reconstruction.
\end{firewall}

\begin{definition}[Positive resolvent scale card (PRSC)]
\label{def:n9v66-PRSC}
PRSC requires:
\begin{enumerate}
\item exact kernel removal and the telescope
\eqref{eq:n9v66-full-telescope};
\item QWC or another mass-length off-diagonal estimate;
\item local first and second metric derivatives including the moving
physical projector;
\item an absolute marked-kernel estimate when a polymer norm is used;
\item explicit UV and retained IR remainders;
\item summability of nonlinear debits against the logarithmic scale
measure;
\item regulator symmetry equivariance; and
\item a separate reflection-positivity proof for the chosen scale
integration.
\end{enumerate}
\end{definition}

\begin{theorem}[PRSC supplies a quasi-local alternative to sharp annuli]
\label{thm:n9v66-PRSC}
Under PRSC, the massless physical covariance is decomposed into positive
massive quasi-local slices with exact metric derivatives, plus declared
UV and IR remainders.  Every integrated slice has a genuine mass floor;
the infrared remainder is not inverted; and the total nonlinear debit
is tested against the actual logarithmic scale measure rather than a
per-slice bound alone.
\end{theorem}

\begin{proof}
Use \Cref{thm:n9v66-telescope,prop:n9v66-slice-norm,
lem:n9v66-mass-locality,prop:n9v66-resolvent-derivatives,
prop:n9v66-remainder-bounds,prop:n9v66-log-count}.  PRSC items 4 and 8
supply the two properties not implied by operator resolvent algebra.
\end{proof}

\begin{assessment}[Progress at ninth-series V66]
\label{ass:n9v66-status}
The sharp-annulus bookkeeping now has a positive quasi-local resolvent
counterpart.  The construction exposes three distinct remaining rows:
absolute kernel normalization for polymers, logarithmic accumulation of
nonlinear scale debits, and reflection positivity of the sliced
integration.  None is hidden by the exact telescope.
\end{assessment}

\begin{programme}[Ninth-series V67: reflection test for rational slices]
\label{prog:n9v66-V67}
V67 will analyze the Osterwalder--Schrader reflection form of massive
resolvents and their differences.  It will give a spectral/Stieltjes
criterion for a rational covariance to be reflection positive and test
whether the positive slices $C_j$ pass it.  If they fail, the programme
will go upstream to a heat-kernel or Markov-time decomposition.
\end{programme}

\begin{firewall}[Claim boundary after ninth-series V66]
\label{fire:n9v66-boundary}
V66 proves the positive resolvent algebra, energy norms, derivative
formulas, and conditional quasi-locality.  It does not prove the
absolute polymer row, scale-debit summability for the physical
trajectory, reflection positivity of individual slices, chiral phase
or anomaly control, closure of the infrared sector, OS reconstruction,
or the full SM--Einstein continuum.
\end{firewall}

\begin{theorem}[Ninth-series V66 result]
\label{thm:n9v66-result}
After exact kernel removal, the massless Einstein covariance has the
positive resolvent telescope \eqref{eq:n9v66-full-telescope}.  A
geometric slice has energy norm at most $(1-q)/(1+q)$ and localization
length $O(m_j^{-1})$ under QWC.  The exact metric derivatives are
\eqref{eq:n9v66-R-first}--\eqref{eq:n9v66-R-second}; the infrared
remainder stays explicit; and per-slice smallness must be integrated over
the logarithmic scale count \eqref{eq:n9v66-log-count}.
\end{theorem}

\begin{proof}
Combine \Cref{thm:n9v66-telescope,prop:n9v66-slice-norm,
prop:n9v66-log-count,lem:n9v66-mass-locality,
prop:n9v66-resolvent-derivatives,prop:n9v66-remainder-bounds,
thm:n9v66-PRSC}.
\end{proof}

\paragraph{Parent-version record.}
V66 was formed from
\texttt{Renewal\_SM\_Einstein\_Continuum\_Research\_Notes\_V65.tex}.
The V65 parent had $24\,352$ logical source lines, $877\,499$ bytes, and
SHA--256
\[
 \texttt{3E9A3152985189C0E31B708D85336F34B7EEF3836DB3FB43C417348876D9B788}.
\]

% =============================================================================
% END APPENDED NINTH-SERIES V66 MATERIAL
% =============================================================================
% =============================================================================
% BEGIN APPENDED NINTH-SERIES V67 MATERIAL
% =============================================================================

\section{Ninth-series V67: reflection test and a spatial-shell repair}
\label{sec:v9v67}

\subsection{Reflection positivity for a Gaussian covariance}

V66 constructed positive resolvent differences but left their
Osterwalder--Schrader status open.  Let time reflection be
$\vartheta(t,\mathbf x)=(-t,\mathbf x)$ and let $\Theta f=f\circ\vartheta$
with the appropriate tensor conjugation.  A covariance $C$ is reflection
positive if
\begin{equation}
 (f,f)_{\rm OS}:=\langle\Theta f,Cf\rangle\ge0
 \label{eq:n9v67-OS-form}
\end{equation}
for every test function supported at positive time.

For a translation-invariant scalar time covariance at fixed spatial
mode, write
\begin{equation}
 \widehat C(p_0)=F(p_0^2).
 \label{eq:n9v67-time-symbol}
\end{equation}

\begin{theorem}[Rational time-spectral criterion]
\label{thm:n9v67-rational-criterion}
Suppose
\begin{equation}
 F(x)=\sum_{r=1}^N\frac{a_r}{x+E_r^2},
 \qquad 0<E_1<\cdots<E_N,
 \label{eq:n9v67-rational}
\end{equation}
with real residues.  The covariance is reflection positive if and only
if
\begin{equation}
 a_r\ge0\quad\hbox{for every }r.
 \label{eq:n9v67-positive-residues}
\end{equation}
\end{theorem}

\begin{proof}
Fourier inversion in time gives, for $t,s>0$,
\begin{equation}
 C(-t-s)=\sum_{r=1}^N\frac{a_r}{2E_r}e^{-E_r(t+s)}.
 \label{eq:n9v67-Hankel-kernel}
\end{equation}
Hence the OS form is
\begin{equation}
 (f,f)_{\rm OS}
 =\sum_{r=1}^N\frac{a_r}{2E_r}
 \left|\int_0^\infty e^{-E_rt}f(t)\,dt\right|^2,
 \label{eq:n9v67-OS-squares}
\end{equation}
which is nonnegative when all residues are.  Conversely, if $a_k<0$,
choose $N$ distinct positive times $t_i$.  The matrix
$(e^{-E_rt_i})_{r,i}$ is nonsingular after a generic arbitrarily small
perturbation of the $t_i$, because the exponentials with distinct
$E_r$ are linearly independent.  A finite linear combination of narrow
test bumps near the $t_i$ can therefore make all Laplace values in
\eqref{eq:n9v67-OS-squares} vanish except the $k$th.  The OS form is then
negative.  Approximation removes the narrow-bump idealization.
\end{proof}

\begin{remark}[Stieltjes form]
\label{rem:n9v67-Stieltjes}
The nonrational analogue is a positive time-spectral representation
\begin{equation}
 F(p_0^2)=\int_0^\infty
 \frac{2E}{p_0^2+E^2}\,d\rho(E),
 \qquad d\rho\ge0,
 \label{eq:n9v67-Stieltjes}
\end{equation}
equivalently a completely monotone positive-time kernel
$C(t)=\int e^{-Et}d\rho(E)$.  V67 needs only the rational criterion for
the explicit slices.
\end{remark}

\subsection{Resolvent differences fail the OS test}

Consider the free reflection-compatible operator
\begin{equation}
 A=-\partial_0^2+B,
 \qquad B\ge0,
 \label{eq:n9v67-product-operator}
\end{equation}
and fix a spatial eigenvalue $b\ge0$.  A V66 slice has time symbol
\begin{equation}
 F_j(p_0^2)
 =\frac1{p_0^2+b+m_j^2}
  -\frac1{p_0^2+b+m_{j-1}^2},
 \qquad m_j<m_{j-1}.
 \label{eq:n9v67-resolvent-slice-symbol}
\end{equation}

\begin{theorem}[Positive resolvent slices are not reflection positive]
\label{thm:n9v67-resolvent-fails}
Every nonzero slice \eqref{eq:n9v67-resolvent-slice-symbol} is a positive
operator covariance but fails reflection positivity.  Its two time
residues are $+1$ at
$E_j=\sqrt{b+m_j^2}$ and $-1$ at
$E_{j-1}=\sqrt{b+m_{j-1}^2}$.
\end{theorem}

\begin{proof}
The scalar function in \eqref{eq:n9v67-resolvent-slice-symbol} is
nonnegative for real $p_0$, proving ordinary covariance positivity.
Its displayed partial fractions have one negative residue.  Apply
\Cref{thm:n9v67-rational-criterion}.
\end{proof}

\begin{proposition}[Explicit two-time negative test]
\label{prop:n9v67-explicit-negative}
Let $0<E_j<E_{j-1}$ and choose $0<t_1<t_2$.  For coefficients
\begin{equation}
 c_1=e^{-E_jt_2},
 \qquad c_2=-e^{-E_jt_1},
 \label{eq:n9v67-test-coefficients}
\end{equation}
the positive-residue Laplace value vanishes, while
\begin{equation}
 c_1e^{-E_{j-1}t_1}+c_2e^{-E_{j-1}t_2}\ne0.
 \label{eq:n9v67-negative-mode-survives}
\end{equation}
Therefore a pair of narrow positive-time bumps with these coefficients
has strictly negative OS norm for the slice.
\end{proposition}

\begin{proof}
The first Laplace value is
$c_1e^{-E_jt_1}+c_2e^{-E_jt_2}=0$.  If the second vanished, division by
the nonzero exponentials would give
$e^{-(E_{j-1}-E_j)(t_1-t_2)}=1$, impossible because both differences are
nonzero.  Insert the two values in \eqref{eq:n9v67-OS-squares}.
\end{proof}

\begin{decision}[Status of the V66 telescope]
\label{dec:n9v67-resolvent-status}
The V66 resolvent slices remain valid for positive Gaussian comparison,
relative forms, quasi-local inverse estimates, and Feshbach bookkeeping.
They must not be used as independently integrated OS-positive RG layers.
\end{decision}

\subsection{Reflection-preserving spatial sandwich decomposition}

The failure came from subtracting time poles.  Instead decompose only
the spatial operator.  Let $B>0$ on its declared physical subspace and
choose bounded real functions $\psi_j$ such that
\begin{equation}
 \sum_{j=0}^J\psi_j(B)^2=I.
 \label{eq:n9v67-spatial-square-partition}
\end{equation}
Define
\begin{equation}
 Q_j=\psi_j(B),
 \qquad
 C_j^{\rm sp}=Q_j(-\partial_0^2+B)^{-1}Q_j.
 \label{eq:n9v67-spatial-slices}
\end{equation}

\begin{theorem}[Exact reflection-positive spatial decomposition]
\label{thm:n9v67-spatial-RP}
Assume $Q_j$ acts only on spatial and internal variables and commutes
with time reflection.  Then
\begin{equation}
 C=(-\partial_0^2+B)^{-1}=\sum_{j=0}^JC_j^{\rm sp},
 \label{eq:n9v67-spatial-sum}
\end{equation}
every $C_j^{\rm sp}$ is a positive covariance, and every slice is
reflection positive.
\end{theorem}

\begin{proof}
Since $Q_j$ commutes with $C$, the square partition gives
\begin{equation}
 \sum_jQ_jCQ_j=C\sum_jQ_j^2=C.
 \end{equation}
Ordinary positivity follows from
$\langle f,Q_jCQ_jf\rangle\ge0$.  The time kernel is
\begin{equation}
 C_j^{\rm sp}(t,s)
 =Q_j\frac{e^{-|t-s|\sqrt B}}{2\sqrt B}Q_j.
 \label{eq:n9v67-spatial-time-kernel}
\end{equation}
For $f$ supported at $t>0$, reflection replaces $|t-s|$ by $t+s$, and
\begin{align}
 \langle\Theta f,C_j^{\rm sp}f\rangle
 &=\frac12\left\lVert
  \int_0^\infty B^{-1/4}e^{-t\sqrt B}Q_jf(t)\,dt
 \right\rVert^2\ge0.
 \label{eq:n9v67-spatial-OS-square}
\end{align}
\end{proof}

\begin{corollary}[Independent OS-positive spatial shells]
\label{cor:n9v67-independent-shells}
Independent Gaussian fields with covariances $C_j^{\rm sp}$ sum to the
full free field and each has a nonnegative OS pre-Hilbert form.  If
$\psi_j$ is supported where
$\kappa_j^2\le B\le4\kappa_j^2$, its time kernel decays at least as
$e^{-\kappa_j|t-s|}$.
\end{corollary}

\begin{proof}
Covariances add by \eqref{eq:n9v67-spatial-sum}; reflection positivity
is \Cref{thm:n9v67-spatial-RP}.  On the support of $Q_j$,
$\sqrt B\ge\kappa_j$ in \eqref{eq:n9v67-spatial-time-kernel}.
\end{proof}

\subsection{Infrared and locality conditions}

\begin{firewall}[Spatial zero mode]
\label{fire:n9v67-spatial-zero}
The formula $B^{-1/4}$ in \eqref{eq:n9v67-spatial-OS-square} requires
the zero spatial mode to be removed, pinned, compactified, or retained
in a separately constructed infrared sector.  A spatial partition does
not normalize the massless constant mode.
\end{firewall}

\begin{proposition}[Spatial quasi-locality interface]
\label{prop:n9v67-spatial-locality}
Suppose $B$ is a finite-range uniformly elliptic spatial operator and
the complete filtered Poisson multiplier satisfies, for $\tau\ge0$,
\begin{equation}
 \left\lVert P_xQ_j\frac{e^{-\tau\sqrt B}}{2\sqrt B}Q_jP_y\right\rVert
 \le A_j\kappa_j^{-1}
 e^{-c\kappa_jd_{\rm sp}(x,y)}e^{-\kappa_j\tau}.
 \label{eq:n9v67-filter-locality}
\end{equation}
Then $C_j^{\rm sp}$ is spatially quasi-local at scale
$\kappa_j^{-1}$ and exponentially local in time at rate $\kappa_j$.
Products with fixed-range marked insertions retain any slightly weakened
exponential rate after collar convolution.
\end{proposition}

\begin{proof}
Set $\tau=|t-s|$ in
\eqref{eq:n9v67-spatial-time-kernel} and apply
\eqref{eq:n9v67-filter-locality}.  Fixed-range insertions shift the
spatial distance by their collar radii; convolution of exponentially
weighted rows preserves every strictly smaller exponential rate.
\end{proof}
\begin{firewall}[Smooth does not automatically mean exponential]
\label{fire:n9v67-smooth-locality}
A $C^\infty$ compactly supported spectral multiplier generally gives
rapid off-diagonal decay, but an exponential estimate such as
\eqref{eq:n9v67-filter-locality} needs analytic functional calculus or a
specific heat-kernel construction.  The multiplier must be chosen with
the locality norm in view.
\end{firewall}

\begin{firewall}[Preferred reflection structure]
\label{fire:n9v67-preferred-time}
The product form $A=-\partial_0^2+B$ and the spatial sandwich use the
reflection time selected for OS reconstruction.  A general fluctuating
metric need not admit this global split.  The regulator must supply a
reflection-compatible foliation or transfer operator before
\Cref{thm:n9v67-spatial-RP} applies.
\end{firewall}

\subsection{Metric derivatives of a spatial filter}

If $B=B(g)$, both the propagator and $Q_j=\psi_j(B)$ vary with the slow
metric.  Their exact first derivative is
\begin{equation}
 D_gC_j^{\rm sp}[h]
 =(D_gQ_j[h])CQ_j+Q_j(D_gC[h])Q_j+Q_jC(D_gQ_j[h]).
 \label{eq:n9v67-slice-derivative}
\end{equation}

\begin{proposition}[Exact partition derivative ledger]
\label{prop:n9v67-partition-derivative}
Differentiating \eqref{eq:n9v67-spatial-square-partition} gives
\begin{equation}
 \sum_j\bigl((D_gQ_j[h])Q_j+Q_jD_gQ_j[h]\bigr)=0.
 \label{eq:n9v67-Q-cancellation}
\end{equation}
Differentiating the exact covariance sum gives the stronger ledger
\begin{equation}
 \sum_j\bigl((D_gQ_j[h])CQ_j+Q_jC(D_gQ_j[h])\bigr)
 =D_gC[h]-\sum_jQ_j(D_gC[h])Q_j.
 \label{eq:n9v67-moving-filter-ledger}
\end{equation}
Thus the moving-filter rows supply the exact repartition correction; they
vanish in the total only when $D_gC[h]$ commutes with every $Q_j$.
They remain real marked interactions at an individual scale.
\end{proposition}

\begin{proof}
The first formula is the derivative of $\sum_jQ_j^2=I$.  For every
metric, functional calculus gives the exact identity
$\sum_jQ_jCQ_j=C$.  Differentiate it and move the middle terms
$\sum_jQ_j(D_gC[h])Q_j$ to the right, obtaining
\eqref{eq:n9v67-moving-filter-ledger}.  If $D_gC[h]$ commutes with all
$Q_j$, that sum is $D_gC[h]\sum_jQ_j^2=D_gC[h]$.
\end{proof}
\begin{firewall}[Noncommuting metric derivative]
\label{fire:n9v67-noncommuting-derivative}
Although $Q_j$ commutes with $C$ at a fixed metric, $D_gQ_j[h]$ need not
commute with either.  It must be computed by divided differences,
resolvent calculus, or a heat-kernel formula; replacing it by
$\psi_j'(B)D_gB[h]$ is valid only in a commuting variation.
\end{firewall}

\subsection{Reflection-positive spatial-shell card}

\begin{definition}[Reflection-positive spatial-shell card (RPSSC)]
\label{def:n9v67-RPSSC}
RPSSC requires:
\begin{enumerate}
\item a reflection-compatible physical free operator
$-\partial_0^2+B$ with declared spatial zero-mode owner;
\item a real square partition \eqref{eq:n9v67-spatial-square-partition};
\item slice covariances formed by spatial sandwiching, not time-pole
subtraction;
\item a proved spatial off-diagonal norm for every filter;
\item exact first and second moving-filter metric derivatives;
\item marked interaction and determinant rows inside the scale CPRAC
budget;
\item a retained infrared spatial shell; and
\item compatibility of the interacting integration with reflection and
the exact gauge/constraint sector.
\end{enumerate}
\end{definition}

\begin{theorem}[RPSSC repairs the V66 reflection failure]
\label{thm:n9v67-RPSSC}
Under RPSSC, the free physical covariance is an exact sum of
OS-positive Gaussian spatial shells.  Each noninfrared shell has a time
gap and the declared spatial locality bound.  The V66 resolvent
differences may still control forms and inverses, but the actual
reflection-positive Gaussian integrations use
\eqref{eq:n9v67-spatial-slices}.
\end{theorem}

\begin{proof}
Use \Cref{thm:n9v67-resolvent-fails,dec:n9v67-resolvent-status,
thm:n9v67-spatial-RP,cor:n9v67-independent-shells,
prop:n9v67-spatial-locality,prop:n9v67-partition-derivative}.
\end{proof}

\begin{assessment}[Progress at ninth-series V67]
\label{ass:n9v67-status}
The reflection question has a definite answer.  Positive resolvent
differences fail OS positivity because their time-spectral measures have
a negative residue.  A spatial square partition repairs the problem:
it preserves the positive time spectral measure and gives an exact
sum of OS-positive shells, at the cost of a preferred reflection split
and moving-filter locality estimates.
\end{assessment}

\begin{programme}[Ninth-series V68: analytic spatial filters and divided differences]
\label{prog:n9v67-V68}
V68 will choose an explicit analytic square partition, prove its
spatial off-diagonal decay from the local operator $B$, and derive first
and second metric derivatives by resolvent divided differences.  The
target is the moving-filter row left open in RPSSC.
\end{programme}

\begin{firewall}[Claim boundary after ninth-series V67]
\label{fire:n9v67-boundary}
V67 proves the free Gaussian reflection test and spatial-shell repair.
It does not prove an exponentially local analytic square partition for
the physical Einstein operator, control the moving-filter derivatives,
preserve reflection positivity after the full interacting chiral
integration, close the infrared sector, reconstruct the OS theory, or
prove the full SM--Einstein continuum.
\end{firewall}

\begin{theorem}[Ninth-series V67 result]
\label{thm:n9v67-result}
Every nonzero V66 resolvent difference fails reflection positivity by
the negative-residue criterion
\eqref{eq:n9v67-positive-residues}.  For a reflection-compatible free
operator $-\partial_0^2+B$, the spatial sandwich covariances
$\psi_j(B)C\psi_j(B)$ form an exact OS-positive decomposition whenever
$\sum_j\psi_j(B)^2=I$.  Their use requires separate spatial locality and
moving-filter estimates.
\end{theorem}

\begin{proof}
Combine \Cref{thm:n9v67-rational-criterion,
thm:n9v67-resolvent-fails,prop:n9v67-explicit-negative,
thm:n9v67-spatial-RP,thm:n9v67-RPSSC}.
\end{proof}

\paragraph{Parent-version record.}
V67 was formed from
\texttt{Renewal\_SM\_Einstein\_Continuum\_Research\_Notes\_V66.tex}.
The V66 parent had $24\,797$ logical source lines, $892\,904$ bytes, and
SHA--256
\[
 \texttt{221C9FDC3C52AD57FDD2CD99CCA39D1991393600A426FEBB78BFC5BB0FD441C6}.
\]

% =============================================================================
% END APPENDED NINTH-SERIES V67 MATERIAL
% =============================================================================
% =============================================================================
% BEGIN APPENDED NINTH-SERIES V68 MATERIAL
% =============================================================================

\section{Ninth-series V68: smooth spatial shells and moving-filter bounds}
\label{sec:v9v68}

\subsection{The analytic exact-gap obstruction}

V67 repaired reflection positivity by decomposing only the spatial
operator $B$.  Its programme asked for an analytic square partition with
an exact spectral gap.  Those two requirements are incompatible.

\begin{lemma}[No nontrivial analytic exact high-pass filter]
\label{lem:n9v68-analytic-gap}
Let $f$ be real analytic on a connected open neighbourhood of
$[0,\infty)$.  If there is a $\lambda_0>0$ such that
\begin{equation}
 f(\lambda)=0\qquad(0\le\lambda\le\lambda_0),
 \label{eq:n9v68-analytic-zero-interval}
\end{equation}
then $f$ vanishes identically.
\end{lemma}

\begin{proof}
The zero set contains an interval and hence has an accumulation point in
the analytic domain.  The identity theorem gives $f\equiv0$ on the
connected domain.
\end{proof}

\begin{corollary}[Analytic filters leak into the massless time spectrum]
\label{cor:n9v68-analytic-leak}
A nonzero analytic spatial filter cannot both vanish below
$\lambda_0$ and give a strict time gap in the V67 covariance
$f(B)(-\partial_0^2+B)^{-1}f(B)$.  If $f$ has only a finite-order zero
at $0$, the filtered equal-time symbol retains a fractional threshold
power and generally has an algebraic spatial tail.
\end{corollary}

\begin{proof}
The first claim is \Cref{lem:n9v68-analytic-gap}.  If
$f(\lambda)=c\lambda^r+O(\lambda^{r+1})$, then the equal-time multiplier
is
\begin{equation}
 \frac{f(\lambda)^2}{2\sqrt\lambda}
 =\frac{c^2}{2}\lambda^{2r-1/2}
  +O(\lambda^{2r+1/2}),
 \label{eq:n9v68-fractional-threshold}
\end{equation}
which is nonanalytic at the massless threshold.  Exponential kernel
decay would give analytic continuation of the Fourier multiplier to a
strip, contradicting this branch behaviour in the flat model.
\end{proof}

\begin{example}[First-order leakage in three spatial dimensions]
\label{ex:n9v68-first-order-tail}
For $B=-\Delta_{\mathbf x}$ in three dimensions and $r=1$,
\eqref{eq:n9v68-fractional-threshold} behaves as
$|\mathbf p|^3$.  After a smooth ultraviolet cutoff, its nonlocal kernel
tail has homogeneity $|\mathbf x|^{-6}$.  It has several finite moments,
but it is not exponentially local.
\end{example}

\begin{decision}[Filter regularity selected]
\label{dec:n9v68-filter-choice}
The reflection-positive shell construction will use a
$C^\infty$ exact-gap square partition.  It gives arbitrary prescribed
polynomial locality order.  Exponential spatial locality is not claimed
for an exact massless spectral shell.
\end{decision}

\subsection{An exact smooth square partition}

Choose $\vartheta\in C^\infty([0,\infty);[0,\pi/2])$ such that
\begin{equation}
 \vartheta(t)=0\quad(t\le1),
 \qquad
 \vartheta(t)=\pi/2\quad(t\ge4).
 \label{eq:n9v68-angle-cutoff}
\end{equation}
Let $\kappa_0>\kappa_1>\cdots>\kappa_J>0$ and set
\begin{equation}
 h_j(\lambda)=\sin\vartheta(\lambda/\kappa_j^2),
 \qquad
 l_j(\lambda)=\cos\vartheta(\lambda/\kappa_j^2).
 \label{eq:n9v68-high-low}
\end{equation}
Define
\begin{align}
 \psi_0&=h_0,
 \label{eq:n9v68-psi0}\\
 \psi_j&=\left(\prod_{r=0}^{j-1}l_r\right)h_j,
 \qquad 1\le j\le J,
 \label{eq:n9v68-psij}\\
 \psi_{\rm IR}&=\prod_{r=0}^Jl_r.
 \label{eq:n9v68-psiIR}
\end{align}

\begin{theorem}[Smooth exact square partition]
\label{thm:n9v68-square-partition}
The functions in \eqref{eq:n9v68-psi0}--
\eqref{eq:n9v68-psiIR} are real and smooth and satisfy
\begin{equation}
 \sum_{j=0}^J\psi_j(\lambda)^2
 +\psi_{\rm IR}(\lambda)^2=1
 \qquad(\lambda\ge0).
 \label{eq:n9v68-square-partition}
\end{equation}
For $j\ge1$,
\begin{equation}
 \operatorname{supp}\psi_j
 \subset[\kappa_j^2,4\kappa_{j-1}^2],
 \label{eq:n9v68-annular-support}
\end{equation}
while $\psi_{\rm IR}$ is supported below $4\kappa_J^2$ and equals one
below $\kappa_J^2$.
\end{theorem}

\begin{proof}
Since $h_j^2+l_j^2=1$, repeated use of
\begin{equation}
 1=h_0^2+l_0^2
 =h_0^2+l_0^2h_1^2+l_0^2l_1^2
 \end{equation}
telescopes to \eqref{eq:n9v68-square-partition}.  Flatness of
$\vartheta$ at the two constant regions makes every factor smooth.
The high factor $h_j$ vanishes below $\kappa_j^2$, while the preceding
factor $l_{j-1}$ vanishes above $4\kappa_{j-1}^2$.  The infrared support
and plateau follow in the same way.
\end{proof}

\begin{corollary}[Exact OS-positive shell covariance]
\label{cor:n9v68-OS-shells}
With $Q_j=\psi_j(B)$ and $Q_{\rm IR}=\psi_{\rm IR}(B)$,
\begin{equation}
 C=\sum_{j=0}^JQ_jCQ_j+Q_{\rm IR}CQ_{\rm IR},
 \qquad C=(-\partial_0^2+B)^{-1},
 \label{eq:n9v68-OS-decomposition}
\end{equation}
and every term is reflection positive.  Each $j\ge1$ shell has time
decay at least $e^{-\kappa_j|t-s|}$.
\end{corollary}

\begin{proof}
Apply \Cref{thm:n9v67-spatial-RP} and
\eqref{eq:n9v68-square-partition}.  The lower support bound in
\eqref{eq:n9v68-annular-support} gives the time decay through
\eqref{eq:n9v67-spatial-time-kernel}.
\end{proof}

\subsection{Flat spatial off-diagonal bounds}

Let the spatial dimension be $s$ and first take
$B=-\Delta_{\mathbf x}$ on a flat torus or Euclidean lattice away from
the Brillouin boundary.  For a noninfrared shell define
\begin{equation}
 F_{j,\tau}(\lambda)
 =\frac{\psi_j(\lambda)^2}{2\sqrt\lambda}
 e^{-\tau\sqrt\lambda}.
 \label{eq:n9v68-filtered-Poisson}
\end{equation}

\begin{theorem}[Arbitrary-order spatial locality]
\label{thm:n9v68-flat-locality}
Assume a geometric scale ratio
$c_q\kappa_j\le\kappa_{j-1}\le C_q\kappa_j$.  For every integer
$N\ge0$, the kernel of $F_{j,\tau}(-\Delta)$ satisfies
\begin{equation}
 |K_{j,\tau}(\mathbf x,\mathbf y)|
 \le C_N\kappa_j^{s-1}
 e^{-c\kappa_j\tau}
 (1+\kappa_j|\mathbf x-\mathbf y|)^{-N},
 \label{eq:n9v68-flat-kernel}
\end{equation}
with constants uniform in $j$ away from the ultraviolet lattice edge.
\end{theorem}

\begin{proof}
In Fourier variables the multiplier is
$F_{j,\tau}(|\mathbf p|^2)$ and is supported where
$\kappa_j\le|\mathbf p|\le C\kappa_j$.  Rescale
$\mathbf p=\kappa_j\mathbf q$.  The factor
$1/(2|\mathbf p|)$ gives $\kappa_j^{-1}$ and Fourier volume gives
$\kappa_j^s$.  The rescaled multiplier and all its derivatives are
uniformly bounded by a polynomial in $\kappa_j\tau$ times
$e^{-\kappa_j\tau}$, which is bounded by
$C_Ne^{-c\kappa_j\tau}$.  Integrate by parts $N$ times in a Fourier
direction parallel to $\mathbf x-\mathbf y$ to obtain
\eqref{eq:n9v68-flat-kernel}.
\end{proof}

\begin{corollary}[Moment order selected by smoothness]
\label{cor:n9v68-moments}
For any prescribed spatial moment $q\ge0$, choose $N>s+q$ in
\eqref{eq:n9v68-flat-kernel}.  Then the absolute spatial $q$th moment is
finite and has scale
\begin{equation}
 \int |\mathbf x-\mathbf y|^q
 |K_{j,\tau}(\mathbf x,\mathbf y)|,d\mathbf y
 \le C_{N,q}\kappa_j^{-q-1}e^{-c\kappa_j\tau}.
 \label{eq:n9v68-shell-moment}
\end{equation}
\end{corollary}

\begin{proof}
Insert \eqref{eq:n9v68-flat-kernel}, use polar coordinates, and rescale
$r$ by $\kappa_j$.
\end{proof}

\begin{firewall}[Uniform elliptic backgrounds need a multiplier theorem]
\label{fire:n9v68-variable-background}
The Fourier proof is exact on a flat translation-invariant background.
For a variable spatial metric, the corresponding claim requires a
uniform off-diagonal spectral-multiplier theorem for $B(g)$, including
boundary and constraint-reduction constants.  Ellipticity alone is not
a substitute for that theorem.
\end{firewall}

\subsection{Divided differences for moving filters}

Let $B=B(g)$ be a self-adjoint finite-regulator matrix in a fixed
physical-bundle trivialization, and let
\begin{equation}
 H_h=D_gB[h],
 \qquad H_{h,k}=D_g^2B[h,k].
 \label{eq:n9v68-B-marks}
\end{equation}
For a smooth compactly supported $f$, choose an almost-analytic
extension $\widetilde f$.

\begin{theorem}[Exact first and second filter derivatives]
\label{thm:n9v68-filter-derivatives}
The first derivative is
\begin{equation}
 D_gf(B)[h]
 =\frac1\pi\int_{\mathbb C}
 \overline\partial\widetilde f(z)
 (z-B)^{-1}H_h(z-B)^{-1},d^2z.
 \label{eq:n9v68-filter-first}
\end{equation}
The second derivative is
\begin{align}
 D_g^2f(B)[h,k]
 &=\frac1\pi\int_{\mathbb C}
 \overline\partial\widetilde f(z)
 \Bigl[
 R_zH_kR_zH_hR_z+R_zH_hR_zH_kR_z
 +R_zH_{h,k}R_z
 \Bigr]d^2z,
 \label{eq:n9v68-filter-second}\\
 R_z&=(z-B)^{-1}.
\end{align}
These formulas include all noncommuting divided-difference rows.
\end{theorem}

\begin{proof}
The Helffer--Sj\"ostrand formula is
\begin{equation}
 f(B)=\frac1\pi\int_{\mathbb C}
 \overline\partial\widetilde f(z)(z-B)^{-1},d^2z.
 \end{equation}
Differentiate under the finite-regulator integral and use
$D_gR_z[h]=R_zH_hR_z$.  A second differentiation yields the two
orderings of $H_h,H_k$ and the direct $H_{h,k}$ term.
\end{proof}

\begin{corollary}[Scale of the moving-filter rows]
\label{cor:n9v68-filter-scaling}
For a uniformly scaled shell profile
$f_j(\lambda)=f(\lambda/\kappa_j^2)$, there are profile constants
$C_1,C_2$ such that
\begin{align}
 \lVert D_gf_j(B)[h]\rVert
 &\le C_1\kappa_j^{-2}\lVert H_h\rVert,
 \label{eq:n9v68-Df-bound}\\
 \lVert D_g^2f_j(B)[h,k]\rVert
 &\le C_2\kappa_j^{-4}
       \lVert H_h\rVert\lVert H_k\rVert
   +C_1\kappa_j^{-2}\lVert H_{h,k}\rVert.
 \label{eq:n9v68-D2f-bound}
\end{align}
The same bounds hold for the finitely overlapping products defining
$\psi_j$, with constants depending only on the geometric scale ratio.
\end{corollary}

\begin{proof}
Rescale $z=\kappa_j^2\zeta$ in
\eqref{eq:n9v68-filter-first}--
\eqref{eq:n9v68-filter-second}.  Each resolvent contributes
$\kappa_j^{-2}$, the area element contributes $\kappa_j^4$, and the
almost-analytic density has the compensating profile scaling.  This
gives one net $\kappa_j^{-2}$ per first variation and
$\kappa_j^{-4}$ for the product of two first variations.  The direct
second variation has only the first-derivative scale.  The product rule
and finite overlap give the last statement.
\end{proof}

\begin{corollary}[Dimensionless weak-field mark]
\label{cor:n9v68-dimensionless-mark}
If on the $j$th shell
\begin{equation}
 \lVert H_h\rVert\le\epsilon_{1,j}\kappa_j^2,
 \qquad
 \lVert H_{h,k}\rVert\le\epsilon_{2,j}\kappa_j^2,
 \label{eq:n9v68-relative-B-marks}
\end{equation}
then
\begin{align}
 \lVert D_gQ_j[h]\rVert&\le C\epsilon_{1,j},
 \label{eq:n9v68-DQ-relative}\\
 \lVert D_g^2Q_j[h,k]\rVert
 &\le C(\epsilon_{1,j}^2+epsilon_{2,j}).
 \label{eq:n9v68-D2Q-relative}
\end{align}
Thus the moving-filter debit is small only when the metric derivative of
$B$ is small relative to the shell energy.
\end{corollary}

\begin{proof}
Insert \eqref{eq:n9v68-relative-B-marks} into
\eqref{eq:n9v68-Df-bound}--\eqref{eq:n9v68-D2f-bound} and use the product
construction of $Q_j$.
\end{proof}

\begin{firewall}[The relative shell mark is not automatic]
\label{fire:n9v68-relative-mark}
For a general metric variation, $D_gB[h]$ is a principal second-order
operator and need not have small norm relative to $\kappa_j^2$.
Equation \eqref{eq:n9v68-relative-B-marks} must follow from the chosen
slow/fast metric split, a small fluctuation chart, or a form-relative
estimate.  Smooth filtering alone does not provide it.
\end{firewall}

\subsection{Smooth spatial-shell filter card}

\begin{definition}[Smooth spatial-shell filter card (SSFC)]
\label{def:n9v68-SSFC}
SSFC requires:
\begin{enumerate}
\item the exact smooth square partition
\eqref{eq:n9v68-square-partition};
\item the OS-positive sandwich covariance
\eqref{eq:n9v68-OS-decomposition};
\item a chosen polynomial locality order exceeding every moment and
cluster summability order used later;
\item a uniform spectral-multiplier theorem on the accepted spatial
metric charts;
\item exact divided-difference derivatives
\eqref{eq:n9v68-filter-first}--\eqref{eq:n9v68-filter-second};
\item a scale-relative bound for $D_gB$ and $D_g^2B$;
\item an explicit infrared shell and ultraviolet lattice-edge owner;
and
\item preservation of the exact physical constraint and reflection
structure.
\end{enumerate}
\end{definition}

\begin{theorem}[SSFC closes the free moving-filter row]
\label{thm:n9v68-SSFC}
Under SSFC, the free physical covariance is decomposed into exact
OS-positive shells with enough polynomial spatial locality for every
declared finite moment.  The first and second moving-filter metric rows
are given by the noncommuting formulas
\eqref{eq:n9v68-filter-first}--\eqref{eq:n9v68-filter-second} and obey
the dimensionless bounds
\eqref{eq:n9v68-DQ-relative}--\eqref{eq:n9v68-D2Q-relative}.
\end{theorem}

\begin{proof}
Use \Cref{thm:n9v68-square-partition,cor:n9v68-OS-shells,
thm:n9v68-flat-locality,cor:n9v68-moments,
thm:n9v68-filter-derivatives,cor:n9v68-dimensionless-mark}; SSFC item 4
extends the flat multiplier bound to the accepted variable backgrounds.
\end{proof}

\begin{assessment}[Progress at ninth-series V68]
\label{ass:n9v68-status}
The attempted analytic exact-gap filter has been rejected by a precise
identity-theorem obstruction.  Smooth exact-gap filters preserve OS
positivity and provide arbitrary polynomial locality order.  Their
moving metric rows are now exact divided differences, with smallness
reduced to the scale-relative derivatives of the physical spatial
operator.
\end{assessment}

\begin{programme}[Ninth-series V69: one interacting reflection step]
\label{prog:n9v68-V69}
V69 will integrate one SSFC Gaussian shell in the presence of a local
reflection-factorized interaction.  It will prove the precise condition
under which reflection positivity passes to the retained measure and
will track the connected metric marks and the shell CPRAC debit.
\end{programme}

\begin{firewall}[Claim boundary after ninth-series V68]
\label{fire:n9v68-boundary}
V68 proves the analytic obstruction, smooth square partition, flat
off-diagonal estimate, and divided-difference calculus.  It does not
prove the needed uniform multiplier theorem on fluctuating physical
metrics, the scale-relative metric marks, interacting reflection
positivity, chiral phase or anomaly control, infrared closure, OS
reconstruction, or the full SM--Einstein continuum.
\end{firewall}

\begin{theorem}[Ninth-series V68 result]
\label{thm:n9v68-result}
No nonzero analytic spatial filter has an exact low-energy gap.  The
smooth construction \eqref{eq:n9v68-psi0}--
\eqref{eq:n9v68-psiIR} instead gives an exact square partition and
OS-positive shells with arbitrary polynomial locality order.  Their
noncommuting moving-filter derivatives are exactly
\eqref{eq:n9v68-filter-first}--\eqref{eq:n9v68-filter-second} and are
dimensionless under the relative shell bounds
\eqref{eq:n9v68-relative-B-marks}.
\end{theorem}

\begin{proof}
Combine \Cref{lem:n9v68-analytic-gap,cor:n9v68-analytic-leak,
thm:n9v68-square-partition,cor:n9v68-OS-shells,
thm:n9v68-flat-locality,thm:n9v68-filter-derivatives,
thm:n9v68-SSFC}.
\end{proof}

\paragraph{Parent-version record.}
V68 was formed from
\texttt{Renewal\_SM\_Einstein\_Continuum\_Research\_Notes\_V67.tex}.
The V67 parent had $25\,197$ logical source lines, $907\,876$ bytes, and
SHA--256
\[
 \texttt{B041D31AB77AC7AFE9889F6E6EFC4AA4999F1EB127855093A4D8518B643CAF4F}.
\]

% =============================================================================
% END APPENDED NINTH-SERIES V68 MATERIAL
% =============================================================================
% =============================================================================
% BEGIN APPENDED NINTH-SERIES V69 MATERIAL
% =============================================================================

\section{Ninth-series V69: one exact interacting reflection step}
\label{sec:v9v69}

\subsection{The reflection cone}

V68 constructed OS-positive free spatial shells.  Interacting
reflection positivity is a statement about the exact weighted measure,
not about positivity of its covariance or of a truncated effective
action.  Let $\Theta$ be an antilinear reflection involution and let
${\cal A}_+$ be the algebra of bounded observables supported at positive
time.  A probability measure $\mu$ is reflection positive if
\begin{equation}
 \mathbb E_\mu[(\Theta F)F]\ge0
 \qquad(F\in{\cal A}_+).
 \label{eq:n9v69-RP-definition}
\end{equation}

Define the algebraic reflection cone
\begin{equation}
 {\cal K}_+=left\{
  \sum_{r=1}^N(\Theta G_r)G_r:
  N<\infty,\ G_r\in{\cal A}_+
 \right\},
 \label{eq:n9v69-reflection-cone}
\end{equation}
and use its $L^1(\mu)$ closure when limits are needed.

\begin{lemma}[The reflection cone is multiplicative]
\label{lem:n9v69-cone-product}
If $W_1,W_2\in{\cal K}_+$, then $W_1W_2\in{\cal K}_+$.
In particular, $(\Theta G)G\in{\cal K}_+$.
\end{lemma}

\begin{proof}
Write $W_1=\sum_r(\Theta F_r)F_r$ and
$W_2=\sum_s(\Theta G_s)G_s$.  Since $\Theta$ reverses complex scalars
and preserves products of commuting configuration observables,
\begin{equation}
 W_1W_2=\sum_{r,s}\Theta(F_rG_s)(F_rG_s).
\end{equation}
\end{proof}

\begin{theorem}[Reflection-cone weighting]
\label{thm:n9v69-cone-weighting}
Let $\mu$ be reflection positive.  Suppose $W\ge0$ is integrable,
$0<Z_W=\mathbb E_\mu W<\infty$, and $W$ belongs to the
$L^1$ reflection-cone closure with convergence sufficient for the
quadratic forms below.  Then
\begin{equation}
 d\mu_W=Z_W^{-1}W\,d\mu
 \label{eq:n9v69-weighted-measure}
\end{equation}
is reflection positive.
\end{theorem}

\begin{proof}
For an algebraic cone element
$W=\sum_r(\Theta G_r)G_r$ and $F\in{\cal A}_+$,
\begin{equation}
 \mathbb E_\mu[(\Theta F)FW]
 =\sum_r\mathbb E_\mu[Theta(FG_r)(FG_r)]\ge0.
 \label{eq:n9v69-weighted-square}
\end{equation}
Pass to the declared $L^1$ quadratic-form limit and divide by $Z_W$.
\end{proof}

\begin{corollary}[Positive/negative-time factorization]
\label{cor:n9v69-factorized-action}
Suppose
\begin{equation}
 V=V_++\Theta V_++V_0,
 \label{eq:n9v69-action-split}
\end{equation}
where $V_+$ is positive-time measurable and
$e^{-V_0/\hbar}$ belongs to the reflection-cone closure.  If the full
Boltzmann factor is nonnegative and integrable, then
\begin{equation}
 e^{-V/\hbar}
 =\Theta(e^{-V_+/\hbar})e^{-V_+/\hbar}e^{-V_0/\hbar}
 \label{eq:n9v69-density-factor}
\end{equation}
preserves reflection positivity.
\end{corollary}

\begin{proof}
Use \Cref{lem:n9v69-cone-product,thm:n9v69-cone-weighting}.
\end{proof}

\begin{firewall}[Nonnegative is weaker than reflection positive]
\label{fire:n9v69-nonnegative}
A pointwise nonnegative reflection-invariant density need not lie in
${\cal K}_+$.  Cross-plane plaquette, hopping, Yukawa, and determinant
factors require a positive character expansion, transfer-kernel
factorization, or another explicit cone proof.  Ordinary positivity is
not enough.
\end{firewall}

\subsection{Products and exact shell marginalization}

Let $\mu_{<}$ be a reflection-positive retained Gaussian measure and
let $\mu_j$ be one V68 shell Gaussian measure.

\begin{lemma}[Product of reflection-positive measures]
\label{lem:n9v69-product-RP}
The product $\mu_{<}\otimes\mu_j$ is reflection positive under the
diagonal reflection.
\end{lemma}

\begin{proof}
For finite sums $F=\sum_aF_a^{<}F_a^j$, the OS form is the quadratic
form of the Schur product of the two positive Gram matrices
\begin{equation}
 M_{ab}^{<}=\mathbb E_{\mu_<}[(\Theta F_a^{<})F_b^{<}],
 \qquad
 M_{ab}^j=\mathbb E_{\mu_j}[(\Theta F_a^j)F_b^j].
 \end{equation}
The Schur product theorem makes it nonnegative.  Density extends the
result to the completed positive-time algebra.
\end{proof}

Let $\psi$ denote the retained field and $\zeta_j$ the shell field.  Set
\begin{equation}
 W_j(\psi,\zeta_j)=
 \exp[-V_j(\psi,\zeta_j)/\hbar].
 \label{eq:n9v69-joint-density}
\end{equation}

\begin{theorem}[One exact interacting shell preserves OS positivity]
\label{thm:n9v69-one-shell-RP}
Assume $W_j\ge0$ is integrable and lies in the reflection-cone closure
of the joint product measure.  Then
\begin{equation}
 d\mu_j^{\rm joint}
 =Z_j^{-1}W_j(\psi,\zeta_j)
 d\mu_< (\psi)d\mu_j(\zeta_j)
 \label{eq:n9v69-joint-measure}
\end{equation}
is reflection positive.  Its exact retained marginal
\begin{equation}
 d\mu_<^{\rm eff}(\psi)
 =Z_j^{-1}R_j(\psi)d\mu_< (\psi),
 \qquad
 R_j(\psi)=\int W_j(\psi,\zeta_j)d\mu_j(\zeta_j),
 \label{eq:n9v69-retained-marginal}
\end{equation}
is also reflection positive.
\end{theorem}

\begin{proof}
The product measure is reflection positive by
\Cref{lem:n9v69-product-RP}; cone weighting gives the joint statement.
For $F$ depending only on $\psi$,
\begin{equation}
 \mathbb E_{\mu_<^{\rm eff}}[(\Theta F)F]
 =\mathbb E_{\mu_j^{\rm joint}}[(\Theta F)F]\ge0,
 \end{equation}
which proves the marginal statement.  Fubini gives the density in
\eqref{eq:n9v69-retained-marginal}.
\end{proof}

\begin{corollary}[Successive exact shell marginals]
\label{cor:n9v69-successive-shells}
If at every step the new shell product and its complete interaction
density satisfy the hypotheses of
\Cref{thm:n9v69-one-shell-RP}, reflection positivity is preserved by
induction through any finite number of exact shell integrations.
\end{corollary}

\begin{proof}
Use the retained marginal from one step as the reflection-positive base
measure for the next and repeat the theorem.
\end{proof}

\begin{firewall}[The effective density need not display a simple cone]
\label{fire:n9v69-effective-cone}
The proof establishes reflection positivity of the retained marginal.
It does not assert that $R_j(\psi)$ admits a short local factorization of
the form \eqref{eq:n9v69-density-factor}.  The cone structure may be
encoded nonlocally in the exact marginal.
\end{firewall}

\begin{firewall}[Truncating the connected logarithm]
\label{fire:n9v69-truncation}
A KP or cumulant expansion may estimate
$U_j=-\hbar\log R_j$, but a finite truncation generally changes
$R_j$ and need not remain in the reflection cone.  Reflection positivity
belongs to the exact marginal; expansions are used to prove convergence
and locality without replacing it.
\end{firewall}

\subsection{Exact shell metric marks}

Let the shell covariance and interaction depend on a reflection-
compatible slow metric $g$.  Relative to a fixed reference
disintegration, let
\begin{equation}
 b_{j,h}=D_g\log d\mu_{j,g}[h]
 \label{eq:n9v69-shell-score}
\end{equation}
denote the complete centred shell score, including the V68 moving-filter
rows.  Define
\begin{equation}
 X_{j,h}=-\hbar^{-1}D_gV_j[h]+b_{j,h}.
 \label{eq:n9v69-complete-shell-mark}
\end{equation}
Let $\nu_{j,\psi}$ be the normalized conditional shell law at fixed
retained field.

\begin{theorem}[Exact interacting-shell Hessian]
\label{thm:n9v69-shell-Hessian}
Under a second-order differentiation domain,
\begin{align}
 U_j(\psi;g)&=-\hbar\log R_j(\psi;g),
 \label{eq:n9v69-Uj}\\
 D_gU_j[h]&=-\hbar\mathbb E_{\nu_{j,\psi}}X_{j,h},
 \label{eq:n9v69-Uj-first}\\
 D_g^2U_j[h,k]
 &=-\hbar\mathbb E_{\nu_{j,\psi}}D_gX_{j,h}[k]
   -\hbar\operatorname{Cov}_{\nu_{j,\psi}}
       (X_{j,h},X_{j,k}).
 \label{eq:n9v69-Uj-second}
\end{align}
Every covariance, reference, interaction, and moving-filter row occurs
once in the complete mark.
\end{theorem}

\begin{proof}
Differentiate the exact conditional partition in
\eqref{eq:n9v69-retained-marginal}.  The logarithmic derivative of its
complete integrand is \eqref{eq:n9v69-complete-shell-mark}.  The
normalized-expectation derivative is the mean derivative plus
covariance, as in V56 and V61.
\end{proof}

\begin{firewall}[Differentiating an inequality]
\label{fire:n9v69-differentiate-RP}
Reflection positivity of $\mu_g$ does not imply that the first or second
metric derivative of every OS quadratic form has a fixed sign.
Equation \eqref{eq:n9v69-Uj-second} is controlled by relative-form and
cluster estimates, not by differentiating the inequality
\eqref{eq:n9v69-RP-definition}.
\end{firewall}

\subsection{One-step CPRAC debit}

Let ${\mathfrak a}_{<}$ be the retained positive physical metric form
after the local RZMOC owners are assigned.  Suppose the direct row and
the complete covariance row obey
\begin{align}
 -\hbar\mathbb E D_gX_{j,h}[h]
 &\ge-\delta_j{\mathfrak a}_{<}(h),
 \label{eq:n9v69-direct-debit}\\
 \hbar\operatorname{Var}(X_{j,h})
 &\le c_j{\mathfrak a}_{<}(h).
 \label{eq:n9v69-covariance-debit}
\end{align}

\begin{proposition}[Aggregate shell margin]
\label{prop:n9v69-shell-margin}
If the incoming retained Hessian satisfies
\begin{equation}
 {\mathfrak h}_{<}(h)
 \ge(1-\rho_{<}){\mathfrak a}_{<}(h),
 \label{eq:n9v69-incoming-margin}
\end{equation}
then after the exact shell integration,
\begin{equation}
 {\mathfrak h}_{<}^{\rm new}(h)
 \ge(1-\rho_{<}-\delta_j-c_j)
 {\mathfrak a}_{<}(h).
 \label{eq:n9v69-new-margin}
\end{equation}
The one-step strict margin is preserved if
$\rho_{<}+\delta_j+c_j<1$.
\end{proposition}

\begin{proof}
Add \eqref{eq:n9v69-Uj-second} to the incoming Hessian and use
\eqref{eq:n9v69-direct-debit}--
\eqref{eq:n9v69-covariance-debit}.
\end{proof}

\begin{proposition}[Marked row criterion for $c_j$]
\label{prop:n9v69-marked-row}
Decompose $X_{j,h}=\sum_xX_{j,x}(h_x)$ and let $A_x$ be a coercive local
reference on the shell-tested retained subspace.  If
\begin{equation}
 \sup_x\sum_y
 \left\lVert
 A_x^{-1/2}\operatorname{Cov}(X_{j,x},X_{j,y})A_y^{-1/2}
 \right\rVert
 \le r_j,
 \label{eq:n9v69-marked-row}
\end{equation}
then one may take $c_j=\hbar r_j$.  On the massless sector the same row
must instead pass the V63 $H^{-1}$ or gradient-factorization test.
\end{proposition}

\begin{proof}
Apply the V57 block Schur estimate to the covariance form and multiply
by $\hbar$.  The final sentence is the V63 constant-mode obstruction.
\end{proof}

\begin{firewall}[Per-shell margins must be summable]
\label{fire:n9v69-margin-sum}
Even if every $\delta_j+c_j$ is small, an infinite sequence of bounds
of the form \eqref{eq:n9v69-new-margin} closes only when the accumulated
debits are summable or are absorbed into running positive owners.  V66's
logarithmic scale count remains in force.
\end{firewall}

\subsection{Dynamical metric obstruction}

The preceding product theorem assumes a shell Gaussian measure
independent of the retained field.  In the SM--Einstein problem the
spatial operator $B(g)$, and hence $C_j^{\rm sp}(g)$, depends on the
retained metric.

\begin{proposition}[Pointwise conditional RP is insufficient]
\label{prop:n9v69-conditional-insufficient}
Suppose for every fixed $g$ the conditional shell measure
$\mu_j(d\zeta\mid g)$ is reflection positive.  This fact alone does not
imply that
\begin{equation}
 \mu_<(dg)\mu_j(d\zeta\mid g)
 \label{eq:n9v69-conditional-joint}
\end{equation}
is jointly reflection positive, because an OS test couples
$g_+$ to $\Theta g_+$ while the two conditional covariances are evaluated
at different full retained configurations.
\end{proposition}

\begin{proof}
The obstruction already appears in a two-label finite reflection system.
Let the positive-time retained label be $a\in\{1,2\}$, the reflected
negative-time label be $b$, and give a reflected pair $(a,b)$ the
nonnegative symmetric weight
\begin{equation}
 K=\begin{pmatrix}1&2\\2&1\end{pmatrix}.
 \label{eq:n9v69-two-label-kernel}
\end{equation}
Every diagonal conditional test is positive because $K_{aa}=1$.  The
joint OS form of a positive-time function $f$ is proportional to
$\sum_{a,b}\overline{f(b)}K_{ba}f(a)$.  For $f=(1,-1)$ this equals
$-2<0$, since $K$ has eigenvalue $-1$.  Thus positivity of all diagonal
conditional kernels does not imply positivity of the off-diagonal joint
kernel.  A metric-conditioned shell disintegration has exactly this
logical form, with retained half-configurations replacing the two
labels.
\end{proof}

\begin{definition}[Joint reflection-positive shell kernel]
\label{def:n9v69-joint-kernel}
A metric-dependent shell kernel is jointly reflection positive if its
joint retained/shell reference measure already satisfies
\eqref{eq:n9v69-RP-definition}.  This may be proved by a transfer-matrix
construction, a reflection-cone density relative to a fixed product
reference, or a completely positive boundary kernel.
\end{definition}

\begin{firewall}[Chiral phase]
\label{fire:n9v69-chiral-phase}
A complex chiral determinant is not a nonnegative density in the
reflection cone.  Conjugate pairing can restore ordinary positivity,
but reflection-cone membership still needs a transfer or character
proof.  The V53--V56 analytic phase branch controls complex expectations;
it does not load \Cref{thm:n9v69-cone-weighting}.
\end{firewall}

\subsection{Interacting reflection shell card}

\begin{definition}[Interacting reflection-positive shell card (IRPSC)]
\label{def:n9v69-IRPSC}
IRPSC requires:
\begin{enumerate}
\item an OS-positive retained measure and SSFC shell covariance;
\item a jointly reflection-positive shell kernel when the covariance
depends on the retained metric;
\item a nonnegative complete interaction density in the reflection-cone
closure;
\item exact marginalization without truncating the retained density;
\item the complete metric mark \eqref{eq:n9v69-complete-shell-mark};
\item direct and covariance debits satisfying
\eqref{eq:n9v69-direct-debit}--\eqref{eq:n9v69-covariance-debit};
\item summability or running ownership of those debits across scales;
and
\item separate chiral phase, anomaly, gauge, and boundary-sector
control.
\end{enumerate}
\end{definition}

\begin{theorem}[IRPSC performs one admissible interacting scale step]
\label{thm:n9v69-IRPSC}
Under IRPSC, exact integration of one spatial Gaussian shell preserves
reflection positivity of the retained measure.  Its induced metric
source and Hessian are exactly
\eqref{eq:n9v69-Uj-first}--\eqref{eq:n9v69-Uj-second}, and the retained
strict relative margin is updated by
\eqref{eq:n9v69-new-margin}.  No truncated action is used to justify the
OS statement.
\end{theorem}

\begin{proof}
Use \Cref{thm:n9v69-one-shell-RP,thm:n9v69-shell-Hessian,
prop:n9v69-shell-margin,prop:n9v69-marked-row}.  IRPSC items 2, 7, and 8
supply the conditional, scale-summation, and physical-sector hypotheses
not implied by a single fixed-background step.
\end{proof}

\begin{assessment}[Progress at ninth-series V69]
\label{ass:n9v69-status}
One exact interacting shell step now has separate OS and coercivity
gates.  Reflection positivity follows from joint cone membership and
exact marginalization.  The metric Hessian follows from the complete
mean/covariance mark and consumes an explicit CPRAC debit.  A
metric-dependent conditional Gaussian introduces a new joint-kernel
obligation that fixed-background shell positivity cannot discharge.
\end{assessment}

\begin{programme}[Ninth-series V70: compulsory companion audit]
\label{prog:n9v69-V70}
V70 will inspect the newest Yang--Mills and Navier--Stokes notes and
compare their current transfer, determinant, or boundary positivity
mechanisms with IRPSC.  V71 will continue with the joint
metric-dependent shell kernel selected by that audit.
\end{programme}

\begin{firewall}[Claim boundary after ninth-series V69]
\label{fire:n9v69-boundary}
V69 proves the abstract exact interacting-shell preservation theorem.
It does not establish reflection-cone membership for the full chiral
Standard Model interaction, joint positivity of a metric-dependent
shell kernel, summability of scale debits, anomaly control, infrared
closure, OS reconstruction, or the full SM--Einstein continuum.
\end{firewall}

\begin{theorem}[Ninth-series V69 result]
\label{thm:n9v69-result}
An exact SSFC shell integration preserves OS positivity when its joint
interaction density belongs to the reflection-cone closure; the exact
retained marginal inherits positivity.  Its metric Hessian is the
complete direct row minus the response covariance in
\eqref{eq:n9v69-Uj-second}, with relative debit
$\delta_j+c_j$.  Pointwise reflection positivity of a shell conditioned
on a dynamical metric is insufficient; a jointly positive kernel is
required.
\end{theorem}

\begin{proof}
Combine \Cref{thm:n9v69-cone-weighting,
thm:n9v69-one-shell-RP,thm:n9v69-shell-Hessian,
prop:n9v69-shell-margin,prop:n9v69-conditional-insufficient,
thm:n9v69-IRPSC}.
\end{proof}

\paragraph{Parent-version record.}
V69 was formed from
\texttt{Renewal\_SM\_Einstein\_Continuum\_Research\_Notes\_V68.tex}.
The V68 parent had $25\,631$ logical source lines, $923\,364$ bytes, and
SHA--256
\[
 \texttt{1FB259B63F4797EF4C6D432B42174C2834C176DF912138573192E5D462E56983}.
\]

% =============================================================================
% END APPENDED NINTH-SERIES V69 MATERIAL
% =============================================================================
% =============================================================================
% BEGIN APPENDED NINTH-SERIES V70 MATERIAL
% =============================================================================

\section{Ninth-series V70: compulsory audit and calibrated coarse directions}
\label{sec:v9v70}

\subsection{Companion files inspected}

This is the compulsory audit five versions after V65.  The newest
Yang--Mills note was V99.  The Navier--Stokes note remained V26 with the
same digest as at V60 and V65.

\begin{center}
\begin{tabular}{|p{0.43\textwidth}|r|r|}
\hline
File & bytes & lines\\
\hline
\texttt{Renewal\_Yang\_Mills\_Mass\_Gap\_Research\_Notes\_V99.tex}
& $1\,625\,638$ & $45\,516$\\
\texttt{Renewal\_NS\_Stability\_Research\_Notes\_V26.tex}
& $278\,283$ & $5\,319$\\
\hline
\end{tabular}
\end{center}
Their SHA--256 digests were
\begin{align}
 &\texttt{7CA9D687D343116056B357B4B3C8245E4B0FC6CB0F01BFBCA0FFA4C21B0A6421},
 \label{eq:n9v70-YM-hash}\\
 &\texttt{82C887F8C2C69E4F28FAD7C3DC8DE5B9099CB5A4BF431EBA1FFF3E91935978AD}.
 \label{eq:n9v70-NS-hash}
\end{align}

\subsection{Audit finding: the full coarse map cannot contract}

Yang--Mills V99 identifies an exact eigenvalue-one sector in a coarse
map: an interaction already measurable in the retained variable passes
through conditional expectation unchanged.  The proper target is a
strict contraction only on a calibrated stable complement, with the
running marginal/relevant coordinates treated separately.  This
applies before any model-specific estimate.

Let $B:\Omega\to\Sigma$ be an exact block map, let
$\kappa(d\omega\mid s)$ be its reference disintegration, and define the
coarse action map
\begin{equation}
 ({\cal T}S)(s)
 =-\hbar\log\int
 e^{-S(\omega)/\hbar}\kappa(d\omega\mid s).
 \label{eq:n9v70-coarse-map}
\end{equation}

\begin{theorem}[Exact retained pullback identity]
\label{thm:n9v70-pullback}
For every coarse functional $R$ for which the integrals exist,
\begin{equation}
 {\cal T}(S+R\circ B)={\cal T}S+R.
 \label{eq:n9v70-pullback}
\end{equation}
Consequently
\begin{equation}
 D{\cal T}_S[R\circ B]=R.
 \label{eq:n9v70-unit-differential}
\end{equation}
Any norm that assigns the same nonzero size to $R\circ B$ and $R$
therefore gives $\lVert D{\cal T}_S\rVert\ge1$.
\end{theorem}

\begin{proof}
On the fibre $B\omega=s$, the pullback $R(B\omega)$ equals the constant
$R(s)$.  Factor $e^{-R(s)/\hbar}$ out of the integral in
\eqref{eq:n9v70-coarse-map} and take the logarithm.  Differentiate the
resulting affine identity.
\end{proof}

\begin{firewall}[No complete-map strict contraction]
\label{fire:n9v70-full-contraction}
The obstruction is exact and independent of weak coupling, clustering,
or block size.  A proof that estimates the complete coarse differential
by a constant below one has omitted a retained pullback direction or has
changed its norm between scales without recording the calibration.
\end{firewall}

\subsection{SM--Einstein marginal and relevant owners}

At scale $j$, let ${\mathfrak B}_j$ be a Banach space of complete local
and quasi-local interactions modulo constants.  Choose a bounded
projection
\begin{equation}
 M_j:{\mathfrak B}_j\to{\mathfrak M}_j,
 \qquad Q_j=I-M_j.
 \label{eq:n9v70-MQ}
\end{equation}
The declared space ${\mathfrak M}_j$ must contain every retained
pullback direction relevant to the trajectory.  At minimum its ledger
includes the cosmological/vacuum and Newton rows from V58, gauge kinetic
couplings, Higgs mass and quartic rows, Yukawa matrices, admissible field
normalizations, and any symmetry-allowed relevant or marginal operator
generated by the regulator.  This is a ledger, not a claim that the list
is already closed under the physical RG.

\begin{definition}[Calibrated SM--Einstein slice]
\label{def:n9v70-calibrated-slice}
For prescribed running coordinates $m_j\in{\mathfrak M}_j$, define
\begin{equation}
 \Sigma_j(m_j)=\{S\in{\mathfrak B}_j:M_jS=m_j\}.
 \label{eq:n9v70-calibrated-slice}
\end{equation}
The exact coarse map respects the calibrated trajectory if
\begin{equation}
 {\cal T}_j\Sigma_j(m_j)
 \subseteq\Sigma_{j+1}(m_{j+1}).
 \label{eq:n9v70-calibration-invariance}
\end{equation}
\end{definition}

\begin{firewall}[Calibration is not tuning by notation]
\label{fire:n9v70-calibration-tuning}
The sequence $m_j$ must be constructed from renormalization conditions
and the exact coarse map.  Declaring a slice does not prove that the
physical bare parameters enter it, that the generated marginal list is
complete, or that the trajectory has a continuum limit.
\end{firewall}

\begin{theorem}[Calibrated stable-sector contraction]
\label{thm:n9v70-stable-contraction}
Let ${\cal D}_j\subset\Sigma_j(m_j)$ be convex.  Suppose
\eqref{eq:n9v70-calibration-invariance} holds and
\begin{equation}
 \sup_{S\in{\cal D}_j}
 \lVert Q_{j+1}D{\cal T}_{j,S}Q_j\rVert
 \le\kappa_j<1.
 \label{eq:n9v70-stable-derivative}
\end{equation}
Then for $S,S'\in{\cal D}_j$,
\begin{equation}
 \lVert Q_{j+1}({\cal T}_jS-{cal T}_jS')\rVert_{j+1}
 \le\kappa_j\lVert Q_j(S-S')\rVert_j.
 \label{eq:n9v70-stable-contraction}
\end{equation}
\end{theorem}

\begin{proof}
Since $S,S'$ lie in the same calibrated slice,
$S-S'=Q_j(S-S')$.  The segment
$S_t=S'+t(S-S')$ stays in ${\cal D}_j$.  Integrate the Fr\'echet
derivative:
\begin{equation}
 Q_{j+1}({\cal T}_jS-{cal T}_jS')
 =\int_0^1Q_{j+1}D{\cal T}_{j,S_t}Q_j(S-S')\,dt,
\end{equation}
and apply \eqref{eq:n9v70-stable-derivative}.
\end{proof}

\begin{corollary}[Cofinal stable Duhamel bound]
\label{cor:n9v70-discrete-Duhamel}
If stable differences satisfy
\begin{equation}
 d_{j+1}\le\kappa d_j+\epsilon_j,
 \qquad 0\le\kappa<1,
 \label{eq:n9v70-defect-recursion}
\end{equation}
then for $n>m$,
\begin{equation}
 d_n\le\kappa^{n-m}d_m
 +\sum_{r=m}^{n-1}\kappa^{n-1-r}\epsilon_r.
 \label{eq:n9v70-discrete-Duhamel}
\end{equation}
Thus a uniformly bounded defect inserted a cofinally increasing number
of scales before $n$ has vanishing terminal effect, provided the nearby
defect tail also tends to zero.
\end{corollary}

\begin{proof}
Iterate \eqref{eq:n9v70-defect-recursion}.  The final statement follows
by splitting the sum into distant and fixed terminal tails.
\end{proof}

\subsection{Principal Gaussian differential and its unit sector}

For a positive quadratic block Hessian
\begin{equation}
 H=\begin{pmatrix}A&B\\B^*&C\end{pmatrix},
 \qquad C>0,
 \label{eq:n9v70-quadratic-block}
\end{equation}
elimination gives
\begin{equation}
 {\mathscr S}(H)=A-BC^{-1}B^*.
 \label{eq:n9v70-Schur-map}
\end{equation}

\begin{proposition}[Differential of the principal Schur row]
\label{prop:n9v70-Schur-differential}
For $\dot H=(\dot A,\dot B,\dot C)$,
\begin{equation}
 D{\mathscr S}_H[\dot H]
 =\dot A-dot BC^{-1}B^*-BC^{-1}\dot B^*
  +BC^{-1}\dot C C^{-1}B^*.
 \label{eq:n9v70-DS}
\end{equation}
If $\dot B=\dot C=0$, then
\begin{equation}
 D{\mathscr S}_H[(\dot A,0,0)]=\dot A.
 \label{eq:n9v70-Gaussian-unit}
\end{equation}
The retained quadratic row is therefore another exact unit sector and
belongs in $M_j$.
\end{proposition}

\begin{proof}
Differentiate \eqref{eq:n9v70-Schur-map} and use
$D(C^{-1})[\dot C]=-C^{-1}\dot C C^{-1}$.  The special case is
immediate.
\end{proof}

\begin{corollary}[Stable Gaussian candidate]
\label{cor:n9v70-Gaussian-stable}
With
\begin{equation}
 \theta=\max\{\lVert BC^{-1}\rVert,
                    \lVert C^{-1}B^*\rVert\},
 \end{equation}
the nonunit part obeys
\begin{equation}
 \lVert D{\mathscr S}_H[0,\dot B,\dot C]\rVert
 \le2\theta\lVert\dot B\rVert+
      \theta^2\lVert\dot C\rVert.
 \label{eq:n9v70-Gaussian-stable-bound}
\end{equation}
It is a strict stable candidate only after the actual block geometry
and energy norm make the right-hand operator below one.
\end{corollary}

\begin{proof}
Apply submultiplicativity to \eqref{eq:n9v70-DS}.
\end{proof}

\subsection{Interaction with reflection-positive shell integration}

The V69 shell covariance and cone density can depend strongly on the
running calibrated coordinates.  Their stable variation is the object
that should be small.

\begin{proposition}[Marginal motion versus stable motion]
\label{prop:n9v70-marginal-stable}
Let the exact one-shell map be written on a calibrated chart as
\begin{equation}
 S=m+q,
 \qquad m=M_jS,quad q=Q_jS.
 \label{eq:n9v70-mq-coordinates}
\end{equation}
Then the differential decomposes exactly as
\begin{equation}
 D{\cal T}_{j,S}[\dot m+dot q]
 =D_m{\cal T}_{j,S}[\dot m]+D_q{\cal T}_{j,S}[\dot q].
 \label{eq:n9v70-map-split}
\end{equation}
The first term determines the running reference family and need not be
contractive.  Only
$Q_{j+1}D_q{\cal T}_{j,S}Q_j$ enters
\eqref{eq:n9v70-stable-derivative}.
\end{proposition}

\begin{proof}
This is linearity of the Fr\'echet differential with the complementary
projections \eqref{eq:n9v70-MQ}.  The contraction theorem itself
removes the $M_j$ difference by restricting to a calibrated slice.
\end{proof}

\begin{decision}[Reflection kernel acquisition after the audit]
\label{dec:n9v70-reflection-acquisition}
V71 will construct a jointly reflection-positive transfer kernel for a
family of shell references indexed by the calibrated coordinates $m$.
Metric and coupling motion along $m$ will be built into that exact
kernel.  Stable interaction differences will then be tested by the V69
cone and CPRAC bounds and the V70 projected contraction.
\end{decision}

\begin{firewall}[Running couplings must stay inside the positivity domain]
\label{fire:n9v70-positivity-domain}
Calibration removes marginal directions from a contraction norm; it
does not exempt them from reflection positivity, stability, or
normalizability.  The full running path $m_j$ must remain inside the
domain where the transfer kernel, Higgs reference, gauge character
weights, and fermion owner are well defined.
\end{firewall}

\subsection{Audit integration card}

\begin{definition}[Calibrated SM--Einstein coarse card (CSECC)]
\label{def:n9v70-CSECC}
CSECC requires:
\begin{enumerate}
\item the exact coarse map and pullback identity
\eqref{eq:n9v70-pullback};
\item a complete declared relevant/marginal projection $M_j$;
\item a prescribed running path $m_j$ satisfying
\eqref{eq:n9v70-calibration-invariance};
\item a strict principal Gaussian bound on the stable complement;
\item a support-resolved nonlinear stable correction inside the
remaining margin;
\item exact determinant, covariance, chart, bad-block, phase, and
moving-filter ownership;
\item compatibility of the running path with the V69 reflection cone;
and
\item the cofinal defect estimate
\eqref{eq:n9v70-discrete-Duhamel}.
\end{enumerate}
\end{definition}

\begin{theorem}[CSECC gives the correct noncircular contraction target]
\label{thm:n9v70-CSECC}
Under CSECC, the exact coarse map follows its prescribed marginal and
relevant trajectory while contracting only the aligned stable
difference.  Stable regulator defects propagate by
\eqref{eq:n9v70-discrete-Duhamel}; the exact pullback sector is never
incorrectly assigned a norm below one.
\end{theorem}

\begin{proof}
Use \Cref{thm:n9v70-pullback,thm:n9v70-stable-contraction,
cor:n9v70-discrete-Duhamel,prop:n9v70-Schur-differential,
prop:n9v70-marginal-stable}.  The remaining CSECC items ensure the
projected derivative is the complete physical derivative.
\end{proof}

\begin{assessment}[Direction after the V70 audit]
\label{ass:n9v70-direction}
The full-map contraction target has been removed.  The SM--Einstein
trajectory must carry exact running cosmological, Newton, gauge, Higgs,
Yukawa, and generated marginal owners, while only the calibrated stable
complement contracts.  This also clarifies V69: a metric-dependent shell
reference belongs to the running family rather than to a perturbative
stable error.
\end{assessment}

\begin{programme}[Ninth-series V71: calibrated joint transfer kernel]
\label{prog:n9v70-V71}
V71 will give a finite-time-lattice transfer-kernel criterion for joint
reflection positivity of the retained metric and one matter/gravity
shell.  It will distinguish positive diagonal conditional kernels from
a positive two-boundary kernel and integrate the calibrated parameter
dependence exactly.
\end{programme}

\begin{firewall}[Claim boundary after ninth-series V70]
\label{fire:n9v70-boundary}
The audit imports no unproved companion hypothesis.  V70 proves the
pullback obstruction and calibrated contraction theorem, but it does
not construct the physical running trajectory, close the stable
constant, build the joint transfer kernel, control the chiral phase or
anomaly, close the infrared theory, reconstruct the OS continuum, or
prove the full SM--Einstein continuum.
\end{firewall}

\begin{theorem}[Ninth-series V70 result]
\label{thm:n9v70-result}
The exact coarse map has the retained pullback identity
\eqref{eq:n9v70-pullback} and therefore cannot contract strictly on its
complete interaction space.  After all running relevant and marginal
owners are placed in $M_j$, the projected derivative criterion
\eqref{eq:n9v70-stable-derivative} gives the correct stable contraction
and cofinal defect bound.  The calibrated running family must also
remain inside the joint reflection-positive domain.
\end{theorem}

\begin{proof}
Combine \Cref{thm:n9v70-pullback,
thm:n9v70-stable-contraction,cor:n9v70-discrete-Duhamel,
prop:n9v70-Schur-differential,thm:n9v70-CSECC}.
\end{proof}

\paragraph{Parent-version record.}
V70 was formed from
\texttt{Renewal\_SM\_Einstein\_Continuum\_Research\_Notes\_V69.tex}.
The V69 parent had $26\,102$ logical source lines, $940\,539$ bytes, and
SHA--256
\[
 \texttt{95349B1F2214BEDF6A603CDDE2FE7371CF878DCDE9DEBCC929A021C0C80C5B12}.
\]

% =============================================================================
% END APPENDED NINTH-SERIES V70 MATERIAL
% =============================================================================
% =============================================================================
% BEGIN APPENDED NINTH-SERIES V71 MATERIAL
% =============================================================================

\section{Ninth-series V71: calibrated joint transfer kernels}
\label{sec:v9v71}

\subsection{Boundary-state positivity}

V69 proved that pointwise reflection positivity of a shell conditioned
on a metric does not imply joint positivity.  V70 placed the running
relevant and marginal coordinates in a calibrated family.  The correct
finite-time-lattice object is therefore a transfer kernel on the joint
boundary state.

Let
\begin{equation}
 X=X_{\rm met}\times X_{\rm shell}\times X_{\rm phys}
 \label{eq:n9v71-joint-state}
\end{equation}
contain the retained spatial metric, the shell boundary field, and all
physical gauge/Higgs labels needed at one time slice.  Let $\nu_m$ be a
positive reference measure at calibrated coordinate $m$ and put
${\cal H}_m=L^2(X,\nu_m)$.

\begin{definition}[Positive joint transfer kernel]
\label{def:n9v71-positive-kernel}
A measurable symmetric kernel $K_m(x,y)$ is a positive joint transfer
kernel if it defines a positive self-adjoint operator $T_m$ on
${\cal H}_m$:
\begin{equation}
 \int\overline{f(x)}K_m(x,y)f(y)
 \,d\nu_m(x)d\nu_m(y)\ge0
 \qquad(f\in{\cal H}_m),
 \label{eq:n9v71-kernel-positive}
\end{equation}
and if $K_m(x,y)\ge0$ almost everywhere when an ordinary bosonic path
probability is required.
\end{definition}

\begin{proposition}[Gram criterion]
\label{prop:n9v71-Gram}
Suppose there is a measure space $(Y,\sigma)$ and a measurable feature
map $\Phi_m:Y\times X\to\mathbb C$ such that
\begin{equation}
 K_m(x,y)=\int_Y
 \overline{\Phi_m(\xi,x)}\Phi_m(\xi,y)d\sigma(\xi),
 \label{eq:n9v71-Gram-factor}
\end{equation}
with the required integrability.  Then $K_m$ is a positive operator
kernel.  Pointwise nonnegativity remains a separate condition.
\end{proposition}

\begin{proof}
Fubini gives
\begin{equation}
 \langle f,T_mf\rangle
 =\int_Y\left|
   \int_X\Phi_m(\xi,y)f(y)d\nu_m(y)
  \right|^2d\sigma(\xi)\ge0.
\end{equation}
The Gram representation can have complex oscillating entries, so it
does not imply pointwise nonnegativity.
\end{proof}

\begin{corollary}[Heat-semigroup criterion]
\label{cor:n9v71-heat-kernel}
If $H_m$ is self-adjoint and bounded below on ${\cal H}_m$, then for
$a_t>0$
\begin{equation}
 T_m=e^{-a_t(H_m-E_m)}\ge0
 \label{eq:n9v71-heat-transfer}
\end{equation}
for any $E_m\le\inf\operatorname{spec}H_m$.  If its semigroup is
positivity preserving and has a kernel, that kernel is nonnegative and
is a positive joint transfer kernel.
\end{corollary}

\begin{proof}
The spectral multiplier $e^{-a_t(\lambda-E_m)}$ is nonnegative.
Positivity preservation sends nonnegative functions to nonnegative
functions and yields a nonnegative kernel under the stated kernel
hypothesis.
\end{proof}

\subsection{Transfer construction of OS positivity}

Let $T_m\ge0$ be a bounded joint transfer operator and let
$\Omega_m\in{\cal H}_m$ be a boundary vector.  Positive-time cylinder
observables act by multiplication operators $M_f$ interspersed with
powers of $T_m$.

\begin{theorem}[Joint transfer operator implies reflection positivity]
\label{thm:n9v71-transfer-RP}
Consider a reflection-symmetric finite time-lattice path functional for
which the only coupling across the reflection bond is $K_m$.  Contract
the positive half-path, its boundary weight, and a positive-time
cylinder observable $F$ into a boundary vector
\begin{equation}
 v_F(x)=\int F(\omega_+)A_m(\omega_+,x)d\omega_+\in{\cal H}_m,
 \label{eq:n9v71-positive-cylinder}
\end{equation}
where the negative half amplitude is the reflected complex conjugate.
Then
\begin{equation}
 \langle\Theta F,F\rangle_{\rm E}
 =Z_m^{-1}\langle v_F,T_mv_F\rangle
 =Z_m^{-1}\lVert T_m^{1/2}v_F\rVert^2\ge0.
 \label{eq:n9v71-transfer-square}
\end{equation}
Hence a positive joint transfer kernel gives OS positivity of the
finite-time joint metric--shell functional.
\end{theorem}

\begin{proof}
Reflection symmetry identifies the negative-half amplitude with
$\overline{v_F(x)}$.  Integrating the two halves leaves
$\int\overline{v_F(x)}K_m(x,y)v_F(y)d\nu_m(x)d\nu_m(y)$, which is
$\langle v_F,T_mv_F\rangle$.  Positivity of $T_m$ gives its square root
and the last equality.
\end{proof}
\begin{remark}[Why the metric is part of the state]
\label{rem:n9v71-metric-state}
The off-diagonal kernel $K_m((g,z),(g',z'))$ couples two boundary
metrics.  Its positivity controls the matrix of cross terms missing
from pointwise conditional kernels $K_g(z,z')$.  The two-label
counterexample in V69 is excluded precisely by
\eqref{eq:n9v71-kernel-positive}.
\end{remark}

\subsection{Half-potentials and cross-plane interactions}

Let $V_{+,m}$ be a real multiplication operator supported on one time
half and define
\begin{equation}
 T_m^{V}=e^{-V_{+,m}/(2\hbar)}T_m
         e^{-V_{+,m}/(2\hbar)}.
 \label{eq:n9v71-half-potential-transfer}
\end{equation}

\begin{proposition}[Half-potential stability]
\label{prop:n9v71-half-potential}
If $T_m\ge0$ and the multipliers in
\eqref{eq:n9v71-half-potential-transfer} are defined on the form domain,
then $T_m^V\ge0$.  If both factors are positivity preserving, the
resulting kernel is nonnegative.
\end{proposition}

\begin{proof}
For every $f$,
\begin{equation}
 \langle f,T_m^Vf\rangle
 =\left\langle e^{-V_{+,m}/(2\hbar)}f,
 T_m e^{-V_{+,m}/(2\hbar)}f\right\rangle\ge0.
\end{equation}
Positivity preservation is stable under composition with a nonnegative
multiplication operator.
\end{proof}

\begin{firewall}[A cross-plane interaction is not a half-potential]
\label{fire:n9v71-cross-plane}
Terms coupling the two boundary slices cannot be assigned to
$V_{+,m}$ and its reflection.  Their kernel must itself be positive as
an operator, for example through a nonnegative character expansion or a
Gram factorization.  Splitting such a term in half without a kernel
proof is circular.
\end{firewall}

\begin{proposition}[Orthogonal physical projection]
\label{prop:n9v71-physical-projection}
Let $P_m$ be an orthogonal Gauss/constraint projection commuting with
$T_m$.  Then
\begin{equation}
 T_m^{\rm phys}=P_mT_mP_m\ge0
 \label{eq:n9v71-projected-transfer}
\end{equation}
on $P_m{\cal H}_m$.  If $P_m$ is realized by averaging a
positivity-preserving group action, it also preserves nonnegative
kernels.
\end{proposition}

\begin{proof}
$\langle f,P_mT_mP_mf\rangle=\langle P_mf,T_mP_mf\rangle\ge0$.
Haar averaging of positivity-preserving transformations preserves
pointwise positivity.
\end{proof}

\begin{firewall}[Faddeev--Popov replacement]
\label{fire:n9v71-FP}
An indefinite ghost determinant is not an orthogonal projection of the
form \eqref{eq:n9v71-projected-transfer}.  A gauge-fixed representation
requires its own graded reflection proof and cannot inherit positivity
from the physical Haar projector by notation.
\end{firewall}

\subsection{Exact shell trace at one transfer step}

Suppose for the moment
\begin{equation}
 {\cal H}_m={\cal H}_{{\rm ret},m}\otimes{\cal H}_{{\rm sh},m}
 \label{eq:n9v71-H-tensor}
\end{equation}
and $T_m\ge0$ is trace class in the shell factor.

\begin{theorem}[Positive partial trace]
\label{thm:n9v71-partial-trace}
The one-step retained operator
\begin{equation}
 T_{{\rm ret},m}=\operatorname{Tr}_{\rm sh}T_m
 \label{eq:n9v71-retained-transfer}
\end{equation}
is positive.  If $T_m$ has a nonnegative joint kernel, then
\begin{equation}
 K_{{\rm ret},m}(g,g')
 =\int K_m((g,z),(g',z))d\nu_{{\rm sh},m}(z)\ge0.
 \label{eq:n9v71-retained-kernel}
\end{equation}
\end{theorem}

\begin{proof}
For an orthonormal shell basis $(e_\alpha)$ and retained $u$,
\begin{equation}
 \langle u,T_{{\rm ret},m}u\rangle
 =\sum_\alpha
 \langle u\otimes e_\alpha,T_m(u\otimes e_\alpha)\rangle\ge0.
\end{equation}
The kernel formula is the diagonal shell trace and is nonnegative term
by term.
\end{proof}

\begin{firewall}[Full-history marginal need not remain one-step Markov]
\label{fire:n9v71-non-Markov}
Integrating a shell field over an entire time history preserves OS
positivity by V69, but the resulting retained action can be nonlocal in
time.  It need not equal a product of powers of the one-step partial
trace \eqref{eq:n9v71-retained-transfer}.  A local retained transfer
operator requires a Markov dilation, an enlarged boundary state, or a
separate memory-decay construction.
\end{firewall}

\subsection{Calibrated families}

Let $m\mapsto H_m$, $P_m$, and $V_{+,m}$ be the running family selected
by V70.  Define
\begin{equation}
 T_m^{\rm phys,V}
 =P_m e^{-V_{+,m}/(2\hbar)}
 e^{-a_t(H_m-E_m)}
 e^{-V_{+,m}/(2\hbar)}P_m.
 \label{eq:n9v71-calibrated-transfer}
\end{equation}

\begin{theorem}[Calibrated positivity domain]
\label{thm:n9v71-calibrated-domain}
Assume that for every $m$ in a connected domain ${\cal P}$:
\begin{enumerate}
\item $H_m$ is self-adjoint, uniformly lower bounded, and has a
positivity-preserving semigroup;
\item $P_m$ is an orthogonal positivity-preserving physical projection;
\item the half-potential is real and form-admissible; and
\item every remaining cross-plane factor is incorporated by a proved positive sandwich or positive-kernel Schur product.
\end{enumerate}
Then \eqref{eq:n9v71-calibrated-transfer} is a positive joint transfer
operator for every $m\in{\cal P}$.  A calibrated trajectory
$m_j\in{\cal P}$ may move without being contractive while preserving
the transfer positivity gate at every scale.
\end{theorem}

\begin{proof}
Apply \Cref{cor:n9v71-heat-kernel,
prop:n9v71-half-potential,prop:n9v71-physical-projection} and use the assumed positive sandwich or positive-kernel Schur product for the cross-plane factors.  Positivity is asserted pointwise
along the parameter path, not inferred from a contraction estimate.
\end{proof}

\begin{firewall}[Operator positivity is not open under arbitrary perturbation]
\label{fire:n9v71-positive-domain}
If $T_m$ has spectrum touching zero, a small operator-norm perturbation
need not remain positive.  The domain ${\cal P}$ must be constructed
from a positive factorization such as
\eqref{eq:n9v71-calibrated-transfer}, or from a uniform positive floor;
continuity in $m$ alone is insufficient.
\end{firewall}

\subsection{Duhamel marks along the running family}

Assume a fixed Hilbert-bundle trivialization and write
\begin{equation}
 \dot H=D_mH_m[\dot m].
 \label{eq:n9v71-H-dot}
\end{equation}

\begin{proposition}[Exact calibrated transfer derivative]
\label{prop:n9v71-transfer-derivative}
Ignoring the displayed half-potential and projection factors for one
line, the semigroup derivative is
\begin{equation}
 D_me^{-a_tH_m}[\dot m]
 =-\int_0^{a_t}
 e^{-(a_t-s)H_m}\dot H e^{-sH_m}\,ds.
 \label{eq:n9v71-transfer-Duhamel}
\end{equation}
The derivative of \eqref{eq:n9v71-calibrated-transfer} additionally
contains the two half-potential derivatives, two $D_mP_m$ rows, and the
ground-energy derivative.  These are running-family owners, not stable
contraction errors.
\end{proposition}

\begin{proof}
Differentiate the semigroup by Duhamel's formula, then apply the product
rule to \eqref{eq:n9v71-calibrated-transfer}.
\end{proof}

\begin{proposition}[Stable cone perturbations]
\label{prop:n9v71-stable-cone}
Fix $m\in{\cal P}$.  If a stable interaction perturbation has density
$W_q$ in the joint reflection-cone closure and its V69 complete marked
rows obey the projected derivative bound of V70, then exact weighting
preserves OS positivity and its stable difference is controlled by the
calibrated contraction constant.  Motion of $m$ is excluded from that
constant.
\end{proposition}

\begin{proof}
Use V69 reflection-cone weighting for $W_q$ and V70 stable contraction
on the common slice $m$.  The pullback identity prevents including
$D_mT_m$ in a strict complete-map norm.
\end{proof}

\subsection{Fermion and anomaly boundaries}

\begin{firewall}[Fermionic transfer positivity]
\label{fire:n9v71-fermion}
The bosonic nonnegative-kernel criterion does not directly cover
Grassmann variables.  A chiral Standard Model transfer construction
must prove the appropriate graded reflection involution and positivity
of the fermionic boundary form before its determinant is integrated
out.  A nonnegative conjugate-paired determinant is still not by itself
a joint transfer-kernel factorization.
\end{firewall}

\begin{firewall}[Anomalous physical projection]
\label{fire:n9v71-anomaly}
If the fermion measure changes under the gauge transformation used to
define $P_m$, the projected transfer operator acquires a cocycle and
\Cref{prop:n9v71-physical-projection} no longer represents the physical
theory.  Regulator-level anomaly cancellation remains prior to joint
transfer positivity.
\end{firewall}

\subsection{Calibrated joint transfer card}

\begin{definition}[Calibrated joint transfer card (CJTC)]
\label{def:n9v71-CJTC}
CJTC requires:
\begin{enumerate}
\item a joint boundary state containing the retained metric and all
shell boundary labels;
\item a positive self-adjoint, positivity-preserving transfer operator
or Gram kernel;
\item a physical orthogonal projection with anomaly-free covariance;
\item positive cross-plane kernels and admissible half-potentials;
\item a calibrated parameter domain ${\cal P}$ containing the entire
running trajectory;
\item exact Duhamel, moving-projection, and ground-energy derivatives;
\item a positive partial trace or an enlarged state for shell memory;
\item stable reflection-cone perturbations inside the V69 CPRAC and V70
contraction margins; and
\item a graded replacement for the physical chiral fermion sector.
\end{enumerate}
\end{definition}

\begin{theorem}[CJTC closes the abstract joint-kernel row]
\label{thm:n9v71-CJTC}
Under CJTC, the metric-dependent shell reference is jointly
OS-positive, rather than merely conditionally positive at fixed metric.
Exact one-step shell tracing preserves transfer positivity, and exact
full-history marginalization preserves OS positivity.  The calibrated
parameter motion is incorporated in the reference family; only aligned
stable perturbations consume the projected contraction margin.
\end{theorem}

\begin{proof}
Use \Cref{thm:n9v71-transfer-RP,
prop:n9v71-half-potential,prop:n9v71-physical-projection,
thm:n9v71-partial-trace,thm:n9v71-calibrated-domain,
prop:n9v71-stable-cone}.  CJTC item 9 supplies the fermionic statement
not contained in the bosonic proof.
\end{proof}

\begin{assessment}[Progress at ninth-series V71]
\label{ass:n9v71-status}
The conditional-metric obstruction from V69 now has a precise repair:
quantize the metric and shell on the same boundary state and prove
positivity of their joint transfer operator.  Calibrated coupling motion
is exact reference motion, while stable cone perturbations remain the
contractive part.  The main upstream task is to construct a closed
joint Hamiltonian form whose semigroup is positivity preserving.
\end{assessment}

\begin{programme}[Ninth-series V72: joint Dirichlet-form construction]
\label{prog:n9v71-V72}
V72 will build a joint metric--bosonic-shell quadratic form, prove
closedness by KLMN under the existing relative margins, and state a
Markov/Beurling--Deny criterion for positivity preservation of its
semigroup.  It will isolate which gauge, conformal, and fermionic rows
prevent a direct scalar Dirichlet-form argument.
\end{programme}

\begin{firewall}[Claim boundary after ninth-series V71]
\label{fire:n9v71-boundary}
V71 proves the joint transfer, partial-trace, and calibrated-family
criteria.  It does not construct the physical joint Hamiltonian,
establish its Markov property, prove chiral fermion reflection
positivity or anomaly cancellation, control all stable rows, close the
infrared sector, reconstruct the OS continuum, or prove the full
SM--Einstein continuum.
\end{firewall}

\begin{theorem}[Ninth-series V71 result]
\label{thm:n9v71-result}
A metric-dependent shell step is jointly reflection positive if its
boundary evolution is generated by a positive self-adjoint transfer
operator on the combined metric--shell state, with positive physical
projection and cross-plane kernels.  Its one-step shell partial trace is
positive, while a full-history marginal may require memory.  Calibrated
running directions remain in the exact transfer family and stable cone
perturbations alone enter the V70 contraction.
\end{theorem}

\begin{proof}
Combine \Cref{prop:n9v71-Gram,cor:n9v71-heat-kernel,
thm:n9v71-transfer-RP,thm:n9v71-partial-trace,
thm:n9v71-calibrated-domain,thm:n9v71-CJTC}.
\end{proof}

\paragraph{Parent-version record.}
V71 was formed from
\texttt{Renewal\_SM\_Einstein\_Continuum\_Research\_Notes\_V70.tex}.
The V70 parent had $26\,487$ logical source lines, $954\,369$ bytes, and
SHA--256
\[
 \texttt{AF5EBAD857B28F42F20D84258BE0FB66F63330FB181541F15B336497451EC3E4}.
\]

% =============================================================================
% END APPENDED NINTH-SERIES V71 MATERIAL
% =============================================================================
% =============================================================================
% BEGIN APPENDED NINTH-SERIES V72 MATERIAL
% =============================================================================

\section{Ninth-series V72: joint Dirichlet forms and the conformal obstruction}
\label{sec:v9v72}

\subsection{Closed forms as a source of positive transfer kernels}

V71 reduced joint reflection positivity to a positive self-adjoint,
positivity-preserving transfer semigroup on the combined boundary
state.  A standard construction begins with a symmetric Dirichlet form.
Let $(X,\mu)$ be the finite-regulator joint metric--bosonic-shell state
space and let ${\cal H}=L^2(X,\mu)$.

\begin{definition}[Dirichlet reference form]
\label{def:n9v72-Dirichlet}
A densely defined closed nonnegative symmetric form
$({\cal E}_0,{\cal D}({\cal E}_0))$ is Dirichlet if for every real
$F\in{\cal D}({\cal E}_0)$ and every normal contraction
$\eta:\mathbb R\to\mathbb R$,
\begin{equation}
 \eta(0)=0,
 \qquad |\eta(a)-\eta(b)|\le|a-b|,
 \qquad
 {\cal E}_0(\eta\circ F)\le{\cal E}_0(F).
 \label{eq:n9v72-Markov-property}
\end{equation}
\end{definition}

\begin{theorem}[Dirichlet form to positive transfer]
\label{thm:n9v72-Dirichlet-transfer}
Let $H_0\ge0$ be the self-adjoint operator represented by
${\cal E}_0$.  Then $e^{-tH_0}$ is positivity preserving and
$L^\infty$-contractive for $t\ge0$.  If it has an integral kernel, that
kernel is nonnegative and defines a V71 positive joint transfer kernel.
\end{theorem}

\begin{proof}
The Markov property of a closed symmetric form is equivalent to the
sub-Markov property of its associated semigroup.  Positivity
preservation follows by applying the normal contractions
$F\mapsto F_+$ and $F\mapsto F\wedge1$.  A positivity-preserving
integral operator has a nonnegative kernel up to null sets.
\end{proof}

\subsection{KLMN loads lower-order negative rows}

Let ${\mathfrak r}$ be a symmetric form on ${\cal D}({\cal E}_0)$ with
negative part satisfying
\begin{equation}
 {\mathfrak r}_-(F)
 \le a{\cal E}_0(F)+b\lVert F\rVert_2^2,
 \qquad 0\le a<1.
 \label{eq:n9v72-KLMN-bound}
\end{equation}

\begin{theorem}[KLMN joint Hamiltonian]
\label{thm:n9v72-KLMN}
The form
\begin{equation}
 {\cal E}={\cal E}_0+{\mathfrak r}
 \label{eq:n9v72-perturbed-form}
\end{equation}
is closed and bounded below and defines a unique self-adjoint
Hamiltonian $H$.  Moreover,
\begin{equation}
 {\cal E}(F)
 \ge(1-a){\cal E}_0(F)-b\lVert F\rVert_2^2.
 \label{eq:n9v72-KLMN-lower}
\end{equation}
After the ground shift $H+b$, the transfer operator is positive as a
self-adjoint operator.
\end{theorem}

\begin{proof}
The strict relative bound in \eqref{eq:n9v72-KLMN-bound} is the KLMN
hypothesis.  It gives closedness, semiboundedness, uniqueness, and
\eqref{eq:n9v72-KLMN-lower}.  Spectral calculus makes
$e^{-t(H+b)}$ a positive operator.
\end{proof}

\begin{firewall}[KLMN positivity versus positivity preservation]
\label{fire:n9v72-KLMN-positive}
KLMN proves operator positivity of the transfer after a ground shift.
It does not prove that the semigroup maps nonnegative functions to
nonnegative functions.  A general nonlocal, complex, magnetic, or
matrix-valued relatively bounded form can violate the Beurling--Deny
criterion.
\end{firewall}

\begin{proposition}[Real potential perturbations]
\label{prop:n9v72-real-potential}
Let ${\mathfrak r}(F)=\int V|F|^2d\mu$ with real $V$.  If
$V_-$ satisfies \eqref{eq:n9v72-KLMN-bound}, then the semigroup of
${\cal E}_0+V$ is positivity preserving.  If $V\ge0$, the perturbed form
is again Dirichlet and its semigroup is sub-Markovian.
\end{proposition}

\begin{proof}
For bounded truncations $V^{(N)}$, the Trotter product
\begin{equation}
 e^{-t(H_0+V^{(N)})}
 =\operatorname*{s-lim}_{k\to\infty}
 \left(e^{-tH_0/k}e^{-tV^{(N)}/k}\right)^k
 \label{eq:n9v72-Trotter-potential}
\end{equation}
is positivity preserving because both factors are.  KLMN form
convergence passes this property to $V$.  When $V\ge0$, multiplication
by $e^{-tV}$ is also $L^\infty$-contractive, giving the Markov property.
\end{proof}

\subsection{Positive principal block and its Schur margin}

Write joint differential variables schematically as metric coordinates
$q$ and bosonic shell coordinates $z$.  A local principal form has
coefficient matrix
\begin{equation}
 {cal G}=
 \begin{pmatrix}
 G&C\\ C^*&K
 \end{pmatrix},
 \qquad K>0.
 \label{eq:n9v72-principal-block}
\end{equation}

\begin{proposition}[Principal Schur criterion]
\label{prop:n9v72-principal-Schur}
The block ${\cal G}$ is nonnegative if and only if
\begin{equation}
 G-CK^{-1}C^*\ge0.
 \label{eq:n9v72-principal-Schur}
\end{equation}
If $G>0$ and
\begin{equation}
 \lVert G^{-1/2}CK^{-1/2}\rVert\le\gamma<1,
 \label{eq:n9v72-principal-cross}
\end{equation}
then
\begin{equation}
 G-CK^{-1}C^*\ge(1-\gamma^2)G.
 \label{eq:n9v72-principal-margin}
\end{equation}
\end{proposition}

\begin{proof}
Complete the square:
\begin{align}
 \langle(u,v),{\cal G}(u,v)\rangle
 &=\langle v+K^{-1}C^*u,
          K(v+K^{-1}C^*u)\rangle\\
 &\quad+langle u,(G-CK^{-1}C^*)u\rangle.
\end{align}
This proves the equivalence.  Equation
\eqref{eq:n9v72-principal-cross} gives
$CK^{-1}C^*\le\gamma^2G$.
\end{proof}

\begin{firewall}[Lower-order absorption cannot repair a bad principal symbol]
\label{fire:n9v72-principal-symbol}
If ${\cal G}$ has a negative principal direction, no KLMN estimate
against the same nonnegative principal form can make it a Dirichlet
form.  The negative contribution grows at the same derivative order and
must be removed by constraint reduction, contour choice, or a new
positive reference before lower-order CPRAC rows are loaded.
\end{firewall}

\subsection{The Einstein conformal direction}

At a spatial point, consider the DeWitt-type quadratic form on symmetric
metric variations
\begin{equation}
 {cal G}_\lambda(h,h)
 =|h|_g^2-\lambda(\operatorname{tr}_gh)^2.
 \label{eq:n9v72-DeWitt-form}
\end{equation}

\begin{lemma}[Conformal negative direction]
\label{lem:n9v72-conformal-negative}
For a pure trace variation $h=\varphi g$ in spatial dimension $d$,
\begin{equation}
 {cal G}_\lambda(\varphi g,\varphi g)
 =d(1-\lambda d)\varphi^2.
 \label{eq:n9v72-conformal-value}
\end{equation}
It is negative whenever $\lambda>1/d$.  The general-relativistic value
$\lambda=1/(d-1)$ lies in this range for $d>1$.
\end{lemma}

\begin{proof}
$|\varphi g|_g^2=d\varphi^2$ and
$\operatorname{tr}_g(\varphi g)=d\varphi$.  Substitute in
\eqref{eq:n9v72-DeWitt-form}; compare $1/(d-1)$ with $1/d$.
\end{proof}

\begin{theorem}[Scalar Dirichlet-form no-go for the unreduced Einstein metric]
\label{thm:n9v72-conformal-no-go}
If the joint boundary Hamiltonian uses the real unreduced Einstein
DeWitt form with $\lambda=1/(d-1)$ as its metric principal diffusion
coefficient, then its scalar quadratic form is not nonnegative and
cannot satisfy \Cref{def:n9v72-Dirichlet}.  Lower-order Higgs, gauge,
vacuum, determinant, or counterterm rows satisfying KLMN bounds do not
alter this conclusion.
\end{theorem}

\begin{proof}
Choose test functions oscillating only in the local conformal coordinate.
Their principal energy has the negative coefficient in
\eqref{eq:n9v72-conformal-value} and grows quadratically with frequency.
Every strictly lower-order KLMN term grows more slowly or is bounded by
a coefficient below the positive reference; it cannot reverse the sign
of an arbitrarily high-frequency principal direction.  Hence no
nonnegative Dirichlet principal form exists on the unreduced real
coordinate space.
\end{proof}

\begin{assessment}[Meaning of the conformal no-go]
\label{ass:n9v72-conformal-meaning}
The no-go concerns the direct scalar heat-kernel construction of the
real unreduced gravitational configuration metric.  It does not rule
out a physical constraint reduction that removes the conformal
direction, a rigorously controlled contour representation, a transfer
construction in different variables, or a positive ultraviolet
completion.  Each alternative changes the form and needs its own OS and
Einstein-source proof.
\end{assessment}

\subsection{A conditional positive joint bosonic form}

Let ${\cal H}_{\rm red}$ be a declared physical metric space on which a
positive principal form ${\cal E}_{\rm met}$ has been proved.  Let
${\cal E}_{\rm YM}$ be the compact gauge heat-kernel form after Gauss
projection and let ${\cal E}_{\rm H}$ contain the positive Higgs kinetic
and quartic reference.  Put
\begin{equation}
 {\cal E}_0={\cal E}_{\rm met}+{\cal E}_{\rm YM}+{\cal E}_{\rm H}
              +{\cal E}_{\rm cross}.
 \label{eq:n9v72-joint-reference}
\end{equation}

\begin{theorem}[Conditional joint bosonic transfer]
\label{thm:n9v72-joint-bosonic}
Assume:
\begin{enumerate}
\item the three diagonal forms are closed Dirichlet forms on a common
dense tensor core;
\item the principal cross coefficient obeys
\eqref{eq:n9v72-principal-cross} with $\gamma<1$;
\item lower-order real negative rows obey
\eqref{eq:n9v72-KLMN-bound}; and
\item all nonpotential rows satisfy the Beurling--Deny locality and
reality conditions.
\end{enumerate}
Then the closed joint form is semibounded and its transfer semigroup is
positivity preserving.  It gives the bosonic part of a V71 joint
transfer kernel.
\end{theorem}

\begin{proof}
The principal Schur estimate
\eqref{eq:n9v72-principal-margin} makes the joint principal form closed
and nonnegative on the common core.  The strict lower-order bound loads
KLMN.  The real potential part preserves positivity by
\Cref{prop:n9v72-real-potential}; item 4 supplies the Markovian locality
for the remaining diffusion rows.  Apply
\Cref{thm:n9v72-Dirichlet-transfer} after the ground shift.
\end{proof}

\begin{firewall}[A form-relative row need not be Markovian]
\label{fire:n9v72-relative-Markov}
CPRAC controls closedness and inversion.  It does not automatically
verify the normal-contraction inequality
\eqref{eq:n9v72-Markov-property}.  Moving projectors, Berry connections,
imaginary phase rows, and matrix-valued first-order couplings must be
tested separately for positivity preservation.
\end{firewall}

\subsection{Constraint projection and positivity}

\begin{proposition}[Invariant positive subspace]
\label{prop:n9v72-invariant-subspace}
Let $P$ be an orthogonal positivity-preserving projection such that
$P{\cal D}({\cal E})\subset{\cal D}({\cal E})$ and
\begin{equation}
 {\cal E}(PF,(I-P)G)=0.
 \label{eq:n9v72-reducing-form}
\end{equation}
Then $P$ reduces $H$, and the physical semigroup
$Pe^{-tH}P$ is positivity preserving on $P{\cal H}$.
\end{proposition}

\begin{proof}
Equation \eqref{eq:n9v72-reducing-form} makes the two subspaces
orthogonal for both the Hilbert and form inner products, so the
represented operator and its functional calculus reduce.  Composition
of the positivity-preserving $P$ and semigroup preserves the positive
cone.
\end{proof}

\begin{firewall}[Solving the Hamiltonian constraint]
\label{fire:n9v72-Hamiltonian-constraint}
The gravitational Hamiltonian constraint is not automatically an
orthogonal positivity-preserving projection satisfying
\eqref{eq:n9v72-reducing-form}.  Treating it as one would beg the central
problem.  V72's conditional bosonic theorem begins only after a positive
physical metric form or an equivalent joint transfer construction has
been supplied.
\end{firewall}

\subsection{Fermionic type separation}

\begin{firewall}[Fermions are not scalar Dirichlet variables]
\label{fire:n9v72-fermion-type}
Fermionic kinetic terms are first-order and Grassmann-valued.  They do
not enter the scalar Markov form
\eqref{eq:n9v72-joint-reference}.  Their transfer positivity is a graded
Hilbert-space statement, followed by determinant or Pfaffian ownership.
The scalar KLMN theorem cannot prove chiral reflection positivity or
anomaly cancellation.
\end{firewall}

\begin{definition}[Joint Dirichlet transfer card (JDTC)]
\label{def:n9v72-JDTC}
JDTC requires:
\begin{enumerate}
\item a positive physical metric principal form, with the conformal
direction explicitly removed or reconstructed;
\item closed Dirichlet gauge and Higgs references;
\item a strict principal Schur margin for metric--matter cross diffusion;
\item KLMN bounds for all real lower-order negative rows;
\item Beurling--Deny tests for every nonpotential perturbation;
\item a positivity-preserving physical projection reducing the form;
\item a calibrated parameter domain maintaining all margins; and
\item a separate graded chiral-fermion transfer and anomaly proof.
\end{enumerate}
\end{definition}

\begin{theorem}[JDTC supplies the joint semigroup row of CJTC]
\label{thm:n9v72-JDTC}
Under JDTC, the bosonic metric--gauge--Higgs boundary theory has a closed
semibounded Hamiltonian whose ground-normalized semigroup is
positivity preserving on the physical subspace.  It therefore supplies
the bosonic joint transfer kernel required by V71.  The conclusion is
conditional on the positive physical treatment of the Einstein
conformal direction and does not include the chiral fermion sector.
\end{theorem}

\begin{proof}
Use \Cref{prop:n9v72-principal-Schur,
thm:n9v72-KLMN,prop:n9v72-real-potential,
thm:n9v72-joint-bosonic,prop:n9v72-invariant-subspace}.  JDTC items 1
and 8 are explicit inputs because of
\Cref{thm:n9v72-conformal-no-go,fire:n9v72-fermion-type}.
\end{proof}

\begin{assessment}[Progress at ninth-series V72]
\label{ass:n9v72-status}
The abstract joint transfer requirement now has a closed-form route for
the bosonic sector: a positive principal Schur block, KLMN lower rows,
and Beurling--Deny positivity.  The same analysis exposes a genuine
upstream obstruction.  The real unreduced Einstein DeWitt metric has a
negative conformal principal direction, so it cannot be inserted into a
scalar Dirichlet heat kernel and repaired later by lower-order rows.
\end{assessment}

\begin{programme}[Ninth-series V73: conformal-mode alternatives]
\label{prog:n9v72-V73}
V73 will compare three precise upstream routes: fixed-volume/unimodular
reduction, a positive higher-derivative conformal reference with
cofinal tuning, and complex contour rotation.  It will test each against
the weak Einstein source, OS positivity, anomaly ownership, and recovery
of the full metric continuum, rejecting any route that merely deletes a
physical equation.
\end{programme}

\begin{firewall}[Claim boundary after ninth-series V72]
\label{fire:n9v72-boundary}
V72 proves the conditional Dirichlet/KLMN construction and the
unreduced conformal no-go.  It does not solve the conformal factor or
Hamiltonian constraint, construct the physical chiral transfer,
establish anomaly cancellation, sum all scale debits, close the
infrared theory, reconstruct the OS continuum, or prove the full
SM--Einstein continuum.
\end{firewall}

\begin{theorem}[Ninth-series V72 result]
\label{thm:n9v72-result}
A positive physical metric principal form, strict metric--matter Schur
margin, KLMN lower-row bound, and Beurling--Deny property produce the
bosonic joint transfer semigroup required by CJTC.  The unreduced real
Einstein DeWitt form fails this route because its pure conformal
coefficient is $d(1-d/(d-1))<0$; lower-order relative bounds cannot
repair that principal sign.
\end{theorem}

\begin{proof}
Combine \Cref{thm:n9v72-Dirichlet-transfer,
prop:n9v72-principal-Schur,lem:n9v72-conformal-negative,
thm:n9v72-conformal-no-go,thm:n9v72-joint-bosonic,
thm:n9v72-JDTC}.
\end{proof}

\paragraph{Parent-version record.}
V72 was formed from
\texttt{Renewal\_SM\_Einstein\_Continuum\_Research\_Notes\_V71.tex}.
The V71 parent had $26\,934$ logical source lines, $971\,389$ bytes, and
SHA--256
\[
 \texttt{3CE49506BC39A7B4B093640E39AC73D1A63593E111637BD0B511B52B3E31E2C3}.
\]

% =============================================================================
% END APPENDED NINTH-SERIES V72 MATERIAL
% =============================================================================
% =============================================================================
% BEGIN APPENDED NINTH-SERIES V73 MATERIAL
% =============================================================================

\section{Ninth-series V73: conformal-mode alternatives and rejection tests}
\label{sec:v9v73}

\subsection{Conformal decomposition and the required tests}

V72 proved that the real unreduced DeWitt principal form is indefinite
in the conformal direction.  Write a spacetime metric locally as
\begin{equation}
 g=e^{2\varphi}\widehat g,
 \qquad
 \det\widehat g=\det\bar g,
 \label{eq:n9v73-conformal-split}
\end{equation}
so that variations split into a trace-free part and
$2\dot\varphi g$.  A viable treatment must satisfy four independent
tests:
\begin{enumerate}
\item a positive physical principal transfer form;
\item recovery of the trace Einstein equation and cosmological owner;
\item compatibility with the weak tensor-source topology of V23--V27;
and
\item reflection positivity or an explicit replacement sufficient for
OS reconstruction.
\end{enumerate}

\subsection{Global fixed volume does not remove the local instability}

\begin{proposition}[Fixed total volume removes only one mode]
\label{prop:n9v73-global-volume}
At first order, the fixed-volume condition is
\begin{equation}
 \int_M\operatorname{tr}_g h\,dV_g=0.
 \label{eq:n9v73-global-volume-condition}
\end{equation}
It does not remove local conformal variations.  Indeed, for every
nonzero smooth $\varphi$ with $\int_M\varphi,dV_g=0$, the variation
$h=\varphi g$ satisfies \eqref{eq:n9v73-global-volume-condition} but has
negative DeWitt norm at the general-relativistic parameter.
\end{proposition}

\begin{proof}
For $h=\varphi g$,
$\operatorname{tr}_gh=n\varphi$, so the mean-zero condition gives
\eqref{eq:n9v73-global-volume-condition}.  V72 gives
\begin{equation}
 {\cal G}_{1/(n-1)}(h,h)
 =n\left(1-\frac n{n-1}\right)\varphi^2<0.
\end{equation}
There are infinitely many mean-zero choices of $\varphi$.
\end{proof}

\begin{firewall}[Canonical fixed volume versus pointwise unimodularity]
\label{fire:n9v73-volume-unimodular}
Fixing the single number $\operatorname{Vol}(M)$ is not the pointwise
condition $\det g=\det\bar g$.  The former leaves all mean-zero
conformal modes, while the latter removes the local trace tangent.  The
two constructions must not be identified.
\end{firewall}

\subsection{Pointwise unimodular reduction}

On the pointwise slice
\begin{equation}
 {\cal U}_{\bar g}=\{g:dV_g=dV_{\bar g}\},
 \label{eq:n9v73-unimodular-slice}
\end{equation}
the tangent condition is
\begin{equation}
 \operatorname{tr}_gh=0.
 \label{eq:n9v73-tracefree-tangent}
\end{equation}

\begin{proposition}[Principal sign on the unimodular tangent]
\label{prop:n9v73-unimodular-sign}
For every $h$ satisfying \eqref{eq:n9v73-tracefree-tangent},
\begin{equation}
 {\cal G}_\lambda(h,h)=|h|_g^2\ge0
 \label{eq:n9v73-unimodular-positive}
\end{equation}
for every $\lambda$.  Thus the algebraic conformal negative direction
from V72 is absent on this tangent space.
\end{proposition}

\begin{proof}
The trace term in \eqref{eq:n9v72-DeWitt-form} vanishes.
\end{proof}

The restricted Euler equation is the trace-free Einstein equation.  Let
$E_{\mu\nu}$ denote the complete weak Einstein residual, including the
renormalized Standard Model source and all local owners.

\begin{theorem}[Trace recovery from conservation]
\label{thm:n9v73-trace-recovery}
Let $M$ be connected and suppose, in the distributional sense,
\begin{equation}
 E_{\mu\nu}-\frac1n(\operatorname{tr}_gE)g_{\mu\nu}=0,
 \qquad
 \nabla^\mu E_{\mu\nu}=0.
 \label{eq:n9v73-tracefree-conserved}
\end{equation}
Then there is a constant $c$ such that
\begin{equation}
 E_{\mu\nu}=c,g_{\mu\nu}.
 \label{eq:n9v73-E-constant}
\end{equation}
Thus the full Einstein equation is recovered up to one cosmological
integration constant on each connected component.
\end{theorem}

\begin{proof}
The first equation gives $E_{\mu\nu}=\tau g_{\mu\nu}/n$ with
$\tau=\operatorname{tr}_gE$ a distribution.  Metric compatibility and
the second equation give $\nabla_\nu\tau=0$.  A distribution with zero
derivative on a connected manifold is constant.  Substitute back.
\end{proof}

\begin{corollary}[Weak measure-source compatibility]
\label{cor:n9v73-weak-source}
The conclusion remains valid when $E$ initially has measure-valued
coefficients, provided its trace and covariant divergence are defined
as distributions.  Equation \eqref{eq:n9v73-tracefree-conserved} then
forces the trace measure to be a constant multiple of $dV_g$.
\end{corollary}

\begin{proof}
The proof of \Cref{thm:n9v73-trace-recovery} is distributional and shows
that the trace distribution is constant; such a distribution is the
declared constant density relative to the metric volume.
\end{proof}

\begin{firewall}[Conservation is a physical hypothesis]
\label{fire:n9v73-conservation}
The trace-recovery theorem needs the complete renormalized Ward identity.
A gravitational or gauge anomaly producing
$\nabla^\mu E_{\mu\nu}\ne0$ invalidates the integration-constant
argument.  A trace anomaly is compatible with the theorem only when it
is included in a conserved renormalized source.
\end{firewall}

\begin{firewall}[Equation recovery is not quantum equivalence]
\label{fire:n9v73-unimodular-quantum}
Theorem \ref{thm:n9v73-trace-recovery} recovers the classical or weak
field equation up to $c$.  It does not prove equality of off-shell
Einstein and pointwise-unimodular path measures, equality of their
correlation functions, or recovery of local determinant fluctuations.
The slice measure, its coarea determinant, and its joint transfer kernel
remain separate obligations.
\end{firewall}

\begin{proposition}[Global-volume Laplace relation]
\label{prop:n9v73-volume-Laplace}
If a rigorous coarea disintegration by total volume exists with positive
fixed-volume partition density $Z(V)$, then the cosmological canonical
partition function satisfies
\begin{equation}
 Z(\Lambda)=\int_0^\infty
 e^{-\Lambda V/\hbar}Z(V)\,dV.
 \label{eq:n9v73-volume-Laplace}
\end{equation}
This relation concerns global volume sectors and does not prove
equivalence with the pointwise slice \eqref{eq:n9v73-unimodular-slice}.
\end{proposition}

\begin{proof}
Disintegrate the measure by $V=\operatorname{Vol}(g)$; the cosmological
term is constant and equal to $\Lambda V$ on each fibre.  Tonelli gives
\eqref{eq:n9v73-volume-Laplace}.  Apply
\Cref{fire:n9v73-volume-unimodular} for the final statement.
\end{proof}

\subsection{Positive higher-derivative reference}

Around a flat background, a positive $R^2$ row gives a conformal
quadratic term of fourth derivative order.  Abstract its competition
with the Einstein conformal row as
\begin{equation}
 q_{\alpha,\beta}(\varphi)
 =\alpha\lVert\Delta\varphi\rVert_2^2
  -\beta\lVert\nabla\varphi\rVert_2^2,
 \qquad \alpha,\beta>0.
 \label{eq:n9v73-p4-p2}
\end{equation}

\begin{theorem}[Higher derivatives do not cure the cofinal infrared sign]
\label{thm:n9v73-higher-IR}
On a Laplacian eigenmode with eigenvalue $\lambda>0$,
\begin{equation}
 q_{\alpha,\beta}=(\alpha\lambda^2-\beta\lambda)
 \lVert\varphi_\lambda\rVert_2^2.
 \label{eq:n9v73-p4-p2-symbol}
\end{equation}
It is negative for $0<\lambda<\beta/\alpha$.  On growing tori,
$\lambda_1\downarrow0$, so no fixed $\alpha$ makes
\eqref{eq:n9v73-p4-p2} nonnegative uniformly in volume.
\end{theorem}

\begin{proof}
The eigenmode calculation gives
\eqref{eq:n9v73-p4-p2-symbol}.  The displayed polynomial equals
$\lambda(\alpha\lambda-\beta)$, which has the stated sign.  Growing
tori supply positive eigenvalues arbitrarily close to zero.
\end{proof}

\begin{corollary}[Semibounded but not positive]
\label{cor:n9v73-higher-semibounded}
For every eigenvalue,
\begin{equation}
 \alpha\lambda^2-\beta\lambda
 \ge-\frac{\beta^2}{4\alpha}.
 \label{eq:n9v73-interpolation}
\end{equation}
Thus the fourth-order form can be ground-shifted and treated by KLMN,
but this does not provide a nonnegative massless principal form.
\end{corollary}

\begin{proof}
Complete the square
$\alpha(\lambda-\beta/(2\alpha))^2-\beta^2/(4\alpha)$.
\end{proof}

\begin{theorem}[Biharmonic semigroup is not positivity preserving]
\label{thm:n9v73-biharmonic-not-Markov}
On the real line, the semigroup with Fourier multiplier
$e^{-t p^4}$ does not have a nonnegative probability kernel for
$t>0$.  Hence a positive fourth-order quadratic form is not a scalar
Dirichlet diffusion merely because its operator is nonnegative.
\end{theorem}

\begin{proof}
Assume a nonnegative integrable kernel $k_t$ of mass one existed.  Its
characteristic function would be $e^{-tp^4}$.  Then
\begin{equation}
 \frac{1-e^{-tp^4}}{p^2}
 =\int_{\mathbb R}\frac{1-\cos(px)}{p^2}k_t(x)\,dx.
 \end{equation}
The left side tends to zero.  Fatou's lemma and
$(1-\cos(px))/p^2\to x^2/2$ give
$\int x^2k_t(x)dx=0$.  Thus $k_t$ is concentrated at zero, whose
characteristic function is one, a contradiction.
\end{proof}

\begin{proposition}[Fourth-order propagator residue obstruction]
\label{prop:n9v73-fourth-OS}
For $0<m_1<m_2$,
\begin{equation}
 \frac1{(p_0^2+m_1^2)(p_0^2+m_2^2)}
 =\frac1{m_2^2-m_1^2}
 \left(\frac1{p_0^2+m_1^2}-
       \frac1{p_0^2+m_2^2}\right).
 \label{eq:n9v73-fourth-partial-fraction}
\end{equation}
The negative residue violates the V67 rational OS criterion.  Thus a
generic factored positive fourth-order covariance is not reflection
positive.
\end{proposition}

\begin{proof}
The partial fraction identity is direct.  Apply
\Cref{thm:n9v67-rational-criterion}.
\end{proof}

\begin{assessment}[Higher-derivative verdict]
\label{ass:n9v73-higher-verdict}
A positive $p^4$ row can regularize ultraviolet estimates and make the
conformal operator semibounded.  It neither removes the cofinal
infrared negative $p^2$ band nor supplies a positivity-preserving or
OS-positive scalar transfer.  It is therefore not a standalone solution
of the conformal problem.
\end{assessment}

\subsection{Complex conformal contour}

The formal rotation $\varphi=i\chi$ changes a negative Gaussian
quadratic form into a decaying one.  The zero-dimensional model already
shows the positivity cost.

\begin{proposition}[Gaussian contour sign]
\label{prop:n9v73-contour-sign}
The real integral
\begin{equation}
 \int_{\mathbb R}e^{+a\varphi^2/(2\hbar)}d\varphi,
 \qquad a>0,
 \label{eq:n9v73-divergent-Gaussian}
\end{equation}
diverges.  On the contour $\varphi=i\chi$, its normalized formal second
moment is
\begin{equation}
 \langle\varphi^2\rangle_{i\mathbb R}
 =-\langle\chi^2\rangle_{\mathbb R}
 =-\frac\hbar a<0.
 \label{eq:n9v73-negative-moment}
\end{equation}
Therefore the rotated conformal coordinate cannot be interpreted as a
real reflection-positive scalar observable.
\end{proposition}

\begin{proof}
The divergence is immediate.  Substitute $\varphi=i\chi$; after
normalization the contour Jacobian cancels, while
$\varphi^2=-\chi^2$ and the decaying Gaussian has variance $\hbar/a$.
A real scalar observable with negative squared norm violates positivity.
\end{proof}

\begin{firewall}[Contour deformation needs a homology theorem]
\label{fire:n9v73-contour-homology}
For the nonlinear metric action, a contour rotation is valid only if the
integrand is holomorphic in a domain connecting the contours, boundary
arcs vanish, determinant zeros and branch cuts are avoided, and the
chosen contour is compatible with gauge and diffeomorphism orbits.
The chiral determinant and metric degeneracy make these substantive
conditions.
\end{firewall}

\begin{firewall}[Physical observables after rotation]
\label{fire:n9v73-contour-observables}
One may hope that the conformal coordinate is removed from the physical
OS algebra by constraints.  That is a constraint-reconstruction
argument, not scalar reflection positivity of the rotated metric.  It
must prove positivity of the remaining observable algebra and recover
the weak Einstein trace equation.
\end{firewall}

\subsection{Comparison and acquisition decision}

\begin{center}
\small
\begin{tabular}{|p{0.18\textwidth}|p{0.21\textwidth}|p{0.21\textwidth}|p{0.25\textwidth}|}
\hline
Route & Principal sign & Einstein trace & OS/measure status\\
\hline
Global fixed volume & local negative modes remain & one global mode
fixed & does not solve V72 obstruction\\
\hline
Pointwise unimodular & trace-free principal form positive & recovered
up to cosmological constant under conservation & slice coarea and
quantum equivalence open\\
\hline
Positive $R^2$ reference & UV positive, IR negative band & retained in
formal action & fourth-order Markov and OS tests fail generically\\
\hline
Complex contour & rotated Gaussian decays & analytic continuation
required & no positive real conformal observable measure\\
\hline
\end{tabular}
\end{center}

\begin{decision}[Conformal acquisition order]
\label{dec:n9v73-acquisition}
The next constructive route is constraint reduction rather than a
standalone fourth-order or contour repair.  V74 will treat the conformal
factor as a variable solved by the Hamiltonian/York constraint, derive
its exact implicit derivatives and Schur complement, and test whether a
positive reduced physical form results.  Pointwise unimodular trace
recovery remains a comparison check, not an asserted quantum
equivalence.
\end{decision}

\begin{definition}[Conformal-mode comparison card (CMCC)]
\label{def:n9v73-CMCC}
CMCC requires for any proposed route:
\begin{enumerate}
\item removal of the V72 negative principal direction without a false
volume-uniform gap;
\item recovery of the complete weak trace Einstein equation and the V58
cosmological owner;
\item a positive joint transfer or a proved physical replacement;
\item regulator-level diffeomorphism/gauge Ward identities and anomaly
control;
\item compatibility with the chiral determinant and metric chart
measure;
\item cofinal recovery of all physical metric observables; and
\item no double counting between constraint, contour, and counterterm
rows.
\end{enumerate}
\end{definition}

\begin{theorem}[V73 conformal comparison]
\label{thm:n9v73-CMCC}
Global fixed volume fails CMCC item 1.  A fixed positive higher-
derivative reference fails the cofinal infrared and generic OS tests.
A naive imaginary conformal contour fails positivity for the real
conformal observable.  Pointwise unimodular reduction removes the
algebraic negative direction and recovers the field equation up to a
cosmological constant under a conserved Ward identity, but its quantum
measure equivalence remains unproved.  Constraint reduction is therefore
the next nonrejected acquisition.
\end{theorem}

\begin{proof}
Use \Cref{prop:n9v73-global-volume,
prop:n9v73-unimodular-sign,thm:n9v73-trace-recovery,
thm:n9v73-higher-IR,thm:n9v73-biharmonic-not-Markov,
prop:n9v73-fourth-OS,prop:n9v73-contour-sign}.
\end{proof}

\begin{assessment}[Progress at ninth-series V73]
\label{ass:n9v73-status}
The conformal alternatives now have explicit rejection tests.  Neither
global volume fixing, a positive $R^2$ ultraviolet row, nor a naive
complex rotation supplies the full positive Einstein transfer.
Pointwise trace-free reduction has a rigorous weak equation-recovery
theorem but not a quantum equivalence theorem.  The programme therefore
moves to solving the conformal constraint and inspecting the reduced
Schur form.
\end{assessment}

\begin{programme}[Ninth-series V74: implicit conformal constraint]
\label{prog:n9v73-V74}
V74 will formulate a finite-regulator York/Lichnerowicz constraint
$F(\varphi,y)=0$, prove local solution and first/second response formulas
under a coercive conformal linearization, and compute the exact reduced
Hessian.  It will identify the gap and sign conditions needed for a
positive nonlocal physical metric form.
\end{programme}

\begin{firewall}[Claim boundary after ninth-series V73]
\label{fire:n9v73-boundary}
V73 rejects two standalone repairs and proves conditional trace recovery
for the unimodular equation.  It does not solve the Hamiltonian
constraint, establish quantum equivalence of a reduced measure, build
the chiral joint transfer, prove anomaly cancellation, close the
infrared sector, reconstruct the OS continuum, or prove the full
SM--Einstein continuum.
\end{firewall}

\begin{theorem}[Ninth-series V73 result]
\label{thm:n9v73-result}
The conformal principal obstruction cannot be removed by global volume
fixing or by a fixed positive fourth-order row with uniform OS control.
Pointwise unimodular reduction makes the DeWitt tangent positive and,
with distributional conservation, recovers the full weak Einstein
equation up to one cosmological constant.  A complex contour does not
provide a positive real conformal observable measure.  The remaining
constructive target is the exact constraint-reduced physical form.
\end{theorem}

\begin{proof}
Combine \Cref{prop:n9v73-global-volume,
prop:n9v73-unimodular-sign,thm:n9v73-trace-recovery,
thm:n9v73-higher-IR,prop:n9v73-fourth-OS,
prop:n9v73-contour-sign,thm:n9v73-CMCC}.
\end{proof}

\paragraph{Parent-version record.}
V73 was formed from
\texttt{Renewal\_SM\_Einstein\_Continuum\_Research\_Notes\_V72.tex}.
The V72 parent had $27\,348$ logical source lines, $987\,286$ bytes, and
SHA--256
\[
 \texttt{1DA15B47505FE54216D55B2031E7B8A6E59D955B92F5D0BCA3894DCF553A4E4C}.
\]

% =============================================================================
% END APPENDED NINTH-SERIES V73 MATERIAL
% =============================================================================
% =============================================================================
% BEGIN APPENDED NINTH-SERIES V74 MATERIAL
% =============================================================================

\section{Ninth-series V74: implicit conformal constraint and reduced Hessian}
\label{sec:v9v74}

\subsection{Finite-regulator constraint branch}

Let $Y_n$ be the finite-dimensional retained physical data at regulator
$n$ and let $X_n$ be the real conformal-variable space.  Write the
regulated Hamiltonian/York constraint as
\begin{equation}
 F_n(\varphi,y)=0,
 \qquad F_n:X_n\times Y_n\to X_n.
 \label{eq:n9v74-constraint}
\end{equation}
Fix a solution $(\bar\varphi,\bar y)$ and put
\begin{equation}
 L=D_\varphi F_n(\bar\varphi,\bar y).
 \label{eq:n9v74-linearized-constraint}
\end{equation}

\begin{theorem}[Local conformal branch]
\label{thm:n9v74-IFT}
If $L$ is invertible, there are neighbourhoods of $\bar y$ and
$\bar\varphi$ and a unique $C^2$ map $\varphi(y)$ such that
\begin{equation}
 F_n(\varphi(y),y)=0,
 \qquad \varphi(\bar y)=\bar\varphi.
 \label{eq:n9v74-local-branch}
\end{equation}
For tangent directions $h,k\in Y_n$,
\begin{align}
 D_y\varphi[h]
 &=-L^{-1}D_yF[h],
 \label{eq:n9v74-phi-first}\\
 D_y^2\varphi[h,k]
 &=-L^{-1}\Bigl(
 D_y^2F[h,k]
 +D_yD_\varphi F[h,D_y\varphi[k]]
 +D_\varphi D_yF[D_y\varphi[h],k]\\
 &\hspace{7.5em}
 +D_\varphi^2F[D_y\varphi[h],D_y\varphi[k]]
 \Bigr).
 \label{eq:n9v74-phi-second}
\end{align}
All derivatives on the right are evaluated at the background solution.
\end{theorem}

\begin{proof}
The finite-dimensional implicit-function theorem gives the branch.
Differentiate \eqref{eq:n9v74-local-branch} once:
$LD_y\varphi[h]+D_yF[h]=0$.  Differentiate that identity in direction
$k$, retain both mixed terms, and solve with $L^{-1}$.
\end{proof}

\begin{corollary}[Quantitative response]
\label{cor:n9v74-response-bound}
If $L\ge\lambda I$ and
\begin{equation}
 \lVert D_yF[h]\rVert\le b_1\lVert h\rVert,
 \label{eq:n9v74-F-first-bound}
\end{equation}
then
\begin{equation}
 \lVert D_y\varphi[h]\rVert
 \le\lambda^{-1}b_1\lVert h\rVert.
 \label{eq:n9v74-phi-first-bound}
\end{equation}
If the four second-derivative rows in
\eqref{eq:n9v74-phi-second} have bounds $b_{20},b_{11},b_{11}',b_{02}$
in the corresponding product norms, then
\begin{equation}
 \lVert D_y^2\varphi[h,k]\rVert
 \le\lambda^{-1}
 \left[b_{20}+(b_{11}+b_{11}')\lambda^{-1}b_1
       +b_{02}\lambda^{-2}b_1^2\right]
 \lVert h\rVert\lVert k\rVert.
 \label{eq:n9v74-phi-second-bound}
\end{equation}
\end{corollary}

\begin{proof}
$\lVert L^{-1}\rVert\le\lambda^{-1}$.  Apply this to
\eqref{eq:n9v74-phi-first} and then to each term of
\eqref{eq:n9v74-phi-second}.
\end{proof}

\subsection{A semilinear York model and its floor}

A standard finite-regulator conformal constraint has the schematic
form
\begin{equation}
 F(\varphi,y)
 =-c_d\Delta_y\varphi+R_y\varphi
  -A_y\varphi^{-p}+B_y\varphi^q,
 \qquad c_d,p,q>0,quad\varphi>0.
 \label{eq:n9v74-semi-York}
\end{equation}
Its linearization is
\begin{equation}
 L=-c_d\Delta_y+R_y+pA_y\bar\varphi^{-p-1}
       +qB_y\bar\varphi^{q-1}.
 \label{eq:n9v74-York-linearization}
\end{equation}

\begin{proposition}[Coercive branch criterion]
\label{prop:n9v74-York-floor}
Suppose the discrete Laplacian is nonnegative and
\begin{equation}
 R_y+pA_y\bar\varphi^{-p-1}
       +qB_y\bar\varphi^{q-1}\ge\lambda>0
 \label{eq:n9v74-York-potential-floor}
\end{equation}
pointwise.  Then $L\ge\lambda I$, the local branch of
\Cref{thm:n9v74-IFT} exists, and it remains positive after shrinking the
$Y_n$ neighbourhood whenever $\bar\varphi$ has a positive minimum.
\end{proposition}

\begin{proof}
The quadratic form of $-c_d\Delta_y$ is nonnegative, so
\eqref{eq:n9v74-York-potential-floor} gives the operator floor.  Apply
the implicit-function theorem.  Continuity of the finite-dimensional
solution map in a norm controlling pointwise values preserves half the
positive minimum on a small neighbourhood.
\end{proof}

\begin{proposition}[Uniqueness in a monotone interval]
\label{prop:n9v74-monotone-uniqueness}
Suppose two positive solutions $\varphi_1,\varphi_2$ lie in an interval
on which the derivative of the nonlinear potential in
\eqref{eq:n9v74-semi-York} is at least $\lambda>0$.  Then
$\varphi_1=\varphi_2$.
\end{proposition}

\begin{proof}
Subtract the two equations and test with
$w=\varphi_1-\varphi_2$.  The Laplacian contributes
$c_d\lVert\nabla w\rVert_2^2$.  The mean-value theorem makes the
potential contribution at least $\lambda\lVert w\rVert_2^2$.  Their sum
can vanish only for $w=0$.
\end{proof}

\begin{firewall}[The floor is not automatic]
\label{fire:n9v74-floor}
The coefficients $R_y,A_y,B_y$ contain curvature, transverse momentum,
mean curvature, cosmological, and matter energy rows with convention-
dependent signs.  Equation \eqref{eq:n9v74-York-potential-floor} must be
proved on the accepted physical sector; it cannot be inferred from the
ellipticity of $-\Delta$ alone.
\end{firewall}

\subsection{Quasi-locality of the solved conformal response}

Assume $L$ has fixed interaction range $R_0$ and a uniform off-diagonal
row in a block decomposition.

\begin{theorem}[Gapped constraint response is quasi-local]
\label{thm:n9v74-response-locality}
If $L\ge\lambda I$ with regulator-uniform $\lambda>0$, then its inverse
has the V65 weighted decay
\begin{equation}
 \lVert P_xL^{-1}P_y\rVert
 \le C_Le^{-m_Ld(x,y)}.
 \label{eq:n9v74-L-inverse-decay}
\end{equation}
If $D_yF[h_x]$ has a fixed collar around $x$, then the kernel of
$D_y\varphi[h_x]$ has the same exponential decay up to that collar.
The second response in \eqref{eq:n9v74-phi-second} is a finite
convolution of these kernels and remains exponentially summable at every
strictly smaller rate.
\end{theorem}

\begin{proof}
Apply \Cref{lem:n9v65-weighted-inverse} to $L$.  Formula
\eqref{eq:n9v74-phi-first} is one inverse followed by a local insertion.
Every term in \eqref{eq:n9v74-phi-second} is an outer $L^{-1}$ followed
by local first/second derivative insertions and at most two first
responses.  Convolution of exponentially weighted rows preserves any
strictly smaller exponent.
\end{proof}

\begin{firewall}[Constraint criticality]
\label{fire:n9v74-criticality}
If the smallest eigenvalue of $L$ tends to zero, the conformal response
can become long range and the branch may bifurcate.  A pseudo-inverse
does not give the uniform bound \eqref{eq:n9v74-L-inverse-decay}.  The
closing mode must be calibrated as a marginal/infrared variable or the
programme must change branches explicitly.
\end{firewall}

\subsection{Exact reduced-action calculus}

Let $S(\varphi,y)$ be a twice differentiable finite-regulator action and
define
\begin{equation}
 S_{\rm red}(y)=S(\varphi(y),y).
 \label{eq:n9v74-reduced-action}
\end{equation}

\begin{proposition}[General pullback Hessian]
\label{prop:n9v74-general-pullback}
Writing $\varphi_h=D_y\varphi[h]$ and
$\varphi_{hk}=D_y^2\varphi[h,k]$, one has
\begin{align}
 D_yS_{\rm red}[h]
 &=S_y[h]+S_\varphi[\varphi_h],
 \label{eq:n9v74-Sred-first}\\
 D_y^2S_{\rm red}[h,k]
 &=S_{yy}[h,k]+S_{y\varphi}[h,\varphi_k]
   +S_{\varphi y}[\varphi_h,k]
   +S_{\varphi\varphi}[\varphi_h,\varphi_k]
   +S_\varphi[\varphi_{hk}].
 \label{eq:n9v74-Sred-second}
\end{align}
\end{proposition}

\begin{proof}
Apply the first- and second-order chain rules to
\eqref{eq:n9v74-reduced-action}.
\end{proof}

\begin{firewall}[Constraint equation versus stationary elimination]
\label{fire:n9v74-constraint-stationary}
The Hamiltonian constraint $F=0$ arises from lapse variation and need
not equal $S_\varphi=0$ for the Euclidean boundary action under study.
In that case the last term and both mixed terms in
\eqref{eq:n9v74-Sred-second} must be retained.  Calling the result a
Schur complement without proving $F=S_\varphi$ is incorrect.
\end{firewall}

\begin{theorem}[Stationary constraint Schur complement]
\label{thm:n9v74-stationary-Schur}
Suppose on the branch
\begin{equation}
 F=S_\varphi=0,
 \qquad L=S_{\varphi\varphi}>0.
 \label{eq:n9v74-stationary-hyp}
\end{equation}
Put $A=S_{yy}$ and $B=S_{y\varphi}$.  Then
\begin{align}
 D_yS_{\rm red}[h]&=S_y[h],
 \label{eq:n9v74-stationary-first}\\
 D_y^2S_{\rm red}&=A-BL^{-1}B^*.
 \label{eq:n9v74-reduced-Schur}
\end{align}
If $A>0$ and
\begin{equation}
 \lVert A^{-1/2}BL^{-1/2}\rVert\le\gamma<1,
 \label{eq:n9v74-reduced-margin-hyp}
\end{equation}
then
\begin{equation}
 D_y^2S_{\rm red}\ge(1-\gamma^2)A.
 \label{eq:n9v74-reduced-margin}
\end{equation}
\end{theorem}

\begin{proof}
$S_\varphi=0$ removes the reaction term in
\eqref{eq:n9v74-Sred-first}--\eqref{eq:n9v74-Sred-second}.  Differentiating
$S_\varphi=0$ gives $\varphi_h=-L^{-1}B^*h$.  Substitute this into the
remaining Hessian terms to obtain
\eqref{eq:n9v74-reduced-Schur}.  The norm hypothesis gives
$BL^{-1}B^*\le\gamma^2A$.
\end{proof}

\begin{firewall}[Positive constraint linearization is not enough]
\label{fire:n9v74-L-not-enough}
The floor $L>0$ proves branch stability and locality.  The reduced
physical Hessian is positive only if the cross coefficient satisfies
\eqref{eq:n9v74-reduced-margin-hyp}, or if the full nonstationary
pullback \eqref{eq:n9v74-Sred-second} has a separate positive estimate.
\end{firewall}

\subsection{The coarea determinant}

At a regular isolated constraint root, elimination with a delta density
has an exact Jacobian.

\begin{proposition}[Finite-regulator constraint coarea]
\label{prop:n9v74-coarea}
Assume $\varphi(y)$ is the unique root in the integration chart.  Then
for a test density $G$,
\begin{equation}
 \int_{X_n}G(\varphi,y)\delta(F(\varphi,y))\,d\varphi
 =\frac{G(\varphi(y),y)}{|\det L(y)|}.
 \label{eq:n9v74-coarea}
\end{equation}
Thus the reduced effective action contains
\begin{equation}
 S_{\rm Jac}(y)=\hbar\log|\det L(y)|
 \label{eq:n9v74-Jac-action}
\end{equation}
with the sign shown for an $e^{-S/\hbar}$ density.
\end{proposition}

\begin{proof}
Use the multivariable delta-function change of variables
$u=F(\varphi,y)$ near the regular root.  Its Jacobian is
$|\det D_\varphi F|$.  Rewriting $|\det L|^{-1}$ as an exponential gives
\eqref{eq:n9v74-Jac-action}.
\end{proof}

\begin{corollary}[Jacobian response owner]
\label{cor:n9v74-Jac-response}
If $L>0$, then
\begin{align}
 D_yS_{\rm Jac}[h]
 &=\hbar\operatorname{Tr}(L^{-1}D_yL[h]),
 \label{eq:n9v74-Jac-first}\\
 D_y^2S_{\rm Jac}[h,k]
 &=\hbar\operatorname{Tr}\left(
 L^{-1}D_y^2L[h,k]
 -L^{-1}D_yL[k]L^{-1}D_yL[h]
 \right).
 \label{eq:n9v74-Jac-second}
\end{align}
It is an eliminated determinant row of the V65 type and must be counted
once.
\end{corollary}

\begin{proof}
Differentiate $\log\det L$ and use
$D(L^{-1})=-L^{-1}(DL)L^{-1}$.
\end{proof}

\begin{firewall}[Multiple roots and determinant sign]
\label{fire:n9v74-multiple-roots}
With several regular roots, \eqref{eq:n9v74-coarea} is a sum over
branches with absolute Jacobians.  At a singular root the formula fails,
and across a determinant-sign change no single smooth logarithmic owner
exists.  Branch selection is physical data, not a local notation choice.
\end{firewall}

\subsection{Reflection-compatible deterministic reduction}

Let $\Phi(y)=(\varphi(y),y)$ be the reconstructed metric map.

\begin{theorem}[Reflection-equivariant pushforward]
\label{thm:n9v74-reflection-pushforward}
Let $\mu$ be an OS-positive measure on retained data $y$.  Suppose
\begin{equation}
 \Phi(\Theta y)=\Theta\Phi(y)
 \label{eq:n9v74-reflection-equivariant}
\end{equation}
and composition with $\Phi$ sends positive-time observables to
positive-time observables.  Then the pushforward
$\Phi_\#\mu$ is OS-positive on the reconstructed physical metric
observable algebra.
\end{theorem}

\begin{proof}
For a positive-time observable $O$ of the reconstructed field,
\begin{align}
 \mathbb E_{\Phi_\#\mu}[(\Theta O)O]
 &=\mathbb E_\mu[Theta(O\circ\Phi)(O\circ\Phi)]\ge0.
\end{align}
The equality uses \eqref{eq:n9v74-reflection-equivariant}; the last
inequality is OS positivity of $\mu$.
\end{proof}

\begin{corollary}[Spatial constraint solving preserves time support]
\label{cor:n9v74-spatial-support}
If $F=0$ is solved independently on each spatial time slice and the
branch is reflection equivariant, spatial nonlocality of $L^{-1}$ does
not mix positive and negative times.  The hypothesis of
\Cref{thm:n9v74-reflection-pushforward} is then satisfied.
\end{corollary}

\begin{proof}
The value of $\Phi(y)$ at time $t>0$ depends only on retained data at the
same positive time slice.  Reflection covariance follows from uniqueness
of the branch and covariance of $F$.
\end{proof}

\begin{firewall}[Spacetime-elliptic solving can cross the reflection plane]
\label{fire:n9v74-time-nonlocal}
If the constraint branch is obtained from an elliptic inverse on the
full spacetime, a positive-time reconstructed observable may depend on
negative-time retained data.  Then deterministic pushforward does not
preserve the positive-time algebra and
\Cref{thm:n9v74-reflection-pushforward} does not apply.
\end{firewall}

\subsection{Conformal constraint reduction card}

\begin{definition}[Conformal constraint reduction card (CCRC)]
\label{def:n9v74-CCRC}
CCRC requires:
\begin{enumerate}
\item a real positive conformal branch solving the exact regulated
Hamiltonian/York constraint;
\item a regulator-uniform coercive linearization or an explicit retained
critical mode;
\item the first and second implicit responses
\eqref{eq:n9v74-phi-first}--\eqref{eq:n9v74-phi-second};
\item the full pullback Hessian, using a Schur complement only under
\eqref{eq:n9v74-stationary-hyp};
\item a strict reduced physical margin;
\item the coarea determinant and its V65 response owner;
\item reflection-equivariant spatial-slice reconstruction;
\item recovery of the trace equation by the V73 conservation theorem;
and
\item branch, anomaly, chiral, and cutoff compatibility.
\end{enumerate}
\end{definition}

\begin{theorem}[CCRC supplies a conditional positive conformal reduction]
\label{thm:n9v74-CCRC}
Under CCRC, the conformal variable is eliminated on a unique real branch
with controlled first and second responses.  The resulting physical
metric Hessian is the exact pullback
\eqref{eq:n9v74-Sred-second}, with the positive Schur estimate
\eqref{eq:n9v74-reduced-margin} in the stationary case.  The coarea
determinant is explicit, and spatial reflection-equivariant
reconstruction preserves OS positivity of the retained measure.
\end{theorem}

\begin{proof}
Use \Cref{thm:n9v74-IFT,thm:n9v74-response-locality,
prop:n9v74-general-pullback,thm:n9v74-stationary-Schur,
prop:n9v74-coarea,thm:n9v74-reflection-pushforward}.  CCRC item 8 invokes
\Cref{thm:n9v73-trace-recovery}.
\end{proof}

\begin{assessment}[Progress at ninth-series V74]
\label{ass:n9v74-status}
Constraint reduction now has exact local formulas and explicit gates.
A coercive conformal linearization gives a unique quasi-local branch;
the reduced Hessian is a Schur complement only for stationary
elimination; the delta constraint contributes its own determinant; and
spatial-slice reconstruction can preserve OS positivity.  Existence of
a regulator-uniform physical branch and its strict reduced margin remain
the hard inputs.
\end{assessment}

\begin{programme}[Ninth-series V75: compulsory companion audit]
\label{prog:n9v74-V75}
V75 will inspect the newest Yang--Mills and Navier--Stokes notes and
compare their current stable-sector, constraint, inverse-locality, and
branch mechanisms with CCRC.  V76 will continue after the audit.
\end{programme}

\begin{firewall}[Claim boundary after ninth-series V74]
\label{fire:n9v74-boundary}
V74 proves the implicit, Hessian, determinant, and reflection
compilers.  It does not prove a global York branch for physical SM data,
a uniform conformal floor or reduced margin, equivalence to the full
Einstein quantum measure, chiral transfer positivity, anomaly
cancellation, infrared closure, OS reconstruction, or the full
SM--Einstein continuum.
\end{firewall}

\begin{theorem}[Ninth-series V74 result]
\label{thm:n9v74-result}
An invertible conformal constraint linearization gives the exact
responses \eqref{eq:n9v74-phi-first}--
\eqref{eq:n9v74-phi-second}.  Stationary elimination produces the
reduced Schur Hessian $A-BL^{-1}B^*$ and margin
$1-\gamma^2$; delta-function elimination contributes
$\hbar\log|\det L|$; and a reflection-equivariant spatial constraint
pushforward preserves the retained OS form.  All conclusions remain
conditional on a real uniform physical branch.
\end{theorem}

\begin{proof}
Combine \Cref{thm:n9v74-IFT,thm:n9v74-response-locality,
thm:n9v74-stationary-Schur,prop:n9v74-coarea,
thm:n9v74-reflection-pushforward,thm:n9v74-CCRC}.
\end{proof}

\paragraph{Parent-version record.}
V74 was formed from
\texttt{Renewal\_SM\_Einstein\_Continuum\_Research\_Notes\_V73.tex}.
The V73 parent had $27\,807$ logical source lines, $1\,004\,754$ bytes,
and SHA--256
\[
 \texttt{E5E18100C64D5465E58FFF12DF16907347342FAB096D3445394CDA4E8D5F7E3F}.
\]

% =============================================================================
% END APPENDED NINTH-SERIES V74 MATERIAL
% =============================================================================

% =============================================================================
% BEGIN APPENDED NINTH-SERIES V75 MATERIAL
% =============================================================================

\section{Ninth-series V75: companion audit and minimum-energy constraint splitting}
\label{sec:v9v75}

\subsection{Companion audit at the compulsory fifth-version gate}

This is the compulsory companion inspection announced in
\Cref{prog:n9v74-V75}.  The files actually inspected were
\begin{align*}
 &\texttt{Renewal\_Yang\_Mills\_Mass\_Gap\_Research\_Notes\_V2.tex},\\
 &\texttt{Renewal\_Yang\_Mills\_Mass\_Gap\_Research\_Notes\_V100\_f7.tex},\\
 &\texttt{Renewal\_NS\_Stability\_Research\_Notes\_V26.tex}.
\end{align*}
Their byte counts and SHA--256 digests at inspection time were,
respectively,
\begin{center}
\begin{tabular}{c|r|l}
file & bytes & SHA--256\\ \hline
YM V2 & $52\,322$ &
\texttt{E1A288C922997FFDD9FE506837311DAE7DEA6F100E9769228B877BD0E352E2AF}\\
YM V100\_f7 & $1\,649\,875$ &
\texttt{816BF7CEA62C101B9826EE351952E54DC6D90C33F9D5911EF1156E1BF01B0723}\\
NS V26 & $278\,283$ &
\texttt{82C887F8C2C69E4F28FAD7C3DC8DE5B9099CB5A4BF431EBA1FFF3E91935978AD}
\end{tabular}
\end{center}

The Yang--Mills archive supplies an exact minimum-energy interpolation
and energy-orthogonal Gaussian retained/detail decomposition.  Its new
V2 supplies an exact shellable-loop plaquette word, a cubic BCH
substitution, and, under stated assembly symmetries, a quadratic symbol
split into the Yang--Mills and topological marginal owners plus an
$O(|k|^2)$ row.  It explicitly leaves the global equivariant filling
and complete stable-norm comparison open.  NS V26 is unchanged at this
inspection and its rank-one Oseen derivative estimates add no new
constraint branch mechanism.

\begin{decision}[Audit transfer]
\label{dec:n9v75-audit-transfer}
The minimum-energy interpolation is transferred below to the
constraint-reduced metric sector.  The transfer is made only at the
finite-regulator linearized level where it is exact.  The YM V2
curvature localization will later be used as a conditional input to
the matter-source rows of the York map; its unproved global assembly
and norm hypotheses are not imported as conclusions.  No new NS row is
imported.
\end{decision}

\subsection{Minimum-energy right inverse}

Let $X$ and $Y$ be finite-dimensional real Hilbert spaces.  Let
$H:X\to X$ be self-adjoint and strictly positive, and let
$Q:X\to Y$ be surjective.  Define
\begin{equation}
 C:=QH^{-1}Q^*,\qquad
 H_Y:=C^{-1},\qquad
 J:=H^{-1}Q^*C^{-1}.
 \label{eq:n9v75-C-HY-J}
\end{equation}

\begin{lemma}[Coarse covariance is invertible]
\label{lem:n9v75-C-positive}
$C$ is strictly positive on $Y$ and hence all operators in
\eqref{eq:n9v75-C-HY-J} are well defined.
\end{lemma}

\begin{proof}
For $y\ne0$,
\[
 \langle y,Cy\rangle_Y
 =\lVert H^{-1/2}Q^*y\rVert_X^2.
\]
Surjectivity of $Q$ implies injectivity of $Q^*$, so the last member is
strictly positive.
\end{proof}

\begin{theorem}[Exact minimum-energy lift]
\label{thm:n9v75-minimum-lift}
The map $J$ is a right inverse of $Q$ and satisfies
\begin{equation}
 QJ=I_Y,
 \qquad
 \langle z,HJy\rangle_X=0
 \quad(z\in\ker Q,\ y\in Y).
 \label{eq:n9v75-right-orthogonal}
\end{equation}
Consequently $Jy$ is the unique minimizer of
$\frac12\langle x,Hx\rangle_X$ subject to $Qx=y$, and
\begin{equation}
 \min_{Qx=y}\frac12\langle x,Hx\rangle_X
 =\frac12\langle y,H_Yy\rangle_Y.
 \label{eq:n9v75-minimum-value}
\end{equation}
\end{theorem}

\begin{proof}
Equation \eqref{eq:n9v75-C-HY-J} gives
$QJ=QH^{-1}Q^*C^{-1}=I_Y$ and $HJ=Q^*C^{-1}$.  Thus
\[
 \langle z,HJy\rangle_X
 =\langle Qz,C^{-1}y\rangle_Y=0
\]
for $z\in\ker Q$.  Every point of the affine fibre $Qx=y$ has the
unique form $x=Jy+z$ with $z\in\ker Q$.  Orthogonality gives
\begin{align*}
 \langle Jy+z,H(Jy+z)\rangle_X
 &=\langle Jy,HJy\rangle_X+\langle z,Hz\rangle_X\\
 &=\langle y,C^{-1}y\rangle_Y+\langle z,Hz\rangle_X.
\end{align*}
Strict positivity of $H$ proves uniqueness and
\eqref{eq:n9v75-minimum-value}.
\end{proof}

\begin{corollary}[Energy projection]
\label{cor:n9v75-energy-projection}
Put $P=JQ$ and $D=I_X-P$.  Then
\begin{equation}
 P^2=P,\qquad D^2=D,\qquad \operatorname{ran}D=\ker Q,
 \qquad HP=P^*H,
 \label{eq:n9v75-projections}
\end{equation}
and every $x\in X$ has the exact orthogonal split
\begin{equation}
 \langle x,Hx\rangle_X
 =\langle Qx,H_YQx\rangle_Y+\langle Dx,HDx\rangle_X.
 \label{eq:n9v75-energy-Pythagoras}
\end{equation}
\end{corollary}

\begin{proof}
$QJ=I$ gives $P^2=P$, $D^2=D$, $QD=0$, and $Dz=z$ for
$z\in\ker Q$.  Moreover
$HP=Q^*C^{-1}Q=P^*H$.  Apply
\Cref{thm:n9v75-minimum-lift} to $x=JQx+Dx$.
\end{proof}

\subsection{Gaussian disintegration without a mixed transfer}

Let $\gamma_H$ be the centred Gaussian probability measure on $X$
with precision $H$, and let $E=\ker Q$.  Denote by $H_E$ the
restriction of the quadratic form of $H$ to $E$.

\begin{theorem}[Exact retained/detail Gaussian factorization]
\label{thm:n9v75-Gaussian-factorization}
Under $\gamma_H$, the variables
\begin{equation}
 y=Qx,\qquad z=Dx
 \label{eq:n9v75-yz}
\end{equation}
are independent centred Gaussians.  Their precisions are $H_Y$ and
$H_E$, respectively, and
\begin{equation}
 x=Jy+z.
 \label{eq:n9v75-Gaussian-reconstruction}
\end{equation}
The conditional mean and covariance are
\begin{align}
 \mathbb E[x\mid Qx=y]&=Jy,
 \label{eq:n9v75-conditional-mean}\\
 \operatorname{Cov}(x\mid Qx)
 &=H^{-1}-H^{-1}Q^*C^{-1}QH^{-1}.
 \label{eq:n9v75-conditional-covariance}
\end{align}
In particular,
\begin{equation}
 \mathbb E[Dx\mid Qx]=0,
 \qquad
 \mathbb E[\langle u,Qx\rangle\langle v,Dx\rangle]=0
 \label{eq:n9v75-zero-linear-transfer}
\end{equation}
for all compatible $u,v$.
\end{theorem}

\begin{proof}
The linear map $Y\oplus E\to X$, $(y,z)\mapsto Jy+z$, is bijective:
surjectivity follows from $x=JQx+Dx$, and applying $Q$ proves
injectivity.  Its Jacobian is a nonzero constant.  The exponent splits
by \eqref{eq:n9v75-energy-Pythagoras}, so the joint density in $(y,z)$
is the product of the $H_Y$ and $H_E$ Gaussian densities.  This proves
independence and \eqref{eq:n9v75-conditional-mean}.  Since
$J\operatorname{Cov}(y)J^*=JCJ^*$, the detail covariance is
\begin{align*}
 H^{-1}-JCJ^*
 &=H^{-1}-H^{-1}Q^*C^{-1}QH^{-1},
\end{align*}
which proves \eqref{eq:n9v75-conditional-covariance}.  The two
identities in \eqref{eq:n9v75-zero-linear-transfer} follow from the
centering and independence of $z$.
\end{proof}

\begin{firewall}[The zero statement is conditional centering]
\label{fire:n9v75-zero-transfer}
Equation \eqref{eq:n9v75-zero-linear-transfer} does not say that an
interacting detail integration is zero or that nonlinear functions of
$z$ disappear.  It says exactly that the centred Gaussian detail has
zero conditional first moment and no retained/detail linear covariance.
Interactions generate cumulants which require the stable estimates of
the coarse recursion.
\end{firewall}

\subsection{Identification with the stationary Schur lift}

Write $X=X_d\oplus Y$, let $Q(\xi,y)=y$, and represent $H$ as
\begin{equation}
 H=
 \begin{pmatrix}
  L&B^*\\ B&A
 \end{pmatrix},
 \qquad L>0.
 \label{eq:n9v75-block-H}
\end{equation}
Define the Schur operator
\begin{equation}
 S:=A-BL^{-1}B^*.
 \label{eq:n9v75-S}
\end{equation}

\begin{proposition}[Block formula and equivalence of margins]
\label{prop:n9v75-block-formula}
The following are equivalent:
\begin{enumerate}
\item $H>0$;
\item $S>0$.
\end{enumerate}
When they hold,
\begin{equation}
 QH^{-1}Q^*=S^{-1},
 \qquad
 Jy=\binom{-L^{-1}B^*y}{y},
 \qquad
 H_Y=S.
 \label{eq:n9v75-block-J}
\end{equation}
Thus the minimum-energy coarse precision is exactly the stationary
Schur complement of \Cref{thm:n9v74-stationary-Schur}.
\end{proposition}

\begin{proof}
For $\xi\in X_d$ and $y\in Y$, completing the square gives
\begin{equation}
 \left\langle\binom{\xi}{y},
 H\binom{\xi}{y}\right\rangle
 =\langle \xi+L^{-1}B^*y,
 L(\xi+L^{-1}B^*y)\rangle+\langle y,Sy\rangle.
 \label{eq:n9v75-square-completion}
\end{equation}
This proves the equivalence.  The unique fibre minimizer is obtained
by setting the square to zero, giving the formula for $J$ and minimum
energy $\langle y,Sy\rangle$.  Uniqueness in
\Cref{thm:n9v75-minimum-lift} then gives $H_Y=S$ and hence
$QH^{-1}Q^*=S^{-1}$.
\end{proof}

\begin{corollary}[V74 cross-margin is the exact Gaussian gate]
\label{cor:n9v75-cross-margin}
If $A>0$ and
\begin{equation}
 \lVert A^{-1/2}BL^{-1/2}\rVert\le\gamma<1,
 \label{eq:n9v75-gamma}
\end{equation}
then $S\ge(1-\gamma^2)A>0$.  Consequently $H>0$, the Gaussian
factorization of \Cref{thm:n9v75-Gaussian-factorization} exists, and
its retained precision obeys the same lower bound.
\end{corollary}

\begin{proof}
Inequality \eqref{eq:n9v75-gamma} is equivalent to
$BL^{-1}B^*\le\gamma^2A$.  Use
\eqref{eq:n9v75-S} and \Cref{prop:n9v75-block-formula}.
\end{proof}

\begin{firewall}[Constraint slope versus energy lift]
\label{fire:n9v75-slope-distinction}
For the V74 constraint equation $F(\varphi,y)=0$, the tangent slope is
\begin{equation}
 R_F:=D_y\varphi=-L_F^{-1}F_y,
 \qquad L_F=D_\varphi F.
 \label{eq:n9v75-constraint-slope}
\end{equation}
The minimum-energy slope obtained from an action Hessian is
$R_*=-S_{\varphi\varphi}^{-1}S_{\varphi y}$.  These are equal only
after a compatibility statement relating the constraint to stationary
elimination.  Positivity of $L_F$ alone does not identify them.
\end{firewall}

\subsection{Gauge-source transfer from the new YM note}

The YM V2 symbol result can enter the Einstein constraint only through
an explicit source map.  Let $\rho_{\rm YM}(F)$ denote the regulated
gauge contribution to $F(\varphi,y)$, and suppose the declared assembly
hypotheses of that note hold.  At quadratic order write
\begin{equation}
 \rho_{\rm YM}^{(2)}
 =c_{\rm YM}\operatorname{Tr}F^2
  +c_\theta\operatorname{Tr}(F\widetilde F)
  +\rho_{\rm der}^{(2)},
 \qquad
 \lVert\widehat\rho_{\rm der}^{(2)}(k)\rVert\le C_2|k|^2.
 \label{eq:n9v75-YM-source-split}
\end{equation}

\begin{proposition}[Ownership consequence of the YM symbol split]
\label{prop:n9v75-YM-ownership}
Under \eqref{eq:n9v75-YM-source-split}, the first two terms must be
assigned to the running gauge/topological reference coefficients in the
calibrated unit sector of V70.  The last term has two additional coarse
momenta and is eligible for the stable sector after the missing YM
complete-norm comparison is proved.  Counting all three terms as a
small stable perturbation is invalid.
\end{proposition}

\begin{proof}
The first two symbols are nonzero at $k=0$ and hence survive the
four-dimensional canonical dilation as marginal owners.  The final
symbol vanishes to second order at $k=0$ and has the dimension-six
scaling predicted by the V70 calibrated decomposition.  Eligibility
is conditional because a pointwise Fourier estimate does not by itself
control the complete polymer/oscillation norm.
\end{proof}

\begin{firewall}[Audit hypotheses remain visible]
\label{fire:n9v75-YM-hypotheses}
The exact shellable-loop Stokes and BCH formulas do not prove global
equivariant filling ownership, treatment of torons and cohomology, or
the complete norm estimate.  Equation \eqref{eq:n9v75-YM-source-split}
is therefore a conditional transfer, not a proof that the full compact
gauge source already lies in the SM--Einstein stable slice.
\end{firewall}

\subsection{Audit ledger and next acquisition}

\begin{theorem}[Ninth-series V75 result]
\label{thm:n9v75-result}
At every finite regulator, a positive precision $H$ and surjective
retained map $Q$ determine the unique minimum-energy right inverse
$J=H^{-1}Q^*(QH^{-1}Q^*)^{-1}$.  It yields the exact energy Pythagoras
identity \eqref{eq:n9v75-energy-Pythagoras}, independent Gaussian
retained/detail variables, and zero conditional centred detail.  In a
detail/retained block, the retained precision is exactly the stationary
Schur complement.  The actual York tangent need not be this lift; its
defect is the next required quantity.
\end{theorem}

\begin{proof}
Combine \Cref{thm:n9v75-minimum-lift,
cor:n9v75-energy-projection,thm:n9v75-Gaussian-factorization,
prop:n9v75-block-formula,fire:n9v75-slope-distinction}.
\end{proof}

\begin{assessment}[Progress at ninth-series V75]
\label{ass:n9v75-status}
The companion audit has supplied an exact geometry for the Gaussian
physical/detail split and has sharpened the V74 Schur margin into the
coarse covariance criterion $QH^{-1}Q^*=S^{-1}$.  It also identifies
the unresolved mismatch between constraint solving and stationary
energy interpolation.  The new YM marginal/stable split fixes the
ownership order for the gauge source but remains conditional on its
global assembly and complete-norm gates.
\end{assessment}

\begin{programme}[Ninth-series V76: constraint--stationarity defect]
\label{prog:n9v75-V76}
V76 will continue immediately after this audit.  It will derive an
exact square-completion formula for the difference between the York
constraint slope $R_F$ and the stationary slope $R_*$, express that
difference through the defect $F-S_\varphi$, and bound the nonstationary
second-response reaction term.  The goal is a checkable positive
reduced-Hessian card without assuming $F=S_\varphi$.
\end{programme}

\begin{firewall}[Claim boundary after ninth-series V75]
\label{fire:n9v75-boundary}
V75 does not prove a global York branch, equality of the constraint and
stationary slopes, the YM global filling/norm gates, an interacting
zero transfer, a uniform continuum Schur margin, chiral transfer
positivity, anomaly cancellation, regulator independence, OS
reconstruction, or the full SM--Einstein continuum.
\end{firewall}

\paragraph{Parent-version record.}
V75 was formed from
\texttt{Renewal\_SM\_Einstein\_Continuum\_Research\_Notes\_V74.tex}.
The V74 parent had $28\,292$ counted source lines, $1\,021\,925$ bytes,
and SHA--256
\[
 \texttt{D68A759CB66DCEFBB6B3A92ADCA7B45DA4BE33BC47FBC606E0E1AF2AF7DD6E16}.
\]

% =============================================================================
% END APPENDED NINTH-SERIES V75 MATERIAL
% =============================================================================

% =============================================================================
% BEGIN APPENDED NINTH-SERIES V76 MATERIAL
% =============================================================================

\section{Ninth-series V76: the constraint--stationarity defect}
\label{sec:v9v76}

V75 separated two maps which had previously been easy to conflate: the
minimum-energy lift selected by an action Hessian and the tangent lift
selected by the York constraint.  This note computes their difference
exactly.  The computation shows that a nonstationary tangent slope does
not itself reduce the graph Hessian.  It contributes a positive square.
The only new indefinite quadratic row is the reaction of the
off-stationary action gradient against the curvature of the constraint
branch.

\subsection{Finite-regulator defect calculus}

Let $X_d$ be the eliminated detail space and $Y$ the retained physical
space.  Let $S:X_d\times Y\to\mathbb R$ be $C^3$ and let
$F:X_d\times Y\to X_d$ be $C^2$, after use of the fixed
finite-regulator Hilbert identification of $X_d$ and its dual.  Suppose
\begin{equation}
 F(\varphi(y),y)=0
 \label{eq:n9v76-constraint-branch}
\end{equation}
on a $C^2$ branch.  At a point of this branch define
\begin{align}
 D&:=S_{\varphi\varphi},&
 B^*&:=S_{\varphi y},&
 A&:=S_{yy},
 \label{eq:n9v76-action-blocks}\\
 L&:=F_\varphi,&
 G&:=F_y.
 \label{eq:n9v76-constraint-blocks}
\end{align}
Assume throughout this section that $D=D^*>0$ and that $L$ is
invertible.  Put
\begin{equation}
 S_0:=A-BD^{-1}B^*,
 \qquad
 R_*:=-D^{-1}B^*,
 \qquad
 R_F:=-L^{-1}G.
 \label{eq:n9v76-Schur-slopes}
\end{equation}
The first operator is the stationary Schur Hessian, while $R_*$ and
$R_F$ are the stationary and constraint tangent slopes.

\begin{definition}[Constraint--stationarity defect]
\label{def:n9v76-defect}
The vector-valued defect is
\begin{equation}
 \mathcal E(\varphi,y)
 :=F(\varphi,y)-\nabla_\varphi S(\varphi,y).
 \label{eq:n9v76-defect}
\end{equation}
Its first derivatives obey
\begin{equation}
 L=D+\mathcal E_\varphi,
 \qquad
 G=B^*+\mathcal E_y.
 \label{eq:n9v76-defect-derivatives}
\end{equation}
\end{definition}

\begin{lemma}[Exact tangent-defect identity]
\label{lem:n9v76-tangent-defect}
Let $M:=R_F-R_*$.  Then
\begin{equation}
 M
 =L^{-1}\bigl(
 \mathcal E_\varphi D^{-1}B^*-\mathcal E_y
 \bigr).
 \label{eq:n9v76-M}
\end{equation}
In particular, $\mathcal E=0$ as a map in a neighbourhood implies
$M=0$ there.
\end{lemma}

\begin{proof}
Using \eqref{eq:n9v76-defect-derivatives} and
$D^{-1}-L^{-1}=L^{-1}(L-D)D^{-1}$ gives
\begin{align*}
 M
 &=-L^{-1}(B^*+\mathcal E_y)+D^{-1}B^*\\
 &=L^{-1}\mathcal E_\varphi D^{-1}B^*
   -L^{-1}\mathcal E_y.
\end{align*}
\end{proof}

\begin{firewall}[A value defect and a slope defect are different]
\label{fire:n9v76-value-slope}
The equality $\mathcal E=0$ at one point does not imply $M=0$ there.
Formula \eqref{eq:n9v76-M} depends on
$\mathcal E_\varphi$ and $\mathcal E_y$.  Conversely $M=0$ does not
make the action stationary on the constraint branch; stationarity also
requires $\nabla_\varphi S=0$, which by
\eqref{eq:n9v76-constraint-branch} is equivalent to
$\mathcal E=0$ at that point.
\end{firewall}

\subsection{The graph square and the reaction row}

For $h\in Y$, the branch tangent is
\begin{equation}
 T_Fh=(R_Fh,h)\in X_d\oplus Y.
 \label{eq:n9v76-TF}
\end{equation}

\begin{theorem}[Exact graph-Hessian square]
\label{thm:n9v76-graph-square}
The action Hessian restricted to the constraint tangent graph is
\begin{equation}
 T_F^*D^2S\,T_F
 =S_0+M^*DM.
 \label{eq:n9v76-graph-square}
\end{equation}
Consequently, if $S_0>0$, every constraint slope has graph Hessian at
least $S_0$; equality holds on a vector $h$ precisely when $Mh=0$.
\end{theorem}

\begin{proof}
The Hessian block matrix is
\[
 \begin{pmatrix}D&B^*\\B&A\end{pmatrix}.
\]
Substituting $R_F=R_*+M=-D^{-1}B^*+M$ into
\[
 A+BR_F+R_F^*B^*+R_F^*DR_F
\]
cancels both terms linear in $M$ and leaves
$A-BD^{-1}B^*+M^*DM$.  Strict positivity of $D$ proves the last
statement.
\end{proof}

\begin{firewall}[The graph Hessian is not yet the pullback Hessian]
\label{fire:n9v76-graph-not-pullback}
When $\nabla_\varphi S\ne0$, the second derivative of
$S(\varphi(y),y)$ also contains the second fundamental-form reaction
$S_\varphi[D_y^2\varphi]$.  Omitting it would repeat the error excluded
by \Cref{fire:n9v74-constraint-stationary}.
\end{firewall}

Define the symmetric bilinear reaction form
\begin{equation}
 \mathcal R[h,k]
 :=S_\varphi[\varphi_{hk}],
 \qquad
 \varphi_{hk}:=D_y^2\varphi[h,k].
 \label{eq:n9v76-reaction}
\end{equation}

\begin{theorem}[Exact nonstationary reduced Hessian]
\label{thm:n9v76-reduced-Hessian}
The full reduced action $S_{\rm red}(y)=S(\varphi(y),y)$ satisfies
\begin{equation}
 D_y^2S_{\rm red}
 =S_0+M^*DM+\mathcal R.
 \label{eq:n9v76-full-reduced}
\end{equation}
On the constraint branch,
$S_\varphi=-\mathcal E$ and hence
\begin{equation}
 \mathcal R[h,k]
 =-\langle\mathcal E,\varphi_{hk}\rangle_{X_d}.
 \label{eq:n9v76-reaction-defect}
\end{equation}
\end{theorem}

\begin{proof}
The general chain rule
\eqref{eq:n9v74-Sred-second} is the graph Hessian plus
$S_\varphi[\varphi_{hk}]$.  Apply
\Cref{thm:n9v76-graph-square}.  Finally,
$F=0$ and \eqref{eq:n9v76-defect} give
$\nabla_\varphi S=-\mathcal E$, proving
\eqref{eq:n9v76-reaction-defect}.
\end{proof}

\begin{corollary}[Sharp relative reaction gate]
\label{cor:n9v76-reaction-gate}
Suppose $S_0>0$ and the self-adjoint operator representing
$\mathcal R$ satisfies
\begin{equation}
 \eta:=\lVert S_0^{-1/2}\mathcal R S_0^{-1/2}\rVert<1.
 \label{eq:n9v76-eta}
\end{equation}
Then
\begin{equation}
 D_y^2S_{\rm red}
 \ge (1-\eta)S_0+M^*DM
 \ge(1-\eta)S_0>0.
 \label{eq:n9v76-positive-reduced}
\end{equation}
\end{corollary}

\begin{proof}
The definition of $\eta$ gives
$\mathcal R\ge-\eta S_0$.  Insert this in
\eqref{eq:n9v76-full-reduced} and use $D>0$.
\end{proof}

\subsection{A computable bound for the reaction gate}

Differentiating the constraint twice gives
\begin{equation}
 \varphi_{hk}=-L^{-1}\mathcal N_F[h,k],
 \label{eq:n9v76-second-response}
\end{equation}
where
\begin{align}
 \mathcal N_F[h,k]
 :={}&F_{yy}[h,k]+F_{y\varphi}[h,R_Fk]
       +F_{\varphi y}[R_Fh,k]
       +F_{\varphi\varphi}[R_Fh,R_Fk].
 \label{eq:n9v76-NF}
\end{align}

\begin{proposition}[Product bound for the reaction debit]
\label{prop:n9v76-reaction-bound}
Assume constants $\sigma,\nu\ge0$ satisfy
\begin{align}
 \lVert \mathcal E^*L^{-1}\rVert&\le\sigma,
 \label{eq:n9v76-sigma}\\
 \lVert\mathcal N_F[h,h]\rVert_{X_d}
 &\le\nu\langle h,S_0h\rangle_Y
 \quad(h\in Y).
 \label{eq:n9v76-nu}
\end{align}
Then
\begin{equation}
 |\mathcal R[h,h]|
 \le\sigma\nu\langle h,S_0h\rangle_Y,
 \qquad
 \eta\le\sigma\nu.
 \label{eq:n9v76-reaction-product}
\end{equation}
In particular $\sigma\nu<1$ is a sufficient strict reduced-Hessian
gate.
\end{proposition}

\begin{proof}
Equations \eqref{eq:n9v76-reaction-defect} and
\eqref{eq:n9v76-second-response} give
\[
 \mathcal R[h,h]
 =\langle\mathcal E,L^{-1}\mathcal N_F[h,h]\rangle.
\]
Cauchy--Schwarz in the declared norms and
\eqref{eq:n9v76-sigma}--\eqref{eq:n9v76-nu} prove the diagonal bound.
For a symmetric bilinear form, the least constant in
$|\mathcal R[h,h]|\le c\langle h,S_0h\rangle$ equals
$\lVert S_0^{-1/2}\mathcal R S_0^{-1/2}\rVert$, proving the norm
claim.  Apply \Cref{cor:n9v76-reaction-gate}.
\end{proof}

\begin{corollary}[Stationary elimination is recovered exactly]
\label{cor:n9v76-stationary-recovered}
If $F=\nabla_\varphi S$ as maps on a neighbourhood, then
\begin{equation}
 \mathcal E=0,\qquad M=0,\qquad\mathcal R=0,
 \qquad D_y^2S_{\rm red}=S_0.
 \label{eq:n9v76-stationary-limit}
\end{equation}
\end{corollary}

\begin{proof}
All derivatives of $\mathcal E$ vanish.  Use
\Cref{lem:n9v76-tangent-defect,thm:n9v76-reduced-Hessian}.
\end{proof}

\subsection{Weighted locality of the slope defect}

For a block metric $d$ and $\mu>0$, write
\begin{equation}
 \lVert T\rVert_\mu
 :=\sup_x\sum_y e^{\mu d(x,y)}
 \lVert P_xTP_y\rVert.
 \label{eq:n9v76-weighted-norm}
\end{equation}
This norm is submultiplicative whenever the block products are
defined.

\begin{theorem}[Quasi-local tangent defect]
\label{thm:n9v76-M-locality}
Suppose the five weighted norms on the right below are finite.  Then
\begin{equation}
 \lVert M\rVert_\mu
 \le\lVert L^{-1}\rVert_\mu
 \left(
  \lVert\mathcal E_\varphi\rVert_\mu
  \lVert D^{-1}\rVert_\mu
  \lVert B^*\rVert_\mu
  +\lVert\mathcal E_y\rVert_\mu
 \right).
 \label{eq:n9v76-M-locality}
\end{equation}
Thus uniform gaps for $L$ and $D$, together with the finite-range
hypotheses of V65 and V74, make the constraint--stationarity slope
defect exponentially summable.
\end{theorem}

\begin{proof}
Take the weighted norm of the exact identity
\eqref{eq:n9v76-M} and use submultiplicativity.  Exponential decay of
the two inverses follows from the weighted-conjugation estimate of V65
under the stated gap and finite-range hypotheses.
\end{proof}

\begin{firewall}[A closing gap is still a retained mode]
\label{fire:n9v76-closing-gap}
Small local coefficients in $\mathcal E$ do not compensate for a
diverging $\lVert L^{-1}\rVert_\mu$ or
$\lVert D^{-1}\rVert_\mu$.  A closing conformal or stationary-detail
mode must be moved into the retained infrared sector, exactly as in the
V74 criticality firewall.
\end{firewall}

\subsection{Determinant ownership}

There are two distinct reductions in play.  Gaussian stationary
integration eliminates a freely integrated detail.  Constraint coarea
restricts the measure to a root of $F$.  Their determinant powers are
different.

\begin{proposition}[Stationary Gaussian determinant]
\label{prop:n9v76-Gaussian-determinant}
If the block Hessian
$H=\left(\begin{smallmatrix}D&B^*\\B&A\end{smallmatrix}\right)$ is
positive, then
\begin{equation}
 \det H=(\det D)(\det S_0),
 \label{eq:n9v76-block-determinant}
\end{equation}
and integration of the quadratic detail $\xi$ at fixed $y$ gives
\begin{equation}
 \int_{X_d}
 e^{-\frac1{2\hbar}\langle(\xi,y),H(\xi,y)\rangle}\,d\xi
 =(2\pi\hbar)^{\dim X_d/2}(\det D)^{-1/2}
 e^{-\frac1{2\hbar}\langle y,S_0y\rangle}.
 \label{eq:n9v76-Gaussian-integral}
\end{equation}
\end{proposition}

\begin{proof}
Use the triangular change of variables
$\xi'=\xi+D^{-1}B^*y$, which has unit Jacobian, and the square
completion \eqref{eq:n9v75-square-completion}.  The Gaussian integral
gives \eqref{eq:n9v76-Gaussian-integral}; comparison of normalizations
gives \eqref{eq:n9v76-block-determinant}.
\end{proof}

\begin{proposition}[Constraint determinant is a separate owner]
\label{prop:n9v76-constraint-determinant}
For an isolated regular root of $F(\varphi,y)=0$, constraint coarea
gives
\begin{equation}
 \int_{X_d}G(\varphi,y)\delta(F(\varphi,y))\,d\varphi
 =\frac{G(\varphi(y),y)}{|\det L|}.
 \label{eq:n9v76-coarea-again}
\end{equation}
The factor $|\det L|^{-1}$ cannot be replaced by
$(\det D)^{-1/2}$ unless an independent derivation proves an equality
of the two measures, including determinant powers and ghosts.
\end{proposition}

\begin{proof}
Equation \eqref{eq:n9v76-coarea-again} is the local delta-function
change of variables of \Cref{prop:n9v74-coarea}.  The Gaussian factor
in \eqref{eq:n9v76-Gaussian-integral} comes instead from integrating
over every detail value with a quadratic weight.  One operation uses a
delta Jacobian to a full inverse power and the other a Gaussian
normalization to a half inverse power; they are different operations.
\end{proof}

\begin{firewall}[No determinant double counting]
\label{fire:n9v76-determinant-double-count}
A construction must declare whether the conformal detail is integrated,
constrained, gauge-fixed with a Faddeev--Popov factor, or reconstructed
from retained data.  Multiplying both determinant factors merely
because both formulas are available double counts the eliminated row.
Any cancellation with ghosts or lapse integration must be derived in
the same regulated measure.
\end{firewall}

\subsection{Constraint--stationarity compatibility card}

\begin{definition}[CSCC]
\label{def:n9v76-CSCC}
The constraint--stationarity compatibility card requires:
\begin{enumerate}
\item a real $C^2$ constraint branch and an invertible $L=F_\varphi$;
\item a positive stationary detail Hessian $D=S_{\varphi\varphi}$;
\item a positive stationary Schur operator $S_0$;
\item the exact defect data
      $\mathcal E,\mathcal E_\varphi,\mathcal E_y$;
\item the slope identity \eqref{eq:n9v76-M} and a weighted-locality
      bound for $M$;
\item the exact branch curvature \eqref{eq:n9v76-second-response};
\item a reaction margin $\eta<1$, or the sufficient product gate
      $\sigma\nu<1$;
\item exactly one declared determinant/ghost ownership rule;
\item spatial reflection compatibility of the reconstructed branch;
      and
\item explicit retention of every closing $L$, $D$, or $S_0$ mode.
\end{enumerate}
\end{definition}

\begin{theorem}[CSCC gives a strict local physical Hessian]
\label{thm:n9v76-CSCC}
If CSCC holds uniformly over a finite-regulator chart, then the
constraint pullback has the exact Hessian
\eqref{eq:n9v76-full-reduced}, is strictly positive by
\eqref{eq:n9v76-positive-reduced}, and has an exponentially summable
tangent defect by \eqref{eq:n9v76-M-locality}.  The chosen determinant
owner is explicit and counted once.
\end{theorem}

\begin{proof}
Items 1--7 allow application of
\Cref{thm:n9v76-reduced-Hessian,cor:n9v76-reaction-gate,
thm:n9v76-M-locality}.  Items 8--10 impose the measure,
reflection-support, and critical-mode gates required by
\Cref{prop:n9v76-constraint-determinant,
thm:n9v74-reflection-pushforward,fire:n9v76-closing-gap}.
\end{proof}

\subsection{Effect of the YM audit on CSCC}

\begin{proposition}[Marginal source calibration precedes a small-defect claim]
\label{prop:n9v76-YM-before-defect}
Suppose the YM source contribution to $\mathcal E$ has the conditional
split \eqref{eq:n9v75-YM-source-split}.  A uniform smallness estimate
for $\mathcal E$, $\mathcal E_y$, or $\mathcal E_\varphi$ can be sought
only after the $\operatorname{Tr}F^2$ and
$\operatorname{Tr}(F\widetilde F)$ coefficients have been absorbed into
the calibrated running reference constraint.  After that calibration,
the $O(|k|^2)$ quadratic row is eligible for the stable defect norm,
conditional on the YM complete-norm gate.
\end{proposition}

\begin{proof}
The two marginal symbols do not shrink under four-dimensional
canonical dilation, so leaving them in $\mathcal E$ cannot yield a
cofinal small-defect estimate.  Calibration removes their declared
reference projection.  The residual symbol has two extra momenta by
\eqref{eq:n9v75-YM-source-split}; its conversion to the stable norm is
exactly the open YM PLC norm comparison and therefore remains a stated
hypothesis.
\end{proof}

\begin{theorem}[Ninth-series V76 result]
\label{thm:n9v76-result}
For a nonstationary constraint branch, the reduced Hessian has the
exact decomposition
\[
 D_y^2S_{\rm red}
 =S_0+M^*DM+\mathcal R,
 \qquad
 M=L^{-1}(\mathcal E_\varphi D^{-1}B^*-\mathcal E_y).
\]
The slope mismatch is positive in the graph Hessian.  All possible
new loss is isolated in the reaction
$\mathcal R=S_\varphi[D_y^2\varphi]$, which is strictly controlled by
either \eqref{eq:n9v76-eta} or the computable gate
$\sigma\nu<1$.  The Gaussian and coarea determinants are distinct
measure owners and cannot be multiplied without a regulated
derivation.
\end{theorem}

\begin{proof}
Use \Cref{lem:n9v76-tangent-defect,thm:n9v76-graph-square,
thm:n9v76-reduced-Hessian,prop:n9v76-reaction-bound,
prop:n9v76-Gaussian-determinant,
prop:n9v76-constraint-determinant}.
\end{proof}

\begin{assessment}[Progress at ninth-series V76]
\label{ass:n9v76-status}
The central V74 ambiguity has been converted into two measured objects:
the first-derivative slope defect $M$ and the second-derivative reaction
$\mathcal R$.  The former is a positive square and is quasi-local under
the two gap hypotheses.  The latter is the sole nonstationary debit and
has an explicit product bound.  The next bottleneck is upstream:
construct a physical regulated Einstein--SM constraint/action pair for
which the defect and reaction gates are uniform after all marginal
matter owners are calibrated.
\end{assessment}

\begin{programme}[Ninth-series V77: upstream source compatibility]
\label{prog:n9v76-V77}
V77 will go upstream to the source map.  It will formulate a
stress-energy Ward-defect identity relating the regulated constraint,
metric action gradient, gauge fixing, and anomaly rows.  The target is
to identify which parts of $\mathcal E$ vanish by exact diffeomorphism
covariance, which are calibrated marginal owners, and which may enter
the stable reaction bound.
\end{programme}

\begin{firewall}[Claim boundary after ninth-series V76]
\label{fire:n9v76-boundary}
V76 proves finite-regulator identities and sufficient gates.  It does
not prove that the physical Einstein constraint equals a Euclidean
action gradient, a uniform bound on $\mathcal E$ or $\mathcal R$, the
YM complete-norm gate, a global York branch, chiral transfer positivity,
anomaly cancellation, infrared closure, regulator independence, OS
reconstruction, or the full SM--Einstein continuum.
\end{firewall}

\paragraph{Parent-version record.}
V76 was formed from
\texttt{Renewal\_SM\_Einstein\_Continuum\_Research\_Notes\_V75.tex}.
The V75 parent had $28\,687$ counted source lines, $1\,036\,100$ bytes,
and SHA--256
\[
 \texttt{3D07FECED15E683D1D800E58D97A0379320A2446543F64CB0E3EC55164300868}.
\]

% =============================================================================
% END APPENDED NINTH-SERIES V76 MATERIAL
% =============================================================================

% =============================================================================
% BEGIN APPENDED NINTH-SERIES V77 MATERIAL
% =============================================================================

\section{Ninth-series V77: upstream Ward and source compatibility}
\label{sec:v9v77}

V76 reduced positivity of a nonstationary York pullback to the size of
the constraint--stationarity defect $\mathcal E$ and its derivatives.
The present note goes upstream and identifies what that defect compares.
The Hamiltonian constraint is a lapse derivative, whereas the
stationary lift uses a conformal derivative.  A map between those two
cotangent spaces is indispensable.  Diffeomorphism Ward identities
control gauge projections of the Euler derivatives; without a rank
hypothesis they do not turn the lapse equation into the conformal
stationarity equation.

\subsection{Typed lapse-to-conformal comparison}

At a fixed regulator let the field chart be
\begin{equation}
 \mathcal Q=N\oplus X_d\oplus Y,
 \qquad q=(n,\varphi,y),
 \label{eq:n9v77-field-chart}
\end{equation}
where $N$ is the declared scalar lapse/multiplier sector, $X_d$ is the
eliminated conformal sector, and $Y$ contains the remaining fields.  Let
$\Gamma:\mathcal Q\to\mathbb R$ be $C^3$.  Its scalar constraint is
\begin{equation}
 C(q):=\nabla_n\Gamma(q)\in N.
 \label{eq:n9v77-C}
\end{equation}
Let $U(q):N\to X_d$ be a $C^2$ isomorphism on the declared scalar
constraint sector and set
\begin{equation}
 F(q):=U(q)C(q),
 \qquad
 \mathcal E(q):=F(q)-\nabla_\varphi\Gamma(q).
 \label{eq:n9v77-F-E}
\end{equation}
The isomorphism hypothesis makes $F=0$ equivalent to $C=0$ in this
sector.  It is a finite-regulator hypothesis to be constructed, not a
formal identification of lapse and conformal indices.

\begin{proposition}[Exact source-defect derivatives]
\label{prop:n9v77-defect-derivatives}
For a variation $h\in T_q\mathcal Q$,
\begin{align}
 D\mathcal E[h]
 &=(DU[h])C+U\Gamma_{nh}-\Gamma_{\varphi h},
 \label{eq:n9v77-DE}\\
 D^2\mathcal E[h,k]
 &=(D^2U[h,k])C+(DU[h])\Gamma_{nk}
 +(DU[k])\Gamma_{nh}\\
 &\quad+U\Gamma_{nhk}-\Gamma_{\varphi hk}.
 \label{eq:n9v77-D2E}
\end{align}
On the constraint surface $C=0$ these reduce, in the $\varphi$ and
$y$ directions, to
\begin{align}
 \mathcal E&=-\nabla_\varphi\Gamma,
 \label{eq:n9v77-E-on-C}\\
 \mathcal E_\varphi&=U\Gamma_{n\varphi}-\Gamma_{\varphi\varphi},
 \label{eq:n9v77-Ephi-on-C}\\
 \mathcal E_y&=U\Gamma_{ny}-\Gamma_{\varphi y}.
 \label{eq:n9v77-Ey-on-C}
\end{align}
\end{proposition}

\begin{proof}
Apply the product rule once and twice to
$U\nabla_n\Gamma-
\nabla_\varphi\Gamma$.  On $C=0$ every undifferentiated factor $C$
vanishes, giving the last three formulas.
\end{proof}

\begin{corollary}[Exact compatibility conditions]
\label{cor:n9v77-compatibility}
At a constrained point, simultaneous vanishing of
$\mathcal E$, $\mathcal E_\varphi$, and $\mathcal E_y$ is equivalent
to
\begin{equation}
 \Gamma_\varphi=0,
 \qquad
 U\Gamma_{n\varphi}=\Gamma_{\varphi\varphi},
 \qquad
 U\Gamma_{ny}=\Gamma_{\varphi y}.
 \label{eq:n9v77-compatibility}
\end{equation}
These are a value equation and two mixed-Hessian equations; the lapse
constraint alone proves none of them.
\end{corollary}

\begin{proof}
Use \eqref{eq:n9v77-E-on-C}--\eqref{eq:n9v77-Ey-on-C}.
\end{proof}

\begin{firewall}[The map $U$ is mathematical data]
\label{fire:n9v77-U-data}
Multiplying a constraint by an unspecified symbol does not construct
$U$.  Its domain, range, zero modes, boundary conditions, covariance,
weighted locality, and regulator-uniform inverse bounds must be stated.
If $U$ has a kernel, $F=0$ can discard lapse constraints.  If its
inverse becomes singular, the discarded direction is a retained
critical owner.
\end{firewall}

\subsection{Reference compatibility and perturbation ownership}

Let $(\Gamma_0,U_0)$ be a reference pair satisfying the off-shell
identity
\begin{equation}
 U_0\nabla_n\Gamma_0=\nabla_\varphi\Gamma_0
 \label{eq:n9v77-reference-compatibility}
\end{equation}
on a declared chart.  Write
\begin{equation}
 \Gamma=\Gamma_0+I,
 \qquad U=U_0+V.
 \label{eq:n9v77-reference-split}
\end{equation}

\begin{lemma}[Exact perturbative defect identity]
\label{lem:n9v77-reference-defect}
The full defect is
\begin{equation}
 \mathcal E
 =U_0I_n-I_\varphi+V(\Gamma_{0,n}+I_n).
 \label{eq:n9v77-reference-defect}
\end{equation}
On the full constraint surface $C=\Gamma_n=0$, it becomes
\begin{equation}
 \mathcal E=U_0I_n-I_\varphi.
 \label{eq:n9v77-reference-on-shell}
\end{equation}
\end{lemma}

\begin{proof}
Expand $(U_0+V)(\Gamma_{0,n}+I_n)
-(\Gamma_{0,\varphi}+I_\varphi)$ and cancel
\eqref{eq:n9v77-reference-compatibility}.  On the constraint surface
the vector multiplying $V$ is zero.
\end{proof}

\begin{firewall}[A compatible reference pair is not automatic]
\label{fire:n9v77-reference-not-automatic}
Equation \eqref{eq:n9v77-reference-compatibility} is stronger than a
Ward identity and can fail even for a diffeomorphism-invariant action.
It must be proved for the chosen lapse parametrization, conformal
coordinate, boundary terms, and gauge fixing.  V77 uses it only as a
conditional compiler for separating calibrated reference owners from
the stable perturbation $I$.
\end{firewall}

\begin{proposition}[Marginal owners cannot remain in $I$]
\label{prop:n9v77-marginal-calibration}
Suppose $I=I_{\rm marg}+I_{\rm st}$ and a component of
$U_0I_{\rm marg,n}-I_{\rm marg,\varphi}$ is invariant under the
coarse dilation.  Then a cofinal estimate
$\lVert\mathcal E\rVert\to0$ is impossible unless that component is
zero or is absorbed into a running compatible reference pair.  After
that calibration, \eqref{eq:n9v77-reference-on-shell} assigns only
$U_0I_{{\rm st},n}-I_{{\rm st},\varphi}$ to the stable defect row.
\end{proposition}

\begin{proof}
An invariant component has the same nonzero norm at every rescaled
stage, contradicting convergence of the full defect to zero.  Moving it
into the reference pair removes it from $I$ but preserves the exact
identity \eqref{eq:n9v77-reference-defect}.  This is the calibrated
projection mechanism of V70.
\end{proof}

The YM V2 owners in \eqref{eq:n9v75-YM-source-split} are an immediate
instance: the $\operatorname{Tr}F^2$ and
$\operatorname{Tr}(F\widetilde F)$ coefficients belong in
$\Gamma_0$, while the two-momentum row may enter $I_{\rm st}$ only
after its complete-norm gate is proved.

\subsection{Finite-regulator Ward identity}

Let $\mathfrak g$ be a finite-dimensional regulated gauge-parameter
space and
\begin{equation}
 R(q):\mathfrak g\to T_q\mathcal Q
 \label{eq:n9v77-generator}
\end{equation}
the infinitesimal gauge/diffeomorphism generator.  In the splitting
\eqref{eq:n9v77-field-chart}, write
$R=(R_n,R_\varphi,R_y)$.  Define the Ward residual
\begin{equation}
 \mathscr W(q):=R(q)^*\nabla\Gamma(q)\in\mathfrak g.
 \label{eq:n9v77-Ward-residual}
\end{equation}

\begin{lemma}[Infinitesimal invariance]
\label{lem:n9v77-invariance}
If $\Gamma$ is invariant along every generated orbit, then
\begin{equation}
 \mathscr W=0.
 \label{eq:n9v77-Ward-zero}
\end{equation}
More generally, the exact component identity is
\begin{equation}
 R_n^*\Gamma_n+R_\varphi^*\Gamma_\varphi
 +R_y^*\Gamma_y=\mathscr W.
 \label{eq:n9v77-component-Ward}
\end{equation}
\end{lemma}

\begin{proof}
For $\xi\in\mathfrak g$, infinitesimal invariance gives
$D\Gamma[R\xi]=0$.  By the Hilbert adjoint this is
$\langle R^*\nabla\Gamma,\xi\rangle=0$ for every $\xi$.  The component
formula is the definition after splitting the gradient and generator.
\end{proof}

\begin{theorem}[What the Ward identity controls on the constraint surface]
\label{thm:n9v77-Ward-projection}
On $C=\Gamma_n=0$,
\begin{equation}
 R_\varphi^*\Gamma_\varphi
 =\mathscr W-R_y^*\Gamma_y.
 \label{eq:n9v77-Ward-on-C}
\end{equation}
If, in addition,
\begin{equation}
 \lVert R_\varphi^*v\rVert_{\mathfrak g}
 \ge r\lVert v\rVert_{X_d}
 \quad(v\in X_d)
 \label{eq:n9v77-rank-bound}
\end{equation}
for some $r>0$, then
\begin{equation}
 \lVert\mathcal E\rVert_{X_d}
 =\lVert\Gamma_\varphi\rVert_{X_d}
 \le r^{-1}\lVert\mathscr W-R_y^*\Gamma_y\rVert_{\mathfrak g}.
 \label{eq:n9v77-E-Ward-bound}
\end{equation}
Without \eqref{eq:n9v77-rank-bound}, the Ward identity controls only
the displayed gauge projection of $\Gamma_\varphi$.
\end{theorem}

\begin{proof}
Set $\Gamma_n=0$ in \eqref{eq:n9v77-component-Ward}.  Apply the lower
bound \eqref{eq:n9v77-rank-bound} to
$v=\Gamma_\varphi$ and then use
\eqref{eq:n9v77-E-on-C}.
\end{proof}

\begin{firewall}[Ward covariance does not imply conformal stationarity]
\label{fire:n9v77-Ward-not-stationarity}
The rank condition \eqref{eq:n9v77-rank-bound} says that the regulated
gauge generator is surjective onto the eliminated conformal sector.
This is a substantial geometric condition.  If a conformal direction
is not a gauge direction, $R_\varphi^*$ can have a kernel and
$\mathscr W=0$ leaves that component of $\Gamma_\varphi$ uncontrolled.
Thus diffeomorphism invariance alone does not prove
$\mathcal E=0$.
\end{firewall}

\subsection{Differentiated Ward identity and the gauge zero-mode no-go}

\begin{proposition}[Exact differentiated Ward identity]
\label{prop:n9v77-differentiated-Ward}
Let $H_\Gamma=D^2\Gamma$.  For every variation $h$,
\begin{equation}
 D\mathscr W[h]
 =(DR[h])^*\nabla\Gamma+R^*H_\Gamma h.
 \label{eq:n9v77-DW}
\end{equation}
At a stationary point $\nabla\Gamma=0$ this reduces to
\begin{equation}
 R^*H_\Gamma=D\mathscr W.
 \label{eq:n9v77-DW-stationary}
\end{equation}
\end{proposition}

\begin{proof}
Differentiate the product $R^*\nabla\Gamma$.
\end{proof}

\begin{corollary}[Unfixed invariant Hessians have gauge null vectors]
\label{cor:n9v77-gauge-null}
If $\mathscr W=0$ identically near a stationary point, then
\begin{equation}
 R^*H_\Gamma=0,
 \qquad H_\Gamma R=0.
 \label{eq:n9v77-gauge-null}
\end{equation}
Hence a nontrivial unreduced gauge-invariant Hessian cannot be strictly
positive on the full field space.
\end{corollary}

\begin{proof}
Equation \eqref{eq:n9v77-DW-stationary} gives the first identity.
Take adjoints and use self-adjointness of the Hessian for the second.
Any nonzero vector in the range of $R$ is a null vector.
\end{proof}

This exact no-go reinforces the V72--V76 acquisition order: positivity
must be stated on a gauge-reduced physical tangent space, or after a
declared gauge-fixing term has lifted the gauge normal directions.

\subsection{Gauge fixing, regulator breaking, and anomaly ownership}

Let $\chi:\mathcal Q\to Z$ be a $C^2$ gauge condition and
\begin{equation}
 \Gamma_{\rm gf}(q)
 =\frac1{2\alpha}\lVert\chi(q)\rVert_Z^2,
 \qquad\alpha>0.
 \label{eq:n9v77-gauge-fixing}
\end{equation}

\begin{proposition}[Exact gauge-fixing Ward row]
\label{prop:n9v77-gf-Ward}
The gauge-fixing contribution to the Ward residual is
\begin{equation}
 \mathscr W_{\rm gf}
 =\alpha^{-1}(D\chi\,R)^*\chi.
 \label{eq:n9v77-gf-Ward}
\end{equation}
It vanishes on the gauge slice $\chi=0$, but its derivative there is
\begin{equation}
 D\mathscr W_{\rm gf}[h]
 =\alpha^{-1}(D\chi\,R)^*D\chi[h].
 \label{eq:n9v77-Dgf-Ward}
\end{equation}
Thus a zero value Ward row can still contribute a nonzero positive
gauge-normal Hessian row.
\end{proposition}

\begin{proof}
$\nabla\Gamma_{\rm gf}=\alpha^{-1}D\chi^*\chi$, so applying $R^*$
gives \eqref{eq:n9v77-gf-Ward}.  Differentiate.  Every derivative
hitting $D\chi R$ is multiplied by $\chi$ and vanishes on the gauge
slice, leaving \eqref{eq:n9v77-Dgf-Ward}.
\end{proof}

At the effective finite-regulator level the Ward ledger has the form
\begin{equation}
 \mathscr W
 =\mathscr W_{\rm gf}+\mathscr W_{\rm reg}
  +\mathscr A+\mathscr W_{\rm bdry},
 \label{eq:n9v77-Ward-ledger}
\end{equation}
where the last three terms denote, respectively, explicit cutoff
breaking, the regulated measure/anomaly row, and boundary breaking.

\begin{definition}[Ward ownership rule]
\label{def:n9v77-Ward-ownership}
Every term in \eqref{eq:n9v77-Ward-ledger} must be assigned exactly one
of the following statuses:
\begin{enumerate}
\item an exact row retained in the reference action and constraint;
\item a BRST/ghost partner with an identity proved in the same
      regulator;
\item a calibrated marginal coefficient;
\item a stable row with a norm tending to zero; or
\item a genuine anomaly or boundary observable retained in the target
      theory.
\end{enumerate}
\end{definition}

\begin{firewall}[Anomaly cancellation must be regulator compatible]
\label{fire:n9v77-anomaly}
Formal cancellation of continuum representation traces does not by
itself set $\mathscr A=0$ in
\eqref{eq:n9v77-Ward-ledger}.  The regulated chiral determinant phase,
measure Jacobian, boundary sector, and limiting Ward identity must be
constructed together.  V77 makes no anomaly-cancellation claim.
\end{firewall}

\subsection{Weak diffeomorphism identity and trace recovery}

The conservation hypothesis used in V73 also has an exact Ward origin.
Let $M$ be a compact manifold without boundary, or take compactly
supported test vector fields away from the boundary.  Let
$E_g^{\mu\nu}$ and $E_\psi$ be the metric and matter Euler derivatives
of a differentiable action.

\begin{theorem}[Weak Bianchi--Ward identity]
\label{thm:n9v77-weak-Bianchi}
If the action is invariant under diffeomorphisms generated by every
smooth test vector field $\xi$, then
\begin{equation}
 \int_M E_g^{\mu\nu}(\nabla_\mu\xi_\nu+
 \nabla_\nu\xi_\mu)\,dV
 +\langle E_\psi,\mathcal L_\xi\psi\rangle=0.
 \label{eq:n9v77-weak-Ward}
\end{equation}
On the matter shell $E_\psi=0$, symmetry of $E_g$ implies
\begin{equation}
 \nabla_\mu E_g^{\mu\nu}=0
 \label{eq:n9v77-conservation}
\end{equation}
in the distributional sense.  With Ward residual
$\mathscr A(\xi)$, the right side of
\eqref{eq:n9v77-weak-Ward} is $\mathscr A(\xi)$ and the divergence has
the corresponding distributional defect.
\end{theorem}

\begin{proof}
The infinitesimal field variations are
$\delta g=\mathcal L_\xi g$ and
$\delta\psi=\mathcal L_\xi\psi$.  Diffeomorphism invariance makes the
first variation zero and gives \eqref{eq:n9v77-weak-Ward}.  On the
matter shell, symmetry yields twice
$\int E_g^{\mu\nu}\nabla_\mu\xi_\nu,dV$.  Distributional integration
by parts makes this
$-2\langle\nabla_\mu E_g^{\mu\nu},\xi_\nu\rangle$.  Vanishing for every
test $\xi$ proves \eqref{eq:n9v77-conservation}.  The residual version
is the same calculation with a nonzero right side.
\end{proof}

\begin{corollary}[Exact input required by V73]
\label{cor:n9v77-V73-input}
The conservation premise in the V73 trace-recovery theorem follows on
the matter shell only after the limiting diffeomorphism Ward residual
vanishes.  If a residual survives, V73 yields trace recovery only after
its divergence defect is included; it cannot be silently discarded.
\end{corollary}

\begin{proof}
Apply \Cref{thm:n9v77-weak-Bianchi} before invoking
\Cref{thm:n9v73-trace-recovery}.
\end{proof}

\subsection{Ward--source compatibility card}

\begin{definition}[WSCC]
\label{def:n9v77-WSCC}
The Ward--source compatibility card requires:
\begin{enumerate}
\item a typed lapse sector $N$, conformal sector $X_d$, and retained
      sector $Y$;
\item a covariant isomorphism $U:N\to X_d$ with uniform locality and
      inverse control;
\item the exact source identities
      \eqref{eq:n9v77-E-on-C}--\eqref{eq:n9v77-Ey-on-C};
\item either a proved compatible reference pair or explicit ownership
      of its failure;
\item a finite-regulator generator $R$ and Ward ledger
      \eqref{eq:n9v77-Ward-ledger};
\item a projection/rank statement before using Ward control on
      $\Gamma_\varphi$;
\item differentiated gauge-fixing, regulator, anomaly, ghost, and
      boundary rows;
\item calibration of all noncontracting source owners, including the
      YM marginal pair;
\item a stable complete-norm bound on the remaining value and
      first/second derivative defects; and
\item a vanishing limiting diffeomorphism residual sufficient for the
      V73 conservation premise.
\end{enumerate}
\end{definition}

\begin{theorem}[WSCC supplies the inputs to the V76 reaction gate]
\label{thm:n9v77-WSCC}
If WSCC holds uniformly and the matter Euler equation is included in
the retained stationary system, then the V76 quantities
$\mathcal E$, $\mathcal E_\varphi$, and $\mathcal E_y$ are typed,
owned, and bounded in the same regulated norm.  The Ward identity
controls the conformal value defect only on the range certified by the
rank statement.  Vanishing of the limiting Ward residual supplies the
distributional conservation input needed for trace recovery.
\end{theorem}

\begin{proof}
Items 1--4 give the exact formulas of
\Cref{prop:n9v77-defect-derivatives,lem:n9v77-reference-defect}.
Items 5--7 and \Cref{thm:n9v77-Ward-projection,
prop:n9v77-gf-Ward} determine which gauge projections are controlled.
Items 8--9 place the remaining terms in the calibrated stable norm used
by V76.  Item 10 and \Cref{thm:n9v77-weak-Bianchi} give conservation.
\end{proof}

\begin{theorem}[Ninth-series V77 result]
\label{thm:n9v77-result}
For a lapse constraint $C=\Gamma_n$ and an explicit isomorphism $U$ to
the conformal sector, the V76 defect is exactly
$\mathcal E=UC-\Gamma_\varphi$.  On $C=0$, its first derivatives are
$U\Gamma_{n\varphi}-\Gamma_{\varphi\varphi}$ and
$U\Gamma_{ny}-\Gamma_{\varphi y}$.  A Ward identity bounds only the
gauge projection $R_\varphi^*\Gamma_\varphi$ unless the quantitative
rank condition \eqref{eq:n9v77-rank-bound} holds.  At stationary points
an invariant unfixed Hessian has exact gauge null vectors, while gauge
fixing contributes the differentiated row
\eqref{eq:n9v77-Dgf-Ward}.  Diffeomorphism conservation needed for V73
follows only after the matter equation and limiting Ward residual are
controlled.
\end{theorem}

\begin{proof}
Combine \Cref{prop:n9v77-defect-derivatives,
thm:n9v77-Ward-projection,cor:n9v77-gauge-null,
prop:n9v77-gf-Ward,thm:n9v77-weak-Bianchi}.
\end{proof}

\begin{assessment}[Progress at ninth-series V77]
\label{ass:n9v77-status}
The upstream obstruction is now explicit.  Constraint--stationarity
compatibility requires a constructed lapse-to-conformal map and two
mixed-Hessian identities; covariance alone supplies only projected Ward
control.  The regulator, gauge-fixing, anomaly, and boundary terms have
distinct derivative ledgers.  This prevents a false proof of a small
reaction debit from a formal Bianchi identity and identifies the exact
data a regulated SM--Einstein source construction must deliver.
\end{assessment}

\begin{programme}[Ninth-series V78: projected physical quotient]
\label{prog:n9v77-V78}
V78 will construct the finite-regulator physical quotient associated
with $R$, prove the Moore--Penrose projector formulas under a uniform
Faddeev--Popov gap, and express the Ward-controlled and transverse
components of $\mathcal E$ separately.  This will replace the strong
surjectivity condition \eqref{eq:n9v77-rank-bound} by an exact
gauge/physical decomposition and expose the genuinely physical defect
that must be estimated.
\end{programme}

\begin{firewall}[Claim boundary after ninth-series V77]
\label{fire:n9v77-boundary}
V77 does not construct the physical lapse-to-conformal map $U$, prove a
compatible reference pair, close the YM stable norm, establish the
finite-regulator chiral Ward identity or anomaly cancellation, prove a
uniform Faddeev--Popov gap, bound the physical transverse defect, close
the infrared limit, prove regulator independence or OS reconstruction,
or construct the full SM--Einstein continuum.
\end{firewall}

\paragraph{Parent-version record.}
V77 was formed from
\texttt{Renewal\_SM\_Einstein\_Continuum\_Research\_Notes\_V76.tex}.
The V76 parent had $29\,210$ counted source lines, $1\,053\,839$ bytes,
and SHA--256
\[
 \texttt{730482CD8F89352F7EC9E4DB3E68ACD5722553C0A85B89D967EE1F339AC5C7DC}.
\]

% =============================================================================
% END APPENDED NINTH-SERIES V77 MATERIAL
% =============================================================================

% =============================================================================
% BEGIN APPENDED NINTH-SERIES V78 MATERIAL
% =============================================================================

\section{Ninth-series V78: the projected physical defect}
\label{sec:v9v78}

The rank condition in \eqref{eq:n9v77-rank-bound} is unnecessarily
strong when the conformal chart contains both gauge and physical
directions.  This note replaces it by an exact orthogonal quotient.  A
Ward identity determines the gauge component of the V76 defect.  The
orthogonal component is a genuinely physical source defect and must be
estimated by the dynamics.  Moving-projector terms are included in the
derivative ledger.

\subsection{Gauge Gram operator and Moore--Penrose quotient}

Let $\mathfrak g$ and $X_d$ be finite-dimensional real Hilbert spaces,
and let
\begin{equation}
 A:\mathfrak g\to X_d
 \label{eq:n9v78-A}
\end{equation}
be the conformal component $R_\varphi$ of the regulated gauge
generator at a fixed field.  Stabilizers form $\ker A$.  Put
\begin{equation}
 \mathfrak g_0:=(\ker A)^\perp,
 \qquad A_0:=A|_{\mathfrak g_0},
 \qquad G:=A_0^*A_0.
 \label{eq:n9v78-G}
\end{equation}

\begin{assumption}[Uniform gauge Gram gap]
\label{ass:n9v78-Gram-gap}
There is $m>0$ such that
\begin{equation}
 G\ge m^2 I_{\mathfrak g_0}.
 \label{eq:n9v78-Gram-gap}
\end{equation}
Equivalently,
$\lVert A_0\xi\rVert\ge m\lVert\xi\rVert$ on the nonstabilizer
parameter space.
\end{assumption}

\begin{theorem}[Exact gauge and physical projectors]
\label{thm:n9v78-projectors}
Under \Cref{ass:n9v78-Gram-gap}, define
\begin{equation}
 \Pi_g:=A_0G^{-1}A_0^*,
 \qquad
 \Pi_p:=I_{X_d}-\Pi_g.
 \label{eq:n9v78-projectors}
\end{equation}
Then
\begin{align}
 &\Pi_g^2=\Pi_g=\Pi_g^*,
 &&\operatorname{ran}\Pi_g=\operatorname{ran}A,
 \label{eq:n9v78-Pg}\\
 &\Pi_p^2=\Pi_p=\Pi_p^*,
 &&\operatorname{ran}\Pi_p=\ker A^*.
 \label{eq:n9v78-Pp}
\end{align}
Thus
\begin{equation}
 X_d=\operatorname{ran}A\mathbin{\mathop\oplus^{\perp}}\ker A^*.
 \label{eq:n9v78-quotient-split}
\end{equation}
The operator $A_0G^{-1}A_0^*$ is the Moore--Penrose range projector and
is independent of the choice of orthonormal coordinates on
$\mathfrak g_0$.
\end{theorem}

\begin{proof}
$G>0$ by \eqref{eq:n9v78-Gram-gap}.  Direct multiplication gives
\[
 \Pi_g^2=A_0G^{-1}(A_0^*A_0)G^{-1}A_0^*=\Pi_g,
\]
and self-adjointness is immediate.  Its range is contained in
$\operatorname{ran}A_0=\operatorname{ran}A$, while
$\Pi_gA_0=A_0$ proves the reverse inclusion.  Orthogonal-complement
duality gives
$(\operatorname{ran}A)^perp=\ker A^*$ and proves the claims for
$\Pi_p$.  The range and orthogonality characterization is intrinsic,
so it is unchanged by an orthogonal reparametrization of
$\mathfrak g_0$.
\end{proof}

\begin{firewall}[Stabilizers are not inverted]
\label{fire:n9v78-stabilizers}
The inverse in \eqref{eq:n9v78-projectors} acts only on
$\mathfrak g_0$.  Global gauge stabilizers, reducible connections, and
Killing directions lie in $\ker A$ and cannot be assigned an artificial
inverse.  They require a separate group-volume, orbit-type, or retained
zero-mode treatment.
\end{firewall}

\subsection{Exact Ward reconstruction of the gauge defect}

Return to the V77 field chart.  Set
\begin{equation}
 b:=\mathscr W-R_n^*\Gamma_n-R_y^*\Gamma_y.
 \label{eq:n9v78-b}
\end{equation}
The component Ward identity
\eqref{eq:n9v77-component-Ward} is
\begin{equation}
 A^*\Gamma_\varphi=b.
 \label{eq:n9v78-Astar-gradient}
\end{equation}

\begin{theorem}[Gauge/physical source-defect split]
\label{thm:n9v78-defect-split}
On the constraint surface $\Gamma_n=0$, define
\begin{equation}
 \mathcal E_g:=\Pi_g\mathcal E,
 \qquad
 \mathcal E_p:=\Pi_p\mathcal E.
 \label{eq:n9v78-Egp}
\end{equation}
Then
\begin{align}
 \mathcal E_g
 &=-A_0G^{-1}P_0b,
 \label{eq:n9v78-Eg}\\
 \mathcal E_p
 &=-\Pi_p\Gamma_\varphi,
 \label{eq:n9v78-Ep}\\
 \mathcal E&=\mathcal E_g+\mathcal E_p,
 \qquad
 \lVert\mathcal E\rVert^2
 =\lVert\mathcal E_g\rVert^2+\lVert\mathcal E_p\rVert^2.
 \label{eq:n9v78-E-Pythagoras}
\end{align}
Here $P_0:\mathfrak g\to\mathfrak g_0$ is the orthogonal projection.
Moreover,
\begin{equation}
 \lVert\mathcal E_g\rVert
 \le m^{-1}\lVert P_0b\rVert.
 \label{eq:n9v78-Eg-bound}
\end{equation}
No Ward identity determines $\mathcal E_p$.
\end{theorem}

\begin{proof}
On the constraint surface,
$\mathcal E=-\Gamma_\varphi$ by
\eqref{eq:n9v77-E-on-C}.  Since $A$ vanishes on its stabilizer space,
$A^*\Gamma_\varphi=P_0b$.  Apply $A_0G^{-1}$ to
\eqref{eq:n9v78-Astar-gradient} and use
\eqref{eq:n9v78-projectors}; this proves
\eqref{eq:n9v78-Eg}.  The physical formula and Pythagoras identity
follow from \Cref{thm:n9v78-projectors}.  The singular values of
$A_0G^{-1}$ are the reciprocals of the singular values of $A_0$, so
\eqref{eq:n9v78-Gram-gap} gives
$\lVert A_0G^{-1}\rVert\le m^{-1}$.
\end{proof}

\begin{corollary}[Replacement for the V77 surjectivity gate]
\label{cor:n9v78-replace-rank}
If $\mathcal E_p=0$, then the Ward residual and retained Euler
derivative give the full estimate
\begin{equation}
 \lVert\mathcal E\rVert
 \le m^{-1}\lVert P_0(\mathscr W-R_y^*\Gamma_y)\rVert
 \label{eq:n9v78-full-Ward-bound}
\end{equation}
on the constraint surface.  In general the exact bound is
\begin{equation}
 \lVert\mathcal E\rVert^2
 \le m^{-2}\lVert P_0(\mathscr W-R_y^*\Gamma_y)\rVert^2
 +\lVert\mathcal E_p\rVert^2.
 \label{eq:n9v78-Ward-physical-bound}
\end{equation}
\end{corollary}

\begin{proof}
Use \eqref{eq:n9v78-E-Pythagoras} and
\eqref{eq:n9v78-Eg-bound}, with $\Gamma_n=0$.
\end{proof}

\begin{firewall}[The physical defect is the real dynamical input]
\label{fire:n9v78-physical-defect}
Setting the Ward residual, matter Euler derivative, and gauge defect to
zero still leaves $\mathcal E_p=-\Pi_p\Gamma_\varphi$.  This is the
component normal to the gauge orbit.  It must vanish by a physical
equation or be bounded in the calibrated stable norm; covariance cannot
remove it.
\end{firewall}

\subsection{Moving-projector derivatives}

Suppose $A=A(q)$ has locally constant rank, so the stabilizer dimension
does not jump in the chart.  Choose a smooth local trivialization of
$\mathfrak g_0$ and write a dot for a directional derivative.

\begin{lemma}[Derivative of the range projector]
\label{lem:n9v78-DPi}
The derivative of $\Pi_g=A_0G^{-1}A_0^*$ is
\begin{equation}
 \dot\Pi_g
 =\Pi_p\dot A_0G^{-1}A_0^*
  +A_0G^{-1}\dot A_0^*\Pi_p.
 \label{eq:n9v78-DPi}
\end{equation}
Consequently
\begin{equation}
 \lVert\dot\Pi_g\rVert
 =\lVert\dot\Pi_p\rVert
 \le 2m^{-1}\lVert\dot A_0\rVert.
 \label{eq:n9v78-DPi-bound}
\end{equation}
\end{lemma}

\begin{proof}
Differentiate $G^{-1}$ using
$\dot G^{-1}=-G^{-1}\dot G G^{-1}$ and
$\dot G=\dot A_0^*A_0+A_0^*\dot A_0$.  Substitution into the product
rule for $A_0G^{-1}A_0^*$ and collection of terms gives
\eqref{eq:n9v78-DPi}.  Now
$\lVert A_0G^{-1}\rVert=
\lVert G^{-1}A_0^*\rVert\le m^{-1}$ and both projectors have norm one,
which proves the bound.
\end{proof}

\begin{proposition}[Exact derivative ledger for the physical defect]
\label{prop:n9v78-DEp}
For any variation $h$,
\begin{align}
 D\mathcal E_p[h]
 &=(D\Pi_p[h])\mathcal E+\Pi_pD\mathcal E[h],
 \label{eq:n9v78-DEp}\\
 D^2\mathcal E_p[h,k]
 &=(D^2\Pi_p[h,k])\mathcal E
 +(D\Pi_p[h])D\mathcal E[k]\\
 &\quad +(D\Pi_p[k])D\mathcal E[h]
 +\Pi_pD^2\mathcal E[h,k].
 \label{eq:n9v78-D2Ep}
\end{align}
The same formulas hold with $p$ replaced by $g$.
\end{proposition}

\begin{proof}
Differentiate the product $\Pi_p\mathcal E$ once and twice.
\end{proof}

\begin{corollary}[First physical-defect bound]
\label{cor:n9v78-DEp-bound}
Under \Cref{ass:n9v78-Gram-gap},
\begin{equation}
 \lVert D\mathcal E_p[h]\rVert
 \le 2m^{-1}\lVert DA_0[h]\rVert\lVert\mathcal E\rVert
 +\lVert D\mathcal E[h]\rVert.
 \label{eq:n9v78-DEp-bound}
\end{equation}
Thus pointwise smallness of $\mathcal E_p$ alone is insufficient for
the V76 slope and reaction gates; motion of the quotient must also be
controlled.
\end{corollary}

\begin{proof}
Apply \eqref{eq:n9v78-DPi-bound} to
\eqref{eq:n9v78-DEp} and use $\lVert\Pi_p\rVert=1$.
\end{proof}

\begin{proposition}[Second projector bound]
\label{prop:n9v78-D2Pi}
On a constant-rank chart satisfying
\eqref{eq:n9v78-Gram-gap}, there is a universal polynomial $P$ such
that
\begin{equation}
 \lVert D^2\Pi_p[h,k]\rVert
 \le P\bigl(m^{-1},\lVert A_0\rVert,
 \lVert DA_0\rVert,\lVert D^2A_0\rVert\bigr)
 \lVert h\rVert\lVert k\rVert.
 \label{eq:n9v78-D2Pi-bound}
\end{equation}
One may take $P$ to be the sum of the operator norms of the finitely
many terms obtained by differentiating
\eqref{eq:n9v78-DPi}; every term contains at most two factors of
$G^{-1}$ after using the collected form.
\end{proposition}

\begin{proof}
Differentiate \eqref{eq:n9v78-DPi}.  Replace each derivative of
$G^{-1}$ by $-G^{-1}(DG)G^{-1}$ and each derivative of $G$ by the
product rule applied to $A_0^*A_0$.  The gap gives
$\lVert G^{-1}\rVert\le m^{-2}$ and
$\lVert A_0G^{-1}\rVert,
\lVert G^{-1}A_0^*\rVert\le m^{-1}$.  Submultiplicativity bounds every
resulting term by the stated polynomial.  There are finitely many terms
and their coefficients are universal.
\end{proof}

\begin{firewall}[Rank-changing strata]
\label{fire:n9v78-rank-change}
At a rank change the projector need not be differentiable and the
constants in \eqref{eq:n9v78-DPi-bound}--
\eqref{eq:n9v78-D2Pi-bound} can diverge.  Orbit-type changes must be
covered by separate strata with compatible boundary ownership; a
single smooth quotient chart cannot cross them silently.
\end{firewall}

\subsection{Weighted locality of the quotient}

Let the fields and gauge parameters carry compatible block metrics and
use the weighted norm \eqref{eq:n9v76-weighted-norm}.

\begin{theorem}[Quasi-local physical projector]
\label{thm:n9v78-projector-locality}
Suppose $A_0$ and $A_0^*$ have finite weighted norms and the Gram
inverse has a finite weighted norm.  Then
\begin{align}
 \lVert\Pi_g\rVert_\mu
 &\le\lVert A_0\rVert_\mu
 \lVert G^{-1}\rVert_\mu
 \lVert A_0^*\rVert_\mu,
 \label{eq:n9v78-Pg-local}\\
 \lVert\Pi_p\rVert_\mu
 &\le1+\lVert\Pi_g\rVert_\mu.
 \label{eq:n9v78-Pp-local}
\end{align}
If $G$ is finite range and has a regulator-uniform spectral gap, the
V65 weighted inverse argument supplies such a bound at a sufficiently
small positive $\mu$.
\end{theorem}

\begin{proof}
Use \eqref{eq:n9v78-projectors} and submultiplicativity of the weighted
norm.  The last statement is the gapped finite-range inverse theorem
already used in V65, V74, and V76.
\end{proof}

\begin{firewall}[Gribov and infrared gates]
\label{fire:n9v78-Gribov}
A local Gram gap is not a global gauge slice.  Approaching a Gribov
horizon, reducible orbit, or infrared gauge zero mode can close the gap
and destroy uniform locality.  The quotient is therefore a chartwise
compiler; global orbit ownership remains a separate gate.
\end{firewall}

\subsection{Insertion into the V76 reaction gate}

Let $\ell_L:=\lVert L^{-1}\rVert$ in the norm used for
\eqref{eq:n9v76-sigma}.  Orthogonality gives the elementary estimate
\begin{equation}
 \lVert\mathcal E^*L^{-1}\rVert
 \le\ell_L\bigl(
 \lVert\mathcal E_g\rVert+\lVert\mathcal E_p\rVert
 \bigr).
 \label{eq:n9v78-sigma-split}
\end{equation}

\begin{proposition}[Ward/physical reaction bound]
\label{prop:n9v78-reaction-bound}
On the constraint surface, suppose
\eqref{eq:n9v76-nu} holds and put
$b_0=P_0(\mathscr W-R_y^*\Gamma_y)$.  Then the V76 reaction constant
obeys
\begin{equation}
 \eta
 \le\nu\ell_L
 \left(m^{-1}\lVert b_0\rVert+
 \lVert\mathcal E_p\rVert\right).
 \label{eq:n9v78-eta-split}
\end{equation}
Consequently
\begin{equation}
 \nu\ell_L
 \left(m^{-1}\lVert b_0\rVert+
 \lVert\mathcal E_p\rVert\right)<1
 \label{eq:n9v78-reaction-gate}
\end{equation}
is a sufficient strict reaction gate.
\end{proposition}

\begin{proof}
Use \eqref{eq:n9v78-Eg-bound} in
\eqref{eq:n9v78-sigma-split}, and then apply
\Cref{prop:n9v76-reaction-bound}.
\end{proof}

This formula identifies the acquisition order.  Gauge-fixing,
regulator, anomaly, boundary, and retained Euler rows enter $b_0$.
The transverse term $\mathcal E_p$ requires a physical compatibility
estimate.  Both are amplified by the conformal inverse and branch
curvature factors $\ell_L$ and $\nu$.

\subsection{Physical quotient card}

\begin{definition}[PQC]
\label{def:n9v78-PQC}
The physical quotient card requires:
\begin{enumerate}
\item a constant-rank orbit-type chart and explicit stabilizer sector;
\item a uniform nonstabilizer Gram gap $G\ge m^2$;
\item the exact projectors \eqref{eq:n9v78-projectors};
\item weighted locality of $G^{-1}$ and both projectors;
\item the Ward reconstruction \eqref{eq:n9v78-Eg};
\item a separately owned physical defect $\mathcal E_p$;
\item first and second moving-projector derivative rows;
\item the gauge/regulator/anomaly/boundary ledger for $b_0$;
\item the strict combined reaction gate
      \eqref{eq:n9v78-reaction-gate}; and
\item compatible gluing or retained boundary data at orbit-type and
      Gribov strata.
\end{enumerate}
\end{definition}

\begin{theorem}[PQC gives a type-correct V76 input]
\label{thm:n9v78-PQC}
Under PQC, the V76 source defect has an orthogonal decomposition into a
Ward-controlled gauge part and an explicitly retained physical part.
Both components and their moving-projector derivatives are bounded in
the declared local norm.  If \eqref{eq:n9v78-reaction-gate} holds, the
nonstationary reduced Hessian remains strictly positive by
\eqref{eq:n9v76-positive-reduced}.
\end{theorem}

\begin{proof}
Use \Cref{thm:n9v78-defect-split,lem:n9v78-DPi,
prop:n9v78-DEp,prop:n9v78-D2Pi,
thm:n9v78-projector-locality,prop:n9v78-reaction-bound}, followed by
\Cref{cor:n9v76-reaction-gate}.
\end{proof}

\begin{theorem}[Ninth-series V78 result]
\label{thm:n9v78-result}
After removing stabilizers, the Gram inverse
$G^{-1}=(A_0^*A_0)^{-1}$ gives exact orthogonal gauge and physical
projectors.  On the lapse-constraint surface the gauge part of the V76
defect is
\[
 \mathcal E_g=-A_0G^{-1}P_0(
 \mathscr W-R_y^*\Gamma_y),
\]
with norm at most $m^{-1}$ times its Ward source, whereas
$\mathcal E_p=-\Pi_p\Gamma_\varphi$ is a genuinely physical input.
The projector derivative is the off-diagonal formula
\eqref{eq:n9v78-DPi}; hence quotient motion is counted in the first and
second defect derivatives.  The combined sufficient reaction gate is
\eqref{eq:n9v78-reaction-gate}.
\end{theorem}

\begin{proof}
Combine \Cref{thm:n9v78-projectors,thm:n9v78-defect-split,
lem:n9v78-DPi,prop:n9v78-DEp,prop:n9v78-reaction-bound}.
\end{proof}

\begin{assessment}[Progress at ninth-series V78]
\label{ass:n9v78-status}
The strong V77 surjectivity assumption has been replaced by a chartwise
physical quotient.  Ward control now enters only where it is valid, and
the remaining physical conformal defect is exposed rather than hidden.
The quotient is local under a Gram gap and its motion has an exact
derivative ledger.  The next upstream question is how the physical
defect is constrained by the Hamiltonian propagation algebra and the
retained matter equations.
\end{assessment}

\begin{programme}[Ninth-series V79: constraint propagation]
\label{prog:n9v78-V79}
V79 will derive a finite-regulator constraint-propagation/Duhamel
identity.  It will separate homogeneous propagation of an initial
constraint defect from forcing by Ward, regulator, boundary, and matter
Euler residuals, and will test which physical component of
$\mathcal E_p$ can be controlled without assuming conformal
stationarity.
\end{programme}

\begin{firewall}[Claim boundary after ninth-series V78]
\label{fire:n9v78-boundary}
V78 does not prove a uniform Gram gap, a global gauge slice, Gribov
closure, smallness of $\mathcal E_p$, the YM complete-norm gate,
constraint propagation, chiral anomaly cancellation, infrared closure,
regulator independence, OS reconstruction, or the full SM--Einstein
continuum.
\end{firewall}

\paragraph{Parent-version record.}
V78 was formed from
\texttt{Renewal\_SM\_Einstein\_Continuum\_Research\_Notes\_V77.tex}.
The V77 parent had $29\,754$ counted source lines, $1\,073\,901$ bytes,
and SHA--256
\[
 \texttt{1279269C51D77FB7917DD68A4C1FE8D67840024135E26621949C8F75A46D5410}.
\]

% =============================================================================
% END APPENDED NINTH-SERIES V78 MATERIAL
% =============================================================================

% =============================================================================
% BEGIN APPENDED NINTH-SERIES V79 MATERIAL
% =============================================================================

\section{Ninth-series V79: constraint propagation and physical observability}
\label{sec:v9v79}

V78 isolated the physical conformal defect
$\mathcal E_p=-\Pi_p\Gamma_\varphi$, which is invisible to the Ward
projection.  Constraint propagation is a possible upstream control
mechanism, but only if this defect enters the propagation equation
through an injective observation operator.  The present note proves the
exact Duhamel compiler and the corresponding observability criterion.
It also proves that propagation of $C=0$ alone cannot control an
unobserved component of $\mathcal E_p$.

\subsection{Off-shell propagation identity}

Let $\mathcal Q$ be a finite-regulator field space, let
$C:\mathcal Q\to\mathcal H_C$ be the regulated constraint map, and
let $V:\mathcal Q\to T\mathcal Q$ be the declared evolution vector
field.  Let
\begin{equation}
 e_p(q):=\Pi_p(q)\nabla_\varphi\Gamma(q)in\mathcal H_p
 \label{eq:n9v79-ep}
\end{equation}
be the physical conformal Euler residual from V78, and let
$e_r(q)\in\mathcal H_r$ collect the other retained Euler residuals.

\begin{definition}[Constraint propagation identity]
\label{def:n9v79-CPI}
A regulated constraint propagation identity is an off-shell identity
\begin{equation}
 DC(q)[V(q)]
 =K(q)C(q)+B_p(q)e_p(q)+B_r(q)e_r(q)+w(q),
 \label{eq:n9v79-CPI}
\end{equation}
where $K$ acts on $\mathcal H_C$, $B_p$ and $B_r$ are typed source
maps, and $w$ contains the explicit Ward, cutoff, anomaly, gauge-fixing,
and boundary breaking rows not already included in the Euler residuals.
\end{definition}

\begin{firewall}[Propagation is an identity to be proved]
\label{fire:n9v79-CPI-not-formal}
Equation \eqref{eq:n9v79-CPI} does not follow from naming $C$ a
constraint.  It must be derived from the regulated Bianchi/BRST
identity, including lapse and shift conventions, boundary terms,
gauge-fixing contributions, and every regulator breaking.  V79 proves
the consequences of a typed identity; it does not assert that the
physical regulator already supplies it.
\end{firewall}

Let a differentiable path $q(t)$ obey the inexact evolution equation
\begin{equation}
 \dot q(t)=V(q(t))+r_V(t),
 \label{eq:n9v79-inexact-flow}
\end{equation}
and put $c(t)=C(q(t))$.

\begin{lemma}[Exact propagated forcing]
\label{lem:n9v79-forcing}
If \eqref{eq:n9v79-CPI} holds along the path, then
\begin{equation}
 \dot c(t)=K(t)c(t)+f(t),
 \label{eq:n9v79-linear-C}
\end{equation}
with
\begin{equation}
 f(t)=B_p(t)e_p(t)+B_r(t)e_r(t)+w(t)
      +DC(q(t))[r_V(t)].
 \label{eq:n9v79-f}
\end{equation}
\end{lemma}

\begin{proof}
The chain rule gives
$\dot c=DC[V]+DC[r_V]$.  Substitute
\eqref{eq:n9v79-CPI}.
\end{proof}

\subsection{Duhamel propagation and sharp stability bounds}

Let $U(t,s)$ be the fundamental solution of
\begin{equation}
 \partial_tU(t,s)=K(t)U(t,s),
 \qquad U(s,s)=I.
 \label{eq:n9v79-fundamental}
\end{equation}

\begin{theorem}[Exact constraint Duhamel formula]
\label{thm:n9v79-Duhamel}
For $t\ge t_0$, every solution of
\eqref{eq:n9v79-linear-C} satisfies
\begin{equation}
 c(t)=U(t,t_0)c(t_0)+\int_{t_0}^tU(t,s)f(s)\,ds.
 \label{eq:n9v79-Duhamel}
\end{equation}
Conversely the right side is the unique solution with the prescribed
initial constraint.
\end{theorem}

\begin{proof}
Differentiate $U(t,s)^{-1}c(t)$ with respect to $t$, or directly
differentiate the right side using
\eqref{eq:n9v79-fundamental}.  It solves
\eqref{eq:n9v79-linear-C} and has value $c(t_0)$ at $t_0$.
Finite-dimensional linear ODE uniqueness proves the converse.
\end{proof}

\begin{proposition}[Logarithmic-norm bound]
\label{prop:n9v79-log-bound}
Suppose the symmetric part of $K$ obeys
\begin{equation}
 \frac12(K(t)+K(t)^*)\le\omega(t)I.
 \label{eq:n9v79-log-hyp}
\end{equation}
Then
\begin{equation}
 \lVert U(t,s)\rVert
 \le\exp\left(\int_s^t\omega(\tau)\,d\tau\right),
 \label{eq:n9v79-U-bound}
\end{equation}
and
\begin{align}
 \lVert c(t)\rVert
 &\le e^{\int_{t_0}^t\omega}
       \lVert c(t_0)\rVert\\
 &\quad+\int_{t_0}^t
 e^{\int_s^t\omega}\lVert f(s)\rVert\,ds.
 \label{eq:n9v79-c-bound}
\end{align}
\end{proposition}

\begin{proof}
For a homogeneous solution $u$, differentiate its squared norm:
\[
 \frac12\frac d{dt}\lVert u\rVert^2
 =\left\langle u,\frac{K+K^*}{2}u\right\rangle
 \le\omega\lVert u\rVert^2.
\]
Gronwall gives \eqref{eq:n9v79-U-bound}.  Insert it into
\eqref{eq:n9v79-Duhamel} and use the triangle inequality.
\end{proof}

\begin{corollary}[Exact propagation of the zero constraint]
\label{cor:n9v79-zero-propagation}
If $c(t_0)=0$ and $f(t)=0$ for all $t$, then $c(t)=0$ for all later
times on which the identity is defined.  A negative bound
$\omega\le-\kappa<0$ additionally damps nonzero initial violations:
\begin{equation}
 \lVert c(t)\rVert
 \le e^{-\kappa(t-t_0)}\lVert c(t_0)\rVert
 +\int_{t_0}^te^{-\kappa(t-s)}\lVert f(s)\rVert\,ds.
 \label{eq:n9v79-damped}
\end{equation}
\end{corollary}

\begin{proof}
Use \Cref{thm:n9v79-Duhamel,prop:n9v79-log-bound}.
\end{proof}

\begin{firewall}[Preservation is not attraction]
\label{fire:n9v79-preservation-not-attraction}
The homogeneous identity preserves an initially zero constraint even
when $\omega=0$ or $\omega>0$.  It attracts nearby violations only
under a dissipative estimate such as $\omega\le-\kappa$.  Hamiltonian
constraint propagation is often neutral rather than dissipative, so a
strict decay claim cannot be inserted without proof.
\end{firewall}

\subsection{What propagation can observe}

Assume now that the path remains on the exact constraint surface,
\begin{equation}
 c(t)=0.
 \label{eq:n9v79-c-zero}
\end{equation}
Then \eqref{eq:n9v79-linear-C} gives the algebraic source relation
\begin{equation}
 B_p(t)e_p(t)
 =-g(t),
 \qquad
 g:=B_re_r+w+DC[r_V].
 \label{eq:n9v79-observation-equation}
\end{equation}

\begin{theorem}[Pointwise physical observability]
\label{thm:n9v79-observability}
Suppose there is $\beta>0$ such that
\begin{equation}
 \lVert B_p(t)v\rVert_{\mathcal H_C}
 \ge\beta\lVert v\rVert_{\mathcal H_p}
 \quad(v\in\mathcal H_p)
 \label{eq:n9v79-observability-gap}
\end{equation}
uniformly along the constrained path.  Then
\begin{equation}
 \lVert e_p(t)\rVert
 \le\beta^{-1}\left(
  \lVert B_r(t)e_r(t)\rVert+
  \lVert w(t)\rVert+
  \lVert DC(q(t))[r_V(t)]\rVert
 \right).
 \label{eq:n9v79-ep-bound}
\end{equation}
Consequently a vanishing retained Euler, Ward-breaking, and evolution
residual forces $e_p=0$ whenever the observability gap remains open.
\end{theorem}

\begin{proof}
Apply \eqref{eq:n9v79-observability-gap} to
\eqref{eq:n9v79-observation-equation} and use the triangle inequality.
\end{proof}

\begin{theorem}[Kernel obstruction]
\label{thm:n9v79-kernel-obstruction}
If $\ker B_p(t)$ contains a nonzero vector, the constraint propagation
identity at that time cannot bound the full physical residual.  More
precisely, for $0\ne v\in\ker B_p(t)$ and every scalar $a$,
\begin{equation}
 B_p(e_p+av)=B_pe_p,
 \label{eq:n9v79-invisible}
\end{equation}
while $\lVert e_p+av\rVert$ is unbounded as $|a|\to\infty$.
\end{theorem}

\begin{proof}
Linearity and $B_pv=0$ prove the first identity.  The reverse triangle
inequality gives
$\lVert e_p+av\rVert\ge |a|\lVert v\rVert-
\lVert e_p\rVert$.
\end{proof}

\begin{firewall}[No control from $C=0$ without observability]
\label{fire:n9v79-no-observability}
Constraint preservation says that the total forcing in
\eqref{eq:n9v79-f} vanishes on an exact constrained trajectory.  It
does not identify the individual sources, and it cannot see
$\ker B_p$.  Therefore using $C=0$ to conclude
$\mathcal E_p=0$ without a lower bound such as
\eqref{eq:n9v79-observability-gap} is circular.
\end{firewall}

\subsection{Observable and silent physical sectors}

Remove the kernel of $B_p$ exactly as V78 removed gauge stabilizers.
Set
\begin{equation}
 \mathcal H_{p,o}:=(\ker B_p)^\perp,
 \qquad
 \mathcal H_{p,s}:=\ker B_p,
 \label{eq:n9v79-observable-silent}
\end{equation}
and let $P_o,P_s$ be their orthogonal projectors.

\begin{proposition}[Moore--Penrose observable reconstruction]
\label{prop:n9v79-observable-reconstruction}
If $B_{p,o}:=B_p|_{\mathcal H_{p,o}}$ obeys
$B_{p,o}^*B_{p,o}\ge\beta^2I$, then
\begin{equation}
 P_oe_p
 =-(B_{p,o}^*B_{p,o})^{-1}B_{p,o}^*g,
 \label{eq:n9v79-eo}
\end{equation}
and
\begin{equation}
 \lVert P_oe_p\rVert\le\beta^{-1}\lVert g\rVert.
 \label{eq:n9v79-eo-bound}
\end{equation}
The silent component $P_se_p$ remains an independent physical owner.
\end{proposition}

\begin{proof}
Apply $B_{p,o}^*$ to
$B_{p,o}P_oe_p=-g$, invert the positive Gram operator, and use the
least-singular-value bound.  Since $B_pP_se_p=0$, the propagation
equation contains no information about the silent component.
\end{proof}

\begin{corollary}[Refined V78 physical defect]
\label{cor:n9v79-refined-defect}
On the exact constraint surface, the V78 physical source defect
$\mathcal E_p=-e_p$ splits as
\begin{equation}
 \mathcal E_p=\mathcal E_{p,o}+\mathcal E_{p,s},
 \qquad
 \lVert\mathcal E_{p,o}\rVert\le\beta^{-1}\lVert g\rVert,
 \label{eq:n9v79-Epsplit}
\end{equation}
where $\mathcal E_{p,s}=-P_se_p$ must be controlled by an additional
physical equation, calibration, or retained infrared variable.
\end{corollary}

\begin{proof}
Use \Cref{prop:n9v79-observable-reconstruction} and
$\mathcal E_p=-\Pi_p\Gamma_\varphi=-e_p$.
\end{proof}

\subsection{Insertion into the reaction margin}

Let
\begin{equation}
 b_0=P_0(\mathscr W-R_y^*\Gamma_y)
 \label{eq:n9v79-b0}
\end{equation}
be the V78 gauge Ward source.  Combining the two orthogonal splits
gives
\begin{equation}
 \mathcal E=mathcal E_g+mathcal E_{p,o}
 +\mathcal E_{p,s}.
 \label{eq:n9v79-three-E}
\end{equation}

\begin{proposition}[Ward--propagation reaction gate]
\label{prop:n9v79-combined-gate}
Under the V78 Gram gap $m$, the observable gap $\beta$, and the V76
branch bounds $\nu$ and $\ell_L$, the reaction constant satisfies
\begin{equation}
 \eta
 \le\nu\ell_L\left(
 m^{-1}\lVert b_0\rVert+
 \beta^{-1}\lVert g\rVert+
 \lVert\mathcal E_{p,s}\rVert
 \right).
 \label{eq:n9v79-eta}
\end{equation}
Therefore strict inequality of the right side with $1$ is a sufficient
positive reduced-Hessian gate.
\end{proposition}

\begin{proof}
Use \eqref{eq:n9v78-Eg-bound} for the gauge part,
\eqref{eq:n9v79-Epsplit} for the observable physical part, and the
triangle inequality in \eqref{eq:n9v78-sigma-split}.  Apply
\Cref{prop:n9v76-reaction-bound}.
\end{proof}

\begin{firewall}[The silent sector is not a remainder by definition]
\label{fire:n9v79-silent-not-small}
Membership in $\ker B_p$ gives no small factor.  A silent physical mode
may be relevant, marginal, topological, or unstable.  It can enter the
stable norm only after an independent scaling and locality estimate;
otherwise it must be calibrated or retained.
\end{firewall}

\subsection{Discrete cofinal propagation}

For compatibility with the exact coarse recursion, consider
\begin{equation}
 c_{j+1}=T_jc_j+f_j.
 \label{eq:n9v79-discrete}
\end{equation}

\begin{proposition}[Discrete Duhamel formula]
\label{prop:n9v79-discrete-Duhamel}
For $n>j_0$,
\begin{equation}
 c_n=T_{n-1}\cdots T_{j_0}c_{j_0}
 +\sum_{r=j_0}^{n-1}T_{n-1}\cdots T_{r+1}f_r,
 \label{eq:n9v79-discrete-Duhamel}
\end{equation}
where the product is the identity when $r=n-1$.  If
$\lVert T_j\rVert\le\rho<1$, then
\begin{equation}
 \lVert c_n\rVert
 \le\rho^{n-j_0}\lVert c_{j_0}\rVert
 +\sum_{r=j_0}^{n-1}\rho^{n-1-r}\lVert f_r\rVert.
 \label{eq:n9v79-discrete-bound}
\end{equation}
\end{proposition}

\begin{proof}
Iterate \eqref{eq:n9v79-discrete}; induction gives
\eqref{eq:n9v79-discrete-Duhamel}.  Submultiplicativity gives the
bound.
\end{proof}

\begin{firewall}[Unit propagation sectors require calibration]
\label{fire:n9v79-unit-sector}
If a constraint owner has an eigenvalue-one coarse pullback, the strict
$\rho<1$ estimate fails exactly as in V70.  That component must be
included in the moving reference constraint or treated as retained
data.  Stable contraction applies only after this calibrated projection.
\end{firewall}

\subsection{Constraint propagation card}

\begin{definition}[CPC]
\label{def:n9v79-CPC}
The constraint propagation card requires:
\begin{enumerate}
\item an off-shell regulated identity \eqref{eq:n9v79-CPI};
\item a declared evolution residual $r_V$ and the complete forcing
      \eqref{eq:n9v79-f};
\item a homogeneous propagator bound in the physical norm;
\item exact initial constraint ownership;
\item separate Ward, cutoff, anomaly, gauge-fixing, boundary, and
      retained Euler forcing rows;
\item the observable/silent split
      \eqref{eq:n9v79-observable-silent};
\item a uniform Gram gap on the observable sector;
\item independent ownership of the silent physical sector;
\item insertion of all three defect components into
      \eqref{eq:n9v79-eta}; and
\item calibrated removal of unit propagation sectors in the cofinal
      recursion.
\end{enumerate}
\end{definition}

\begin{theorem}[CPC controls exactly the observable physical defect]
\label{thm:n9v79-CPC}
Under CPC, an initially vanishing constraint obeys the exact Duhamel
formula \eqref{eq:n9v79-Duhamel}.  On an exact constrained trajectory,
the observable component of the physical conformal defect is bounded by
\eqref{eq:n9v79-Epsplit}.  The silent component remains explicit.  If
the right side of \eqref{eq:n9v79-eta} is less than one, the V76
reduced-Hessian margin remains strict.
\end{theorem}

\begin{proof}
Use \Cref{thm:n9v79-Duhamel,
prop:n9v79-observable-reconstruction,
cor:n9v79-refined-defect,prop:n9v79-combined-gate}.
\end{proof}

\begin{firewall}[Lorentzian propagation is not OS reconstruction]
\label{fire:n9v79-not-OS}
The parameter $t$ in this compiler labels the declared constraint
evolution.  A Lorentzian propagation identity does not prove Euclidean
reflection positivity, existence of a positive transfer kernel, or OS
reconstruction.  Those remain the separate V67--V72 transfer gates.
\end{firewall}

\begin{theorem}[Ninth-series V79 result]
\label{thm:n9v79-result}
A typed off-shell propagation identity gives the exact forced equation
$\dot c=Kc+f$ and Duhamel formula
\eqref{eq:n9v79-Duhamel}.  It preserves $c=0$ when every forcing row
vanishes.  On an exact constrained path it controls the physical
conformal residual only through $B_p$: the observable component has the
Moore--Penrose formula \eqref{eq:n9v79-eo}, while
$\ker B_p$ is rigorously invisible.  Combining this split with the V78
Ward quotient yields the reaction gate \eqref{eq:n9v79-eta}.
\end{theorem}

\begin{proof}
Combine \Cref{lem:n9v79-forcing,thm:n9v79-Duhamel,
thm:n9v79-observability,thm:n9v79-kernel-obstruction,
prop:n9v79-observable-reconstruction,prop:n9v79-combined-gate}.
\end{proof}

\begin{assessment}[Progress at ninth-series V79]
\label{ass:n9v79-status}
Constraint preservation and physical stationarity are now separated.
The former is a Duhamel problem; the latter follows from propagation
only on a quantitatively observable sector.  The remaining silent
physical residual is an explicit owner rather than an assumed zero.
This identifies a concrete target for the next companion audit: compare
the newest YM and NS stable norms and propagation mechanisms with the
observable/silent split before choosing the V81 acquisition.
\end{assessment}

\begin{programme}[Ninth-series V80: compulsory companion audit]
\label{prog:n9v79-V80}
V80 will inspect the newest Yang--Mills and Navier--Stokes notes, audit
their current propagation, marginal calibration, and complete-norm
results against CPC, and set the next direction.  V81 will continue
after the audit.
\end{programme}

\begin{firewall}[Claim boundary after ninth-series V79]
\label{fire:n9v79-boundary}
V79 does not derive the physical regulated propagation identity, prove
the observable gap, control the silent physical sector, close the YM
norm gate, establish anomaly cancellation, a global York branch,
infrared closure, regulator independence, OS reconstruction, or the
full SM--Einstein continuum.
\end{firewall}

\paragraph{Parent-version record.}
V79 was formed from
\texttt{Renewal\_SM\_Einstein\_Continuum\_Research\_Notes\_V78.tex}.
The V78 parent had $30\,245$ counted source lines, $1\,090\,338$ bytes,
and SHA--256
\[
 \texttt{1C52C16F9F7B52E2E5F4EFA8F85C5EE41116F63D25177749D1F22E6A83E312A7}.
\]

% =============================================================================
% END APPENDED NINTH-SERIES V79 MATERIAL
% =============================================================================

% =============================================================================
% BEGIN APPENDED NINTH-SERIES V80 MATERIAL
% =============================================================================

\section{Ninth-series V80: companion audit and joint calibration rank}
\label{sec:v9v80}

\subsection{Compulsory companion audit}

The newest companion notes present at the V80 gate were inspected
read-only:
\begin{align*}
 &\texttt{Renewal\_Yang\_Mills\_Mass\_Gap\_Research\_Notes\_V5.tex},\\
 &\texttt{Renewal\_NS\_Stability\_Research\_Notes\_V26.tex}.
\end{align*}
Their inspection records were
\begin{center}
\begin{tabular}{c|r|l}
file & bytes & SHA--256\\ \hline
YM V5 & $103\,087$ &
\texttt{09320D6A0C1E99A0ADC763EC593052942AA33F6ACF6CEDA520CD2B8113C18570}\\
NS V26 & $278\,283$ &
\texttt{82C887F8C2C69E4F28FAD7C3DC8DE5B9099CB5A4BF431EBA1FFF3E91935978AD}
\end{tabular}
\end{center}

YM V5 proves an exact implicit chart on a prescribed output marginal
level set.  Its stable derivative is
$C-D_{\rm ms}A^{-1}B_{\rm sm}$, and its two centered leakage blocks
give a quadratic calibration debit when complete conditional covariance
bounds hold.  It also requires the entire finite-cutoff quadratic
commutant, including anisotropic owners, to be calibrated; continuum
rotational restoration is left open.  NS V26 is unchanged.  Its exact
rank-one Oseen-kernel norms concern a Navier--Stokes tensor input class
and supply no SM constraint-propagation identity or calibration row.

\begin{decision}[V80 audit transfer]
\label{dec:n9v80-transfer}
The YM level-set mechanism is transferred to a joint SM calibration
map containing both the full finite-cutoff marginal ledger and a
finite-dimensional calibratable part of the V79 silent physical defect.
The transfer keeps the inverse-Jacobian and leakage gates explicit.
The NS constants are not imported.  V81 will continue after this audit
by differentiating the resulting moving reference through constraint
propagation.
\end{decision}

\subsection{Joint output conditions}

At a fixed cutoff let $P,Q,Z$ and $W$ be finite-dimensional real
Hilbert spaces.  Here $p\in P$ denotes running reference parameters,
$q\in Q$ a stable interaction, and
\begin{equation}
 \mathcal F:P\times Q\to Z
 \label{eq:n9v80-F}
\end{equation}
collects the output conditions to be held fixed.  In the application,
write
\begin{equation}
 Z=Z_{\rm marg}\oplus Z_{\rm phys},
 \qquad
 \mathcal F(p,q)=
 \binom{\beta(p,q)-\beta_*}{P_c\mathcal E_{p,s}(p,q)}.
 \label{eq:n9v80-joint-F}
\end{equation}
The first row contains every noncontracting finite-cutoff coefficient.
The second contains a declared finite-dimensional projection $P_c$ of
the silent physical defect from V79.  Let
\begin{equation}
 \mathcal R:P\times Q\to W
 \label{eq:n9v80-R}
\end{equation}
be any remaining output, such as the stable interaction together with
$(I-P_c)\mathcal E_{p,s}$.

\begin{theorem}[Joint calibration implicit chart]
\label{thm:n9v80-joint-IFT}
Suppose $\mathcal F$ is $C^1$ near $(p_0,q_0)$,
$\mathcal F(p_0,q_0)=0$, and
\begin{equation}
 A_0:=D_p\mathcal F(p_0,q_0):P\to Z
 \label{eq:n9v80-A0}
\end{equation}
is an isomorphism.  Then there are neighbourhoods and a unique $C^1$
map
\begin{equation}
 p=\mathfrak k(q),
 \qquad
 \mathfrak k(q_0)=p_0,
 \qquad
 \mathcal F(\mathfrak k(q),q)=0.
 \label{eq:n9v80-k}
\end{equation}
At every point of the chart,
\begin{equation}
 D\mathfrak k=-A^{-1}B,
 \qquad
 A:=D_p\mathcal F,
 \quad B:=D_q\mathcal F.
 \label{eq:n9v80-Dk}
\end{equation}
\end{theorem}

\begin{proof}
Apply the finite-dimensional implicit-function theorem to
$\mathcal F(p,q)=0$.  Differentiating
$\mathcal F(\mathfrak k(q),q)=0$ gives
$A D\mathfrak k+B=0$, proving \eqref{eq:n9v80-Dk}.
\end{proof}

\begin{theorem}[Joint calibration Schur differential]
\label{thm:n9v80-Schur}
Define the aligned output
\begin{equation}
 \mathcal R_{\rm al}(q):=\mathcal R(\mathfrak k(q),q).
 \label{eq:n9v80-Ral}
\end{equation}
Then
\begin{equation}
 D\mathcal R_{\rm al}
 =D_q\mathcal R-D_p\mathcal R\,A^{-1}B.
 \label{eq:n9v80-Schur}
\end{equation}
If on a convex aligned chart
\begin{equation}
 \lVert D_q\mathcal R\rVert\le c,
 \quad \lVert D_p\mathcal R\rVert\le d,
 \quad \lVert A^{-1}\rVert\le a,
 \quad \lVert B\rVert\le b,
 \label{eq:n9v80-four-bounds}
\end{equation}
then
\begin{equation}
 \lVert\mathcal R_{\rm al}(q)-
 \mathcal R_{\rm al}(q')\rVert
 \le(c+dab)\lVert q-q'\rVert.
 \label{eq:n9v80-aligned-bound}
\end{equation}
\end{theorem}

\begin{proof}
The chain rule and \eqref{eq:n9v80-Dk} give
\eqref{eq:n9v80-Schur}.  Submultiplicativity bounds its norm by
$c+dab$; integrate along the segment from $q'$ to $q$.
\end{proof}

\begin{firewall}[Calibration changes the derivative]
\label{fire:n9v80-naive-derivative}
The derivative at fixed reference parameter is $D_q\mathcal R$.
The derivative along fixed output conditions is
\eqref{eq:n9v80-Schur}.  Dropping the second term loses the response of
the running cosmological, Newton, gauge, Higgs, Yukawa, anisotropic, and
physical-defect coordinates.
\end{firewall}

\subsection{The calibration rank obstruction}

The square isomorphism in \Cref{thm:n9v80-joint-IFT} is not a notation
choice.  It encodes the number and independence of the available
renormalization conditions.

\begin{theorem}[Finite-dimensional rank obstruction]
\label{thm:n9v80-rank-obstruction}
Let $r=\dim P$ and $s=\dim Z$.
\begin{enumerate}
\item If $r<s$, no derivative $D_p\mathcal F:P\to Z$ is surjective.
Thus arbitrary nearby joint output conditions cannot be solved by
varying $p$.
\item If $r=s$ but $D_p\mathcal F$ is singular, no local unique
$C^1$ calibration follows from the inverse-function mechanism.
\item If $r>s$ and $D_p\mathcal F$ is surjective, the local solution
set has at least $r-s$ parameter directions.  A unique calibration
requires an additional transverse slice.
\end{enumerate}
\end{theorem}

\begin{proof}
The rank of a linear map $P\to Z$ is at most $r$, proving the first
claim.  The derivative of a locally unique $C^1$ inverse must be an
inverse of $D_p\mathcal F$, proving the second in the stated
inverse-function setting.  For the third, the submersion theorem makes
each regular level set a manifold of dimension $r-s$; choosing a
complement to the kernel supplies a transverse slice.
\end{proof}

\begin{corollary}[Only finitely many silent conditions can be calibrated]
\label{cor:n9v80-finite-silent}
After the full marginal ledger consumes $r_{\rm marg}$ independent
parameter directions, at most
$r-r_{\rm marg}$ additional independent components of
$P_c\mathcal E_{p,s}$ can be imposed by the same parameter vector.
Every remaining silent physical component requires a dynamical stable
estimate, an additional physical equation, a new running parameter, or
explicit retention.
\end{corollary}

\begin{proof}
Apply \Cref{thm:n9v80-rank-obstruction} to the direct-sum target
\eqref{eq:n9v80-joint-F}.
\end{proof}

\begin{firewall}[A functional defect cannot be tuned away by a short list]
\label{fire:n9v80-functional-defect}
If the silent physical defect contains an unrestricted local function
or infinitely many independent cutoff modes, finitely many SM running
couplings cannot set it identically to zero.  Only a finite projection
can enter \eqref{eq:n9v80-joint-F}; the complement must be proved
irrelevant, constrained by another equation, or retained.
\end{firewall}

\subsection{Rectangular calibration with a declared slice}

The underdetermined case can still be made exact.  Suppose
$A_0=D_p\mathcal F(p_0,q_0)$ is surjective.  Set
\begin{equation}
 P_0:=(\ker A_0)^\perp,
 \qquad
 J_0:=A_0^*(A_0A_0^*)^{-1}:Z\to P_0.
 \label{eq:n9v80-J0}
\end{equation}

\begin{proposition}[Minimum-norm calibration slice]
\label{prop:n9v80-minimum-slice}
$A_0J_0=I_Z$, and $J_0z$ is the unique minimum-norm solution of
$A_0u=z$.  Restricting $p-p_0$ to $P_0$, the derivative is an
isomorphism and the implicit calibration chart exists uniquely on that
slice.  Its base derivative is
\begin{equation}
 D\mathfrak k(q_0)=-J_0B_0.
 \label{eq:n9v80-Dk-rectangular}
\end{equation}
\end{proposition}

\begin{proof}
Surjectivity makes $A_0A_0^*>0$, and direct multiplication gives
$A_0J_0=I$.  Every solution is $J_0z+v$ with $v\in\ker A_0$, and the
two summands are orthogonal.  Hence $J_0z$ is the unique minimum-norm
solution.  The restriction $A_0|_{P_0}$ is bijective with inverse
$J_0$, so the implicit-function theorem applies on the affine slice.
Differentiation gives \eqref{eq:n9v80-Dk-rectangular}.
\end{proof}

This is the calibration analogue of the V75 minimum-energy lift and
the V78 range projector.  Its kernel directions are genuine scheme or
reference freedoms and must be fixed once, not reselected silently at
each scale.

\subsection{Centered leakage and the joint debit}

Decompose the aligned output derivative into four blocks at a base
point:
\begin{equation}
 A=D_p\mathcal F,
 \quad B=D_q\mathcal F,
 \quad C=D_q\mathcal R,
 \quad D=D_p\mathcal R.
 \label{eq:n9v80-ABCD}
\end{equation}
Then \eqref{eq:n9v80-Schur} is $C-DA^{-1}B$.

\begin{proposition}[Quadratic joint leakage debit]
\label{prop:n9v80-quadratic-debit}
Assume at $q=0$ that
\begin{equation}
 B(p_0,0)=0,
 \qquad D(p_0,0)=0.
 \label{eq:n9v80-centred}
\end{equation}
Suppose on $\lVert q\rVert\le r$,
\begin{equation}
 \lVert B\rVert\le L_B\lVert q\rVert,
 \qquad
 \lVert D\rVert\le L_D\lVert q\rVert,
 \qquad
 \lVert A^{-1}\rVert\le a.
 \label{eq:n9v80-leakage-bounds}
\end{equation}
Then
\begin{equation}
 \lVert DA^{-1}B\rVert
 \le aL_BL_D\lVert q\rVert^2.
 \label{eq:n9v80-quadratic-debit}
\end{equation}
If additionally $\lVert C\rVert\le\kappa_0<1$, the aligned output is a
strict contraction whenever
\begin{equation}
 \kappa_0+aL_BL_Dr^2<1.
 \label{eq:n9v80-joint-contraction}
\end{equation}
\end{proposition}

\begin{proof}
Multiply the three bounds in
\eqref{eq:n9v80-leakage-bounds}.  The derivative formula
\eqref{eq:n9v80-Schur} and the triangle inequality give the contraction
constant.
\end{proof}

\begin{firewall}[Gaussian centering is operator-specific]
\label{fire:n9v80-centering}
The V75 energy-orthogonal Gaussian split can prove
\eqref{eq:n9v80-centred} only when the selected condition and output
insertions are orthogonal for that exact Gaussian and the projections
are matched to it.  Interacting conditional cumulants generate $B$ and
$D$; complete covariance bounds, not formal centering, must prove
\eqref{eq:n9v80-leakage-bounds}.
\end{firewall}

\subsection{Full finite-cutoff owner space}

Let $Z_{\rm phys}$ denote the desired continuum-invariant marginal
coordinates and let $Z_{\rm an}$ contain finite-cutoff marginal
operators allowed by the regulator symmetry but not by the target
continuum symmetry.  The actual condition space is
\begin{equation}
 Z_{\rm marg}=Z_{\rm phys}\oplus Z_{\rm an}.
 \label{eq:n9v80-anisotropic-split}
\end{equation}

\begin{proposition}[Finite-cutoff completeness requirement]
\label{prop:n9v80-full-owner}
Every noncontracting eigenvector of the exact linearized coarse map must
belong to $Z_{\rm marg}$, whether or not it is invariant under the
target continuum group.  Omitting an anisotropic noncontracting owner
places it in the purported stable space and prevents a strict stable
contraction.  Vanishing of its calibrated trajectory in the continuum
requires a separate symmetry-restoration estimate.
\end{proposition}

\begin{proof}
If an omitted direction $v$ has coarse multiplier of modulus at least
one, the derivative norm of the alleged stable map is at least one on
$v$.  This contradicts strict contraction.  Adding $v$ to the
calibration target removes that contradiction but supplies no estimate
that its selected level tends to zero.
\end{proof}

For SM--Einstein this ledger includes all cutoff-allowed cosmological,
Einstein, curvature-squared, gauge, Higgs, Yukawa, gauge-fixing,
anisotropic, and relevant boundary owners, modulo exact symmetries and
redundancies proved at the same cutoff.

\begin{firewall}[Continuum symmetry cannot be assumed at finite cutoff]
\label{fire:n9v80-symmetry-restoration}
Hypercubic, foliation, or gauge-fixed regulator symmetry can permit
more marginal tensors than the continuum Lorentz/diffeomorphism
commutant.  Calibrating only the desired continuum coefficients assumes
the restoration theorem that the programme must prove.
\end{firewall}

\subsection{Aligned silent-defect bound}

Let the residual silent output after calibration be
\begin{equation}
 e_{\rm rem}(q):=(I-P_c)\mathcal E_{p,s}(
 \mathfrak k(q),q).
 \label{eq:n9v80-erem}
\end{equation}

\begin{proposition}[Calibrated residual estimate]
\label{prop:n9v80-erem-bound}
Assume $e_{\rm rem}(q_0)=0$ and, on a convex aligned ball,
\begin{equation}
 \left\lVert
 D_qe_{\rm rem}-D_pe_{\rm rem}A^{-1}B
 \right\rVert\le\kappa_e.
 \label{eq:n9v80-ke}
\end{equation}
Then
\begin{equation}
 \lVert e_{\rm rem}(q)\rVert
 \le\kappa_e\lVert q-q_0\rVert.
 \label{eq:n9v80-erem-bound}
\end{equation}
Inserting this into V79 replaces the silent term in
\eqref{eq:n9v79-eta} and gives the sufficient reaction gate
\begin{equation}
 \nu\ell_L\left(
 m^{-1}\lVert b_0\rVert+
 \beta^{-1}\lVert g\rVert+
 \kappa_e\lVert q-q_0\rVert
 \right)<1.
 \label{eq:n9v80-calibrated-reaction}
\end{equation}
\end{proposition}

\begin{proof}
The derivative in \eqref{eq:n9v80-ke} is exactly the aligned derivative
from \Cref{thm:n9v80-Schur}.  Integrate it along the segment.  Then use
\Cref{prop:n9v79-combined-gate}.
\end{proof}

\subsection{Joint calibration card}

\begin{definition}[JCC]
\label{def:n9v80-JCC}
The joint calibration card requires:
\begin{enumerate}
\item the complete finite-cutoff noncontracting owner space, including
      anisotropic and boundary coordinates;
\item a declared finite-dimensional calibratable projection of the
      V79 silent physical defect;
\item enough independent running parameters to satisfy the rank
      condition of \Cref{thm:n9v80-rank-obstruction};
\item an invertible square Jacobian or a surjective Jacobian with a
      fixed transverse slice;
\item a uniform inverse bound for that Jacobian;
\item the exact aligned derivative \eqref{eq:n9v80-Schur};
\item complete conditional covariance bounds for both leakage blocks;
\item stable control of the uncalibrated silent remainder;
\item the strict reaction gate
      \eqref{eq:n9v80-calibrated-reaction}; and
\item a separate cofinal symmetry-restoration theorem for unwanted
      finite-cutoff marginal trajectories.
\end{enumerate}
\end{definition}

\begin{theorem}[JCC closes the finite-dimensional calibration step]
\label{thm:n9v80-JCC}
Under JCC, the prescribed marginal levels and calibratable silent
physical conditions determine a local moving reference chart.  Its
remaining output derivative is the Schur differential
\eqref{eq:n9v80-Schur}; centered leakage contributes only the quadratic
debit \eqref{eq:n9v80-quadratic-debit}; and the uncalibrated silent
remainder enters the explicit reaction gate
\eqref{eq:n9v80-calibrated-reaction}.
\end{theorem}

\begin{proof}
Use \Cref{thm:n9v80-joint-IFT,thm:n9v80-Schur,
prop:n9v80-quadratic-debit,prop:n9v80-erem-bound}.  The rank and full
owner requirements are supplied by
\Cref{thm:n9v80-rank-obstruction,prop:n9v80-full-owner}.
\end{proof}

\begin{theorem}[Ninth-series V80 result]
\label{thm:n9v80-result}
The compulsory audit transfers the YM V5 marginal-level mechanism to a
joint SM condition map.  An invertible calibration Jacobian gives
$D\mathfrak k=-A^{-1}B$ and the exact aligned derivative
$D_q\mathcal R-D_p\mathcal R A^{-1}B$.  Centered leakage has a quadratic
debit.  The rank obstruction proves that only as many independent
silent physical conditions can be calibrated as there are unused
running parameter directions, while all other components require
dynamical stable control or retention.
\end{theorem}

\begin{proof}
Combine \Cref{thm:n9v80-joint-IFT,thm:n9v80-Schur,
thm:n9v80-rank-obstruction,prop:n9v80-quadratic-debit,
prop:n9v80-erem-bound}.
\end{proof}

\begin{assessment}[Progress at ninth-series V80]
\label{ass:n9v80-status}
The V79 silent physical sector is no longer a single undifferentiated
remainder.  Its finite calibratable projection, its stable complement,
and the rank needed to separate them are explicit.  The YM audit also
forces the finite-cutoff anisotropic owners into the reference space.
The next step is to propagate constraints along this moving chart and
include the connection term generated by $D\mathfrak k$.
\end{assessment}

\begin{programme}[Ninth-series V81: moving-reference propagation]
\label{prog:n9v80-V81}
V81 will continue after the audit.  It will differentiate the constraint
and Ward data along $p=\mathfrak k(q)$, derive the exact connection
forcing caused by reference motion, and prove a discrete Duhamel bound
which counts that forcing once.  The target is a coupled calibration--
propagation gate compatible with V79 and JCC.
\end{programme}

\begin{firewall}[Claim boundary after ninth-series V80]
\label{fire:n9v80-boundary}
V80 does not prove the model-specific calibration Jacobian invertible,
the complete covariance bounds, rotational/diffeomorphism restoration,
smallness of the uncalibrated physical defect, the regulated constraint
propagation identity, anomaly cancellation, infrared closure,
regulator independence, OS reconstruction, or the full SM--Einstein
continuum.
\end{firewall}

\paragraph{Parent-version record.}
V80 was formed from
\texttt{Renewal\_SM\_Einstein\_Continuum\_Research\_Notes\_V79.tex}.
The V79 parent had $30\,737$ counted source lines, $1\,107\,011$ bytes,
and SHA--256
\[
 \texttt{47EBBA6BD0A9933481DA48150938EFD89B5943D4CF27547541FB75794887F425}.
\]

% =============================================================================
% END APPENDED NINTH-SERIES V80 MATERIAL
% =============================================================================

% =============================================================================
% BEGIN APPENDED NINTH-SERIES V81 MATERIAL
% =============================================================================

\section{Ninth-series V81: moving-reference constraint propagation}
\label{sec:v9v81}

V80 constructs a local reference chart $p=\mathfrak k(q)$ by fixing
joint marginal and calibratable physical output conditions.  Constraint
propagation along that chart is not obtained by differentiating at
fixed $p$.  Reference motion contributes an exact connection term.  The
same term appears either inside the aligned derivative or as an
explicit forcing, but never both.

\subsection{Calibration with a moving target}

Let
\begin{equation}
 \mathcal Z:P\times Q\to Z
 \label{eq:n9v81-Z}
\end{equation}
be a $C^2$ joint output map, and let $z_*(s)\in Z$ be a prescribed
scale- or time-dependent target.  A calibrated path satisfies
\begin{equation}
 \mathcal Z(p(s),q(s))=z_*(s).
 \label{eq:n9v81-level}
\end{equation}
Set
\begin{equation}
 A:=D_p\mathcal Z,
 \qquad B:=D_q\mathcal Z,
 \label{eq:n9v81-AB}
\end{equation}
and assume $A$ is invertible on the declared calibration slice.

\begin{lemma}[Exact reference velocity]
\label{lem:n9v81-reference-velocity}
Along every calibrated differentiable path,
\begin{equation}
 \dot p=A^{-1}(\dot z_*-B\dot q).
 \label{eq:n9v81-pdot}
\end{equation}
For a fixed target this reduces to
$\dot p=-A^{-1}B\dot q$.
\end{lemma}

\begin{proof}
Differentiate \eqref{eq:n9v81-level} and solve
$A\dot p+B\dot q=\dot z_*$ for $\dot p$.
\end{proof}

\begin{corollary}[Reference-speed bound]
\label{cor:n9v81-reference-speed}
If $\lVert A^{-1}\rVert\le a$ and
$\lVert B\rVert\le b$, then
\begin{equation}
 \lVert\dot p\rVert
 \le a\bigl(\lVert\dot z_*\rVert+b\lVert\dot q\rVert\bigr).
 \label{eq:n9v81-pdot-bound}
\end{equation}
\end{corollary}

\begin{proof}
Take norms in \eqref{eq:n9v81-pdot}.
\end{proof}

\subsection{The calibration connection}

Let $H:P\times Q\to E$ be any $C^1$ regulated quantity: a constraint,
Ward residual, projector, branch operator, or reaction coefficient.

\begin{definition}[Calibration covariant derivative]
\label{def:n9v81-covariant-derivative}
For a fixed output target define
\begin{equation}
 \nabla_q^{\rm cal}H
 :=D_qH-D_pH\,A^{-1}B.
 \label{eq:n9v81-cal-derivative}
\end{equation}
The operator
\begin{equation}
 \mathcal A_H:=D_pH\,A^{-1}B
 \label{eq:n9v81-connection}
\end{equation}
is the calibration connection acting on $H$.
\end{definition}

\begin{theorem}[Exact moving-reference chain rule]
\label{thm:n9v81-chain-rule}
Along \eqref{eq:n9v81-level},
\begin{equation}
 \frac d{ds}H(p(s),q(s))
 =\nabla_q^{\rm cal}H\,\dot q
 +D_pH\,A^{-1}\dot z_*.
 \label{eq:n9v81-moving-chain}
\end{equation}
In particular, at a fixed target the derivative is precisely
$\nabla_q^{\rm cal}H\,\dot q$.
\end{theorem}

\begin{proof}
The ordinary chain rule gives $D_pH\dot p+D_qH\dot q$.
Insert \eqref{eq:n9v81-pdot} and collect the two terms.
\end{proof}

\begin{corollary}[Connection bound]
\label{cor:n9v81-connection-bound}
If
$\lVert D_pH\rVert\le h_p$,
$\lVert A^{-1}\rVert\le a$, and
$\lVert B\rVert\le b$, then
\begin{equation}
 \lVert\mathcal A_H\rVert\le h_pab
 \label{eq:n9v81-connection-bound}
\end{equation}
and
\begin{equation}
 \left\lVert\frac d{ds}H\right\rVert
 \le(\lVert D_qH\rVert+h_pab)\lVert\dot q\rVert
 +h_pa\lVert\dot z_*\rVert.
 \label{eq:n9v81-chain-bound}
\end{equation}
\end{corollary}

\begin{proof}
Use submultiplicativity in
\eqref{eq:n9v81-connection}--\eqref{eq:n9v81-moving-chain}.
\end{proof}

\begin{firewall}[Coordinate connection, not new physics]
\label{fire:n9v81-coordinate-connection}
The connection term records motion of the chosen reference chart.  It
is not an additional interaction.  If one works directly with the
aligned function $H(\mathfrak k(q),q)$, the connection is already
inside its derivative.  Adding it again as a forcing double counts the
same reference response.
\end{firewall}

\subsection{Moving-reference constraint equation}

Let $C(p,q)$ be the regulated constraint.  Suppose that at fixed
reference parameter its declared propagation law is
\begin{equation}
 D_qC[V_q]
 =KC+B_pe_p+B_re_r+w,
 \label{eq:n9v81-fixed-propagation}
\end{equation}
where all quantities may depend on $(p,q)$.  Let an inexact calibrated
path obey
\begin{equation}
 \dot q=V_q+r_q.
 \label{eq:n9v81-q-flow}
\end{equation}

\begin{theorem}[Exact moving-reference propagation]
\label{thm:n9v81-moving-propagation}
Along the calibrated path, $c(s)=C(p(s),q(s))$ satisfies
\begin{equation}
 \dot c=Kc+f_{\rm dyn}+f_{\rm cal},
 \label{eq:n9v81-moving-C}
\end{equation}
where
\begin{align}
 f_{\rm dyn}
 &:=B_pe_p+B_re_r+w+D_qC[r_q],
 \label{eq:n9v81-fdyn}\\
 f_{\rm cal}
 &:=D_pC\,A^{-1}(\dot z_*-BV_q-Br_q).
 \label{eq:n9v81-fcal}
\end{align}
Equivalently,
\begin{equation}
 \dot c
 =\nabla_q^{\rm cal}C[V_q+r_q]
 +D_pC\,A^{-1}\dot z_*.
 \label{eq:n9v81-covariant-C}
\end{equation}
The two representations are identical.
\end{theorem}

\begin{proof}
The full chain rule is
$\dot c=D_qC[V_q+r_q]+D_pC\dot p$.  Use
\eqref{eq:n9v81-fixed-propagation} for $D_qC[V_q]$ and
\eqref{eq:n9v81-pdot} for $\dot p$, giving
\eqref{eq:n9v81-moving-C}--\eqref{eq:n9v81-fcal}.  Alternatively apply
\Cref{thm:n9v81-chain-rule} directly to $H=C$, proving
\eqref{eq:n9v81-covariant-C}.
\end{proof}

\begin{corollary}[Calibration-forcing bound]
\label{cor:n9v81-fcal-bound}
If
\begin{equation}
 \lVert D_pC\rVert\le c_p,
 \quad\lVert A^{-1}\rVert\le a,
 \quad\lVert B\rVert\le b,
 \label{eq:n9v81-fcal-hyp}
\end{equation}
then
\begin{equation}
 \lVert f_{\rm cal}\rVert
 \le c_pa\left(
 \lVert\dot z_*\rVert+b\lVert V_q\rVert
 +b\lVert r_q\rVert\right).
 \label{eq:n9v81-fcal-bound}
\end{equation}
\end{corollary}

\begin{proof}
Take norms in \eqref{eq:n9v81-fcal}.
\end{proof}

\begin{firewall}[A fixed-reference identity is not chart covariant]
\label{fire:n9v81-fixed-not-moving}
Even when the fixed-$p$ forcing in
\eqref{eq:n9v81-fixed-propagation} vanishes, a moving calibrated
reference can generate $f_{\rm cal}$.  Exact propagation of $c=0$
requires either cancellation of that term in the covariant identity or
its inclusion in the forcing ledger.  Ignoring reference motion is not
a proof of constraint preservation.
\end{firewall}

\subsection{Duhamel estimate with reference motion}

Let $U(s,r)$ be the fundamental solution for $K$ as in V79.

\begin{corollary}[Moving-reference Duhamel formula]
\label{cor:n9v81-Duhamel}
Equation \eqref{eq:n9v81-moving-C} has the exact solution
\begin{equation}
 c(s)=U(s,s_0)c(s_0)
 +\int_{s_0}^sU(s,r)
  \bigl(f_{\rm dyn}(r)+f_{\rm cal}(r)\bigr)\,dr.
 \label{eq:n9v81-Duhamel}
\end{equation}
If the logarithmic-norm bound
\eqref{eq:n9v79-log-hyp} holds, then
\begin{align}
 \lVert c(s)\rVert
 &\le e^{\int_{s_0}^s\omega}\lVert c(s_0)\rVert\\
 &\quad+\int_{s_0}^se^{\int_r^s\omega}
 \bigl(\lVert f_{\rm dyn}(r)\rVert+
       \lVert f_{\rm cal}(r)\rVert\bigr)\,dr.
 \label{eq:n9v81-Duhamel-bound}
\end{align}
\end{corollary}

\begin{proof}
Apply \Cref{thm:n9v79-Duhamel,prop:n9v79-log-bound} to the total
forcing in \eqref{eq:n9v81-moving-C}.
\end{proof}

\begin{proposition}[Cofinal integrability gate]
\label{prop:n9v81-integrability}
Suppose $\omega\le0$, $c(s_0)=0$, and
\begin{equation}
 \int_{s_0}^{\infty}
 \bigl(\lVert f_{\rm dyn}(r)\rVert+
       \lVert f_{\rm cal}(r)\rVert\bigr)\,dr<\infty.
 \label{eq:n9v81-integrable-forcing}
\end{equation}
Then $c(s)$ remains uniformly bounded.  If in addition
$\omega\le-\kappa<0$ and the total forcing tends to zero, then
\begin{equation}
 \lVert c(s)\rVert\longrightarrow0.
 \label{eq:n9v81-c-to-zero}
\end{equation}
\end{proposition}

\begin{proof}
For $\omega\le0$, \eqref{eq:n9v81-Duhamel-bound} is bounded by the
integral in \eqref{eq:n9v81-integrable-forcing}.  In the damped case,
split the convolution at a large $R$: the early part is suppressed by
$e^{-\kappa(s-R)}$, and the late part is bounded by the supremum of the
forcing after $R$ times $\kappa^{-1}$.  Send first $s$ and then $R$ to
infinity.
\end{proof}

\begin{firewall}[Neutral propagation retains accumulated connection error]
\label{fire:n9v81-neutral}
When $\omega=0$, integrable forcing yields boundedness but need not
yield $c(s)\to0$; its integral can converge to a nonzero constraint
offset.  A neutral constraint sector therefore requires exact
cancellation, a terminal calibration condition, or a separate
vanishing-tail argument.
\end{firewall}

\subsection{Discrete scale connection}

At two successive cutoff levels let a fixed-reference coarse step
produce the provisional data
\begin{equation}
 (\widetilde p_{j+1},q_{j+1})
 \label{eq:n9v81-provisional}
\end{equation}
from $(p_j,q_j)$.  The output calibration selects
$p_{j+1}=\mathfrak k_{j+1}(q_{j+1})$.  Put
\begin{equation}
 \delta p_j:=p_{j+1}-\widetilde p_{j+1}.
 \label{eq:n9v81-deltap}
\end{equation}

\begin{lemma}[Exact finite connection increment]
\label{lem:n9v81-finite-connection}
For the next-scale constraint $C_{j+1}$,
\begin{equation}
 f_{{\rm cal},j}
 :=C_{j+1}(p_{j+1},q_{j+1})
  -C_{j+1}(\widetilde p_{j+1},q_{j+1})
 =\overline{D_pC}_{j+1}\,\delta p_j,
 \label{eq:n9v81-finite-connection}
\end{equation}
where
\begin{equation}
 \overline{D_pC}_{j+1}
 :=\int_0^1D_pC_{j+1}(
 \widetilde p_{j+1}+t\delta p_j,q_{j+1})\,dt.
 \label{eq:n9v81-average-DpC}
\end{equation}
Hence
\begin{equation}
 \lVert f_{{\rm cal},j}\rVert
 \le L_{C,j}\lVert\delta p_j\rVert
 \label{eq:n9v81-finite-connection-bound}
\end{equation}
if $\lVert D_pC_{j+1}\rVert\le L_{C,j}$ on the segment.
\end{lemma}

\begin{proof}
Apply the fundamental theorem of calculus to the line segment in
parameter space and then take norms.
\end{proof}

Suppose the fixed-reference constraint recursion is
\begin{equation}
 \widetilde c_{j+1}=T_jc_j+f_{{\rm dyn},j}.
 \label{eq:n9v81-fixed-discrete}
\end{equation}

\begin{theorem}[Exact calibrated discrete recursion]
\label{thm:n9v81-discrete-recursion}
The calibrated constraint satisfies
\begin{equation}
 c_{j+1}=T_jc_j+f_{{\rm dyn},j}+f_{{\rm cal},j}.
 \label{eq:n9v81-calibrated-discrete}
\end{equation}
For $n>j_0$,
\begin{align}
 c_n={}&T_{n-1}\cdots T_{j_0}c_{j_0}\\
 &+\sum_{r=j_0}^{n-1}T_{n-1}\cdots T_{r+1}
   (f_{{\rm dyn},r}+f_{{\rm cal},r}).
 \label{eq:n9v81-discrete-Duhamel}
\end{align}
\end{theorem}

\begin{proof}
Add \eqref{eq:n9v81-finite-connection} to
\eqref{eq:n9v81-fixed-discrete}; iteration gives the Duhamel sum as in
\Cref{prop:n9v79-discrete-Duhamel}.
\end{proof}

\begin{proposition}[Two-rate convolution]
\label{prop:n9v81-two-rate}
Assume, after removal of calibrated unit sectors,
\begin{equation}
 \lVert T_j\rVert\le\sigma<1,
 \quad\lVert f_{{\rm dyn},j}\rVert\le\epsilon_j,
 \quad\lVert f_{{\rm cal},j}\rVert\le K\rho^{j-j_0},
 \quad 0<\rho<1.
 \label{eq:n9v81-two-rate-hyp}
\end{equation}
Then, writing $N=n-j_0$,
\begin{align}
 \lVert c_n\rVert
 &\le\sigma^N\lVert c_{j_0}\rVert
 +\sum_{r=j_0}^{n-1}\sigma^{n-1-r}\epsilon_r
 +K\,\Xi_N(\sigma,\rho),
 \label{eq:n9v81-two-rate}\\
 \Xi_N(\sigma,\rho)
 &:={}
 \begin{cases}
 \dfrac{\sigma^N-\rho^N}{\sigma-\rho},&\sigma\ne\rho,\\
 N\sigma^{N-1},&\sigma=\rho.
 \end{cases}
 \label{eq:n9v81-Xi}
\end{align}
In both cases $\Xi_N(\sigma,\rho)\to0$.
\end{proposition}

\begin{proof}
Apply norms to \eqref{eq:n9v81-discrete-Duhamel}.  The connection sum
is
$K\sum_{k=0}^{N-1}\sigma^{N-1-k}\rho^k$.  Evaluate the finite
geometric convolution; when the two rates coincide it has $N$ equal
terms.  Exponential decay dominates the linear factor in the resonant
case.
\end{proof}

\begin{firewall}[Reference mismatch must be measured against the transported reference]
\label{fire:n9v81-wrong-increment}
The correct increment is
$p_{j+1}-\widetilde p_{j+1}$, not merely $p_{j+1}-p_j$.
The fixed-reference coarse step may itself transport marginal
coefficients.  Using the latter difference counts physical beta-flow as
an error and can create a spurious connection forcing.
\end{firewall}

\subsection{Second calibration response}

The V76 reaction gate uses second derivatives, so the moving reference
requires a second response as well.  For a fixed target let
$p=\mathfrak k(q)$ and write
$p_h=D\mathfrak k[h]$.

\begin{proposition}[Exact second calibration response]
\label{prop:n9v81-D2k}
For $h,k\in Q$,
\begin{align}
 D^2\mathfrak k[h,k]
 =-A^{-1}\bigl(&\mathcal Z_{qq}[h,k]
 +\mathcal Z_{pq}[p_h,k]
 +\mathcal Z_{qp}[h,p_k]\\
 &+\mathcal Z_{pp}[p_h,p_k]\bigr),
 \label{eq:n9v81-D2k}
\end{align}
where $p_h=-A^{-1}Bh$ and $p_k=-A^{-1}Bk$.
\end{proposition}

\begin{proof}
Differentiate
$\mathcal Z(\mathfrak k(q),q)=z_*$ twice.  The term containing the
second derivative of $\mathfrak k$ is $A D^2\mathfrak k[h,k]$; move all
other chain-rule terms to the opposite side and apply $A^{-1}$.
\end{proof}

\begin{corollary}[Uniform second-response bound]
\label{cor:n9v81-D2k-bound}
If $\lVert A^{-1}\rVert\le a$,
$\lVert B\rVert\le b$, and all four second derivative blocks in
\eqref{eq:n9v81-D2k} have norm at most $M_2$, then
\begin{equation}
 \lVert D^2\mathfrak k[h,k]\rVert
 \le aM_2(1+2ab+a^2b^2)
 \lVert h\rVert\lVert k\rVert.
 \label{eq:n9v81-D2k-bound}
\end{equation}
\end{corollary}

\begin{proof}
$\lVert p_h\rVert\le ab\lVert h\rVert$ and similarly for $p_k$.
Apply these bounds term by term in \eqref{eq:n9v81-D2k}.
\end{proof}

\subsection{Ward and reaction data on the moving chart}

For the Ward residual and physical defect, the exact aligned derivatives
are
\begin{align}
 D\mathscr W_{\rm al}
 &=\mathscr W_q-\mathscr W_pA^{-1}B,
 \label{eq:n9v81-W-aligned}\\
 D\mathcal E_{p,\rm al}
 &=(\mathcal E_p)_q-(\mathcal E_p)_pA^{-1}B.
 \label{eq:n9v81-E-aligned}
\end{align}
Their second derivatives contain $D^2\mathfrak k$ from
\eqref{eq:n9v81-D2k}.

\begin{proposition}[Moving-chart reaction sufficiency]
\label{prop:n9v81-moving-reaction}
Suppose the V80 joint chart obeys uniform bounds on
$A^{-1},B,D^2\mathfrak k$, and the aligned Ward, observable, and silent
defect rows satisfy the V79 estimates with common constants
$m,\beta,\nu,\ell_L$.  Then the sufficient strict reaction gate remains
\begin{equation}
 \nu\ell_L\left(
 m^{-1}\lVert b_{0,\rm al}\rVert+
 \beta^{-1}\lVert g_{\rm al}\rVert+
 \lVert\mathcal E_{p,s,\rm al}\rVert
 \right)<1,
 \label{eq:n9v81-moving-reaction}
\end{equation}
provided every first and second derivative entering these three aligned
rows is computed with
\eqref{eq:n9v81-cal-derivative} and
\eqref{eq:n9v81-D2k}.
\end{proposition}

\begin{proof}
At each point of the chart, the orthogonal decompositions and algebraic
V79 reaction estimate are unchanged.  The chain rules
\eqref{eq:n9v81-W-aligned}--\eqref{eq:n9v81-E-aligned} and
\Cref{prop:n9v81-D2k} supply the derivatives required in the V76
$\nu$ bound.  Uniformity of the stated constants permits application of
\Cref{prop:n9v79-combined-gate} pointwise along the chart.
\end{proof}

\begin{firewall}[Pointwise recalibration can destroy locality]
\label{fire:n9v81-locality}
Even if each fixed-$p$ operator is local, the maps $A^{-1}$,
$D\mathfrak k$, and $D^2\mathfrak k$ can be nonlocal or singular.  The
moving chart is admissible for the continuum programme only when these
responses satisfy the same weighted locality and collar bounds as the
constraint and physical quotient.
\end{firewall}

\subsection{Calibration--propagation card}

\begin{definition}[CPCal]
\label{def:n9v81-CPCal}
The calibration--propagation card requires:
\begin{enumerate}
\item a JCC chart with uniform $A^{-1}$, first-response, and
      second-response bounds;
\item an explicit target trajectory $z_*(s)$ or target sequence;
\item the exact reference velocity \eqref{eq:n9v81-pdot};
\item a fixed-reference constraint identity with complete dynamic
      forcing;
\item exactly one representation of the calibration connection;
\item weighted locality of $A^{-1}$ and the connection operators;
\item a homogeneous constraint propagator bound;
\item cofinal summability or strict contraction of the connection
      forcing;
\item aligned Ward and physical-defect derivatives through second
      order; and
\item the uniform moving reaction gate
      \eqref{eq:n9v81-moving-reaction}.
\end{enumerate}
\end{definition}

\begin{theorem}[CPCal gives coupled calibrated propagation]
\label{thm:n9v81-CPCal}
Under CPCal, constraint evolution along the moving reference has the
exact continuous Duhamel formula \eqref{eq:n9v81-Duhamel} and the exact
discrete recursion \eqref{eq:n9v81-discrete-Duhamel}.  The reference
connection is counted once, its cofinal contribution is bounded by
\eqref{eq:n9v81-two-rate}, and the second calibration response enters
the V76 reaction ledger explicitly.  If the strict reaction and
propagation gates hold, the local reduced Hessian stays positive and
the calibrated stable constraint error tends to zero.
\end{theorem}

\begin{proof}
Use \Cref{thm:n9v81-moving-propagation,cor:n9v81-Duhamel,
thm:n9v81-discrete-recursion,prop:n9v81-two-rate,
prop:n9v81-D2k,prop:n9v81-moving-reaction}.  The last convergence claim
uses the strict contraction hypotheses in CPCal and does not apply to
the neutral-sector firewall \ref{fire:n9v81-neutral}.
\end{proof}

\begin{theorem}[Ninth-series V81 result]
\label{thm:n9v81-result}
Along a joint output level set, the reference velocity is
$\dot p=A^{-1}(\dot z_*-B\dot q)$ and every regulated quantity has the
exact aligned derivative
$D_qH-D_pHA^{-1}B$.  Applied to the constraint, this produces the
calibration forcing \eqref{eq:n9v81-fcal} and the Duhamel formula
\eqref{eq:n9v81-Duhamel}.  Across discrete cutoff steps, the exact
connection increment is the line integral
\eqref{eq:n9v81-finite-connection}; two contracting rates give the
explicit convolution \eqref{eq:n9v81-Xi}.  The second calibration
response is \eqref{eq:n9v81-D2k} and must be included in the reaction
bound.
\end{theorem}

\begin{proof}
Combine \Cref{lem:n9v81-reference-velocity,
thm:n9v81-chain-rule,thm:n9v81-moving-propagation,
lem:n9v81-finite-connection,prop:n9v81-two-rate,
prop:n9v81-D2k}.
\end{proof}

\begin{assessment}[Progress at ninth-series V81]
\label{ass:n9v81-status}
Calibration and constraint propagation now form one exact system.  The
moving reference produces a measurable connection forcing rather than
an implicit change of scheme, and its discrete convolution decays when
both the stable constraint sector and reference mismatch contract.  A
neutral constraint sector still requires exact cancellation.  The next
bottleneck is to construct a regulator in which the finite-cutoff Ward
identity, calibration chart, and reflection-positive transfer use the
same gauge/ghost measure.
\end{assessment}

\begin{programme}[Ninth-series V82: common regulated measure]
\label{prog:n9v81-V82}
V82 will formulate a single-measure compatibility theorem for
constraint coarea, gauge fixing, ghost determinants, joint calibration,
and OS transfer.  It will identify which determinant factors are
Jacobians, which arise from Gaussian integration, and which enter the
graded fermionic measure, with no double counting.
\end{programme}

\begin{firewall}[Claim boundary after ninth-series V81]
\label{fire:n9v81-boundary}
V81 does not construct the model-specific calibration chart, prove
weighted locality of its inverse, derive the physical Ward/constraint
identity, establish exact neutral-sector cancellation, build the common
gauge/ghost measure, prove anomaly cancellation, infrared closure,
regulator independence, OS reconstruction, or the full SM--Einstein
continuum.
\end{firewall}

\paragraph{Parent-version record.}
V81 was formed from
\texttt{Renewal\_SM\_Einstein\_Continuum\_Research\_Notes\_V80.tex}.
The V80 parent had $31\,232$ counted source lines, $1\,124\,892$ bytes,
and SHA--256
\[
 \texttt{CCC6AF206C459EA1FCC8E12D5A59AC902C7B9093FAAA6A1DB665EF8E44634E89}.
\]

% =============================================================================
% END APPENDED NINTH-SERIES V81 MATERIAL
% =============================================================================

% =============================================================================
% BEGIN APPENDED NINTH-SERIES V82 MATERIAL
% =============================================================================

\section{Ninth-series V82: one regulated measure and the coupled Jacobian}
\label{sec:v9v82}

V74 and V76 distinguished constraint coarea from stationary Gaussian
integration.  V77--V81 added gauge fixing, ghosts, Ward rows, and a
moving reference.  These pieces must now be placed in one regulated
measure.  The first obstruction is already finite dimensional: when a
constraint and a gauge condition have cross derivatives, their
Jacobian is a block determinant.  The correct gauge operator after
constraint solving is a Schur complement.

\subsection{Combined constraint and gauge coordinates}

Let $U$, $A$, and $Y$ be finite-dimensional oriented real Hilbert
spaces.  Coordinates are $(u,a,y)$, where $u\in U$ is eliminated by a
constraint, $a\in A$ parametrizes a local gauge orbit, and $y\in Y$ is
retained.  Let
\begin{equation}
 F:U\times A\times Y\to U,
 \qquad
 \chi:U\times A\times Y\to A
 \label{eq:n9v82-F-chi}
\end{equation}
be $C^2$.  At a simultaneous regular root define
\begin{equation}
 L:=F_u,
 \quad R:=F_a,
 \quad S:=\chi_u,
 \quad M:=\chi_a,
 \quad
 \mathbb J:=
 \begin{pmatrix}L&R\\S&M\end{pmatrix}.
 \label{eq:n9v82-block-J}
\end{equation}

\begin{theorem}[Constraint-first gauge Jacobian]
\label{thm:n9v82-Schur-Jacobian}
Suppose $L$ is invertible.  The implicit constraint root
$u=u(a,y)$ satisfies
\begin{equation}
 u_a=-L^{-1}R.
 \label{eq:n9v82-ua}
\end{equation}
The derivative of the restricted gauge condition
$\chi_{\rm red}(a,y):=\chi(u(a,y),a,y)$ is
\begin{equation}
 M_{\rm red}:=D_a\chi_{\rm red}
 =M-SL^{-1}R.
 \label{eq:n9v82-Mred}
\end{equation}
Moreover,
\begin{equation}
 \det\mathbb J=(\det L)(\det M_{\rm red}).
 \label{eq:n9v82-block-det}
\end{equation}
Thus the Faddeev--Popov operator on the constrained branch is
$M_{\rm red}$, not $M$, unless the cross term $SL^{-1}R$ vanishes.
\end{theorem}

\begin{proof}
Differentiate $F(u(a,y),a,y)=0$ with respect to $a$ to obtain
$Lu_a+R=0$.  The chain rule for $\chi_{\rm red}$ then gives
\eqref{eq:n9v82-Mred}.  Finally multiply $\mathbb J$ on the left by
the unit-determinant block matrix
\[
 \begin{pmatrix}I&0\\-SL^{-1}&I\end{pmatrix}
\]
to obtain the upper triangular block matrix
$\left(\begin{smallmatrix}L&R\\0&M-SL^{-1}R\end{smallmatrix}\right)$.
Taking determinants proves \eqref{eq:n9v82-block-det}.
\end{proof}

\begin{corollary}[Regularity equivalence]
\label{cor:n9v82-regularity}
Under $L$ invertible, the combined map $(u,a)\mapsto(F,\chi)$ is
regular if and only if $M_{\rm red}$ is invertible.
\end{corollary}

\begin{proof}
Use \eqref{eq:n9v82-block-det}.
\end{proof}

\begin{firewall}[Separate determinants require triangularity]
\label{fire:n9v82-separate-determinants}
The product $(\det L)(\det M)$ is generally not the combined
Jacobian.  Replacing $M_{\rm red}$ by $M$ requires a proof that
$SL^{-1}R=0$, or a controlled determinant comparison.  Constraint
solving can change the gauge orbit derivative even at one finite
cutoff.
\end{firewall}

\subsection{Double coarea and the three normalization choices}

Assume for simplicity that a chart contains one simultaneous root
$(u(a,y),a(y))$ for each $y$ and that $\mathbb J$ is invertible there.

\begin{theorem}[Exact double-delta reduction]
\label{thm:n9v82-double-delta}
For every compactly supported test density $H$ in the chart,
\begin{equation}
 \int H(u,a,y)\delta(F(u,a,y))\delta(\chi(u,a,y))
 \,du\,da
 =\frac{H(u(y),a(y),y)}{|\det\mathbb J(y)|}.
 \label{eq:n9v82-double-delta}
\end{equation}
Equivalently, reducing first by $F$ and then by $\chi_{\rm red}$ gives
the factors
\begin{equation}
 |\det L|^{-1}|\det M_{\rm red}|^{-1}.
 \label{eq:n9v82-two-inverse-factors}
\end{equation}
\end{theorem}

\begin{proof}
Use $(F,\chi)$ as local coordinates in $(u,a)$; the absolute
change-of-variables Jacobian is $|\det\mathbb J|$.  The second form
follows from \eqref{eq:n9v82-block-det}.
\end{proof}

\begin{proposition}[Gauge quotient and direct pushforward choices]
\label{prop:n9v82-three-choices}
Starting from \eqref{eq:n9v82-double-delta}, the following are distinct
regulated measures.
\begin{enumerate}
\item The raw double-delta measure retains the full inverse factor
$|\det\mathbb J|^{-1}$.
\item Inserting the positive local gauge quotient factor
$|\det M_{\rm red}|$ cancels only the gauge-coordinate coarea and leaves
$|\det L|^{-1}$.
\item Inserting both $|\det M_{\rm red}|$ and $|\det L|$ produces the
direct deterministic pushforward density $H(u(y),a(y),y)$.
\end{enumerate}
None of these choices is an algebraic consequence of the others; the
starting measure determines which one is correct.
\end{proposition}

\begin{proof}
Multiply the right side of
\eqref{eq:n9v82-two-inverse-factors} by the stated inserted factors.
\end{proof}

\begin{firewall}[Absolute value and orientation]
\label{fire:n9v82-absolute-value}
Bosonic delta coarea produces absolute determinants.  Oriented
Faddeev--Popov and ghost formulas produce signed or complex
determinants.  They coincide with the positive quotient factor only on
a chart where the determinant sign and phase are fixed and the chosen
orientation is compatible.  Crossing a zero is a singular orbit or
branch event, not a harmless change of notation.
\end{firewall}

\subsection{Gaussian, ghost, and fermion factors}

The remaining standard determinant mechanisms have different powers
and algebraic types.

\begin{proposition}[Finite-dimensional determinant dictionary]
\label{prop:n9v82-dictionary}
Let $D_b>0$ act on an $n$-dimensional real bosonic space, let
$M_g$ act on an $r$-dimensional ghost space, and let $D_f$ act on a
finite complex fermion space.  With fixed orientation conventions,
\begin{align}
 \int_{\mathbb R^n}e^{-\frac12\langle x,D_bx\rangle}\,dx
 &=(2\pi)^{n/2}(\det D_b)^{-1/2},
 \label{eq:n9v82-boson-det}\\
 \int e^{-\bar cM_gc}\,d\bar c\,dc
 &=\det M_g,
 \label{eq:n9v82-ghost-det}\\
 \int e^{-\bar\psi D_f\psi}\,d\bar\psi\,d\psi
 &=\det D_f.
 \label{eq:n9v82-fermion-det}
\end{align}
The first is a positive inverse square root.  The latter two are
Berezin determinants and need not be positive real numbers.
\end{proposition}

\begin{proof}
Orthogonally diagonalize $D_b$ and multiply the one-dimensional
Gaussian integrals.  Expand each Grassmann exponential; only the term
containing every generator survives Berezin integration, and its
coefficient is the determinant.
\end{proof}

\begin{firewall}[Mutually exclusive elimination mechanisms]
\label{fire:n9v82-exclusive}
If the same conformal variable $u$ is fixed by $\delta(F)$, it is not
also freely integrated as the Gaussian in
\eqref{eq:n9v82-boson-det}.  Multiplying
$|\det L|^{-1}$ by $(\det D_b)^{-1/2}$ for that same eliminated row
double counts unless the regulated construction contains two genuinely
independent variables.  Ghost and fermion determinants likewise require
their declared Berezin fields in the common measure.
\end{firewall}

\begin{firewall}[Graded positivity is separate]
\label{fire:n9v82-graded}
Equations \eqref{eq:n9v82-ghost-det} and
\eqref{eq:n9v82-fermion-det} do not define positive scalar weights in
general.  Fermionic OS positivity is a graded statement about the
reflection algebra and covariance.  It cannot be replaced by positivity
of a determinant, especially for a chiral phase.
\end{firewall}

\subsection{Aligned determinant response}

Let $K(p,q)$ be any invertible matrix family of constant determinant
sign on a V80 calibration chart $p=\mathfrak k(q)$.  Put
\begin{equation}
 \dot K_{\rm al}[h]
 :=K_q[h]-K_p[A^{-1}Bh].
 \label{eq:n9v82-K-aligned}
\end{equation}

\begin{theorem}[First aligned log-determinant response]
\label{thm:n9v82-aligned-logdet}
For a fixed calibration target,
\begin{equation}
 D_q\log|\det K(\mathfrak k(q),q)|[h]
 =\operatorname{Tr}\bigl(K^{-1}\dot K_{\rm al}[h]\bigr).
 \label{eq:n9v82-aligned-logdet}
\end{equation}
For the combined constraint/gauge Jacobian this response can be written
equivalently as
\begin{equation}
 D\log|\det\mathbb J|
 =D\log|\det L|+D\log|\det M_{\rm red}|,
 \label{eq:n9v82-split-response}
\end{equation}
provided every derivative is the same aligned derivative.
\end{theorem}

\begin{proof}
Jacobi's formula gives $D\log|\det K|=\operatorname{Tr}(K^{-1}DK)$
on a constant-sign chart.  Insert the V81 aligned derivative
\eqref{eq:n9v81-cal-derivative}.  The split form follows by
differentiating \eqref{eq:n9v82-block-det}.
\end{proof}

\begin{proposition}[Second aligned determinant response]
\label{prop:n9v82-second-logdet}
Let $K_h$ and $K_{hk}$ denote the first and second derivatives of the
composite $K(\mathfrak k(q),q)$, with
$D\mathfrak k$ and $D^2\mathfrak k$ given by V80--V81.  Then
\begin{equation}
 D^2\log|\det K|[h,k]
 =\operatorname{Tr}\bigl(
 K^{-1}K_{hk}-K^{-1}K_kK^{-1}K_h
 \bigr).
 \label{eq:n9v82-second-logdet}
\end{equation}
\end{proposition}

\begin{proof}
Differentiate \eqref{eq:n9v82-aligned-logdet} and use
$D(K^{-1})[k]=-K^{-1}K_kK^{-1}$.
\end{proof}

\begin{firewall}[Connection response is part of the determinant owner]
\label{fire:n9v82-det-connection}
The terms containing $A^{-1}B$ and $D^2\mathfrak k$ belong to the same
log-determinant owner as the fixed-reference trace.  Adding a separate
``calibration determinant correction'' after using
\eqref{eq:n9v82-aligned-logdet} would count the connection twice.
\end{firewall}

\subsection{Reflection-compatible scalar reweighting}

Even a positive determinant weight need not preserve OS positivity.
The sufficient structure is reflection factorization.  Let $\mu$ be an
OS-positive bosonic measure, let $\Theta$ be reflection, and let
$\mathcal A_+$ be the positive-time observable algebra.

\begin{theorem}[Reflection-factorized reweighting]
\label{thm:n9v82-OS-reweight}
Suppose a nonnegative integrable weight has the form
\begin{equation}
 w=(\Theta a)a
 \label{eq:n9v82-factorized-weight}
\end{equation}
for some $a\in\mathcal A_+$, and
$0<\int w\,d\mu<\infty$.  Then the normalized measure
\begin{equation}
 d\mu_w:=\frac{w\,d\mu}{\int w\,d\mu}
 \label{eq:n9v82-muw}
\end{equation}
is OS positive.
\end{theorem}

\begin{proof}
For $F\in\mathcal A_+$,
\begin{equation}
 \int(\Theta F)F\,d\mu_w
 =\frac{\int\Theta(aF)(aF)\,d\mu}{\int(\Theta a)a\,d\mu}
 \ge0.
\end{equation}
The numerator is nonnegative by OS positivity of $\mu$, and the
denominator is positive by hypothesis.
\end{proof}

\begin{corollary}[Spatial-slice determinant criterion]
\label{cor:n9v82-slice-factorization}
If a positive scalar Jacobian weight is a product of independent
spatial-slice factors paired by time reflection, with any reflection-
plane factor admitting a positive square root in the boundary algebra,
then it has the form \eqref{eq:n9v82-factorized-weight} and preserves
the bosonic OS form.
\end{corollary}

\begin{proof}
Take $a$ to be the product of all positive-time slice factors and the
positive square root of the reflection-plane factor.  Reflection gives
the negative-time product.
\end{proof}

\begin{firewall}[Positivity of a determinant is insufficient]
\label{fire:n9v82-positive-not-OS}
A general positive nonlocal function $w$ need not be
reflection-factorized, so it need not preserve the OS cone.  For the
constraint and gauge Jacobians, factorization or a positive joint
transfer-kernel representation must be proved.  Absolute values alone
do not supply that proof.
\end{firewall}

\subsection{One-measure ownership card}

\begin{definition}[OMOC]
\label{def:n9v82-OMOC}
The one-measure ownership card requires:
\begin{enumerate}
\item one explicit finite-regulator field, ghost, and multiplier
      measure;
\item a declared order or combined treatment of $F=0$ and $\chi=0$;
\item the block Jacobian $\mathbb J$ and constrained gauge operator
      $M_{\rm red}$;
\item a choice among the raw coarea, gauge quotient, and direct
      pushforward normalizations;
\item Gaussian determinants only for independently integrated bosonic
      variables;
\item oriented ghost and fermion determinants with their phases and
      zero modes retained;
\item first and second aligned responses of every determinant;
\item exactly one owner for every Jacobian and calibration connection;
\item reflection factorization or a positive joint transfer kernel for
      every bosonic scalar reweighting; and
\item a graded reflection/anomaly theorem for the ghost and chiral
      fermion sectors before the continuum limit.
\end{enumerate}
\end{definition}

\begin{theorem}[OMOC prevents determinant double counting]
\label{thm:n9v82-OMOC}
Under OMOC, sequential constraint and gauge reduction agrees with the
combined Jacobian formula \eqref{eq:n9v82-block-det}; each Gaussian,
coarea, ghost, and fermion determinant is tied to a distinct integration
operation; and moving-reference responses are the aligned traces
\eqref{eq:n9v82-aligned-logdet}--
\eqref{eq:n9v82-second-logdet}.  If the bosonic weights satisfy
\eqref{eq:n9v82-factorized-weight}, their insertion preserves the
bosonic OS cone.
\end{theorem}

\begin{proof}
Use \Cref{thm:n9v82-Schur-Jacobian,thm:n9v82-double-delta,
prop:n9v82-three-choices,prop:n9v82-dictionary,
thm:n9v82-aligned-logdet,prop:n9v82-second-logdet,
thm:n9v82-OS-reweight}.  The graded sector remains subject to item 10
and is not inferred from the bosonic theorem.
\end{proof}

\begin{theorem}[Ninth-series V82 result]
\label{thm:n9v82-result}
At a simultaneous regular constraint/gauge root, the exact combined
Jacobian factors as
\[
 \det\mathbb J=\det L\,
 \det(M-SL^{-1}R).
\]
Thus constraint solving changes the gauge operator by a Schur term.
Double-delta coarea, gauge quotient, direct pushforward, Gaussian
integration, ghost integration, and fermion integration carry distinct
determinant powers and signs.  Along the V80 calibration chart their
responses use the V81 aligned derivative.  A positive bosonic
determinant reweighting preserves OS positivity under the explicit
reflection factorization \eqref{eq:n9v82-factorized-weight}, not from
positivity alone.
\end{theorem}

\begin{proof}
Combine \Cref{thm:n9v82-Schur-Jacobian,
thm:n9v82-double-delta,prop:n9v82-dictionary,
thm:n9v82-aligned-logdet,thm:n9v82-OS-reweight}.
\end{proof}

\begin{assessment}[Progress at ninth-series V82]
\label{ass:n9v82-status}
The measure ledger now has a single finite-regulator algebra.  Its new
substantive point is the constrained Faddeev--Popov Schur operator
$M_{\rm red}$; separate lapse and gauge determinants are valid only
under triangularity or after this correction.  The remaining hard
step is the graded one: construct a reflection-compatible gauge/ghost/
chiral boundary transfer and prove its anomaly and zero-mode identities
uniformly with the same regulator.
\end{assessment}

\begin{programme}[Ninth-series V83: graded boundary transfer]
\label{prog:n9v82-V83}
V83 will formulate the finite-dimensional Grassmann reflection form,
derive the Schur complement for integrating an interior fermion block,
and state a checkable criterion under which gauge ghosts cancel only the
declared orbit Jacobian while the physical chiral determinant retains
its separate phase and anomaly owner.
\end{programme}

\begin{firewall}[Claim boundary after ninth-series V82]
\label{fire:n9v82-boundary}
V82 does not prove global gauge-slice regularity, constant determinant
sign, reflection factorization of the physical Jacobians, graded OS
positivity, chiral anomaly cancellation, uniform determinant response,
infrared closure, regulator independence, OS reconstruction, or the
full SM--Einstein continuum.
\end{firewall}

\paragraph{Parent-version record.}
V82 was formed from
\texttt{Renewal\_SM\_Einstein\_Continuum\_Research\_Notes\_V81.tex}.
The V81 parent had $31\,849$ counted source lines, $1\,144\,944$ bytes,
and SHA--256
\[
 \texttt{3AE57FBCA594C0D612653E5F8071E901E18A7DDA09FBA755D6AA7AF2A9681737}.
\]

% =============================================================================
% END APPENDED NINTH-SERIES V82 MATERIAL
% =============================================================================

% =============================================================================
% BEGIN APPENDED NINTH-SERIES V83 MATERIAL
% =============================================================================

\section{Ninth-series V83: graded Schur transfer and ghost ownership}
\label{sec:v9v83}

V82 placed the bosonic constraint and gauge Jacobians in one regulated
measure but left the graded sector open.  This note works entirely at a
finite cutoff.  It proves the exact Grassmann Schur complement, gives a
necessary and sufficient positivity criterion for the resulting
finite-dimensional Gaussian reflection form, and identifies the sign
left over when an oriented ghost determinant is paired with an absolute
bosonic gauge coarea factor.

\subsection{Interior Grassmann elimination}

Let $V_b$ and $V_i$ be finite-dimensional complex spaces of boundary
and interior fermion variables.  With independent Grassmann generators
$\psi_b,\bar\psi_b,\psi_i,\bar\psi_i$, let
\begin{equation}
 D=
 \begin{pmatrix}
  D_{bb}&D_{bi}\\D_{ib}&D_{ii}
 \end{pmatrix}
 \label{eq:n9v83-D-block}
\end{equation}
and suppose $D_{ii}$ is invertible.  Define
\begin{equation}
 D_{\partial}:=D_{bb}-D_{bi}D_{ii}^{-1}D_{ib}.
 \label{eq:n9v83-D-boundary}
\end{equation}

\begin{theorem}[Exact fermionic Schur integration]
\label{thm:n9v83-fermion-Schur}
With a fixed Berezin orientation,
\begin{align}
 &\int
 \exp\left[-
 \begin{pmatrix}\bar\psi_b&\bar\psi_i\end{pmatrix}
 D
 \binom{\psi_b}{\psi_i}
 \right]d\bar\psi_i\,d\psi_i
 \notag\\
 &\qquad
 =\det D_{ii}\,
   \exp(-\bar\psi_bD_{\partial}\psi_b).
 \label{eq:n9v83-Schur-integral}
\end{align}
Moreover,
\begin{equation}
 \det D=(\det D_{ii})(\det D_{\partial}).
 \label{eq:n9v83-Schur-det}
\end{equation}
\end{theorem}

\begin{proof}
Make the unit-Berezinian triangular translation
\begin{equation}
 \psi_i'=\psi_i+D_{ii}^{-1}D_{ib}\psi_b,
 \qquad
 \bar\psi_i'=\bar\psi_i+\bar\psi_bD_{bi}D_{ii}^{-1}.
 \label{eq:n9v83-Grassmann-shift}
\end{equation}
Direct expansion gives
\begin{equation}
 \bar\psi D\psi
 =\bar\psi_i'D_{ii}\psi_i'
  +\bar\psi_bD_{\partial}\psi_b.
\end{equation}
The interior Berezin integral is $\det D_{ii}$, proving
\eqref{eq:n9v83-Schur-integral}.  The determinant identity follows from
the same block triangular elimination, or by integrating all variables
and comparing with $\det D$.
\end{proof}

\begin{corollary}[Boundary covariance]
\label{cor:n9v83-boundary-covariance}
If $D$ and $D_{\partial}$ are invertible, the boundary block of the
full covariance is
\begin{equation}
 P_bD^{-1}P_b^*=D_{\partial}^{-1}.
 \label{eq:n9v83-boundary-covariance}
\end{equation}
\end{corollary}

\begin{proof}
Solve
$D\binom{x_b}{x_i}=\binom{f_b}{0}$.  The second block gives
$x_i=-D_{ii}^{-1}D_{ib}x_b$; the first becomes
$D_{\partial}x_b=f_b$.
\end{proof}

\begin{firewall}[The Schur operator need not inherit positivity]
\label{fire:n9v83-Schur-not-positive}
The algebra above requires only invertibility of $D_{ii}$.  A Euclidean
Dirac operator is generally not a positive self-adjoint matrix, and
$D_{\partial}$ need not be positive.  Positivity of
$\det D_{ii}$ or $\det D_{\partial}$ does not prove positivity of the
graded reflection form.
\end{firewall}

\subsection{A finite exterior-power reflection criterion}

Let $V_+$ be a finite-dimensional complex one-particle space and
\begin{equation}
 \mathcal F_+:=\bigoplus_{k=0}^{\dim V_+}\Lambda^kV_+
 \label{eq:n9v83-Fock}
\end{equation}
its fermionic Fock space with the standard exterior-power inner
product.  Fix the graded reflection convention so that the reflected
Gaussian pairing associated with a matrix $K:V_+\to V_+$ is
\begin{equation}
 \mathfrak B_K(F,G)
 :=\sum_{k=0}^{\dim V_+}
 \langle f_k,(\Lambda^kK)g_k\rangle_{\Lambda^kV_+},
 \label{eq:n9v83-BK}
\end{equation}
where $F\leftrightarrow(f_0,\ldots,f_n)$ and similarly for $G$.
Equivalently, on degree-$k$ monomials the matrix element is the
determinant of the corresponding $k\times k$ submatrix of $K$.  The
fixed reflection twist is part of this definition and removes the
order-reversal sign ambiguity.

\begin{theorem}[Graded Gaussian reflection criterion]
\label{thm:n9v83-graded-criterion}
The form $\mathfrak B_K$ is positive semidefinite,
\begin{equation}
 \mathfrak B_K(F,F)\ge0\quad(F\in\mathcal F_+),
 \label{eq:n9v83-graded-positive}
\end{equation}
if and only if $K$ is positive semidefinite on $V_+$.
\end{theorem}

\begin{proof}
If $K\ge0$, choose an orthonormal eigenbasis with eigenvalues
$\lambda_j\ge0$.  On the exterior basis indexed by subsets $I$,
$\Lambda^kK$ has eigenvalues $\prod_{j\in I}\lambda_j$, all
nonnegative.  Every summand in \eqref{eq:n9v83-BK} is therefore
positive semidefinite.  Conversely, restrict
\eqref{eq:n9v83-graded-positive} to the degree-one sector; it becomes
$\langle v,Kv\rangle\ge0$ for every $v\in V_+$.
\end{proof}

\begin{corollary}[Strict and degenerate sectors]
\label{cor:n9v83-strict-degenerate}
If $K>0$, the reflection form is strictly positive on every exterior
sector.  If $K$ has a kernel, every wedge containing a vector in that
kernel is null.  The physical Hilbert space is obtained only after
quotienting these null vectors and completing.
\end{corollary}

\begin{proof}
Use the eigenvalue products in the proof of
\Cref{thm:n9v83-graded-criterion}.
\end{proof}

\begin{proposition}[Positive fermionic transfer from one-particle data]
\label{prop:n9v83-second-quantization}
If $0\le T_1\le I$ on $V_+$, then its fermionic second quantization
\begin{equation}
 \Gamma_f(T_1):=\bigoplus_{k=0}^{\dim V_+}\Lambda^kT_1
 \label{eq:n9v83-second-quantization}
\end{equation}
is a positive contraction on $\mathcal F_+$.  If $W$ is any bounded
operator on $\mathcal F_+$, then
\begin{equation}
 W^*\Gamma_f(T_1)W\ge0.
 \label{eq:n9v83-sandwich}
\end{equation}
\end{proposition}

\begin{proof}
In an eigenbasis of $T_1$, every eigenvalue of
$\Lambda^kT_1$ is a product of numbers in $[0,1]$, hence lies in
$[0,1]$.  The sandwich inequality follows from
$\langle v,W^*\Gamma_f(T_1)Wv\rangle=
\langle Wv,\Gamma_f(T_1)Wv\rangle\ge0$.
\end{proof}

\begin{firewall}[A boundary Schur complement is not yet a transfer]
\label{fire:n9v83-Schur-not-transfer}
To apply \Cref{thm:n9v83-graded-criterion}, one must derive the
cross-reflection matrix $K$ from the boundary covariance
$D_{\partial}^{-1}$ with the declared reflection and chirality maps.
The identity \eqref{eq:n9v83-boundary-covariance} alone does not show
$K\ge0$ or exhibit a one-particle contraction $T_1$.
\end{firewall}

\subsection{Zero modes and retained Grassmann variables}

Suppose now that $D_{ii}$ is singular and let
$Z_i=\ker D_{ii}$.  Choose a complementary subspace $Z_i^\perp$ only
after declaring the regulator inner product.

\begin{proposition}[Zero-mode saturation]
\label{prop:n9v83-zero-modes}
If $D_{ii}$ has a nonzero right or left zero mode, the uninserted
interior Berezin integral vanishes:
\begin{equation}
 \int e^{-\bar\psi_iD_{ii}\psi_i}
 \,d\bar\psi_i\,d\psi_i=\det D_{ii}=0.
 \label{eq:n9v83-zero-integral}
\end{equation}
A nonzero correlation requires insertions which saturate every
unpaired zero-mode generator.  Consequently $D_{ii}^{-1}$ in
\eqref{eq:n9v83-D-boundary} cannot be replaced silently by a
pseudoinverse; the zero modes must be retained as explicit boundary,
topological, or infrared Grassmann variables.
\end{proposition}

\begin{proof}
Choose bases whose first generator is a right or left zero direction.
The exponent contains no term capable of supplying the corresponding
top-degree Grassmann factor, so the coefficient selected by Berezin
integration is zero.  Insertions can supply it, but then it remains an
explicit observable rather than part of an inverse covariance.
\end{proof}

\begin{firewall}[Index sectors]
\label{fire:n9v83-index}
For a chiral operator, unequal left and right zero-mode dimensions are
index data.  Removing them to make a square invertible interior block
changes the determinant line and its gauge transformation.  They must
remain in the anomaly and topological-sector ledger.
\end{firewall}

\subsection{Exact ghost cancellation ledger}

Return to the constrained gauge operator
\begin{equation}
 M_{\rm red}=M-SL^{-1}R
 \label{eq:n9v83-Mred}
\end{equation}
from V82.  Bosonic gauge-coordinate coarea contributes
$|\det M_{\rm red}|^{-1}$, whereas a complex ghost pair contributes
$\det M_{\rm red}$.

\begin{theorem}[Oriented ghost/coarea product]
\label{thm:n9v83-ghost-product}
On an invertible real gauge chart,
\begin{equation}
 \frac{\det M_{\rm red}}{|\det M_{\rm red}|}
 =\operatorname{sgn}\det M_{\rm red}
 \label{eq:n9v83-ghost-sign}
\end{equation}
is the exact product of the ghost determinant and the bosonic gauge
coarea factor.  The product equals $+1$ precisely on an orientation
chart where $\det M_{\rm red}>0$.  For a complex operator the residual
is the phase
$\det M_{\rm red}/|\det M_{\rm red}|$.
\end{theorem}

\begin{proof}
Multiply the two factors.  The stated sign and phase characterizations
are the definitions of the real sign and complex phase of a nonzero
determinant.
\end{proof}

\begin{corollary}[No second Faddeev--Popov insertion]
\label{cor:n9v83-no-second-FP}
If the ghost Berezin integral already supplies
$\det M_{\rm red}$, inserting another determinant with the bosonic
delta gauge condition produces an extra owner.  Exact cancellation of
the gauge-coordinate inverse uses one ghost determinant, not two.
\end{corollary}

\begin{proof}
The exact product is \eqref{eq:n9v83-ghost-sign}; multiplying by another
determinant changes it by that nonconstant factor.
\end{proof}

\begin{firewall}[Gribov sign]
\label{fire:n9v83-Gribov-sign}
The sign in \eqref{eq:n9v83-ghost-sign} cannot change without a zero of
$M_{\rm red}$.  Such a zero is the local Gribov/orbit singularity where
the coarea chart and ghost inverse both fail.  Taking an absolute value
to hide the sign changes the oriented ghost theory and requires a new
BRST and locality analysis.
\end{firewall}

\subsection{Physical fermions are not ghost owners}

Let $D_{\rm ch}$ be the regulated physical chiral operator.  Its
Berezin integral gives $\det D_{\rm ch}$ on its determinant line.

\begin{proposition}[Separation of determinant lines]
\label{prop:n9v83-determinant-lines}
The product
\begin{equation}
 (\det M_{\rm red})(\det D_{\rm ch})
 \label{eq:n9v83-separate-lines}
\end{equation}
contains two independent determinant owners unless an explicit
intertwining isomorphism identifies their domain, codomain, symmetry
action, boundary condition, and spectrum.  Equality of matrix size or
opposite formal powers is insufficient for cancellation.
\end{proposition}

\begin{proof}
A determinant is a top exterior-power map on the determinant line of
its operator.  Multiplication pairs two such lines but does not identify
either with the inverse of the other.  Cancellation requires an
isomorphism under which one operator is the inverse owner of the other;
the listed data are necessary for that identification.
\end{proof}

For an invertible differentiable family, the infinitesimal phase and
Ward response is governed by
\begin{equation}
 D\log\det D_{\rm ch}[h]
 =\operatorname{Tr}(D_{\rm ch}^{-1}DD_{\rm ch}[h]),
 \label{eq:n9v83-chiral-response}
\end{equation}
with the logarithm understood only on a chosen local determinant-line
trivialization.

\begin{firewall}[Local trace versus global anomaly]
\label{fire:n9v83-global-anomaly}
Equation \eqref{eq:n9v83-chiral-response} is local on the invertible
chart.  It does not trivialize the determinant line around a loop,
control spectral flow through zero, or prove cancellation of local or
global gauge and gravitational anomalies.  Those require a regulator-
compatible statement for the complete SM representation.
\end{firewall}

\subsection{Ghost reflection and the physical quotient}

Ghosts are Grassmann fields, but the physical scalar product is not the
unrestricted ghost Fock product.

\begin{proposition}[Conditional BRST descent]
\label{prop:n9v83-BRST-descent}
Let $\mathcal K$ be a finite-dimensional graded state space with a
nilpotent operator $Q$, $Q^2=0$.  Suppose a transfer $T$ obeys
\begin{equation}
 TQ=QT
 \label{eq:n9v83-TQ}
\end{equation}
and a positive semidefinite form on $\ker Q$ has radical containing
$\operatorname{ran}Q$.  Then $T$ induces a well-defined operator on the
cohomology
\begin{equation}
 H_Q=\ker Q/\operatorname{ran}Q,
 \label{eq:n9v83-cohomology}
\end{equation}
and the form descends to a positive semidefinite form there.
\end{proposition}

\begin{proof}
Commutation maps closed vectors to closed vectors and exact vectors to
exact vectors, so $T$ is well defined on equivalence classes.  If
$v$ is replaced by $v+Qw$, all added inner-product terms vanish because
the exact subspace lies in the radical.  Positivity is inherited from
the representative in $\ker Q$.
\end{proof}

\begin{firewall}[BRST hypotheses remain open]
\label{fire:n9v83-BRST-open}
V83 does not construct a nilpotent finite-cutoff SM--Einstein BRST
charge, prove \eqref{eq:n9v83-TQ}, or identify the exact subspace with
the radical of a reflection form.  The proposition is a gate, not a
claim that ghost positivity has been solved.
\end{firewall}

\subsection{Graded boundary transfer card}

\begin{definition}[GBTC]
\label{def:n9v83-GBTC}
The graded boundary transfer card requires:
\begin{enumerate}
\item a finite Grassmann algebra and fixed Berezin orientation;
\item an invertible interior block after all zero modes are retained;
\item the exact boundary Schur operator $D_{\partial}$;
\item the reflected one-particle cross matrix $K$ derived from
      $D_{\partial}^{-1}$;
\item $K\ge0$, or an equivalent positive one-particle transfer
      representation;
\item a fixed positive-orientation chart for $M_{\rm red}$ or an
      explicit residual sign/phase owner;
\item exactly one ghost determinant for the gauge coarea cancellation;
\item a separate determinant line for every physical chiral species;
\item regulator-compatible local and global anomaly cancellation; and
\item a BRST cohomology descent or another construction of the physical
      graded Hilbert space.
\end{enumerate}
\end{definition}

\begin{theorem}[GBTC supplies a conditional graded transfer]
\label{thm:n9v83-GBTC}
Under GBTC, interior fermion integration produces the exact boundary
operator \eqref{eq:n9v83-D-boundary}; its reflected Gaussian form is
positive by \Cref{thm:n9v83-graded-criterion}; ghost integration cancels
only the declared oriented gauge-coordinate factor; and the physical
chiral determinant and its anomaly remain separately owned.  If the
BRST hypotheses hold, the transfer descends to a positive semidefinite
physical cohomology form.
\end{theorem}

\begin{proof}
Use \Cref{thm:n9v83-fermion-Schur,
thm:n9v83-graded-criterion,thm:n9v83-ghost-product,
prop:n9v83-determinant-lines,prop:n9v83-BRST-descent}.
\end{proof}

\begin{theorem}[Ninth-series V83 result]
\label{thm:n9v83-result}
Finite interior Grassmann integration gives
$D_{\partial}=D_{bb}-D_{bi}D_{ii}^{-1}D_{ib}$ and the factor
$\det D_{ii}$.  The induced finite Gaussian reflection form is positive
exactly when its reflected one-particle cross matrix $K$ is positive;
this is equivalent to positivity of every exterior power.  A ghost
determinant paired with bosonic gauge coarea leaves the exact
orientation sign or phase \eqref{eq:n9v83-ghost-sign}.  Physical chiral
determinants inhabit separate determinant lines and cannot be used as
undeclared ghost cancellations.
\end{theorem}

\begin{proof}
Combine \Cref{thm:n9v83-fermion-Schur,
thm:n9v83-graded-criterion,thm:n9v83-ghost-product,
prop:n9v83-determinant-lines}.
\end{proof}

\begin{assessment}[Progress at ninth-series V83]
\label{ass:n9v83-status}
The graded finite-cutoff algebra is now explicit.  The remaining
positivity question has been reduced to the spectrum of a reflected
one-particle boundary matrix, while zero modes, ghost orientation, and
the physical chiral determinant are kept as separate owners.  The next
step is to determine when Schur elimination preserves the one-particle
reflection cone and to quantify the debit when the interior/boundary
coupling is perturbed.
\end{assessment}

\begin{programme}[Ninth-series V84: reflected Schur margin]
\label{prog:n9v83-V84}
V84 will derive a reflected one-particle Schur-complement margin, bound
its change through resolvent identities, and isolate the zero-crossing
event which forces a retained boundary mode.  The target is a robust
finite-cutoff graded transfer gate suitable for the moving calibration
chart.
\end{programme}

\begin{firewall}[Claim boundary after ninth-series V83]
\label{fire:n9v83-boundary}
V83 does not prove positivity of the physical reflected cross matrix,
a uniform interior fermion gap, absence of determinant-line holonomy,
BRST nilpotence, chiral anomaly cancellation, Gribov closure, infrared
closure, regulator independence, OS reconstruction, or the full
SM--Einstein continuum.
\end{firewall}

\paragraph{Parent-version record.}
V83 was formed from
\texttt{Renewal\_SM\_Einstein\_Continuum\_Research\_Notes\_V82.tex}.
The V82 parent had $32\,288$ counted source lines, $1\,161\,118$ bytes,
and SHA--256
\[
 \texttt{711FEB2E7C68DAEF248BBC51C2BD0C77C5DA1D2098FA0D34A7471E644A23440A}.
\]

% =============================================================================
% END APPENDED NINTH-SERIES V83 MATERIAL
% =============================================================================

% =============================================================================
% BEGIN APPENDED NINTH-SERIES V84 MATERIAL
% =============================================================================

\section{Ninth-series V84: reflected Schur margins and closing modes}
\label{sec:v9v84}

V83 gives an exact criterion: the finite fermionic Gaussian reflection
form is positive precisely when its reflected one-particle cross matrix
$K$ is positive.  This note derives the dependence of $K$ on interior
Schur elimination, proves a quantitative perturbation margin, and
specifies what must happen when the margin closes.  All estimates remain
finite-cutoff statements.

\subsection{Differentiating the interior Schur operator}

Rename the blocks of the V83 fermion operator as
\begin{equation}
 D=
 \begin{pmatrix}
  A&B\\C&E
 \end{pmatrix},
 \qquad
 S:=A-BE^{-1}C,
 \label{eq:n9v84-blocks}
\end{equation}
where $E$ is the interior block and $S=D_{\partial}$ is the boundary
Schur operator.

\begin{lemma}[First Schur response]
\label{lem:n9v84-DS}
For every differentiable variation,
\begin{equation}
 \dot S
 =\dot A-\dot B E^{-1}C-BE^{-1}\dot C
  +BE^{-1}\dot E E^{-1}C.
 \label{eq:n9v84-DS}
\end{equation}
If
\begin{equation}
 \lVert E^{-1}\rVert\le e^{-1},
 \label{eq:n9v84-E-gap}
\end{equation}
then
\begin{align}
 \lVert\dot S\rVert
 \le{}&\lVert\dot A\rVert
 +e^{-1}\lVert\dot B\rVert\lVert C\rVert
 +e^{-1}\lVert B\rVert\lVert\dot C\rVert\\
 &+e^{-2}\lVert B\rVert\lVert\dot E\rVert
       \lVert C\rVert.
 \label{eq:n9v84-DS-bound}
\end{align}
\end{lemma}

\begin{proof}
Differentiate $A-BE^{-1}C$ and use
$\dot E^{-1}=-E^{-1}\dot E E^{-1}$.  The norm estimate follows by
submultiplicativity.
\end{proof}

\begin{proposition}[Finite interior perturbation]
\label{prop:n9v84-finite-Schur}
Let primed blocks be $A'=A+\Delta A$ and similarly for $B,C,E$.
Suppose
\begin{equation}
 \eta_E:=\lVert E^{-1}\Delta E\rVert<1.
 \label{eq:n9v84-etaE}
\end{equation}
Then $E'$ is invertible,
\begin{align}
 \lVert E'^{-1}\rVert
 &\le\frac{\lVert E^{-1}\rVert}{1-\eta_E},
 \label{eq:n9v84-Eprime-inverse}\\
 \lVert E'^{-1}-E^{-1}\rVert
 &\le\frac{\lVert E^{-1}\rVert^2\lVert\Delta E\rVert}
 {1-\eta_E},
 \label{eq:n9v84-E-resolvent}
\end{align}
and the boundary perturbation obeys
\begin{align}
 \lVert S'-S\rVert
 \le{}&\lVert\Delta A\rVert
 +\lVert\Delta B\rVert\lVert E'^{-1}\rVert
   (\lVert C\rVert+\lVert\Delta C\rVert)\\
 &+\lVert B\rVert\lVert E'^{-1}-E^{-1}\rVert
   \lVert C\rVert
 +\lVert B\rVert\lVert E'^{-1}\rVert
   \lVert\Delta C\rVert.
 \label{eq:n9v84-S-finite-bound}
\end{align}
\end{proposition}

\begin{proof}
$E'=E(I+E^{-1}\Delta E)$ and the Neumann series give
\eqref{eq:n9v84-Eprime-inverse}.  The resolvent identity
\[
 E'^{-1}-E^{-1}=-E'^{-1}\Delta E E^{-1}
\]
gives \eqref{eq:n9v84-E-resolvent}.  Finally expand
\begin{align*}
 S'-S={}&\Delta A-Delta B E'^{-1}(C+\Delta C)\\
 &-B(E'^{-1}-E^{-1})C-BE'^{-1}\Delta C
\end{align*}
and take norms.
\end{proof}

\begin{firewall}[Interior invertibility is a physical gate]
\label{fire:n9v84-interior-gap}
The factor $(1-\eta_E)^{-1}$ diverges at an interior zero mode.  Such a
mode cannot remain in the eliminated block.  It must be retained as in
\Cref{prop:n9v83-zero-modes}, with its determinant-line and boundary
quantum numbers intact.
\end{firewall}

\subsection{Boundary covariance and reflection extraction}

Assume $S$ is invertible and put
\begin{equation}
 C_{\partial}:=S^{-1}.
 \label{eq:n9v84-Cpartial}
\end{equation}
Let
\begin{equation}
 \mathfrak R_q:\operatorname{End}(V_b)\to
 \operatorname{End}(V_+)
 \label{eq:n9v84-R}
\end{equation}
be the declared linear extraction which restricts the boundary
covariance to opposite sides of the reflection plane and applies the
fixed spin/chirality reflection map.  Define
\begin{equation}
 K(q):=\mathfrak R_q(C_{\partial}(q)).
 \label{eq:n9v84-K}
\end{equation}
The subscript allows the reflection frame or boundary identification to
move with the background.

\begin{lemma}[Reflected covariance response]
\label{lem:n9v84-DK}
The exact derivative is
\begin{equation}
 \dot K
 =(\dot{\mathfrak R})(C_{\partial})
 -\mathfrak R(C_{\partial}\dot S C_{\partial}).
 \label{eq:n9v84-DK}
\end{equation}
Hence, if
$\lVert C_{\partial}\rVert\le s^{-1}$,
$\lVert\mathfrak R\rVert\le r_0$, and
$\lVert\dot{\mathfrak R}\rVert\le r_1$, then
\begin{equation}
 \lVert\dot K\rVert
 \le r_1s^{-1}+r_0s^{-2}\lVert\dot S\rVert.
 \label{eq:n9v84-DK-bound}
\end{equation}
\end{lemma}

\begin{proof}
Differentiate \eqref{eq:n9v84-K} and use
$\dot C_{\partial}=-C_{\partial}\dot S C_{\partial}$.
The bound follows by submultiplicativity.
\end{proof}

\begin{proposition}[Finite reflected-covariance perturbation]
\label{prop:n9v84-finite-K}
Suppose
\begin{equation}
 \eta_S:=\lVert S^{-1}(S'-S)\rVert<1.
 \label{eq:n9v84-etaS}
\end{equation}
Then
\begin{equation}
 \lVert S'^{-1}-S^{-1}\rVert
 \le\frac{\lVert S^{-1}\rVert^2\lVert S'-S\rVert}
 {1-\eta_S}.
 \label{eq:n9v84-S-resolvent}
\end{equation}
If the extraction also changes from $\mathfrak R$ to
$\mathfrak R'$, then
\begin{align}
 \lVert K'-K\rVert
 \le{}&\lVert\mathfrak R'-\mathfrak R\rVert
       \lVert S'^{-1}\rVert\\
 &+\lVert\mathfrak R\rVert
       \lVert S'^{-1}-S^{-1}\rVert.
 \label{eq:n9v84-K-finite-bound}
\end{align}
\end{proposition}

\begin{proof}
Apply the resolvent argument of
\Cref{prop:n9v84-finite-Schur} to $S'=S+(S'-S)$.  Add and subtract
$\mathfrak R(S'^{-1})$ in $K'-K$ and take norms.
\end{proof}

\begin{firewall}[Moving reflection data]
\label{fire:n9v84-moving-reflection}
The term $(\dot{\mathfrak R})(C_{\partial})$ is present when the metric,
normal, spin frame, chirality map, or boundary subspace changes.  A
fixed-flat-background cross-covariance estimate does not control this
term on a moving Einstein chart.
\end{firewall}

\subsection{The strict reflected margin}

The order relation below applies only after reflection covariance has
made $K$ self-adjoint.

\begin{theorem}[Robust one-particle reflection margin]
\label{thm:n9v84-margin}
Let $K_0=K_0^*$ satisfy
\begin{equation}
 K_0\ge\kappa_0I,
 \qquad\kappa_0>0.
 \label{eq:n9v84-K0-margin}
\end{equation}
If $K=K^*$ and
\begin{equation}
 \lVert K-K_0\rVert\le\delta<\kappa_0,
 \label{eq:n9v84-K-delta}
\end{equation}
then
\begin{equation}
 K\ge(\kappa_0-\delta)I>0.
 \label{eq:n9v84-K-margin}
\end{equation}
Consequently the complete fermionic Fock reflection form has strictly
positive exterior-power operators, with degree-$k$ lower bound
$(\kappa_0-\delta)^k$.
\end{theorem}

\begin{proof}
For every unit vector $v$,
\[
 \langle v,Kv\rangle
 \ge\langle v,K_0v\rangle-lVert K-K_0\rVert
 \ge\kappa_0-\delta.
\]
The exterior-power eigenvalues are products of $k$ eigenvalues of $K$,
as in \Cref{thm:n9v83-graded-criterion}.
\end{proof}

\begin{corollary}[Pathwise margin]
\label{cor:n9v84-path-margin}
Let $K(q_0)\ge\kappa_0I$, let every $K(q)$ be self-adjoint on a convex
ball, and suppose
\begin{equation}
 \lVert DK(q)\rVert\le L_K.
 \label{eq:n9v84-LK}
\end{equation}
Then
\begin{equation}
 K(q)\ge
 \bigl(\kappa_0-L_K\lVert q-q_0\rVert\bigr)I.
 \label{eq:n9v84-path-margin}
\end{equation}
The reflection form remains strict on the ball of radius
$r<\kappa_0/L_K$.
\end{corollary}

\begin{proof}
Integrate $DK$ along the line segment and apply
\Cref{thm:n9v84-margin}.
\end{proof}

\begin{firewall}[Hermiticity is exact, not a small debit]
\label{fire:n9v84-Hermitian}
If $K\ne K^*$, the quadratic expression
$\langle v,Kv\rangle$ need not be real.  Bounding the anti-Hermitian
part by a small number does not prove reflection positivity.  Exact
reflection covariance must make the physical cross matrix Hermitian,
or the reflection structure itself must be corrected.
\end{firewall}

\subsection{Moving-calibration margin}

Along the V80--V81 chart $p=\mathfrak k(q)$, every dot in
\eqref{eq:n9v84-DS} and \eqref{eq:n9v84-DK} must be the aligned
derivative
\begin{equation}
 \dot H_{\rm al}[h]=H_q[h]-H_p[A^{-1}Bh].
 \label{eq:n9v84-aligned-dot}
\end{equation}

\begin{proposition}[Calibrated reflected debit]
\label{prop:n9v84-calibrated-debit}
Assume uniform gaps
$\lVert E^{-1}\rVert\le e^{-1}$ and
$\lVert S^{-1}\rVert\le s^{-1}$, and bounds
\begin{equation}
 \lVert A^{-1}\rVert\le a,
 \qquad\lVert B\rVert\le b.
 \label{eq:n9v84-cal-AB}
\end{equation}
If the fixed-reference derivatives satisfy
$\lVert H_q\rVert\le h_q$ and
$\lVert H_p\rVert\le h_p$ for every block and reflection map occurring
in \eqref{eq:n9v84-DK}, then the aligned derivatives obey
\begin{equation}
 \lVert\dot H_{\rm al}\rVert\le h_q+h_pab.
 \label{eq:n9v84-aligned-block-bound}
\end{equation}
Substitution into \eqref{eq:n9v84-DS-bound} and
\eqref{eq:n9v84-DK-bound} gives an explicit finite constant
$L_K^{\rm al}$.  The strict calibrated reflection gate is
\begin{equation}
 L_K^{\rm al}r<\kappa_0.
 \label{eq:n9v84-calibrated-margin}
\end{equation}
\end{proposition}

\begin{proof}
Equation \eqref{eq:n9v84-aligned-block-bound} is the triangle inequality
applied to \eqref{eq:n9v84-aligned-dot}.  The Schur and covariance
response lemmas express $DK$ as a finite sum of products of the listed
blocks and inverses, so their displayed bounds produce
$L_K^{\rm al}$.  Apply \Cref{cor:n9v84-path-margin}.
\end{proof}

\begin{firewall}[Two inverse gaps]
\label{fire:n9v84-two-gaps}
The interior gap $e$ and boundary Schur gap $s$ are independent.  The
interior determinant can remain nonzero while the boundary effective
operator approaches zero, or conversely.  A uniform graded transfer
estimate requires both, plus the reflected margin $\kappa_0$.
\end{firewall}

\subsection{Closing reflected modes}

Let $K=K^*$ be nonnegative.  For a threshold $\epsilon>0$, define the
spectral projectors
\begin{equation}
 P_{\rm crit}:=\mathbf1_{[0,\epsilon)}(K),
 \qquad
 P_{\rm gap}:=I-P_{\rm crit}.
 \label{eq:n9v84-critical-projectors}
\end{equation}

\begin{proposition}[Exact retained-mode split]
\label{prop:n9v84-retained-modes}
The projectors in \eqref{eq:n9v84-critical-projectors} commute with
$K$ and satisfy
\begin{equation}
 P_{\rm gap}KP_{\rm gap}\ge\epsilon P_{\rm gap}.
 \label{eq:n9v84-gap-sector}
\end{equation}
Every null or closing reflected one-particle mode lies in
$P_{\rm crit}V_+$ and must be retained in the boundary state.  On the
exterior algebra there is the canonical factorization
\begin{equation}
 \Lambda(V_+)cong
 \Lambda(P_{\rm crit}V_+)\widehat\otimes
 \Lambda(P_{\rm gap}V_+),
 \label{eq:n9v84-Fock-factorization}
\end{equation}
and the second factor has reflection margin at least
$\epsilon^k$ in degree $k$.
\end{proposition}

\begin{proof}
The spectral theorem gives commutation and
\eqref{eq:n9v84-gap-sector}.  Exterior algebra over an orthogonal direct
sum is the graded tensor product of the two exterior algebras.  Apply
\Cref{thm:n9v84-margin} to the gapped restriction.
\end{proof}

\begin{theorem}[First-loss criterion]
\label{thm:n9v84-first-loss}
Let $K(t)=K(t)^*$ be continuous and suppose $K(0)>0$.  If positivity
fails at some later time, then at the first failure time $t_*$,
\begin{equation}
 \lambda_{\min}(K(t_*))=0.
 \label{eq:n9v84-first-zero}
\end{equation}
Thus strict positivity cannot change sign without a boundary null mode.
\end{theorem}

\begin{proof}
The smallest eigenvalue of a finite-dimensional self-adjoint matrix is
continuous in operator norm.  It starts positive.  At the boundary of
the open positive cone it equals zero.
\end{proof}

\begin{firewall}[A null reflection mode is not discarded]
\label{fire:n9v84-null-mode}
Quotienting a null vector is legitimate only after the complete OS
form is constructed and its radical is shown invariant under the
transfer and observables.  Before that proof, a closing mode is retained
as boundary/infrared data; it is not removed by a pseudoinverse.
\end{firewall}

\subsection{Reflected Schur margin card}

\begin{definition}[RSMC]
\label{def:n9v84-RSMC}
The reflected Schur margin card requires:
\begin{enumerate}
\item an invertible interior fermion block with gap $e$ after explicit
      zero-mode retention;
\item the exact Schur derivative \eqref{eq:n9v84-DS};
\item an invertible boundary Schur operator with gap $s$;
\item a typed reflected extraction $\mathfrak R$ including its metric
      and frame derivatives;
\item exact Hermiticity of the cross matrix $K$;
\item a reference margin $K_0\ge\kappa_0I$;
\item the calibrated derivative debit $L_K^{\rm al}$;
\item the strict gate \eqref{eq:n9v84-calibrated-margin};
\item spectral retention of every mode below the declared threshold;
      and
\item compatibility of the retained critical Fock factor with ghosts,
      chirality, anomalies, and the interacting boundary transfer.
\end{enumerate}
\end{definition}

\begin{theorem}[RSMC supplies robust finite-cutoff graded positivity]
\label{thm:n9v84-RSMC}
Under RSMC, interior elimination gives a differentiable boundary
covariance, the reflected one-particle matrix remains strictly positive
on the calibrated ball, and the induced Gaussian Fock reflection form
is positive by V83.  If the margin approaches the threshold, the exact
split \eqref{eq:n9v84-Fock-factorization} retains all closing modes and
preserves a uniform margin on the complementary factor.
\end{theorem}

\begin{proof}
Use \Cref{lem:n9v84-DS,lem:n9v84-DK,
thm:n9v84-margin,prop:n9v84-calibrated-debit,
prop:n9v84-retained-modes,thm:n9v83-graded-criterion}.
\end{proof}

\begin{theorem}[Ninth-series V84 result]
\label{thm:n9v84-result}
The boundary fermion Schur operator has response
\eqref{eq:n9v84-DS}, its reflected cross covariance has response
\eqref{eq:n9v84-DK}, and a reference margin
$K_0\ge\kappa_0I$ survives every self-adjoint perturbation of norm less
than $\kappa_0$.  Along the moving calibration chart the connection
changes every block derivative by \eqref{eq:n9v84-aligned-dot}.  A loss
of strict positivity must pass through a zero reflected mode, which is
then retained by \eqref{eq:n9v84-critical-projectors} rather than
inverted.
\end{theorem}

\begin{proof}
Combine \Cref{lem:n9v84-DS,lem:n9v84-DK,
thm:n9v84-margin,prop:n9v84-calibrated-debit,
thm:n9v84-first-loss,prop:n9v84-retained-modes}.
\end{proof}

\begin{assessment}[Progress at ninth-series V84]
\label{ass:n9v84-status}
The graded Gaussian gate is now quantitative.  Three separate gaps are
visible: the interior inverse, the boundary Schur inverse, and the
reflected one-particle margin.  Calibration and moving reflection data
have explicit derivative debits.  The next compulsory companion audit
will test whether the newest YM or NS notes supply usable complete-norm
or critical-mode mechanisms before the interacting graded transfer is
attempted.
\end{assessment}

\begin{programme}[Ninth-series V85: compulsory companion audit]
\label{prog:n9v84-V85}
V85 will inspect the newest YM and NS notes, compare their current
stable norms, calibrated inverse gates, and critical-mode ownership with
RSMC, and select the interacting graded acquisition.  V86 will continue
after the audit.
\end{programme}

\begin{firewall}[Claim boundary after ninth-series V84]
\label{fire:n9v84-boundary}
V84 does not prove the physical reflected matrix Hermitian or positive,
uniform fermion gaps, locality of the reflected extraction, interacting
graded positivity, BRST descent, anomaly cancellation, infrared
closure, regulator independence, OS reconstruction, or the full
SM--Einstein continuum.
\end{firewall}

\paragraph{Parent-version record.}
V84 was formed from
\texttt{Renewal\_SM\_Einstein\_Continuum\_Research\_Notes\_V83.tex}.
The V83 parent had $32\,754$ counted source lines, $1\,178\,598$ bytes,
and SHA--256
\[
 \texttt{7260A5307F15F1E8EC3741EA115EDA5206CE7721B1DAA480CAB24A0AA24514F9}.
\]

% =============================================================================
% END APPENDED NINTH-SERIES V84 MATERIAL
% =============================================================================

% =============================================================================
% BEGIN APPENDED NINTH-SERIES V85 MATERIAL
% =============================================================================

\section{Ninth-series V85: companion audit and a positive reference graph}
\label{sec:v9v85}

\subsection{Compulsory companion audit}

The current companion files at this gate were inspected read-only:
\begin{align*}
 &\texttt{Renewal\_Yang\_Mills\_Mass\_Gap\_Research\_Notes\_V10.tex},\\
 &\texttt{Renewal\_Yang\_Mills\_Mass\_Gap\_Research\_Notes\_V11.tex},\\
 &\texttt{Renewal\_NS\_Stability\_Research\_Notes\_V26.tex}.
\end{align*}
The exact records were
\begin{center}
\begin{tabular}{c|r|l}
file & bytes & SHA--256\\ \hline
YM V10 & $181\,163$ &
\texttt{8F78A8449AB4B8F7D4EA0E241DA49BF7E1FEE8057C0677C27BEAF223E0D5E40C}\\
YM V11 & $181\,163$ &
\texttt{8F78A8449AB4B8F7D4EA0E241DA49BF7E1FEE8057C0677C27BEAF223E0D5E40C}\\
NS V26 & $278\,283$ &
\texttt{82C887F8C2C69E4F28FAD7C3DC8DE5B9099CB5A4BF431EBA1FFF3E91935978AD}
\end{tabular}
\end{center}
YM V10 and V11 were byte-identical at inspection time; therefore V11
contained no mathematical continuation beyond the substantive V10
endpoint.  V10 proves an explicit second-moment stencil bound for a
quadratic pullback residue and emphasizes that the residue can be
$O(1)$ in the insertion-scale stable norm even though canonical
dilation later contracts it.  It absorbs this term into a unique moving
stable reference graph.  NS V26 remains unchanged and its rank-one
Oseen-kernel constants do not address the reflected fermion cone.

\begin{decision}[V85 audit transfer]
\label{dec:n9v85-transfer}
The moving-reference graph mechanism is transferred to the
self-adjoint reflected one-particle sector.  An $O(1)$ Schur/rechart
residue will be admitted only if it fits inside an invariant contraction
ball that stays strictly within the positive cone.  The YM stencil
constant itself is not imported because the SM fermion interpolation
stencil has not yet been constructed.  No NS result is imported.
\end{decision}

\subsection{Why irrelevant does not mean insertion-small}

Let $X_{\rm sa}$ be a finite-dimensional normed space of self-adjoint
one-particle boundary operators.  Suppose a linear coarse dilation
$\mathcal D$ satisfies
\begin{equation}
 \lVert\mathcal D x\rVert\le\rho_D\lVert x\rVert,
 \qquad0<\rho_D<1,
 \label{eq:n9v85-dilation}
\end{equation}
on a moment-zero subspace.

\begin{proposition}[Insertion-size obstruction]
\label{prop:n9v85-insertion-size}
The contraction \eqref{eq:n9v85-dilation} places no upper bound smaller
than one on $\lVert x\rVert$ before $\mathcal D$ is applied.  In
particular an $O(1)$ pullback residue can be canonically irrelevant but
still violate the V84 gate
$\lVert x\rVert<\kappa_0$ at its insertion scale.
\end{proposition}

\begin{proof}
For any nonzero $x$ in the stable subspace and any scalar $a$,
\eqref{eq:n9v85-dilation} gives
$\lVert\mathcal D(ax)\rVert\le\rho_D|a|\lVert x\rVert$, while
$\lVert ax\rVert=|a|\lVert x\rVert$ is arbitrarily large.  The
post-dilation ratio does not bound the input radius.
\end{proof}

\begin{firewall}[No power-counting positivity proof]
\label{fire:n9v85-power-counting}
Two additional momenta can make a row contract after rescaling without
making the reflected cross matrix positive before rescaling.  The
positive-cone margin must be checked at every insertion, or the full
$O(1)$ row must be incorporated in the reference operator.
\end{firewall}

\subsection{The reflected positive reference graph}

Let $\mathcal M$ be a compact domain of calibrated marginal parameters.
For each $m\in\mathcal M$, let $K_0(m)=K_0(m)^*$ be a base reflected
operator satisfying
\begin{equation}
 K_0(m)\ge\kappa I
 \quad(m\in\mathcal M)
 \label{eq:n9v85-base-margin}
\end{equation}
for some $\kappa>0$.  Let $B_r\subset X_{\rm sa}$ be the closed norm
ball of radius $r$.  Abstract one complete rescaled stable step by
\begin{equation}
 \Phi_m:B_r\to X_{\rm sa}
 \label{eq:n9v85-Phi}
\end{equation}
and let $c(m)\in X_{\rm sa}$ be the full insertion-scale Schur/rechart
forcing, including any nonvanishing quadratic residue.

\begin{theorem}[Positive reference graph]
\label{thm:n9v85-positive-graph}
Suppose, uniformly in $m$,
\begin{align}
 \lVert\Phi_m(x)-\Phi_m(x')\rVert
 &\le\rho\lVert x-x'\rVert,
 \qquad0\le\rho<1,
 \label{eq:n9v85-contraction}\\
 \lVert\Phi_m(0)\rVert+\lVert c(m)\rVert
 &\le(1-\rho)r,
 \label{eq:n9v85-ball-gate}\\
 r&<\kappa.
 \label{eq:n9v85-positive-ball}
\end{align}
Then there is a unique $h(m)\in B_r$ satisfying
\begin{equation}
 h(m)=\Phi_m(h(m))+c(m).
 \label{eq:n9v85-graph}
\end{equation}
The full reference operator
\begin{equation}
 K_*(m):=K_0(m)+h(m)
 \label{eq:n9v85-Kstar}
\end{equation}
obeys the uniform strict margin
\begin{equation}
 K_*(m)\ge(\kappa-r)I>0.
 \label{eq:n9v85-Kstar-margin}
\end{equation}
\end{theorem}

\begin{proof}
For $x\in B_r$,
\begin{align*}
 \lVert\Phi_m(x)+c(m)\rVert
 &\le\rho\lVert x\rVert+
 \lVert\Phi_m(0)\rVert+\lVert c(m)\rVert\le r.
\end{align*}
Thus $x\mapsto\Phi_m(x)+c(m)$ maps the complete ball into itself and
is a contraction.  Banach's theorem gives the unique fixed point.
Since $h=h^*$ and $\lVert h\rVert\le r$, one has
$h\ge-rI$.  Combine with \eqref{eq:n9v85-base-margin}.
\end{proof}

\begin{corollary}[Fock margin on the graph]
\label{cor:n9v85-Fock-margin}
Under \Cref{thm:n9v85-positive-graph}, the degree-$k$ Gaussian
reflection operator $\Lambda^kK_*(m)$ has lower bound
$(\kappa-r)^k$.  Hence the reference Gaussian Fock reflection form is
strictly positive at every finite cutoff.
\end{corollary}

\begin{proof}
Apply \Cref{thm:n9v83-graded-criterion,thm:n9v84-margin}.
\end{proof}

\begin{firewall}[The invariant ball is the real gate]
\label{fire:n9v85-ball-real-gate}
The theorem does not require $c(m)\to0$, but it does require the
quantitative invariant-ball condition \eqref{eq:n9v85-ball-gate} and
$r<\kappa$.  Calling $c$ irrelevant supplies neither inequality.
\end{firewall}

\subsection{Lipschitz dependence on the marginal path}

Assume
\begin{align}
 \lVert\Phi_m(x)-\Phi_{m'}(x)\rVert
 &\le L_\Phi|m-m'|,
 \label{eq:n9v85-Phi-m}\\
 \lVert c(m)-c(m')\rVert
 &\le L_c|m-m'|.
 \label{eq:n9v85-c-m}
\end{align}

\begin{proposition}[Reference-graph response]
\label{prop:n9v85-graph-response}
The fixed graph satisfies
\begin{equation}
 \lVert h(m)-h(m')\rVert
 \le\frac{L_\Phi+L_c}{1-\rho}|m-m'|.
 \label{eq:n9v85-h-Lipschitz}
\end{equation}
If also
$\lVert K_0(m)-K_0(m')\rVert\le L_0|m-m'|$, then
\begin{equation}
 \lVert K_*(m)-K_*(m')\rVert
 \le\left(L_0+\frac{L_\Phi+L_c}{1-\rho}\right)|m-m'|.
 \label{eq:n9v85-Kstar-Lipschitz}
\end{equation}
\end{proposition}

\begin{proof}
Subtract the two fixed-point equations and add and subtract
$\Phi_m(h(m'))$:
\begin{align*}
 \lVert h(m)-h(m')\rVert
 &\le\rho\lVert h(m)-h(m')\rVert
 +(L_\Phi+L_c)|m-m'|.
\end{align*}
Rearrange.  The second bound follows from
\eqref{eq:n9v85-Kstar} and the triangle inequality.
\end{proof}

This response is part of the V81 calibration connection.  If the
determinant or reflected covariance is differentiated along $K_*(m)$,
$Dh$ is already present and must not be added as a second graph
correction.

\subsection{Tracking the moving positive graph}

Let the exact recharted stable recursion be
\begin{equation}
 x_{j+1}=\Phi_{m_{j+1}}(x_j)+c(m_{j+1})+d_j,
 \label{eq:n9v85-recursion}
\end{equation}
where $d_j$ is the declared interaction, regulator, and chart defect.
Set
\begin{equation}
 e_j:=x_j-h(m_j).
 \label{eq:n9v85-ej}
\end{equation}

\begin{theorem}[Positive-graph tracking]
\label{thm:n9v85-tracking}
Under the preceding hypotheses,
\begin{equation}
 \lVert e_{j+1}\rVert
 \le\rho\lVert e_j\rVert
 +\rho L_h|m_{j+1}-m_j|+\lVert d_j\rVert,
 \qquad
 L_h:=\frac{L_\Phi+L_c}{1-\rho}.
 \label{eq:n9v85-tracking}
\end{equation}
If $|m_{j+1}-m_j|\to0$ and $\lVert d_j\rVert\to0$, then
$e_j\to0$.  If eventually
\begin{equation}
 \lVert e_j\rVert<\kappa-r,
 \label{eq:n9v85-track-positive}
\end{equation}
then the full operator $K_*(m_j)+e_j$ remains strictly positive.
\end{theorem}

\begin{proof}
Subtract the fixed-point equation for $h(m_{j+1})$ from
\eqref{eq:n9v85-recursion}.  The $c(m_{j+1})$ terms cancel, and
contraction gives
\[
 \lVert e_{j+1}\rVert
 \le\rho\lVert x_j-h(m_{j+1})\rVert+\lVert d_j\rVert.
\]
Insert $h(m_j)$ and use \eqref{eq:n9v85-h-Lipschitz}.  Iterating the
result shows convergence because a geometric convolution of a bounded
sequence tending to zero also tends to zero.  Finally
$K_*(m_j)\ge(\kappa-r)I$ and
$e_j\ge-\lVert e_j\rVert I$ prove positivity under
\eqref{eq:n9v85-track-positive}.
\end{proof}

\begin{firewall}[The physical marginal path remains an input]
\label{fire:n9v85-marginal-input}
The fixed graph is conditional on $m_j$.  It does not construct the SM
beta functions, prove convergence of the marginal increments, or show
that unwanted finite-cutoff anisotropic coefficients vanish.  Those are
the V80 joint-calibration and symmetry-restoration gates.
\end{firewall}

\subsection{Stencil transfer remains conditional}

If a future SM fermion block interpolation has finite polyphase symbol
$\widehat Q_f(k)$ and a self-adjoint reference cross tensor $E_f$, its
quadratic pullback is
\begin{equation}
 \widehat K_f(k)=\widehat Q_f(k)^*E_f\widehat Q_f(k).
 \label{eq:n9v85-fermion-stencil}
\end{equation}
Subtracting $\widehat K_f(0)$ gives a moment-zero residue.  Under
inversion symmetry it has vanishing first moment and hence an
$O(|k|^2)$ symbol by the same elementary Taylor estimate used in the YM
audit.

\begin{firewall}[No borrowed YM stencil constant]
\label{fire:n9v85-no-borrowed-stencil}
The YM constant $\Xi_Q(\mu)$ depends on its own curvature block stencil,
aliases, and commutant.  The SM fermion/gravity boundary interpolation
has different spin, chirality, gauge, and metric indices.  Its stencil
and weighted moment constant must be computed independently before
\eqref{eq:n9v85-fermion-stencil} can feed
\eqref{eq:n9v85-ball-gate}.
\end{firewall}

\subsection{Positive reference graph card}

\begin{definition}[PRGC]
\label{def:n9v85-PRGC}
The positive reference graph card requires:
\begin{enumerate}
\item a self-adjoint base reflected operator with uniform margin
      $\kappa$;
\item the complete insertion-scale residue $c(m)$, including every
      $O(1)$ quadratic rechart row;
\item a complete stable map with contraction factor $\rho<1$;
\item the invariant-ball gate \eqref{eq:n9v85-ball-gate};
\item the positive-cone gate $r<\kappa$;
\item Lipschitz control of the graph along the calibrated marginal path;
\item a complete defect $d_j$ and vanishing forcing tail;
\item the tracking positivity condition
      \eqref{eq:n9v85-track-positive};
\item connection and determinant responses counted once; and
\item the model-specific fermion/gravity polyphase stencil and its
      weighted moment bound.
\end{enumerate}
\end{definition}

\begin{theorem}[PRGC absorbs an insertion-scale residue]
\label{thm:n9v85-PRGC}
Under PRGC, the full $O(1)$ reflected Schur/rechart residue is absorbed
into a unique self-adjoint reference graph $h(m)$ lying strictly inside
the positive cone.  Deviations track this graph by
\eqref{eq:n9v85-tracking}; if the forcing tails vanish, the reflected
one-particle operators are eventually positive and their Gaussian Fock
forms satisfy the V83 criterion.
\end{theorem}

\begin{proof}
Use \Cref{thm:n9v85-positive-graph,
cor:n9v85-Fock-margin,prop:n9v85-graph-response,
thm:n9v85-tracking,thm:n9v83-graded-criterion}.
\end{proof}

\begin{theorem}[Ninth-series V85 result]
\label{thm:n9v85-result}
The compulsory audit identifies an $O(1)$ insertion-size obstruction
which canonical irrelevance does not remove.  A complete contracting
stable map and invariant ball produce the unique reference graph
$h(m)=\Phi_m(h(m))+c(m)$.  If the ball radius satisfies $r<\kappa$, the
full reflected reference has margin $\kappa-r$, and deviations track it
under \eqref{eq:n9v85-tracking}.  The construction requires an
independent SM fermion/gravity stencil constant; the YM value is not
transferred.
\end{theorem}

\begin{proof}
Combine \Cref{prop:n9v85-insertion-size,
thm:n9v85-positive-graph,thm:n9v85-tracking,
fire:n9v85-no-borrowed-stencil}.
\end{proof}

\begin{assessment}[Progress at ninth-series V85]
\label{ass:n9v85-status}
The reflected margin no longer depends on pretending every irrelevant
row is insertion-small.  Large but controlled quadratic residues can
become part of a moving positive reference.  The next step is to treat
interactions: determine a cone-preserving bound for the exact boundary
weight around this reference and distinguish positive transfer
sandwiches from truncated cumulant approximations.
\end{assessment}

\begin{programme}[Ninth-series V86: interacting graded cone]
\label{prog:n9v85-V86}
V86 will continue after the audit.  It will prove a finite-dimensional
cone-preserving sandwich theorem for even boundary interactions, derive
a relative perturbation margin for the transfer operator, and show why
a finite connected-log truncation does not preserve graded reflection
positivity without an exact positive representation.
\end{programme}

\begin{firewall}[Claim boundary after ninth-series V85]
\label{fire:n9v85-boundary}
V85 does not prove the model-specific invariant ball, physical
fermion/gravity stencil bound, interacting graded positivity, BRST
descent, chiral anomaly cancellation, symmetry restoration, infrared
closure, regulator independence, OS reconstruction, or the full
SM--Einstein continuum.
\end{firewall}

\paragraph{Parent-version record.}
V85 was formed from
\texttt{Renewal\_SM\_Einstein\_Continuum\_Research\_Notes\_V84.tex}.
The V84 parent had $33\,245$ counted source lines, $1\,194\,520$ bytes,
and SHA--256
\[
 \texttt{639D3D540BA0D5A05DCEBFCEF829FC0934D748DB97E1501A58D87ACAD92FB31A}.
\]

% =============================================================================
% END APPENDED NINTH-SERIES V85 MATERIAL
% =============================================================================

% =============================================================================
% BEGIN APPENDED NINTH-SERIES V86 MATERIAL
% =============================================================================

\section{Ninth-series V86: interacting graded cones and the volume-margin obstruction}
\label{sec:v9v86}

V83--V85 construct a positive Gaussian reference transfer on a finite
fermionic boundary Fock space.  Interactions must preserve that cone.
This note proves exact sandwich and channel criteria, gives the sharp
relative perturbation debit, and identifies a thermodynamic obstruction:
a uniform one-particle lower bound need not give a volume-uniform lower
bound on the full Fock transfer.  Cone preservation, rather than a
global inverse margin, is the quantity which can survive large volume.

\subsection{Exact positive sandwiches}

Let $\mathcal H$ be a finite-dimensional Hilbert space, graded or
ungraded, and let $T_0=T_0^*\ge0$ be a reference transfer.

\begin{theorem}[Boundary sandwich criterion]
\label{thm:n9v86-sandwich}
For every operator $W$ on $\mathcal H$,
\begin{equation}
 T_W:=W^*T_0W\ge0.
 \label{eq:n9v86-sandwich}
\end{equation}
If $T_0\ge\tau I$ and $W^*W\ge w^2I$, then
\begin{equation}
 T_W\ge\tau w^2I.
 \label{eq:n9v86-sandwich-margin}
\end{equation}
In particular, if $V=V^*$ and $W=e^{-V/2}$, then
\begin{equation}
 T_W\ge\tau e^{-\lVert V\rVert}I.
 \label{eq:n9v86-exp-margin}
\end{equation}
\end{theorem}

\begin{proof}
For every $f$,
$\langle f,T_Wf\rangle=
\langle Wf,T_0Wf\rangle\ge0$.  Under the two lower bounds this is at
least $\tau\lVert Wf\rVert^2\ge\tau w^2\lVert f\rVert^2$.
If $V=V^*$, spectral calculus gives
$W^*W=e^{-V}\ge e^{-\lVert V\rVert}I$.
\end{proof}

The formula is the operator form of splitting an even boundary
interaction into reflected half-weights.  It does not require $W$ to
commute with $T_0$.

\begin{corollary}[Positive reflection-plane factor]
\label{cor:n9v86-plane-factor}
If $C=C^*\ge0$, then
\begin{equation}
 W^*CW\ge0.
 \label{eq:n9v86-plane-factor}
\end{equation}
Thus a reflection-plane interaction is admissible whenever its exact
boundary operator is positive and the positive/negative-time pieces
enter through paired half-weights.
\end{corollary}

\begin{proof}
Apply \Cref{thm:n9v86-sandwich} with $T_0=C$.
\end{proof}

\begin{firewall}[Cross-plane interactions need a positive centre]
\label{fire:n9v86-cross-plane}
An interaction joining positive and negative times cannot in general be
assigned to $W$ on one side.  It must be represented by a positive
central operator $C$, a positive transfer kernel, or a sum of positive
channels.  Merely writing half of its action on each side does not prove
such a representation.
\end{firewall}

\subsection{Positive channel sums and exact auxiliary labels}

\begin{theorem}[Positive channel closure]
\label{thm:n9v86-channels}
Let $T_\alpha\ge0$ and let $W_\alpha$ be arbitrary operators.  For
nonnegative weights $p_\alpha$,
\begin{equation}
 T:=\sum_{\alpha}p_\alpha W_\alpha^*T_\alpha W_\alpha
 \label{eq:n9v86-channel-sum}
\end{equation}
is positive whenever the finite sum, or norm-convergent infinite sum,
exists.  The same conclusion holds for a positive measure integral of
channels.
\end{theorem}

\begin{proof}
For every $f$, each term
$p_\alpha\langle W_\alpha f,T_\alpha W_\alpha f\rangle$ is
nonnegative.  Sum or take the monotone/norm limit.
\end{proof}

\begin{corollary}[Auxiliary-label representation]
\label{cor:n9v86-auxiliary-label}
If an exact local decomposition introduces a label $\alpha$ with a
positive conditional weight and each conditional boundary transfer is a
positive sandwich, retaining and then summing the label preserves the
reflection cone.  Integrating the label at the action level is optional;
positivity is already visible in \eqref{eq:n9v86-channel-sum}.
\end{corollary}

\begin{proof}
Use \Cref{thm:n9v86-channels}.
\end{proof}

This provides the safe use of the exact auxiliary-label idea seen in
the companion YM decomposition: the label is an interacting owner until
the positive channel sum has been established.

\begin{firewall}[Log-sum-exp hides the cone]
\label{fire:n9v86-log-sum}
Replacing \eqref{eq:n9v86-channel-sum} by the scalar effective action
$-\log\sum_\alpha e^{-S_\alpha}$ can hide the operator sandwiches and
their boundary labels.  Equality of scalar partition weights does not
by itself reconstruct a positive boundary transfer kernel.
\end{firewall}

\subsection{Positive partial trace}

Let $\mathcal H=\mathcal H_b\otimes\mathcal H_i$ and let
$\operatorname{Tr}_i$ denote the ordinary partial trace.

\begin{proposition}[Interior state elimination preserves positivity]
\label{prop:n9v86-partial-trace}
If $T\ge0$ on $\mathcal H_b\otimes\mathcal H_i$, then
\begin{equation}
 \operatorname{Tr}_iT\ge0
 \label{eq:n9v86-partial-trace}
\end{equation}
on $\mathcal H_b$.
\end{proposition}

\begin{proof}
For an orthonormal basis $(e_j)$ of $\mathcal H_i$ and $f\in\mathcal H_b$,
\begin{equation}
 \langle f,(\operatorname{Tr}_iT)f\rangle
 =\sum_j\langle f\otimes e_j,T(f\otimes e_j)\rangle\ge0.
\end{equation}
\end{proof}

\begin{firewall}[Berezin elimination must match the transfer representation]
\label{fire:n9v86-Berezin-partial}
The Grassmann Schur integral of V83 and the Hilbert-space partial trace
are two representations of interior elimination only after the coherent-
state kernel, reflection twist, and normalization are matched.  Their
formal similarity is not an identification.  GBTC must supply the
intertwiner.
\end{firewall}

\subsection{A relative interacting margin}

Positive sandwich structure is sufficient and exact.  When an
interaction is available only as an additive transfer perturbation, a
relative form bound gives a second route.

\begin{theorem}[Relative transfer margin]
\label{thm:n9v86-relative-margin}
Let $T_0>0$ and let $E=E^*$.  If
\begin{equation}
 \delta:=\lVert T_0^{-1/2}ET_0^{-1/2}\rVert<1,
 \label{eq:n9v86-relative-delta}
\end{equation}
then
\begin{equation}
 (1-\delta)T_0\le T_0+E\le(1+\delta)T_0.
 \label{eq:n9v86-relative-order}
\end{equation}
In particular $T_0+E>0$.
\end{theorem}

\begin{proof}
The self-adjoint operator
$R=T_0^{-1/2}ET_0^{-1/2}$ satisfies
$-\delta I\le R\le\delta I$.  Conjugate
$(1-\delta)I\le I+R\le(1+\delta)I$ by $T_0^{1/2}$.
\end{proof}

\begin{corollary}[Absolute debit from a strict finite-volume margin]
\label{cor:n9v86-absolute-debit}
If $T_0\ge\tau I$ and $E=E^*$, then
\begin{equation}
 \lVert E\rVert<\tau
 \label{eq:n9v86-absolute-gate}
\end{equation}
implies positivity.  The relative constant obeys
$\delta\le\lVert E\rVert/\tau$.
\end{corollary}

\begin{proof}
$\lVert T_0^{-1/2}\rVert^2\le\tau^{-1}$; apply
\Cref{thm:n9v86-relative-margin}.
\end{proof}

\begin{firewall}[Self-adjointness is again exact]
\label{fire:n9v86-selfadjoint}
The relative norm of a non-self-adjoint $E$ can be small while
$T_0+E$ fails to define a Hermitian reflection form.  Reflection
covariance must first prove $E=E^*$ in the physical boundary
representation.
\end{firewall}

\subsection{The Fock-volume margin obstruction}

Let $K>0$ on an $n$-dimensional one-particle space, with eigenvalues
$\lambda_1,\ldots,\lambda_n$.  Its Gaussian Fock transfer is
$\Gamma_f(K)=\bigoplus_k\Lambda^kK$.

\begin{theorem}[Exact many-mode minimum]
\label{thm:n9v86-Fock-minimum}
The smallest eigenvalue of $\Gamma_f(K)$ is
\begin{equation}
 \lambda_{\min}(\Gamma_f(K))
 =\prod_{j=1}^n\min\{1,\lambda_j\}.
 \label{eq:n9v86-Fock-minimum}
\end{equation}
Consequently, if $K\ge\kappa I$, then
\begin{equation}
 \Gamma_f(K)\ge\min\{1,\kappa\}^n I,
 \label{eq:n9v86-Fock-lower}
\end{equation}
and this lower bound decays exponentially in $n$ whenever
$0<\kappa<1$.
\end{theorem}

\begin{proof}
In the occupation-number basis indexed by subsets
$I\subset\{1,\ldots,n\}$, the eigenvalue is
$\prod_{j\in I}\lambda_j$.  To minimize it, include exactly those
indices with $\lambda_j<1$.  This gives
\eqref{eq:n9v86-Fock-minimum}.  The lower bound follows from
$\lambda_j\ge\kappa$.
\end{proof}

\begin{corollary}[No volume-uniform absolute perturbation gate]
\label{cor:n9v86-no-global-gate}
A one-particle margin $K\ge\kappa I$ with fixed
$0<\kappa<1$ does not yield a positive lower bound for
$\Gamma_f(K)$ uniform in the number of boundary modes.  Therefore the
absolute gate \eqref{eq:n9v86-absolute-gate} can require exponentially
small global perturbations as the volume grows.
\end{corollary}

\begin{proof}
Let $n\to\infty$ in \eqref{eq:n9v86-Fock-lower}; equality occurs for
$K=\kappa I$.
\end{proof}

\begin{firewall}[Strict finite-volume margin is the wrong continuum norm]
\label{fire:n9v86-wrong-global-norm}
The thermodynamic OS property is nonnegativity of all local reflected
observables, not a volume-uniform lower eigenvalue for the complete Fock
operator.  A proof based only on the inverse of the global transfer
margin will generally lose exponentially with volume.  Exact cone
sandwiches, positive channel sums, or uniform local relative forms are
required.
\end{firewall}

\subsection{Local tensor products and cone survival}

For disjoint boundary blocks $X$, suppose $T_X\ge0$.

\begin{proposition}[Product-cone stability]
\label{prop:n9v86-product-cone}
For every finite family $\Lambda$,
\begin{equation}
 T_\Lambda:=\bigotimes_{X\in\Lambda}T_X\ge0.
 \label{eq:n9v86-product-positive}
\end{equation}
If each local interaction is inserted as a sandwich
$W_X^*T_XW_X$, their tensor product remains positive.  This conclusion
does not require the product of the local strict margins to stay bounded
away from zero.
\end{proposition}

\begin{proof}
Write $T_X=S_X^*S_X$.  Then
$T_\Lambda=(\bigotimes_XS_X)^*(\bigotimes_XS_X)\ge0$.  Replace $S_X$
by $S_XW_X$ for the sandwiched statement.
\end{proof}

\begin{proposition}[Local channel assembly]
\label{prop:n9v86-local-channels}
Suppose every finite-volume interacting transfer is a finite or
norm-convergent sum
\begin{equation}
 T_\Lambda=\sum_\alpha p_{\Lambda,\alpha}
 W_{\Lambda,\alpha}^*T_{0,\Lambda,\alpha}
 W_{\Lambda,\alpha},
 \qquad p_{\Lambda,\alpha}\ge0,
 \label{eq:n9v86-local-channel-assembly}
\end{equation}
with $T_{0,\Lambda,\alpha}\ge0$.  Then every finite-volume OS form is
nonnegative.  If the local Schwinger functions converge and the limit
commutes with reflection, their limiting OS quadratic forms remain
nonnegative.
\end{proposition}

\begin{proof}
Finite-volume positivity follows from
\Cref{thm:n9v86-channels}.  For a fixed local observable, its OS value
is a nonnegative real number at each volume; convergence preserves the
closed half-line $[0,\infty)$.
\end{proof}

\begin{firewall}[Existence of the limit is separate]
\label{fire:n9v86-limit-existence}
The last proposition transfers positivity through an already proved
local limit.  It does not prove tightness, convergence of Schwinger
functions, regulator independence, clustering, or reconstruction.
Those remain continuum gates.
\end{firewall}

\subsection{Why finite cumulant truncation is not a cone proof}

The exact conditional integration may be written as an infinite
connected expansion, but positivity is a nonlinear closed-cone
property.

\begin{proposition}[Polynomial truncation counterexample]
\label{prop:n9v86-truncation-counterexample}
The positive transfer family
\begin{equation}
 T(t)=e^{-tH},
 \qquad H=\begin{pmatrix}0&0\\0&1\end{pmatrix},
 \qquad t\ge0,
 \label{eq:n9v86-positive-family}
\end{equation}
has first-order connected/Duhamel polynomial approximation
\begin{equation}
 T^{[1]}(t)=I-tH
 =\begin{pmatrix}1&0\\0&1-t\end{pmatrix},
 \label{eq:n9v86-first-truncation}
\end{equation}
which is not positive for $t>1$.  Hence finite truncation of an exact
positive expansion does not preserve the positive cone without a
separate relative-margin estimate.
\end{proposition}

\begin{proof}
$T(t)$ has eigenvalues $1$ and $e^{-t}$, both positive.  The truncated
matrix has eigenvalue $1-t$, which is negative for $t>1$.
\end{proof}

\begin{firewall}[Exponentiating a truncated scalar action is not enough]
\label{fire:n9v86-truncated-action}
Even when an exponentiated truncated scalar action is pointwise
positive, OS positivity concerns the reflected boundary kernel, not
pointwise density.  The truncation must be reorganized into
\eqref{eq:n9v86-sandwich}, \eqref{eq:n9v86-channel-sum}, or satisfy the
relative operator gate \eqref{eq:n9v86-relative-delta}.
\end{firewall}

\subsection{Interacting graded cone card}

\begin{definition}[IGCC]
\label{def:n9v86-IGCC}
The interacting graded cone card requires:
\begin{enumerate}
\item a V85 positive Gaussian reference graph;
\item exact reflection covariance and self-adjoint boundary operators;
\item every half-space interaction represented by a boundary sandwich;
\item every cross-plane interaction represented by a positive centre or
      positive channel sum;
\item an exact coherent-state identification between Grassmann Schur
      integration and the Hilbert transfer;
\item local rather than global-volume inverse margins for any additive
      remainder;
\item retained zero modes and determinant-line phases;
\item a complete stable norm which bounds all non-sandwich remainders;
\item finite-volume positive channel assembly; and
\item convergence of local reflected observables before passing the cone
      inequality to the limit.
\end{enumerate}
\end{definition}

\begin{theorem}[IGCC preserves the interacting finite-volume cone]
\label{thm:n9v86-IGCC}
Under IGCC, every finite-volume interacting graded transfer is positive.
Exact sandwich and channel pieces preserve the cone without using a
global strict margin; additive local remainders preserve it under the
relative gate \eqref{eq:n9v86-relative-delta}.  If local reflected
Schwinger functions converge, their limiting OS forms are nonnegative.
\end{theorem}

\begin{proof}
Use \Cref{thm:n9v86-sandwich,thm:n9v86-channels,
prop:n9v86-partial-trace,thm:n9v86-relative-margin,
prop:n9v86-local-channels}.  The volume-margin obstruction
\Cref{thm:n9v86-Fock-minimum} explains why IGCC uses local cone
representations rather than a global inverse bound.
\end{proof}

\begin{theorem}[Ninth-series V86 result]
\label{thm:n9v86-result}
Exact boundary sandwiches, positive channel sums, and positive partial
traces preserve the interacting graded cone.  An additive self-adjoint
remainder is admissible under the relative bound
$\lVert T_0^{-1/2}ET_0^{-1/2}\rVert<1$.  However, for an $n$-mode
Gaussian Fock transfer the exact minimum is
$\prod_j\min\{1,\lambda_j(K)\}$, so a fixed one-particle margin below one
decays exponentially with volume.  The continuum strategy must
therefore use exact local cone representations and convergence of local
OS forms, not a volume-uniform inverse of the full Fock transfer.
\end{theorem}

\begin{proof}
Combine \Cref{thm:n9v86-sandwich,thm:n9v86-channels,
thm:n9v86-relative-margin,thm:n9v86-Fock-minimum,
prop:n9v86-local-channels}.
\end{proof}

\begin{assessment}[Progress at ninth-series V86]
\label{ass:n9v86-status}
The interacting graded problem now has a viable volume-stable target:
an exact local positive-channel representation around the moving
reference graph.  The global-margin route is ruled out quantitatively.
The next bottleneck is the continuum limit itself: obtain tight local
Schwinger bounds, compatibility between regulators, and a Cauchy
estimate strong enough to pass the complete SM--Einstein Ward and OS
identities to a common limit.
\end{assessment}

\begin{programme}[Ninth-series V87: regulator comparison]
\label{prog:n9v86-V87}
V87 will formulate a two-regulator coupling and telescoping criterion
for local Schwinger functions.  It will track calibrated reference,
constraint, determinant, Ward, and positive-channel differences in one
norm and identify the summability needed for a regulator-independent
local limit.
\end{programme}

\begin{firewall}[Claim boundary after ninth-series V86]
\label{fire:n9v86-boundary}
V86 does not construct the physical positive-channel representation,
prove the coherent-state/BRST identification, close the complete stable
norm, establish chiral anomaly cancellation, prove tightness or local
Schwinger convergence, regulator independence, OS reconstruction, or
the full SM--Einstein continuum.
\end{firewall}

\paragraph{Parent-version record.}
V86 was formed from
\texttt{Renewal\_SM\_Einstein\_Continuum\_Research\_Notes\_V85.tex}.
The V85 parent had $33\,636$ counted source lines, $1\,208\,625$ bytes,
and SHA--256
\[
 \texttt{E35D2B4F82346963F04578ACB59A8DF9FE7FE42B467AEF645A8A692DB5BD04FD}.
\]

% =============================================================================
% END APPENDED NINTH-SERIES V86 MATERIAL
% =============================================================================

% =============================================================================
% BEGIN APPENDED NINTH-SERIES V87 MATERIAL
% =============================================================================

\section{Ninth-series V87: regulator telescopes for local Schwinger data}
\label{sec:v9v87}

V86 identifies local reflection functionals, rather than a global Fock
inverse margin, as the correct thermodynamic objects.  This note gives a
two-regulator comparison theorem for those functionals.  It separates
seven discrepancy rows, proves the telescoping criterion for a local
limit, and states the additional cross-tower condition needed for
regulator independence.  Existence of the required estimates remains a
model-specific gate.

\subsection{Compatible local algebras}

Let $\mathcal O$ range over bounded spacetime regions.  For every
regulator index $n$ let $\mathcal A_n(\mathcal O)$ be a finite
$\mathbb Z_2$-graded observable algebra and let
\begin{equation}
 \omega_n:\mathcal A_n(\mathcal O)\to\mathbb C
 \label{eq:n9v87-omega-n}
\end{equation}
be its normalized regulated functional.  Let
\begin{equation}
 I_n^{n+1}:\mathcal A_n(\mathcal O)	o
 \mathcal A_{n+1}(\mathcal O^{+r_n})
 \label{eq:n9v87-embedding}
\end{equation}
be a unital parity-preserving reconstruction map, with a declared collar
$r_n$.  Write
$I_n^m=I_{m-1}^m\cdots I_n^{n+1}$.

\begin{definition}[Local comparison seminorm]
\label{def:n9v87-comparison-seminorm}
For a normed test class $\mathcal B_n(\mathcal O)\subset
\mathcal A_n(\mathcal O)$, define
\begin{equation}
 d_n(\mathcal O)
 :=\sup_{0\ne F\in\mathcal B_n(\mathcal O)}
 \frac{|\omega_{n+1}(I_n^{n+1}F)-\omega_n(F)|}
 {\lVert F\rVert_{\mathcal B_n}}.
 \label{eq:n9v87-dn}
\end{equation}
\end{definition}

\begin{theorem}[Local telescoping criterion]
\label{thm:n9v87-local-telescope}
Suppose the embeddings are norm-nonexpanding on the transported test
classes and
\begin{equation}
 \sum_{n=n_0}^{\infty}d_n(\mathcal O^{+R_n})<\infty,
 \label{eq:n9v87-summable-d}
\end{equation}
where $R_n$ contains the accumulated collars needed to compare a fixed
initial observable.  Then for every
$F\in\mathcal B_{n_0}(\mathcal O)$ the sequence
\begin{equation}
 \omega_m(I_{n_0}^mF)
 \label{eq:n9v87-expectation-sequence}
\end{equation}
is Cauchy.  More precisely, for $M>N\ge n_0$,
\begin{equation}
 |\omega_M(I_{n_0}^MF)-\omega_N(I_{n_0}^NF)|
 \le\lVert F\rVert_{\mathcal B_{n_0}}
 \sum_{n=N}^{M-1}d_n(\mathcal O^{+R_n}).
 \label{eq:n9v87-tail-bound}
\end{equation}
\end{theorem}

\begin{proof}
Insert the intermediate values
$\omega_{n+1}(I_{n_0}^{n+1}F)$ and
$\omega_n(I_{n_0}^{n}F)$ and telescope from $N$ to $M-1$.
The transported test norm is no larger than the initial norm, so each
difference is bounded by the corresponding $d_n$.  Summability makes
the tail tend to zero.
\end{proof}

\begin{firewall}[Collars must remain local]
\label{fire:n9v87-collars}
If the accumulated reconstruction collar of a fixed observable fills
the whole volume before the comparison tail is controlled, the estimate
is not local and cannot establish a thermodynamic limit.  Finite-range
or summably decaying reconstruction must keep $R_n$ compatible with the
chosen continuum support.
\end{firewall}

\subsection{The seven-row hybrid ledger}

For one step, insert intermediate normalized functionals which change
one construction component at a time: calibrated reference, stable
interaction, constraint reduction, determinant/ghost density, Ward
counterterm, positive-channel representation, and reconstruction.
Denote them
\begin{equation}
 \omega_n^{(0)},\omega_n^{(1)},\ldots,\omega_n^{(7)},
 \quad
 \omega_n^{(0)}=\omega_n,
 \quad
 \omega_n^{(7)}=\omega_{n+1}\circ I_n^{n+1}.
 \label{eq:n9v87-hybrids}
\end{equation}

\begin{proposition}[Exact hybrid comparison]
\label{prop:n9v87-hybrid}
If
\begin{equation}
 |\omega_n^{(j)}(F)-\omega_n^{(j-1)}(F)|
 \le\epsilon_n^{(j)}\lVert F\rVert_{\mathcal B_n},
 \qquad1\le j\le7,
 \label{eq:n9v87-hybrid-row}
\end{equation}
then
\begin{equation}
 d_n\le
 \epsilon_n^{\rm ref}+epsilon_n^{\rm st}
 +\epsilon_n^{\rm con}+epsilon_n^{\rm det}
 +\epsilon_n^{\rm Ward}+epsilon_n^{\rm chan}
 +\epsilon_n^{\rm rec}.
 \label{eq:n9v87-seven-ledger}
\end{equation}
\end{proposition}

\begin{proof}
Telescope the seven hybrid differences and use the triangle inequality.
\end{proof}

The rows in \eqref{eq:n9v87-seven-ledger} have fixed meanings:
\begin{enumerate}
\item $\epsilon^{\rm ref}$ is the V80--V85 moving-reference mismatch,
measured against the transported reference;
\item $\epsilon^{\rm st}$ is the complete stable interaction difference;
\item $\epsilon^{\rm con}$ is the constraint-branch/coarea difference;
\item $\epsilon^{\rm det}$ contains the aligned V82 Jacobian, ghost, and
physical fermion determinant responses, each once;
\item $\epsilon^{\rm Ward}$ is the gauge/diffeomorphism/anomaly residual
and counterterm difference;
\item $\epsilon^{\rm chan}$ compares exact positive boundary-channel
representations; and
\item $\epsilon^{\rm rec}$ is the observable reconstruction error.
\end{enumerate}

\begin{firewall}[A small total difference does not identify its owner]
\label{fire:n9v87-owner-cancellation}
Opposite errors in two hybrid rows can cancel in a direct comparison.
Such cancellation is unusable for Ward, anomaly, positivity, or
universality unless an identity proves it.  The hybrid path keeps each
owner bounded before summation.
\end{firewall}

\subsection{Couplings for bosonic polynomial observables}

For unbounded bosonic fields, bounded-observable comparison is not
enough.  Let $X_{n,j}=\Phi_n(f_j)$ be reconstructed smeared fields on a
coupling of two successive regulators.

\begin{theorem}[Coupled moment comparison]
\label{thm:n9v87-moment-coupling}
Fix $k\ge1$.  Suppose a coupling $\pi_n$ of the reconstructed regulator
$n$ and $n+1$ fields satisfies, for $1\le j\le k$,
\begin{align}
 \lVert X_{n+1,j}-X_{n,j}\rVert_{L^k(\pi_n)}
 &\le\delta_n(f_j),
 \label{eq:n9v87-coupled-difference}\\
 \max\{\lVert X_{n,j}\rVert_{L^k(\pi_n)},
       \lVert X_{n+1,j}\rVert_{L^k(\pi_n)}\}
 &\le M_j.
 \label{eq:n9v87-uniform-moment}
\end{align}
Then
\begin{align}
 &\left|
 \mathbb E\prod_{j=1}^kX_{n+1,j}
 -\mathbb E\prod_{j=1}^kX_{n,j}
 \right|\\
 &\qquad\le
 \sum_{j=1}^k\delta_n(f_j)
 \prod_{\ell\ne j}M_\ell.
 \label{eq:n9v87-product-bound}
\end{align}
Thus summability of the right side gives a Cauchy sequence of smeared
$k$-point functions.
\end{theorem}

\begin{proof}
Use the exact product telescope
\begin{align*}
 \prod_{j=1}^kx_j-\prod_{j=1}^ky_j
 =\sum_{j=1}^k(x_j-y_j)
 \left(\prod_{\ell<j}x_\ell\right)
 \left(\prod_{\ell>j}y_\ell\right).
\end{align*}
Apply H\"older with $k$ factors to each expectation and use
\eqref{eq:n9v87-coupled-difference}--
\eqref{eq:n9v87-uniform-moment}.
\end{proof}

\begin{firewall}[Weak convergence without uniform integrability]
\label{fire:n9v87-uniform-integrability}
A coupling which makes smeared fields close in probability does not
control polynomial Schwinger functions without uniform moments or
uniform integrability.  Rare bad fields can carry a nonvanishing moment
tail even when their probability tends to zero.
\end{firewall}

\subsection{Determinant-density comparison}

The determinant row often changes a normalized scalar density.  The
following elementary bound converts an aligned log-determinant estimate
into a bounded-observable comparison.

\begin{proposition}[Normalized exponential reweighting]
\label{prop:n9v87-reweighting}
Let $\mu$ be a probability measure and let
\begin{equation}
 d\nu=Z^{-1}e^{-V}\,d\mu,
 \qquad Z=\int e^{-V}\,d\mu.
 \label{eq:n9v87-reweight}
\end{equation}
If $|V|\le d$, then for every bounded $F$,
\begin{equation}
 |\nu(F)-\mu(F)|
 \le(e^{2d}-1)\lVert F\rVert_\infty.
 \label{eq:n9v87-reweight-bound}
\end{equation}
\end{proposition}

\begin{proof}
$e^{-d}\le Z\le e^d$, so the density
$r=Z^{-1}e^{-V}$ lies in $[e^{-2d},e^{2d}]$.  Hence
$|r-1|\le e^{2d}-1$ and
$|\nu(F)-\mu(F)|=|\mu((r-1)F)|$ has the stated bound.
\end{proof}

\begin{corollary}[Aligned log-determinant row]
\label{cor:n9v87-logdet-row}
If the total change of the correctly owned aligned log-determinants
between two regulators is bounded pointwise by $d_n$, then its bounded-
observable hybrid row satisfies
\begin{equation}
 \epsilon_n^{\rm det}\le e^{2d_n}-1.
 \label{eq:n9v87-det-row}
\end{equation}
Summability of $d_n$ and boundedness of its initial terms imply
summability of the tail of \eqref{eq:n9v87-det-row}.
\end{corollary}

\begin{proof}
Use \Cref{prop:n9v87-reweighting}.  Since
$e^{2d}-1\le 2e^{2d_*}d$ for $0\le d\le d_*$, a summable tail of
$d_n$ gives a summable exponential tail.
\end{proof}

\begin{firewall}[Complex determinant phases are not scalar reweightings]
\label{fire:n9v87-complex-phase}
The proposition applies to real bounded changes of a positive density.
Chiral and oriented ghost phases require comparison of graded
functionals or determinant-line parallel transport; taking their
absolute values changes the theory.
\end{firewall}

\subsection{Positive-channel comparison}

Let $\rho_n^{\mathcal O}$ be a normalized local boundary density matrix
representing the V86 positive channel functional in region $\mathcal O$.

\begin{proposition}[Trace-norm channel row]
\label{prop:n9v87-channel-row}
For every bounded local operator $F$,
\begin{equation}
 |\operatorname{Tr}[(\rho_{n+1}^{\mathcal O}
 -\rho_n^{\mathcal O})F]|
 \le\lVert\rho_{n+1}^{\mathcal O}-\rho_n^{\mathcal O}\rVert_1
 \lVert F\rVert.
 \label{eq:n9v87-trace-row}
\end{equation}
Hence one may take
$\epsilon_n^{\rm chan}$ to be the local trace-norm difference after
transport to a common boundary space.
\end{proposition}

\begin{proof}
This is trace/operator norm duality.
\end{proof}

\begin{firewall}[Common boundary space]
\label{fire:n9v87-common-boundary}
Density matrices at different cutoffs cannot be subtracted until the
one-particle reconstruction, retained zero modes, chirality maps, and
graded tensor order have been identified.  That identification is part
of $I_n^{n+1}$ and its error belongs to
$\epsilon_n^{\rm rec}$.
\end{firewall}

\subsection{Distributional Schwinger limits}

Let $\mathcal D$ be a nuclear test-function space and let
$S_n^{(k)}$ be the regulated smeared $k$-point form, multilinear on
$\mathcal D^k$.

\begin{theorem}[Continuous multilinear limit]
\label{thm:n9v87-distribution-limit}
Suppose for every $f_1,\ldots,f_k$ the sequence
$S_n^{(k)}(f_1,\ldots,f_k)$ is Cauchy and there are a continuous
seminorm $p$ and a constant $C_k$, independent of $n$, such that
\begin{equation}
 |S_n^{(k)}(f_1,\ldots,f_k)|
 \le C_k\prod_{j=1}^kp(f_j).
 \label{eq:n9v87-uniform-distribution-bound}
\end{equation}
Then the pointwise limit $S^{(k)}$ is a continuous multilinear form
with the same bound and therefore defines a distributional $k$-point
kernel.
\end{theorem}

\begin{proof}
Multilinearity passes through pointwise limits.  Taking the limit in
\eqref{eq:n9v87-uniform-distribution-bound} gives the same bound for
$S^{(k)}$, which is the standard continuity criterion for a multilinear
form.  The Schwartz kernel theorem identifies it with a distribution on
the product test space.
\end{proof}

\begin{firewall}[All orders and growth remain necessary]
\label{fire:n9v87-all-orders}
Convergence at each fixed order does not alone construct a probability
or graded quantum field.  One also needs consistency, permutation/
graded symmetry, suitable growth or characteristic-functional
continuity, and compatibility of all orders with one state.
\end{firewall}

\subsection{Passage of reflection positivity and Ward identities}

Assume each embedding intertwines reflection on local observables,
\begin{equation}
 I_n^{n+1}\Theta_n=\Theta_{n+1}I_n^{n+1}.
 \label{eq:n9v87-reflection-intertwining}
\end{equation}

\begin{theorem}[Closed-cone passage to the limit]
\label{thm:n9v87-OS-limit}
Suppose \Cref{thm:n9v87-local-telescope} defines a limiting functional
$\omega$ and every $\omega_n$ is reflection positive on the transported
positive-time algebra.  Then
\begin{equation}
 \omega((\Theta F)F)\ge0
 \label{eq:n9v87-limit-OS}
\end{equation}
for every local positive-time $F$ in the inductive test algebra.
\end{theorem}

\begin{proof}
By \eqref{eq:n9v87-reflection-intertwining}, the regulated numbers
$\omega_n((\Theta F)F)$ are nonnegative.  They converge by the local
telescoping theorem.  The closed half-line $[0,\infty)$ contains the
limit.
\end{proof}

Let $\delta_n$ be a regulated symmetry derivation and $\delta$ the
candidate limiting derivation.

\begin{theorem}[Ward identity passage]
\label{thm:n9v87-Ward-limit}
Suppose for every local $F$,
\begin{align}
 |\omega_n(\delta_nF)|&\le\eta_n\lVert F\rVert_{\mathcal B},
 &\eta_n&\longrightarrow0,
 \label{eq:n9v87-Ward-residual}\\
 \lVert\delta_nF-\delta F\rVert_{\mathcal B}&\longrightarrow0,
 \label{eq:n9v87-generator-convergence}
\end{align}
after common reconstruction, and the functionals are uniformly bounded
on $\mathcal B$.  Then
\begin{equation}
 \omega(\delta F)=0.
 \label{eq:n9v87-limit-Ward}
\end{equation}
\end{theorem}

\begin{proof}
Add and subtract $\omega_n(\delta_nF)$ and use
\begin{align*}
 |\omega(\delta F)|
 \le{}&|\omega(\delta F)-\omega_n(\delta F)|
 +|\omega_n((\delta-\delta_n)F)|
 +|\omega_n(\delta_nF)|.
\end{align*}
The local functional convergence, generator convergence, and residual
bound send the three terms to zero.
\end{proof}

\begin{firewall}[Anomaly is the nonvanishing limit]
\label{fire:n9v87-anomaly-limit}
If $\omega_n(\delta_nF)$ converges to a nonzero functional, the result is
an anomaly, not a failed estimate to be omitted.  Cancellation requires
the complete SM fermion, ghost, boundary, and regulator ledger to make
the combined residual vanish.
\end{firewall}

\subsection{Cross-tower regulator independence}

Convergence of one regulator tower does not identify its limit with a
second tower.

\begin{theorem}[Cofinal cross-comparison criterion]
\label{thm:n9v87-cross-tower}
Let $\omega_n$ and $\widetilde\omega_m$ be two locally Cauchy regulator
towers with limits $\omega$ and $\widetilde\omega$.  Suppose that for
every local $F$ there are cofinal indices $n_j,m_j$ and common
reconstructions such that
\begin{equation}
 |\omega_{n_j}(F)-\widetilde\omega_{m_j}(F)|
 \le\chi_j\lVert F\rVert_{\mathcal B},
 \qquad\chi_j\longrightarrow0.
 \label{eq:n9v87-cross-comparison}
\end{equation}
Then $\omega(F)=\widetilde\omega(F)$ for every such $F$.
\end{theorem}

\begin{proof}
Use the triangle inequality
\begin{align*}
 |\omega(F)-\widetilde\omega(F)|
 \le{}&|\omega(F)-\omega_{n_j}(F)|
 +|\omega_{n_j}(F)-\widetilde\omega_{m_j}(F)|\\
 &+|\widetilde\omega_{m_j}(F)-\widetilde\omega(F)|
\end{align*}
and send $j\to\infty$.
\end{proof}

\begin{firewall}[Universality is a cross-comparison theorem]
\label{fire:n9v87-universality}
Two separately convergent lattice or spectral regulators can converge
to different theories.  Matching formal continuum actions or a finite
set of renormalized couplings does not prove equality of all local
functionals.  A common refinement or cross-tower estimate such as
\eqref{eq:n9v87-cross-comparison} is required.
\end{firewall}

\subsection{Regulator comparison card}

\begin{definition}[RCC]
\label{def:n9v87-RCC}
The regulator comparison card requires:
\begin{enumerate}
\item local graded algebras and reflection-compatible reconstruction
      maps with controlled collars;
\item the seven hybrid rows \eqref{eq:n9v87-seven-ledger};
\item summability of their local one-step errors;
\item uniform moment bounds for every bosonic polynomial degree used;
\item common boundary identifications and trace-norm channel control;
\item uniform distributional seminorm bounds at every Schwinger order;
\item consistency and graded symmetry of the limiting forms;
\item vanishing complete Ward/constraint residuals and convergence of
      their generators;
\item finite-volume reflection positivity in the exact V86 channel
      representation; and
\item cofinal cross-tower comparisons for regulator independence.
\end{enumerate}
\end{definition}

\begin{theorem}[RCC gives a regulator-independent local OS functional]
\label{thm:n9v87-RCC}
Under RCC, every local observable has a regulator-independent limiting
functional.  The limiting Schwinger forms are distributions, the local
OS inequalities pass to the limit, and every symmetry whose complete
Ward residual vanishes satisfies its limiting Ward identity.  This
conclusion is at the level of local functionals; Hilbert-space
reconstruction, clustering, and a mass spectrum remain separate.
\end{theorem}

\begin{proof}
Use \Cref{prop:n9v87-hybrid,thm:n9v87-local-telescope,
thm:n9v87-moment-coupling,thm:n9v87-distribution-limit,
thm:n9v87-OS-limit,thm:n9v87-Ward-limit,
thm:n9v87-cross-tower}.
\end{proof}

\begin{theorem}[Ninth-series V87 result]
\label{thm:n9v87-result}
Summable local one-step discrepancies imply Cauchy convergence of all
transported bounded observables.  Uniform coupled moments extend the
result to polynomial Schwinger functions.  The exact hybrid ledger
splits the comparison into reference, stable, constraint, determinant,
Ward, positive-channel, and reconstruction rows.  Reflection positivity
and vanishing Ward identities pass through the local limit, while
regulator independence additionally requires the cross-tower criterion
\eqref{eq:n9v87-cross-comparison}.
\end{theorem}

\begin{proof}
Combine \Cref{thm:n9v87-local-telescope,
prop:n9v87-hybrid,thm:n9v87-moment-coupling,
thm:n9v87-OS-limit,thm:n9v87-Ward-limit,
thm:n9v87-cross-tower}.
\end{proof}

\begin{assessment}[Progress at ninth-series V87]
\label{ass:n9v87-status}
The continuum target is now expressed as a summable comparison problem
for local observables, compatible with the V86 volume-margin
obstruction.  Regulator independence is correctly separated from
single-tower convergence.  The next step is to convert the limiting OS
functional into a Hilbert space and transfer semigroup while retaining
the gauge/BRST quotient, then state the exact clustering and spectral
conditions needed for the particle interpretation.
\end{assessment}

\begin{programme}[Ninth-series V88: OS reconstruction compiler]
\label{prog:n9v87-V88}
V88 will construct the pre-Hilbert quotient from the limiting local OS
form, prove that reflection-compatible time translations descend to a
contraction semigroup under the required hypotheses, and isolate the
additional continuity and clustering gates.  No spectral claim will be
made without those inputs.
\end{programme}

\begin{firewall}[Claim boundary after ninth-series V87]
\label{fire:n9v87-boundary}
V87 does not prove the seven discrepancy rows summable, all-order moment
bounds, consistency of a limiting state, cross-regulator comparison,
anomaly cancellation, Hilbert reconstruction, clustering, a physical
spectrum, or the full SM--Einstein continuum.
\end{firewall}

\paragraph{Parent-version record.}
V87 was formed from
\texttt{Renewal\_SM\_Einstein\_Continuum\_Research\_Notes\_V86.tex}.
The V86 parent had $34\,095$ counted source lines, $1\,225\,388$ bytes,
and SHA--256
\[
 \texttt{839C39EAF32DF10EDA55E23D33869BBDB78FDB59CA8EC87C47CCDB5C357EDD07}.
\]

% =============================================================================
% END APPENDED NINTH-SERIES V87 MATERIAL
% =============================================================================

% =============================================================================
% BEGIN APPENDED NINTH-SERIES V88 MATERIAL
% =============================================================================

\section{Ninth-series V88: OS Hilbert and transfer reconstruction}
\label{sec:v9v88}

V87 gives a conditional regulator-independent local functional with
reflection positivity and limiting Ward identities.  This note proves
the algebraic Hilbert reconstruction and states the exact additional
hypotheses needed for a positive-energy time-translation semigroup.  It
separates OS positivity, semigroup contractivity, clustering, a spectral
gap, field-operator domains, and the BRST/constraint quotient.

\subsection{The OS pre-Hilbert quotient}

Let $\mathcal A$ be a unital $\mathbb Z_2$-graded complex algebra, let
$\mathcal A_+\subset\mathcal A$ be the positive-time subalgebra, and
let $\Theta$ be the fixed antilinear reflected graded involution.  Let
$\omega:\mathcal A\to\mathbb C$ be normalized and suppose the
sesquilinear form
\begin{equation}
 (F,G)_{\rm OS}:=\omega((\Theta F)G),
 \qquad F,G\in\mathcal A_+,
 \label{eq:n9v88-OS-form}
\end{equation}
is Hermitian and positive semidefinite.

\begin{lemma}[Cauchy--Schwarz for the OS form]
\label{lem:n9v88-CS}
For all $F,G\in\mathcal A_+$,
\begin{equation}
 |(F,G)_{\rm OS}|^2
 \le(F,F)_{\rm OS}(G,G)_{\rm OS}.
 \label{eq:n9v88-CS}
\end{equation}
\end{lemma}

\begin{proof}
If $(G,G)_{\rm OS}>0$, positivity of
$(F+zG,F+zG)_{\rm OS}$ for
$z=-(G,F)_{\rm OS}/(G,G)_{\rm OS}$ gives the inequality.  If
$(G,G)_{\rm OS}=0$, positivity for every real and imaginary scalar
multiple of $G$ forces $(F,G)_{\rm OS}=0$.
\end{proof}

Define the null space
\begin{equation}
 \mathcal N:=\{F\in\mathcal A_+:(F,F)_{\rm OS}=0\}.
 \label{eq:n9v88-null}
\end{equation}

\begin{theorem}[OS Hilbert reconstruction]
\label{thm:n9v88-Hilbert}
$\mathcal N$ is a linear subspace and is orthogonal to all of
$\mathcal A_+$.  Therefore
\begin{equation}
 \mathcal H_0:=\mathcal A_+/\mathcal N,
 \qquad
 \langle[F],[G]\rangle:=(F,G)_{\rm OS}
 \label{eq:n9v88-H0}
\end{equation}
is a well-defined inner-product space.  Its completion
$\mathcal H_{\rm OS}$ is a Hilbert space, and
\begin{equation}
 \Omega:=[1]
 \label{eq:n9v88-vacuum}
\end{equation}
has norm one.
\end{theorem}

\begin{proof}
By \Cref{lem:n9v88-CS}, every null vector has zero pairing with every
vector.  This proves linearity of $\mathcal N$ and independence of the
representatives in \eqref{eq:n9v88-H0}.  The quotient form is strictly
positive by construction.  Complete it.  Finally
$\lVert\Omega\rVert^2=\omega(1)=1$.
\end{proof}

\begin{firewall}[The null space is essential]
\label{fire:n9v88-null-essential}
Reflection positivity is generally semidefinite.  Gauge, equation-of-
motion, zero-mode, or boundary relations can create null vectors.
Treating the pre-Hilbert form as strictly positive before quotienting
can make operators and time translations depend on the chosen
representative.
\end{firewall}

\subsection{Time translations}

Let $\tau_t$ be Euclidean time translations on local observables.  For
$t\ge0$, assume $\tau_t\mathcal A_+\subset\mathcal A_+$ on the chosen
domain.

\begin{definition}[OS translation gates]
\label{def:n9v88-translation-gates}
The translation gates are:
\begin{enumerate}
\item $\tau_{t+s}=\tau_t\tau_s$ and $\tau_0=I$;
\item null preservation,
      $\tau_t\mathcal N\subset\mathcal N$;
\item OS contraction,
\begin{equation}
 \lVert[\tau_tF]\rVert_{\rm OS}
 \le\lVert[F]\rVert_{\rm OS};
 \label{eq:n9v88-OS-contraction}
\end{equation}
\item OS symmetry,
\begin{equation}
 (F,\tau_tG)_{\rm OS}=(\tau_tF,G)_{\rm OS};
 \label{eq:n9v88-OS-symmetry}
\end{equation}
and
\item strong continuity at zero on a dense local domain.
\end{enumerate}
\end{definition}

\begin{theorem}[Positive-energy transfer semigroup]
\label{thm:n9v88-semigroup}
Under the translation gates,
\begin{equation}
 T(t)[F]:=[\tau_tF]
 \label{eq:n9v88-T}
\end{equation}
extends uniquely to a strongly continuous self-adjoint contraction
semigroup on $\mathcal H_{\rm OS}$.  There is a unique self-adjoint
operator $H\ge0$ such that
\begin{equation}
 T(t)=e^{-tH}.
 \label{eq:n9v88-H}
\end{equation}
If $\omega$ is time-translation invariant, then
$T(t)\Omega=\Omega$ and $H\Omega=0$.
\end{theorem}

\begin{proof}
Null preservation makes \eqref{eq:n9v88-T} well defined on the
quotient.  The contraction gate extends it continuously to the
completion, the algebraic semigroup law survives the extension, and
the continuity gate gives strong continuity.  OS symmetry makes every
$T(t)$ self-adjoint.  The spectral theorem for strongly continuous
self-adjoint contraction semigroups yields a self-adjoint nonnegative
generator $H$ with \eqref{eq:n9v88-H}.  Translation invariance gives
$[\tau_t1]=[1]$; differentiation at zero gives $H\Omega=0$.
\end{proof}

\begin{firewall}[Reflection positivity does not alone prove contraction]
\label{fire:n9v88-contraction-open}
The inequality \eqref{eq:n9v88-OS-contraction} requires the appropriate
Euclidean translation/reflection structure, often formulated through
reflection positivity at translated planes or an already constructed
positive transfer.  Positivity of the form at one plane alone does not
justify inserting this inequality.
\end{firewall}

\begin{proposition}[Symmetry from covariance]
\label{prop:n9v88-symmetry-covariance}
Suppose $\omega\circ\tau_t=\omega$ and
\begin{equation}
 \Theta\tau_t=\tau_{-t}\Theta.
 \label{eq:n9v88-reflection-translation}
\end{equation}
Whenever all products are in the common domain,
\eqref{eq:n9v88-OS-symmetry} holds.
\end{proposition}

\begin{proof}
Use covariance and translate the second expression:
\begin{align*}
 (\tau_tF,G)_{\rm OS}
 &=\omega((\Theta\tau_tF)G)
 =\omega((\tau_{-t}\Theta F)G)\\
 &=\omega((\Theta F)(\tau_tG))
 =(F,\tau_tG)_{\rm OS}.
\end{align*}
\end{proof}

\subsection{Spatial symmetries}

Let $g$ be a Euclidean transformation preserving the positive
half-space and commuting with the reflection structure.

\begin{proposition}[Unitary spatial reconstruction]
\label{prop:n9v88-spatial-unitary}
If $\omega\circ\alpha_g=\omega$,
$\alpha_g\mathcal A_+=\mathcal A_+$, and
$\Theta\alpha_g=\alpha_g\Theta$, then
\begin{equation}
 U(g)[F]:=[\alpha_gF]
 \label{eq:n9v88-Ug}
\end{equation}
is a well-defined unitary on $\mathcal H_{\rm OS}$.
\end{proposition}

\begin{proof}
The hypotheses give
$(\alpha_gF,\alpha_gG)_{\rm OS}=(F,G)_{\rm OS}$, so null vectors are
preserved and the quotient map is isometric.  Apply the inverse
$\alpha_{g^{-1}}$ for surjectivity, then extend to the completion.
\end{proof}

\begin{firewall}[Full Poincar\'e reconstruction needs more]
\label{fire:n9v88-Poincare}
The Hamiltonian and spatial symmetry operators above do not by
themselves construct Lorentz boosts or prove the spectrum condition for
the full energy--momentum vector.  That step requires full Euclidean
covariance, regularity, analytic continuation, and the complete OS
axioms on the limiting Schwinger functions.
\end{firewall}

\subsection{Bounded observable representation and field domains}

Let $B\in\mathcal A_+$ be such that left multiplication preserves the
null space:
\begin{equation}
 B\mathcal N\subset\mathcal N.
 \label{eq:n9v88-null-multiplier}
\end{equation}

\begin{proposition}[Null-compatible observable]
\label{prop:n9v88-observable}
If, in addition,
\begin{equation}
 \lVert[BF]\rVert_{\rm OS}
 \le C_B\lVert[F]\rVert_{\rm OS},
 \label{eq:n9v88-B-bound}
\end{equation}
then
\begin{equation}
 \pi(B)[F]:=[BF]
 \label{eq:n9v88-piB}
\end{equation}
defines a bounded operator with
$\lVert\pi(B)\rVert\le C_B$.
\end{proposition}

\begin{proof}
Equation \eqref{eq:n9v88-null-multiplier} makes the map independent of
the representative, and \eqref{eq:n9v88-B-bound} extends it continuously
to the Hilbert completion.
\end{proof}

\begin{firewall}[Unbounded fields require closability]
\label{fire:n9v88-unbounded-fields}
Smeared metric, gauge, Higgs, and fermion fields are generally
unbounded.  Their polynomial action on local equivalence classes needs
a common dense invariant domain, null compatibility, adjoint relations,
and closability.  Moment convergence alone does not establish these
operator-domain properties.
\end{firewall}

\subsection{Constraint and Ward descent}

One-point vanishing of a constraint is weaker than vanishing of its
operator insertion.  Let $C$ be a regulated-to-limit local constraint
insertion acting on the reconstructed domain.

\begin{theorem}[Weak constraint operator criterion]
\label{thm:n9v88-constraint-operator}
Suppose for all local positive-time $F,G$ in a dense invariant domain,
\begin{equation}
 \omega((\Theta F)CG)=0.
 \label{eq:n9v88-constraint-matrix-elements}
\end{equation}
If multiplication by $C$ is null compatible, then the induced operator
$\pi(C)$ vanishes on that domain.
\end{theorem}

\begin{proof}
Equation \eqref{eq:n9v88-constraint-matrix-elements} says
$\langle[F],\pi(C)[G]\rangle=0$ for every vector in a dense set.
Therefore $\pi(C)[G]=0$ for every $G$ in the domain.
\end{proof}

\begin{firewall}[Expectation zero is insufficient]
\label{fire:n9v88-one-point}
The identity $\omega(C)=0$ gives only the vacuum expectation of the
constraint.  It does not imply \eqref{eq:n9v88-constraint-matrix-elements}
and cannot establish the physical constraint operator equation.
The V87 Ward/constraint residual must vanish with arbitrary local
insertions.
\end{firewall}

Let $Q$ be a densely defined nilpotent BRST operator on a graded dense
domain, $Q^2=0$.

\begin{proposition}[BRST and transfer compatibility]
\label{prop:n9v88-BRST-transfer}
Suppose $Q$ preserves the OS null space, the exact subspace
$\operatorname{ran}Q$ is null in the closed subspace $\ker Q$, and
\begin{equation}
 T(t)Q=QT(t).
 \label{eq:n9v88-TQ}
\end{equation}
Then the physical cohomology
\begin{equation}
 \mathcal H_{\rm phys}
 :=\overline{\ker Q/\operatorname{ran}Q}
 \label{eq:n9v88-Hphys}
\end{equation}
inherits a positive inner product and a contraction semigroup.  Its
generator is the restriction/quotient of $H$.
\end{proposition}

\begin{proof}
The radical hypothesis makes the inner product independent of the
cohomology representative.  Commutation sends closed vectors to closed
vectors and exact vectors to exact vectors, so $T(t)$ descends.  Its
contraction and semigroup properties survive the quotient and
completion; the generator is induced by strong continuity.
\end{proof}

\begin{firewall}[Order of quotients]
\label{fire:n9v88-quotient-order}
The OS null quotient and BRST cohomology can be interchanged only after
the two null structures and $Q$ are shown compatible.  Ghost Berezin
integration, an oriented determinant cancellation, and formal
nilpotence do not by themselves prove the Hilbert-space hypotheses of
\Cref{prop:n9v88-BRST-transfer}.
\end{firewall}

\subsection{Clustering, the vacuum, and a mass gap}

Let
\begin{equation}
 P_0:=\mathbf1_{\{0\}}(H)
 \label{eq:n9v88-P0}
\end{equation}
be the spectral projection onto the zero-energy space.

\begin{theorem}[Long-time transfer limit]
\label{thm:n9v88-long-time}
For every $f\in\mathcal H_{\rm OS}$,
\begin{equation}
 \lim_{t\to\infty}T(t)f=P_0f
 \label{eq:n9v88-long-time}
\end{equation}
strongly.  Hence connected time-separated correlations cluster to the
projection onto $\ker H$.  They factor through the vacuum vector alone
precisely when $P_0=|\Omega\rangle\langle\Omega|$ on the cyclic sector.
\end{theorem}

\begin{proof}
By the spectral theorem,
$T(t)-P_0$ is multiplication by $e^{-t\lambda}$ for $\lambda>0$ and by
zero at $\lambda=0$.  Dominated convergence gives strong convergence.
The factorization statement is the characterization of a rank-one
zero-energy projection in the cyclic sector.
\end{proof}

\begin{theorem}[Spectral-gap equivalence]
\label{thm:n9v88-gap-equivalence}
Let $P_\perp=I-P_0$.  For $m>0$, the following are equivalent:
\begin{enumerate}
\item $\operatorname{spec}(H|_{P_\perp\mathcal H})\subset[m,\infty)$;
\item for every $t\ge0$,
\begin{equation}
 \lVert T(t)P_\perp\rVert\le e^{-mt};
 \label{eq:n9v88-gap-decay}
\end{equation}
\item for all $f,g\in P_\perp\mathcal H$,
\begin{equation}
 |\langle f,T(t)g\rangle|
 \le e^{-mt}\lVert f\rVert\lVert g\rVert.
 \label{eq:n9v88-correlation-gap}
\end{equation}
\end{enumerate}
\end{theorem}

\begin{proof}
The spectral theorem gives (1)$\Rightarrow$(2), and operator norm
duality gives (2)$\Leftrightarrow$(3).  If (2) holds but the spectral
measure has support below $m$, then
$\sup_{\lambda\in\operatorname{spec}(H|_{P_\perp})}e^{-t\lambda}$
exceeds $e^{-mt}$ for some $t$, a contradiction.  Thus (2) implies
(1).
\end{proof}

\begin{firewall}[Clustering is not a mass gap]
\label{fire:n9v88-cluster-not-gap}
Strong convergence $T(t)\to P_0$ follows without a positive spectral
gap and can be arbitrarily slow.  A mass gap requires the uniform
exponential estimate \eqref{eq:n9v88-gap-decay}, after removal of all
vacuum, gauge, topological, and massless physical sectors appropriate to
the claimed theorem.
\end{firewall}

For the full SM--Einstein continuum, massless gravitons and the photon
are expected physical sectors rather than errors.  A global positive
gap above the vacuum is therefore not the target.  Spectral statements
must be sector-specific, with the YM colour-singlet gap distinguished
from massless gravitational and electromagnetic branches.

\subsection{OS reconstruction card}

\begin{definition}[OSRC]
\label{def:n9v88-OSRC}
The OS reconstruction card requires:
\begin{enumerate}
\item an RCC limiting local graded functional with Hermitian reflection
      form;
\item the OS null quotient \eqref{eq:n9v88-H0};
\item reflection-compatible time translations preserving the positive
      algebra and null space;
\item contraction, symmetry, and strong continuity of those
      translations;
\item null-compatible bounded observables and closable smeared-field
      domains;
\item all-insertion constraint and Ward identities;
\item a finite-cutoff-to-limit BRST charge satisfying the Hilbert-space
      descent hypotheses;
\item full Euclidean covariance and regularity for analytic
      continuation;
\item uniqueness/clustering only where a vacuum claim is made; and
\item sector-specific exponential estimates for every claimed mass
      gap.
\end{enumerate}
\end{definition}

\begin{theorem}[OSRC gives a conditional physical transfer theory]
\label{thm:n9v88-OSRC}
Under OSRC, the limiting local reflection form reconstructs a Hilbert
space, vacuum vector, and nonnegative Hamiltonian with
$T(t)=e^{-tH}$.  Constraint insertions vanish as operators when
\eqref{eq:n9v88-constraint-matrix-elements} holds, and the compatible
BRST quotient carries the physical contraction semigroup.  Vacuum
uniqueness and sector gaps follow only from their respective clustering
and exponential spectral gates.
\end{theorem}

\begin{proof}
Use \Cref{thm:n9v88-Hilbert,thm:n9v88-semigroup,
thm:n9v88-constraint-operator,prop:n9v88-BRST-transfer,
thm:n9v88-long-time,thm:n9v88-gap-equivalence}.
\end{proof}

\begin{theorem}[Ninth-series V88 result]
\label{thm:n9v88-result}
A Hermitian positive limiting reflection form gives the exact quotient
$\mathcal H_0=\mathcal A_+/\mathcal N$.  Under explicit null-
preservation, contraction, symmetry, and continuity gates, Euclidean
time translations reconstruct $e^{-tH}$ with $H\ge0$.  Constraint
vanishing requires all matrix elements, BRST descent requires
compatibility with the OS null space, and a spectral gap is equivalent
to uniform exponential decay on the orthogonal spectral sector.
Ordinary clustering alone does not prove such a gap.
\end{theorem}

\begin{proof}
Combine \Cref{thm:n9v88-Hilbert,thm:n9v88-semigroup,
thm:n9v88-constraint-operator,prop:n9v88-BRST-transfer,
thm:n9v88-gap-equivalence}.
\end{proof}

\begin{assessment}[Progress at ninth-series V88]
\label{ass:n9v88-status}
The local limit of V87 now has a precise reconstruction target.  The
main remaining analytic gaps are no longer hidden inside the phrase
``OS reconstruction'': time-translation contraction, unbounded-field
domains, all-insertion constraints, BRST compatibility, Euclidean
analytic continuation, and sector-specific spectral estimates are
separate gates.
\end{assessment}

\begin{programme}[Ninth-series V89: Euclidean continuation and sectors]
\label{prog:n9v88-V89}
V89 will organize the limiting Schwinger axioms needed beyond the
half-space Hilbert construction, prove passage of graded symmetry and
Euclidean covariance through the RCC comparison, and formulate the
sector decomposition that keeps photon and graviton branches massless
while allowing a colour-singlet YM gap claim.
\end{programme}

\begin{firewall}[Claim boundary after ninth-series V88]
\label{fire:n9v88-boundary}
V88 does not prove the RCC hypotheses, time-translation contraction,
field closability, a limiting BRST charge, anomaly cancellation, full
Euclidean covariance, analytic continuation, clustering, any spectral
gap, or the full SM--Einstein continuum.
\end{firewall}

\paragraph{Parent-version record.}
V88 was formed from
\texttt{Renewal\_SM\_Einstein\_Continuum\_Research\_Notes\_V87.tex}.
The V87 parent had $34\,635$ counted source lines, $1\,244\,803$ bytes,
and SHA--256
\[
 \texttt{0C66A277EDF5164ED3DC72EB7B8CCF5395C01AB35C110119A42ABB927E2B557B}.
\]

% =============================================================================
% END APPENDED NINTH-SERIES V88 MATERIAL
% =============================================================================

% =============================================================================
% BEGIN APPENDED NINTH-SERIES V89 MATERIAL
% =============================================================================

\section{Ninth-series V89: Euclidean identities and spectral sectors}
\label{sec:v9v89}

V88 reconstructs a conditional Hamiltonian from the limiting
half-space functional.  This note proves how graded symmetry and
Euclidean covariance pass through the V87 comparison, derives the
two-point spectral Laplace representation, and formulates sector gaps
without contradicting the massless photon and graviton branches of the
full coupled theory.

\subsection{Graded permutation symmetry in the limit}

Let $S_n^{(k)}(f_1,\ldots,f_k)$ be the regulated multilinear Schwinger
form, with homogeneous test fields of parities $|f_j|\in\{0,1\}$.

\begin{proposition}[Passage of graded symmetry]
\label{prop:n9v89-graded-limit}
Suppose for every adjacent transposition,
\begin{equation}
 S_n^{(k)}(\ldots,f_i,f_{i+1},\ldots)
 =(-1)^{|f_i||f_{i+1}|}
 S_n^{(k)}(\ldots,f_{i+1},f_i,\ldots)
 \label{eq:n9v89-graded-symmetry}
\end{equation}
and both sides converge pointwise under the common reconstruction of
V87.  Then the limiting form satisfies the same identity and hence the
full Koszul permutation rule.
\end{proposition}

\begin{proof}
Take the limit in \eqref{eq:n9v89-graded-symmetry}.  Adjacent
transpositions generate the permutation group, and multiplication of
their parity signs gives the Koszul sign.
\end{proof}

\begin{firewall}[Regulator ordering conventions]
\label{fire:n9v89-ordering}
The common reconstruction must preserve the ordering of Grassmann
generators and the reflection twist.  A cutoff-dependent reordering can
insert signs into the comparison map; these belong to the V87
reconstruction row and cannot be ignored after taking the limit.
\end{firewall}

\subsection{Euclidean covariance with a restoration defect}

Let $G_E$ be the target Euclidean group or a connected subgroup on
which all test functions remain in the declared charts.  Write
$g\cdot f$ for the tensor/spinor action on test fields.  Define the
regulated covariance residual
\begin{equation}
 \mathfrak C_{n,g}^{(k)}(f_1,\ldots,f_k)
 :=S_n^{(k)}(g\cdot f_1,\ldots,g\cdot f_k)
   -S_n^{(k)}(f_1,\ldots,f_k).
 \label{eq:n9v89-covariance-residual}
\end{equation}

\begin{theorem}[Covariance restoration criterion]
\label{thm:n9v89-covariance-limit}
Suppose $S_n^{(k)}\to S^{(k)}$ in the V87 distributional sense and,
for every compact $K\subset G_E$ and bounded test set $B\subset
\mathcal D^k$,
\begin{equation}
 \sup_{g\in K}\sup_{(f_1,\ldots,f_k)\in B}
 |\mathfrak C_{n,g}^{(k)}(f_1,\ldots,f_k)|
 \longrightarrow0.
 \label{eq:n9v89-covariance-vanish}
\end{equation}
Then the limiting Schwinger form is $G_E$-covariant:
\begin{equation}
 S^{(k)}(g\cdot f_1,\ldots,g\cdot f_k)
 =S^{(k)}(f_1,\ldots,f_k).
 \label{eq:n9v89-limit-covariance}
\end{equation}
The convergence is locally uniform in $g$ on the tested bounded sets.
\end{theorem}

\begin{proof}
For fixed $g$ and test tuple, add and subtract the two regulated forms:
\begin{align*}
 &S^{(k)}(g\cdot f)-S^{(k)}(f)\\
 &=[S^{(k)}(g\cdot f)-S_n^{(k)}(g\cdot f)]
   +\mathfrak C_{n,g}^{(k)}(f)
   +[S_n^{(k)}(f)-S^{(k)}(f)].
\end{align*}
The first and third terms vanish by distributional convergence and the
middle by \eqref{eq:n9v89-covariance-vanish}.  Uniform versions of the
same three bounds give local uniformity.
\end{proof}

\begin{corollary}[Exact subgroup plus broken generators]
\label{cor:n9v89-subgroup-restoration}
If the regulator preserves a subgroup $G_n\subset G_E$ exactly, only
transformations outside that subgroup require estimates in
\eqref{eq:n9v89-covariance-vanish}.  Every finite-cutoff anisotropic
marginal owner from V80 must nevertheless be calibrated or proved to
vanish; exact covariance of a small subgroup does not remove it.
\end{corollary}

\begin{proof}
For $g\in G_n$ the residual is zero.  The second statement follows from
\Cref{prop:n9v80-full-owner}: a noncontracting anisotropic direction is
not controlled by stable convergence.
\end{proof}

\begin{firewall}[Pointwise restoration is too weak for generators]
\label{fire:n9v89-generator-restoration}
Pointwise convergence in $g$ can prove invariance at each fixed group
element but need not justify differentiation at the identity.  Angular
momentum, boosts after continuation, and Ward generators require the
locally uniform or differentiable control in
\eqref{eq:n9v89-covariance-vanish}.
\end{firewall}

\subsection{Two-point spectral representation}

Assume the V88 translation gates and let $H\ge0$ be the reconstructed
Hamiltonian.  For $f,g\in\mathcal H_{\rm OS}$ define
\begin{equation}
 C_{f,g}(t):=\langle f,e^{-tH}g\rangle,
 \qquad t\ge0.
 \label{eq:n9v89-Cfg}
\end{equation}

\begin{theorem}[Laplace spectral measure]
\label{thm:n9v89-Laplace}
There is a finite complex measure
\begin{equation}
 \mu_{f,g}(B):=\langle f,E_H(B)g\rangle
 \label{eq:n9v89-mu}
\end{equation}
on $[0,\infty)$ such that
\begin{equation}
 C_{f,g}(t)=\int_{[0,\infty)}e^{-t\lambda}
 \,d\mu_{f,g}(\lambda).
 \label{eq:n9v89-Laplace}
\end{equation}
For $f=g$, $\mu_{f,f}$ is positive.  The function
\begin{equation}
 C_{f,g}(z)=\int e^{-z\lambda}\,d\mu_{f,g}(\lambda)
 \label{eq:n9v89-holomorphic}
\end{equation}
is holomorphic for $\operatorname{Re}z>0$ and bounded there by
$\lVert f\rVert\lVert g\rVert$.
\end{theorem}

\begin{proof}
Apply the spectral theorem to $H$.  The total variation of
$\mu_{f,g}$ is at most $\lVert f\rVert\lVert g\rVert$.  The Laplace
formula follows from functional calculus.  On every closed right
half-plane $\operatorname{Re}z\ge\epsilon>0$, differentiation under the
integral is justified by the bounded functions
$\lambda^ke^{-\epsilon\lambda}$, proving holomorphy.  The bound follows
from $|e^{-z\lambda}|\le1$.
\end{proof}

\begin{corollary}[Gap and two-point decay]
\label{cor:n9v89-gap-decay}
If $f,g$ have spectral support in $[m,\infty)$, then
\begin{equation}
 |C_{f,g}(t)|\le e^{-mt}\lVert f\rVert\lVert g\rVert.
 \label{eq:n9v89-two-point-decay}
\end{equation}
Conversely, a uniform bound of this form for every $f,g$ in a closed
$H$-invariant subspace implies that the spectrum of $H$ on that
subspace lies in $[m,\infty)$.
\end{corollary}

\begin{proof}
The first statement follows from
\eqref{eq:n9v89-Laplace}; the converse is
\Cref{thm:n9v88-gap-equivalence} restricted to the subspace.
\end{proof}

\begin{firewall}[Higher-point continuation]
\label{fire:n9v89-higher-continuation}
The two-point Laplace representation follows from the reconstructed
semigroup.  Analytic continuation of all time-ordered $k$-point
functions requires compatible operator domains, energy bounds, graded
locality, and tube analyticity.  It cannot be inferred from
\eqref{eq:n9v89-holomorphic} alone.
\end{firewall}

\subsection{Operator-generated spectral sectors}

Let $\mathcal O_s$ be a declared family of physical local observables
after the OS and BRST quotients.  Define its vacuum-generated subspace
\begin{equation}
 \mathcal H_s:=overline{\operatorname{span}}
 \{f(H)O\Omega: O\in\mathcal O_s,
 \ f\in C_c([0,\infty))\}.
 \label{eq:n9v89-Hs}
\end{equation}
This definition makes $\mathcal H_s$ invariant under bounded functions
of $H$.

\begin{definition}[Observable-sector gap]
\label{def:n9v89-sector-gap}
The family $\mathcal O_s$ has a Hamiltonian gap $m_s>0$ above its
zero-energy component if
\begin{equation}
 E_H((0,m_s))\mathcal H_s=\{0\}.
 \label{eq:n9v89-sector-gap}
\end{equation}
\end{definition}

\begin{proposition}[Sector gap criterion]
\label{prop:n9v89-sector-gap}
The gap \eqref{eq:n9v89-sector-gap} holds if and only if
\begin{equation}
 \lVert e^{-tH}(I-P_0)|_{\mathcal H_s}\rVert
 \le e^{-m_st}
 \quad(t\ge0).
 \label{eq:n9v89-sector-decay}
\end{equation}
Equivalently all connected two-point functions of vectors in
$\mathcal H_s$ obey the corresponding uniform exponential bound.
\end{proposition}

\begin{proof}
Apply \Cref{thm:n9v88-gap-equivalence} to the invariant subspace
$\mathcal H_s$.
\end{proof}

\begin{firewall}[An operator label is not a superselection sector]
\label{fire:n9v89-label-not-sector}
Calling $O$ a ``Yang--Mills observable'' does not make
$\mathcal H_s$ invariant under the full coupled dynamics.  Gauge,
Higgs, fermion, and gravitational interactions can mix states with the
same exact quantum numbers.  The definition
\eqref{eq:n9v89-Hs} includes all spectral states actually reached by the
operator family.
\end{firewall}

\subsection{Massless contamination and the coupled-theory no-go}

\begin{theorem}[Massless-overlap obstruction]
\label{thm:n9v89-massless-obstruction}
Let $O$ be a physical observable and put
$v=(I-P_0)O\Omega$.  If
\begin{equation}
 E_H((0,\epsilon))v\ne0
 \quad\text{for every }\epsilon>0,
 \label{eq:n9v89-low-overlap}
\end{equation}
then no positive $m$ satisfies an exponential bound
\begin{equation}
 \lVert e^{-tH}v\rVert\le Ce^{-mt}
 \quad(t\ge0)
 \label{eq:n9v89-false-gap}
\end{equation}
with finite $C$ independent of $t$.
\end{theorem}

\begin{proof}
If \eqref{eq:n9v89-false-gap} held, the spectral theorem applied to
$\lVert e^{-tH}v\rVert^2$ would force the spectral measure of $v$ to be
supported in $[m,\infty)$: any positive measure in $(0,m)$ decays more
slowly than $e^{-2mt}$.  This contradicts
\eqref{eq:n9v89-low-overlap}.
\end{proof}

\begin{corollary}[No global SM--Einstein vacuum gap]
\label{cor:n9v89-no-global-gap}
If the reconstructed physical theory contains photon or graviton states
with energy arbitrarily close to zero and nonzero local physical
overlap, then the full orthogonal complement of the vacuum has no
strict Hamiltonian gap.  A gap claim must exclude or project away these
massless branches by an exact invariant-sector construction.
\end{corollary}

\begin{proof}
Choose an observable vector with the massless spectral overlap and apply
\Cref{thm:n9v89-massless-obstruction}.
\end{proof}

\begin{firewall}[Gravity can contaminate a nominal YM channel]
\label{fire:n9v89-gravity-contamination}
In a coupled theory, a gauge-invariant colour-singlet operator can share
all exact quantum numbers with multiphoton or multigraviton states.
Unless an exact commuting projection or selection rule removes those
states, its spectral measure may extend to zero.  The isolated pure-YM
mass-gap statement does not transfer automatically to the fully coupled
SM--Einstein spectrum.
\end{firewall}

\subsection{Exact invariant projections}

Let $P_s$ be an orthogonal projection on the physical Hilbert space.

\begin{proposition}[Projected spectral sector]
\label{prop:n9v89-projected-sector}
If
\begin{equation}
 [P_s,H]=0,
 \qquad P_s\Omega=0,
 \label{eq:n9v89-projector-commute}
\end{equation}
then $P_s\mathcal H$ is $H$-invariant and has a gap $m_s$ precisely when
\begin{equation}
 P_sE_H((0,m_s))=0.
 \label{eq:n9v89-projected-gap}
\end{equation}
If $P_s$ is built from an exact conserved charge, different spectral
values of that charge give orthogonal invariant sectors.
\end{proposition}

\begin{proof}
Commutation with $H$ implies commutation with every spectral projection
of $H$, so the range is invariant and
\eqref{eq:n9v89-projected-gap} is exactly the spectral gap definition.
Spectral projections of a self-adjoint conserved charge are mutually
orthogonal and commute with $H$.
\end{proof}

\begin{firewall}[Colour is constrained, not a physical global charge]
\label{fire:n9v89-colour-charge}
The local nonabelian gauge constraint does not supply a physical colour
superselection projection by the elementary conserved-charge argument.
The physical algebra is gauge invariant after BRST/constraint descent.
A Yang--Mills gap in the coupled theory must therefore be formulated by
the spectral support of an appropriate gauge-invariant observable
family or another rigorously constructed invariant subspace.
\end{firewall}

\subsection{Mass and momentum}

If the full OS continuation reconstructs commuting self-adjoint spatial
momenta $\mathbf P$ with the relativistic spectrum condition, define
\begin{equation}
 M^2:=H^2-\mathbf P^2
 \label{eq:n9v89-mass-operator}
\end{equation}
by joint functional calculus.

\begin{proposition}[Hamiltonian gap versus invariant mass]
\label{prop:n9v89-H-vs-mass}
A lower bound $H\ge m$ on a sector implies only an energy gap in the
chosen frame.  Identification of $m$ as a particle mass requires the
joint energy--momentum representation and spectral information about
$M^2$.  Conversely a massive relativistic branch has
$H(\mathbf p)=\sqrt{m^2+|\mathbf p|^2}$ and therefore energy at least
$m$, but multiparticle thresholds and massless contamination must be
analyzed in the same joint spectrum.
\end{proposition}

\begin{proof}
The first statement follows because $H$ alone contains no momentum
information.  The second is the relativistic dispersion relation on a
mass-$m$ spectral branch; joint functional calculus distinguishes it
from other branches and thresholds.
\end{proof}

\subsection{Euclidean and spectral sector card}

\begin{definition}[ESSC]
\label{def:n9v89-ESSC}
The Euclidean and spectral sector card requires:
\begin{enumerate}
\item all-order V87 distributional Schwinger limits;
\item exact limiting graded symmetry;
\item locally uniform vanishing of every Euclidean covariance defect;
\item the V88 Hilbert and time-translation gates;
\item common invariant domains for higher-point analytic continuation;
\item full Euclidean covariance and the OS regularity/growth axioms;
\item a physical BRST/constraint quotient before sectors are named;
\item operator-generated or exactly projected invariant sectors;
\item explicit removal or accounting of photon and graviton massless
      spectral overlap; and
\item a joint energy--momentum analysis for every particle-mass claim.
\end{enumerate}
\end{definition}

\begin{theorem}[ESSC yields type-correct spectral claims]
\label{thm:n9v89-ESSC}
Under ESSC, graded symmetry and Euclidean covariance pass to the
regulator-independent Schwinger functions.  The reconstructed
two-point functions have the spectral representation
\eqref{eq:n9v89-Laplace}.  Every claimed gap is attached to an
operator-generated or invariant projected subspace and is equivalent to
the exponential estimate \eqref{eq:n9v89-sector-decay}.  The full
SM--Einstein Hilbert space has no global gap when a physical photon or
graviton branch reaches zero energy.
\end{theorem}

\begin{proof}
Use \Cref{prop:n9v89-graded-limit,
thm:n9v89-covariance-limit,thm:n9v89-Laplace,
prop:n9v89-sector-gap,thm:n9v89-massless-obstruction,
prop:n9v89-projected-sector}.
\end{proof}

\begin{theorem}[Ninth-series V89 result]
\label{thm:n9v89-result}
Graded permutation symmetry passes directly through the V87 local
limit, while Euclidean covariance requires the uniform restoration
estimate \eqref{eq:n9v89-covariance-vanish}.  The V88 Hamiltonian gives
the positive spectral Laplace measure \eqref{eq:n9v89-Laplace}.  A
sector gap is an absence of spectral support in $(0,m)$ on a precisely
defined invariant or operator-generated subspace.  Massless photon or
graviton overlap obstructs a global gap and can also contaminate a
nominal Yang--Mills colour-singlet channel unless an exact selection
rule removes it.
\end{theorem}

\begin{proof}
Combine \Cref{prop:n9v89-graded-limit,
thm:n9v89-covariance-limit,thm:n9v89-Laplace,
thm:n9v89-massless-obstruction,prop:n9v89-projected-sector}.
\end{proof}

\begin{assessment}[Progress at ninth-series V89]
\label{ass:n9v89-status}
The continuum spectral target is now compatible with the actual coupled
field content.  Covariance restoration is a quantitative regulator
row, analytic two-point continuation follows from the reconstructed
semigroup, and gap claims are sector-specific.  The next compulsory
audit will compare the newest YM definition of its gap and stable
reference with this massless-contamination constraint before the final
ten-version acquisition block is chosen.
\end{assessment}

\begin{programme}[Ninth-series V90: compulsory companion audit]
\label{prog:n9v89-V90}
V90 will inspect the newest YM and NS notes, compare their current
reference-graph, contraction, and spectral targets with ESSC, and select
the V91 direction.  V91 will continue after the audit.
\end{programme}

\begin{firewall}[Claim boundary after ninth-series V89]
\label{fire:n9v89-boundary}
V89 does not prove covariance restoration, higher-point analytic
continuation, a physical BRST quotient, clustering, any sector gap,
decoupling from photon/graviton continua, the joint energy--momentum
spectrum, regulator independence, or the full SM--Einstein continuum.
\end{firewall}

\paragraph{Parent-version record.}
V89 was formed from
\texttt{Renewal\_SM\_Einstein\_Continuum\_Research\_Notes\_V88.tex}.
The V88 parent had $35\,129$ counted source lines, $1\,262\,364$ bytes,
and SHA--256
\[
 \texttt{006B3352C2072A5E33A2AEC7320F37D1882CD1D5FA29CB980F60F1BDEED3679A}.
\]

% =============================================================================
% END APPENDED NINTH-SERIES V89 MATERIAL
% =============================================================================

% =============================================================================
% BEGIN APPENDED NINTH-SERIES V90 MATERIAL
% =============================================================================

\section{Ninth-series V90: companion audit and reconstruction-lift defects}
\label{sec:v9v90}

\subsection{Compulsory companion audit}

The newest companion files present at this boundary were inspected
read-only:
\begin{align*}
 &\texttt{Renewal\_Yang\_Mills\_Mass\_Gap\_Research\_Notes\_V14.tex},\\
 &\texttt{Renewal\_Yang\_Mills\_Mass\_Gap\_Research\_Notes\_V15.tex},\\
 &\texttt{Renewal\_NS\_Stability\_Research\_Notes\_V26.tex}.
\end{align*}
Their exact records were
\begin{center}
\begin{tabular}{c|r|l}
file & bytes & SHA--256\\ \hline
YM V14 & $237\,592$ &
\texttt{56FEE5EE73F12EDBF70AC1CD8CDB0320E6EA1D1839D38D6EFD83E399510C8E7B}\\
YM V15 & $237\,592$ &
\texttt{56FEE5EE73F12EDBF70AC1CD8CDB0320E6EA1D1839D38D6EFD83E399510C8E7B}\\
NS V26 & $278\,283$ &
\texttt{82C887F8C2C69E4F28FAD7C3DC8DE5B9099CB5A4BF431EBA1FFF3E91935978AD}
\end{tabular}
\end{center}
YM V14 and V15 were byte-identical at inspection time, so the
substantive endpoint was V14.  It replaces an equal-face energy-lift
assumption by the exact vertical correction $M=-A^{-1}K$, bounds that
correction by a relative singular value, propagates it through tensor
powers, and gives a contraction margin tolerant of a nonzero lift
defect.  Its remaining input is a finite certified physical Wilson
matrix.  NS V26 remains unchanged and supplies no regulator-lift or
Schwinger comparison row.

\begin{decision}[V90 audit transfer]
\label{dec:n9v90-transfer}
The tensor-power lift-defect estimate is transferred to the V87
reconstruction row.  A summable physical-versus-reference lift defect
will imply summable reconstruction errors at every fixed Schwinger
order.  The YM numerical matrices and face constants are not imported.
Reflection, chirality, covariance, and constraint intertwining defects
remain separate SM estimates.  No NS result is imported.
\end{decision}

\subsection{Two reconstruction lifts}

Let $X$ and $Y$ be finite-dimensional normed field spaces at two
successive regulators.  Let
\begin{equation}
 J_0:Y\to X,
 \qquad J:Y\to X,
 \qquad M:=J-J_0
 \label{eq:n9v90-lifts}
\end{equation}
be a reference reconstruction and the physical reconstruction.  Put
\begin{equation}
 j_\#:=\max\{\lVert J\rVert,\lVert J_0\rVert\},
 \qquad m_*:=\lVert M\rVert.
 \label{eq:n9v90-jm}
\end{equation}

\begin{lemma}[Multilinear lift telescope]
\label{lem:n9v90-tensor-telescope}
For $k\ge1$,
\begin{equation}
 J^{\otimes k}-J_0^{\otimes k}
 =\sum_{r=1}^k
 J^{\otimes(r-1)}\otimes M\otimes
 J_0^{\otimes(k-r)},
 \label{eq:n9v90-tensor-identity}
\end{equation}
and therefore
\begin{equation}
 \lVert J^{\otimes k}-J_0^{\otimes k}\rVert
 \le k j_\#^{k-1}m_*.
 \label{eq:n9v90-tensor-bound}
\end{equation}
\end{lemma}

\begin{proof}
Insert $J-J_0=M$ one tensor position at a time; adjacent intermediate
tensor products telescope.  Bound each of the $k$ terms by
$j_\#^{k-1}m_*$.
\end{proof}

\begin{proposition}[Schwinger reconstruction row]
\label{prop:n9v90-Schwinger-row}
Let $S_X^{(k)}$ be a continuous $k$-linear form on $X$ with
\begin{equation}
 |S_X^{(k)}(x_1,\ldots,x_k)|
 \le C_k\prod_{r=1}^k\lVert x_r\rVert.
 \label{eq:n9v90-Sk-bound}
\end{equation}
Then
\begin{align}
 &|S_X^{(k)}(Jy_1,\ldots,Jy_k)
 -S_X^{(k)}(J_0y_1,\ldots,J_0y_k)|\\
 &\qquad\le
 C_k k j_\#^{k-1}m_*
 \prod_{r=1}^k\lVert y_r\rVert.
 \label{eq:n9v90-Schwinger-row}
\end{align}
\end{proposition}

\begin{proof}
View $S_X^{(k)}$ as a bounded linear functional on the projective
$k$-fold tensor product and apply
\eqref{eq:n9v90-tensor-bound} to
$y_1\otimes\cdots\otimes y_k$.
\end{proof}

\begin{corollary}[Summable lift criterion at fixed order]
\label{cor:n9v90-summable-lift}
For a regulator tower, suppose at order $k$ that
\begin{equation}
 \sup_n C_{n,k}<\infty,
 \qquad
 \sup_n j_{\#,n}\le J_*<\infty,
 \qquad
 \sum_nm_{*,n}<\infty.
 \label{eq:n9v90-summable-hyp}
\end{equation}
Then the physical-versus-reference lift contribution to the V87
reconstruction row is summable:
\begin{equation}
 \sum_n\epsilon_{n,k}^{\rm lift}
 \le kJ_*^{k-1}\sup_nC_{n,k}
 \sum_nm_{*,n}<\infty.
 \label{eq:n9v90-summable-row}
\end{equation}
\end{corollary}

\begin{proof}
Apply \Cref{prop:n9v90-Schwinger-row} at each step and sum.
\end{proof}

\begin{firewall}[Fixed order is not an all-order generating bound]
\label{fire:n9v90-fixed-order}
The factor $kJ_*^{k-1}C_{n,k}$ may grow rapidly with $k$.  Summability
for every fixed order does not give uniform convergence of the
Schwinger generating series.  A characteristic-functional argument
requires its own equicontinuity and positivity estimates.
\end{firewall}

\subsection{Energy-selected lift defect}

Let $Q:X\to Y$ be surjective, $QJ_0=I$, and let $H>0$ be a physical
quadratic form on $X$.  Set $E=\ker Q$ and write the $H$ blocks relative
to $X=E\oplus J_0Y$ as
\begin{equation}
 A:=H|_{E\times E},
 \qquad K:=P_E^*HJ_0,
 \qquad D:=J_0^*HJ_0.
 \label{eq:n9v90-AKD}
\end{equation}

\begin{proposition}[Physical reconstruction correction]
\label{prop:n9v90-energy-lift}
The unique $H$-minimum lift is
\begin{equation}
 J_H=J_0-A^{-1}K,
 \qquad M_H:=-A^{-1}K:Y\to E.
 \label{eq:n9v90-JH}
\end{equation}
With
\begin{equation}
 \eta_H:=\lVert A^{-1/2}KD^{-1/2}\rVert<1,
 \label{eq:n9v90-etaH}
\end{equation}
and bounds $A\ge\lambda_AI$, $D\le\Lambda_DI$, one has
\begin{equation}
 \lVert M_H\rVert
 \le\eta_H\sqrt{\Lambda_D/\lambda_A}.
 \label{eq:n9v90-MH-bound}
\end{equation}
\end{proposition}

\begin{proof}
Completing the square gives
\begin{align*}
 \langle z+J_0y,H(z+J_0y)\rangle
 &=\langle z+A^{-1}Ky,A(z+A^{-1}Ky)\rangle\\
 &\quad+\langle y,(D-K^*A^{-1}K)y\rangle.
\end{align*}
The minimizer is $z=-A^{-1}Ky$.  Positivity of the Schur complement
implies \eqref{eq:n9v90-etaH}; writing
$M_H=-A^{-1/2}(A^{-1/2}KD^{-1/2})D^{1/2}$ gives the bound.
\end{proof}

This is the V75 energy lift in reference coordinates and the mechanism
audited in YM V14.  Its value for SM fields depends on the actual
gauge-, constraint-, zero-mode-, and boundary-reduced precision $H$.

\begin{firewall}[Positive finite matrices need a uniform certificate]
\label{fire:n9v90-uniform-certificate}
Finite-dimensional positivity gives $\eta_H<1$ at one cutoff but no
uniform separation from one and no summability of
\eqref{eq:n9v90-MH-bound}.  RCC requires interval or analytic bounds on
the actual sequence of reduced matrices.
\end{firewall}

\subsection{Intertwining defects}

Let $R_Y$ and $R_X$ denote a reflection, symmetry, constraint, or BRST
operator at the two regulators.  Define the lift intertwining defect
\begin{equation}
 \Delta_R:=R_XJ-JR_Y.
 \label{eq:n9v90-DeltaR}
\end{equation}

\begin{proposition}[Intertwining transfer bound]
\label{prop:n9v90-intertwining}
If $J=J_0+M$, then
\begin{equation}
 \Delta_R=(R_XJ_0-J_0R_Y)+(R_XM-MR_Y)
 \label{eq:n9v90-DeltaR-split}
\end{equation}
and
\begin{equation}
 \lVert\Delta_R\rVert
 \le\lVert R_XJ_0-J_0R_Y\rVert
 +(\lVert R_X\rVert+\lVert R_Y\rVert)m_*.
 \label{eq:n9v90-DeltaR-bound}
\end{equation}
\end{proposition}

\begin{proof}
Expand \eqref{eq:n9v90-DeltaR} using $J=J_0+M$ and apply the triangle
inequality.
\end{proof}

\begin{corollary}[Exact symmetry of the physical lift]
\label{cor:n9v90-exact-intertwining}
If $H$, $Q$, and $J_0$ are equivariant and the energy-minimizing lift is
unique, then $R_XJ_H=J_HR_Y$.  Thus $\Delta_R=0$ exactly for that
symmetry.
\end{corollary}

\begin{proof}
For fixed $y$, both $R_XJ_Hy$ and $J_HR_Yy$ lie over $R_Yy$ and have
the same $H$ energy by equivariance.  Uniqueness of the fibre minimizer
makes them equal.
\end{proof}

\begin{firewall}[Approximate intertwining is not exact OS covariance]
\label{fire:n9v90-approximate-reflection}
A small reflection defect may contribute a small V87 comparison row,
but each finite-cutoff OS inequality used in V86 requires an exact
reflection structure.  One must either choose an exactly equivariant
physical lift or construct positivity in a common exact reflection
representation before taking a small-defect limit.
\end{firewall}

\subsection{Robust comparison contraction}

Let $L_0$ be the ideal typed one-step comparison operator on the stable
rows and let $L=L_0+\mathcal E_H$ include the physical lift defect.

\begin{proposition}[Lift-defect tolerance]
\label{prop:n9v90-contraction-tolerance}
If in a declared weighted type norm
\begin{equation}
 \lVert L_0\rVert\le q_0<1,
 \qquad
 \lVert\mathcal E_H\rVert\le\epsilon_H<1-q_0,
 \label{eq:n9v90-contraction-hyp}
\end{equation}
then
\begin{equation}
 \lVert L\rVert\le q_0+\epsilon_H<1.
 \label{eq:n9v90-contraction}
\end{equation}
If nonlinear and marked corrections have norm
$\epsilon_{\rm nl}$, the complete strict gate is
\begin{equation}
 q_0+\epsilon_H+\epsilon_{\rm nl}<1.
 \label{eq:n9v90-complete-gate}
\end{equation}
\end{proposition}

\begin{proof}
Apply the triangle inequality once and then include the nonlinear
operator as a second perturbation.
\end{proof}

\begin{firewall}[Cycles survive norm reweighting]
\label{fire:n9v90-cycles}
Typed weights can make a strictly triangular degree-lowering block
small.  If the physical lift defect creates reverse arrows and a cycle,
its spectral contribution belongs in $\epsilon_H$.  No equivalent norm
can turn an operator of spectral radius at least one into a contraction.
\end{firewall}

\subsection{Characteristic-functional lift comparison}

Let $X$ now be a random finite-regulator bosonic field and define
\begin{equation}
 Z_J(f):=\mathbb E\,e^{i\langle Jf,X\rangle}.
 \label{eq:n9v90-ZJ}
\end{equation}

\begin{proposition}[Characteristic lift bound]
\label{prop:n9v90-characteristic-bound}
If $\mathbb E\lVert X\rVert<\infty$, then
\begin{equation}
 |Z_J(f)-Z_{J_0}(f)|
 \le m_*\lVert f\rVert\,\mathbb E\lVert X\rVert.
 \label{eq:n9v90-characteristic-bound}
\end{equation}
\end{proposition}

\begin{proof}
Use $|e^{iu}-e^{iv}|\le|u-v|$ and
\begin{align*}
 |\langle(J-J_0)f,X\rangle|
 \le m_*\lVert f\rVert\lVert X\rVert,
\end{align*}
then take expectations.
\end{proof}

\begin{corollary}[Summable characteristic reconstruction row]
\label{cor:n9v90-characteristic-sum}
If the reconstructed first moments are uniformly bounded and
$\sum_nm_{*,n}<\infty$, then the physical-versus-reference lift part of
the characteristic-functional comparison is summable uniformly on
bounded test sets.
\end{corollary}

\begin{proof}
Sum \eqref{eq:n9v90-characteristic-bound} and take the supremum over a
bounded set of $f$.
\end{proof}

\subsection{Reconstruction lift-defect card}

\begin{definition}[RLDC]
\label{def:n9v90-RLDC}
The reconstruction lift-defect card requires:
\begin{enumerate}
\item reference and physical reconstruction maps on a common typed
      field space;
\item the actual reduced energy matrix and exact minimum-energy lift;
\item a uniform relative singular-value certificate for the lift
      defect;
\item summability of $m_{*,n}$ or a moving reference graph which absorbs
      its nonvanishing part;
\item uniform fixed-order Schwinger bounds for
      \eqref{eq:n9v90-summable-row};
\item characteristic-functional equicontinuity for the all-order
      construction;
\item exact reflection intertwining at every cutoff;
\item bounded covariance, constraint, and BRST intertwining defects with
      their own owners;
\item the complete contraction gate \eqref{eq:n9v90-complete-gate}; and
\item zero-mode retention whenever an energy or Schur gap closes.
\end{enumerate}
\end{definition}

\begin{theorem}[RLDC closes the lift part of RCC]
\label{thm:n9v90-RLDC}
Under RLDC, the physical reconstruction lift contributes a summable
V87 comparison row at every fixed Schwinger order and on bounded
characteristic-functional test sets.  Exact reflection covariance is
preserved, other intertwining defects remain explicitly bounded, and
the stable comparison operator stays contractive under
\eqref{eq:n9v90-complete-gate}.
\end{theorem}

\begin{proof}
Use \Cref{cor:n9v90-summable-lift,
cor:n9v90-characteristic-sum,cor:n9v90-exact-intertwining,
prop:n9v90-intertwining,prop:n9v90-contraction-tolerance}.
\end{proof}

\begin{theorem}[Ninth-series V90 result]
\label{thm:n9v90-result}
The compulsory audit transfers the exact lift-defect mechanism to the
SM regulator comparison.  If $J-J_0=M$, every $k$-point reconstruction
row is bounded by $C_kkj_\#^{k-1}\lVert M\rVert$, and summable lift
defects yield summable fixed-order errors.  An energy-selected physical
lift has correction $-A^{-1}K$ and the explicit relative bound
\eqref{eq:n9v90-MH-bound}.  Exact equivariance plus uniqueness gives
exact symmetry intertwining; all non-equivariant defects and reverse
transport cycles remain separate quantitative gates.
\end{theorem}

\begin{proof}
Combine \Cref{lem:n9v90-tensor-telescope,
prop:n9v90-Schwinger-row,prop:n9v90-energy-lift,
prop:n9v90-intertwining,cor:n9v90-exact-intertwining,
prop:n9v90-contraction-tolerance}.
\end{proof}

\begin{assessment}[Progress at ninth-series V90]
\label{ass:n9v90-status}
The reconstruction row in RCC now has a computable physical source: the
vertical energy-lift correction.  Fixed-order and characteristic-
functional comparisons are separated, and reflection equivariance is
an exact uniqueness consequence when the reduced energy data are
equivariant.  The next version will use characteristic functionals to
address the all-order consistency and measure-existence gap left open in
V87--V89.
\end{assessment}

\begin{programme}[Ninth-series V91: characteristic functional]
\label{prog:n9v90-V91}
V91 will continue after the audit.  It will formulate the positive-
definiteness, normalization, equicontinuity, and projective consistency
conditions under which the limiting bosonic characteristic functional
defines a measure on distributions, then attach the finite Grassmann
coefficient functional without treating it as a positive scalar
measure.
\end{programme}

\begin{firewall}[Claim boundary after ninth-series V90]
\label{fire:n9v90-boundary}
V90 does not provide the physical reduced SM energy matrices, certify a
uniform lift gap, prove summability of the reconstruction defect,
construct the limiting characteristic measure, close the graded
functional, establish anomaly cancellation, clustering, spectral
sectors, or the full SM--Einstein continuum.
\end{firewall}

\paragraph{Parent-version record.}
V90 was formed from
\texttt{Renewal\_SM\_Einstein\_Continuum\_Research\_Notes\_V89.tex}.
The V89 parent had $35\,582$ counted source lines, $1\,279\,534$ bytes,
and SHA--256
\[
 \texttt{84BE7FFA9091F9B5DE53C7115363FF8C4F1097FAF7188E3B9BC4605D0E892C6D}.
\]

% =============================================================================
% END APPENDED NINTH-SERIES V90 MATERIAL
% =============================================================================

% =============================================================================
% BEGIN APPENDED NINTH-SERIES V91 MATERIAL
% =============================================================================

\section{Ninth-series V91: characteristic functionals and the graded local state}
\label{sec:v9v91}

V87 proves convergence of each fixed Schwinger order under summable
comparison rows, but a collection of moments need not come from one
measure.  V90 therefore redirects the bosonic construction through the
characteristic functional.  This note proves passage of normalization
and positive definiteness, isolates the equicontinuity needed for
Minlos tightness, recovers moments under explicit integrability, and
attaches the local Grassmann coefficients as a graded functional rather
than as a scalar probability measure.

\subsection{Finite-regulator characteristic functionals}

Let $\mathcal D$ be a real nuclear test-function space containing all
bosonic SM--Einstein species after the declared gauge, constraint, and
coordinate reconstruction.  Let $\mathcal D'$ be its continuous dual.
At regulator $n$, let $\mu_n$ be a probability measure on a finite-
dimensional reconstruction space embedded in $\mathcal D'$, and define
\begin{equation}
 Z_n(f):=\int_{\mathcal D'}e^{i\phi(f)}\,d\mu_n(\phi),
 \qquad f\in\mathcal D.
 \label{eq:n9v91-Zn}
\end{equation}

\begin{lemma}[Regulated characteristic properties]
\label{lem:n9v91-characteristic-properties}
For every $n$,
\begin{align}
 Z_n(0)&=1,
 \label{eq:n9v91-normalization}\\
 Z_n(-f)&=\overline{Z_n(f)},
 \label{eq:n9v91-reality}
\end{align}
and for every $f_1,\ldots,f_N\in\mathcal D$ and
$c_1,\ldots,c_N\in\mathbb C$,
\begin{equation}
 \sum_{r,s=1}^N\overline{c_r}c_sZ_n(f_s-f_r)\ge0.
 \label{eq:n9v91-positive-definite}
\end{equation}
\end{lemma}

\begin{proof}
The first two identities follow directly from
\eqref{eq:n9v91-Zn}.  For the last,
\begin{align*}
 \sum_{r,s}\overline{c_r}c_sZ_n(f_s-f_r)
 &=\int\left|\sum_sc_se^{i\phi(f_s)}\right|^2d\mu_n(\phi)\ge0.
\end{align*}
\end{proof}

\begin{firewall}[The bosonic body measure is an input]
\label{fire:n9v91-bosonic-body}
Equation \eqref{eq:n9v91-Zn} requires a positive bosonic probability
measure.  An oriented ghost determinant or chiral determinant phase
cannot be absorbed into $\mu_n$ unless the V82--V86 measure and graded
reflection analysis has produced a positive scalar body and separate
integrable graded coefficients.
\end{firewall}

\subsection{Passage to a positive-definite limit}

Assume $Z_n(f)$ is Cauchy for every $f$, as supplied conditionally by
the V87 hybrid comparison and V90 lift row, and set
\begin{equation}
 Z(f):=\lim_{n\to\infty}Z_n(f).
 \label{eq:n9v91-Z-limit}
\end{equation}

\begin{proposition}[Closed positive-definite limit]
\label{prop:n9v91-positive-limit}
The limit satisfies $Z(0)=1$, $Z(-f)=\overline{Z(f)}$, and
\begin{equation}
 \sum_{r,s=1}^N\overline{c_r}c_sZ(f_s-f_r)\ge0
 \label{eq:n9v91-limit-positive}
\end{equation}
for every finite test family.
\end{proposition}

\begin{proof}
The identities pass through scalar limits.  For fixed finite test and
coefficient families, the left side of
\eqref{eq:n9v91-positive-definite} is a finite sum of convergent complex
numbers and is nonnegative real for every $n$.  Its limit remains in
the closed half-line $[0,\infty)$.
\end{proof}

Pointwise convergence does not by itself make $Z$ continuous at zero.
The missing input is uniform equicontinuity.

\begin{theorem}[Equicontinuity transfer]
\label{thm:n9v91-equicontinuity}
Suppose that for every $\epsilon>0$ there is a neighbourhood
$U_\epsilon$ of zero in $\mathcal D$ such that
\begin{equation}
 \sup_n\sup_{f\in U_\epsilon}|Z_n(f)-1|<\epsilon.
 \label{eq:n9v91-uniform-equicontinuity}
\end{equation}
Then $Z$ is continuous at zero.  Positive definiteness then implies
continuity on all of $\mathcal D$.
\end{theorem}

\begin{proof}
Take $n\to\infty$ in
\eqref{eq:n9v91-uniform-equicontinuity} to get
$|Z(f)-1|\le\epsilon$ on $U_\epsilon$.  For a normalized positive-
definite function, the $3\times3$ positivity matrix for
$0,f,g$ gives the standard inequality
\begin{equation}
 |Z(f)-Z(g)|^2\le2\bigl(1-\operatorname{Re}Z(f-g)\bigr).
 \label{eq:n9v91-PD-continuity}
\end{equation}
Thus continuity at zero implies continuity at every $g$.
\end{proof}

\begin{proposition}[Quadratic equicontinuity criterion]
\label{prop:n9v91-quadratic-equicontinuity}
If there is a continuous seminorm $p$ on $\mathcal D$ and $C<\infty$
such that
\begin{equation}
 1-\operatorname{Re}Z_n(f)\le Cp(f)^2
 \quad(n\ge n_0,\ f\in\mathcal D),
 \label{eq:n9v91-quadratic-bound}
\end{equation}
then \eqref{eq:n9v91-uniform-equicontinuity} holds for the tail.
Finite initial regulators can be added by intersecting finitely many
continuity neighbourhoods.
\end{proposition}

\begin{proof}
The elementary inequality for characteristic functions
\begin{equation}
 |Z_n(f)-1|^2\le2(1-\operatorname{Re}Z_n(f))
\end{equation}
is \eqref{eq:n9v91-PD-continuity} with $g=0$.  Thus
$|Z_n(f)-1|\le\sqrt{2C}\,p(f)$ uniformly on the tail.
Choose a sufficiently small $p$-ball and intersect with the finitely
many initial continuity neighbourhoods.
\end{proof}

\begin{firewall}[Pointwise characteristic convergence is not tightness]
\label{fire:n9v91-pointwise-not-tight}
Without \eqref{eq:n9v91-uniform-equicontinuity}, probability mass can
escape into increasingly singular field directions while every tested
subsequence appears pointwise bounded.  Continuity at zero is the
nuclear-space tightness gate in the Minlos construction.
\end{firewall}

\subsection{Finite-dimensional restrictions and projective consistency}

For a finite-dimensional subspace $E\subset\mathcal D$, let
$Z_E=Z|_E$.

\begin{theorem}[Consistent finite-dimensional marginals]
\label{thm:n9v91-finite-marginals}
If $Z$ is normalized, continuous, and positive definite, then for every
finite-dimensional $E$ there is a unique probability measure $\mu_E$
on $E'$ satisfying
\begin{equation}
 Z_E(f)=\int_{E'}e^{i\ell(f)}\,d\mu_E(\ell).
 \label{eq:n9v91-finite-Bochner}
\end{equation}
If $E\subset F$, the pushforward of $\mu_F$ under restriction
$F'\to E'$ equals $\mu_E$.
\end{theorem}

\begin{proof}
The finite-dimensional Bochner theorem applied after choosing a basis
of $E$ gives existence and uniqueness.  If $E\subset F$, the
pushforward of $\mu_F$ has characteristic function $Z_F|_E=Z_E$.
Uniqueness in the Bochner theorem makes it $\mu_E$.
\end{proof}

\begin{theorem}[Minlos measure]
\label{thm:n9v91-Minlos}
Let $\mathcal D$ be nuclear.  If $Z:\mathcal D\to\mathbb C$ is
normalized, continuous at zero, and positive definite, then there is a
unique Radon probability measure $\mu$ on $\mathcal D'$ with its
cylindrical Borel sigma algebra such that
\begin{equation}
 Z(f)=\int_{\mathcal D'}e^{i\phi(f)}\,d\mu(\phi).
 \label{eq:n9v91-Minlos}
\end{equation}
Its every finite-dimensional marginal is the measure of
\Cref{thm:n9v91-finite-marginals}.
\end{theorem}

\begin{proof}
The consistent Bochner measures from
\Cref{thm:n9v91-finite-marginals} define a cylindrical probability
measure.  The Minlos theorem states that continuity at zero on a nuclear
space makes this cylindrical measure tight and countably additive on
$\mathcal D'$, giving a Radon measure.  Equality
\eqref{eq:n9v91-Minlos} follows on every one-dimensional projection.
Characteristic-function uniqueness on all finite-dimensional
projections proves uniqueness of $\mu$.
\end{proof}

\begin{corollary}[Single-tower bosonic measure construction]
\label{cor:n9v91-bosonic-measure}
If the V87--V90 comparison makes $Z_n$ pointwise Cauchy and the
equicontinuity gate \eqref{eq:n9v91-uniform-equicontinuity} holds, then
the limiting $Z$ defines the unique bosonic distribution-valued measure
\eqref{eq:n9v91-Minlos}.
\end{corollary}

\begin{proof}
Use \Cref{prop:n9v91-positive-limit,
thm:n9v91-equicontinuity,thm:n9v91-Minlos}.
\end{proof}

\begin{firewall}[Nuclearity and support]
\label{fire:n9v91-nuclearity}
The topology on $\mathcal D$ determines both nuclearity and the support
space $\mathcal D'$.  A bound in an unspecified norm does not select a
continuum distribution space.  Metric, gauge, Higgs, and matter species
may require a product of distinct test-function topologies, whose
nuclearity and covariance must be checked.
\end{firewall}

\subsection{Recovery of Schwinger moments}

Let $f_1,\ldots,f_k\in\mathcal D$ and suppose
\begin{equation}
 \int\prod_{j=1}^k(1+|\phi(f_j)|)\,d\mu(\phi)<\infty.
 \label{eq:n9v91-moment-integrability}
\end{equation}

\begin{proposition}[Moment derivatives]
\label{prop:n9v91-moment-derivatives}
Under \eqref{eq:n9v91-moment-integrability},
\begin{equation}
 D^kZ(0)[f_1,\ldots,f_k]
 =i^k\int\prod_{j=1}^k\phi(f_j)\,d\mu(\phi).
 \label{eq:n9v91-moment-derivatives}
\end{equation}
\end{proposition}

\begin{proof}
Differentiate
$\exp(i\sum_jt_j\phi(f_j))$ once in each $t_j$.  The absolute value of
the resulting integrand is bounded by
$\prod_j|\phi(f_j)|$, integrable by
\eqref{eq:n9v91-moment-integrability}.  Dominated convergence permits
all derivatives to pass through the integral at $t=0$.
\end{proof}

\begin{corollary}[Agreement with fixed-order limits]
\label{cor:n9v91-moment-agreement}
If the regulated $k$-point functions converge as in V87 and their
products are uniformly integrable under a common coupling, then their
limit equals the derivative in
\eqref{eq:n9v91-moment-derivatives}.
\end{corollary}

\begin{proof}
Uniform integrability permits passage of the regulated moments to the
limiting measure, while \Cref{prop:n9v91-moment-derivatives} identifies
the latter with the derivative of $Z$.
\end{proof}

\begin{firewall}[Moments need not determine the measure]
\label{fire:n9v91-moment-indeterminate}
Even all finite moments can be indeterminate without growth conditions.
The characteristic functional determines $\mu$ uniquely; the moment
sequence is a derived object only where differentiability and
integrability are proved.
\end{firewall}

\subsection{Support on constraint equations}

Let $C_f(\phi)$ be a measurable smeared constraint residual for test
label $f$.

\begin{proposition}[Quadratic support criterion]
\label{prop:n9v91-support}
If
\begin{equation}
 \int|C_f(\phi)|^2\,d\mu(\phi)=0,
 \label{eq:n9v91-C-square-zero}
\end{equation}
then $C_f(\phi)=0$ for $\mu$-almost every $\phi$.  If this holds on a
countable dense test set and $f\mapsto C_f(\phi)$ is almost surely
continuous, then the constraint vanishes distributionally for
$\mu$-almost every field.
\end{proposition}

\begin{proof}
A nonnegative integrable function with zero integral vanishes almost
surely.  Intersect the countably many full-measure sets.  On their
intersection, continuity extends vanishing from the dense set to every
test function.
\end{proof}

\begin{firewall}[Ward expectation versus support]
\label{fire:n9v91-Ward-support}
A signed Ward identity $\int C_f\,d\mu=0$ can result from cancellation
between positive and negative values and does not imply
\eqref{eq:n9v91-C-square-zero}.  Measure support on the constraint
surface requires a nonnegative quadratic criterion or the stronger
all-insertion operator identity of V88.
\end{firewall}

\subsection{Regulator independence of the bosonic measure}

\begin{proposition}[Characteristic universality]
\label{prop:n9v91-characteristic-universality}
Let two regulator towers satisfy the V87 cross-comparison criterion for
the bounded Weyl observables $e^{i\phi(f)}$.  Then their limiting
characteristic functionals agree on all $f\in\mathcal D$.  If both
satisfy the Minlos hypotheses on the same nuclear topology, their
limiting measures are equal.
\end{proposition}

\begin{proof}
The cross-comparison theorem gives equality of the two limiting Weyl
expectations, hence of the characteristic functionals.  Uniqueness in
\Cref{thm:n9v91-Minlos} gives equality of measures.
\end{proof}

\subsection{The finite local Grassmann functional}

Let $V_{\mathcal O}$ be the finite-dimensional fermion/ghost test space
for a bounded region $\mathcal O$, and set
\begin{equation}
 \mathfrak A_{\mathcal O}
 :=\mathcal B_{\mathcal O}\widehat\otimes
 \Lambda(V_{\mathcal O}),
 \label{eq:n9v91-local-graded-algebra}
\end{equation}
where $\mathcal B_{\mathcal O}$ is a bosonic cylindrical test algebra on
$\mathcal D'$.

Choose an exterior basis $(\theta_I)$.  Every element is uniquely
\begin{equation}
 F(\phi,\theta)=\sum_IF_I(\phi)\theta_I.
 \label{eq:n9v91-F-expansion}
\end{equation}

\begin{theorem}[Graded coefficient functional]
\label{thm:n9v91-graded-functional}
Suppose for every $I$ there is a coefficient
$a_I\in L^1(\mu)$, with only finitely many indices in each local
algebra.  Then
\begin{equation}
 \Omega_{\mathcal O}(F)
 :=\sum_I\int_{\mathcal D'}F_I(\phi)a_I(\phi)
 \,d\mu(\phi)
 \label{eq:n9v91-Omega-local}
\end{equation}
defines a continuous linear functional on bounded coefficient vectors,
with
\begin{equation}
 |\Omega_{\mathcal O}(F)|
 \le\sum_I\lVert F_I\rVert_\infty\lVert a_I\rVert_{L^1(\mu)}.
 \label{eq:n9v91-Omega-bound}
\end{equation}
If these functionals agree under inclusions
$\mathfrak A_{\mathcal O_1}\hookrightarrow
\mathfrak A_{\mathcal O_2}$, they define a unique functional on the
inductive union of local graded algebras.
\end{theorem}

\begin{proof}
Finite summation and H\"older's $L^\infty$--$L^1$ inequality give
\eqref{eq:n9v91-Omega-bound}.  Compatibility makes the value independent
of the chosen containing region, which defines the inductive-limit
functional uniquely.
\end{proof}

The coefficient functions $a_I$ contain the finite Grassmann Schur,
ghost, physical chiral, boundary, and interaction-channel data of
V82--V86.

\begin{firewall}[No Grassmann probability measure]
\label{fire:n9v91-no-Grassmann-measure}
The exterior variables in \eqref{eq:n9v91-F-expansion} are nilpotent
algebra generators.  The coefficients $a_I$ define a graded linear
functional, not a positive scalar probability distribution on
anticommuting numbers.  Graded OS positivity remains the boundary-
transfer condition of V83--V86.
\end{firewall}

\begin{proposition}[Passage of the graded OS form]
\label{prop:n9v91-graded-OS-limit}
Suppose each local $\Omega_{n,\mathcal O}$ is graded reflection positive,
its coefficient functionals converge under the V87 comparison, and the
reflection/order maps are exactly intertwined.  Then the limiting
$\Omega_{\mathcal O}$ satisfies
\begin{equation}
 \Omega_{\mathcal O}((\Theta F)F)\ge0
 \label{eq:n9v91-graded-OS}
\end{equation}
for every positive-time local $F$.
\end{proposition}

\begin{proof}
The regulated left sides are nonnegative real numbers and converge by
coefficient-functional convergence.  The nonnegative half-line is
closed.
\end{proof}

\begin{firewall}[Local coefficient consistency is substantive]
\label{fire:n9v91-coefficient-consistency}
Integrating out a fermion mode can change all lower-degree coefficients
through a Schur determinant.  Compatibility under enlarging
$\mathcal O$ must be proved from the exact graded channel/partial-trace
construction; copying finite-region coefficients does not establish it.
\end{firewall}

\subsection{Characteristic and graded state card}

\begin{definition}[CGSC]
\label{def:n9v91-CGSC}
The characteristic and graded state card requires:
\begin{enumerate}
\item a positive bosonic body measure at every regulator;
\item pointwise Cauchy convergence of $Z_n$ from the complete V87
      ledger;
\item uniform equicontinuity at zero on a declared nuclear test space;
\item the limiting positive-definiteness
      \eqref{eq:n9v91-limit-positive};
\item uniform integrability for every recovered moment;
\item cross-tower characteristic comparison for universality;
\item quadratic or all-insertion constraint support identities;
\item local $L^1(\mu)$ Grassmann coefficient families;
\item compatibility of those coefficients under local inclusions; and
\item exact graded reflection positivity and Ward/anomaly identities for
      the inductive local functional.
\end{enumerate}
\end{definition}

\begin{theorem}[CGSC constructs the local continuum state]
\label{thm:n9v91-CGSC}
Under CGSC, the limiting bosonic characteristic functional defines a
unique regulator-independent Radon probability measure on
$\mathcal D'$.  Its integrable moments agree with the V87 Schwinger
limits.  The compatible Grassmann coefficients define a unique local
graded functional over that bosonic body, and its finite-cutoff graded
OS inequalities pass to the limit.
\end{theorem}

\begin{proof}
Use \Cref{prop:n9v91-positive-limit,
thm:n9v91-equicontinuity,thm:n9v91-Minlos,
prop:n9v91-moment-derivatives,
prop:n9v91-characteristic-universality,
thm:n9v91-graded-functional,prop:n9v91-graded-OS-limit}.
\end{proof}

\begin{theorem}[Ninth-series V91 result]
\label{thm:n9v91-result}
Pointwise limits of the regulated characteristic functionals preserve
normalization and positive definiteness.  Uniform equicontinuity at zero
makes the limit continuous, and the nuclear Minlos theorem then gives a
unique bosonic measure on distributions.  Moments are recovered only
under explicit integrability.  Fermion and ghost variables remain a
compatible local Grassmann coefficient functional over the bosonic
body; graded OS positivity passes as a closed-cone limit and is not
replaced by a fictitious Grassmann probability measure.
\end{theorem}

\begin{proof}
Combine \Cref{prop:n9v91-positive-limit,
thm:n9v91-equicontinuity,thm:n9v91-Minlos,
prop:n9v91-moment-derivatives,thm:n9v91-graded-functional,
prop:n9v91-graded-OS-limit}.
\end{proof}

\begin{assessment}[Progress at ninth-series V91]
\label{ass:n9v91-status}
The all-order consistency gap left by fixed-order Schwinger convergence
now has an exact target: a continuous positive-definite characteristic
functional on a specified nuclear test space, plus a compatible local
graded coefficient system.  The next bottleneck is quantitative
equicontinuity.  It must be derived from uniform smeared-field tails in
the full calibrated SM--Einstein measure, including bad-field and
constraint-branch contributions.
\end{assessment}

\begin{programme}[Ninth-series V92: equicontinuity from tails]
\label{prog:n9v91-V92}
V92 will derive characteristic equicontinuity from uniform quadratic
and exponential smeared-field bounds, split good and bad fields, and
state a cofinal tail card which is stable under the V87 hybrid
comparison and the V90 physical lift.
\end{programme}

\begin{firewall}[Claim boundary after ninth-series V91]
\label{fire:n9v91-boundary}
V91 does not prove a positive bosonic body for the complete chiral
theory, equicontinuity, Minlos tightness for the physical regulator,
uniform moment integrability, constraint support, consistency of the
graded coefficients, anomaly cancellation, clustering, spectral
sectors, or the full SM--Einstein continuum.
\end{firewall}

\paragraph{Parent-version record.}
V91 was formed from
\texttt{Renewal\_SM\_Einstein\_Continuum\_Research\_Notes\_V90.tex}.
The V90 parent had $36\,021$ counted source lines, $1\,294\,433$ bytes,
and SHA--256
\[
 \texttt{FCC17136DB6D178862A1F6B7636032B4561DE94F2BF322F6C62B9225425F405C}.
\]

% =============================================================================
% END APPENDED NINTH-SERIES V91 MATERIAL
% =============================================================================

% =============================================================================
% BEGIN APPENDED NINTH-SERIES V92 MATERIAL
% =============================================================================

\section{Uniform random-seminorm tails and characteristic equicontinuity}
\label{sec:v9v92}

\subsection{Purpose and the point left open in V91}

V91 reduced the bosonic measure-reconstruction problem to continuity at the
origin of the limiting characteristic functional on a nuclear test space.
The purpose of the present note is to supply a regulator-level criterion for
that continuity which survives the three operations already forced by the
programme: removal of bad-field regions, replacement of a reference measure
by a hybrid measure, and passage through the physical reconstruction lift.
The central object is a random seminorm dominating all smeared fields.  Its
tail must be uniform in the complete cofinal regulator, not merely at each
fixed cutoff.

Throughout this note, $E$ is a real locally convex test space, $p$ is a
continuous seminorm on $E$, $I$ is the directed regulator set, and
$(\Omega_i,\mathcal F_i,\mu_i)$ is a probability space for each $i\in I$.
The real random linear functional at regulator $i$ is denoted by
$\Phi_i:E\to L^0(\mu_i)$ and
\[
 Z_i(f):=\int_{\Omega_i}e^{\mathrm i\Phi_i(f)}\,d\mu_i.
\]
No common sample space is assumed.

\begin{lemma}[The elementary cosine estimate]
\label{lem:n9v92-cosine}
For every real random variable $Y$,
\[
 0\leq 1-\Re\mathbb E e^{\mathrm iY}
 \leq \mathbb E\min\!\left\{2,\frac{Y^2}{2}\right\}.
\]
Consequently, if $|\Phi_i(f)|\leq X_i p(f)$ almost surely for a nonnegative
random variable $X_i$, then for every $R>0$,
\begin{equation}
 1-\Re Z_i(f)
 \leq \frac{R^2p(f)^2}{2}+2\mu_i\{X_i>R\}.
 \label{eq:n9v92-truncated-cosine}
\end{equation}
\end{lemma}

\begin{proof}
The pointwise inequalities $0\leq1-\cos y\leq\min\{2,y^2/2\}$ give the
first assertion after integration.  On $\{X_i\leq R\}$ use
$|\Phi_i(f)|\leq Rp(f)$; on the complement use $1-\cos y\leq2$.
\end{proof}

\begin{definition}[Uniform random-seminorm tail card]
\label{def:n9v92-tail-card}
The family $\{\Phi_i\}_{i\in I}$ satisfies the $p$-tail card if there are
nonnegative random variables $X_i$ and a decreasing function
$a:[0,\infty)\to[0,1]$ such that
\begin{equation}
 |\Phi_i(f)|\leq X_i p(f)\quad\mu_i\text{-a.s. for every }f\in E,
 \qquad
 \sup_{i\in I}\mu_i\{X_i>R\}\leq a(R),
 \qquad
 a(R)\longrightarrow0.
 \label{eq:n9v92-tail-card}
\end{equation}
The domination is required on one full-measure event on which it holds for
all $f$ in a fixed countable dense subspace, followed by continuous
extension.  This convention avoids an uncountable intersection of null
sets.
\end{definition}

\begin{theorem}[Uniform tails imply characteristic equicontinuity]
\label{thm:n9v92-tail-equicontinuity}
If the $p$-tail card holds, then the family $\{Z_i\}_{i\in I}$ is
equicontinuous at the origin of $E$.  More precisely,
\begin{equation}
 \sup_{i\in I}\bigl(1-\Re Z_i(f)\bigr)
 \leq \inf_{R>0}\left\{\frac{R^2p(f)^2}{2}+2a(R)\right\}.
 \label{eq:n9v92-uniform-modulus}
\end{equation}
Moreover,
\begin{equation}
 \sup_{i\in I}|Z_i(f)-1|^2
 \leq 2\sup_{i\in I}\bigl(1-\Re Z_i(f)\bigr),
 \label{eq:n9v92-modulus-square}
\end{equation}
so the usual complex-valued equicontinuity follows.
\end{theorem}

\begin{proof}
Taking the supremum in \eqref{eq:n9v92-truncated-cosine} and then the
infimum in $R$ proves \eqref{eq:n9v92-uniform-modulus}.  Given
$\varepsilon>0$, first choose $R$ with $2a(R)<\varepsilon/2$ and then a
$p$-neighbourhood on which $R^2p(f)^2/2<\varepsilon/2$.  For the second
claim, Cauchy--Schwarz and $|e^{\mathrm i y}-1|^2=2(1-\cos y)$ give
\[
 |Z_i(f)-1|^2
 \leq\mathbb E_i|e^{\mathrm i\Phi_i(f)}-1|^2
 =2(1-\Re Z_i(f)).
\]
\end{proof}

\begin{remark}[Why a cutoff-independent bad probability is insufficient]
\label{rem:n9v92-fixed-bad-event}
An estimate
\[
 1-\Re Z_i(f)\leq \tfrac12 A^2p(f)^2+2\delta
\]
with fixed $\delta>0$ does not prove continuity at $f=0$: its right-hand
side does not tend to zero.  The tail card repairs this defect because the
threshold $R$ is chosen after the target accuracy and before $f$ is sent to
zero.  Thus ``good field with high probability'' must mean a uniformly
tight family of field seminorms, rather than one predetermined good set.
\end{remark}

\subsection{Quadratic and exponential sufficient conditions}

\begin{proposition}[Uniform second moments]
\label{prop:n9v92-second-moment}
Suppose
\begin{equation}
 \sup_{i\in I}\mathbb E_i|\Phi_i(f)|^2\leq C p(f)^2
 \qquad(f\in E).
 \label{eq:n9v92-quadratic-form-bound}
\end{equation}
Then
\begin{equation}
 1-\Re Z_i(f)\leq \frac C2p(f)^2,
 \qquad
 |Z_i(f)-1|\leq \sqrt C\,p(f),
 \label{eq:n9v92-quadratic-modulus}
\end{equation}
uniformly in $i$.  In particular this is precisely the quadratic
equicontinuity input used in V91.
\end{proposition}

\begin{proof}
Apply Lemma~\ref{lem:n9v92-cosine} without truncation and then
\eqref{eq:n9v92-modulus-square}.
\end{proof}

\begin{proposition}[Exponential random-seminorm card]
\label{prop:n9v92-exponential-card}
Assume $|\Phi_i(f)|\leq X_i p(f)$ and, for some $q,\alpha>0$ and
$M<\infty$,
\begin{equation}
 \sup_{i\in I}\mathbb E_i e^{\alpha X_i^q}\leq M.
 \label{eq:n9v92-exponential-seminorm}
\end{equation}
Then the tail card holds with
$a(R)=\min\{1,Me^{-\alpha R^q}\}$.  For every $k>0$,
\begin{equation}
 \sup_i\mathbb E_i X_i^k
 \leq M\alpha^{-k/q}\Gamma\!\left(1+\frac{k}{q}\right).
 \label{eq:n9v92-all-moments}
\end{equation}
In particular, V91's quadratic condition holds with
$C=M\alpha^{-2/q}\Gamma(1+2/q)$.
\end{proposition}

\begin{proof}
Markov's inequality gives
$\mu_i\{X_i>R\}\leq Me^{-\alpha R^q}$.  The layer-cake formula and this
bound give
\[
 \mathbb E_iX_i^k
 =k\int_0^\infty r^{k-1}\mu_i\{X_i>r\}\,dr
 \leq Mk\int_0^\infty r^{k-1}e^{-\alpha r^q}\,dr,
\]
and the change of variable $u=\alpha r^q$ evaluates the last expression as
the right-hand side of \eqref{eq:n9v92-all-moments}.  Finally
$|\Phi_i(f)|^2\leq X_i^2p(f)^2$.
\end{proof}

\begin{proposition}[Sub-Gaussian smeared fields]
\label{prop:n9v92-subgaussian}
Suppose every $\Phi_i(f)$ is centred and
\begin{equation}
 \mathbb E_i e^{\lambda\Phi_i(f)}
 \leq \exp\!\left(\frac{C\lambda^2p(f)^2}{2}\right)
 \qquad(\lambda\in\mathbb R, f\in E, i\in I).
 \label{eq:n9v92-subgaussian-mgf}
\end{equation}
Then \eqref{eq:n9v92-quadratic-form-bound} holds with the same $C$.
\end{proposition}

\begin{proof}
For fixed $i,f$, apply \eqref{eq:n9v92-subgaussian-mgf} with $\lambda$ and
$-\lambda$, add, and use
$\cosh x\geq1+x^2/2$.  This gives
\[
 1+\frac{\lambda^2}{2}\mathbb E_i\Phi_i(f)^2
 \leq \mathbb E_i\cosh(\lambda\Phi_i(f))
 \leq e^{C\lambda^2p(f)^2/2}.
\]
Divide by $\lambda^2/2$ and let $\lambda\to0$.
\end{proof}

\begin{remark}[Centres and one-point functions]
\label{rem:n9v92-centres}
If $\Phi_i(f)=m_i(f)+\widetilde\Phi_i(f)$, with the centred part satisfying
Proposition~\ref{prop:n9v92-subgaussian} and
$|m_i(f)|\leq C_1p(f)$ uniformly, then
\[
 |\Phi_i(f)|^2\leq2|\widetilde\Phi_i(f)|^2+2C_1^2p(f)^2
\]
gives the required quadratic bound.  Thus a nonzero Higgs expectation or
other one-point function is harmless only after its continuity is bounded
uniformly in the regulators.
\end{remark}

\subsection{Stability under the hybrid comparison}

The V87 comparison passes among several regulated measures.  The following
two lemmas identify sufficient hypotheses for transporting the tail card;
neither hypothesis is automatic from partition-function convergence.

\begin{lemma}[Bounded-density transport]
\label{lem:n9v92-bounded-density}
Let $\nu_i=r_i\mu_i$ be probability measures and suppose
$0\leq r_i\leq K$ almost surely, uniformly in $i$.  If $X_i$ has
$\mu_i$-tail $a$, then
\begin{equation}
 \nu_i\{X_i>R\}\leq K a(R).
 \label{eq:n9v92-bounded-density-tail}
\end{equation}
Thus the tail card is preserved.
\end{lemma}

\begin{proof}
Integrate $r_i\mathbf1_{\{X_i>R\}}$ and use the essential bound.
\end{proof}

\begin{lemma}[$L^s$ density transport]
\label{lem:n9v92-lp-density}
Let $1<s\leq\infty$, let $s'=s/(s-1)$, and suppose
\begin{equation}
 \nu_i=r_i\mu_i,qquad \sup_i\|r_i\|_{L^s(\mu_i)}\leq K_s.
 \label{eq:n9v92-lp-density-hyp}
\end{equation}
Then
\begin{equation}
 \nu_i\{X_i>R\}
 \leq K_s a(R)^{1/s'}.
 \label{eq:n9v92-lp-density-tail}
\end{equation}
For $s=\infty$ this agrees with
Lemma~\ref{lem:n9v92-bounded-density}.
\end{lemma}

\begin{proof}
Apply H\"older's inequality to
$\int r_i\mathbf1_{\{X_i>R\}}d\mu_i$.
\end{proof}

\begin{proposition}[Coupling transport]
\label{prop:n9v92-coupling-transport}
Suppose $\Phi_i$ and $\widehat\Phi_i$ are defined on a coupling and
\begin{equation}
 |\widehat\Phi_i(f)|\leq |\Phi_i(f)|+Y_i p(f).
 \label{eq:n9v92-coupling-domination}
\end{equation}
If $|\Phi_i(f)|\leq X_i p(f)$ and both $X_i$ and $Y_i$ are uniformly tight,
then $\widehat\Phi_i$ satisfies the tail card.  Quantitatively,
\begin{equation}
 \mathbb P\{X_i+Y_i>R\}
 \leq \mathbb P\{X_i>R/2\}+\mathbb P\{Y_i>R/2\}.
 \label{eq:n9v92-sum-tail}
\end{equation}
\end{proposition}

\begin{proof}
The domination holds with random seminorm $X_i+Y_i$, and the displayed
union bound tends to zero uniformly.
\end{proof}

\begin{firewall}[Determinant weights]
\label{fire:n9v92-determinant-weight}
The normalized determinant-density comparison of V87 does not by itself
give \eqref{eq:n9v92-lp-density-hyp}.  A determinant can have small mean
error while developing large peaks on small sets.  The continuum proof must
therefore provide either a uniform $L^s$, $s>1$, bound for each normalized
bosonic density in the hybrid chain, or a direct coupled tail estimate.  The
chiral determinant-line phase is still outside this positive-density
argument.
\end{firewall}

\subsection{Stability under physical lifts and reconstructed branches}

Let $E_{\rm phys}$ and $E_{\rm ref}$ be locally convex test spaces.  At
regulator $i$, let $J_i:E_{\rm phys}\to E_{\rm ref}$ be the physical lift
of V90 and define
\[
 \Phi_i^{\rm phys}(f):=\Phi_i^{\rm ref}(J_if).
\]

\begin{proposition}[Uniformly bounded physical lift]
\label{prop:n9v92-lift-tail}
Suppose
\begin{equation}
 |\Phi_i^{\rm ref}(g)|\leq X_i p_{\rm ref}(g),
 \qquad
 p_{\rm ref}(J_if)\leq C_Jp_{\rm phys}(f)
 \label{eq:n9v92-lift-bounds}
\end{equation}
uniformly in $i$.  If $X_i$ has uniform tail $a$, then the physical fields
satisfy the $p_{\rm phys}$-tail card with random seminorm $C_JX_i$ and tail
$a(R/C_J)$.  The same statement holds for the quadratic and exponential
cards with the corresponding powers of $C_J$.
\end{proposition}

\begin{proof}
Substitution in \eqref{eq:n9v92-lift-bounds} gives
$|\Phi_i^{\rm phys}(f)|\leq C_JX_ip_{\rm phys}(f)$.  The remaining claims
follow by rescaling.
\end{proof}

\begin{corollary}[V90 relative lift debit]
\label{cor:n9v92-relative-lift}
In the notation of V90, suppose
$J_i=J_{0,i}-A_i^{-1}K_i$ and there is a common pair of seminorms such that
\begin{equation}
 p_{\rm ref}(J_{0,i}f)\leq C_0p_{\rm phys}(f),
 \qquad
 p_{\rm ref}(A_i^{-1}K_if)\leq\eta_i p_{\rm phys}(f),
 \qquad
 \sup_i\eta_i\leq\eta_\#<\infty.
 \label{eq:n9v92-v90-lift-bound}
\end{equation}
Then Proposition~\ref{prop:n9v92-lift-tail} applies with
$C_J=C_0+\eta_\#$.
\end{corollary}

\begin{proof}
Use the triangle inequality for $p_{\rm ref}$.
\end{proof}

\begin{remark}[Closing singular values]
\label{rem:n9v92-closing-lift}
If $\|A_i^{-1}K_i\|$ diverges, reference-field tightness need not imply
physical-field tightness.  This is the probabilistic form of the V84/V90
closing-mode warning.  Such modes must be retained as explicit variables,
or the numerator must be proved to vanish at a rate that keeps
\eqref{eq:n9v92-v90-lift-bound} uniform.
\end{remark}

For nonlinear constraint reconstruction, let $y$ denote independent
coordinates, let $u=\Psi_i(y)$ be the reconstructed coordinates, and let
$\mathcal O_i(f;y,\Psi_i(y))$ be the full smeared observable.

\begin{proposition}[Constraint-branch tail transfer]
\label{prop:n9v92-constraint-branch}
Suppose there are random variables $X_i,Y_i\geq0$ and constants $b,c\geq0$
independent of $i$ such that
\begin{align}
 |\mathcal O_i(f;y,0)|&\leq X_i p(f),
 \label{eq:n9v92-independent-tail}\\
 |\mathcal O_i(f;y,\Psi_i(y))-\mathcal O_i(f;y,0)|
 &\leq (bY_i+c)p(f).
 \label{eq:n9v92-reconstruction-tail}
\end{align}
If $X_i$ and $Y_i$ are uniformly tight, then the reconstructed full
observable satisfies the tail card.
\end{proposition}

\begin{proof}
The full observable is dominated by
$[X_i+bY_i+c]p(f)$.  Uniform tightness follows from two applications of the
union bound, or from Proposition~\ref{prop:n9v92-coupling-transport}.
\end{proof}

\begin{remark}[Where the implicit-chart constants enter]
\label{rem:n9v92-implicit-chart}
In a local implicit chart, $Y_i$ may be taken to dominate the independent
field norm, while $b$ contains the uniform inverse-Jacobian and response
constants.  The proposition is useful only after the chart radius, inverse
gap, and leakage estimates are uniform on the chosen cofinal family.  A
pointwise implicit-function theorem at each regulator is insufficient.
\end{remark}

\subsection{Several species and the cofinal card}

The full SM--Einstein field contains finitely many bosonic species at a
fixed geometric presentation, as well as derivatives and local coordinate
components.  A single Hilbertian seminorm on the nuclear test space can
absorb finitely many component seminorms.

\begin{lemma}[Finite species assembly]
\label{lem:n9v92-finite-species}
Let $\Phi_i^{(a)}$ for $1\leq a\leq N$ obey
\[
 |\Phi_i^{(a)}(f_a)|\leq X_i^{(a)}p_a(f_a)
\]
and suppose each $X_i^{(a)}$ is uniformly tight.  Define
\[
 p(f)^2:=\sum_{a=1}^Np_a(f_a)^2,
 \qquad X_i:=\left(\sum_{a=1}^N(X_i^{(a)})^2\right)^{1/2}.
\]
Then $|\sum_a\Phi_i^{(a)}(f_a)|\leq X_ip(f)$ and $X_i$ is uniformly tight.
\end{lemma}

\begin{proof}
Cauchy--Schwarz gives the domination.  If $X_i>R$, then some
$X_i^{(a)}>R/\sqrt N$; the union bound proves uniform tightness.
\end{proof}

\begin{definition}[Cofinal characteristic-tail card]
\label{def:n9v92-cofinal-card}
A cofinal regulator family passes the characteristic-tail card if the
following data are recorded with constants independent of the full
regulator tuple:
\begin{enumerate}
\item a nuclear physical test space $E_{\rm phys}$ and one continuous
Hilbertian seminorm $p_{\rm phys}$;
\item a positive normalized bosonic body measure at every regulator;
\item random seminorms $X_i$ satisfying
$|\Phi_i^{\rm phys}(f)|\leq X_ip_{\rm phys}(f)$ and one uniform tail
$a(R)\to0$;
\item for every hybrid density change, either a uniform $L^s$, $s>1$,
bound as in Lemma~\ref{lem:n9v92-lp-density}, or a direct coupling estimate;
\item a uniform physical-lift bound as in
Proposition~\ref{prop:n9v92-lift-tail}, with every closing mode retained;
\item a uniform constraint-reconstruction estimate as in
Proposition~\ref{prop:n9v92-constraint-branch};
\item compatible local Grassmann coefficient bounds, kept algebraically
separate from the positive bosonic body.
\end{enumerate}
\end{definition}

\begin{theorem}[Tail card closes V91's equicontinuity hypothesis]
\label{thm:n9v92-closes-v91}
Assume the cofinal characteristic-tail card and suppose the physical
characteristic functionals converge pointwise on $E_{\rm phys}$:
$Z_i(f)\to Z(f)$.  Then $Z$ is continuous at the origin, normalized,
positive definite, and satisfies $Z(-f)=\overline{Z(f)}$.  Hence, if
$E_{\rm phys}$ is nuclear, the conditional Minlos theorem of V91 produces a
unique Radon probability measure on $E_{\rm phys}'$ with characteristic
functional $Z$.
\end{theorem}

\begin{proof}
Theorem~\ref{thm:n9v92-tail-equicontinuity} gives an $i$-independent
modulus at the origin.  Passing to the pointwise limit preserves this
modulus and therefore continuity.  Normalization, reality, and positive
definiteness pass to pointwise limits by the finite-matrix argument of V91.
Minlos then applies under the stated nuclearity hypothesis.
\end{proof}

\begin{corollary}[Cross-tower stability]
\label{cor:n9v92-cross-tower}
Suppose two cofinal regulator towers are joined by the V87 hybrid chain and
each comparison step satisfies one of the transport hypotheses in
Definition~\ref{def:n9v92-cofinal-card}, with common final tail modulus.
If their physical characteristic functionals have the same pointwise
limit, then they reconstruct the same Radon probability measure on
$E_{\rm phys}'$.
\end{corollary}

\begin{proof}
Both limits are continuous by
Theorem~\ref{thm:n9v92-closes-v91}; uniqueness in Minlos identifies the
measures because their characteristic functionals agree.
\end{proof}

\subsection{Consequences, remaining debit, and next step}

V92 turns an analytic field-tail estimate into the exact continuity input
needed by the V91 reconstruction theorem.  It also locates three failure
modes that a genuine SM--Einstein construction must exclude: concentration
of normalized determinant weights, divergence of the physical lift on
closing modes, and loss of the implicit-chart inverse gap.  These are now
explicit ledger rows rather than hidden uses of dominated convergence.

\begin{assessment}[Status after ninth-series V92]
\label{assess:n9v92}
The characteristic-functional route is conditionally closed from a
uniform cofinal random-seminorm tail card to a Radon bosonic measure on the
distribution space.  The programme has not yet derived that tail card from
the interacting regulated SM--Einstein action.  In particular, a coercive
local estimate must control the conformal gravitational direction, gauge
zero modes, Higgs/Yukawa interactions, normalized determinant densities,
and the reconstructed constraint branch with constants uniform in volume,
mesh, shell, and calibration parameters.
\end{assessment}

\begin{programme}[Ninth-series V93: coercive local tail mechanism]
\label{prog:n9v92-V93}
V93 will formulate a local coercive-block criterion yielding the random
seminorm tail card.  It will distinguish stable positive directions from
the Euclidean conformal and gauge-orbit directions, prove an exponential
tail bound under a block energy inequality, and identify exactly which
SM--Einstein terms can and cannot supply the required coercivity.
\end{programme}

\begin{firewall}[Claim boundary after ninth-series V92]
\label{fire:n9v92-boundary}
V92 does not prove positivity of the complete chiral bosonic body, a
uniform interacting exponential moment, control of the Euclidean
conformal factor, an anomaly-free determinant-line trivialization,
constraint support, graded coefficient consistency, clustering, a
sectorial mass statement, or the full SM--Einstein continuum.
\end{firewall}

\paragraph{Parent-version record.}
V92 was formed from
\texttt{Renewal\_SM\_Einstein\_Continuum\_Research\_Notes\_V91.tex}.
The V91 parent had $36\,535$ validator-counted source lines,
$1\,313\,788$ bytes, and SHA--256
\[
 \texttt{0157730F1902BFCFE596A19BE7E3C0226A133E581BCDD58D27191B43D0BA3991}.
\]

% =============================================================================
% END APPENDED NINTH-SERIES V92 MATERIAL
% =============================================================================

% =============================================================================
% BEGIN APPENDED NINTH-SERIES V93 MATERIAL
% =============================================================================

\section{From local coercive energy to the cofinal field-tail card}
\label{sec:v9v93}

\subsection{The action-to-probability conversion}

V92 isolated a uniform random-seminorm tail as the missing probabilistic
input for Minlos reconstruction.  A lower bound for a bare action is not by
itself such an input, because normalization can cancel the apparent gain.
This note begins with the exact normalization hypothesis needed to convert
coercivity into a uniform exponential moment.

For each regulator $i$, let $\rho_i$ be a reference probability measure,
let $V_i$ be a measurable interaction, and suppose
\[
 0<Z_i:=\int e^{-V_i}\,d\rho_i<\infty,
 \qquad d\mu_i:=Z_i^{-1}e^{-V_i}\,d\rho_i.
\]
Let $\mathcal E_i\geq0$ be a measurable block energy.

\begin{theorem}[Normalized coercivity lemma]
\label{thm:n9v93-normalized-coercivity}
Assume that there are constants $c>0$, $z_*>0$ and real numbers $C_i$ such
that
\begin{equation}
 V_i\geq c\mathcal E_i-C_i
 \quad\rho_i\text{-a.s.},
 \qquad
 Z_i e^{-C_i}\geq z_*.
 \label{eq:n9v93-normalized-coercivity}
\end{equation}
Then, for every $0\leq\alpha\leq c$,
\begin{equation}
 \mathbb E_{\mu_i}e^{\alpha\mathcal E_i}
 \leq z_*^{-1}
 \label{eq:n9v93-energy-exponential}
\end{equation}
uniformly in $i$.
\end{theorem}

\begin{proof}
The action bound gives
\[
 \int e^{\alpha\mathcal E_i-V_i}\,d\rho_i
 \leq e^{C_i}\int e^{-(c-\alpha)\mathcal E_i}\,d\rho_i
 \leq e^{C_i},
\]
because $\rho_i$ is a probability measure and $c-\alpha\geq0$.  Divide by
$Z_i$ and use the second inequality in
\eqref{eq:n9v93-normalized-coercivity}.
\end{proof}

\begin{remark}[The extensive-constant cancellation]
\label{rem:n9v93-extensive-cancellation}
In a volume-$|\Lambda_i|$ system one often has $C_i=C|\Lambda_i|$.
Discarding this term and asking for a volume-independent lower bound on
$Z_i$ is usually false.  The invariant datum is the normalized combination
$Z_ie^{-C_i}$.  Thus a legitimate coercivity card contains both an action
lower bound and a matching free-energy lower bound.
\end{remark}

\begin{corollary}[Field seminorm from block energy]
\label{cor:n9v93-field-from-energy}
Suppose, in addition, that a random seminorm $X_i$ dominating the physical
smeared field satisfies the deterministic reconstruction inequality
\begin{equation}
 |\Phi_i(f)|\leq X_i p(f),
 \qquad
 X_i^q\leq A\mathcal E_i+B
 \label{eq:n9v93-reconstruction-inequality}
\end{equation}
with $q,A>0$ and $B\geq0$ independent of $i$.  Then
\begin{equation}
 \sup_i\mathbb E_{\mu_i}
 \exp\!\left(\frac{c}{A}X_i^q\right)
 \leq e^{cB/A}z_*^{-1}.
 \label{eq:n9v93-field-exponential}
\end{equation}
Consequently the V92 exponential random-seminorm card holds.
\end{corollary}

\begin{proof}
From \eqref{eq:n9v93-reconstruction-inequality},
$cX_i^q/A\leq c\mathcal E_i+cB/A$.  Apply
Theorem~\ref{thm:n9v93-normalized-coercivity} with $\alpha=c$.
\end{proof}

\begin{remark}[Scaling the energy]
\label{rem:n9v93-energy-scaling}
If the deterministic estimate reads $X_i^q\leq A_i\mathcal E_i+B_i$, the
conclusion is useful only when $\sup_iA_i<\infty$ and
$\sup_iB_i<\infty$.  This is where the mesh weights, block overlap, volume
normalization, and Sobolev reconstruction constants enter.  A coercive
finite-dimensional action at every fixed lattice does not imply the
cofinal tail card.
\end{remark}

\subsection{A local block criterion}

Let $\mathcal B_i$ be a finite block cover.  Associate to each block a
nonnegative local energy $e_{i,B}$ and set
\[
 \mathcal E_i:=\sum_{B\in\mathcal B_i}w_{i,B}e_{i,B},
 \qquad w_{i,B}>0.
\]
The next result records a standard bounded-overlap mechanism without
assuming a translation-invariant lattice.

\begin{proposition}[Bounded-overlap reconstruction]
\label{prop:n9v93-bounded-overlap}
Suppose the field is represented by local variables $u_{i,B}$ and there
are local test seminorms $p_{i,B}$ such that
\begin{align}
 |\Phi_i(f)|
 &\leq\sum_{B\in\mathcal B_i}|u_{i,B}|p_{i,B}(f),
 \label{eq:n9v93-local-pairing}\\
 \sum_{B\in\mathcal B_i}w_{i,B}^{-2/q}
 p_{i,B}(f)^2
 &\leq C_p^2p(f)^2,
 \label{eq:n9v93-test-embedding}\\
 \sum_{B\in\mathcal B_i}w_{i,B}^{2/q}|u_{i,B}|^2
 &\leq C_e^2\bigl(1+\mathcal E_i\bigr)^{2/q}
 \label{eq:n9v93-field-energy-embedding}
\end{align}
with constants independent of $i$.  Then
\begin{equation}
 |\Phi_i(f)|\leq C_pC_e(1+\mathcal E_i)^{1/q}p(f).
 \label{eq:n9v93-block-random-seminorm}
\end{equation}
Thus one may take
$X_i=C_pC_e(1+\mathcal E_i)^{1/q}$ in
Corollary~\ref{cor:n9v93-field-from-energy}.
\end{proposition}

\begin{proof}
Insert
$|u_{i,B}|p_{i,B}(f)=
 [w_{i,B}^{1/q}|u_{i,B}|]
 [w_{i,B}^{-1/q}p_{i,B}(f)]$
in \eqref{eq:n9v93-local-pairing} and apply Cauchy--Schwarz.  The two
factors are bounded by \eqref{eq:n9v93-test-embedding} and
\eqref{eq:n9v93-field-energy-embedding}.
\end{proof}

\begin{remark}[Why local estimates are preferable]
\label{rem:n9v93-local-cone}
The V86 full-Fock margin can decay exponentially with volume.  The block
criterion does not ask for one global strict positivity margin.  It asks
for uniform local energy and reconstruction constants plus bounded
overlap.  This is compatible with passing to arbitrary compactly supported
test functions and then using the nuclear topology.
\end{remark}

\begin{proposition}[Good blocks and rare defects]
\label{prop:n9v93-good-bad-blocks}
Suppose $\Phi_i=\Phi_i^{\rm g}+\Phi_i^{\rm d}$, where the good part obeys
\eqref{eq:n9v93-block-random-seminorm} with energy $\mathcal E_i$, and
\begin{equation}
 |\Phi_i^{\rm d}(f)|\leq Y_i p(f),
 \qquad
 \sup_i\mu_i\{Y_i>R\}\leq b(R),
 \qquad b(R)\longrightarrow0.
 \label{eq:n9v93-defect-tail}
\end{equation}
If \eqref{eq:n9v93-energy-exponential} holds, then the complete field
satisfies the V92 tail card.
\end{proposition}

\begin{proof}
The good random seminorm has an exponential tail by Markov's inequality;
the defect seminorm is uniformly tight by hypothesis.  Their sum is
uniformly tight by the V92 coupling-transport proposition.
\end{proof}

\subsection{Stable bosonic directions in the regulated theory}

The preceding statements are abstract.  The point of the next ledger is to
separate actual positive terms from directions on which positivity has not
been obtained.

\begin{proposition}[Quartic absorption for the Higgs amplitude]
\label{prop:n9v93-higgs-absorption}
Let $\lambda>0$ and $m^2\in\mathbb R$.  For every $0<\varepsilon<\lambda$
there is $C_{\varepsilon,m,\lambda}<\infty$ such that
\begin{equation}
 \lambda r^4+m^2r^2
 \geq (\lambda-\varepsilon)r^4-C_{\varepsilon,m,\lambda}
 \qquad(r\geq0).
 \label{eq:n9v93-higgs-quartic}
\end{equation}
One may take
$C_{\varepsilon,m,\lambda}=(m_-^2)^2/(4\varepsilon)$, where
$m_-^2:=\max\{-m^2,0\}$.
\end{proposition}

\begin{proof}
If $m^2\geq0$ the assertion is immediate.  Otherwise Young's inequality
$m_-^2r^2\leq\varepsilon r^4+(m_-^2)^2/(4\varepsilon)$ gives the claim.
\end{proof}

\begin{corollary}[Positive matter energy at fixed geometry]
\label{cor:n9v93-matter-positive}
At a fixed positive Euclidean metric, the standard bosonic terms
\[
 \int\left(\frac14|F|^2+|D H|^2+\lambda|H|^4+m^2|H|^2\right)d\mathrm{vol}
\]
bound the gauge curvature, the covariant Higgs gradient, and the Higgs
quartic norm from below, up to a volume term.  They do not control gauge
orbit directions, flat connections, metric degeneracy, or the conformal
gravitational direction.
\end{corollary}

\begin{proof}
All displayed terms except the possibly negative quadratic Higgs term are
nonnegative.  Apply Proposition~\ref{prop:n9v93-higgs-absorption}
pointwise and integrate.  The list of uncontrolled directions follows
because the energy depends on curvature or covariant derivatives rather
than on a gauge representative, and because the metric has so far been
held fixed.
\end{proof}

\begin{remark}[Gauge orbits]
\label{rem:n9v93-gauge-orbits}
On a finite lattice with compact gauge group, link variables already range
over a compact set, but this compactness does not provide a continuum
Sobolev bound for a gauge potential.  Such a bound requires a gauge slice,
a quotient formulation, or gauge-invariant observables.  The V79
observable/silent split and the V82 coarea Jacobian identify the additional
inverse and determinant estimates required by a slice.
\end{remark}

\subsection{The Euclidean conformal obstruction}

We now make precise why the ordinary Euclidean Einstein--Hilbert action
cannot supply the positive block energy required above.  This is an
upstream obstruction, not a missing estimate.

\begin{theorem}[Conformal-factor unboundedness in four dimensions]
\label{thm:n9v93-conformal-unboundedness}
Let $(M,\bar g)$ be a closed four-dimensional Riemannian manifold and
consider
\begin{equation}
 S_{\rm EH}^{E}(g)
 :=-\frac1{16\pi G}\int_M(R_g-2\Lambda)\,d\mathrm{vol}_g,
 \qquad G>0.
 \label{eq:n9v93-euclidean-EH}
\end{equation}
For $g=e^{2\sigma}\bar g$,
\begin{equation}
 S_{\rm EH}^{E}(e^{2\sigma}\bar g)
 =-\frac1{16\pi G}\int_M e^{2\sigma}
 \bigl(\bar R+6|\bar\nabla\sigma|^2\bigr)d\mathrm{vol}_{\bar g}
 +\frac{\Lambda}{8\pi G}\int_Me^{4\sigma}d\mathrm{vol}_{\bar g}.
 \label{eq:n9v93-conformal-action}
\end{equation}
If $M$ admits a nonconstant smooth function, then
$S_{\rm EH}^{E}$ is unbounded below along a sequence of smooth conformal
metrics with conformal factors bounded in $L^\infty$.
\end{theorem}

\begin{proof}
The conformal scalar-curvature formula is
\[
 R_{e^{2\sigma}\bar g}
 =e^{-2\sigma}
 \bigl(\bar R-6\bar\Delta\sigma-6|\bar\nabla\sigma|^2\bigr).
\]
Multiplying by $d\mathrm{vol}_{e^{2\sigma}\bar g}=e^{4\sigma}
d\mathrm{vol}_{\bar g}$ and integrating the Laplacian term by parts gives
\[
 \int_Me^{2\sigma}(-6\bar\Delta\sigma)d\mathrm{vol}_{\bar g}
 =12\int_Me^{2\sigma}|\bar\nabla\sigma|^2d\mathrm{vol}_{\bar g},
\]
which proves \eqref{eq:n9v93-conformal-action}.

Choose a coordinate ball and a nonzero compactly supported smooth cutoff
$\chi$ inside it.  In one coordinate $x^1$, set
$\sigma_k=a\chi(x)\sin(kx^1)$ with fixed sufficiently small $a>0$.
Then $\|\sigma_k\|_\infty\leq a\|\chi\|_\infty$, so the scalar-curvature
and cosmological terms without derivatives remain uniformly bounded.
On a smaller ball on which $\chi$ is nonzero, the leading derivative is
$ak\chi\cos(kx^1)$; hence
\[
 \int_Me^{2\sigma_k}|\bar\nabla\sigma_k|^2d\mathrm{vol}_{\bar g}
 \longrightarrow\infty
\]
at order $k^2$.  Its coefficient in
\eqref{eq:n9v93-conformal-action} is negative.  Therefore
$S_{\rm EH}^{E}(e^{2\sigma_k}\bar g)\to-\infty$.
\end{proof}

\begin{corollary}[No direct positive-density coercivity from Einstein--Hilbert]
\label{cor:n9v93-no-direct-EH-coercivity}
For the sign convention \eqref{eq:n9v93-euclidean-EH}, no lower bound of
the form
\[
 S_{\rm EH}^{E}(g)\geq c\|\nabla\sigma\|_{L^2}^2-C
\]
with $c>0$ can hold on all smooth conformal metrics, even with a fixed
$L^\infty$ bound on $\sigma$.  Consequently the normalized coercivity
lemma cannot be applied to the unmodified Euclidean Einstein--Hilbert
density in the conformal direction.
\end{corollary}

\begin{proof}
The sequence in Theorem~\ref{thm:n9v93-conformal-unboundedness} has a
right-hand side tending to $+\infty$ and a left-hand side tending to
$-\infty$.
\end{proof}

\begin{firewall}[What matter terms do not repair]
\label{fire:n9v93-matter-conformal}
At bounded matter fields, the usual local SM action does not produce a
positive term growing like the high-frequency conformal kinetic energy in
Theorem~\ref{thm:n9v93-conformal-unboundedness}.  Thus its presence does
not by itself repair the conformal instability.  A proposed repair must be
stated explicitly: for example a justified contour deformation, a
different positive gravitational regulator with a proved removal theorem,
a canonical/constraint construction, or a Lorentzian route followed by a
proved Euclidean reconstruction.  These alternatives are not equivalent
without further work.
\end{firewall}

\subsection{The complete coercive-tail ledger}

\begin{definition}[SM--Einstein coercive-tail card]
\label{def:n9v93-coercive-tail-card}
A cofinal regulated construction passes the coercive-tail card on a local
physical test space if it supplies, uniformly in the complete regulator:
\begin{enumerate}
\item a positive normalized bosonic body $d\mu_i=Z_i^{-1}e^{-V_i}d\rho_i$;
\item a nonnegative local block energy $\mathcal E_i$ and normalized
coercivity constants satisfying
\eqref{eq:n9v93-normalized-coercivity};
\item a mesh- and volume-uniform deterministic reconstruction inequality
\eqref{eq:n9v93-reconstruction-inequality};
\item a treatment of gauge-orbit and flat directions compatible with the
V79 constraint split and V82 coarea factor;
\item a conformal-gravity prescription that replaces the failed direct
Einstein--Hilbert lower bound and is compatible with reflection positivity;
\item uniform $L^s$, $s>1$, control or direct tail transport for the
normalized determinant weights of V82--V87;
\item uniform physical-lift and reconstructed-branch bounds from V92,
with closing modes retained;
\item a separate algebraic estimate for the Grassmann coefficients and a
trivialized anomaly-free determinant line.
\end{enumerate}
\end{definition}

\begin{theorem}[Coercive-tail implication chain]
\label{thm:n9v93-implication-chain}
If a cofinal family passes Definition~\ref{def:n9v93-coercive-tail-card},
then its physical bosonic characteristic functionals are equicontinuous at
the origin.  Every pointwise limit therefore reconstructs a Radon
probability measure on the physical distribution space under the nuclearity
hypothesis of V91.
\end{theorem}

\begin{proof}
The normalized coercivity lemma and the deterministic reconstruction
inequality give the exponential random-seminorm card by
Corollary~\ref{cor:n9v93-field-from-energy}.  The remaining ledger rows
transport that card through hybrid measures, lifts, and constraint
reconstruction by V92.  V92 then gives equicontinuity and invokes the V91
Minlos theorem.
\end{proof}

\begin{assessment}[Status after ninth-series V93]
\label{assess:n9v93}
V93 proves a rigorous action-to-tail implication and an explicit local
block route to the V92 card.  For positive gauge-curvature and Higgs
directions, standard Euclidean terms have the correct sign after quartic
absorption.  The ordinary Euclidean Einstein--Hilbert conformal direction
has the wrong sign and is unbounded below.  Therefore the direct positive
Gibbs-measure route cannot be completed until a gravitational prescription
is chosen and proved compatible with the OS and cofinal comparison ledgers.
\end{assessment}

\begin{programme}[Ninth-series V94: conformal prescriptions and OS compatibility]
\label{prog:n9v93-V94}
V94 will go upstream to compare the admissible gravitational routes at the
level of exact mathematical obligations.  It will analyze conformal
contour rotation, positive higher-derivative stabilization, and
constraint-reduced variables, prove the finite-dimensional contour lemma
that is actually available, and identify which steps fail to preserve a
positive bosonic body or reflection positivity.
\end{programme}

\begin{firewall}[Claim boundary after ninth-series V93]
\label{fire:n9v93-boundary}
V93 does not construct a conformal contour, prove a positive gravitational
measure, remove a stabilizing regulator, establish reflection positivity
for gravity, control determinant densities, trivialize the chiral
determinant line, prove the interacting tail card, or construct the full
SM--Einstein continuum.
\end{firewall}

\paragraph{Parent-version record.}
V93 was formed from
\texttt{Renewal\_SM\_Einstein\_Continuum\_Research\_Notes\_V92.tex}.
The V92 parent had $37\,050$ validator-counted source lines,
$1\,333\,510$ bytes, and SHA--256
\[
 \texttt{9A5B3F8855FB443835B8D101C6FE7AFDA1CCED32AB14DDAAEEC71175482D92C8}.
\]

% =============================================================================
% END APPENDED NINTH-SERIES V93 MATERIAL
% =============================================================================

% =============================================================================
% BEGIN APPENDED NINTH-SERIES V94 MATERIAL
% =============================================================================

\section{Conformal prescriptions tested against measure and OS positivity}
\label{sec:v9v94}

\subsection{Decision problem created by V93}

V93 proved that the real Euclidean Einstein--Hilbert action is unbounded
below along bounded-amplitude, high-frequency conformal factors.  The
purpose of V94 is to test three common responses at the level of exact
finite-dimensional or Gaussian statements: rotate the conformal contour,
add a positive higher-derivative term, or eliminate the unstable variable
by a constraint reduction.  Convergence, probability positivity, and
reflection positivity are distinct requirements; satisfying one does not
imply the others.

\subsection{What conformal contour rotation actually produces}

Consider first one unstable quadratic coordinate with formal weight
$e^{a x^2/2}$, $a>0$.  Its real-axis integral diverges.  On the imaginary
axis $z=\mathrm i y$ the integral converges, but its normalized transform is
not the characteristic function of a real random variable.

\begin{lemma}[Finite-dimensional conformal rotation]
\label{lem:n9v94-contour-rotation}
Let $a>0$ and let $F$ be entire with
$|F(\mathrm i y)|\leq C e^{\beta y^2}$ for some $\beta<a/2$.  Then
\begin{equation}
 \frac{\displaystyle\int_{\mathrm i\mathbb R}
 e^{az^2/2}F(z)\,dz}
 {\displaystyle\int_{\mathrm i\mathbb R}e^{az^2/2}\,dz}
 =\frac{\displaystyle\int_{\mathbb R}e^{-ay^2/2}F(\mathrm i y)\,dy}
 {\displaystyle\int_{\mathbb R}e^{-ay^2/2}\,dy}.
 \label{eq:n9v94-contour-expectation}
\end{equation}
Both sides converge absolutely.  For $F_t(z)=e^{\mathrm i t z}$,
\begin{equation}
 Z_{\rm rot}(t)=e^{t^2/(2a)}.
 \label{eq:n9v94-rotated-transform}
\end{equation}
\end{lemma}

\begin{proof}
Parametrize the oriented contour by $z=\mathrm i y$, so $dz=\mathrm i dy$.
The common factor $\mathrm i$ cancels between numerator and denominator,
and the growth assumption gives absolute convergence.  For $F_t$ the
right-hand side is the Gaussian moment-generating function of $-y$, which
equals \eqref{eq:n9v94-rotated-transform}.
\end{proof}

\begin{proposition}[Failure of real probability positivity]
\label{prop:n9v94-contour-not-characteristic}
The function $Z_{\rm rot}$ in
\eqref{eq:n9v94-rotated-transform} is not the characteristic function of a
probability measure on $\mathbb R$.  More generally, a rotated conformal
coordinate cannot be inserted as a real random distribution merely by
calling the convergent contour integral a probability measure.
\end{proposition}

\begin{proof}
Every probability characteristic function $Z$ satisfies
$|Z(t)|\leq Z(0)=1$, by the triangle inequality.  But
$Z_{\rm rot}(t)>1$ for $t\neq0$.  Equivalently, its second derivative at the
origin has the wrong sign for $Z''(0)=-\mathbb E X^2$.
\end{proof}

\begin{remark}[Scope of the obstruction]
\label{rem:n9v94-contour-scope}
Contour rotation can define analytic continuations of selected holomorphic
observables.  Proposition~\ref{prop:n9v94-contour-not-characteristic}
shows that it does not directly supply the positive bosonic body required
by the V91--V93 Minlos route for a real metric.  To use the contour route,
one must specify a real physical observable algebra on which positivity is
recovered and prove OS positivity there; convergence of the contour alone
does not do this.
\end{remark}

The same calculation applies to a finite vector of unstable modes.

\begin{corollary}[Several rotated directions]
\label{cor:n9v94-multidimensional-rotation}
Let $A$ be a positive definite $m\times m$ matrix and rotate
$z=\mathrm i y$ in the formal weight $e^{z^TAz/2}$.  Its normalized
linear transform is
\begin{equation}
 Z_{\rm rot}(t)=\exp\!\left(\frac12t^TA^{-1}t\right),
 \label{eq:n9v94-vector-rotated-transform}
\end{equation}
which is not a probability characteristic function unless $m=0$.
\end{corollary}

\begin{proof}
The multivariate Gaussian moment-generating formula proves the identity;
the modulus exceeds one for every nonzero $t$.
\end{proof}

\subsection{Positive higher derivatives and removal of the stabilizer}

Let $-\Delta$ have positive eigenvalues $\lambda$ on the nonconstant scalar
modes of a compact background.  The simplest quadratic model of an
$R^2$-type conformal stabilizer is
\begin{equation}
 Q_{\varepsilon,b}(\sigma)
 =\frac12\sum_\lambda
 \bigl(\varepsilon\lambda^2-a\lambda+b\bigr)|\sigma_\lambda|^2,
 \qquad a>0,quad\varepsilon>0.
 \label{eq:n9v94-stabilized-quadratic}
\end{equation}

\begin{lemma}[Spectral positivity threshold]
\label{lem:n9v94-stabilizer-threshold}
For $b=0$, the form \eqref{eq:n9v94-stabilized-quadratic} is nonnegative on
all nonconstant modes if and only if
\begin{equation}
 \varepsilon\lambda_1\geq a,
 \label{eq:n9v94-epsilon-threshold}
\end{equation}
where $\lambda_1$ is the first positive eigenvalue.  Hence a family with
fixed $a$ cannot remain nonnegative while $\varepsilon\to0$.
For general fixed $b$, the minimum of
$\varepsilon\lambda^2-a\lambda+b$ over $\lambda\geq0$ is
$b-a^2/(4\varepsilon)$; it tends to $-\infty$ as
$\varepsilon\downarrow0$.
\end{lemma}

\begin{proof}
For $b=0$ the coefficient is
$\lambda(\varepsilon\lambda-a)$ and is nonnegative for every
$\lambda\geq\lambda_1$ exactly under
\eqref{eq:n9v94-epsilon-threshold}.  For general $b$, complete the square:
\[
 \varepsilon\lambda^2-a\lambda+b
 =\varepsilon\left(\lambda-\frac{a}{2\varepsilon}\right)^2
 +b-\frac{a^2}{4\varepsilon}.
\]
\end{proof}

\begin{corollary}[Collapse of uniform coercivity]
\label{cor:n9v94-coercivity-collapse}
The stabilizer in \eqref{eq:n9v94-stabilized-quadratic} cannot both be
removed with $\varepsilon\to0$ and provide a regulator-independent positive
coercivity constant for all conformal modes.  Therefore V93's exponential
tail constant cannot be obtained uniformly from this quadratic term alone
along such a removal.
\end{corollary}

\begin{proof}
A regulator-independent coercivity constant would in particular make all
quadratic coefficients uniformly positive.  Lemma~\ref{lem:n9v94-stabilizer-threshold}
excludes this as $\varepsilon\to0$.
\end{proof}

Positive higher derivatives also face an independent OS test.  The next
one-dimensional time calculation supplies an explicit negative reflected
norm.

\begin{theorem}[A fourth-order Gaussian reflection obstruction]
\label{thm:n9v94-fourth-order-OS}
Fix $m>0$ and $\varepsilon>0$.  The translation-invariant Gaussian
covariance with Fourier transform
\begin{equation}
 \widehat C_\varepsilon(p)
 =\frac1{(p^2+m^2)\,[1+\varepsilon(p^2+m^2)]}
 \label{eq:n9v94-fourth-order-covariance}
\end{equation}
is not reflection positive on the positive time half-line.
\end{theorem}

\begin{proof}
Set $M^2=m^2+\varepsilon^{-1}$.  Partial fractions give
\[
 \widehat C_\varepsilon(p)
 =\frac1{p^2+m^2}-\frac1{p^2+M^2}.
\]
After Fourier transformation in time,
\[
 C_\varepsilon(t+s)
 =\frac{e^{-m(t+s)}}{2m}-\frac{e^{-M(t+s)}}{2M}
 \qquad(t,s>0).
\]
The reflected quadratic form of a real test function $f$ supported in
positive time is therefore
\begin{equation}
 \mathcal Q(f)
 =\frac1{2m}\left|\int_0^\infty f(t)e^{-mt}dt\right|^2
 -\frac1{2M}\left|\int_0^\infty f(t)e^{-Mt}dt\right|^2.
 \label{eq:n9v94-negative-reflection-form}
\end{equation}
Choose two nonnegative smooth bumps with disjoint small supports near
distinct positive times.  A nontrivial linear combination can be chosen
so that its $e^{-mt}$ moment vanishes.  Since $e^{-mt}$ and $e^{-Mt}$ are
not proportional on two distinct times, choose the bumps sufficiently
narrow that the $e^{-Mt}$ moment of that combination is nonzero.  Then
\eqref{eq:n9v94-negative-reflection-form} is strictly negative.
\end{proof}

\begin{remark}[Interpretation]
\label{rem:n9v94-higher-derivative-interpretation}
The negative residue of the heavy pole is the reflected negative norm.
The theorem does not state that every possible gravitational
higher-curvature regulator violates OS positivity.  It proves that
positivity of a higher-derivative Euclidean quadratic form is not enough;
the induced covariance must separately pass the V83--V88 reflection test.
\end{remark}

\subsection{Constraint reduction: the algebraic gain and analytic debit}

An unstable variable can sometimes be removed by a constraint rather than
integrated.  The finite-dimensional calculation is elementary but
clarifies the distinction.

\begin{proposition}[Linear constraint reduction of an indefinite form]
\label{prop:n9v94-linear-constraint-reduction}
Let $H_x,H_y$ be finite-dimensional real Hilbert spaces, let
$A=A^*$ on $H_x$, let $C=C^*>0$ on $H_y$, and let
$B:H_x\to H_y$.  Define
\begin{equation}
 Q(x,y)=\frac12\langle Ax,x\rangle
 -\frac12\langle Cy,y\rangle+\langle Bx,y\rangle.
 \label{eq:n9v94-indefinite-block-form}
\end{equation}
The $y$-stationarity constraint is $Cy=Bx$.  On its graph,
\begin{equation}
 Q_{\rm red}(x):=Q(x,C^{-1}Bx)
 =\frac12\langle(A+B^*C^{-1}B)x,x\rangle.
 \label{eq:n9v94-reduced-positive-form}
\end{equation}
Thus $Q_{\rm red}$ is positive if
$A+B^*C^{-1}B>0$, even though the real Gaussian integral of
$e^{-Q(x,y)}$ over $y$ diverges.
\end{proposition}

\begin{proof}
Substitute $y=C^{-1}Bx$ into \eqref{eq:n9v94-indefinite-block-form}:
the last two terms combine to
$\frac12\langle C^{-1}Bx,Bx\rangle$.  For fixed $x$, the exponent
$-Q(x,y)$ grows like $\langle Cy,y\rangle/2$ along $|y|\to\infty$, so its
real $y$ integral diverges.
\end{proof}

\begin{remark}[Relation to the gravitational constraints]
\label{rem:n9v94-gravity-constraint-relation}
Proposition~\ref{prop:n9v94-linear-constraint-reduction} is a model of the
possible gain from solving lapse, shift, or conformal constraint variables
before constructing a measure.  It is not a proof that the nonlinear
Einstein constraint set has a global positive chart.  Such a proof needs
existence, uniqueness, regularity, uniform inverse bounds, control of
multiple branches, and compatibility with the V80--V81 moving calibration.
\end{remark}

\begin{theorem}[Conditional reduced-variable route]
\label{thm:n9v94-reduced-route}
Suppose a regulated gravitational construction supplies:
\begin{enumerate}
\item a measurable reduced variable space $X_i$ and a reconstruction map
$R_i:X_i\to\{\text{constraint-satisfying metrics and matter data}\}$;
\item a positive probability measure $\nu_i$ on $X_i$ satisfying the V93
normalized coercive-tail card;
\item uniform local bounds for $R_i$ and its physical lift as required in
V80, V81, V90, and V92;
\item a reflected positive local algebra on $X_i$ whose reconstructed
Schwinger functions obey the V86--V88 OS hypotheses.
\end{enumerate}
Then the reconstructed physical characteristic functionals are
equicontinuous and every pointwise cofinal limit has a positive bosonic
Minlos measure and OS Hilbert-space reconstruction.
\end{theorem}

\begin{proof}
Items 1--3 transport the V93 random-seminorm tail through reconstruction by
the V92 constraint-branch and lift theorems, giving equicontinuity and
Minlos reconstruction.  Item 4 passes reflection positivity to the limit
by V87 and invokes the quotient and semigroup construction of V88.
\end{proof}

\begin{firewall}[Locality and covariance after reduction]
\label{fire:n9v94-reduced-locality}
A constraint solution is generally nonlocal because it involves an inverse
elliptic operator.  Positivity of the reduced measure therefore does not
automatically imply Euclidean locality, Markov locality, or covariance of
the reconstructed metric observables.  These must be proved through the
intertwining estimates and symmetry-restoration rows of V89--V90.
\end{firewall}

\subsection{Route comparison and programme decision}

\begin{definition}[Gravitational-route acceptance card]
\label{def:n9v94-route-card}
A proposed treatment of the conformal direction is acceptable for the
present SM--Einstein programme only if it provides all of the following:
\begin{enumerate}
\item a normalized positive bosonic body for real physical observables, or
a reconstructed positive physical algebra with an explicit replacement for
that body;
\item regulator-uniform local field tails of the V92--V93 type;
\item reflection positivity on a local positive-time algebra;
\item uniform constraint and calibration reconstruction, including closing
modes;
\item cofinal removal of auxiliary contours or stabilizers without collapse
of the preceding constants;
\item Euclidean symmetry restoration and a route to the V88 semigroup.
\end{enumerate}
\end{definition}

The finite-dimensional conclusions of this note give the following exact
status.  Imaginary contour rotation supplies convergence but fails the real
probability characteristic test on the rotated coordinate.  A simple
positive fourth-order stabilizer cannot be removed with uniform quadratic
coercivity and can introduce a negative reflected norm.  Constraint
reduction can turn an indefinite block form into a positive reduced form,
but transfers the burden to global constraint charts, nonlocal
reconstruction, and OS compatibility.  Of these three routes, only the
last is presently aligned with the already developed physical-lift and
constraint ledgers, and even that route remains conditional.

\begin{assessment}[Status after ninth-series V94]
\label{assess:n9v94}
V94 rules out two shortcuts.  A convergent imaginary conformal contour is
not a real probability measure, and a positive fourth-order Euclidean
quadratic form need not be reflection positive or uniformly removable.
The constraint-reduced route has a rigorous conditional implication chain,
but the nonlinear reduced gravitational measure and its uniform local
chart have not been constructed.
\end{assessment}

\begin{programme}[Ninth-series V95: reduced gravitational chart audit]
\label{prog:n9v94-V95}
V95 will formulate the nonlinear reduced-chart theorem needed by
Theorem~\ref{thm:n9v94-reduced-route}: a quantitative implicit constraint
solve, its coarea density, branch exclusion, and the induced OS locality
debit.  Because V95 is a compulsory five-version boundary, it will first
inspect the newest YM and NS notes and transfer only genuinely new,
compatible estimates.
\end{programme}

\begin{firewall}[Claim boundary after ninth-series V94]
\label{fire:n9v94-boundary}
V94 does not construct a nonlinear Einstein constraint chart, a positive
reduced gravitational measure, an OS-positive gravitational regulator, a
uniform determinant density, an anomaly-free chiral determinant line, or
the full SM--Einstein continuum.
\end{firewall}

\paragraph{Parent-version record.}
V94 was formed from
\texttt{Renewal\_SM\_Einstein\_Continuum\_Research\_Notes\_V93.tex}.
The V93 parent had $37\,463$ validator-counted source lines,
$1\,350\,089$ bytes, and SHA--256
\[
 \texttt{8344192803E08ECC01BE09B61601DBFD6B9135FA8B737B2FE51899792BDD1704}.
\]

% =============================================================================
% END APPENDED NINTH-SERIES V94 MATERIAL
% =============================================================================

% =============================================================================
% BEGIN APPENDED NINTH-SERIES V95 MATERIAL
% =============================================================================

\section{Quantitative reduced constraint charts and the V95 companion audit}
\label{sec:v9v95}

\subsection{Compulsory companion audit}

Before continuing the reduced gravitational route chosen in V94, the
newest Yang--Mills and Navier--Stokes notes in the shared
\texttt{latex\_notes} directory were inspected read-only.  At the instant
of inspection the companion records were
\begin{align*}
 &\texttt{Renewal\_Yang\_Mills\_Mass\_Gap\_Research\_Notes\_V18.tex},
 &&302\,194\text{ bytes},\\
 &\texttt{Renewal\_Yang\_Mills\_Mass\_Gap\_Research\_Notes\_V19.tex},
 &&302\,194\text{ bytes},\\
 &\texttt{Renewal\_NS\_Stability\_Research\_Notes\_V26.tex},
 &&278\,283\text{ bytes}.
\end{align*}
The two YM files were byte-identical, with SHA--256
\[
 \texttt{34EB88AB7B46DC8491E653DDC8D96EC94B749AB3E15B8C47E102EA5D6412180A},
\]
so their live mathematical content ended at eighth-series V18.  The NS V26
digest was
\[
 \texttt{82C887F8C2C69E4F28FAD7C3DC8DE5B9099CB5A4BF431EBA1FFF3E91935978AD}.
\]

The YM result relevant here is exact conditional Gaussian matching.  For a
fine precision $H$, a surjective coarse map $B$, coarse precision
$S=(BH^{-1}B^*)^{-1}$, and conditional lift
$J=H^{-1}B^*S$, the whitened map $H^{1/2}JS^{-1/2}$ is an isometry.
Consequently conditional expectation is contractive on every
covariance-matched Wick chaos even when $J$ is not a small perturbation of a
simple coordinate lift.  The YM note correctly records localization of
this matched map as the remaining bottleneck.

The NS V26 result computes derivative norms on the rank-one tensor inputs
actually generated by the nonlinearity.  The restricted norm equals the
full norm for one derivative but is strictly smaller for part of the
second-derivative range.  Its transferable principle is to estimate a
constraint Hessian on tangent-graph inputs rather than automatically use
its full bilinear operator norm.  None of the NS numerical kernel constants
is a gravitational constant and none is imported.

\begin{assessment}[V95 audit decision]
\label{assess:n9v95-audit}
The reduced gravitational chart will be measured in its exact induced
energy metric, following the YM conditional-isometry result, and its
quadratic remainder will use a structured tangent-input norm, following the
NS rank-one refinement.  The YM locality warning is retained as a separate
OS debit.  No YM lattice constant, NS Oseen-kernel constant, mass-gap claim,
or regularity claim is transferred.
\end{assessment}

\subsection{A quantitative nonlinear constraint chart}

Let $X,Y,Z$ be finite-dimensional real normed spaces with
$\dim Y=\dim Z$, and let
\[
 \mathcal C:U_X\times U_Y\longrightarrow Z
\]
be $C^2$.  The independent variables are $x\in X$, the variables to be
reconstructed are $y\in Y$, and the constraint is
$\mathcal C(x,y)=0$.  Fix a bounded invertible map $A:Y\to Z$.

\begin{theorem}[Uniform contraction chart]
\label{thm:n9v95-contraction-chart}
Let $r_x,r_y>0$ and $0\leq\kappa<1$.  Suppose, for
$\|x\|\leq r_x$ and $\|y\|\leq r_y$,
\begin{align}
 \|I_Y-A^{-1}D_y\mathcal C(x,y)\|&\leq\kappa,
 \label{eq:n9v95-Dy-close}\\
 \|A^{-1}\mathcal C(x,0)\|&\leq(1-\kappa)r_y.
 \label{eq:n9v95-centre-residual}
\end{align}
Then for every $\|x\|\leq r_x$ there is a unique
$\psi(x)\in\overline B_Y(0,r_y)$ such that
$\mathcal C(x,\psi(x))=0$.  The map $\psi$ is $C^2$ on the interior,
\begin{align}
 \|\psi(x)\|&\leq
 \frac{\|A^{-1}\mathcal C(x,0)\|}{1-\kappa},
 \label{eq:n9v95-psi-bound}\\
 D\psi(x)&=-[D_y\mathcal C(x,\psi(x))]^{-1}
 D_x\mathcal C(x,\psi(x)),
 \label{eq:n9v95-psi-derivative}\\
 \|[D_y\mathcal C(x,\psi(x))]^{-1}\|
 &\leq\frac{\|A^{-1}\|}{1-\kappa}.
 \label{eq:n9v95-inverse-bound}
\end{align}
Uniqueness is asserted inside the declared $y$-ball; roots outside that
ball are not excluded.
\end{theorem}

\begin{proof}
For fixed $x$, define
$T_x(y)=y-A^{-1}\mathcal C(x,y)$.  By the mean-value formula and
\eqref{eq:n9v95-Dy-close},
$\|T_x(y)-T_x(y')\|\leq\kappa\|y-y'\|$.  Moreover,
\[
 \|T_x(y)\|
 \leq\|T_x(y)-T_x(0)\|+\|T_x(0)\|
 \leq\kappa r_y+\|A^{-1}\mathcal C(x,0)\|
 \leq r_y.
\]
Thus $T_x$ maps the closed ball to itself and is a contraction.  Banach's
fixed-point theorem gives the unique fixed point there, which is exactly a
constraint root.  Comparing the fixed point to zero gives
\eqref{eq:n9v95-psi-bound}.

The bound \eqref{eq:n9v95-Dy-close} says
$A^{-1}D_y\mathcal C=I-E$ with $\|E\|\leq\kappa$; the Neumann series gives
invertibility and \eqref{eq:n9v95-inverse-bound}.  The ordinary implicit
function theorem then gives $C^2$ regularity, and differentiating
$\mathcal C(x,\psi(x))=0$ gives
\eqref{eq:n9v95-psi-derivative}.
\end{proof}

\begin{corollary}[First response bound]
\label{cor:n9v95-first-response}
If $\|D_x\mathcal C\|\leq L_x$ on the chart, then
\begin{equation}
 \|D\psi(x)\|\leq
 L_\psi:=\frac{\|A^{-1}\|L_x}{1-\kappa}.
 \label{eq:n9v95-first-response-bound}
\end{equation}
\end{corollary}

\begin{proof}
Combine \eqref{eq:n9v95-psi-derivative} and
\eqref{eq:n9v95-inverse-bound}.
\end{proof}

\subsection{Second response on the actual tangent graph}

Write $J_xh=(h,D\psi(x)h)$ for the tangent graph lift.  Define the
structured constraint-Hessian norm
\begin{equation}
 \|D^2\mathcal C\|_{J_x}
 :=\sup_{h,k\neq0}
 \frac{\|D^2\mathcal C(x,\psi(x))[J_xh,J_xk]\|}
 {\|h\|\,\|k\|}.
 \label{eq:n9v95-structured-Hessian}
\end{equation}
This is at most
$\|D^2\mathcal C\|\|J_x\|^2$, but it can be smaller because the two
arguments lie on the tangent graph.

\begin{proposition}[Exact second response]
\label{prop:n9v95-second-response}
On the chart of Theorem~\ref{thm:n9v95-contraction-chart},
\begin{equation}
 D^2\psi(x)[h,k]
 =-[D_y\mathcal C]^{-1}
 D^2\mathcal C[J_xh,J_xk],
 \label{eq:n9v95-second-response}
\end{equation}
where all derivatives of $\mathcal C$ are evaluated at
$(x,\psi(x))$.  Hence
\begin{equation}
 \|D^2\psi(x)\|
 \leq\frac{\|A^{-1}\|}{1-\kappa}
 \|D^2\mathcal C\|_{J_x}.
 \label{eq:n9v95-second-response-bound}
\end{equation}
\end{proposition}

\begin{proof}
Differentiate
$D\mathcal C(x,\psi(x))[J_xh]=0$ in direction $k$.  The derivative of
the first differential is
$D^2\mathcal C[J_xh,J_xk]$, while the derivative of the graph vector
contributes $D_y\mathcal C\,D^2\psi[h,k]$.  Solving gives
\eqref{eq:n9v95-second-response}; use
\eqref{eq:n9v95-inverse-bound} for the estimate.
\end{proof}

\begin{remark}[Transfer from the NS audit]
\label{rem:n9v95-NS-transfer}
The numerator in \eqref{eq:n9v95-second-response} never receives arbitrary
vectors in $X\oplus Y$: it receives $J_xh,J_xk$.  Replacing
\eqref{eq:n9v95-structured-Hessian} by the full Hessian norm can therefore
lose a genuine geometric factor, just as the NS rank-one calculation can
improve the second-derivative norm.  Any improvement here must nevertheless
be computed for the Einstein constraint tensor; the NS numerical gain does
not transfer.
\end{remark}

\subsection{Exact delta-constraint density and its variation}

Fix Lebesgue densities on $X,Y,Z$ compatible with chosen orientations.
The absolute determinant below is independent of the orientation choice.

\begin{theorem}[Local delta reduction]
\label{thm:n9v95-delta-reduction}
Under Theorem~\ref{thm:n9v95-contraction-chart}, suppose that the declared
$y$-ball contains the only roots included in the regulated integration
domain.  For every continuous compactly supported $F$ in the chart,
\begin{equation}
 \int_X\int_Y F(x,y)\,\delta(\mathcal C(x,y))\,dy\,dx
 =\int_X
 \frac{F(x,\psi(x))}
 {|\det D_y\mathcal C(x,\psi(x))|}\,dx.
 \label{eq:n9v95-delta-reduction}
\end{equation}
Thus a positive ambient density $e^{-S(x,y)}$ induces the positive reduced
density
\begin{equation}
 e^{-S(x,\psi(x))}
 |\det D_y\mathcal C(x,\psi(x))|^{-1}dx.
 \label{eq:n9v95-reduced-density}
\end{equation}
\end{theorem}

\begin{proof}
For fixed $x$, the inverse-function theorem makes
$y\mapsto z=\mathcal C(x,y)$ a local diffeomorphism near its unique root.
The change-of-variables formula gives
\[
 \int F(x,y)\delta(\mathcal C(x,y))dy
 =\int F(x,y_x(z))\delta(z)
 |\det D_y\mathcal C(x,y_x(z))|^{-1}dz,
\]
which evaluates at $z=0$.  Integrate over $x$.
\end{proof}

\begin{remark}[Delta density versus induced surface measure]
\label{rem:n9v95-delta-versus-surface}
The density in \eqref{eq:n9v95-delta-reduction} is the one defined by the
ambient delta function and the chosen constraint coordinates.  Hausdorff
measure on the graph instead contains
$\det(I+D\psi^*D\psi)^{1/2}dx$.  Confusing these two measures loses a
Jacobian.  A Faddeev--Popov or coarea determinant must be combined with the
specific measure actually used, as already required by V82.
\end{remark}

\begin{proposition}[Aligned logarithmic Jacobian response]
\label{prop:n9v95-logdet-response}
Let $K(x)=D_y\mathcal C(x,\psi(x))$.  Along any direction $h\in X$,
\begin{equation}
 D\log|\det K(x)|[h]
 =\operatorname{tr}\!\left(
 K(x)^{-1}
 \bigl[D_xD_y\mathcal C[h]+D_y^2\mathcal C[D\psi(x)h]\bigr]
 \right).
 \label{eq:n9v95-logdet-response}
\end{equation}
Consequently,
\begin{equation}
 |D\log|\det K|[h]|
 \leq \|K^{-1}\|_{\rm op}
 \|D_xD_y\mathcal C[h]+D_y^2\mathcal C[D\psi h]\|_{\rm tr},
 \label{eq:n9v95-logdet-trace-bound}
\end{equation}
where $\|\cdot\|_{\rm tr}$ is the trace norm after identifying $Y$ and $Z$.
\end{proposition}

\begin{proof}
Jacobi's formula gives
$D\log|\det K|=\operatorname{tr}(K^{-1}DK)$ while the determinant stays
nonzero.  The chain rule gives the displayed derivative of $K$, and
$|\operatorname{tr}(PQ)|\leq\|P\|_{\rm op}\|Q\|_{\rm tr}$ gives the bound.
\end{proof}

\begin{firewall}[Dimension and locality of the determinant row]
\label{fire:n9v95-determinant-extensivity}
A bound obtained from
$|\operatorname{tr}T|\leq(\dim Y)\|T\|$ is extensive and does not yield a
cofinal local density estimate.  The trace in
\eqref{eq:n9v95-logdet-response} must be localized into uniformly summable
blocks or renormalized as a density.  It must then pass the V92 $L^s$
reweighting card.  Pointwise invertibility of $K$ alone is insufficient.
\end{firewall}

\subsection{Matched energy geometry of the nonlinear lift}

Let $H_x$ be a positive definite quadratic form on $X\oplus Y$ at a point
of the constraint graph.  The tangent lift is $J_x:X\to X\oplus Y$ and its
induced precision is
\begin{equation}
 S_x:=J_x^*H_xJ_x.
 \label{eq:n9v95-induced-precision}
\end{equation}

\begin{theorem}[Exact tangent energy isometry]
\label{thm:n9v95-tangent-isometry}
If $J_x$ is injective and $H_x>0$, then $S_x>0$ and
\begin{equation}
 R_x:=H_x^{1/2}J_xS_x^{-1/2}
 \quad\text{satisfies}\quad R_x^*R_x=I_X.
 \label{eq:n9v95-tangent-isometry}
\end{equation}
Thus
\begin{equation}
 \|J_xh\|_{H_x}=\|h\|_{S_x},
 \qquad
 \|(R_x^*)^{\otimes n}\|\leq1
 \quad(n\geq1).
 \label{eq:n9v95-tensor-isometry}
\end{equation}
No small ambient-norm bound on $D\psi$ is required for these identities.
\end{theorem}

\begin{proof}
By definition,
$R_x^*R_x=S_x^{-1/2}J_x^*H_xJ_xS_x^{-1/2}=I$.  The first equality in
\eqref{eq:n9v95-tensor-isometry} is
$\langle H_xJ_xh,J_xh\rangle=\langle S_xh,h\rangle$; the tensor bound
follows from $\|R_x^*\|=1$.
\end{proof}

\begin{assessment}[Transfer from the YM audit]
\label{assess:n9v95-YM-transfer}
Theorem~\ref{thm:n9v95-tangent-isometry} is the nonlinear tangent analogue
of the YM exact conditional isometry.  It removes a false requirement that
the physical constraint lift be perturbatively close to a simple ambient
lift.  It does not remove the need to compare $H_x,S_x$ between nearby
points, nor does it prove that the square roots or the graph lift are local.
For polymer and OS estimates, quadratic-form energy frames should be used
whenever possible so that nonlocal square roots are not inserted merely to
write an isometry.
\end{assessment}

\subsection{Branch exhaustion and OS locality}

The local theorem does not show that every relevant configuration belongs
to one chart.  A cofinal construction must make this omission probabilistic
and uniform.

\begin{definition}[Cofinal branch-exhaustion card]
\label{def:n9v95-branch-card}
There are nested chart domains $\mathfrak U_i(R)$ and a decreasing function
$a_{\rm br}(R)\to0$ such that
\begin{equation}
 \sup_i\mu_i\{\text{configuration is outside }\mathfrak U_i(R)\}
 \leq a_{\rm br}(R),
 \label{eq:n9v95-branch-tail}
\end{equation}
and on each $\mathfrak U_i(R)$ the constraint roots are exhausted by a
finite set of labelled charts with compatible overlap densities.  If the
number of branches is not one, the branch label is retained as a discrete
physical variable until equivalence or suppression is proved.
\end{definition}

\begin{proposition}[Branch tail preserves the V92 card]
\label{prop:n9v95-branch-tail-transfer}
Suppose the reconstructed random seminorm is bounded by $X_i(R)$ inside
$\mathfrak U_i(R)$ and
\begin{equation}
 \sup_i\mu_i\{X_i(R)>L\}\leq a_R(L),
 \qquad a_R(L)\longrightarrow0
 \quad(L\to\infty).
 \label{eq:n9v95-inside-tail}
\end{equation}
If there is a diagonal choice $R(L)\to\infty$ for which
$a_{R(L)}(L)+a_{\rm br}(R(L))\to0$, then the complete reconstructed field
satisfies the V92 tail card.
\end{proposition}

\begin{proof}
The event that the full seminorm exceeds $L$ is contained in the union of
the chart-exit event and the inside-chart exceedance event.  Apply the union
bound and the stated diagonal choice.
\end{proof}

Let $\theta$ be time reflection and let $\mathcal A_+$ denote the algebra
generated by independent variables supported at positive time.

\begin{proposition}[A sufficient support condition for reduced OS algebras]
\label{prop:n9v95-OS-support}
Suppose the reconstruction $R_i$ satisfies
\begin{equation}
 R_i\theta=\theta R_i,
 \qquad
 F\in\mathcal A_+^{\rm full}Longrightarrow F\circ R_i\in\mathcal A_+^{\rm red}.
 \label{eq:n9v95-OS-support-condition}
\end{equation}
If the reduced measure $\nu_i$ is reflection positive on
$\mathcal A_+^{\rm red}$, then the reconstructed full Schwinger functional
is reflection positive on $\mathcal A_+^{\rm full}$.
\end{proposition}

\begin{proof}
For $F\in\mathcal A_+^{\rm full}$, reflection covariance gives
$\theta(F\circ R_i)=(\theta F)\circ R_i$.  Therefore
\[
 \int\overline{\theta F(R_i x)}F(R_i x)d\nu_i(x)
 =\int\overline{\theta(F\circ R_i)(x)}(F\circ R_i)(x)d\nu_i(x)
 \geq0.
\]
\end{proof}

\begin{firewall}[Elliptic reconstruction is generally two-sided]
\label{fire:n9v95-elliptic-OS}
Solving an elliptic Einstein constraint on a full Euclidean time slice or
spacetime generally produces an inverse kernel with support on both sides
of the reflection plane.  Then the second implication in
\eqref{eq:n9v95-OS-support-condition} fails: a positive-time full
observable can pull back to negative-time reduced variables.  Exponential
off-diagonal decay only makes this a small leakage; it does not give exact
finite-regulator OS positivity.  One must retain boundary variables, use a
one-sided domain decomposition, or prove positivity by a different exact
factorization.
\end{firewall}

\subsection{Reduced gravitational chart card}

\begin{definition}[Quantitative reduced-gravity chart card]
\label{def:n9v95-reduced-chart-card}
At every cofinal regulator the construction must provide:
\begin{enumerate}
\item contraction-chart constants $A,\kappa,r_x,r_y$ uniform in the
declared local norm, with closing modes retained;
\item first and structured second response bounds from
Corollary~\ref{cor:n9v95-first-response} and
Proposition~\ref{prop:n9v95-second-response};
\item the exact delta/coarea density and a localized logarithmic determinant
bound which passes the V92 reweighting card;
\item positive induced energy forms $H_x,S_x$ and controlled variation of
their quadratic-form frames;
\item the branch-exhaustion tail of
Definition~\ref{def:n9v95-branch-card};
\item an exact reflection-compatible domain decomposition satisfying
Proposition~\ref{prop:n9v95-OS-support}, or an alternative direct OS
factorization;
\item the V80--V81 calibration connection and V90 physical lift, evaluated
on the same charts.
\end{enumerate}
\end{definition}

\begin{theorem}[Conditional reduced-route closure]
\label{thm:n9v95-conditional-route-closure}
Assume the reduced-variable probability measures pass the V93
coercive-tail card and the reconstruction passes
Definition~\ref{def:n9v95-reduced-chart-card}.  Assume also pointwise
convergence of all physical characteristic functionals and local Schwinger
functions.  Then the limit has a Radon bosonic measure on the physical
distribution space and a positive OS pre-Hilbert reconstruction.  The
result is independent of the chosen cofinal tower whenever the V87 hybrid
comparisons and V92 tail transports identify the limiting characteristic
and Schwinger functionals.
\end{theorem}

\begin{proof}
V93 gives the random-seminorm tail on reduced variables.  The response,
branch, determinant, and lift rows transport it to physical observables by
V92 and Proposition~\ref{prop:n9v95-branch-tail-transfer}.  V91--V92 give
the Radon measure.  Exact finite-regulator reflection positivity follows
from Proposition~\ref{prop:n9v95-OS-support}; V87 passes it to the limit and
V88 forms the OS quotient.  The last assertion follows from the V87
cross-tower comparison and uniqueness of characteristic reconstruction.
\end{proof}

\begin{assessment}[Status after ninth-series V95]
\label{assess:n9v95}
V95 supplies the quantitative nonlinear chart, exact response formulas,
delta density, matched tangent isometry, branch-tail criterion, and a
sufficient OS support theorem required by the V94 reduced route.  The
active upstream bottleneck is now exact reflection-compatible localization
of an elliptic constraint solve.  Ordinary global elliptic inversion is
two-sided in Euclidean time and does not preserve the positive-time
algebra.
\end{assessment}

\begin{programme}[Ninth-series V96: boundary-retained domain decomposition]
\label{prog:n9v95-V96}
Guided by the YM localization warning, V96 will keep reflection-plane data
as explicit variables and split the constraint solve into two half-space
problems.  It will prove the finite-dimensional conditional-independence
and Schur-complement factorization which preserves reflection positivity,
state the nonlinear domain-decomposition card, and quantify the gluing
debit without assuming that an elliptic inverse is one-sided.
\end{programme}

\begin{firewall}[Claim boundary after ninth-series V95]
\label{fire:n9v95-boundary}
V95 does not construct a global Einstein constraint atlas, prove branch
exhaustion, localize the coarea determinant, produce a reflection-compatible
nonlinear domain decomposition, control the chiral determinant line,
derive the interacting tail card, or construct the full SM--Einstein
continuum.
\end{firewall}

\paragraph{Parent-version record.}
V95 was formed from
\texttt{Renewal\_SM\_Einstein\_Continuum\_Research\_Notes\_V94.tex}.
The V94 parent had $37\,821$ validator-counted source lines,
$1\,365\,230$ bytes, and SHA--256
\[
 \texttt{BBDE2ED51094B2DDA45121E1CDB933E46ABCD9A0F8BCB1B9CB053462273B6E21}.
\]

% =============================================================================
% END APPENDED NINTH-SERIES V95 MATERIAL
% =============================================================================

% =============================================================================
% BEGIN APPENDED NINTH-SERIES V96 MATERIAL
% =============================================================================

\section{Boundary-retained domain decomposition and reflected squares}
\label{sec:v9v96}

\subsection{Replacing a two-sided inverse by conditional half-space solves}

V95 showed that an elliptic constraint reconstruction generally fails to
preserve the positive-time algebra when written as one global inverse.  The
correct object for OS positivity is instead a domain decomposition: retain
all reflection-plane data as boundary variables and solve the two interiors
conditionally and separately.  The resulting reflected square is exact and
does not require an elliptic Green function to have one-sided support.

Let $X_+$ and $X_-$ be measurable interior spaces, let $B$ be a boundary
space, and let $\theta:X_+\to X_-$ be a measurable bijection.  Reflection
acts by
\[
 \Theta(x_+,x_-,b)=(\theta^{-1}x_-,\theta x_+,b).
\]

\begin{theorem}[Conditional reflected-square criterion]
\label{thm:n9v96-reflected-square}
Let $\nu$ be a positive finite measure on $B$, and for $\nu$-almost every
$b$ let $K_b$ be a positive finite measure on $X_+$.  Suppose
\begin{equation}
 0<N:=\int_B K_b(X_+)^2\,d\nu(b)<\infty.
 \label{eq:n9v96-normalization}
\end{equation}
Define a probability measure $\mu$ on $X_+\times X_-\times B$ by
\begin{equation}
 \int G\,d\mu
 :=N^{-1}\int_B\int_{X_+}\int_{X_+}
 G(x_+,\theta y_+,b)
 \,dK_b(x_+)dK_b(y_+)d\nu(b).
 \label{eq:n9v96-factorized-measure}
\end{equation}
Then $\mu$ is reflection invariant and reflection positive.  More exactly,
for every bounded measurable $F=F(x_+,b)$,
\begin{equation}
 \int \overline{F(\Theta q)}F(q)\,d\mu(q)
 =N^{-1}\int_B
 \left|\int_{X_+}F(x_+,b)dK_b(x_+)\right|^2d\nu(b)
 \geq0.
 \label{eq:n9v96-reflected-square}
\end{equation}
\end{theorem}

\begin{proof}
Reflection invariance follows by exchanging the two $K_b$ integrations.
For $q=(x_+,\theta y_+,b)$ one has
$F(q)=F(x_+,b)$ and
$F(\Theta q)=F(y_+,b)$.  Fubini's theorem turns the two interior integrals
into the squared modulus in \eqref{eq:n9v96-reflected-square}.
\end{proof}

\begin{remark}[Boundary variables belong to the positive algebra]
\label{rem:n9v96-boundary-algebra}
The positive-time algebra may depend on $b$, because reflection fixes the
boundary.  Its OS norm is the $L^2(\nu)$ norm of the conditional amplitude
$b\mapsto\int F(x_+,b)dK_b(x_+)$.  Eliminating $b$ before establishing this
square can hide the positivity mechanism and create a nonlocal cross-plane
kernel.
\end{remark}

\subsection{The Gaussian Schur realization}

The reflected-square criterion is realized exactly by local positive
Gaussian actions.  Let $H_+,H_B$ be finite-dimensional real Hilbert spaces,
let $A=A^*>0$ on $H_+$, let $D=D^*$ on $H_B$, and let
$C:H_+\to H_B$.  Use reflected copies $x,y\in H_+$ and boundary
$b\in H_B$.

\begin{theorem}[Boundary-retained Gaussian factorization]
\label{thm:n9v96-gaussian-factorization}
Consider
\begin{equation}
 Q(x,y,b)
 =\frac12\langle Ax,x\rangle+\langle Cx,b\rangle
 +\frac12\langle Ay,y\rangle+\langle Cy,b\rangle
 +\frac12\langle Db,b\rangle.
 \label{eq:n9v96-gaussian-action}
\end{equation}
Set
\begin{equation}
 S_B:=D-2CA^{-1}C^*.
 \label{eq:n9v96-boundary-Schur}
\end{equation}
Then the full quadratic form is positive definite if and only if
$S_B>0$.  In that case the normalized Gaussian measure
$e^{-Q}dx\,dy\,db$ has the factorization
\begin{equation}
 e^{-Q}dx\,dy\,db
 =e^{-\langle S_Bb,b\rangle/2}db
 \prod_{u\in\{x,y\}}
 e^{-\langle A(u+A^{-1}C^*b),u+A^{-1}C^*b\rangle/2}du,
 \label{eq:n9v96-gaussian-square-factorization}
\end{equation}
up to a positive constant.  Hence it satisfies
Theorem~\ref{thm:n9v96-reflected-square} and is reflection positive.
\end{theorem}

\begin{proof}
Complete the square separately in $x$ and $y$:
\[
 \frac12\langle Au,u\rangle+\langle Cu,b\rangle
 =\frac12\langle A(u+A^{-1}C^*b),u+A^{-1}C^*b\rangle
 -\frac12\langle CA^{-1}C^*b,b\rangle.
\]
Adding the two copies gives
\eqref{eq:n9v96-gaussian-square-factorization}.  Since $A>0$, the Schur
complement criterion says $Q>0$ exactly when $S_B>0$.  The two conditional
interior kernels at fixed $b$ are identical, so
Theorem~\ref{thm:n9v96-reflected-square} applies.
\end{proof}

\begin{corollary}[Boundary covariance and conditional means]
\label{cor:n9v96-gaussian-conditionals}
Under Theorem~\ref{thm:n9v96-gaussian-factorization}, the boundary
covariance is $S_B^{-1}$ and, conditioned on $b$, the two interiors are
independent Gaussians with covariance $A^{-1}$ and mean
$-A^{-1}C^*b$.  Thus the global elliptic coupling between the halves is
carried entirely by the retained boundary variable.
\end{corollary}

\begin{proof}
Read the covariance and means from
\eqref{eq:n9v96-gaussian-square-factorization}.
\end{proof}

\begin{remark}[Local cut sets]
\label{rem:n9v96-local-cut}
For a finite-range discretization, enlarge the reflection plane to a
boundary layer containing every cell touched by an interaction crossing
the plane.  Removing that layer disconnects the positive and negative
interiors, so no direct $x$--$y$ block remains.  The theorem then applies
without estimating decay of a global inverse.  The layer thickness is the
interaction range and must remain controlled under the regulator flow.
\end{remark}

\subsection{Nonlinear constraints on the two halves}

Let $u_\pm$ denote independent interior variables and $v_\pm$ variables to
be reconstructed.  Retain boundary data $b$.  On the positive half assume
a constraint
\[
 \mathcal C_+(u_+,v_+;b)=0
\]
with a V95 contraction chart
$v_+=\psi_+(u_+;b)$.  Define the negative-half constraint and chart by
reflection:
\[
 \mathcal C_-:=\theta\mathcal C_+\theta^{-1},
 \qquad \psi_-:=\theta\psi_+\theta^{-1}.
\]

\begin{theorem}[Nonlinear reflected constraint factorization]
\label{thm:n9v96-nonlinear-factorization}
Suppose:
\begin{enumerate}
\item at each retained $b$, the positive and negative constraint roots are
exhausted by the reflected V95 charts above;
\item the action decomposes as
\begin{equation}
 S=S_B(b)+S_+(u_+,v_+;b)
 +S_+(\theta^{-1}u_-,\theta^{-1}v_-;b);
 \label{eq:n9v96-action-decomposition}
\end{equation}
\item the reference density and constraint delta functions factor over the
two interiors at fixed $b$;
\item all coarea and gauge-slice factors on the negative half are reflected
copies of positive factors, and their product with the boundary density is
nonnegative.
\end{enumerate}
Then reduction of $v_\pm$ by the V95 delta formula gives
\begin{equation}
 d\mu=N^{-1}d\nu(b)\,dK_b(u_+)\,dK_b(\theta^{-1}u_-),
 \label{eq:n9v96-nonlinear-product}
\end{equation}
with positive $K_b$ and $\nu$.  Consequently the reconstructed theory is
reflection positive on observables depending on the positive interior and
the retained boundary.
\end{theorem}

\begin{proof}
Apply the V95 delta-reduction theorem separately to
$\mathcal C_+$ and $\mathcal C_-$.  The action decomposition makes the two
Boltzmann factors identical after reflection.  The coarea determinants are
absolute determinants and likewise form reflected copies.  All remaining
boundary-only factors are absorbed into $d\nu$.  The assumed nonnegativity
makes these genuine measures, and
Theorem~\ref{thm:n9v96-reflected-square} proves the result.
\end{proof}

\begin{firewall}[The gluing equations are not optional]
\label{fire:n9v96-gluing}
The retained datum $b$ must include every trace needed to glue the two
half-solutions into one admissible configuration: metric and matter traces,
the necessary normal data, boundary gauge frames or edge modes, and any
constraint flux.  If a further gluing equation couples $u_+$ directly to
$u_-$ after $b$ is fixed, conditional independence fails.  The cure is to
enlarge $b$ until the coupling factors through it, not to discard the
equation.
\end{firewall}

\begin{proposition}[Tail assembly through the boundary]
\label{prop:n9v96-tail-assembly}
Assume, under the probability version of
\eqref{eq:n9v96-nonlinear-product}, that
\begin{align}
 |\Phi_B(f)|&\leq X_Bp(f),
 \label{eq:n9v96-boundary-tail}\\
 |\Phi_+(f)|&\leq X_+(b)p(f),
 \label{eq:n9v96-half-tail}
\end{align}
and the negative field is the reflected copy.  If $X_B$ is uniformly tight
under the boundary marginal and
\begin{equation}
 \sup_{i,b}K_{i,b}\{X_+(b)>R\}/K_{i,b}(X_+)
 \leq a_+(R)\longrightarrow0,
 \label{eq:n9v96-uniform-conditional-tail}
\end{equation}
then the full smeared field satisfies the V92 tail card.
\end{proposition}

\begin{proof}
The full random seminorm is bounded by $X_B+X_++X_-$.  Conditional on $b$,
the two interior tails obey the same uniform bound.  Integrating the union
bound over the boundary marginal gives
\[
 \mu\{X_B+X_++X_->R\}
 \leq \mu\{X_B>R/3\}+2a_+(R/3),
\]
which tends to zero uniformly.
\end{proof}

\subsection{Approximate factorizations and the limit}

Exact factorization is the clean finite-regulator route.  For a regulator
whose cross-plane leakage is nonzero, one may still pass positivity to the
limit if the negative part is quantitatively vanishing on a common dense
algebra.

\begin{lemma}[Vanishing reflected negativity]
\label{lem:n9v96-vanishing-negativity}
Let $\mathcal A_+$ be a complex vector space with seminorm $\|\cdot\|_*$.
Suppose regulated reflected forms $q_i$ satisfy
\begin{equation}
 q_i(F,F)\geq-\varepsilon_i\|F\|_*^2,
 \qquad \varepsilon_i\longrightarrow0,
 \label{eq:n9v96-approximate-OS}
\end{equation}
and $q_i(F,F)\to q(F,F)$ for every $F\in\mathcal A_+$.  Then
$q(F,F)\geq0$ for all $F\in\mathcal A_+$.
\end{lemma}

\begin{proof}
Take the limit inferior in \eqref{eq:n9v96-approximate-OS} for fixed $F$.
\end{proof}

\begin{proposition}[Trace-norm control of a reflected kernel debit]
\label{prop:n9v96-kernel-debit}
Let $q_i^0(F,F)\geq0$ be an exactly factorized reflected form represented on
a Hilbert completion by a positive trace-class operator $T_i^0$.  Let the
actual form be represented by a self-adjoint trace-class operator
$T_i=T_i^0+E_i$.  Then
\begin{equation}
 q_i(F,F)\geq-\|E_i\|_{\rm op}\|F\|^2
 \geq-\|E_i\|_1\|F\|^2.
 \label{eq:n9v96-trace-debit}
\end{equation}
Hence $\|E_i\|_1\to0$ supplies
Lemma~\ref{lem:n9v96-vanishing-negativity}.
\end{proposition}

\begin{proof}
$\langle F,T_iF\rangle
=\langle F,T_i^0F\rangle+\langle F,E_iF\rangle
\geq-\|E_i\|_{\rm op}\|F\|^2$, and
$\|E_i\|_{\rm op}\leq\|E_i\|_1$.
\end{proof}

\begin{remark}[Small action error is not yet a small reflected operator]
\label{rem:n9v96-action-versus-kernel}
An $L^1$ difference between scalar Boltzmann densities does not
automatically give the trace-norm estimate in
Proposition~\ref{prop:n9v96-kernel-debit} on an unbounded observable
algebra.  The comparison must be proved in the reflected kernel norm or on
a uniformly bounded dense algebra, followed by closure.  This is the same
distinction between scalar reweighting and positive-channel control made in
V82 and V87.
\end{remark}

\subsection{Gauge, ghost, and fermion boundary data}

\begin{proposition}[Gauge quotient with retained edge modes]
\label{prop:n9v96-gauge-edge-modes}
Suppose the gauge slice on each half is fixed relative to a boundary gauge
frame $g_B$, and all gauge transformations nontrivial on the cut are
retained in $g_B$.  If the two half-space Faddeev--Popov/coarea densities
are nonnegative reflected copies and the boundary Haar density is positive,
then the bosonic gauge-reduced measure has the factorization
\eqref{eq:n9v96-nonlinear-product}.
\end{proposition}

\begin{proof}
At fixed $g_B$, the two based gauge groups act independently on the two
interiors.  Their slice changes of variables therefore factor.  Reflection
identifies the two absolute Jacobians, and the remaining Haar integration
is boundary-only, so it is part of $d\nu$.
\end{proof}

\begin{firewall}[Orientation signs and chiral phases]
\label{fire:n9v96-boundary-phases}
Proposition~\ref{prop:n9v96-gauge-edge-modes} assumes nonnegative bosonic
slice densities.  An oriented ghost determinant can carry a sign, and a
chiral fermion determinant is a section of a determinant line rather than
a positive scalar.  Reflected pairing may produce an absolute square only
after anomaly cancellation, compatible boundary conditions, and a proved
trivialization.  The absolute coarea determinant cannot be silently
substituted for those data.
\end{firewall}

For fermions, retain the cut spinors and apply the boundary Schur reduction
of V83 separately on the two halves.  If the induced one-particle reflected
boundary matrix is positive, the exterior-power theorem of V83 gives the
graded reflected square.  This fermionic condition is independent of the
bosonic conditional factorization and both must hold.

\subsection{The domain-decomposition card}

\begin{definition}[Reflection-compatible domain-decomposition card]
\label{def:n9v96-domain-card}
A regulated reduced SM--Einstein measure passes the domain-decomposition
card if it provides:
\begin{enumerate}
\item a finite-width reflection boundary layer carrying every interaction,
constraint, gauge, and fermion trace that crosses the cut;
\item V95 contraction charts on each half at fixed boundary data, related
exactly by reflection;
\item an action, reference density, coarea factor, and gauge quotient which
factor as reflected positive half-kernels times a positive boundary density;
\item either exact finite-regulator factorization or a reflected-operator
debit tending to zero as in
Proposition~\ref{prop:n9v96-kernel-debit};
\item uniform conditional half-space and boundary field-tail estimates as
in Proposition~\ref{prop:n9v96-tail-assembly};
\item a boundary fermion/ghost treatment satisfying the V83 graded
one-particle criterion and the determinant-line requirements;
\item gluing consistency and branch exhaustion with the boundary label
retained until uniqueness is proved.
\end{enumerate}
\end{definition}

\begin{theorem}[Boundary-retained implication]
\label{thm:n9v96-boundary-implication}
If a cofinal family passes
Definition~\ref{def:n9v96-domain-card}, its limiting local Schwinger
functional is reflection positive and its bosonic characteristic functional
is continuous at the origin.  Under the V87 convergence, V91 nuclearity,
and V88 time-translation hypotheses, it therefore yields a Radon bosonic
measure and an OS Hilbert space with a nonnegative Hamiltonian.
\end{theorem}

\begin{proof}
Exact positivity follows from
Theorem~\ref{thm:n9v96-reflected-square}; in the approximate case it follows
from Lemma~\ref{lem:n9v96-vanishing-negativity}.  The conditional tail
assembly gives the V92 card and hence characteristic continuity.  Apply the
V87 passage to the limit, the V91--V92 Minlos reconstruction, and the V88 OS
quotient and semigroup theorem.
\end{proof}

\begin{assessment}[Status after ninth-series V96]
\label{assess:n9v96}
V96 resolves the formal OS-locality problem of a two-sided elliptic inverse
at the structural level: retain the reflection boundary, solve the two
halves conditionally, and obtain an exact reflected square.  It supplies
Gaussian and nonlinear factorization theorems, tail assembly, and a
vanishing-negativity alternative.  What remains is to verify this card for
a concrete regulated Einstein constraint system with gauge and chiral
matter, including a positive boundary Schur form and uniform local chart
constants.
\end{assessment}

\begin{programme}[Ninth-series V97: linearized Einstein boundary symbol]
\label{prog:n9v96-V97}
V97 will test the card on the linearized Euclidean Einstein constraints in
a flat half-space model.  It will separate gauge, transverse-traceless,
constraint, conformal, and boundary modes; compute the half-space
Dirichlet-to-Neumann quadratic symbol; and determine which retained
boundary channels are positive and which still inherit the conformal sign.
\end{programme}

\begin{firewall}[Claim boundary after ninth-series V96]
\label{fire:n9v96-boundary}
V96 does not prove that the Einstein constraint action admits the required
positive half-kernels, compute its boundary symbol, establish anomaly-free
fermion gluing, derive uniform nonlinear chart constants, construct the
interacting measure, or complete the SM--Einstein continuum.
\end{firewall}

\paragraph{Parent-version record.}
V96 was formed from
\texttt{Renewal\_SM\_Einstein\_Continuum\_Research\_Notes\_V95.tex}.
The V95 parent had $38\,320$ validator-counted source lines,
$1\,384\,654$ bytes, and SHA--256
\[
 \texttt{17878F216C92E06FB1E458254BD67FA53A04C562805252597E9AF5891D6C820B}.
\]

% =============================================================================
% END APPENDED NINTH-SERIES V96 MATERIAL
% =============================================================================

% =============================================================================
% BEGIN APPENDED NINTH-SERIES V97 MATERIAL
% =============================================================================

\section{Linearized Einstein half-space symbol and the conformal boundary channel}
\label{sec:v9v97}

\subsection{Quadratic form and two invariant test sectors}

V96 showed that a local positive action factorizes across a retained
reflection boundary.  Boundary retention does not, however, change the sign
of an interior quadratic form.  This note computes the relevant sign in the
flat linearized Einstein model and then identifies the additional role of
the Hamiltonian constraint.

Use the Euclidean Einstein--Hilbert convention of V93,
\[
 S_E(g)=-\frac1{16\pi G}\int R_g\,d^4x,
\]
and expand $g_{\mu\nu}=\delta_{\mu\nu}+h_{\mu\nu}$ about flat space.  Up to
a boundary term, the quadratic Fierz--Pauli form is
\begin{equation}
 S_E^{(2)}(h)=\frac1{64\pi G}\int
 \left(
 \partial_\lambda h_{\mu\nu}\partial_\lambda h_{\mu\nu}
 -2\partial_\mu h_{\mu\nu}\partial_\lambda h_{\lambda\nu}
 +2\partial_\mu h_{\mu\nu}\partial_\nu h
 -\partial_\lambda h\partial_\lambda h
 \right)d^4x,
 \label{eq:n9v97-FP-form}
\end{equation}
where $h=h_{\mu\mu}$.  The formula is used only on compactly supported
fields or with the declared boundary term included.

\begin{proposition}[TT and conformal restrictions]
\label{prop:n9v97-sector-signs}
On transverse-traceless fields,
\begin{equation}
 \partial_\mu h_{\mu\nu}=0,qquad h=0,
 \qquad
 S_E^{(2)}(h)=\frac1{64\pi G}\int|\partial h|^2d^4x\geq0.
 \label{eq:n9v97-TT-positive}
\end{equation}
On pure conformal fields $h_{\mu\nu}=2\sigma\delta_{\mu\nu}$,
\begin{equation}
 S_E^{(2)}(2\sigma\delta)
 =-\frac3{8\pi G}\int|\partial\sigma|^2d^4x\leq0.
 \label{eq:n9v97-conformal-negative}
\end{equation}
\end{proposition}

\begin{proof}
The TT formula follows immediately from
\eqref{eq:n9v97-FP-form}.  In the conformal sector,
$h=8\sigma$ and
$\partial_\mu h_{\mu\nu}=2\partial_\nu\sigma$.  The four terms in the
integrand of \eqref{eq:n9v97-FP-form} have respective coefficients
$16,-8,32,-64$ multiplying $|\partial\sigma|^2$.  Their sum is $-24$,
which gives \eqref{eq:n9v97-conformal-negative}.
\end{proof}

\begin{remark}[Gauge directions]
\label{rem:n9v97-gauge-directions}
Infinitesimal diffeomorphism directions
$h_{\mu\nu}=\partial_\mu\xi_\nu+\partial_\nu\xi_\mu$ are null directions of
the gauge-invariant quadratic form modulo boundary terms.  They must be
quotiented, fixed with their Faddeev--Popov factor, or retained as boundary
gauge data.  Adding a positive gauge-fixing term controls these null
directions but does not change the negative conformal restriction
\eqref{eq:n9v97-conformal-negative}.
\end{remark}

\subsection{The half-space Dirichlet-to-Neumann calculation}

Let $t>0$ be Euclidean time and $x\in\mathbb R^3$ tangential to the
reflection plane.  For a tangential Fourier frequency $k\neq0$, put
$\omega=|k|$.

\begin{lemma}[Scalar half-space action]
\label{lem:n9v97-scalar-DN}
Let $b(k)$ be Schwartz and let
\begin{equation}
 u(t,k)=e^{-|k|t}b(k).
 \label{eq:n9v97-harmonic-extension}
\end{equation}
Then $(-\partial_t^2+|k|^2)u=0$, $u(0,k)=b(k)$, and, with a fixed unitary
Fourier convention,
\begin{equation}
 \int_0^\infty\int_{\mathbb R^3}
 \bigl(|\partial_tu|^2+|\nabla_xu|^2\bigr)dx\,dt
 =\int_{\mathbb R^3}|k|\,|b(k)|^2dk.
 \label{eq:n9v97-DN-form}
\end{equation}
Thus the Dirichlet-to-Neumann operator is $|D|=(-\Delta_x)^{1/2}$.
\end{lemma}

\begin{proof}
The differential equation and boundary value are immediate.  Plancherel's
theorem makes the integrand at frequency $k$ equal
$2|k|^2e^{-2|k|t}|b(k)|^2$.  Its $t$ integral is
$|k||b(k)|^2$.
\end{proof}

\begin{theorem}[Linearized Einstein boundary symbol]
\label{thm:n9v97-Einstein-boundary-symbol}
Consider decaying half-space solutions in either of the following sectors.
\begin{enumerate}
\item $h_{0\mu}=0$ and the tangential boundary tensor $b_{ab}(k)$ is
transverse and trace free:
$k_ab_{ab}=0$, $b_{aa}=0$.
\item $h_{\mu\nu}=2\sigma\delta_{\mu\nu}$ with boundary value $s(k)$.
\end{enumerate}
The corresponding on-shell half-space quadratic forms are
\begin{align}
 Q_+^{\rm TT}(b)
 &=\frac1{64\pi G}\int |k|\,|b_{ab}(k)|^2dk,
 \label{eq:n9v97-TT-boundary-form}\\
 Q_+^{\rm conf}(s)
 &=-\frac3{8\pi G}\int |k|\,|s(k)|^2dk.
 \label{eq:n9v97-conformal-boundary-form}
\end{align}
Hence the TT Dirichlet-to-Neumann symbol is positive, while the conformal
symbol has the same magnitude sign defect as the bulk action.
\end{theorem}

\begin{proof}
Extend every boundary component by
\eqref{eq:n9v97-harmonic-extension}.  The declared tensor conditions make
the extension four-dimensionally transverse and trace free, so
\eqref{eq:n9v97-TT-positive} and Lemma~\ref{lem:n9v97-scalar-DN} give
\eqref{eq:n9v97-TT-boundary-form}.  The conformal result follows from
\eqref{eq:n9v97-conformal-negative} and the same lemma.
\end{proof}

\begin{corollary}[Boundary retention alone does not repair the sign]
\label{cor:n9v97-retention-not-repair}
If the conformal boundary value $s$ is integrated as a real Gaussian
variable with precision read from
\eqref{eq:n9v97-conformal-boundary-form}, the integral diverges.  Thus the
positive boundary Schur hypothesis of the V96 Gaussian factorization fails
in this channel.
\end{corollary}

\begin{proof}
The exponent $-Q_+^{\rm conf}$ is positive and quadratic in every nonzero
mode $s(k)$, so its real Gaussian integral diverges.
\end{proof}

\begin{remark}[Infrared modes]
\label{rem:n9v97-infrared}
At $k=0$, the decaying harmonic extension used above does not exist for a
nonzero constant boundary value and the symbol $|k|$ vanishes.  Constant
volume, toron-like, and global constraint modes must therefore be retained
as separate infrared variables.  They cannot be placed inside the inverse
$|D|^{-1}$ used below.
\end{remark}

\subsection{What the linearized spatial constraint removes}

The four-dimensional Euclidean conformal mode and the conformal coordinate
used in canonical constraint data should not be identified without a
calculation.  Nevertheless, the linearized vacuum Hamiltonian constraint
shows the mechanism by which nonzero scalar boundary data can be removed
from the physical slice.

Let $\gamma_{ij}=\delta_{ij}+q_{ij}$ be a perturbation of a flat
three-dimensional slice.  Its linearized scalar curvature is
\begin{equation}
 \delta R^{(3)}(q)=\partial_i\partial_jq_{ij}-\Delta\operatorname{tr}q.
 \label{eq:n9v97-linearized-scalar-curvature}
\end{equation}

\begin{proposition}[Vacuum scalar constraint in transverse gauge]
\label{prop:n9v97-vacuum-constraint}
Impose the transverse gauge $\partial_jq_{ij}=0$ and the vacuum linearized
Hamiltonian constraint $\delta R^{(3)}(q)=0$.  Then every nonzero Fourier
mode is trace free.  Hence the nonzero modes of $q$ are TT.
\end{proposition}

\begin{proof}
In transverse gauge,
\eqref{eq:n9v97-linearized-scalar-curvature} becomes
$-\Delta\operatorname{tr}q=0$.  At $k\neq0$ this says
$|k|^2\operatorname{tr}\widehat q(k)=0$, so the trace vanishes.  Together
with transversality this is the TT condition.
\end{proof}

\begin{corollary}[Sourced trace and inverse-order debit]
\label{cor:n9v97-sourced-trace}
If the linearized constraint is
$\delta R^{(3)}(q)=16\pi G\rho$ in transverse gauge, then
\begin{equation}
 \operatorname{tr}\widehat q(k)
 =16\pi G|k|^{-2}\widehat\rho(k)
 \qquad(k\neq0).
 \label{eq:n9v97-sourced-trace}
\end{equation}
Thus the scalar metric datum is determined rather than freely integrated,
but its reconstruction loses two derivatives and is infrared singular.
\end{corollary}

\begin{proof}
Fourier transform the sourced version of
\eqref{eq:n9v97-linearized-scalar-curvature} after imposing transversality.
\end{proof}

\begin{firewall}[Four-dimensional and canonical reductions]
\label{fire:n9v97-two-conformal-notions}
Proposition~\ref{prop:n9v97-vacuum-constraint} does not by itself remove the
four-dimensional conformal factor from a covariant Euclidean path integral.
It concerns canonical boundary data after a spatial gauge choice.  To use
it in the V96 domain decomposition, lapse, shift, extrinsic curvature,
boundary conditions, the momentum constraints, and the coarea density must
all be included in one reflected reduction.
\end{firewall}

\subsection{A positive reduced boundary Schur mechanism}

The correct algebraic possibility is to constrain the negative scalar
channel instead of integrating it.  The following operator version of the
V94 block calculation is the relevant boundary test.

\begin{theorem}[Constraint elimination of a negative boundary channel]
\label{thm:n9v97-negative-channel-elimination}
Let $\mathcal H_m,\mathcal H_s$ be real Hilbert spaces.  Let
$A=A^*\geq a_0I>0$ on $\mathcal H_m$, let
$C=C^*\geq c_0I>0$ on $\mathcal H_s$, and let
$B:\mathcal H_m\to\mathcal H_s$ be bounded.  Consider
\begin{equation}
 Q(m,s)=\frac12\langle Am,m\rangle
 -\frac12\langle Cs,s\rangle+\langle Bm,s\rangle.
 \label{eq:n9v97-boundary-indefinite-block}
\end{equation}
Do not integrate $s$; impose its stationarity constraint
$Cs=Bm$.  Then
\begin{equation}
 Q_{\rm red}(m)
 =\frac12\left\langle
 (A+B^*C^{-1}B)m,m\right\rangle
 \geq\frac{a_0}{2}\|m\|^2.
 \label{eq:n9v97-boundary-reduced-positive}
\end{equation}
The graph lift $m\mapsto(m,C^{-1}Bm)$ is an exact isometry between the
induced energy metric and the ambient reduced graph metric in the sense of
V95.
\end{theorem}

\begin{proof}
Substitution gives
\[
 -\frac12\langle C^{-1}Bm,Bm\rangle
 +\langle Bm,C^{-1}Bm\rangle
 =\frac12\langle B^*C^{-1}Bm,m\rangle.
\]
The added operator is positive.  The isometry is the V95 tangent-isometry
identity applied to this linear graph.
\end{proof}

\begin{remark}[Order of the conformal boundary operator]
\label{rem:n9v97-conformal-order}
For the half-space conformal form,
$C$ is a positive multiple of $|D|$ on nonzero modes, even though it enters
$Q$ with a minus sign.  If a source coupling has symbol $B(k)$, the induced
positive term has symbol $B(k)^*|k|^{-1}B(k)$.  Uniform tails therefore
depend on the low-frequency vanishing order of $B(k)$ and on separate
control of the retained zero mode.
\end{remark}

\begin{firewall}[The constraint must be the correct one]
\label{fire:n9v97-correct-constraint}
The positivity in
Theorem~\ref{thm:n9v97-negative-channel-elimination} uses exactly the
stationarity equation $Cs=Bm$ of the same quadratic form.  An independently
imposed equation $Ls=Bm$ gives a different reduced form and need not be
positive.  A concrete Einstein regulator must prove that its lapse,
Hamiltonian, and boundary equations produce the required compatible Schur
relation after gauge reduction.
\end{firewall}

\subsection{Linearized boundary acceptance card}

\begin{definition}[Einstein boundary-symbol card]
\label{def:n9v97-boundary-symbol-card}
At each nonzero tangential momentum, a regulated linearized gravitational
boundary reduction must provide:
\begin{enumerate}
\item a gauge quotient or slice separating pure gauge modes and retaining
boundary gauge transformations;
\item a positive TT Dirichlet-to-Neumann block comparable to $|k|$;
\item an explicit scalar constraint eliminating every negative conformal
block by a compatible Schur relation of
Theorem~\ref{thm:n9v97-negative-channel-elimination};
\item uniform bounds for the scalar response $C^{-1}B$ in the physical
source and test seminorms;
\item separate variables and calibration conditions for $k=0$ and every
closing mode;
\item reflected equality of the two half-space symbols and positive gluing
density on the retained boundary;
\item the matching matter, gauge, ghost, and fermion boundary blocks of the
V96 domain-decomposition card.
\end{enumerate}
\end{definition}

\begin{theorem}[Linearized conditional conclusion]
\label{thm:n9v97-linearized-conclusion}
If a finite regulator passes
Definition~\ref{def:n9v97-boundary-symbol-card}, its reduced nonzero-mode
bosonic boundary quadratic form is positive and the two reflected halves
give an OS-positive Gaussian functional.  If the comparison constants and
zero-mode tails are uniform, the resulting characteristic functionals pass
the V92 tail card.
\end{theorem}

\begin{proof}
The TT block is positive by
Theorem~\ref{thm:n9v97-Einstein-boundary-symbol}.  The constrained scalar
block is positive by
Theorem~\ref{thm:n9v97-negative-channel-elimination}; pure gauge modes have
been removed and zero modes retained separately.  V96 then writes the two
halves as a reflected Gaussian square.  Uniform positive comparison of the
precision with the declared physical Sobolev weights gives the V93 Gaussian
energy moment and hence the V92 tail card.
\end{proof}

\begin{assessment}[Status after ninth-series V97]
\label{assess:n9v97}
V97 computes the flat half-space symbols.  The TT boundary form is positive
and proportional to $|k|$; the freely integrated conformal boundary form is
negative with the same symbol.  Boundary retention therefore does not cure
the conformal problem.  At the linearized canonical level, the vacuum
Hamiltonian constraint removes the nonzero trace in transverse gauge, and
a compatible stationary Schur constraint turns a negative scalar block
into a positive reduced contribution.  The full coupled Einstein constraint
must still be shown to have exactly this compatibility.
\end{assessment}

\begin{programme}[Ninth-series V98: sourced boundary Schur and infrared ledger]
\label{prog:n9v97-V98}
V98 will analyze the coupled source operator $B$ in the scalar boundary
Schur term, derive sharp Fourier-order conditions for bounded response and
field tails, separate neutral and charged zero modes, and state the finite
set of matter and gauge source estimates needed before nonlinear
perturbation of the boundary chart can begin.
\end{programme}

\begin{firewall}[Claim boundary after ninth-series V97]
\label{fire:n9v97-boundary}
V97 is a flat linearized calculation.  It does not prove the nonlinear
Einstein constraint Schur identity, positive-energy control with Standard
Model sources, zero-mode tightness, a uniform interacting boundary measure,
chiral determinant-line gluing, or the full SM--Einstein continuum.
\end{firewall}

\paragraph{Parent-version record.}
V97 was formed from
\texttt{Renewal\_SM\_Einstein\_Continuum\_Research\_Notes\_V96.tex}.
The V96 parent had $38\,737$ validator-counted source lines,
$1\,401\,439$ bytes, and SHA--256
\[
 \texttt{C5DF5D591799D535CB9F7E36680B91C46762E15D77E04E71FBE1EA9A3CE95A7F}.
\]

% =============================================================================
% END APPENDED NINTH-SERIES V97 MATERIAL
% =============================================================================

% =============================================================================
% BEGIN APPENDED NINTH-SERIES V98 MATERIAL
% =============================================================================

\section{Sourced conformal Schur response and the infrared ledger}
\label{sec:v9v98}

\subsection{Fourier order of the reduced scalar channel}

V97 showed that the nonzero conformal boundary precision has magnitude
$C_0=c|D|$, $c>0$, but enters the unreduced Einstein quadratic form with a
negative sign.  If a compatible constraint is
$C_0s=Bm$, eliminating $s$ gives the positive Schur term
$B^*C_0^{-1}B$.  The purpose of V98 is to record the exact derivative and
infrared costs of this statement.

For a Schwartz function on $\mathbb R^3$ with Fourier transform
$\widehat f$, use the homogeneous seminorm
\begin{equation}
 \|f\|_{\dot H^r}^2
 :=\int_{\mathbb R^3}|k|^{2r}|\widehat f(k)|^2dk
 \label{eq:n9v98-homogeneous-norm}
\end{equation}
whenever the integral is finite.  Zero modes are excluded from homogeneous
inverse powers and will be treated separately.

\begin{theorem}[Exact order of a scalar boundary Schur term]
\label{thm:n9v98-order}
Let $C_0=c|D|$, $c>0$, and let $B=b|D|^\beta$ on a scalar source $m$, with
$b\in\mathbb R$ and $\beta\in\mathbb R$.  On nonzero modes the constraint
$C_0s=Bm$ has solution
\begin{equation}
 s=\frac bc|D|^{\beta-1}m.
 \label{eq:n9v98-response}
\end{equation}
Moreover,
\begin{align}
 \|s\|_{\dot H^r}
 &=\frac{|b|}{c}\|m\|_{\dot H^{r+\beta-1}},
 \label{eq:n9v98-response-Sobolev}\\
 \langle Bm,C_0^{-1}Bm\rangle
 &=\frac{b^2}{c}\|m\|_{\dot H^{\beta-1/2}}^2.
 \label{eq:n9v98-Schur-Sobolev}
\end{align}
The reduced contribution is nonnegative, but it is infrared singular when
$\beta<1/2$.
\end{theorem}

\begin{proof}
At each $k\neq0$ the response multiplier is
$(b/c)|k|^{\beta-1}$, proving
\eqref{eq:n9v98-response} and
\eqref{eq:n9v98-response-Sobolev}.  The Schur multiplier is
$b^2|k|^{2\beta}/(c|k|)$, which gives
\eqref{eq:n9v98-Schur-Sobolev}.  Its sign is nonnegative because
$C_0^{-1}>0$.
\end{proof}

\begin{corollary}[Hamiltonian-source order]
\label{cor:n9v98-Hamiltonian-order}
The flat linearized spatial constraint of V97 has
$|D|^2s=\gamma\rho$ for a constant $\gamma$.  Rewriting it as
$C_0s=B\rho$ with $C_0=c|D|$ requires
\begin{equation}
 B=c\gamma|D|^{-1}.
 \label{eq:n9v98-Hamiltonian-B}
\end{equation}
Thus $\beta=-1$,
\begin{align}
 s&=\gamma|D|^{-2}\rho,
 \label{eq:n9v98-Poisson-response}\\
 \langle B\rho,C_0^{-1}B\rho\rangle
 &=c\gamma^2\|\rho\|_{\dot H^{-3/2}}^2.
 \label{eq:n9v98-Hamiltonian-Schur}
\end{align}
The boundary stationary source is one inverse derivative rougher than the
bulk constraint source.
\end{corollary}

\begin{proof}
The identity $c|D|s=c\gamma|D|^{-1}\rho$ is equivalent to the constraint on
nonzero modes.  Apply Theorem~\ref{thm:n9v98-order} with $b=c\gamma$ and
$\beta=-1$.
\end{proof}

\begin{remark}[Compatibility is a substantive condition]
\label{rem:n9v98-compatibility}
The boundary variation of a harmonic conformal half-space action naturally
has order $|D|s$.  The spatial Hamiltonian constraint has order
$|D|^2s$.  Equation \eqref{eq:n9v98-Hamiltonian-B} shows the nonlocal source
operator required to make them the same stationary Schur equation.  A
complete derivation must obtain this operator from lapse and boundary
reduction; it cannot be inserted solely to make the sign positive.
\end{remark}

\subsection{Mean modes, neutrality, and calibrated backgrounds}

The norm in \eqref{eq:n9v98-Hamiltonian-Schur} is not automatically finite
near $k=0$.

\begin{lemma}[Neutrality gives one infrared zero]
\label{lem:n9v98-neutrality}
Let $\rho\in L^1(\mathbb R^3)$ satisfy
\begin{equation}
 \int_{\mathbb R^3}\rho(x)dx=0,
 \qquad
 \int_{\mathbb R^3}|x|\,|\rho(x)|dx<\infty.
 \label{eq:n9v98-neutrality-hyp}
\end{equation}
Then
\begin{equation}
 |\widehat\rho(k)|
 \leq |k|\int|x|\,|\rho(x)|dx.
 \label{eq:n9v98-neutrality-bound}
\end{equation}
Consequently the low-frequency part of
$\|\rho\|_{\dot H^{-3/2}}$ is finite if $\widehat\rho$ is otherwise locally
bounded.
\end{lemma}

\begin{proof}
Use the zero mean to write
\[
 \widehat\rho(k)=\int(e^{-\mathrm i k\cdot x}-1)\rho(x)dx
\]
and $|e^{-\mathrm i u}-1|\leq|u|$.  For $|k|<1$, the integrand of the
squared homogeneous norm is bounded by a constant times $|k|^{-1}$;
including the three-dimensional radial measure gives
$\int_0^1r\,dr<\infty$.
\end{proof}

\begin{lemma}[Nonzero mean produces a logarithmic obstruction]
\label{lem:n9v98-nonzero-mean}
If $\rho\in L^1(\mathbb R^3)$ and
$\widehat\rho(0)\neq0$, then
$\rho\notin\dot H^{-3/2}$.
\end{lemma}

\begin{proof}
Continuity of $\widehat\rho$ gives
$|\widehat\rho(k)|\geq c_0>0$ on a ball about zero.  Therefore
\[
 \int_{|k|<\varepsilon}|k|^{-3}|\widehat\rho(k)|^2dk
 \geq4\pi c_0^2\int_0^\varepsilon\frac{dr}{r}=\infty.
\]
\end{proof}

For Standard Model energy density, literal neutrality is generally the
wrong condition.  The mean source must instead be carried by a background
Einstein equation or retained global mode.

\begin{definition}[Mean/fluctuation calibration]
\label{def:n9v98-mean-calibration}
On a compact regulated spatial region, split
\begin{equation}
 \rho=\bar\rho+\rho_0,
 \qquad
 \bar\rho:=|\Sigma|^{-1}\int_\Sigma\rho,
 \qquad
 \int_\Sigma\rho_0=0.
 \label{eq:n9v98-mean-split}
\end{equation}
The mean $\bar\rho$ is assigned to a retained finite-dimensional background
coordinate $z$ satisfying a calibrated background Einstein equation.  Only
$\rho_0$ is placed in the inverse Laplacian.  The calibration response of
$z$ is governed by the V80--V81 connection and must remain uniformly
invertible.
\end{definition}

\begin{proposition}[Finite-volume inverse after mean removal]
\label{prop:n9v98-finite-volume-inverse}
Let $-\Delta_0$ be the Laplacian restricted to mean-zero functions on a
compact connected slice.  Then $-\Delta_0$ is positive and invertible, and
the sourced trace is
\begin{equation}
 s_0=\gamma(-\Delta_0)^{-1}\rho_0.
 \label{eq:n9v98-mean-zero-inverse}
\end{equation}
The bound
\begin{equation}
 \|s_0\|_{L^2}
 \leq\gamma\lambda_1^{-1}\|\rho_0\|_{L^2}
 \label{eq:n9v98-Poincare-response}
\end{equation}
depends on the first positive eigenvalue $\lambda_1$ and therefore is not
volume-uniform on expanding flat tori.
\end{proposition}

\begin{proof}
The spectral theorem on the orthogonal complement of constants gives
invertibility and \eqref{eq:n9v98-mean-zero-inverse}.  The inverse norm is
$\lambda_1^{-1}$, proving \eqref{eq:n9v98-Poincare-response}.
\end{proof}

\begin{remark}[Thermodynamic versus local norms]
\label{rem:n9v98-local-norms}
The loss in \eqref{eq:n9v98-Poincare-response} is unavoidable in a global
$L^2$ norm.  A continuum construction must instead use local negative
Sobolev norms, screened variables, background response, or retained
infrared coordinates.  This is parallel to the YM audit's separation of
localized energy frames from global matched-chaos contraction.
\end{remark}

\subsection{Stability under a perturbed scalar precision}

The flat operator $C_0$ will be perturbed by background curvature,
calibration, and interactions.  Relative form control is the correct
criterion.

\begin{theorem}[Relative Schur stability]
\label{thm:n9v98-relative-Schur}
Let $C_0>0$ on a Hilbert space and let $E=E^*$ satisfy
\begin{equation}
 \|C_0^{-1/2}EC_0^{-1/2}\|\leq\eta<1.
 \label{eq:n9v98-relative-perturbation}
\end{equation}
Set $C=C_0+E$.  Then
\begin{equation}
 (1-\eta)C_0\preceq C\preceq(1+\eta)C_0
 \label{eq:n9v98-form-comparison}
\end{equation}
and
\begin{equation}
 (1+\eta)^{-1}C_0^{-1}
 \preceq C^{-1}
 \preceq(1-\eta)^{-1}C_0^{-1}.
 \label{eq:n9v98-inverse-comparison}
\end{equation}
For every source map $B$,
\begin{equation}
 \frac1{1+\eta}B^*C_0^{-1}B
 \preceq B^*C^{-1}B
 \preceq\frac1{1-\eta}B^*C_0^{-1}B.
 \label{eq:n9v98-Schur-comparison}
\end{equation}
\end{theorem}

\begin{proof}
Write
$C=C_0^{1/2}(I+K)C_0^{1/2}$ with
$K=C_0^{-1/2}EC_0^{-1/2}$ and $-\eta I\preceq K\preceq\eta I$.
This proves \eqref{eq:n9v98-form-comparison}.  Invert the positive operator
inequality to obtain \eqref{eq:n9v98-inverse-comparison}, and sandwich it
between $B^*$ and $B$.
\end{proof}

\begin{corollary}[Response debit]
\label{cor:n9v98-response-debit}
Under Theorem~\ref{thm:n9v98-relative-Schur},
\begin{equation}
 \|C_0^{1/2}C^{-1}B m\|
 \leq\frac1{1-\eta}\|C_0^{-1/2}Bm\|.
 \label{eq:n9v98-response-debit}
\end{equation}
\end{corollary}

\begin{proof}
$C^{-1}=C_0^{-1/2}(I+K)^{-1}C_0^{-1/2}$ and
$\|(I+K)^{-1}\|\leq(1-\eta)^{-1}$.
\end{proof}

\begin{firewall}[Absolute perturbations miss the infrared]
\label{fire:n9v98-absolute-perturbation}
An operator bound $\|E\|\leq\varepsilon$ does not imply
\eqref{eq:n9v98-relative-perturbation} when $C_0=c|D|$ has spectrum closing
at zero.  Low modes must be retained or the perturbation must vanish with
the same Fourier order.  This is the boundary version of the closing-mode
rule in V84, V90, and V92.
\end{firewall}

\subsection{Several sources and source-tail transport}

Let $m=(m_1,\ldots,m_N)$ collect energy, momentum, Higgs, gauge, and other
source channels, and write $B m=\sum_aB_am_a$.

\begin{proposition}[Positive source Gram matrix]
\label{prop:n9v98-source-Gram}
The block operator
\begin{equation}
 \mathcal G_{ab}:=B_a^*C^{-1}B_b
 \label{eq:n9v98-source-Gram}
\end{equation}
is positive semidefinite on the direct sum of source spaces.  Its cross
terms satisfy
\begin{equation}
 |\langle B_am_a,C^{-1}B_bm_b\rangle|
 \leq
 \langle B_am_a,C^{-1}B_am_a\rangle^{1/2}
 \langle B_bm_b,C^{-1}B_bm_b\rangle^{1/2}.
 \label{eq:n9v98-source-cross}
\end{equation}
\end{proposition}

\begin{proof}
$\mathcal G=B^*C^{-1}B=(C^{-1/2}B)^*(C^{-1/2}B)\succeq0$.
The cross estimate is Cauchy--Schwarz in the $C^{-1}$ inner product.
\end{proof}

\begin{remark}[Rank is fixed by the constraint channel]
\label{rem:n9v98-source-rank}
Many matter sources may feed one scalar constraint, but the resulting Gram
matrix has rank at most the scalar constraint space.  Treating every source
as an independent positive scalar Schur direction double counts the same
constraint.  The complete block $B^*C^{-1}B$ must be retained.
\end{remark}

\begin{proposition}[Quadratic-source exponential tail]
\label{prop:n9v98-quadratic-source-tail}
Let $X_i\geq0$ satisfy
\begin{equation}
 \sup_i\mathbb E e^{\alpha X_i^2}\leq M.
 \label{eq:n9v98-field-subgaussian-norm}
\end{equation}
Suppose a nonlinear source seminorm obeys
\begin{equation}
 Y_i\leq C(1+X_i^2).
 \label{eq:n9v98-quadratic-source-bound}
\end{equation}
Then for $R>C$,
\begin{equation}
 \sup_i\mathbb P\{Y_i>R\}
 \leq M\exp\!\left[-\frac\alpha C(R-C)\right].
 \label{eq:n9v98-source-tail}
\end{equation}
If a reconstructed scalar seminorm is bounded linearly by $Y_i$ in a
declared negative Sobolev source norm, it is uniformly tight and therefore
passes the V92 tail card.
\end{proposition}

\begin{proof}
$Y_i>R$ implies $X_i^2>(R-C)/C$.  Apply Markov's inequality to
$e^{\alpha X_i^2}$.  A bounded linear response preserves the resulting
tail after rescaling.
\end{proof}

\begin{firewall}[Composite fields require renormalization]
\label{fire:n9v98-composite-renormalization}
The pointwise Standard Model stress tensor is a composite field.  The bound
\eqref{eq:n9v98-quadratic-source-bound} is not automatic in the continuum
and generally requires regulated products, counterterms, and a uniform
composite-operator norm.  V98 proves tail transport after such a bound; it
does not construct the renormalized stress tensor.
\end{firewall}

\subsection{Infrared and source acceptance card}

\begin{definition}[Sourced conformal boundary card]
\label{def:n9v98-source-card}
A cofinal reduced Einstein boundary construction passes the source card if
it supplies:
\begin{enumerate}
\item an exact derivation of the compatible stationary constraint
$Cs=Bm$ from the regulated action, lapse, and boundary terms;
\item a mean/fluctuation split assigning every zero mode to a finite
calibrated background variable rather than to $C^{-1}$;
\item source-space bounds for $C^{-1}B$ and $B^*C^{-1}B$ in local or
homogeneous Sobolev norms, uniform in volume and mesh;
\item relative precision control
\eqref{eq:n9v98-relative-perturbation} on all inverted modes;
\item a complete source Gram block, including cross terms, with no channel
double counting;
\item uniform composite-source tails of the form
\eqref{eq:n9v98-source-tail};
\item separate Gauss-law charge/flux, momentum-constraint, and global
geometric ledgers, with their boundary edge modes retained;
\item compatibility with the V95 calibration chart and V96 reflected
factorization.
\end{enumerate}
\end{definition}

\begin{theorem}[Source-card implication]
\label{thm:n9v98-source-card-implication}
If the TT block passes the V97 boundary-symbol card and the scalar block
passes Definition~\ref{def:n9v98-source-card}, the complete reduced
linearized bosonic boundary form is positive on all inverted modes.  If the
declared source tails and calibrated zero-mode tails are uniform, the
reconstructed metric fields satisfy the V92 characteristic-tail card.
\end{theorem}

\begin{proof}
TT positivity is V97.  Scalar positivity and its stability follow from the
V97 constraint-elimination theorem and
Theorem~\ref{thm:n9v98-relative-Schur}.  Response bounds transport the
source tails to the scalar metric, while the zero modes are controlled in
their finite calibration chart.  Finite species assembly and constraint
reconstruction are V92.
\end{proof}

\begin{assessment}[Status after ninth-series V98]
\label{assess:n9v98}
V98 computes the exact sourced boundary orders.  The Hamiltonian-source
trace response loses two derivatives and its positive stationary Schur term
requires an $\dot H^{-3/2}$ source norm.  A nonzero mean violates that norm,
so the mean energy must move a retained background coordinate rather than
enter the inverse.  Relative form control stabilizes the nonzero-mode
Schur operator, and sub-Gaussian field norms yield exponential tails for a
quadratically bounded regulated source.  The missing inputs are the exact
action/constraint compatibility and the renormalized stress-tensor bound.
\end{assessment}

\begin{programme}[Ninth-series V99: nonlinear sourced Schur chart]
\label{prog:n9v98-V99}
V99 will perturb the sourced scalar equation nonlinearly in the weighted
spaces fixed here.  It will prove a contraction theorem with a retained mean
coordinate, derive first and second source responses, and give a quantitative
positivity margin for the reduced Hessian under relative nonlinear
remainders.
\end{programme}

\begin{firewall}[Claim boundary after ninth-series V98]
\label{fire:n9v98-boundary}
V98 does not derive the full Einstein boundary constraint from a regulator,
renormalize the Standard Model stress tensor, control nonneutral infinite-
volume sources, prove nonlinear Schur positivity, establish determinant-line
gluing, or construct the full SM--Einstein continuum.
\end{firewall}

\paragraph{Parent-version record.}
V98 was formed from
\texttt{Renewal\_SM\_Einstein\_Continuum\_Research\_Notes\_V97.tex}.
The V97 parent had $39\,103$ validator-counted source lines,
$1\,416\,170$ bytes, and SHA--256
\[
 \texttt{62BFACE21D331372AA008B7AD5AFDCE815F08BEB14F35FFFBB4C9B956E3D6084}.
\]

% =============================================================================
% END APPENDED NINTH-SERIES V98 MATERIAL
% =============================================================================

% =============================================================================
% BEGIN APPENDED NINTH-SERIES V99 MATERIAL
% =============================================================================

\section{Nonlinear sourced scalar chart and reduced-Hessian margin}
\label{sec:v9v99}

\subsection{Weighted spaces after removal of the mean mode}

V98 identified the correct nonzero-mode scalar precision $C_0$ and showed
that ordinary operator norms become singular at low momentum.  The natural
fixed-point norm is instead the energy norm of $C_0$.  Let $S$ be a real
Hilbert space on which $C_0$ is positive and injective after all zero and
closing modes have been retained outside the inverse.  Define
\begin{equation}
 \|s\|_S:=\|C_0^{1/2}s\|,
 \qquad
 \|j\|_{S^*}:=\|C_0^{-1/2}j\|.
 \label{eq:n9v99-energy-dual-norms}
\end{equation}
Then $C_0^{-1}:S^*\to S$ has norm one.  The variables denoted by $m$ below
include matter and gauge sources together with the finite calibrated
background coordinates of V80--V81 and V98; the inverted scalar $s$ has
zero mean.

Consider the nonlinear sourced equation
\begin{equation}
 \mathcal F(s,m):=C_0s-Bm-\mathcal N(s,m)=0,
 \label{eq:n9v99-nonlinear-constraint}
\end{equation}
where $B:M\to S^*$ is bounded in the declared source norm.

\begin{theorem}[Weighted nonlinear scalar chart]
\label{thm:n9v99-weighted-chart}
Fix $r_s,r_m>0$ and $0\leq\kappa<1$.  Suppose, for
$\|s\|_S\leq r_s$ and $\|m\|_M\leq r_m$,
\begin{align}
 \|\mathcal N(s,m)-\mathcal N(t,m)\|_{S^*}
 &\leq\kappa\|s-t\|_S,
 \label{eq:n9v99-N-Lipschitz}\\
 \|Bm+\mathcal N(0,m)\|_{S^*}
 &\leq(1-\kappa)r_s.
 \label{eq:n9v99-centre-bound}
\end{align}
Then for every $\|m\|_M\leq r_m$ there is a unique
$s=\psi(m)$ in the closed $S$-ball of radius $r_s$ solving
\eqref{eq:n9v99-nonlinear-constraint}.  Moreover,
\begin{equation}
 \|\psi(m)\|_S
 \leq\frac{\|Bm+\mathcal N(0,m)\|_{S^*}}{1-\kappa}.
 \label{eq:n9v99-solution-bound}
\end{equation}
\end{theorem}

\begin{proof}
Define
\[
 T_m(s):=C_0^{-1}[Bm+\mathcal N(s,m)].
\]
Because $C_0^{-1}:S^*\to S$ has norm one,
\eqref{eq:n9v99-N-Lipschitz} makes $T_m$ a contraction with constant
$\kappa$.  The centre bound and
$\|T_m(s)-T_m(0)\|_S\leq\kappa r_s$ show that it maps the closed ball into
itself.  Banach's theorem gives the unique fixed point.  Comparing the fixed
point with zero gives \eqref{eq:n9v99-solution-bound}.
\end{proof}

\begin{corollary}[Relative precision perturbations are included]
\label{cor:n9v99-relative-precision}
If
\begin{equation}
 \mathcal N(s,m)=-E(m)s+\mathcal N_{\rm nl}(s,m)
 \label{eq:n9v99-linear-remainder-split}
\end{equation}
and
\begin{equation}
 \|C_0^{-1/2}E(m)C_0^{-1/2}\|
 +\operatorname{Lip}_{s,S\to S^*}(\mathcal N_{\rm nl})
 \leq\kappa<1,
 \label{eq:n9v99-relative-Lipschitz}
\end{equation}
then Theorem~\ref{thm:n9v99-weighted-chart} applies.  Thus the V98 relative
Schur perturbation and the nonlinear chart use the same dimensionless
smallness parameter.
\end{corollary}

\begin{proof}
In the norms \eqref{eq:n9v99-energy-dual-norms}, the Lipschitz norm of
$s\mapsto-E(m)s$ is exactly
$\|C_0^{-1/2}E(m)C_0^{-1/2}\|$.  Add the nonlinear Lipschitz constant.
\end{proof}

\subsection{First and second sourced responses}

Assume now that $\mathcal N$ is $C^2$ on the chart.  Put
\begin{equation}
 A(s,m):=C_0-D_s\mathcal N(s,m).
 \label{eq:n9v99-linearized-scalar-operator}
\end{equation}

\begin{lemma}[Uniform inverse in the energy frame]
\label{lem:n9v99-energy-inverse}
Under \eqref{eq:n9v99-N-Lipschitz},
\begin{equation}
 \|A(s,m)^{-1}\|_{S^*\to S}\leq\frac1{1-\kappa}.
 \label{eq:n9v99-energy-inverse}
\end{equation}
\end{lemma}

\begin{proof}
Factor
$A=C_0^{1/2}(I-K)C_0^{1/2}$ with
$K=C_0^{-1/2}D_s\mathcal N C_0^{-1/2}$ and $\|K\|\leq\kappa$.
The Neumann series for $(I-K)^{-1}$ gives the bound.
\end{proof}

\begin{proposition}[Exact first source response]
\label{prop:n9v99-first-response}
On the solution graph,
\begin{equation}
 D\psi(m)
 =A(\psi(m),m)^{-1}
 [B+D_m\mathcal N(\psi(m),m)].
 \label{eq:n9v99-first-response}
\end{equation}
Hence
\begin{equation}
 \|D\psi(m)\|_{M\to S}
 \leq\frac{
 \|B+D_m\mathcal N(\psi(m),m)\|_{M\to S^*}}
 {1-\kappa}.
 \label{eq:n9v99-first-response-bound}
\end{equation}
\end{proposition}

\begin{proof}
Differentiate
$C_0\psi-Bm-\mathcal N(\psi,m)=0$ and solve with
Lemma~\ref{lem:n9v99-energy-inverse}.
\end{proof}

For $h\in M$, write
\begin{equation}
 J_mh:=(D\psi(m)h,h)\in S\oplus M.
 \label{eq:n9v99-source-graph-lift}
\end{equation}

\begin{proposition}[Exact structured second response]
\label{prop:n9v99-second-response}
If $B$ is linear, then
\begin{equation}
 D^2\psi(m)[h,k]
 =A(\psi(m),m)^{-1}
 D^2\mathcal N(\psi(m),m)[J_mh,J_mk].
 \label{eq:n9v99-second-response}
\end{equation}
Consequently, with
\begin{equation}
 L_2(m):=sup_{h,k\neq0}
 \frac{\|D^2\mathcal N[J_mh,J_mk]\|_{S^*}}
 {\|h\|_M\|k\|_M},
 \label{eq:n9v99-structured-N-Hessian}
\end{equation}
one has
\begin{equation}
 \|D^2\psi(m)\|\leq\frac{L_2(m)}{1-\kappa}.
 \label{eq:n9v99-second-response-bound}
\end{equation}
\end{proposition}

\begin{proof}
Differentiate the first variation of
\eqref{eq:n9v99-nonlinear-constraint} once more.  The left side contributes
$A D^2\psi[h,k]$; the nonlinear chain rule contributes
$D^2\mathcal N[J_mh,J_mk]$.  Solve using
Lemma~\ref{lem:n9v99-energy-inverse}.
\end{proof}

\begin{remark}[Actual-input norm]
\label{rem:n9v99-actual-input}
The quantity $L_2$ is the V95 structured tangent-input norm.  It can be
strictly smaller than the full Hessian norm because the nonlinear remainder
receives only source-graph inputs.  Any claimed gain must be established for
the regulated Einstein--matter constraint, just as the audited NS note
computed its gain on the actual rank-one Oseen inputs.
\end{remark}

\subsection{When constraint reduction preserves positivity}

Let $\mathscr S(m,s)$ be a $C^2$ reduced-boundary action before elimination
of $s$, and let $s=\psi(m)$ be any $C^2$ constraint graph.  At a graph point
write its Hessian in blocks
\[
 H=
 \begin{pmatrix}
 H_{mm}&H_{ms}\\
 H_{sm}&H_{ss}
 \end{pmatrix},
 \qquad L:=D\psi(m),
 \qquad g:=D_s\mathscr S(m,\psi(m)).
\]

\begin{lemma}[Exact graph Hessian]
\label{lem:n9v99-graph-Hessian}
For $\mathscr S_{\rm red}(m)=\mathscr S(m,\psi(m))$,
\begin{align}
 D^2\mathscr S_{\rm red}[h,k]
 ={}&H_{mm}[h,k]+H_{ms}[h,Lk]+H_{sm}[Lh,k]
 +H_{ss}[Lh,Lk]
 \notag\\
 &\quad+g[D^2\psi[h,k]].
 \label{eq:n9v99-graph-Hessian}
\end{align}
\end{lemma}

\begin{proof}
Differentiate the first derivative
$D\mathscr S_{\rm red}[h]=D_m\mathscr S[h]+D_s\mathscr S[Lh]$ and apply the
chain rule.
\end{proof}

\begin{theorem}[Reduced-Hessian positivity margin]
\label{thm:n9v99-reduced-Hessian-margin}
Assume
\begin{equation}
 H_{mm}\succeq aI,
 \qquad
 -MI\preceq H_{ss}\preceq-cI<0
 \label{eq:n9v99-Hessian-signs}
\end{equation}
for $a,c,M>0$.  Define the stationary response
\begin{equation}
 L_0:=-H_{ss}^{-1}H_{sm}.
 \label{eq:n9v99-stationary-response}
\end{equation}
If
\begin{equation}
 \|L-L_0\|\leq\delta_L,
 \qquad
 \|g\|\,\|D^2\psi\|\leq\delta_g,
 \label{eq:n9v99-mismatch-debits}
\end{equation}
then
\begin{equation}
 D^2\mathscr S_{\rm red}
 \succeq(a-M\delta_L^2-\delta_g)I.
 \label{eq:n9v99-positivity-margin}
\end{equation}
In particular the reduced Hessian is positive if
$M\delta_L^2+\delta_g<a$.
\end{theorem}

\begin{proof}
Complete the square in the first four terms of
\eqref{eq:n9v99-graph-Hessian}:
\begin{align*}
 H_{mm}+H_{ms}L+L^*H_{sm}+L^*H_{ss}L
 ={}&H_{mm}-H_{ms}H_{ss}^{-1}H_{sm}\\
 &+(L-L_0)^*H_{ss}(L-L_0).
\end{align*}
Because $H_{ss}<0$, the stationary Schur correction
$-H_{ms}H_{ss}^{-1}H_{sm}$ is nonnegative.  The mismatch term is bounded
below by $-M\delta_L^2I$, and the final term in
\eqref{eq:n9v99-graph-Hessian} is bounded below by $-\delta_gI$.
\end{proof}

\begin{corollary}[Exact stationary constraint]
\label{cor:n9v99-exact-stationarity}
If the constraint equation is exactly
$D_s\mathscr S(m,s)=0$, then on its graph $g=0$ and
$L=L_0$.  Under \eqref{eq:n9v99-Hessian-signs},
\begin{equation}
 D^2\mathscr S_{\rm red}
 =H_{mm}-H_{ms}H_{ss}^{-1}H_{sm}
 \succeq aI.
 \label{eq:n9v99-exact-Schur-positive}
\end{equation}
\end{corollary}

\begin{proof}
Differentiate $D_s\mathscr S=0$ to obtain
$H_{sm}+H_{ss}L=0$.  Apply
Theorem~\ref{thm:n9v99-reduced-Hessian-margin} with zero debits.
\end{proof}

\begin{remark}[Why the compatibility debit is quadratic]
\label{rem:n9v99-quadratic-mismatch}
After completing the square, the response mismatch costs
$M\delta_L^2$, not a first-order term.  This is the nonlinear boundary
analogue of centering at the exact conditional lift in the audited YM note.
The separate $\delta_g$ term remains first order in the stationarity defect
and second response; it vanishes only for the correct action-derived
constraint.
\end{remark}

\begin{firewall}[A positive Hessian is local information]
\label{fire:n9v99-local-Hessian}
Theorem~\ref{thm:n9v99-reduced-Hessian-margin} proves strict convexity only
inside the declared chart.  It does not bound the action globally, exclude
other branches, or supply the normalized free-energy lower bound of V93.
Those require the V95 branch-exhaustion card and an independent large-field
coercivity estimate.
\end{firewall}

\subsection{Tail transport through the nonlinear source chart}

\begin{proposition}[Polynomial response tail]
\label{prop:n9v99-polynomial-tail}
Suppose on the branch-exhausting chart
\begin{equation}
 \|Bm+\mathcal N(0,m)\|_{S^*}
 \leq a_1X+a_2X^2+a_0
 \label{eq:n9v99-polynomial-source}
\end{equation}
for a random source seminorm $X\geq0$.  Then
\begin{equation}
 \|\psi(m)\|_S
 \leq\frac{a_1X+a_2X^2+a_0}{1-\kappa}.
 \label{eq:n9v99-polynomial-response}
\end{equation}
If $X$ has a uniform exponential-square moment, the family
$\|\psi(m)\|_S$ is uniformly tight and has an exponential tail in its
level.  Adding a uniformly tight retained mean coordinate preserves the
V92 random-seminorm tail card.
\end{proposition}

\begin{proof}
The deterministic bound is
Theorem~\ref{thm:n9v99-weighted-chart}.  If
$\sup_i\mathbb E e^{\alpha X_i^2}<\infty$, the event that the quadratic
polynomial exceeds $R$ forces $X_i^2\geq cR-C$ for constants independent of
$i$; Markov's inequality gives an exponential bound.  The sum with the mean
coordinate is tight by the V92 coupling lemma.
\end{proof}

\begin{proposition}[Chart-exit control from the source tail]
\label{prop:n9v99-chart-exit}
Suppose the contraction and positivity hypotheses hold whenever
$X\leq R$, with radii $r_s(R),r_m(R)$ satisfying
\eqref{eq:n9v99-centre-bound}.  If
$\sup_i\mathbb P\{X_i>R\}\leq a(R)\to0$, then the failure probability of
the declared scalar chart is at most $a(R)$, apart from separately recorded
calibration and determinant failures.
\end{proposition}

\begin{proof}
On $\{X\leq R\}$ all hypotheses hold by assumption, so a chart failure is
contained in $\{X>R\}$ or in one of the explicitly excluded ledger events.
\end{proof}

\subsection{Nonlinear sourced-Schur card}

\begin{definition}[Nonlinear sourced scalar card]
\label{def:n9v99-nonlinear-card}
A cofinal boundary construction passes the nonlinear sourced scalar card if
it provides:
\begin{enumerate}
\item a mean/closing-mode split and energy-dual spaces
\eqref{eq:n9v99-energy-dual-norms};
\item a uniform contraction constant $\kappa<1$ and centre bound
\eqref{eq:n9v99-centre-bound};
\item structured first and second responses from
Propositions~\ref{prop:n9v99-first-response} and
\ref{prop:n9v99-second-response};
\item an action-derived stationarity identity, or explicit mismatch debits
$\delta_L,\delta_g$ satisfying
$M\delta_L^2+\delta_g<a$;
\item branch-exit, determinant-density, and calibration-failure tails
tending to zero uniformly;
\item a normalized large-field coercivity/free-energy pair as in V93;
\item compatibility with the reflected half-space factorization of V96 and
the composite-source renormalization row of V98.
\end{enumerate}
\end{definition}

\begin{theorem}[Nonlinear conditional boundary closure]
\label{thm:n9v99-conditional-closure}
If the TT, gauge, matter, and boundary sectors satisfy their V96--V98 cards
and the scalar sector satisfies
Definition~\ref{def:n9v99-nonlinear-card}, then the reduced bosonic action
has a positive local Hessian on the good charts, the reconstructed scalar
field has the V92 tail card, and the limiting characteristic and reflected
functionals satisfy the conditional V91 and V88 reconstruction conclusions.
\end{theorem}

\begin{proof}
The positivity margin is
Theorem~\ref{thm:n9v99-reduced-Hessian-margin}; the scalar tail is
Proposition~\ref{prop:n9v99-polynomial-tail} plus branch exhaustion.  V96
supplies reflected positivity, V92 characteristic equicontinuity, V91
Minlos reconstruction, and V88 the OS quotient, under the remaining cards
listed in the hypothesis.
\end{proof}

\begin{assessment}[Status after ninth-series V99]
\label{assess:n9v99}
V99 constructs the nonlinear mean-zero scalar chart in the correct energy
and dual-energy spaces, derives its exact first and structured second
responses, and proves a quantitative reduced-Hessian margin.  Centering on
the stationary response makes the compatibility loss quadratic; a residual
stationarity defect remains a separate debit.  The main unresolved
mathematical input is now an action-derived regulated Einstein constraint
whose nonlinear remainder is uniformly contractive in these weighted
spaces and whose Standard Model source is renormalized with the required
tail.
\end{assessment}

\begin{programme}[Ninth-series V100: audit and completion ledger]
\label{prog:n9v99-V100}
V100 will perform the compulsory YM/NS audit, consolidate V91--V99 into one
implication and obstruction ledger, and identify the first unresolved
finite theorem for the next series.  It will then be frozen as the ninth
archive and followed by a compact new V1 containing context, the major
results, and the next attack.
\end{programme}

\begin{firewall}[Claim boundary after ninth-series V99]
\label{fire:n9v99-boundary}
V99 does not derive the nonlinear regulated Einstein constraint, prove its
weighted contraction, renormalize and bound the Standard Model stress
tensor, obtain the normalized global coercivity card, trivialize the chiral
determinant line, or construct the full SM--Einstein continuum.
\end{firewall}

\paragraph{Parent-version record.}
V99 was formed from
\texttt{Renewal\_SM\_Einstein\_Continuum\_Research\_Notes\_V98.tex}.
The V98 parent had $39\,533$ validator-counted source lines,
$1\,431\,676$ bytes, and SHA--256
\[
 \texttt{52B8CEAEEDBD2BD0DE058940F306CFD5BECC6FA6ABB654CE2EFFA59F4B3ED12C}.
\]

% =============================================================================
% END APPENDED NINTH-SERIES V99 MATERIAL
% =============================================================================

% =============================================================================
% BEGIN APPENDED NINTH-SERIES V100 MATERIAL
% =============================================================================

\section{V100 audit, Ward-relative boundary control, and ninth-series ledger}
\label{sec:v9v100}

\subsection{Compulsory companion audit}

The current companion notes were inspected read-only before the V100
consolidation.  At the audit instant the files were
\begin{align*}
 &\texttt{Renewal\_Yang\_Mills\_Mass\_Gap\_Research\_Notes\_V22.tex},
 &&349\,604\text{ bytes},\\
 &\texttt{Renewal\_Yang\_Mills\_Mass\_Gap\_Research\_Notes\_V23.tex},
 &&349\,604\text{ bytes},\\
 &\texttt{Renewal\_NS\_Stability\_Research\_Notes\_V26.tex},
 &&278\,283\text{ bytes}.
\end{align*}
The two YM files were byte-identical, with SHA--256
\[
 \texttt{72EDE7370D6BC47B04B35984E3F5D25C9AE7A848B57D8C1FE0BA14720D802EEB},
\]
and their content ended at eighth-series V22.  The unchanged NS digest was
\[
 \texttt{82C887F8C2C69E4F28FAD7C3DC8DE5B9099CB5A4BF431EBA1FFF3E91935978AD}.
\]

Since the V95 audit, YM V19--V22 has proved a localized two-norm contraction
from a vertical inverse gap, exact completed-square and covariance-rechart
identities, an explicit Wilson Hessian Lipschitz stock, a differentiated
gauge Ward identity away from critical backgrounds, and a second-moment
relative covariance bound at stationary translation-invariant backgrounds.
Two points transfer.  First, a massless covariance perturbation must be
controlled relatively in the precision metric, not absolutely.  Second,
the failure of gauge annihilation away from a critical background is fixed
by the action gradient and belongs to the moving reference connection, not
to a fictitious mass term.  The numerical Wilson incidence constants do not
transfer to gravity.

The NS V26 file has not changed.  Its restriction of derivative norms to
the actual nonlinear input cone remains incorporated in the structured
Hessians of V95 and V99; no further NS statement is imported.

\begin{assessment}[V100 audit decision]
\label{assess:n9v100-audit}
The next gravitational finite theorem will use a covariant quotient Hessian
at the moving calibrated background.  Its stationary boundary remainder
will be tested by a Ward-relative moment estimate, while the exact
connection-gradient term is assigned to the V81 moving reference ledger.
This changes the direction from seeking an absolute small boundary kernel
to proving the correct zero-momentum cancellation.
\end{assessment}

\subsection{A boundary Ward-relative moment lemma}

The following abstract transfer is stated for the order-one boundary
precision found in V97.  It is simpler than the order-two bulk YM estimate
but uses the same cancellation mechanism.

\begin{theorem}[Even second-moment boundary remainder]
\label{thm:n9v100-boundary-moment}
Let $L(k)$ be a self-adjoint matrix symbol for $k\in\mathbb R^d$ with
\begin{equation}
 L(0)=0,
 \qquad L(-k)=L(k),
 \qquad
 \sup_{|q|\leq K}\|D^2L(q)\|\leq M_2.
 \label{eq:n9v100-moment-hyp}
\end{equation}
Then, for $|k|\leq K$,
\begin{equation}
 \|L(k)\|\leq\frac{M_2}{2}|k|^2.
 \label{eq:n9v100-quadratic-vanishing}
\end{equation}
If a positive boundary precision satisfies
\begin{equation}
 C_0(k)\succeq c|k|I
 \qquad(k\neq0),
 \label{eq:n9v100-order-one-floor}
\end{equation}
then
\begin{equation}
 \|C_0(k)^{-1/2}L(k)C_0(k)^{-1/2}\|
 \leq\frac{M_2}{2c}|k|.
 \label{eq:n9v100-relative-boundary}
\end{equation}
In particular the relative perturbation vanishes at zero momentum despite
the closing spectrum of $C_0$.
\end{theorem}

\begin{proof}
Evenness gives $DL(0)=0$.  Taylor's formula along the segment $tk$ gives
\[
 L(k)=\int_0^1(1-t)D^2L(tk)[k,k]dt,
\]
which proves \eqref{eq:n9v100-quadratic-vanishing}.  The precision floor
implies $\|C_0(k)^{-1/2}\|^2\leq(c|k|)^{-1}$; sandwich the preceding bound.
\end{proof}

\begin{proposition}[Kernel second moment supplies $M_2$]
\label{prop:n9v100-kernel-moment}
Let a translation-invariant lattice kernel have symbol
\begin{equation}
 L(k)=\sum_xe^{-\mathrm i k\cdot x}\ell(x)
 \label{eq:n9v100-kernel-symbol}
\end{equation}
and finite second moment
\begin{equation}
 \mathfrak m_2(\ell):=\sum_x|x|^2\|\ell(x)\|<\infty.
 \label{eq:n9v100-kernel-second-moment}
\end{equation}
Then $\sup_k\|D^2L(k)\|\leq\mathfrak m_2(\ell)$.  If
$\sum_xe^{\mu|x|}\|\ell(x)\|\leq K_\mu$, then
\begin{equation}
 \mathfrak m_2(\ell)\leq2\mu^{-2}K_\mu.
 \label{eq:n9v100-exponential-to-moment}
\end{equation}
\end{proposition}

\begin{proof}
Differentiate \eqref{eq:n9v100-kernel-symbol} twice and take the absolute
sum.  The elementary inequality $r^2\leq2\mu^{-2}e^{\mu r}$ gives the final
bound.
\end{proof}

\begin{remark}[How $L(0)=0$ is to be proved]
\label{rem:n9v100-Ward-zero}
For a gauge or diffeomorphism invariant action at a critical background,
the differentiated Ward identity makes the quotient Hessian annihilate the
infinitesimal symmetry directions.  Combined with the spanning low-momentum
symbol of those directions, this can force $L(0)=0$.  Away from criticality,
the differentiated symmetry vector contributes a bilinear form linear in
the action gradient.  That term must be subtracted by the moving calibration
connection before Theorem~\ref{thm:n9v100-boundary-moment} is applied.
The precise gravitational connection term remains to be derived.
\end{remark}

\begin{corollary}[Perturbed boundary Schur stability]
\label{cor:n9v100-boundary-Schur-stability}
If
\begin{equation}
 \sup_{0<|k|\leq K}
 \frac{M_2|k|}{2c}\leq\eta<1,
 \label{eq:n9v100-relative-gate}
\end{equation}
then $C=C_0+L$ satisfies the V98 relative Schur bounds with parameter
$\eta$ on the declared carrier.  Low modes outside that carrier must be
retained rather than estimated by enlarging $\eta$.
\end{corollary}

\begin{proof}
Theorem~\ref{thm:n9v100-boundary-moment} gives the relative perturbation
hypothesis of the V98 relative Schur theorem.
\end{proof}

\begin{firewall}[Stationarity and covariance]
\label{fire:n9v100-stationarity}
Evenness and a second moment do not imply $L(0)=0$.  Diffeomorphism
invariance at a noncritical background does not give exact Hessian
annihilation without its connection-gradient correction.  Likewise,
translation invariance is lost on a general curved background.  A curved
version must use localized pseudodifferential principal symbols and
commutator remainders, not silently reuse the Fourier theorem.
\end{firewall}

\subsection{Ninth-series implication ledger}

The results V91--V99 can now be placed in one chain.

\begin{theorem}[Ninth-series conditional implication chain]
\label{thm:n9v100-implication-chain}
Suppose a cofinal regulated SM--Einstein construction supplies all of the
following:
\begin{enumerate}
\item pointwise convergence of physical characteristic and local Schwinger
functionals;
\item the V93 normalized coercivity/free-energy pair and a uniform local
reconstruction inequality;
\item a positive real bosonic body and V92 $L^s$ or coupled transport across
every determinant and hybrid-measure reweighting;
\item the V95 quantitative constraint atlas, coarea density, branch
exhaustion, and matched physical lift;
\item the V96 boundary-retained reflected factorization, including gauge
edge modes and the V83 graded fermion criterion;
\item the V97 positive TT boundary symbol and action-compatible elimination
of the negative conformal channel;
\item the V98 calibrated mean split, source Gram estimate, relative Schur
control, and renormalized composite-source tail;
\item the V99 nonlinear weighted contraction and reduced-Hessian margin;
\item anomaly cancellation with a compatible determinant-line
trivialization, plus the V87 symmetry and comparison ledgers.
\end{enumerate}
Then every cofinal pointwise limit has a Radon bosonic measure on the
physical distribution space, a consistent graded coefficient functional,
reflection-positive Schwinger functions, and the V88 OS Hilbert-space and
nonnegative-Hamiltonian reconstruction.  Two towers satisfying the common
V87 comparison card produce the same reconstructed theory.
\end{theorem}

\begin{proof}
Items 2, 3, 7, and 8 give the V92 random-seminorm tail, hence characteristic
equicontinuity.  V91--V92 and item 1 give the Radon Minlos measure.  Items 4
and 5 give the exact reflected-square or vanishing-negativity passage; V87
passes OS positivity and symmetry to the limit, and V88 performs the
quotient and semigroup reconstruction.  Items 6--8 ensure that the
gravitational scalar reduction is positive and uniformly reconstructible.
Item 9 supplies the missing graded global consistency.  Cross-tower
identity follows from convergence of the same characteristic and Schwinger
functionals and their uniqueness statements.
\end{proof}

\begin{assessment}[Results genuinely closed in V91--V100]
\label{assess:n9v100-closed}
The ninth series closes the following abstract mathematics:
\begin{enumerate}
\item characteristic equicontinuity from uniformly tight random seminorms,
including exponential, hybrid-density, physical-lift, and branch variants;
\item normalized action coercivity plus deterministic reconstruction implies
the tail card;
\item the Euclidean Einstein--Hilbert conformal action is unbounded below;
imaginary contour rotation is not a real characteristic functional; and a
simple fourth-order stabilizer can violate OS positivity;
\item a quantitative nonlinear constraint chart with exact first and
structured second responses, delta density, log-determinant response, and
matched tangent energy isometry;
\item exact reflected-square positivity from boundary-retained domain
decomposition, with nonlinear, gauge, tail, and approximate variants;
\item the flat half-space TT and conformal Dirichlet-to-Neumann signs;
\item the sourced scalar Fourier orders, mean-mode obstruction, relative
Schur stability, and nonlinear weighted response;
\item the reduced-Hessian positivity margin and the Ward-relative
second-moment boundary estimate.
\end{enumerate}
Each item is a proved implication or obstruction.  None asserts that the
physical regulator supplies all hypotheses.
\end{assessment}

\begin{assessment}[Open physical construction rows]
\label{assess:n9v100-open}
The full continuum remains open because the manuscript does not yet contain:
\begin{enumerate}
\item a concrete reflection-positive regulated Einstein constraint action
whose conformal equation is the stationary Schur constraint used above;
\item uniform nonlinear weighted estimates for that constraint on a curved,
moving calibrated background;
\item a constructed and uniformly bounded renormalized Standard Model
stress tensor in the required negative Sobolev source norm;
\item a normalized global coercivity and free-energy theorem including all
large-field, determinant, gauge, and branch sectors;
\item an anomaly-free chiral determinant-line gluing theorem compatible with
the reflection cut;
\item completed RG marginal trajectories, Euclidean symmetry restoration,
clustering, and the remaining sector reconstructions.
\end{enumerate}
\end{assessment}

\subsection{First theorem for the next series}

\begin{definition}[Covariant stationary-quotient card]
\label{def:n9v100-CSQ}
The next finite target, CSQ, consists of a regulated calibrated background
$\bar\phi$ and:
\begin{enumerate}
\item the exact differentiated diffeomorphism and internal-gauge Ward
identities for the full gauge-fixed/coarea reduced action;
\item an explicit connection-gradient bilinear $\Gamma(D S)$ such that the
quotient Hessian
$H^{\rm q}=D^2S-\Gamma(D S)$ annihilates all infinitesimal gauge directions;
\item a reflected half-space boundary Schur remainder
$L=H^{\rm q}_{\rm eff}-H^{\rm q}_{0,\rm eff}$ with a uniformly summable
second kernel moment;
\item a zero-symbol proof from the Ward identity and reflection parity;
\item the order-one relative gate
\eqref{eq:n9v100-relative-gate}, with closing modes retained;
\item assignment of $\Gamma(D S)$ to the V81 moving reference rather than to
the stable covariance.
\end{enumerate}
\end{definition}

\begin{theorem}[What CSQ would close]
\label{thm:n9v100-CSQ-closes}
If CSQ holds with relative parameter $\eta<1$, the perturbed nonzero-mode
boundary precision remains positive and its inverse and scalar Schur term
obey the V98 comparison bounds.  This supplies the linearized precision row
of the V99 nonlinear chart without an absolute spectral gap.
\end{theorem}

\begin{proof}
Items 2--4 permit
Theorem~\ref{thm:n9v100-boundary-moment}; item 5 gives the relative
perturbation norm below one.  The V98 relative Schur theorem then proves the
precision, inverse, and Schur comparisons.  Item 6 keeps the nonstationary
debit out of the stable covariance.
\end{proof}

\begin{assessment}[Status after ninth-series V100]
\label{assess:n9v100}
The ninth series has converted the gravitational positivity problem from an
ill-posed integral over a negative conformal direction into a precise
conditional reduced-boundary programme.  The first unresolved finite target
is CSQ, followed by the renormalized composite-source and global coercivity
rows.  The full SM--Einstein continuum is not constructed.
\end{assessment}

\begin{programme}[Tenth-series V1 restart]
\label{prog:n9v100-next}
This V100 file will be frozen with suffix
\texttt{\_f9}.  A compact tenth-series V1 will restate the context, summarize
the major proved results without copying the full legacy text, preserve the
claim boundary, and begin CSQ with the exact covariant Hessian identity on a
finite-dimensional gauge manifold.
\end{programme}

\begin{firewall}[Claim boundary after ninth-series V100]
\label{fire:n9v100-boundary}
V100 proves the audit transfer, the boundary moment theorem, and the
conditional implication ledger.  It does not prove CSQ, construct the
regulated Einstein constraint measure, bound the renormalized SM stress
tensor, trivialize the chiral determinant line, close the RG trajectories,
or construct the full SM--Einstein continuum.
\end{firewall}

\paragraph{Parent-version record.}
V100 was formed from
\texttt{Renewal\_SM\_Einstein\_Continuum\_Research\_Notes\_V99.tex}.
The V99 parent had $39\,970$ validator-counted source lines,
$1\,446\,520$ bytes, and SHA--256
\[
 \texttt{00812D529E8C3C9A3A00B7EA62EC062B7F14E4374D2B0CB01AC4721FFFF2A812}.
\]

% =============================================================================
% END APPENDED NINTH-SERIES V100 MATERIAL
% =============================================================================
\end{document}
